% Gerado por tools/tei2latex.py a partir de entrada.xml. Regenerar, não editar.

%%>>> parte-1

%%>>> art-28

%%>>> t204

a) Se a relação entre a linguagem e as idéas não fosse \emph{convencional}, ou \emph{contingente}, mas \emph{natural}, ou \emph{necessaria}, as palavras \pg{333} \pg{334}  não poderiam estar sujeitas a variações, mas deveriam conservar sempre a mesma significação; porque o signal \pg{335}  \emph{natural} é uniforme e constante, como a natureza. Ora não é assim; porque as palavras estão sujeitas a variações, na mesma lingua, como é por todos sabido. Logo a relação entre a linguagem e as idéas é \emph{convencional}.

b) Se a relação entre a linguagem e as idéas fosse \emph{natural}, á mesma palavra deveria despertar a mesma idéa, e palavras diversas ideas diversas; porque, posta a causa, deve seguir-se necessariamente o effeito relativo. Ora não acontece assim; porque a mesma palavra desperta, ás vezes, idéas diversas em diversas pessoas, constituidas nas mesmas circumstancias physiologieas, como, pelo contrario, palavras diversas despertam sempre a mesma idéa. Logo a relação entre a linguagem e as idéas é \emph{convencional} \footnote{As idéas e, em geral, os actos internos, só os percebe o sujeito, que os forma ou produz; os outros só podem percebel-os por meio de \emph{sgnaes}. --- \emph{Signal} é uma coisa que leva ao conhecimento de outra coisa, com a qual está associada. Divide-se em \emph{natural} e \emph{convencional}, conforme a relação entre o signal e a coisa significada deriva da \emph{natureza} ou da convenção. --- A. vos é um \emph{signal}. D'ella serve-se o homem --- não só para emittir sons \emph{inarticulados}, que exprimem de algum modo as proprias affeições (o que fazem tambem os animaes irracionaes) --- mas tambem para emittir \emph{sons articulados}, isto é, para formar \emph{palavras}, seguindo a direcção da intelligencia e o imperio da vontade (o que é exclusivamente proprio do homem). --- Ha uma outra linguagem constituida por \emph{gestos convencionaes}, e chama-se \emph{mimica}; --- e uma outra composta de signaes \emph{graphicos}, e é a \emph{escriptura}. Mas a denominação de \emph{linguagem}, dada á \emph{mimica} e á \emph{escriptura}, é menos propria. --- Ha differentes \emph{typos} de linguagem, conforme os differentes modos de juntar os sons na formação das palavras. Os principaes typos, aos quaes se reduzem todas as formas de linguagem, são tres: o das linguas \emph{monosyllabicas}, que são formadas por sons destacados e de uma só syllaba immutavel, significativa de uma idéa especial (\emph{raiz}) --- o das linguas \emph{agglutinantes}, em que domina uma fusão de monosyllabos \emph{radícaes}, cada um dos quaes conserva a propria significação, --- o das linguas \emph{flecivas}, que admittem mudanças nas proprias raizes e em que as relações do discurso são significadas pela união de prefixos ou suffixos. --- O \emph{terceiro} typo, que provavelmente é uma derivação do \emph{segundo}, é o mais perfeito, porque é mais apto para representar a immensa variedade de ideas, que o homem pode e quer manifestar. --- Muitos escriptores, entre os quaes \emph{Mas Müller} e \emph{Curtéus}, opinam que a lingua fallada pelo primeiro homem fosse o \emph{monosyllabica}.\par A conclusão, exposta no texto, é evidente. A palavra \emph{burro} é uma só, e todavia para um italiano significa a \emph{manteiga}, e para um portuguez denota o asno. --- Como, pelo contrario, as palavras \emph{ensis} e \emph{gladius} são diversas, e todavia significam uma e e mesma coisa --- a \emph{espada}. Diga-se o mesmo com relação a todos os \emph{synonymos}. Concluimos que a linguagem, embora possa e deva dizer-se \emph{natural} ao homem, emquanto este é \emph{naturalmente} dotado de orgãos vocaes e de instincto e razão para servir-se d'elles, comtudo\pgn{336}é essencialmente \emph{convencional}, emquanto a relação entre a palavra e a idea depende de uma \emph{convenção}, ou \emph{accordo}, estabelecido e confirmado pelo uso. --- Quando dizemos, porém, que uma palavra foi ligada a uma idéa por uma convenção, --- não entendemos que os homens, antes de começarem a fallar, se tenham reunido em congresso e tenham proposto uma certa lingua, como, ás vezes, têm praticado, tratando-se de significar novos objectos ou novas idéas, --- mas entendemos que o uso estabeleceu e firmou essa convenção, embora o uso a tenha, ás vezes, mudado. --- Não se oppõe ao que deixamos dito o facto de algumas palavras terem uma significação \emph{natural}, como são as que reproduzem o som das coisas (por. ex. \emph{w/ulado}). São excepções; porque são muito poucos os objectos, que possam ser significados de um modo \emph{émitativo}, nem o são em \emph{todas} as linguas. ]\par (i) A \emph{palavra} é essencialmente um \emph{signal} e um signal de \emph{ideas}. Por isso, se não quizermos acceitar os absurdos do \emph{nominalismo}, devemos admittir que a idéa precede natural e essencialmente a palavra, como a coisa significada precede o signal. --- Nenhum vocabulo pode despertar uma idéa no nosso espirito, se n\&o conhecermos antecedentemente a sua significação, isto é, se não possuirmos a idea, que elle exprime. A palavra \emph{externa é effeito} da \emph{interna}, e por isso é \emph{posterior} a esta. Assim como eu não posso fallar com 0s outros sem que primeiramente tenha fallado commigo, assim tambem ninguem pode perceber a linguagem dos outros,sem que primeiramente tenham apprehendido a fallar comsigo. Diz S. Thomaz: « Planum est quod illud quod significatur interius in anima existens \emph{prius} est quam ipsum verbum voce prolatum, utpote \emph{causa} eius existens » (\emph{De diff}. \emph{verbi div}. \emph{et hum}.).\par De resto, \& propria experiencia, tanto \emph{interna} como \emph{externa}, attesta que a idéa é independente da palavra. --- Todos estão convencidos de que, para}.

%%>>> t205

\pg{336} \tese{205}{A linguagem suppõe as idéas.}

á) \textasciicircum{} linguagem compõe-se da signaes \emph{arbitrarios}, que a vontade humana, e não a natureza, assumiu para significar as coisas. Ora a vontade, attribuindo uma certa significação a uma, certa palavra, por si indifferente para exprimir esta ou aquella coisa, segue necessariamente a intelligencia, e d'esta recebe a idéa, que se quer exprimir, e por isso a idéa deve ser presupposta á palavra; assim, quando os homens estabeleceram que a palavra \textsc{Deus} signifieasse o supremo Auctor de todas as creaturas, deviam já ter a idéa desse Ente nobilissimo. Logo a linguagem suppõe as idéas.

a) Se a palavra não suppuzesse a idéa, mas esta nascesse d'aquella, seguir-se-hia que quem ignorasse uma coisa havia de percebel-a ao ouvir a palavra correspondente. Ora esta consequencia é falsa; pois quem ignora uma coisa não chega a conhecel-a pelo facto de ouvir a palavra, que a denota; assim quem não sabe o que é a \emph{electricidade}, não a percebe nem depois de ouvir essa palavra. Logo a linguagem suppõe as idéas \footnote{\pgn{337}ter idéas claras e distinctas, relativas ás coisas materiaes e tambem ás espirituaes, não é preciso mais do que a união das proprias coisas com a intelligenecia por meio das respectivas especies intelligiveis, sem auxilio da palavra. A nossa intelligencia, levada pelo desejo de saber e estimulada pela acção dos objectos, percebe primeiramente a essencia das coisas materiaes e depois, servindo-se da synthese e da analyse, eleva-se ao conhecimento de outros objectos, até ao conhecimento de Deus. --- A experiencia \emph{interna} é confirmada pela \emph{externa}. Os \emph{surdos-mudos} adquirem e possuem muitos conhecimentos relativos á ordem material e espiritual, e todavia não ouvem e não fallam. Muitas vezes percebemos uma coisa e não encontramos o vocabulo correspondente; temos a \emph{idéa}, mas falta-nos o \emph{termo}: o que mostra que a idea não nasce da palavra.\par \emph{De Bonald} (1754-1840), tradicionalista, sustentou --- que a palavra e a idéa são duas coisas inseparaveis, --- que a palavra produz a idéa, --- e que sem palavras não ha idéas. --- Parece incrivel que um homem tão illustrado tivesse cahido num erro tão grosseiro. E, na verdade, de que modo poderia a palavra produzir a idéa? Quando se falla, ou percebemos a palavra, ou não. Sea percebemos, é porque já temos a idéa significada, e assim a palavra não produz a idéa. Se não a percebemos, ella não produz idéa nenhuma.}.

%%>>> t206

\pg{337} \tese{206}{A linguagem não é necessaria, mas é utilissima.}

a) Não é necessaria. --- É um corollario da proposição antecedente. Se a linguagem suppõe as idéas, é claro que a palavra não é necessaria para a consecução das verdades, tanto da ordem material como da ordem metaphysica e moral.

b) Mas é utilissima, --- E utilissima \emph{para quem falla} e \emph{para quem escuta}. --- \emph{Para quem falla}; porque a palavra é o \emph{symbolo sensivel} para a formação e fixação das idéas, emquanto é sobretudo no phantasma da \emph{palavra} que a intelligencia conhece e reconhece o proprio objecto. --- \emph{Para quem escuta} ; porque a linguagem é o meio \emph{externo} e \emph{social}, pelo qual o homem manifesta os seus sentimentos e transmitte aos autros as proprias ideas, e assim se torna o vinculo da convivencia humana e o principal factor de todo o progresso \footnote{Ninguem pode contestar a grandissima \emph{utilidade} da linguagem para a nossa vida \emph{intellectual} e \emph{social}. --- A mossa intelligencia, como vimos, para que possa pensar e tornar a pensar, precisa de imagens ou representações sensiveis. E' a imaginação que as forma. Muitas vezes o trabalho desta faculdade é assaz intenso, sobretudo quando se trata de procurar ou variar imagens ou representações, que tenham relação com objectos abstractos ou immateriaes. Por isso, as imagens mais aptas para o conhecimento são as que são mais claras, mais faceis, mais simples e mais flexiveis. Ora a imagem sensivel da \emph{palavra}, entre todas as imagens, --- é a mais simples e clara, porque tem a sua significação determinada, --- é a mais facil e flexivel, porque, conservando-se na memoria, pode ser representada sem difficuldade, e é capaz de significar todas as coisas, mesmo as \emph{immalteriaes}. E' por isso que, ás vezes, faltam-nos todas as imagens sensiveis de objectos abstrusos, mas não nos falta a imagem da \emph{palavra}, que os exprime. --- Uti-\par 22 Vor. II\pgn{338}lissima para a vida \emph{intellectual}, a palavra não o é menos para a vida \emph{social}. O homem nasceu para 2 sociedade, e precisa de signaes sensiveis para estreitar e conservar os laços, que o prendem aos seus semelhantes. Ora o signal mais apto é a \emph{palavra}... A utilidade da palavra vê-se mais claramente, quando consideramos a difficuldade, que se encontra numa conversa com uma pessoa muda ou ignorante da lingua... Mas a palavra é tambem utilissima para o completo desenvolvimento intellectual. E' verdade que ha um duplice meio para apprehender as sciencias: a investigação propria e o magisterio alheio; mas é tambem verdade que ninguem tem chegado á celebridade, unicamente pelo esforço proprio. A investigação e o descobrimento de verdades, relativas a uma sciencia, exige tantas e tão grandes qualidades intellectuaes, moraes e physicas, que difficilmente se encontram reunidas numa só pessoa. E não fallamos nas circumstancias da vida, que difficultam, muitas vezes, se não tornam impossivel, um estudo serio e constante, necessario para adquirir uma certa sciencia. Os homens mais sabios de todos os tempos deveram sempre e muito a seus mestres... Diga-se o mesmo com relação ao estudo das artes, á pratica do commercio, á apprendizagem de todo e qualquer ramo do conhecimento e da actividade humana.}.

%%>>> t207

\pg{338} \tese{207}{O homem não inventou, mas podia inventar a linguagem.}

a) O homem não inventou a linguagem. --- É a questão de facto. Todas as tradições religiosas dos povos, sobretudo a tradição christã, attestam que o homem foi creado adulto num estado de civilisação e de felicidade, collocado na sociedade domestica, constituido chefe e mestre de todo o genero humano, e dotado do dom da \emph{palavra}. --- Com effeito, a Sagrada Escriptura, que, mesmo prescindindo de todo o caracter sobrenatural, é sempre o mais antigo e o mais auctorisado documento historico, diz que o primeiro homem \emph{fallou desde o principio}; o que não teria podido fazer, se a invenção da linguagem tivesse sido deixada á sua industria e ao seu estudo \footnote{Conta o livro do \emph{Genesis} (Il, 19) que Deus apresentou os animaes ao primeiro homem, para que este lhes impuzesse os nomes, e que Adão effectivamente deu a cada um o nome. Este trecho mostra que o primeiro homem fallou, e fallou desde o principio, embora não resolva \emph{explicitamente} a questão relativa á \emph{origem} da linguagem. Mas é facil deduzir do contexto que á lingua, fallada por Adão, não era o effeito dos seus esforços, mas um dom de Deus, emquanto Deus ou infundiu no homem a sciencia da linguagem, ou lhe concedeu a faculdade de encontrar immediata e facilmente os termos significativos dos seus pensamentos. E esta deducção esta em perfeita harmonia com a sapientissima Providencia divina. Tendo sido destinado ao altissimo officio de chefe e mestre do genero humano, devia o homem, \emph{desde o principio}, possuir todos os meios necessarios para o cumprimento d'esse officio. Ora um dos meios mais indispensaveis era o poder de \emph{fallar}, e de fallar \emph{immediatamente}. Quanto tempo e quantos esforços teriam sido necessarios a Adão, occupado em tantas e tão graves coisas, para procurar e formar uma lingua! --- Esta opinião --- que Deus infundiu no homem a sciencia da \emph{linguagem} --- não só é commum entre os philosophos christãos,.\pgn{339}mas é defendida por muitos incredulos e racionalistas, como \emph{Rousseau}, \emph{Humboldt}, \emph{Herder}. \emph{Rousseau} chama a linguagem um \emph{dom do ceu}. --- Deus, provavelmente, revelou as \emph{raizes primitivas} da linguagem. Em seguida, sobretudo depois da dispersão do genero humano, a linguagem devia desenvolver-se em conformidade com as condições psychicas, physiologicas, geographicas, etnographicas, e com um certo convencionalismo e arte, dependente sempre da natural estructura do mechanismo vocal e das circumstancias locaes. E' neste sentido que se pode e deve acceitar a hypothese da \emph{elaboração progresiva}, ou do \emph{desenvolvimento successivo} da linguagem. Por isso a opinião que a primeira lingua não fosse a \emph{flexiva}, mas a \emph{monosyllabica}, encontra o appoio de muitos sabios; e este appoio funda-se não só no exame das linguas mais antigas, que têm raizes communs, --- mas tambem no das linguas \emph{polosyllabicas}, que manifestam o producto da elaboração progressiva, devida á variação das condições ethnicas nos diversos periodos da civilisação humana.\par *\par Muitos escriptores rejeitam uma extraordinaria intervenção divina com relação á \emph{origem} da linguagem, e, para explicar o \emph{modo} por que o homem começou a fallar, recorrem a varias e differentes hypotheses. As principaes são \emph{quatro}: a \emph{evolucionista}, a da \emph{elaboração progressiva}, a da \emph{instituição convencional}, e a do \emph{instincto philologico}. De cada uma faremos uma breve exposição com a respectiva critica.\par \emph{a}) \emph{Hypothese evolucionista}. E' uma invenção de \emph{Darwin}, \emph{Spencer}, e admitte que os homens, barbaros e salvagens no periodo primitivo, passaram, depois de uma longa serie de seculos, da linguagem \emph{natural} (gestos, sons emocionaes, ete.) para a \emph{convencional}. Na evolução succederam-se tres formas glotticas: a \emph{cinetica} ou \emph{mimica} (gestos, movimentos imitativos das coisas), a \emph{phonica} (sons organicos por meio de palavras de harmonia imitativa), \& \emph{glottica perfeita}, que depois, pela reflexão e pela cultura, se tornou artificial e deu origem ás varias linguas. Na passagem da forma \emph{phonica} para a \emph{gloitica}. \emph{perfeita} houve tres periodos: o \emph{grito}, \emph{a vocalisagüo}, a \emph{articulação}.\par Esta explicação da \emph{origem} da linguagem não pode admittir-se. Porquanto, funda-se no \emph{transformismo}, que, como veremos, é falso. --- Além d'isso, suppõe que o primeiro estado do homem fosse o de \emph{selvagem}, sem o vinculo social da palavra. Ora, se isso fosse verdade, as linguas primitivas deviam ser muito imperfeitas, e os selvagens, que ainda ha no mundo, deveriam fallar uma lingua muito pobre e imperfeita, e essa lingua deveria revelar esse primitivo estado de barbarie. Mas não é assim. A historia da glottologia reconhece --- que as linguas primitivas eram relativamente perfeitas, --- que a lingua, de que se serve o selvagem, á muitas vezes um prodigio de riqueza e de logica, e deve ser considerada como um monumento vivo e indelevel da civilisagão primitiva (Cf. \emph{De Maistre}, Soirées, t. I, p. 121). Entre a linguagem e a intelligencia ha aquella mesma relação, que existe entre o effeito e a causa, entre o signal e a coisa significada. Por isso a linguagem de um povo é um criterio infallivel para se averiguar o seu estado primitivo e as suas affinidades com os outros povos. Um povo que possue\pgn{340}uma lingua rica e logica não pode ser senão um povo \emph{degenerado} e \emph{decahido} de um estado florescente de cultura intellectual. Esta opinião é dividida por muitos sabios, como \emph{Mas Müller}, \emph{De Nadallac}, e outros. \emph{Horacio Hale} cita, para exemplo, a lingua dos australianos, considerados como uma das raças mais inferiores, e diz que ella é um prodigio de riqueza grammatical e de perfeição logica. Nas questões de anthropologia, como nas de biologia, não ha só \emph{progresso}, ha tambem \emph{regresso}, e o estudo d'este é, ás vezes, mais instructivo que o d'aquelle.\par 5) \emph{Hypothese da elaboração progressiva}. E' admittida por \emph{Rabier}, \emph{Witney}. Embora não acceite a passagem do homem do estado selvagem para o civilisado, comtudo chega a conclusões quasi semelhantes ás dos evolucionistas. Reconhecendo no homem a intelligencia e a capacidade de formar uma linguagem, sustenta que effectivamente o homem, sem auxilios extraordinarios do Creador, chegou a formar a linguagem \emph{artificial} por uma \emph{progressiva elaboração} da linguagem \emph{natural}.\par Não negamos que, na formação da linguagem, tenha havido uma elaboração progressiva, mas devemos reconhecer que essa elaboração suppõe uma revelação primitiva. Sem esta revelação, o homem não teria chegado 2 formar a linguagem senão depois de muito tempo. Ora o homem devia, desde o principio, possuir a sciencia da linguagem, porque, desde o principio, viveu em sociedade e precisou da linguagem.\par \emph{c}) \emph{Hypothese da instituição convencional}. Attribue-se a \emph{Democrito}. Segundo esta hypothese, um homem de genio, depois de um primitivo periodo de vida selvagem, inventou um systema de sons articulados, que, em seguida, se tornou commum por uma especie de \emph{convenção}.\par Tambem esta é uma hypothese arbitraria. O homem foi civilisado desde o principio. Não se pode conceber um ser racional e naturalmente sociavel, como é o homem, destituido, por muito tempo, de um meio tão necessario e comtudo tão natural, como é a palavra. --- Além d'isso, não se concebe como um homem selvagem tenha podido, só por si, formar um systema linguistico relativamente perfeito e impol-o aos outros selvagens.\par \emph{d}) \emph{Hypothese do instincto philologico}. E' \& opinião de \emph{Mac Müller}, \emph{Reid}, \emph{Jouffroy}, \emph{Garnier}, \emph{Renan}, que attribuem a origem da linguagem a um instincto especial --- o \emph{instincto philologico}, --- pelo qual as mesmas ideas suggerem, naturalmente e em todos os homens, as mesmas palavras.\par Nem esta hypothese pode admittir-se. Está averiguado que as linguas nos primeiros periodos historicos apparecem já como o producto, não do \emph{instincto}, mas de um grande trabalho reflexivo. --- Além d'isso, esse supposto instincto, sendo \emph{natural}, deveria ser \emph{uniforme} em todo o genero humano e, por isso, produzir uma \emph{unica} lingua. Ora não é assim. --- Ha mais. Essa lingua, produzida pelo instincto, deveria ser não só unica, mas tambem \emph{invariavel}, porque invariavel é o instincto \emph{natural}. Ora quem ignora a con- : tinua variabilidade das linguas? --- Finalmente, esse instincto; sendo \emph{natural}, deveria produzir o seu effeito em todos os homens, dotados do apparelho vocal, devidamente desenvolvido. Mas não produz; porquanto os \emph{surdos-mudos} são dotados de orgãos glottologicos, e todavia nao fallam.\par e}.

b) Mas podia inventar a linguagem, --- Para que o homem podesse inventar a linguagem, tres requisitos eram necessarios e sufficientes: á \emph{faculdade} de \emph{pensar} e \emph{generalisar}, \emph{á faculdade} \pg{340}  de produzir sons \emph{articulados}, e a \emph{tendencia} para communicar as proprias \emph{tidéas} e os proprios \emph{sentimentos}. Ora o homem é dotado d'esse triplice requisito; pensa e generalisa, produz \pg{341}  sons articulados, tem uma grande tendencia para communicar aos outros o que se passa no intimo do seu espirito. Logo o homem podia inventar a linguagem \footnote{Resolvida a questão \emph{historica}, resolve-se agora a questão \emph{philosophica}. Podia o homem, sem uma intervenção extraordinaria de Deus, inventar a linguagem? Como se vê, fallamos --- não do \emph{facto}, mas da \emph{possibilidade theorica}. Posta assim a questão, respondemas que Deus podia crear o homem, dotal-o da faculdade e tendencia para entender e fallar, deixando-lhe, comtudo, a liberdade e o cuidado de formar uma linguagem. Nem isso era impossivel para o homem, desde que se lhe concedesse, na sua primeira origem, um conveniente desenvolvimento, necessario para tal formação. De facto, o homem, vendo uma coisa sensivel e material, podia indical-a aos outros por meio de gestos, e em seguida, lovado pela mesma tentencia para fallar, podia produzir uma palavra articulada com o fim de significar essa mesma coisa, primeiramente indicada com o gesto. Formadas as palavras, significativas das coisas materiaes, podia servir-se d'ellas para exprimir as coisas immateriaes, que com as materiaes tivessem uma certa analogia. Todavia devemos admittir que o trabalho teria sido difficil, longo e confuso, e o resultado incerto, inconstante e imperfeito.}. \textbar{}

%%>>> art-29

\section*{Artigo V.\\ Volições, seu objecto e divisão.}
\addcontentsline{toc}{section}{Artigo V. --- Volições, seu objecto e divisão.}

%%>>> t208

\pg{339} \tese{208}{Volições.}
\emph{Volições} são os actos, pelos quaes a nossa vontade tende para o bem \emph{immaterial}, percebido por meio da \emph{intelligencia}. --- Assim como, depois de uma percepção \emph{sensitiva}, se produz no homem uma inclinação para o bem \emph{sensivel}; assim tambem, depois de uma percepção \emph{intellectual}, nasce em nós uma tendencia para o bem \emph{immaterial} \footnote{O acto da \emph{vontade}, pois, para ser \emph{voluntario}, deve ser-precedido pelo acto da intelligencia. --- Todavia a \emph{volição} e a \emph{percepção} são actos distinctos e diversos, como distinctas e diversas são as faculdades, de que derivam. --- D'ahi se vê o erro de \emph{Fouillée}, que inventou a theoria das \emph{idéas-forças}.\par A \emph{volicüo}, ou o \emph{voluntario}, differe do \emph{espontaneo}. Por quanto, o \emph{espontanen} não exige a percepção da intelligencia nem a deliberação da vontade; ao passo que a \emph{volição} não. se concebe sem esses dois actos. Às operações dos animaes são \emph{espontaneas}, mas não são \emph{voluntarias}.\par O bem \emph{universal} é o objecto \emph{proprio} e \emph{directo} da vontade: --- \emph{proprio}, porque e \emph{immaterial}, como \emph{immaterial} é a vontade; --- \emph{directo}, porque, sendo \emph{adequado}, é desejado em \emph{si} e \emph{por si}. --- O bem \emph{particular}, embora seja objecto \emph{proprio} da vontade, todavia não é \emph{directo}, porque a vontade tende para elle, não como para \emph{fim}, mas como para \emph{meio}; tal é a \emph{sciencia}, \emph{á virtude}, etc. Os bens \emph{materiaes} não são objectos \emph{proprios} da vontade, mas podem ser \emph{iadirectos}, emquanto a vontade os deseja para as faculdades \emph{organicas},}.

%%>>> t209

\tese{209}{Objecto das volições.}
O objecto das \emph{volições} é o mesmo que o da \emph{vontade}: \emph{o bem}. --- O bem é \emph{universal} ou \emph{particular}, \pg{342}  conforme é dotado de todas as perfeições, ou só possue uma perfeição finita e relativa á sua especie.

%%>>> t210

\tese{210}{Divisão das voligües.}
As \emph{volicoes} dividem-se --- em \emph{necessarias} e \emph{livres}, --- \emph{elicitas} e \emph{imperadas}.

a) Voligóes necessarias e livres, --- \emph{Necessarias} são as que a vontade não póde deixar de produzir; taes são as \emph{volicoes}, que têm por objecto o bem \emph{universal}. --- \emph{Livres} são as que a vontade produz, mas que podia deixar de produzir; taes são as \emph{voligoes}, que se referem a um bem \emph{particular} \footnote{Tendo provado (n.º 130 e 131) que a vontade não é livre na sua tendencia para o bem \emph{universal}, mas é livre na sua tendencia para o bem \emph{particular}, é evidente que em nós se dão volições \emph{necessarias} e \emph{livres}.}.

b) Voligóes elicitas e imperadas. --- -\emph{E/icitas} são as que derivam \emph{immediatamente} da vontade; como são os actos de \emph{querer}, de amar, de \emph{consentir}, etc. --- \emph{Imperadas} são as que derivam \emph{immediatamente} de uma outra faculdade, \emph{movida} pela \emph{vontade} ; como são os actos de \emph{escrever}, de \emph{passear}, ete. \footnote{O acto \emph{imperado é voluntario} e por isso \emph{imputavel}; porque é a vontade que o executa, embora \emph{mediante} uma outra faculdade, e porque o effeito mais depende da causa motora do que do instrumento.}. --- 9M. \textbf{Voligües} e sentimentos. --- As \emph{volicoes} são acompanhadas por algumas \emph{emoções}, agradaveis ou desagradaveis, chamadas vulgarmente \emph{sentimentos}. --- Os \emph{sentimentos}, pois, são \emph{emoções}, \emph{agradaveis ou desagradaveis}, \emph{que affectam a vontade na sua tendencia para o bem ou na sua fuga do mal} \footnote{A palavra --- \emph{sentimento} --- tem muitas e diversas significações.\par Uma vezes denota o consenso unanime dos homens em admittir, sem exame nem raciocinio, certas verdades primitivas, especulativas ou praticas; e então corresponde ao chamado \emph{senso commun}.\par Outras vezes significa a \emph{emoção produzida na alma por uma percepção intellectual}. Assim dizemos: \emph{sentimento} da \emph{verdade}, do \emph{bem}, do \emph{bello}, etc. --- O sentimento, tomado nesta significação, divide-se em \emph{esthetico}, \emph{intellectual}, \emph{moral} ou \emph{religioso}, segundo a emoção deriva da percepção do \emph{bello}, da \emph{verdade}, do \emph{bem} ou das \emph{coisas religiosas}.\par Outras vezes não se distingue das \emph{paixões}, e significa a \emph{emoção}, \emph{agradavel} ou \emph{desagradavel}, \emph{resultante da apprehensão do bem ou do mal}. D'ahi a divisão da sensibilidade \emph{interna} em \emph{sentimentos} de tres especies, a saber: \emph{pessoaes}, \emph{sociaes} e \emph{racionaes}, conforme se referem ás \emph{pessoas}, \emph{á sociedade}, ou êncerram um conceito da \emph{razão}. --- Por isso podem chamar-se \emph{sentimentos} --- não só as emoções que nascem da apprehensão do bem ou do mal meramente sensivel, --- mas tambem as que affectam ao mesmo tempo o appetite \emph{sensitivo} e o \emph{racional}, como o amor da \emph{familia}, \emph{á inveja}, etc.\par \emph{A} definição, que apresentamos no texto, é a que corresponde á realidade. O sentimento não se refere a uma faculdade \emph{perceptiva}, mas a uma faculdade \emph{appetitiva}, --- Esta faculdade \emph{appetitiva} não é a \emph{sensitiva}, mas é á \emph{intellectual},\pgn{343}chamada commummente \emph{vontade} ; porquanto, se os actos do appetite \emph{sensitivo} têm por objecto um bem \emph{sensivel}, os \emph{sentimentos} têm por objecto um bem \emph{immaterial}. --- Os \emph{sentimentos}, pois, não são phenomenos meramente passivos; porque são modificações da vontade, que tende para o seu bem.\par Existem, pois, muitas differenças entre as \emph{paicües} e os \emph{sentimentos}. As \emph{paixões} têm por objecto um bem \emph{sensivel}, residem no appetite \emph{sensitivo}, e produzem sempre uma mudança organica; ao passo que os \emph{sentimentos} têm por objecto um bem \emph{tinmmaterial}, residem na \emph{vontade}, e ordinariamente não produzem mudança organica.\par Dissemos que os \emph{sentimentos} não produzem \emph{ordinariamente} mudança\par organica; porque a emoção da \emph{vontade}, se for vehemente, reflecte-se e encontra um echo nas faculdades \emph{sensitivas}; de modo que nesse caso a vontade e o appetite \emph{sensitivo} se unem no amor do bem e no odio do mal.}.

%%>>> t219

\pg{343} \tese{219}{Divisão dos sentimentos.}
Os \emph{sentimentos} São \emph{onze}, a saber: \emph{amor} e \emph{odio}, \emph{desejo} e \emph{abominação}, \emph{alegria} e \emph{tristeza}, \emph{esperança} e \emph{desespero}, \emph{audacia} e \emph{medo}, e, finalmente, \emph{tra}. Com effeito, os \emph{sentimentos} devem ser tantos, quantas são as \emph{paéaoes} ; pois os mesmos movimentos, que o bem e o mal \emph{sensivel} excitam no appetite \emph{sensitivo}, são excitados pelo bem e pelo mal \emph{immaterial} no appetite \emph{intellectual}, que é a \emph{vontade} \footnote{Os actos ordenados da vontade, repetidos frequentemente, produzem os \emph{habitos moraes}: a \emph{prudencia}, a \emph{justiça}, \emph{á temperança}, e á \emph{fortaleza}.\par A \emph{prudencia}, considerada como virtude \emph{moral}, é o habito de uma sabia conducta. E' a forma de todas as virtudes \emph{moraes}. --- A \emph{justiça} é o habito de dar a cada um o que lhe pertence. --- \emph{A temperanca} é o habito de regular e ordenar o appetite \emph{sensiítivo} e as \emph{paixões}. --- A \emph{fortaleza} é o habito de luctar contra as difficuldades da vida e os obstaculos contrarios á virtude.\par Estas noções ácerca dos actos da \emph{vontade}, que muito em resumo temos aqui exposto, serão convenientemente desenvolvidas na \emph{Moral}.}.

=------ ------

\textbf{ORIGEM} E \textbf{DESTINO DO HOMEM}

213. \textbf{Divisão} da \textbf{secção} quarta. --- O conhecimento da \emph{essencia}, das \emph{faculdades} e das \emph{operações} de um ente, leva ao conhecimento da sua \emph{origem} e do seu \emph{destino}. Por isso, tendo considerado a \emph{essencia}, as \emph{faculdades} e as \emph{operações} do homem, resta-nos estudar a sua \emph{origem} e o seu \emph{destino}. Dividimos a secção \emph{quarta} em \emph{dois} capitulos. No \emph{primeiro} capitulo tratamos da \emph{origem} do homem; no \emph{segundo}, do seu \emph{destino}.

%%>>> cap-5

\chapter*{CAPITULO PRIMEIRO\\[0.3em]\large ORIGEM DO HOMEM}
\addcontentsline{toc}{chapter}{CAPITULO PRIMEIRO --- ORIGEM DO HOMEM}

{\small\noindent\textsc{Summario}: --- Origem do corpo humano. --- Origem da alma humana. --- Origem .do composto humano.\par}

%%>>> art-30

\section*{Artigo I.\\ Origem do corpo humano.}
\addcontentsline{toc}{section}{Artigo I. --- Origem do corpo humano.}

\pg{344} 214. O corpo humano \textbf{teve} uma origem. --- O \emph{corpo humano} não só não existiu sempre, mas é muito mais recente que os viventes inferiores. Teve, pois, uma origem.

E d'onde derivou o \emph{corpo humano}?

O \emph{corpo humano} não derivou do concurso de causas fortuitas, nem da evolução de uma especie inferior, mas, na sua origem primitiva, foi formado \emph{Immediatamente} por \textsc{Deus}, e, no estado actual, se propaga por \emph{geração} \footnote{A \emph{paleontologia} demonstra que o apparecimento do homem sobre a terra é posterior ao dos vegetaes e dos animaes. Os restos do homem apparecem nas camadas mais superficiaes.}.

215. O corpo \textbf{humano} não derivou do concurso de \textbf{causas} fortuitas. --- Um effeito, que revela ordem e finalidade, não póde derivar do concurso de causas \emph{fortuitas} ; pois as causas \emph{fortuitas}, na sua operação, não formam \emph{planos} a seguir, nem se propõem \emph{fins} a alcançar; aliás não seriam \emph{fortuitas}. Ora o \emph{corpo humano} \pg{345}  é um effeito, que revela ordem e finalidade; pois a \emph{anatomia} e a \emph{physiologia} ensinam que elle é um todo completo, resultando de diversos orgãos, harmonieamente dispostos e capazes de produzir as respectivas operações. Logo o \emph{corpo humano} não derivou do concurso de causas \emph{fortuitas} \footnote{A \emph{ordem} e a \emph{finalidade} não podem conceber-se sem uma intelligencia. O \emph{acaso}, sendo cego, não póde conhecer a relação entre os meios e o fim, --- conhecimento que é indispensavel para ordenar varios seres e escolher os meios proporcionados e efficazes.\par Consulte-se o que na \emph{Cosmologia geral} (cap. 11, a1t. 111) dissemos Acerca da \emph{finalidade} dos seres em geral e do \emph{organismo animal} em particular.}.

216. O corpo humano não derivou \textbf{da evolução de uma especie inferior} \footnote{O \emph{transformismo}, admittindo que a derivação de uma especie de outra é a lei reguladora de toda a natureza organica, não podia deixar de concluir que o homem deriva da evolução de uma especie inferior. \emph{Dórwin} diz; u A origem do homem não é differente da dos outros animaes; o homem, como os mammiferos, deriva de uma forma desconhecida e inferior 7.\par Se tivermos a curiosidade de averiguar quem é este nosso pretendido antepassado, o escriptor inglez nos ensinará que « foi um quadrupede piloso, fornecido de cauda e de orelhas ponteagudas, que vivia no antigo mundo, habitando provavelmente sobre as arvores ». --- \emph{J/aeckel} (\emph{La création naturelle}) faz uma descripção muito minuciosa do \emph{homo primigenius}, indicando a côr do cabello e da pelle, o comprimento dos braços e das pernas, a forma do craneo, e a transformação do grito animal em sons articulados.\par \emph{Büchner} é mais leal neste ponto. Não descreve o /\emph{omo primigenius}, porque o não conhece. Diz assim: « O nosso antepassado, isto é, o homem-macaco ou o macaco-homem é até hoje desconhecido: chegará um dia em que elle encontrar-se-ha nos tropicos, verdadeira patria dos grandes macacos anthropoides viventes, ou escavando na Africa, ou fazendo pesquizas nas ilhas do Archipelago da \emph{Malasia}. E, se nunca se encontrar, resignar-nos-hemos ao destino, lamentando a extraordinaria imperfeição dos indieios geologicos e as lacunas, que abrem na sciencia os logares submergidos » (\emph{O homem considerado segundo os resultados da sciencia}, pag. 2).}.

a) Se o corpo humano tivesse derivado de uma especie inferior, as \emph{especies} esteriam sujeitas á mutabilidade. Ora a razão e a experiencia provam que as \emph{especies} não estão, nem podem estar sujeitas á mutabilidade. Logo o corpo humano não derivou da evolução de uma especie inferior \footnote{\emph{Haeckel} e \emph{Huxley} afirmam que a descendencia do homem de uma serie de formas animaes inferiores é uma deducção legitima e necessaria do systema de \emph{Darwin}. Não contradizemos a affirmacao. Mas de nada vale que uma deduoção seja legitima, quando o principio é falso. Porisso, tendo demonstrado, na \emph{Cosmologia especial}, a falsidade do \emph{transformismo}, podemos concluir que é tambem falsa a hypothese da descendencia do homem de uma especie inferior, --- hypothese que está baseada na supposta lei da \emph{transformação} das especies,}.

\pg{346} b) Se o corpo humano tivesse derivado da evolução da especie inferior, que mais se approxima da especie humana e que é o \emph{simio}, deveriam encontrar-se, ao menos no estado de \emph{fosseis}, as formas intermedias entre o corpo do homem e o corpo do simio; pois a especie inferior, mudando-se para a superior, passaria por varias formas intermedias. Ora estas formas não se encontram, nem no estado de fosseis. Logo o corpo humano não derivou da evolução de uma especie inferior \footnote{Os \emph{darwinistas} dizem qué o homem, no periodo \emph{quaternario}, sé approximava manifestamente do typo \emph{s/miano}, e que por isso o homem \emph{quaternario} é uma das formas intermedias entre o macaco e o homem actual.\par Respondemos que as descobertas da \emph{paleontologia} são contrarias aos \emph{darwinistas}. Citaremos algumas. --- Em \emph{Cro-Mognon} (\emph{Franca}) foram encontrados varios craneos, e estes incontestavelmente pertencem a uma raça intelligente e dotada de \emph{grande cerebro}, como confessa o proprio \emph{Broca} (\emph{Les troglodytes de la Vésére} --- na \emph{Revue scientifique}, 16 nov. 1812, pag. 478). --- Em \emph{Sténgenas} (Suecia) encontraram-se craneos, que eram ainda mais antigos que os de \emph{Cro-Mognon}, cuja capacidade é \emph{sensivelmente superior d dos craneos actuaes}. (Cf. E.-T. Hamy, \emph{Précis de paleontologie humaine}, pag. 190, Paris, 1870). --- Em \emph{Mentone} (\emph{Italia}), no anno de 1872, foi encontrado um esqueleto humano bem conservado e tão antigo que remonta á \emph{edade da pedra}. Como póde ver-se no museu de \emph{Paris}, os seus caracteres são os seguintes: \emph{boa estatura}, \emph{cranêo sensivelmente desenvolvido} e \emph{bello angulo facial}. \emph{Riviére} diz a respeito do esqueleto de \emph{Mentone}: u O esqueleto, de que se trata, não apresenta \emph{nenhum} caracter que possa, \emph{em qualquer ponto}, approximalo do macaco; os craneos humanos, com que tem mais analogia, são os craneos encontrados em \emph{Cro}- \emph{Mognon} » (\emph{Comptes rendus de l'Académie des sciences}, 29 abril 1802).\par Os \emph{darwinistas} appellam para o famoso craneo de \emph{Veanderthal}, e dizem que ha uma grande semelhança entre os arcos superciliares d'esse craneo e a crista frontal do gorilla. --- Este recurso dos \emph{darwinistas} é inutil. O sabio \emph{Prüner-Bey} não só demonstrou que não ha nenhuma approximação possivel entre os arcos superciliares d'aquelle craneo e a crista frontal do gorilla, mas tambem reconheceu que o volume do \emph{cerebro} do homem de \emph{Neanderthal é superior ao volume medio do cerebro do homem moderno}, \emph{e que toda a superficie d'aquelle encephalo}, \emph{sem excepção alguma}, \emph{tem a conformação do typo humano}. (\emph{Congrés international d}' \emph{anthropologie}, \emph{tenu} á Paris, en 1861).\par O proprio \emph{Huxley} diz a este respeito: « Não devemos considerar os ossos encontrados no valle de \emph{Neander} como restos de um ser medio entre o macaco e o homem. Provam, quando muito, a existencia de um homem, cujo craneo teve uma leve semelhança com o dos macacos n.\par Muitos naturalistas sustentam a nossa these, dizendo que o homem \emph{quaternario} não se approximava do macaco, e por isso não é uma das formas intermedias entre o proprio macaco e o homem actual. Citaremos apenas os testemunhos de \emph{Quatrefages}, \emph{Virchow}, \emph{Joly} e \emph{Mortillet}.\par \emph{Quatrefages}, depois de ter comparado os mais antigos ossos humanos com o esqueleto do homem moderno, conelue d'este modo:\par « Dolichocephalo ou brachyocephalo, grande ou pequeno, orthognato ou prognato, o homem \emph{quaternario} é sempre homem na accepção completa da\par ,\pgn{347}palavra. Sempre que foi possivel um juizo ácerca dos seus restos, encontrou-se nelle o pó e á mão, que caraoterisam a nossa especie; a columna vertebral mostrou a duplice curvatura, a que \emph{Laurence} dava tão grande importancia e que \emph{Serres} considerava como o attributo do reino humano. Quanto mais se estudar, tanto mais estaremos certos de que cada esqueleto, grande ou pequeno, traz comsigo, na sua forma e proporções, o attestado da sua origem.\par \& Todos os ossos da cabeça humana moderna encontram-se nas cabeças fosseis e apresentam as mesmas relações. Quer se considerem isoladamente, quer no seu conjuncto, os ossos fosseis correspondem aos modernos. A enorme arcada supereiliar do homem de \emph{ANeanderthal} não póde disfarçar o caracter inteiramente humano d'este craneo excepcional.\par \& Em todas as raças fosseis encontra-se o caracter essencialmente humano: o craneo predomina sobre a face. Nellas, como em nós, a caixa craneana, destinada a encerrar o cerebro, conserva sempre uma capacidade egual a dos nossos dias.\par « Podemos, pois, com certeza applicar ao homem fossil, que conhecemos, as palavras de \emph{Huxley}. Nas epochas \emph{quaternarias}, como no periodo actual, u \emph{nenhum ser intermediario preenche a lacuna}, \emph{que separa o homem do troglodyte} ». E' impossivel negar a existencia d'este abysmo \emph{n}. (\emph{L'espéce humaine}).\par \emph{Wirchow}, num discurso dirigido aos naturalistas allemaes, reunidos no Congresso de Munich, em 22 de setembro de 1818, disse: « Tenho-me dedicado de um modo especial ao estudo da \emph{Antropologia}; mas devo declarar que todo o adiantamento positivo, que temos feito na Anthropologia prehistorica, nos tornou menos possivel o poder provar o nexo do homem com algum animal inferior. A Anthropologia estuda presentemente a questão Acerca dos homens fosseis... Pois bem, se estudarmos os restos certos e mais antigos dos homens fosseis, que na serie de descendencia ou de ascendencia devem approximar-se mais dos primitivos antepassados, encontramos invariavelmente um \emph{homem}, tal qual existe hoje. No fim de tudo, \emph{devemos reconhecer que falta o typo fossil de um desenvolvimento inferior}. ... Até agora tem sido impossivel encontrar uma transição immediata entre o homem e o macaco, porque semelhante successão realmente não existe 7.\par \emph{Joly}, livre pensador e patrono da \emph{geração espontanea}, não hesita em dizer o seguinte: « Até \emph{hoje}, \emph{o estudo dos craneos e dos outros restos humanos prehistoricos não nos auctoriza a pensar}, \emph{e muito menos a sustentar que o homem primitivo}, \emph{e por conseguinte o homem actual}, \emph{tenha uma origem simiana} n (\emph{Revue scientifique} du 18 janvier 1879).\par \emph{Mortillet}, defensor ardente do \emph{transformismo}, depois de enumerar todos os ossos fosseis de macacos, que se poderam encontrar nos dois continentes, diz: « Poderemos nós entre estes restos reconhecer alguns vestigios de um \emph{anthropopitheco}, de um precursor do homem? Não o creio » (\emph{Le Préhistorique}).\par Do mesmo modo se exprimem outros sabios naturalistas insuspeitos, como Hamy (\emph{Précis de Paléontologie hwmaine}), --- Nadaillae (\emph{Les premiers hommes et les temps prélistoriques}), --- o mesmo Virchow (\emph{ArcAéologie celtique et gauloise}) --- Le Bon (na \emph{Revue scientifique}, 17 déc. 1881).\par De resto o proprio \emph{Darwin} confessa que se não encontraram intermediarios\pgn{348}, nem sequer no estado de fosseis, entre o homem e o macaco. Mas, se \emph{Darwin} não encontrou esses intermediarios, como póde dizer que existem ? Não é este o methodo \emph{positivo}, que os \emph{transformistas} querem seguir.\par \emph{Lyel}, sequaz do \emph{darwinismo}, espera encontrar \emph{centenas de quadrumanos fosseis}, para estabelecer a continuidade entre o homem e o Macaco, mas no entretanto pede que \emph{não nos deixemos impressionar}, \emph{em nossos Juizos}, \emph{por esta falta de provas} (\emph{L'anciennetê de Phomme}, pag. 550-551). --- Respondemos a \emph{Lye}! que elle, com o seu pedido, se colloca fóra do verdadeiro methodo scientifico; pois é antiscientifico enunciar uma theoria sem poder adduzir as provas. A falta d'aquellas provas é um elemento indispensayel na apreciação do \emph{transformismo}. Pela nossa vez pedimos aos \emph{darwinistas}, que, antes de apregoarem que entre o homem e o macaco existem formas intermediarias, encontrem as mesmas formas. Uma theoria de facto só com os factos póde e deve ser demonstrada.}.

c) De dois seres, caracterisados por differenças profundas e radieaes, um não pode derivar da evolução de outro; pois que todo o ente é \emph{substancialmente} semelhante a outro, de que deriva \pg{347} \pg{348} . Ora \textbf{entre} o corpo \emph{humano} e o do animal, que mais se approxima do homem e que é o \emph{simo}, existem differencas radicaes e profundas, tanto \emph{an«tomicas} como \emph{physiologicas}. Logo o corpo humano não derivou da evolução de uma especie inferior \footnote{Os \emph{transformistas}, fundados no principio de que seres substancialmente semelhantes entre si devem ter uma origem commum, empregam todos os esforcos para demonstrar que entre o bruto e o homem ha uma grandissima e completa semelhança e d'ahi concluir que um e outro tiveram uma e a mesma origem. Se devessemos acreditar nas suas palavras, não haveria sciencia natural, que n\&o subministrasse argumentos em favor d'aquella profunda semelhança; a anatomia, a physiologia, a embryologia, apresentariam factos mais que sufficientes para provar a sua theoria.\par Não só aeceitamos, mas até defendemos o principio em que se funda o raciocinio dos transformistas --- que dois seres, substancialmente semelhantes entre si, têm uma origem commum; porque, se elles são substancialmente semelhantes entre si, pertencem á mesma especie, e, como taes, procedem necessariamente da mesma origem.\par Mas deveremos acceitar egualmente a applicacio que fazem d'aquelle principio? Deveremos por ventura admittir que o homem e o bruto sao perfeitamente semelhantes nos seus elementos substanciaes, e não differem senão nas condições accidentaes? Nào; e vamos demonstral-o.\par Neste artigo mostraremos as divergencias, que existem entre o homem e o animal, considerados no seu \emph{organismo}, e, no seguinte, notaremos as differencas entre os dois seres, em relação ás faculdades \emph{perceptivas}.\par O animal que mais se approxima do homem, \emph{em relação ao organismo}, é o \emph{simio} ou \emph{macaco}; €, entre os macacos, os mais semelhantes ao homem são os \emph{anthropoides} ou \emph{anthropomorphos}. Contam-se quatro especies de macacos \emph{anthropoides}: o \emph{gibbon} e o \emph{orangolango} na Asia orienta], o \emph{chimpanzé} e o \emph{gorilla} na Africa occidental.\par As nossas investigações, portanto, limitam-se ao homens e aos \emph{antlropoides}.\par Antes de tudo reconhecemos que entre o \emph{organismo} do homem e o do macaco, considerados na sua \emph{estructura geral}, existe uma analogia incontestavel. Tém os mesmos orgãos, os mesmo apparelhos, o mesmo systema nervoso, com os seus annexos e relações.\par Mas esta analogia do \emph{organismo}, que approxima o homem do macaco,\pgn{349}encontra-se em todos os mammiferos, \emph{e não prova que o homem deriva do Macaco}, \emph{assim como não prova que os mammiferos derivam uns dos outros}. Para concluir que o homem deriva do macaco, os \emph{dariwinistas} deveriam provar que o homem é \emph{semelhante} ao macaco, não só na \emph{estructura geral do organismo}, mas tambem nas particularidades \emph{especificas} do mesmo organismo. Mas isto é que nunca poderá ser provado.\par Aquellas analogias \emph{geraes}, que existem entre o homem e o Macaco, são largamente compensadas e excedidas por divergencias tão caracteristicas e profundas, que, assim como necessariamente nos determinam a reconhecer nelles uma diversidade de especie, assim tambem nos levam a attribuir-lhes uma diversidade de origem.\par Examinando \emph{pAysiologicamente} o homem e O macaco, vemos que entre elles existem profundas divergencias principalmente com relação á \emph{posição}, ao \emph{craneo}, ao \emph{cerebro}, \emph{ás extremidades}, ao \emph{pello} e ás \emph{forças}.\par a) O /\emph{liomem}, e só elle, caminha: a cabeça direita e levantada, a conformação do esqueleto, a acção regular dos musculos, o peso equilibrado das visceras, --- tudo demonstra que a nossa \emph{posição é vertical}, e que a natureza nos formou para estarmos direitos. --- O \emph{macaco}, pelo contrario, é trepador; \& sua cabeça inclinada para a terra, a conformação do organismo, a disposição dos ossos, a funcção dos musculos, --- tudo isto indica que á sua \emph{posição é horizontal}, e que a \emph{vertical} é para elle violenta e nunca permanente.\par Esta differença de \emph{posição} no homem e no macaco faz com que elles sejam dois typos perfeitamente distinctos um do outro.\par Ouçamos a este respeito alguns dos melhores naturalistas.\par \emph{Quatrefages} (\emph{Rapport sur les progrês de l'anthropologie}, pag. 244, Paris 1867) diz: « O. homem \emph{caminha}, e caminha sobre os membros posteriores; todos os macacos \emph{trepam}. Nos dois grupos, todo o apparelho locomotor tem a marca d'estes differentes destinos; \emph{os dois typos são perfeitamente distinctos} n. --- \emph{Duvernoy}, \emph{Gratiolet} e \emph{Alis}, citados por \emph{Quatrefages} (obr. cit, pag. 245), affirmam que « 0 \emph{typo-m«acaco}, \emph{aperfeicoando-se}, \emph{ndo perde o seu caracter fundamental e fica sempre perfeitamente distincto do typo humano} s. --- \emph{Godron} (\emph{De l'espéce et des races}) diz o seguinte: « Entre todos os seres da creacüo, só o homem está organisado para a posição \emph{vertical}, só elle caminha naturalmente direito; é este um caracter \emph{essencial}, \emph{que o separa de todos os animaes} .\par Se, pois, \emph{o} homem e o macaco são typos perfeitamente \emph{distinctos}, não podem ter uma origem commum, ou não podem derivar um do outro.\par \emph{Darwin} responde que os caracteres, relativos á posição \emph{vertical} do homem, são o simples resultado da \emph{adaptação} e por isso insufficientes para separar o homem do macaco (\emph{Origem do homem}, vol. l, pag. 790).\par A resposta de \emph{Darwin} não só não destroe a força do nosso argumento, mas é absurda. --- Não destroe a força do nosso argumento. Por quanto, aquelles caracteres, embora fossem o resultado de uma \emph{adaptação} a um novo modo de existir, comtudo nao deixariam de estabelecer entre o homem eo macaco uma \emph{differenca profunda} e \emph{essencial}. --- E' adsurda. Com effeito,\pgn{350}aquella \emph{adaptação}, não podendo conceber-se sem uma \emph{finalidade}, exige uma \emph{causa} dotada de \emph{intelligencia}. Ora, no systema de \emph{Darwin}, a \emph{adaptação} não tem a sua razão sufficiente numa causa dotada de \emph{intelligencia}; pois que é effeito da \emph{selecção}, e esta é o resultado \emph{cego} e \emph{fatal} dos factos.\par 5b) O craneo do homem distingue-se do craneo do macaco tanto na \emph{capacidade}, como na \emph{conformação}.\par \emph{a}) Distingue-se na \emph{capacidade}. Por quanto, o \emph{craneo} do homem tem uma \emph{capacidade}, que é tres vezes superior á do \emph{craneo} do macaco. \emph{Godron}, seguindo o methodo de \emph{Tiedemann}, encheu de areia o \emph{craneo} de um \emph{Europeu} e o de um \emph{orangotango}, e viu que aquelle continha litros 1,58, ao passo que este continha apenas litros 0,44. (\emph{De Vespêce}, t. IT. --- Note-se que a \emph{capacidade} do craneo é quasi a mesma em todas as raças humanas.\par \emph{b}) Distingue-se na \emph{conformação}. Na verdade, o \emph{craneo} do homem é \emph{vertical} na sua \emph{face}. e, se é \emph{obliquo} em certas raças, especialmente na \emph{negra}, esta \emph{obliquidade} é notavelmente menos sensivel que nos \emph{macacos}. Tal conformação aprecia-se pelo \emph{angulo facial}. Por quanto, o \emph{angulo facial}, no \emph{homem}, está comprehendido entre 70º e 85º; ao passo que nos macacos anthropoides não passa de 40º. --- Exceptua-se o \emph{gibbon}, cujo \emph{angulo facial} póde attingir 60"; mas é este animal que, entre todos os macacos, sob outros aspectos, mais diverge do homem.\par Oucamos o que a este respeito dizem \emph{Aeby}, \emph{Bichoff}, \emph{Huxley}.\par \emph{Car}. \emph{Aeby}, naturalista de \emph{Berna}, no seu livro --- \emph{Les formes du cráne de l'homme et des singes}, depois de apresentar medidas e comparações, sob todos os aspectos possiveis, dos craneos de todas as raças humanas e dos craneos dos macacos e dos mammiferos, escreve o seguinte: « D'estas comparações, que fizemos, resulta que a \emph{distancia}, \emph{que separa o homem do macaco mais perfeito}, \emph{é mais consideravel que a que separa os macacos}, \emph{uns dos outros}; e por conseguinte não hesitamos um instante em sustentar que o typo do craneo humano se distingue inteiramente do typo simiano, e que especialmente os macacos \emph{antliropomorphos se unem}, \emph{sob todos os aspectos}, \emph{com os seus alliados naturaes e mesmo com os mamnmiferos inferiores}, \emph{por um modo incomparavelmente mais estreito do que com o homem} » (pag. \emph{V}1).\par \emph{Bichoff}, depois de ter tambem estudado os craneos do homem e do macaco, conelue que \emph{existe uma falta absoluta de factos para estabelecer}, \emph{ou mesmo para explicar} a \emph{passagem do macaco para o homem}.\par \emph{Huxley} diz: « As differencas entre o craneo d'um homem e o d'um gorilla são enormes; mesmo as que existem entre o homem e o macaco da ordem mais elevada são consideraveis; cada osso particular do gorilla apresenta signaes, que o distinguem do osso correspondente do corpo humano ».\par Os \emph{darwinistas} respondem que o craneo de um macaco novo se approxima muito do craneo do homem.\par A resposta dos \emph{darwinistas} é inutil. Para a comparação ser producente, é necessario que seja feita entre os craneos de um homem \emph{adulto} e de um macaco \emph{adulto}. Se assim se fizer, o resultado é muito diverso. Diz o citado\pgn{351}\emph{Aeby} (pag. 87): « Não póde negar-se que na edade juvenil ha uma ligeira approximacüo de typos; mas esta approximação não destroe a proposição estabelecida com relação á \emph{edade adulta}: que o \emph{craneo humano á inteiramente diverso do craneo simiano}. \emph{O intervallo entre o homem e o macaco é sempre incomparavelmente maior que o que separa este ultimo do resto dos animaes} n.\par 6) O \emph{cerebro} do homem differe sensivelmente do cerebro do macaco, tanto em relação ao \emph{volume}, como em relação ao \emph{desenvolvimento}.\par \emph{a}) O \emph{volume} do cerebro humano é, relativamente á massa do corpo, muito mais consideravel que o volume do cerebro dos mammiferos, e quasi tres vezes mais que o dos anthropoides.\par \emph{Huxley}, sem negar o facto, pretende que a mesma differenca, existente entre o cerebro do homem e o dos animaes, existe entre os cerebros das diversas racas humanas.\par A esta difficuldade responde um transformista allemão, \emph{Schaafhausen}, com as seguintes palavras: « A assercüo, formulada por \emph{Huxley}, de que os homens, mesmo no peso do cerebro, differem mais entre si, do que dos macacos, é erronea.... O cerebro do \emph{Australiano} excede, em volume, duas ou ires vezes o do gorilla, ao passo que o cerebro regularmente desenvolvido de um \emph{Europeu} não excede senão de um quinto o do mesmo \emph{Australiano} n (Cit. por \emph{Hamard}, nas \emph{Quest}. \emph{scient}. de \emph{Bruxelles}, julho 1818).\par \emph{b}) O cerebro do homem \emph{desenvolve-se} d'uma maneira opposta Aquella por que \emph{se desenvolve} o cerebro do macaco. Por quanto, no homem desenvolve-se e aperfeicóa-se primeiro o lóbo anterior e depois o lateral; no macaco, porém, forma-se primeiramente o lóbo lateral e depois o anterior. Ora é impossivel, como observa \emph{Qauatrefages}, que sejam completamente semelhantes, e que derivem um de outro, dois seres que se formam d'um modo contrario, porque a historia natural reconheceu a seguinte lei permanente, que não admitte excepção --- que o semelhante se desenvolve de um modo semelhante. (\emph{Conf}. \emph{pop}. \emph{acerca} da \emph{origem do homem}).\par 4) O \emph{homem é bimano}; as suas mãos regulares e perfeitas são orgão do tacto, verdadeiros compassos naturaes para medir, admiraveis instrumentos para manifestar os pensamentos da intelligencia e para communicar as disposições da vontade; e os pés sustentam todo o peso do corpo. --- O macaco é \emph{quadrupede}; os seus pés anteriores, que têm forma de garras, são propriamente orgão de suspensão, e os posteriores não podem sustentar o peso do corpo senão por breve tempo e com grande difficuldade. Diz \emph{Quatrefages}: « O exame do organismo em geral e das extremidades em particular mostra as differencas essenciaes, que existem entre a forma e a disposição do homem e á do bruto; o que torna inadmissivel a opinião de que um derive de outro v (\emph{Rapport}, etc., pag. 247).\par €) Os macacos, e quasi todos os mammiferos, têm o corpo, principalmente no dorso, coberto de pello, o que não acontece em relação ao homem.\par Este facto não só constitue uma enorme distancia entre o homem e o macaco, mas é absolutamente inexplicavel no systema \emph{darwinista}. Na verdade, sob o aspecto physico, a ausencia de pello, para proteger o corpo contra o frio e a chuya, constitue o homem, relativamente ào macaco, num estado\pgn{352}de inferioridade, que não póde explicar-se pela \emph{selecção natural}. Por quanto, segundo \emph{Darwin}, \& especie de macacos, que perdeu o pello, devia achar-se inferior ás outras, e, como tal, não podia ser aproveitada pela \emph{selecção}, mas devia desapparecor. Todavia não foi assim. Não só não desappareceu, mas na sua evolução produziu uma especie mais perfeita --- o homem.\par Os \emph{darwinistas} dizem: O pello desappareceu, porque se tornou inutil.\par Respondemos que o pello, se era inutil para alguns macacos e para os homens que vivem em climas quentes, não era inutil para os que vivem em climas frigidissimos. Por isso a razão do desapparecimento do pello não póde ser a que os adversarios apontam.\par Insistem: Os vestidos fizeram desapparecer o pello do homem.\par Os \emph{darwinistas} tomam o effeito pela causa. O pello não desappareceu, porque usamos de vestido; mas usamos de vestido, porque o pello não existe. --- Além d'isso, se fosse verdade o que elles dizem, a cabeça, que está habitualmente coberta, não devia ter cabello.\par f) O corpo humano, considerado physicamente, é mais fraco que o corpo dos animaes. Ora o homem, se fosse uma evolução mais aperfeiçoada dos \&nimaes, não seria mais fraco do que estes; pois que a evolução far-se-hia por obra da \emph{selecção}, e, esta, segundo \emph{Darwis}, não só não diminue, mas augmenta a perfeição recebida das especies passadas.\par Os naturalistas reconhecem muitas outras dessemelhanças entre o homem e o macaco, --- dessemelhanças na distribuição dos dentes, na disposição das entranhas, no comprimento das mãos e dos pés, --- em todas as partes, tanto nas mais notaveis como nas mais pequenas de seus organismos, que nos levam a negar-lhes uma descendencia commum.\par Escutemos o resumo, que os melhores naturalistas fazem das differenças entre e macaco e o homem. Citaremos \emph{Priüner-BDey}, \emph{Bermeister} e \emph{Luci}.\par \emph{Prüner-Bey}, na sessão da Sociedade de Anthropologia, em 18 de novembro de 1569, pronunciou as seguintes palavras: « São as differenças que impressionam o espirito. --- O macaco está coberto de um vestido de pello, que falta ao homem. E' este um caracter anatomico importantissimo, que compelle o homem a supprir pela industria o que a natureza lhe recusa. --- O macaco tem um dente canino, que lhe serve de arma; o homem não o possue, e, para a sua defeza, teme de empregar armas aperfeiçoadas. --- A' medida que toda a construcção ossea do macaco denota um animal quadrupede e trepador; a do homem mostra-o adaptado para caminhar. --- O systema muscular offerece os mesmos caracteres. \emph{Gratiolet} demonstrou que, em relação á circulação do sangue, o homem, comparado com o macaco, é um ser arterial. --- O estado das visceras em alguns macacos, como no gorilla e no chimpanzé, revela um animal herbivoro perfeitamente caracterisado. --- O estudo dos craneos é tambem concludente. Todos os simios apresentam a mesma face e egualmente distineta da humana. No homem o lóbo frontal predomina sobre o occipital ; no macaco nota-se o opposto. --- O systema dentario do macaco revela um herbivoro, armado para a sua defeza; o homem é omnivoro e só a sua industria lhe subministra armas. Não ha transição possivel entre o macaco e o homem, a menos que as leis do desenvolvimento não sejam imaginarias ».\par \emph{Burmeister} escreve: « O homem distingue-se do macaco, emquanto á\pgn{353}estructura do corpo, por um desenvolvimento maior do cerebro, por um esqueleto destinado a facilitar o andar em posição vertical, por um desenvolvimento maior do pelvis e por uma differenca radical na disposição das duas extremidades; porque, no homem, só a extremidade anterior é uma verdadeira mão, a extremidade posterior, não; ao passo que no macaco é tudo o contrario, as extremidades posteriores são verdadeiras mãos, e as anteriores assemelham-se antes a pés ».\par \emph{Luci} diz o seguinte: « Em uma grande serie de macacos, do mundo antigo e moderno, notei que todos os ossos das extremidades anteriores e posteriores não confirmam a opinião de \emph{Huxley}, quando disse que as differencas anatomieas entre o homem e o gorilla são menos sensiveis do que entre o gorilla e os outros macacos inferiores » (\emph{Archivo de Anthrop}., v. vi, pag. 16).\par Depois d'estes testemunhos da sciencia e dos melhores naturalistas ácerca da differenca entre o homem e o macaco, vê-se a falsidade da proposição de \emph{Huxley} (\emph{Demonstration}. \emph{de la place de l'homme dans la nature}), sustentada por \emph{Darwin} (\emph{Origem do homem}, pag. 191), que \emph{ua differenga entre o homem e os macacos superiores é menor que a que existe entre os macacos superiores e os inferiores do mesmo grupo} n. 1\par *\par Examinemos agora os principaes argumentos dos \emph{transformistas}.\par a) Dizem: A \emph{anatomia} ensina que o homem é formado segundo o mesmo plano e modelo do macaco. Os ossos do esqueleto, os musculos, nervos, veias, entranhas do homem, são semelhantes aos correspondentes ossos do esqueleto, aos musculos, nervos, veias, entranhas do macaco.\par Concedemos que a estructura \emph{geral}! do corpo do homem tem uma certa affinidade com a do corpo do macaco. Mas que revelará esta \emph{generica} semelhança entre o homem e o macaco? Manifestará que estes seres têm uma origem commum? Não; pois, se assim fosse, todos os mammiferos deveriam derivar uns dos outros, visto que todos convêm na estructura \emph{geral} do organismo. Para que o homem e o macaco tivessem uma origem commum, era necessario que convissem entre si, não só na estructura \emph{generica}, mas tambem na \emph{especifica}; pois a razão e a experiencia demonstram que dois seres têm uma origem commum, quando pertencentes á mesma \emph{especie}. Ora entre o organismo do homem e o do macaco existem divergencias radicaes e profundas, que fazem com que sejam dois typos \emph{especificamente} diversos, e que por isso um não possa derivar e não derive do outro.\par \emph{Aeby} diz: « Examinando o macaco e o homem, deprehendemos não só que o plano fundamental lhes é commum com todos os vertebrados, mas tambem \emph{que sobre este plano estão levantados edificios completamente differentes}. Na verdade, a sua conformação não concorda senão raramente, mesmo num ponto isolado; \emph{o accordo é quasi sempre appavente}; \emph{em relação} ao \emph{todo}, \emph{não têm nada de commum entre si}. \emph{Em toda} a \emph{serie dos mammiferos não se encontra uma distancia}, \emph{que possa comparar-se}, \emph{nem mesmo de longe}, \emph{com a que separa o homem do macaco} » (pag. 82).\par pb) Continuam: A \emph{physiologia} demonstra --- que o bruto, á semelhança do homem, vê com os olhos, ouve com os ouvidos, gosta com a lingua, ca-\par 23 Vor. LI.\pgn{354}minha com os pés, --- recebe do homem e lhe communica a hydrophobia e a variola --- que os macacos tambem padecem apoplexias, ophtalmias, inflammações intestinaes, gostam de chá, café e bebidas espirituosas, --- que o animal nasce como o homem, e finalmente morre como elle.\par Este argumento refuta-se como o precedente. \emph{O} homem é \emph{animal}, como o é o macaco, e por isso deve existir entre elles uma certa affinidade com relação á vida vegetativa e sensitiva. Mas esta affinidade é apenas \emph{generica}, porque \emph{animal é genero}, e, se manifesta o nexo e a harmonia existente entre os seres creados, não demonstra que um descenda do outro.\par o) Raplicam: A \emph{embryologia} põe em evidencia que o principio e o desenvolvimento successivo do embryao do homem são identicos ao principio e desenvolvimento do dos outros vertebrados, que na escala zoologiea estão no grau immediatamente inferior ao homem.\par Respondemos que o embryão do homem e o do macaco, ainda que no seu começo e até a um certo ponto do desenvolvimento apresentem uma \emph{semelhança generica}, todavia não são \emph{realmente identicos}. Porquanto, se fossem \emph{identicos} na \emph{realidade}, se entre elles não existisse uma differenca \emph{especifica}, não haveria nenhuma repugnaneia em admittir que do embryão de um cão se formasse um homem, e do embryão \emph{do} homem sahisse um gato. Quem poderá sustentar semelhante extravagancia? Se, portanto, os embryóes do homem e do macaco divergem no seu termo, devem egualmente divergir no seu começo e desenvolvimento e conter em si, ao menos \emph{virtualmente}, as differencas especificas, que tão claramente apparecem no resultado.\par Esta objecg\&o dos \emph{transformistas} baséa-se na theoria da \emph{unidade de typo}, segundo a qual o embryao de uma especie superior percorreria, no seu desenvolvimento, todas as formas dos animaes inferiores; de modo que as especies inferiores seriam differentes pausas do mesmo embryo.\par A theoria da \emph{unidade de typo} é falsa. Porquanto, embora a nossa intelligencia, considerando o que as especies têm de \emph{commum}, forme o typo de \emph{animal}, todavia esse typo é apenas \emph{logico} e \emph{subjectivo} e não é \emph{real} ou \emph{objectivo}, pois que, na \emph{realidade}, as especies são \emph{essencialmente} diversas entre si. \emph{Müller} escreve: « Sustentaya-se ha pouco tempo, com muita seriedade, que o feto humano, antes de chegar ao estado perfeito, percorre successivamente os diversos graus de evolução, nos quaes permanecem fixos, durante toda a vida, os animaes das classes inferiores. \emph{Esta hypothese ndo tem o minimo fundamento}, como demonstrou \emph{Büer}. O embryão humano nunca é semelhante a um radiado, a um mollusco, a um verme. O \emph{desenho} da formação d'estes animaes é \emph{inteiramente differente} do dos vertebrados. O homem poderia, quando muito, ser semelhante a estes ultimos; porque elle tambem é vertebrado, e o seu organismo é constituido (\emph{nos caracteres mais geraes}) segundo o typo commum a esta grande divisão do reino animal. Mas elle nunca é semelhante a um peixe, a um reptil, \emph{a} uma aye, etc. A analogia, existente entre o homem e um reptil, não excede a que existé entre todos os animaes vertebrados » (\emph{Manual de physiologia}, t. III, 7128). --- Com \emph{Miller} concordam outros sabios, como Flourens (\emph{Histoire des travaux de Cuvier}, p. 274), Milne Edwards (\emph{Lecons sur la phys}. \emph{et l'anat}. \emph{comp}., t. I, pag. 38).\par Os \emph{darwinistas} adduzem em defeza da sua opinião outros argumentos,\pgn{355}tirados da existencia dos \emph{orgãos rudimentares}, do \emph{atavismo}, --- argumentos que já foram expostos e refutados na \emph{Cosmologia especial}.\par Apezar de todos os desmentidos com que a razão e a experiencia condemuam o \emph{transformismo}, os materialistas não se cançam de apregoar esse systema como o unico e legitimo resultado da sciencia.\par \emph{Haechel} diz: « Os nossos antepassados, especialmente o velho \emph{AwmphAioxo}, merecem mais respeito e consideração do que toda a cafila dos santos inuteis, a quem as nossas nações altamente civilisadas dedicam templos e procissões. Aquelles realmente são a carne. da nossa carne, o sangue do nosso sangue.... O veneravel ancião \emph{Amphioxo}, ainda sem craneo, encontra-se na raia do reino dos vertebrados. Os mammiferos primitivos vão subindo successivamente até ao \emph{Homo Sapiens}. Todas as bibliothecas da sciencia ecclesiastica e da falsa philosophia nada valem; nós illuminamos esses erros com o sol da historia do desenvolvimento animal .\par O fanatismo dos materialistas não podia ser mais ridiculo!}.

\pg{349} 217. 0 \textbf{corpo humano}, na sua \textbf{origem primitiva}, foi \textbf{formado immediatamente} por \textbf{Deus}.

a) O corpo humano, na sua origem primitiva, ou derivou do concurso de causas fortuitas, --- ou foi o resultado da evolução \pg{350}  de uma especie inferior, --- ou foi formado \emph{immediatamente} por \textsc{Deus}. A disjuneção é completa. Ora o corpo humano, na sua origem primitiva, não derivou do concurso de causas \emph{fortuitas}, nem foi o resultado da evolução de uma especie inferior. Logo foi formado \emph{immediatamente} por \textsc{Deus}.

b) O corpo do \emph{primeiro} homem, revelando em si perfeição, ordem e finalidade, foi formado immediatamente por uma causa \emph{intelligente}; aliás o effeito seria superior á causa, ou dar-se-hia \pg{351} \pg{352}  um effeito sem causa. Esta causa intelligente não póde ser finita; porque uma causa finita, sendo um ente \emph{particular}, não \pg{353}  póde formar uma especie, que é um ente \emph{commum}, aliás o effeito seria mais nobre do que a sua causa. Logo deve ser \pg{354}  a causa infinita, que é \textsc{Deus}; porque só Elle, por ser o principio \emph{universalissimo}, pode produzir entes \emph{universaes} e \emph{communs}, \pg{355}  como são as \emph{especies}. Logo o corpo humano, na sua origem primitiva, foi formado \emph{immediatamente} por \textsc{Deus} \footnote{Quando dizemos que o corpo humano, na sua origem primitiva, foi formado \emph{immediatamente} por \textsc{Deus}, queremos entender que, para a formação do corpo do primeiro homem, não bastou que Dzvus tivesse creado a materia e as forças da materia, mas que foi necessaria uma acção \emph{directa} e \emph{especial}, distineta do \emph{concurso} que Elle, como Causa primeira, presta aos actos das causas segundas, pela qual formou da materia preexistente o corpo humano.\par Por isso a formação do corpo do \emph{primeiro} homem não foi uma \emph{creação}; pois a \emph{creação}, sendo a produccáo do \emph{nada}, não presuppõe materia alguma. --- Todavia a formação do corpo do primeiro homem, embora não fosse uma creacüo, foi comtudo uma operação só propria de \textsc{Deus}, que não podia ser exercida pelas forças naturaes das causas segundas. Por quanto, as causas segundas, se podem produzir um ser de uma especie, quando esta especie já existe, não a podem produzir, quando a especie não está ainda formada. Ora não existindo, na primeira formação das coisas, um corpo humano, de que podesse derivar, por via de geração, um outro semelhante na especie, foi necessario que o corpo do primeiro homem fosse formado immediatamente por Deus (\emph{Sum}. \emph{TA}., p. I, q. 91, a. 2). A razão d'isso encontramol-a nas seguintes palavras de S. Thomaz: « A geração natural de toda e qualquer especie faz-se por meio de uma determinada materia. Ora a materia, de que é gerado naturalmente o homem, é o semen humano. Portanto um individuo da especie humana não pode ser naturalmente gerado de qualquer outra materia. E, como só Deus, que é auctor da natureza, pode produzir as coisas fóra da ordem da propria natureza, só Deus podia formar o corpo do primeiro homem » (1\emph{b}., q. 92, a. 4).\par Nem se diga que o corpo humano, se tivesse sido formado immediatamente por Deus, não poderia estar sujeito a tantas imperfeições, sendo Deus um agente dotado de infinita sabedoria e de omnipotencia. O corpo humano é perfeito, não \emph{absoluta} mas \emph{relativamente}, isto é, é perfeito com relação ao fim \emph{prosimo}, que Deus lhe determinou e que consiste em servir ás operações da alma racional. Considerado sob este ponto de vista, o corpo humano é perfeito, E se elle é corruptivel e mortal, essa disposição é uma consequencia da materia, de que é composto. Deus podia formar um corpo \emph{naturalmente}\par *\pgn{356}incorruptivel e immortal: mas devemos advertir que, na formação das coisas, o criterio não é o que Deus pode fazer, mas é o que convem á natureza das proprias coisas. --- Dizemos \emph{naturalmente}; porque \emph{sobrenaturalmente}, isto ê, por uma graça não devida ao homem, Deus tinha dado o remedio efficaz contra as doenças e a morte. Veio o peccado e, com o peccado, o doença, a morte. Mas a doença e a morte, se são penas do peccado, são tambem disposições e consequencias da natureza humana, isto é, que naturalmente derivam dos principios constitutivos do corpo humano.}.

\pg{356} 218..0 corpo humano, no estado actual, propaga-se por \textbf{geração}. --- Um ser vivo propaga-se por \emph{geração}, quando por tal forma deriva de um outro ser vivo, que ambos convêm na mesma natureza ou especie. Ora a experiencia attesta que, no estado actual, o corpo do homem deriva de outro homem. Logo o corpo humano, no estado actual, propaga-se por \emph{geração}.

\textsc{Deus}, quando formou \emph{immediatamente} o corpo do primeiro homem, deu-lhe a virtude ou o poder de transmittir a outros a sua semelhança e assim perpetuar a propria vida \footnote{Todos os tratados de \emph{Astoria natural} apresentam noções ácerca da faculdade, que os seres vivos possuem, de reproduzir e multiplicar a propria substancia. --- S. Thomaz nota a differenca entre a formação do corpo primeiro homem, e a formação do corpo dos seus descendentes. 'l'onde-se proposto a difficuldade que, assim como os nossos primeiros paes foram produzidos immediatamente por Deus, assim tambem deviam ser produzidos os seus descendentes, porque a identidade da especie exige a identidade do modo da producção, o Angelico Doutor responde e diz que nas primeiras producções se tratava de instituir a natureza, e por isso aquelles produeções deviam derivar immediatamente de Deus, mas, depois da instituição da natureza, esta podia produzir os seus effeitos por um operação natural. Eis as palavras do S. Doutor: « In primis operibus natura instituenda erat, et ideo oportebat illa opera immediata a principio supernaturali esse: sed postmodum natura instituta in proprios effectus per naturalem operationem pervenire poterat, et hoc de perfectione ipsius erat » (2 \emph{Dist}. XX; q. I, a. 1 ad 4).}.

%%>>> art-31

\section*{Artigo II.\\ Origem da alma Irunaaaana,.}
\addcontentsline{toc}{section}{Artigo II. --- Origem da alma Irunaaaana,.}

%%>>> t219-2

\tese{219}{A alma humana tem uma origem.}
Todo o ente \emph{contingente} tem uma origem; pois que só o ente \emph{necessario} é eterno. Ora a alma humana é um ser \emph{contingente}; porque não só não existiu sempre, mas tambem podia deixar de existir. Logo a alma humana tem uma origem.

Neste artigo provamos que a alma não é produzida pela \emph{evolução de uma especie inferior}, nem por \emph{geração}, nem por \pg{357}  \emph{emanação da substancia divina}, mas que é \emph{creada} immediatamente por \textsc{Deus}, no momento da sua união com o corpo \footnote{Sao \emph{quatro}, pois, os principaes systemas ácerca da \emph{origem} da alma humana: o \emph{evolucionismo}, o \emph{generacianismo}, o \emph{emanatismo} e o \emph{creacianismo}.\par A questão ácerca da \emph{origem} da \emph{alma humana} é duplice; por quanto póde referir-se --- ou á alma do \emph{primeiro} homens, ---- ou á alma dos homens que vão apparecendo continuamente sobre a terra. Todavia tratamos estas duas questões, como se fossem uma só; pois que ambas se resolvem com o mesmo principio, baseado na propria natureza da alma.\par D'esta questão trata S. Thomaz --- na \emph{Sum}. \emph{Th}., p. I, q. 90, a. 1, 2; q. 118, a. 9, --- nas q. d. \emph{De pot}., q. 9, a. 9; --- no 2 \emph{Sent}. dist. 17, 18; --- na \emph{Sum}. c. \emph{Gentes}, l. II, c. 85 e seg.}.

%%>>> t220

\tese{220}{A alma humana não é produzida pela evolução de uma especie inferior.}

a) Uma especie inferior, embora podesse na sua evolução produzir a alma humana, todavia não o poderia fazer, sem produzir ao mesmo tempo o corpo; pois que o agente finito não produz uma forma, senão modificando a materia. Ora o corpo humano não é produzido por uma especie inferior, que se desenvolva. Logo nem a alma humana é produzida pela evolução de uma especie inferior.

b) A alma humana, como demonstrámos, é \emph{espiritual}. Ora um ente \emph{espiritual} não é produzido pela evolução de uma especie inferior; pois o resultado da evolução de uma substancia \emph{material}, como é toda a especie inferior ao homem, é \emph{materia}; aliás o effeito seria superior á causa. Logo a alma humana não é produzida pela evolução de uma especie inferior \footnote{O homem e o animal possuem uma \emph{alma}, que é forma \emph{substancial} do seu organismo e principio \emph{primeiro} de todas as suas operações.\par Ora operações ha no homem, como as da \emph{ntelligencia} e da \emph{vontade}, que, tendo por objecto o \emph{abstracto}, o \emph{universal}, o \emph{immaterial}, não podem derivar senão de uma faculdade \emph{espiritual}, isto é, independente da materia; visto que as faculdades \emph{organicas}, ou dependentes da materia, só têm por objecto o \emph{concreto}, o \emph{singular}, o \emph{material}. Logo a \emph{ntelligencia} e á \emph{vontade} são faculdades \emph{espirituaes}, ou independentes da materia. E, como a faculdade não póde exceder em perfeição a substancia de que deriva, a alma, que é o principio de que resultam a \emph{intelligencia.e} a \emph{vontade}, é uma substancia \emph{espiritual}, isto é, independente da materia.\par Mas as operações do animal, pelo facto de terem \emph{todas} um objecto \emph{concreto}, \emph{singular} e \emph{material}, mostram que /\emph{odas} as faculdades do mesmo animal são \emph{materiaes} e \emph{organicas}. Sendo todas essas faculdades \emph{materiaes} e \emph{organicas}, é claro que tambem a alma do animal, da qual brotam, é \emph{material}; porque pela natureza dos effeitos conhecemos a natureza da causa.\par A \emph{alma} do homem, pois, é \emph{espiritual}; a do \emph{animal é material}.\par Se a alma do homem é \emph{espiritual} e a do animal é \emph{material}, segue-se que ellas são \emph{especificamente} diversas entre si; pois a essencia da alma do homem\pgn{358}possue um elemento, de que não é dotada a alma do animal e que é a \emph{espiritualidade}; ora todo o elemento \emph{essencial} produz diversidade na \emph{especie}.\par Sendo a alma do homem \emph{especificamente} diversa da do animal e superior a ella, é claro que \emph{a aima humana} não póde derivar da evolução do \emph{animal}; porque o individuo de uma especie não deriva do individuo de uma especie diversa, e porque o menos não póde dar o mais.\par D'ahi resulta que, embora o \emph{corpo humano} fosse em tudo e por tudo semelhante ao corpo do animal, todavia o /;\emph{ome}; nunca poderia derivar do \emph{animal}. Por quanto, o homem, \emph{por causa da sua alma espiritual}, é constituido numa \emph{especie diversa} da do animal: ora, como dissemos, o individuo de uma especie não póde derivar do individuo de outra especie.\par Alguns £\emph{ransformistas}, como Filippi (\emph{Revue de cours scientifiques}, année 1864, pag. 468, Paris), na sua pretencao de estabelecer affinidades genealogicas entre o homem e o macaco, prescindem dos caracteres \emph{psychologicos} d'este dois seres, e só attendem aos seus caracteres \emph{anatomicos}.\par Este processo é \emph{falso} e \emph{illogico}. --- E' \emph{falso}; porque o homem e o macaco devem ser considerados como são em si mesmos, e por isso tanto sob o aspecto \emph{anatomico} como sob o \emph{psychologico}, se quizermos determinar a sua origem. --- E' \emph{illogico}: porque os \emph{transform/istas}, pelo facto de admittirem que a evolução \emph{psychologica} dos seres vivos corresponde á sua evolução \emph{anatomica}, não podem deixar de considerar no homem e no macaco a duplice serie dos phenomenos --- \emph{psychologicos} e \emph{anatomicos}.\par Mais coherente foi \emph{Darwin}, encarando a questão e procurando mostrar que o homem e o macaco, relativamente ás faculdades \emph{intellectuaes}, pouco ou nada divergem entre si. \textasciicircum{}\par Examinemos os principaes argumentos de \emph{Darwin} e dos seus sequazes.\par a) \emph{Darwin} diz: A \emph{psychologia animal} demonstra que o animal, como o homem, possue muitos e variados conhecimentos; experimenta dor e prazer, amor e odio, sympathia e inveja, esperança e desespero. Alguma differenca, que existe entre a intelligencia do homem e a do animal, é apenas accidental; pois não é maior que a que existe entre o homem civilisado e o selvagem. Logo o animal é dotado de intelligencia e de razão, como é o homem, e por isso este deriva d'aquelle.\par Respondemos concedendo que o animal é dotado de \emph{conhecimento} e de \emph{appetite}. Mas este \emph{conhecimento} e este \emph{appetite} não sahem da esphera do \emph{sensivel}, do \emph{concreto} e do \emph{material}; e por isso, se manifestam que o homem convém com o animal no \emph{genero}, manifestam tambem que aquelle differe d'este na \emph{especie}. Por quanto, a differenca, que existe entre a percepção \emph{intellectual} e a \emph{sensitiva}, não é apenas gradual ou accidental, mas é especifica e essencial; pois trata-se de actividades completamente distinctas e dotadas de natureza inteiramente diversa.\par Todavia do facto de o homem convir com o animal no \emph{genero}, n\&o se segue que aquelle deriva d'este; porque, se assim fosse, a planta deveria derivar da pedra (pois ambas convém no \emph{genero} de \emph{substancia}), e o animal deveria derivar da planta (porque estas substancias estão sujeitas ao mesmo \emph{genero} de \emph{vivente}) --- Repetimos que, para concluirem para a descendencia\pgn{359}commum do homem e do animal, os \emph{transformistas} deveriam demonstrar que estes dois seres não só convêm no \emph{genero}, mas tambem na \emph{differença especifica}. Ora nunca poderão demonstrar que o animal é \emph{racional}, como o é o homem.\par Negamos que entre o homem e o animal exista menor differença que entre o homem civilisado e o selvagem. Por quanto, a differença entre o homem civilisado e o selvagem é apenas \emph{accidental}; ao passo que entre o homem e o animal a differenca é \emph{essencial}. O selvagem póde ser civilisado; o animal nunca poderá ser intelligente.\par ) \emph{Darwin} acerescenta: O animal é dotado de intelligencio; porque emprega os meios para os fins, opéra com prudencia, tem consciencia das suas acções, falla, abstrahe, julga, raciocina, possue o sentimento do bello, tem fé e religião, e não é destituido de sentimento moral.\par Todas estas asserções ou sao falsas, ou nada provam.\par \emph{a}) E' falso que o animal possua intelligencia, embora num grau menos elevado que o homem; porque, se possuisse inteligencia, deveria conceber ideas universaes. Ora o animal, como demonstrámos na \emph{Cosmologia}, não tem nem póde ter ideas universaes.\par \emph{db}) E' certo que o animal emprega os meios para os fins e parece operar com prudencia. Estas operações emanam do \emph{instincto}, que foi dado aos animaes pelo sapientissimo Auctor da natureza. Mas o \emph{instincto} deriva de uma faculdade \emph{sensitiva}, que é á força \emph{estimativa}; pois que os animaes, por não terem ideas universaes, não conhecem a \emph{relação} entre o fim e os meios nem o \emph{clivo} das suas operações.\par \emph{c}) E' falso que o animal seja dotado de \emph{consciencia} ou \emph{reflexdo}; pois estes actos emanam da inteligencia, de que é absolutamente destituido o animal. Certos actos d'elle, que parecem revelar \emph{consciencia} ou \emph{reflerão} (como quando o cão deixa um alimento que appetece, porque teme o castigo que o ameaça, fazendo como uma \emph{comparação} entre o prazer proveniente da comida e a dor produzida pelo castigo) derivam do \emph{sentido commuwm}, da memoria \emph{sensitiva} e da \emph{força estimativa}, que são todas faculdades \emph{organicas} e proprias dos animaes superiores.\par \emph{d}) E' falso que o animal possua o dom da \emph{palavra}. Não insistimos neste ponto, pois já foi tratado na \emph{Cosmologia especial}.\par \emph{e}) E' falso que o animal \emph{abstraia}, \emph{julgue e raciocine}. Porquanto, se o animal abstrahisse, julgasse e raciocinasse, deveria revelal.o nas suas operações. Ora as operações do animal, sendo todas \emph{concretas} e \emph{singulares}, não revelam que elle abstraia, julgue ou raciocine. Alguns actos, que parecem encerrar um \emph{juizo} ou um \emph{raciocinio}, derivam da \emph{força estimativa}.\par \emph{f}) E' falso que o animal-possua o sentimento do \emph{bello}. Porquanto, esse sentimento suppõe necessariamente a intelligencia, quer o bello seja \emph{espiritual}, quer seja \emph{material}. A impressio agradavel, que a bellezza \emph{material} possa produzir no animal, não póde chamar-se sentimento do bello.\par 9) E' falso que o animal tenha fé e \emph{religião}. \emph{Darwin} opina que a fée a religião são o resultado da illusio dos que \emph{sonham}, e como os animaes tambem \emph{sonham}, tambem elles tem fé e \emph{religião}. --- Este discurso é absurdo. A fé e a \emph{religião} não podem derivar dos sonhos. Se não tivessemos nenhum\pgn{360}conhecimento de \textsc{DEus}, quando estamos acordados, não poderiamos sonhar com Elle, quando estamos a dormir; visto que o sonho só póde despertar na phantasia imagens de objectos \emph{anteriormente} percebidos.\par \emph{h}) E' falso que os animaes possuam o \emph{sentimento moral}. Porquanto (supposto que ao \emph{sentimento moral} se dê a significação que lhe dão os homens de criterio são) devemos dizer que não ha \emph{sentimento moral}, onde não ha \emph{liberdade}, e onde fallecem principios \emph{universaes} e \emph{invariaveis}. Ora o animal é destituido de \emph{liberdade}, nem possue principios \emph{universaes} e \emph{immutaceis}.\par 6) Os \emph{transformistas} objectam: Ha muitos factos, que demonstram que o animal é dotado de intelligencia e de razão.\par Respondemos que a maior parte d'esses factos foram inventados ou mal interpretados pelos escriptores ou exploradores, que os contaram.\par \emph{Paulo Gereais} diz: « Nunca será demasiada a desconfiança ácerca da facilidade, com que certos observadores menos cuidadosos tem concedido a muitos animaes sentimento e raciocinio, que só existem, na maior parte dos casos, no espirito dos observadores. Casos houve em que os antigos viajantes abusaram da confiança, que nelles tinham os seus contemporaneos. E cumpre dizel-o, muitas vezes, como tambem hoje, o gosto do publico por tudo quanto \textbar{} é exaggerado, ou devido só á imaginação dos escriptores, tem levado alguns d'estes a publicar, como verdades, muitas pretendidas observações, que na realidade não passam de contos imaginosos e divertidos. E por isso indispensavel ter o maior cuidado na escolha, que deve fazer-se entre os pormenores publicados ácerca dos animaes ». --- O mesmo dizem \emph{Cuvier} e \emph{Chaillu}.\par A consequencia, que d'alguns factos certos e averiguados póde legitimamente deduzir-se, é que em alguns animaes o grau do conhecimento \emph{sensitivo} é mais elevado, e ha nelles uma semelhança, ainda que remotissima, da intelligencia, --- semelhança impressa no sentido da \emph{estimativa}; porque os seres mais perfeitos de um grau inferior participam, de algum modo, das prerogativas dos seres do grau immediatamente superior.\par Com 0s mesmos criterios, apresentados aqui o na \emph{Cosmologia especial}, refutam-se todas as difficuldades, que os \emph{transformistas} adduzem para demonstrar que os animaes são dotados de intelligencia, embora não desenvolvida, --- difficuldades que --- ou nada provam, porque são baseadas sobre falsos suppostos, pois confundem a \emph{idea} com o \emph{phantasma}, o \emph{pensamento} com a \emph{sensação}, a. \emph{intelligencia} com o \emph{instincto} --- ou, se provam alguma coisa, provam que o animal é dotado de faculdades \emph{sensitivas}.\par a) Os \emph{transformistas} insistem: Nas camadas geologicas, nas cavernas, nos antigos tumulos encontraram-se instrumentos de \emph{pedra lascada}, de \emph{pedra polida}, de \emph{bronse} e de \emph{ferro}. Os instrumentos de \emph{pedra lascada} s\&o obras do \emph{homem primitivo}, porque são tão grosseiros, que suppõem uma intelligencia proxima da dos animaes; os da \emph{pedra polida} mostram o homem jà um pouco mais adeantado, e os de \emph{bronse} e de \emph{ferro} manifestam o homem como é no estado actual. Ora estes factos demonstram que o homem progrediu necessaria mas lentamente, o que este progresso lento não pôde deixar de ser o resultado de uma evolução, cujo ponto de partida é o estado selvagem e quasi bestial. ; :\par E' certo o facto da descoberta d'aquelles instrumentos, mas é falso ---\pgn{361}4) que os instrumentos de \emph{pedra lascada} sejam obra do \emph{homem primitivo}, --- \emph{b}) que esses instrumentos sejam tão grosseiros, que revelem uma intelligencia proxima da dos animaes, --- \emph{c}) que o homem progrediu lenta e necessariamente, --- \emph{d}) e que o ponto de partida da evolução fosse o estado selvagem.\par \emph{d}) E' falso que os instrumentos de \emph{pedra lascada} sejam obra do \emph{homem primitivo}. Na verdade, a historia demonstra que, desde a mais remota antiguidade, no \emph{Egypto} e na \emph{Asia menor}, se conhecia o uso dos metaes. --- O sabio \emph{Chabas}, no sen livro intitulado --- \emph{Éludes sur l'antiquité historique} (pag. 860), diz que, no seculo XVII antes de Jesus Christo, os Egy pelos « \emph{conheciam}, \emph{havia Jd} 2000 \emph{annos}, \emph{o uso de todos os metaes e tinham todos os habitos de um luco desenvolvido pela riqueza} s». --- \emph{M}. \emph{Schlieman} diz que, escavando nas ruinas de \emph{Hissarlih}, na antiga \emph{T'roada}, na profundidade de 10 ou 15 metros, encontrou, ao mesmo tempo e \emph{na mesma camada}, objectos de bronze e de prata, misturados com bellas louças de barro. Estas camadas, por conterem os restos de tres civilisacgóes, pertenciam á mais remota antiguidade.\par 5) E' falso que os instrumentos de \emph{pedra lascada} demonstrem que o \emph{homem primitivo} tivesse uma intelligencia proxima das faculdades perceptivas dos animaes, e é falso pelas seguintes razões:\par 1.º) O homem da edade da \emph{pedra}, como vimos no artigo antecedente, emquanto á constituição physica e especialmente emquanto á \emph{capacidade} e \emph{forma} do \emph{craneo}, era semelhante ao homem moderno. Se era semelhante emquanto á \emph{constituição physica}, a analogia exige que lhe fosse tambem semelhante na intelligencia.\par 2. Além d'isso, a grosseria dos instrumentos, de que o homem primitivo se teria servido, não é o nivel do desenvolvimento da intelligencia ou da razão; aliás deveriamos dizer que os escriptores dos seculos do \emph{oz}70 da nossa litteratura não tinham a intelligencia desenvolvida, visto que os instrumentos de \emph{mechanica} não estavam nesse tempo tão aperfeiçoados como hoje.\par 3.º) Finalmente, conhecemos nós todos os instrumentos, de que os homens primitivos se serviam? Não. O tempo nos conservou apenas os de \emph{pedra}, porque não sujeitos a alterações. Mas quantos outros instrumentos devia o homem fabricar com materiaes mais frageis e corruptiveis, que não chegaram até nós? Logo os instrumentos de \emph{pedra lascada} não podem dar uma idea adequada da intelligencia dos homens primitivos.\par c) E' falsa a theoria dos \emph{transformistas} \&cerca do \emph{progresso} da humanidade. E' certo que o homem é essencialmente perfectivel e que se aperfeiçoa e progride; mas esta \emph{lei} de progresso não é \emph{absoluta} nem \emph{necessaria}, não é \emph{universal} nem se verifica no sentido dos \emph{transformistas}. Se appellarmos para a historia, esta nos dirá que a \emph{Persia}, a \emph{Syria}, a Asia \emph{menor}, a \emph{Africa do norte}, depois de se terem elevado ao apogeo da civilisacüo e da gloria, descahiram, precepitando-se na. obscuridade e tornando-se quasi barbaras.\par \emph{d}) E' falso que o estado primitivo do homem fosse o selvagem e quasi bestial, e é falso pelas seguintes razàües:\par 1.º) A tradição de todos os povos, não só civilisados mas tambem barbaros, diz que o estado primitivo do homem não foi um estado selvagem, mas civilisado, e que a primeira edade do genero humano não foi uma edade de \emph{pedra}, mas uma edade de ouro. --- A tradição da \emph{Grecia} foi representada\pgn{362}principalmente por \emph{Platão}, \emph{Hesiodo}, Arato, \emph{Cratilo}, \emph{Telecides}, \emph{Ferecrates}, \emph{Dicearcho}, diseipulo de \emph{Aristoteles}. \emph{Platão}, na sua \emph{Politica} e no livro III da \emph{Republica}, diz que \emph{os primeiros homens}, \emph{creados e dirigidos por Deus}, \emph{viviam} num \emph{estado de felicidade}, e accrescenta que attingia estas noticias, ácerca da condição primitiva do homem, na tradição dos antepassados. --- A tradição dos \emph{Romanos} não era menos viva nem menos certa. Virgilio, na \emph{Georgica}, falla num antigo reino de \emph{Saturno}, numa antiga edade de civilisacao e de felicidade. \emph{Ovidio}, na sua \emph{Metarmophose}, celebra \emph{a edade d'ouro}, \emph{que foi primeira} a \emph{amanhecer}, \emph{reinando então uma eterna primavera}, --- Os \emph{mongoes}. como attesta \emph{Bergman}, diziam: « \emph{Os nossos primeiros paes viviam num estado feliz}; a \emph{sua felicidade durou pouco tempo}; \emph{perderam-na por culpa propria}. \emph{Do solo brotavam e cresciam fructos doces e brancos}; \emph{o aspecto d'elles seduziu os olhos do homem}; \emph{comeu e acabou a bemaventuranca} ». --- Os \emph{persas}, como attesta \emph{Creuser}, tinham nos seus livros estas palavras: « \emph{O primeiro homem e a primeira mulher}, \emph{Meschia} e \emph{Meschiani}, \emph{foram felizes na sua origem}, \emph{perfeitos e sujeitos a Orimusd}, \emph{seu Deus e seu auctor} ». --- Os \emph{chineses}, os \emph{indios}, os \emph{americanos}, e, finalmente, todos os povos, cujas tradições são conhecidas, admittiam que a primeira edade do mundo foi uma edade de \emph{ouro}.\par 2.) Se o estado primitivo do homem fosse o selvagem, o homem nunca se teria civilisado. Porquanto, embora um, povo selvagem possa, \emph{por si} só, sahir do seu estado de barbarie, comtudo é um \emph{facto historico} que nenhum povo sahiu jamais de tal estado sem o auxilio externo de outro povo já civilisado. Poderiam os \emph{transformistas} citar um exemplo de um povo, que, \emph{unicamente pelos seus esforços}, se elevasse do estado selvagem ao estado de civilisacao? Não. O que sabemos é que o Japão recebeu da China a civilisação, Roma da Grecia, e esta do Egypto e dos povos orientaes. Se, nos nossos tempos, alguns selvagens abrandaram seus costumes e receberam a luz da civilisação, esta mudança é devida ao contacto com povos civilisados e sobre tudo á abnegação e ao heroismo dos missionarios catholicos.\par A. \emph{archeologia prehistorica} confirma as nossas asserções. A. \emph{Bertrand}, na sua \emph{Archéologie celtique et gauloise}, escreve: « Em geologia póde haver uma lei immutayel na successão dos terrenos de toda a crusta do globo, terrenos primarios, segundarios, terciarios e quaternarios, com subdivisões egual e nitidamente distinctas; mas nas agglomerações humanas, na successão das camadas da civilisação, não existe semelhante lei. Seria um erro gravissimo julgar que todas as racas humanas tém passado necessariamente pelas mesmas phases de desenvolvimento e percorrido toda a serie dos estados sociaes, que a theoria quer impor-Ihes » (pag. 46).\par « Alguns philosophos theoricos pretenderam que o homem fôra condemnado em toda a parte a passar successivamente, e como por uma lei da sua propria natureza, do estado de cacador nomada ao de pastor, em seguida ao de agricultor, antes de attingir a perfeição do estado social. Taes theoria\& têm sido até hoje desmentidas pelos factos, pelo menos emquanto á Europa. As \emph{primeiras gerações de homens}, \emph{entregues} a \emph{si mesmas}, \emph{não souberam em parte alguma das nossas regiões ultrapassar o limite certo}, \emph{que pela Providencia foi assignado ao desenvolvimento das suas faculdades isoladas}. Por duas vezes differentes, na Gallia, foram novos grupos humanos, que fizeram surgir\pgn{363}do seu somno as populações anteriores, com as quaes elles parece terém-se fundido posteriormente, communicando-]hes novas aptidões, mas talvez tambem participando d'ellas » (pag. 79).\par Mas admittamos que se verifieasse aquella \emph{lei} de progresso lento e continuo, formulada pelos \emph{transformistas}. Seguir-se-hia por ventura que o homem deriva do macaco? Não, com certeza. A unica consequencia, que se poderia legitimamente deduzir dos \emph{progressos realizados} pelo homem, é que neste existe uma faculdade espiritual ---- a razão; que o homem, desde a sua apparição sobre a terra, era já immensamente distante do animal, e por isso não derivou pela evolução d'este. Porquanto, o \emph{progresso} não póde conceber-se sem o uso da razão; como, pelo contrario, o estado \emph{estacionario} dos animaes prova que elles são destituidos de razão.}.

%%>>> t221

\tese{221}{A alma humana não é produzida por geração.}
Se a alma humana fosse produzida por \emph{geração}, deveria ser produzida ou \pg{358} \pg{359}  por geração \emph{corporea}, ou por geração \emph{espiritual}. Ora a nossa alma não póde ser produzida por \emph{geração corporea}, nem por \emph{geração espiritual}.

a) Não póde ser produzida por \emph{geração corporea}. Uma força \emph{material} não póde produzir uma substancia \emph{espiritual} ou \emph{immaterial} \pg{361} ; pois o effeito não póde ser superior á causa. Ora á força, que opéra na geração \emph{corporea}, é e não pode deixar de \pg{362}  ser \emph{material}, emquamto a alma humana, por ser dotada de faculdades e operações \emph{zmmateriaes} ou \emph{espirituaes}, é \emph{immaterial}

\emph{p} ou \emph{espiritual}. Logo a alma humana não póde ser produzida por \emph{geração corporea} \footnote{A conclusão é evidente. O producto de uma energia corporea ou é corpo, ou depende do corpo. Ora a alma humana não é corpo, porque é simples, nem depende do corpo, porque é subsistente ou espiritual. Logo a alma humana não é o producto de uma energia corporea. Diz S. Thomaz: u Impossibile est virtutem activam, quae est in materia, extendere suam actionem ad producendum immaterialem effectum. Manifestum est autem quod principium intellectivum in homine est principium transcendens materiam; habet enim operationem, in qua non communicat corpus. Et ideo impossibile est quod virtus, quae est in semine, sit productiva intellectivi principii » (\emph{Sum}. \emph{Th}., p. I, q. 118, a. 2), --- Em menos palavras: A origem de uma coisa corresponde ao seu modo de operar. Ora a nossa alma é independente do corpo na sua operação. Logo é independente do corpo na sua origem (Cf. \emph{C}. \emph{Gent}., l. II, c. 86, n. 1).}.

b) Não póde ser produzida por \emph{geração espiritual}, A alma poderia ser produzida por \emph{geração espiritual}, se a alma dos paes --- ou se dividisse em \emph{partes}, --- ou transmittisse uma \emph{materia espiritual}, --- ou tivesse o poder de \emph{criam} outras almas. Ora estas tres hypotheses são impossiveis. --- I? impossivel a \emph{primeira}; pois a alma dos paes, sendo \emph{simples}, não tem partes, em que possa dividir-se. --- E' impossivel a \emph{segunda}; pois aquella materia, embora existisse, não poderia formar a alma; aliás esta seria \emph{composta}. --- E' impossivel a \emph{terceira}; porque a \emph{creação}, sendo a producção do \emph{nada}, não é acção propria das creaturas, que só operam num sujeito preexistente, transformando-o. Logo a alma humana não pode ser produzida por \emph{geração espiritual} \footnote{Desenvolvamos as partes da \emph{menor} do argumento. --- A alma dos filhos não pode ser uma \emph{parte} da alma dos paes. Nem se diga que a alma dos filhos, se não é uma \emph{parte}, é um \emph{producto} da alma dos paes. Porquanto, correspondendo o \emph{ser} de uma coisa á sua \emph{origem}, se a alma depende de um sujeito na sua \emph{producção}, depende tambem de um sujeito no seu \emph{ser}, e por isso não é subsistente. Diz S. Thomaz: « Ponere animam intellectivam a generante causari, nihil aliud est quam ponere eam non subsistentem, et\pgn{364}per consequens corrumpi cum corpore » (\emph{Sum}. \emph{Th}., p. I, q. 118, a. 2). --- A alma dos filhos não pode ser o producto de uma \emph{materia espiritual}, transmittida pelos paes. Porquanto, se assim fosse, não só a alma dos paes, mas tambem a dos filhos seria \emph{composta}, e não \emph{simples}. Seria \emph{composta} a alma dos \emph{paes}, porque tal materia \emph{espiritual}, transmittida por elles, seria uma \emph{parte} da alma, e por isso esta seria divisivel e composta. Seria \emph{composta} a alma dos filhos, porque resultaria constituida d'essa materia ou sujeito \emph{espiritual} e da \emph{forma}, pois o termo da geração é sempre o \emph{composto}; e d'este modo a nossa alma, numericamente \emph{una}, seria acto ou forma de materias diversas, isto é, da materia espiritual e da materia corporea, o que é absurdo, porque um e o mesmo acto não pode ser recebido em potencias de genero diverso (Cf. 2 \emph{Sent}. dist. 17, q. 1, a. 2). --- A alma dos filhos não pode derivar da alma dos paes por \emph{creação}. A produccáo do \emph{nada} é um acto exclusivamente proprio de Deus, que, assim como é independente na existencia, assim tambem é independente na operação, de modo que não exige um sujeito preexistente, sobre o qual exerça a sua actividade. Nem se diga, com \emph{Avicenua} e com \emph{Pedro Lombardo}, que Deus communica a sua força creativa á alma dos paes. Não pode communicar-se, diz-S. Thomaz, um poder infinito a uma creatura, dotada de natureza finita. Ora a creação exige um poder infinito, porque entre o ser e o não ser a distancia é infinita. Logo não pode communiecar-se a uma creatura o poder creativo » (\emph{In IV Sent}., dist. 5, a. 3).\par O systema refutado é o \emph{generacianismo}. Chama-se assim, porque admitte que por \emph{geração} se produz todo o homem --- \emph{alma} e \emph{corpo}.\par O \emph{generacianismo} é duplice: \emph{corporeo} e \emph{espiritual}.\par O \emph{generacianismo corporeo}, ensinado por \emph{Apollinario} e, talvez, por \emph{Tertulliano}, sustenta que a alma humana é produzida pela força plastica do germen, que forma o corpo.\par O \emph{generacianismo espiritual} admitte que a \emph{alma} dos filhos deriva da \emph{alma} dos paes. Esta theoria, como dizemos no texto, pode significar --- ou que a alma dos filhos é uma \emph{parte} da alma dos paes, --- ou que é formada por uma materia \emph{espiritual}, transmittida pelos paes, como pensavam alguns escriptores catholicos no tempo de 8. Agostinho, --- ou que é \emph{creada} pela propria alma dos paes, como ensinou ha poucos annos, depois de \emph{Avicenna} e de \emph{P}. \emph{Lombardo}, o dr. \emph{Frohschammer}.\textasciicircum{}\par \emph{Antonio IRosmini} reproduziu, sob um outro aspecto, o systema, do \emph{generacianismo}, Disse que a alma dos filhos, \emph{emquanto sensitiva}, deriva por \emph{geração}, e que, em seguida, se transforma em \emph{intellectual} pela apparicao da \emph{idea de ente}, que \textsc{Deus} lhe manifesta. D'este modo, segundo \emph{Rosmini}, na geração do individuo humano, concorrem dois agentes: o homem com a geração, \textsc{Deus} com a manifestação da sua luz; o homem forma o animal, \textsc{Deus} cria a alma intellectual no mesmo instante em que é produzido o \emph{tomem animal}, iluminando-a com o esplendor do seu rosto e communieando-l]he alguma coisa de si mesmo, --- o \emph{ente} ideal, que é a luz das creaturas intellectuaes. (\emph{Antropologia}, lib. IV, c. V, n. 819. --- \emph{Psicologia}, v. I, lib. IV, e. 98). --- -\pgn{365}A opinião de \emph{Rosmini} nem sequer tem o merito da novidade, porque foi defendida por alguns antigos philosophos, e é exposta por S. Thomaz com as seguintes palavras: « Alii dicunt, quod illa eadem anima, quae \emph{primo} fuit \emph{vegetativa} tantum, postmodum per actionem virtutis, quae est in semine, perdueitur ad hoe, ut 2\emph{psa eadem}. fiat \emph{sensitiva}, et tandem ad hoc ut \emph{ipsa eadem}. \emph{fiat intellectiva}, non quidem per virtutem activam seminis, sed per virtutem superioris agentis, scilicet. \emph{Dei}, \emph{deforis illustrantis} » (\emph{Sum}. \emph{Th}., p. I, q. 118, a. 2 ad 2).\par O systema de \emph{Rosmini} é falso pelas seguintes razões:\par \emph{a}) E' uma renovação do \emph{generacianismo}. Por quanto, o ser da alma é unico e semplicissiBTo, e por isso se a alma deriva por \emph{geração}, emquanto \emph{sensitiva}, deve tambem derivar por \emph{geração}, emquanto é \emph{intellectual}; o que o proprio \emph{Rosmini} concede (\emph{Cf}. \emph{Psicologia}, loc. cit.). Mas se a alma humana deriva por \emph{geração}, não é \emph{creada}, não é \emph{espiritual}, nem \emph{immortal}.\par \emph{b}) \emph{Rosmini} diz que a alma \emph{sensitiva}, contemplando o ente \emph{ideal}, que \textsc{Deus} lhe manifesta, se transforma em \emph{intellectual}. Ora isto é absurdo. Por quanto a alma \emph{sensitiva}, como \emph{tal}, não póde perceber o \emph{ente ideal}; visto que um ente \emph{idea}! só póde ser percebido por uma forma intellectual ou espiritual. Por isso a visão do \emph{ente ideal}, longe de transformar a alma sensitiva em intellectual, suppõe-na ja \emph{intellectual}, quando aquelle \emph{ente} lhe apparece.\par c) Embora a visão do ente ideal tivesse o poder de transformar a alma de \emph{sensitiva} em \emph{intellectual}, todavia a alma \emph{sensitiva} não poderia estar sujeita a uma transformação. Por quanto a.transformação exige composição no ser que se transforma; ora a alma sensitiva não é composta, mas é \emph{simples}. --- O proprio \emph{Rosmini} confessa que a alma sensitiva, tornando-se racional, fica substancialmente a mesma, e não muda de natureza. E então fica \emph{formalmente} sensitiva, e não é \emph{espiritual}.\par \emph{d}) Finalmente, como argumenta S. Thomaz contra os sequazes d'esta opinião, a acção de \textsc{Deus}, illuminaudo a alma sensitiva, deve produzir um effeito. Ora este effeito ou é subsistente, ou não. Se é \emph{subsistente}, é uma coisa distincta da alma sensitiva, que não é subsistente; e assim teremos varios principios vitaes no homem, o que demonstrámos ser falso. So não é \emph{subsistente}, é uma perfeição \emph{accidental} da alma \emph{sensitiva}; e assim a alma humana, sendo intelligente só \emph{accidentalmente} e não por \emph{natureza}, não será \emph{espiritual}, nem differir\& senão \emph{accidentalimente} da alma dos animaes, e por isso corromper-se-ha na corrupção do composto (\emph{Jb}. --- Esta doutrina de \emph{Rosmini} foi condemnada pela Egreja.\par Objectam: A nossa alma é \emph{sensitiva}. Ora toda a alma \emph{sensitiva} deriva por \emph{geração}. Logo a nossa alma deriva por geração.\par Resposta. A alma humana é \emph{sensitiva}, mas só \emph{virtualmente}, porque \emph{formalmente} é intellectual. Sendo \emph{formalmente} intellectual, não póde ter origem por \emph{geração}; pois tudo o que é gerado é \emph{material}. (Cf. c. \emph{Gentes}, 1. 2, c. 89).}.

222. A \textbf{alma humana não} é \textbf{produzida} por \textbf{emanação da substan}-

HM \textbf{cia divina}.

á) Se a alma humana fosse produzida por \emph{emanação} da substancia divina, seria --- ou uma \emph{parte}, --- ou uma \emph{propriedade} ou \emph{manifestação} da Divindade. Ora não póde ser \emph{parte} \pg{365}  da Divindade; porque esta, por ser simplieissima, é indivisivel. Nem é uma \emph{propriedade} ou \emph{manifestação} da mesma Divindade; pois que a nossa alma não é um accidente, nem um modo, nem um phenomeno de uma outra substancia, mas é uma substancia distincta de todas as outras. Logo a alma humana não é produzida por \emph{emanação} da substancia divina.

b) Se a alma fosse produzida por \emph{emanação} da substancia divina, --- ou a mesma alma teria todas as perfeições divinas, --- ou \textsc{Deus} estaria sujeito aos defeitos, a que está sujeito o homem. Ora estas consequencias são absurdas. Logo a alma humana não é produzida por \emph{emanação} da substancia divina \footnote{O \emph{emanatismo} foi ensinado, --- na antiguidade, por \emph{PytAagoras} e pelos \emph{Staicos} --- e, nos tempos modernos, por \emph{Poiret} e por todos os \emph{pantheistas}.\par Não nos demoramos na refutação do \emph{emanatismo}; porquanto este systema é apenas uma \emph{applicação} do \emph{pantheismo}, e o \emph{pantheismo} foi exposto e refutado na \emph{Cosmologia geral}, cap. I, art. IV. (\emph{Of}. e. \emph{Gentes}, 1. 2, c. 85).}.

2238. A \textbf{alma humana} é \textbf{creada immediatamente} por \textbf{Deus}.

a) A alma humana, sendo effeito, deve ter uma causa. Esta causa ou é finita ou infinita. Ora não é uma causa finita; porque, como vimos, a nossa alma não é produzida por evolução de uma especie inferior, nem por \emph{geração}. Logo é a causa infinita, é \textsc{Deus}. Ora \textsc{Deus} tira a nossa alma --- ou da propria substancia divina, --- ou de uma materia preexistente, --- ou do \emph{nada}. Mas não a tira da propria substancia divina; pois essa \emph{emanação}, como vimos, é absurda. Nem a tira de uma materia preexistente; porque a nossa alma é \emph{espiritual}. Logo a tira do \emph{nada}. Tirar do \emph{nada é criar}. Logo a alma humana é \emph{creada} immediatamente por \textsc{Deus}.

b) Uma forma, que não depende da materia na sua \emph{existencia}, não depende della na sua \emph{origem}; pois que a origem é o começo da existencia, e todo o ser possue no seu começo a mesma natureza, que conserva na sua duração. Ora a alma humana não depende da materia na sua existencia, pois que é \emph{espiritual}. Logo nem depende da materia na sua origem. Se não depende da materia na sua origem, é tirada do \emph{nada}, \emph{é creada}. Ora a \emph{creação} é uma acção exclusivamente propria de \textsc{Deus}; pois os agentes finitos não podem produzir as formas, senão mudando a materia preexistente. Logo a alma humana é \emph{creada} immediatamente por \textsc{Deus}.

Dizemos que alina humana é creada \emph{Zmmedtatamente} por \textsc{Deus}, para signifiearmos que \textsc{Deus}, na creação da nossa alma, não se serve dos agentes finitos, nem como de \emph{causas instrumentaes} ; porque o agente finito presuppõe sempre um sujeito em que opére, e a \emph{creação} é a producção do \emph{nada} \footnote{Cf. \emph{Sum}. \emph{Th}, p. I, q. 90, aa. 2 e 8; c. \emph{Gentes}, l. II, c. 81. --- Como todo o agente imprime a propria semelhança no eífeito, a nossa alma é semelhante a Deus, como o são todos os seres que existem. --- Mas a semelhança da nossa alma com Deus é mais perfeita que a das creaturas corporeas;\pgn{367}porque, emquanto a semelhança d'estas consiste num simples vestigio, a semelhança d'aquella é uma verdadeira \emph{imagem}, embora infinitamente distante do original. --- A nossa alma é \emph{imagem} de Deus --- 1.º) porque é incorporea, indivisivel, espiritual, immorta!, como é Deus, --- 2.º) porque é dotada de intelligencia, vontade e memoria, --- 8.º) porque é capaz de sabedoria, de virtude, da graça e de felicidade --- 4.º) porque domina e preside a todos 0s animaes, --- 5.9) porque, assim como Deus Padre, conhecendo-se a si mesmo pela intelligencia, gera o Verbo, que é seu Filho, e, amando o seu Verbo, dà origem ao Espirito Santo, assim tambem o homem, conhecendo-se a si mesmo, produz o verbo intellectual, que o exprime, e, revendo-se no verbo, produz um acto de amor, e por isso representa, de algum. modo, o mysterio da SS. Trindade. (Cf. \emph{Sum}. \emph{Th}., p. 1, q. 93). --- A nobilissima origem do nosso corpo e da nossa alma vemol-a indicada claramente na Sagrada Escriptura. Quando trata da origem das plantas e dos animaes irracionaes, a B. Escriptura diz que Deus mandou ás causas segundas que os produzissem; mas quando falla na origem do primeiro homem, claramente indica qua foi Deus quem formou immediatamente o corpo de Adão e produziu do nada a alma humana: --- \emph{Formavit Deus hominem de limo terrae}, \emph{et inspiravit}. \emph{in faciem eius spiraculum vitae n}. (\emph{Gen}., IL, 7).\par Obiectam: A theoria da creacáo \emph{immediata} das almas repugna á grandeza e magestade de \textsc{Deus}, pois suppõe que \textsc{Deus} esteja sujeito ás creaturas.\par Respondemos que a creacüo immediata da nossa alma não só não torna Drs sujeito ás creaturas, mas é um testemunho da sua Omnipotencia. --- Não torna DEkus sujeito ás creaturas; pois, quando cria as almas humanas, \textsc{Deus} segue a ordem, que Elle proprio estabeleceu. --- E' testemunho da sua Omnipotencia; pois a creacüo, sendo a producção do \emph{nada}, só compete Àquelle que, na existencia e na operação, não depende de outro ente.\par Pergunta-se: As almas humanas, quando são creadas por \textsc{DEUS}, são eguaes ou deseguaes entre si em perfeição?\par A pergunta, como se vé, não se refere \& perfeição \emph{accidental}, mas á \emph{substancial}. --- Não se refere á perfeição \emph{accidental}; pois que as almas, no momento em que são creadas, ainda não possuem as perfeições, que adquirem mais tarde (como a \emph{sciencia}, a \emph{virtude}, etc), e pelas quaes se tornam \emph{accidentalmente} deseguaes. --- Mas refere-se á perfeição \emph{substancial}; pois é esta perfeição que as almas possuem no momento da sua creação.\par Ora a perfeição \emph{substancial} póde ser \emph{especifica} ou \emph{entitativa}, --- emquanto é um elemento \emph{essencial}, proprio de uma especie e não de outra, --- ou emquanto se considera nos seus graus de intensidade.\par Todos os philosophos concordam em relação á perfeição \emph{especifica} das almas, porque todos admittem que as almas são \emph{especificamente} eguaes; mas divergem com relação a sua perfeição \emph{entitativa}.\par Uns dizem que as almas humanas, quando são creadas por \textsc{DEus}, possuem todas a mesma perfeição \emph{entitativa}, isto é, sao \emph{entitativamente} eguaes, e que, se divergem entre si, isto é, se uma é mais intelligente que a outra, esta\pgn{368}diversidade deriva do organismo, que, como vimos, concorre, pelo menos como \emph{condição}, em \emph{todas} as operações da alma.\par Outros philosophos dizem que as almas humanas, quando são creadas por \textsc{Deus}, são \emph{entitativamente} deseguaes entre si, de modo que, embora todas possuam as mesmas perfeições \emph{especificas}, todavia umas possuem-nas num grau mais \emph{intenso} que as outras. Por quanto, devendo a forma ser sempre proporcionada á meteria, é claro que, quanto mais perfeito for o organismo, tanto Mais perfeita deverá ser a sua forma. Ora a experiencia attesta que entre os organismos humanos um é mais perfeito que outro. Logo alma, que é destinada a informar o organismo mais perfeito, deve ser \emph{entitativamente} mais perfeita que uma outra, destinada a informar um organismo menos perfeito. D'ahi a distinegao \emph{numerica} de uma alma de outra, --- distineção proveniente das suas relações aos diversos organismos.\par Ambas estas opiniões são provaveis; pois baseam-se em boas razões e têm defensores entre os melhores philosophos. --- Todavia preferimos a segunda; porque \textsc{DEus}, quando cria a nossa alma, não a cria \emph{fóra} do organismo, mas no \emph{proprio organismo}, e, por isso, proporcionada a este; visto que a forma \emph{substancial} deve ser proporcionada á materia. E' a doutrina de S. Thomaz. -Ào quesito --- se \emph{todas as almas são eguaes na sua creagüo} --- o Angelico Doutor responde \emph{negativamente}, e dà a seguinte razão: Todas as almas racionaes pertencem á mesma especie, e differem unicamente no numero. Esta diversidade deriva da materia; e por isso, como as almas não são compostas de materia, toda a diversidade e distineção de graus nas almas deve ser causada pela diversidade do corpo; de modo que, quanto mais perfeita é a compleicio do corpo, tanto mais nobre deve'ser a alma, visto que á coisa recebida deve ser sempre proporcionada ao sujeito, que a recebe ». (In 2 \emph{Sent}., dist. 32, q. 2, a. 3), --- A mesma coisa o S. Doutor ensina na Summa \emph{Theol}., dizendo: « Quanto corpus est melius dispositum, tanto meliorem sortitur animam... Cujus ratio est quia actus et forma recipitur in materia, secundum materiae capacitatem » (p. I, q. 85, a. 7). Como se vé, trata-se de desegualdade \emph{entitativa} e \emph{intrinseca}, proveniente da desegualdade dos organismos humanos. --- Como adverte S. Thomaz, á diversidade material, que é causa da diversidade das almas, não é a \emph{especifica}, que existe entre as varias partes do organismo, \emph{especificamente diversas} (a \emph{cabeça}, o \emph{thoraw}, as \emph{pernas}), e que é commum a todos os individuos da mesma especie, mas éa \emph{individual}, que existe entre os varios individuos e não pertence á especie. A primeira é causada pelas formas, que fazem com que o corpo seja dotado de taes partes: a segunda, pelo contrario, produz a diversidade nas formas. (In 2. \emph{Sent}. dist. 32, q. 2, a. 8 ad 6).}.

%%>>> t224

\pg{360} \pg{363} \pg{364} \pg{366} \pg{367} \tese{224}{A alma é creada no momento da sua união com o corpo.}

a) A alma humana é forma \emph{substancial} do corpo, e com elle constitue uma substancia ou natureza completa. Ora a ordem \pg{368}  natural exige que a forma substancial de um composto só seja produzida, quando é formado o proprio composto. Logo a alma humana é creada no momento da sua união com o corpo.

b) A alma humana é substancia \emph{incompleta}; porque é parte, embora nobilissima, da natureza humana. Ora nenhuma natureza é creada por \textsc{Deus} num estado incompleto e imperfeito. Logo a alma é creada no momento da sua união com o corpo \footnote{Apresentamos, em confirmação, da nossa these, este outro argumento. --- As almas humanas, embora convenham entre si na \emph{especie}, todavia differem\pgn{369} no \emph{numero}, pois que a alma não é uma e a mesma em todos os homens, mas a alma de um individuo é distincta da de outro individuo. Esta distincção numerica deriva --- ou da propria alma, --- ou do corpo. Ora não póde derivar da alma; pois toda a differenca na \emph{forma} produz differenca na \emph{especie}. Logo deriva do \emph{corpo}, pois a materia, de que o corpo é composto, é divisivel em partes distinctas. Se a distineção \emph{numerica} das almas deriva do corpo, é claro que não podem ser creadas antes da sua união com o corpo e fora d'elle; pois todo o ente, que existe na realidade, é \emph{individuo}, \emph{singular}, e por isso \emph{numericamente} distincto dos autros entes. Logo a alma humana é creada no momento da sua união com o corpo e no corpo.\par A. doutrina, que sustentamos, foi ensinada por \emph{Aristoteles}, \emph{Santo Agostinho}, \emph{Santo Thomas}, e por todos os philosophos christãos. \textsc{DEus}, dizem, elles, \emph{infunde a alma}, \emph{criando-a}; \emph{e cria} « \emph{alma}, \emph{infundindo-a}. --- Cf. \emph{Sum}. \emph{Th}., p. I, q. 90, a. 4; q. 118, \& 8; \emph{c}. \emph{Gentes}, 1. II, c. 83; \emph{De Pot}., q. 3, a. 10.\par D'esta doutrina afastaram-so \emph{Platão} e \emph{Leibnitz}.\par a) \emph{Platão} disse que as almas humanas, antes da sua união com os corpos, existiam nas estrellas; e, em seguida, por faltas commettidas, foram expulsas do ceu e encarceradas nos corpos, de que hão de sahir, depois da expiação, para voltarem para a sua antiga mansão. --- Alguns sequazes de \emph{Origenes} seguiram esta theoria.\par A opinião de \emph{Platão} é falsa pelas seguintes razões:\par a) Se as almas humanas \emph{preexistissem} aos corpos, seriam substancias \emph{completas}; pois tudo o que é produzido por \textsc{Deus}, e não tem relação com um composto, é substancia \emph{completa}; e por isso, sendo a alma uma substancia \emph{completa}, a sua união com o corpo seria \emph{accidental}. Ora a união da alma com o corpo, como demonstrámos, não é \emph{accidental}, mas \emph{substancial}.\par \emph{b}) Na doutrina de \emph{Platão}, a união da alma com o corpo não só seria \emph{accidental}, mas \emph{violenta} e \emph{contraria d}. \emph{natureza}; porque a alma, por si, não teria relação nenhuma com o corpo, mas estaria enearcerada nelle de um modo repugnante a todas as suas inclinações. E assim o homem, a mais nobre das creaturas visiveis, não seria uma natureza modelada pela intelligencia divina, mas representaria apenas o castigo infligido a um espirito mau.\par \emph{c}) Se a alma estivesse no corpo, como num carcere, \emph{o} homem nada desejaria mais do que a morte; pois sô a morte poderia restituir-Ihe a liberdade perdida. Ora o homem nada mais teme do que a morte.\par \emph{d}) Finalmente a \emph{preewistencia} das almas é uma affirmacaão gratuita e contraria ao testemunho da consciencia.\par b) \emph{Leibnits} pensou --- que as nossas almas foram creadas no principio do tempo e encerradas em certos corpusculos organicos, tendo apenas a sensibilidade, --- e »ue, quando esses corpusculos se desenvolvem para formar\par 24 Vor. II.\pgn{370}os corpos humanos, tambem as almas se desenvolvem e de sensitivas se tornam intellectuaes.\par A opinião de \emph{Leibnitz} é inacceitavel pelas seguintes razões:\par a) E' contraria é razão. A alma é forma substancial do corpo. Ora nenhuma forma pode ser produzida fóra do proprio sujeito. Logo a alma deve ser creada no proprio corpo. Diz S. Thomaz: « Actus proprius fit in potentia propria. Cum ergo anima sit proprius actus corporis, anima producta est in corpore » (\emph{Sum}. \emph{Th}., p. L, q. 90, a. 4).\par \emph{b}) Admitte que a alma de \emph{sensitiva} se transforma em \emph{intellectual}, --- transformação que, como demonstrámos contra \emph{Rosmini}, é absurda.\par \emph{c}) E' contraria é experiencia, a qual attesta que, pela geração, não se desenvolve um organismo já existente, mas se produz um novo organismo.\par A alma humana, pois, é creada no momento da sua união com o corpo.\par Mas poderá determinar-se esse momento?\par Os philosophos não estão de accordo com relação a este ponto. --- Uns dizem que a alma humana é creada por \textsc{Deus} no \emph{primeiro} momento da \emph{conceição} do corpo. --- Outros affirmam que a alma humana só é creada por \textsc{Deus}, quando o corpo está sufficiente e devidamente organisado. E a doutrina de S. Thomaz (\emph{Sum}. \emph{Th}., p. I, q. 118, a. 2; c. \emph{Gentes}. l. I1, c. 89).\par Seguimos a segunda opinião por estas razões:\par 1.) A \emph{materia} não póde receber a \emph{forma}, se não estiver convenientemente disposta em relação á mesma forma. Ora o corpo esta convenientemente disposto em relação á alma, quando organisado; pois que a alma é forma \emph{substancial} do corpo \emph{organisado}. Logo a alma é creada, quando o corpo está sufficiente e devidamente organisado.\par 2." Na geração do homem, devemos conceder ás causas segundas, que são os principios geradores, tudo o que se attribue aos principios geradores dos outros seres vivos, inferiores ao homem. Ora aos outros seres vivos, inferiores ao homem, attribuimos o poder de formar os organismos, sem uma \emph{especial} intervenção da Causa primeira. Logo os principios geradores humanos têm o poder de formar o organismo, e por isso a nossa alma é creada e unida ao corpo, só depois d'este chegar á devida organisação.\par 3.º) A experiencia demonstra que as disposições organicas dos paes se transmittem aos filhos, \emph{e} produzem uma certa semelhança entre estes e aquelles. Ora esta transmissão indica que á causa efficiente do organismo é a força do germen, e não é a alma humana. --- E' á força do germen; pois este, contendo \emph{virtualmente} as qualidades do organismo gerador, e operando segundo o impulso que este Ihe communiocou, tende para reproduzir no filho a semelhança dos paes. --- Não é a alma humana; pois esta, sendo creada immediatamente por \textsc{Deus} e não derivando da substancia dos paes, nao possue aquella tendencia. Logo o organismo humano é formado pela força do germen, e por isso a nossa alma é creada, quando o corpo está sufficiente e devidamente organisado.\par Diremos, pois, que o embryão, antes de formar o corpo humano, passa por tres differentes phases. Na primeira phase, o embryão é apenas dotado\pgn{371}de vida \emph{vegetetiva}, proveniente de uma forma substancial \emph{vegetativa}, que é produzida pela força do germen. Na segunda, é dotado tambem de vida \emph{sensitiva}, proveniente da \emph{forma} substancial \emph{sensitiva}, que é produzida pela força do mesmo germen, e que exerce tambem as funcções da forma \emph{vegetativa}. Na terceira, o corpo já devidamente organisado, é informado pela alma \emph{intellectual}, que é creada immediatamente por \textsc{DEus} e que exerce as funcções das duas formas inferiores, que são a \emph{vegetativa} e a \emph{sensitiva}.\par Quando se produz a forma \emph{sensitiva}, acaba a \emph{vegetativa}; e, quando é creada a alma \emph{intellectual}, acaba a forma \emph{sensitiva}; porque a forma superior contém \emph{virtualmente} a inferior, e porque a materia não póde ter ao mesmo tempo mais que uma forma \emph{substancial},\par Esta successão de formas \emph{substanciaes} no embryão humano é muito conforme com o modo por que a natureza opéra; pois é proprio d' esta passar \emph{gradualmente} do imperfeito e indeterminado para o perfeito e determinado.\par Nem se diga que a nossa theoria é uma reproducção da dos \emph{transformistas}, a qual admitte que uma especie superior passa, no seu desenvolvimento, por todas as especies inferiores. Porquanto, os \emph{transformistas} suppõem que uma especie superior, no seu desenvolvimento, é \emph{indeterminada} e por isso constitue tantas especies inferiores, quantas são as suas differentes pausas; ao passo que, no nosso systema, o embryão humano é \emph{determinado desde o seu principio}, e, embora adquira as formas \emph{vegetativa} e \emph{sensitiva}, todavia nunca é \emph{planta}, nunca é uma \emph{especie} de \emph{animal irracional}, mas é sempre o embryão humano, que passa pelas formas inferiores ató chegar ao seu, completo desenvolvimento.\par Alguns escriptores não admittem esta successão de formas \emph{substanciaes} no embryão humano, porque temem ser desmentidos pela \emph{physiologia}.\par O receio d'esses escriptores é vão. A \emph{p/Aysiologia} nunca os desmentirá. Por quanto, esta sciencia só diz que o embryão, desde o começo do seu desenvolvimento, é vivo. Ninguem o nega. Mas a questão é outra. Sendo a vida do embryão, no seu começo, \emph{vegetativa} (pois manifesta-se apenas por uma crescente multiplicação de cellulas), resta saber se essa vida deriva de uma forma \emph{vegetativa}, ou da alma \emph{intellectual}. \emph{A phiysiologia}, pois, nada póde dizer relativamente a esta questão.\par \emph{Fredault}, no seu tratado \emph{Physiologie générale}, diz que alma, devendo fabricar-se o organismo, adaptado ás proprias operações, deve ser creada e unida ao corpo no primeiro momento da conceicio, antes da formação do proprio organismo.\par Resposta. A alma é causa \emph{formal} do organismo, e não causa \emph{efficiente}. Ora não é a causa \emph{formal} a que forma o sujeito; mas o sujeito é formado e predisposto por uma energia diversa, que é verdadeira causa efficiente, e só quando esta disposição tem chegado ao seu ultimo complemento, o sujeito se torna capaz de ser actuado pela causa \emph{formal}. \emph{A forma} não é o \emph{principio}, mas é o \emph{termo} da geração. Diz S. Thomaz: « Generatio non sequitur, sed praecedit formam substantialem » (\emph{De pot}.,, q. 3, a. 12). \emph{A}. geração tende para a forma, como para o seu termo, e essa tendencia continua até a organisacüo do sujeito; de modo que a geração comprehende todo o processo previo, necessario e sufficiente para a formação do organismo. A causa \emph{efficiente}\pgn{372}, que construe o organismo, é a \emph{energia formativa}, communicada ao semen pela alma do principio gerador. Essa energia, que é encerrada no semen, é causa efficiente \emph{instrumental}; e a alma do principio gerador, que communica essa energia ao semen, é causa efficiente \emph{principal}. Por isso diz S. Thomaz: « In semine non est anima, sed virtus animae » (\emph{De pot}., q. 3, a. 9 ad 9). --- Isto, e só isto, explica a semelhança entre paes e filhos nas feições e nas qualidades do organismo e, por meio d'este, nas qualidades intellectuaes e moraes (\emph{Sum}. \emph{Th}., p. 1, q. 85, a. 7). Porquanto, sendo, como o proprio \emph{Fredault} observa, o organismo do sujeito gerado uma reproducção, mais ou menos perfeita, do organismo do principio gerador, como poderia verificar-se essa conformidade, essa semelhança, se a causa efficiente da elaboração do organismo não fosse uma energia derivada do principio gerador, mas fosse a alma racional do sujeito gerado? E pode sequer conceber-se que a alma, no formação do organismo, opere como um architecto, delineando um plano, que reproduza o organismo dos paes, e executando-o com tenacidade de vontade? Não, por certo; porque a alma, como dissemos, é causa \emph{formal}, \emph{e} não \emph{efficiente}, do corpo, com que está unida. --- Notamos ------ que essa energia \emph{formativa} acaba depois de ter produzido o seu effeito, que é a organisacao do corpo (\emph{Sum}. \emph{Th}., p. I, q. 118, a. 1 ad 4), --- e que a alma, informando o corpo succede não a essa energia, mas ao previo principio activo, que exercia o officio de forma n'aquelle organismo. --- Elaborado pela energia \emph{formativa} e unido com a alma racional, o organismo desenvolve-se e augmenta sob a acção da propria alma.\par Coneluimos con S. Thomaz: « Anima intellectiva creatura Deo \emph{in fine generationis hwmanae}, quae simul est et sensitiva et nutritiva, corruptis formis praeexistentibus » (\emph{Sum}. \emph{Th}., p. I, q. 118, a. 2ad 2). --- Mas, quando é que a geração chegou ao seu \emph{termo}? Os physiologistas não estão de accordo sobre este punto e nunca o estarão. Por isso nunca poderá determinar-se o momento, no qual \& alma é creada e unida ao corpo. --- Em tal estado de coisas, a Egreja, Mãe carinhosa, que segue sempre a via mais segura, quando se trata da salvação de seus filhos, ordena sabiamente que se administre o santo Baptismo ao feto que dé signaes de vida, embora não se saiba se essa vida deriva da alma espiritual ou de uma forma inferior.}.

%%>>> art-32

\section*{Artigo III.\\ Origem do composto humano.}
\addcontentsline{toc}{section}{Artigo III. --- Origem do composto humano.}

{\small\noindent\textsc{Summario}: --- Fim do homem. --- Immortalidade da alma humana. --- Estado e operações da alma separada.\par}

%%>>> t225

\tese{225}{O composto humano tem uma origem.}
O \emph{composto humano}, ou o \emph{homem}, resultando da união da alma \emph{intellectual} e \pg{369} \pg{370}  do corpo \emph{organico}, deve ser posterior a estes dois elementos componentes; pois que o todo é sempre posterior ás suas \pg{371}  partes. Ora a alma e o corpo, como demonstrámos, têm uma origem. Logo o \emph{composto humano} tem uma origem.

A apparição do homem sobre a terra é um facto \emph{relativamente} recente, como se demonstra principalmente pelas ultimas descobertas da \emph{paleontologia} \footnote{As ultimas descobertas da \emph{paleontologia} vieram confirmar a antiga tradição que a apparioüo do homem sobre a terra remonta a uma epocha, comprehendida entre 5.000 e 10.000 annos.\par Todavia os \emph{transformistas}, baseando-se no principio de que a evolução de uma especie inferior para a superior é muitissimo lenta e por isso exige um grandissimo numero de annos e de seculos, concluem que a existencia do homem sobre a terra remonta a uma antiguidade, comprehendida entre os 20.000 e os 800.000 annos.\par \emph{Haeckel} diz: « As innumeras e importantes descobertas, feitas nestes ultimos annos, attestam um facto \emph{capital}: que a existencia do genero humano remonta \emph{certamente} a mais de 20.000 annos. Mas mais de 100.000 e talvez centenares de milhares de annos ja passaram depois da origem do homem sobre a terra... Embora fizessemos remontar a existencia e a dispersão do\pgn{373}homem sobre a terra a 200.000 annos, a um numero qualquer de centenares de milhares de annos, tudo isto seria um espaço de tempo infinitamente breve em relação á incommensuravel duração, exigida pela gradual evolução da longa serie dos antepassados do homem » (\emph{Hist}. \emph{de la creat}. \emph{nat}.).\par \emph{Darwin} diz: « A remotissima antiguidade do homem foi posta em evidencia por homens insignes, começando por Bucher de Perthes, e \emph{esta} é a \emph{base necessaria para se comprehender a sua origem} » (\emph{A origem do homem}).\par Do mesmo modo se exprimem outros \emph{transformistas}, como \emph{Zaberowshi}- \emph{Moindron}, W. \emph{de Fonvielle}, \emph{V}. \emph{de Meunier}, \emph{Vogt}, \emph{Rolle}, \emph{Büchner}, \emph{Broca}, \emph{Conestrini}, \emph{Bunsen}, \emph{Beverley Randolph}, etc.\par *\par Os argumentos, que os \emph{transformistas} apresentam para provar a remotissima antiguidade do homem, reduzem-se a cinco classes, e são o \emph{paleontologico}, o \emph{anthropologico}, o \emph{geologico}, o \emph{archeologico} e o \emph{chronologico}. --- De cada um d'estes argumentos faremos uma breve critica.\par a) O argumento \emph{paleontologico} é o seguinte: O homem primitivo viveu ao mesmo tempo, em que viviam certos animaes já extinctos. Ora a extincção \emph{successiva} d'esses animaes, sendo muito lenta, só poude effectuar-se num incommensurayel espaço de tempo. Logo a existencia do homem sobre a terra é muito mais antiga do que narram as historias,e por isso o \emph{Homem prehistorico} é um facto incontestavel.\par A \emph{maior} do syllogismo é verdadeira, mas a \emph{menor} é falsa, e por isso é falsa tambem a \emph{consequencia}.\par \emph{a}) A \emph{maior} é verdadeira. Por quanto, é certo que hoje não vivem sobre \& terra muitas especies de animaes, que viveram outr'ora. Entre essas especies contam-se o \emph{ursus spelaeus}, \emph{á hyaena spelaea}, \& \emph{tigris spelaea}, o \emph{elephas primigenius}, o \emph{rhinoceros tichorhinus}, o \emph{hyppopotamus maior}, o \emph{alces hibernicus}. Mesmo em nossos dias extinguiram-se duas aves --- o \emph{didus} e o \emph{dinornis}, que existiam no seculo XVII, encontrando-se aquelle nas ilhas \emph{AMauricia}, \emph{Bourbon} e \emph{Rodrigo}, e este na \emph{Nova Zelandia}.\par 5) E' falsa a \emph{menor}, isto é, que a extinccao d'aquelles animaes devesse ser muito \emph{lenta} e \emph{successiva}, e por isso só podesse effectuar-se num incommensuravel espaço de tempo. Como podem os \emph{transformistas} provar que aquelle desapparecimeuto foi \emph{lento} e \emph{successivo}, e que não foi obra de um cataclismo repentino, ao menos em relação á maior parte dos animaes?\par Tudo leva a admittir, como verdadeira, esta segunda parte da \emph{disjunctiva}.\par O diluvio, cuja existencia é attestada --- não só pela Biblia Sagrada, que, embora se prescinda da sua inspiração divina, é o livro mais antigo, mais authentico e mais digno de fé, --- mas tambem pela constante e.universal tradição de todos os povos e pelas descobertas geologicas, o diluvio, dizemos, é a verdadeira causa d'aquelle desapparecimento.\par O grande \emph{Cuvier} escreve: « Penso com \emph{Deluc} e \emph{Dolomieu} que, se alguma coisa em geologia se acha indubitayelmente justificada, é que a superficie do nosso globo foi victima de uma grande e subita revolução, cuja data não póde ir além de cinco a seis mil annos. Esta revolução abaixou e fez desapparecer os paizes, onde morayam os homens e os animaes hoje desconhecidos;\pgn{374}e, pelo contrario, seccou o fundo do mar, formando o terreno, que presentemente habitamos. O pequeno numero de homens, que escapou áquella destruição, occupou o novo terreno e propagou-se nelle » (\emph{Discours sur les révolutions du globe}). Do mesmo modo fallam \emph{Lambert} (\emph{Le déluge Mosaique}), \emph{Gainet} (\emph{Accord de la Bible et de la géologie}), e outros. --- As recentes descobertas, feitas nas grutas de \emph{Engis} de \emph{Neanderthal}, mostram abertamente que aquelles restos de homens e de animaes extinctos, que se encontraram misturados, foram alli lançados pela violencia das aguas.\par Mas, embora aquellas especies tivessem desappareeido \emph{successivamente}, é certo que para isso era mais que sufficiente o espaço de \emph{oito} ou \emph{des} mil annos, o que os sabios geralmente admittem; tanto mais que as especies extinctas são relativamente poucas, e o seu desapparecimento, embora fosse \emph{successivo} emquanto ao \emph{modo} por que se operava, podia ser \emph{simwulianeo} emquanto ao \emph{tempo} em que se operava.\par Os \emph{transformistas} respondem que só para o desapparecimento de uma especie eram necessarios muitos seculos. --- Mas esta resposta é gratuita, pois que não adduzem prova alguma que a confirme.\par b) O argumento \emph{anthropologico} é o seguinte: Entre os fosseis encontraram-se craneos humanos, cuja estructura é tão differente da dos nossos craneos, que manifesta o typo de um homem muito inferior, tanto nas qualidades physicas como nas intellectuaes, ao homem moderno. Ora, para que o homem se elevasse d'aquelle infimo estado ao grau perfeito de \emph{humanidade}, em que se encontra hoje, devia passar muitissimo tempo, --- pelo menos o espaço de 100.000 annos; pois a evolução é lentissima. Logo existiram verdadeiramente \emph{homens prehistoricos}.\par Este argumento é falso na \emph{maior}, na \emph{menor} e na \emph{conclusão}.\par \emph{a}) E' falso na \emph{maior}. Porquanto, como acima demonstrámos, é falso que os craneos dos homens antigos sejam \emph{especificamente} diversos dos dos homens modernos. Todos os caracteres, que os transformistas apontam como \emph{prehistoricos}, são \emph{historicos} e \emph{historico-contemporaneos}, como provaram os melhores naturalistas, \emph{Quatrefages}, \emph{Caruel}, \emph{Calori} e muitos outros.\par 5) E' falso na \emph{menor}. Com effeito, demonstrámos que o homem não passou nem podia passar do estado selvagem para o estado de civilisacüo.\par \emph{c}) E' falso na \emph{conclusão}. Porquanto, uma assergão, que se funda sobre suppostos falsos, não póde deixar de ser falsa.\par o) O argumento \emph{geologico} é o seguinte: A \emph{geologia}, nas suas descobertas, encontrou, nas mais profundas camadas, restos de corpo humano. Ora, para \& formação d' essas camadas, foram necessarios centenares de milhares de annos. Logo a existencia do homem remonta a uma antiguidade de centenares de milhares de annos.\par Antes de tudo devemos advertir que, nas mais antigas e profundas camadas, não se encontram fosseis humanos. O homem não podia viver no periodo \emph{glacial}, como confessam \emph{Lyel} e \emph{Agassis}. Sô na epocha \emph{quaternaria}, ou postpliocena, se encontram restos humanos em estado de \emph{fosseis}.\par Ora a espessura e a natureza dos materiaes, que formam as camadas, não indicam que a existencia do homem sobre a terra conte centenares de milhares de annos. Porquanto, a geologia, se descobre o numero e o nome\pgn{375}de muitos elementos, que compõem a superficie terrestre, não determina nem póde determinar a edade da formação de cada estrato, e, por conseguinte, não póde marcar o tempo em que viveram os homens, cujos restos se encontram nos mesmos estratos.\par Dizem: Embora não possamos determinar exactamente o tempo, em que se formaram as camadas, onde se encontram restos humanos, comtudo devemos dizer que esse tempo abrangeu grandissimo numero de seculos.\par Respondemos que tudo isto é gratuito.\par Quem sabe se-os sedimentos ou depositos são effeitos de uma acção lentissima, ou de uma revolução repentina? Os geologos narram que, ha poucos annos, o rio Pó (\emph{Italia}) arrastou, em breve tempo, ao mar cerca de 4.000 metros cubicos de terra, --- que o \emph{Ganges} arrebatou ás montanhas do \emph{Himalaya} tanta materia, quanta bastaria para reconstruir sessenta vezes a mais alta Pyramide do Egypto, --- que são frequentes as elevações e depressões de terreno, devidas aos terremotos e outras causas.\par E, sendo isto verdade, poderão os transformistas provar que as camadas da epocha \emph{quaternaria} não foram formadas por uma inundação ou por uma repentina elevação de terreno, tanto mais que as camadas d'aquella epocha apresentam apenas uma altura de á a 12 metros, e que a actividade dos agentes exteriores na epocha \emph{quaternaria} era extraordinaria? Estamos convencidos de que nunca o poderão provar.\par Em confirmação do que temos dicto, citaremos o testemunho de \emph{Lapparent}, cuja competencia é reconhecida. Este sabio \emph{geologo} escreve:\par « A parte da era moderna, que se denomina \emph{epocha quaternaria}, é assignalada pela apparicüo do homem sobre o globo. Depois d'este grande facto, o mundo organico deixou de se enriquecer de novas especies, mas muitas formas desappareceram ou emigraram. A epocha \emph{quaternaria} distingue-se dos tempos actuaes por uma actividade extraordinaria de precipitações atmosphericas..... Embora tão proxima da nossa, esta epocha é ainda hoje um mysterio para nós. A successão dos depositos é muito escura. A ausencia ou \& escassez de restos organicos tornam sensivelmente delicada e difficil a determinação da edade relativa....\par « A epocha \emph{quaternaria} está caracterizada por uma actividade inteiramente extraordinaria dos agentes exteriores. Julgar, pelo que se passa hoje, do tempo que foi necessario para desembaraçar certos valles, entulhados de depositos moveis, seria esquecer que as precipitações atmosphericas eram, pelo menos, dez ou vinte vezes mais abundantes do que em nossos dias...\par « A sciencia não póde ainda encontrar um chronometro, para medir o tempo passado, mesmo no periodo immediatamente anterior ao nosso... Para nós, basta ter provado até que ponto são destituidos de base rigorosa todos esses calculos, que distribuem generosamente centenares e milhares de seculos entre as diversas phases da \emph{epocha quaternaria}. Nos factos geologicos da \emph{epocha quaternaria} nada se encontra, que possa motivar as apreciações de certos auctores » (\emph{Traité de géologie}, pag. 1114-1119).\par d) O argumento \emph{archeologico} é o seguinte. A \emph{archeologia}, revelando com as-suas descobertas as obras dos homens primitivos, mostra o tempo ou a data da existencia d'elles sobre a terra. Esse tempo póde dividir-se em tres\pgn{376}epochas: da \emph{pedra}, do \emph{bronze} e do \emph{ferro}. Ora, para o homem passar de uma epocha para outra, deviam decorrer inanumerayeis seculos, pois a evolução da civilisação humana foi lentissima. --- \emph{Molot} e \emph{Canestrini} dão á epocha do \emph{bronze} uma antiguidade de 28 a 42 seculos, e á da \emph{pedra} uma antiguidade de 47 a 70 seculos.\par Respondemos a este. argumento dos \emph{transformistas}, provando ser falso --- que o uso da \emph{pedra} seja o signal de uma antiquissima existencia do homem sobre a \emph{terra}, --- e que o uso da \emph{pedra}, do \emph{bronse} e do \emph{ferro} indique tres epochas \emph{successivas} e \emph{distinctas}.\par a) E' falso que o uso da \emph{pedra} seja signal da antiquissima existencia do homem sobre a terra. Porquanto, se assim fosse, se o uso da \emph{pedra} fosse o caracter distinctivo dos tempos \emph{preAistoricos}, esse uso teria desapparecido nos tempos \emph{Historicos}. Ora não é assim. A historia attesta que os povos historicos, como os lsraelitas (\emph{Ex}., IV, 25 --- \emph{Jos}., V, 2), os Egypeios (Herodoto, II, 86 --- Diodoro, I, 91), os Romanos (Livio, I, 54), os Indios (Aldovrandi, \emph{Mus}. \emph{Metall}. os Persas (Herod., \emph{Polymnia}) os antigos Allemães (Diod. Sic. \emph{Rei}. \emph{ant}. l. IV), os actuaes habitantes do \emph{Kamtscialha}, da \emph{Nova Hollanda}, das ilhas \emph{Marianas}, de \emph{Sandwich}, da \emph{Oceania}, e de muitos outros paizes, faziam e fazem uso da \emph{pedra}.\par De resto, o proprio \emph{B/chner}, não sabemos com quanta coherencia, sustenta a nossa opinião com as seguintes palavras: « Ào passo que o \emph{bronze} foi prompta e completamente substituido pelo \emph{ferro}, as armas de \emph{pedra} duraram muito tempo, e o seu uso, como dissemos, prolongou-se até aos tempos \emph{historicos}. Às ultimas armas de \emph{pedra} foram certamente lavradas com instrumentos de \emph{ferro}. Lé-se que os Bretóes combateram com estas armas contra \emph{Guilherme}, o conquistador ». (\emph{O homem considerado segundo a sciencia}, \emph{p}. 1).\par Nem se diga que nos tempos \emph{Aistoricos} o homem só se serviu da pedra \emph{polida} e que nos tempos \emph{prehistoricos} só fez uso da pedra \emph{lascada}. Por quanto, nos mesmos logares e na mesma camada, encontraram-se e encontram-se misturados, uns com os outros, instrumentos informes e perfeitos de. pedra \emph{lascada} e de pedra \emph{polida}; o que não indica uma successão chronologica dos differentes typos de \emph{pedra}.\par Coneluimos que o uso da \emph{pedra} não é o signal, que manifeste a \emph{antéquissima} existencia do homem e a sua rudez primitiva.\par \emph{b}) E' falso que o uso da \emph{pedra}, do \emph{bronze} e do \emph{ferro}, indique tres epochas \emph{successivas} e \emph{distinctas}. Por quanto, como dissemos, desde a mais remota antiguidade, em que, segundo os \emph{transformistas}, só devia conhecer-se a \emph{pedra}, conheciam-se, ao mesmo tempo, \emph{o bronse} e o \emph{ferro}, e fazia-se uso d'estes metaes. Assim \emph{Z'ubalcain}, oitavo descendente de \emph{Adão}, era habil artista em \emph{bronze} e em \emph{ferro}; os \emph{phenicios} tambem serviam-se do \emph{bronze} e do \emph{ferro}; os \emph{gregos}, na guerra de Troia, tinham armas de \emph{bronse} e de \emph{ferro}; as armas do exercito de \emph{Xerves}, como diz Herodoto, eram tambem de \emph{bronze} e de \emph{ferro}.\par Se, pois, nos tempos \emph{Astoricos}, o uso da \emph{pedra}, do \emph{bronse} e do \emph{ferro}, não foi \emph{successivo} mas \emph{simultaneo}, devemos concluir --- que a theoria das tres epochas, \emph{successivas} e \emph{distinctas}, não só não tem base solida, mas é desmentida pelos factos, --- e que a longuissima serie de seculos, em queo homem teria passado de uma epocha para outra, é um mero producto da imaginação.\par \textasciicircum{}}.

996. O \textbf{composto humano não deriva} do \textbf{concurso de causas fortuitas}. --- Se o \emph{composto humano} derivasse do concurso de causas \pg{374}  fortuitas, tambem d'esse concurso deveriam derivar a alma e o corpo. Ora o consequente é falso; porque as causas fortuitas, \pg{375}  --- por serem a negação de toda a \emph{ordem} e de toda a \emph{finalidade}, não podem produzir o nosso corpo, em que tudo é harmonia \pg{376} , tudo está disposto para um fim certo e determinado, --- e, por serem finitas e limitadas, não podem \emph{criar} uma substancia \pg{377}  \emph{espiritual}, como é a nossa alma. Logo o \emph{composto humano} não deriva do concurso de causas fortuitas \footnote{\emph{Du Bois-Reimond}, que é tido por uma celebridade scientifica, escreve:\par « \emph{Imaginemos} que todos os atomos, que constituiam Cesar num determinado momento da sua vida, por exemplo, no Rubicaão, sejam com arte mechanica collocados no seu logar e que nelle se imprima á velocidade e a direcção, que tinham naquelle momento. \emph{Para nós}, reappareceria Cesar em alma e corpo. O Cesar artificial teria, desde o primeiro momento da sua formação, as mesmas sensações, os mesmos desejos, os mesmos pensamentos, que o seu modelo tinha no Rubição, a sua memoria seria cheia das mesmas imagens; seria dotado das mesmas faculdades, herdadas por geração ou adquiridas. \emph{Imaginemos} que o mesmo trabalho mechanico se renove ao mesmo tempo, em varios pontos, com os mesmos atomos de \emph{carbone}, de \emph{Hydrogenio}, etc. Estes novos Cesares distinguir-se-hiam entre si, só por causa do logar, onde foram formados » (\emph{Les bornes de la philosophie naturelle}).\par Aquelle « \emph{para nós} » dispensa qualquer commentario.}.

297. O \textbf{composto humano não deriva da evolução de} uma \textbf{especie inferior}. --- O \emph{composto humano} resulta da união da alma e do corpo. Ora a nossa alma e o nosso corpo não derivam da evolução de uma especie inferior. Logo o \emph{composto humano} não deriva da evolução de uma especie inferior.

e) O argumento \emph{chronologico} é o seguinte: Os annaes dos povos do Oriente e da Chaldêa, a historia da India e da China, provam que a existencia do homem sobre a terra remonta a centenares de milhares de annos.

Respondemos ás diversas partes do argumento.

\emph{a}) Com relação aos \emph{Egypcios} e á sua pretendida antiguidade, \emph{Champollion} diz: « Demonstro que nenhum monumento egypcio é realmente anterior ao \&uno 2.000 antes da era christà ». --- Nem os adversarios appellem para os famosos zodiacos de \emph{Denderah} e de \emph{Esneh}. Por quanto, --- com relação ao zodiaco de \emph{Denderah}, insignes sabios, como \emph{Laland}, \emph{Hamilton}, \emph{Delambre}, \emph{Biot}, \emph{Chamgpollion} e outros, demonstraram que elle não passaya da epocha dos \emph{Ptolomeus} (820 a 805 annos antes de Christo), --- e com relação ao zodiaco de \emph{Esneh}, \emph{Latronne} e \emph{Bankes} provaram que elle foi construido nos primeiros annos de \emph{Antonino Pio}, imperador romano.

\emph{b}) Com relação ás observações astronomicas dos \emph{CAaldeus}, quo mostrariam que o homem existe ha mais que 470.000 annos, diremos, com \emph{Laplace}, que essas observações não excedem de 800 annos a era christa.

c) Com relação aos \emph{Indios}, \emph{William Jones}, seguido por \emph{Heeren}, afirma que a origem d'aquelle imperio se deve collocar cerca de 2.000 annos antes da era vulgar. --- Nem recorram ás famosas tabuas astronomicas dos mesmos \emph{indios}, para provar a antiguidade d'esse povo. \emph{A/aprot}, nas suas \emph{Memorias da} Asia, escreve: « Às tabuas astronomicas dos \emph{indios}, \& que se attribuia uma prodigiosa antiguidade, foram feitas no seculo VII da era vulgar, e posteriormente referidas, por caleulos, a uma epocha anterior ».

\emph{d}) Com relação aos \emph{Chinezes}, \emph{Klaprot} nega que haja certeza historica, Acerca d'esse povo, anterior ao anno (32 antes da era christã, --- \emph{Abel Remusat} admitte a \emph{certeza} nos factos historicos até ao anno 2.260, e reconhece só \emph{probabilidade} nas tradições que remontam ao anno 2.637 antes de Christo.

%%>>> t228

\pg{372} \pg{373} \pg{378} \tese{228}{O "primeiro composto humano foi formado immediatamente por Deus.}
Se \textsc{Deus} formou \emph{immediatamente} o corpo do \emph{primeiro} homem, e \emph{criou} a sua alma, deve dizer-se que o \emph{primeiro} composto humano foi formado \emph{Immediatamente} por \textsc{Deus}. Ora \textsc{Deus} formou \emph{mmediatamente} o corpo do \emph{primeiro} homem, e \emph{criou} á sua alma. Logo o \emph{primeiro} composto humano foi formado \emph{immediatamente} por \textsc{Deus} \footnote{\emph{Mivart}, seguido por alguns \emph{transfurmistas}, admitte que o corpo do \emph{primeiro} homem derivou do macaco; mas que neste corpo, que uma evolução progressiva, \emph{devida ao simples jogo de causas naturaes}, tornou proporcionado para o homem, \textsc{Deus} infundiu uma alma racional.\par Esta opinião é inacceitavel pelas seguintes razões:\par \emph{a}) Admitte que as especies se transformam, a inferior na superior. Ora tal transformação, como demonstrámos, é falsa.\par 5) Attribue a transformação do corpo do macaco em corpo humano ao \emph{simples jogo das causas naturaes}. Ora esta. transformação é absurda: porque o organismo humano apresenta, desde o seu desenvolvimento, uma \emph{finalidade}, isto é, uma adaptação de cada orgão ao seu fim respectivo. Esta \emph{finalidade} não pode provir do \emph{acaso}, do \emph{simples jogo de causas naturaes}.\par \emph{c}) Se fosse verdadeira a theoria de \emph{Mivart}, deveriam encontrar-se os typos \emph{intermediarios} entre os macacos, de que teria derivado o corpo do homem, e o proprio homem, --- typos que deviam ser semelbantes ao corpo humano, mas que não deviam ser dotados de razão. Ora esses typos nunca se encontraram. Pelo contrario, onde a sciencia encontrou o homem \emph{anatonico} fossil, encontrou tambem vestigios da vida racional.\par \emph{d}) Aquelle grupo de animaes, que teriam sido antepassados \emph{émmediatos} do corpo humano, seria uma verdadeira anomalia na natureza viva. Por quanto, o organismo d'aquelles animaes suppõe-se adaptado para as funcções da vida intellectual, mas destituido de vida intellectual. Ora entre o organismo e as faculdades \emph{psychicas} deve existir uma proporção; porque a nossa vida \emph{psychica} depende, de algum modo, do organismo.}.

%%>>> t229

\tese{229}{0 composto humano, no estado actual, deriva de outro homem e de Deus, mas por modo diverso.}

a) Deriva de outro homem e de \textsc{Deus}. Porquanto, o homem, concorrendo para a formação do corpo, concorre para a producção do \emph{composto humano}; e \textsc{Deus}, ereando a alma e infundindo-a no corpo, produz o mesmo composto.

b) Por \emph{modo diverso}. Com effeito, --- na \emph{creação} da alma, \textsc{Deus} opéra como Causa \emph{primeira} e \emph{unica}, e o homem não concorre nem sequer como causa \emph{instrumental} ; --- na producção do \emph{corpo}, \textsc{Deus} opéra como Causa \emph{primeira}, e o homem como causa \emph{segunda}, em virtude da actividade, que elle recebeu do Auctor da natureza \footnote{Respondamos a algumas objecções dos adversarios.\par a) Dizem: O filho diz-se \emph{gerado} pelo pae. Ora o filho não se poderia\pgn{379}dizer gerado, se o pae não gerasse a alma e o corpo do filho. Logo o pae gera tambem a alma do filho.\par Resposta. O \emph{termo} da geração não é o corpo, nem a alma, mas é o \emph{composto}; e é do \emph{termo} que tira a sua denominação todo e qualquer movimento. Logo quem concorre, como causa efficiente, para a geração do composto humano, pode chamar-se e é verdadeiro \emph{pae}. Ora as pessoas, ás quaes damos o nome de \emph{paes}, concorrem, como causas efficientes, para a geração do composto humano, porque, embora só Deus produza a alma, os paes comtudo dispõem proximamente a materia para a recepção da alma, e como que trazem a alma para o corpo, e por isso são causa efficiente da união da alma com o corpo, isto é, da geração do composto. E' isto basta para que possa dizer-se que os paes geram o ser dos seus filhos, ainda que não tenham gerado a alma. Portanto pode dizer-se com razão que gera o /\emph{iomem} quem concorre para a formação do composto humano, embora não gere a alma, assim como dizemos que mata o \emph{homem} quem concorre para destruir esse composto, embora não mate a alma. --- Em poucas palavras: « A \emph{geração} tem por termo o \emph{composto}. Ora o homem concorre para a formação do composto. Logo o homem \emph{gera} o homem, embora não gere a alma. Diz S. Thomaz: « Homo \emph{generat} sibi simile, in quantum per virtutem seminis ejus disponitur materia ad susceptionem talis formae » (\emph{Sum}. \emph{T}77., p. I, q. 118, a. 2 ad 4).\par 5) Insistem: E' um facto que os filhos se parecem sempre com os paes não só nas qualidades physicas e physiologicas, mas tambem nas qualidades intellectuaes e moraes. Ora esta semelhança nas qualidades intellectuaes e moraes seria inexplicavel, se a alma dos filhos não fosse gerada pelos paes.\par Resposta. Não é verdade que os filhos se parecem sempre com os paes tambem nas qualidades intellectuaes e moraes. Mas, embora assim fosse, não poderiamos d'ahi concluir que a alma dos filhos deriva dos paes. Na verdade, ninguem ignora que a vida vegetativa e a sensitiva dependem do organismo physico, e que a propria vida psychica ou intellectual, embora independente intrinsecamente do corpo, comtudo suppõe e exige o concurso dos orgão corporeos. Ora o corpo dos filhos é gerado pelos paes. Não admira, pois, se os filhos se parecem com os paes --- não só nas qualidades physicas e physiologicas, --- mas tambem nas intellectuaes e moraes. E dá-se esta semelhança --- não porque os paes produzem a alma dos filhos, --- mas porque. produzem o corpo, que, como dissemos, concorre tambem, de algum modo, para a nossa vida psychica ou intellectual.\par €) Acerescentam : Um só effeito não pode ser produzido por varios agentes. Ora o homem é um só ente. Logo o homem não pode derivar, ao mesmo tempo, dos paes e de Deus.\par Resposta. Um e o mesmo effeito não pode ser produzido por varios agentes, quando estes não são coordenados entre si, mas pode ser produzido, quando entre os agentes existe uma coordenação. Ora Deus e os paes são agentes coordenados; porque Deus opera como Causa primeira, e os paes como causas segundas. Os paes subministram a materia e a preparam gradualmente até á ultima disposição, que tem por termo a infusão da alma,\par n\pgn{380}e param ahi, porque a infusão da alma excede a força creada. Mas, se param as causas segundas, não para a Causa primeira, a qual continua só e produz o effeito, que as causas secundas não podiam produzir, isto ó, crea a alma humana. --- E, sendo a energia creada para Deus, como o instrumento é para o supremo e principal agente, devemos dizer que o corpo humano é formado pela força de Deus, como agente primario, e pela energia do homem, como agente secundario, ao passo que a alma é produzida unicamente pela força divina e é produzida na materia, previamente disposta pela energia humana (Cf. S. Th., \emph{c}. \emph{Gentes}, l. 2, c. 89). Esta coordenação da Causa primeira e da causa segunda explica a unidade admiravel do homem.\par \emph{d}) Replicam: Dois agentes, embora coordenados, devem. produzir um duplice effeito. Ora o homem tem um só \emph{ser}. Logo o homem não pode derivar de Deus e de outro homem.\par Resposta. Embora Deus e o homem concorram, cada um na sua ordem, para a producção do homem, comtudo não ha no homem gerado um duplice ser, --- um que derive do principio gerador e outro do principio creador, mas ha apenas um só \emph{ser}; porque o Creador dà o ser á alma no corpo, e o principio gerador dispõe o corpo para que possa participar d'esse ser pela sua união com a alma. Diz S. Thomaz: « In homine non est duplex esse, quia non est sic intelligendum corpus esse \emph{separatím} a generante et animam a creante, quasi corpori acquiratur esse separatim a generante et separatim animae a creante; sed quia \emph{Creans dat esse animae in corpore}, \emph{et generans disponit corpus ad hoc}, \emph{quod hujus esse sit particeps per animam sibi wunitam} » (\emph{De pot}., q. 9, a. 9 ad 20).}.

%%>>> t230

\pg{379} \tese{230}{Todos os homens pertencem á mesma especie.}

r

a) Todos os homens pertencerão á mesma especie, se todos possuirem os mesmos elementos \emph{essenciaes} e differirem apenas nos caracteres \emph{accidentaes}; porque só os elementos \emph{essenciaes} \pg{380}  constituem a \emph{especie}. Ora todos os homens possuem os mesmos elementos \emph{essenciaes} (\emph{animalidade} e \emph{racionalidade}), e só differem nos caracteres \emph{accidentaes} (\emph{conformação do craneo}, \emph{cor}, ete.). Logo todos os homens pertencem á mesma \emph{especie} \footnote{Proecuraremos demonstrar que \emph{todos} os homens \emph{convém} entre si nos elementos \emph{essenciaes}, e que só \emph{divergem} nos caracteres \emph{accidentaes}.\par a) Todos os homens \emph{convém} entre si nos elementos \emph{essenciaes}. Os elementos \emph{essenciaes}, que constituem o homem, são a \emph{animalidade} e a \emph{racionalidade}. Ora todos os homens serão dotados de \emph{animalidade} e \emph{racionalidade}, se todos possuirem, emquanto ao \emph{essencial}, os mesmos caracteres, tanto \emph{anatomicos} e \emph{physiologicos} (provenientes da \emph{animalidade}). como \emph{psychologicos} e moraes (provenientes da \emph{racionalidade}). Mas effectivamente todos os homens possuem, emquanto ao \emph{essencial}, os mesmos caracteres, tanto \emph{anatomicos} e \emph{physiologicos}, como \emph{psychologicos} e \emph{moraes}.\par \emph{a}) Todos os homens possuem os mesmos caracteres \emph{azatomicos} e \emph{physiologicos}. Por quanto, todos os homens, seja qual for a sua raça, têm a mesma organisacüo corporea (\emph{emquanto ao essencial}), a mesma posição, \emph{o} mesmo numero de dentes, o mesmo contorno craneano. O cerebro de um negro do Congo não differe do de um europeu; as dobras de todos os cerebros humanos estão ligadas ao mesmo typo fundamental. Em todas as racas o organismo passa pelas mesmas phases, está sujeito ás mesmas doenças e ás mesmas vicissitudes e circumstancias, que costumam dar-se exclusivamente entre os seres da mesma especie.\par \emph{T'lourens}, na Academia de Paris, pronunciou as seguintes palavras: « Os\pgn{381}homens, seja qual for a sua casta, brancos ou negros, vermelhos ou amarellos, todos têm proporcionalmente a mesma capacidade de craneo. Nem o cerebro apresenta differenças; pelo contrario, o cerebro humano differe do do orangotango na forma, no peso, no volume, nas cireumyoluções ».\par \emph{b}) Todos os homens possuem os mesmos caracteres \emph{psychologicos} e moraes. Por quanto, todas as raças --- têm \emph{intelligencia}, capaz de perceber o \emph{immaterial} e de se adiantar na civilisação, --- são dotadas de \emph{vontade}, capaz de praticar o bem e de fugir do mal, --- e possuem a idea de \textsc{Deus}, de \emph{verdade}, de \emph{justiça}, de \emph{religião}, de \emph{honestidade}, e o espirito de \emph{sociabilidade}. Ora isto demonstra que todos os homens pertencem á mesma especie.\par b) As divergencias, que existem entre os homens, são apenas \emph{accidentaes}. Na verdade, as principaes differencas, que se notam entre os homens e formam diversas raças, são a cor e a \emph{conformação do craneo}. Ora as differencas na \emph{cor} e na \emph{conformação} do \emph{craneo} são apenas \emph{accidentaes}.\par \emph{a}) A differenca na \emph{cor} será apenas um caracter \emph{accidental}, se a propria \emph{cor} depender de uma causa \emph{accidental} e \emph{ewtrinseca}: por quanto os caracteres \emph{essenciaes} do ente derivam de uma causa \emph{substancial} e \emph{intrinseca}. Ora a \emph{cor} depende sobretudo de uma causa \emph{accidental} e \emph{extrinseca}, que é o \emph{clima}. Na verdade, os descendentes das raças brancas, que vivem na Africa, tornam-se pretos; as colonias portuguezas das costas africanas apresentam muitos exemplos d'esta mudança de \emph{cor}. Os arabes são mais claros nos paizes mais frescos. O europeu, que se aclimata no Egypto, em pouco tempo adquire a cor pardacenta, e na Abyssinia a bronzeada.\par O proprio \emph{Bwrmeister} (\emph{Historia da creacdo}, pag. 507) reconhece que o sol e o clima exercem uma grande influencia na \emph{cor} da pelle. Eis as suas palavras: « Póde dizer-se que alguns africanos se tornam brancos debaixo da zona temperada, estando muitas gerações sob a acção dos raios solares mais obliquos; mas nunca se tornam tão brancos como os europeus. Por outro lado, os povos brancos tornam-se morenos sob os tropicos, mas nunca se tornam negros na Africa, nem vermelhos na America ».\par Objectam: Se a \emph{cor} dependesse do \emph{clima}, o europeu, que vive na Africa, deveria tornar-se preto, --- e este, vivendo na Europa, deveria tornar-se branco. Ora não acontece assim. Logo a cor não depende do \emph{clima}.\par Resposta. Antes de tudo, ha muitos factos para provar que o europeu, vivendo na Africa, depois de muitas gerações, se torna completamente preto, e vice-versa. --- Mas, embora \emph{Hoje} não se desse nenhum d'esses factos, não poderia dizer-so que o \emph{clima} não influiu, nos tempos passados, na \emph{cor} e por isso na formação das diversas raças. Por quanto, \emph{hoje} que as raças estão formadas e as influencias climatericas são menos violentas, a cor depende mais da \emph{raça} do que do \emph{clima}. Mas, no desenvolvimento da humanidade, quando as raças não estavam ainda formadas, \& \emph{cor} dependeu sobretudo do \emph{clima}. --- \emph{Wuttke} (\emph{Historia do paganismo}, I, pag. 30) diz: « A conclusão que, contra a unidade da especie humana, se tira do facto de o negro se não tornar branco e reciprocamente, é arbitraria e tem tanto valor como esta outra: assim como o homem, na edade adulta, não muda a physionomia, assim tambem não a muda, quando é criança ». --- \emph{Burdach} escreve: « Na\pgn{382}edade infantil, os homens não podiam ter adquirido ainda toda a sua personalidade, estayam mais dependentes da natureza e sujeitos ás influencias do mundo exterior; mas agora que as forças naturaes não actuam com a mesma intensidade, como anteriormente, e todas as relações tomaram um caracter mais estavel, o homem tornou-se mais independente, e já não está tão sujeito ás influencias do clima ».\par Pergunta-se: Qual é a causa \emph{intrinseca} da \emph{cor} da pelle?\par Os physiologistas não concordam em determinar a causa \emph{intrinseca} da \emph{cor} da pelle. Segundo alguns, a cor negra da pelle é uma consequencia de uma substancia colorida, que se interpõe nas cellulas. Segundo outros, provém de uma segregação de figado, que nos paizes quentes é eliminado pela pelle sob forma combustivel. \emph{Foissac} (\emph{Influencia do clima sobre} o \emph{homem}) affirma que a cor escura da pelle é consequencia da alimentação vegetal e de substancias em que predomina o carbone. \emph{Pruner} diz que o sangue do negro é mais escuro por causa do excesso do carbone nelle existente.\par \emph{b}) A \emph{conformação do craneo} tambem é um caracter \emph{accidental}. Com effeito, um caracter, que depende de causas \emph{extrinsecas} e \emph{accidentaes}, não póde deixar de ser \emph{accidental}. Ora \& \emph{conformação do craneo}, que nas differentes racas varia apenas entre 85º e 70º, depende de causas \emph{extrinsecas}, que são o \emph{clima}, o \emph{desenvolvimento da intelligencia}, \emph{á alimentação}, etc. --- Depende do \emph{clima}; pois as observações de \emph{Bür} demonstram que os povos, que vivem nas praias do mar, tém o craneo mais achatado que o dos habitantes das montanhas. --- Depende do \emph{desenvolvimento da intelligencia}, pois uma vida intellectual, muito activa e desenvolvida, reflecte-se ordinariamente sobre o orgão, que \emph{acompanha} o pensamento e que é o cerebro. Com o desenvolvimento do cerebro, e por isso de todo o systema nervoso, modifica-se todo o organismo. \emph{Prichard} affirma explicitamente que os \emph{Habitos} e o \emph{desenvolvimento} da \emph{intelligencia} exercem grande influencia na constituição physica do homem, e sobre tudo na forma do craneo e do rosto. --- Depende da \emph{alimentação}. \emph{Buffon} diz que uma alimentação grosseira faz degenerar os povos, e que os que vivem na miseria são mal constituidos. \emph{Br} encontrou, entre os tartaros, vestigios da influencia, que o modo de viver e a \emph{alimentação} exercem na forma do craneo e do rosto. Logo a \emph{conformação do craneo} é um caracter \emph{accidental}.\par \emph{c}) Não só a cor e a \emph{conformação} do \emph{craneo}, mas toda e qualquer differenca que existe entre as diversas raças é apenas \emph{accidental}. --- Porquanto não ha propriedade de uma certa raça, que insensivelmente não possa ser adquirida por individuos de qualquer outra raca. Ora isto mostra que aquellas propriedades são \emph{accidentaes}; pois que não podemos adquirir propriedades \emph{essenciaes}. --- De mais, os individuos que differem na cor, muitas vezes convém entre si na \emph{conformação do} craneo ou na qualidade do \emph{cabello}, e vice-versa. Logo estes caracteres são \emph{accidentaes}; aliás dois individuos, ao mesmo tempo, pertenceriam a especies diversas. --- Finalmente, os caracteres, proprios dos individuos de uma raca, encontram-se ás vezes em individuos de outras racas. Sendo assim, é evidente que aquelles caracteres são \emph{accidentaes}; pois os caracteres \emph{essenciaes}, proprios dos individuos de uma especie, não podem passar para os individuos de uma outra especie.}.

b) Todos os homens perteneerão á mesma \emph{especie}, se as diversas raças, eruzando-se, se propagarem em linha recta ao indefinido; porquanto, a fecundidade continua ou indefinida, \pg{382}  de que certos individuos são dotados, é signal de que elles pertencem á mesma \emph{especie}, ao passo que o \emph{hybridismo} mostra que os individuos divergem na \emph{especie}. Ora todas as diversas \pg{383}  raças humanas, eruzando-se, propagam-se em linha recta ao indefinido. Logo todos os homens pertencem á mesma \emph{especie} \footnote{Todos os physiologistas apontam a fecundidade continua como o caracter da especie e a cansa da sua fixidez. \emph{Flourens} diz: « A especie é a continuada fecundidade; é este o caracter decisivo » (\emph{Journal des savants}, 1868, pag. 628). O mesmo dizem \emph{Buffon}, \emph{Jussieu}, \emph{Candolle}, \emph{Geoffroy Saint}- \emph{Hillaire}, \emph{Chevreuil}, \emph{Leunis}, \emph{Müller}, etc.\par Os naturalistas reconhecem a fecundidade continua entre os individuos das diversas raças humanas. \emph{Wilbrand} (\emph{Descendencia do genero humano de} "\emph{m só par}) diz que é da união entre individuos das raças mais diversas, que resulta a prole mais formosa, mais forte, mais energica e fecunda.\par A eschola \emph{americana}, representada por \emph{Morton} e \emph{Agassis}, pretendeu demonstrar que os mesticos, que nascem dos individuos de differentes racas, não têm fecundidade poerduravel. Mas essa asserção foi desmentida por \emph{Bachman}, Waitz, \emph{Scherser}, \emph{Quatrefages}, \emph{Wagner}, e por outros naturalistas, que encontraram nas populações do \emph{Mexico}, das \emph{Philippinas}, do \emph{Nicaragua} e \emph{Paraguay}, do \emph{Novo Canadá}, das ilhas \emph{Pitcairn}, e d'outros paizes, muitos exemplos de fecundidade perduravel nos mestiços.\par De resto os melhores physiologistas, alguns dos quaes muito insuspeitos, sustentam a doutrina --- que todos os homens, seja qual for a sua raça, pertencem á mesma \emph{especie}.\par \emph{Humboldt} diz: Affirmando a unidade do genero humano, rejeito como lamentavel e indigna a hypothese das raças superiores e inferiores. Ha raças mais civilisadas, mais illustradas, nobilitadas pelo estudo e civilisação; mas uma não é por si mais nobre do que outra » (\emph{Kosmos}, I, pag. 385).\par \emph{Owen}, depois de ter comparado as diversas raças humanas, conclue: « Em consequencia d'isto, cheguei á conclusão que o homem é uma especie unica, e que as differenças peculiares das raças indicam apenas variedades ».}.

Coneluimos, pois, que todos os individuos humanos, pelo facto de possuirem todos os mesmos elementos \emph{essenciaes} e só differirem nos caracteres \emph{accidentaes}, pertencem á mesma especie.

Obiectam: As raças humanas differem tanto entre si, que não podem deixar de constituir especies diversas.

Resposta. As racas humanas differem menos entre si, do que differem as diversas raças de animaes, as quaes todavia pertencem á mesma especie. Mas, apezar de todas as differencas, o certo é que entre as racas humanas nunca se poude fixar um limite certo e determinado; d'ahi a divergencia entre os anthropologistas na classificação das raças humanas. Se não existe um limite certo e determinado, que distinga uma raça de outra, é signal que estas raças não constituem especies diversas; pois entre especies diversas ha um limite certo e determinado. \emph{Herper} (\emph{Ideas para a philosophia da historia da humanidade} VII, 2) escreve: « Sobre a terra não ha quatro nem cinco raças, nem variedades exclusivas. Às cores confundem-se por tal modo, que formam tantas sombras e graduações de um só grande quadro, que se extende \& todos os espacos da terra ». \emph{Humboldt} (\emph{Kosmos}, I, pag. 377) diz: « Se consideramos apenas as discrepancias extremas das raças, somos inclinados a consideral-as como especies diversas; mas em favor da unidade da especie humana fallam as multiplices variedades da pelle, da cor e da estructura do craneo, que os recentes estudos geographicos têm trazido a luz ».

%%>>> t231

\pg{381} \pg{384} \tese{231}{Todos os homens poderam derivar da mesma familia ou do mesmo par.}
Esta proposição é um corollario da precedente. --- Se todos os homens, por serem dotados dos mesmos elementos \emph{essenciaes}, pertencem á mesma especie, é claro que todos poderam derivar da mesma familia ou do mesmo par; visto que a descendencia dos individuos da mesma especie em linha recta é continua ou indefinida \footnote{A \emph{possibilidade} da origem commum de todos os homens não é o \emph{facto} d'essa origem. São duas questões distinctas; a solução da primeira pertence á \emph{scencia}, a da segunda á \emph{historia}. Se os homens \emph{poderam} derivar de uma só familia, não se segue que \emph{de facto} tenham derivado.\par Ora --- \emph{a}) a \emph{sciencia} effectivamente prova a \emph{possibilidade} da descendencia de todos os homens de um só par; --- 0) e a \emph{Historia} prova o \emph{facto} que todos os homens descenderam do mesmo par.\par a) A \emph{sciencia} demonstra a \emph{possibilidade} da descendencia de todos os homens de um só par. Porquanto, como dizemos no texto, se todos os homens pertencem á mesma especie, é claro que todos podem ter derivado da mesma familia ou do mesmo par; pois a especie é fecunda em linha recta ao indefinido. \emph{Schaafhausen} (\emph{Diario da reunião dos materialistas allemaes}, 1861, pag. 42) diz: « A possibilidade da descendencia de todos os homens de um só par deve ser admittida ». O mesmo dizem \emph{Giebel}, \emph{Waits}, \emph{Perty}.\par \emph{b}) A \emph{historia} prova o \emph{facto} que todos os homens descendem do mesmo par. Por quanto, a \emph{Biblia}, que é a historia mais antiga, mais authentica e mais digna de fé, affirma que todos os homens derivam de um só par; e esta affinmação não contradiz as descobertas certas da \emph{scencia}, pois esta, como vimos, prova a \emph{possibilidade} da origem commum de todo o genero humano. As affirmações da \emph{Biblia} encontram apoio --- nas tradições ácerca do genero humano, que são conhecidas, e que fallam de um homem e de uma mulher, de que derivaram todos os homens, --- e na \emph{glottologia}, a qual reconhece, mesmo \emph{pAilologicamente}, a origem commum d'aquellas linguas, que pareciam pertencer a um typo inteiramente diverso; ora a unidade de origem da\pgn{385}linguagem indica e prova a unidade de origem do genero humano (Cf. Trombetti, \emph{L'unità d'origine del linguaggio}, Bologna, 1905).\par \emph{Miller} escreve: « Se estabelecermos a questão, partindo das sciencias naturaes, não resta outra explicação possivel, senão a admissão pura e simples da descendencia de todos de um só par. A historia natural não procura se verdadeiramente todos os homens derivam de uma só familia; tudo o que ella póde attestar é que á humanidade se acha num estado, como se todos derivassem de um só par, e esta conclusão é muito importante para a unidade da especie ». :\par \&\par As principaes objecções, que os adversarios apresentam contra a descendencia de todo o genero humano de um só par, são tiradas --- da \emph{linguistica}, --- da habitação da \emph{America}, antes da sua descoberta, --- e da existencia das diversas racas humanas.\par \emph{a}) Dizem: Se todos os homens derivassem de uma só familia, não existiriam entre as varias linguas as profundidas divergencias, que existem.\par A. esta objecção respondem \emph{Humboldt}, \emph{Pott}, \emph{Müller} e \emph{Heitzig}.\par \emph{Humboldt} diz o seguinte: « Por mais isoladas que parecam certas linguas, e extranhos os sons, as articulações e as propriedades, todas têm entre si analogia, e as suas affinidades serão tanto melhor conhecidas, quanto mais progredir o estudo philosophico da historia dos povos e da linguistica ».\par \emph{Pott} escreve: « E' certo que, á medida que progride a sciencia das linguas, um grande numero das mesmas, que ainda hoje parece serem isoladas e não terem parentesco com as outras, serão incluidas nos grupos maiores, e estes serão cada vez menos numerosos. A \emph{linguistica} não se oppõe á doutrina, que admitte a descendencia de todos os homens e de todos os povos de um só par ».\par \emph{M}. \emph{Müller} e \emph{Heitsig} procuraram demonstrar que, tanto em relação á forma como em relação á materia, existe uma certa affinidade entre os varios grupos das linguas \emph{turanas}, \emph{semiticas} e \emph{arianas}.\par \emph{b}) Insistem: Se todos os homens derivassem de um só par, seria inexplicavel o facto de se terem encontrado homens na America.\par Resposta. Embora não possamos determinar o modo, por que os homens passaram á America, é certo que o podiam fazer, visto que os dois hemispherios são quasi contiguos (Lyell, \emph{Principes de géologie}, v. 2, pag. 472), e, como observou \emph{Humboldt} (\emph{Histoire de la géographie}, \emph{etc}.), para passar da Asia para a America bastava uma viagem de trinta e seis horas por mar.\par E não só podiam passar, mas ha provas de que effectivamente passaram á America. Estas provas são certas memorias.e certos ritos, encontrados naquelle paiz pelos primeiros exploradores, como a tradição da immigracao de seus antepassados, \emph{a} memoria de um diluvio e da conservação de uma só familia, a cireumcisüo entre os Mexicanos, a observancia do sabbado, a\par 25 Vor. II.\pgn{386}celebração do jubileu, e muitas outras coisas, attestando que esses povos tém a origem commum com as outras nações.\par c) Replicam: Se todos os homens derivassem da mesma familia, não haveria tantas raças humanas.\par Respondemos que a descendencia de todos os homens de uma só familia e a existencia das raças são duas coisas conciliaveis. Por quanto, como provámos, a diversidade das raças provém de causas \emph{accidentaes} e \emph{extrinsecas}, como são o \emph{clima}, o \emph{desenvolvimento da intelligencia}, a \emph{alimentação}, etc.}.

O proprio \emph{Burmeister} diz: « Todas as nações da terra pertencem, no sentido historico, em quanto á natureza, a uma só especie, e as suas differencas podem apenas ser consideradas como variedades ».

\emph{J}. \emph{Müller} escreve: « As \emph{raças humanas são formas de uma especie}, pois são fecundas entre si e reproduzem-se por meio da geração. Não são \emph{especies} differentes de um \emph{genero}; porque, se o fossem, seriam estereis os seus descendentes respectivos » (\emph{Phys}. \emph{de l'homme}).

\emph{Quatrefages} diz: « Os grupos humanos, seja qual for a sua differenca, real ou apparente, são apenas raças de uma e mesma especie, e não de especies distinctas... Entre os verdadeiros sabios, poucos deixarão de admittir este ponto de partida. Elles hão de concluir, como os grandes homens, de quem sou discipulo, conto \emph{Linneu}, \emph{Buffon}, \emph{Lamarch}, \emph{Cuvier}, \emph{Geoffroy}, \emph{Humboldt}, \emph{Müller}, \emph{que todos os homens pertencem d mesma especie}, \emph{que só existe uma unica especie de homens} » (\emph{L'éspéce humaine}, concl. du liv. I).

Seguem a mesma doutrina os mais celebres naturalistas e ethnographos: \emph{Blumembach}, os dois Wagner, \emph{Burdach}, \emph{Blainville}, \emph{Flourens}, \emph{Gratiolet}, \emph{Wilbrand}, \emph{Waitz}, \emph{Prichard}, \emph{Buffon}, \emph{Cuvier}, 'etc.

939. \textbf{Os homens} constituem um reino \textbf{especial}. --- Os homens constituirão um reino \emph{especial}, se possuirem um principio de operações \emph{caracteristicas}, diverso e superior ao principio de que são dotados os vegetaes e os animaes; porque são esses principios que distinguem os varios reinos da natureza. Ora os homens, pela sua alma \emph{racional}, possuem um principio de operações \emph{caracteristicas}, pertencente a uma ordem diversa e \pg{386}  superior á ordem, á qual pertencem os vegetaes e os animaes. Logo os homens constituem um reino \emph{especial} \footnote{\emph{Reino}, em \emph{physiologia}, indica o supremo genero dos entes, que conyém numa commum e caracteristica ordem de operações, embora possam differir com relação a outras notas, mesmo essenciaes. A denominação, como se vê, é \emph{metaphorica}, mas é bem applicada; porque, assim como os reinos politicos tiram a sua distincção das leis fundamentaes proprias e dos principios da vida civil e social, que pertencem a cada um e derivam do chefe respectivo, assim tambem os reinos da natureza distinguem-se pelas suas Operações caracteristicas, e, por isso, pelo principio radical, de que essas operações emanam. Portanto, um grupo de entes, que possuem uma nota especial e caracteristica, constitue um reino proprio ou especial, e este reino é distincto de um outro reino, que é constituido por outro grupo de entes, dotados de outra nota especial e caracteristica. Assim, as plantas, porque possuem a vida \emph{vegetativa}, constituem um reino especial, distincto do reino \emph{animal}, constituido pelos seres dotados tambem de vida \emph{sensitiva}. Ora o homem, --- embora pelo seu \emph{corpo organico} não constitua um reino \emph{especial} ou \emph{proprio}, porque, se se considera somente o organismo humano, a differenca entre o homem e o macaco é menos sensivel do que a que existe entre os macacos e os zoophytos (e todavia os macacos e os zoophytos pertencem ao mesmo reino), --- comtudo pela sua \emph{alma espiritual} acha-se collocado numa ordem immensamente superior á ordem, em que se encontram as plantas e os animaes, A distancia, que separa da vida sensitiva a intellectual, é muito mais sensivel do que a distancia, que separa da vida vegetativa a sensitiva; porque a vida vegetativa e a sensitiva pertencem ambas á \emph{mesma} ordem, que é a \emph{material}, ao passo que a vida sensitiva e a intellectual pertencem a ordens \emph{diversas}, e a distancia entre duas coisas da \emph{mesma} ordem é sempre menos sensivel do que a distancia entre duas coisas, pertencentes a ordens diversas. Por isso, se os \emph{brutos} não se podem reduzir ao reino \emph{vegetal}, mas constituem um reino proprio, muito menos pode reduzir-se o homem ao reino \emph{animal}, mas elle deve constituir um reino \emph{proprio}, que é o \emph{reino humano}.\par Muitos naturalistas negam esta verdade e dizem que o homem pertence ao reino \emph{animal}, embora constitua um typo ou uma familia especial. Muitos outros, porém, entre os quaes \emph{Bonnat}, \emph{Geoffroy St}. \emph{Hilaire}, \emph{Quatrefages}, \emph{Tiedemann}, \emph{Godron}, acceitam e defendem a nossa these.\par Obiectam: Os proprios escholasticos affirmam que o homem é \emph{animal} e que no \emph{genero} convem com os brutos. Logo o homem não pode constituir um reino proprio e especial.\par Resposta. O homem é \emph{animal}, mas é animal \emph{racional}, e por isso, se convem com o bruto no elemento \emph{generico}, differe d'elle no elemento \emph{especifico}, que é a razão. Differindo do bruto no elemento \emph{especifico}, o homem\pgn{387}constitue um reino proprio e distincto do reino dos brutos. Do mesmo modo, o bruto, embora convenha com a planta no genero de \emph{substancia viva}, comtudo differe no elemento especifico, que é a \emph{sensibilidade}, e por isso constitue um reino proprio e distincto do reino dos vegetaes. --- Nem se diga que o homem, por estar subordinado a um \emph{genero}, não pode constituir um reino proprio. Por quanto, --- se é verdade que elle não pode constituir um reino proprio pelo facto de estar subordinado a um \emph{genero}, quando a palavra \emph{genero} se toma \emph{physiologicamente}, porque, neste caso, o \emph{genero}, por ser uma das ultimas divisões do \emph{reino}, \emph{é inferior} a este e por isso não constitue uin reino proprio, --- não é verdade que o homem não possa constituir um reino especiale proprio pelo facto de estar subordinado a um genero, se a palavra \emph{genero} se toma \emph{philosophicamente}, como a tomam os escholasticos, emquanto significa o elemento determinavel e perfectivel da essencia, porque então, o \emph{genero} tanto pode estar subordinado ao \emph{reino} (assim ha muitos generos sob o reino animal, vegetal, mineral) como pode ter subordinados muitos reinos (assim o genero \emph{substancia} comprehende debaixo de si todos os reinos da natureza). --- Como se vé, a difficuldade funda-se num equivoco.}.

DESTINO DO HOMEM

%%>>> art-33

\section*{Artigo I.\\ Fim do homem.}
\addcontentsline{toc}{section}{Artigo I. --- Fim do homem.}

%%>>> t233

\pg{385} \pg{387} \tese{233}{O homem foi creado para um fim.}
Todo o ser, produzido por uma causa intelligente, é destinado para um fim proporcionado á sua natureza. Ora o homem é um ser creado por \textsc{Deus}, que é a propria Sabedoria. Logo o homem foi creado para um fim \footnote{Em todas as operações dos agentes racionaes é necessario distinguir um duplice fim: o fim do \emph{agente} e o fim da \emph{obra}. O fim do \emph{agente} é o motivo, que leva este a produzir a obra. O fim da \emph{obra} é o destino, que o agente dà á propria obra. Estes dois fins podem ser distinctos um do outro, e podem tambem identificar-se. --- Na creação do homem temos este duplice fim. --- O fim do \emph{agente} é o motivo, que levou \textsc{Deus} a crear o homem. Este fim é a gloria de \textsc{DEus}; pois \textsc{DEus} não póde deixar de ser o centro dos louvores e das homenagens das creaturas. --- O fim da \emph{obra} é o destino, que \textsc{Deus} deu ao homem. E' d'este fim que tratamos no presente artigo.}.

%%>>> t234

\tese{234}{0 fim do homem é a felicidade.}
O fim de um ser é aquelle, para o qual o mesmo ser tende sempre e necessariamente \pg{388} . Ora o homem tende sempre e necessariamente para a \emph{felicidade}. Logo o fim do homem é a \emph{felicidade} \footnote{O desejo da \emph{felicidade} é uma lei natural. Não ha, nem póde haver ninguem, que não deseje ser feliz. Procura-se a felicidade sempre e em toda a parte. Todos os nossos actos, mesmo os que serão causa de remorsos e de inquietações, derivam do desejo da felicidade. Por quanto, todo o fim \emph{inferior} ou \emph{proximo} dos actos humanos está immediata ou mediatamente subordinado ao fim \emph{supremo} ou \emph{ultimo}, que é a \emph{felicidade}.}.

%%>>> t235

\tese{235}{A felicidade encontra-se no bem.}
A felicidade é o fim constante e necessario de todas as nossas tendencias e operações. Ora o fim de toda a tendencia e operação é o \emph{bem}; pois que todo o objecto, emquanto se refere á vontade, é e diz-se \emph{bem}. Logo a felicidade encontra-se no \emph{bem} \footnote{A \emph{felicidade} é uma coisa que attrahe. Ora o \emph{bem}, como tal, attrahe; e o mal, como tal, repelle. Logo a felicidade encontra-se no \emph{bem}. --- Se o homem commette o \emph{mal}, commette-o sempre no intuito de apagar a sua sede de felicidade. Mas, então, não abraça o \emph{mal}, como tal, mas abraça-o emquanto lhe \emph{parece} um \emph{bem}.}.

%%>>> t236

\tese{236}{O bem, em que se encontra a felicidade, deve ser universal.}
O bem, em que se encontra a felicidade, deve socegar completamente todas as aspirações da alma, de modo que esta nada mais tenha a desejar; porque a \emph{felicidade é um estado de perfeição}, \emph{constituido pela posse de todos os bens}. Ora o bem, que socega completamente todas as aspirações da alma, só pode ser o bem \emph{universal}; pois só este encerra todo o bem e exclue todo o mal. Logo o bem, em que se encontra a felicidade, deve ser \emph{universal} \footnote{O bem \emph{universal}, em que se encontra a felicidade, deve ser um e o mesmo para todos os homens. Por quanto, sendo o fim \emph{ultimo} de um ser proporcionado á sua natureza, é claro que, se todos os homens são dotados da mesma natureza especifica, todos elles devem tender para um e o mesmo bem \emph{universal}. --- O bem \emph{universal}, cuja posse forma a felicidade do homem, não é uma entidade \emph{abstracta} ou \emph{imaginaria}, mas é \emph{concreta} e \emph{real}; pois que uma entidade \emph{abstracta} ou \emph{imaginaria} não pode exercer no espirito de \emph{todos} os homens uma influencia constante e irresistivel e Deus não pode ser o auctor de um desejo vão e irrisorio.\par O bem \emph{universal}, pois, existe na realidade. Mas onde encontral-o? Nos bens \emph{finitos}, ou no bem \emph{infinito}? --- Eis a questão, que vamos resolver.}.

%%>>> t237

\tese{237}{A felicidade não se encontra nos bens finitos.}
O objecto da felicidade é o bem \emph{universal}; porque só este sacia completamente as tendencias ou aspirações da alma. Ora os bens \emph{finitos}, que são as \emph{riquezas}, os \emph{prazeres}, a \emph{gloria}, a \emph{sciencia}, \emph{á virtude}, não podem saciar completamente as aspirações da nossa alma; pois que, sendo bens \emph{particulares}, não contêm todos \pg{389}  os bens, nem excluem todos os males. Logo á felicidade não se encontra nos bens \emph{finitos} \footnote{Os systemas, que collocaram o objecto da nossa felicidade nos bens \emph{finitos}, podem reduzir-se a tres. --- Uns, como \emph{Epicuro}, seguido pelos sensualistas e materialistas, collocaram o objecto da nossa felicidade nos bens do \emph{corpo}, que são as \emph{riquezas} e os \emph{prazeres}; --- outros, como \emph{Zeno}, collocaram-no nos bens da \emph{alma}, que são a \emph{gloria}, \emph{á sciencia}, a \emph{virtude}; --- outros, como \emph{Aristoteles}, constituiram-no na somma dos bens da \emph{alma} e do \emph{corpo}.\par Estes tres systemas são falsos, porque o objecto da nossa felicidade não consiste nos bens do \emph{corpo}, --- nem nos bens da \emph{alma}, --- nem na \emph{somma} dos bens da \emph{alma} e do \emph{corpo}.\par a) O objecto da felicidade não consiste nos bens do \emph{corpo}. Na verdade, os bens do \emph{corpo} são os \emph{prazeres} e as \emph{riquezas}. Ora nem os \emph{prazeres corporeos}, nem as \emph{riquezas} são o objecto da nossa felicidade.\par \emph{a}) Não os \emph{prazeres corporeos}. Por quanto estes, se podem ser o fim das faculdades \emph{sensitivas}, não são o fim das faculdades \emph{intellectuaes}, pois o objecto deve ser proporcionado á respectiva faculdade. Se os \emph{prazeres corporeos} não são o fim das faculdades \emph{intellectuaes}, não podem ser o fim do homem; pois o homem é tal pelas suas faculdades superiores, que são as \emph{intellectuaes}. ---- De mais, o objecto da \emph{felicidade} deve saciar completamente as aspirações da alma. Ora a experiencia attesta que os \emph{prazeres corporeos}, longe de saciarem a alma e apagarem-lhe a sede da felicidade, são a causa de muitas inquietações e remorsos. Logo o objecto da felicidade humana não póde consistir nos prazeres \emph{corporeos}.\par \emph{b}) Não as \emph{riquezas}. Com effeito, as \emph{riquezas} não aperfeicoam, nem tornam felizes as faculdades \emph{superiores} do homem; não se desejam como \emph{fim}, mas como \emph{meios} para o fim, pois servem para sustentar a vida, para adquirir \emph{gloria} ou \emph{praser}; não dão socego ao possuidor, porque quanto mais se possuem, tanto mais se desejam, e são a causa de muitos cuidados e anciedades, e a occasiüo de muitos vicios e discordias.\par p) O objecto da felicidade não é constituido pelos bens da \emph{alma}. Os bens da alma são principalmente a \emph{gloria}, \& \emph{sciencia} e a \emph{virtude}. Ora nenhum d'estes bens é objecto da felicidade.\par a) Não é a \emph{gloria}. Porquanto a \emph{gloria}, sendo uma coisa \emph{ewtrinseca}, \emph{instavel}, \emph{dependente da vontade dos outros}, \emph{fonte de muitas anciedades e cuidados}, não póde satisfazer e saciar completamente as tendencias da alma.\par \emph{b}) Não é a \emph{sciencia}. Com effeito, a \emph{sciencia}, embora seja um grande bem, comtudo não é o bem \emph{supremo} ou \emph{universal}, porque não só não encerra todos os outros bens, mas está unida com muitos males, trabalhos e.imperfeições, e por isso não póde tornar o homem completamente feliz.\par \emph{c}) Não é a \emph{virtude}. Por quanto, a \emph{virtude} é um meio para a felicidade, mas não é a felicidade; pois esta é o premio da \emph{virtude}. --- Além disso, \& \emph{virtude}, exigindo muitos sacrificios, não sacia as tendencias da alma.\par 0) O objecto da felicidade não é a somma de todos os bens da alma e do \emph{corpo}. Antes de tudo, é impossivel que o homem possua, ao mesmo tempo, \emph{todos} os bens \emph{finitos}, tanto da \emph{alma}, como do \emph{corpo}; porque muitos não dependem da nossa vontade, o muitos excluem-se mutuamente, assim a\pgn{390}\emph{virtude}, que é um bem da \emph{alma}, não póde exercer-se sem incommodos e sacrificios. --- Mas, ainda que o homem possuisse, ao mesmo tempo, \emph{todos} os bens \emph{finitos}, nem por isso seria feliz. Porquanto, a posse de todos esses bens, sendo acompanhada do medo da perda e de muitos trabalhos, não póde satisfazer completamente o desejo de felicidade (Cf. \emph{Sum}. \emph{Th}., 1-2, q. 2).\par Pergunta-se: Se a felicidade não se encontra nos bens \emph{finitos}, porque é que tendemos para elles?\par Resposta. A nossa vontade tende para os bens particulares e limitados, por isso mesmo que tende para o bem \emph{universal} e \emph{illimitado}. Por quanto, tendendo para o bem \emph{universal} e \emph{illimitado}, a vontade sente um impulso e uma attracção para todos os objectos que participam do bem \emph{universal}. A tendencia, pois, para os bens \emph{limitados} e \emph{particulares} é muito conforme com a indole da vontade. Todo o mal consiste em fazer dos bens \emph{finitos} o nosso fim \emph{ultimo}, quando elles são apenas \emph{meios} para a consecução do bem universal. D'ahi a regra pratica --- que devemos servirnos ou abster-nos dos bens \emph{finitos}, emquanto nos dirigem para o fim ou nos afastam d'elle; pois que é o fim que determina o emprego dos meios.\par *\par Com os mesmos criterios refutam-se algumas theorias ácerca do objecto da felicidade, apresentadas nestes ultimos annos. Reduzem-se a tres: a primeira é dos \emph{materialistas}, a segunda é de \emph{Hegel}, a terceira de \emph{Kant}.\par a) Os \emph{materialistas} sustentam que o objecto da felicidade se encontra no proprio homem.\par Esta opinião é absurda. Se o objecto da felicidade se encontrasse no proprio homem, este deveria bastar a si mesmo. Ora a consciencia attesta que o homem não basta a si mesmo, aliás todos seriam felizes. --- O homem deseja a verdade, o bem. Mas elle não é nem a verdade, nem o bem. Logo o objecto da felicidade não se encontra no proprio homem.\par \emph{b}) \emph{Hegel} diz que o fim ultimo do homem é a \emph{evolução continua}, \emph{universal e illimitada da humanidade}.\par Esta opinião tambem é absurda. Por quanto, a \emph{evolução continua}, \emph{universal e illimitada da humanidade} é apenas um corollario do Systema pantheistico de \emph{Hegel} ácerca da \emph{Idea}; e já not\&mos as falsidades o contradicções d'esse systema. --- De mais, aquella \emph{evolução da humanidade}, imaginada por \emph{Zegel}, é contraria á experiencia e á razão. E' contraria á experiencia, a qual attesta que muitos povos, depois de terem progredido no caminho da civilisação, retrocederam. E' contraria á razão; porque, sendo as acções do homem livres, a humanidade é capaz de progredir e de retroceder. --- Finalmente, que felicidade póde encontrar o homem no progresso da humanidade?\par c) \emph{Kant} disse que a verdadeira felicidade não consiste na posse do bem universal, mas num \emph{progresso continuo e indefinido para} a \emph{felicidade}, \emph{que todavia nunca se alcança}.\par A opinião de \emph{Kant} é falsa pelas seguintes razões.\par 1.º) A felicidade perfeita consiste na posse de todo o bem e na quietação de todos os desejos. Ora o \emph{progresso continuo} e \emph{indefinido paru a felicidade}\pgn{391}não é a posse do bem, mas é uma tendencia para elle; não é quietação, é movimento. Logo esse \emph{progresso} não é a verdadeira felicidade.\par 9.) Esse \emph{progresso para a felicidade} suppõe uma serie real de bens finitos, contingentes e relativos, sem um bem infinito, necessario e absoluto. Ora a supposicaão é absurda. Por quanto, se não existisse um ente infinito, necessario e absoluto, não poderia existir o ente finito, contingente e relativo, como provámos na \emph{Ontologia}.}.

%%>>> t238

\tese{238}{A felicidade encontra-se unicamente no bem infinito, que é Deus.}
A felicidade encontra-se no bem \emph{universal}, isto é, no \pg{390}  bem, que é só bem e por isso sacia inteiramente todas as tendencias e aspirações da alma. Ora o bem \emph{universal}, isto é, o bem, que é só bem e por isso sacia completamente as nossas \pg{391}  aspirações, é só o bem \emph{infinito}, que é \textsc{Deus}. Logo a felicidade encontra-se unicamente no bem \emph{infinito}, que é \textsc{Deus} \footnote{Quanta verdade e quanto sentimento encerram estas palavras de Santo Agostinho: « Senhor, ereaste-nos para Vós; eo nosso coração estará sempre inquieto, emquanto não descancar em Vós! » (Cf. \emph{Sum}. \emph{Th}., 1-2, p. 8, a. 1).\par Objectam: Entre o homem e o objecto da sua felicidade deve existir uma proporção. Ora nenhuma proporção existe entre \textsc{DEus} e o homem. Logo \textsc{Deus} não é objecto da nossa felicidade.\par Resposta. E' certo que entre o homem e o objecto da sua felicidade deve existir uma proporção; mas não é necessario que esta proporção seja de \emph{entidade}, basta- que seja de \emph{ordem}. Ora, se entre \textsc{Deus} e o homem não existe uma proporção de \emph{entidade}, pois \textsc{Deus} é infinito e o homem é finito, existe uma proporção de \emph{ordem}, emquanto o homem, tendendo com a sua intelligencia e vontade para a verdade e para o bem \emph{universal}, tende naturalmente para \textsc{Deus}, em que essa verdade e esse bem está concretisado.\par Insistem: Faculdades finitas, como são as do homem, não podem attingir um objecto infinito, como é \textsc{Deus}.\par Respondemos com uma distincção. E' verdade que faculdades finitas não podem attingir um objecto infinito por um modo \emph{adequado}; mas é falso que o não possam attingir por um modo \emph{inadequado}, pois, para isso, basta uma proporção de \emph{ordem}. Portanto a felicidade do homem, embora seja \emph{infinita} e \emph{increada} em relação ao seu \emph{objecto}, que 6 \textsc{Deus}, é \emph{finita} e \emph{creada} em relação á \emph{posse} do mesmo \textsc{Deus}. --- Em poucas palavras: uma faculdade \emph{finita} pode attingir um objecto \emph{infinito}, embora por um modo \emph{finito}.}.

239. \emph{A} felicidade do homem consiste no conhecimento e \textbf{no amor de Deus}. :

a) Sendo \textsc{Deus} o objecto da nossa felicidade, é claro que só na posse d'Elle a nossa alma póde encontrar o descanço e o termo de todas as suas tendencias. Mas \textsc{Deus}, por ser um espirito purissimo, só póde ser possuido pelas nossas faculdades espirituaes, que são a intelligencia e a vontade, emquanto uma o conhece e a outra o ama. Logo a felicidade do homem consiste no conhecimento e no amor de \textsc{Deus}.

b) \textsc{Deus} é o objecto da nossa felicidade, emquanto somos \emph{homens}. Ora somos \emph{homens}, isto é, distinguimo-nos de todas as creaturas visiveis, enquanto possuimos intelligencia e vontade. Logo \textsc{Deus} é o objecto da nossa felicidade, emquanto sacia a intelligencia e a vontade. Ora a intelligencia sacia-se no conhecimento de \textsc{Deus}, e a vontade no amor d'Elle. Logo \pg{392}  devemos concluir que a felicidade do homem consiste no conhecimento e no amor de \textsc{Deus} \footnote{(Cf. \emph{Sum}. \emph{Th}., 1-2, q. 3, aa. 2, 3). Consistindo a \emph{felicidade} no conhecimento e no amor de \textsc{Deus}, é evidente que haverá tantas especies de \emph{felicidade}, quantos forem os diversos modos por que conhecemos e amamos a \textsc{Deus}. Ora os diversos modos por que conhecemos e amamos a \textsc{Deus} são dois. Porquanto, --- ou podemos conhecer a \textsc{Deus} por uma visão \emph{intuitiva} e \emph{immediata}, por forma que \textsc{Deus} se veja como é em si mesmo e se ame com um amor correspondente; --- ou podemos conhecer a \textsc{Deus} por uma percepção \emph{abstracta} e \emph{mediata}, tirada das coisas creadas, e amal-o com um amor correspondente. O conhecimento \emph{intuitivo} e \emph{immediato} de \textsc{Deus}, por ser superior ás exigencias e condições da nossa natureza, constitue a \emph{felicidade sobrenatural}; o seu conhecimento \emph{abstracto} e «\emph{ediato}, por ser proporcionado á nossa natureza, forma a felicidade \emph{natural}. A felicidade \emph{sobrenatural} contém \emph{eminentemente} a \emph{natural}. --- Ao homem só era devida a felicidade \emph{natural}; mas \textsc{DEUS}, na sua infinita Bondade e Misericordia, quiz destinal-o á felicidade \emph{sobrenatural}. --- Devemos advertir que, embora para a felicidade do homem concorram os actos da intelligencia e os da vontade, todavia esta felicidade consiste \emph{radical} e \emph{essencialmente} no acto da \emph{intelligencia}, isto é, na contemplação de \textsc{Deus}. Por quanto a felicidade consiste na \emph{posse} de Dus. Ora \textsc{Deus} possue-se com o acto da intelligencia e não com o da vontade; pois que a intelligencia \emph{attrahe} a si o objecto, ao passo que a vontade \emph{tende} para o objecto \emph{anteriormente} percebido (\emph{Sum}. \emph{Th}., 1-2, q. 8, a. 4).}.

%%>>> t240

\tese{240}{Na vida presente a felicidade do homem não é perfeita.}
A felicidade do homem consiste no conhecimento e no amor de \textsc{Deus}. Ora este conhecimento e este amor de \textsc{Deus}, na vida presente, não podem saciar completamente todas as aspirações da alma; pois andam misturados com muitos males e perigos, incertezas e afflieções. Logo, na vida presente, a felicidade do homem não é perfeita \footnote{O conhecimento e o amor de \textsc{Deus}, se não tornam o homem \emph{perfeitamente} feliz sobre a terra, dão-lhe comtudo mais consolação do que todos os bens da terra. Só um Bem \emph{espiritual} e \emph{infinito} póde contentar a nossa alma, que é \emph{espiritual} e suspira pelo \emph{infinito} (\emph{Sum}. \emph{Th}., 1-2, q. 5, a. 3).}.

241. A felicidade do homem será perfeita na vida futura, --- O homem tende para a felicidade, que exelue todo o mal, toda a incerteza, toda a affieção; e este desejo, sendo muito natural e legitimo, não póde deixar de derivar do proprio Auctor da natureza. Derivando de \textsc{Deus}, o desejo da felicidade deve ser completamente saciado no conhecimento e no amor do proprio \textsc{Deus}. Ora não é completamente saeiado na vida presente. Logo será saciado na vida futura. Logo a felicidade do homem será perfeita na vida futura \footnote{Por isso a nossa alma, embora separada do corpo, póde ser completamente feliz; porque \textsc{Deus}, que é o unico objecto da nossa felicidade, é attingido\pgn{393} pela intelligencia e pela vontade, e estas faculdades estão sempre e necessariamente unidas com a alma (\emph{Sum}. \emph{TA}., 1-2, q. 4, a. 5).\par Se o nosso fim \emph{ultimo} é a \emph{posse} do Bem infinito, na vida futura, por meio do conhecimento e do amor, é certo que o nosso fim \emph{prowimo}, aqui sobre a terra, deve consistir em conhecermos, amarmos, e servirmos a \textsc{Deus}. Os proprios philosophos pagãos synthetisavam a regra e o. fim dos actos humanos naquella celebre sentença: \emph{Sequere} DEUM.}.

%%>>> art-34

\section*{Artigo II.\\ Immortalidade da alma humana.}
\addcontentsline{toc}{section}{Artigo II. --- Immortalidade da alma humana.}

%%>>> t249

\pg{393} \tese{249}{Immortalidade.}
\emph{Immortalidade} é a propriedade pela qual uma substancia \emph{viva} não está sujeita á morte, mas é destinada para viver eternamente. --- A \emph{immortalidade}, pois, não compete ás substancias ?\emph{morganicas} \footnote{A. \emph{émmortalidade} póde ser \emph{essencial}, \emph{natural} e \emph{gratuita}. --- \emph{Essencial} é a que se identifica com a mesma essencia do ente. Esta \emph{inmortalidade} é só propria de \textsc{DEus}; pois.a \emph{im?nortalidade} das creaturas é \emph{contingente}, como contingente é o seu ser, e por isso exige o concurso continuo da Causa creadora. --- \emph{Natural} é a que é exigida pela mesma natureza do ser. Esta exigencia resulta da. conveniencia e harmonia, que \textsc{Deus} estabeleceu nas creaturas. --- \emph{Gratuita} é a que \textsc{DEus} póde conceder á sua creatura, sem que a natureza d'esta o exija. --- A \emph{émmortalidade}, que é propria da nossa alma, é a \emph{natural}. --- S. Thomaz trata d'este assumpto principalmente --- na \emph{Sum}. \emph{Th}., p. I, q. 75, a. 6; q. 104, aa. 8 o 4; --- na \emph{Sum}. c. \emph{Gentes}, l. II, cc. 30, 55, 80; --- nas \emph{QQ}. \emph{DD}. \emph{De Anima}, a. 14; \emph{De pot}., q. 5, aa. 9 e 4; \emph{De verit}., q. 19, a. 1; --- 2 \emph{Sent}., dist. 19, q. 1, a. 1.}.

243. Immortalidade intrinseca e \textbf{extrinseca}. --- A \emph{immortalidade} pode ser \emph{intrinseca} e \emph{extrinseca}.

a) A \emph{intrinseca} é uma propriedade, pela qual uma substancia \emph{viva} não contém em si nenhum principio de corrupção, de modo que, em virtude da sua propria natureza, é destinada para viver eternamente. --- Por isso a immortalidade \emph{Intrinseca} costuma chamar-se \emph{Zncorruptibilidade}.

b) A \emph{extrinseca} é uma propriedade, pela qual uma substancia \emph{viva}, \emph{intrinsecamente immortal} ou \emph{incorruptivel}, não póde ser ou não é destruida por um agente exterior. --- A immortalidade \emph{extrinseca} diz-se 2\emph{ndestructibilidade}.

244. A \textbf{alma humana} é immortal. --- Uma substancia \emph{viva é immortal}, se for tal não só \emph{intrinseca} mas tambem \emph{extrinsecamente}. Ora a alma humana é \emph{immortal} não só \emph{intrinseca} mas tambem \emph{eztrinsecamente}. Logo a alma humana é \emph{immortal} \footnote{Para que possa concluir-se que uma substancia \emph{viva é immortal}, não basta provar que ella não contém em si mesma principio algum de corrupção,\pgn{394}mas é necessario provar tambem que nenhuma força exterior tem o poder ou á vontade de a destruir.\par Negam a \emph{íimmortalidade} da alma humana --- os \emph{materialistas}, afirmando que a nossa alma, por se identificar com a materia ou com as forças da materia, acaba com o corpo, --- os sequazes da \emph{palingenesis} e, em geral, os \emph{pantheistas}, sustentando que a alma fica absorvida pela essencia divina e por isso perde a propria individualidade, --- todos os que, para se entregarem mais livremente ás desordens, negam uma vida futura.}.

245. A \textbf{alma} humana é \textbf{intrinsecamente immortal}.

a), Uma substancia \emph{viva é intrinsecamente} immortal, quando não acaba --- nem por dissolução dos seus elementos (como acabam os \emph{compostos}), --- nem pela destruição do sujeito, em que existe e opéra (como acabam as formas \emph{mão-subsistentes}). Ora a alma humana não acaba --- nem por dissolução dos elementos, porque é \emph{simples}, --- nem por destruição do organismo, porque é \emph{espiritual}, isto é, independente da materia. Logo a alma humana é \emph{intrinsecamente} immortal \footnote{Este argumento é concludente. Tudo o que deixa de existir, deixa de existir --- ou porque se dissolvem os seus elementos, assim acabam todas as substancias compostas de \emph{materia} e de \emph{forma}, --- ou porque lhes falta o sujeito em que e com que operavam e existiam, assim acabam as almas dos animaes, pois estas, não tendo uma operação propria e independente da materia, não podem operar e por isso não podem existir fóra do organismo, e naturalmente acabam. Mas nenhuma d'estas hypotheses se verifica em relação á nossa alma. Esta não póde acabar --- nem como acabam os \emph{compostos}, pois a alma, como vimos, não é composta de \emph{materia} e de \emph{forma}, mas é essencialmente \emph{simples}, como o são todas as formas, --- nem como acabam as almas dos animaes, e todas as formas \emph{ndo-subsistentes}, pois \emph{a} alma humana, como provámos, tem operações \emph{proprias}, subsistencia \emph{propria}, ser \emph{proprio}, e este \emph{ser} constitue a sua propria \emph{essencia}; de modo que o \emph{ser} não se pode separar da forma subsistente, sem que essa forma se separe desi mesma. Logo a alma humana, por ser forma simples e subsistente, não pode estar sujeita á corrupção, e é por si immortal. Diz S. Thomaz: « \emph{Dupliciter} aliquid corrumpitur: uno modo \emph{per se}, alio modo \emph{per accidens}.... Quod \emph{per se} habet esse, non potest generari vel corrumpi nisi \emph{per se}. Quae vero \emph{non subsistunt}, ut accidentia et formae \emph{materiales}, dicuntur fieri et corrumpi per generationem et corruptionem \emph{compositorum}. Ostensum est autem quod animae brutorum \emph{non sunt per se subsistentes}, sed \emph{sola} anima humana. Unde animae brutorum corrumpuntur, corruptis corporibus; anima autem humana non posset corrumpi, nisi \emph{per se} corrumperetur. Quod quidem est impossibile non solum de ipsa, sed de quolibet \emph{swubsistente}, quod est \emph{forma tantum}. Manifestum est enim quod id quod secundum se convenit alicui, est inseparabile \&b ipso. \emph{Esse} autem \emph{per se} convenit formae, quae est actus... Impossibile est autem quod forma separetur a seipsa; unde impossibile est quod forma \emph{subsistens} desinat esse » (\emph{Sum}. \emph{Th}. p. I, q. 15, a. 6).\par Os adversarios pretendem destruir este argumento com varias objecções. Oitemos e confutemos as principaes.\par Dizem; A alma humana é \emph{forma} do corpo. Logo não pode perseverar\pgn{395}na existencia depois da corrupção do proprio corpo, porque acaba o fim pela qual foi creada.\par Resposta. E' verdade que a alma é \emph{forma} do corpo. Mas a sua união com o corpo não é o fim \emph{ultimo} da sua creação; é apenas o fim \emph{proximo}. Não é o fim \emph{ultimo}, porque o fim \emph{ultimo} pelo qual foi creada a alma é \& felicidade eterna, a qual consiste no conhecimento e no amor de Deus e que se alcança sem o concurso do corpo. Mas é o fim \emph{proximo}, pois é um \emph{meio} para conseguirmos, com os actos voluntarios da nossa vida mortal,o mesmo fim \emph{ultimo}. --- Além d'isso, \emph{a} alma é \emph{forma} do corpo, mas é uma forma \emph{subsistente}. Por isso, se, emquanto é forma, é para ella uma coisa natural a união com o corpo, todavia, emquanto é \emph{subsistente}, é para ella uma coisa natural a duração na existencia depois da dissolução do composto.\par Continuam: A alma nasce, cresce e envelhece com o corpo. Logo deve tambem morrer com o corpo; pois duas coisas, que são semelhantes no começo e na continuação da sua existencia, devem ser semelhantes no termo.\par Resposta. Negamos que as vicissitudes da alma correspondam ás do corpo. Por quanto, o corpo e a alma são entes produzidos; mas o corpo é produzido por \emph{geração}, e a alma pela \emph{creacdo}. --- O corpo e a alma crescem; pras o corpo cresce pela \emph{assimilação} da materia, e a alma cresce pela \emph{acquisicdo} de novas ideas, que não têm extensão nem materia. --- A alma não enfraquece, quando o corpo perde as suas forças; pois \emph{ordinariamente} nunca a alma é mais forte e vigorosa, como quando o corpo perde o vigor juvenil; e se, algumas vezes, a alma se resente das alterações do organismo, isso não mostra que uma e outra substancia são dotadas da mesma natureza, mas que estão unidas em unidade de \emph{pessoa}. --- Sendo tão dessemelhantes no começo e na continuação da sua existencia, a alma e o corpo hão de ser dessemelhantes com relação ao \emph{destino}; pois o corpo, por ser composto e material, resolve-se nos seus elementos, ao passo que a alma, por ser \emph{espiritual}, não póde corromper-se, mas deve subsistir em si.}.

b) Um desejo, que deriva da natureza e não se oppõe a um bem de ordem superior, não póde ser frustrado; pois, em tal caso, deriva do proprio \textsc{Deus}, e por isso deve ser satisfeito. Ora a nossa alma tem um vivissimo desejo de viver eternamente, e esse desejo deriva da natureza (pois a vontade deseja viver sempre, porque a intelligencia percebe uma vida sem fim) e não se oppõe a um bem de ordem superior. Logo o desejo de viver eternamente não póde ser frustrado. Logo a alma, considerada na sua \emph{intrinseca} natureza, é \emph{immortal} \footnote{O desejo de viver eternamente não constitue propriamente a \emph{immortalidade} da alma, mas é um \emph{signal} d'ella. Por isso este segundo argumento é \emph{a-posteriori}. Diz S. Thomaz: « Potest autem huius rei accipi \emph{signum} ex hoc quod unumquodque naturaliter suo modo esse desiderat. Desiderium autem in rebus cognoscentibus sequitur cognitionem. Sensus autem non cognoscit \emph{esse}, nisi \emph{sub hic et nunc}. Sed intellectus apprehendit esse \emph{absolute}, et \emph{secundum} omne \emph{tempus}. Unde omne habens intellectum \emph{naturaliter} desiderat \emph{esse semper}. Naturale autem desiderium non potest esse inane. Omnis igitur intellectualis substantia est incorruptibilis » (\emph{Sum}. \emph{T}7., p. L, q. V5, a. 6). --- \emph{A} razão é apodictica. A tendencia; nos seres cognoscitivos, segue o proprio\pgn{396}conhecimento. A alma sensitiva não apprehende a existencia senão determinada pelo \emph{espaço} e pelo \emph{tempo}, isto é, emquanto se realiza \emph{Aic et nunc}, e nao deseja senão essa existencia, assim percebida. Mas a alma intellectual, cujo objecto é o \emph{universal}, apprehende a existencia, isenta de toda a limitação do tempo e do espaço, e por isso perpetua, e para ella tende irresistivelmente.\par Objectam: Ha tantos desejos \emph{naturaes}, que todavia são \emph{énefficazes}. Ora um d'elles é o desejo da \emph{Immortalidade}.\par Respondemos que a \emph{znefficacia} póde dar-se com relação aos desejos \emph{condicionaes} e \emph{secundarios}, mas não com relação a um desejo \emph{absoluto} e \emph{radical}, como é o da \emph{immortalidade}; pois um desejo, que não depende de condições e que é a raiz de todos os outros desejos, mostra o proprio destino do ente, e não póde ser vão, mas deve ficar satisfeito.\par Insistem: Os animaes desejam a \emph{vida}, e todavia a sua alma não é \emph{immortal}.\par Resposta. Os animaes desejam a \emph{conservação} da vida, mas não podem desejar uma vida \emph{futura} e \emph{immortal}. Por quanto, como dissemos, o desejo corresponde á percepção; ora o animal, não possuindo ideas \emph{abstractas} ou \emph{universaes}, não póde \emph{perceber} e por isso não póde desejar uma vida \emph{futura} e \emph{immortal},}.

246. A \textbf{alma humana} é \textbf{extrinsecamente immortal}. --- A alma humana, sendo \emph{espiritual}, só póde acabar por \emph{aniquilação}, isto é, pela redueção ao \emph{nada}. Sendo a \emph{aniquilação} um acto que excede a força de toda a creatura (pois esta, assim como não póde tirar do \emph{nada} ou \emph{criar}, assim tambem não póde reduzir ao \emph{nada} ou \emph{aniquilar}), segue-se que só \textsc{Deus} póde \emph{aniquilar} a nossa alma. Ora \textsc{Deus} nunca \emph{aniquilará} a nossa alma; porque, na ordem actual das coisas, esta aniquilação repugna á \emph{Sabedoria}, \emph{á Bondade}, \emph{á Justiça} e á \emph{Santidade} de \textsc{Deus}.

a) Repugna á \emph{Sabedoria} de \textsc{Deus}. --- A \emph{Sabedoria} de \textsc{Deus} exige que a sua economia, no governo do mundo, se conforme com a natureza de cada creatura. Ora tal conformidade não se daria, se \textsc{Deus} aniquilasse a nossa alma, depois de a ter creado \emph{intrinsecamente} incorruptivel ou immortal \footnote{Se a alma humana, por ser subsistente, não acaba com a dissolução ou corrupção do composto, ella continuará a existir sempre, a não ser que seja aniquilada por Deus (Cfr. \emph{Sum}. \emph{c}. \emph{Gentes}, l. II, c. 55). Esta aniquilação não seria o effeito de uma acção \emph{positiva} de Deus, pois esta não pode ter por termo o nada, mas seria o effeito da cessação da influencia divina, de que as creaturas precisam para a sua duração na existencia. Ora Deus, auctor da natureza, governa as coisas em conformidade com a natureza e exigencias d'ellas, e por isso a nenhuma tira o que lhe compete essencialmente. Se Deus fizesse d'outro modo, metter-se-hia em contradicção comsigo mesmo; porque operaria contra a ordem estabelecida por Elle, quando formou os archetypos das essencias productiveis. Por isso, sera alma humana exige, pela sua natureza, \& conservação na existencia, é certo que Deus, embora, pela sua potencia \emph{absoluta}, possa destruir a nossa alma, comtudo, pela sua potencia \emph{ordenada}, isto é, pela sua potencia considerada em relação á ordem natural,\pgn{397}por Elle livremente estabelecida, não a destruirá. Diz S. Thomaz: « Ex hae sola potentia (primi agentis) nihil potest dici corruptibile; tum quia res dicuntur necessariae et contingentes secundum potentiam quae est in eis, et non secundum potentiam Dei; tum etiam quia Deus, qui est institutor naturae, non subtrahit rebus id quod est proprium naturis earum. Ostensum est autem quod proprium naturis intellectualibus est quod sint perpetuae. Unde hoe eis à Deo non subtrahitur » (\emph{C}. \emph{Gentes}, l. II, c. 55; Cf. \emph{De pot}., q. 5, a. 4). --- Para se comprehender a força d'esta razão, advertimos que por dois modos pode o homem conhecer o que Deus fará no mundo: ou pela \emph{revelação}, ou pela \emph{rasdo}, investigando o que a essencia de cada ente exige. No primeiro caso, temos a fé; no segundo, a \emph{sciencia}.}.

b) Repugna á \emph{Bondade} de \textsc{Deus}. --- E' certo que o homem sente uma tendencia irresistivel para a felicidade \emph{perfeita}, a qual consiste na posse \emph{eterna} de todos os bens e na exclusão \emph{eterna} de todos os males, --- tendencia, que, sendo natural, deriva do proprio Auctor da natureza; e tambem é certo que a felicidade \emph{perfeita} se não encontra na vida presente. Logo devemos dizer --- ou que \textsc{Deus} deu ao nosso espirito uma tendencia e um desejo que, por nunca ficar satisfeito, só nos devia servir de tormento, --- ou que existe uma vida eterna, em que a alma consiga a sua felicidade \emph{perfeita}. Ora repugna 4 \emph{Bondade} de \textsc{Deus} excitar um desejo, que nunca seja satisfeito. Logo deve existir uma vida \emph{eterna}, e por isso \textsc{Deus} nunca \emph{aniquilará} a nossa alma \footnote{Dizemos que a felicidade \emph{perfeita} traz comsigo a idea da \emph{eternidade}; porque a incerteza ou o medo de que o bem possuido possa acabar já inquieta e torna triste e infeliz a alma.}.

c) Repugna á \emph{Justiça} de \textsc{Deus}. --- A \emph{Justiça} de \textsc{Deus} exige que os bons recebam a recompensa \emph{adequada} ás suas virtudes, e os maus soffram o castigo \emph{proporcionado} ás suas eulpas. Ora a experiencia attesta que, na vida presente, muitos actos bons ou não são recompensados, ou não recebem um premio \emph{adequado}, e que muitos actos maus ou não são castigados, ou não o são como merecem. Logo deve existir uma vida futura, em que o homem receba o premio ou o castigo adequado dos seus actos; e por isso \textsc{Deus} não \emph{aniquila} a nossa alma \footnote{O proprio \emph{Rousseau} dizia: « Embora não existissem outras provas em favor da immortalidade da alma, senão os triumphos dos malvados e as oppressóes: dos bons, só isso tirar-me-hia toda a duvida. Esta nota tão dissonante na harmonia universal obrigar-me-hia a concluir que, com a morte, não acaba o homem todo, antes tudo reentra na ordem \emph{rn}.\par \emph{Cousin} escreve: « O mundo, em que vivemos, não nos mostra a perfeita\pgn{398}harmonia, que deve existir entre a felicidade e a virtude; e, como esta harmonia é necessaria e deve existir, é necessario concluir que existe um mundo, em que a harmonia se restabeleca ».\par Objectam: Para recompensar a virtude e castigar o vicio, não é necessaria uma vida futura; sobre a terra, os actos humanos tem a sua sancção na paz, que deriva da virtude, e no remorso, causado pelo vicio.\par Respondemos que, embora a virtude e o vicio recebam nesta vida o respectivo premio ou castigo, \emph{consistente} na paz ou no remorso da consciencia, todavia esta sancção não é \emph{adequada}; pois que, nem a paz, que os justos gosam, exclue todos os males, nem o remorso, que agita os impios, exclue todos os bens. --- Além d'isso, ha muitos actos virtuosos e muitos crimes, que nao recebem recompensa ou castigo sobre a terra. Que recompensa recebe aqui um soldado que morre pela sua patria, ou um christão que morre pelo seu \textsc{DEus}? --- Finalmente, a paz, que os bons gosam, e o remorso, que atormenta os maus, não só não excluem mas suppõem a \emph{émmortalidade} da alma. Porquanto, nem os justos gosariam, nem os maus seriam lacerados pelos remorsos, se não estivessem convencidos de que existe uma vida futura e immortal em que todos hão de receber o premio ou o castigo, conforme a qualidade das suas obras.}.

d) Repugna á \emph{Santidade} de \textsc{Deus}. --- A \emph{Santidade} de \textsc{Deus} exige que se proponha um meio efficaz para induzir os homens \pg{398}  a praticar o bem e a fugir do mal. Ora, não haveria meio algum efficaz, se não existisse uma vida futura, em que a alma gose ou soffra eternamente pelo bem que fez ou pelo mal que praticou. Logo a \emph{Santidade} de \textsc{Deus} exige que a nossa alma não seja \emph{aniquilada} \footnote{Uma vida eterna, feliz ou infeliz, é o unico meio \emph{efficaz} para garantir a ordem moral e social. Qualquer outro meio não actua tão efficazmente sobre a vontade, para a determinar á pratica do bem e á fuga do mal.\par Objectam: O argumento, que se basóa nos attributos divinos, prova que a alma humana não acaba com a dissolução do organismo, mas não prova que deva durar \emph{eternamente}, isto é, que seja \emph{immortal}; aliás deveria admittir-se a repugnante theoria ácerca da eternidade do premio e do castigo.\par Resposta. O argumento, tirado dos atributos divinos, prova não só que a alma sobrevive ao corpo, mas tambem que é \emph{immortal}. Por quanto, accommodando-se a \emph{Sabedoria} de \textsc{DEUs}, no governo do mundo, á natureza das suas creaturas, é claro que, se a nossa alma é \emph{intrinsecamente} immortal, \textsc{Deus} a sustentará eternamente na existencia.\par Negamos que a theoria ácerca da eternidade dos premios e dos castigos seja absurda ou repugnante; pois é muito conforme com a nossa razão. Na verdade, é justo que a felicidade, que \textsc{Deus} dará aos bons, seja eterna; porque, se não fosse eterna, não seria \emph{perfeita}, não excluiria todo o mal. --- Tambem é justo que o castigo, que \textsc{DEus} infligirà aos maus, seja eterno. Por quanto, quem morre em peccado mortal, não tendo já tempo para se arrepender, é inimigo de \textsc{Deus} por toda a eternidade, e, como tal, merece \emph{eternamente} o castigo da sua rebelliao. Diz S. Thomaz: « Culpa, quae est irreparabilis, de se habet quod perpetuo durat; et ideo debetur ei poena aeterna » (\emph{Sum}. \emph{Th}, 1-9, q. 81, a. 4 ad 3). --- De mais, o peceado grave, sendo offensa de uma Bondade infinita, contém em si uma certa malicia infinita; pois a gravidade da offensa mede-se pela qualidade da pessoa offendida. Contendo uma malicia infinita,\pgn{399}o peceado grave deve receber um castigo infinito, --- \emph{infinito} não na \emph{intensidade} (pois a creatura não é capaz de soffrer um castigo \emph{infinitamente} intenso), mas na \emph{duração}, soffrendo por toda a eternidade.\par Examinemos agora o valor de algumas provas em favor da \emph{immortalidade} da alma. Limitamo-nos a citar as de \emph{Descartes}, \emph{Galluppi}, \emph{Rosmini}.\par a) \emph{Descartes} adduziu esta, prova: A alma humana é uma substancia. Ora toda a substancia é indestruetivel. Logo a alma humana é indestructivel.\par Esta argumentação é viciosa. A alma humana não é \emph{incorruptivel} pelo facto de ser \emph{substancia}; pois ha tantos seres, que são substancias, e todavia são corruptiveis, como são os corpos. Para demonstrar que a nossa alma é \emph{incorruptivel}, é necessario demonstrar que é substancia \emph{espiritual}.\par \emph{b}) \emph{Galluppi} estabeleceu o seguinte \emph{sorites}: A alma humana é simples. O que é simples não tem partes. O que não tem partes não acaba por dissolução de partes. O que não acaba por dissolução de partes é naturalmente indestructivel. Logo a alma humana é naturalmente indestructivel,\par Nem esta razão é sufficiente. Porquanto, do facto de um ente ser simples, e não acabar pela dissolução dos elementos, não se segue que seja indestruetivel; porque póde acabar, por lhe faltar o sujeito em que existia. Assim todas as formas \emph{substanciaes não-subsistentes} são \emph{simples} e todavia acabam, --- não pela dissolução dos elementos, que ellas não têm, --- mas por acabar o sujeito em que existiam, porque não se sustentam a si mesmas. Logo a \emph{émmortalidade} da alma deve fundar-se numa razão differente.\par \emph{c}) \emph{Rosmini} disse que o \emph{ente universal}, mostrando-se á alma \emph{sensitiva}, não só a torna \emph{intellectual}, mas tambem \emph{immortal} (\emph{Psicologia}, \emph{v}. I, lib. V).\par Esta opinião é falsa pelas seguintes razões:\par 1.º) Suppõe que a alma \emph{sensitiva} do homem se transforma em \emph{intellectual}, pela manifestação do ente \emph{universal}. Ora esta supposição, como provámos no capitulo anterior, é falsa e absurda.\par 2.) A \emph{immortalidade} é um dote \emph{intrinseco} e \emph{essencial} da nossa alma, e por isso deve deduzir-se da intima natureza d'ella. Ora, se a apparição do ente \emph{universal} desse á nossa alma a \emph{immortalidade}, este dote seria apenas \emph{extrinseco} e \emph{accidental}.\par 3.º) Segundo \emph{Rosmini}, o ente \emph{universal}, emquanto \emph{informa} a alma, communiea-lhe a \emph{émmortalidade}. Ora isto é impossivel. Por quanto, esse ente \emph{universal} --- ou é \emph{ideal}, --- ou é \emph{real}. Se é \emph{ideal}, communicará \& alma propriedades \emph{ideaes}, mas não a \emph{immortalidade}, que é uma propriedade \emph{real}. Se é \emph{real}, é \textsc{Deus}; mas \textsc{Deus} não póde ser \emph{forma} de nenhum ente creado.}.

247. A \textbf{immortalidade da alma} humana e o consenso \textbf{da humanidade}. --- A convicção da \emph{immortalidade} da alma humana está tão profundamente gravada no seio da humanidade, que sempre e em toda a parte se manifestou como crença \emph{pratica}, \emph{publica} e \emph{efficaz}, saneeionada, sobretudo, pelo culto aos mortos e pelas apotheoses aos heroes e aos benemeritos da patria. Ora esta convicção tão intima e tão universa] da humanidade só póde ser o resultado da voz da propria natureza, a qual, sendo \pg{400}  egual em todos os homens, em todos desperta o mesmo sentimento e o mesmo desejo de uma vida sem fim \footnote{Poucas verdades ha que tenham em seu favor o consenso unanime e constante de todos os povos, como a. \emph{immortalidade} da nossa alma.\par \emph{Cicero} diz: « O \emph{consenso de todas as nações}, assim como nos leva a reconhecer a existencia dos deuses, assim tambem nos obriga a admittir á sobrevivencia das nossas almas ».\par \emph{Bossuet} escreve: u As historias antigas e modernas attestam que a idea de uma vida immortal se encontra confusa nos povos selvagens, mas clara e distincta nos povos, que têm um exacto conhecimento de \textsc{DEus}; pois estes vêem que a creatura racional póde viver eternamente feliz, admirando as grandezas de \textsc{DEus}, os planos da sua sabedoria e a belleza das suas obras ».\par \emph{Lamennais} (\emph{Essais sur lindif}.) diz: « A immortalidade da alma, dogma capital, de que ninguem, como diz Celso, póde afastar-se, foi sempre uma crença universal do genero humano, como confessam os mais fogosos adversarios do Christianismo. \emph{Voltaire} e \emph{Boligbroke} concordam nisso 7.\par Este consenso unanime traduz-se pelo \emph{culto} aos \emph{nortos}, pelos \emph{cuidados} no \emph{futuro}, e pelo \emph{modo de fallar}. ---- Traduz-se pelo \emph{culto aos mortos}. Todos os povos professam respeito pelos seus mortos, e este respeito é mais ou menos profundo, segundo o maior ou menor grau de civilisação. Os \emph{egypcios}, os \emph{gregos} e os \emph{romanos} distinguiam-se nisso entre todos. --- T'raduz-se pelo \emph{cuidado no futuro}, porque o homem pensa em perpetuar a sua memoria, em preparar ao seu nome uma aureola eterna de gloria. --- Traduz-se pelo no \emph{modo de fallar}, pois o amigo promette ao seu amigo uma fidelidade \emph{eterna}, uma memoria \emph{eterna}, etc. \emph{Poisson} (\emph{Entretiens sur la} vie \emph{future}) diz: « Escutai o homem; as suas promessas são immutaveis, as suas affeições perpetuas, os seus amores eternos; as suas glorias seculares; falla como opéra, como um ser immortal ». \emph{P}. \emph{Janet} (\emph{La famille}) escreve: « A eternidade é tão natural ao amor, que este nada sabe pedir ou conceder, sem prometter a eternidade ».\par Ora este consenso unanime e constante não póde derivar dos prejuizos, das paixões, da ignorancia, ou do erro (pois estas causas não são universaes nem constantes), mas, como dissemos no texto, é produzido pela propria voz da natureza, a qual não poderia dizer que á alma é immortal, se esta verdadeiramente o não fosse. Com razão disse \emph{Cicero}: « Se o consenso unanime de todos é a voz da natureza, e se effectivamente todos os povos estão de accordo em que com a morte não acaba tudo, é forgoso que acceitemos este principio » (\emph{Tuscul}., 1, 15).\par Não póde negar-se que a crença na \emph{immortalidade} da alma humana apparece, em muitas religiões, misturada com muitos erros e superstições. Mas estes erros e estas superstições variam segundo os tempos e os povos, ao passo que o fundo d'essa crença ficou sempre illeso.\par Objectam: O \emph{Buddhismo}, que conta quasi 400.000.000 de sequazes, admitte a \emph{aniquilação} de toda a personalidade, pois \emph{Nirvána} é considerado como a aniquilação das almas completamente purificadas.\par Resposta. Antes de tudo advertimos que não ha religião mais incoherente e mutavel que o \emph{Buddhismo}, tendo-se alliado com todos os systemas religiosos, desde o monotheismo até á mais grosseira idolatria. --- Em seguida, diremos que, se alguns metaphysicos do meio dia da India consideram \emph{Nirvána} como a aniquilação das almas, a maior parte dos \emph{Buddhistas} pensam que\pgn{401}as almas, entrando em \emph{Nirvána}, se purificam de todos os males e imperfeições, mas conservam a propria identidade e individualidade. E' a opinião de \emph{Max Müller}, \emph{Barthélemy S}. \emph{Hilaire}, \emph{Barth}, \emph{Tele}, etc.}.

%%>>> art-35

\section*{Artigo III.\\ Condição e operações da alma separada.}
\addcontentsline{toc}{section}{Artigo III. --- Condição e operações da alma separada.}

\pg{394} \pg{395} \pg{396} \pg{397} \pg{399} \pg{401} 248. A \textbf{alma separada} é sempre individua. --- A alma humana, embora separada, tem comtudo uma relação \emph{essencial} ao seu corpo e, por essa relação, distingue-se das outras almas da mesma especie e por isso é \emph{individua}. Ora o que é \emph{essencial} a uma coisa, não pode separar-se d'ella. Logo a alma humana, mesmo separada do corpo, conserva sempre uma ordem essencial ao seu corpo, pela qual se distingue de todas as outras almas humanas, e por isso é sempre \emph{individua} \footnote{Como dissemos no cap. III da secção I d'este tratado, o principio que torna \emph{individua} a alma humana e a distingue das outras, é a \emph{essencial} relação ao seu corpo; de modo que a alma de Pedro é de Pedro e não é de Paulo, porque se refere essencialmente 20 corpo de Pedro e não ao de Paulo. --- As almas são todas eguaes na \emph{especie}, isto é, não se distinguem umas das outras pelos seus elementos \emph{especificos}, mas distinguem-se pela sua relação aos seus corpos. Esta relação não é uma propriedade \emph{accidental} da alma, nem é uma consequencia da distincção, como pretendem os \emph{Scotistas}, mas é o principio e a razão da distincção, é o elemento \emph{essencial} que constitue \emph{humana} e \emph{individua} a nossa alma. --- Por isso a nossa alma é composta, na sua propria essencia, de dois elementos, --- um \emph{commum}, e outro \emph{proprio}. O elemento \emph{commum} é que ella é \emph{acto primeiro do corpo humano}; o elemento \emph{proprio} é que \emph{esta} alma é acto d'este corpo, e não de outro. Esta \emph{essencial} relação de cada alma a um certo e determinado corpo constitue a \emph{differença} entre as almas, e é o fundamento da sua distineção e multiplicação \emph{numerica}. Portanto, a. individuação da alma depende da sua essencial relação ao corpo, e não da sua união actual com o corpo; porque o \emph{ser} da alma não depende da materia, mas é communicado á materia. Com razão, pois, dissemos que o principio da \emph{individuação} é a propria \emph{essencia} da alma, emquanto se refere a um certo e determinado corpo. A essencia é o \emph{principio}, á causa da individuação; o corpo é o \emph{receptivo}, o \emph{termo} d'essa individuação, porque é ao corpo que a alma se refere e é no corpo que a alma. encontra o complemento na sua especie. --- Ora essa relação, pelo facto de constituir o elemento \emph{proprio} da essencia da alma, deve permanecer e permanece na propria alma, mesmo quando separada do corpo, e, permanecendo, deve produzir o seu effeito, que é o de distinguir uma alma de todas as outras da mesma especie, isto é, de tornal-a \emph{individua}. --- S. Thomaz ensina esta doutrina em varios pontos das suas obras, sobretudo --- na \emph{Sum}. \emph{Th}., p. I, q. 76, a. 2, --- na \emph{Sum}. \emph{c}. \emph{Gentes}., l. IL. c. 81; nas QQ. DD. \emph{De anima}, a. 9, --- in I \emph{Sent}. dist. 8, q. 5, a. 9 ad 6. Cf. Caiet. in op. \emph{de ente et essencia}, c. 6, q. 19. --- Limitamo-nos a citar as seguintes palavras do Angelico Doutor: « Diversitas animarum se-\par 26 Vor. I.\par É\pgn{402}paratarum est secundum diversam commensurationem ad corpora. Haee enim anima est commensurata huic corpori, et non illi, illa autem alii... Hujusmodi commensurationes remanent in animabus, etiam pereuntibus corporibus » (Sum. c. \emph{Gent}., 1. II, c. 81, n. 2).\par Mas, se a alma, separada do corpo, não está sujeita a alguma mudança radical na sua. essencia, o corpo, depois de abandonado pela alma, soffre uma transformação substancial. O corpo morto deixa de ser corpo \emph{Jano}, embora conserve, por algum tempo, a forma exterior de homem. Na verdade, o nosso corpo não é /\emph{nmmano} senão pela alma espiritual, que o colloca nessa determinada especie e o distingue dos corpos das outras especies; porque é a forma que dá a especie á materia. Quando, pois, a alma abandona o corpo, este perde a sua especie e soffre uma mudança substancial. Deixando de ser humano, o nosso corpo torna a fazer parte das substancias corporeas. --- Todavia o corpo humano, separado da alma, se não permanece o mesmo pela identidade da \emph{especie}, e por isso do \emph{numero}, permanece o mesmo pela identidade da \emph{materia}, a qual conserva sempre uma certa relação, embora \emph{extrinseca}, com o principio formal, que um dia a vivificou. Diz S. Thomaz: « Corpus mortuum alicuius sancti non est idem numero quod primo fuit dum viveret, propter diversitatem formae, quae est anima; est tamen idem identitate \emph{maleriae}, quae est iterum suae formae unienda » (\emph{Sum}. \emph{Th}., p. 8, q. 25, a. 6 ad 3). --- D'onde se vé o erro de \emph{Rosmini}, condemnado pela Egreja (prop. 22), quando diz que, pela virtude divina, pode ser separada do corpo animado a alma intellectual, permanecendo sempre o corpo na sua qualidade de \emph{animal}, pelo facto de permanecer nelle o principio animal. E' a alma que faz com que o nosso corpo seja \emph{animal} (\emph{animal} vem da \emph{anima}), e por isso, afastando-se do corpo a alma intellectual, afasta-se tambem o principio \emph{animal}, e o corpo não pode permanecer na sua qualidade de \emph{animal}.}.

%%>>> t249-2

\pg{402} \tese{249}{O estado de separação é connatural á alma.}
A alma

t humana, por ser subsistente, é naturalmente incorruptivel, de modo que a sua existencia deve ser eterna. Mas se a alma deve ter uma existencia eterna, esta existencia não a pode ter sempre no corpo, porque o corpo é naturalmente corruptivel. Logo a alma deve ter naturalmente a existencia fóra do corpo, quando este se corrompe, e por isso o estado de separação deve dizer-se e é \emph{connatural} á alma \footnote{A conclusão é evidente. A alma humana é \emph{forma subsistente}. Emquanto é \emph{forma}, ella diz ordem ao corpo; emquanto é \emph{subsistente}, é immortal. \textsc{Duas} coisas, portanto, convêm naturalmente á alma: a união com o corpo e a immortalidade. --- Mas o corpo, com que se une a alma, por ser composto de elementos contrarios, é naturalmente corruptivel, e por isso tambem corruptivel é a união, e um dia deve acabar. Logo se é natural a união da alma com o corpo, tambem é natural a dissolução d'essa união, e por isso é natural o estado de separação. --- O estado da união é \emph{natural}, emquanto deve existir \emph{por algum tempo}; mas não é \emph{natural}, emquanto deva existir \emph{sempre}. Logo pode chamar-se \emph{natural} tanto o estado de \emph{união}, como o estado de \emph{separação}, não ao mesmo tempo, mas successivamente. --- A alma tende \emph{naturalmente} tanto para um estado, como para o outro. Na verdade, o que ella deseja primeira e principalmente, com um transporte irresistivel,\pgn{403}é uma existencia eternamente feliz, e, em seguida, sabendo de dever existir agora no corpo e depois fóra do corpo, tende naturalmente para esse duplice estado, visto não poder viver sempre no estado de \emph{união}. --- Sendo \emph{connatural}, o estado de \emph{separação} não se pode chamar \emph{violento}. Porquanto, o estado \emph{violento} é o que repugna á \emph{principal} tendencia de um sujeito; ora a \emph{principal} tendencia da alma é para uma existencia eternamente feliz, sendo para ella uma coisa muito \emph{secundaria} a existencia no corpo, ou fóra do corpo. --- Todavia, devemos conceder que o estado de \emph{união} é mais natural que o de \emph{separação}; porque, no estado de \emph{separação}, a alma não pode, como veremos, exercer \emph{todas} as operações, que são proprias da especie humana e que ella exerce no estado de \emph{união}.\par Se o estado de separação é \emph{natural} á alma, a resurreição, isto é, a reconstituição do composto humano não é \emph{devida} ao homem. Na verdade, o corpo não é uma parte, um elemento da \emph{essencia da alma}, mas é um elemento, uma parte do \emph{composto humano}. E, se é essencial á alma que \emph{possa unir-se} ao corpo, não lhe é essencial que lhe seja \emph{actualmente unida}; aliás, deixando o corpo, devia immediatamente perder a existencia. Diz S. Thomaz: « Corpus non est de \emph{essentia} animae; sed anima ex \emph{natura suae essentiae} habet quod sit corpori \emph{unibilis} » (\emph{Sum}. \emph{Th}., p. I, q. 75, a. 7 ad 3). Nem podia ser d'outro modo. A alma humana --- não é uma forma \emph{material}, que tem o seu \emph{ser} no composto, e por isso, dissolvido o composto, perde o proprio ser, --- mas é uma forma \emph{subsistente} ou \emph{espiritual}, que tem o \emph{ser} por si e em si, de modo que não recebe do corpo o ser, mas; pelo contrario, communica o ser ao corpo, e por isso, separada do corpo, conserva o proprio ser e, com o ser, as proprias operações, embora não possa chamar-se perfeita na sua natureza, emquanto lhe é natural a união com o corpo. Diz o Angelico: « Anima subsistit in suo esse: unde relinquitur quod post separationem a corpore perfectum esse habeat; unde perfectam operationem habere potest, licet non habeat perfectam naturam » (\emph{Sum}. \emph{Th}. 1-2, q. 4, a. 5 ad 2). Portanto a alma humana pode, na vida futura, existir e operar separada do corpo; e permaneceria separada para sempre, se Deus, na sua infinita bondade, não tivesse concedido ao homem o beneficio, absolutamente gratuito, da resurreição; a qual, por isso, não pertence á ordem da \emph{natureza}, mas á da \emph{graça}, \emph{e} não pode chamar-se \emph{natural}, senão em quanto é \emph{conveniente} á natureza humana. Diz S. Thomaz: « Nullum principium \emph{activum}. resurrectionis est in materia...; unde resurrectio, \emph{simpliciter} loquendo, est \emph{miraculosa}, non \emph{naturalis}, nisi \emph{secundum quid} » (\emph{Sum}. \emph{Th}., \emph{Suppl}., q. '(5, \&, 9).}.

%%>>> t250

\pg{403} \tese{250}{A alma separada conserva um desejo innato de se reunir ao corpo.}
E' certo que a alma, quando é creada e está unida ao corpo, tem um desejo innato, isto é, uma aptidão natural, uma ordem essencial para constituir o composto. Ora esta ordem, esta aptidão dura sempre na alma, mesmo quando é separada do corpo; porque nem então perde o caracter de substancia \emph{incompleta} e de forma \emph{informante}, e, por quanto depende della, nunca teria desamparado o corpo. Logo a alma separada conserva um desejo innato de se reunir ao corpo \footnote{Diz S. Thomaz: « Secundum se convenit animae corpori uniri, sicut secundum se convenit corpori levi esse sursum. Et sicut corpus leve manet\pgn{404}quidem leve, cum a loco proprio fuerit separatum, cum aptitudine tamen et inclinatione ad proprium locum, ita anima humana manet in suo esse, cum fuerit à corpore separata, \emph{Aabens aptitudinem et inclinationem naturalem ad corporis wnionem} » (\emph{Sum}. \emph{Th}., p. I, q. 76, a. 1 ad 6). --- Este desejo innato da alma não é \emph{desordenado}, ainda que a resurreigio não seja \emph{devida} ao homem; porque um desejo não é \emph{desordenado}, quando o seu \emph{termo}, que, neste caso, é a \emph{união} da alma com o corpo, seja \emph{naturalmente} possivel, embora, depois da separação, seja impossivel \emph{per accidens}, isto é, não pela condição da alma, mas pela corrupção do composto. --- Nem este desejo innato constitue a alma separada num estado de violencia; porque, como dissemos, o estado de \emph{violencia} é o que repugna \& \emph{principal} tendencia de uma natureza, e \& \emph{principal} tendencia da alma não é a sua reunião com o corpo, mas é o desejo de viver eternamente feliz no modo, que Ihe fór \emph{naturalmente possivel}. Além d'isso, é necessario advertir que a alma separada, guiada somente pela luz da razão, --- ou julga impossivel a sua reunião com o corpo, --- ou, se a julga possivel, deve julgal-a possivel com o mesmo corpo miseravel e corruptivel; e por isso o seu desejo --- ou não é efficaz nem absoluto, porque diz respeito a uma coisa impossivel, --- ou não é tal que a sua irrealisação possa constituir a alma num estado de \emph{violencia}.\par A existencia d'esse desejo 7»1\emph{ato}, se não mostra a \emph{necessidade} da resurreição, mostra comtudo a sua \emph{conveniencia}. Na verdade, a alma, separada do corpo, acha-se num estado menos conforme com a sua natureza, a qual tende para a união. Ora parece pouco conveniente que dure eternamente um estado, menos conforme com a natureza das coisas. --- De mais, o modo, por que a alma separada percebe, é, como veremos, \emph{preternatural} e, num certo sentido, menos \emph{perfeito} que o modo por que percebe no estado de união. Ora parece menos conveniente que o conhecimento da alma seja sempre \emph{imperfeito} e \emph{preternatural} na vida futura, na qual ella deve \emph{aaturalmente} encontrar a perfeita felicidade. Finalmente, é conveniente que o corpo, assim como, na terra, foi o instrumento das obras virtuosas ou viciosas da alma, assim tambem, na vida futura, participe da sorte da mesma alma. (Cf. S. Th. \emph{Comp}. \emph{theol}. \emph{verit}. c. 151; \emph{c}. \emph{Gentes}, 1. IV, c. 19).\par Pars refutar as objecções, que se fazem contra a resurreicüo dos corpos, basta advertir --- 1.9) que, na resurreicüo, os corpos não reassumem todos os elementos da vida toda; aliás, renovando-se os elementos cada 8 ou 10 annos, “teriamos um corpo monstruoso; --- 2.9) que os corpos \emph{simples} ou \emph{elementares}, embora passem por muitas combinações e transformações, são indestruetiveis; --- 3.º) que o poder de Deus é infinito, e, se poude formar do lodo o corpo humano, pode reconstituil-o com elementos existentes.}.

%%>>> t251

\pg{404} \tese{251}{A alma, separada do corpo, produz algumas operações.}
Todo o ente, que conserva a propria substancialidade, produz algumas operações; pois a operação segue necessariamente a existencia. Ora a alma humana, separada do corpo, conserva a propria substancialidade; porque ella, como provámos, sendo \emph{simples}, não se resolve em partes, e, sendo \emph{espiritual}, não precisa de um sujeito, em que exista e opére. Logo o alma, separada do corpo, produz algumas operações \footnote{A existencia e a operação são duas coisas tão intrinsecamente ligadas\pgn{405}entre si, que uma não póde estar sem outra. Assim como a operação suppõe necessariamente um ser existente, assim o ser existente necessariamente opéra. Um ente, que não opera, não existe. --- Isto verifica-se sobretudo nos seres vivos, nos quaes \emph{viver} é \emph{ser}: \emph{vivere viventibus est esse}. Por isso, assim como repugna que um ser se separe de si mesmo, assim tambem repugna que a vida se separe da nossa alma.}.

%%>>> t252

\pg{405} \tese{252}{A alma, separada do corpo, não produz as operações organicas.}
Com effeito, a alma humana, embora seja o \emph{principio} e a \emph{raiz} das faculdades \emph{organicas}, que são as \emph{vegetativas} e as \emph{sensitivas}, comtudo não é o seu \emph{sujeito}; porque, como vimos, não é a alma só que vegeta e sente, mas é o \emph{composto} de alma e de corpo. Se o \emph{sujeito} das faculdades \emph{organicas} é o \emph{composto humano}, é claro que a alma, separada do corpo, não produz as operações \emph{organicas} \footnote{As faculdades \emph{organicas}, se não ficam na alma \emph{actualmente}, ficam nella \emph{virtual} ou \emph{radicalmente}. Na verdade, a alma é o \emph{principio primeiro}, de que as faculdades emanam para o \emph{composto}; porque a forma \emph{substancial} é a fonte de toda a actividade, de que os seres \emph{compostos} são dotados. Ora a alma, separada do corpo, por não mudar de natureza, não deixa de ser o \emph{principio primeiro} das faculdades \emph{organicas}; e por isso, assim como conserva sempre a capacidade de tornar a informar o proprio corpo, assim tambem conserva a capacidade de, \emph{quando reunida ao corpo}, reproduzir nelle as faculdades \emph{vegetativas} e \emph{sensitivas}, E", pois, muito exacto dizer que as faculdades \emph{organicas} não ficam \emph{actual} mas \emph{radicalmente} na alma separada, tirando a semelhança da arvore, cuja \emph{rais}, embora não possua \emph{actualmente} o tronco e os ramos, todavia tem a força de produzir estes e aquelle.\par Diz S. Thomaz: « Quaedam potentiae sunt in \emph{conjuncto}, sicut in \emph{subjecto}, sieut omnes potentiae \emph{sensitivae} partis et \emph{nutritivae}. Destructo autem \emph{subjecto}, non potest \emph{accidens} remanere. Unde, corrupto conjuncto, non manent hujusmodi potentiae \emph{acfw}, sed \emph{virtute} tantum manent in anima, sicut in \emph{príncipio} vel \emph{radice}. Et sic falsum est quod quidam dicunt, hujusmodi potentias in anima remanere, etiam corpore corrupto; et multo falsius quod dicunt, etiam actus harum potentiarum remanere in anima separata; quia talium potentiarum nulla est actio nisi per organum corporeum ». (\emph{Sum}. \emph{Th}., p. I, q. 77, a. 8). --- Por isso a alma separada não vé com os olhos corporeos, nem cheira, nem produz os actos da imaginação ou da memoria sensitiva, nem gosa ou se entristece pelo appetite sensitivo, embora recorde o passado pela memoria intellectual e gose ou se entristeca pelo appetite intellectual. --- \$8. Thomaz explica noutro logar (\emph{De anima}, a. 19) o sentido d'aquella expressão --- que as faculdades \emph{sensitivas} e \emph{nutritivas} permanecem na alma separada em \emph{raiz} ou em \emph{virtude} e diz: « Hujusmodi potentiae dicuntur in anima separata remanere, \emph{ut in radice}; non quia sint \emph{actu} in ipsa, sed quia anima separata est talis \emph{virtutis}, ut si uniatur corpori, iterum potest causare has potentias, sicut et vitam ».}.

%%>>> t253

\tese{253}{A alma, separada do corpo, produz as operações intellectuaes.}

a) A alma, separada do corpo, produz algumas operações. Ora não produz as operações \emph{organicas}. Logo produz as operações \pg{406}  \emph{intellectuaes}; aliás não produziria operação alguma, porque todas as operações humanas ou são \emph{organicas} ou \emph{intellectuaes}.

b) As faculdades \emph{intellectuaes}, que são a \emph{intelligencia} e a \emph{vontade}, n\&o só derivam da alma, como do seu \emph{principio}, mas tambem residem na alma, como em seu \emph{sujeito}; pois que propriedades \emph{espirituaes} só podem residir num sujeito \emph{espiritual}. Se as faculdades \emph{zntellectwaes} residem na alma, como em seu \emph{sujeito}, é evidente que ella, separada do corpo, continua a produzir as operações \emph{intellectuaes} \footnote{Portanto a alma separada \emph{entende e quer}. --- O conhecimento \emph{intellectual} é o producto de uma faculdade \emph{espiritual} ou \emph{inorganica}. E, se agora a nossa intelligencia depende das faculdades sensitivas, esta dependencia é apenas \emph{objectiva}, emquanto as faculdades \emph{sensitivas} apresentam o \emph{objecto} do conhecimento. Por isso a alma separada, conservando a propria natureza, continua a produzir as operações \emph{intellectuaes}, e, tendo mudado o modo de existir, s muda o modo de operar, recebendo o \emph{objecto} dos conhecimentos, não das faculdades \emph{sensitivas}, mas do Auctor da natureza, como explicaremos d'aqui a pouco. --- E' claro que a alma, separada do corpo, \emph{raciocina}, \emph{reflecte} sobre si mesma, e \emph{recorda} os conhecimentos \emph{passados}; pois á \emph{razão}, a \emph{consciencia} e a \emph{memoria intellectual} são, como dissemos, diversas funcções de uma e mesma faculdade, que é a \emph{intelligencia}.\par S. Thomaz é neste ponto, como em todos os outros, de uma clareza e precisão admiravel Antes de tudo ensina que a \emph{intelligencia} e a \emph{vontade} permanecem na alma separada: « \emph{Omnes} potentiae comparantur ad animam solam sicut ad \emph{principium}. Sed \emph{quaedam} potentiae comparantur ad animam solam sieut ad \emph{subjectum}, ut \emph{intellectus} et \emph{voluntas}; et hujusmodi potentiae necesse est quod maneant in anima, corpore destructo » (\emph{Sum}. \emph{T}\textasciicircum{}., p. I, q. 77, a. 8). Em seguida adverte que a alma separada não só entende, mas é mais livre é prompta no exercicio das suas funcções intellectuaes: « Anima separata est quodammodo \emph{liberior} ad intelligendum, in quantum per gravedinem et occupationem corporis a puritate intelligentiae impeditur » (\emph{Sum}. \emph{Th}., p. I, q. 89, a. 2 ad 1). E dá a prova d'essa afirmação, dizendo --- que o entendimento compete á alma não emquanto é forma do corpo, mas emquanto é espirito, --- que o mesmo entendimento não se exerce pelo concurso intrinseco de algum orgão corporeo, mas pela exclusiva energia do espirito, --- e que por isso á actividade e a perspicacia da intelligencia devem estar na razão \emph{directa} do seu afastamento da materia: « Quantum ad actum intelligibilem habet (anima) aliquam perfectionem, quam habere non potest dum est corpori unita. Nec hoc est inconveniens; quia operatio intellectualis competit animae, secundum quod supergreditur corporis proportionem: intellectus enim non est actus alicujus organi corporalis » (\emph{De anima}, a. 17 ad 1).\par \emph{Rosmini} affirma que a alma separada se acha reduzida a \emph{um estado de trevas}, \emph{de somno perpetuo}, \emph{destituida de reflexão sobre si mesma}, \emph{destituida de toda á consciencia}, \emph{como se não existisse} (prop. 28). A alma \emph{existe}, diz\pgn{407}este escriptor, mas \emph{não vive}: « A alma, destituida de \emph{sentimento} (em que, na opinião d'elle, consiste a vida), e concentrada na intuição da idéa, é destituida da sua propria vida » (\emph{Teodicea}, p. 641). --- E' curioso este \emph{Jtosmini}! Be a alma existe, vive; se não vive, não existe; porque, como dissemos, para os entes vivos, \emph{ser é viver}. Se a alma não vive, é morta, o por isso é mortal.}.

%%>>> t254

\tese{254}{A alma, separada do corpo, conhece por meio das especies intelligiveis.}
O conhecimento exige a união do objecto conhecido com a faculdade cognoscitiva. Logo a nossa alma, separada \pg{407}  do corpo, não póde conhecer objecto algum, sem que este se lhe una. Ora o objecto não póde unir-se á alma por meio da propria substancia. Logo deve unir-se por meio de uma \emph{imagem} ou especie \emph{Intelligivel}. Logo a alma, separada do corpo, conhece por meio das especies \emph{Intelligiveis}.

As \emph{especies intelligiveis}, por meio das quaes a alma separada conhece, são as que ella formou durante a vida mortal e as que \textsc{Deus} lhe infunde na intelligencia depois da morte \footnote{Advirta-se que só fallamos no conhecimento \emph{natural} da alma separada, isto é; no conhecimento que compete a toda a alma separada, \emph{segundo a evigencia da sua natureza}. Por quanto, o meio \emph{ideal} do conhecimento \emph{sobrenatural} dos \emph{bemaventurados} não é a especie \emph{intelligivel} do objecto conhecido, mas é a propria Essencia de Deus, contemplada \emph{intuitivamente}.\par A alma, separada do corpo, entende. Mas qual é o meio \emph{ideal}, pelo qual ella descobre os objectos, aos quaes se extende a sua energia intellectual? E' porventura, a propria Essencia divina? Seria excellente esse meio, mas não é proporcionado á nossa capacidade. Exclusa esta hypothese, que é a dos \emph{ontologistas}, devemos dizer que o meio ideal do conhecimento natural da alma separada consiste nas formas ideaes, representativas dos objectos, isto é, nas \emph{especies intelligiveis}, que, actuando na intelligencia, determinam esta faculdade \& respectiva percepção.\par E d'onde derivam para a alma separada as \emph{especies intelligiveis}? Os Anjos receberam immediatamente de Deus taes especies desde o principio da sua creação; e este meio de conhecimento é muito connatural a esses Espiritos, que são substancias absolutamente immateriaes (Cf. \emph{Sum}. \emph{Th}, p. l, q. 55, a. 2). --- Mas a alma humana separada não é Anjo, e, d'outra parte, já não pode extrahir dos \emph{phantasmas} as especies intelligiveis, pois, como vimos, a alma separada não pode exercer os actos das faculdades organicas. Como poderá, pois, a nossa alma adquirir as especies, indispensaveis para o conhecimento intellectual? S. Thomaz resolve a questão. Diz elle que a alma, separada do corpo, se não se muda emquanto á \emph{natureza}, deve mudar-se e se muda emquanto ao \emph{modo de operar}; porque o modo de operar é proporcionado ao modo de existir. Ora, se para a alma, unida com o corpo, é connatural que entenda pela especies intelligiveis, abstrahidas dos phantasmas, --- para a alma, separada do corpo, é connatural que entenda pelas especies intelligiveis, infusas por Deus, como entendem os Anjos (\emph{Sum}. \emph{Th}., p. L, q. 89, a. 1). --- Este novo modo de conhecer é connatural á nossa alma, não absolutamente emquanto é alma, mas emquanto é alma separada do corpo. Diz S. Thomaz: « Haeo coguitio est sibi naturalis, non simpliciter, sed in quantum est separata » (\emph{De anima}, a. 1V ad 2).\par Aqui levanta-se uma gravissima difficuldade. Se a alma, separada do corpo, recebe de \textsc{DEus} as especies \emph{intelligiveis}, então a sua união com o corpo é prejudicial para ella; pois é mais perfeito ' conhecer as coisas por\pgn{408}meio de especies \emph{intelligiveis}, infusas por \textsc{Deus}, do que por meio de especies \emph{intelligiveis}, abstrahidas dos phantasmas.\par S. Thomaz responde a essa diffieuldade na \emph{Sum}. \emph{Th}. p. 1, q. 89, a. 1. Diz o S. Doutor que, embora, \emph{absolutamente} fallando, o conhecimento por meio de especies \emph{ntelligiveis}, infusas por \textsc{Deus}, seja mais perfeito que o conhecimento por meio de especies \emph{intelligiveis}, abstrahidas dos \emph{phanthasmas}, todavia, \emph{relativamente} \& nossa alma, este segundo modo de conhecer é mais perfeito do que o primeiro. Na verdade, a intelligencia, para entender, precisa de uma \emph{forma} ou \emph{especie espiritual}, que represente o objecto entendido. --- Se a, intelligencia é divina, essa forma \emph{espiritual} é unica, pois é a propria essencia divina; visto que \textsc{DEus}, contemplando a sua essencia, vé todas as coisas. --- Se a intelligencia é finita, as formas ou especies são \emph{accidentes} e diversas entre si em numero e perfeição. Por quanto, as especies \emph{intelligiveis}, infusas nas intelligencias \emph{superiores}, são poucas mas dotadas de uma representação mais extensa ou universal; ao passo que as especies \emph{intelligiveis}, infusas nas intelligencias \emph{inferiores}, sao mais numerosas mas menos uni versaes. --- D'onde se segue que, se as intelligencias \emph{inferiores} recebessem as \emph{formas} ou \emph{especies intelligiveis}, que são proprias das intelligencias \emph{superiores}, a sua sciencia seria confusa, porque não chegariam a distinguir todos os objectos, representados por aquellas especies. Isto acontece continuamente entre nós; os homens rudes e ignorantes não descobrem numa idea \emph{universal} tudo o que nella descobrem os homens cultos e disciplinados, mas é necessario apontar-lhes \emph{cada uma} das verdades, contidas nessa idea \emph{universal}. Ora \& intelligencia da alma humana é infima entre todas as intelligencias; porque, a nossa intelligencia, embora \emph{possa} entender sempre, comtudo não entende sempre \emph{actualmente}, como entendem as intelligencias \emph{separadas} da materia. Sendo a nossa intelligencia infima entre todas, é claro que, se ella, desde o principio, tivesse recebido de \textsc{Deus} as especies \emph{iztelligiveis}, como as receberam as intelligencias separadas, os nossos conhecimentos seriam imperfeitos e confusos, porque aquellas especies teriam excedido e como que deslumbrado a força natural da nossa alma. Para que, pois, os nossos conhecimentos fossem distinctos e especificos, era necessario que a nossa alma pouco a pouco recolhesse e abstrahisse dos objectos visiveis as especies \emph{intelligiveis}. Mas os objectos visiveis não podem ser percebidos senão pelas faculdades \emph{sensiítivas}, e estas residem, como em seu \emph{sujeito}, no \emph{composto} de alma e de corpo. Logo, para a perfeição da nossa vida \emph{intellectual}, foi necessario que a alma fosse unida ao corpo.\par D'estas doutrinas resulta que, embora a nossa alma, separada do corpo, seja mais livre, prompta e efficaz no exercicio das funcções intellectuaes, comtudo os seus conhecimentos são menos determinados e mais confusos do que os que adquire na sua união com o corpo; não porque as especies \emph{intelligiveis}, infusas por \textsc{DEus}, sejam indeterminadas e imperfeitas, mas porque, como dissemos, excedem a capacidade natural da nossa alma, a qual não póde descobrir e distinguir por ellas tudo o que as intelligencias superiores descobrem e distinguem. --- "Todavia esses conhecimentos, provenientes da infusão das especies \emph{intelligiveis}, tornam-se mais determinados e distinctos, e por isso\pgn{409}mais perfeitos, não só por causa das antigas especies \emph{intelligiveis} (que a alma formou, quando unida ao corpo)e do habito de distinguir o particular no universal, mas tambem porque \textsc{Deus}, como premio das boa sobras, infunde na intelligencia da alma justa novos e particulares conhecimentos, que a alma não póde adquirir com as suas forças naturaes.\par A alma, separada do corpo, entende --- não só por meio das especies \emph{intelligiveis}, infusas por \textsc{DEus}, --- mas tambem por meio das especies, formadas durante a vida mortal e conservadas na memoria \emph{intellectual}, que, por ser uma faculdade \emph{immaterial}, permanece na alma separada, --- e tambem vendo as substancias separadas. Diz S. Thomaz: « Anima post mortem tribus modis intelligit: uno modo per species, quas recepit a rebus, dum erat in corpore; alio modo per species in ipsa sui separatione a corpore sibi divinitus infusas; tertio modo videndo substantias separatas » (\emph{De verit}, a. 1; Cf. \emph{Sum}. \emph{Th}. p. I, q. 79, aa. 5 e 6).\par Pergunta-se: A intelligencia \emph{agente} permanece na alma \emph{separada}?\par Respondemos que permanece. Por quanto, sendo uma faculdade espiritual e uia propriedade da alma, a intelligencia \emph{agente} não pode corromper-se nem separar-se da alma. E se a sua luz não se reflecte nem nas especies, infusas por \textsc{Deus}, nem nas especies adquiridas na vida mortal, porque estas especies são por si luzes, reflecte-se, comtudo, e concentra-se na propria essencia da alma, \& qual, unida com o corpo, era intelligivel \emph{em potencia}, e, separada do corpo, se torna intelligivel \emph{em acto}, por si mesma, sem o intermedio de uma especie, que a represente. Esta illuminação da alma pela intelligencia \emph{agente} é necessaria. Por quanto, embora a alma seja uma substancia \emph{espiritual} e, como tal, não careca de illuminação, comtudo o seu ser, diz S. Thomaz, é o proprio ser do composto: « animae humanae remanet esse compositi post corporis destructionem » (\emph{Sum}. \emph{TA}., 1-2, q. 4, a. 5 ad 2), e conserva sempre uma relação ao corpo; e, por isso, não é totalmente livre de qualquer potencialidade e imperfeição, não tem aquelle grau de pureza necessario para a intelligibilidade, e, para a remoção d'essa potencialidade ou imperfeição, precisa de ser illuminada, tornando-se intelligivel \emph{em acto}. --- Além d'esse, a intelligencia \emph{agente} não produz outro acto por lhe faltar a materia, sobre a qual possa exercer a sua energia, até a resurreicüo dos corpos.}.

%%>>> t255

\tese{255}{A alma, separada do corpo, conhece as causas immateriaes do mundo corporeo.}
A alma, separada do corpo, deve possuir os conhecimentos, que são proporcionados á sua natureza. Ora o \pg{408} \pg{409}  conhecimento das causas \emph{immateriaes} do mundo corporeo é proporcionado á natureza da nossa alma. Logo a alma, separada do corpo, conhece as causas \emph{Zmmateriaes} do mundo corporeo \footnote{A alma separada, participando do modo de existir proprio das substancias angelicas, participa tambem do modo de conhecer proprio d'essas substancias. Ora as substancias angelicas percebem pelas especies infusas e percebem as \emph{causas immateriaes}, ou as \emph{razões intelligiveis} do mundo corporeo. Diz S. Thomaz: « Assim como das mãos de Deus sahiram as coisas para que subsistissem na sua natureza, assim tambem da Sabedoria infinita derivaram para as substancias intellectuaes as formas das coisas, pelas quaes (formas) percebessem as coisas » (\emph{De anima}, a. 18). Por isso na intelligencia angelica encontra-se da um modo \emph{immaterial} o que de um modo \emph{material} se encontra no mundo corporeo. Estas formas ideaes por si representam as \emph{razões intelligiveis}, isto é, as \emph{essencias} das coisas, o não se extendem aos \emph{individuos} senão \emph{per accidens}, emquanto os individuos são necessarios para a realisação concreta d'essas formas. Por quanto, a ordem ideal deve corresponder á ordem\pgn{410}real; ora o que a natureza tem em vista, na ordem real das coisas, é á \emph{essencia}, é \& \emph{especie}, e não se extende aos \emph{individuos} senão emquanto são necessarios para a existencia concreta da mesma \emph{essencia}, ou \emph{especie}. Be, pois, a alma separada entende como entendem e o que entendem as substancias angelicas, deve ella conhecer pelas especies infusas e deve conhecer as causas \emph{immateriaes} do mundo corporeo. --- Todavia o conhecimento, que 2 alma separada tem d'essas causas, não é tão clero e distincto, como é o dos Anjos. Na verdade, as especies infusas, representativas das causas immateriaes do mundo corporeo, não são, como dissemos, inteiramente proporcionadas á intelligencia humana, habituada a recolher das proprias coisas materiaes as especies intelligiveis; e por isso o conhecimento, que a alma separada tem dos particulares, contidos nas especies infusas, não é claro e distincto, mas generico e indeterminado (Cf. S. Thom. /\emph{oc}. \emph{cit}.).}.

%%>>> t256

\pg{410} \tese{256}{A alma, separada do corpo, conhece os objectos materiaes na sua singularidade.}
Por quanto a alma, separada do corpo, percebe por meio das especies \emph{enfe/ligiveis}, que são infusas por Dzvus, e que são uma .semelhança, embora remotissima, das ideas arehetypas da intelligencia divina. Ora estas ideas arehetypas representam as coisas \emph{materiaes}, como são em si, isto é, na sua \emph{singularidade}. Logo a alma, separada do corpo, conhece os objectos \emph{materiaes} na sua \emph{singularidade} \footnote{A nossa alma, no estado de união com o corpo, não pode, como dissemos, conhecer os objectos materiaes na sua singularidade; porque, agora, ella entende por meio de especies intelligiveis, abstrahidas dos phantasmas e despidas das condições materiaes, e por isso representativas do universal e do immaterial,e não do singular e do material. Mas, depois da separação, a nossa alma não entende pelas especies abstrahidas dos phantasmas, mas pelas especies, infusas por Deus. Ora estas especies, infusas por Deus, representam o \emph{singular} e o \emph{material}, como tal; e por isso a alma, separada do corpo, percebe os objectos \emph{materiaes} na sua \emph{singularidade}. Oucamos a magnifieca theoria de S. Thomaz a esta respeito. As formas ou especies ideaes, pelas quaes a intelligencia conhece, podem ser --- ou meramente \emph{representativas} --- ou \emph{productivas}, conforme são \emph{copias} ou \emph{typos} de um objecto. As primeiras informam a intelligencia de quem considera um objecto jà existente; as segundas pertencem ao artista, que executa uma obra sob a influencia do modelo concebido. Ora, assim como a forma ideal, que é \emph{copia} de objecto, leva ao conhecimento d'elle na medida e no modo por que o representa; assim tambem a forma ideal, que é /\emph{ypo} do objecto, leva ao conhecimento d'este na medida e no modo por que o produz. Por isso, o artista, que só dà a uma obra a forma, e não a materia, pela idéa e na idéa só pode ver \& sua obra em ordem á forma, e não em ordem á materia. E, como nenhum artista creado é capaz de produzir a materia, mas suppõe-na preexistente, segue-se que elle, no seu conceito typico, não conhece a sua obra senão de um modo \emph{universal}; porque, em virtude d'esse conceito, só percebe a forma, capaz de ser repetida muitas vezes e em muitos logares. A concretisação \emph{individual}, que a obra recebe da materia e que a determina a um certo tempo e a um certo logar, é conhecida pelo artista mas só pela percepção sensitiva, --- Pelo contrario, o artista increado, Deus, pela sua intelligencia\pgn{411} não só produz nas coisas a \emph{forma}, que é principio de \emph{universalidade}, mas tambem a \emph{materia}, que é principio de \emph{singularidade}, e por isso, nas idéas archetypas da sua arte, conhece, ao mesmo tempo, o universal e o singular, com respeito a todas as coisas creadas. Ora é d' esta arte divina que derivam para as substancias intellectuaes as formas ou especies, pelas quaes ellas percebem por um modo, que lhe é connatural; porque, como dissemos, assim como da arte divina procedem as coisas materiaes, para que subsistam nas proprias naturezas, assim tambem da mesma arte derivam para as substancias intellectuaes as formas intelligiveis, pelas quaes as coisas são conhecidas por essas substancias como são produzidas por Deus, e por isso na sua universalidade e singularidade, emquanto á forma e emquanto á materia. E, visto a nossa alma, depois de separada do corpo, entender como e o que entendem as substancias angelicas, concluimos que ella conhece, e de um modo \emph{directo}, isto é, pela especie \emph{propria}, os objectos \emph{materiaes} na sua \emph{singularidade} (Cf. \emph{De anima}, a. 20).\par Todavia, para a alma, separada do corpo, poder conhecer os objectos \emph{materiaes} na sua \emph{singularidade}, não basta que \textsc{Deus} infunda as especies \emph{intelligiveis} ; porque a nossa alma, tendo uma natural capacidade para entender pelas formas \emph{universaes}, abstrahidas dos phantasmas, e não pelas formas \emph{infusas}, não póde distinguir nestas as especies infimas, e, muito menos, os objectos individuos e singulares; mas é necessario que essas formas \emph{infusas} recebam uma determinação e concretisação --- ou por meio dos conhecimentos, que foram adquiridos na vida, e que se referem a objectos \emph{materiaes} e \emph{singulares}, --- ou pelas affeioóes a objectos determinados, como aos parentes, amigos, etc., --- ou pela ordem que a alma conserva para o proprio corpo, --- ou pela Providencia de \textsc{Deus}, que póde dirigir os pensamentos da alma para um objecto particular (Cf. S. Th. \emph{De anima}, a. 90).\par Pergunta-se: A alma separada conhece os acontecimentos humanos?\par Respondemos que a alma separada não póde, na ordem da \emph{natureza}, conhecer os acontecimentos humanos, se \textsc{Deus} e os outros espiritos lh'os não manifestarem; porque, sahindo a alma do corpo, acabam as suas communicações com o mundo material e só faz parte do mundo espiritual (\emph{Sum}. \emph{TR}. p: 10:89, a 8y}.

%%>>> t257

\pg{411} \tese{257}{A alma, separada do corpo, conhece as substancias espirituaes.}
A alma humana, separada do corpo, deve conhecer os objectos, \emph{actualmente} intelligiveis e proporcionados á sua faculdade cognitiva. Ora os objectos \emph{actualmente} intelligiveis e proporcionados á faculdade cognitiva da alma separada são as substancias \emph{espirituaes}. Logo a alma, separada do corpo, conhece as substancias \emph{espirituaes}. --- Essas substancias são a \emph{essencia} da propria \emph{alma}, os outros \emph{espiritos}, e \textsc{Deus} \footnote{A alma separada conhece a sua \emph{essencia}, os outros \emph{espiritos} e \textsc{Deus}.\par a) Conhece a sua \emph{essencia}. --- A alma separada deve conhecer tudo o que é intelligivel em \emph{acto}, que é objecto \emph{proporcionado}, e que lhe está unido pela propria substancia ou por uma sua especie. Ora a essencia da alma, com relação á propria alma, ---- é intelligivel \emph{em acto}, porque é illuminada pela intelligencia \emph{agente}, --- 6 objecto \emph{proporcionado}, porque é \emph{identico} ao principio perceptivo, --- e está unida com esse principio perceptivo, porque,\pgn{412}como dissemos, é uma coisa só com elle. Logo a alma separada conhece a sua essencia. --- E conhece-a \emph{por si mesma}, sem necessidade de uma especie ou forma, que a represente. Na verdade, a especie ou forma, se é necessaria, quando o objecto é \emph{distincto} do principio cognoscente, não é necessaria, quando um se identifica com o outro, como no nosso caso.\par Este conhecimento, que a alma tem de si mesma, é \emph{intuitivo} e \emph{adequado}.\par E' \emph{intuitivo}. --- Antes de mais nada, notamos que o conhecimento pode ser \emph{intuitivo} ou \emph{abstractivo}. E' \emph{intuitivo}, --- se é causado pela acção do objecto, quer \emph{immediata}, isto é, proveniente directamente do \emph{objecto}, quer \emph{mediata}, isto é, proveniente directamente da \emph{especie propria} do objecto, --- e se attinge o mesmo objecto, realmente presente. E' \emph{abstractivo}, --- se ou não é causado pela acção (immediata ou mediata) do objecto, ou não tem por termo directo o mesmo objecto, realmente presente, --- mas attinge directamente propriedades ou manifestações do objecto, e só indirectamente o proprio objecto. --- São \emph{duas}, pois, as condições necessarias e sufficientes para o conhecimento \emph{intuitivo}. \emph{A primeira} é que seja causado pela acção \emph{immediata} do objecto, ou pela especie \emph{propria} d'elle, isto é, pela especie, que gera um conceito \emph{proprio}, que exprime o objecto \emph{como é em si mesmo} (e não por meio de negações ou analogias), embora nao represente sempre a \emph{essencia}. \emph{A segunda} é que perceba \emph{directamente} o proprio objecto \emph{presente}; de modo que não pode chamar-se \emph{intuitivo} o conhecimento de um objecto, que é conhecido noutro objecto, anteriormente percebido. --- Feita esta distincção, devemos dizer que o conhecimento, que a alma tem agora da sua \emph{essencia}, é apenas \emph{abstractivo} ou \emph{arguitivo}. Na verdade, a nossa alma, no estado da união, não se percebe a si mesma senão reflectindo sobre os seus actos perceptivos e vitaes (S. Th. \emph{De mente}, a. 8). Por quanto, sendo forma do organismo, ella tem agora as olhos virados para o mundo corporeo, e, quando se concentra em si mesma, vê-se actuada por muitos pensamentos, relativos aos objectos sensiveis e por isso não pode perceber a propria essencia senão atravez esses pensamentos (\emph{Sum}. \emph{Th}., p. I, q. 87, a. 1). Além d'isso, a alma, no estado de união, pelo facto de attingir o intelligivel não em si mesmo mas no sensivel, percebe a substancia só por meio dos phenomenos, por que esta se manifesta e de que é principio e fundamento, e por isso só pela existencia e qualidade dos seus actos pode perceber a propria existencia e a qualidade do seu ser. Logo a percepção, que a nossa alma tem agora de si mesma, é \emph{abstractiva} ou \emph{arguitiva}.... Mas, separada do corpo, a alma muda o seu modo de operar, como muda o seu modo de existir. Então já não olha para o mundo corporeo, mas para os objectos intelligiveis, como são em si mesmos. Ora, entre os objectos intelligiveis, o mais intimo e o mais proporcionado é á sua \emph{essencia}. Logo a alma separada tem conhecimento da sua \emph{essencia}. Este conhecimento, sendo causado pela acção \emph{immediata}, que a alma, \emph{actual} e \emph{realmente} presente, exerce sobre si mesma, é \emph{intuitivo} (\emph{De anima}, a. 17).\par E' \emph{adequado}. Se o conhecimento, que na vida presente temos da nossa alma, é \emph{abstractivo}, é tambem e necessariamente \emph{inadaquado}; porque o que se conhece de um modo \emph{abstractivo} conhece-se de um modo \emph{parcial}, e não como\pgn{413}é \emph{em si mesmo}. --- Mas o conhecimento, que a alma separada tem da sua essencia, por ser \emph{intuitivo}, \emph{é adequado}. De facto, diz-se \emph{adequado} o conhecimento.de um objecto, quando a forma ou especie, que determina o principio cognoscente, representa \emph{adequadamente} o proprio objecto, do modo que nada se encontra no objecto, que a forma não exprima. Ora a forma, que determina a alma ao conhecimento da sua \emph{essencia}, \emph{é a} mesma \emph{essencia}. Logo 9 conhecimento, que a alma separada tem da sua \emph{essencia}, \emph{é adequado}.\par Como possa a nossa alma ter um conhecimento \emph{intuitivo} da sua \emph{essencia} depois da separação, ao passo que o não pode ter no estado de união, explica-o o Card. Gaetano no seu Commentario \& \emph{Summa Theologica}. A. alma humana, diz o eminente theologo, pela sua propria natureza, é intelligivel \emph{em potencia}. De facto, a alma é, pela sua natureza, forma do corpo, e por isso segue a condição das formas materiaes. Ora estas são intelligiveis \emph{em potencia}; porque, sendo unidas com a materia, não podem actuar immediatamente, e por si, na intelligencia, nem exercer o munus da especie intencional. Logo tambem a nossa alma, no estado de união, é só intelligivel em \emph{potencia}, e não é conhecida por si mesma senão atravez os phantasmos e de um modo abstractivo e inadequado. --- Separada, porém, do corpo, á nossa alma pode receber e recebe directamente sobre a sua essencia a luz da intelligencia \emph{agente}; de modo que ella, que era intelligivel em. \emph{potencia}, se torna; em virtude d'esta illuminacàáo, immediatamente intelligivel \emph{em acto}, sem que a intelligendia precise do auxilio de outra especie » (In I p., q. 89, a. 2).\par b) Conhece os outros \emph{espiritos}. --- E isto é muito natural. Fazendo parte do mundo \emph{espiritual}, a alma separada deve entrar em communieacüo com as substancias, que o compõem, e tal communicacio não pode conceber-se sem um mutuo conhecimento.\par Esta substancias \emph{espirituaes} são as outras \emph{almas separadas}, e os \emph{Anjos}. --- Das outras \emph{almas separadas} a alma tem um conhecimento \emph{adequado} ; pois a especie intelligivel, que representa uma outra alma, imprime na alma cognoscente tudo o que essa especie representa, visto que entre a alma conhecida e a cognoscente, por serem \emph{especificamente} eguaes, ha uma proporção perfeita, e tudo é recebido em conformidade com o modo e com a aptidão do sujeito, que recebe. --- Dos Anjos a alma separada tem um conhecimento \emph{inadequado} ; porque a especie intelligivel, que representa o Anjo, não imprime na alma separada tudo o que ella (especie) representa, pois o \emph{Anjo é essencialmente} superior \& alma humana, e, como dissemos, \emph{quidquid recipitur per modum recipientis recipitur}. --- E' claro que fallamos sempre no conhecimento de ordem \emph{natural}; porque não pode dizer-se o mesmo do conhecimento da ordem \emph{sobrenatural}, que a alma adquire pela \emph{luz da gloria} (\emph{Sum}. \emph{Th}., p. I, q. 89, a. 2).\par A alma separada conhece as outras almas e os Anjos --- não só por \emph{si mesma}, mas tambem, e sobretudo, pelas especies infusas por Deus. Estas especies são \emph{necessarias}, Por quanto, a alma separada representa por si mesma a essencia \emph{especifica} das outras almas, mas não representa a existencia e as notas \emph{individuaes}, que estão fóra da essencia e que variam em cada individuo\pgn{414}. --- Se a alma não representa tudo o que se encontra nas outras almas, muito menos representa tudo o que se encontra nos Anjos. Na verdade, a alma separada, se representa os predicados \emph{communs} e \emph{genericos}, nos quaes convêm todas as substancias \emph{espirituaes}, não representa a differença \emph{especifica}, que constitue o Anjo e o distingue da alma humana.\par Entrando em communicação com as outras almas e com os Anjos, a alma separada deve poder \emph{fallar} com essas creaturas. Mas de que modo podem fallar os espiritos? S. Thomaz ensina que o espirito falla, dirigindo e manifestando, por um acto da vontade, o proprio pensamento a outro espirito. O homem falla a outro homem, dirigindo-lhe a palavra \emph{oral}; o espirito falla a outro espirito, dirigindo-lhe a palavra \emph{mental} (\emph{Sum}. \emph{Th}., p. I, q. 107, a. 1). O espirito que falla não infunde nem produz coisa alguma no espirito, com quem falla; mas este, por meio da especie, que possue, conhece o espirito que falla e o pensamento que lhe dirige (\emph{De verit}., q. 9, a. 6 ad 4).\par 0) Conhece \textsc{Deus}. E conhece-o como o conhecem os Anjos, embora, sempre pelos \emph{effeitos}, e não \emph{intuitivamente}. --- S. Thomaz distingue \emph{tres} modos de conhecer uma coisa pelos seus effeitos. O \emph{primeiro} é quando o effeito é um \emph{meio} para o conhecimento da causa, de modo que do conhecimento de um se passa ao conhecimento da outra; e isto verifica-se na vida presente. O \emph{segundo} é quando no effeito \emph{se} vé a causa, como na copia o typo; assim vemos uma pessoa no espelho, onde se reflecte a sua imagem. Este segundo modo differe do primeiro; porque no 1º modo temos dois conhecimentos distinctos e consecutivos, isto é, do effeito e da causa, ao passo que no 2º modo ha um só acto perceptivo, pelo qual se percebe, ao mesmo tempo, o effeito e a causa. O \emph{terceiro} é quando a semelhança da causa, que resplandece no effeito, é a forma intelligivel, pela qual o mesmo effeito conhece a sua causa. Ora o \emph{primeiro} d'estes modos não compete á alma separada no conhecimento que ella tem de Deus; porque, em tal caso, o seu conhecimento seria \emph{discursivo}, e por isso não semelhante ao dos Anjos. Mas competem-lhe o \emph{segundo} e o \emph{terceiro}, emquanto conhece Deus nas outras substancias espirituaes e em si mesma. --- Este conhecimento não é, não pode ser \emph{adequado}, porque nenhuma creatura pode representar Deus como é em si mesmo e nenhuma intelligencia finita é capaz de \emph{comprehender} o infinito, --- mas é muito superior ao que temos agora. Diz S. Thomaz: « A alma separada conhece que Deus é, que Deus é \& causa de todo o creado, que é livre de todo o limite e imperfeição, que excede infinitamente na perfeição o que existe ou pode ser concebido pela inteligencia creada: --- verdades todas que de algum modo á nossa alma conhece tambem no seu estado actual » (\emph{C}. \emph{Gentes}, l. III, c. 49). --- Na ordem \emph{natural}, \emph{se existisse}, a alma separada conheceria Deus tambem pelas especies intelligiveis, infusas immediatamente pelo proprio Deus, conforme os meritos das obras bóas, praticadas na vida mortal. Dizemos --- se \emph{existisse}; porque o homem, pela divina bondade, foi elevado á ordem \emph{sobrenatural}, a qual contem eminentemente a ordem \emph{natural}, e, neste estado de elevação, tal infusão de especies intelligiveis é inutil, por serem as almas dos justos admittidas á visão \emph{intuitiva} de Deus.}.

%%>>> t258

\tese{258}{A alma separada é immutavel na vontade.}
A alma foi unida ao corpo, para que, nelle e com elle, se ordenasse e dirigisse para o seu ultimo fim, que só na vida futura pode ser possuido perfeitamente. Logo a vida presente é um estado de \pg{412} \pg{413}  \emph{via} e de \emph{provação}, e a vida futura é um estado de \emph{termo} e de \emph{recompensa}; e a morte, sendo a dissolução do composto humano, marca o fim de um estado e o começo do outro. Se, no estado de \emph{via} e de \emph{provação}, a alma pode estar sujeita a \pg{414}  perpetuas mudanças, não o pode no. estado de \emph{termo} e de \emph{recompensa}, e esta potencia ou impotencia funda-se na propria natureza desses dois estados. D' onde se segue que a alma será sempre, por toda a eternidade, naquella disposição, rela-

tivamente ao objecto escolhido como fim ultimo, na qual se

encontrou no momento da separação. Logo a alma separada

é immutavel na sua vontade \footnote{A alma está unida ao corpo. Tal união é util ao corpo, que da alma\par recebe toda a perfeição e o proprio ser. Mas esta utilidade não é o fim prin-\par cipal da união; porque a forma não é para a materia, mas a materia é para a forma. Portanto o bem e a vantagem da alma é o que o Creador primeira 9 principalmente teve em vista na formação do composto humano. O bem 9 á vantagem, que a alma pode receber da sua união com o corpo, consiste nisto que ella, no corpo e com o corpo, podo e deve naturalmente alcançar a perfeição conveniente. Ora, sendo a perfeição o bem do ente e tendo o bem a razão de fim, é evidente que a alma foi unida ao corpo para, nelle e com elle, alcançar o seu ultimo fino. Logo, separada do corpo, a alma suspende € termina o seu movimento para o ultimo fim, e permanece eternamente naquella disposição para com esse fim, que teve no momento da morte. --- Nem podia ser d'outro modo. A nossa vontade segue sempre a intelligencia. Dá-se agora este facto, que, assim como a intelligencia, embora adhira sempre aos primeiros principios, comtudo passa de uma conclusão para outra diversa, e até contraria, assim tambem a vontade, embora tenda sempre para a sua felicidade, comtudo muda constantemente com relação ao objecto concreto, onde julgamos encontrar essa felicidade, porque não ha na terra um bem, que seja o bem infinito, o nosso ultimo fim. Mas, separada a alma do corpo, a intelligencia conhece immediatamente que Deus é o Primeiro, o unico, o infinito Bem, em que se encontra a verdadeira felicidade e que é, por isso, 9 nosso ultimo fim. Conhecendo Deus como o \emph{ultimo} fim, a intelligencia não pode procurar outro fim fóra do Deus, e por isso fixa-se immutavelmente nas suas relações com elle, e, em conformidade com essas relações, a vontade produz os seus actos. Para os bons Deus é o ultimo fim, por elles eternamente possuido; para os maus Deus é o ultimo fim, por elles eternamente perdido. A esta diversidade de conhecimento na intelligencia corresponde una diversidade de appetite na vontade. Por isso, á vontade dos bons não quer senão o que Deus quer, e portanto não quer senão o bem, o bem em tudo e tudo no bem; e a vontade dos maus quer sempre o mal, o mal em tudo e tudo no mal, não porque o mal possa ser ou seja o objecto da vontade, mas porque tudo o que elles querem, querem-no para um fim mau. A . vida dos justos é uma vida de luz, de amor, de alegria; a vida dos maus é uma vida de trevae, de odio, de desespero (Cf. c. \emph{Gentes}, 1. IV, c. 95),\par Agora digamos alguma coisa especial com relação \& immutabilidade da vontade dos \emph{justos}. Que a vontade dos justos, depois da morte, seja immutavel no bem é certissimo. Porquanto, a nossa alma, na vida, sempre é mutavel com relação aos varios bens, porque está sempre \emph{em}. \emph{potencia} com relação ao seu fim ultimo. Mas, alcançando este seu fim ultimo, já não pode estar sujeita a mudanças, desejando outro bem fóra d'esse fim; porque a posse do fim ultimo sacia completamente as tendencias da vontade, de modo que lhe não deixa a liberdade de tender para outro bem. --- Além disso, quem possue tudo, nada pode desejar fóra d'esse tudo. Ora quem possue o bem infinito, tem tudo, e nada pode desejar, que repugne ao bem, em que se encontra a felicidade perfeita. --- E' claro que tudo o que deixamos dito se refere á ordem \emph{natural}, Na ordem \emph{sobrenatural}, a que nos elevou a divina\pgn{416}Bondade, e que é a unica ordem que existe, as almas dos justos, sendo admittidas á visão intuitiva de Deus, amam a Bondade divina com um amor \emph{necessario}, e por isso vontade d'elles é eternamente immutavel no bem. (Cf. c. \emph{Gentes}, 1. IV, c. 92).\par Be os \emph{justos} têm a vontade immutavel no \emph{bem}, os \emph{maus} tem-na no \emph{mal}. Estes detestam as suas culpas, que foram causa de uma pena eterna, e detestam-nas porque não perderam a tendencia para a felicidade; mas não têm, não podem ter uma verdadeira contricção, porque a vontade d'elles é inteiramente mã, por estarem inteiramente separados do centro de toda a rectidão. A vontade dos reprobos detesta a pena mas não a culpa. --- Nem esta immutabilidade da vontade cessará depois da resurreição. Na reunião da alma com o corpo, á alma não será disposta segundo as exigencias do corpo, mas este será disposto segundo as exigencias d'aquella; o corpo participará da sorte da alma, e será seu docil instrumento (C£. c. \emph{Gentes}, 1. IV, co. 93, 95) --- \emph{E}! esta uma terrivel verdade, mas muito conforme com a razão. Oucamos um auctor insuspeito, \emph{Mivart}. Fallando na condemnacao dos reprobos ás penas eternas, diz assim: « Quaes serão as provaveis consequencias d'esse desastre? A razão, entregue a si mesma, muito pouco pode dizer Acerca da alma separada. Não temos razões para provar que a alma possa então desfazer o que fez durante a vida; antes o contrario. Nem essa affirmação pode encontrar apoio na doutrina da conservação das energias; e não ha nada que possa contradizer a esserc\&o da Egreia, quando ensina que o estado, em que se encontra a alma no fim da provação d'esta vida, é immutayel. Se isto é assim, o homem, que morre no estado de aversão relativamente á Luz suprema e ao Bem infinito, deve permanecer sempre nesse estado com todas as suas possiveis consequencias. Poderão dizer que não é necessario que essas consequencias sejam eternas. Mas, se a causa é immutavel como podem mudar as consequencias? Além d' isso, devemos . considerar o que é proprio da Eternidade, quando o tempo tiver acabado. E' possivel que, se podessemos entender o que é a Eternidade, a hypothese relativa a uma mudança na condição da alma nos apparecesse absurda. Ha mais: tal mudança não nos parece compativel com a divina Justiça. Porquanto, todo e qualquer castigo temporal, por quanto comprido, a que succedesse uma bemaventurança eterna, estabeleceria, praticamente, uma egualdade na vida de todos os homens; pois uma multidão de seculos é nada em comparação com a eternidade. A sciencia, pelo menos, não apoia a opinião que admitte a possibilidade de uma mudança nas consequencias de uma acção uma vez realisada » (\emph{Happiness in Hell}, XIX. Cent. n. 190, p. 909).}.

%%>>> t259

\tese{259}{A alma separada é dotada de livre arbitrio.}
O \emph{livre arbitrio} é uma propriedade da nossa vontade. Ora a vontade permanece na alma separada. Logo a alma separada é dotada de \emph{livre arbitrio}. --- Nem a existencia do \emph{livre arbitrio} na alma separada é incompativel com a sua \emph{immutabilidade} ; porque, emquanto a immutabilidade diz respeito ao \emph{fim}, a \emph{eleição}, que é o acto proprio do livre arbitrio, diz respeito aos \emph{meios}. --- Immutavel no mal, a vontade dos reprobos é livre só para o mal, para arrastar-se, na caligem do espirito do mal, no lodo dos vicios: é a liberdade dos escravos. Immutavel no bem, a \pg{415} \pg{416} \pg{417}  vontade dos justos, é livre para o bem, para escolher este ou aquelle bem, para se perder felizmente, sob a luz serena do summo Bem, no pélago de todos os prazeres: é a liberdade dos filhos de \textsc{Deus} \footnote{A \emph{iémmutabilidade} da vontade não é incompativel com o \emph{livre arbitrio}, cujo acto é a \emph{eleição}. De facto, a \emph{eleição}, como dissemos, não se refere ao \emph{fim}, mas aos \emph{meios}. A vontade adhere natural e necessariamente ao seu ultimo fim, como a intelligencia adhere necessariamente aos primeiros principios. E, assim como a natureza é o principio eo fundamento de todas as propriedades de um ente, assim tambem o desejo do ultimo fim é o fundamento e o principio de todos os outros desejos, porque todo o movimento procede de um immovel (\emph{Sum}. \emph{Th}., p. I, q. 82, a. 1). Por isso, se nao é incompativel com o livre arbitrio a immobilidade da vontade, com a qual agora desejamos a felicidade em geral, nem é incompativel com o mesmo livre arbitrio a immobilidade da vontade, com a qual esta faculdade adherirá a um determinado objecto, como ao seu fim ultimo. E, assim como nos é inherente á natureza, pela qual desejamos nesta vida a felicidade em geral, assim tambem na vida futura nos será inherente, como uma segunda natureza, aquella disposição especial, que nos levou a desejar este ou aquelle bem, como o nosso ultimo fim, e, em virtude d'essa disposição, durará eternamente na vontade a tendencia para o bem desejado. Portanto, esta disposição não é contraria ao livre arbitrio, como não é contraria ao livre arbitrio de Deus a tendencia da sua vontade para o seu unico objecto proporcionado, que é a propria Essencia infinita. --- Esta liberdade, que permanece na alma separada, é a de \emph{exercicio}; porque a liberdade de \emph{contrariedade}, que não pertence á essencia do livre arbitrio e que é uma especie de excrescencia morbida, permittida por Deus no tempo da \emph{provação}, não pode permanecer na vida futura, na qual, como dissemos, a vontade é eternamente immutavel no bem ou no mal. --- E', pois, livre a alma dos justos, como é livre a dos reprobos. Mas entro uma e outra liberdade a differenga é enorme. A vontade dos justos, isenta de todo o mal, é livre com relação a todos os bens; a vontade dos reprobos, privada de todo o bem, é livre para todos os males, (C£. c. \emph{Gentes}, l. IV, c. 95).\par 27 i Nor. I1.}.

%%>>> t260

\tese{260}{Corollario pratico.}
A nossa alma, d'aqui a algum tempo, separar-se-ha do corpo, e entrará na eternidade. Todo o nosso empenho deve consistir em que ella se ache, no momento da separação, revestida de graça divina; para que, apresentando-se assim no tribunal divino, possa receber a immareescivel coroa de gloria, que N. S. Jesus \textsc{CurISTO} nos mereceu pela sua dolorosa Paixão e Morte.

\textbf{METAPHYSICA}

%%>>> parte-2

\part*{PARTE QUARTA\\[0.3em]\large THEODICEA}
\addcontentsline{toc}{part}{PARTE QUARTA --- THEODICEA}

\pg{418} \pg{419} PROLOGO

1. \textbf{Theodicea}. --- \emph{Theodicea} é a parte da Metaphysica, que \emph{trata de} \textsc{Deus}, \emph{emquanto é conhecido pela luz da razão}. --- A \emph{TTheodicea} costuma chamar-se tambem \emph{Theologia natural} \footnote{A palavra \emph{theodicea} (\emph{defesa de} Deus), embora pela sua etymologia possa denotar uma parte do tratado de \textsc{Deus} (e neste sentido foi empregada por \emph{Leibnits}, que a formou), todavia hoje toma-se para significar \emph{toda} a sciencia que trata de \textsc{Deus}, servindo-se da luz natural da razão.}.

2. Objecto da Theodicea. --- O objecto da \emph{Theodicea} é \textsc{Deus}, \emph{emquanto conhecido pela luz da razão}. \textsc{Deus} é o Ser necessario, infinitamente perfeito, creador do mundo. --- A \emph{Theodicea}, pois, distingue-se da Theologia \emph{revelada}; porque esta trata de \textsc{Deus}, \emph{emquanto conhecido pela luz da} Fé \footnote{Lembramos que o criterio para distinguir as sciencias não só se baséa no \emph{objecto}, que a\&tingem, mas tambem no \emph{modo}, por que o attingem. D'ahi \& distincção entre a "Theologia \emph{revelada} e a \emph{natural}.}.

3. Excellencia da Theodicea. --- A excellencia das sciencias \emph{especulativas} depende principalmente da excellencia dos seus objectos. Ora o objecto da \emph{Theodicea} é o mais excellente entre os objectos de todas as outras sciencias \emph{racionaes}; porque \textsc{Deus} é o Ser infinitamente perfeito. Logo a \emph{Theodicea} é a mais excellente entre todas as sciencias \emph{racZonaes} \footnote{A excellencia e sublimidade da \emph{Theodicea} é causa da sua difficuldade. \textsc{Deus}, sendo o ente infinito, excede infinitamente a capacidade da nossa intelligencia, que é essencialmente finita. Sabemos que Deus existe, que é o principio e o-fim de todas as coisas, que contém todas as perfeições possiveis e immaginaveis num grau purissimo e infinito; mas esta idea não representa \emph{adequadamente} a essencia de \textsc{Deus}. Todavia um conhecimento\pgn{420}\emph{inadequado} de \textsc{Deus} aperfeiçoa e ennobrece a nossa intelligencia, mais do que o conhecimento \emph{adequado} de todos os seres creados.\par D'ahi se vé que a \emph{Theodicea} é altissima \emph{sabedoria}. Porquanto, a \emph{sabedoria} consiste no conhecimento da causa \emph{absolutamente suprema} dos entes. Uma causa póde ser \emph{relativa} ou \emph{absolutamente suprema}. E' \emph{relativamente} suprema, quando é a mais alta só num certo genero de causas; pois a \emph{causa} divide-se em \emph{material}, \emph{formal}, \emph{efficiente} e \emph{final}. E' \emph{absolutamente} suprema, quando é mais alta do que toda e qualquer outra causa, de modo que todas as outras causas dependem d'ella. Por isso, se o conhecimento das causas \emph{relativamente} supremas é \emph{sabedoria}, o conhecimento da causa \emph{absolutamente} suprema é \emph{altissima sabedoria}. Ora a \emph{Theodicea} trata da causa \emph{absolutamente} suprema, que é \textsc{Deus}. Logo a \emph{Teodicea é altissima sabedoria}.\par Dizemos que a \emph{Theodicea} é mais excellente que todas as outras sciencias \emph{racionaes}, pois a \emph{Sagrada Theologia}, que é mais elevada do que a \emph{Theodicea}, não se guia pela luz da \emph{razão}, mas pela luz da Fé. --- A S. Theologia é mais elevada que a "Theodicea --- pela \emph{certeza} dos \emph{conhecimentos}, pela \emph{extensão} do \emph{objecto}, e pela \emph{excellencia} do \emph{fim}, que é a visão intuitiva de Deus.}.

\pg{420} 4. \textbf{Utilidade da Theodicea}. --- O estudo da \emph{Theodicea} é summamente util para o \emph{individuo} e para a \emph{sociedade}.

a) B? util para o \emph{individuo}, Por quanto, a nossa \emph{intelligencia}, conhecendo \textsc{Deus}, descobre a Verdade infinita, de que emanam todas as verdades, e a Causa primeira, que produziu todas as creaturas; e a nossa vontade, illuminada e guiada por esse conhecimento, une-se a Dxzvs por amor, e nesta união encontra o deseanco nas suas tendencias, a paz nas luctas, á força nas diffieuldades, o allivio nas dores e miserias da vida. j

b) B' util para a \emph{sociedade}. Na verdade, sendo a sociedade a collecção dos individuos, é claro que, se estes conhecerem e amarem a \textsc{Deus}, aquella terá tranquillidade e socego, e, fortalecida pela harmonia e honestidade dos cidadãos, progredirá num caminho de civilisação e de gloria não caduea \footnote{\textsc{DEus} é o principio e o fim de todas as coisas. Logo ninguem póde prescindir d'Elle. --- Se prescindirmos de \textsc{Deus} na vida \emph{particular}, encontraremos o erro, o vicio, e, por consequencia, o remorso. --- Se prescindirmos de \textsc{Deus} na vida \emph{social}, teremos \emph{o} despotismo, o roubo, a anarchia. --- Se prescindirmos de \textsc{Deus} nas sciencias, os nossos conhecimentos serão \emph{imperfeitos}, porque não baseados na razão suprema de todas as coisas. Com razão \emph{Bacon} dizia: « Quem considerar a dependencia, a serie e o encadeamento das causas, facilmente se convencerá de que, segundo a mythologia dos poetas, o primeiro annel da cadóa natural está preso aos pés do throno de Jupiter » (\emph{De dign}. \emph{et augm}. \emph{scient}., l. I, n. 5).\par Por isso a \emph{Theodicea} tem um caracter essencialmente real e scientifico, porque encerra a explicação de todos os problemas, a resolução de todas as ineognitas, o primeiro e mais forte elo da grande cadóa dos seres, --- Aos \emph{materialistas} e \emph{positivistas}, que excluem do quadro das sciencias a \emph{Theodicea}, offerecemos um magnifico trecho de um escriptor muito conhecido d'elles --- \emph{Vacherot}. E o seguinte: « O que é esta \emph{Theologia} aristotelica, senão o remate\pgn{421}de uma \emph{psychologia} e de uma \emph{physica}, fundadas na experiencia e na analyse da realidade? Onde está o methodo \emph{a-priori}? Onde está a mera especulação neste vasto systema, cuja base é tão solida como elevado é o seu vertice? Se alli a natureza é explicada pela alma, e a alma por \textsc{Deus}, quem levou o espirito a este supremo principio de explicação, senão a physica primeiramente, e depois á psychologia? Isto manifesta a razão por que a doutrina de Aristoteles ficou firme sobre as ruinas de tantas especulações antigas e modernas, como um d'esses monumentos cuja base solida e profunda desafia a acção do tempo e os golpes da tempestade... \textsc{Deus} é a maior palavra das linguas humanas. Nenhuma poude esquecel-a. Todas, ainda as mais barbaras, celebraram-na com maior ou menor precisão, profundidade e pureza. E' o problema por excellencia das mais sublimes philosophias » (\emph{Le Nouveau Spiritualisme}, pags. 164 e 288).}.

\pg{421} 5. \textbf{Divisão} da \textbf{Theodicea}. --- Dividimos a \emph{Theodicea} em \emph{quatro} capitulos. No \emph{primeiro} capitulo tratamos da \emph{existencia} de \textsc{Deus}; --- no \emph{segundo}, da sua \emph{essencia} e dos seus \emph{aítributos}; --- no \emph{terceiro}, das suas \emph{faculdades e operações àmmanentes} ; --- no \emph{quarto}, das suas \emph{operações transeuntes} \footnote{Ninguem se admire se o nosso tratado da \emph{Theodicea} é relativamente breve. Eserevemos sobretudo para estudantes, que se destinam ao estado ecclesiastico e que devem applicarse, mais tarde, a um estudo completo da S. Theologia. Por isso, para ganhar tempo e evitar repetições, julgamos conveniente apresentar em resumo as principaes verdades, que se referem a Deus e se conhecem pela luz natural da razão, e depois tratar, com o devido desenvolvimento, as outras partes da Philosophia racional.}.

%%>>> cap-6

\chapter*{CAPITULO PRIMEIRO\\[0.3em]\large EXISTENCIA DE DEUS}
\addcontentsline{toc}{chapter}{CAPITULO PRIMEIRO --- EXISTENCIA DE DEUS}

{\small\noindent\textsc{Summario}: --- Demonstrabilidade da existencia de \textsc{Deus}. --- Demonstração da existencia de \textsc{Deus}. --- A existencia de \textsc{Deus} e o atheismo.\par}

%%>>> art-36

\section*{Artigo I.\\ Demonstrabilidade da existencia de Deus.}
\addcontentsline{toc}{section}{Artigo I. --- Demonstrabilidade da existencia de Deus.}

6. A \textbf{existencia de Deus} póde ser \textbf{demonstrada pela razão humana}. --- A razão humana póde deduzir da existencia dos effeitos a existencia da causa. Ora todos os seres sensiveis são effeitos. Logo a razão póde deduzir da existencia dos seres sensiveis a existencia da sua causa, que é \textsc{Deus} \footnote{Os systemas, que contestam a nossa proposição sao --- o (\emph{r«dicionalismo}, admittindo que a existencia de \textsc{Deus} só pode ser conhecida pela\pgn{422}\emph{revelação} ou pela \emph{tradição}, --- o \emph{Kantismo}, sustentando que percebemos os \emph{phenomenos} ou as \emph{manifestações}, mas não os \emph{noumenos} ou as \emph{substancias}, --- o \emph{positivismo} e o \emph{sensualismo}, affirmando que a razão humana só attinge objectos sensiveis, --- o \emph{agnosticismo}, ensinando que toda a nossa sciencia se resume e concentra na esphera dos phenomenos e que por isso o \emph{ultraphenomenico} é \emph{incognoscivel}, --- o \emph{immanentismo}, negando á razão a possibilidade de demonstrar a existencia de Deus e attribuindo-a á \emph{consciencia}.\par Por se sufficientemente conhecida a natureza dos primeiros systemas apontados, falta-nos agora explicar a natureza do ultimo, que é o \emph{ómmanentismo}. Segundo as pretenções d'este systema, a razão não pode demonstrar a existencia de \textsc{Deus}, como não pode demonstrar nenhuma verdade religiosa e moral, porque esta demonstração não pode ser verificada pelos dados da experiencia. As verdades religiosas e moraes não se provam com os raciocinios, como se provam as verdades mathematicas; devem brotar do fundo da nossa alma, do intimo da nossa consciencia. Se, pois, quizermos encontral-as, devemos reflectir sobre os nossos actos, investigar as exigencias do nosso espirito, sondar as necessidades do nosso coração, escutar a voz da consciencia, onde residem os dictames religiosos e moraes. Para esse trabalho é indispensavel uma vontade bóa e recta. Por quanto, os sentimentos religiosos e moraes, embora se encontrem em todos os homens, mesmo nos incredulos, comtudo só nos homens de bóa vontade e de vida irreprehensivel se encontram num estado \emph{explicito}, \emph{claro} e \emph{consciente}, --- Esta é a chamada \emph{philosophia nova}, ou philosophia da \emph{immanencia}, da \emph{acção}, do \emph{voluntarismo}, do \emph{pragmatismo}, que tem por fundadores \emph{Renouvier}, \emph{Boutrouz}, \emph{Bergson}, \emph{Mach}, \emph{William James}. E' uma emanação da doutrina de \emph{Kent}, que, depois de ter negado todo o valor á razão \emph{especulativa}, quando esta investiga o \emph{sobresensivel}, ensina que a razão \emph{pratica} deve admittir a existencia de mesmo \emph{sobresensivel}, como postulado ou presupposto necessario para a ordem moral --- Esta theoria tem sido perfilhada pelos \emph{modernistas}. Dizem elles que a existencia de Deus não pode ser demonstrada por meio dos factos \emph{externos}, isto é, pela existencia e natureza dos seres \emph{creados}, mas sô por meio dos factos \emph{internos} e \emph{subjeclivos}. A razão não pode demonstrar a existencia de Deus, ou não o pode fazer sem a cooperação da vontade, do sentimento e dos outros factores \emph{subjectivos}, que \emph{supprem a deficencia dos argumentos racionaes}.\par Tendo refutado na \emph{Anthropologia} o \emph{tradicionalismo}, o \emph{Kantismo}, o \emph{positipismo} e o \emph{sensualismo}, resta-nos refutar o \emph{agnosticismo} e o \emph{immanentismo}.\par \emph{A} pretenção do \emph{agnosticismo}, que ensina ser \emph{incognoscivel} para a razão a existencia de Deus, é uma pretenção absurda. Que as sciencias \emph{eoperimentaes} não reconheçam senão os \emph{p/enomenos}, e que para ellas o \emph{sobr'esensivel} seja \emph{incognoscivel}, concedemol-o de bom grado. Mas seguir-se-ha d'ahi que, alem do sensivel, não existe mais nada, e que, atraz do phenomeno, não existe o nowmeno? Não, com certeza. Se o sobresensivel não é objecto das sciencias experimentaes, é objecto de outra sciencia mais elevada. Os udversarios imitam nisso o celebre D. Ferrante, nos \emph{promessos esposos} de Manzoni, que negava a existencia da peste em Milão, só porque esta não cabia em nenhuma das suas categorias! Se a existencia de Deus não é perceptivel pelos\pgn{423}sentidos, é perceptivel pela razão. E a razão não prescinde dos dados da experiencia sensivel, mas funda-se nelles, explica-os e completa-os. Por quanto, no fundo de cada facto sensivel está sempre uma questão \emph{metaphysica}, que diz respeito á \emph{origem} e á \emph{natureza} do proprio facto e que só a razão resolve. --- A Metaphysica, repetimos, não é um tecido de syllogismos forjados \emph{a-priori}; 6 \& investigação das causas \emph{ultimas}, de que dependem os factos. E, como essas causas devem ser e são reaes, porque reaes são 08 factos, tambem real deve ser e é a Metaphysica, que tem por seu objecto as causas \emph{ultimas} dos entes. --- Coneluimos que \& existencia de Deus, se é \emph{incognoscivel} pelos sentidos, não é \emph{incognoscivel} pela razão.\par Não menos absurda é a pretengüo dos \emph{émmanentistas}. Dizem estes --- que Deus não se conhece pela razão, mas pela \emph{consciencia}, --- que é na consciencia que se manifesta o divino, --- que é no intimo da alma que se encontra Deus. --- Mas, como a consciencia é \emph{psychologica} ou \emph{moral}, conforme refere a existencia das nossas operações intellectuaes, ou indica a \emph{conveniencia} d'essas operações com as regras dos costumes, d'ahi a necessidade de perguntarmos aos \emph{immanentistas} se a consciencia, de que fallam, é a \emph{psychologica}, ou é a \emph{moral}? --- Be entendem fallar da consciencia \emph{psychologica}, enganam-se redondamente. Esta consciencia revela-nos as nossas operações, as nossas necessidades e desejos, a existencia do eu, a realidade nossa, mas não a realidade de Deus. Diz-nos, é certo, á consciencia que temos a aspiração, o desejo do infinito, de Deus; mas á \emph{aspiração} do infinito não é a \emph{presença} do infinito. Se o infinito se sentisse na nossa consciencia, não haveria aspiração ou desejo; o que se possue, não se deseja. O infinito está fóra e sobre nós; o divino não é \emph{immanente}, é \emph{transcendente}. Fundando-se nestas aspirações ou desejos, a razão, pelo principio de \emph{causalidade}, eleva-se a Deus... Admittimos que Deus está em nós; mas está como causa creadora o conservadora do nosso ser; e esta presença de Deus, como causa do nosso ser, não se percebe pela consciencia, mas pelo raciocinio, emquanto os seres finitos suppõem e exigem a existencia de um ser infinito... Se teimarmos em procurar a Deus \emph{unicamente} com os dados da consciencia, nunca o encontraremos, a não ser que se admitta --- ou que o ser divino se identifica com o nosso ser, como dizem os \emph{pantheistas}, --- ou que o ser divino, embora distincto do nosso, está sempre presente á nossa intelligencia, como objecto de percepção immediata e directa, como affirmam os \emph{ontologíistas}, --- ou que o ser divino é representado e unido á nossa intelligencia por uma idéa ou especie \emph{innata}, como querem os \emph{énnatistas}. --- Não menos grosseiro é o engano dos \emph{immanentistas}, se ou quando dizem --- que é \& consciencia \emph{moral} que revela o \emph{divino}, --- que Deus vive em nós, como suprema lei moral, --- que a consciencia \emph{moral} é o echo da voz de Deus, é o \emph{proprio Deus que falla}. Por quanto, se sentimos que os dictames da lei moral nos \emph{obrigam}, sentimos tambem que essa \emph{obrigação} não deriva da nossa vontade, mas nos é imposta por um superior, que está sobre e fóra de nós; pois ninguem pode ser subdito e superior de si mesmo. Logo da existencia da lei moral podemos o devemos deduzir a existencia de um legislador supremo... Alem d'isso, \& consciencia, sendo a regra \emph{proxima} dos nossos actos, é o \emph{echo} da voz de\pgn{424}Deus, que é a regra \emph{ultima}. E' o echo da voz de Deus; mas não é o proprio Deus; porque o echo de uma voz distingue-se da propria voz, e da pessoa que a emitte. O echo da lei moral só pode servir de base para a conclusão da existencia de um supremo legislador. --- Não podemos tambem negar que a vontade tem o seu logar na demonstração da existencia de Deus. Mas é preciso determinar com exactidão o papel da mesma faculdade nesta demonstração. A vontade pode entrar --- ou para levar a intelligencia a \emph{considerar} os argumentos, que provam a existencia de Deus, --- ou para levar a intelligencia a \emph{admittir cegamente} a existencia de Deus. Na primeira hypothese, o influxo da vontade na intelligencia é apenas \emph{extrinseco} e \emph{mediato}, e isto verifica-se em todos os conhecimentos; na segunda, o influxo é \emph{intrinseco} e \emph{immediato}, e isto verifica-se no acto de Fé, que não é determinado pela \emph{evidencia} da verdade proposta, mas pelo imperio da vontade. Ora é claro que a vontade entra na demonstração da existencia de Deus --- não emquanto leva a intelligencia a \emph{admittir cegamente} essa existencia, mas emquanto a determina a considerar os argumentos, que \emph{prova} tal verdade; de modo que o assenso da intelligencia é devido á força dos argumentos, e não ao imperio da vontade. Esta nunca poderá fazer com que seja verdadeira uma proposição, que a intelligencia julgou falsa. Dizer, pois, com os \emph{émmanentistas}, --- que o sentimento deve substituir a razão, --- que a emoção é um conhecimento, --- que a disposição subjectiva deve ter e tem o logar da verdade objectiva, é blasonar um erro tão absurdo quão pueril.\par A. Fé confirma as conclusões da razão. O Concilio Vaticano emanou a seguinte definição: « Si quis dixerit Drum unum et verum, Creatorem et Dominum nostrum, \emph{per ea quae facta sunt}, \emph{naturali rationis lumine certo cognosci non posse}, anathema sit». --- O Concilio baseou a sua definição no Livro da \emph{Sabedoria} (\emph{c}. 19, vv. 1-5) e na Epistola de S. Paulo aos Romanos (c. 1, vv. 20-21). --- Estas auectoridades devem ser recebidas, ao menos pelos \emph{tradicionalistas}, que reconhecem a \emph{revelação} .como unico criterio de verdade.}. -

7. A \textbf{existencia de Deus deve ser demonstrada}. --- A existencia de um ente deve ser demonstrada, quando ella não é percebida \emph{directa} e \emph{immediatamente} --- nem pelas faculdades \emph{sensitivas}, --- \pg{422} \pg{423}  nem pela \emph{intelligencia}. Ora a existencia de \textsc{Deus} não é percebida \emph{directa} e \emph{immediatamente} --- nem pelas faculdades \emph{sensitivas}, porque \textsc{Deus} não é um ente sensivel, --- nem pela \emph{intelligencia}, \pg{424}  pois, como demonstrámos na \emph{Anthropologia}, a visão \emph{intuitiva} de \textsc{Deus} excede a natural capacidade da nossa intelligencia. Logo á existencia de \textsc{Deus} deve ser demonstrada \footnote{Os que negam a necessidade de demonstrar a existencia de \textsc{Deus} são --- a) os \emph{sentimentalistas}, sustentando que conhecemos \emph{immediatamente} a \textsc{DEus} por meio de um sentido arcano e divino. --- 5) os \emph{ontologistas}, admittindo uma \emph{zntuicüo immediata} ou \emph{directa} de \textsc{Deus}.\par Estas duas opiniões são falsas. -\par \emph{a}) E' falsa a dos \emph{sentimentalistas}; pois a existencia d'esse sentido arcano e divino é contraria ao testemunho da consciencia.\par 5) E' falsa a dos \emph{ontologistas}; pois o \emph{ontologismo}, como vimos, € falso.\par Aqui apenas accrescentaremos que a S. Egreja, Mestra infallivel de verdade, declarou suspeitas e perigosas, entre outras, estas proposições: « Immediata Dk: cognitio, habitualis saltem, intellectui humano essentialis est, ita ut sine ea nihil cognoscere possit: siquidem est ipsum lumen intel lectuale ». --- « Esse illud, quod in omnibus et sine quo nihil intelligimus, est esse divinum ». (Decr. S. Officii, die 18 Septemb. a. 1861).}.

8. \emph{A} \textbf{existencia} de \textbf{Deus} não se demonstra \textbf{a-priori}. --- A demonstração diz-se \emph{a-priori}, quando a existencia de um ser se prova \pg{425}  pela sua \emph{causa efficiente}. Ora a existencia de \textsc{Deus} não se prova pela sua \emph{causa efficiente}; porque \textsc{Deus} é a Causa primeira, e por isso absolutamente independente de toda e qualquer causa. Logo a existencia de \textsc{Deus} não se demonstra \emph{a-priori}.

9. A \textbf{existencia} de \textbf{Deus} não \textbf{se demonstra a-simultaneo}. --- A demonstração é \emph{a-simultaneo}, quando a existencia de um ente se prova pela analyse da idea ou noção do proprio ente. Ora este processo, applicado á existencia de \textsc{Deus}, é vieioso; visto que, para uma coisa existir na realidade, não basta que exista no pensamento, pois deduzir da ordem \emph{dia}? a ordem \emph{real} é tirar uma consequencia mais extensa que as premissas. Logo a existencia de \textsc{Deus} não se demonstra \emph{a-simultaneo} \footnote{Esta proposição --- \textsc{DEus} \emph{existe} --- é \emph{immediatamente} evidente \emph{em si} mesma, mas não é \emph{immediatamente} evidente \emph{em relação a} nós.\par a) E' \emph{immediatamente} evidente \emph{em si mesma}. Uma proposição é \emph{immediatamente} evidente \emph{em si mesma}, quando o predicado convém \emph{essencialmente} ao sujeito. Ora a \emph{existencia} convém \emph{essencialmente} a \textsc{DEus}; porque \textsc{Deus} é o ser \emph{necessario}, \emph{e} não póde deixar de existir. Logo a proposição --- \textsc{Deus} \emph{existe} --- é \emph{immediatamente} evidente em si mesma.\par \emph{b}) Não é \emph{immediatamente} evidente \emph{em relação a nós}. Uma proposição é \emph{immediatamente} evidente em relação a nós, quando directamente e sem necessidade de raciocinio conhecemos a essencia do sujeito e do predicado; assim a proposição --- \emph{o todo é maior que} a \emph{parte} --- é \emph{immediatamente} evidente \emph{em relação} a \emph{nós}, pois basta enunciar as palavras \emph{todo} e \emph{parte}, para logo percebermos que o /\emph{odo} é maior que a \emph{parte}. Ora a essencia de \textsc{Deus}, não sendo por nós conhecida senão depois de um raciocinio, e ainda assim muito imperfeitamente, não se nos manifesta logo á primeira vista como exigindo a existencia. Logo a proposição --- Dus \emph{existe} --- não é \emph{immediatamente} evidente \emph{em relação a nós}. --- Cf. S. Thom., --- \emph{Sum}. \emph{Th}., p. I, q. 2,a.1; --- c. \emph{Gentes}., 1. I, cc. 10, 11; 1 \emph{Sent}. Dist. 8, q. 1, a. 2; \emph{De verit}., q. 10, a. 12.\par Os mais celebres argumentos \emph{a-simultaneo} em favor da \emph{ecistencia} de \textsc{DEUS} são os de \emph{Santo Anselmo}, \emph{Descartes} e \emph{Leibnits}.\par a) \emph{Santo} Anselmo raciocina d'este modo: Uma proposição é \emph{immediatamente} evidente \emph{em relação a nós}, quando a sua verdade se nos manifesta pela simples enunciação dos termos. Ora a verdade d'esta proposição --- \textsc{Deus} \emph{existe} --- torna-se-nos manifesta pela simples enunciação dos termos. Com effeito, quando se profere o nome de \textsc{Deus}, immediatamente entendemos que \textsc{Deus} é o Ser maior que se póde imaginar, e que por isso deve existir na \emph{realidade}; porque, se existisse unicamente na intelligencia, não seria o Ser maior que se póde imaginar, pois superior a Elle seria o ser, que existisse na intelligencia e tambem na \emph{realidade}.\par Critica. Ainda que todos admittissem que \textsc{DEus} é o ser maior que se póde imaginar, todavia d'ahi só poderia concluir-se que \textsc{Deus} existe na \emph{intelligencia}, mas não na \emph{realidade}; porque, como dissemos, deduzir da ordem \emph{ideal} a ordem \emph{real} é tirar uma consequencia, que não está incluida nas premissas\pgn{426}. --- Para provar que \textsc{Deus} existe tambem na \emph{realidade}, deve demonstrar-se que existe na \emph{realidade} um ser, que é o maior que se póde imaginar. --- Cf. S. Thom., c. \emph{Gentes}, 1. I, c. 11.\par 1) \emph{Descartes} diz: Todo o attribuito, que está contido na idea \emph{clara} e \emph{distincta} de um sujeito, deve affirmar-se do mesmo sujeito. Ora na idea \emph{clara} e \emph{distincta} de \textsc{Deus} está contida a existencia \emph{real}; porque \textsc{Deus} é infinitamente perfeito, e a existencia \emph{real} é uma perfeição. Logo deve affirmar-se de \textsc{Deus} a existencia \emph{real}, isto é, \textsc{DEus} existe \emph{realmente}.\par Critica. O argumento de \emph{Descartes} tem o mesmo defeito que o de \emph{Santo} Anselmo; pois passa da ordem \emph{ideal} para a \emph{real}. Na verdade, tudo o que esta contido na idea \emph{clara} e \emph{distincta} de um sujeito deve attribuir-se ao sujeito na mesma \emph{ordem}, em que está o sujeito; de modo que, se o sujeito estiver na ordem \emph{real}, o predicado tambem estará na ordem \emph{real}, --- mas, se o sujeito estiver na ordem \emph{ideal}, o predicado tambem existirá necessariamente na ordem \emph{ideal}. Logo na \emph{idea} de \textsc{DEus} está contida a existencia \emph{ideal}, e não a \emph{real}. --- Lembramos que, para \emph{Descartes}, a idea do \emph{infinito é innata} (o que é falso), e que a verdade da proposição maior do seu syllogismo se baséa na veracidade do proprio \textsc{DEus} (o que constitue um \emph{circulo vicioso}).\par 0) \emph{Leibnitz} apresenta este argumento: O ente infinitamente perfeito é possivel. Se é possivel, existe na \emph{realidade}; porque á existencia \emph{real}, por ser uma perfeição, pertence necessariamente á essencia do Ente infinitamente perfeito. Logo o Ente infinitamente perfeito existe na \emph{realidade}.\par Critica. Este argumento é inconcludente pelo mesmo defeito, que apontámos nos dois antecedentes. Do facto de \emph{concebermos}, como \emph{possivel}, o Ente infinitamente perfeito, cuja essencia exige a existencia, só resulta que nao podemos \emph{conceber}, como \emph{possivel}, esse ente, sem que o \emph{concebamos} como existente; mas não se segue que a realidade deva corresponder ao que \emph{concebemos} na intelligencia.}.

10. A \textbf{existencia} de \textbf{Deus demonstra-se a-posteriori}. --- A demonstração \emph{a-posteriori} faz-se, quando, pela existencia do effeito, \pg{426}  se prova a existencia da causa respectiva. Ora a existencia de \textsc{Deus} demonstra-se pela existencia das creaturas; porque estas, por serem \emph{effeitos}, suppoóem uma causa proporcionada, que é \textsc{Deus}. Logo a existencia de \textsc{Deus} demonstra-se \emph{a-posteriori} \footnote{Cf. S. Thom., --- \emph{Sum}. \emph{Th}., p. I, q. 2, a. 2; --- \emph{c}. \emph{Gentes}, l. I, c. 12.\par A demonstração da existencia de Deus basea-se no \emph{principio de causalidade}. E' por isso que os \emph{materialistas}, os \emph{positivistas} e todos os \emph{incredulos}, no intuito de negar esta verdade fundamental, ou, ao menos, de prescindir d'ella, atacam a noção de causa e esse principio de \emph{causalidade}. --- Não repetimos tudo o que na \emph{Ontologia} escrevemos em defeza do dito principio. Limitamo-nos a observar que o principio de \emph{causalidade} é uma verdade, que a intelligencia descobre e formula á vista das mudanças, dos phenomenos, dos factos, que se succedem, uns aos outros, no mundo visivel. Ora quando um phenomeno começa, percebe-se immediatamente que está relacionado com um outro, de que deriva. Esta relação não é apenas de \emph{successdo}, mas de \emph{dependencia}. Tal dependencia é uma coiea real, objectiva, existente nos proprios seres creados; a intelligencia \emph{descobre-a}, mas \emph{não a inventa}. Por isso o principio de \emph{causalidade} tem un valor real e objectivo; porque, embora este juizo --- \emph{todo o effeito tem a sua causa} --- seja um producto da razão, comtudo a rasdo \emph{ndo o produz senão fundada nos dados da experiencia sensivel}. A experiencia apresenta os factos; a razão descobre o nexo, que os liga. ---\pgn{427}Advirta-se que o principio de \emph{causalidade} não se applica a todo o ser que \emph{existe}, mas a todo o ser que começa a existir, e por isso a Causa \emph{primeira} não tem causa.\par Os \emph{Kantistas} não negam o principio de \emph{causalidade}, mas dizem que elle só tem valor, quando se trata de dois phenomenos, ambos sujeitos á observação sensivel, porque, em tal caso, podemos averiguar a dependencia de um phenomeno do outro, --- mas não tem valor, quando se trata de duas realidades, uma das quaes não está sujeita á experiencia, porque, em tal caso, não podemos conhecer a \emph{dependencia} ou a \emph{origem} de uma realidade da outra, e por isso nem a razão de \emph{causa} e de \emph{ejfeíto}; o que se verifica no nosso caso, porque, se o mundo está sujeito á experiencia, Deus escapa a toda a experiencia. Por isso não podemos basear-nos no principio de \emph{causalidade}, para deduzir da existencia do mundo a existencia de Deus.\par Respondemos que, para o principio de \emph{causalidade} ter todo o seu valor, não é necessario que ambas as realidades, que se dizem proceder uma da outra, estejam sujeitas á experiencia sensivel. Basta que um phenomeno, um facto, presente aos nossos sentidos, manifeste um caracter de mutabilidade e de contingencia, para que possamos e devamos concluir que esse phenomeno \emph{depende} necessariamente de um antecedente. Pouco importa, para O caso, que esse antecedente seja visivel ou invisivel, natural ou sobrenatural; elle existe, deve existir, porque existe o phenomeno, que não tem em si a razão sufficiente da sua existencia. Ninguem viu a força da \emph{gravitação}; e, comtudo, todos a admittem, porque se revela na queda dos graves. Ninguem viuo talento militar de Napoleão; mas todos o admittem, porque se manifestou nos seus grandes planos de guerra. Deus é invisivel na sua essencia, mas revela-se nas coisas creadas, que suppõem e exigem uma força infinita. Portanto, ou devemos negar o principio de \emph{causalidade}, e temos a completa derrocada de todas as sciencias, --- ou devemos admittir que, em força d'esse principio, podemos e devemos elevar-nos ao conhecimento da existencia de Deus. --- Nem se diga que uma percepção mais profunda e completa dos phenomenos naturaes pode fazer com que se prescinda de uma causa primeira. Os phenomenos naturaes, ainda que podessem ser percebidos em toda a sua comprehensão, nunca podem perder os \emph{essenciaes} caracteres de \emph{contingencia} e de \emph{mutabilidade}, e por isso da sua \emph{dependencia} de uma causa primeira. Pelo contrario: quanto mais se conhece o \emph{movimento}, a \emph{larmonia}, a \emph{finalidade} dos seres visiveis, tanto mais necessaria apparece a «xistencia de um supremo e omnipotente Ordenador.\par Objectam: Todos os nossos conhecimentos são \emph{relativos} ao mundo \emph{visive}? o á intelligencia \emph{finita} do homem. Logo não podemos conhecer a \textsc{Deus}, que se suppõe \emph{invisivel} e \emph{infinito}.\par Resposta. E' certo que todos os nossos conhecimentos, \emph{que adquirimos}\pgn{428}\emph{de um modo directo e immediato}, se referem ao mundo visivel. Mas d'ahi não se segue que a nossa razão não possa, nem de um modo \emph{indirecto} e \emph{mediato}, adquirir outros conhecimentos, relativos ao mundo invisivel. Os caracteres de contingencia e de mutabilidade, que observamos no mundo visivel, levam-nos necessariamente a admittir a existencia de um Ente invisivel, necessario e immutavel. --- Como tambem é certo que todos os conhecimentos são proporcionados á capacidade da intelligencia cognoscente e ao modo por que os objectos se tornam presentes á mesma intelligencia, de modo que só apprehendemos as notas do objecto, que actuaram na faculdade, ficando incognosciveis as que escaparam ao alcance d'ella. Mas esta \emph{relatividade}, se prova que não podemos \emph{comprehender} \& Deus, porque, a intelligibilidade de Deus excede infinitamente a capacidade da intelligencia creada, não prova que não possamos, \emph{de algum modo}, conhecer o omnipotente Creador de todas as coisas.\par .Devemos, porém, advertir que, para os sequazes do \emph{agnosticismo}, \emph{relatividade} significa \emph{subjectivismo}; e por isso quando dizem que os nossos conhecimentos são \emph{relativos}, entendem que são \emph{subjeclivos}, em quanto nos fazem conhecer as coisas não como são em si mesmas, mas como são no nosso espirito, ás quaes nada corresponde na realidade. Não nos detemos na refutação d'este erro, a qual já está feita. Só acerescentaremos que, se todos os nossos conhecimentos sao \emph{swbjectivos}, são \emph{subjectivos} não só os que se referem ao mundo invisivel, mas tambem 03s que se referem ao visivel, e por isso, se não existe Deus, nem existe coisa alguma no mundo. --- Nem se diga que os conhecimentos relativos ao mundo \emph{visivel} são \emph{objectivos}, porque podem ser verificados pela experiencia, pelos factos: o que se não pode fazer com os conhecimentos relativos ao mundo \emph{invisivel}, a Deus. Negamos que os nossos conhecimentos relativos a Deus não se possam verificar. pelos factos. Na verdade, como provamos nós que Deus existe? Provamol-o pela origem e natureza do mundo, do movimento, da vida, da ordem, da finalidade; porque a origem e a natureza d'estas coisas suppõem a existencia de uma causa extramundana. Ora o mundo, o movimento, a vida, a ordem, a finalidade, são factos, e factos grandiosos, admiraveis, que actuam nos nossos sentidos e despertam a admiração na nossa intelligencia. Portanto o conhecimento . de Deus é tão \emph{objectivo}, como \emph{objectivos} são os factos, que nos servem de degrau para nos elevarmos a Elle.\par Continuam: A existencia de Deus não pode ser demonstrada pelos effeitos. Na verdade, entre a causa e o effeito deve existir uma proporção; ora nenhuma proporção pode existir entre as creaturas, que são seres finitos, e Deus, que é o Ser infinito.\par Resposta. Demonstra-se á causa pelos effeitos por dois modos. Se os effeitos possuem a mesma natureza da causa (0 que se verifica nas causas \emph{univocas}) então não só levam ao conhecimento da \emph{existencia}, mas tambem da \emph{natureza} da cansa, de modo que adquirimos um conhecimento \emph{perfeito} da causa. Se os effeitos não têm a mesma natureza da causa (\emph{o} que se verifica nas causas \emph{equivocas}), então só leyam ao conhecimento da \emph{ewístencia} da causa, e não da \emph{natureza}. e por isso o conhecimento da causa é \emph{imperfeito}\pgn{429}, e o grau d'esta imperfeição é proporcionado ao grau de dessemelhanca entre os effeitos e a causa. Por isso, \emph{todo} o effeito tem com a sua causa uma proporção de \emph{ordem}, emquanto é produzido por ella; mas nem \emph{todo} o effeito tem com a causa uma proporção de \emph{natureza}. Appliquemos estes principios ao nosso caso. Os seres visiveis, sendo finitos e contingentes, não exprimem toda a virtude da sua causa, não possuem á mesma \emph{natureza} de Deus, e por isso não nos levam ao conhecimento \emph{adequado} da \emph{natureza} de Deus. Mas, sendo \emph{effeitos}, e effeitos que só podem derivar de um Ente necessario e omnipotente, levam certamente ao conhecimento da \emph{existencia} d'esse Ente. --- Todavia, embora o conhecimento, que temos da \emph{natureza} de \textsc{Deus}, não seja \emph{adequado}, mas \emph{imperfeito}, comtudo é \emph{positivo}; pois as coisas creadas são semelhanças \emph{positivas}, embora \emph{imperfeitas}, das divinas perfeições, visto que todo o agente tende a imprimir no effeito a propria \emph{semelhança}.}.

%%>>> art-37

\section*{Artigo II.\\ Demonstração da existencia de Deus.}
\addcontentsline{toc}{section}{Artigo II. --- Demonstração da existencia de Deus.}

11. \textbf{Especies de argumentos para} a \textbf{demonstração da existencia de Deus}. --- A existencia de \textsc{Deus} demonstra-se \emph{a-posteriori}. Ora os argumentos, que demonstram \emph{a-posteriori} a existencia de \pg{427} \pg{428}  \textsc{Deus}, são de tres especies: \emph{metaphysicos}, \emph{physicos} e \emph{moraes}. Os \emph{metaphysicos} basé6am-se na propria \emph{natureza} das coisas, --- os \emph{physicos}, na \emph{ordem} que existe no mundo e, em particular, nas \pg{429}  \emph{operações} dos agentes \emph{Zrraciomaes}, --- os \emph{moraes}, na existencia da lei \emph{moral} e no \emph{consenso} de todos os povos \footnote{Na classificação dos argumentos, que demonstram a existencia de \textsc{Deus}, em \emph{methaphysicos}, \emph{physicos}, \emph{e moraes}, conforme se baséam na \emph{essencia} das coisas, ou na \emph{ordem} existente no mundo e nas \emph{operações} dos agentes \emph{irracionaes}, ou na lei \emph{moral} e no \emph{consenso} dos povos, seguimos a doutrina dos mais sensatos auctores.\par Alguns dividem esses argumentos em \emph{cosmologicos}, \emph{psychologicos} e \emph{moraes}; conforme os factos, em que se fundam, são \emph{cosmologicos} (baseados nas leis \emph{physicas}), \emph{psychologicos} (baseados nas leis \emph{intellectuaes}), ou \emph{moraes} (baseados nas leis \emph{moraes}). --- Esta divisão tambem é acceitavel.\par \emph{Kant}, dizendo que o espirito humano se eleva a \textsc{DEus} por meio de um d'estes tres principios: de \emph{identidade}, de \emph{causalidade} ou de \emph{finalidade}, reduz as provas a tres classes: \emph{ontologicas}, \emph{cosmologicas} e \emph{physico-teleologicas}, conforme se baseam no principio de \emph{identidade}, de \emph{causalidade} ou de \emph{finalidade}. --- A. doutrina de \emph{Kant} é inacceitavel. O principio, em que se baseam todas as provas em favor da existencia de \textsc{Deus}, é um só: o \emph{principio de causalidade}, e a distineção das provas deriva da natureza dos \emph{factos}, aos quaes o dito principio se applica. O principio de \emph{identidade} não prova coisa alguma, como o proprio \emph{Kant} confessa; o principio de \emph{finalidade}, longe de ser a base do argumento \emph{fteleologíco}, é apenas um corollario d'esse argumento.\par A. existencia de \textsc{DEus} é uma verdade tão necessaria e fundamental, que recorre em todos os grandes problemas da philosophia. Nem admira. Nào encontrando nos agentes finitos a causa \emph{sufficiente} dos seres, a razão deve elevar-se a uma causa \emph{suprema} e \emph{infinita}. \emph{A origem} do \emph{mundo} (\emph{Cosm}., \emph{ger}.\pgn{430}Cap. I, art. V), á \emph{finalidade} nas operações dos seres irracionaes (\emph{Cosm}. \emph{ger}., Cap. II, art. III a ordem \emph{physica} com as suas \emph{leis} (\emph{Cosm}. \emph{ger}., Cap. II, art. IV e V), a \emph{origem} e o \emph{desenvolvimento da vida} (\emph{Cosm}. \emph{esp}., Cap. II, art. III), á \emph{origem} da \emph{alma} e do \emph{corpo} humano (\emph{Anthr}., Sec. IV, Cap. D, tudo o que nos cerca, com os seus caracteres de \emph{contingencia} e de \emph{mutabilidade}, proclama a existencia de uma Causa infinitamente sabia, poderosa e boa --- a existencia de \textsc{DEus}.\par Escutámos esta voz maravilhosa da creação e podemos dizer, com immensa satisfação da nossa alma, que procurámos sempre pôr em relevo a racionabilidade e a necessidade da existencia de um Ente supremo, e mostrar a fatuidade, os erros e as contradieções dos adversarios.\par Mas isso não obsta a que façamos aqui, como no logar proprio, uma synthese do que dissemos em todo este trabalho, indieando os principaes argumentos em defeza da \emph{existencia} de Deus. --- Cf. \emph{Sum}. \emph{Th}., p. I, q. 2, a. 8; --- \emph{c}. \emph{Gentes}, l. I, c. 18.}.

. 12. Argumentos \textbf{metaphysicos}. --- Os argumentos \emph{metaphysicos}, em favor da \emph{existencia} de \textsc{Deus}, baséam-se em que, --- a) se no mundo existe o movimento ou mudança, deve existir um motor \emph{emmovel}, --- \emph{b}) se existem as causas \emph{segundas}, deve existir a causa \emph{primeira}, --- c) se existem seres \emph{contingentes}, deve existir um ser \emph{necessario}, --- \emph{d}) se existem varios e diversos graus de perfeição, deve existir um ser infinitamente perfeito, --- e) se existem verdades \emph{absolutas} e \emph{necessarias}, deve existir uma intelligencia \emph{absoluta} e \emph{necessaria}.

\pg{430} a) As substancias creadas estão sujeitas a continuos movimentos ou mudanças. Ora os movimentos, sendo passagens da \emph{potencia} para o \emph{acto}, devem ser determinados --- n\&o pelas proprias substancias que se movem, pois ninguem dá o que não tem, --- mas por um motor distincto, pois o ente, dotado apenas de \emph{potencia}, não adquire o \emph{acto}, se lo não communiea um outro ente, que já o possue. Este motor --- ou é \emph{movido}, --- ou é \emph{immovel}. Se é \emph{immovel}, é \textsc{Deus}. Se é \emph{movido}, subiremos de motor em motor, ató encontrarmos necessariamente um motor \emph{immovel}; porque, se não existisse este motor \emph{immovel}, não poderiam existir os motores \emph{moveis}, que não têm em si a razão do seu movimento, e por isso não haveria movimento nas substancias creadas. Logo existe \textsc{Deus} \footnote{Aqui a palavra \emph{movimento} toma-se em sentido mais lato, emquanto exprime n\&o só o \emph{movimento local}, mas toda e qualquer \emph{mudança} ou \emph{passagem da potencia para o acto}. D'este modo o argumento exposto equivale a este: O mundo é \emph{mutavel}. Ora todo o ente \emph{mutavel} depende do ente \emph{immutavel}. Logo existe o ente \emph{immutavel}. O ente \emph{immutavel} é \textsc{Deus}. Logo existe \textsc{DEus}. --- Ora, encontrando-se a \emph{mudança} não só nos corpos mas tambem nos espiritos, o argumento póde desdobrar-se; pois, baseando-nos tanto no movimento dos \emph{corpos} como no dos \emph{espiritos}, podemos elevar-nos ao descobrimento de uma causa, que move tudo, sendo ella \emph{immovel} ou \emph{immutavel}.\par Eis os dois argumentos:\par a) A experiencia attesta que as substancias materiaes, por exemplo, os planetas, estão sujeitas a continuos movimentos. Ora todo o corpo que se move, é movido por um agente; pois o corpo, sendo por si indifferente tanto para a quietação como para o movimento, não póde determinar-se por si ao movimento. Logo ou devemos admittir uma serie infinita de motores \emph{moveis}, ou deve existir um motor \emph{immovel}, que seja a causa de todo o movimento dos corpos. Ora uma serie infinita de motores \emph{moveis} repugna; porque numa serie \emph{infinita} de motores \emph{moveis} não haveria um \emph{primeiro} motor, e assim não existiria uma causa proporcionada ou adequada de todos os movimentos\pgn{431}nos corpos. Logo deve existir um motor \emph{immovel}. Este motor \emph{immovel} deve ser infinitamente perfeito; pois, se o hão fosse, se lhe faltasse alguma perfeição, poderia tender para ella, e assim passaria da, potencia ao acto, \emph{e} não seria \emph{imnmovel}. O ser infinitamente perfeito é \textsc{Deus}. Logo existe \textsc{Deus}.\par Este argumento, muito popular, pode desenvolver-se \emph{scientificamente} da maneira seguinte. Existe no mundo o movimento. Para a explicação da existencia do movimento adduzem-se duas hypotheses: a \emph{mechanica} e \& \emph{dynamica}. Ora nenhuma d'estas hypotheses pode, na explicação do movimento, prescindir da intervenção de um primeiro Motor. --- Não pode prescindir \& hypothese \emph{mechanica}. Na verdade, para os sequazes d'esta hypothese, a materia é \emph{meramente passiva}, de modo que os corpos são movidos por um impulso extrinseco, mas elles não movem, e só communicam aos outros corpos o movimento recebido. Se todos os corpos, de que se compõe o universo, são meramente passivos, é necessario que a força impulsiva do movimento venha de fóra, e proceda de um primeiro Motor, que não seja movido. Nem se diga que o movimento da materia teve um começo \emph{necessario}, ou que é \emph{eterno}; porque um \emph{começo necessario} é uma contradicção (o que começa ê necessariamente contingente), \emph{e} um movimento \emph{eterno} da materia, embora não envolvesse repugnancia, não deixaria de exigir um primeiro motor, desde que se supponha a materia meramente \emph{passiva}. --- Não pode prescindir a hypothese \emph{dynamica}. Com effeito, para os sequazes d'esta hypothese, \& materia não é meramente passiva, mas é essencialmente dotada de actividade, em virtude da qual um corpo pode, em certas circumstancias, mover outro corpo, \emph{attrahendo-o}. E' certo que o mundo inteiro se move pela \emph{lei da attracção}; mas a \emph{attracção} não é o \emph{primeiro} motor, nem é o motor \emph{sufficiente} do nosso systema solar. Nào é o \emph{primeiro} motor; porquanto, a força de \emph{attrac}- “\emph{ção}, de que é dotada a materia, não está sempre e necessariamente em acto, mas, pelo contrario, tende para o repouso das partes numa unica massa, e por isso é indifferente para o acto e para a potencia, e, não contendo em si a razão da passagem da potencia para o acto, exige e suppõe uma causa que a determine e mova. Não é o motor \emph{sufficiente} do nosso systema solar; porquanto este systema deve ser mantido em equilibrio e regulado nos seus movimentos por duas forças --- uma \emph{centripeta} (sem a qual não haveria unidade no systema) e outra \emph{centrifuga} (sem a qual todos os corpos confundir-se-hiam numa unica massa) e se a força \emph{centripeta}, da qual deriva a attracgüo, pode dizer-se essencial á materia, não se pode dizer essencial a força \emph{centrifuga}, a qual por isso deve ser accrescentada e unida por um motor extranho e superior á propria materia.\par Dizem que este argumento, se prova a existencia de um Ente superior ao mundo e causa de todo o movimento, não prova a existencia de um Deus tão perfeito, como o concebem os espiritualistas. --- Respondemos que o nosso argumento, embora não provasse senão a \emph{existencia} de um Ser supremo, já teria alcançado o seu fim, que é o de demonstrar a mesma \emph{existencia}, e não a \emph{natureza} do Ente supremo. São duas questões distinctas. --- Todavia devemos admittir que este argumento, que prova a necessidade de um \emph{primeiro Motor}, contém, \emph{em germe}, a razão de todos os attributos da Divindade.\pgn{432}Na verdade, se a materia, por ser indifferente ao repouso e ao movimento, tem a razão do seu movimento não em si mas numa causa exterior, segue-se que tal causa existe fóra da materia, e por isso é immaterial, intelligente, eterna, necessaria, e tambem infinita e unica. Porquanto, o primeiro Motor, não podendo ser movido, é \emph{todo} «\emph{cto}, sem mistura de potencialidade, é acto \emph{purissimo}, isto é, infinitamente perfeito e é unico.\par Os melhores naturalistas e philosophos modernos reconhecem esta verdade. Citaremos apenas \emph{Couchy}, \emph{Clarke}, \emph{B}. \emph{S}. \emph{Hilaire}, \emph{P}. \emph{Janet}. No fim, ouviremos um insuspeito, \emph{Rousseau}.\par \emph{Couchy} diz: « O unico ser, que póde produzir as forças physicas, é o Ente necessario. A força é a expressão da sua vontade. As diversas forças do equilibrio e movimento são causas secundarias. A gravitação universal, a attraecüo dos corpos na razão inversa do quadrato das distancias, que os separam, é uma lei que exprime a vontade divina » (\emph{Comptes rendus hebd}. \emph{de lAcad}. \emph{des sciences}, V. XXI, Paris, 1845).\par \emph{Clarke}, discipulo de \emph{Newton}, fallando da \emph{attracção}, de que depende todo o movimento dos corpos celestes, escreve: « Ella póde ser o effeito de um impulso, mas não póde ser material, ó antes o effeito de uma causa immaterial » (\emph{Physica de Rohault}, trad. por \emph{Clarke}, II, c. XI).\par \emph{Barthélemy Saint-Hilaire} diz: « Concebe-se que a materia siga hoje as leis infalliveis, que a regem, e conserve uma regularidade eterna; mas foi necessario um impulso primitivo, que ordenou tudo para a inexhauriyel serie dos tempos. Até as anomalias, que offerece o admiravel systema dos ceus, attestam a presença immanente e indefectivel d'Aquelle que os fez. O sobrenatural está em tudo... Só é necessario que a sciencia se resigne a resolver certos problemas por outros meios, que não sejam os de uma experiencia impossivel; e o da origem de todas as coisas é um d'esses problemas, a que só se renuncia por timidez, pensando que se pratica uma reserva prudente. A \emph{questão da origem é inevitavel}, e de nada valeria querer illudil-a ». (\emph{Journal des Savants}, 1862, pag. 607).\par \emph{P}. \emph{Janet} diz: « A experiencia não mostra senão que um corpo é movido por outro, e este por outro, e isto ao infinito. Ora, se fóra d'esta cadêa não existe uma causa motora, teremos uma serie de movimentos sem causa e sem razão » (\emph{Le matérialisme contemporain}, 1875, pag. 58).\par \emph{Rousseau} escreve: « Vejo a materia ora em movimento e ora em repouso; d'ahi coneluo que nem o repouso nem o movimento lhe são essenciaes. Mas o movimento, sendo uma acção, mostra ser effeito de uma causa... Conceber a materia productora do movimento é conceber um effeito destituido de. causa... Não é evidente que, se o movimento fosse essencial á materia, um não poderia separar-se da outra, mas estaria sempre unido com ella, conservando-se sempre no mesmo grau, sempre o mesmo em cada uma das partes, não poderia augmentar nem diminuir, e não se poderia conceber a materia em repouso?!... Quando me dizem que o movimento não é \emph{essencial} \& materia, mas \emph{necessario}, pretendem embaracar-me com palavras, que poderiam facilmente confutar-se, se tivessem algum sentido. Porque o movimento da materia --- ou deriva da propria materia, e então lhe é essencial\pgn{433} --- ou deriva de uma causa extrinseca, e já não é necessario á propria materia, senão emquanto a causa motora actua nella » (\emph{Emilio}, t. III).\par Na \emph{Cosmologia geral} tambem desenvolvemos este ponto.\par b) No mundo não só existe o movimento \emph{local}, mas tambem o movimento \emph{espiritual}; e este, como aquelle, suppõe e exige necessariamente a existencia de um primeiro motor immovel. Na verdade, se reflectirmos em nós mesmos, “percebemos que a nossa alma tambem se move, emquanto passa continuamente da potencia para o acto. O acto é um complemento, uma addição que aperfeiçõa a potencia. Ora este acto não pode ser e não é extrahido do fundo da potencia, porque ninguem dá o que não tem; mas deve provir de um agente extrinseco, de um motor, que possua \emph{actualmente} o que falta \& potencia e que seja capaz de lh'o communicar. Este motor é a propria alma, que se move a si mesma emquanto uma faculdade, que está em \emph{acto}, move uma outra, que esta \emph{em potencia}. Todavia a alma não pode ser o \emph{primeiro} motor. O primeiro acto da faculdade, que move todas as outras, não pode ser o effeito do impulso das outras faculdades, que estão em potencia, mas deve ser o effeito de um motor distincto e superior \& propria faculdade. Este motor --- ou é o primeiro Motor, --- ou é um motor intermedio, movido, finalmente, por esse primeiro Motor immovel.\par Este argumento, que se refere ao movimento da alma humana em geral, conserva toda a sua força, quando se refere ao movimento da \emph{intelligencia} e da \emph{vontade}.\par a) Dá-se o movimento na :\emph{ntelligencia}. Ainda que a intelligencia mova a vontade, e esta mova aquella, comtudo, por ser absurdo o processo ao infinito, devemos admittir que o primeiro acto da intelligencia precedeu o primeiro acto da vontade. Diz S. Thomaz: « Non oportet procedere in infinitum, sed statur in intellectu, sicut in primo. Omnem enim voluntatis motum necesse est quod praecedat apprehensio; sed non omnem apprehensionem praecedit motus voluntatis » (\emph{Sum}. \emph{Th}., p. I, q. 82, a. 4, ad 8). Mas de que modo passa a intelligencia da potencia para o acto? Espontaneamente? Não; porque faltaria a razão sufficiente da passagem. Pela acção do objecto intelligivel? Tambem não; porque, sendo este objecto formado pela intelligencia agente, a acção do objecto é \emph{posterior} á da intelligencia. A. causa motora deve ser superior ao objecto e a sujeito, ao intelligivel e ao inteligente; e esta causa é Deus. Aristoteles, ainda que pagão, exprime esta verdade nas seguintes magnificas palavras: « Quod est principium motus in anima? Patet vero, quemadmodum in universo est Deus, sic etiam in anima; moventur enim omnia ab inexistente in nobis Numine. Intellectus autem principium non intellectus sed praestantius quidpiam est. Quid autem et scientia praestantius est et ratione, nisi Deus? (\emph{Moral}. \emph{Eudem}., l. 7, c. 14). D'onde se vê a verdade do aphorismo: \emph{Eu penso}; \emph{logo Deus existe}.\par \emph{b}) Não só o: acto da intelligencia, mas tambem o da \emph{vontade} exige e suppõe a existencia de um primeiro Motor. A vontade é movida necessariamente pelo bem absoluto. Esse movimento provoca uma tendencia \emph{wuniversal}, \emph{necessaria}, \emph{indestructivel}. Ora um effeito, dotado de semelhantes qualidades, como é essa tendencia, não pode derivar de uma causa imaginaria ou phantastica, mas deve derivar de uma causa real e objectiva, que deve\par 28 Vor. II.\pgn{434}identificar-se conv o proprio bem absoluto; porque todo o effeito suppõe uma causa proporcionada. Ora o bem absoluto é Deus. Logo Deus existe.\par Como se vê, o principio --- \emph{tudo o que se move é movido por um motor distincto e}, \emph{em ultima analyse}, \emph{immovel} --- é muito fecundo e serve para a intelligencia e resolução de innumeraveis problemas.}.

b) No mundo existem muitas e diversas \emph{causas efficientes}. Estas causas estão por tal modo subordinadas entre si, que formam uma serie, na qual uma depende de outra, a inferior \pg{431} \pg{432}  é \emph{effeito} da superior. Ora, nesta serie de causas \emph{eficientes}, deve existir uma causa \emph{primeira}, que seja só \emph{causa} e não \emph{effeito}, que seja causa de todas as causas; visto que, se não existisse uma \pg{433}  causa \emph{primeira}, não poderiam existir, como existem, as causas \emph{segundas}, pois as causas \emph{segundas} são \emph{effeitos}, e todo o effeito \pg{434}  exige necessariamente uma causa proporcionada. A causa \emph{primeira} é \textsc{Deus}. Logo existe \textsc{Deus} \footnote{Os \emph{materialistas}, para destruirem a força d'este argumento, recorrem a uma serie \emph{infinita} de causas \emph{efficientes}, das quaes uma produza outra, sem sermos obrigados a admittir uma \emph{primeira} causa \emph{efficiente}.\par O recurso é vão. Porquanto, essa serie infinita, embora fosse \emph{possivel}. estaria disposta ou em linha \emph{recta}, ou em linha \emph{circular}. Ambas as hypotheses são absurdas. --- E' absurda a \emph{primeira} ; porque, se estivesse disposta em linha \emph{recta}, de modo que a causa inferior fosse produzida pela superior, essa serie não teria uma causa \emph{primeira} (porque o infinito não tem principio). Não tendo uma causa \emph{primeira}, não teria a \emph{segunda}, nem a \emph{terceira}, nem causa alguma; visto que numa serie de causas, dispostas em linha \emph{recta}, a inferior não póde existir sem a superior. Logo deve admittir-se, fóra d'esta serio, uma causa primeira. --- E' absurda a segunda hypothese. Na verdade, se a serie \emph{infinita} das causas \emph{efficientes} estivesse disposta em linha \emph{circular}, um e o mesmo ente seria causa e effeito de outro ente. Na. verdade, supponhamos uma linha \emph{circular}, formada por a, b, o, d... E' evidente que, sendo \& causa de b, b de c, o de d, d de a, d seria causa e effeito (ainda que remoto) de a; o que é absurdo.\par Os adversarios dizem que a \emph{primeira} causa \emph{efficiente} é a \emph{materia}.\par Esta asserção tambem é absurda. Por quanto, como provámos, a materia não é nem póde ser eterna, mas foi produzida, e por isso não póde ser a causa \emph{primeira} de todas as coisas.\par Note-se a differenca entre este e o precedente argumento. Aqui, \textsc{Deus} considera-se como a causa \emph{efficiente} de tudo o que existe; além, \textsc{Deus} con- Siderou-se tambem como causa \emph{final} de todas as coisas, emquanto tudo tende para \textsc{DEus}, como Ente infinitamente intelligivel e amavel.}.

c) Os entes visiveis são todos \emph{contingentes}, isto é, existem mas podiam deixar de existir. Se são \emph{contingentes}, foram produzidos por outro ente; porque o \emph{contingente}, podendo existir e deixar de existir, não contém em si a razão sufficiente da propria existencia. Este outro ente, que produziu o \emph{contingente}, ou é tambem \emph{contingente}, ou é \emph{necessario}. Se é \emph{necessario}, É \textsc{Deus}. Se é \emph{contingente}, foi produzido por outro, e este por outro, até se encontrar inevitavelmente o ente \emph{necessario} ; porque, se não existisse o ente \emph{necessario}, não poderia existir nenhum ente \emph{contingente}, pois repugna uma serie de entes \emph{dependentes} sem um \emph{primeiro}, de que todos dependam. Logo existe \textsc{Deus} \footnote{Em poucas palavras: O mundo é \emph{contingente}. Ora o ente \emph{contingente} depende do ente \emph{necessario}. Logo existe um ente \emph{necessario}. O ente necessario é \textsc{Deus}. Logo existe \textsc{Deus}.\par \emph{Vazherot} afirma muitas vezes (sem nunca o demonstrar) que o mundo,\pgn{435}se é \emph{contingente} na \emph{apparencia}, \emph{é necessario} e \emph{absoluto} na \emph{realidade}. --- Esta affirmação é absurda. O mundo, como demonstrámos na \emph{Cosmologia geral}, \emph{é essencialmente contingente} no seu conjuncto e em cada uma das suas partes.\par \emph{Kant} diz que a prova, deduzida da idea do ser \emph{necessario}, não tem valor; porque, como a prova \emph{ontologica}, passa da ordem \emph{ideal} para a \emph{real}, apoiando-se num principio meramente \emph{subjectivo}, que é o de \emph{causalidade}.\par Rejeitamos com \emph{Kant} a prova \emph{ontologica} ácerca da existencia de \textsc{DEus}, mas negamos que a prova \emph{cosmologica} se reduza á \emph{ontologica}. Na verdade, a prova \emph{cosmologica} não consiste, como pensa \emph{Kant}, em deduzir da \emph{idea} de um ser \emph{necessario} a sua existencia \emph{real}, mas em deduzir da existencia \emph{real} dos seres \emph{contingentes} a existencia \emph{real} do ser \emph{necessario}, que é Dxgus. --- Além d'isso, negamos que o principio de \emph{causalidade} seja meramente \emph{subjectivo}; elle é tambem, e primeiramente, \emph{objectivo}, porque não é lei do nosso espirito, senão porque é lei das coisas.}.

\pg{435} d) Os entes creados, embora possuam todos uma perfeição \emph{finita}, não á possuem todos no mesmo grau, mas um é \emph{mais} perfeito, e outro o é \emph{menos}. Isto mostra que a perfeição das creaturas é \emph{participada}; porque, se o não fosse, não haveria razão por que numa creatura a perfeição fosse \emph{maior}, e noutra \emph{menor}. Sendo \emph{participada}, a perfeição das creaturas deve derivar de um ente, e de um ente, que a possua pela \emph{sua propria essencia}. Ora a perfeição, que um ente possue pela sua \emph{propria essencia}, deve ser \emph{infinita}; porque, se fosse \emph{finita}, seria \emph{participada}. Logo existe um ente \emph{infinitamente} perfeito, de que derivaram, como de um oceano immenso, as muitas e diversas perfeições das creaturas. O ente \emph{infinitamente} perfeito é \textsc{Deus}. Logo existe \textsc{Deus} \footnote{Em poucas palavras: Os entes creados são todos \emph{finitos}. Ora a existencia dos entes \emph{finitos} suppõe necessariamente a existencia do ente \emph{infinito}. Logo existe o ente \emph{infinito}. O ente \emph{infinito} é \textsc{DEus}. Logo existe \textsc{Deus}.\par Objectam: A diversidade das perfeições, que se encontram nos seres finitos, não demonstra a existencia de um ente infinito.\par \emph{A} objecção já está refutada no texto. Entre os seres finitos, um é \emph{mais} ou \emph{menos} perfeito que o outro. Ora uma coisa diz-se \emph{mais} ou menos perfeita, emquanto \emph{mais} ou \emph{menos} se approxima de um ente \emph{absolutamente perfeito}: assim uma photographia diz-se \emph{mais} ou \emph{menos} perfeita, emquanto \emph{mais} ou menos se approxima do original, que é uma coisa \emph{absoluta}. Logo deve existir um ente \emph{absoluta} e \emph{infinitamente} perfeito, que seja a causa das varias e \emph{limitadas} perfeições, que se encontram nas creaturas. --- O \emph{relativo}, que se não refira ao \emph{absoluto}, é uma contradicção.}.

e) A intelligencia humana conhece verdades \emph{absolutas} e

\emph{necessarias}, tanto da ordem \emph{especulativa} como da \emph{pratica}. Taes verdades devem ter uma \emph{razão sufficiente}. Esta \emph{razão sufficiente} não se encontra nos \emph{seres creados}, nem na propria \emph{intelligencia humana}. Não se encontra nos \emph{seres creados}; pois estes, sendo \emph{relativos} e \emph{contingentes}, não podem produzir o \emph{absoluto} e O \pg{436}  \emph{necessario}. Nem se encontra na :\emph{ntelligencia creada}; pois esta não forma as verdades \emph{absolutas} e \emph{necessarias}, descobre-as apenas e recebe-as como regra do seu pensamento. Logo a \emph{razão sufficiente} d'aquellas verdades deve encontrar-se numa \emph{intelli}- \emph{Jencia absoluta} e \emph{necessaria}; pois ellas são entidades \emph{ideaes}, \emph{absolutas} e \emph{necessarias}. A intelligencia \emph{absoluta} e \emph{necessaria} é a de \textsc{Deus}. Logo existe \textsc{Deus} \footnote{E' certo que, além da ordem \emph{physica}, composta de substancias materiaes, existe uma ordem \emph{logica}, composta de verdades \emph{necessarias}, \emph{absolutas}; por quanto estas verdades não podem ser formadas nem mudadas pela nossa inteligencia. Esta faculdade limita-se a descobril-as, e descobre-as como independentes de toda a experiencia e de todo o pensamento creado. Logo ou devemos negar a verdade do principio \emph{de} razão \emph{sufficiente}, ou deve existir uma intelligencia \emph{necessaria} e \emph{absoluta}, em que estas verdades encontram o seu fundamento. Ora o principio \emph{de razão sufficiente} é verdadeiro. Logo deve existir uma intelligencia \emph{absoluta} e \emph{necessaria}, que é a de \textsc{Deus}.\par Nem se diga que, neste argumento, passamos da ordem \emph{ideal} para a \emph{real}. Por quanto, as verdades \emph{absolutas} e \emph{necessarias} são \emph{objectivas} ; pois, como dissemos, ellas são verdades independentemente do nosso pensamento, não podemos formal-as nem mudal-as, são sempre o que são, embora não pensemos nellas e as contestemos. Sendo \emph{objectivas}, devem encontrar na \emph{wvealidade} um fundamento. Este fundamento não se encontra nas coisas mutaveis, nem na intelligencia creada. Logo deve encontrar-se num Ente \emph{necessario} e \emph{eterno}, em \textsc{Deus}. :}.

13. \textbf{Argumentos physicos}. --- Os argumentos \emph{physicos} basóam-se em que, se existe uma ordem maravilhosa no mundo e nas operações dos agentes \emph{irracionaes}, deve existir um ordenador infinitamente sabio e poderoso.

a) Uma ordem maravilhosa e constante existe entre as muitas e diversas substancias, que compõem o universo. Esta ordem exige uma causa infinitamente sabia e poderosa; pois SÓ essa causa, --- por ser \emph{infinitamente sabia}, podia conhecer a natureza de todas as substancias e o posto de cada uma d'ellas no plano da ereac\&o, e escolher as leis mais aptas para a constancia e regularidade da mesma ordem; --- e, por ser \emph{infinitamente poderosa}, podia executar esse plano. A causa infinitamente \emph{sabía} e \emph{poderosa} é \textsc{Deus}. Logo existe \textsc{Deus} \footnote{Não insistimos neste ponto; porque na \emph{Cosmologia geral} provámos que a ordem \emph{physica}, com as suas \emph{leis}, é effeito de uma intelligencia \emph{infinita}.\par \emph{Kant} diz: A ordem \emph{physica} é effeito de uma intelligencia, mas não de uma intelligencia \emph{infinita}. Logo não demonstra a existencia de \textsc{DEus}.\par O que dizemos no texto mostra a falsidade da affirmação de \emph{Kant}. --- Além d'isso, a inteligencia, que estabeleceu a ordem no mundo, só podia\pgn{437}ser a do Ente, que \emph{criou} o proprio mundo. Ora o ente, que \emph{criou} o mundo, deve ter um poder \emph{infinito}. O ente, dotado de poder \emph{infinito}, deve ter uma natureza \emph{infinita} e por isso uma intelligencia \emph{infinita}.\par Os materialistas objectam: Esta ordem e regularidade, que existe em todo o universo e em cada substancia, não é effeito de uma causa intelligente e deliberada, mas é obra do \emph{acaso}.\par Resposta. O \emph{acaso}, sendo \emph{privação} de intelligencia e deliberação, é \emph{nada}. Ora é absurdo dizer que o \emph{nada} produziu a \emph{realidade}. Logo a ordem, que existe no mundo, sendo uma \emph{realidade}, não foi produzida pelo acaso, mas por um Ente real, infinitamente sabio e poderoso.}.

b) Não só a ordem, que existe no universo, mas tambem a que existe nas operações de cada agente \emph{irracional}, proclama \pg{437}  a existencia de uma Intelligencia soberana e infinita. Na verdade, todo o agente \emph{irracional}, nas suas operações, tende para um fim, empregando para isso os meios mais efficazes, e exeeutando um plano com maravilhosa exactidão, constancia e uniformidade. Este fim, esta escolha de meios e este plano, devem ser effeitos de uma intelligencia \emph{infinita}; porque só uma intelligencia \emph{infinita} podia determinar a cada agente \emph{tirra}= \emph{cional} um fim proporcionado, modelar planos complicados e ao mesmo tempo perfeitos, e, para os realizar, servir-se da forças cegas da natureza. A intelligencia \emph{infinita} é propria de \textsc{Deus}. Logo existe \textsc{Deus} \footnote{Tambem este ponto ácerca da \emph{finalidade} nas operações dos agentes \emph{irracionaes} foi sufficientemente tratado na \emph{Cosmologia geral}.}:

14. Argumentos \textbf{moraes}. --- Os argumentos \emph{moraes} fundam-se em que, --- a) se existe uma lei \emph{moral}, deve existir um legislador, --- 25) se todos os povos admittiram a existencia de um Ente supremo, Este deve existir.

a) A consciencia attesta-nos que temos gravada no nosso coração uma lei moral, á qual não nos podemos subtrahir, e que o cumprimento d'essa lei produz tranquillidade, assim como a sua transgressão causa remorso. Ora a existencia d'esta lei, como a respectiva sancção, suppõe a existencia de um legislador. Este legislador não póde ser o proprio homem; aliás seria inexplicavel o sentimento da nossa dependencia de um poder superior. Logo deve ser um Ente superior ao homem e dotado de virtude ou poder \emph{infinito}; porque só um ente infinitamente poderoso podia gravar no coração dos homens essa lei. O Ente superior, de que o homem depende e a que deve obedecer, é \textsc{Deus}. Logo existe \textsc{Deus} \footnote{Este argumento, apresentado tambem por \emph{Kant} para a demonstração da existencia de \textsc{Deus}, será mais desenvolvido na \emph{Moral}.\par Aqui apenas adyertimos que a existencia da \emph{lei moral} considera-se como um \emph{facto experimental}, attestado pela consciencia. Provado o facto da existencia da \emph{lei}, legitimamente deduz-se a existencia de um \emph{legislador}.}.

\pg{438} b) Todos os homens, em todos os tempos e em todos os logares, reconheceram e admittiram a existencia de um Ente supremo. Este consenso unanime, constante e universal, não póde ser effeito do \emph{erro}; pois o \emph{erro} é causa de discordias, é mutavel e particular. Logo é effeito da \emph{verdade}; porque só a \emph{verdade} é unanime, constante e universal, e só ella pode le- 'var os homens a mortifiear as paixões e renunciar á propria vontade, para se sujeitarem á vontade de um Ente supremo. Logo existe \textsc{Deus} \footnote{O consenso do genero humano, relativamente á \emph{existencia} de um Ente supremo, não podia ser mais constante, nem mais universal. Esta crença num Ente supremo podia ser e foi desfigurada por muitos erros; mas conservou sempre o fundo intacto. Muitos povos ignoraram e ignoram ainda a \emph{essencia} de \textsc{DEus}, mas nenhum ignorou ou ignora a sua \emph{existencia}.\par \emph{Cicero} disse: « Não existe povo algum, por inculto e selvagem, que não tenha fé em \textsc{DEus}, embora lhe não conheça a essencia » (\emph{De nat}. \emph{Deor}. I, 17).\par \emph{Quatrefages} escreve: « Obrigado pelo meu ensino a passar revista a todas as raças humanas, procurei o \emph{atheismo} tanto entre os povos mais selvagens, como entre os mais civilisados. Não \emph{o encontrei em parte alguma}, a \emph{não ser nalgum individuo ou eschola muito limitada}, como se viu na Europa no seculo passado e como se vê ainda no presente » (\emph{fespéce humaine}, pag. 955, Paris, 1871)... E' noutro logar não se exprime menos claramente: « Pouco a pouco fez-se a luz, e os povos da \emph{Australia} e da \emph{Melanesia}, os \emph{Boschimans}, os \emph{Ottentotes}, os \emph{Cafres}, os \emph{Bechuanas}, foram eliminados do numero dos povos atheus e reconhecidos como religiosos » (\emph{Rapport}, pag. 410).\par De resto, os \emph{Anaaes da propagação da} Né provam, com razões claras e baseadas na experiencia, que povos atheus não passam de um mytho.\par Objectam: A crença de todos os povos na existencia de um Ente supremo derivou da \emph{ignorancia}, ou da astucia dos \emph{sacerdotes}, ou da prepotencia dos \emph{legisladores}, ou do \emph{medo}, ou dos \emph{prejuizos}.\par Resposta: A crença unanime de todos os povos em \textsc{Deus} não derivou --- a) nem da \emph{ignorancia}, --- \emph{b}) nem da astucia dos \emph{sacerdotes}, --- \emph{c}) nem da prepotencia dos \emph{legisladores}, --- \emph{d}) nem do \emph{medo}, --- \emph{e}) nem dos \emph{prejuizos}.\par a) Não derivou da \emph{ignorancia}. Por quanto, o estudo das causas naturaes é o meio mais facil e efficaz para se conhecer a existencia de \textsc{DEus}; visto que essas causas, sendo \emph{secundarias}, \emph{contingentes}, \emph{imperfeitas}, exigem uma causa \emph{primeira}, \emph{necessaria}, \emph{perfeitissima}, que é \textsc{Deus}.\par A experiencia confirma a these. Na verdade, os maiores naturalistas são os mais convencidos e enthusiastas adoradores de \textsc{Deus}. --- Citaremos apenas os testemunhos de \emph{Linneu} e \emph{Kepler}.\par \emph{Linneu} escreve: « Quando acordei, \textsc{Deus} \emph{eterno}, \emph{inmenso}, \emph{omnisciente}, \emph{omnipotente}, acabava de passar. Vi-o ao longe, e fiquei profundamente maravilhado. Segui os vestigios de seus passos, atravez das obras da creação, e por toda a parte, mesmo nas mais pequenas coisas, que poder! que sabedoria! que ineffavel perfeição! O sol e todo o systema planetario, immenso, inealeulavel, appareceram-me suspensos pelo Primeiro Motor, Causa das\pgn{439}causas, Governador e Conservador do universo. Todas as coisas são, pois, testemunho da sabedoria divina; a belleza, a harmonia, as justas proporções das creaturas proclamam o poder do grande \textsc{Deus} » (\emph{Syst}. \emph{de la Nat}. p. 10).\par \emph{Kepler} conclue a sua obra, intitulada « \emph{Os quatro livros da harmonia celeste} », por estas sublimes palavras: « E agora só me resta levantar as mãos e os olhos para ao ceu, e dirigir ao Auctor de toda a luz a minha humilde oração: O' Vós, que, pelos esplendores sublimes que derramastes sobre toda a natureza, elevais os nossos desejos até á divina luz da vossa graça, sede bemdito! Senhor e Creador de todas as consolações, que experimentei no extasis, a que me elevou a contemplação da obra das vossas mãos, eis acabado este livro, que contém o fructo do meu trabalho e de toda a intelligencia, que me destes. Proclamei deante dos homens toda a grandeza das vossas obras, mostrando-lhes, na medida das minhas forças, a sua immensa perfeição. Esforcei-me por chegar até á verdade, por a conhecer o mais perfeitamente possivel, e, se me escapou alguma coisa indigna de Vós, fazei que a conheça para a eliminar. Ter-me-hei deixado seduzir pela presumpção, na presença da belleza maravilhosa das vossas obras? Terei eu procurado a minha gloria entre os homens, levantando este monumento, que só devia ser consagrado á vossa honra? Ah! se assim foi, tende piedade e misericordia de mim, e concedei-me a graça de que o meu trabalho não produza mal algum, mas concorra para a glorificação do vosso Nome e para a salvação das almas ».\par Os mesmos sentimentos encontram-se nos escriptos dos mais celebres naturalistas, como \emph{Copernico}, \emph{Galileu}, \emph{Newton}, \emph{Jussieu}, \emph{Cuvier}, \emph{Agassis}," \emph{Biot}, \emph{Buffon}, \emph{Lavoisier}. \emph{Brongniart}, \emph{Blainville}, \emph{Bertholet}, \emph{Gay-Lussac}, \emph{Elie de Beaumont}, \emph{Thénard}, \emph{Réaumur}, \emph{Barande}, \emph{Milne-Edwards}, \emph{Dumas}, \emph{Chevreul}. \emph{Liebig}, \emph{Flourens}, \emph{Quatrefages}, \emph{Pasteur}, etc. Nem admira; porque, como diz \emph{Herschel}, ü quanto mais se extende o campo da Physica, tanto mais numerosas e\pgn{440}inconcussas se tornam as demonstrações da existencia eterna de uma intelligencia creadora e omnipotente; os geologos, os astronomos, os naturalistas. trouxeram todos a sua pedra para este grande templo da sciencia, --- templo levantado ao proprio \textsc{DEUS} n. --- Não nos cancaremos de repetir com \emph{Bacon}; u A pouca sciencia afasta de \textsc{Deus}, a muita conduz a Elle ».\par b) Não derivou da astucia dos \emph{sacerdotes}; pois a instituição dos \emph{sacerdotes é effeito}, e não \emph{causa}, da crença universal na Divindade.\par 0) Não derivou da prepotencia dos \emph{legisladores}. Os legisladores basearam as leis na propria crença religiosa dos seus povos. --- Além d'isso, os povos não teriam acceitado, ao menos em toda a parte, uma crença, que é contraria a todas as paixões. --- Finalmente, as leis variam em todos os povos, mas a crença em \textsc{Deus} foi sempre a mesma por todu a parte.\par 4) Não derivou do \emph{medo}. Na verdade, esse \emph{medo} ou era conforme com a razão, ou não. Se não era conforme com a razão, não podia produzir a crença em \textsc{DEus}; pois essa crença é universal e constante, e um effeito universal e constante não póde derivar de uma causa particular e mutavel, como é o \emph{medo}. Se era conforme com a razão, esse \emph{medo} já suppunha a crenca em \textsc{Deus}; pois ninguem tem \emph{medo} de um ente, que se não conhece como existente. --- A verdade é que o \emph{medo}, longe de produzir a crença em \textsc{Deus}, deu sempre occasião ao \emph{atheismo}, pois os libertinos, \emph{pelo facto de temer o castigo}, esforcam-se por negar a existencia de \textsc{Deus}.\par €) Não derivou dos \emph{prejuízos}; pois os \emph{prejuízos} variam segundo a diversidade dos povos, ao passo que a crença em \textsc{DEus} é universal e constante.}.

\pg{439} EI

%%>>> art-38

\section*{Artigo III.\\ A. existencia de Deus e o atheismo.}
\addcontentsline{toc}{section}{Artigo III. --- A. existencia de Deus e o atheismo.}

15. \textbf{Atheismo}. --- \emph{Atheismo} é o erro dos que não reconhecem ou negam a \emph{existencia} de \textsc{Deus}.

16. \textbf{Divisão} do \textbf{atheismo}. --- O \emph{atheismo} divide-se em \emph{pratico} e \emph{theoretico}. --- O \emph{pratico} é o erro dos que, embora affirmem a existencia de \textsc{Deus} com a \emph{razão}, negam-na com os \emph{factos}, pois vivem, como se \textsc{Deus} não existisse. --- O \emph{theoretico} É o erro dos que negam com a \emph{razão} a existencia de \textsc{Deus}.

17. \textbf{Divisão} do atheismo theoretico. --- O atheismo \emph{theoretico} póde ser \emph{negativo} ou \emph{positivo}. --- O \emph{negativo} é o erro dos que não admittem a existencia de \textsc{Deus}, porque o não conhecem. --- O \emph{positivo} é o erro dos que, tendo conhecido a \textsc{Deus}, estão \emph{séria} e \emph{sinceramente} convencidos de que Elle não existe.

\pg{440} 18. O \textbf{atheismo pratico} é possivel. --- O atheismo \emph{pratico} não é uma depravação da \emph{intelligencia}, mas sim da \emph{vontade}. Ora a nossa vontade está sujeita á depravação, e, para se entregar mais livremente ás desordens, póde deixar de pensar em \textsc{Deus} e de Lhe prestar as homenagens da sua piedade e obediencia. --- Este \emph{atheismo} não só é \emph{possivel}, mas existe realmente \footnote{\emph{Bayle} diz: « Os que não servem a \textsc{Deus}, e que antes o offendem com uma vida cheia de crimes e vicios, são uma especie de atheus, é o atheismo dos libertinos e devassos; porque effectivamente vivem, como se não houvesse Deus, sem medo da sua Justiça, sem gratidão para com a sua Bondade, sem respeito para o seu Nome, sem obediencia ás suas leis, e, ainda que podessem tirar \textsc{DEus} do seu throno para o esmagar sob os seus pés e aniquilar para sempre, não viveriam com mais licença e ousadia. Sao atheus pelas suas acções; conhecem a \textsc{Deus} com a lingua, mas negam-no com os costumes; fallam como crentes, mas operam como incredulos » (\emph{Pensées div}., t. IV, pag. 91). E}.

19. O \textbf{atheismo} theoretico-negativo não é possivel nos homens \textbf{adultos}. --- A existencia de \textsc{Deus}, embora nos não seja \emph{immediatamente} evidente, todavia descobre-se logo que começamos a applicar os primeiros principios da razão. Por quanto, sendo natural e facillimo deduzir da existencia dos effeitos a existencia da causa respectiva, é claro que, contemplando o universo, não podemos deixar de reconhecer uma causa, que o \pg{441}  produziu e formou. Logo o atheismo \emph{theoretico-negativo} é impossivel nos homens \emph{adultos} \footnote{A historia contém exemplos de homens que, isolados, desde crianças, de toda á convivencia social, chegaram, só pela força da razão, ao conhecimento de uma causa productora do mundo.}.

20. O \textbf{atheismo} theoretico-positivo é impossivel. --- E' impossivel que o homem esteja \emph{séria} e \emph{sinceramente} convencido da tal- .sidade de uma coisa, attestada pela voz da natureza. Ora a existencia de \textsc{Deus} é attestada pela voz da natureza; pois é admittida por todos os homens, em todos os tempos, em todos os logares, e só o que deriva da natureza se encontra em todos os homens, em todos os tempos, em todos os logares. Logo é impossivel que o homem esteja \emph{séria} e \emph{sinceramente} convencido de que não existe \textsc{Deus}. --- A experiencia confirma a these; pois os que se dizem \emph{atheus} são \emph{scepticos} ou \emph{pantheistas} \footnote{O nosso seculo foi chamado \emph{era do atheismo}. E' certo que alguns infelizes empregam esforços para destruir, se fosse possivel, a propria idea de \textsc{Deus}. \emph{Fouerbach} diz: « Só o homem é o verdadeiro salvador! Só o homem é o nosso \textsc{DEus}, o nosso Juiz, o nosso Redemptor ». \emph{Rénan} escreve: « \textsc{Deus}, a Providencia, a immortalidade, não passam de palavras velhas e pesadas ». A loja maconica de Liege, em 1865, dizia: « Não só devemos collocar-nos acima de todas as religiões, mas acima de toda a crença num Deus ».\par Mas as pretendidas \emph{convicções} d'esses desvairados não passam do \emph{duvidas}... O espirito nunca poderá sinceramente convencer-se de que \textsc{Deus} não existe e, pelo menos com as suas duvidas, protestará sempre contra as estultas affirmações do coração. Já nos seus tempos \emph{Seneca} dizia: « Mentem os que dizem estar convencidos de que \textsc{Deus} não existe; pois, se ás vezes o dizem, todavia de noite, quando estão sós, começam a duvidar ».\par Estas \emph{duvidas} ácerca da existencia de \textsc{DEus} derivam da \emph{superficialidade} do pensamento e, sobretudo, da \emph{corrupção} da vontade.\par As \emph{duvidas} derivam da \emph{superficialidade} do pensamento. A este respeito, \emph{Caro} escreve: « Ha entre nós uma classe de espiritos scientificamente desorganisados, habeis para comprehender os phenomenos e as leis, nas diversas ordens da realidade, na natureza ou na historia, mas incapazes, por causa de uma inercia systematica, que se toma por uma força, de se elevar á concepção da \emph{causa}, que é a razão ultima das leis naturaes, assim como as leis são a razão dos factos » (\emph{L'idée de Dieu}, 1818, pag. 42).\par Mas as \emph{duvidas} derivam principalmente da \emph{corrupção} da vontade. Não é na \emph{intelligencia}, que o atheu nega a existencia de \textsc{DEus} ou duvida d'ella, mas é no \emph{coração}. O homem, cuja vida é pura e ordenada, nem sequer duvida da existencia de \textsc{Deus}. \emph{Rousseau} dizia; « Conservai a vossa alma num estado de desejar que exista \textsc{Deus}, e nunca duvidareis da sua existencia ». \emph{Za Bruyére} escreve: « Desejava ouvir a um homem sabio, moderado, casto /\emph{e} justo, que \textsc{Deus} não existe ou que a nossa alma não é immortal; elle á fallaria sem interesse. Mas esse homem não se encontra s.}.

21. \textbf{Consequencias} do \textbf{atheismo}. --- O \emph{atheismo} é a destruição da \emph{sciencia}, da \emph{sociedade} e da \emph{natureza humana}.

\pg{442} a) E' a destruição da \emph{sciencia}. Na verdade, a sciencia é o conhecimento das \emph{causas} dos seres. Ora, se \textsc{Deus} não existisse, os problemas mais fundamentaes seriam inexplieaveis. Porquanto, --- na ordem \emph{physica}, seria inexplicavel a origem do mundo, da vida, da alma humana, a uniformidade e constancia da ordem e das leis physicas; --- na ordem \emph{logica}, faltaria a base da \emph{possibilidade} das coisas, da \emph{eternidade} e \emph{necessidade} das essencias, da \emph{immutabilidade} da verdade; --- na ordem \emph{moral}, seria inexplieavel a existencia da \emph{lei moral}, a virtude não teria o premio, nem o vicio o castigo, e o bem não differiria do mal.

b) E' a destruição da \emph{sociedade}. Porquanto, uma sociedade não póde subsistir, sem que haja fidelidade e respeito nos subditos, desinteresse e justiça no principe, amor e concordia em todos. Mas uma sociedade sem \textsc{Deus} é o reino da anarchia, tendo por direito a força material, e por lei o despotismo \footnote{\emph{Voltaire} diz (\emph{Dict}. \emph{phiL}, a. \emph{athée}): « Eu não quereria ser subdito de um principe atheu, que, se encontrasse o seu interesse em me esmigalhar, com certeza esmigalhar-me-hia. Eu não quereria, se fosse soberano, ter cortezaos atheus, pois, se tivessem interesse em me enyenenar, deveria tomar, todos os dias, contra-venenos. E', pois, absolutamente necessario, para os principes e para os povos, que a idea de um Ente supremo, creador, governador, remunerador, seja gravada profundamente em todos os espiritos 7.}.

c) E' a destruição da \emph{natureza humana}. Com effeito, se \textsc{Deus} não existisse, a nossa alma não seria immortal, e assim o homem seria a mais desgraçada das creaturas; porque, depois de ter soffrido tantas dores, depois de ter suspirado tanto pela felicidade sem nunca a encontrar, o \emph{nada} seria a recompensa dos seus soffrimentos e o descanco das suas tendencias \footnote{Estas consequencias do \emph{atheismo} são tambem provas, ainda que \emph{indirectas}, da nossa these. Comtudo, \emph{Stuart Mill} diz que não temos prova alguma, nem \emph{seriamente provavel}, para demonstrar a existencia de \textsc{Deus}!\par --- e s+e---}.

/

\pg{443} \textbf{ESSENCIA} E \textbf{ATTRIBUTOS DE DEUS}

\textsc{SUuMMARIO}: --- Essencia de \textsc{Deus}, seu constitutivo e attributos. --- Attributos

de \textsc{Deus} negativos. --- Attributos de \textsc{Deus} positivos.

e

%%>>> art-39

\section*{Artigo I.\\ Essencia de Deus, seu constitutivo e attributos.}
\addcontentsline{toc}{section}{Artigo I. --- Essencia de Deus, seu constitutivo e attributos.}

22. \textbf{Essencia de Deus}. --- \emph{HEssencia} de \textsc{Deus} é o que faz com que Elle seja \textsc{Deus}, e se distinga de todos os outros seres.

23. \textbf{Essencia physica} e metaphysica de \textbf{Deus}. --- A essencia \emph{physica} de \textsc{Deus} é constituida por todas as perfeigoes, de que a Divindade é dotada. A \emph{metaphysica} é a perfeição fundamental, que \emph{se concebe} como raiz e fonte de todas as outras perfeições divinas, e pela qual \textsc{Deus} se distingue das creaturas. --- \emph{A} essencia \emph{physica} representa \textsc{Deus} como é em si mesmo; a \emph{metaphysica} exprime \textsc{Deus} como é em relação ao nosso modo de entender \footnote{Lembramos que a essencia --- é \emph{physica}, quando os seus elementos constitutivos são \emph{reaes} e \emph{concretos} (assim o \emph{corpo organico} e á \emph{alma racional} constituem a essencia \emph{physica} do homem), --- é \emph{metaphysica}, quando os seus elementos constitutivos são \emph{logicos} e \emph{abstractos}, denotando o \emph{genero} e a \emph{differença} (assim a \emph{animalidade} e a \emph{racionalidade} constituem a essencia \emph{metaphysica} do homem). --- Da essencia brotam as \emph{propriedades}. Estas distinguem-se \emph{realmente} da essencia \emph{pliysica}, porque não a constituem mas resultam d'ella (assim a intelligencia e a vontade são realmente distinctas da essencia \emph{physica} da nossa alma), --- mas podem não distinguir-se \emph{realmente} da essencia \emph{metaphysica}, porque, ainda que não possam considerar-se como elementos \emph{constitutivos} da essencia, comtudo podem identificar-se com esses elementos (assim a \emph{immortalidade}, embora não possa considerar-se como um elemento constitutivo da essencia \emph{metaphysica} da nossa alma, comtudo não se distingue \emph{realmente} da mesma essencia). E' claro que não podemos fallar da essencia de \textsc{Deus}, como fallamos da essencia das creaturas. Em Deus, a essencia \emph{physica} significa e abrange toda a perfeição, que compete ao Ente infinito, sem sombra de imperfeição, --- e a essencia \emph{metaphysica} exelue toda e,qualquer composição. Por isso o \emph{constitutivo} da essencia \emph{metaphysica} de Deus não é uma coisa realmente distincta das outras propriedades, e a sua relação com estas deve, mais ou menos, conceber-se como se concebe a relação do ente, considerado na sua razão communissima, com as suas propriedades \emph{íranscendentaes}.\par Dizemos que a essencia \emph{metaphysica} exprime \textsc{Deus} como Elle é em \emph{relação} ao \emph{nosso modo de entender}. Por: quanto, considerando os seres creados,\pgn{444}encontramos em cada um d'elles uma perfeição, que é a primeira a ser conheocida, que constitue o mesmo ser numa determinada especie, distinguindo-o de todos os outros, e que é a raiz e a fonte de todas as outras perfeições. E, como subimos ao conhecimento de \textsc{Deus} por meio das creaturas, procuramos encontrar em \textsc{DEus} essa perfeição \emph{primeira}, \emph{caracteristica e radical}; embora na essencia de \textsc{DEus}, como é em si mesma, não se encontre attributo algum, que seja anterior aos outros attributos e de que estes fluam. --- Este modo de entender a essencia de \textsc{Deus} é o resultado da limitação da nossa intelligencia; não podemos entender o infinito senão por partes, e, ainda assim, com muita imperfeição. Todavia, a propria intelligencia emenda ou explica o seu modo de ver; pois que, depois de ter considerado um attributo divino, como fundamento dos outros, prova que, em \textsc{Deus}, não só não ha distincção alguma entre a essencia e os attributos, mas nem sequer entre os mesmos attributos.}.

\pg{444} 24. A \textbf{essencia physica de Deus consiste na} infinidade.

a) Em \textsc{Deus}, o ser (ou a existencia) não é recebido nem subsiste na essencia, mas é \emph{subsistente por si}; visto que o ser de \textsc{Deus} não se distingue da sua essencia, e o que não se distingue de uma coisa, não pode ser recebido nem subsistir nella. Ora o ente, eujo ser não é recebido nem subsiste na essencia, é necessariamente infinito; porque o ser é limitado

\emph{Z} ou finito só porque e quando é recebido numa certa e determinada essencia, Logo Deus é infinito, isto é, a essencia \emph{physica} de \textsc{Deus} consiste na \emph{infinidade} \footnote{Desenvolvamos este argumento. Uma substancia, como dissemos na \emph{Ontologia}, pode \emph{subsistir}, ou \emph{existir em si}, por dois modos: ou porque é a existencia, ou porque \emph{tem} a existencia. No primeiro caso, a existencia (ou o ser) não é recebida na essencia, ou natureza, como na sua potencia proporcionada, e por isso é acto puro e subsistente por si; --- no segundo, a existencia (ou o ser) é recebida na essencia, como acto e principio de subsistencia, e por isso não é acto puro, nem subsiste por si, mas subsiste na essencia, como na sua potencia. Ora, em Deus, a essencia e o ser, ou a existencia, não se distinguem \emph{realmente}, porque o ser não é accidental, mas é essencial á essencia divina; porque esta é \emph{essencialmente} existente. Se o ser não se distingue \emph{realmente} da essencia divina, Deus, assim como é a sua essencia, isto é, a sua Divindade, assim tambem é o seu proprio ser, ou a sua propria existencia. Mas o ser, que não se distingue da essencia, não é recebido na propria essencia; porque o ser, que é recebido na essencia, deve necessariamente distinguir-se d'esta, como o acto se distingue necessariamente da potencia (assim a côr, recebida numa parede, distingue-se da propria parede). O ser, que não é recebido numa essencia, não é limitado ou determinado a uma certa especie; visto que o ser emtanto se limita ou se determina a uma certa especie, emquanto é recebido numa certa potencia ou natureza (assim o ser \emph{umano} é limitado ou determinado a uma certa especie, e não a uma outra, porque é recebido na \emph{essencia} ou natureza \emph{humana}). Portanto, o ser, que não é recebido numa essencia, não é limitado ou cireumscripto, mas é illimitado ou infinito. E, como o \emph{ser}, por ser \emph{acto},\pgn{445}indica e é \emph{perfeição}, segue-se que o ser subsistente por si é a perfeição subsistente, illimitada, infinita; de modo que não ha nem pode conceber-se entitade, perfeição, attributo, que esse ser não possua, assim como a alvura, se fosse subsistente em si, havia de conter tudo o que tem razão de alvor. Deus, portanto, sendo o ser subsistente por si, é infinito na sua perfeição, E' o raciocinio do Doutor Angelico, quando diz: « Deus est ipsum esse per se subsistens; ex quo oportet quod totam perfectionem essendi in se contineat. Manifestum est eonim quod si aliquod calidum non habeat totam perfectionem ealidi, hoc ideo est, quia calor non participatur secundum perfectam rationem; sed, si calor esset per se subsistens, non posset ei aliquid deesse de virtute caloris. Unde cum Deus sit ipsum esse subsistens, nihil de perfectione essendi potest ei deesse. Omnium autem perfectiones pertinent ad perfectionem essendi; secundum hoc enim aliqua perfecta sunt quod aliquo modo esse habeant. Unde sequitur quod nullius rei perfectio Deo desit » (\emph{Sum}. \emph{TAeolog}., p. I, q. 4, a. 2. --- Mais adiante accrescenta: « Cum esse divinum non sit esse receptum in aliquo; sed ipse sit suum esse subsistens, manifestum est quod ipse Deus sit infinitus et perfectus » (\emph{Jb}., q. 7, a. 1) E noutro logar: « Deus est infinitus: quod sic videri potest. Esse enim hominis terminatum est ad hominis speciem, quia est receptum in natura speciei humanae, et simile est de esse equi, vel cujuslibet creaturae. Esse autem Dei, quum non sit in aliquo receptum, sed sit esse purum, non limitatur ad aliquem modum perfectionis essendi, sed totum esse in se habet » (\emph{De Pot}., q. 1, a. 2).}. ,

\pg{445} b) A essencia \emph{physica} de \textsc{Deus} é constituida por todas as perfeições, de que a Divindade é dotada. Ora as perfeicoes, de que a Divindade é dotada, são \emph{infinitas} ; visto que \textsc{Deus}, sendo \emph{immutavel} e \emph{necessario}, deve possuir todas as possiveis perfeições; pois, se Lhe faltasse alguma, seria, com relação a esta, \emph{mutavel} e \emph{contingente}, \emph{e} por isso dependeria de uma causa. Logo á essencia \emph{physica} de \textsc{Deus} consiste na \emph{infinidade}.

c) Se as perfeicoes, que constituem a essencia \emph{physica} de \textsc{Deus}, fossem \emph{finitas} ou \emph{limitadas}, essa limitação deveria provir --- ou de uma causa \emph{extrinseca}, --- ou da indole da propria essencia divina, em que as perfeicoes são recebidas. Ora não pode provir --- nem de uma causa \emph{eztrinseca}, porque Dus, sendo a Causa \emph{primeira}, não depende de outra causa, --- nem da indole da propria essencia divina, porque esta, sendo perfeita por \emph{si mesma}, e não por \emph{participação}, não pode deixar de possuir \emph{a} plenitude da perfeição. Logo as perfeições, que constituem a essencia \emph{physica} de \textsc{Deus}, são \emph{infinitas}. Logo a essencia \emph{physica} de \textsc{Deus} consiste na \emph{infinidade} \footnote{A \emph{infinidade} de \textsc{Deus} \emph{indirectamente} demonstra-se pelo facto de se não encontrar nº Elle nenhuma das causas, que limitam as creaturas. Na verdade, as creaturas, que são compostas de \emph{materia} e \emph{forma}, são limitadas, porque á \emph{forma} limita á \emph{materia}, quando a determina a constituir uma certa especie e não outra, --- e a \emph{materia} limita a \emph{forma}, tornando-a \emph{sin}-,\pgn{446}\emph{gular} e \emph{individua}. Ora \textsc{Deus} nao pode ser limitado d'este modo: porque, como provaremos d'aqui a pouco, não é composto da \emph{materia} e \emph{forma}. --- As creaturas todas pertencem a um certo \emph{genero} e a uma certa \emph{especie}, e por isso são limitadas, porque não possuem as perfeições proprias dos ouires generos e das outras, especies. Ora \textsc{Deus}, sendo o ser \emph{necessario}, \emph{immutavel} e \emph{independente}, transcende todos os generos e todas as especies. --- As creaturas todas são compostas de \emph{essencia} e de \emph{ser}, ou \emph{existencia}, e por isso são limitadas, porque o \emph{ser}, sendo recebido numa essencia \emph{finita}, como é a essencia da creatura, não pode deixar de ser \emph{finito}. Ora, como dissemos, o ser divino, \emph{existindo} por assim o exigir a sua \emph{essencia}, subsiste por \emph{si} e não é recebido numa essencia, que o limite. --- As creaturas são limitadas pelo \emph{espaço} e pelo \emph{tempo}. \textsc{DEus}, por ser \emph{immenso}, não é limitado pelo \emph{espaço}, e, por ser \emph{eterno}, não é determinado pelo \emph{tempo}. Logo \textsc{Deus} é \emph{infinito}. --- S. Thomaz prova a mesma verdade, dizendo que Deus é \emph{infinito}; porque a razão da limitação num ente deriva ou da \emph{causa} do proprio ente ou da \emph{essencia}, em que o ser é recebido, mas Deus não tem causa e o seu ser não é recebido na essencia, mas, como provámos, é subsistente por si. Eis as suas palavras: « Ipsum esse, absolute consideratum, infinitum est... Si igitur alieujus esse sit finitum, oportet quod limitetur esse illud per aliquid aliud, quod sit aliqualiter causa illius esse, vel \emph{receptivum} ejus. Sed esse divini non potest esse aliqua causa, quia ipse necesse est esse per seipsum: nec esse ejus est receptivum, quum ipse sit suum esse, Ergo esse suum est infinitum, et ipse infinitus » (C. \emph{Gent}., lib. I, c. 48).\par Os erros ácerca da \emph{infinidade} de \textsc{Deus} podem reduzir-se a tres.\par a) Uns escriptores, não entendendo por \emph{infinito} senão a \emph{quantidade}, disseram que \textsc{Deus} não é infinito, porque uma \emph{quantidade infinita} repugna.\par Esta opinião refuta-se, negando o \emph{supposto}. A \emph{infinidade}, que se attribue a \textsc{Deus}, não é de \emph{quantidade}, mas sim de \emph{perfeição}.\par \emph{b}) Outros, como \emph{Locke}, \emph{Hobbes}, \emph{Robinet}, dizem que Deus é \emph{infinito}, não \emph{absoluta} mas \emph{relativamente}; pois faltam-lhe muitas perfeições.\par Esta opinião é absurda; porque \textsc{DEus}, como provámos, é \emph{essencialmente} infinito nas suas adoraveis perfeições. j\par c) Outros, como \emph{Hegel}, seguido por \emph{Vacherot}, \emph{Renan} e \emph{Taine}, não fazendo \textbar{} distincção entre a idea meramente \emph{subjectiva} do ente em geral e o Ente necessario, nem entre o conceito de \emph{infinidade} e o Infinito real, dizem que \textsc{DEus} não é um ente \emph{real} nem \emph{pessoal}, mas é a \emph{categoria do ideal}.\par Como é evidente, estas affirmações encerram o \emph{pantheismo}. Um \textsc{Deus} ideal, formado por nós, não é \textsc{Deus}. --- Não nos demoraremos na refutação do \emph{systema} de \emph{Hegel}; pois o refutámos na \emph{Cosmologia geral}, tratando do \emph{pantheismo}. Aqui apenas accrescentaremos que as theorias de \emph{Hegel}, servilmente copiadas ou imitadas por muitos escriptores, são tidas, na propria Allemanha, naquella conta que merecem. \emph{Schopenhauer} diz --- que \emph{Hegel} u é uma \emph{cabeça mediocre},\pgn{447}\emph{que por todos os meios quiz passar por um grande philosopho} », --- e que a philosophia de \emph{Hegel}, como a de \emph{Fichte} e de \emph{Schelling}, \emph{é} «\emph{um} minimum \emph{de pensamento diluido in quinhentas paginas de phraseologia nauseabunda} n...\par Davus, sendo infinito nas suas perfeições, deve possuir todas as perfeições, espalhadas nas creaturas; pois a causa \emph{efficiente} deve possuir toda a perfeição, que se encontra no effeito. --- E' o processo da \emph{causalidade}.\par Mas esta posse é diversa, segundo a natureza das perfeições. Por quanto as perfeições podem ser \emph{absolutas} ou \emph{miwtas}. +\par \emph{Absolutas} são as perfeições, que no seu conceito não encerram imperfeição alguma, nem se oppõem a perfeições mais elevadas, como são a \emph{vida}, a \emph{belleza}, a \emph{bondade}, a \emph{sabedoria}, etc. Estas perfeições attribuem-se a \textsc{Deus} no proprio conceito, que exprimem, mas por um modo mais \emph{eminente}, porque em \textsc{Deus} se encontram por um modo infinitamente mais perfeito e sublime do modo, por que se encontram nas creaturas, pois \textsc{DEus} não é \emph{bom}, \emph{santo}, \emph{amavel}, como é o homem. Por isso os nomes, com que exprimimos as perfeições \emph{absolutas} de \textsc{Deus}, não são \emph{univocos} nem \emph{equivocos}, mas \emph{analogos}. (\emph{Sum}. \emph{Th}., p. I, q. 19, a. 5). --- E' o processo da \emph{eminencia}.\par \emph{Mixtas} são as perfeições, que, embora no seu conceito exprimam uma entidade positiva, todavia encerram alguma imperfeição e se oppõem a uma perfeição mais elevada; tal é a \emph{corporeidade}, que, embora signifique uma coisa positiva, todavia encerra alguma imperfeição (pois o \emph{corpo}, por ser \emph{composto}, \emph{é dependente}, \emph{potencial}, \emph{finito}, ete.) e se oppõe a uma perfeição mais elevada, que é a \emph{espiritualidade} (pois o \emph{espirito} é mais nobre que o \emph{corpo}). Estas perfeicgóes não se attribuem a \textsc{Deus} no seu proprio conceito, mas só \emph{eminentemente}, emquanto attribuimos a \textsc{Deus} toda a entidade, que essas perfeições encerram, e afastamos d'Elle toda a imperfeição, de que as perfeições estão cercadas. --- E' o processo da \emph{exclusão}.\par D'ahi se vê que o processo da \emph{exclusão} é, para nós, o mais natural, para nos elevarmos ao conhecimento de \textsc{Deus}. Por quanto, devemos afastar de \textsc{Deus} --- não só as perfeições creadas, que no seu proprio conceito indicam alguma imperfeição, --- mas tambem todo o limite de que estão cercadas as perfeições creadas, que no seu conceito não encerram alguma imperfeição. D'este modo mais sabemos o que \textsc{Deus} não é, do que o que Elle é.\par Todavia, este processo, --- além de manifestar.a perfeição infinita de \textsc{DEus}, pois o Ente, de que se afasta toda a imperfeição, é perfeitissimo, --- subminisbra-nos um verdadeiro conhecimento de Dus, pois tanto melhor conheceremos um ente, quanto mais o distinguirmos de todos os outros seres. --- \textsc{DEUS}, pois, não é absolutamente \emph{i5cognoscivel}, como pretende \emph{H}. \emph{Spencer}.}.

25. \emph{A} \textbf{essencia} metaphysica de \textbf{Deus} consiste na subsistencia do proprio ser. --- A essencia \emph{metaphysica} de \textsc{Deus} consiste naquella perfeição, --- que se concebe como raiz e fonte de todas as outras perfeicoes divinas, e pela qual \textsc{Deus} se distingue das creaturas. Ora a \emph{subsistencia do proprio ser} é a perfeição \pg{446} \pg{447}  --- que se concebe como raiz e fonte de todas as outras perfeições divinas, porque, como provámos, Deus possue todas as perfeições, possiveis e imaginaveis, pelo facto de ser o \emph{proprio ser subsistente}, --- e pela qual se distingue das creaturas, visto que Deus é o Ente infinito, porque é o proprio ser subsistente, e as creaturas são finitas, porque \emph{têm} o ser recebido numa determinada essencia. Logo a essencia \emph{metaphysica} de \textsc{Deus} consiste na \emph{subsistencia do proprio ser} \footnote{A conclusão é certa, e é uma consequencia dos principios expostos no numero anterior. O ser divino; pelo facto de não ser recebido em nenhuma\pgn{448}essencia ou potencia, não tem limites, encerra uma perfeição absoluta sem sombra de imperfeição. Repetimos con S. Thomaz: « \emph{Ew hoc quod Deus est ipsum esse subsistens}, \emph{oportet quod totam}. \emph{perfectionem essendi in se contineat} » (\emph{Sum.-Th}., p. I, q. 4, a. 2). Sendo subsistente, o ser divino distingue-se das creaturas, que não \emph{são} o ser, mas \emph{têm} o ser, recebido numa determinada potencia ou essencia e por isso limitado a uma corta especie. Como o ser ó \emph{acto}, o ser que não é recebido numa potencia, mas que é acto \emph{puro}, é necessariamente infinito ou illiimitado; ao passo que o ser, que é recebido numa potencia, é limitado por ella. Por isso, não só pode dizer-se que o \emph{constitutivo} da essencia divina consiste em que \textsc{Deus} \emph{do ser subsistente por si}, mas tambem pode dizer-se que esse \emph{constitutivo} consiste om que \textsc{Deus} é \emph{acto purissimo}, \emph{não recebido em nenhuma potencia}. As expressões são synonimas. --- Concluimos --- que a razão formal e inürinseca, pela qual Deus é unico, eterno, infinito, é porque Elle é o proprio ser subsistente, é acto purissimo, cuja essencia é constituida pelo proprio ser, --- e que, pelo contrario, a razão, pela qual as creaturas são multiplicaveis, imperfeitas, finitas, é porque têm o ser recebido na essencia, e por isso são compostas de potencia e de acto, de essencia e de ser.\par Sendo Deus o proprio ser subsistente, o nome mais proprio de Deus é \textsc{JEHOYAH}, que significa: Qui EST? --- \textsc{AQUELLE QUE} É. Por quanto, os nomes significam a forma ou essencia das coisas. Ora, em Deus, a essencia identifica-se com o ser, isto é, o ser é á sua propria essencia, de modo que á pergunta: \emph{o que é Deus}? podemos responder: \emph{Deus é o proprio ser} --- \emph{Deus est ipsum esse}. Logo o nome \emph{Qui est} é muito proprio de Deus. Com effeito,. de nenhuma creatura podemos dizer que ella é o proprio ser; porque o \emph{ser}, nas \emph{creaturas}, não é a \emph{essencia}, nem é um \emph{predicado essencial}, mas é distincto da essencia, embora seja recebido nella; e por isso os nomes, que significam as coisas creadas, exprimem a essencia, e não e ser, assim \emph{tomem} exprime a essencia de um individuo, isto é, \emph{a humanidade}, mas não exprime o \emph{ser} ou a \emph{existencia} d'esse individuo, que poderia não existir. Escutemos S. Thomaz: « Cum in omni quod est, sit considerare \emph{quidditatem suam}, per quam subsistit in natura determinata, et \emph{esse suum}, per quod dicitur de eo quod est in actu, hoc nomen \emph{res} imponitur rei a quidditate sua, hoc nomen \emph{qui est} vel \emph{ens} imponitur ab ipso actu essendi. Cum autem ita sit quod in qualibet re creata, essentia sua differat a suo esse, res illa proprie denominatur a quidditate sua et non ab actu essendi, sicut homo sb humanitate. In Deo autem ipsum esse suum est sua quidditas; et ideo nomen quod sumitur ab esse, proprie nominat ipsum, et est proprium nomen ejus; sicut proprium nomen hominis quod sumitur a quidditate sua » (I \emph{Sent}, Dist. 8, q. 1, a. 1).\par E não pode deixar de ser assim. Quando digo: \emph{Pedro é homem}, attribue 2 esse individuo duas coisas: o \emph{ser}, ou a \emph{existencia}, e a \emph{natureza humana}, isto é, attribue-lhe o ser limitado e determinado a uma certa especie; mas quando digo: \emph{Deus é aquelle que} é, quero affrmar que Deus não é este ou aquelle ente, mas que ó o ser, todo o ser, o proprio ser, o ser infinito, de modo que tudo o que tem razão de ser deve encontrar-se e se encontra\pgn{449}em Deus. Eis o sentido do nome \emph{qui est}, que o proprio Deus manifestou a Moyses. Ló-se no \emph{Exodo}; « Ait Moyses ad Deum: Ecce ego vadam ad filios Israel, et dicam eis: Deus patrum vestrorum misit me ad vos. Si dixérint mihi: quod est nomen ejus? quid dicam eis? Dixit Deus ad Moysen: Eco \emph{sum qui sux}, Ait: Sic dices filiis. Israel: Qui EST misit me ad vos » (Cap. III, vv. 183, 14). --- E' este um nome mais universal e menos determinado que os outros; mas por isso mesmo é o mais proprio. Sendo a essencia divina infinitamente superior á intelligencia creada, o nome mais proprio de Deus é o que melhor traduz essa transcendencia e indefinibilidade, e exclue as imperfeitas determinações das creaturas. Diz S. Thomaz: « Quia divina comprehendi a nobis non possunt, conveuientius a nobis significantur per nomina.communia, quae indefinite aliquid significant, quam per nomina specialia, quae definite rei species exprimunt; unde hoc nomen \emph{qui est}, quod secundum Damascenum significat substantiae pelagus infinitum, convenientissimum nomen dicitur esse » (\emph{De pot}., q. 10, a. 1 ad 9).\par Nem se diga que o nome \emph{qui est} não é o mais proprio de Deus, porque não o distingue das creaturas, pois Deus é e as creaturas tambem \emph{são}. Por quanto, ha uma grande differenca entre o \emph{ser} de Deus e o \emph{ser} das creaturas. O ser de Deus, por ser subsistente em si, identifica-se com a essencia, e por isso se torna individuo e distincto de todos os outros entes; ao passo que o ser das creaturas, por não ser subsistente em si, não se identifica com a essencia, mas-é recebido e subsiste na essencia, e por isso se torna individuo e distincto por causa da mesma essencia. Portanto o nome \emph{qui est} é exclusivamente proprio de Deus, porque só em Deus o ser ó a mesma essencia, e o nome, como dissemos, exprime a essencia. Oucamos de novo S. Thomaz: « Ipsum esse Dei distinguitur et individuatur a quolibet alio esse per hoc ipsum, quod est esse per se subsistens, et non adveniens alicui naturae, quae sit aliud ab ipso esse. Omne autem aliud esse, quod non est subsistens, oportet quod individuetur per naturam et substantiam, quae in tali esse subsistit. Et in eis verum est, quod esse hujus est aliud ab esse illius per hoc quod est alterius naturae » (\emph{De pot}., q. 7, a. 2 ad 5).\par Todos os philosophos concordam em que em \textsc{Deus}, considerado segundo o nosso modo de entender, existe um attributo constitutivo da essencia divina, mas divergem na determinação d'aquelle attributo.\par Ha quatro outras opiniões a este respeito.\par a) Uns dizem que o attributo \emph{constitutivo} da essencia divina é a somma \emph{actual} de todas as perfeições possiveis.\par \emph{b}) Outros, como \emph{Scoto}, sustentam que a essencia divina não é consti- Guida pela somma \emph{actual}, mas pela \emph{exigencia radical} de todas as perfeições.\par c) Outros, seguindo \emph{Billuart}, afirmam que o \emph{constitutivo} da essencia divina é a \emph{intelligencia}.\par \emph{d}) Outros finalmente ensinam que o constitutivo da essencia de \textsc{Deus} é a \emph{aseidade} (\emph{ens a se}), isto é, aquella propriedade que faz com que \textsc{Deus} exista \emph{por si}, sem receber de outra causa o ser que possue.\par 29 : Vor. II.}.

26. \textbf{Especies} -de \textbf{attributos} divinos, \textbf{derivados da essencia}. --- Os attributos divinos, derivados da essencia, dividem-se em \emph{quiescentes} e \emph{operativos}. --- Os \emph{quiescentes} são os que se referem 4 \pg{448} \pg{449}  \emph{essencia}, emquanto é uma \emph{entidade}. Os \emph{operativos} sio os que se referem á \emph{essencia}, emquanto é \emph{principio de operações}, e comprehendem as \emph{faculdades}, com as \emph{operações} respectivas.

\pg{450} Os attributos \emph{quiescentes} dividem-se em \emph{negativos} e \emph{positivos}. Os \emph{negativos} são os que \emph{negam} de \textsc{Deus} as imperfeições das creaturas. Os \emph{positivos} são os que \emph{affirmam} de \textsc{Deus} as perfeições creadas, mas isentas de todo o limite.

Neste artigo e no seguinte tratamos dos attributos \emph{qu?escentes}, \emph{negativos} e \emph{positivos} ; nos capitulos seguintes trataremos dos attributos \emph{operativos} \footnote{A. nossa intelligencia, não podendo, por ser finita, comprehender a Essencia divina, que excede, pela sua perfeição infinita, a capacidade de toda a creatura, procura exprimil-a, de algum modo, por muitos e differentes conceitos, A esta pluralidade e diversidade de conceitos corresponde, da parte de Deus, não uma pluralidade ou diversidade real de perfeigóes, mas uma plena e infinita perfeição, que equivale a todas as perfeições e que as contém de um modo eminente. --- A esta pluralidade e diversidade de conceitos, que exprimem a unica e simplicissima Essencia divina, correspondem varios nomes, que por isso não são \emph{synonymos}, porque, embora todos exprimam a mesma Essencia, comtudo não a exprimem todos sob o mesmo aspecto, mas um exprime-a sob um aspecto e outro sob outro aspecto, e sempre de um modo \emph{deficiente} e \emph{parcial} (Cf. \emph{Sum}. \emph{Th}., p. I, q. 18, a. 4).}.

Todas estas opiniões são inacceitaveis.

a) E' inaeceitavel á \emph{primeira}. Por quanto, a \emph{infinidade} das perfeições constitue a essencia \emph{physica} de \textsc{DEUS}, mas não a \emph{metaphysica}; e aqui procuramos o \emph{constitutivo} da essencia \emph{metaphysica}: de \textsc{Deus}.

\emph{b}) E' inacceitavel a \emph{segunda}. Com effeito, aquella \emph{existencia radical}, longe de constituir a essencia de \textsc{Deus}, suppõe-na já constituida. .

c) E' inacceitavel a \emph{terceira}. Na verdade, a \emph{intelligencia} não é á \emph{primeira} coisa, que conhecemos de \textsc{Deus}, nem o distingue de todos os seres creados; pois ha muitas creaturas dotadas de intelligencia. --- Além d'isso, não podemos deduzir da \emph{Intelligencia} muitos attributos de Drxus, como a \emph{eternidade}, a \emph{immensidade}, \emph{á immutabilidade}, ec.

\emph{d}) E' inacceitavel a \emph{quarta}. Porquanto, a \emph{aseidade}, embora distinga Deus das creaturas, comtudo não representa o intrinseco, proximo e formal principio, que se concebe como constitutivo da essencia divina, e pelo qual Deus se distingue de todas as outras coisas. Na verdade, a independencia (\emph{esse} a se), embora exclua a causa efficiente, todavia não se concebe como uma coisa \emph{intrinseca} ao proprio ente, nem como a raiz de todos os attributos divinos, porque podemos sempre perguntar: porque Deus é \emph{ens a se}? E a resposta ' deve ser esta: porque \textsc{DEus} é o ser subsistente por si, e o acto purissimo. Por isso, este attributo, embora não seja o primeiro na ordem do \emph{conhecimento}, é o primeiro na ordem da \emph{realidade}, porque é o que se concebe como uma coisa \emph{Znfrinseca} e como a raiz e a fonte de todos os outros attributos divinos. --- Notamos que alguns philosophos e theologos tomam estas duas expressões --- \emph{Deus est ens per se subsistens} e \emph{Deus est ens a se} --- como equivalentes, e, dizendo que o constitutivo da essencia divina é a \emph{aseidade}, entendem que é a \emph{subsistencia do ser divino}. Mas, rigorosamente fallando, ha uma differenca entre as duas expressões.

%%>>> art-40

\section*{Artigo II.\\ Attributos de Deus negativos.}
\addcontentsline{toc}{section}{Artigo II. --- Attributos de Deus negativos.}

\pg{451} 27. Attributos de \textbf{Deus} negativos. --- Attributos \emph{negativos} São, como dissemos, os que \emph{negam} de \textsc{Deus} as imperfeições proprias das creaturas. --- Os principaes são: a \emph{simpicidade}, a \emph{immutabilidade}, a \emph{eternidade}, \emph{á inmensidade} e á \emph{unidade} \footnote{Estes attributos são \emph{negativos} em quanto ao \emph{termo}, mas não em quanto á \emph{significação}, porque significam uma perfeição pura e simplicissima, de que a essencia de \textsc{Deus} é dotada. --- A \emph{infinidade} tambem colloca-se entre os attributos \emph{negativos}; mas d'ella fallámos no artigo precedente.}.

28. Simplicidade de \textbf{Deus}. --- A \emph{simplicidade} é o attributo, que exelue de \textsc{Deus} toda a \emph{composição} e \emph{divisibilidade} \footnote{\emph{Simplicidade}, em geral, é a \emph{negação de partes}; porque todos chamam \emph{simples} o ente, que não tem partes. A \emph{simplicidade} é uma especie de \emph{unidade}, á qual accrescenta a \emph{indivisibilidade}. \emph{Uno} é o ente, que é \emph{indiviso}. Mas o ente uno pode ter partes, e, pela divisão, tornar-se multiplice; por isso, se exclue a multiplicidade \emph{actual}, não exclue a multiplicidade \emph{potencial}. Pelo contrario, o \emph{simples}, por não ter partes, exclue a multiplicidade \emph{actual} e a \emph{potencial}. --- A \emph{simplicidade}, considerada em si, não representa nem uma perfeição, nem uma imperfeição. Ás vezes o ente \emph{simples} é mais perfeito que o \emph{composto}, ás vezes é menos perfeito. E' mais perfeito, quando representa um \emph{acto} ou uma \emph{forma} completa e subsistente, dotada de toda a perfeição, propria da sua especie; assim o Anjo é mais perfeito que o homem. E' menos perfeito, quando representa uma \emph{parte} de um todo; assim a alma humana é menos perfeita que o composto humano. Diz 8. Thomaz: « Simplicitas \emph{per se} non est causa nobilitatis, sed perfectio; unde ubi perfecta bonitas in uno simplici invenitur, simplex est nobilius, quam compositum: quando autem e contrario simplex est imperfectum, tunc compositum est nobilius, quam simplex, sicut homo nobilior est terra » (4 Sent. \emph{Dist}. 11, q. 2, a. 1 ad 1).\par A? \emph{simplicidade} oppõe-se diametralmente a \emph{composição}. --- Esta pode ser \emph{real}, ou \emph{de rusão}. E' \emph{real}, quando a coisa é por si composta de partes, independentemente da nossa reflexão. E' \emph{de razão}, quando a nossa intelligencia, pela sua limitação natural, exprime por varios e diversos conceitos uma coisa, que por si é simplicissima. A composição \emph{real} denota imperfeição na coisa, e por isso não pode admittir-se em Deus; a \emph{de razão} denota apenas imperfeição na nossa intelligencia, e é compativel com a simplicidade de Deus (Cf. \emph{C}. \emph{Gent}., 1. II, c. 14, --- \emph{De ver}., q. 5, a. 2 ad 2). --- \emph{A} composição \emph{real} subdivide-se em \emph{physica}, \emph{metaphysica} e \emph{logica}. A \emph{physica} é propria do ente, que consta de diversos elementos \emph{physicos}, como são a \emph{materia} e a \emph{forma}. \emph{A metaphysica} é propria do ente, que consta de partes \emph{metaphysicas}, como são a \emph{potencia} e o \emph{acto}, \emph{á essencia} e o \emph{supposto}, \emph{á essência} e á \emph{existencia}, \& \emph{essencia} e os \emph{attributos}, \& \emph{substancia} e o \emph{accidente}. \emph{A logica} é propria do ente, que consta de elementos \emph{logicos}, como são o \emph{genero} e a \emph{differença}. --- Ninguem se admire se reduzimos a composição \emph{logica á real}. A. composição \emph{logica}, se é \emph{de} razão emquanto ao seu conceito \emph{formal}, porque\pgn{452}os elementos são productos da intelligencia, é, comtudo, \emph{real} emquanto ao \emph{fundamento}, pois presuppõe uma composição \emph{real} no ente composto de \emph{genero} e de \emph{differenca}, visto que o \emph{genero} se funda nà \emph{potencia}, e a \emph{differenca}, no \emph{acto}. --- Provaremos que Deus é \emph{simplieissimo}, se provarmos que Elle exclue de si toda e qualquer especie de \emph{composição}.}.

\pg{452} 29. \textbf{Deus} é \textbf{simplicissimo}.

a) Todo o ente composto é \emph{finito}; porque resulta de partes, e estas são essencialmente finitas. Ora \textsc{Deus} é \emph{infinito}. Logo \textsc{Deus} não é \emph{composto}, mas \emph{simplicissimo}.

b) O ente \emph{composto} não só é \emph{posterior}, ao menos por \emph{natureza}, ás partes componentes, mas tambem depende de uma causa \emph{extrinseca}, que reune as mesmas partes componentes. Ora Dszus é o ente \emph{primeiro}, a causa independente de toda a causa. Logo \textsc{Deus} não é \emph{composto}, mas \emph{simplicissimo} \footnote{Cf. S. Thom., \emph{Sum Th}. p. I, q. 8, a. €; --- c. \emph{Gent}., 1. T, e. 16 e seg., --- \emph{De pot}. q. 7, a. 1. --- Depois de termos demonstrado em geral que em \textsc{Deus} não ha composição alguma, demonstraremos em particular que em \textsc{Deus} não ha composição \emph{physica}, nem \emph{metaphysica}, nem \emph{logica}.\par a) Em \textsc{Deus} não ha composição \emph{physica}. ---- Em Dus haveria composição \emph{physica}, se elle fosse \emph{corpo}, ou \emph{forma corporea}, ou \emph{composto} de \emph{materia} e \emph{forma}. Ora \textsc{Deus} não é \emph{corpo}, nem \emph{forma} de corpo, nem \emph{composto}.\par \emph{a}) \textsc{Deus} não é \emph{corpo}. O corpo não é a mais nobre entre as substancias. Porquanto, o corpo \emph{vivo} é mais nobre do que o \emph{não-vivo}. Ora o corpo \emph{vivo} é mais nobre do que o \emph{não-vivo}, não emquanto é \emph{corpo} (aliás todos os corpos seriam eguaes em \emph{nobresa}), mas emquanto é \emph{vivo}. Logo o principio, de que deriva a \emph{vida} para o corpo, é mais nobre do que o \emph{corpo}, como tal. Ora \textsc{Deus} é o ser nobilissimo. Logo \textsc{Deus} não é \emph{corpo} (\emph{Sum}. \emph{Th}. p. I, q. 8, a. 1).\par \emph{b}) \textsc{Deus} não é \emph{forma} nem \emph{materia} de algum corpo. --- Não é \emph{forma}. Porquanto, a \emph{forma corporea}, além de ser \emph{finita} e por isso \emph{creada} e \emph{contingente}, depende, pelo menos \emph{e«trinsecamente} (como a alma humana), da materia. Ora Deus é infinito, Causa primeira, Ente necessario e absolutamente independente. Logo não pode ser \emph{forma} de algum corpo. --- Muito menos pode ser \emph{materia}; pois esta é um ser \emph{potencial}, \emph{finito}, \emph{contingente}, e por isso dotada de propriedades contrarias ás de Deus (\emph{Sum}. \emph{Th}., p. T, q. 8 a. 8).\par c) \textsc{Deus} não é \emph{composto} de \emph{materia} e \emph{forma}. Com effeito, todo o \emph{composto}, além de ser \emph{finito}, \emph{é dependente} não só da causa \emph{efficiente}, que reune as partes, mas tambem da \emph{forma}, de que recebe as perfeigóes. Ora Deus não depende de nenhuma causa, por ser a Causa primeira, e possue todas as perfeições não por \emph{participação}, mas por \emph{essencia}. (\emph{Sum}. \emph{Tu}, ipSDq:9, 8: 2) --- E\textasciicircum{}, pois, absurdo o \emph{antAropomorphismo}, que attribue a \textsc{Deus} a \emph{forma humana}.\par D'isto resulta que \textsc{Deus} é \emph{purissimo espirito}. \emph{Espirito} é uma substancia independente da \emph{materia}. Ora a substancia divina exclue toda e qualquer mistura de \emph{materia} e de \emph{potencia}. Logo \textsc{Deus} é \emph{purissimo espirito}.\par Totavia, do facto de \textsc{Deus} ser \emph{purissimo espirito}, e não ter \emph{corpo}, não póle concluir-se (como concluem os \emph{pantheistas}) que Elle não é \emph{pessoa}. Por quanto, \emph{pessoa} é a \emph{substancia individua de natureza racional}. Ora \textsc{Deus} é \emph{substancia}, subsistente em si e por si, e esta \emph{substancia} é infinitamente \emph{intelligente}. Logo \textsc{Deus} é \emph{pessoa}. --- E não só é \emph{pessoa}, mas possue a \emph{personalidade}\pgn{453} no grau mais elevado, que póde imaginar-se; porque a \emph{substancia} de \textsc{Deus} é \emph{perfeitissima}, e a sua intelligencia é \emph{infinita}. --- A materia não é sempre elemento necessario da \emph{personalidade}; pelo contrario, sendo a \emph{razão} o elemento principal da \emph{pessoa}, é claro que, quanto mais a razão for independente da materia, tanto mais nobre será a \emph{pessoa}.\par b) Em \textsc{Deus} não ha composição \emph{metaphysica}. --- Na verdade, \textsc{DEUS} não é composto de \emph{potencia} e \emph{acto}, de \emph{essencia} e \emph{supposto}, de \emph{essencia} e \emph{existencia}, de \emph{essencia} e \emph{attributos}, de \emph{substancia} e \emph{accidentes}.\par 4) \textsc{Deus} não é composto de \emph{potencia} e \emph{acto}. Todo o ente, composto de \emph{potencia} e \emph{acto}, \emph{é mutavel}. Ora \textsc{Deus} é \emph{immutavel}. Logo não é composto de \emph{potencia} e \emph{acto}. --- Por isso \textsc{Deus} é \emph{acto purissimo}, contendo \emph{todas} as \emph{possiveis} perfeições num grau \emph{infinito}; porque, se não possuisse alguma perfeição ou fosse capaz de augmento, estaria em. \emph{potencia} para aquella perfeição ou augmento de perfeição, e por isso seria \emph{composto} de \emph{potencia} e \emph{acto}.\par \emph{b}) \textsc{Deus} não é composto de \emph{essencia} e de \emph{supposto}. Um ente é composto de \emph{essencia} e de \emph{supposto}, quando os elementos, que estão contidos neste, differem, na \emph{especie} ou no \emph{numero}, dos que estão contidos naquella; de modo que não se póde dizer que o \emph{supposto} é a \emph{essencia}. Assim é homem é composto de \emph{essencia} e de \emph{supposto}; porque, além de \emph{todos} os elementos \emph{especificos}, que estão contidos na \emph{humanidade} e que são a \emph{animalidade} e á \emph{racionalidade}, o \emph{supposto humano} contém outros elementos, que são as \emph{propriedades individuaes} e que não estão contidas no conceito \emph{especifico} de \emph{humanidade}; de modo que não podemos dizer que \emph{Pedro} é a \emph{humanidade}, mas só que \emph{Pedro} possue a \emph{humanidade}. Ora as \emph{propriedades individuaes} derivam da \emph{materia}. Logo nos \emph{espiritos puros}, em que não ha \emph{materia}, tudo o que está no \emph{supposto}, encontra-se na \emph{essencia}; de modo que esta póde attribuir-se áquelle por \emph{identidade}, dizendo que o \emph{supposto} é a propria \emph{essencia}. \textsc{DEus} é \emph{purissimo espirito}. Logo, em \textsc{DEus}, não ha distineccüo entre \emph{supposto} e \emph{essencia}, e por isso \textsc{DEUS} não é composto de \emph{supposto} e \emph{essencia} ; de modo que podemos dizer que \textsc{Deus} é a \emph{Divindade} (\emph{Sum}. \emph{Th}., p. I, q. 9, a. 8; --- c. \emph{Gent}. 1. I, c. 21).\par \emph{c}) \textsc{Deus} não é composto de \emph{essencia} e de \emph{existencia}. Um ente, composto de \emph{essencia} e de \emph{existencia}, isto é, cuja \emph{essencia} se distingue \emph{realmente} da \emph{evistencia}, não existe \emph{essencialmente}, por assim o exigir a sua essencia, mas recebeu o \emph{ser} de outra causa. Ora \textsc{Deus} existe \emph{essencialmente}, não recebendo de outra causa o ser que possue. Logo \textsc{Deus} não é composto de \emph{essencia} e de \emph{existencia}. --- Além disso, a \emph{essencia} refere-se \& \emph{existencia}, como a \emph{potencia} ao \emph{acto}. Ora \textsc{Deus} não é composto de \emph{potencia} e \emph{acto} (\emph{Sum}. \emph{Th}., p. I, q. 8, 8. 4; --- c. \emph{Gentes}, l. T, c. 22).\par \emph{d}) \textsc{Deus} não é composto de \emph{essencia} e \emph{attributos}. A. Essencia divina, sendo o proprio ser, possue a plenitude da perfeição, e por isso não ha. attributo algum, que Ella não contenha. Ora a \emph{essencia} de \textsc{Deus}, se se distinguisse \emph{realmente} dos seus attributos, aperfeicoar-se-hia pelos mesmos attributos, e já não possuiria a plenitude da perfeição. Logo em DEus não ha distincção \emph{real} entre a \emph{essencia} e os \emph{attributos}. Logo \textsc{Deus} não é composto de \emph{essencia} e de \emph{attributos}. --- Alem d'isso, quando os attributos são \emph{realmente} distinctos da essencia, esta refere-se iquelles, como a potencia ao acto. Ora a Essencia divina, por ser o proprio ser subsistente, é acto purissimo, sem sombra de\pgn{454}potencialidade. Diz S. Thomaz: « Perfectio divini \emph{esse} non attenditur secundum aliquod \emph{additum} supra ipsum; sed quia ipsum secundum \emph{seipsum} perfectum est. Bonitas igitur Dei non est aliquid additum suae substantiae, \emph{sed} sua \emph{substantia est sua bonitas} » (\emph{C Gentes}, 1. T, c. 88) --- Se, em \textsc{Deus}, os \&ttributos são uma e a mesma coisa com a essencia, é claro que, assim como dizemos que \textsc{Deus} é a \emph{Divindade}, assim tambem devemos dizer que Elle é á \emph{Verdade}, a \emph{Belleza}, \& \emph{Bondade}, a \emph{Misericordia}, etc.\par Se, em Deus, os attributos não se distinguem \emph{realmente} da \emph{essencia}, tambem não se distinguem \emph{realmente} entre Sj, e por isso, na realidade, um identifica-se com outro. E' um corollario da these precedente. Se cada attributo divino se identifica com a Essencia, é evidente que um não se distingue realmente do outro, porque uma coisa não se distingue realmente de si mesma. A Essencia divina encerra, na sua absoluta simplicidade, todas as perfeições possiveis; e só a imperfeição e limitação da nossa intelligencia considera como derivadas da essencia e distinctas entre si as propriedades, que, na realidade, são uma e a mesma coisa. --- E claro que isto se refere aos attributos \emph{absolutos} de Deus, e não aos que se dizem \emph{relativos ad intra} e que são oppostos entre si, como são o \emph{Paternidade}, a \emph{Filiação} e a \emph{Processão} ; porque estes \textbf{atgeibutos} são realmente distinctos e oppostos e constituem as tres Pessoas da Santissima Trindade --- Panne, Figo e EsPmrro Sawro. Diz S. Thomaz: « Omnia in Deo sunt unum \emph{re}, praeter ingenerationem, generationem et processionem, quae constituunt personas re distinctas » (T Sent. \emph{Dist}. 2, q. 1, a. 8 ad 6.\par Mas, se os attributos divinos não se distinguem \emph{realmente} um do outro, porque cada um se identifica com a Essencia divina, comtudo distinguem-se \emph{logicamente}, porque o conceito de um não se identifica com o conceito de outro; assim o conceito de \emph{justiça} não é o conceito de \emph{misericordia}. E esta distineção \emph{logica} não é arbitraria, mas tem o seu \emph{fundamento real} na propria Essencia divina. Porquanto, a cada um d'esses diversos conceitos, relativos a attributos diversos, corresponde, na sua plenitude simplicissima, a unica Essencia divina. Se assim não fosse, isto é, se de nenhum modo existisse em Deus o fundamento d'essa distinecüo, os nomes \emph{absolutos} e \emph{positivos}, que exprimem as divinas perfeições, seriam \emph{synonymos}; o que é falso; pois esses nomes, embora exprimam todos a Essencia divina, exprimem-na comtudo sob differentes aspectos; assim a \emph{Justiça} exprime a Essencia divina emquanto dá a cada um o que merece,a \emph{misericordia} exprime a mesma Essencia divina emquanto tem dó dos nossos males. E' a doutrina de S. Thomaz. Diz assim: « In Deo est sapientia, bonitas, et huj usmodi, quorum quodlibet est ipsa divina essentia, et ita sunt unum re. Et quia unumquodque eorum est in Deo secundum verissimam sui rationem, et ratio sapientiae non est ratio bonitatis in quantum hujusmodi, relinquitur quod sunt diversa ratione non tantum ex parte ipsius ratiocinantis, sed ex proprietate ipsius rei; et inde est quod Ipse non est causa rerum omnino aequivoca, cum secundum formam suam producat effectus similes non univoce sed analogice » (I \emph{Sent}. Dist. 2, q. 1, a. 2; cf. \emph{Sum}. \emph{Th}, P. I, q. 13, a. 2). E mais adeante: « Pluralitas nominum non tantum est ex parte intellectus nostri, formantis diversas conceptiones de Deo; sed etiam ex parte ipsius Dei, inquantum scilicet. est\pgn{455}aliquid in Deo correspondens omnibus istis conceptionibus, scilicet plena et omnimoda ejus perfectio » (I \emph{Sent}. Dist. 2 q. 1, \&. 8).\par €) \textsc{Deus} não é composto de \emph{substancia} e \emph{accidentes}. O \emph{accidente} refere-se á \emph{substancia}, como o \emph{acto á potencia}. Ora \textsc{DEus} é \emph{acto purissimo}. --- Além d'isso, o \emph{accidente} aperfeiçoa a \emph{substancia}. Ora a \emph{substancia divina}, possuindo \& plenitude da perfeição, não pode receber novas perfeições. Logo \textsc{Deus} não é composto de \emph{substancia} e de \emph{accidentes}.\par c) Em \textsc{Deus} não existe composição \emph{logica}, isto é, de \emph{genero} e \emph{differença}. Por quanto o \emph{genero} refere-se á \emph{differenga}, como a \emph{potencia} ao \emph{acto}. Logo \textsc{Deus}, por ser \emph{acto purissimo}, não é composto de \emph{genero} e de \emph{differenca}. --- Além d'isso, o ente, composto de \emph{genero} e de \emph{differenca}, é \emph{finito}, porque só possue a perfeição propria de uma certa especie de entes. Ora Dus é \emph{infinito}. --- Mais, a composição de \emph{genero} e de \emph{differenca}, embora seja \emph{intencional}, comtudo suppõe uma composição \emph{real} de \emph{potencia} e de \emph{acto} ; pois o \emph{genero} funda-se na \emph{potencia}, e a \emph{differença} no \emph{acto}. Ora em Deus não ha composição \emph{real}, e por isso não pode haver composição \emph{logica}. --- Se Deus podesse estar subordinado a um \emph{genero}, este seria o \emph{ente} ou a \emph{substancia}. Ora não pode ser o \emph{ente}; porque este não se attribue a Deus e ás creaturas em significação «\emph{nivoca}; visto que, quando se attribue a Deus, \emph{ente} significa o proprio ser subsistente em si, e, quando se attribue ás creaturas, significa o ser recebido num sujeito ou potencia, qué é a essencia. Nem pode ser a \emph{substancia}; porque a \emph{substancia}, que é \emph{genero}, não é o proprio ser, aliás toda a substancia Seria o proprio ser, isto é, o seu ser não seria produzido por outra causa; e por isso Deus, que é o proprio ser, não pode estar sujeito ao genero de \emph{substancia}. --- Todavia, \emph{analogicamente e} num sentido lato, podemos dizer que \textsc{Deus} é \emph{substancia}, emquanto esta significa um ente \emph{subsistente}. --- D'ahi se vê que não podemos formar uma \emph{rigorosa} definição de \textsc{Deus}, porque uma definição \emph{rigorosa} é composta de \emph{genero} e de \emph{differença}. Mas nem por isso confundimos \textsc{Deus} com as creaturas. \textsc{Deus} distingue-se das creaturas --- não por differenças especificas e componentes, --- mas pela pureza absoluta, pela substancia simplicissima do proprio ser (\emph{Sum}. \emph{Th}. p. 1 q: 3 a. 5).}.

\pg{453} 30. \textbf{Immutabilidade de Deus}. --- \emph{Immutabilidade} é o attributo, que afasta de \textsc{Deus} toda a \emph{mudança}. --- A \emph{mudança} pode ser

\pg{454} duplice: \emph{physica} ou \emph{moral}. A mudança \emph{physica} consiste na adquisição ou na perda de alguma entidade ou perfeição; a

\pg{455} \emph{moral} consiste na passagem da \emph{vontade} de um para outro acto,

de um para outro proposito \footnote{S. Thomaz trata d'isto --- na \emph{Sum}, \emph{Th}., p. I, q. 9; --- q. 14, a. 15; q. 19, a. 7; --- \emph{c}. \emph{Gentes}, l. T, c. 18 e seg.; --- 1 \emph{Sent}. Dist. 8, q. 8, a. 1 e 9,}.

31. \textbf{Deus} é immultavel.

a) \textsc{Deus} é \emph{absolutamente} immutavel. Todo o ente, que se muda de algum modo, é composto de \emph{potencia} e de \emph{acto}. Ora \textsc{Deus} é acto simplieissimo e purissimo, sem sombra de potencialidade. Logo \textsc{Deus} é \emph{absolutamente} immutavel.

b) \textsc{Deus} é \emph{physicamente} immutavel. Todo o ente, que se muda \emph{physicamente}, ou adquire uma perfeição que não tinha, ou perde uma perfeição que tinha. Ora \textsc{Deus}, sendo o Ente \emph{infinito e necessario}, possue todas as perfeicóes possiveis, e não póde perdel-as, aliás perderia a essencia, com que as perfeições se identificam. Logo \textsc{Deus} é \emph{physicamente} immutavel.

\pg{456} c) \textsc{Deus} é \emph{moralmente} immutavel. Com effeito, a mudança \emph{moral}, ou de \emph{vontade}, deriva de tres causas, --- ou da mudança \emph{physica} do sujeito, que, por isso, começa a desejar uma coisa, pela qual tinha repugnancia, --- ou da mudança do \emph{conhecimento}, emquanto a intelligencia descobre uma causa, que antes ignorava, e assim determina a vontade a mudar de proposito, --- ou da leviandade e inconstancia da propria \emph{vontade}, que, sem motivo sufficiente, passa de um proposito para outro. Ora nenhuma d'estas tres causas pode admittir-se em Dus, que, sendo infinitamente perfeito, não está sujeito a mudanças Physicas, nem á ignorancia ou inconstancia. Logo \textsc{Deus} é \emph{moralmente} immutavel \footnote{D'onde resulta que \textsc{Deus}, quando quer ou permitte effeitos diversos ou contrarios entre si, não muda de vontade; porque uma coisa é mudar de vontade e outra coisa é querer a mudança no effeito, e um agente póde, com o mesmo decreto, prescrever que presentemente se faça uma coisa, e que depois se faça o contrario. --- Como tambem, quando criou o mundo, \textsc{DEus} não mudou de vontade, porque não mudou de essencia (visto que em \textsc{Deus}, como provaremos, as \emph{operações} se identificam com a \emph{essencia}), mas, desde a eternidade, decretou que, em um certo tempo, existisse o mundo. Diz S, Thomaz: « Aliud est mutare voluntatem, et aliud est velle aliquarum rerum mutationem. Potest enim aliquis, eádem voluntate immobiliter permanente, velle quod nunc fiat hoe, et postea fiat contrarium. Sed tunc voluntas mutaretur, si aliquis inciperet velle quod prius non voluit. Quod quidem accidere non potest, nisi praesupposita mutatione vel ex parte cognitionis, vel circa dispositionem substantiae ipsius volentis » (\emph{Sum}. \emph{Th}, p. I, q. 19, as Entende-se egualmente como é que Deus, sem alguma mudança na sua vontade, possa agora amar e amanhã deixar de amar uma e a mesma pessoa, que hoje é santa e amanhã é peccadora; porque a mudança não se deu na vontade de Deus, mas só no homem, que antes era digno de amor, e depois se tornou digno de odio. Diz S. Agostinho: « Oum iratus malis dicitur (Deus), et placidus bonis, illi mutantur, non Ipse: sicut lux infirmis oculis aspera, firmis lenis est, ipsorum scilicet mutatione, non sua » (\emph{De Trin}., 1. V, c. 16).\par Objectam: Deus é livre com relação ás creaturas. A liberdade suppõe a indifferença; de modo que o acto, que se diz livre, pode ser e pode não ser realizado. Ora se Deus pode realizar e pode deixar do realizar um acto, é evidente que, quando o realiza, a sua vontade se muda. :\par Resposta. Tudo o que Deus quer fóra da sua essencia, não só o quer desde a eternidade, mas tambem com o mesmo simplicissimo acto, com que ama a si mesmo; pois, em Deus, não ha nem pode haver composição ou vicissitude de actos. O acto da divina vontade, sendo simplicissimo, refere-se ao objecto de um modo proporcionado á condição do mesmo objecto. Logo Deus ama a si mesmo como \emph{im}, e ama as outras coisas como \emph{meios} para o mesmo fim. E, como estes meios não têm uma relação necessaria com o fim, Deus pode querel-os e deixar de querel-os; e, ainda que ame a si mesmo \emph{necessariamente}, comtudo ama ou quer as outras coisas \emph{contingente} e \emph{livremente}, porque assim o exige a diversa qualidade dos objectos. --- Por isso\pgn{457}Deus quer verdadeiramente as creaturas, quando as quer; e não as quer verdadeiramente, quando não as quer; mas o querer e o não querer de Deus não suppõe uma mudança da parte de Deus, mas da parte das creaturas, que se referem \emph{realmente} ao proprio Deus. --- Para a comprehens\&o d'esta verdade, lembramos que a causa efficiente creada não se muda, porque opera; mas muda-se, porque, operando, passa necessariamente da potencia para o acto; e por isso Deus, que é acto purissimo, não se muda quando opera. --- Em poucas palavras: a liberdade, no seu conceito abstractissimo, exprime indifferenca não para com a propria \emph{entidade}, mas .para com o \emph{termo}, de modo que a vontade possa querer ou deixar de querer um objecto; mas não exige que este querer ou deixar de querer seja expressão de actos \emph{entitativamente} diversos, se não quando ella se considera na vontade creada, a qual, pela sua volição, passa da potencia para o acto.\par \emph{Renan}, seguido pelos defensores do \emph{progresso indefinido}, distingue um duplice \textsc{Deus}, --- o \textsc{Deus} \emph{ideal} e o \textsc{Deus} \emph{real}. O \textsc{Deus} \emph{ideal} é, segundo esse eseriptor, immutavel e eterno; o \textsc{Deus} \emph{real}, que se identifica com a \emph{universalidade das coisas ewistentes}, \emph{é mutavel} e capaz de progresso indefinido.\par Estas impias e absurdas affirmações, que encerram o mais crasso \emph{pantheismo}, não merecem ser novamente refutadas, pois já o foram no capitulo antecedente, quando provámos que existe um \textsc{Deus} \emph{real} e \emph{verdadeiro}.}.

\pg{457} 82. \textbf{Eternidade de Deus}. --- A \emph{eternidade} é o attributo, que exclue de \textsc{Deus} o \emph{principio}, o \emph{fim}, e toda a \emph{successão} \footnote{\emph{Boecio} dà á seguinte definição da \emph{eternidade}: « A \emph{eternidade é a posse total e perfeita de uma vida sem fim} ». A nalysemos esta definição. A \emph{eternidade} --- é \emph{posse}, porque é uma coisa estavel, que se conserva com socego e não póde ser tirada por nenhuma força extranha; --- é posse \emph{total}, porque a duração de \textsc{Deus} é sempre completa e indivisivel, e distingue-se da duração do ser finito, a qual não é \emph{total} mas \emph{successiva} ; --- \emph{perfeita}, porque a \emph{eternidade} é o \emph{presente}, que dura sempre, e por isso é diversa do tempo \emph{presente} dos seres finitos, que, embora, por não ter partes, seja \emph{total}, comtudo não é uma coisa absolutamente perfeita (pois não é perfeição possuir uma coisa apenas por um instante, e o tempo \emph{presente} das coisas finitas é um instante); --- é posse da uma \emph{vida}; porque a duração refere-se principalmente no \emph{acto} dos seres nobilissimos, que é \emph{viver}. --- de uma vida \emph{sem fim}, isto é, que não principiou nem ha de acabar. --- A \emph{posse} sem \emph{principio} e sem \emph{fim} é o elemento \emph{material} da \emph{eternidade}; a \emph{stimultaneidade} da posse é o elemento \emph{formal}.\par \emph{A eternidade}, pois, consiste na exclusão do \emph{principio}, do \emph{fim}, e de toda a \emph{successdo}, e identifica-se com a propria substancia de \textsc{Deus}, emquanto Esta não teve \emph{principio}, nem terá \emph{fim}, e não está sujeita a \emph{successão}.\par Por isso o \emph{tempo} não é uma parte da \emph{eternidade}, porque a \emph{eternidade} não tem \emph{partes}; mas a \emph{eternidade} corresponde e assiste a todos os seculos, que se vão succedendo, assim como um \emph{centro} immovel corresponde, por meio de linhas, a todos os pontos de uma esphera, que se mova em volta d'elle.\par D'onde se vé o erro --- tanto de \emph{Loche}, quando disse que a \emph{eternidade} é constituida pela infinita multiplicação do \emph{tempo}, --- como de \emph{Clarke}, quando opinou que a \emph{eternidade} é o proprio tempo; considerado em \emph{abstracto}.\par S. Thomaz trata d'este argumento --- na \emph{Summ}, \emph{Theol}. p. T, q. 10; --- \emph{c}. \emph{Gentes}, l. I, c. 15; --- I \emph{Sent}. Dist, 8, q. 2, a. 1, e Dist. 19, q.2, a. 1 e 9.}

\pg{458} 33. \textbf{Deus} é \textbf{eterno}. --- Um ente e é diz-se \emph{eterno}, quando não tem \emph{principio}, nem \emph{fim}, e não está sujeito a \emph{successão}. Ora \textsc{Deus} não tem \emph{principio} nem \emph{fim}, e não está sujeito á \emph{successão}. --- Não tem \emph{principio} nem \emph{fim}. Na verdade, \textsc{Deus}, sendo o Ente \emph{necessario}, \emph{immutavel} e \emph{independente}, existiu sempre e não póde deixar de existir; pois, se tivesse principiado a existir, quando não existia, ou podesse deixar de existir, seria \emph{contingente}, \emph{mutavel} e \emph{dependente}. --- Não está sujeito a \emph{successão}. Porquanto, ' a \emph{successão} indica \emph{mudança} e \emph{imperfeição}, e \textsc{Deus} é \emph{immutavel} e infinitamente \emph{perfeito}. Logo \textsc{Deus} é \emph{eterno} \footnote{Dizemos no texto que toda a \emph{successão} indica \emph{imperfeição}; porquanto o ente, que admitte \emph{successão} nos seus estados ou qualidades, adquire perfeições que não tinha, ou perde as que tinha, e por consequencia é finito ou limitado nas suas perfeições.\par \textsc{Deus} não só é \emph{eterno}, mas é a sua \emph{eternidade}; porque a eternidade não é uma coisa distincta de Deus, mas é o seu proprio ser uniforme; e por isso Deus, assim como é a sua essencia, assim tambem é a sua eternidade. Diz S. Thomaz: « Ratio aeternitatis consequitur immutabilitatem, sicut ratio temporis consequitur motum. Unde cum Deus sit maxime immutabilis, sibi maxime competit esse aeternum. Nec solum est aeternus, sed est sua aeternitas; cum tamen nulla alia res sit sua duratio, quia non est suum \emph{esse}. Deus autem est suum esse uniforme. Unde: sicut est sua essentia, ita sua aeternitas » (\emph{Sum}. \emph{Th}, p. I, q. 10, a. 2).\par Como a \emph{eternidade} consiste na perfeita immutabilidade, e todas as creaturas são mutayeis, segue-se que a nenhuma creatura compete a eternidade, no sentido rigoroso da palavra, e por isso só Deus é eterno. --- Todavia costumamos dizer que algumas creaturas participam da eternidade, como participam da immutabilidade no seu ser ou nos seus actos. Mas trata-se sempre de uma eternidade impropria e imperfeita (\emph{Sum}. \emph{Th}., P- 1, q. 10, a. 3).}.

34. \textbf{Immensidade de Deus}. --- \emph{Zmmensidade} é o attributo, pelo qual \textsc{Deus} está presente \emph{actualmente} em todas as coisas \emph{emistentes} e \emph{virtualmente} em todas as coisas \emph{possiveis} \footnote{Quando dizemos que \textsc{Deus} está \emph{virtualmente} presente em todas as coisas \emph{possiveis}, entendemos que \textsc{Deus} tem a \emph{virtude} ou o poder de estar em todas as coisas, que não existem, mas que podem existir. ' E' necessario distinguir a \emph{immensidade} da \emph{ubiquidade}. --- A \emph{immensidade} é a presença divina, \emph{actual} ou \emph{virtual}, em todas as coisas, tanto \emph{existentes} como \emph{possiveis}; a \emph{ubiquidade} é a \emph{actual} presença de \textsc{Deus} nas coisas \emph{ezístentes}. --- A \emph{immensidade} é uma coisa eterna, absoluta e infinita, porque se extende a todo o possivel, e prescinde da \emph{actual} existencia das coisas; a \emph{ubiquidade} é uma coisa temporanea, relativa, não é rigorosamente infinita, porque suppõe a existencia do mundo, e porque as coisas creadas não são infinitas. --- A \emph{ubiquidade} está contida na \emph{immensidade}. --- Por isso \textsc{Deus},\pgn{459}embora não esteja presente nas coisas finitas, senão depois da creação, todavia é \emph{immenso} desde a eternidade; porque a propria Essencia de \textsc{Deus} exige que Elle possa e deva estar presente em todas as coisas, \emph{existentes} e \emph{possiveis}. --- S. Thomaz trata d'este assumpto --- na \emph{Sum}. \emph{Th}., p. I, q. 8, --- c. \emph{Gentes}, 1. III, c. 68; --- I \emph{Sent}. Dist. 91, q. 3, a. 1.}.

35. \textbf{Deus} é immenso. --- \textsc{Deus} será \emph{immenso}, se estiver presente --- não só em todas as coisas \emph{creadas} e em todos os logares \emph{existentes}, --- mas tambem em todas as coisas e logares \pg{459}  \emph{possiveis}. Ora \textsc{Deus} --- está presente em todas as coisas \emph{creadas}, porque Elle produziu e conserva todas as creaturas, e um agente não póde operar onde não existe; --- está presente em todos os logares \emph{existentes}, porque todo o logar é uma coisa creada; --- póde estar presente em todas as coisas e logares \emph{possiveis}, porque \textsc{Deus} deve estar em todo o ser, a que se extende o seu poder (pois, em \textsc{Deus}, os attributos identificam-se com a substancia), e o poder de \textsc{Deus} extende-se a todo o \emph{possivel}. Logo \textsc{Deus} é \emph{emmenso} \footnote{\textsc{Deus} está em todos os logare com \emph{poder}, emquanto tudo\par ;\par 7 gares --- o seu \emph{poder}, Lhe obedece, --- com a sua \emph{presença}, emquanto Elle vé tudo o que se faz e pensa, --- com a sua \emph{essencia}, emquanto produz e conserva o ser das creaturas (Cf. \emph{Sum}. \emph{Th}, p. I, q. 8, a. 9).\par Pergunta-se: Onde estava \textsc{DEus}, antes da creacüo do mundo? e agora como está Elle no mundo?\par Antes da creacüo, \textsc{Deus} estava em si mesmo. E, quando o mundo foi creado, \textsc{Deus} não deixou de estar em si mesmo, nem se mudou pelo facto de estar presente em todas as coisas; assim como não se muda o sol, quando se mudam as coisas, que elle illumina. --- E, agora, \textsc{DEus} está no mundo, n\&o diffundido por toda a parte, como a luz ou o ar, mas de um modo ineffavel, que excede toda a \emph{imaginução}. Uma semelhança, embora infinitamente \emph{inadequada}, da presença de \textsc{Deus} nas coisas é a nossa alma, a qual, ainda que \emph{simples} e \emph{espiritual}, está toda em todo o organismo e toda em cada parte d'elle. Dizemos --- infinitamente \emph{inadequada}, porque a nossa alma é limitada ao nosso corpo, ao passo que \textsc{Deus} não é limitado a este mundo.\par Muitos negaram ou desvirtuaram a \emph{Zmmensidade} de \textsc{DEus}.\par \emph{a}) Uns, como os \emph{Manicheus} e \emph{Descartes}, negaram a \emph{ubiquidade} de \textsc{Deus}; pois aquelles diziam que só as coisas \emph{espirituaes} estão sujeitas ao \emph{poder} de \textsc{Deus}, e \emph{Descartes} affirmou que \textsc{Deus} está em todas as coisas com o seu \emph{poder} e \emph{presença}, mas não com a sua \emph{essencia} ou \emph{substancia}.\par Ambas essas affirmações são falsas. --- E' falsa a dos \emph{Manicheus}; porque o poder de \textsc{DEus} é infinito e extende-se “a todas as creaturas, espirituaes e corporeas, --- E' falsa a de \emph{Descartes}; porque, identificando-se em \textsc{Deus} as operações com a essencia, onde existem aquellas, deve existir esta.\par \emph{b}) Outros, seguindo as theorias de \emph{Newton} e \emph{Clarcke}, disseram que a \emph{immensidade divina} é o espaço infinito.\par Esta opinião foi refutada na \emph{Ontologia}.\par \emph{c}) \emph{Spinoza} disse que a \emph{émunensidade} não é o attributo de um Ente individuo e pessoal, mas é a \emph{collecção} de todas as coisas.\par Estas affermações contêm o \emph{pantheismo}, que já foi refutado.\par \emph{d}) \emph{Bayle} diz que \textsc{Deus}, sendo \emph{simplicissimo}, não póde ser \emph{immenso}; porque a \emph{simplicidade} e a \emph{iémmensidade} são attributos contradictorios.\par Respondemos que a \emph{simplicidade} não só não é contraria á \emph{immensidade},\pgn{460}mas é uma condição indispensavel para ella; porque \textsc{Deus} não seria \emph{immenso}, se não fosse \emph{simplicissimo}, pois o \emph{composto}, resultando de partes necessariamente finitas e limitadas, não póde ser 2\emph{mmenso}.}.

\pg{460} 36. \textbf{Unidade de Deus}. --- \emph{Unidade} é o attributo, pelo qual a Essencia divina não póde dividir-se, multiplicar-se e constituir varios e diversos deuses \footnote{Quando, pois, dizemos que \textsc{Deus} é \emph{uno}, entendemos que é uno não só \emph{especifica} mas tambem \emph{numericamente}, e por isso é \emph{unico}.\par S. Thomaz trata d'este assumpto --- na \emph{Sum Th}., p. I, q. 1, a.8 e 4; --- c. \emph{Gentes}, l. I, c. 42; --- \emph{De pot}., q: 8, a. 6; ----I \emph{Sent}. Dist. 2, a. 1, --- II \emph{Sent}. Dist. 91, q. 1, a. 1.}.

37. \textbf{Deus} é uno.

£) \textsc{Deus} é o Ente infinitamente perfeito. Ora o ente infinitamente perfeito deve ser \emph{uno}. Porquanto, se fossem varios, deveria haver entre elles uma distineção. Esta distinegüo só poderia dar-se, emquanto um fosse dotado de uma perfeição, que faltasse a outro. Mas, neste caso, aquelle, a quem faltasse uma perfeição, não seria \textsc{Deus}; mas só o seria Aquelle, que fosse infinitamente perfeito. Logo \textsc{Deus} é \emph{uno}.

b) Uma forma multiplica-se, emquanto é recebida em muitos e diversos sujeitos; assim a \emph{humanidade} multiplica-se, emquanto é recebida em muitos e diversos homens. Se, porém, a forma não for recebida em sujeito algum, mas subsistir em si mesma, não poderá multiplicar-se, e existirá um unico individuo d'aquella especie; assim, se a \emph{Humanidade} fosse subsistente em si e por si, existiria um unico homem. Ora a Divindade não é recebida em sujeito algum (porque, em \textsc{Deus}, infinitamente \emph{simples}, a \emph{essencia} identifica-se com o \emph{supposto}), mas é subsistente em si e por si. Logo a Divindade não póde multiplicar-se e formar varios suppostos. Logo \textsc{Deus} é \emph{uno} \footnote{Nunca devemos esquecer que a \emph{subsistencia do ser divino} é, segundo o nosso modo de entender, a razão e a raiz de todos os \&ttribütos divinos. Ser das creaturas é simples como é simples uma parte ou uina coisa incompleta, porque é uma forma recebida num sujeito e individualisada pela natureza, que subsiste nesse ser, e por isso é multiplicavel e molti- Plica-se conforme a multidão das creaturas. Pelo contrario, o ser de Deus é simples, como é simples o acto e a perfeição, porque é o ser por si subsistente e por si individualisado e distincto de todo e qualquer outro ser, e por isso não é multiplieavel, mas é \emph{unico}.}.

c) Se existissem varios deuses, \emph{especificamente} eguaes, mas \emph{numericamente} distinctos, seriam necessariamente compostos de \emph{natureza divina}, \emph{commum} a todos, e de \emph{notas} ou \emph{attributos}, \emph{proprios} de cada um (aliás seria impossivel a egualdade \emph{especifica} a e distineção \emph{mumerica}). Ora \textsc{Deus}, sendo \emph{simplicissimo}, exclue toda \pg{461}  e qualquer composição. Logo não podem existir varios deuses. Logo \textsc{Deus} é \emph{uno} \footnote{Os systemas ou crenças contrarias \& \emph{unidade} de \textsc{Deus} foram o \emph{dualismo} e o \emph{polytheismo}.\par a) O \emph{dualismo} admittia que existem \emph{dois principios infinitos}, \emph{necessarios}, e \emph{supremos} --- um \emph{bom} e outro \emph{mau}, e que do \emph{principio bom} derivam todos os \emph{bens} e do \emph{mau} todos os \emph{males}.\par O \emph{dualismo} foi chamado tambem \emph{Manicheismo}, porque foi \emph{Manes} ou \emph{Manicheu} quem primeiramente o ensinou. --- Nos tempos modernos foi resuscitado pelo sceptico \emph{Bayle}, ensinando que o \emph{dualismo}, embora-seja absurdo \emph{a-priori}, comtudo póde provar-se \emph{a-posteriori}.\par O \emph{dualismo} é falso-pelas seguintes razões:\par \emph{a}) \emph{Baséa-se num falso fundamento}. Porquanto, admitte que existe um principio inteiramente mau; porque, aliás, não poderia explicar-se a existencia do mal. Ora esta razão é falsa; porque o mal, derivando das causas creadas, é uma consequencia do \emph{bem limitado}.\par \emph{b}) \emph{Um principio inteiramente mau não póde existir}. Porquanto, esse principio --- ou é privado de toda a perfeição e entidade, --- ou falta-lhe apenas a bondade \emph{moral}, --- Na primeira hypothese, um ente inteiramente mau seria \emph{nada} e o \emph{nada} não existe. --- Na segunda, um principio inteiramente mau seria, ao mesmo tempo, infinito e finito. Seria \emph{infinito}, porque o systema suppõe-no tal. Seria \emph{finito}, porque faltar-lhe-hia uma perfeição.\par €) \emph{O dualismo perverte a moral}. Com effeito, se todo o \emph{bem} devesse attribuir-se a um principio \emph{bom}, e todo o \emph{mal} a um principio \emph{mau}, o homem não poderia dizer-se e não seria causa do bem nem do mal, e por isso não seria merecedor de premio nem de castigo.\par \emph{d}) \emph{Esta hypothese não explica a existencia do bem e do mal}. Porquanto --- ou estes dois principios seriam eguaes, --- ou o principio \emph{bom} seria mais forte que o \emph{mau}, --- ou este mais forte que aquelle. Na primeira hypothese, as forças dos dois principios paralysar-se-hiam, e assim não existiria, sobre a terra, nem o bem nem o mal. Na segunda, não haveria mal algum. Na terceira, não haveria bem algum. Ora estas consequencias são falsas.\par \emph{Bayle} diz que os dois principios fizeram um pacto, para se não transtornarem mutuamente nas suas obras.\par Mas esta explicação é absurda. Se assim fosse, esses dois principios seriam fracos, e nenhum d'elles seria infinitamente perfeito. --- De mais, esse pacto ou seria bom ou mau. Se fosse bom, não poderia ser acceite pelo principio \emph{mau}; pois este, sendo inteiramente \emph{mau}, não póde fazer nem tolerar bem algum. Se fosse mau, não poderia ser acceite pelo principio \emph{bom} ; pois este, segundo o \emph{dualismo}, não póde fazer nem tolerar mal algum.\par b) O \emph{polytheismo} é o erro que admitte varios \emph{deuses}.\par Este erro, que derivou sobre tudo da corrupção da vontade, não foi constante nem universal, --- Não foi \emph{constante}; porque não existia antes do \emph{diluvio}, \emph{e} porque a primeira crença dos povos foi o \emph{monotheismo} (Cf. Hooke, \emph{Princip}. \emph{vélig}. \emph{natural}. \emph{et révél},, --- Frank, \emph{Études Orientales}). --- Não foi \emph{universal}; porque, além de muitos philosophos (como \emph{Anaxagoras}, \emph{Socrates}, \emph{Platão}, \emph{Aristoteles}, \emph{Cicero}), alguns povos, como o \emph{d'Zsrael}, professavam o \emph{monotheismo}, e muitos outros, embora -admittissem deuses, comtudo reconheciam \emph{Jupiter}, como pae e senhor de todos os deuses.}.

%%>>> art-41

\section*{Artigo III.}
\addcontentsline{toc}{section}{Artigo III.}

\pg{462} \textbf{Attributos de Deus positivos},

38. Attributos de \textbf{Deus} positivos. --- Attributos de \textsc{Deus} \emph{positivos} são os que \emph{affrmam} de \textsc{Deus} as perfeições dos seres creados, mas isentas de toda a imperfeição ou de todo o limite. --- Os principaes são: a \emph{verdade}, a \emph{bondade} e a \emph{belleza} \footnote{Depois dos attributos \emph{negativos} de \textsc{Deus}, devemos considerar os attributos \emph{positivos}, E primeiramente consideraremos os attributos \emph{transcendentaes}, que \emph{analogicamente} convêm a Dzvus e ás creaturas. Como dissemos na Ontologia, attributos \emph{transcendentaes} são os que convêm a \emph{todo o ente}, emquanto \emph{tal}, e que, embora não sejam realidades distinctas do proprio ente, comtudo são explicações mais claras do mesmo, porque exprimem o que o ente não exprime. --- São quatro: \emph{unidade}, \emph{verdade}, \emph{bondade}, \emph{bellezza}. --- Tendo tratado da \emph{unidade}, resta-nos tratar dos outros attributos.}.

39. \textbf{Verdade} de \textbf{Deus}. --- \emph{Verdade} de \textsc{Deus} é o attributo, que consiste na \emph{conformidade}, ou \emph{adequação}, entre o acto da intelligencia divina e o ser divino \footnote{S. Thomaz trata d' este assumpto --- na \emph{Sum}. \emph{Th}., p. T q. Ignem c. \emph{Gentes}, 1. I, c. 60 e seg., --- \emph{De verit}., q. I, --- I \emph{Sent}. Dist. 19, q.5, a.a. 9 e 8.}.

40. \textbf{Deus} é verdadeiro. --- Um ente é \emph{verdadeiro}, quando ha conformidade entre o seu ser e o acto da intelligencia, á qual de algum modo se refere. Ora entre o ser de \textsc{Deus} e o acto da sua intelligenecia não só ha \emph{conformidade}, mas absoluta \emph{identidade}; porque, em \textsc{Deus}, á substancia e os actos das faculdades são uma e a mesma coisa. Logo \textsc{Deus} é \emph{verdadeiro} \footnote{A verdade, por ser uma adequação entre a intelligencia e o objecto, deve encontrar-se em \textsc{Deus}, cuja intelligencia é uma e a mesma coisa com \& essencia. Na sua essencia, como veremos, \textsc{DEus} conhece todas as coisas. --- Nem se diga que a verdade, consistindo \emph{na composição} e na \emph{divisão}, não pode encontrar-se em \textsc{Deus}, cuja intelligencia é \emph{simplicissima}, como é a sua essencia. O principio --- a \emph{verdade consiste na composição e na divisão} --- verifica-se unicamente com relação á verdade, que é propria da intelligencia \emph{humana}, e não pode applicar-se á intelligencia \emph{divina}. Com effeito, quando se diz que a verdade consiste na adequação entre a intelligencia e o objecto, --- não devemos entender que o \emph{acto} da intelligencia deve por tal modo egualar o \emph{objecto}, que, se este fór composto e material, tambem composto e material deva ser aquelle (o que é absurdo).--- mas devemos entender que o \emph{verbo} da intelligencia deve por tal forma corresponder ao seu \emph{objecto}, que tudo o que o \emph{verbo} exprime se deve encontrar realmente no \emph{objecto}, o qual pode ser simples ou composto, espiritual ou material. Portanto a \emph{qualidade}, isto é, a \emph{simplicidade} ou a \emph{composição} do \emph{acto intellectual} não entra na razão ou essencia da \emph{verdade}, mas refere-se exclusivamente ao principio intelligente; e por isso, se a, intelligencia do homem, por ser\pgn{463}a menos perfeita, não pode alcançar a verdade se não por meio da \emph{composição} e da \emph{divisão} (juizo \emph{affirmativo} ou \emph{negativo}), a intelligencia de \textsc{Deus}, por ser infinitamente perfeita e simplicissinia, como a sua essencia, com um acto unico e simplicissimo alcança toda a vendade, quer esta se refira á essencia divina, quer se refira ás coisas creadas, simples ou compostas, espirituaes ou materiaes. (\emph{Cf}. \emph{C}. \emph{Gentes}, l. I. c. 59).}.

\pg{463} 41. \textbf{Deus} é a \textbf{verdade}. --- Consistindo a verdade na conformidade entre o acto da intelligencia e o ser, \textsc{DEUS} será a propria \emph{verdade}, se fôr o seu entender e o seu proprio ser. Ora \textsc{Deus} é, effectivamente, o seu entender e o seu proprio ser. Logo \textsc{Deus} é a \emph{verdade}. --- \textsc{Deus} é a verdade \emph{summa} e \emph{infinita}, porque summa e infinitamente perspicaz é a sua intelligencia, e summa e. infinitamente intelligivel é o seu ser \footnote{Em \textsc{Deus} não só ha a verdade, mas Elle é a propria verdade. Por quanto, tudo o que está em Deus não é \emph{participado} nem \emph{accidental}, mas é \emph{essencial} e \emph{substancial}, é o proprio ser divino, é o proprio Deus. Logo Deus, assim como é á sua essencia, o seu ser, assim tambem é a sua verdade. Diz S. Thomaz: « De Deo nihil \emph{participative} dici potest, cum sit suum esse, quod nihil participat. Sed veritas est in Deo. Si igitur non dicatur \emph{participative}, oportet quod dicatur \emph{essentialiter}. \textsc{DEUS} ergo est sua veritas ». (\emph{C}. \emph{Gentes}, l. I, c. 60). --- A verdade divina é tão pura, que não tolléra sombra de falsidade ou de engano. « A verdade, diz o mesmo Angelico, não é compativel com a falsidade, como a candura não é compativel com a côr preta. Ora Deus não só é verdadeiro, mas é a propria verdade. Logo em Deus não pode haver falsidade » (\emph{Ib}., c. 61).}.

42. \textbf{Deus} é a \textbf{verdade} na ordem ontologica e na ordem logica.

a) \textsc{Deus} É A VERDADE NA ORDEM ONTOLOGICA. --- \textsc{Deus} será a verdade na ordem \emph{ontologica}, se a sua essencia ou substancia se conformar com a sua intelligencia; porque a verdade \emph{ontologica} consiste na conformidade do objecto com a intelligencia. Ora a essencia de \textsc{Deus}, não só se conforma com a sua intelligencia, mas identifica-se com ella, e \textsc{Deus} é a sua essencia. Logo Deus é a verdade na ordem \emph{ontologica}.

b) \textsc{Deus} É A VERDADE NA ORDEM LoGICA. --- \textsc{Deus} será a verdade na ordem \emph{logica}, se a sua intelligencia se conformar com a sua essencia; porque a verdade \emph{logica} consiste na conformidade da intelligencia com o objecto. Ora, como dissemos, á intelligencia de \textsc{Deus} não só: se conforma com a sua essencia, mas identifica-se com ella, e \textsc{Deus} é a sua intelligencia. Logo \textsc{Deus} é á verdade na ordem \emph{logica} \footnote{A verdade, como dissemos, pode considerar-se no \emph{objecto} e na \emph{intelligencia}. Considerada no \emph{objecto}, é a conformidade d'este com a intelligencia (\emph{veritas rei}); --- considerada na intelligencia, é \& conformidade d'esta com o objecto (\emph{veritas intellectus}) Em Deus ha a verdade do \emph{objecto} ou \emph{ontologica}, e ha a verdade da \emph{intelligencia} ou \emph{logica}. Diz 8. Thomaz: « Veritas\pgn{464}invenitur \emph{in intellectu}, secundum quod apprehendit \emph{rem ut est}; et \emph{in} re, secundum quod habet esse \emph{conformabile intellectui}. Hoc autem maxime invenitur in Deo, Nam esse suum mon solum est conforme suo intellectui, sed etiam est ipsum suum intelligere.., et ipse est suum esse et intelligere. Unde sequitur quod non solum in ipso sit veritas, sed quod ipse sit ipsa summa veritas » (\emph{Sum}. \emph{Th}., p. I, q. 16, a. 5). \textbar{}}.

\pg{464} 43. A verdade ontologica de \textbf{Deus} é a \textbf{causa efficiente da verdade} ontologica das \textbf{creaturas}. --- A verdade \emph{ontologica} de Dzvus será a causa \emph{efficiente} da verdade \emph{ontologica} das creaturas, se o ser d'estas for effeito e expressão da Essencia divina. Ora o ser das creaturas, tendo sido produzido por \textsc{Deus} e em conformidade com as idéas arehetypas, que são imitações diversas da Bsseneia infinita, é, na sua physica realidade, expressão e effeito da mesma Essencia. Logo a verdade \emph{ontologica} de \textsc{Deus} é á causa \emph{efficiente} da verdade \emph{ontologica} das creaturas \footnote{A Essencia divina é uma simplicissima unidade, mas é imitavel e participavel, por infinitos e diversos modos, pelas creaturas. « Cada creatura, diz S. Thomaz, possue a propria especie, emquanto de algum modo participa da semelhança da Essencia divina ». (\emph{Sum}. \emph{Th}, p. 1, q. 15, a. 2) Por isso cada realidade creada é uma \emph{expressão}, ou \emph{semelhança}, da unica Essencia divina. ---- E' uma \emph{expressão}, e tambem um \emph{effeito}. Diz o mesmo santo Doutor: « Um ente, que possue o ser \emph{por essencia}, é necessariamente causa de todos os entes, que o possuem por \emph{participação}... Mas Deus possue o ser pela sua essencia, porque é o proprio ser; e todos os outros entes \emph{so} o possuem por participação, porque o ente, que é o proprio ser, é unico. Portanto, Deus é causa do ser de todos os outros entes » (\emph{C}. \emph{Gentes}, 1. II, e. 15). --- Sendo uma expressão e um effeito da Essencia divina, o ser das creaturas não pode deixar de corresponder ás ideas archetypas da intelligencia divina. Por isso, o ser das creaturas refere-se \emph{necessariamente} á intelligencia divina, emquanto não pode deixar de corresponder ás ideas archetypas, e só \emph{accidentalmente} á intelligencia/ creada, emquanto \emph{pode} ser conhecido por ella; e a verdade \emph{ontologica} das creaturas é \emph{necessaria} com relação a Deus, e é \emph{contingente} com relação a nós. --- Como tambem essa verdade é \emph{unica} e \emph{multipla}. E' \emph{unica} com relação a Deus: porque unica é a Essencia divina, que pode ser e é diversamente imitada e participada. E' \emph{multipla} com relação ás creaturas; porque, consistindo a verdade \emph{ontologica} na conformidade do ente com a intelligencia, são tantas as verdades, quantas são as entidades.}.

44. A verdade logica de \textbf{Deus} é a \textbf{causa exemplar da verdade} logica da inteligencia humana. --- A verdade \emph{logica} da intelligencia humana depende das coisas, emquanto a nossa intelligenoeia, para alcançar e possuir essa verdade, deve conformar-se com as proprias coisas. Ora conformar-se com as coisas é conformar-se com as ideas archetypas da intelligencia divina, de que as coisas dependem essencialmente e que as coisas exprimem na sua realidade. Logo as ideas da nossa intelligencia emtanto \pg{465}  são dotadas de verdade \emph{logica}, emquanto imitam as idéas archetypas divinas, isto é, a verdade \emph{logica} de \textsc{Deus}. Logo a verdade \emph{logica} de \textsc{Deus} é a causa \emph{exemplar} da verdade \emph{logica} da intelligencia humana \footnote{« A verdade divina, diz S. Thomaz, é a medida de toda a verdade. Por quanto, a verdade da nossa intelligencia depende e é medida pela coisa, que está fóra da alma; pois dizemos que a nossa intelligencia é verdadeira, emquanto se conforma com a coisa. Mas a verdade da coisa é medida pela intelligencia divina, que é a causa de todos os seres creados, como a verdade das coisas artificiaes é medida pela arte do artista. Sendo, pois, \textsc{Deus} a primeira intelligencia e o primeiro intelligivel, é forçoso concluir que a sua verdade é a medida da verdade da intelligencia creada. Por isso a verdade divina é a \emph{primeira} verdade » (\emph{C}. \emph{Gentes}, l. 1, c. 62). --- Portanto a verdade \emph{logica}, propria da nossa intelligencia, pode dizer-se, ao mesmo tempo, \emph{unica} e \emph{multiplice}. E' \emph{unica} com relação a \textsc{Deus}, emquanto as nossas ideas são \& expressão ou a semelhança, ainda que imperfeita, da unica e substancial verdade da intelligencia divina. E' \emph{multiplice} com relação ás \emph{formas} da intelligencia humana, que existem em muitas e distinctas intelligencias e representam muitos e diversos objectos.}.

45. A \textbf{verdade} divina é a \textbf{causa} final de toda a \textbf{verdade creada}. --- \textsc{Deus} não só é o principio, mas tambem é o fim de toda a creação; de modo que toda a perfeição creada não existe senão para imitar e revelar, embora de um modo inadequado, a correspondente perfeição divina. Ora a verdade creada é um seffeito, e a sua existencia não pode ter e não tem outro fim, senão o de imitar e manifestar a verdade divina. Logo a verdade divina é a causa \emph{final} de toda a verdade creada \footnote{\emph{A} verdade creada é, como dissemos, \emph{ontologica} e \emph{logica}, conforme se considera no \emph{objecto} ou na \emph{intelligencia}. Ora --- a verdade \emph{ontologica}, que é o proprio \emph{objecto} emquanto se conforma com a intelligencia, imita e manifesta a verdade \emph{ontologica} de Dzvus, porque todo o ente participado revela o ente por essencia, do qual deriva; --- e a verdade \emph{logica}, que é o verbo da intelligencia emquanto se conforma com o objecto, revela a verdade da intelligencia divina. Por isso a verdade creada, emquanto imita e revela a verdade divina, tende para esta, como para o seu fim natural; e d'este modo a verdade divina é a causa \emph{final} da verdade creada.}.

46. \textbf{Bondade} de \textbf{Deus}. --- \emph{Bondade} de \textsc{Deus} é o attributo, que consiste na \emph{conveniencia} da infinita perfeição de \textsc{Deus} com a sua santissima vontade. --- A. \emph{bondade}, pois, é a propria entidade, a propria perfeição do ente, emquanto \emph{convem} á vontade, e attrahe-a ou pode attrahil-a \footnote{S. Thomaz trata da \emph{bondade} de \textsc{DEUS} --- na \emph{Sum}. \emph{Th}. p. l, q. 6 --- c. \emph{Gentes}, l. I, c. 37 e seg. --- \emph{De verit}. q. 21.\par 30 Vor. IL.}.

47. \textbf{Deus} é bom. --- Um ente é e diz-se \emph{dom}, emquanto é dotado de perfeição, que tem a força de attrahir a vontade.

\pg{466} Ora Deus contém em si toda a perfeição, pela qual é objecto de um amor sem limites da parte da sua vontade; pois todo o ente não pode deixar de amar a propria perfeição. Logo \textsc{Deus} é \emph{bom}. --- \textsc{Deus} é \emph{summa} e \emph{infinitamente bom}, porque \emph{summa} e \emph{infinita} é a sua entidade ou perfeição \footnote{Deus é \emph{bom}, e \emph{summa} e \emph{infinitamente} bom, porque é \emph{perfeito}, e summa e \emph{infinitamente} perfeito. A bondade está na razão \emph{directa} da perfeição; porque a bondade nao é senão a perfeição, emquanto se refere á vontade. Diz o P. Lessio; « Imagina uma infinita luz, um infinito poder, uma infinita sabedoria, uma infinita belleza, uma infinita doçura, uma infinita alegria, uma infinita riqueza, uma infinita gloria, uma infinita majestade, uma infinita santitade. Pois bem, a Bondade divina contém todas estas perfeigóes como um oceano infinito de todos os bens. E não só contem todo o bem, mas contém-no de um modo infinitamente elevado; porque não o contem por muitas e diversas formas ou qualidades, mas por uma simplicissima e elevadissima forma, que é a sua essencia. Este modo de conter as perfeigóes, torna estas infinitamente mais excellentes, do que se existissem distinctas uma de outra. (\emph{De perfectione divina}, l. 7, e. 1, n. 2).}.

48. \textbf{Deus} é a \textbf{bondade}. --- A 2\emph{ondade} de \textsc{Deus} é a sua perfeição; e por isso tal é a natureza da sua bondade, qual é a natureza da sua perfeição. Ora a perfeição, em \textsc{Deus}, não é uma coisa accidental, addieionada á essencia, mas é uma coisa substancial, é a divina essencia; porque toda a perfeição é um ente, e a essencia divina é o proprio ser. Logo \textsc{Deus}, assim como é a sua essencia, assim tambem é a sua \emph{bondade} \footnote{Se \textsc{Deus} não fosse a sua bondade, isto é, se não fosse bom pela sua \emph{essencia}, seria bom \emph{por participação}. Mas então, como o ente, que é bom \emph{por participação}, presuppõe outro ente, anterior a si, do qual recebeu a bondade, Deus não seria a Causa primeira, o Motor immoyel,o Ente absoluto e infinito: o que é absurdo. Logo Deus é bom pela sua essencia, isto é, é a propria bondade. --- Se Deus é a propria Bondade, é claro que exclue absolutamente todo e qualquer \emph{mal}, que é opposto ao \emph{bem}. Alem d'isso, se o mal podesse existir em Deus, existiria como uma coisa \emph{essencial}, seria a propria \emph{essencia} de Deus, porque tudo o que está em Deus é o proprio Deus: o que é absurdo, porque o mal, sendo privação do bem, é nada, e o nada não pode constituir nenhuma essencia, muito menos a perfeitissima essencia divina. (Cf. c. \emph{Gentes}, l. T, c. 38 e 89).}.

49. A \textbf{bondade} divina é a causa efficiente de toda a bondade \textbf{creada}. --- A bondade divina será a causa \emph{efficiente} de toda a bondade creada, se as perfeicoóes, de que são dotadas as creaturas, derivarem da unica e simplicissima bondade de \textsc{Deus}. Ora é certo que todas as perfeigoes, de que são dotadas as creaturas, derivam da unica e simplicissima bondade de \textsc{Deus}; porque as perfeigóes das creaturas são finitas, contingentes, mutaveis, e por isso derivam da perfeição divina, que é infinita \pg{467} , necessaria; immutavel. Logo a bondade divina. é a causa \emph{efficiente} de toda a bondade creada \footnote{Toda a bondade das creaturas deriva, como da causa \emph{efficiente}, da bondade de Deus. Dizem os escholasticos: \emph{quod est maximum in wnoquoque genere est causa aliorum quae sunt in illo genere}. --- Mas que differenca entre a bondade de Deus e a das creaturas! A bondade de Deus é imparticipada e fonte inexhaurivel de todo o bem; a bondade das creaturas é participada e deriva de Deus. --- A bondade de Deus é illimitada, porque não ha quem a possa limitar; a bondade das creaturas é limitada, porque é recebida numa essencia finita. --- A bondade de Deus á substancial e identifica-se com a sua essencia; a bondade das creaturas é uma forma accidental e addicionada á essencia. --- A bondade de Deus é immutavel, não pode estar sujeita ao augmento nem á diminuição; a bondade das creaturas é capaz de augmento até ao infinito e de diminuição até ao nada. Por causa d'esta excellencia da bondade divina sobre a creada disse \textsc{JESUS CHRISTO}: « Só Deus é bom: \emph{Nemo bonus nisi solus Deus} » (Luc., XVIII, 19).}.

50. A \textbf{bondade} divina é a \textbf{causa exemplar da bondade creada}. --- O ente, que possue uma determinada qualidade por participação, não se diz dotado d'essa qualidade, se não emquanto manifesta uma certa \emph{semelhança} com o ente, que possue a mesma qualidade por essencia; assim como o ferro não se diz igneo, senão emquanto tem uma certa semelhança com o fogo. Ora \textsc{Deus} é bom por essencia, e todos os seres creados são bons por participação. Logo todo o ente creado diz-se \emph{bom}, emquanto tem uma certa semelhança com a bondade de Deus. Logo a bondade divina é a causa \emph{exemplar} da bondade creada \footnote{Cf. \emph{c}. \emph{Gentes}, l. 1, c. 40. --- A bondade divina é a causa \emph{exemplar}, mas não é a causa \emph{formal} da bondade creada, porque a bondade das creaturas não é uma e a mesma coisa com a bondade divina. A bondade creada é \emph{distincta} da bondade divina. --- Por isso a bondade creada é \emph{unica} e \emph{multiplice}. E) \emph{unica}, emquanto se refere a Deus, como ao seu primeiro principio \emph{efficiente} e \emph{exemplar}. E' \emph{multiplice}, emquanto a bondade ou perfeição creada se acha multiplieada tantas vezes, quantas são as creaturas.\par Advertimos que as creaturas não participam da \emph{bondade} de \textsc{Deus} pelo mesmo modo, por que participam da sua \emph{verdade}. Com effeito, as creaturas participam da \emph{bondade} de \textsc{Deus} --- não só \emph{extrinsecamente}, emquanto derivam de Dxzus, --- mas tambem \emph{intrinsecamente}, emquanto possuem em si mesmas a forma de bondade, que é uma semelhança, embora infinitamente remota e deficiente, da propria bondade de \textsc{Deus}; ao passo que ellas se dizem \emph{overdadeiras}, --- não \emph{intrinsecamente}, pois a verdade reside, como no seu proprio sujeito, na \emph{ntelligencia}, e não no \emph{objecto}, --- mas apenas \emph{ewtrinsecamente}, emquanto os seres creados, por isso mesmo que são creados, se referem, e de um modo \emph{necessario}, á intelligencia divina. (Cf. Sum. \emph{Th}., p. I, q. 6, a. 4).}.

51. A \textbf{bondade} divina é a \textbf{causa} final da bondade \textbf{creada}. --- A bondade divina será a causa \emph{final} da bondade creada, se \textsc{Deus}, infundindo a bondade nas creaturas, tiver tido o fim de

\pg{468} ------= EE E q AE Rs ia ie ap

468 METAPHYSICA --- PARTE QUARTA --- THEOD!CEA communicar e manifestar a sua bondade infinita. Ora tal foi, effectivamente, o fim do Creador. Logo a bondade divina é a causa \emph{final} da bondade creada \footnote{Não pode duvidar-se de que Deus creou o mundo com o fim de communicar a outros seres a propria bondade. Porquanto, todo o agente, quando opéra, tem em vista a propria bondade ou perfeição --- para a adquirir, se não a possue, como o agente creado, --- e para a communicar, se a possue, como o agente increado.}.

52. \textbf{Belleza} de \textbf{Deus}. --- \emph{Belleza} de \textsc{Deus} é o attributo, que consiste na infinidade das perfeições divinas, reduzidas a uma unica e simplicissima unidade, e sobremaneira agradaveis a quem as contempla \footnote{Como dissemos na \emph{Ontologia}, \& \emph{belleza é a conveniente unidade na variedade}, \emph{que agrada a quem} a vé: \emph{pulchritudo est multitudinis debitae debita unitas}, \emph{quae visa placet}. Na belleza devemos distinguir dois elementos --- o elemento \emph{objectivo} e o \emph{subjectivo}. O elemento \emph{objectivo}, que constitue propriamente a essencia da belleza, é á \emph{conveniente unidade na variedade}, isto é, a devida multiplicidade reduzida á devida unidade; o elemento \emph{subjectivo}, que não é propriamente \emph{constitutivo} mas \emph{effeito} da belleza, é o \emph{agrado}, que a propria belleza produz em quem a contempla. --- O elemento \emph{objectivo é composto}. \emph{A multiplicidade} ou \emph{variedade} é á \emph{materia} da belleza; a \emph{unidade} é a \emph{forma}, como diz S. Agostinho: « Omnis pulchritudinis \emph{forma} est \emph{unitas} » (\emph{Ep}. 18, al. 68, ad Coelest. n. 2). --- O elemento \emph{subjectivo} indica que a belleza, como a verdade e a bondade, é uma coisa \emph{relativa}. A verdade refere-se á faculdade cognoscitiva, a bondade á appetitiva, a belleza refere-se á cognoscitiva e á appetitiva; porque todos dizem que \emph{pulchrum}, \emph{est quod visum placet}, ora uma coisa \emph{vé-se} pela faculdade cognoscitiva, e \emph{agrada} á appetitiva. --- A belleza é um attributo \emph{transcendental}; porque é proprio de todo o ente, emquanto \emph{tal}. Todo o ente é \emph{multiplo} (real ou logicamente) e \emph{uno}, e a sua belleza é maior ou menor, conforme o maior ou menor numero de partes e conforme a maior ou menor harmonia e união, que existe entre ellas.}.

53. \textbf{Deus} é \textbf{bello}. --- Consistindo a belleza na multiplicidade das perfeicoes devidas, reduzida á conveniente unidade, \textsc{Deus} será \emph{bello}, se possuir todas as perfeicoes, que são proprias do Ente increado, e se todas essas perfeigoes forem reduzidas á devida unidade; de modo que essa unidade em tanta variedade possa agradar a quem a contempla. Ora \textsc{Deus} possue perfeições infinitas, e todas essas perfeigoes constituem a unica e simplicissima Essencia divina, e tal unidade numa variedade infinita não pode deixar de agradar immensamente a quem a contempla. Logo Dszvs é \emph{bello} \footnote{A conclusão é evidente, e prova-se tambem com argumentos \emph{a-posteriori}. O mundo, creado por Deus, é dotado de belleza, e é por isso que foi chamado \emph{cosmos} pelos gregos, e \emph{mundo} pelos latinos. Ora toda a perfeição do effeito deve preexistir na sua Causa efficiente. --- Além d'isso,\pgn{469}a belleza, que se admira nas creaturas, é mutavel e imperfeita. Ora tudo o que é mutavel e imperfeito suppõe e exige uma causa absolutamente immutavel e perfeita. Logo deve existir um Ente immutayel e perfeitamente bello, de que deriva toda a belleza, espalhada nas creaturas.}.

\pg{469} 54. \textbf{Deus} é a \textbf{sua belleza}. --- Todas as perfeigoes divinas são uma e a mesma coisa com a essencia divina; porque \textsc{Deus} é perfeito pela sua essencia. Ora a belleza é uma perfeição. Logo a belleza é uma e a mesma coisa com a essencia divina; e por isso Dzvus, assim como é a sua essencia, assim tambem é a sua \emph{belleza}. --- Além d'isso, consistindo a belleza divina numa infinita variedade de perfeicoes, reduzidas a uma simplieissima unidade, é evidente que \textsc{Deus} será a sua belleza, se fôr a sua perfeição e a sua unidade. Ora \textsc{Deus} é a sua perfeição e a sua unidade; porque tudo o que está em \textsc{Deus}, é \textsc{Deus}. Logo Dxvus é a sua \emph{belleza} \footnote{\textsc{Deus} é a propria belleza! Quem poderá comprehender e exprimir todo o encanto, toda a suavidade, que esta expressão encerra?... A belleza constitue a propria Essencia de Dzvus!... Os escriptores dizem que a belleza de Deus é \emph{summa}, \emph{infinita}, \emph{incorruptivel}, \emph{àmmutavel}, \emph{amabilissima}; mas essas palavras levam-nos ao perfeito conhecimento da belleza divina? Não. Só \textsc{Deus} contempla e comprehende a sua belleza infinita e eterna; só Dxzus ama a sua belleza com amor infinito e eterno; só \textsc{DEus} é infinitamente bemaventurado na contemplação e no amor da sua belleza... A creatura racional vê um pallido reflexo d'esta belleza infinita nas coisas visiveis, e d'esse reflexo, depurado das imperfeições, a intelligencia forma um ideal de belleza, que é superior a toda a belleza creada e representa, ainda que de um modo muito inadequado e deficiente, a ineffavel belleza divina. Esse ideal é sufficiente para alma se destacar de todos os encantos creados e se unir intimamente, pela fé e pelo amor, ao Summo Bem... Isto é muito, mas não é tudo. Deus quiz, na sua infinita misericordia, elevar o homem a uma ordem sobrenatural, que tem por fim a visão \emph{clara} e: \emph{intuitiva} da propria essencia de Deus. \emph{Videbimus Deum sicuti est}! A nossa intelligencia verá, pela graça de Deus, a propria essencia divina como é em si mesma, e por isso infinitamente bella, a propria belleza! Não será uma visão \emph{comprehensiva} (porque o finito não pode comprehender o infinito), mas será verdadeira \emph{visão}. Esta visão transformará a alma e a tornará semelhante a Deus,e será principio de um amor e de uma felicidade sem limites! Felizes os que, conscientes da sua excelsa dignidade e do seu eterno destino, desprezam as bellezas caducas e miseraveis das creaturas, e collocam as suas esperancas e os seus amores na unica e immutavel Belleza divina!}.

55. A \textbf{belleza} divina é a \textbf{causa efficiente} de \textbf{toda} a \textbf{belleza creada}. --- A belleza divina será a causa \emph{efficiente} de toda a belleza creada, se esta derivar de \textsc{Deus}, emquanto e porque Elle é dotado de belleza infinita. Ora é certo que a belleza creada deriva de \textsc{Deus}, emquanto e porque Elle é dotado de belleza infinita; porque ninguem pode communicar a outros uma \pg{470}  perfeição, senão emquanto e porque a possue. Logo a belleza divina é a causa \emph{efficiente} de toda a belleza creada \footnote{De \textsc{DEus}, como da causa \emph{efficiente}, deriva toda a belleza, como toda e qualquer qualidade das creaturas. --- A belleza creada é multipla e diversa, porque muitos e diversos são os entes creados. São bellos os seres \emph{inorganicos}, sobretudo os erystaes e as pedras preciosas; principalmente a \emph{us}, que é esplendida em si e torna esplendidos os outros corpos, e por isso é imagem e symbolo da Belleza divina. --- Mais bellos são os seres \emph{vivos}; sobretudo o homem, em cuja alma espiritual \textsc{Deus} gravou a propria imagem e semelhança. --- Immensamente, ineffavelmente bella é a SS. Humanidade de N. S. Jesus \textsc{CuRISTO}, pela integridade, regularidade e harmonia das suas feições immaeuladas, por aquelle esplendor de magestade e de graça, com que a Divindade, que nella habita substancialmente, a circumda e glorifica. --- Mas de todas essas bellezas a fonte unica é a Belleza divina.}.

56. A \textbf{heileza} divina é a \textbf{causa exemplar} de toda a \textbf{belleza creada}. --- 'Todo o agente tende a imprimir e imprime, de algum modo, no effeito a semelhança da propria forma, por que opera. Ora Deus opera pela sua essencia, que é bella e é a propria belleza. Logo Deus, produzindo as creaturas, imprime nellas a semelhança da sua belleza. Logo a belleza divina é a causa \emph{exemplar} de toda a belleza creada \footnote{Deve haver uma \emph{analogia}, uma \emph{semelhança} entre a belleza divina e a belleza creada, como ha sempre analogia ou semelhança entre a causa e o effeito. --- Mas a \emph{analogia}, \& \emph{semelhança} não é \emph{egualdade}. \textsc{Deus} nao é causa \emph{univoca} das creaturas; por isso a forma das creaturas não pode ser \emph{especificamente egual} á forma do Creador. Portanto a belleza divina é infinitamente mais pura, mais elevada, mais attrahente. Ora, se a belleza das creaturas, embora imperfeita e caduca, é digna de estima e tanto prende o coração, de quanta estima não é digna e quanto não deve ser amada a infinita e purissima belleza de Deus?!}.

57. A \textbf{belleza} divina é a \textbf{causa} final da \textbf{belleza creada}. --- A belleza divina será a causa \emph{final} da belleza creada, se \textsc{Deus} dotou de belleza as creaturas, para que. estas representassem a sua belleza infinita. Ora a intenção de \textsc{Deus}, na producção do mundo, foi effectivamente a de revelar nas visiveis perfeições das creaturas as proprias perfeições invisiveis. Logo a belleza divina é a causa \emph{final} da belleza creada \footnote{Se a belleza divina é a causa \emph{final} da belleza creada, esta belleza creada não pode nem deve considerar-se como o \emph{termo} e o \emph{fim ultimo} das nossas tendencias e dos nossos desejos, mas unicamente como o \emph{meio} e o \emph{degrau} para nos elevarmos á contemplação e ao amor da belleza divina. O espirito, que assim considera a belleza e, em geral, as perfeições das creaturas, não pode deixar de se unir intimamente ao infinito Bem.\par Concluimos este artigo e confirmamos as nossas considerações com as palavras do nosso Angelico Doutor: « Consideratio creaturarum animos hominum in amorem divinae bonitatis accendit. Quidquid enim bonitatis et perfectionis\pgn{471}in diversis creaturis particulariter distributum est, totum in Deo universaliter est adunatum, sicut in fonte totius bonitatis. Si igitur creaturarum bonitas, pulehritudo et suavitas sic animos hominum allicit; ipsius Der fins bonitatis, rivulis bonitatis in singulis creaturis repertis diligenter comparatus, animos hominum inflammatos totaliter ad se trahet » (C. \emph{Gentes}, 1. TT, c. 2).}.

{\small\noindent\textsc{Summario}: --- Intelligencia de Deus, sua natureza e objecto. --- Vontade de Deus, sua natureza e objecto. --- Vida divina, sua existencia e natureza.\par}

\pg{471} FACULDADES E \textbf{OPERAÇÕES} IMMANENTES \textbf{DE DEUS}

\textbf{Intelligencia de Deus}, \textbf{sua natureza} e \textbf{objecto}.

58. \textbf{Faculdades} de \textbf{Deus}. --- \emph{Faculdades} de \textsc{Deus} são as perfeições ou os attributos divinos, que se referem á \emph{operação}. --- Estas faculdades são duas: a \emph{intelligencia} e a \emph{vontade} \footnote{Depois de termos tratado dos attributos \emph{quiescentes} de \textsc{Deus}, devemos agora tratar dos seus attributos \emph{operativos}, e, por isso, das \emph{faculdades} e \emph{operações}. --- Sendo a operação uma consequencia e uma expressão da natureza, nunca aloancaremos um pleno conhecimento de um ente, se não conhecermos convenientemente a sua operação. Por quanto, pelo modo e pela qualidade da operação conhecemos a natureza e a qualidade da virtude ou faculdade operativa, e essa virtude ou faculdade manifesta a essencia do ente, visto que um ente opera de um certo e determinado modo, pelo facto de possuir uma certa e determinada essencia. --- Ora, ha duas. especies de operações; uma que permanece no agente, e é perfeição d'elle; outra que tem o termo fóra do agente, e é perfeição do effeito externo. A primeira diz-se \emph{émmianente}, porque tem o principio e o termo no mesmo agente; a segunda chama-se /\emph{ransewnte}, porque tem o termo fóra do agente. --- A primeira precede a segunda. Por isso, tratando das operações divinas, primeiramente fallaremos das \emph{immanentes}, e depois das \emph{transeuntes}.}.

59. Intelligencia de \textbf{Deus}. --- \emph{Intelligencia} de \textsc{Deus} é a faculdade, pela qual Elle conhece a sua essencia e todas as coisas. --- Da intelligencia procede a \emph{sciencia}. --- Diremos da \emph{existencia} d'esta faculdade, da sua \emph{natureza} e \emph{objecto} \footnote{8. Thomaz trata d'este assumpto sobretudo --- na \emph{Sum}. \emph{Th}., p. 1, q. 14 --- e. \emph{Gentes}, l. I, c. 44 e seg. --- \emph{De Verit}, q. 9 e 9, --- I \emph{Sent}., Dist. 35.}.

60. \textbf{Deus} é dotado de inteligencia.

a) Toda a substancia \emph{espiritual}, ou independente da \emph{materia}, é dotada de \emph{ntelligencia}. Ora \textsc{Deus} é \emph{purissimo espirito}. Logo \textsc{Deus} é dotado de \emph{inteligencia} \footnote{O grau do conhecimento de um ser é sempre proporcionado ao grau da sua immaterialidade; de modo que, quanto mais um ente se afastar\pgn{472}da materia, tanto mais elevado será o seu conhecimento. Assim a alma das plantas, por estar inteiramente ligada á materia, é destituida de todo e qualquer conhecimento. A alma dos animaes, embora não possa operar sem o organismo, todavia recebe as especies \emph{sensiveis} dos objectos, e por isso os animaes são dotados de conhecimento, embora no infimo grau. \emph{A} alma humana, por ser \emph{espiritual}, isto é, independente da materia, é dotada de conhecimento \emph{ntellectwal}, ou de intelligencia. \textsc{Deus}, sendo um espirito \emph{purissimo}, não pode deixar de ser summamente intelligente. Mas escutemos o Angelico Doutor: « Cognoscentia a non cognoscentibus in hoc distinguuntur, quia non cognoscentia nihil habent nisi formam suam tantum; sed cognoscens natum est habere formam (\emph{intentionalem}) etiam rei alterius: nam species cogniti est in cognoscente. Unde manifestum est quod natura rei non cognoscentis est magis coarctata et limitata. Natura autem rerum cognoscentium habet maiorem amplitudinem et extensionem. Propter quod dicit Philosophus (\emph{De anima}, l. 3, t. 91) quod anima est quodammodo omnia (\emph{non realiter}, \emph{sed intentionaliter}). --- Coarctatio autem formae est per materiam. Unde et supra diximus quod formae, secundum quod sunt magis immateriales, secundum hoc magis accedunt ad quamdam infinitatem. Patet igitur quod \emph{Zmmaterialitas} alicujus rei est \emph{ratio} quod sit \emph{cognoscitiva}; et secundum modum immaterialitatis est modus cognitionis. Unde dicitur (\emph{De anima}, l. II, t. 124) quod plantae non cognoscunt propter suam materialitatem. Sensus autem cognoscitivus est, quia receptivus est specierum sine materia; et intellectus adhuc magis cognoscitivus, quia magis separatus a materia, et immixtus; ut dicitur (\emph{De anima}, l. III. t. 4). Unde eum Deus sit in summo immaterialitatis, sequitur quod ipse sit in summo cognitionis » (\emph{Sum}. \emph{Th}. p. I, q. 14, a. 1).}.

\pg{472} b) \textbf{Deus} é infinitamente perfeito, e por isso não pode faltar-Ihe alguma perfeição \emph{absoluta}. Ora a intelligencia é uma perfeição \emph{absoluta}. Logo \textsc{Deus} é dotado de \emph{intelligencia} \footnote{A \emph{intelligencia}, como tambem a \emph{sciencia}, que é o fructo do exercicio d'esta faculdade, é uma perfeição \emph{absoluta}, e por isso convem e attribue-se a Deus em sentido \emph{positivo} e \emph{proprio}, \emph{e} não num sentido \emph{metaphorico}, como Ihe se attribue a \emph{ira} ou a \emph{misericordia}. De facto, a \emph{intelligencia}, e por isso a \emph{sciencia}, não encerra, no seu conceito formal, imperfeição alguma e não se oppõe a uma perfeição mais elevada, e portanto deve attribuir-se a Deus no sentido \emph{proprio}, no mesmo sentido em que se attribue ao homem e ao Anjo; embora em Deus se encontre de um modo infinitamente superior ao modo por que se encontra nas creaturas, pois em Deus é uma perfeição substancial, que se identifica com a essencia, emquanto nas creaturas é uma perfeição accidental, que se distingue da propria essencia. Logo a \emph{intelligencia}; como tambem a \emph{sciencia}, deve attribuir-se a \textsc{Deus}, ainda que se lhe não possa attribuir a imperfeição, de que se acha cercada nas creaturas. Diz S. Thomaz: « Quia perfectiones procedentes a Deo in ereaturas altiori modo sunt in Deo, oportet quod quandocumque nomen sumptum a quaecumque perfectione creaturae Deo attribuitur, secludatur ab ejus significatione omne illud quod pertinet ad imperfectum modum, qui competit creaturae. Unde scientia non est qualitas in Deo, vel habitus, sed substantia et actus purus » (\emph{Sum}. \emph{Th}, p. T, q. 14, a. 1 ad 1).}.

\pg{473} c) \textsc{Deus} estabeleceu a ordem, que resplandece no mundo. Ora esta ordem é effeito de uma causa infinitamente intelligente. Logo \textsc{Deus} é dotado de \emph{intelligencia} \footnote{Não adduzimos mais argumentos em defeza d'esta proposição, porque é uma consequencia do que dissemos ácerca da \emph{existencia} e da \emph{essencia} de \textsc{Deus}. Todos os povos admittem que o Ente supremo recebe as homenagens dos mortaes, ouve as supplicas dos que recorrem a Elle, julga as acções humanas, e recompensa ou condemna conforme a qualidade d'essas acções, e por isso admittem que \textsc{Deus} é dotado de \emph{intelligencia}.}.

61. A intelligencia de \textbf{Deus} identifica-se com a \textbf{sua essencia}, e é \textbf{acto} purissimo e infinito.

a) Identifica-se com a sua essencia. --- \textsc{Deus} é \emph{simplicissimo}. Ora, se a sua intelligencia não se identificasse com a sua essencia, \textsc{Deus} seria \emph{composto} de \emph{essencia} e de \emph{attributos}: o que é falso. Logo a intelligencia de \textsc{Deus} identifica-se com a sua essencia. --- Por isso \textsc{Deus}, assim como é a sua \emph{essencia}, assim tambem é a sua \emph{intelligencia}.

b) E' acto purissimo. --- Se a intelligencia de \textsc{Deus} não fosse \emph{acto purissimo}, mas fosse \emph{potencial}, isto é, passasse da \emph{potencia} para o \emph{acto}, ella, como á nossa, seria \emph{mutavel} e \emph{composta}. Ora a intelligencia de \textsc{Deus}, identificando-se com a sua essencia, é \emph{immutavel} e \emph{simplicissima}. Logo é \emph{acto purissimo}.

c) E' acto infinito. --- A intelligencia de \textsc{Deus} identifica-se com a essencia e é acto purissimo. Ora a essencia de \textsc{Deus} é infinita, porque \textsc{Deus} é o proprio ser; e um acto purissimo é tambem infinito, porque é sempre a potencia que limita o acto. Logo á intelligencia de \textsc{Deus} é acto \emph{infinito} \footnote{A conclusão é evidente. Deus é simplicissimo; e tudo o que está em \textsc{DEus} identifica-se com \textsc{Deus}, é o proprio \textsc{DEus}. Por isso a intelligencia, como tambem o acto de entender, são a propria essencia divina; porque, se fossem, como em nós, uma perfeição distincta, haveria em \textsc{DEus} composição de \emph{potencia} e de \emph{acto}: o que repugna. « Em \textsc{Deus}, diz S. Thomaz, a intelligencia, o intelligente, o objecto entendido, a especie intelligivel, e o acto de entender são uma e a mesma coisa. D'ahi se segue que, embora seja \textsc{Deus} infinitamente intelligente, nem por isso so admitte qualquer multiplicidade na sua substancia » (\emph{Sum}. \emph{Th}, p. I, q. 14, a. 4).}.

62. \textbf{Deus} conhece tudo pela sua essencia.

a) Se \textsc{Deus} conhecesse as coisas por uma especie intelligivel, distincta da sua essencia, essa especie seria o \emph{principio formal} do conhecimento divino, e \textsc{Deus} seria \emph{dependente} de um outro ente na operação e por isso no ser; porque a operação e o ser, em \textsc{Deus}, são uma e a mesma coisa, Ora \textsc{Deus} é a Causa primeira, absolutamente independente de todo e qualquer ente. Logo \textsc{Deus} conhece tudo pela sua \emph{essencia}.

\pg{474} b) \textsc{Deus} conhecerá tudo pela sua \emph{essencia}, se esta representar \emph{adequadamente} tudo o que Elle entende --- o seu proprio ser e os entes finitos, existentes e possiveis; porque, então, não ha razão sufficiente para procurar outra especie intelligivel. Ora a essencia divina representa \emph{adequadamente} --- o ser divino, porque, em \textsc{Deus}, a essencia e o ser são uma e a mesma coisa, --- e todo o ente finito, existente e possivel, porque a essencia divina é a causa efficiente e exemplar de todas as coisas. Logo \textsc{Deus} conhece tudo pela sua \emph{essencia}.

c) A especie intelligivel, principio formal do conhecimento, refere-se á intelligencia, como o acto á potencia. Portanto essa especie, se fosse, em \textsc{Deus}, distineta da essencia, seria \emph{acto} da. intelligencia divina, a qual, por consequencia, seria \emph{potencia}. Ora não pode haver, em Deus, composição de acto e de potencia. Logo Deus conhece tudo pela sua \emph{essencia} \footnote{Cf. c. \emph{Gentes}, 1. T, c. 41. --- Para a intelligencia do assumpto, lembramos certos principios. O conhecimento consiste nisto que o cognoscente se torna semelhante ao conhecido. Este une-se áquelle por uma forma ou imagem intencional, que se chama \emph{especie intelligivel}, e que actua na intelligencia e lhe representa o objecto. Por isso a perfeição do conhecimento depende inteiramente da perfeição d'esta forma ou especie, a qual, pela sua vez, deve ser proporcionada á virtude ou energia intellectual, que ella determina ao acto. D'ahi se segue --- que, quanto mais perfeita fôr a intelligencia, tanto mais excellente deve ser a especie intelligivel, --- e que, quanto mais excellente fôr a especie, tanto mais perfeita será a sua maneira de representar. Portanto, todo o conhecimento intellectual exige e suppõe tres coisas: a intelligencia cognoscente, o objecto conhecido, e a especie intelligivel. O objecto conhecido chama-se objecto \emph{material} ou \emph{quasi-material} do conhecimento. A razão, pela qual o objecto \emph{material} determina a intelligencia e que é representada ou expressa pela especie intelligivel, diz-se objecto \emph{formal}. Sendo assim, é claro que o objecto \emph{formal}, quando se trata das faculdades \emph{cognoscitivas}, coincide de algum modo com a \emph{especie intelligivel}; porque o objecto \emph{material} não actua na intelligencia senão pela especie e pelo que está representado na especie. Por isso, quando se trata do objecto \emph{formal} de uma intelligencia, procuramos e indicamos a \emph{especie intelligivel}, que actua naquella intelligencia e lhe torna as coisas presentes sob um certo e determinado aspecto. --- O objecto \emph{material} de toda e qualquer intelligencia divide-se em \emph{primario} e \emph{secundario}. \emph{Primario} é o que é proporcionado á intelligencia, é conhecido por si mesmo e por uma especie adequada, e esse conhecimento é a razão pela qual são conhecidas todas as outras coisas; assim o objecto \emph{primario} da nossa intelligencia, no estado presente, é a essencia das coisas materiaes. \emph{Secundario} é o que não é proporcionado á inteligencia, não é conhecido por si mesmo, mas por outra coisa e por uma especie inadequada; assim o objecto \emph{secundario} da nossa intelligencia é constituido por tudo o que podemos conhecer por meio das especies, derivadas das coisas sensiveis ou materiaes.\par Applieando estas noções á intelligencia divina, é evidente que o seu\pgn{475}objecto \emph{formal} não pode ser uma forma ou especie intelligivel, distincta da essencia divina; porque, como dissemos no texto, a intelligencia divina, sendo acto purissimo e infinitamente perfeito, não pode ser aperfeiçoada e completada, em ordem ao acto, por uma especie, distincta d'ella e da essencia. --- Emquanto ao objecto \emph{material}, a inteligencia divina, como vamos provar, comprehende o proprio ser divino e todas as outras coisas, que podem ser de algum modo conhecidas; porque a essencia divina, sendo forma perfeitissima e de infinito valor, representa tudo á intelligencia divina, com a qual é uma e a mesma coisa.\par Portanto, se nos perguntarem se para o acto da intelligencia divina é necessaria a \emph{especie intelligivel}, respondemos que, se a especie \emph{intelligivel} se toma no sentido rigoroso emquanto significa «a semelhança de uma coisa, a qual, á maneira de um principio subjectivo, realmente distincto da intelligencia, completa e determina esta faculdade, por si mesma incompleta e indifferente, ao conhecimento de um determinado objecto v, então devemos dizer que ella não é necessaria nem pode existir \emph{formalmente} em Deus, porque não pode admittir-se que a intelligencia divina seja por si incompleta e deficiente com relação ao seu acto; --- mas, se a \emph{especie intelligivel} se toma num sentido \emph{lato}, emquanto significa ua razão sufficiente pela qual a intelligencia conhece um objecto, isto é, emquanto é a semelhança da coisa, que torna a intelligencia capaz de conhecer a mesma coisa », ella é necessaria para o acto da intelligencia divina e existe \emph{formalmente} um \textsc{DEUS}, porque, se \textsc{DEus} conhece, deve existir nelle a razão sufficiente do seu conhecimento. --- Esta especie, tomada no sentido \emph{lato} e existente em \textsc{Deus}, não é uma coisa \emph{realmente} distincta da essencia de Deus, nem da sua intelligencia, nem do seu acto intellectual, mas identifica-se com cada uma d'estas coisas, como ensina a cada passo S. Thomaz. --- Dissemos que a especie, em \textsc{Deus}, não se distingue \emph{realmente} da \emph{essencia}, nem da \emph{intelligencia}, nem do \emph{acto intellectual}; porque, \emph{logicamente}, distingue-se da \emph{intelligencia} e do \emph{acto intellectual}. Porquanto, a \emph{especie} considera-se como um principio, que \emph{torna} a \emph{intelligencia capaz} de perceber um certo objecto, e por isso \emph{logicamente} distingue-se da intelligencia e do seu acto. --- Mas nem \emph{logicamente} se distingue da \emph{essencia divina}. Por quanto, a especie deve \emph{conter em si} o \emph{objecto cognoscivel}, para que possa determinar a intelligencia á percepção d'elle; ora a essencia divina, contendo em si todas as coisas, pode e deve exercer o munus da dita especie. Diz S. Thomaz: « Species intelligibilis in Dro nihil aliud est quam \emph{essentía sua}, in qua omnes species rerum comprehenduntur » (\emph{Sum}. \emph{Th}., p. I, q. 14, a. 5 ad 3).\par Da doutrina exposta resultam os seguintes corollarios: --- 1.9) \emph{Unica} é a especie intelligivel em \textsc{Deus}, porque unica é a sua essencia, capaz de determinar por si só \& intelligencia divina ao conhecimento de todas as coisas; --- 2.º) Pela mesma unica especie, que é a sua essencia, \textsc{DEus} conhece o proprio ser e todas as coisas, --- o proprio ser pela essencia, que é a mesma coisa com o mesmo ser, --- todas as outras coisas pela essencia, que é uma\pgn{476}semelhança d'essas coisas; --- 8.9) o conhecimento divino não é \emph{especificado} pela especie propria das outras coisas, mas unicamente pela divina essencia, e por isso é sempre infinito (\emph{Sum}. \emph{T}/., p. I, q. 14, a. 5 ad 3).}.

\pg{475} 63. \textbf{Deus conhece perfeitamente} o \textbf{proprio ser}.

a) \textbf{Conhece} o proprio ser, --- A intelligencia de \textsc{Deus}, sendo um acto purissimo e de infinita energia, deve conhecer tudo o que é intelligivel e lhe está presente. Ora o ser divino, não só é intelligivel, porque é immaterial, mas tambem está intimamente \pg{476}  presente á intelligencia, porque é uma e a mesma coisa com ella. Logo \textsc{Deus} conhece o proprio ser \footnote{A conclusão é evidente, e é um corollario, que deriva da natureza do conhecimento intellectual. Este, por ser immaterial e inorganico, é essencialmente capaz de perfeita reflexão, pela qual se conhece a existencia e a natureza do cognoscente. Deveriamos negar á intelligencia divina uma perfeição, que attribuimos á intelligencia humana? Nào, com certeza.}.

b) Conhece perfeitamente o proprio ser. --- A perfeição do conhecimento, que consiste na completa adequação entre a cognoscibilidade do objecto e o proprio acto cognoscitivo, depende de duas condições: que a especie intelligivel --- reproduza fielmente o objecto, --- e se una intimamente com a faculdade. Ora a essencia divina, que faz as vezes da especie intelligivel no conhecimento de \textsc{Deus}, não só exprime fielmente o ser divino e se une intimamente com a intelligencia divina, mas identifica-se com o mesmo ser e com a mesma intelligencia. Logo \textsc{Deus} conhece \emph{perfeitamente} o proprio ser \footnote{Podemos expor o argumento numa outra forma. Para a percepção intellectual são necessarios tres elementos: a faculdade intellectual, o objecto intelligivel, \emph{a} união do objecto com a faculdade; e quanto mais perfeitos forem estes elementos, tanto mais perfeita será a percepção. Ora a intelligencia divina ésinfinitamente perspicaz, --- o ser divino é infinitamente intelligivel, --- a união do ser com a intelligencia não poderia ser mais intima, porque são uma e a mesma coisa. Logo a percepção que \textsc{Deus} tem do proprio ser é infinitamente perfeita. --- Esta percepção chama-se tambem \emph{comprehensão}; porque \emph{comprehender} um objecto é conhecel-o tanto, quanto elle é cognoscivel, de modo que á intelligencia não escape nenhum elemento, nenhuma propriedade, nenhuma efficacia, de que elle é dotado. Ora o conhecimento de Deus é proporcionado á sua cognoscibilidade, isto é, \textsc{DEUS} se conhece a si mesmo tanto, quanto Elle é cognoscivel; porque esse conhecimento e essa cognoscibilidade fundam--se na mesma purissima actualidade do ser divino. Logo \textsc{Deus} \emph{comprehende} o proprio ser. Diz S. Thomaz: « Tunc dicitur aliquid \emph{comprehendi}, quando res cognoscitur ita perfecte, sicut cognoscibilis est... Manifestum est autem quod \textsc{Drus} ita perfecte cognoscit seipsum, sicut perfecte cognoscibilis est... 'lanta est autem virtus DEI in cognoscendo, quanta est actualitas ejus in existendo; quia per hoc quod actu est et ab omni materia et potentia separatus, \textsc{Dkvs} cognoscitivus est. Unde manifestum est quod tantum seipsum cognoscit, quantum cognoscibilis est. Et propter hoc seipsum perfecte \emph{comprehendit} » (\emph{Sum}. \emph{Th}., p. T, q. 14, a. 8).\par Que \textsc{Deus} conheça o proprio ser, como objecto \emph{primario} da sua intelligencia, é fóra de toda a duvida. Objecto \emph{primario} da intelligencia é, como dissemos, o que é proporcionado á intelligencia, é conhecido por si mesmo e por uma especie adequada, e esse conhecimento é a razão pela qual são conhecidas as outras coisas. Ora o ser divino --- 6 o unico objecto proporcionado\pgn{477} á intelligencia divina, porque só esse ser é infinito como infinita é essa intelligencia, --- só elle é conhecido por si mesmo e por uma especie \emph{adequada}, porque só elle está intimamente unido com a intelligencia divina e só elle é representado por uma especie \emph{adequada}, que é a divina essencia, --- e o conhecimento do ser divino é a razão pela qual são conhecidas todas as coisas, porque \textsc{Deus} conhece tudo na sua essencia e pela sua essencia. \textbar{} Logo o ser divino é o objecto \emph{primario} da intelligencia divina.}.

\pg{477} ------------------ ------------------------===--- "

%%>>> cap-7

\chapter*{CAPITULO II\\[0.3em]\large FACULDADES E OPERAÇÕES IMMANENTES DE DEUS 477}
\addcontentsline{toc}{chapter}{CAPITULO II --- FACULDADES E OPERAÇÕES IMMANENTES DE DEUS 477}

64. \textbf{Deus} conhece \textbf{todas as} outras \textbf{coisas}, \textbf{existentes} e \textbf{possiveis}.

4) A intelligencia de Deus, identifieando-se com o proprio ser d'Elle, é infinitamente perfeita. Ora uma intelligencia infinitamente perfeita não pode deixar de conhecer tudo o que é cognoscivel, isto é, tudo o que tem alguma razão de entidade, quer exista quer não exista; porque, se uma coisa só, existente ou possivel, lhe eseapasse, ella seria necessariamente limitada e imperfeita. Logo \textsc{Deus} conhece todas as outras coisas, existentes e possiveis.

b) \textsc{Deus} conhece \emph{perfeitamente} o proprio ser, isto é, conhece-o tanto, quanto é cognoscivel. Ora \textsc{Deus} não conheceria perfeitamente o proprio ser, se não conhecesse tambem todas as outras coisas, existentes e possiveis. Por quanto, um ente não conhece perfeitamente o proprio ser, se não conhece perfeitamente a efficacia e a extensão do seu poder, --- e não conhece perfeitamente a efficacia e a extensão do seu poder, se não conhece todos os effeitos, que esse poder produz ou pode produzir. Mas o poder de \textsc{Deus} extende-se a todas as coisas, existentes e possiveis; porque Elle é a primeira causa efficiente de tudo o que existe ou pode existir. Logo \textsc{Deus} conhece todas as outras coisas, existentes e possiveis \footnote{E' o argumento de S. Thomaz, que diz: « Respondeo dicendum quod necesse est Deum cognoscere alia a se. Manifestum est enim quod seipsum perfecte intelligit; alioquin suum esse non esset perfectum, cum suum esse sit suum intelligere. Si autem perfecte aliquid cognoscitur, necesse est quod virtus ejus perfecte cognoscatur. Virtus autem alicujus rei perfecte cognosci non potest, nisi cognoscantur ea, ad quae virtus se extendit. Unde, cum virtus divina se extendat ad alia, eo quod ipsa est prima Causa effectiva omnium entium, \textbar{} necesse est quod Deus alia a se cognoscat » (\emph{Sum}. \emph{Th}., p. I, q. 14, a. B). \textbar{}\par Do argumento conclue-se --- que \textsc{Deus} conhece todas as outras coisas, existentes e possiveis, --- \emph{não immediatamente}, em \emph{si mesmas}, --- mas em \emph{si mesmo}, isto é, na propria \emph{essencia} como num \emph{meio conhecido}, --- e que por isso todas as coisas constituem o objecto \emph{secundario} da intelligencia divina. Na verdade, o que se conhece immediatamente em si mesmo, e por \textbar{} uma especie adequada ao proprio cognoscivel, é o objecto \emph{primario}; em- \textbar{} quanto o que se conhece mediatamente num outro ente, e por uma especie \textbar{} propria d'esse outro ente, é o objecto \emph{secundario}. Ora Deus conhece todas \textbar{} as outras coisas, existentes e possiveis, pela sua essencia, isto é, por uma \textbar{}\pgn{478}especie, que é adequada unicamente ao seu ser mas excedo immensamente todas as outras coisas, nella \emph{virtualmente contidas}, e por isso conhece o proprio ser como objecto \emph{primario}, e todas as outras coisas como objecto \emph{secundario}. --- E não pode ser d' outro modo. A intelligencia divina não entende as coisas, senão pela sua essencia, que exerce o munus de especie intelligivel. Ora a essencia determina necessariamente a intelligenecia a conhecer \emph{immediata} e \emph{directamente} a mesma essencia, e \emph{mediata} e \emph{indirectamente} as outras coisas, que estão nella virtualmente contidas. Por quanto, \textsc{Deus} só conhece as outras coisas, porque conhece a sua essencia, --- não emquanto é uma essencia, --- mas emquanto pode ser e é diversamente participada e imitada por muitas coisas. Diz S. Thomaz (1. e.): « Dupliciter aliquid cognoscitur: uno modo \emph{in seipso}, alio modo \emph{tn altero}. In seipso quidem cognoscitur aliquid, quando cognoscitur per \emph{speciem}. \emph{adaequatam ipsi cognoscibili}; sicut cum oculus videt hominem per speciem hominis. In alio autem videtur id quod videtur \emph{per speciem continentis}; sicut cum pars videtur in toto per speciem totius, vel cum homo videtur in speculo per speciem speculi, vel quocumque alio modo contingat aliquid in alio videri. Bie igitur dicendum est, quod Deus seipsum videt in seipso, quia seipsum videt per essentiam suam. Alia autem a se videt non in seipsis, sed in se ipso, inquantum essentia sua \emph{continet} similitudinem aliorum ab ipso ».\par Como apparece das noções expostas, o que é objecto \emph{primario} para a intelligencia humana é \emph{secundario} para a intelligencia divina, e o que é \emph{secundario} para nós é \emph{primario} para \textsc{Deus}. Na verdade, o objecto \emph{primario} da nossa inteligencia, no estado presente, é a essencia abstracta das coisas materiaes; porque esta essencia --- é proporcionada á natureza e ao modo de ser da nossa intelligencia, --- 6 conhecida por uma especie adequada, --- e é a razão, pela qual conhecemos todas as outras coisas, e o proprio Deus, que, por isso, é objecto \emph{secundario} da nossa inteligencia. Ma o objecto \emph{primario} da intelligencia divina é a propria essencia divina; porque só ella --- é proporcionada á natureza e ao modo de ser da intelligencia divina, --- é conhecida por uma especie adequada, --- e é a razão pela qual \textsc{Deus} conhece todas as outras coisas, que, por isso, constituem o objecto \emph{secundario} da intelligencia divina. Nós não podemos conhecer a \textsc{Deus}, senão depois de conhecermos as creaturas; \textsc{Deus} não pode conhecer as creaturas, senão \emph{depois} de se ter conhecido a si mesmo. --- Por isso a essencia divina é o meio \emph{in quo} \textsc{Deus} percebe todas as coisas; porque, como dissemos, esse meio \emph{primeira} ou \emph{directamente} é percebido em si mesmo e \emph{depois} leva á percepção de outro objecto, e portanto é, ao mesmo tempo, \emph{objecto} e \emph{meio}. --- E claro que, no caso presente, não se trata de prioridade de \emph{tempo}, mas de \emph{natureza}; porque \textsc{Deus}, com um só acto, que é a sua substancia, vê tudo numa unica coisa, que é a sua essencia (\emph{Sum}. \emph{Th}., p. I, q. 14, a. 7).\par Do que temos dito resulta qué a sciencia divina, emquanto aos objectos entendidos, pode e deve dividir-se em sciencia de \emph{visão} e de \emph{simples intelligencia}. --- Sciencia \emph{de visão} é aquella, pela qual \textsc{Deus} conhece o proprio ser e todas as creaturas \emph{existentes} numa determinada differenca de tempo, --- e por isso passadas, presentes e futuras, Sciencia de \emph{simples intelligencia}\pgn{479}é aquella, pela qual \textsc{DEus} conhece as coisas meramente \emph{possiveis}, que não existem e nunca hão de existir. --- E' a divisão dada por S. Thomaz (\emph{Sum}. \emph{Th}., p. I, q. 14, a. 9), e admittida pela eschola thomista, que a proclama \emph{adequada}, porque entre a sciencia de \emph{visão} e a de \emph{simples intelligencia} nao pode conceber-se outra, sendo \emph{adequada} a divisão do objecto. --- Outros theologos, como veremos, admittem uma terceira especie de sciencia, que chamam \emph{media} e que teria por objecto os \emph{futuros condicionados livres}.\par Para a solução das difficuldades contra a doutrina exposta, indicamos os seguintes principios. --- 19) O objecto conhecido, representado na especie intelligivel, aperfeicóa o cognoscente --- não emquanto ao ser \emph{real} d'elle, --- mas emquanto ao ser \emph{intencional}; porque a especie intelligivel é acto ou principio formal da intelligencia, e todo o acto ou toda a forma aperfeiçõa a potencia. Ora a especie, que se une com a intelligencia divina, não é uma coisa distincta de \textsc{DEus}, mas é a sua propria essencia; e por isso só \textsc{Deus} é a perfeição de si mesmo. (\emph{Sum}. \emph{T}., p. I, q. 14, a. 5 ad 4) --- 209) O meio, pelo qual se conhece uma coisa, deve ser \emph{proporcionado} á propria coisa --- não emquanto ao \emph{ser}, como se o meio e a coisa deyessem ter o mesmo ser real, --- mas emquanto á \emph{representação}, de modo que o meio represente a coisa. Ora a essencia divina, embora exceda infinitamente o ser das creaturas, representa-as perfeitamente, e por isso é um meio proporcionado para o conhecimento d'ellas. (\emph{De} Ver., q. 2, a. 8 ad 4). --- 8º) O conhecimento de \textsc{Deus} pode chamar-se \emph{multiplice} e \emph{diverso} --- emquanto ao \emph{objecto conhecido}, que é multiplice e diverso, --- mas não emquanto ao \emph{principio} ou ao \emph{meio} de conhecimento, que é unico, porque é a propria essencia divina. A \emph{unidade} do conhecimento, se se considera no sujeito cognoscente, depende da \emph{unidade} da forma intelligivel, que, em \textsc{Deus}, é a propria essencia divina. Esta, pela sua perfeição infinita, representa tudo o que tem razão de \emph{ente}, e d'ella deriva a multiplicidade e diversidade dos objectos entendidos. --- 49) A sciencia de \textsc{DEus} é, como veremos, a \emph{causa} das coisas, emquanto \textsc{Deus} produz as creaturas em conformidade com as ideas exemplares da sua intelligencia; mas é causa juntamente com a vontade e com o poder. Por isso, embora \textsc{Deus} conheça as coisas desde a eternidade, nao se segue d'ahi que desde a eternidade existam as coisas. O acto da vontade divina com relação ás creaturas é necessariamente livre, e por isso \textsc{Deus} pode querer e deixar de querer a existencia das creaturas, e pode querel-a nesta ou naquella differenga de tempo. (\emph{Sum Th}., p. I, q. 14, a. 8) --- 59) A infinita perfeição de \textsc{Deus} exige que Elle --- não só conheça tudo o que, de algum modo, é, --- mas tambem que o conheca de um modo perfeitissimo, proporcionado á excellencia do ser divino. Ora é mais perfeito que Deus conheça as coisas na divina essencia, do que na essencia d'ellas. --- Nem por isso este conhecimento é indeterminado ou imperfeito; porque, como diremos, a essencia divina exprime \emph{distincta} e \emph{adequadamente} todas as creaturas, e por isso representa não só os elementos mais \emph{comwumuns} d'ellas, mas tambem os mais \emph{proprios}, até ás ultimas differencas (\emph{Sum}. \emph{T}/., p. I, q. 19, a. 6).\par Tendo provado que \textsc{Deus}, cuja intelligencia e sciencia são uma e a mesma coisa com a substancia, conhece na sua essencia tudo o que é de\pgn{480}algum modo cognoscivel, não seria absolutamente necessaria \& questão relativa a alguns entes particulares --- 8e e como \textsc{DEUS} OS conhece. Todavia; como a nossa intelligencia limitada encontra muita difficuldade em conceber o modo por que a intelligencia divina conhece este ou aquelle ente, julgamos util tratar este assumpto. --- Os objectos, que são O termo do conhecimento divino, podem considerar-se em ordem á \emph{essencia}, ou em ordem á \emph{existencia}. Se se consideram em ordem á \emph{essencia}, podemos perguntar: se e como \textsc{DEUS} conhece os entes \emph{singulares} na sua \emph{materialidade}, e os \emph{males}. Se se consideram em ordem á \emph{existencia}, podemos perguntar: se e como \textsc{Deus} conhece os \emph{futuros necessarios}, os \emph{futuros livres}, quer \emph{absolutos} quer \emph{condicionaes}, e todos os \emph{possiveis}. Vamos responder a cada uma d'estas perguntas.\par a) \textsc{Deus} \emph{conhece os}, \emph{entes singulares na sua materialidade}. \textsc{DEUS} deve conhecer tudo o que Elle produz pela sua arte. Ora a arte divina não \emph{so} produz a \emph{forma}, pela qual o ente é collocado numa certa \emph{especie}, mas tambem produz á \emph{materia}, que é principio de \emph{individuação}. Logo \textsc{Deus} conhece todos OS entes \emph{materiaes}, tanto com relação á \emph{materia}, quanto com relação \& \emph{forma}, e por isso conhece os entes \emph{singulares} na sua \emph{materialidade}. (Cf. \emph{De} Ver, d. 9 a 5) Alem d'isso, é uma perfeição para a nossa intelligencia o conhecimento dos entes \emph{singulares} e \emph{materiaes}. Ora toda \& perfeição da creatura deve encontrar-se, num grau mais elevado, no Creador. Logo \textsc{Deus} conhece os entes \emph{singulares} e \emph{materiaes} de um modo mais perfeito, do que nós. Nós conhecemos, por meio da \emph{intelligencia}, O \emph{universal} e o \emph{immaterial}, e, por meio dos \emph{sentidos}, o \emph{singular} e o \emph{material}; \textsc{DEus} conhece o \emph{universal} e O \emph{singular}, o \emph{immaterial} e o \emph{material} pela mesma faculdade intellectual. --- Diz S. Thomaz: « Cum Deus sit causa rerum Per suam scientiam, intantum se extendit scientia Dei, inquantum 8e extendit ejus causalitas. Unde, cum virtus activa Dei se extendat non solum ad formas, a quibus accipitur ratio universalis, sed etiam usque ad materiam, necesse est quod scientia Dei usque ad singularia se extendat, quae per materiam individuantur. Cum enim sciat alia a se per essentiam suam, inquantum\par est similitudo rerum; velut principium activum earum, necesse est quod\par essentia sua sit principium sufficiens cognoscendi omnia, quae per ipsum\par fiunt, non solum in universali, sed etiam in singulari. Et esset simile de\par scientia artificis, si esset productiva totius rei, et non formae tantum 7?\par (Sum. \emph{Th}., p. I, q. 14, a. 11).\par Para a solução das dificultades, notamos com Sg. Thomaz (\emph{art}. \emph{cit}.): ---\par 19) \emph{A}. nossa intelligencia abstrahe dos principios individuantes a especie\par intelligivel. Por isso 2 especie intelligivel da nossa intelligencia não pode\par representar OS principios individuantes, e não podemos ter um conheoci-\par mento \emph{directo} dos entes singulares e materiaes. Mas a especie intelligivel\par da intelligencia divina, que é \& propria essencia de \textsc{DEus}, não é immate-\par rial por abstracção, mas por si mesma, e é o principio de todos os ele-\par mentos, quer especificos, quer individuaes, que entram mna composição da\par coisa. Por isso \textsc{DEUS} conhece, pela sua essencia, não 80.0 universal, mas\par tambem o singular. --- 90) Ainda que á especie da intelligencia divina não\par contenha as condições materiaes, como á especie recebida nos nossos Sen-\par tidos, comtudo pode representar e representa O immaterial e O material.\pgn{481}Para a perfeição do conhecimento não se requer que a coisa se encontre no cognoscente no seu modo de ser; mas basta que a coisa se conheça como ella é. --- 3º) A materia, embora, por ser pura potencia, seja muito dessemelhante de \textsc{Deus}, que ê puro acto, comtudo, emquanto é, tem uma certa semelhança com o \emph{ser} divino.\par b) \textsc{Deus} \emph{conhece os males}. --- Quem entende perfeitamente uma coisa, deve tambem entender tudo o que a essa coisa pode acontecer. Ora ha \emph{bens}, que podem ser corrumpidos pelos \emph{males}. Logo \textsc{DEUS} não entenderia perfeitamente os \emph{bens}, se não entendesse tambem os \emph{ales}... Mas uma coisa emtanto é cognoscivel, emquanto é. Ora, consistindo o mal na \emph{privação} do bem, e a \emph{privação} sendo \emph{nada}, \emph{o} mal não se conhece em si mesmo, mas no bem, de que é privação; e por isso \textsc{Deus}, pelo facto de conhecer os bens, conhece tambem os males, como pela luz se conhecem as trevas. --- Notamos --- 1º) a sciencia de \textsc{Deus} não é causa do mal, mas é causa do bem, pelo qual conhecemos o mal; --- 2º) o conhecimento de uma coisa por meio de outra é imperfeito, quando aquella coisa é cognoscivel por si. Ora o mal não é cognoscivel por si, porque consiste essencialmente na privação do bem, e por isso o mal não pode ser definido, nem conhecido, senão por meio do bem. Diga-se o mesmo com relação ao \emph{peccado}, e, em geral, a todas as \emph{privações} (\emph{Sum}. \emph{Th}., p. I, q. 14, a. 10).\par c) \textsc{DEus} \emph{conhece os futuros contingentes absolutos}. --- \emph{Futuro} é o ente que actualmente não existe, mas que existirá com certeza numa certa differença de tempo. Distingue-se do \emph{possivel}, que nunca existirá. --- Divide-se o \emph{futuro} em \emph{necessario} e \emph{contingente}. \emph{Necessario} é o que deriva de uma causa \emph{necessaria}, isto é, que não pode deixar de produzir um certo effeito. \emph{Contingente} é o que deriva de uma causa \emph{contingente}, isto é, que pode produzir ou deixar de produzir um certo effeito. --- Se a causa \emph{contingente} é \emph{livre}, o effeito futuro chama-se \emph{contingente livre}. --- O futuro \emph{contingente} subdivide-so em \emph{absoluto} e \emph{condicionado}. O contingente \emph{absoluto} é o que será certamente produzido pela causa contingente. O contingente \emph{condicionado} é o que nunca será produzido pela causa contingente, mas seria por ella produzido, se se verificasse esta ou aquella condição. --- O \emph{futuro contingente condicionado} costuma chamar-se \emph{futurivel}. --- Feitas estas distineções, provamos a nossa proposição com as palavras de S. Thomaz (\emph{Sum}. \emph{Th}., p. I, q. 14, a. 18).\par O \emph{futuro contingente absoluto} pode considerar-se sob dois aspectos --- em \emph{si mesmo} e \emph{na sua causa}. Considerado \emph{em si mesmo}, este contingente não é \emph{futuro}, mas é \emph{presente}, --- nem é \emph{indifferente} para a existencia e para a não existencia, mas é \emph{determinado} para a existencia, e, sob este aspecto, pode ser objecto de um conhecimento certo e infallivel. Considerado na sua causa, este contingente é \emph{futuro} e é \emph{indifferente} para a existencia e para a não existencia, porque a \emph{causa contingente} é indifferente com relação aos dois termos oppostos; e por isso o mesmo contingente não pode ser conlhecido com certeza. Portanto quem conhece o effeito contingente só na respectivaà causa, só pode ter d'elle um conhecimento conjectural. \textsc{Deus}, porém, conhece todos os contingentes, não só emquanto estão contidos nas proprias causas, mas tambem emquanto subsistem actualmente em si mesmos.\par 81 Vor. II.\pgn{482}E, ainda que os contingentes recebam a existencia successivamente, todavia \textsc{Deus} não os conhece successivamente, como nós, mas conhece-os todos \emph{ao mesmo tempo}; porque o seu conhecimento é eterno, como eterno é O seu ser, o a eternidade contem todo o tempo e corresponde a todas as differencas do tempo. Por isso, todas as coisas, que existem no tempo, estão presentes, desde a eternidade, a \textsc{Deus}; não só porque \textsc{Deus} possue as idéas exemplares de todas as coisas, mas tambem porque o seu olhar vé, desde a eternidade, todas as coisas, como ellas existem em si mesmas. E' pois, manifesto que os contingentes são infallivelmente conhecidos por \textsc{Deus}, emquanto estão presentes aos olhos divinos, e comtudo são futuros contingentes, se se compararem com as suas causas proximas.\par \emph{A certeza} da sciencia divina com relação aos futuros contingentes \emph{absolutos} não destroe a \emph{liberdade} dos nossos actos. Na verdade, o conhecimento de \textsc{Deus} é verdadeiro. Ora o conhecimento, para que seja verdadeiro, deve conformar-se com o seu objecto; de modo que se limite a descobrir a natureza do objecto, sem alteral-a. Ora a natureza do futuro \emph{contingente} consiste em que elle possa existir ou deixar de existir. Logo esta indifferenca não pode ser destruida ou alterada pela sciencia divina. --- Além d'isso, a sciencia divina refere-se aos futuros contingentes, como a nossa vista se refere aos objectos presentes. Ora é claro que os objectos, que nos estão presentes, não perdem a contingencia, e não se tornam necessarios pelo facto de serem vistos por nós, mesmo de um modo infallivel. Logo nem os futuros contingentes deixam do ser contingentes pelo facto de serem conhecidos infallivelmente por \textsc{Deus}. Diz S. Thomaz: « Sicut.... noster visus non fallitur unquam, videns contingentia ut sunt praesentia, et tamen ex hoc non removetur, quin illa \emph{contingenter} eveniant; ita \textsc{Deus} infallibiliter videt omnia contingentia, sive quae nobis sunt praesentia, sive quae praeterita, sive quae futura, quia sibi non sunt futura, sed ea inspieit esse tune quando sunt; unde ex hoc non removetur quin \emph{contingenter} eveniant» (\emph{De} Ver., q. 2, a. 12) --- Em poucas palavras: o acto \emph{contingente}, que \textsc{Deus} prevé, realizar-se-ha \emph{infallivelmente}, mas não \emph{necessariamente}. Se-a \emph{necessidade} tira a liberdade, não a tira a \emph{infallibilidade}.\par Objectam: Tudo o que \textsc{Deus} conhece desde a eternidade, como existente, é um ente \emph{necessario}. Ora nenhum futuro \emph{contingente é necessario}. Logo \textsc{Deus} não conhece os futuros contingentes.\par Resposta. Tudo o que \textsc{DEus} conhece desde a eternidade, como existente, quer seja ente \emph{necessario} com relação ás suas causas proximas, quer seja \emph{contingente}, existe \emph{necessariamente}, porque é conhecido por \textsc{Deus} como presente, e não pode deixar de existir o que \textsc{Deus} conhece como existente; aliás haveria engano na sciencia divina. Mas a \emph{necessidade} de existencia, que se attribue a um effeito, é duplice: uma \emph{absoluta}, quando o effeito deriva de causas necessarias; --- outra \emph{condicionada}, quando o effeito deriva de causas contingentes, porque tudo o que é, emquanto á, não \textbf{pode} deixar de ser. Ora esta necessidade \emph{condicionada} deve attribuir-se aos effeitos \emph{contingentes}, cuja existencia é vista por \textsc{Deus}; assim como é \textbf{necessario} que Socrates\pgn{483} escreva, quando escreve, e nós vemos que escreve. Portanto, se com o futuro contingente não é compativel a necessidade \emph{absoluta}, é certamente compativel a necessidade \emph{condicionada}. (\emph{Of}. \emph{S}. \emph{Th}., p. I, q. 14, a. 8 ad 8; \emph{De} Ver. q: 2, a. 12; c. \emph{Gentes}; 1- 1, o. 60).\par d) \textsc{Deus} \emph{conhece} os \emph{futuros contingentes condicionados}. --- A questão limita-se aos \emph{contingentes condicionados livres}. --- Que \textsc{Deus} conheça estes \emph{futuros condicionados livres}, todos os philosophos espiritualistas o admittem. Esses futuros encerram ou constituem uma verdade, uma entidade; e tudo o que é verdade e entidade não pode escapar á sciencia divina. --- \emph{A} difficuldade não está no \emph{facto}, está no \emph{modo} de explicar o facto. Com effeito, não podemos dizer que \textsc{Deus} vé esses futuros \emph{condicionados livres}, como vê os futuros \emph{contingentes} ou os \emph{livres absolutos}; porque estes estão \emph{presentes} aos olhos de \textsc{Deus}, e são vistos emquanto \emph{presentes}, porque existem numa differenca de tempo, e a duração simultanea da eternidade divina abrange todas as differencas do tempo, mas os futuros \emph{condicionados livres} não são e não serão nunca \emph{presentes}. --- Como, pois, vé \textsc{Deus} os futuros \emph{condicionados Livres}? --- Ha duas respostas, dadas por duas escholas, egualmente celebres, --- a \emph{thomista} e a \emph{molinista}. Vamos expôr estas respostas com toda a lealdade e com a clareza possivel em materia tão intricada.\par a) Resposta da eschola \emph{thomista}. --- \textsc{DEus} não vé os futuros \emph{condicionados livres} em si mesmos, na propria actualidade, porque não existem e não existirão em nenhuma differenca de tempo; por isso deve vel-os na causa, que os contém e que teria podido dar-lhes a existencia. Esta causa não pode deixar de ser á propria Causa primeira, principio adequado de tudo o que existe ou que teria podido existir. \textsc{DEus} vé tudo isso nos \emph{seus decretos}.\par A essencia divina, --- emquanto é imitavel ao infinito, é, como veremos, o fundamento de todas as possibilidades necessarias e eternas, --- mas, emquanto a realisação de algumas d'estas possibilidades é decretada pela vontade divina, é o fundamento de todas as existencias contingentes, --- presentes, passadas ou futuras. Por quanto, as coisas --- emtanto são possiveis, emquanto \textsc{Deus} é, --- e, emtanto são existentes, emquanto \textsc{Deus} as quer ou permitte. O conhecimento das primeiras é, como dissemos, o objecto da sciencia de \emph{simples intelligencia}; o conhecimento das segundas é o objecto da sciencia de \emph{visão}. --- A sciencia de \emph{simples intelligencia} extende-se á \emph{essencia} e ás \emph{necessarias consequencias} das coisas. A sciencia de \emph{visão} abrange todas as realidades contingentes, queridas ou permittidas nos decretos eternos, não exceptuados os. \emph{futuros livres}, que permanecem \emph{livres}, não obstante a infallivel efficacia da Causa primeira. Ora estes decretos eternos, infinitamente perfeitos, devem extender-se a todos os casos, não só aos existentes mas tambem aos meramente possiveis, e nelles \textsc{Deus} lé não só tudo o que acontece, mas tambem tudo o que aconteceria, se se verificassem certas condições, que nunca se hão de verificar. Tal é a sciencia dos \emph{futuriveis}.\par Por isso, para que \textsc{DEus} saiba o que fará Pedro, se se achar nestas ou naquellas cirecumstancias, deve apenas consultar o decreto geral da sua Providencia\pgn{484}, que regula o nosso livre arbitrio, e ver as soluções, que nelle estão virtualmente contidas para cada um dos casos.\par O \emph{Cardeal Zigliara} escreve as seguintes palavras: « \textsc{Deus} nihil potest cognoscere per aliquid sibi extrinsecum et a se independens, a quo cognitionem accipiat et in cognoscendo dependeat, sed omnia cognoscit in sua essentia ut in medio \emph{in quo} cognito. Sed essentia divina non est causa futuri, sive absoluti sive conditionati, praecise ut essentia est; secus enim, cum ipsa si\& eodem modo exemplar seu idea omnium possibilium, aut haec omnia dicenda essent futura, aut non esset ratio cur aliqua prae aliis possibilibus ad existentiam pervenirent. Ergo essentia divina est medium in quo Deus videt quaecumque futura ut.est determinata per liberum decretum divinae voluntatis, determinantis alia possibilia aliquando futura, alia vero non; scilicet \textsc{Drus} videt in sua essentia omnia futura, quia libere voluit esse futura: quod est videre futura in suo libero decreto » (\emph{Sum}. \emph{Phil}. TTheol., 1. III, c. 2, a. 9).\par Pp) Resposta da eschola \emph{molinista}. --- A sciencia de \emph{visão} nos decretos divinos e a sciencia de \emph{simples intelligencia} não explicam a presciencia dos futuros \emph{condicionados livres}. E' necessaria admittir em \textsc{Deus} uma terceira sciencia --- a \emph{sciencia media}. --- E' claro que estas divisões da sciencia divina, unica e simplicissima em si mesma, se referem á variedade dos objectos, e significam a nossa maneira de conceber as coisas divinas.\par O que é a \emph{sciencia media}? Ella não é só a sciencia dos futuros condicionados livres. Se assim fosse, não haveria divergencia de opiniões; porque todos admittem que \textsc{DEus} tem o conhecimento d'esses objectos. \emph{A} sciencia \emph{media} é uma explicação systematica do \emph{meio}, em que Deus vê os futuros condicionados livres, --- é a sciencia dos \emph{futwriveis}, independente dos decretos divinos e anterior a elles. Diz o theologo \emph{Tournely}: « scientia media est scientia conditionatorum, independens ab omni decreto absoluto et efficaci, eoque anterior » (\emph{De scientia Dei}, a. 5). Portanto, os effeitos do nosso livre arbitrio são independentes da vontade divina; e, se \textsc{DeUs} os conhece, deve conhecelos anteriormente a todo e qualquer decreto.\par Mas se \textsc{DEus} não conhece os \emph{futuriveis} nos decretos da sua vontade, em que \emph{meio} os conhece? --- \emph{Molina} responde que Deus os conhece nas \emph{causas livres}, que são capazes de produzir esses futuros em certas circumstancias e das quaes Elle tem uma comprehensao perfeita, uma \emph{supercomprehensao}. Diz assim; « \textsc{Deus} tem uma profunda e insondavel comprehensao de cada livre arbitrio; Elle vé claramente o que cada uma das causas livres faria, de modo espontaneo, nestas e naquellas circumstancias, e mesmo na infinidade das circumstancias. Esta vista de Deus é o que nós chamamos \emph{sciencia media} » (\emph{Concord}., q. 14, a. 13, disp. 52).\par Todavia este \emph{meio} é rejeitado hoje por todos os philosophos, incluidos os neo-molinistas. Diz o \emph{Card}. \emph{Maszella}; « Deus nao pode conhecer os \emph{fufuriveis} nas suas causas proximas. Porquanto os actos \emph{livres} não têm um nexo \emph{necessario} com as proprias causas, embora dispostas completa e proximamente\pgn{485} para o acto; pois, se a causa livre não fosse, mesmo nesse momento, indifferente para operar e deixar de operar, já não seria livre. E este argumento é invineivel; porque aqui não se trata de um conhecimento \emph{conjectural}, mas de um conhecimento \emph{infallivel} dos futuriveis, com o qual é metaphnysicamente incompativel o erro, qual é o conhecimento de \textsc{DEus} » (\emph{De gratia}, disp. 3, a. 7).\par Nem vale recorrer á infinita energia da vista divina e á sua \emph{supercomprehensdo}. Quanto mais profundo e penetrante é o olhar de quem contempla a natureza de uma causa livre, tanto mais clara e distinctamente conhece que os effeitos, que ella contém, são ainda indeterminados e meramente possiveis. Logo \textsc{Deus} não pode ver, nas causas livres, as determinações, que ainda não existem; a não ser que se diga que estas determinações existem : o que seria a negação da liberdade. Por isso, esta explicação leva logicamente ou a negar a liberdade humana, ou a reduzir a sciencia divina \& condição de sciencia \emph{conjectural}.\par O \emph{meio}, proposto por \emph{Molina} e rejeitado pelos proprios sequazes, ainda não foi substituido. Geralmente, os \emph{neo-molinistas} afirmam que \textsc{Deus} conhece os futuriveis exclusivamente na contemplação das possibilidades eternas ou das essencias das coisas, independentemente de todo e qualquer decreto. --- Esta explicação não serve para o escopo. Comprehende-se facilmante que a Essencia de Deus possa representar e represente á sua intelligencia todas as possibilidades eternas e necessarias, isto é, todas as essencias das coisas contingentes e todos os effeitos \emph{necessarios}, absolutos ou condicionados, que estão nellas contidos. Mas, se os futuros \emph{livres} estivessem tambem contidos nas essencias das coisas, a futurição d'elles, longe de ser livre, seria absolutamente necessaria, e tudo seria regulado pela fatalidade. Logo \textsc{DEus} não pode ver na essencia das coisas os futuros \emph{livres}, que ahi se não encontram.\par Outros, então, appellam para o \emph{meio das contradictorias}. Todos os futuriveis, dizem elles, podem ser representados sob a forma de proposições \emph{contradictorias}; por exemplo: se Pedro é tentado nestas circumstancias, peccará ou deixará de peccar. Não ha meio termo. E, como estas proposições não podem ser ambas verdadeiras ao mesmo tempo, mas é necessario que uma seja verdadeira, e outra falsa, pela essencia mesma das coisas e anteriormente a todo o decreto divino, é claro que a intelligencia divina, a qual conhece a verdade toda, deve conhecer com certeza absoluta qual d'estas proposições é a verdadeira, e qual é a falsa. --- Mas, se o principio de contradicção é sufficiente para mostrar que as duas contradictorias não se podem realisar simultaneamente, não o é para mostrar qual das duas se realisará em certas e determinadas cireumstancias; pois o livre arbitrio torna-as egualmente possiveis. Se uma das duas fosse determinada para a existencia pelo principio de contradicção, ella seria necessaria e não livre. E, se ambas são \emph{possiveis}, qual é a razão sufficiente, pela qual \textsc{Deus} vé que esta será realisada, e aquella não? Para responder a esta pergunta, não ha outro remedio senão recorrer a \emph{supercomprehensüo} de \emph{Molina}.\par Em vista d'isso, não é para admirar que o \emph{P}. \emph{de Regnon}, excellente philosopho e theologo, conclua a exposição da controversia com estas palavras\pgn{486}: « E' um mysterio, um mysterio insondavel (o da sciencia divina dos \emph{futuriveis}) nenhuma das varias explicações apresentadas satisfaz completamente. Devemos renunciar á explicação do \emph{modo} por que Deus vé os futuriveis. Explicar esta sciencia é obra de \emph{dilettantismo} philosophico! ». (\emph{Bannes} et \emph{Molina} p. 119-115).\par Expuzemos, em resumo, a controversia ácerca da sciencia dos \emph{futwriveis} com a palavras de um philosopho moderno (\emph{Farges}, l'Idée de Dieu, pag. 910-8971) Não temos competencia para emittir um juizo definitivo numa questão, que tem cancado as mais profundas intelligencias. Só dizemos, com o citado escriptor, que, devendo escolher entre o mysterio da \emph{sciencia media} e o mysterio dos \emph{decretos divinos}, preferimos este áquelle. A efficacia da Causa primeira, muito diversa da efficacia das causas segundas, pode mover a nossa natureza, sem fazer-lhe violencia, e produzir tudo em nós, mesmo a nossa liberdade, que não pode; sem sacrilegio, subtrahir-se á universal causalidade de \textsc{DEus}, segundo a expressão de S. Agostinho. Este padre da Egreja, criticando \emph{Cicero} e outros pagãos, que não admitiiam em \textsc{Deus} o conhecimento dos \emph{futuriveis}, diz que a theoria d'elles » \emph{dum vult homines facere liberos}, \emph{facit sacrilegos} » (\emph{De civ}. \emph{Dei}, l. 5, c. 9).\par e) \textsc{Deus} \emph{conhece todos os possiveis}. --- 'l'ratamos dos \emph{meramente possiveis}, isto é, dos entes, que, embora não lhes repugne a existencia, comtudo não a têm e não a terão jamais. --- Estes \emph{possiveis}, como dissemos na \emph{Ontologíia}, têm o seu fundamento na essencia divina. --- Que \textsc{Deus} os conheça é fóra de duvida. A intelligencia divina, sendo infinita, conhece tudo o que é cognoscivel. Ora todos os \emph{possiveis} são cognosciveis nos seus elementos communs e individuaes. Logo \textsc{DEus} conhece-os. --- Além d' isso, \textsc{DEus} conhece e comprehende a sua essencia e potencia infinita. Oomprehendendo a sua essencia infinita, \textsc{DEus} deve conhecer todos os modos por que ella pode ser imitada e representada, --- e, comprehendendo a sua potencia infinita, deve conhecer todos os effeitos, que ella pode produzir. --- Finalmente, \textsc{DEus} pode produzir tudo o que é possivel, e produz tudo pela sua intelligencia e vontade, isto é, emquanto conhece e quer as coisas. Ora quem por tal forma produz ou pode produzir as coisas, não pode deixar de as conhecer. Logo Dzvus conhece todos os possiveis.\par E, como a multidão dos possiveis é infinita, podemos dizer que \textsc{DEus} conhece objectos \emph{infinitos}. Diz S. Thomaz: « \textsc{DEus} conhece perfeitamente o seu poder. Ora o poder não pode ser conhecido perfeitamente, se não forem conhecidas todas as coisas, ás quaes pode extender-se; porque é nellas que deve fundar-se o criterio e a quantidade da sua efficacia. Ora o poder de \textsc{Deus}, sendo infinito, extende-se a objectos infinitos. Logo \textsc{Deus} conhece objectos infinitos. (C. \emph{Gentes}, l. I, c. 69). --- Todavia, como adverte o S. Doutor (l. c.), \textsc{Deus} não conhece objectos infinitos pela sciencia \emph{de visão}, porque objectos \emph{én/initos} não existem \emph{actualmente}, nem existiram, nem existirão; mas conhece objectos infinitos pela sciencia de \emph{simples intelligencia}, porque \textsc{Deus} conhece infinitas coisas, que não são, nem serão, nem foram, e que tadavia existem na potencia da creatura, como tambem conhece infinitas coisas, que existem na sua potencia (que Elle pode produzir), mas que não são, nem serão, nem foram. --- Este conhecimento não é confuso, nem\pgn{487}indeterminado; mas é infinitamente distincto e proprio, porque a essencia divina, que faz o officio da especie intelligivel, é uma \emph{semelhança sufficiente} de todas as coisas, que existem ou podem existir, não sô emquanto aos elementos communs, mas tambem emquanto aos elementos proprios de cada ente. (\emph{Sum}. \emph{Th}., p. 1, q. 14. a. 18).\par Objectam: Se \textsc{Deus} conhecesse infinitas coisas, conheceria um numero infinito. Ora um numero infinito não pode ser conhecido, porque é uma contradicção, um nada. Logo Deus não conhece infinitas coisas. --- Respondemos, \emph{negando} á proposição \emph{maior}. \emph{A multidão}, absolutamente considerada, não pode nem deve confundirse com o \emph{numero}. O \emph{numero} é, sem duvida, uma multidão, mas é uma multidão, medida pela \emph{unidade}, (porque é formado pela addig\&o de unidades), e por isso não pode ser infinito. Mas \textsc{Deus}, conhecendo objectos infinitos, --- não os conhece como se conhece um numero, addicionando partes a partes e unidades a unidades, --- mas conhece-os todos e ao mesmo tempo, e não os conhece em si mesmos mas sua essencia divina, emquanto vó que esta pode ser imitada e participada por uma illimitada multidão de creaturas. Diz S. Thomaz: « \textsc{DEUS} non sic cognoscit infinitum vel infinita, quasi enumerando partem, cum cognoscat omnia simul, non successive » (\emph{Sum}. \emph{Th}. p. I, q. 14, a. 12 ad 1; Cf. c. \emph{Gentes}, 1. I. c. 69). E' falso, pois, que \textsc{Deus}, se conhecesse infinitas coisas, conheceria um numero infinito.}.

\pg{478} 65. \textbf{Deus} tem um conhecimento proprio de todas as coisas.

a) Um ser intelligente tem um conhecimento \emph{proprio} de uma coisa, quando a conhece --- não só nos seus elementos \pg{479}  \emph{communs} e \emph{universaes}, pelos quaes convem com as outras coisas do mesmo genero ou da mesma especie, --- mas tambem nos \pg{480}  seus elementos \emph{proprios} e \emph{individuaes}, pelos quaes se distingue de todas as outras coisas do mesmo genero e da mesma \pg{481}  especie. Ora \textsc{Deus} não pode deixar de conhecer todas as coisas tambem nos seus elementos \emph{proprios} e \emph{individuaes}; porque, se as conhecesse somente nos elementos \emph{communs} e \emph{universaes},

\pg{482} "

482 METAPHYSICA --- PARTE QUARTA --- THEODICEA o seu conhecimento seria imperfeito: o que se não pode admittir em \textsc{Deus}, eujo entender é uma e a mesma coisa com a sua perfeitissima essencia. Logo \textsc{Deus} tem um conhecimento \emph{proprio} de todas as coisas.

b) Dzvus é causa de todas as coisas pela sua intelligencia, e per isso deve conhecel-as como são em si mesmas. Ora todas as coisas são em si mesmas o que são --- não só pela natureza, que lhes é commum com outras da mesma especie, --- \pg{483} \pg{484}  --- --- --- m ------ SENA = v ELT

484 METAPHYSICA --- PARTE QUARTA --- THEODICEA mas tambem pelas ultimas differenças e propriedades individuaes, que as distinguem de todas as outras, embora pertencentes á mesma especie. Logo \textsc{Deus} tem um conhecimento \emph{proprio} de todas as coisas.

c) \textsc{Deus} comprehende perfeitamente o proprio ser, e por isso todas as coisas, que estão contidas na sua potencia \emph{activa}, \pg{485} \pg{486}  isto é, a que a sua potencia \emph{activa} pode extender-se. Ora todas as coisas estão contidas na potencia \emph{activa} de \textsc{Deus} emquanto \pg{487}  ás formas \emph{proprias} e ás differenças \emph{ultimas}, porque \textsc{Deus} é o principio de todo o ente. Logo \textsc{Deus} tem um conhecimento \emph{proprio} de todas as coisas \footnote{Para evitar equivocos, advertimos que o conhecimento se pode chamar \emph{proprio} em dois sentidos, muito diversos, --- ou emquanto a especie, que informa a intelligencia, representa uma coisa com aquellas determinadas perfeições, e não outra, --- ou emquanto a especie é tão perfeita, que representa todas e cada uma das coisas com as perfeições, que lhe são proprias. No primeiro caso, a especie é imperfeitissima; no segundo, é infinitamente perfeita. Neste segundo sentido dizemos que \textsc{Deus} conhece todas as coisas com um conhecimento \emph{proprio}, emquanto as conhece pela sua Essencia infinita, representativa de todas as coisas com as propriedades e perfeições, proprias de cada uma. --- Do mesmo modo, a especie \emph{propria} --- pode representar o \emph{universal}, como são todas as especies abstractas das coisas materiaes, que informam a nossa intelligencia no estado de união e que lhe apresentam o objecto \emph{proprio}, \emph{e}, então, o \emph{proprio} não se oppõe ao \emph{comqum}, mas ao \emph{improprio}, ao \emph{indirecto}, --- ou pode representar o \emph{singular}, o \emph{individuo}, e, então, o \emph{proprio} oppõe-se ao \emph{commum}. Ora a essencia divina não se diz especie \emph{propria} no \emph{primeiro} sentido, mas no \emph{segundo}, emquanto entendemos dizer que a essencia divina, pela sua infinita perfeição, representa as coisas como são em si mesmas, isto é, na sua singularidade.\par O t.º argumento, adduzido no texto, é convincente. Conhecer uma coisa nos seus elementos \emph{communs}, p. e, emquanto é \emph{ente}, é conhecel-a \emph{imperfeitamente}, porque não se conhecem os elementos \emph{proprios}, isto é, as ultimas perfeições, que completam uma coisa e a distinguem das outras. Diz S. Thomaz: « Intelligere aliquid in \emph{comniuni} et non \emph{in speciali}, est \emph{imperfecte} aliquid cognoscere. Unde intellectus noster, dum de potentia in actum reducitur, pertingit prius ad cognitionem \emph{wniversalem} et \emph{confusam} de rebus,\pgn{488}quam ad \emph{propriam} rerum cognitionem, sicut de imperfecto ad perfectum procedens. Si igitur cognitio Dri de rebus aliis a se esset \emph{in universali} tantum, et non \emph{in speciali}, sequeretur quod \emph{ejus intelligere non esset omnibus modis perfectum}, et per consequens nec ejus \emph{esse} » (\emph{Sum}. \emph{Thy} p. I, q. 14, 98,76; of. c. \emph{Gentes}, l. T,-c. 5)\par A mesma doutrina é exposta por S. Thomaz no magnifico trecho seguinte. Diz assim: « Toda a perfeição, que se encontra na creatura, deve preexistir, num grau infinitamente mais elevado, no proprio Creador. Ora perfeição das creaturas não é apenas o que lhes é \emph{commum}, isto é, o ser, mas é tambem o que lhes é \emph{proprio}. E por isso todas as coisas creadas preexistem em \textsc{Deus} não só emquanto aos elementos, pelos quaes convérh com outras, mas tambem emquanto aos elementos, pelos quaes se distinguem entre si Contendo todas as perfeições, a essencia divina refere-se a todas as essencias das coisas, \emph{não como} o \emph{commum ao proprio}, \emph{mas como} um \emph{acto perfeito aos imperfeitos}. Ora é evidente que, por meio de um acto perfeito, podemos conhecer os actos imperfeitos não só com um conhecimento commum, mas tambem com um conhecimento \emph{proprio}; assim quem conhece o \emph{homem} tem um conhecimento \emph{proprio} do \emph{animal}. Portanto, contendo eminentemente a Essencia divina todas as perfeições, de que são dotadas as creaturas, \textsc{DEUS} pode conhecer em si \emph{mesmo} todas as coisas com um conhecimento \emph{proprio}» (\emph{Sum}. \emph{Th}., l. c.).\par Para a solução das varias difficuldades contra a verdade exposta e defendida no texto, fazemos as observações seguintes: 1. O conhecimento exige a união do objecto conhecido com o cognoscente por meio de uma especie ou forma intencional, que actue como principio formal e represente o proprio objecto. E, para que o conhecimento seja \emph{proprio}, basta que a especie represente o objecto tambem emquanto aos elementos \emph{proprios} e \emph{individuaes}, e não é necessario que ella seja por tal modo proporcionada ao objecto, que represente exclusivamente este, e não um outro. Ora a Essencia divina, representando todas as coisas emquanto aos elementos communs e \emph{proprios}, é principio de um conhecimento \emph{proprio}, embora não seja limitada aos objectos, que representa. --- 2.º) Pelo facto de a Essencia divina ser um \emph{unico} meio de conhecimento, não se segue que \textsc{Deus} não possa ter um conhecimento \emph{proprio} e \emph{distincto} de \emph{muitos} objectos; porque nenhum d'estes objectos imita \emph{completamente} e reproduz \emph{todas} as perfeições, que ella encerra, mas cada um reproduz, e de um modo deficientissimo, uma ou outra perfeição. Por isso a unica Essencia divina, excedendo infinitamente todos os objectos representados, pode representar objectos infinitos, sem nunca se exgottar. Se houvesse uma imagem de \textsc{Deus}, que imitasse completamente e exprimisse tudo o que está em \textsc{Deus}, essa imagem não poderia deixar de ser \emph{unica}. (Cf. \emph{De} Ver., q. 2, a. 4 ad 2). --- 8.*) \textsc{Deus} conhece todas as coisas pela sua Essencia, que é causa \emph{efficiente} e \emph{ewtrinscca} de todas as coisas, e por isso mesmo o conhecimento, que \textsc{Deus} tem de todas as coisas, é eminentemente \emph{proprio} e \emph{distincto}, visto ser a Essencia divina o principio de todo o \emph{ser}, de que as coisas são constituidas.}.

\pg{488} 66. A \textbf{sciencia divina} é \textbf{causa das coisas creadas}. --- A sciencia de \textsc{Deus} está para as coisas creadas, como a sciencia do artifice está para as coisas artificiaes. Ora a sciencia do artifice é causa das coisas artificiaes; porque o artifice opera pela sua

\pg{489} \textbf{FEE}--------- .?.".\textbf{so---mnm}--------- = == E

CAPITULO II --- FACULDADES E OPERACÓES IMMANENTES DE DEUS 489 intelligencia. Logo tambem a. sciencia de \textsc{Deus} é causa das coisas creadas. --- A sciencia divina é causa das coisas creadas, emquanto está unida, para o effeito, com o acto da vontade \footnote{:A conclusão é certissima. A verdade e a certeza da sciencia das coisas, que a mesma sciencia não produz, mas suppõe já produzidas, depende das proprias coisas e é por ellas medida. Por isso, se a sciencia de Drzus não fosse causa productiva ou efficiente das coisas, a sua verdade e certeza dependeria das coisas creadas e seria medida por ellas. Ora isto é falso. Deus não conhece as coisas, porque são: mas as coisas são, porque \textsc{DEUS} as conhece. (\emph{Sum}. \emph{Th}., p. I, q. 14, a. 8. --- D'ahi a differenca entre a nossa sciencia e a de \textsc{DEus}. Às coisas naturaes são anteriores á nossa, sciencia, e são a medida d'esía; porque a nossa sciencia é verdadeira, emquanto se adapta e corresponde aos objectos percebidos; --- ao passo que a sciencia divina é anterior ás coisas naturaes,e é a medida d'estas; porque as coisas são verdadeiras, emquanto se adaptam e correspondem aos archetypos divinos. Por isso as coisas creadas são intermedias entre a sciencia divina e a humana (\emph{Sum}. \emph{Th}., l. c. ad 8) --- Mas a sciencia divina é causa das coisas, não emquanto é um simples conhecimento das coisas, mas emquanto está unida com o acto da vontade, isto é, \textsc{DEus} produz as coisas emquanto as conhece e as quer. Com effeito, se a sciencia, só por si, fosse a causa efficiente e sufficiente das coisas, tudo o que \textsc{Deus} conhece existiria: o que é falso. Por isso a sciencia divina, que é causa das coisas e que está unida com a vontade, chama-se propriamente \emph{sciencia de approvacdo}. D'ahi se segue: 1.9 que \textsc{DEus} não é causa do \emph{peccado}; porque \textsc{DEus} não \emph{quer} e não \emph{approva} a deformidade do peccado, mas só \emph{permitte}; --- 2.9) que, embora a sciencia divina seja eterna, comtudo as creaturas não são eternas; porque, repetimos, as coisas derivam da sciencia, por meio da vontade, e pór isso não é necessario que ellas existam desde que são conhecidas, mas devem existir quando o tiver determinado a vontade. (Cf. \emph{De} Ver., q. 2, a. 14 ad 2).\par Extendendo-se a sciencia divina a todas as coisas, existentes e possiveis, e sendo ella productiva dos seres creados, é evidente que na intelligencia divina ha muitas \emph{idéas}, quer a \emph{idéa} se tome em ordem ao \emph{conhecimento} e signifique a semelhança intencional dos objectos percebidos, quer se tome em ordem á \emph{operação} e denote a forma exemplar ou o modelo da coisa, que se deve fazer. Com effeito, \textsc{Deus} é dotado de sciencia perfeitissima, identificada com a sua essencia, e é causa de todas as coisas pela sua intelligencia. Ora seria impossivel tanto a sciencia de coisas diversas, como \emph{a} producção das mesmas pela inteligencia, se em \textsc{Deus} não existisse uma pluralidade de \emph{ideas}, quer a \emph{idéa} signifique a semelhança intencional do objecto, que foi percebido, quer signifique o modelo, segundo o qual se ha de fazer uma coisa. --- Que em \textsc{Deus} existam muitas \emph{idéas}, emquanto estas significam as semelhanças intencionaes dos objectos percebidos, entende-se facilmente. Que em \textsc{Deus} existam muitas \emph{idéas}, emquanto estas significam as formas exemplares, os modelos das coisas creadas, deduz-se da soberana independencia de \textsc{Deus} na operação, como no ser. \textsc{Deus} não é uma causa \emph{necessaria}, que opéra em virtude da sua forma natural, e que por isso sô pode produzir um effeito determinado; mas é a causa infinitamente \emph{livre}, que\pgn{490}opera pela sua intelligencia, e que por isso deve ter e tem as formas exemplares á semelhança das quaes produz todas as coisas. --- Logo todas as coisas, feitas por \textsc{Deus}, tem um duplice ser: um em si mesmas, emquanto são realmente distinctas de \textsc{DEus}, mutaveis, compostas, finitas; outro na arte divina, isto é, nos divinos archetypos ou exemplares, e então não se distinguem realmente de \textsc{Deus}, e constituem a propria vida de \textsc{DEus}. , . Esta \emph{multiplicidade} de idéas não importa \emph{composição}, e por isso não é contraria á \emph{simplicidade} de \textsc{Deus}. Dá-se a \emph{composição}, opposta á \emph{simplicidade}, numa intelligencia, que percebe muitas e diversas coisas, quando o principio determinativo da intelligencia, que é a especie intelligivel, ou é \emph{realmente distincto} da intelligencia, ou é \emph{multiplice}; porque, em ambos os casos, ha composição de \emph{potencia} e \emph{acto}. Ora nada d'isso se verifica com relação á intelligencia divina. Porquanto, a essencia divina, que faz o officio de especie intelligivel, identifica-se com a intelligencia divina e é \emph{simplicissima}, embora possa ser imitada pelas creaturas por modos infinitos. Por isso, entendendo a intelligenciw divina todas as coisas por uma unica especie, --- pela essencia divina, que é semelhança de si mesma e de todas as coisas, a multidão das coisas entendidas não importa composição na propria intelligencia. --- A \emph{variedade} e \emph{multiplicidade} de ideas não diz respeito á divina essencia, mas aos termos ou ás coisas, que a podem imitar. (Cf. \emph{Sum}. \emph{Th}., p. 1, q. 15; e. \emph{Gentes}, l. I, e. 51 e seg.; \emph{De} Ver., q. 8, a. 2).\par Para a solução das difficuldades, tenhamos presentes os seguintes principios: 1.9) Ha uma grande differença entre a semelhança, que as creaturas têm com a natureza divina, e a semelhança, que as mesmas creaturas têm com os divinos archetypos. A primeira semelhança é imperfeitissima, porque as creaturas imitam a natureza divina por um modo infinitamente deficiente e remoto. A segunda é perfeitissima, porque as creaturas representam adequadamente, na sua entidade, os divinos archetypos, em conformidade com os quaes foram produzidas. (Cf. \emph{Sum}. \emph{Th}., p. I, q. 44, a. 8 ad 1). --- 2.9) Se o conhecimento divino se refere \emph{ao cognoscente}, \textsc{DEUS} conhece o proprio ser e as outras coisa pelo \emph{mesmo modo}, isto é, pela sua \emph{essencia}. Se o conhecimento se refere \emph{ds coisas conhecidas}, \textsc{Deus} conhece o proprio ser e as outras coisas por um modo \emph{differente}; porque conhece o proprio ser pela sua \emph{essencia}, emquanto esta o representa por \emph{identidade}, e conhece as outras coisas pela sua \emph{essencia}, emquanto esta as representa por \emph{semelhança}. Por isso é necessario admittir em \textsc{DEus} muitas e diversas idéas, visto ser muitos e diversos os modos, por que pode ser imitada pelas creaturas a unica e simplicissima essencia divina, --- 8.9) Às coisas imitam a essencia divina pelo facto de \textsc{Deus} a conhecer como imitavel; e não vice-versa, isto é, não é pelo facto de as coisas imitarem a essencia divina que \textsc{Deus} a conhece como imitavel. Por isso, na divina intelligencia, que conhece a sua essencia emquanto imitavel, subsistem todos os archetypos, á semelhança dos quaes são feitas todas as coisas. --- 4.º) As idéas, ou os archetypos, emquanto subsistem na divina intelligencia, são uma e a mesma coisa com ella; mas, emquanto se consideram fóra de \textsc{Deus}, são nada; porque, fóra de \textsc{Deus}, só existem entes \emph{redes}.\par Do que deixamos exposto vê-se quão absurdas são as theorias dos \emph{pantheistas} e de alguns \emph{theistas} Acerca da \emph{sciencia divina}. --- \emph{Spinoza}, ensinando\pgn{491}que \textsc{Deus} é \emph{impessoal}, conclue que Elle não é intelligente senão no homem.\par x --- \emph{Fichte}, \emph{Schelling}, \emph{Hegel}, e os sequazes do \emph{progresso indefinido}, afirmam que \textsc{Deus} conhece todas as coisas não \emph{simultanea} mas \emph{successivamente}, em virtude das suas evoluções. \emph{Renan} diz que \textsc{DEus} não chegou ainda a conhecer, de um modo perfeito, o mundo, mas que, estando Elle \emph{in fieri}, talvez chegue a possuir uma sciencia infinita, --- \emph{Cicero}, os \emph{Sociniamos}, e, nos tempos modernos, \emph{Ramiron}, disseram que Deus! não tem conhecimento, pelo menos \emph{certo}, dos \emph{futuros livres}; pois, na sua opinião, a prosienois divina não pode conciliar-se com a liberdade humana.}.

\pg{490} 67. \textbf{Propriedades} da \textbf{sciencia divina}. --- As principaes são \emph{quatro}: \emph{é substancial}, \emph{infinita}, \emph{intuitiva}, \emph{immutavel}.

\pg{491} a) E' \emph{substancial} emquanto ao \emph{sujeito}. A sciencia, em \textsc{Deus}, não é uma coisa accidental, como é em nós, mas é uma e a mesma coisa com a sua substancia. --- Identificando-se com a substancia divina, a' sciencia divina é \emph{eterna} \footnote{Se a sciencia de \textsc{Deus} não se identificasse com a sua substancia, a Essencia de \textsc{DEus} referir-se-hia ao conhecimento como a potencia ao acto, e \textsc{Deus} receberia de uma coisa distincta d'Elle a perfeição. Ora isto é absurdo, porque \textsc{Deus} é acto purissimo e é perfeição de si mesmo. Por isso diz S. Thomaz que «in Deo intellectus, intelligens, et id quod intelligitur, et species intelligibilis, et ipsum \emph{ntelligere}, sunt omnino unum et idem » (\emph{Sum}. \emph{Th}., p. L, dq 14, a. 4, 15).}.

b) E' \emph{infinita} emquanto ao \emph{objecto}. Com effeito, \textsc{Deus} conhece tudo o que é, de algum modo, cognoscivel; e por isso a sua sciencia extende-se a todas as coisas, --- existentes e possiveis, --- passadas, presentes e futuras, --- necessarias e livres, --- absolutas e condicionadas \footnote{« A extensão da sciencia ou do conhecimento de um cognoscente é proporcionado á forma, que é o principio do proprio conhecimento. Assim, a especie \emph{sensivel}, que está nos sentidos, é a semelhança de um só individuo, e por isso não faz conhecer senão aquelle individuo. A especie \emph{intelligivel} da nossa intelligencia é a semelhança do objecto emquanto á \emph{natureza} da \emph{especie}, que pode ser participada por infinitos individuos; e por isso a nossa intelligencia, por meio da especie intelligivel de homem, conhece, de algum modo, uma multidão infinita de homens, --- não emquanto aos elementos \emph{individuaes}, mas emquanto aos elementos \emph{especificos}. A essencia divina, pela qual a intelligencia divina conhece, representa uma multidão infinita de coisas, existentes e possiveis; e por isso a sciencia de \textsc{Deus} é infinita ». (\emph{Sum}. \emph{Th}., p. I, q. 14, a. 12).}.

c) E \emph{intuitiva} emquanto ao \emph{meto}. \textsc{Deus}, conhecendo a sua essencia, conhece nella todas as coisas. Ora \textsc{Deus} conhece a sua essencia com um unico acto \emph{intuitivo}, de infinita energia e extensão. Logo a sciencia divina é \emph{intuitiva} \footnote{No conhecimento \emph{discursivo}, que é opposto ao ?\emph{ntuitivo} e que é sempre acompanhado por um \emph{movimento}, \& inteligencia passa da \emph{potencial} para a \emph{actual} percepção da conclusão, e esta não a conhece pelo mesmo acto, por que conhece as premissas, mas por um novo acto, que succede ao primeiro. Por isso o conhecimento \emph{discursivo} suppõe sempre uma composição e\pgn{492}distincc\&o de potencia e de acto, e um conjuncto de sciencia e de ignorancia. Ora, como nada d'isso pode admittir-se em \textsc{Deus}, concluimos que a sciencia de \textsc{Deus} não é \emph{discursiva}, mas \emph{intuitiva}. (\emph{Sum}. \emph{Th}., p. I, q. 14, a. 70).}.

\pg{492} d) E' \emph{immutavel} emquanto á \emph{natureza}. --- A sciencia divina identifica-se, como dissemos, com a substancia divina, e por isso é dotata dos mesmos caracteres. Ora a substancia divina é immutavel. Logo tambem a sciencia divina é \emph{immutavel} \footnote{A sciencia divina refere-se ás creaturas, não emquanto estas existem em si mesmas, mas emquanto existem em \textsc{DEus}. E, como em \textsc{Deus} existem por um modo \emph{invariavel}, assim a sciencia, que \textsc{DEus} tem das creaturas, é invariavel. --- \textsc{Deus} conhece a \emph{variabilidade} ou \emph{mutabilidade} das coisas; mas a sua sciencia é inyariavel ou immutavel. --- A sciencia de \textsc{Deus} seria variavel só no caso de Elle chegar a conhecer uma coisa, que antes ignorava. Ora isto é impossivel. \textsc{Deus}, na sua eternidade, vé tudo o que existe no tempo e conhece tudo o que é cognoscivel. (\emph{Sum}. \emph{Th}, p. I, q. 14, a. 15).}.

\emph{A}

68. \textbf{Termo} immanente da sciencia diviva. --- O acto da intelligencia divina, pelo qual \textsc{Deus} conhece o proprio ser e todas as coisas, existentes e possiveis, deve ter o seu termo. E, como o acto de entender é formalmente \emph{immanente}, o termo tambem deve ser \emph{immanente}, isto é, deve permanecer no principio, de que deriva. Este termo \emph{immanente}, que é perfeitissima expressão de tudo o que \textsc{Deus} conhece e por isso é imagem semelhantissima do proprio ser divino e de todas as coisas, é e diz-se Verbo. --- O Verbo, em \textsc{Deus}, é unico, infinito, substancial; como unico, infinito, substancial é o acto da intelligencia divina, de que procede \footnote{O \emph{verbo mental} encontra-se em toda a natureza intellectual, e por isso, como no homem, tambem em \textsc{Deus}. Já tivemos occasiüo de indicar as \emph{analogias} entre o verbo humano e o verbo divino. Não insistimos no assumpto. A razão humana chega a entender que, emquanto o verbo humano não é a propria substancia da alma, mas é uma coisa accidental, o verbo divino, ao contrario, nio é uma coisa accidental, mas é o proprio Dus. Porquanto, em \textsc{Deus}, a intelligencia, o acto intellectual, o principio de que o verbo deriva, e o proprio verbo não se distinguem da essencia, mas são uma e a mesma coisa com a essencia. Tal é a conclusão da razão. (\emph{Of}. \emph{Sum}. \emph{Th}., p. I, q. 21, 84; c. \emph{Gent}., 1. IV, c. 16; \emph{De Verit}., q. 4; \emph{De dif}. Verb. \emph{div}. \emph{et hum}.).\par Mas a Fé descobre horisontes invisiveis á razão humana, e diz-nos que o Verbo, em \textsc{Deus}, é uma verdudeira \emph{Pessoa}, subsistente na divina natureza e distincta, pela relação da origem, do Principio, de que deriva e que é tambem \emph{Pessoa}, porque tambem Elle subsiste na mesma natureza divina. Na verdade, chama-se \emph{pessoa} quem subsiste numa natureza intellectual e distingue-se de todos os outros, que subsistem na mesma natureza, tanto quando a natureza é uma e a mesma \emph{especificamente} em todos, o que se verifica nas creaturas, como quando a natureza é uma e a mesma \emph{numericamente}, o que se verifica em \textsc{Deus}.\par O Verbo divino procede do seu Principio por verdadeira \emph{geração}, e por\pgn{493}isso é verdadeiro \emph{Filho} de \textsc{Deus}. A \emph{geração} é \& origem de um ser vivo de outro ser vivo, que communica no primeiro a propria natureza em virtude da sua intrinseca actividade e a propria semelhança em consequencia da indole da origem. Ora todos estes elementos encontram-se na processão do Verbo divino. Na verdade, o Verbo --- procede do seu Principio, como \textsc{Deus} vivo de \textsc{Deus} vivo, --- recebe do seu Principio não uma natureza da mesma especie, mas a propria identica natureza, --- procede em virtude da intrinseca actividade e fecundidade do seu Principio, porque a intelligencia divina possue sempre a energia productora do seu Verbo, --- procede como imagem do seu Principio, por isso mesmo que procede como Verbo; porque o verbo é essencialmente imagem do objecto entendido pela intelligencia, e é imagem da propria intelligencia, quando esta se entende a si mesma, e, como a intelligencia divina produz o seu Verbo entendendo-se e comprehendendo-se a si mesma, por isso o Verbo é essencialmente imagem da mesma intelligencia, de que procede. Logo o Verbo procede de \textsc{DEus} por verdadeira geração, e por geração infinitamente perfeita e superior a toda a geração das creaturas; e assim o Verbo é \emph{Filho} de Dxzvs.\par O Verbo 6 \emph{Filho} de \textsc{Deus}, e é \emph{todo} Filho de \textsc{Deus}, porque a \textsc{Deus} se attribue tudo o que concorre para a geração do Verbo. Para, entendermos isto, consideremos o que se verifica na nossa intelligencia. Quando esta pensa num objecto, concebe o seu verbo, que é uma imagem ideal do proprio objecto. Este verbo mental é filho --- tanto da intelligencia, que o produz, --- como do objecto, que, pela sua especie, fecunda a intelligencia e a determina a produzir o verbo, semelhante a si mesmo. Por isso, quando a intelligencia pensa num objecto, distincto de si mesma, o objecto entendido é como o pae do verbo, concebido pela intelligencia, porque. é principio \emph{activo}, emquanto actua na faculdade, e a intelligencia é como a mãe do verbo, porque é principio \emph{passivo}, emquanto concebe e traz no seu seio o verbo. Mas, quando a intelligencia pensa em si mesma, então ella é, ao mesmo tempo, pae e mãe do verbo, e a ella deve attribuir-se tudo o que concorre para a geração do mesmo verbo. E' o que se verifica em \textsc{Deus}. A intelligencia divina produz o seu Verbo pela comprehens\&o de si mesma, e por isso ella é o unico principio gerador do Verbo.\par O Verbo é todo Filho de \textsc{Deus}, e é Filho \emph{wnigenito}. O Filho de \textsc{Deus} não pode deixar de ser \emph{unico}, tanto se se considera o \emph{acto}, que O gera, como se se considera o termo d'esse acto. --- Se se considera o \emph{acto}, este é \emph{unico}, simplicissimo, purissimo, infinito, e infinitamente intelligente, e por isso comprehende a essencia divina e nella tudo o que é cognoscivel, e por isso \textsc{Deus} exprime pelo Verbo tudo o que Elle comprehende. --- Se se considera o \emph{termo} desse acto, que é o \emph{Verbo}, este é infinitamente perfeito, e por isso é uma expressão adequada da: intelligencia divina, de que procede,\pgn{494}Ora se o \emph{acto} do entender divino é \emph{unico}, e se o Verbo divino exprime \emph{tudo} o que \textsc{Deus} entende, o Verbo deve ser \emph{unico}. A repetição ou multiplicação do Verbo, em \textsc{Deus}, significaria impotencia, imperfeição. Poderá haver, na creação, innumeraveis expressões d'esse Verbo, mais ou menos inadequadas; mas estas, embora multiplicadas ao infinito, não exprimirão nunca tudo o que o unico Verbo de \textsc{Deus} exprime. --- Este Verbo, tão glorioso na sua Divindade e depois tão humilhado na sua Humanidade, é o nosso querido \textsc{JESUS}... Outras noções relativas a este assumpto encontram-se nos tratados de Theologia revelada, (Cf. \emph{Sum}. \emph{Th}., p. I, q. 27; p. III, q. 28; \emph{c}. \emph{Gent}., 1. IV, o: 11,13),}.

%%>>> art-42

\section*{Artigo II.\\ Facullade e operação volitiva immanente.}
\addcontentsline{toc}{section}{Artigo II. --- Facullade e operação volitiva immanente.}

\pg{493} 69. \textbf{Vontade de Deus}. --- \emph{Vontade} de \textsc{Deus} é a faculdade, pela qual Elle ama o bem, percebido pela sua intelligencia.

\pg{494} --- Acto da vontade é o \emph{amor}. --- Diremos da \emph{existencia} d'esta faculdade, da sua \emph{indole} e \emph{objecto} \footnote{S. Thomaz trata da \emph{vontade} de \textsc{DEUS} na --- \emph{Sum}. \emph{Th}, p. T, q. 19; --- c. \emph{Gent}., 1. I, c. 62 e seg.; --- \emph{De} Ver, q. 25; --- I \emph{Sent}., Dist. 45.}.

70. \textbf{Deus} é dotado de \textbf{vontade}.

a) Toda a substancia espiritual, que é dotada de \emph{intelligencia}, é tambem dotada de \emph{vontade} ; porque essa substancia, percebendo pela intelligencia o bem conveniente, não pode deixar de ter uma inclinação natural para esse bem, de modo que tenda para elle, quando o não possue, e descance nelle, quando o possue; e tal inclinação é a \emph{vontade}. Ora Dxvus é dotado de \emph{intelligencia}. Logo é tambem dotado de \emph{vontade}.

b) A forma, considerada pela intelligencia, não determina o agente ao acto externo, se não por meio da vontade, eujo objecto é o fim ou o bem; como pode verificar-se no artifice, que concebeu na intelligencia a forma da casa. Ora \textsc{Deus}, artifice sapientissimo, produz as suas operações \emph{transeuntes} na creação conforme a idéa ou forma exemplar, concebida na sua intelligencia. Logo deve admittir-se em \textsc{Deus} a vontade, que o move a operar. --- Por isso em \textsc{Deus}, que todas as creaturas produziu á imitação da sua essencia, a intelligencia \emph{dirige}, á vontade \emph{impera}, o poder \emph{executa}.

c) \textsc{Deus} é infinitamente perfeito. Ora não o seria, se não fosse dotado de \emph{vontade}; não só porque a vontade é uma perfeição \emph{absoluta}, mas tambem porque, se não possuisse a vontade, nem poderia possuir a \emph{santidade}, a \emph{justiça}, a \emph{bondade}, a \emph{felicidade}, e todos os attributos, que residem nessa faculdade. Logo \textsc{Deus} é dotado de \emph{vontade} \footnote{A vontade, em nós, tem um duplice acto --- \emph{appetecer} o bem não possuido, e \emph{amar} o bem possuido e \emph{gosar} d'elle. O primeiro acto indica \emph{imperfeição}, os outros não. Por isso quando o nome de \emph{vontade} se transfere das creaturas a \textsc{DEUS}, não se applica a \textsc{Deus} senão emquanto significa perfeição, isto a, emquanto \textsc{Deus} ama o que possue, porque possue sempre o bem, que é seu objecto. D'onde se segue que a vontade divina impropria e falsamante se chama \emph{appetite}, se esta palavra se toma no sentido rigoroso,\pgn{495}emquanto denota a tendencia para um bem não possuido. --- O objecto da vontade de \textsc{Deus} é a sua propria bondade, que se identifica com a sua essencia. Por isso a vontade de \textsc{Deus} não é movida por um bem extranho, mas por Elle só, e diz-se \emph{movida}, emquanto se chama \emph{movimento} tambem o acto da intelligencia e da vontade.}.

\pg{495} 71. A vontade de \textbf{Deus} identifica-se com a \textbf{sua essencia}, e é \textbf{acto} purissimo, e infinitamente \textbf{perfeita}.

a) Identifica-se com a sua essencia, --- Se a \emph{vontade} de \textsc{Deus} não se identificasse com a sua \emph{essencia}, Elle seria \emph{composto} de \emph{essencia} e de \emph{attributos}. Ora \textsc{Deus} não é composto de \emph{essencia} e de \emph{attributos}. Logo a \emph{vontade} de \textsc{Deus} identifica-se com a sua \emph{essencia}. --- Por isso \textsc{Deus} \emph{quer} pela sua essencia, como \emph{conhece} pela sua essencia.

b) E' acto purissimo. --- A \emph{vontade} de \textsc{Deus} identifica-se com a sua essencia. Ora a essencia de \textsc{Deus} é \emph{acto purissimo}, excluindo toda a sombra de potencia e de imperfeição. Logo a \emph{vontade} de \textsc{Deus} é \emph{acto purissimo}.

o) E' infinitamente perfeita, --- Uma vontade, que ama uma bondade infinita tanto, quanto esta é amavel, é infinitamente perfeita, Ora a vontade divina, por ser adequada á divina bondade, ama esta bondade tanto, quanto é amavel, isto é, com um amor simplesmente infinito. Logo a divina vontade é \emph{infinitamente perfeita} \footnote{Todas as faculdades de um ente, como tambem os seus actos, seguem a natureza do mesmo ente. --- Em \textsc{Deus}, ente \emph{s/mplicissimo}, a vontade, o acto volitivo, a substancia, são uma e a mesma coisa na \emph{realidade}, havendo apenas entre ellas uma distinecüo de razão. No homem, como nas outras creaturas intellectuaes, a vontade, o acto volitivo, e substancia, são coisas \emph{realmente} distinctas. --- Além d'isso, a vontade de \textsc{Deus} está \emph{sempre em acto}, isto é, \emph{quer sempre}, e permanece sempre immovel e invariavel, embora queira coisas diversas, e livremente se refira aos effeitos, que imitam a sua bondade. A creatura ora quer, ora não quer, passa da potencia para o acto, e está sujeita a mudanças. --- O acto da vontade divina produz a bondade e a communica ás coisas, que ella quer. O acto da vontade creada é produzido e medido pelo bem, que se encontra nas coisas.\par *\par A vontade divina, embora una e simplicissima, divide-se --- emquanto ao \emph{modo}, por que se manifesta, em vontade de \emph{beneplacito} e de \emph{signal}, --- emquanto ao \emph{e/feito}, que produz, em \emph{antecedente} e \emph{consequente}, --- emquanto ao \emph{objecto}, que attinge, em \emph{necessaria} e \emph{livre}.\par \emph{a}) Vontade de \emph{beneplacito} e de \emph{signal}. A primeira é a vontade \emph{propriamente} dita; a segunda é a vontade tomada em sentido \emph{metaphorico}, e consiste num \emph{signal} externo, com que \textsc{Deus} manifesta a sua vontade, assim o \emph{castigo}, que \textsc{Deus} inflige ao peccador, é um \emph{signal} de \emph{ira} na sua vontade.\par \emph{b}) Vontade \emph{antecedente} e \emph{consequente}. Pela primeira, \textsc{Deus}, movido pela\pgn{496}sua benevolencia, quer um effeito, independentemente de toda e qualquer condição; assim \textsc{DEus} quer a salvação da todos os homens. Pela segunda, \textsc{Deus}, suppostas certas cireumstancias, quer um effeito ainda que extranho á sua bondade; assim \textsc{Deus} pune o homem, se elle morrer em peceado.\par \emph{c}) Vontade \emph{necessaria} e \emph{livre}. E' \emph{necessaria}, quando n\&o pode deixar de produzir o seu acto; assim \textsc{DEus} ama a sua bondade necessariamente. E' \emph{livre}, quando pode deixar de produzir o seu acto; assim \textsc{DEUS} ama as coisas creadas livremente. --- Na \emph{supposição} que \textsc{DEus} queira uma coisa, não pode deixar de a querer; porque a vontade de \textsc{Deus} é immutavel. Mas, então, é uma necessidade \emph{Aypothetica}, não \emph{absoluta}.}.

\pg{496} 12. \textbf{Deus quer} ou ama tudo \textbf{pela} sua bondade.

a) Uma volição perfeitissima, como é toda a volição. divina, deve ter um perfeitissimo \emph{objecto formal}, ou \emph{motivo}, \emph{e} um perfeitissimo modo de querer. Ora o perfeitissimo \emph{objecto formal}, ou \emph{motivo}, da volic\&o divina não pode deixar de ser a \emph{Bondade divina}, e o perfeitissimo modo de querer é querer ou amar essa Bondade por si mesma e sobre todas as coisas e em todas as coisas, e querer ou amar as outras coisas por ella e nella. Logo \textsc{Deus} quer ou ama tudo pela sua \emph{bondade}.

b) O objecto \emph{formal} dá a cada volição a perfeição especifica, ou quasi-especifica; porque é pelos objectos \emph{formaes} que se especificam os actos. Ora é impossivel que um acto de \textsc{Deus} seja especificado por uma coisa, distincta da sua essencia, e por isso da sua bondade, que é uma e a mesma coisa com a essencia. Logo é impossivel que o objecto \emph{formal} da volição de \textsc{Deus} seja uma coisa distineta da sua bondade. Logo \textsc{Deus} quer ou ama tudo pela sua \emph{bondade}.

c) O objecto \emph{formal}, como toda a causa \emph{final}, refere-se á vontade, como o motor ao movido. Logo se o objecto \emph{formal}, ou o \emph{motivo}, da volição de \textsc{Deus} fosse uma coisa distineta da sua bondade, \textsc{Deus} seria movido e determinado ao acto por um outro motor. Ora repugna que \textsc{Deus}, primeiro motor e infinitamente independente, receba de uma coisa extranha o motivo e o impulso da sua volição. Logo \textsc{Deus} quer ou ama tudo pela sua \emph{bondade} \footnote{Em todas as faculdades devemos distinguir o objecto \emph{formal} e o objecto \emph{material}. O objecto \emph{material} é o que a faculdade attinge; o \emph{formal} 6 0 \emph{motivo}, ou \emph{aspecto}, pelo qual, ou sob o qual, a faculdade attinge o objecto \emph{material}. Dá-se o mesmo com relação á vontade. O objecto \emph{formal}, ou o \emph{Motivo}, exerce uma verdadeira \emph{causalidade}, actua na vontade, e a leva a praticar o acto. --- Quando se trata da vontade do \textsc{Deus}, devemos tambem admittir um objecto \emph{formal}, um \emph{motivo}, uma \emph{razão sufficiente}, pela qual essa faculdade produz o proprio acto; aliás o acto não seria racional. Mas esse motivo, essa \emph{razão sufficiente}, não exerce na vontade divina uma\pgn{497}verdadeira \emph{causalidade}, porque a idéa de \emph{causalidade} desperta a idéa de um \emph{effeito}, de um \emph{producto}, que começa a existir, quando antes não existia. --- Ora o motivo, \emph{sufficiente} e \emph{supremo}, de todas as volições divinas é a \emph{bondade divina}. \textsc{DEus}, assim como vé tudo com um só acto na sua essencia e pela sua essencia, assim tambem com um só acto quer tudo na sua bondade e pela sua bondade. Só este motivo é adequado á perfeição infinita da sua vontade. (\emph{Cf}. \emph{Sum}. \emph{T}\textasciicircum{}., p. I, q. 19, a. 5). --- \textsc{Deus}, nas suas obras, pode ter em vista qualquer outro motivo, como seria o bem, a felicidade do homem; mas o motivo \emph{supremo}, que se encerra nos outros motivos e a todos confere a \emph{appetibilidade}, é sempre a infinita bondade de \textsc{Deus}, que deseja communicar-se a outros entes.\par Advirta-se que \textsc{DEus}, embora tudo queira e opere pela sua bondade e gloria, comtudo ama as creaturas com perfeitissimo amor de \emph{amizade} e de \emph{benevolencia} (que procura o bem da coisa amada); porque a maxima gloria extrinseca, que \textsc{DEus} pode receber das creaturas, está inseparavel e essencialmente unida ao maximo bem, que \textsc{Deus} pode dar á creatura e que consiste na visão da essencia divina e no correspondente amor beatifico.}.

\pg{497} 73. \textbf{Deus} quer ou ama o proprio ser como fim.

a) Um ente quer ou ama um objecto como \emph{fim}, quando esse objecto é amavel \emph{por si mesmo}, e communica a sua amabilidade a todos os outros objectos, que se tornam, por isso, termos da vontade. Ora Dzuvs, sendo a propria bondade infinita, é infinitamente amavel \emph{por si mesmo}, e todas as outras coisas não são amaveis, senão emquanto participam da bondade divina. Logo \textsc{Deus} quer ou ama o proprio ser como \emph{fim}.

b) Um objecto ama-se como \emph{fim}, quando é proporcionado e adequado á natureza e excellencia da faculdade appetitiva; porque essa proporção e adequação termina e satisfaz por completo a tendencia da faculdade. Ora o ser divino, por ser uma bondade infinita, é o «\emph{nico} objecto proporcionado e adequado á natureza e excellencia da vontade divina; porque esta é uma e a mesma coisa com aquelle. Logo \textsc{Deus} quer ou ama o proprio ser como \emph{fim} \footnote{Esta e as seguintes conclusões são simples corollarios da conclusão precedente --- que \textsc{Deus} quer ou ama tudo pela sua bondade.\par 32 Vor. II.}.

74. \textbf{Deus} quer ou ama as outras coisas como meios.

a) Deus quer ou ama as outras coisas. --- \textsc{Deus}, amando o proprio ser \emph{por si mesmo} e como \emph{fim}, porque esse ser é a bondade infinita, não pode deixar de amar tudo o que é participação ou semelhança d'essa bondade; porque não se ama a si mesmo quem não ama a propria semelhança. Ora as outras coisas são participações ou semelhanças da bondade divina. Logo \textsc{Deus} quer ou ama as outras coisas.

b) Deus quer ou ama as outras coisas como meios. --- \textsc{Deus}, pelo \pg{498}  facto de amar, como \emph{fim}, o proprio ser, que é bondade \emph{infinita}, não póde deixar de amar tudo o que, de algum modo, concorre para a perfeição da mesma bondade. Ora as outras coisas, embora, por serem finitas, não possam concorrer para o augmento da bondade infinita, comtudo, por serem semelhanças d'essa bondade, concorrem para a sua manifestação e glorificação; e por isso têm razão de \emph{meios}. Logo \textsc{Deus} quer ou ama as outras coisas como \emph{meios} \footnote{Se \textsc{DEus} quer e ama o proprio ser como \emph{fim} e as outras coisas como \emph{meios}, é claro que o ser divino é o objecto \emph{primario} da vontade divina, e as outras coisas constituem o objecto \emph{secundario}. De facto, objecto \emph{primario} de uma faculdade é o que é attingido \emph{por si}; o objecto \emph{secundario} é o que não é attingido \emph{por si mesmo}, mas por meio de \emph{outro} objecto. Ora a vontade divina ama o ser divino \emph{por si mesmo}, porque é bondade infinita, e não ama as outras coisas senão porque se ama a si mesmo. --- D'onde se vê que a tendencia para o bem e a propria posse do bem não excluem o amor de outras coisas, quando estas estão para esse bem, como meios para o fim.\par Como o bem se encontra propriamente no ente e a verdade na intelligencia, o objecto \emph{primario} da \emph{intelligencia} e o objecto \emph{primario} da vontade --- nalguma coisa conyém entre si, --- e noutra differem. Convém em que, tanto para uma faculdade como para outra, objecto \emph{primario} é o que é attingido \emph{por} si, e não por outro ou noutro. --- Differem em que, emquanto o objecto \emph{primario} da intelligencia é constituido por uma coisa, diversa segundo a diversidade da intelligencia, o objecto \emph{primario} da vontade não é constituido por uma coisa, diversa segundo a diversidade da vontade. Com effeito, o objecto \emph{primario} da intelligencia divina é o \emph{ser divino}, e o objecto \emph{primario} da inteligencia humana, no estado da união com o corpo, é a \emph{essencia abstracta das coisas materiaes}; emquanto o objecto \emph{primario} tanto da vontade divina como da humana é a \emph{felicidade}, que consiste na posse de \textsc{Deus}, ainda que a nossa vontade tenda para \textsc{Deus} emquanto tende para a felicidade e não emquanto Deus é o verdadeiro objecto da felicidade, porque tendemos para o bem como o concebemos.\par Embora \textsc{Deus} e \& creatura racional queiram alguns bens, como meios para o fim, comtudo a tendencia da vontade creada para esses objectos \emph{secundarios} differe da tendencia da vontade divina para os mesmos objectos. Com effeito, \emph{a} vontade creada quer os meios, para que a levem á consecução da felicidade, que não possue mas que deseja possuir; e por isso os taes meios tem a indole ou o caracter de \emph{bens uteis}. \textsc{DEUS}, porém, não quer os meios, isto é, os objectos \emph{secundarios}, para que o levem ao alcance da felicidade, porque a felicidade de \textsc{Deus} é a sua essencia, que ama e em que se deleita, mas quer esses objectos para que lhes possa communicar a propria bondade, e d'este modo manifeste exteriormente as riquezas d'essa bondade. Portanto, emquanto a creatura se serye das coisas, que são ordenadas ao fim, para vantagem propria, \textsc{DEUS} serve-se de nós para a sua bondade e para a vantagem nossa. Escutemos S. Thomaz: « Sicut Deus cognoscendo essentiam suam, cognoscit omnia, quae sunt ab eo, in quantum sunt similitudo quaedam veritatis ejus; ita etiam volendo vel amando essentiam suam, vult omnia, quae sunt ab eo, in quantum habent similitudinem bonitatis eius.\pgn{499}Unde id quod est volitum \emph{primo} ab eo, est bonitas sua tantum. Alia vero vult in ordine ad bonitatem suam; non autem hoc modo ut pér ea aliquid bonitatis acquirat, sicut nos facimus circa alios bene operando, sed ita quod eis de bonitate sua aliquid largiatur; et ideo liberalitas est quasi proprium ipsius, secundum \emph{Avicennam}; quia ex operatione sua non intendit aliquod Sibi commodum provenire, sed vult bonitatem suam in alios diffundere; et ideo \emph{Augustinus} dieit quod ipse utitur nobis ad bonitatem suam et utilitatem nostram » (I \emph{Sent}. \emph{Dist}. 45, q. 1, a. 2).\par Contra a ultima proposição --- que \textsc{DEus} quer ou ama outras coisas --- levantem-se muitas difficuldades. Apresentamos as principaes.\par Obj. 1.º --- Quem possue o bem, que é objecto \emph{primario} da vontade, não pode desejar outra coisa; porque esse bem constitue o fim ultimo, a felicidade. Ora \textsc{Deus} possue a sua bondade infinita, que é objecto \emph{primario} da sua vontade. Logo não pode amar outras coisas.\par Resposta. O bem, que constitue o objecto da felicidade, não pode ser senão o objecto \emph{primario} da vontade, porque este e só este tem a indole, o caracter de \emph{fim ultimo}. E, como o fim ultimo é um só, é evidente que a vontade, a qual possue aquelle bem, não pode querer outro objecto \emph{primario}. Todavia isso não obsta a que a mesma vontade possa querer outros objectos ou bens \emph{secundarios} por amor do bem \emph{primario}, para o qual estão ordenados como meios para o fim; e, neste caso, os bens \emph{secundarios} constituem o objecto \emph{material} da vontade, porque o objecto \emph{formal} é constituido pelo \emph{fim}, de que deriva toda a razão de \emph{bem}. Por isso \textsc{DEUs}, a quem basta a sua bondade, pode querer, por amor da sua mesma bondade, outras coisas, que a Elle se refiram. Diz S. Thomaz: « Ex hoc quod voluntati divinae sufficit sua bonitas, non sequitur quod nihil aliud velit, sed quod nihil aliud vult nisi ratione suae bonitatis. Sicut etiam intellectus divinus, licet sit perfectus ex hoc ipso quod essentiam divinam cognoscit, tamen in es cognoscit alia ». (\emph{Sum}. \emph{Th}., p. I, q. 19, a. 2 ad 3).\par Obj. 2.º --- \textsc{Deus}, desde a eternidade, possua o seu fim, que é a sua bondade, e por isso não pode querer outras coisas; porque quem quer uma coisa pelo fim, deixa de querel-a, desde que seja alcançado o fim.\par Resposta. \textsc{DEus}, embora desde a eternidade possua o seu fim, que é a sua propria bondade, comtudo pode querer e amar outras coisas, como meios ordenados para o fim. Por quanto, como advertimos, a vontade divina não dispõe e ordena os meios para o fim pelo mesmo motivo, por que os dispõe e ordena a vontade humana. Esta dispõe e ordena os meios para o fim, com o intuito de alcançar um fim ou um bem, que não possue, e por isso, alcançado o fim, cessa o uso e o amor dos meios. A vontade divina, pelo contrario, quer as outras coisas pela sua bondade, emquanto, por amor da sua bondade, confere a outras coisas uma certa bondade, que seja semelhança e participação do bem absoluto, e por isso o amor dos \emph{meios} não só não cessa, quando se possue o /\emph{im}, mas funda-se propriamente na posse do mesmo /\emph{im,ou} no \emph{fim} emquanto possuido. Por isso mesmo que é a bondade infinita, \textsc{Deus} a derrama largamente sobre as outras coisas; e, amando a sua bondade, produz\pgn{500} a bondade nos outros. Diz S. Thomaz: « \textsc{DEus} non ordinat creaturas in finem bonitatis suae, quasi per eas bonitatem suam assequatur, sed ut ipsae creaturae divina operatione similitudinem aliquam divinae bonitatis acquirant, quod esse non posset, nisi eo volente et faciente » (I \emph{Sent}. Dist. 4b, q. 1, a. 2 ad 4).\par Obj. 8.º --- Quando são varios e diversos os \emph{objectos} da vontade, varios tambem e diversos são os \emph{actos} da mesma faculdade. Ora em \textsc{DEUS} não pode haver senão um unico acto, como não ha senão um unico ser. Logo \textsc{Deus} não quer, não pode querer as outras coisas.\par Resposta. A unidade ou multiplicidade dos actos de uma faculdade depende, como dissemos, da unidade ou multiplicidade do objecto \emph{formal}, e não do objecto \emph{material}; porque é pelo objecto \emph{formal} que se especificam e distinguem os actos das faculdades. D'onde se segue que a unidade do acto é compativel com a multiplicidade dos objectos \emph{materiaes}, quando muitos objectos materiaes são attingidos sob um só aspecto \emph{formal}, segundo o qual movem e determinam a faculdade. Por isso, assim como a intelligencia divina percebe todas as coisas por um só acto, porque é determinada pela sua essencia, que é a forma intelligivel, representativa de todas as coisas, assim tambem a divina vontade quer todas as coisas com um só acto, porque é levada a querer pela sua bondade, a qual é a unica razão, por que \textsc{Deus} ama primeiramente o seu ser, e segundariamente as outras coisas. Diz S. Thomaz: «Sicut intelligere divinum est unum, quia multa non videt nisi in uno; ita velle divinum est unum et simplex, quia multa non vult nisi per unum, quod est bonitas sua » (\emph{Sum}. \emph{Th}., P- 1, q. 19, a. 2 ad 4).\par Obj. 4.4 --- Não parece \emph{conveniente} que a infinita Magestade de \textsc{Deus} se abaixe até ás creaturas, amando-as e communieando-Ihes a sua bondade.\par Resposta. Negamos que não \emph{convenha} á bondade divina communicar aos outros os proprios thesouros. Se as creaturas, tão pobres em bens e em generosidade, communicam aos outros os proprios bens, como não deverá admittir-se esse tendencia em Dus, tão rico e tão generoso? Nem devemos receiar pela Magestade de \textsc{Deus}. Elle encontra, na sua infinita sabedoria, o modo de se communicar, conveniente á sua grandeza e dignidade, e do acto da sua generosidade tirará sempre a sua maior gloria. Mas escutemos um magnifico trecho de S. Thomaz: « Res naturalis non solum habet naturalem inclinationem respectu proprii boni, ut acquirat ipsum cum non habet, vel ut quiescat in illo cum habet; sed etiam ut proprium bonum in alia diffundat, secundum quod possibile est. Unde videmus quod omne agens, inquantum est actu et perfectum, facit sibi simile. Unde et hoc pertinet ad rationem bonitatis, ut bonum quod quis habet aliis communicet, secundum quod possibile est. Et hoc praecipue pertinet ad bonitatem divinam, a qua per quamdam similitudinem derivatur omnis perfectio. Unde si res naturales, inquantum perfectae sunt, suum bonum. aliis communicant, multo magis pertinet ad voluntatem divinam ut bonum suum aliis per similitudinem communicet, seeundum quod possibile est (scilicet per similitudinem \emph{imperfectam}. et \emph{analogam}). Sic igitur vult et se et alia: sed se ut finem, alia vero ut ad finem, inquantum condecet divinam bonitatem etiam alia ipsam participare » (Sum. \emph{TA}., p. I, q. 19, a. 2).}.

\pg{499} 75. \textbf{Deus} quer ou \textbf{ama necessariamente} o proprio ser e \textbf{livremente as outras coisas}.

a) \textbf{Deus} quer ou ama necessariamente o proprio ser, --- Nenhuma faculdade pode ser indifferente com relação ao objecto, que \pg{500}  lhe é adequado e que por isso lhe satisfaz todas as tendencias; porque, n'esse caso, a indifferenca seria uma desordem.

\pg{501} Ora o ser divino, sendo a propria bondade ou perfeição, é o unico objecto adequado á vontade divina e satisfaz plenamente todas as tendencias d'ella; porque em tal objecto a vontade encontra tudo o que pode desejar. Logo \textsc{Deus} quer ou ama \emph{necessariamente} o proprio ser \footnote{Este argumento pode tambem propor-se d'este modo. O ser divino, por ser a propria bondade infinita, constitue o objecto da felicidade de \textsc{Deus}. Ora a vontade não pode deixar de amar o objecto, que constitue a propria felicidade e que é conhecido como tal. Logo \textsc{Deus} quer ou ama necessariamente o proprio ser. --- Nem podia ser diversamente. O ser divino, sendo infinitamente perfeito, é infinitamente amavel. E poderia a vontade divina negar-lhe o seu amor? Não. E' certo que a vontade não quer ou não ama sem razão sufficiente. Basta reflectir um instante nas imperfeições dos bens, que nos cercam, para dar á vontade motivos sufficientes para os rejeitar. Mas a intelligencia divina, que penetra com o seu olhar no oceano infinito das amabilidades do ser divino, não encontra o minimo defeito, a mais leve imperfeição, que possa deixar indifferente a vontade. Esta, portanto, não pode deixar de amar o bem infinito, o unico digno d'ella, e de o amar na medida da sua bondade, e por isso infinitamente e sem medida. Em \textsc{Deus} ha uma equação perfeita entre a intelligencia, o amor e a plenitude do ser divino, objecto d'este conhecimento e d'este amor.}.

b) \textbf{Deus} quer ou ama livremente as outras coisas, --- Todo o bem, que --- ou não é \emph{adequado} á vontade, isto é, não enche plenamente a capacidade d'essa faculdade, --- ou não tem um nexo \emph{necessario} com esse objecto adequado, não attrahe \emph{necessariamente} a mesma vontade, mas deixa-a na sua indifferenga ou liberdade. Ora nenhum bem \emph{finito}, --- por isso mesmo que é finito, é objecto \emph{adequado} á divina vontade, que é infinita, --- ou tem um nexo \emph{necessario} com o objecto \emph{adequado} da vontade divina, que é constituido pela propria bondade infinita, pois o infinito é independente do finito. Logo \textsc{Deus} quer ou ama \emph{livremente} as outras coisas \footnote{Alem d'isso, pertence á perfeição da vontade, que ella tenda para o objecto em conformidade com a indole e o merito do mesmo objecto; de modo que deve tender \emph{necessariamente} para o objecto \emph{necessario}, e \emph{livremente} para o objecto \emph{livre} ou \emph{contingente}. Ora o ser divino, por ser a propria bondade infinita, apresenta-se como o objecto summamente bom e necessario, e as outras coisas, que existem fóra de \textsc{Deus}, apresentam-se como bens finitos, e contingente ou livremente amaveis. Logo a vontade divina, sendo infinitamente perfeita, deve amar \emph{necessariamente} o proprio ser divino e \emph{livremente} as.creaturas. --- Ha mais. Tudo o que \textsc{DEUS} ama, ama-o emquanto se ama a si mesmo; porque, sendo ordenado o seu amor, não pode amar coisa alguma. senão emquanto está ligada com o ultimo fim, de que deriva toda a amabilidade para os outros bens. Ora \textsc{Deus} conhece perfeitamente que nenhum amor de coisa finita está necessariamente ligado com o amor do proprio ser; porque a sua bondade é absoluta e perfeitissima, e não precisa\pgn{502} de se associar com quelquer outro bem. Logo é impossivel que \textsc{Deus} queira ou ame necessariamente outros bens, distinctos do proprio bem. --- Finalmente, a propria tendencia necessaria e irresistivel da vontade divina para o bem infinito é a razão da contingencia ou liberdade da sua tendencia para os bens finitos e imperfeitos. A. infinita intelligencia de \textsc{Deus} descobre nestes bens uma multidão de motivos, que, se o podem levar a querel-os, podem-no tambem levar a não querel-os, ficando sempre livre a respeito d'elles. \textsc{DEus} n\&o precisa de nenhum bem, de nenhuma coisa finita, para \& sua perfeita felicidade; e por isso Elle é tao independente no seu amor, como o é no seu ser. Emquanto nós amamos as coisas \emph{porque} são \emph{bóas}; \textsc{DEUS} ama as coisas \emph{para} as \emph{fazer bóas}, e dóar-lhes com o ser essa bondade, que nunca teriam tido sem Elle. Emquanto o nosso amor vai procurando o bem para o abraçar e encontrar na posse d'elle a felicidade, o amor divino \emph{produz} o bem e derrama-o sobre todas as coisas. Se, pois, \textsc{Deus} não depende, de modo algum, das creaturas, mas estas dependem, em tudo e por tudo, de \textsc{Deus}, é claro que o amor de \textsc{Deus} para com as creaturas é plenamente \emph{livre}.\par Mas escutemos S. Thomaz, que prova ambas as partes da nossa proposição. Diz assim: « Voluntas divina \emph{necessariam} habitudinem habet ad bonitatem suam, quae est \emph{proprium} ejus objectum. Unde bonitatem suam Deus \emph{ec necessitate} vult; sicut et voluntas nostra ex necessitate vult beatitudinem, sicut et quaelibet alia potentia habitudinem habet ad proprium et principale objectum, ut visus ad colorem, quia de sui ratione est ut ad illud tendat. Alia autem a se Deus vult in quantum ordinantur ad suam boniiatem, ut in finem. Ea autem quae sunt ad finem, non ex necessitate volumus, volentes finem, nisi sint talia, sine quibus finis esse non potest; sicut volumus cibum, volentes conservationem vitae, et navem, volentes transfreiare. Non sic autem ex necessitate volumus ea, sine quibus esse potest, sicut equum ad ambulandum, quia sine hoc possumus ire; et eadem ratio est in aliis. Unde cum bonitas Dei sit perfecta, et esse possit sine aliis, cum nihil ei perfectionis ex aliis accrescat, sequitur quod alia a se eum velle non sit necessarium \emph{absolute}. Et tamen est necessarium \emph{ec suppositione}; supposito enim quod vult, non potest non velle, quia non potest voluntas eius mutari » (\emph{Sum}. \emph{Th}., p. I, q. 19, a. 3).\par Advirta-se o nexo d'esta proposição com a precedente. Na precedente, demonstrámos que \textsc{Deus}, com o mesmo simplicissimo acto da sua vontade, quer o proprio ser como /\emph{im} e as outras coisas como \emph{meios}. Nesta proposição, provamos que o mesmo simplicissimo acto, --- que, considerado em relação ao seu \emph{principio}, \emph{é necessario}, porque é uma e a mesma coisa com a essencia, a qual é necessaria, --- considerado em relação ao seu \emph{termo}, \emph{é necessario}, emquanto se refere á divina bondade, que é o seu objecto \emph{primario} e a razão pela qual ama todas as coisas, mas é \emph{contingente} ou \emph{livre}, emquanto se refere ás outras coisas, que são o seu objecto \emph{secundario}. --- Por isso esta questão attinge propriamente a \emph{liberdade divina}. Por quanto, se se admitte que o acto da vontade divina tem uma ordem ou relação \emph{necessaria} a \emph{tudo} o que \textsc{Deus} quer, \textsc{Deus} não pode chamar-se \emph{livre}. Se, pelo contrario, o mesmo acto da vontade divina tem uma relação \emph{necessaria} á divina bondade, de\pgn{503}modo que não pode deixar de a amar, e tem uma relação não \emph{necessaria} ás outras coisas, de modo que pode deixar de as amar, então deve admittir-se em \textsc{Deus} uma perfeitissima \emph{liberdade}, que não importa \emph{contingencia} e \emph{potencialidade} da parte do acto da vontade, mas da parte das coisas, queridas por \textsc{Deus}. --- \emph{Que seja assim}, isto é, que o mesmo acto da divina vontade seja \emph{necessario} em ordem á divina bondade, e seja \emph{livre} em ordem ás outras coisas, entende-se facilmente. \emph{Como possa ser assim}, isto é, como é que um e o mesmo acto acto é \emph{necessario} e \emph{livre} em ordem a termos diversos, não se comprehende pela limitação da nossa intelligencia e pela transcendencia das operações divinas. Alguma coisa dissemos a este respeito, tratando da \emph{immutabilidade} de \textsc{Deus}. Accrescentamos uma outra consideração. O acto da vontade divina não é \emph{causado} nem \emph{especificado} pelas coisas queridas, que \& elle se referam de um modo contingente; porque estas não são \emph{necessariamente} ligadas com o fim, isto é, com a bondade divina, que pode subsistir sem ellas. Pelo contrario, o acto livre da vontade creada é \emph{causado} e \emph{especificado} pelo seu \emph{objecto}, porque depende d'este, e por isso a vontade creada não pode, com o mesmo acto, querer e deixar de querer, querer necessariamente e querer livremente. Serdo, pois, o acto das creaturas distincto da vontade, e referindo-se ao seu objecto como o effeito á causa, e multiplicando-se, por isso, conforme a multidão dos objectos e a diversa relação que tem com elles, não admira que o acto da vontade divina, --- que é uma e a “mesma coisa com a propria vontade e que de maneira nenhuma é causado pelos objectos, que quer livremente, --- possa ser um e o mesmo, seja qual fór a variedade dos objectos queridos, e a diversa relação que estes têm com elle. Concluimos que a unica e simplicissima volição divina, por ser infinita, pode corresponder e corresponde a muitas e diversas volições, e é \emph{necessaria} ou \emph{livre}, conforme o seu objecto é o bem infinito ou o bem limitado.\par Contra a segunda parte da proposição --- que \textsc{Deus} ama \emph{livremente} as outras coisas --- levantam muitas objecções. Expomos as principaes.\par Dizem: \textsc{DEus} \emph{necessariamente} conhece tudo o que conhece, porque a sua sciencia se identifica com a essencia. Ora tambem a vontade de \textsc{Deus} identifica-se com a sua essencia. Logo \textsc{DEus} tambem quer \emph{necessariamente} tudo o que quer.\par Resposta. Ha uma grande differenga entre o acto da intelligencia e o da vontade, e esta differenga deriva da diverse natureza do conhecimento e do appetite. Com effeito, o conhecimento faz-se emquanto o conhecido se encontra no cognoscente; por isso \textsc{DEUS} conhece as coisas emquanto existem na sua essencia; e, como as coisas existem \emph{necessariamente} na essencia divina, assim \emph{necessariamente} \textsc{Deus} as conhece. Pelo contrario, a vontade refere-se ás coisas, emquanto e como estas existem em si mesmas, e, como as coisas não são em si necessarias, \textsc{DEUS} não as quer \emph{necessariamente}. Por isso, devemos distinguir a proposição \emph{maior}, e dizer que \textsc{Deus} entende \emph{necessariamente} as coisas, não porque a sua sciencia se identifica com a essencia, mas porque a sciencia refere-se ás coisas conhecidas, emquanto são na essencia, e, como nesta existem \emph{necessariamente}, \emph{necessariamente} \textsc{DEUS} as conhece. Como tambem devemos distinguir a proposição \emph{menor} e dizer que a\pgn{504}vontade de \textsc{Deus} se identifica com a essencia, mas nem sempre se refere ao que está na mesma essencia, porque pode referir-se a coisas, que estão fóra da essencia, e que por isso não são queridas \emph{necessariamente} por \textsc{Deus}. --- Devemos advertir que o acto de entender importa uma perfeição no ser intelligente, e o acto de querer importa uma perfeição na coisa querida, e que, por isso, se \textsc{DEus} não entendesse tudo, não seria absolutamente perfeito, mas, se não quizesse as outras coisas, não deixaria de ser perfeito. --- 8. Thomaz resolve a difficuldade proposta com estas palavras: « Sicut divinum esse in se est necessarium, ita et divinum velle et divinum scire; sed divinum scire habet necessariam habitudinem ad scita, non autem divinum velle ad volita. Quod ideo est, quia scientia habetur de rebus secundum quod sunt in sciente; voluntas autem comparatur ad res secundum quod sunt in se ipsis. Quia igitur omnia alia habent necessarium esse secundum quod sunt in Deo; non autem secundum quod sunt in seipsis habent necessitatem absolutam, ita quod sint per se ipsa necessaria; propter hoc Deus quaecumque scit, ex necessitate scit; non autem quaecumque vult, ex necessitate yult » (\emph{Sum}. \emph{Th}, p. I, q. 19, a. 3 ad 6).\par Continuam: Se \textsc{Deus} não quer \emph{necessariamente} as creaturas, pode deixar de as querer. Ora este consequente é falso, porque quem pode deixar de querer o que quer, tem uma vontade \emph{mutavel}. Logo \textsc{Deus} quer \emph{necessariamente} as outras coisas, isto é, as creaturas.\par Resposta. Distinguimos a proposição \emph{maior}, dizendo --- ser verdade que, se Deus quer \emph{livremente} as creaturas, pode deixar de as querer em sentido \emph{absoluto}, --- mas ser falso que, se \textsc{Deus} quer \emph{livremente} as creaturas, pode deixar de as querer mesmo na \emph{supposição} que as tenha querido. Porquanto, tudo o que \textsc{Deus} quer, tanto \emph{necessaria} como \emph{livremente}, é desde a eternidade que o quer; porque o acto da vontade divina é medido, como o seu ser, pela eternidade. Mas ha esta differença entre a vontade necessaria e \& livre: a vontade necessaria refere-se por tal modo ao seu objecto, que não pode absolutamente deixar de o querer, emquanto a vontade livre refere-se por tal modo ao seu objecto, que absolutamente podia não tel-o querido, embora, se por \emph{supposição} o quiz, não possa deixar de o ter querido ou de o querer. E esta necessidade \emph{Aypothetica} ou de \emph{supposição} chama-se necessidade de \emph{immutabilidade}, porque é propriamente a \emph{immutabilidade} da vontade divina que exige que \textsc{Deus} não possa deixar de querer o que desde a eternidade quiz. Portanto a necessidade /\emph{iypothetica} ou de \emph{supposição} não tem o mesmo caracter em \textsc{DEUS} e nas creaturas; porque nas creaturas suppõe Successão e mudança; o que não suppõe em \textsc{Deus}.\par Insistem: \emph{Convém} á divida bondade communicar aos outros a propria perfeição. Ora \textsc{Deus} não pode deixar de fazer o que \emph{convém} á sua bondade. Logo \textsc{Deus} quer \emph{necessariamente} as outras coisas.\par Resposta. São sobretudo os sequazes de \emph{Rosmini} que propugnam um nexo \emph{necessario}, não \emph{physico} mas \emph{moral}, entre a divina bondade e a existencia das creaturas (Cf. Rosmini, \emph{Teosofia}, v. l) Mas o erro ó manifesto. Ha duas especies de \emph{conveniencia} ou \emph{decencia}: uma, quando o contrario \emph{não convém}; outra, quando tambem o contrario \emph{convém}. Ora não é a \emph{primeira} especie que se verifica na creação, mas é a \emph{segunda}. Com effeito, é \emph{conveniente}\pgn{505} que a bondade divina, sendo infinita, communique aos outros os seus thesouros; mas, se \textsc{Deus} os não quizer communicar, não ha nisso nenhum \emph{inconveniente} para a sua bondade, sendo esta perfeitissima e absolutamente independente de qualquer creatura. Inconveniente e absurdo seria o acto divino, se se admittisse nelle uma \emph{necessidade}, embora \emph{moral}, a respeito das creaturas; porque, em tal caso, Dxzus seria dependente das mesmas creaturas. --- Em poucas palavras: Assim como \emph{convém} que \textsc{Deus} queira outras coisas, porque Elle é a fonte inexhaurivel de todo o bem, assim tambem convém que \textsc{DEus} não queira outras coisas, porque Elle é um bem infinito, suffücientissimo a si mesmo. A \emph{conveniencia} ou \emph{decencia} de \textsc{Deus} é compativel tanto com a hypothese da creação, como com a de \emph{não-creação}.\par *\par Das coisas expostas resultam os seguintes corollarios:\par a) \textsc{Deus} é infinitamente \emph{bemaventurado}. \emph{A bemaventurança}, segundo \emph{Boecio} (\emph{De cons}., 1. IV) é « estado perfeito pela reunião de todos os bens --- \emph{status omnium bonorum aggregatione perfectus} ». --- A \emph{bemaventurança} encerra --- a perfeição absoluta, --- o conhecimento d'esta perfeição, --- a complacencia na sua posse, --- a isenção de todos os males, presentes e futuros, --- e o dominio supremo sobre todas as coisas contingentes. Ora \textsc{DEus} --- contém na sua simplicissima essencia a \emph{plenitude} da perfeição, --- conhece e comprehende \emph{adequadamente} a propria essencia e nella todas as coisas, --- quer e ama \emph{necessariamente} a propria bondade ou perfeição, --- não pode, por ser immutayvel, perder algum bem ou adquirir algum mal, --- tem o dominio supremo e absoluto sobre todas as creaturas. Logo \textsc{Deus} é \emph{bemaventurado}. E é bemaventurado \emph{infinitamente}, porque \emph{infinita} é a sua \emph{perfeição}, \emph{infinito} é o \emph{conhecimento} da sua intelligencia, \emph{infinito} é o \emph{amor} da sua vontade, e por isso é a propria \emph{bemaventuranca}. --- \textsc{DEUS} não só é a bemaventurança de si mesmo, mas é tambem a bemaventurança de todas as creaturas intelligentes. Por quanto, \emph{a} bemaventurança das creaturas não pode deixar de ser uma participação finita da infinita bemaventuranca de \textsc{DEUS}, summo bem e fonte de todo o bem. --- \textsc{Deus} é a bemaventuranoa \emph{objectiva} das creaturas, emquanto é conhecido, amado e possuido (\emph{Sum}. \emph{Th}., p. I, q. 21).\par \emph{b}) \textsc{Deus} é dotado de \emph{liberdade}. E' uma consequencia do que dissemos \&cerca da vontade de \textsc{DEus} em relação ás creaturas. Diz S. Thomaz: « Liberum arbitrium habemus respectu eorum, quae non \emph{necessario} volumus, vel naturali instinctu. Non enim ad liberum arbitrium pertinet quod volumus esse felices, sed ad naturalem instinctum. Unde et alia animalia, quae naturali instinctu moventur ad aliquid, non dicuntur libero arbitrio moveri. Cum igitur Deus ex necessitate suam bonitatem. velit, alia vero non ex necessitate, \emph{respectu illorum}, \emph{quae mon em necessitate vult}, \emph{liberum}. \emph{arbitrium habet} » (\emph{Sum}. \emph{Th}., p. I, q. 19, a. 10). --- Esta liberdade é \emph{perfeitissima}. Com effeito, o maior ou menor grau de perfeição na liberdade consiste na maior ou menor independencia, de que a vontade é dotada com relação a um determinado bem. Ora \textsc{Deus} é infinitamente independente de todos os bens creados, porque de nada precisa para a sua felicidade. Logo a liberdade de \textsc{Deus} é \emph{perfeitissima}. O \emph{poder peccar} não constitue a \emph{essencia} da liberdade,\pgn{506}mas, antes, é um seu \emph{defeito}, \emph{e} por isso não pode encontrar-se em \textsc{Deus}. Diz B. Thomaz: « Quod liberum arbitrium diversa eligere possit, servato ordine finis, hoc pertinet ad \emph{perfectionem} libertatis ejus; sed quod eligat aliquid, divertendo ab ordine finis, hoc pertinet ad \emph{defectum} libertatis » (\emph{Sum}. \emph{Th}., Pp: 1, q. 62, a. 8 ad 8).\par c) \textsc{Deus} não quer o mal \emph{physico} de um modo \emph{directo}. Porquanto, a vontade divina não tem e não pode ter outro objecto a vão ser a sua infinita bondade, e, se ama outras coisas, ama-as emquanto participam da mesma bondade. Ora todo o mal, quer \emph{physico} quer \emph{moral}, \emph{é privação} de algum \emph{bem}, e por isso, não tendo em si uma razão de appetibilidade, não pode ser objecto \emph{directo} da vontade divina. --- Todavia, se \textsc{Deus} não pode querer \emph{directamente} o mal \emph{physico}, pode querel-o \emph{indirectamente}, emquanto esse mal esta \emph{accidentalmente} unido a um \emph{bem commum} ou a um bem de uma \emph{ordem mais elevada}; assim não repugna que \textsc{Deus}, querendo a \emph{justiça}, que é um bem, queira a \emph{pena} do peccador, que é um mal \emph{p/ysico}, e, querendo a ordem do mundo, queira a natural corrupção de muitas coisas. --- Mas o mal \emph{moral}, que é o \emph{peccado}, \textsc{DEUS} não o quer \emph{de modo algum}. O mal \emph{moral} contem a aversão da creatura do seu Creador, para quem, emquanto é Summo Bem, todas as coisas devem estar naturalmente ordenadas. Ora \textsc{DEUS} não pode \emph{querer}, nem \emph{indirectamente}, tudo o que se oppõe á essencial ordem das coisas. Logo \textsc{Deus} não \emph{quer}, nem \emph{indirectamente}, \emph{o} mal \emph{moral}, Não o quer, mas \emph{permitte}, e esta \emph{permissão}, diz S. Thomaz, é um \emph{bem}: « Deus vult \emph{permittere} mala fieri: et hoc est \emph{bonum} » (\emph{Sum}. \emph{Th}., p- T, q. 19, a. 9).}.

16. A \textbf{vontade} divina é \textbf{causa das coisas creadas}. --- Se Dus, na creação, operasse --- não pela deliberação da vontade, --- mas pela. necessidade da natureza, operaria sempre a mesma \pg{502} \pg{503}  coisa e sempre do mesmo modo, e produziria tudo o que pode produzir, a não ser que encontrasse impedimento; porque assim opera quem opera por necessidade de natureza. Ora \pg{504}  Dzvus produziu muitas e diversas coisas, e por muitos e diversos modos, nem produziu tudo o que pode produzir, embora \pg{505}  não possa ser e não seja contrariado por ninguem. Logo \textsc{Deus} não opera pela necessidade da natureza, mas pela deliberação \pg{506}  da sua vontade, Logo a vontade divina é causa das coisas creadas \footnote{Vejam-se outras razões, que traz S. Thomaz sobre este assumpto (\emph{Sum}. \emph{Th}., p. I, q. 19, a. 4), especialmente a \emph{primeira}, com que prova que na ordem das causas efficientes, o agente pela intelligencia e pela vontade deve preceder o agente pela natureza, porque só o agente pela intelligencia e pela vontade pode determinar o fim ao agente pela natureza, e todo o agente pela natureza tende para o seu fim; e por isso \textsc{Deus}, que é o primeiro na ordem dos agentes, deve operar pela inteligencia e pela vontade. --- A verdade d'esta proposição torna-se ainda mais clara, se se considera que a potencia divina, por ser infinita, não é por si determinada para produzir esta ou aquella coisa, e que a sciencia divina, sendo por si indifferente com relação a esta ou aquella especie de seres, não pode determinar a potencia a erear este ou aquelle ser. Logo só a vontade pode determinar tal indifferenca, decretando que, entre possiveis infinitos, se produza este ou aquelle ser. (Cf. I \emph{Sent}., Dist. 45, a. 8).\par Sendo necessario, para a producção das coisas, o intervento da intelligencia e da vontade de \textsc{DEus}, pode perguntar-se: Qual das duas faculdades é á causa \emph{mais proxima} do effeito? --- Respondemos que a causa mais proxima é a \emph{vontade}. A sciencia nada pode fazer por si, sem o imperio da vontade, que tudo determina e move. (Cf. \emph{Sum}. \emph{Th}., P. I, q. 19, a. 4 ad 4).}.

77. \textbf{Propriedades do amor} divino. --- O amor divino é \emph{substancial}, \emph{immutavel}, \emph{efficacissimo}, \emph{benevolo}.

a) Eº \emph{substancial}. Em \textsc{Deus} o amor, como o acto intellectual, identifica-se com o ser. Por isso \textsc{Deus} é \emph{amor}.

\pg{507} b) E' \emph{immutavel}. Identificando-se com o proprio ser, o amor divino é immutavel, como immutavel é o ser.

c) E' \emph{efficacissimo}. O homem ama as coisas, porque são bóas; \textsc{Deus}, pelo contrario, ama as coisas, para as tornar bóas. Por isso, o nosso amor é \emph{effeito} da bondade das creaturas; mas o amor divino é \emph{causa} de toda a bondade, de que são dotados os seres finitos \footnote{Por isso \textsc{DEUS} não ama as creaturas com amor \emph{affectivo}. O amor diz-se \emph{affectivo}, quando se ama um objecto pela perfeição, que elle encerra em si e por si; de modo que o amor é um \emph{effeito} da mesma perfeição. Ora \textsc{Deus} não ama nem pode amar as coisas creadas por alguma perfeição, que ellas possuam em si e por si; porque as creaturas, antes de serem objecto da vontade divina, não tinham em si perfeição alguma. Logo \textsc{Deus} não ama as coisas creadas com amor \emph{affectivo}. --- Se \textsc{DEUS} não ama as creaturas com amor \emph{affectivo}, ama-as com amor \emph{e/fectivo}. O amor é \emph{effectivo}, quando emtanto se ama un objecto, emquanto lhe se confere gratuitamente alguma perfeição. Ora \textsc{Deus} ama as creaturas, emquanto lhes concedeu todas as perfeições, de que são dotadas. Logo \textsc{Deus} ama as coisas creadas com amor \emph{effectivo}. Diz S. Thomaz: « Cum amare nihil aliud sit quam velle bonum alicui, manifestum est quod Deus omnia quae sunt amat, non tamen eo modo sieut nos. Quia enim voluntas nostra non est causa bonitatis rerum, sed ab ea movetur sieut ab objecto; amor noster, quo bonum alicui volumus, non \emph{est causa bonitatis} ipsius; sed e converso bonitas ejus, vel vera vel aestimata, \emph{provocat amorem}, quo ei volumus et bonum conservari quod habet, et addi quod non habet, et ad hoc operamur. \textsc{SED} AMOR Dri EST IN- FUNDENS ET CREANS BONITATEM IN REBUS » (\emph{Sum}. \emph{T}., p. 1, q. 20, a. 2).\par Por isso --- 1.9) \textsc{Deus} ama \emph{todas} as coisas \emph{existentes}. Porquanto, estas, emquanto \emph{existem}, são bóas; porque todo o ente é bom, e toda a perfeição é um bem. Ora as coisas não existem senão pela vontade de \textsc{Deus}. Logo todo o bem, que as creaturas existentes possuem, deriva da vontade de \textsc{Deus}, e por isso a \emph{todas} as creaturas existentes \textsc{Deus} quer algum bem. Ora querer um bem é amar. Logo \textsc{DEus} ama \emph{todas} as creaturas \emph{existentes}. (\emph{Sum}. \emph{Th}. ib.) --- 2.º) \textsc{Deus} mais ama as creaturas, que são \emph{mais perfeitas}. Uma creatura é melhor que outra, quando é dotada de natureza mais nobre, ou de uma bondade mais perfeita e elevada. E, como toda a bondade ou perfeição das creaturas deriva da vontade de \textsc{Deus}, é claro que \textsc{DEUS} mais ama as creaturas mais perfeitas, porque lhes quiz um bem maior, uma perfeição mais elevada. --- Note-se que o amor se diz mais on menos intenso emquanto ao \emph{objecto}, isto é, emquanto ao maior ou menor bem que se quer ou confere; porque, da parte do seu \emph{principio}, que é de \textsc{Deus}, o amor é simplicissimo e não admitte o mais nem o menos, mas é sempre egual (\emph{Zb}., a. 8). --- 8.9) Se considerarmos a indole da \emph{natureza creada}, \textsc{DEUS} mais ama os anjos que os homens, mais os homens que as creaturas inferiores; porque a natureza dos anjos é superior a dos homens, e a dos homens é superior á das outras creaturas inferiores (1/., a. 4). --- 4.º) Se considerarmos o estado \emph{moral} das creaturas \textsc{DEus} mais ama os justos do que os peccadores; porque os justos, além da bondade \emph{p/ysica}, que lhes é commum com os peccadores, possuem a bondade \emph{moral}, de que os outros estão privados (\emph{Jb}. a. 4). --- 5.º) \emph{A}" pergunta\pgn{508}: se \textsc{DEUs} ame os meramente \emph{Possiveis}, --- não emquanto ao \emph{ser}, que elles têm em \textsc{Deus} (porque esse ser, identificando-se com o ser divino, é objecto de um amor \emph{necessario} da vontade divina), --- mas emquanto ao \emph{ser}, que os \emph{possiveis} teriam em si mesmos, se fossem creados por \textsc{Deus}, --- os philosophos e os theologos não dão a mesma resposta. A opinião mais provavel é --- que \textsc{DEus} ama as creaturas meramente \emph{possiveis} (que nunca terão existencia) com \emph{amor} ?\emph{nefficas de simples complacencia}, porque tambem ellas têm uma \emph{certa} entidade (embora \emph{não-actual}), constituida por notas especificas e individuaes, e, se \textsc{DEus} quizesse, poderiam existir em si mesmas, --- 9 que esse amor é \emph{necessario}, porque \textsc{Deus} não pode deixar de amar tudo o quetem uma relação necessaria com a sua essencia e omnipotencia, como é o ente meramente \emph{Possivel},}.

\pg{508} d) E' \emph{benevolo}. \textbf{Deus}, amando-nos desde a eternidade e

sem \textbf{algum merito} nosso, deu-nos a existencia, conserva-nos

e rege-nos com paterna e sabia providencia, e todas as coisas

\textbf{dispõe} para o nosso bem, para a nossa eterna felicidade \footnote{Quaes são os nossos deveres em vista dos innumerayeis beneficios, que o amor divino tem espalhado e espalha continuamente sobre nós? Todos esses deveres, como os beneficios de Dus, estão resumidos e concentrados na celebre formula do Apostolo S. João: « \emph{Nos ergo diligamus} Deum, \emph{quoniam} \textsc{DEUS} \emph{prior dilewit nos} » (I Ep., IV, 19).}.

78. \textbf{Termo immanente} do \textbf{amor divino}. --- O acto da vontade

divina, pelo qual \textsc{Deus} ama o proprio ser e todas as Coisas,

que participam da sua bondade infinita, deve ter o seu termo,

como todo e qualquer acto. E, como o acto da vontade é

formalmente \emph{immanente}, tambem \emph{immanente} deve ser o seu

termo. Este termo \emph{immanente} do amor divino costuma cha-

mar-se \emph{espirito}, --- O \emph{espirito}, em \textsc{Deus}, é unico, infinito, substancial, como unico, infinito, substancial é o acto da vontade divina, de que procede \footnote{\textsc{Deus} é purissimo espirito. Portanto o seu primeiro acto é o \emph{pensamento}. O pensamento de \textsc{Deus} não é, como o nosso, multiplice, vario, imperfeito, mas é um acto simplicissimo, unico, infinitamente perfeito. Fructo e termo d'este simplicissimo acto é 0 \emph{verbo}. --- Mas o espirito não só pensa, ama tambem. Dus, conhecendo-se \& si mesmo no seu verbo, e, vendo-se infinitamente perfeito, não pode deixar de se amar com um amor, proporcionado á sua amabilidade. O termo d'este amor é o \emph{espirito}. O \emph{espirito} é o termo \emph{necessario} da vontade divina; porque o acto d'esta faculdade não é \emph{livre}, mas é \emph{natural}, isto é, nasce da mesma faculdade considerada como \emph{natureza}. O \emph{espirito é unico}, como \emph{unico} é o \emph{verbo}. Uma vontade infinita não pode exhalar senão um unico espirito, como uma intelligencia infinita não pode gerar senão um unico verbo. Uma variedade de termos suppõe uma repetição de actos, e por isso uma limitação nas faculdades operativas.\par A Fé vai mais além, e ensina que, em \textsc{DEus}, o \emph{Espirito} é uma Pessoa, Subsistente na unica essencia divina, e procedente do Pae e do Filho, como de unico Principio. \textsc{Deus} pensa e ama com toda a sua substancia, e toda a substancia Elle transfunde nos termos do seu pensamento e do seu amor.\pgn{509}Como o Pae, entendendo-se a si mesmo, se transfunde todo no Verbo, assim Pae e Filho, amando-se, se transfundem no termo do seu amor, que é o \emph{Espirito Santo}. Assim, o Espirito Santo é o termo substancial e pessoal do mutuo amor do Pae e do Filho. Se a intelligencia do Pae é o principio do Filho, á vontade do Pae e do Filho é o principio do Espirito Santo. --- Procedendo da \emph{intelligencia} do Pae, o Verbo é verdadeiro Filho, porque, como dissemos, a intelligencia, quando entende, torna semelhante a si mesma o objecto entendido, --- semelhança indispensavel para que uma origem possa chamar-se \emph{geração}; --- procedendo da \emph{vontade} do Pae e do Filho, o Espirito Santo não procede por \emph{geração}, porque a vontade, longe de tornar semelhante a si mesma o objecto querido, torna-se semelhante a elle. --- Como o Verbo não dá origem a outro verbo, porque é Elle o Verbo, que procede; assim o Espirito Santo não dà origem a outro Espirito, porque é Elle o Espirito que procede. São tres, portanto, as Pessoas divinas; porque dois são os actos \emph{imnmanentes} de \textsc{Deus} --- o conhecimento e o amor, e dois são os termos d'estes actos. --- O Pae, que não procede, é a \emph{primeira} Pessoa; o Filho, que procede da intelligencia do Pae, é a \emph{segunda}; o Espirito Santo, que procede da vontade do Pae e do Filho, é a \emph{terceira} Pessoa, porque o acto da vontade é posterior, na origem, ao acto da intelligencia. --- A theologia \emph{dogmatica} apresenta noções mais minuciosas e profundas acerca d'este altissimo mysterio da SS. Trindade, á qual seja gloria, benção e amor por toda a eternidade!},

%%>>> art-43

\section*{Artigo III.\\ Vida divina, sua existencia e natureza.}
\addcontentsline{toc}{section}{Artigo III. --- Vida divina, sua existencia e natureza.}

\pg{509} 79. \textbf{Vida divina}. --- \emph{Vida divina} é o attributo, pelo qual \textsc{Deus} produz operações \emph{immanentes}, que têm o \emph{principio} e o \emph{termo} no mesmo sujeito. --- Diremos da sua \emph{existencia} e \emph{natureza} \footnote{A palavra \emph{vida}, embora formada pelo facto de se perceberem, nalguns entes, umas \emph{operações inmmanentes}, que tinham nelles o seu principio e o seu termo, comtudo emprega-se para denotar a \emph{substancia}, que tem \textbf{esse} poder de se mover a si mesma, ou de produzir as operações \emph{immanentes}. Por isso a \emph{vida} toma-se em acto \emph{primeiro} e em acto \emph{segundo}. Tomada em acto \emph{primeiro}, significa a propria \emph{substancia} do ser vivo; d'ahi o adagio: \emph{Vivere viventibus est esse}. Tomada em acto \emph{segundo}, significa a \emph{operação vital}, isto é, a operação propria dos seres vivos. Diz S. Thomaz: « Vita dicitur dupliciter: uno modo ipsum \emph{esse} viventis...; alio modo ipsa \emph{operatio} viventis, secundum quam principium vitae in actu reducitur » (\emph{Sum}. \emph{Th}. 1-9, q. 8, a. 2 ad 1) Por isso a vida é um attributo \emph{substancial}. (Cf. \emph{Sum}. \emph{Th}. p. I, q. 18; --- c. \emph{Gentes}, l. c. 97 e seg.).}.

80. \textbf{Deus} é \textbf{dotado} de vida. \textbar{}

a) Um ser é dotado \emph{de vida}, quando produz ou é capaz de produzir operações \emph{immanentes}. Ora \textsc{Deus}, como provámos, produz operações \emph{immanentes}, que são as da sua \emph{intelligencia} e da sua \emph{vontade}. Logo \textsc{Deus} é dotado de \emph{vida}.

\pg{510} b) O ser divino deve conter, na sua simplicissima unidade,

toda á plenitude da perfeição, e por isso todas as perfeições

\emph{absolutas}. Ora a \emph{vida} é uma perfeição \emph{absoluta}, porque no seu

conceito não encerra imperfeição alguma e não se oppõe a

perfeições mais elevadas. Logo \textsc{Deus} é dotado de \emph{vida} \footnote{E' certo que tudo o que se encontra no effeito deye encontrar-se na causa --- ou de um modo \emph{eminente}, --- ou de um modo \emph{proprio} e \emph{formal}. Por isso, em \textsc{Deus}, causa primeira de todas as creaturas, devem encontrar-se todas as perfeições creadas --- ou de um modo \emph{eminente}, --- ou de um modo \emph{proprio} e \emph{formal}. Ora a \emph{vida}, considerada no seu conceito \emph{essencial}, não encerra senão \emph{perfeição}; Porque \emph{viver} é poder produzir um acto, independentemente de um principio \emph{extrinseco}, e isto é perfeição. --- Ev verdade que a \emph{vida} se faz consistir tambem no \emph{movimento immanente}, e que em \textsc{Deus} não ha nem pode haver \emph{movimento}, porque Elle é Purissimo \emph{acto}, sem sombra de \emph{potencia}, e todo o \emph{movimento} importa uma composição de \emph{potencia} e de \emph{acto}. Mas, quando se diz que á \emph{vida} consiste num \emph{Movimento immanente}, a palavra \emph{movimento} deve tomar-se em sentido \emph{latissimo}, emquanto significa toda e quelquer \emph{operação immanente} no sujeito, quer essa operação seja realmente distincta do sujeito, quer não. Neste sentido, essa definição de \emph{vida} pode applicar-se tambem á vida de \textsc{Deus}, Se, porém, a palavra \emph{movimento} se toma em sentido \emph{rigoroso}, emquanto importa uma passagem da potencia para o acto, \emph{e} uma distineção \emph{real} outre o sujeito operante e a sua operação immanente, então essa definição de \emph{vida}, --- se convem á vida das \emph{creaturas}, que, nas operações, passam da potencia para o acto, e cujo acto é realmente distincto da substancia, --- não pode convir \& vida do \emph{Creador}, que é acto purissimo e simplicissimo, --- Mas, repetimos, o \emph{essencial} para o conceito e natureza da \emph{vida} é que a Operação seja \emph{immanente}, independente de um Principio \emph{extrinseco}. Que essa operação \emph{immanente} se exerça com ou sem uma passagem da Potencia para o acto, com ou sem composição de substancia e de acto, isto é \emph{accidental} Para o conceito e natureza da \emph{vida},}.

81. \textbf{Deus} é a \textbf{sua vida}.

a) A vida, nos entes vivos, é o proprio ser. Ora \textsc{Deus} é

o seu ser. Logo \textsc{Deus} é a sua vida.

b) Entender e querer é viver; porque os actos da intel-

ligeneia e da vontade são vitaes. Ora \textsc{Deus} é o seu entender

€ o seu querer. Logo é o seu viver e a sua vida \footnote{Alem d'isso, se \textsc{Deus} não fosse a sua vida, seria Vivo, não por essencia, mas por Participação, e por isso teria recebido de outro ente superior a vida, de que é dotado. Ora é absurdo dizer que Ser infinito tenha recebido de um ente superior alguma perfeição. Logo \textsc{DEus} é a sua vida. (C. \emph{Gentes}, 1. I, c. 98).}.

82. A \textbf{vida} de \textbf{Deus} é \textbf{perfeitissima}.

a) Consistindo a vida na operação \emph{mmanente}, é claro que

quanto mais \emph{immanente}, isto é, independente de um impulso

exterior, fór a operação, tanto mais perfeita será a vida. Ora

á operação de \textsc{Deus}, identifieando-se com a sua essencia, é \pg{511}  absolutamente independente de todo e qualquer impulso exterior. Logo a vida de \textsc{Deus} é \emph{perfeitissima}.

b) A vida de \textsc{Deus} é uma vida \emph{intellectual}, e por isso será tanto mais perfeita, quanto mais immunes do influxo de principios extranhos forem os actos da sua intelligencia e da sua vontade. Ora os actos da intelligencia e da vontade de \textsc{Deus} são inteiramente immunes de todo o influxo extranho; porque \textsc{Deus} é a propria Verdade e Bondade, e os seus actos, como tambem as faculdades, são uma e a mesma coisa com a essencia. Logo a vida de \textsc{Deus} é \emph{perfeitissima} \footnote{A vida \emph{intellectual} é propria das \emph{substancias espirituaes}, mas sô em \textsc{Deus} ella attinge o mais alto grau de \emph{immanencia}, e por isso de perfeição. Diz S. Thomaz: « Cum vivere dicantur aliqua secundum quod operantur ex seipsis, et non quasi ab aliis mota; quanto perfectius competit hoc alicui, tanto perfectius in eo invenitur vita » (\emph{Sum}. \emph{Th}., p. L, q. 18, a. 8). --- Consideremos a vida \emph{intellectual} do \emph{homem}. A nossa intelligencia move-se a si mesma, forma a especie intelligivel, que é o \emph{principio} do conhecimento, e produz o \emph{verbo}, que é a imagem do objecto entendido e o \emph{termo} do acto intellectual. Tem, pois, em si o \emph{principio} e o \emph{termo} da sua operação, e esta operação é propriamente \emph{imumanente}. Mas, se reflectirmos --- em que a \emph{especie intelligivel} é formada pela intelligencia humana depois de esta receber, por meios dos sentidos, um impulso da parte dos objectos externos, --- e em que o \emph{verbo mental}, embora unido com a intelligencia, é comtudo uma coisa accidental, distincto d'esta faculdade, devemos concluir que o nosso acto intellectual não é tão immanente, não é tão intimo, que não se possa conceber outro acto mais perfeito. --- Acto mais perfeito, infinitamente mais perfeito é o de \textsc{Deus}, porque elle, e só elle, é infinitamente immanente. \textsc{DEus}, absolutamente independente no ser, é absolutamente independente na operação. A. especie intelligivel, que, segundo o nosso modo de entender, move e determina a intelligencia divina ao acto do entendimento, não depende, nem proxima nem remotamente, do impulso de um agente externo, porque é a propria essencia divina; como tambem o \emph{verbo}, que é o termo do acto do entendimento, não é uma coisa accidental, distincta da intelligencia, mas é \& propria intelligencia, é a propria essencia divina. Mais immanente, mais intimo não pode ser o acto da intelligencia divina, e por isso não pode ser mais perfeito. Acto perfeitissimo, vida perfeitissima. --- O que dissemos com relação ao acto da intelligencia divina, deve applicar-se ao acto da vontade. A nossa vontade move-se a si mesma para um bem particular; mas, antes de se mover para um bem particular, ella recebe um impulso e é movida pelo ultimo fim, pela felicidade. Portanto o acto da nossa vontade é \emph{immanente}, mas não é perfeitamente immanente. Pelo contrario, á vontade de \textsc{Deus}, tendendo para a felicidade, não é movida por uma coisa extranha, porque Elle é a sua propria felicidade, o seu ultimo fim. Diz S. Thomaz: « Perfectior modus vivendi est eorum, quae habent intellectum. Haec enim perfectius movent seipsa... Sed quamvis intellectus noster ad aliqua se agat, tamen aliqua sunt ei praestituta a natura; sicut sunt prima principia, cirea quae non potest aliter se babere, et ultimus finis, quem non potest non velle.\pgn{512}Unde, lice\& quantum ad aliquid moveat se, tamen oportet quod quantum ad aliqua ab alio moveatur. Illud igitur, cujus sua natura est ipsum ejus intelligere, et cui id, quod naturaliter habet, non determinatur ab alio, hoc est quod obtinet summum gradum vitae. Tale autem est Deus. Unde in Deo maxime est vita. Unde Philosophus, ostenso quod Deus sit intelligens, concludit quod habeat vitam perfectissimam et sempiternam, quia intellectus ejus est perfectissimus et semper in actu » (\emph{Sum}, \emph{Th}., P- I, q. 18, a. 8).}.

\pg{512} 83. A \textbf{vida} de \textbf{Deus} é sempiterna.

a) Um ser deixa de viver, quando d'elle se separa a vida. Ora a vida não pode separar-se de \textsc{Deus}; porque, em \textsc{Deus}, a vida é o seu ser, e uma coisa não pode separar-se de si mesma, pois a separação consiste na divisão de uma coisa de outra. Logo a vida de \textsc{Deus} é sempiterna.

b) Um ser, cuja vida não é sempiterna, mas tem o seu começo e o seu termo, e está sujeita a successio, é mutavel, composto e dependente de uma causa superior. Ora \textsc{Deus} é immutavel, simplieissimo, causa suprema de todas as coisas. Logo a vida de \textsc{Deus} é sempiterna \footnote{Estes e outros argumentos são de S. Thomaz (\emph{C}. \emph{Gentes}, 1. I, c. 99).}.

84. \textbf{Todas as} coisas creadas são vida em \textbf{Deus}. --- O viver de \textsc{Deus} é o seu entender. Ora, em \textsc{Deus}, o entender, o objecto entendido e a faculdade intellectual são uma e a mesma. coisa. Portanto tudo o que existe em Dus como objecto entendido é o viver de \textsc{Deus}, é a sua vida. Ora todas as coisas creadas, sendo effeitos de \textsc{Deus}, existem nelle como objectos entendidos. Logo todas as coisas creadas são vida em \textsc{Deus} \footnote{O argumento, adduzido no texto, é tambem do Angelico Doutor (\emph{Sum}. \emph{Th}, p. I, q. 18, a. 4. --- Para que evitemos equivocos, S. Thomaz faz as seguintes advertencias. --- 1.9) As creaturas são ou existem em \textsc{DEUS} por dois modos: emquanto são conservadas pelo poder divino (assim dizemos que existe em nós o que esta em nosso poder) e emquanto s\&o conhecidas pela intelligencia divina (assim dizemos que existe em nós uma pedra, cuja especie está na nossa intelligencia) Emquanto são conservadas pelo poder divino, ellas existem em \textsc{DEUS} \emph{como existem nas suas proprias naturezas}. E' neste sentido que se devem tomar as palavras de S. Paulo, quando diz: \emph{In épso} (\emph{Deo}) \emph{vivimus}, \emph{movemur et sumus}; porque o nosso viver, o nosso ser e o nosso movimento são causados por \textsc{DEus}. Emquanto são conhecidas pela intelligencia divina, as creaturas são em \textsc{Deus} pelos exemplares archetypos, que as representam como são em si mesmas e que são a propria essencia divina. Por isso as coisas, emquanto existem em \textsc{Deus} por este segundo modo, são a essencia divina e a vida divina (45. ad L) --- 2.0) Uma coisa feita segundo um modelo deve conformar-se com este, emquanto este é um \emph{modelo}, \emph{e} não emquanto tem \emph{um certo modo de ser}. Porque o \emph{ser} do \emph{modelo} pode ser e é muitas vezes diverso do \emph{ser} da \emph{coisa modelada} ; assim o modelo de uma casa, emquanto existe na inteligencia do artifice, tem um ser immateriale intelligivel, ao passo que, se o considerarmos applicado\pgn{513}na propria casa, que existe fóra da intelligencia, o modelo tem um ser material e sensivel. Portanto as ideas exômplares divinas, se as considerarmos nas coisas creadas, ás quaes são applicadas, não vivem, ou podem não viver; mas, se as considerarmos na intelligencia divina, são a propria vida divina, porque tem o ser divino (45. ad 2).}.

{\small\noindent\textsc{Summario}: --- Creação, sua natureza, principio e fim. --- Conservação, sus divisão e necessidade. --- Concurso, sua natureza e extensão. --- Provi- dencia, sua existencia, extensão e propriedades.\par}

\pg{513} \textbf{OPERAÇÕES TRANSEUNTES DE DEUS}

%%>>> art-44

\section*{Artigo I.\\ Creação, sua natureza, principio e fim.}
\addcontentsline{toc}{section}{Artigo I. --- Creação, sua natureza, principio e fim.}

85. \textbf{Creação}. --- \emph{Creação} é o acto, pelo qual \textsc{Deus} tirou do \emph{nada} todas as coisas. --- Tirar do \emph{nada} o ente é produzil-o na sua \emph{totalidade}, isto é, não só emquanto á \emph{forma}, mas tambem emquanto á \emph{materia} \footnote{Tratamos, neste capitulo, dos principaes actos do amor \emph{efectivo} de Dzvus para com as creaturas, que são: a \emph{creacdo}, 8 conservação, o concurso e a \emph{providencia}. --- Neste artigo occupamo-nos da \emph{creacdo}. (\emph{Cf}. \emph{Sum}. \emph{Tos} p.d; q. 44, 45; \emph{c}. \emph{Gent}, 1. II, c. 15 e seg).\par A \emph{creagüo} é a \emph{produccdo de uma coisa do nada}. Para que, pois, a produecüo de uma coisa possa chamar-se \emph{creação}, são necessarios e sufficientes dois requisitos: 1.º) que \emph{nada preewista} da coisa, que se diz \emph{creada}, a não ser a não-repugnancia para existir, ou a absoluta possibilidade, ou a potencia objectiva, e por isso a creação differe de todas as outras especies de producção, como são a geração, a alteração, etc.; --- 2.9) que a coisa, que se diz creada, passe do não-ser para o ser, não emquanto ella é alguma coisa no termo a \emph{quo}, que é o \emph{não-ser}, mas emquanto, pelo influxo de uma causa proporcionada, recebe o ser, que antes não tinha.\par A \emph{creacüo} diz-se \emph{activa} ou \emph{passiva}, conforme se considera em \textsc{DEus} ou na creatura. --- Considerada em \textsc{Deus}, a \emph{creação} é o acto, pelo qual \textsc{Deus} tira do \emph{nada} as coisas. E' a definição apresentada no texto. Este acto de \textsc{Deus} é a sua propria essencia, emquanto se refere ás creaturas. Para a comprehensüo d'esta verdade, é necessario distinguir um duplice acto --- um \emph{immanente} e outro \emph{transeunte}. \emph{O immanente} é o que permanece no agente, e é uma perfeição do proprio agente, como entender, querer, etc.; o \emph{transeunte} é o que passa para o effeito externo, \emph{e} é uma perfeição do proprio effeito, que existe por elle, como aquecer, fabricar, etc. Este duplice acto compete a \textsc{Deus}. Compete-lhe o primeiro acto, porque \textsc{Deus} entende, quer, -ama, gosa; compete-Ihe o segundo, emquanto produz as coisas, conserva-as e\par 88 VOL. Il.\pgn{514}rege-as. E porque o primeiro acto é perfeição do agente, e o segundo é perfeição do effeito, e o agente é anterior ao effeito e é causa d'elle, é evidente que o primeiro acto é a razão do segundo e naturalmente o procede (\emph{C}. \emph{Gent}., l II, c. 1) --- Mas o \emph{acto} não se attribue a \textsc{Deus} e ás creaturas num sentido \emph{wnivoco}, mas só num sentido \emph{analogico}; e por isso do acto de \textsc{Deus} devemos afastar toda a imperfeição, que se encontra no acto das creaturas. O acto das creaturas sahe do agente e é recebido no paciente, porque é uma coisa \emph{accidental}, e por isso é formalmente \emph{transeunte}; ao passo que o acto divino não sahe de \textsc{Deus} para as creaturas, porque o acto de \textsc{DEUS} é a sua propria \emph{substancia}, e é simplesmente \emph{inmanente}. --- Todavia o acto divino não é tão \emph{imnmanente}, que se não possa chamar, de algum modo, \emph{transeunte}; porque convem com o acto \emph{éransewnte} das creaturas em que produz um effeito fóra do agente. Por isso, o acto divino chama-se \emph{formalmente} immanente e \emph{eminente} ou \emph{virtualmente} transeunte, emquanto, ficando no agente com que se identifica, produz o que o agente creado não pode produzir senão por um acto, que d'elle passa para o paciente. Emquanto \textsc{Deus} entende e quer, produz o effeito. --- Considerada nas creaturas, a \emph{creação} é uma \emph{relação das creaturas com o Creador}, \emph{como com o principio do proprio ser}. Esta relação é \emph{ndo-mutia}; porque, se as creaturas se referem \emph{realmente} a \textsc{DEus}, \textsc{Deus} só \emph{logicamente} refere-se ás creaturas; pois o \emph{fundamento} da relação, que é \emph{o ser tirado do nada}, só é \emph{real} nas creaturas. Comtudo, não devemos d'ahi concluir que \textsc{Deus} só \emph{logicamente} é Senhor das creaturas; porque, para \textsc{DEus} ser \emph{realmente} Senhor das creaturas, basta que estas \emph{realmente} derivem e dependam de \textsc{DEus}. --- Outras noções, relativas á \emph{creação}, foram expostas na \emph{Cosmologia}, tratando da \emph{origem do mundo}. Ahi consideramos a creação mais no aspecto \emph{passivo}, e aqui mais no aspecto \emph{activo}.}.

\pg{514} 86. A \textbf{creação} é o \textbf{acto} proprio de \textbf{Deus}.

a) A \emph{creacdo}, por ser um acto \emph{absoluto}, isto é, independente de toda a materia, deve ser exclusivamente proprio da causa; que possue uma existencia \emph{absoluta}, isto é, independente de todo o ser; porque o modo de operar é proporcionado ao modo de existir. Ora só \textsc{Deus} possue uma existencia \emph{absoluta}. Logo a \emph{creação} é o acto proprio de \textsc{Deus}.

b) Todo o effeito deve ser proporcionado e semelhante á forma do agente; porque todo o agente opera pela sua forma. Ora a forma, por que \textsc{Deus} opera, é a sua essencia, e esta consiste em que \textsc{Deus} é o proprio \emph{Ser} subsistente, contendo em si toda a plenitude e perfeição da entidade. Logo o effeito proporcionado e semelhante a \textsc{Deus} é o proprio \emph{ser} das coisas, Mas o ser das coisas, pelo facto de abranger \emph{todos} os elementos constitutivos, só pode ser produzido por \emph{creação}. Logo a \emph{creação} é o acto proprio de \textsc{Deus} \footnote{Desenvolvamus e esclareçamos este ponto. A Causa primeira, que communica \textbf{alentras} coisas todo o ser, que possuem, e por isso produz o ente emquanto é \textbf{enigi\&opera} pela \emph{intelligencia} com o concurso da \emph{vontade} e da \emph{potencia operativa}. isso, para entendermos, de algum modo, a acto creador\pgn{515} de \textsc{DEus}, devemos fallar da sua \emph{potencia}, ou \emph{energia operativa}, tendo, no capitulo antecedente, fallado da sua intelligencia e vontade.\par A \emph{potencia}, sendo uma perfeição \emph{absoluta}, deve encontrar-se em \textsc{DEUS} no seu conceito \emph{formal}, \emph{e} em toda a extensão da sua perfeição, excluindo todas as imperfeições e limitações, de que se acha cercada nas creaturas. Ora a \emph{potencia}, no seu conceito \emph{formal}, \emph{é um principio efficiente do effeito}. Este conceito, como se vê, importa simplesmente o \emph{acto}, e exclue toda a imperfeição da \emph{potencialidade}. --- Mas a \emph{potencia} attribue-se só \emph{analogicamente} a \textsc{Deus} e ás creaturas. Por isso, nas creaturas, a \emph{potencia} distingue-se \emph{realmente} tanto da \emph{essencia}, como do \emph{acto}, \emph{é principio efficiente} do \emph{acto} e do \emph{effeito}, e está sempre unida com a potencialidade, porque suppõe a passagem da potencia para o acto; ao passo que, em \textsc{Deus}, no qual a \emph{essencia}, a \emph{potencia} e O \emph{acto} são uma e a mesma coisa (pois \textsc{DEus} é o proprio Ser subsistente), a \emph{potencia operativa}, emquanto ás operações \emph{transeuntes}, não é principio efficiente do \emph{acto}, mas unicamente do \emph{e/feito}, e, d'este modo, encontra-se em \textsc{Deus} em toda a sua perfeição. (\emph{Cf}. \emph{Sum}. \emph{Th}., p. I, q. 25, a. 1 ad 8).\par Emquanto ás relações da \emph{potencia operativa}, ou \emph{activa}, de \textsc{DEUS} com a sua \emph{intelligencia} e \emph{vontade}, dizemos que, na producção e governo das creaturas, a intelligencia \emph{dirige}, á vontade \emph{manda} e a potencia \emph{executa}. Por quanto, \textsc{DEUS}, operando fóra da sua essencia não por necessidade da sua natureza, mas pela liberdade do seu juizo, não produz todos.os effeitos, que a sua potencia é capaz de produzir; mas, como convém a um agente independente, que livremente se determina a produzir um effeito preconcebido pela intelligencia, produz, sem alguma mudança de si mesmo, os effeitos externos, e os produz no modo, na ordem e no tempo, que o eterno juizo da sua vontade livremente determinou. Por isso, todos os effeitos attribuem-se, como á causa efficiente, ao imperio da vontade divina, porque a potencia só produz os effeitos, que a vontade imperou. Com razão diz o Psalmista (\emph{Ps}. 148): u \emph{Ipse dicit et facta sunt}; \emph{Ipse mandavit et creata sunt} » (Cf. S. Th., 2 \emph{Sent}. Dist. I, q. 1, a. 5 ad 11).\par A potencia divina, identificando-se com a \emph{essencia}, é simplesmente \emph{infinita} (\emph{Sum}. \emph{Th}., p. I, q. 25, a. 2). --- E pode chamar-se \emph{infinita} tambem com relação aos \emph{e/feitos}, --- não emquanto a potencia divina produza uma multidão actualmente infinita, ou produza um effeito infinitamente perfeito (o que repugna á razão de effeito) --- mas emquanto não pode exgottar-se na produec\&o dos effeitos, e a efficacia da sua acção nas creaturas é multiplice e illimitada. Diz S. Thomaz: « Numquam (Deus) tot effectus facit, quin plures facere possit, nec umquam ita intense operatur, quin intensius operari possit. Intensio autem in operatione divina non est attendenda secundum quod operatio est in operante, quia sic semper est infinita, cum operatio sit divina essentia; sed attendenda est secundum quod attingit effectum. Sic \textbf{enim} a Deo moventur quaedam efficacius, quaedam minus efficaciter » (\emph{De pot}.,\pgn{516}q. 1, a 2). --- D'onde se segue que \textsc{Deus} é \emph{omnipotente}, emquanto pode fazer todo o que é \emph{absolutamente possivel}. \emph{O absolutamente impossivel} é uma \emph{contradicção}, e \& contradiegüo não pode ser feita. (\emph{Sum}. \emph{Th}., p. I, q. 25, a. 8, --- c. \emph{Gent}., l. II, c. 25, --- \emph{De pot}., q. 1, a. 3).\par Uma potencia \emph{infinita} deve, na sua operação, ser independente de toda e qualquer coisa extranha, e por isso, para produzir o effeito, não exige que alguma coisa d'esse effeito preexista ou se presupponha. O acto, que não presuppõe nada da coisa, que por elle é produzida, diz-se \emph{creacdo}. --- A \emph{creação} é o acto proprio, especificamente proprio de \textsc{Deus}; porque só \textsc{Deus}, por ser a Causa universal e o primeiro Principio das coisas, pode, na sua operação, ser independente de toda e qualquer causa anterior, material ou efficiente. Por isso é absurdo dizer que a creatura possa, por si, tirar do \emph{nada} as coisas; porque é absurdo que a creatura possa ser o primeiro principio das coisas.\par Faz-se a questão se \textsc{DEUS} possa communiocar á creatura a virtude ou força de crear. Os escholasticos respondem que essa virtude ou força não pode ser communieada á creatura, e dão as seguintes razões. --- A creação exige uma potencia \emph{infinita}, que opere pelo \emph{ser}, como tal, e não pelo ser \emph{determinado} a uma certa especie; porque o termo da creacüo é o \emph{ser}, e a causa deve ser proporcionada ao effeito. Ora nenhuma creatura pode receber uma qualidade infinita, porque uma qualidade recebida deve accommodar-se á capacidade do sujeito, que a recebe. --- Além d'isso, o acto da creatura é uma coisa \emph{accidental}, que deve ser necessariamente recebido num sujeito para que possa -produzir o effeito, e por isso a creatura não pode fazer nada senão transformando a materia preexistente: o que é diametralmente opposto ao acto -\emph{creador}. Portanto a creação exige uma força, que se identifique com a essencia do agente, que seja acto puro e não seja distincto da propria operação. (\emph{De pot}., q. 8, \emph{a} 4). --- S. Thomaz accrescenta que a creatura não pode intervir, nem como causa \emph{Instrumental}, no acto \emph{creador} de \textsc{DEus}. De facto, o \emph{instrumento} serve para dispór o sujeito, de modo que este possa receber em si o acto do agente principal; ora, como a creação não presuppõe algum «sujeito, a creatura não pode intervir, nem como \emph{instrumento}, no acto creador de \textsc{Deus}. Além d'isso, o acto de uma causa, embora \emph{instrumental}, deve derivar da potencia d'ella; ora, sendo impossivel que da potencia de uma causa -\emph{finita} derive um acto de uma força infinita, qual é o acto necessario para a creação, é tambem impossivel que a creatura seja causa \emph{instrumental} do "ente \emph{creado} (Cf. \emph{De pot}., q. a. 9, 4; \emph{Sum}. \emph{Th}., p. I, q. 45, a. 5).\par Dissemos que, na producção das creaturas, a intelligencia de \textsc{Deus} \emph{di}- -\emph{rige}, a sua vontade \emph{manda} e a sua potencia \emph{executa}. Para entendermos, de algum modo, o acto da intelligencia divina nessa obra, lembremo-nos da nobilissima theoria do Angelico Doutor. A intelligencia divina, diz S. Thomaz, conhece perfeitamente a essencia divina --- acto purissimo e oceano infinito de toda a perfeição, --- e conhece-a não só emquanto é tal essencia, mas tambem emquanto pode ser imitada e participada por infinitos modos, embora todos infinitamente inadequados. Conhecendo assim a essencia divina, a intelligencia forma innumerayeis conceitos, ou archetypos, de coisas possiveis\pgn{517}, cada um dos quaes imita, mais ou menos, a mesma essencia, e, em conformidade com estes archetypos, \textsc{DEus} produz, numa certa differenca de tempo, tudo o que Elle quer produzir. E 8. Agostinho já tinha ensinado que todas as creaturas existiam, desde a eternidade, na arte e na sapiencia da intelligencia divina, pela qual foram depois creadas: « Sapientia Dei, per quam facta sunt omnia, secundum artem continet omnia, antequam fabricet omnia » (In \emph{Joan}, c. I, n. 17; \emph{cf}. \emph{Sum}. \emph{Th}., p. I, q. 15, a. 4).}.

\pg{515} 87. \textbf{Deus} creou \textbf{todas as} coisas.

a) \emph{Crear} é tirar do \emph{nada}. Ora \textsc{Deus} tirou do \emph{nada} todas as coisas; porque todas as coisas são \emph{contingentes} na \emph{totalidade} da sua substancia, isto é, tanto em relação á \emph{materia} como em relação á \emph{forma}, e todo o ente \emph{contingente} é produzido, em ultima analyse, pelo ente \emph{necessario}. Logo \textsc{Deus} \emph{creou} todas as coisas.

\pg{516} b) Se o acto da Causa primeira, na producção das coisas, devesse presuppór alguma materia n\&o-produzida, esta materia possuiria o ser \emph{por essencia}; porque o ente, que não possue o ser \emph{por participação}, possue-o \emph{por essencia}. Ora é impossivel

\emph{Z} \pg{517}  que a materia possua o ser \emph{por essencia}; porque o ente que possue o ser \emph{por essencia é acto purissimo}, emquanto a materia é \emph{mera potencia}. Logo o acto da Causa primeira, na producção das coisas, não presuppõe alguma materia n\&o-produzida, e por isso é um acto \emph{creador}. Logo \textsc{Deus} \emph{creou} todas as coisas \footnote{Outros argumentos encontram-se em S. Thomaz (C. \emph{Gentes}, 1. II, c. 15, --- \emph{De pot}, q. a. 5). --- Contra esta these fazem-se muitas objecções, algumas das quaes foram expostas e\textasciicircum{}resolvidas na \emph{Cosmologia}. Para a solução de outras, que se referem 4 Causa creadora, expomos os seguintes principios: --- 1º) \textsc{Deus} opera pelas ideas da sua intelligencia, e por isso, assim como não tem propriamente a idea da materia, separada da forma, mas \emph{so} a idea do composto subsistente, assim tambem não crea a materia sem a forma, mas crea a materia formada, isto é, o composto subsistente. Todavia, pode dizer-se que \textsc{Deus} tem uma idea da materia, emquanto a materia imita, de algum modo, a essencia divina (\emph{De pot}., q. 3, a. 1 ad 13). --- 2.9) O acto de \textsc{Deus}, na producção das coisas, passa do agente para o paciente --- emquanto, permanecendo na essencia divina, com a qual se identifica, produz fôra de \textsc{DEus} toda a substancia do paciente, --- mas não emquanto, derivando da essencia divina, é recebido no paciente; porque, como dissemos, o acto divino não se distingue da essencia, e esta produz o effeito sem produzir o acto. --- 8.º) A causa universal que é \textsc{Deus}, produz o ente \emph{emquanto ente} --- não no sentido que produz a razão abstracta do ente communissimo e analogico (que n\&o pode existir na realidade) --- mas no sentido que produz o ente na sua \emph{totalidade}, isto é, segundo \emph{todos} os elementos, de que o ente é constituido e que são a materia e a forma.\par Não nos demoramos neste assumpto, pois que foi sufficientemente tratado na \emph{Cosmologia}. --- Mas não podemos deixar de fazer uma breve critica do systema de A. \emph{Rosmini} acerca da \emph{creacdo}.\par Este escriptor, na sua \emph{T/eosophia}, explica a \emph{creacdo} por um triplice acto da intelligencia divina: \emph{abstracção}, \emph{imaginação} e \emph{synthese}. --- A Pessoa divina do \emph{Pae}, contemplando o \emph{Verbo}, \emph{abstrahe} d'Elle o \emph{ser inicial}, em que vé contido virtualmente todo o ser finito. E' a primeira \emph{creacdo}, que dà origem á primeira creatura, isto á, ao \emph{ser inicial}. Este ser não tem uma existencia \emph{real}, mas somente \emph{logica}, e reside primeiramente na intelligencia de \textsc{Deus} e depois na das creaturas. --- A esta \emph{abstracção} segue-se a \emph{imaginação}, pela qual DEus-PAE concebe varios \emph{termos} reaes e limitados, de que se compõe a realidade do universo. --- Finalmente, realiza-se a \emph{synthese}, pela qual \textsc{Deus} une os dois elementos, que são o \emph{ser inicial} e os diversos \emph{termos} d'elle, isto é, os entes \emph{reaes finitos} (\emph{Teos}., v. 1, ]. 2, sez. 4. c. 6, a. 7).\par O systema de \emph{Rosmini} é falso pelas seguintes razões:\par 1.9) Depois de attribuir só 4 Pessoa do \emph{Pae} a creação, que (como toda a acção de \textsc{Deus} \emph{ad extra}) \emph{é commun}. ás tres divinas Pessoas, \emph{Rosmini} reduz\pgn{518}a oreacüo a um puro \emph{idealismo}. Na verdade, as tres operações, de que ella seria constituida e que são a \emph{abstracção}, a. \emph{imaginação} e a \emph{synthese}, são productos da intelligencia e não dão origem a nenhuma realidade. Porquanto, o proprio \emph{Rosmini} ensina que a primeira creatura --- o \emph{ser inicial} --- não tem existencia \emph{real}, mas só \emph{logica}. Se não tem existencia \emph{real} o producto da \emph{abstracção}, nem a podem ter os productos da \emph{imaginação} e da \emph{synthese}. D'este modo, os entes creados não passam de entes \emph{ideaes}.\par 2.9) \emph{Rosmini} admitte como \emph{increado} o ser das creaturas. De facto, o ser das creaturas, na opinião d'esse escriptor, é constituido pelo \emph{ser inicial}, e este \emph{ser inicial} é unico, simplicissimo e eterno, porque é o proprio \emph{ser} do \emph{Verbo}, abstrahido somente da subsistencia. Ora esta opinião, além de conter o \emph{pantheismo}, destroe o conceito da \emph{creação}; porque a creação, como dissemos, tem por termo o proprio \emph{ser}, pois produz a \emph{totalidade} da substancia, Assim ensina S. Thomaz: « Oreatio terminatur ad \emph{esse}, tamquam ad proprium effectum » (\emph{De pot}., q. 8, a. 16 ad 21) E noutro logar: « Cum Deus sit \emph{ipsun esse} per suam essentiam, oportet quod \emph{esse creatum} sit proprius effectus ejus, sicut ignire est proprius effectus ignis » (\emph{Sum}. \emph{Th}., P. Lq. 8, a. 1). --- Goherente com esse principio, \emph{Rosmini} diz que --- o ser das creaturas não é \emph{produzido}, mas é \emph{posto}, --- o \emph{ser das creaturas} não se distingue do \emph{ser divino}, senão por causa dos \emph{termos}, --- só \textsc{Deus} pode crear, porque o \emph{ser inicial}, que dá á creatura a forma de ente, é \emph{alguma coisa} de \textsc{Deus}, e só \textsc{DEus} pode dispór do que lhe pertence (\emph{Teos}., v. 1, p. 969).}.

\pg{518} 88. \textbf{Deus} foi \textbf{livre na creação} do mundo. --- \textsc{Deus} foi \emph{livre} na creacüo do mundo, se a-creação não era exigida --- nem pela essencia divina, --- nem pela essencia do mundo. Ora a creação do mundo não era exigida --- nem pela essencia divina, porque esta é um oceano infinito de perfeições, e, para ser felieissima, não precisa de algum bem creado, --- nem pela essencia do mundo, porque o mundo é contingente, e todo o ente contingente é por si indifferente para a existencia e para a não-existencia. Logo \textsc{Deus} foi \emph{livre} na creação do mundo \footnote{Não insistimos neste ponto, que foi sufficientemente desenvolvido, tratando da \emph{vontade} de \textsc{Deus}. \textsc{Deus} é causa do mundo, --- não \emph{natural}, isto é, que não opera por necessidade da sua natureza, --- mas \emph{livre}, isto é, que opera pelo livre juizo da sua vontade. Por isso, não só não foi necessario que \textsc{DEus} creasse o mundo \emph{desde a eternidade}, mas nem foi necessario que o creasse \emph{no tempo}. Diz S. Thomaz: « Voluntas Dei est causa rerum. Sic ergo aliqua necesse est esse, sicut necesse est Deum velle illa: cum necessitas effeotus ex necessitate causae dependeat. Ostensum est autem quod, absolute loquendo, non est necesse Deum velle aliquid, nisi seipsum. Non est ergo necessarium Deum velle quod mundus fuerit semper. Sed eatenus mundus est, quatenus Deus vult illum esse, cum esse mundi ex voluntate Dei dependeat, sicut ex sua causa. Non est \emph{igitur} necessarium mundum Semper esse» (\emph{Sum}. \emph{Th}., p. 1, q. 46, a 1). --- Lembramos que, quando se diz que Deus creou o mundo \emph{no tempo}, --- não entendemos que o tempo tenha preexistido ao mundo, ou que o tempo seja a medida de toda a duração, mesmo da duração divina (o que é absurdo), --- mas entendemos que Deus, immovel na sua eternidade, que abrange e excede toda a duração\pgn{519}temporanea, coreou o mundo \emph{com o tempo}, emquanto, creado o mundo, começou o \emph{movimento}, de que o tempo é medida. O tempo, que se pensa ter existido antes do mundo, é \emph{Zmaginario}, não \emph{real} (\emph{Of}. \emph{De pot}., q. 9, a. 17).}.

\pg{519} 89. \textbf{Deus creou} o mundo para a sua gloria. --- \textsc{Deus}, sendo infinitamente intelligente, devia propor-se um fim na creação do mundo. Este fim não pode ser senão a communicação e por isso a manifestação das perfeicoes divinas, --- communicação e manifestação, de que resulta a \emph{gloria} de \textsc{Deus}. Logo \textsc{Deus} ereou o mundo para a sua \emph{gloria} \footnote{A \emph{gloria} é o \emph{conhecimento} claro da \emph{bondade} de um ser, a qual é digna de louvor. Por isso define-se: \emph{clara notitia cum laude}. D'onde se segue que a \emph{gloria} é maior ou menor, conforme maior ou menor é a \emph{bondade}, de que um ser é dotado, e conforme maior ou menor é o \emph{conhecimento}, que se tem d'essa bondade. A gloria das creaturas é imperfeita, porque a sua bondade é finita e participada. Só a gloria de \textsc{Deus} é infinitamente perfeita, porque só a sua bondade é essencial e infinita. --- Mas, como a bondade divina só por \textsc{DEus} pode ser conhecida adequadamente e por isso louvada como merece, só a gloria, que \textsc{Deus} recebe de si mesmo, é digna de \textsc{DEus}. A gloria, que a creatura pode e deve dar a \textsc{Deus}, é sempre infinitamente inferior ao merito; porque infinitamente deficiente é o conhecimento, que a creatura tem da bondade divina. --- Ainda que infinitamente inferiores ao merito, os nossos louvores dão ao bom \textsc{Deus} uma satisfação immensa e procuram á nossa alma uma felicidade eterna.\par Esta proposição --- que \textsc{Deus} ereou o mundo para a \emph{sua gloria} --- não é contraria á outra proposição --- que \textsc{DEus} ereou o mundo para communicar a sua bondade. Communicando a sua bondade, \textsc{Deus} entende a sua gloria. Como tambem, dando \emph{gloria} a \textsc{DEus}, as creaturas participam da felicidade divina. \textsc{DEus} não quiz separar do nosso bem a sua gloria; o nosso bem é a sua gloria, e a sua gloria é o nosso bem. Felizes as almas, que reconhecem os beneficios de \textsc{Deus} e Lhe tributam a homenagem da gratidão, fidelidade e da obediencia!}.

%%>>> art-45

\section*{Artigo II.\\ Conservação, sua divisão e necessidade.}
\addcontentsline{toc}{section}{Artigo II. --- Conservação, sua divisão e necessidade.}

90. \textbf{Conservação}. --- \emph{Conservação} é o acto da divina vontade, pelo qual as coisas, uma vez creadas, perseveram na existencia. --- Diz-se --- acto da divina \emph{vontade}; porque \textsc{Deus} não produz as operações \emph{transeuntes.por} necessidade da sua \emph{natureza}, mas pelo \emph{Zere} arbitrio da sua \emph{vontade} \footnote{S. Thomaz trata da \emph{conservação} das coisas creadas --- na \emph{Sum}. \emph{T}., p. L q. 104, a. 1, --- na \emph{C}. \emph{Gentes}, 1. III, c. 65, --- \emph{De pot}. q. 5, a 1.}.

91. \textbf{Conservação} positiva e negativa. --- A conservação \emph{positiva}, ou \emph{directa}, é a que suppõe e exige, da parte de \textsc{Deus}, um \pg{520}  continuo e real influxo, pelo qual as coisas, uma vez creadas, possam perseverar na existencia. A \emph{negativa}, ou \emph{indirecta}, é a que suppõe e exige que \textsc{Deus} --- não anniquile as coisas creadas, --- ou afaste d'ellas tudo o que as possa corromper \footnote{Todos os philosophos admittem que \textsc{DEUS} conserva as coisas creadas; mas alguns, como \emph{Crousat}, \emph{Bayle}, \emph{Galluppi}, sustentam que essa conservação é apenas \emph{negativa}. \textsc{DEus}, na opinião d'elles, conserva o mundo, \emph{emquanto o não destroe}. Mas os escholasticos ensinam que \textsc{Deus} conserva o mundo de um modo \emph{positivo}, e que por isso as coisas não poderiam subsistir, nem por um instante, se \textsc{DEus} lhes subtrahisse o seu influxo e o seu apoio.}.

92. \textbf{Deus conserva todas as coisas} de um modo positivo.

a) O ser de todas as coisas, mesmo depois de creadas, não deixou de ser \emph{contingente}, isto €, indifferente ou indeterminado tanto para a existencia como para a não-existencia; porque esse ser é recebido e limitado na essencia, e não é a propria essencia, mas distingue-se realmente d'ella. Se o ser de todas as coisas, mesmo depois de creadas, é, por si, indifferente ou indeterminado tanto para a existencia como para a não-existencia, não tem em si a razão sufficiente da sua continuação na existencia e por isso precisa do influxo \emph{positivo}, ou \emph{directo}, de uma causa, que constantemente o tire da indifferenca e o conserve. Esta causa, que assim conserva o ser de todas as coisas, é a mesma que o tirou do nada e que é Dzvus; porque o ser das coisas finitas depende sempre e exclusivamente d'Aquelle, que é o \emph{proprio Ser subsistente}. Logo \textsc{Deus} conserva todas as coisas de um modo \emph{positivo} \footnote{Este argumento costuma propor-se por differentes modos; mas funda-se sempre sobre o mesmo \emph{fermo medio} --- o caracter de \emph{contingencia} e de \emph{participação}, que revelam todas as coisas creadas. --- Todo o effeito, \emph{emquanto tal}, depende sempre da causa, que o produziu, de modo que o effeito dura tanto, quanto dura o influxo da causa; porque a razão da existencia do effeito não se encontra no proprio effeito, mas na sua causa. Assim a vida do nosso corpo, derivando da alma, depende da propria alma e dura tanto, quanto dura a união da alma com o corpo. Ora o \emph{ser} de todas as coisas é um \emph{effeito} exclusivo. de \textsc{Deus}. E' um \emph{e/feito}, porque não é o ser de \textsc{DEus}: e por isso ó um ser contingente, participado. Sendo \emph{contingente}, o \emph{ser creado} deve ser sempre determinado, para a existencia, pelo Ente \emph{necessario}; e, sendo \emph{participado}, não é a sua essencia e por isso não pode subsistir por si mesmo, mas deve subsistir por virtude do Ente, que é o proprio Ser subsistente. --- S. Thomaz compara a acção de \textsc{DEUS} nas coisas creadas á acção do sol no ar. Toda a creatura, diz o Santo Doutor, refere-se a \textsc{Deus}, como o ar se refere ao sol, que o illumina. Assim como o sol é luminoso por sua natureza, e o ar se torna luminoso pela participação da /\emph{uz} do sol, e não pela participação da \emph{natureza} do sol, assim tambem \textsc{Deus}, e só Elle, é o Ser pela sua essencia, porque á sua essencia é o seu ser, e toda a creatura\pgn{521} é ser por participação, porque a sua essencia não é o seu ser. Por isso, assim como o ar se torna luminoso ao apparecer do sol, e se cobre de trevas ao desapparecer d'elle, assim tambem a creatura existe emquanto \textsc{Deus} a conserva, e deixaria de existir e voltaria para o seu nada, se \textsc{Deus} lhe subtrahisse o influxo da sua potencia (\emph{Sum}. \emph{T}7., p. I, q. 104, a 1).}.

\pg{521} b) Se as coisas creadas se conservassem na existencia pela sua propria virtude, e não por um influxo positivo e directo do Creador, só poderiam ser \emph{a»niquiladas}, isto é, reduzidas ao nada, por um acto \emph{positivo} da divina vontade. Ora isto é simplesmente absurdo; porque um acto \emph{positivo} de \textsc{Deus} não pode ter por termo o \emph{nada}. Logo o \emph{anniquilagüo} das creaturas, que é \emph{possivel}, só pode conceber-se por um acto \emph{negativo} de \textsc{DEUS}, emquanto \textsc{Deus} lhes \emph{subtrahe} o influxo \emph{positivo} e \emph{directo}, que as conserva na existencia. Logo \textsc{Deus} conserva todas as coisas de um modo \emph{positivo} \footnote{Outros argumentos adduz S. Thomaz nos logares citados. --- Agora expomos e resolvemos as principaes objecções contra a nossa these.\par Obj. 13 --- \textsc{Deus} é mais poderoso do que todo o agente creado. Ora o agente creado communica ao effeito o poder de se conservar no ser, mesmo depois de cessado o acto productor. Logo, com maior razão, \textsc{DEus} pode conferir ás coisas creadas o poder de se conservarem no ser, sem um positivo e continuo influxo da sua infinita energia.\par Resposta. A razão, por que os effeitos do agente creado se conseryam no ser, mesmo depois de cessar o acto productor, não se deve attribuir \& perfeição e efficacia do mesmo agente, mas unicamente á potencia da Causa primeira, da qual o ser de todas as coisas é effeito proprio e especifico. Por quanto, o agente creado não pode operar senão transformando a materia, e por isso o seu acto presuppõe o sujeito constituido no seu ser. D'este modo, a conservação do effeito de um agente creado no proprio ser, longe de revelar a perfeição d'esse agente, revela, pelo contrario, a sua imperfeição, emquanto demonstra --- que o agente creado não pode, na sua operação, prescindir de um sujeito preexistente, --- e que é absolutamente necessaria a acção de \textsc{DEus}, que produza e conserve o ser das coisas. --- A difficuldade teria algum valor, se o agente creado produzisse por tal modo o effeito, que, 20 cessar do acto productor, o effeito permanecesse no proprio ser, independentemente de toda a causa: ora isto é falso. (\emph{De pot}., q. 5, a. 1 ad 4).\par Obj. 2.º --- O agente creado produz a \emph{forma}, e por isso o \emph{ser} das coisas; porque o \emph{ser} deriva da \emph{forma}. Ora, acabado o acto do agente creado, o effeito persevera no ser. Logo tambem, acabado o acto do Creador, que é a causa do \emph{ser} das coisas, estas perseveram no ser. s . Resposta. O agente creado, como dissemos, não opera senão transformando a materia, e a sua efficacia extende-se unicamente ás coisas, que, de algum modo, dependem da materia ou estão unidas com a materia. Por isso, se a efficacia do agente creado chega a produzir a \emph{forma substancial}, esta forma não a produz independentemente de todo o sujeito, como se ella fosse subsistente, mas tira-a da potencia da materia, dependentemente da materia, e por isso o ser, que se produz e que é proprio do composto (embora\pgn{522}derive da forma para o composto), não depende, pelo menos exclusivamente, do agente creado, e pode permanecer e permanece, embora acabe a acção d'esse agente. Mas \textsc{Deus} é a causa total e proporcionada do \emph{ser} das coisas, e por isso esse ser não poderia continuar a existir, se \textsc{Deus} não lhe continuasse o influxo do seu poder.\par Obj. 3.º --- Os substancias \emph{materiaes} são compostas de materia e de forma; ora a materia sustenta a forma, e a forma é o principio do ser. Do mesmo modo, as substancias /\emph{mmateriaes} são \emph{naturalmente indestructiveis}. Logo to- i das as substancias creadas continuam a existir pelos seus elementos \emph{intrinsecos}, e não pelo acção \emph{positiva} de \textsc{Deus}.\par Resposta. As substancias \emph{materiaes} são compostas de materia e de forma. Ninguem o contesta. Mas esses elementos, sendo \emph{contingentes} e \emph{participados}, não podem existir por si mesmos, por virtude propria; mas exigem necessariamente o auxilio de um agente extrinseco, de uma Causa superior, que z os conserve. --- Diga-se o mesmo com relação ás substancias \emph{immateriaes}. Estas são \emph{naturalmente} indestructiveis, mas na supposicio que \textsc{DEUS} conserve os principios essenciaes, de que são compostas. Diz S. Thomaz: « Quamdiu principia essentialia rerum sunt, tamdiu res conservantur in esse; sed et ipsa rerum principia esse desinerent, divina actione cessante » (\emph{De pot}., q. 5, a. 1 ad 3).\par Obj. 4.º --- Todas as coisas creadas, sendo obras de \textsc{Deus}, são perfeitas. Ora não seriam perfeitas, se, para a sua conservação, precisassem do influxo de um agente exterior; porque perfeito é o ente, a que nada falta. Logo as coisas creadas não precisam do influxo divino para a sua conservação.\par Resposta. Todas as coisas creadas são perfeitas --- emquanto têm tudo o que convem á sua natureza, --- mas não emquanto são dotadas de \emph{todas} as perfeições e independentemente de \textsc{Deus}. Ora uma das perfeições, que convem.á natureza das coisas creadas, ó que sejam conservadas por \textsc{Deus}, para que não voltem para o nada, d'onde foram tiradas. Escreve S. Agostinho « Creatoris potentia, et Omnipotentis virtus causa subsistendi est omni creaturae; quae virtus ab eis, quae creata sunt regendis, si aliquando cessaret, simul illorum cessaret species, omnisque natura concideret » (\emph{De Gen}. \emph{ad lit}, \textbar{}. IV, c. 12). Por isso S. Thomaz responde á diffieuldade com estas palavras: « Creaturae Dei sunt perfectae in sua natura et in suo ordine; sed inter alia, quae ad eorum perfectionem requiruntur, hoc etiam est, quod a Deo contineantur in esse» (\emph{De pol}, q. 5, a. 1 ad 1).}.

\pg{522} 93. \textbf{Deus conserva todas as coisas pela continuação} do \textbf{acto creador}. --- O ser das coisas finitas, em todos os instantes da sua duração, é o mesmo que era no primeiro instante, e por isso, para continuar na existencia, exige aquella mesma força, que foi necessaria para o seu começo. Ora, para o começo de todas as coisas, foi necessario o acto \emph{creador} de \textsc{Deus}. Logo, para que as coisas continuem na existencia, é necessaria a

4 continuação do mesmo acto. Logo \textsc{Deus} conserva todas as coisas pela continuação do acto \emph{creador} \footnote{O acto, pelo qual \textsc{Deus} conserva as coisas, não é diverso do acto, com que as creou, mas é a continuação do proprio acto creador. Esta verdade é certissima, --- não só porque em \textsc{Deus} não ha distinecüo ou diversidade de actos --- mas tambem porque, mesmo na esphera das causas finitas, quando\pgn{523}um effeito. 6 conservado pela mesma causa, que o produziu, o acto, que conserva o effeito, não é \emph{realmente} distincto do acto, que o produziu, mas é uma continuação do mesmo; assim o acto, com que o sol continua a illuminar o ar, não é \emph{realmente} distincto do acto, com que o illuminou no primeiro instante, mas é uma \emph{continuação} d'esse primeiro acto. Sendo uma continuação do acto \emph{creador}, o acto \emph{conservador} não dà ás creaturas um novo ser, distincto do ser, que receberam pelo acto \emph{creador}, mas conserva aquelle mesmo ser, que as creaturas receberam na creação. Diz S. Thomaz: 4 Conservatio rerum a Deo non est per aliquam novam actionem, sed per continuationem actionis, qua dat esse; quae quidem actio est sine motu et tempore; sicut etiam conservatio luminis in aére est per continuatum influxum solis » (\emph{Sum}. \emph{T}/., p. I, q. 104, a. 1 ad 4). --- Dissemos que o acto \emph{conservador} não é \emph{realmente} distincto do acto \emph{creador}; porque \emph{logicamente} um é distincto do outro, emquanto o acto \emph{creador} refere-se ao \emph{primeiro} instante da existencia das coisas, e o \emph{conservador} aos instantes \emph{posteriores}, --- o acto \emph{creador} não suppõe a existencia das coisas, o que suppõe o acto \emph{conservador}. --- Com razão, pois, a \emph{conservação} chama-se uma.creação \emph{continua}, emquanto a mesma virtude infinita, que tirou do \emph{nada} as coisas, é necessaria em cada momento para as conservar na existencia, --- e não emquanto Dus anniquile e crie continuamente as coisas, como pensavam \emph{Bagle}, \emph{Bertaud} e outros. Sob o influxo da acção divina, as coisas perseveram no ser, sem que sejam anniquiladas em cada instante, como, sob o influxo do sol, o ar continua a ser luminoso, sem que, em cada instante, se torne escuro.}.

\pg{523} 94. \textbf{Deus não anniquila alguma coisa creada}. --- As intenções de \textsc{Deus} relativamente ás coisas revelam-se na propria natureza das coisas. Ora a natureza das coisas mostra que \textsc{Deus} não anniquila nada do que Elle creou. Por quanto, as coisas creadas --- ou são \emph{emmalteriaes}, --- ou são \emph{materiaes}. Se são \emph{immateriaes}, são por si indestruetiveis, porque não ha nellas a potencia para a n\&o-existencia; se são \emph{materiaes}, são indestructiveis emquanto á \emph{materia}. porque esta é o sujeito de toda a geração e corrupção. Logo \textsc{Deus} não anniquila alguma coisa creada \footnote{\textsc{Deus} ereou as coisas, emquanto livremente lhes communicou o ser, e conserya-as, emquanto influe nellas constantemente o ser. Por isso, assim como podia não communicar-lhes o ser, antes que existissem, assim tambem pode subtrahir-lhes o seu influxo, depois de existirem. Subtrahindo ás coisas creadas o seu influxo, \textsc{Deus} reduziria as coisas ao \emph{nada}; porque, como dissemos, a reduecao das coisas ao \emph{nada} só pode ser effeito de uma acção \emph{negativa} de \textsc{Deus}, e não de uma acção \emph{positiva}, que tem por termo o \emph{ser} e não o \emph{nada} (\emph{Sum}. \emph{Th}., p. Y, q. 104, a. 3). --- Mas \textsc{DEus}, embora \emph{possa} reduzir ao nada as coisas creadas, comtudo é certo que não as reduzirá; e esta vontade divina revela-se nos proprios caracteres das coisas. Diz S. Thomaz: \& Quae facturus est Deus secundum ordinem naturaleni, rebus inditum, considerari possunt ex ipsis rerum naturis » (75. a. 4). A potencia e a bondade de \textsc{Deus} revelam-se mais na conservação do que na anniquilacüo das coisas. --- O homem deve corresponder a esta bondade de \textsc{Deus}. Se \textsc{Deus} consorya em cada instante a nossa vida, é justo que em cada instante a nossa vida seja um hymno de reconhecimento pelo nosso Bemfeitor.}.

%%>>> art-46

\section*{Artigo III.\\ Concurso, sua natureza e extensão.}
\addcontentsline{toc}{section}{Artigo III. --- Concurso, sua natureza e extensão.}

\pg{524} 95. \textbf{Concurso}. --- \emph{Concurso} é o acto, pelo qual \textsc{Deus}, como Causa primeira, move ou auxilia as forças \emph{physicas} das causas segundas, e assim influe em todos os actos \emph{mafturaes} d'essas causas. --- O \emph{concurso}, pois, distingue-se da \emph{conservação}. Esta refere-se á \emph{existencia} das creaturas; aquelle, ás \emph{operações} \footnote{Fallamos, como se vé, do concurso \emph{physico}, com que \textsc{DEus}, como Causa primeira e Auctor da natureza, move ou auxilia as forças das causas segundas, --- e não do concurso \emph{moral}, com que \textsc{DEus} pode determinar a causa segunda ao acto por meio de \emph{conselhos}, \emph{exhortações}, etc. Esse concurso \emph{physico} pertence á ordem \emph{natural}, refere-se ás operações \emph{naturaes} das causas segundas, e por isso distingue-se do auxilio da \emph{graça}, com que \textsc{DEus}, emquanto Áuetor da ordem \emph{sobrenatural}, influe nas operações da vida \emph{sobrenatural}. --- O concurso costuma chamar-se tambem \emph{influxo}, moção, acção de \textsc{Deus} nas obras das creaturas. Devemos, comtudo, advertir que o vocabulo « \emph{concurso} » aqui se toma em sentido \emph{lato}; porque, no seu sentido \emph{rigoroso}, \emph{concurso} significa a acção de dois agentes, cada um dos quaes influe no effeito, mas de uma maneira \emph{parcial} (e por isso esse effeito não depende \emph{inteiramente} nem de um nem de outro) e portanto não pode significar o \emph{influxo} de \textsc{DEus} nas operações das causas segunda, visto que o effeito d'essas operações attribue-se \emph{inteiramente} a \textsc{DEus} e \emph{inteiramente} á causa segunda, ainda que de um modo \emph{diverso}, como veremos.}.

96. \textbf{Especies} do \textbf{concurso}. --- O \emph{concurso} pode ser \emph{immediato} ou \emph{mediato}, \emph{previo} ou \emph{simultaneo}.

a) Concurso immediato ou mediato, --- Eº \emph{immediato}, se \textsc{Deus}, como Causa primeira e universal, opera juntamente com as creaturas, causas segundas e particulares, e com ellas realmente produz o mesmo effeito. --- E' \emph{mediato}, se \textsc{Deus} influe no effeito, produzido pelas creaturas, só emquanto \emph{conserva} nellas as faculdades operativas, de que as dotou.

b) Concurso previo ou simultaneo. --- E' \emph{previo}, se \textsc{Deus} move e appliea as creaturas á operação, de modo que o concurso de \textsc{Deus} \emph{prevém} ou \emph{precede} a determinação das creaturas. --- E' \emph{simultaneo}, se \textsc{Deus} não prevém, mas limita-se a \emph{acompanhar} a operação das creaturas \footnote{Sao as especies \emph{principaes}. Porquanto o \emph{concurso} pode ser tambem \emph{especial} ou \emph{geral}, \emph{intrinseca} ou \emph{ewtrinsecamente efficas}. --- E' \emph{especial}, se \textsc{Deus} para cada operação da creatura concede um determinado auxilio, --- é \emph{geral}, se \textsc{Deus} concede o mesmo auxilio para todas as operações. --- E' \emph{intrinsecamente efficas}, se o auxilio de \textsc{Deus} produz o seu effeito por si, independentemente\pgn{525} de toda a determinação da creatura, --- é \emph{extrinsecamente efficaz}, se o auxilio de \textsc{Deus} não produz o seu effeito por si, mas só dependentemente da previa determinação da creatura.}.

\pg{525} 97. \textbf{Deus} concorre \textbf{immediatamente em todas as operações das} creaturas.

a) Todas as operações das creaturas produzem um effeito; e este effeito contém o \emph{ser}, \emph{determinado e particularisado por algumas perfeições}. Ora, se as \emph{perfeições} que determinam e particularisam o \emph{ser}, sendo coisas \emph{particulares}, constituem o effeito \emph{proprio} das causas \emph{particulares}, que são as creaturas, o \emph{ser}, sendo uma coisa \emph{universal}, é o effeito \emph{proprio} da Causa \emph{univer}= \emph{sal}, que é \textsc{Deus}; porque a ordem dos effeitos deve corresponder á ordem das causas. Mas uma coisa, que está contida no effeito de uma causa \emph{particular}, não poderia ser produzida pela causa \emph{universal}, se a mesma causa \emph{particular} não operasse pela virtude ou energia da causa \emph{universal}, isto é, se a causa \emph{universal} não concorresse 2\emph{mumediatamente} na operação da causa \emph{particular}. Logo \textsc{Deus} concorre \emph{immediatamente} em todas as operações das creaturas \footnote{O argumento é de S. Thomaz. Eis as suas palavras: « Secundum ordinem causarum est ordo effectuum. Primum autem in omnibus effectibus est \emph{esse}; nam omnia alia sunt determinationes ipsius. Igitur \emph{esse} est proprius effectus primi agentis, et omnia alia agunt ipsum, inquantum agunt in virtute primi agentis; secunda autem agentia, quae sunt \emph{qwasé particulantia et determinantia actionem primi agentis}, agunt, sicut proprios effectus, alias perfectiones, quae determinant esse » (\emph{C}. \emph{Gentes}, l. III, c. 66). --- Este argumento funda-se naquella verdade: um effeito \emph{commum} e \emph{universal} deve 'derivar de uma causa \emph{commuwm} e \emph{universal}, como do seu principio adequado e proporcionado. Por quanto, se uma causa \emph{creada} e \emph{particular} podesse produzir um effeito \emph{commum} a varios entes e á propria causa, o effeito seria mais nobre do que a causa, e esta seria causa de si mesma. E' por isso que uma causa, subordinada a uma certa especie, não pode por si produzir a propria especie. Posto este fundamento solidissimo, conclue-se facilmente --- que o \emph{ser}, que está contido em todos os effeitos, provenientes das causas particulares, deve attribuir-se á virtude ou energia da causa \emph{universal}, --- “que as causas particulares operam pela virtude ou energia da propria causa "\emph{universal}, --- e que por isso \textsc{Deus} concorre \emph{immediatamente} em todas as operações das creaturas. E' o que diz S. Thomaz em poucas palayras: « Universaliores effectus oportet reducere in universaliores et priores causas. Inter omnes autem effectus universalissimurn est ipsum esse. Unde oportet "quod. sit proprius effectus primae et universalissimae causae, quae est Deus » (\emph{Sum}. \emph{Th}., p. 1, q. 45, a. 8).}.

b) Todas as creaturas, quando operam, passam da potencia para o acto e adquirem uma nova perfeição. Ora nenhum ente pode passar da potencia para o acto, nem adquirir uma

\pg{526} nova perfeição, sem o concurso de outro ente, que esteja em acto e possua essa perfeição. Logo as creaturas, quando operam, precisam do concurso de um ente, que esteja em acto, para que possa determinar a potencialidade d'ellas, e que possua immensas perfeições, para que as possa communicar aos que não as têm. Esse ente, em ultima analyse, é \textsc{Deus}; porque só Elle é superior a todas as creaturas, e é acto purissimo e infinitamente perfeito. Logo \textsc{Deus} concorre \emph{immediatamente} em todas as operações das creaturas \footnote{Estes e outros argumentos encontram-se em S. Thomaz (\emph{Sum}. \emph{TA}., P. I, q. 105, a. 5, --- \emph{c}. \emph{Gentes}, 1. IIl, c. 66, 67, --- \emph{De pot}., q. 8, a. 7). Ainda que apresentados sob varios aspectos, todos esses argumentos reduzem-se a este: \textsc{Deus} é o ente \emph{por essencia} (\emph{cuja} essencia é o seu ser) e por isso de \textsc{Deus} deve derivar todo o ente \emph{participado}. Ora toda a operação da creatura é uma realidade \emph{participada}, é um novo e limitado grau de entidade. Logo toda a operação da creatura deve derivar de \textsc{Deus}, como uma imperfeita participação da realidade infinita. Mas não poderia essa operação derivar de \textsc{Deus}, se \textsc{Deus} não concorresse \emph{immediatamente} na operação da creatura. Logo \textsc{Deus} concorre \emph{immediatamente} em toda a operação da creatura. --- E', pois, necessario admittir, contra \emph{Durando} e outros poucos, que \textsc{Deus}, além do concurso \emph{mediato}, emquanto concede ás creaturas e lhes conserva a faculdade ou energia de operar, presta tambem um concurso \emph{immediato}, emquanto applica essa faculdade ou energia á operação e emquanto toda a energia creada opera pela propria energia increada, --- De resto, o concurso \emph{immediato} é um corollario da conservação \emph{positiva}. Por quanto, se todo o ente finito e participado depende, na \emph{continuação} da sua existencia, de um concurso \emph{positivo} de \textsc{Deus}, tambem deve depender do mesmo concurso \emph{positivo} ou \emph{immediato} de Dus no seu \emph{começo}; porque, sendo a \emph{continuação} na existencia uma \emph{extensão do proprio começo}, a energia que \emph{conserva} o ser das coisas é a mesma que \emph{produziu}. Logo \textsc{Deus} produz todo o ente finito e participado. E como o effeito da operação das creaturas é um ente finito, deve concluir-se que \textsc{DEUS} concorre em todas as operações das creaturas.\par Demonstrada a \emph{existencia} do \emph{concurso immediato} de \textsc{DEus} em todas as operações das creaturas, devemos fazer algumas considerações, que esclarecem a \emph{natureza} e a \emph{condição} d'esse concurso.\par I. Nào ha nada na operação e no effeito da creatura, que não derive total e immediatamente de \textsc{Deus}. Na verdade, se houvesse nessa operação e nesse effeito alguma coisa, que não derivasse \emph{immediatamente} de \textsc{Deus}, o concurso de \textsc{Deus} seria \emph{mediato}. --- D'onde se Segue que a mesma .unica e indivisiyel operação da creatura procede \emph{inteiramente} e \emph{immediatamente} de \textsc{DEus}. Assim como uma e a mesma é a operação da causa instrumental e da causa principal, e não pode determinar-se uma parte da operação, que seja produzida unicamente pela causa instrumental sem o influxo da principal; assim tambem toda a operação das creaturas deriva da energia natural d'ellas e da energia da Causa primeira, e não ha uma parte só d'essa operação que possa altribuir-se exclusivamente ás causas segundas. \textsc{DEus}, pois, e á creatura são causas \emph{fofaes}, cada uma na sua ordem. Diz S. Thomaz: « Patet quod non sic idem effectus causae naturali e divinae virtuti\pgn{527}attribuitur, quasi \emph{partim} a Deo et \emph{partim} a naturali agente fiat, sed \emph{totus} ab \emph{utroque} secundum \emph{alium modwm}; sicut idem effectus totus attribuitur instrumento, et principali agenti etiam totus » (\emph{C}. \emph{Gentes}, ]. TIT, c. 10). E noutro logar: « In operatione, qua Deus operatur movendo naturam, non operatur natura; sed ipsa naturae operatio est etiam operatio virtutis divinae, sicut operatio instrumenti est per virtutem agentis principalis. Nec impeditur quin natura et Deus ad idem operentur, propter ordinem qui est inter Deum et naturam » (\emph{De pot}, q. 8, a. 7 ad 8). E noutro: « Una actio non procedit a duobus agentibus unius ordinis; sed nihil prohibet quin una ct eadem actio procedat a primo et secundo agente » (\emph{Sum}. \emph{Th}., p. T, q. 105, a. 5 ad 2).\par Il. A acção das creaturas é necessaria --- não pela insufficieneia da energia de \textsc{DEus}, --- mas pela immensidade da sua bondade. Por quanto, \textsc{Deus} teria podido produzir, sô por si, o effeito \emph{todo}, não so emquanto ao elemento uniyersal, que é o \emph{ser}, mas tambem emquanto aos elementos particulares, que \emph{determinam}. o ser. Todavia, sendo infinitamente bom, quiz communicar ás creaturas a propria semelhança, pela qual não só fossem bôas, mas podessem, operando, communicar aos outros a bondade. Mas, como a creatura não pode ter uma perfeita semelhança com o Creador, nem a energia participada egualar a energia infinita, segue-se --- que a creatura, operando, não pode attingir o effeito \emph{universal}, que é proprio da energia infinita e que é o \emph{ser}, --- e que por isso, na produecüo de um effeito, que contem esse ser, á mesma creatura precisa do concurso \emph{immediato} de \textsc{Deus}, isto é, não pode ella operar senão pela energia mesma de \textsc{Deus}. Diz S. Thomaz: « Neque est superfluum, si Deus per se ipsum potest omnes effectus naturales producere, quod per quasdam alias causas producantur. Non enim hoc est ex insufficientia divinae virtutis, sed ex immensitate bonitatis ipsius, per quam suam similitudinem rebus communicare voluit, non solum quantum ad hoc quod essent, sed etiam quantum ad hoc quod aliorum causae essent. His enim duobus modis creaturae communiter omnes divinam similitudinem consequuntur... Per hoc etiam decor ordinis in rebus creatis apparet » (C. \emph{Gentes}, l. ITI, c. 10)... 4 Creatura habet aliquam Dei similitudinem partieipando bonitatem ipsius, in quantum est et agit, non tamen ita quod per similitudinis perfectionem ad aequalitatem perveniat; et ideo sicut imperfectum indiget perfecto, ita virtus naturae in agendo indiget operatione divina » (\emph{De pot}., q. 8, a. 1 ad 10).\par III. Toda a causa creada, \emph{ainda que seja completa na sua ordem}, comtudo, absolutamente fallando, é por si incompleta e insufficiente para a producção de todo e qualquer effeito. Na verdade, se a causa segunda não pode operar sem o concurso \emph{immediato} da Causa primeira, é claro que ella é incompleta e insufficiente com relação ao effeito; visto que uma causa abso- Tutamente completa e adequada pode só por si produzir o effeito. Portanto a operação da creatura, que é tambem a operação de \textsc{Deus}, deve conceber-se como uma entidade, dotada de uma duplice relação de dependencia, --- uma com a propria causa creada, e outra com a Causa primeira e universal, porque de uma e de,outra depende de um modo indivisivel. --- Dissemos que a causa \emph{creada pode ser completa} na sua \emph{ordem}; porque, entre essas causas,\pgn{528}ha umas que podem operar sem o auxilio de outra causa creada e por isso são completas na sua ordem, ao passo outras, para operarem, precisam do auxilio de outra causa e por isso são incompletas e parciaes mesmo na ordem de causas segundas. --- Advirta-se que essa \emph{insuficiencia} das causas segundas não significa que as suas faculdades operativas sejam \emph{imperfeitas}. Essas faculdades são perfeitas, muito perfeitas, mas \emph{na sua ordem}, e a ordem d'ellas exige que os seus actos sejam exercidos dependentemente da Causa primeira. Assim como \textsc{Deus}, quando produz as coisas, não pode fazer com que as coisas existam independentemente d' Elle, assim tambem, quando concede ás creaturas as faculdades operativas, não pode fazer com que essas faculdades, embora perfeitas, se exercem independentemente d' Elle; porque \textsc{Deus} não pode fazer com que o ser participado seja ser por essencia. (\emph{De pot}., q. 8, a. Tad 5).\par IV. O effeito depende mais da Causa primeira do que da causa segunda. Na verdade, a causa segunda, sendo effeito da Causa primeira, recebe da Causa primeira a propria substancia. Recebendo da Causa primeira a substancia, a causa segunda recebe d'ella tambem a potencia ou energia operativa; porque a potencia é uma consequencia da substancia. Logo, se a causa segunda produz, pela sua potencia, um determinado effeito, esta causalidade deriva da Causa primeira. Portanto, a razão de causa, com relação a um certo effeito, compete mais á Causa primeira, do que á causa segunda, e por isso o effeito depende mais d'aquella do que d'esta. (S. Th., \emph{De causis}, l. L par. 6) --- Além d'isso, emquanto a energia da creatura não pode extender-se até á producção da materia ou do sujeito, porque a sua actividade presuppõe sempre a materia ou o sujeito, a energia de \textsc{Deus} extende-se atà á materia, creando-a, e assim o effeito procede de \textsc{Deus} segundo \emph{toda} a sua physica entidade. (\emph{De pot}., q. 3, a. 4. 7), --- D'ahi se segue que toda a operação da causa segunda se funda necessariamente na operação da Causa primeira, não só porque a creatura não pode operar sem o concurso divino, mas tambem porque toda a operação da creatura presuppõe o sujeito ou a materia, creada por \textsc{Deus}. (\emph{De pot}., q. 8, a. 4).\par V. As causas segundas são para a Catisa primeira, como o instrumento á para o agente principal. E' a opinião de S. Thomaz e de muitos outros. Por quanto, assim como o instrumento depende, na sua operação, do agente principal, nem pode operar senão pela força do mesmo agente, de modo que é uma causa essencialmente deficiente e imperfeita com relação ao effeito; assim tambem as causas segundas não podem operar sem o concurso da Causa primeira. --- Assim como una é a acção do agente principal e do instrumento e uno é o effeito; assim tambem una é acção e uno é o effeito da Causa primeira e da causa segunda. --- Assim como o instrumento não produz o effeito, se não fór applicado e movido pelo agente principal; assim tambem a causa segunda não produz o effeito, se não fôr applicada e movida pela Causa primeira. --- Assim como o agente principal se serve de instrumentos diversos, conforme a variedade dos actos e dos effeitos; assim tambem \textsc{Deus} se serve de muitas e diversas causas segundas para a producção de muitos e diversos effeitos. --- Todavia, as creaturas são \emph{instrumentos} --- não\pgn{529}num sentido \emph{rigoroso}, mas só num sentido \emph{analogico} ou \emph{metaphorico}. Por quanto, o instrumento não torna o effeito semelhante a si mesmo, mas ao agente principal, porque o instrumento não opera por virtude ou energia propria, mas por virtude do agente principal, que, operando pela propria virtude e forma, transfunde no effeito essa forma, e portanto a propria semelhança; ao passo que as causas segundas, embora operem sob o influxo de \textsc{Deus}, tornam o effeito semelhante a si mesmas, e não a \textsc{Deus}, porque operam pela propria forma. Além d'isso, o instrumento opera por virtude do agente principal, sem receber d'este um principio \emph{intrinseco} do seu movimento; ao passo que as causas segundas são dotadas de um principio \emph{intrinseco} do seu movimento, e algumas são por tal maneira movidas por \textsc{Deus}, que conservam o dominio dos proprios actos, a liberdade dos seus movimentos. (\emph{De verit}. q. 24, a. 1 ad 5).}.

\pg{527} 98. \textbf{Deus concorre immediatamente} em \textbf{todas as operações da vontade humana}.

a) A volição actual é uma realidade mais nobre e mais perfeita do que a propria faculdade volitiva; porque o acto é \pg{528}  mais perfeito do que a potencia. Ora nenhuma realidade, nenhuma perfeição pode encontrar-se nas creaturas, que não derive de \textsc{Deus}, fonte de toda a perfeição e causa universal de toda a realidade participada. Logo \textsc{Deus} é a causa --- não \pg{529}  só da vontade, que é o principio proximo do acto volitivo, --- mas tambem do acto volitivo, e por isso concorre \emph{Zntmediatamente} em todas as operações da vontade humana \footnote{E' necessario distinguir, sempre \emph{da parte do effeito}, uma acção divina que tem por termo a \emph{nossa vontade}, e uma outra acção que tem por termo o \emph{acto volitivo}. A primeira por tal modo produz e conserva a nossa faculdade volitiva, que á creatura não coopera com a acção divina, que lhe confere tal faculdade. A segunda, pelo contrario, produz por tal modo na nossa vontade o acto volitivo, que a mesma vontade coopera com a acção divina como verdadeira causa, embora \emph{particular}, que não pode operar, senão em virtude da causa \emph{universal}. (\emph{De pot}., q. 9, a. \emph{V}. ad 7).}.

b) \textsc{Deus} por tal modo confere a todos os agentes finitos, quer naturaes quer voluntarios, a faculdade ou potencia de operar, que nenhum d'elles pode servir-se d'essa faculdade, se não é determinado ou estimulado pela virtude ou energia do proprio Creador; porque todo o agente potencial e por participação, como é o agente finito, não pode operar senão pela virtude ou energia do agente, que é acto e opera pela sua essencia. Ora o ente, cuja virtude ou energia determina o acto de um agente, quer natural quer voluntario, é causa não só da faculdade operativa, mas tambem do proprio acto d'esse agente. Logo \textsc{Deus} não só é causa da nossa vontade, mas tambem do nosso acto volitivo, e por isso concorre \emph{immediatamente} em todas as operações da vontade humana \footnote{E' o argumento de S. Thomaz, na Sum. \emph{Contra Gentes}, l. III, c. 89, n. 4. --- Como nas operações de todas as creaturas, \textsc{DEUS} influe nas operações da vontade humana por differentes modos. Primeiramente, \textsc{Deus}, emquanto primeiro motor immovel e causa universal de todas as coisas, move a nossa vontade --- \emph{necessariamente} para o bem \emph{universal}, como para o objecto proporcionado d' essa faculdade e \emph{contingentemente} para os bens \emph{particulares}, que não tem um nexo necessario com o bem universal; porque a inclinação natural de uma faculdade para um termo não pode provir senão da Causa, de que deriva essa faculdade (\emph{Sum}. \emph{Th}., 1-2, q. 9, a. 6). Deus, portanto, move\par 54 Vor. II.\pgn{530}á nossa vontade, emquanto nos deu uma certa \emph{forma}, de que derivam certas faculdades e tendencias. --- Em segundo logar, \textsc{DEUS} move a nossa vontade emquanto ao \emph{objecto}. Com effeito, \emph{a} vontade é movida pelo \emph{bem}, e de um modo \emph{necessario} ou \emph{contingente}, conforme o bem é \emph{universal} ou \emph{particular}. Ora \textsc{Deus} é o bem \emph{universal}, e todo o bem \emph{particular} não pode mover a vontade senão emquanto é uma participação da bondade divina, a qual a o ultimo fim do homem, como é o seu primeiro principio. \textsc{Deus}, portanto, move a vontade humana emquanto esta tende para o objecto, que é o \emph{bem}. (\emph{De malo}, q. 6, a. un, --- \emph{Sum}. \emph{Th}., p. I, q. 105, a. 4. --- Finalmente, \textsc{Deus} move a nossa vontade, emquanto influe para que a vontade \emph{queira actualmente}, e assim é \emph{principio efficiente} do \emph{acto volitico}. E, como \textsc{Deus} não corrompe mas conserya a natureza creada por Elle, segue-se que a acção divina move \emph{diversamante} para a operação os agentes \emph{naturaes} e os agentes \emph{livres}. Move os agentes \emph{naturaes} emquanto lhes communica a energia necessaria para actuarem nos outros seres; move os agentes \emph{livres} emquanto lhes communiea a energia necessaria para se detorminarem ao movimento e para se moverem de modo, que sejam senhores dos seus actos. Por isso, o influxo divino, embora, considerado no seu \emph{principio}, seja um e o mesmo (porque a acção de \textsc{DEus} é a sua essencia), --- comtudo, considerado no seu \emph{termo}, que é a creatura, é multiplice e diverso, emquanto \textsc{Deus} influe nas causas segundas por um modo proporcionado á natureza d'ellas. Diz S. Thomaz: \& Ad providentiam divinam non pertinet naturam rerum corrumpere, sed servare. Unde omnia movet \emph{secundum eorum conditionem}; ita quod ex \emph{causis necessariis} per motionem divinam sequantur \emph{effectus em necessitate}, ex \emph{causis} autem \emph{contingentibus} sequantur \emph{effectus contingentes}. Quia igitur voluntas est activum principium \emph{non determinatum ad unum}, sed indifferenter se habens ad multa; sio Deus ipsam movet, quod non ex necessitate ad unum determinat, sed remanet motus ejus \emph{contingens} et \emph{non necessarius}, nisi in his ad quae \emph{naturaliter} movetur » (\emph{Sum}. \emph{Th}., 1-9, q. 10, a. 4. --- E claro que \textsc{Deus} concorre \emph{immediatamente} nas operações da nossa vontade emquanto é o principio \emph{efficiente} do acto volitivo.\par Para a intelligencia d'esta verdade e para a solução das objecções, devemos fazer as advertencias seguintes:\par l. A nossa vontade não pode ser coagida; já o provámos. O acto elicito da vontade escapa a todo o movimento \emph{violento}. O movimento de um ente é \emph{violento}, quando esse ente é determinado ao movimento por um principio \emph{extrinseco}, que não só não communica ao ente nenhuma inclinação\pgn{531}\emph{natural}, mas contraria essa inclinação. \textsc{DEus}, porém, sendo o Auctor de todas as naturezas, confere a todas a inclinação conveniente e movoe-as por tal modo, que o movimento está em perfeita harmonia com as inclinações d'ellas. Por isso o movimento, por que \textsc{Deus} actua nos entes creados, não a violento, mas é \emph{natural}, isto é, proporcionado á natureza de cada um. Diz S. Thomaz: « Illud quod movetur ab altero dicitur cogi, si moveatur contra inclinationem propriam; sed si moveatur ab alio, \emph{quod}. \emph{ipsi dat propriam inclinationem}, non dicitur cogi... Sie igitur Deus, movendo voluntatem, non cogit ipsam; quia dat ei ejus propriam inclinationem » (\emph{Sum}. \emph{Th}., p. I, q. 105, a. 4 ad 1). --- Advertimos que o movimento, que Deus imprime na nossa vontade, é e diz-se \emph{natural}, --- não emquanto a vontade seja por isso obrigada a operar \emph{necessariamente}, --- mas emquanto esse movimento é proporcionado á \emph{natureza} da vontade, a quel exige que \textsc{DEus} á mova e a mova de modo que lhe conserve a liberdade dos seus actos.\par II. O movimento de uni ente vivo deriva de um principio intrinseco, e, se esse ente é dotado de vontade, é indeterminado, emquanto aos actos voluntarios, tanto para uma coisa como para outra. Mas, se o ente vivo é causa \emph{segunda}, isto é, se é vivo por participação e não por essencia, então não só não exclue, mas exige a acção de um extrinseco principio superior, que é Deus; e como recebeu de \textsc{Deus}, --- Causa primeira e universal, vida por essencia e acto purissimo, --- o ser e a potencia operativa, assim tambem deve receber \emph{actualmente} d'Elle a energia necessaria para realizar a sua operação. Portanto o facto de um ser vivo se mover a si mesmo e, se é livre, de se determinar a si mesmo para uma coisa ou para outra, não exclue o influxo da Causa primeira, que é o principio universal de toda a perfeição e entidade participada, e que faz com que o ser vivo e livre se mova a si mesmo e se mova livremente. Diz S. Thomaz: « Cum aliquid movet seipsum, non excluditur, quin ab alio moveatur, « \emph{quo habet hoc ipsum}, \emph{quod}. \emph{seipsum Movel}; et sic non repugnat libertati, quod Deus est causa actus liberi arbitrii » (\emph{De malo}, q. 3, a. 2 ad 4). \textbar{}\par IIL Quando um ente é movido por outro e não se move de modo algum por sei mesmo, então aquelle movimento não é imputavel ao mesmo ente; porque, para um acto ser imputavel, é preciso que o agente tenha o poder de se determinar a si mesmo, isto é, é preciso que seja livre. Pelo contrario, se o ente, que é movido por \textsc{DEus}, é livre, então o seu acto é imputavel. --- Aqui advertimos que a Causa primeira e universal communica ás causas segundas e particulares, dotadas de liberdade, toda a perfeição, que se encontra no acto, e não a deformidade ou imperfeição; de modo que toda a perfeição do acto deve attribuir-se á Causa primeira como ao agente primario, que influe no effeito mais efficazmente do que o agente secundario, e toda a imperfeição ou todo o defeito deve attribuire á causa segunda, que pode faltar e desviar-se do recto caminho. Diz S. Thomaz: « Quia causa\pgn{532}prima magis influit in effectum quam secunda, ideo quidquid perfectionis est in effectu, \emph{principaliter} reducitur ad primam causam; quod autem: est de defectu, reducendum est in causam secundam, quae non iía efficaciter operatur, sicut causa prima » (\emph{De poL}, q. 9, a. 7 ad 15).\par 1V. A Causa primeira e universal, assim como produziu todas as coisas pela intelligencia e pela vontade, assim tambem pela intelligencia e pela vontade as governa e lhes concede tudo o que lhes é necessario para operarem actualmente conforme a natureza e a condição de cada uma. Por isso \textsc{Deus} não só não impõe a necessidade á nossa vontade, que é um principio activo, indifferente para muitas coisas, mas actua por tal modo nesta faculdade, que ella possa determinar-se mais para uma do que para outra coisa. A moção divina accommoda-se \& indole do agente creado. Diz S. Thomaz: « Voluntas divina non solum se extendit ut aliquid fiat per rem quam movet, sed ut etiam eo modo fiat, quo congruit naturae ipsius. Et ideo magis repugnaret divinae motioni, si voluntas ex necessitate moveretur, quod suae naturae non competit, quam si moveretur libere, prout competit suae naturae » (\emph{Sum}. \emph{T}/., 1-2, q. 10, a. 4 ad 1).\par V. A vontade humana tem o dominio dos seus actos. Esse dominio --- se exclue todo o influxo de uma causa exterior, que determine necessariamente a vontade a praticar um certo acto, --- não exclue o influxo da Causa primeira, que não só não determina necessariamente o acto, mas faz com que o proprio acto seja livre. Diz S. Thomaz: « Voluntas dicitur habere dominium sui actus non per exclusionem causae primae, sed quia causa prima non ita agit in voluntate, u\& eam de necessitate ad unum determinet, sicut determinat naturam; et ideo determinatio actus relinquitur in potestate rationis et voluntatis » (\emph{De pot}., q. 8, a. 7 ad 18; --- Cf. c. \emph{Gentes}; 1, 68).}.

\pg{530} 99. O \textbf{concurso de Deus} é \textbf{tambem previo ás operações das creaturas}.

a) As creaturas, como dissemos, não podem por si mesmas passar da potencia para o acto, mas devem ser reduzidas ao acto por \textsc{Deus}, emquanto \textsc{Deus} influe na potencialidade das causas segundas e as tira da indifferença. Ora este influxo não pode deixar de ser previo ou anterior, na ordem da natureza e da causalidade, ás operações das creaturas; porque o acto do motor deve ser previo ao acto do movel. Logo o concurso de \textsc{Deus} é tambem \emph{previo} ás operações das creaturas.

\pg{531} b) Se o concurso de Dus nas operações das creaturas é exclusivamente \emph{simultaneo}, \emph{e} não é tambem \emph{previo}, a \emph{determinação} d'essas operações é independente de Dzvus, e por isso

4 não é \textsc{Deus} quem determina o acto da creatura, mas, pelo contrario, é a creatura que determina o concurso de Deus. Ora estas consequencias são falsas; porque não pode haver, no \pg{532}  mundo, a mais simples entidade, que escape á universal causalidade de \textsc{Deus}, e porque a creatura, não podendo produzir acto algum sem o concurso de \textsc{Deus}, não pode determinar esse concurso. Logo o concurso de \textsc{Deus} é tambem \emph{previo} ás operações das creaturas \footnote{O concurso de \textsc{Deus} nas operações das creaturas pode ser, como dissemos, \emph{previo} ou \emph{simuwltaneo}. --- E' \emph{previo}, quando precede, pelo menos na ordem da natureza e da causalidade, o operação do agente, ao qual se presta, e por isso é recebido na potencialidade, ou faculdade, do proprio agente, e a determina para o acto. Tal é concurso, ou influxo, do esculptor, quando appliea o cinzel e o move para reproduzir no marmore uma determinada imagem. --- E' \emph{simultaneo}, quando não precede a operação do agente nem é recebido na potencialidade d'este, mas é apenas collateral, parallelo e extrinseco á propria operação do agente, e é recebido no \emph{effeito}, que, d'este modo, se attribue \emph{parcialmente} a um e outro agente. Tal é o concurso, que um cavallo dá a outro na tracção de um carro; porque um cayallo não influe na causalidade do outro, e o movimento do carro não se attribue \emph{totalmente}, mas só \emph{parcialmente} \& um e ao outro. --- Advertimos que o concurso \emph{previo} não exclue o \emph{simultaneo}; este pode ser uma \emph{continuação} d'aquelle. Assim, o esculptor applica o cinzel ao marmore, continua a propria acção previa e primitiva com a acção do instrumento; e, emquanto o cinzel produz com a sua acção o effeito proprio, que é \emph{cortar}, o esculptor tambem produz, por\par n\pgn{533}meio da acção do cinzel, o effeito proprio, que é \emph{reproduzir a imagem}; e, d'este modo, o esculptor e o cinzel são causa do mesmo effeito, embora numa ordem differente.\par Dada esta explicação, perguntamos: Deverá admittir-se o concurso \emph{previo} de \textsc{DEus} nas operações das creaturas, ou bastará o concurso exclusivamente \emph{simultaneo}? --- A esta pergunta têm respondido duas celebres escholas: a dos \emph{thomistas} e a dos \emph{molinistas}, sustentando aquella o concurso \emph{previo}, e esta o exclusivamente \emph{simultaneo}. --- Não temos a presumpção de dar uma solução definitiva \& esta controversia, que tem apaixonado tantos sabios, sobretudo nos seculos XVI e XVII. Limitamo-nos a expor, com muita serenidade e imparcialidade, as theorias d'essas duas escholas, citando os argumentos, em que se firmam, e as objecções, com que os adversarios se combatem.\par \emph{A}) \emph{Theoria do concurso simultaneo}. --- Foi um Jesuita hespanbol, \emph{Luiz Molina}, quem propoz esta theoria. Diz assim: « A acção de \textsc{DEus} e a acção da creatura podem comparar-se á acção de dois homens, que puxam uma barca. O movimento todo deriva de um e de outro, embora nenhum dos dois seja a causa total d'elle » (\emph{Concord}., q. 14, a. 13, dis. 26). Todavia, a acção de \textsc{Deus} e a acção da creatura não são duas acções distinctas. A acção é uma só, produzida contemporaneamente por duas causas \emph{parciaes}, que, unidas, formam uma causa \emph{completa}.\par Como se vê, nesta theoria, \textsc{DEUS} não opera \emph{na} creatura, mas com a creatura; não é \emph{motor}, \emph{é cooperador}; não \emph{dirige} o acto da creatura, limita-se a \emph{acompanhal-o}; não é \& causa \emph{primeira} do acto, mas \emph{concorre} para elle de um modo \emph{parallelo}; não dá ás causas segundas o impulso ao acto, mas dá, em união com ellas, o ser ao effeito; não é o senhor absoluto dos agentes creados, é o servo sempre prompto a prestar-lhes auxilio. Eis o fraco da theoria: dizem os adversarios. --- Por quanto, se \textsc{Deus} e a creatura não são agentes \emph{subordinados}, mas apenas \emph{unidos} ou \emph{associados}, sem que um dependa de outro, a acção da creatura escapa á causalidade divina e não se poderá dizer que toda a entidade, a qual se encontra na creatura, deriva de \textsc{Deus}. --- Ha mais: A acção de \textsc{Deus} subordina-se, nesta theoria, á acção da creatura. Com effeito, o concurso de \textsc{Deus}, que se suppõe indifferente para este ou aquelle acto, recebe da creatura a propria determinação, sem a qual \textsc{Deus} não poderia operar. Diz \emph{Molina}: « Para nós o concurso divino não determina a vontade ao consenso. Pelo contrario, é o influxo particular do livre arbitrio, que determina o concurso divino ao acto, emquanto a vontade se resolve a querer ou a não querer, a querer isto e não aquillo » (\emph{Concord}., q. 28, a. 4 e 5, dis. 1). \textsc{Deus}, pois, para operar, deve esperar que o homem o determine para este ou aquelle acto. Mas, como poderá a creatura determinar o concurso divino, se nada pode ella operar sem esse concurso? D'este modo, a liberdade humana, que se pretende salvar, esta reduzida á impotencia.\par Em vista d'estas difficuldades, alguns sequazes de \emph{Molina} abandonaram a theoria do concurso \emph{simultaneo}, embora lhe conservassem o nome, e apresentaram uma theoria, que pode chamar-se --- a theoria da \emph{premoção indeterminada}. A! comparação de dois homens, que puxam uma barca, substituiram a de um navio impellido pelo vento para um porto indeterminado.\pgn{534}O seu raciocinio é o seguinte. O vento impelle o navio, por. ex. para o levante, mas não determina o porto. Este é determinado pelo piloto, que governa e dirige o navio. Ora o vento é a imagem da moção divina, que prevem a acção das creaturas e as move para um termo geral; o piloto é a vontade livre, que determina este ou aquelle acto. Fundam esta theoria na auctoridade de S. Thomaz, quando diz --- que \textsc{DEus}, como motor universal, move a vontade do homem para o objecto universal da vontade, que ó o bem, --- e que o homem, pela razão, se determina para este ou aquelle bem particular, verdadeiro ou apparente (\emph{Sum}. \emph{Th}., 1-2, q. 9, a. 6 ad 8; \emph{De pot}. q. 1, a. 4 ad 8).\par A theoria da premoção \emph{indeterminada} é superior \& do \emph{concurso simultaneo}. Admittindo uma premoção da nossa vontade pela acção divina, ella explica a passagem d'aquella faculdade da potencia para o acto. Todavia, nem esta theoria reconhece uma \emph{total} dependencia do acto da creatura com relação ao acto de \textsc{Deus}. O vento explica o movimento do navio, mas não a sua direcção: esta depende exclusivamente do piloto. Do mesmo modo, se o movimento da vontade humana para o bem em geral depende da moção divina, a \emph{direcção} d'esse movimento para este ou aquelle bem em particular depende da propria vontade. Portanto admitte-se uma realidade positiva, uma verdadeira perfeição (tal é essa direcção), que escapa á causalidade divina, que não deriva da fonte de todo o bem e de toda a perfeição. --- Alem d'isso, tambem esta theoria admitte que o concurso divino, indeterminado por si mesmo, é determinado e especificado pela activa influencia do livre arbitrio humano. Mas como poderá a creatura determinar e especificar com o seu acto o concurso divino, se ella precisa do concurso divino para praticar todo e qualquer acto? --- Por estas e outras razões, muitos philosophos e theologos acceitam a theoria dos \emph{tAomistas}.\par \emph{B}) \emph{Theoria do concurso previo}. --- Esta theoria é mais synthelica e grandiosa, \textsc{Deus} é o motor primeiro o universal, e nada pode escapar á sua acção e causalidade soberana. Diz S. Thomaz: « Quantumceumque natura corporalis vel spiritualis ponatur perfecta, non potest in suum actum procedere, nisi moveatur a Deo » (\emph{Sum}. \emph{T}7., 1-9, q. 109, a. 1). --- « Deus comparatur ad omnia sicut movens ad motum, et activum ad id quod est in potentia » (C. \emph{Gent}., III, 100). Deus não move, como movem as causas segundas, por um modo exterior e violento; a sua moção realiza-se no intimo dos seres, e adapta-se á natureza e ás disposições de cada um. Diz o mesmo Angelico: « Deus cuilibet agenti adest interius quasi in ipso agens, dum ipsum ad agendum moyet » (\emph{Comp}. \emph{T}/., c. 120). Por isso a moção divina, longe de ser determinada e especificada pelo concurso activo das creaturas, determina-so e especifica-se por si mesma, conforme a natureza e as disposições do ser, em que é recebida, fazendo com que o ser necessario produza um acto necessario e o ser livre um acto livre. Na efficacia universal e transcendente da causalidade divina encontra-se a razão sufficiente de todos os actos das creaturas, mesmo dos actos livres. Assim como, emquanto é o Ser primeiro, \textsc{Deus} influe continuamente em todas as coisas o \emph{ser}, que é a ultima actualidade na linha do ser, e por este influxo physico e preveniente confere a\pgn{535}existencia, assim tambem, emquanto é a Causa primeira, Elle influe em todos os agentes o \emph{acto}, que é a ultima actualidade na linha da acção, e applica activamente as faculdades operativas aos seus termos connaturaes. Por isso \textsc{DEus} não é apenas o \emph{companheiro} dos agentes finitos; é o \emph{senhor} e \emph{moderador}, que se serve d'elles, como de instrumentos, para a execução das obras da sua providencia; é a Causa primeira, que actua immediatamente, não nos effeitos, mas nas causas segundas, infundindo-lhes força ou vigor, completando-as o applicando-as ao acto; é o agente universal, cuja influencia abrange a natureza inteira, todos os agentes de cada especie, todas as energias de cada agente. --- Esta influencia da Causa primeira, que determina e applica ao acto as causas segundas, ó chamada \emph{premocádo physica}. E' uma \emph{premoção}, porque a moção da Causa primeira \emph{precede}, pela sua natureza e causalidade, o acto das causas segundas, --- e é uma premoção \emph{physica}, porque ó o acto de uma Causa efficiente, que move e applica activamente a causa segunda ao acto connatural. A \emph{premocdo} chama-se tambem \emph{predeterminagdo}, no mesmo sentido e pelo mesmo motivo.\par Os \emph{thomistas} provam a sua theoria com muitos argumentos, baseados nos principios do Angelico Doutor. Expomos os principaes.\par \emph{a}) \textsc{DEUS} é o primeiro motor. As creaturas não podem operar senão emquanto são movidas e applicadas por \textsc{Deus}; não só porque toda a passagem da potencia para o acto deve ter e tem a sua razão sufficiente na Causa primeira, mas tambem porque um motor segundo deve ser necessariamente movido pelo motor primeiro. Esta moção de \textsc{Deus} --- é \emph{physica}, porque deriva d'Elle, como da primeira Causa efficiente, e influe physicamente em todas as creaturas, --- e é \emph{preveniente}, isto é, anterior ao acto da creatura, porque a moção do motor precede, na ordem da natureza e da causalidade, a moção do movel: « \emph{Motio moventis praecedit motum mobilis}, \emph{ratione}. \emph{et causa} » (\emph{C}. \emph{Gent}., III, 149). Ora uma moção physica e preveniente é uma \emph{premocdo physica}. Logo deve admittir-se esta \emph{premoção}. --- Accrescente-se que, se a causa segunda não é applicada ao acto por uma moção actual e preveniente, mas só recebe de \textsc{DEus} o auxilio de um concurso \emph{simultaneo}, nem ella nem o seu effeito dependem de \textsc{Deus}, como da Causa verdadeiramente primeira. Uma causa não é nem pode dizer-se \emph{primeira} com relação a outra, se não emquanto influe actualmente nella.\par \emph{D}) As causas segundas podem operar e influir no effeito, mas não operam e não influem sempre, e, quando operam, passam da potencia para o acto. Este acto, ou este influxo actual, é uma entidade, uma perfeição, distincta da potencia. Sendo uma entidade, uma perfeição, o acto da creatura deve derivar de \textsc{Deus}, fonte unica de toda a entidade e perfeição. Mas \textsc{Deus} não communica ás causas segundas o poder de influir actualmente no effeito, senão emquanto Elle exerce nessas causas um influxo proprio. Este influxo divino deve preceder naturalmente o influxo das creaturas, como a causa precede o effeito, \emph{e} por isso é uma verdadeira \emph{premocüo physica}. =--- Nem se diga que o influxo actual, que as causas segundas exercem no effeito, deriva de \textsc{Deus}, emquanto \textsc{Deus} concedeu a essas causas as faculdades, de que o influxo deriva. Porquanto, \textsc{Deus}, dando ás causas segundas\pgn{536}as faculdades ou energias operativas, deu-lhes o \emph{poder} de operar, mas não lhes deu o \emph{exercicio}, o \emph{influxo actual} d'esse poder. Quer Elle que, em relação a esse \emph{exercicio}, \& causa segunda dependa da Causa primeira, como o ser segundo depende do Ser primeiro.\par c) E' certo que uma grande parte das causas creadas são por si indifferentes e contingentes, como são os seres dotados de liberdade. Ora uma causa indifferente e contingente, que pode produzir varios actos mas que por si não é determinada para nenhum d'elles, ficará sempre no seu estado de potencialidade e não produzirá acto algum, se não fôr determinada. Esta determinação --- não pode provir da propria \emph{causa}, que é por si indifferente, --- nem do \emph{acto}, que é um effeito, e um effeito nao pode determinar a causa, --- nem do \emph{objecto}, porque este, quando se refere a uma causa \emph{livre}, excerce uma moção moral e não physica, deixa a causa na sua indifferenca, e, se \emph{especifica} o acto, não determina o \emph{exercicio} d'esse acto. Logo a determinação não pode provir, em ultima analyse, senão da Causa primeira por uma moção anterior ao acto da causa segunda, isto é, por uma \emph{premoção physica}. (\emph{Sum}. \emph{Th}, 1-2, q. 9, a. 4).\par \emph{d}) Em todas as faculdades operativas devemos distinguir dois estados --- um \emph{habitual} e outro \emph{actual}. Encontra-se a faculdade no estado \emph{habitual}, ou em \emph{acto primeiro}, quando \emph{pode} ou é \emph{capaz} de operar, mas actualmente não opera; --- encontra-se no estado \emph{actual}, ou em \emph{acto segundo}, quando actualmente opera. Logo uma faculdade operativa, para que possa actualmente operar, deve passar do estado habitual para o actual, do acto primeiro para o segundo. Mas não pode passar por si mesma, pela propria energia; porque está em potencia para o acto, e um ente em potencia só pode ser movido por um ente em acto, e porque o estado habitual é menos perfeito que o actual e nenhum ente pode dar a si mesmo uma perfeição superior \& que possue. Logo a faculdade deve passar da potencia para o acto pela energia de um motor, que seja acto purissimo, porque só este pode dar ás forças creadas a ultima actualidade. Diz S. Thomaz: « Omne quod quandoque est agens in actu, et quandoque in potentia, indiget moveri ab aliquo movente » (\emph{Sum}. \emph{T}., 1-2, q. 9, a. 4). --- « Oportet enim ut id quod est in potentia reducatur in actum per aliquid quod est in actu » (\emph{l}. c., a. 1). Este motor, que completa as faculdades das causas segundas e que é acto purissimo, é o primeiro agente, é \textsc{DEus}. Diz o mesmo S\&S. Doutor: « Complementum virtutis agentis secundi est ex virtute agentis primi » (\emph{C}. \emph{Gent}.. III, 66). Este complemento, esta ultima actualidade é communicada por \textsc{Deus} por uma moção \emph{fransitoria}, que só nElle tem uma existencia fixa (\emph{De pot}., q. 8, a. 7). Tal moção \emph{transitoria} não pode ser um concurso \emph{simultaneo}; porque o concurso \emph{simultaneo} suppõe o agente no estado \emph{actual}, e não no \emph{habitual}; logo deve ser \emph{previa}, deve ser uma \emph{premocéáo}.\par €) Os defensores do \emph{concurso previo} appellam tambem para a auctoridade do \emph{Catecismo do Concilio de Trento}, que diz: « Non solum Deus universa quae sunt providentia sua tuetur atque administrat, verum etiam .\emph{quae} moventur \emph{et agunt aliquid}, intima virtute, ad motum atque operationem ita \emph{zmpellit}, ut, quamvis causarum secundarum efficientiam \emph{non impediat}, \emph{praeveniat} tamen; cum ejus oceultissima vis ad singula pertineat, et,\pgn{537}quemadmodum Sapiens testatur (\emph{Sap}. VIII, 1), \emph{attingat} a \emph{fine usque ad finem fortiter et omnia disponat suaviter} » (P. I, in I art. Symb., n. 22). --- Trata-se --- de \emph{premoção}, porque \textsc{Deus} \emph{impelle} ao acto e \emph{prevem} a efficiencia das causas segundas, --- e de premoção 7/\emph{ijjsica}, e não de premocüo \emph{moral}, porque o \emph{Catecismo} falla de uma \emph{moção} commum a \emph{todos} os agentes, e do uma moção que deriva da omnipotencia divina e consiste numa \emph{virtude occulta} ou \emph{energia intima}, ora a premocao \emph{moral} --- não é \emph{commum} a \emph{todos} os agentes (porque só convem aos agentes \emph{l/vres}) --- não é attributo proprio da omnipotencia divina, porque o anjo e o homem. podem mover moralmente a nossa vontade, --- nem é uma energia occulta ou intima, como é evidente. --- Os \emph{thomistas} adduzem a auctoridade do \emph{Catecismo}, não para elevarem a \emph{premocdo physica} a um artigo de Fé, mas para mostrarem a solidez da sua theoria, recebida e considerada pela Egreja como a melhor explicação da omnipotencia divina e para refutarem os que disseram que a \emph{premoção physica} tinha sido condemnada no Concilio de Trento.\par Os defensores do concurso exclusivamente \emph{sémwultaneo} adduzem varios trechos de S. Thomaz, com que pretendem demonstrar que o Angelico Doutor é contrario a \emph{premocdo physica}. --- Os \emph{thomistas} respondem que esses trechos, rectamente interpretados, nada provam contra a \emph{premocádo}, e, para essa recta interpretação, apresentam os seguintes criterios: --- 1.9) Quando S. Thomaz affirma que \textsc{Deus} \emph{não predetermina} os nossos actos livres, ou que os actos da nossa vontade \emph{não derivam} de um principio externo \emph{determinante}, o S. Doutor refere-se a predeterminacaüo \emph{natural}, que impõe ás coisas um certo e determinado modo de operar e exclue toda e qualquer liberdade de escolha, mas não exclue a predeterminação propria de um ser livre, que livremente é applicado ao seu acto (\emph{Cf}. \emph{Sum}. \emph{Th}., p. I, q. 82, a. 1 ad 1). --- 2.9) Quando S. Thomaz diz que « \emph{nossa vontade se move a si mesma}, \emph{que se determina e é senhora dos seus actos}, ete., attribue á vontade o que é proprio d'ella, como \emph{agente provimo}, mas não exclue, antes suppõe e exige a moção e a applieacio do \emph{primeiro agente} (\emph{Sum}. \emph{Th}. p. I, q. 88, a. 1 ad 8). --- 3.º) Quando S. Thomaz ensina que \textsc{Deus} move a vontade, sem deixar nella impressão alguma, não excluo com isso a moção divina, mas apenas um novo estado habitual, uma nova qualidade permanente: porque, quando \textsc{Deus} influe o vigor, a energia na creatura, para ella operar actualmente, a acção divina ma creatura é \emph{fransítoria}, e por isso não produz um novo estado ou uma nova qualidade (\emph{De pot}., q. 3, a. €; --- \emph{De verit}., q. 99, a. 8). --- 4.º) Quando S. Thomaz sustenta que \textsc{Deus} move a vontade para o bem \emph{universal} e que a propria vontade se move por si mesma para este ou aquelle bem \emph{particular}, refere-so o 8. Doutor áquella \emph{primitiva moção especial}, pela qual \textsc{Deus} imprime na vontade humana uma tendencia \emph{necessaria} para o bem \emph{universal}, para a \emph{felicidade}, mas não exclue a divina moção \emph{geral}, que dá á nossa vontade e a todos os agentes o vigor, a energia indispensavel para o exercicio do acto; porque ensina expressamente, em varios logares das suas obras, que /\emph{odo} o acto da nossa vontade deriva --- não só da propria vontade, como do agente \emph{immediato} --- mas tambem, e sobretudo, de \textsc{DEus}, como do primeiro agente (\emph{Sum}. \emph{Th}., 1-9, q. 109, a. 1; --- \emph{De verit}. q. 22, a. 8).}.

100. O \textbf{concurso previo de Deus deixa intacta} a \textbf{liberdade da vontade humana}. --- Sendo o homem causa segunda, como o é todo e qualquer agente creado, para que a liberdade da sua vontade \pg{533} \pg{534}  fique intacta, não se requer que a vontade não seja movida de modo algum pela Causa primeira, mas requer-se e basta que essa moção não tire á vontade a sua indifferenca, o dominio

74 \pg{535}  de seus actos. Ora a moção divina, sendo \emph{natural} e \emph{intrinseca} ás faculdades das creaturas, não só não destroe o que é proprio d' essas faculdades, mas presta-lhes o auxilio, para que \pg{536}  possam passar da potencia para o acto e operar em conformidade com a sua natureza; de modo que, movendo uma \pg{537}  faculdade \emph{necessaria}, \textsc{Deus} faz com que esta produza um acto \emph{necessario}, e, movendo uma faculdade \emph{livre}, faz com que esta \pg{538}  produza um acto \emph{livre}. Logo o concurso \emph{previo} de \textsc{Deus} deixa intacta a liberdade da vontade humana \footnote{E' o argumento de S. Thomaz, que diz: « Liberum arbitrum est causa sui motus, quia homo per liberum arbitrium seipsum movet ad agendum. Non tamen hoc est de necessitate libertatis quod sit prima causa sui id quod liberum est: sieut nec ad hoo quod aliquid sit eausa alterius, requiritur quod sit prima causa ejus. Deus igitur est prima causa movens et \emph{naturales} causas, et \emph{voluntarias}. Et sicut \emph{naturalibus} causis, movendo eas, non aufert quin actus earum sint \emph{naturales}; ita movendo causas \emph{voluntarias}, non aufert quin actiones earum sint \emph{voluntariae}, \emph{sed potius hoc in eis opevatur}; \emph{operatur enim}, \emph{in unoquoque secundum ejus proprietatem} » (\emph{Sum}. \emph{Th}., P. 1, q. 88, a. 1 ad 8).\par Esta conclusão está contida nas conclusões precedentes. E') certo que \textsc{Deus} \emph{premove} a nossa vontade; é certo que a nossa vontade é \emph{livre}; logo deve ser tambem certo que a premocüo divina deixa intacta a liberdade da nossa vontade. --- Alguns philosophos e theologos reconhecem que \textsc{Deus} \emph{move} a vontade humana, e que a vontade, movida por \textsc{DEus}, permanece livre; mas não admittem a \emph{premordo}, porque dizem que esta é incompativel com a liberdade. Mas não tém razão. Já vimos que, segundo S. Thomaz, toda a \emph{moção} é uma \emph{premoção}; porque a moção do motor \emph{precede} necessariamente, na ordem da natureza e da causalidade, a moção do movel. Logo, se a \emph{premocdo} é incompativel com a liberdade, tambem é incompativel a moção; e, se não é incompativel esta, nem o pode ser aquella. O concurso \emph{simultaneo} não é verdadeira \emph{moção}; quem corre com outro, nem por isso o move. --- Os adversarios da \emph{premoção} esforgam-se para demonstrar que a nossa vontade, sob o impulso \emph{previo} de \textsc{Deus}, não pode ser \emph{livre}; mas todas as razões, que adduzem, fundam-se no supposto que \textsc{Deus} opere, como operam as creaturas. Porque uma causa finita e extrinseca, predeterminando physicamente a nossa vontade, tira-lhe a indifferenga e por isso a liberdade, concluem que se \textsc{Deus} predetermina physicamente ao acto a nossa vontade, tira-lhe tambem a indifferenca e a liberdade. Ora tal supposto é falso. Entre a causalidade creada e a causalidade divina poderá haver e ha uma \emph{analogia}, mas não ha uma \emph{egualdade}. A creatura não pode mover \emph{natuvalmente} á nossa vontade (porque esta, sendo faculdade de uma substancia espiritual, creada por \textsc{Deus}, escapa ao dominio das causas segundas), e por isso pode exercer nella uma violencia, e tirar-lhe de algum modo a liberdade; mas \textsc{Deus}, sendo o Creador e o Senhor da nossa vontade, pode movela \emph{naturalmente}, dando-lhe \emph{uma inclinação} para a coisa, para a qual a move. D'este modo, a moção divina é preveniente e infallivel, e comtudo a vontade humana conserva sempre a sua liberdade (\emph{Sum}. \emph{Th}., p. 1, q. 105, a. 4 ad 1; 1-2, q. 9, a. 6).\par Algumas diffieuldades foram já resolvidas. Vamos resolver ainda uma ou outra das que se apresentam com mais força.\par \emph{a}) Dizem: Se a nossa vontade é \emph{premovida} por \textsc{Deus}, ella não se determina a si mesma ao acto. Ora a essencia da liberdade consiste precisamente nesta auto-determinação. Logo a premoção é incompativel com a liberdade.\par Resposta. A premoção divina não só não é contraria a esta auto-determinação, mas é ella mesma que faz com que a vontade se determine a si mesma. Sem a premoção divina, a nossa vontade não poderia exercer nenhum\pgn{539} acto. E' a premoção, que, como o decreto divino, determina que uma causa \emph{livre} exerça, num certo tempo, um certo acto e o exerça \emph{livremente}. \textsc{Deus} é a causa universalissima e effieacissima. A nossa vontade, como diz S. Agostinho (\emph{De corrept}. \emph{et gratia} c. 14, n. 45), está mais nas mãos de \textsc{DEus} do que nas nossas; e por isso, continua S. Thomaz (\emph{De ver}., a. 99 a. 8), se a nossa vontade pode, pela virtude propria, passar de um acto para outro diverso, e até contrario, muito mais pode realizar essa passagem pela virtude divina, de que deriva e a que está subordinada toda a virtude creada. Mas \textsc{Deus}, repetimos, actua na nossa vontade, respeitando a dignidade d'ella e inclinando-a intrinsecamente para a coisa, para a qual Elle a premove, de modo que á nossa vontade produz o seu acto, porque assim quer.\par \emph{b}) Continuam: Se \textsc{Deus} \emph{premove} a nossa vontade, esta deve operar \emph{necessariamente}. Ora a necessidade oppóe-se á liberdade. Logo a premoção é incompativel com a liberdade.\par Resposta. Se Deus premove a nossa vontade, a operação d'esta faculdade é \emph{infallivel}, mas não é \emph{necessaria}. E' \emph{infallivel}; porque a premoc\&o de \textsc{Deus}, sendo effioacissima, á seguida pela operação da nossa vontade, e a vontade, quando opera, não pode deixar de operar. Mas não \& \emph{necessaria}; porque, como dissemos, a premoção determina a vontade em harmonia com a natureza d'esta faculdade, e por isso não só não lhe tira a liberdade, mas faz com que ella possa operar \emph{livremente}, como se fosse inteiramente independente de todo e qualquer concurso. Diz S. Thomaz: « Deus movet quidem voluntatem \emph{émmutabiliter} propter \emph{efficaciam virtutis moventis}, quae deficere non potest; sed propter naturam voluntatis motae, quae indifferenter se habet ad diversa, non \emph{inducitur necessitas}, \emph{sed manet liberlas}; sicut etiam in omnibus providentia divina \emph{infallibiliter} operatur; et tamen a causis contingentibus proveniunt effectus contingenter, inquantum Deus omnia movet \emph{proportionabiliter}, \emph{unwmquodque secundum suum modum} » (\emph{De malo}, q. 6, a. u. ad 3).\par c) Replicam: Se a premoção de \textsc{Deus} e a operação da nossa vontade estão infallivelmente unidas, é evidente que a nossa vontade, assim como não tem dominio sobre a premoção divina, assim tambem não o tem sobre a operação propria. Logo a premoção e a liberdade são incompativeis.\par Resposta. E' certo que a nossa vontade não tem dominio sobre o \emph{principio} da premocao, que é o proprio \textsc{DEus}, mas é falso que o não tenha sobre o \emph{termo}, que é o acto da mesma vontade. A nossa vontade, sob a premocao divina, pode suspender, interromper, revogar o proprio acto, sem prejuizo da efficacia da mesma premoção; porque o proprio poder de revogar ou\pgn{540}interromper o acto, que a premoção tinha começado, não pode deixar de .ser effeito da premogao. E', pois, sob a causalidade de \textsc{Deus}, motor primeiro, que 1 2 vontade passa livremente de um acto, ou de um objecto, para outro acto, ou outro objecto diverso e até contrario. --- Nem isto se oppõe á infallibilidade da Providencia divina. Quando a nossa vontade, dispondo livremente de si mesma, interrompe um acto, para o qual tinha sido determinada pela premoção, o novo acto não se realiza sem uma moção preveniente de \textsc{Deus}, á qual nenhum acto das creaturas, embora livre, pode escapar; de modo que as causas livres concorrem livremente para a execução dos divinos decretes.\par \emph{d}) Insistem: Se Deus premoye a nossa vontade para todos e cada um dos actos, deve premovel-a tambem para o peccado, e assim \textsc{Deus} é o auctor do peccado. Ora isto é absurdo e impio.\par Já respondemos a esta difficuldade. O peccado é um \emph{acto defeituoso}. Emquanto \emph{acto}, é uma perfeição e por isso deriva de \textsc{Deus}, como do agente primeiro. Emquanto \emph{defeituoso}, é uma imperfeição, que deve attribuir-se á deficiencia das causas segundas. Assim, o \emph{andar torto} de um coxo, --- emquanto é \emph{andar}, é uma perfeição e deriva da alma, fonte de toda a energia vital, --- emquanto é \emph{torto}, é imperfeição e deriva do defeito da perna. Diz S. Thomaz: « Deus est causa omnis actionis, in quantum est actio. Peccatum nominat \emph{ess} et \emph{actionem} cum quodam \emph{defectu}; defectus autem ille est ex causa creata, scilicet libero arbitrio, in quantum deficit ab ordine primi agentis, scilicet Dei. Unde defectus non reducitur in Deum, sicut in causam, sed in liberum arbitrium; sicut defectus claudicationis reducitur in tibiam curvam, sicut in causam, non autem in virtutem motivam, a qua tamen causatur quidquid est motionis in claudicatione. Et secundum hoc Deus est causa «\emph{actus} peccati, non tamen est causa \emph{peccati}, quia non est causa hujus, quod actus sit cum defectu » (\emph{Sum}. \emph{Th}, 1-2, q. 79, a. 2). --- Do resto, esta difficuldade pode oppor-se tambem contra a theoria do concurso \emph{simultaneo}, porque tambem essa theoria admitte que \textsc{DEus} \emph{concorre immediatamente} em todos os nossos actos, e, para o caso, pouco importa que esse concurso seja \emph{previo} ou \emph{simultaneo}.\par Concluimos que, embora a theoria do concurso \emph{simultaneo pareça} explicar ou salvar a liberdade humana melhor do que a do concurso \emph{previo}, comtudo é certo, e todos o reconhecem, que a theoria do concurso \emph{previo} favorece e salva a universalidade e a efficacia da causalidade divina melhor que a outra. Dissemos --- embora \emph{pareça}, porque, na \emph{realidade}, a theoria do concurso \emph{previo}, mais do que a do concurso \emph{simultaneo}, favorece não só os direitos de Deus, mas tambem os da nossa liberdade. Por quanto, a perfeição da creatura consiste na sua união com \textsc{DEus}: e quanto mais intima fôr esta união, tanto mais elevada será a sua perfeição. Ora a união nossa com \textsc{Deus} em nenhum systema se affirma mais intima, como no systema de concurso \emph{previo}, o qual propugna que a nossa vontade não passa da liberdade habitual para a liberdade actual, se não porque ou emquanto é movida pela primeira e increada liberdade.}.

%%>>> art-47

\section*{Artigo IV.}
\addcontentsline{toc}{section}{Artigo IV.}

\pg{539} \textbf{Providencia}, \textbf{sua existencia}, \textbf{extensão}

e \textbf{propriedades}.

101. \textbf{Providencia}. --- \emph{Providencia} é o acto, pelo qual \textsc{Deus} \emph{ordena} e \emph{dirige} todas as coisas para o fim. --- A Providencia, como se vé, importa duas coisas: a \emph{ordenação}, ou \emph{disposição}, \pg{540}  das coisas para o fim, e a \emph{execução} d'essa ordenação. A \emph{ordenação}, por ser acto da intelligencia divina, é \emph{eterna}; a \emph{execução}, por ser realisada no tempo, é \emph{temporanea} \footnote{\textsc{Deus}, como vimos, \emph{creou} todas as coisas; depois de as ter creado, \emph{conserva-as} na existencia, e \emph{concorre} em todos os actos, que ellas produzem.\pgn{541}Mas não basta que \textsc{DEus} concorra em todos os actos das creaturas; é necessario tambem que \emph{regule} esses actos e os \emph{dirija} para o fim, que estabeleceu ás mesmas creaturas. Esta direcção é obra da \emph{providencia}. --- A providencia, como dizemos no texto, abrange duas coisas: a \emph{ordenação} ou \emph{disposição} de todas as creaturas para o fim, e a \emph{execução} d'essa ordenação. A \emph{ordenação}, ou \emph{disposição}, é um acto da \emph{intelligencia} divina, a qual concebe a \emph{idea} ou a razão da ordem das coisas para o fim; e esta ordenação chama-se propriamente \emph{providencia}, que foi por isso definida por \emph{Boecio}: « \emph{divina ratio in summo omnium Principe constituta}, \emph{quae cuncta disponit} » (\emph{De cons}., l. IV, pr. 6) A \emph{execução} accrescenta ao acto da intelligencia, que estabelece a ordem, o acto da \emph{vontade}, que manda a execução da ordem, e o acto da \emph{potencia}, que executa o mandato da vontade e diz-se \emph{governo}, \emph{direcção}. --- Embora a Providencia divina, considerada no seu complexo, importe uma acção concorde da intelligencia, da vontade e da potencia, comtudo, considerada no seu conceito \emph{formal}, é um acto proprio da \emph{intelligencia}; porque \emph{governar} é acto de \emph{prudencia} e de \emph{sabedoria}, e estas qualidades residem na intelligencia, como no seu sujeito. De facto, é proprio da intelligencia conhecer a relação entre os meios e o fim, e determinar os meios mais efficazes para a consecução do fim. --- A Providencia suppõe nas creaturas uma absoluta dependencia de \textsc{Deus}, um fim proprio de cada uma, e um fim commum e unico, para o qual estão todas ordenadas. --- Emquanto a \emph{ordenação} é uma coisa \emph{eterna}, porque eterno é acto da intelligencia que concebeu a ordem, a execução ou o \emph{governo} é uma coisa \emph{temporanea}, porque temporaneas são as creaturas, que Deus dirige para o fim. --- S. Thomaz resumiu estas noções nas seguintes palavras: « \emph{AdWf'ovidentiae} curam duo pertinent: scilicet \emph{ratio ordinis}, quae dicitur \emph{providentia} et \emph{dispositio}; et \emph{executio ordinis}, quae dicitur \emph{gubernatio}. Quorum \emph{primum} est \emph{aeternum}; \emph{secundum}, \emph{temporale} » (\emph{Sum}. \emph{Th}. p. I, q. 22, \emph{a} 1 ad 2). --- Quando o \emph{fim} e os \emph{meios} para o fim \emph{excedem} as forças e as exigencias das creaturas, a providencia é \emph{sobrenatural}; quando o fim e os meios são \emph{proporcionados} ás forças e ás exigencias das creaturas, a providencia é \emph{natural}. Da providencia \emph{sobrenatural} trata a Theologia revelada; aqui occupamo-nos da \emph{natural}. --- S. Thomaz desenvolve esta questão --- na \emph{Sum}. \emph{Th}., p. I, q. 22, --- \emph{c}. \emph{Gent}., 1. III, e. 64, --- \emph{De} ver., q. 5.}.

\pg{541} 102. \textbf{Deus ordena todas as coisas para} um fim.

a) Todo o agente, dotado de intelligencia, não só opera para um fim, mas tambem dispõe ou ordena todas as suas obras para esse fim. Ora \textsc{Deus} é o agente dotado de intelligencia infinita. Logo \textsc{Deus} ordena todas as coisas para um fim \footnote{O fim, que levou \textsc{Deus} a crear a mundo, é a sua bondade, e o fim, que \textsc{Deus} estabeleceu ás suas creaturas, é tambem a sua bondade. --- As creaturas têm tambem um fim \emph{proximo}, que é a perfeição ou a bondade de cada uma; mas o seu fim \emph{supremo} ou \emph{remoto} é sempre a bondade divina. --- "Tendendo para o fim \emph{proximo}, as creaturas tendem para o fim supremo.}.

b) Conforme demonstram a experiencia e a razão, cada uma das coisas creadas está ordenada ou disposta para um fim particular e proximo, que é a propria perfeição, e todas ellas estão ordenadas ou dispostas para um fim universal e supremo, \pg{542}  que é a .participação e a manifestação da perfeição divina. Ora esta ordenação ou disposição, sendo um bem e um bem das creaturas, deve derivar de \textsc{Deus}, fonte de todo o bem e creador de todas as coisas, e por isso deve existir na intelligencia divina a razão, o plano d'essa ordenação ou disposição. Logo \textsc{Deus} ordena todas as coisas para um fim \footnote{Diz S. Thomaz: « Omne bonum, quod est in rebus, a Deo creatum est. In rebus autem creatis invenitur bonum, non solum quantum ad substantiam rerum, sed etiam quantum ad ordinem earum in finem, et praecipue in finem ultimum, qui est bonitas divina. Hoc igitur bonum ordinis in rebus creatis existens a Deo creatum est. Cum autem Deus sit causa rerum per suum intellectum, et sic cujuslibet sui effectus, necesse est quod ratio ordinis rerum in finem in mente divina praeexistat. Ratio autem ordinandorum in finem proprie \emph{providentia} est » (\emph{Sum}. \emph{Th}., Pp: 1, q. 22, a. 1). --- Não nos demoramos neste ponto, que foi tratado e sufficientemente desenyolvido na \emph{Cosmologia geral} e no I capitulo da \emph{T/eodicéa}.}.

103. \textbf{Deus} dirige \textbf{todas as coisas para} o fim.

a) Sendo proprio de um agente perfeito produzir obras relativamente perfeitas, convem á infinita Bondade de \textsc{Deus} governar todas as coisas e dirigil-as para o mais alto grau de perfeição, a que podem chegar. Ora o mais alto grau de perfeição, a que podem chegar as coisas creadas, consiste na consecução do fim. Logo \textsc{Deus} dirige todas as coisas para o fim \footnote{Diz S. Thomaz: « Cum optimi sit optima producere, non convenit summae Dei bonitati quod res productas ad perfectum non perducat. UItima autem perfectio uniuscujusque est in consecutione finis. Unde ad divinam bonitatem pertinet, ut sicut produxit res in esse, ita etiam eas ad finem perducat: quod est \emph{gubernare} » (\emph{Sum}. \emph{Th}., p. I, q. 108, a. 1).}.

b) Se considerarmos as coisas creadas, veremos que todas ellas, embora diversas, não só estão ordenadas para um fim commum, mas tambem tendem para esse fim e o conseguem pelo uso dos meios mais simples e effieazes. Ora a razão sufficiente d'essa tendencia e d'essa consecução não pode encontrar-se nas energias naturaes das proprias coisas, que são geralmente cegas e se encontram todas no estado de potencialidade. Logo tal razão deve encontrar-se na Suprema vontade de \textsc{Deus}, que, executando a ordem por Elle estabelecida, governa e dirige todas as coisas para o fim \footnote{Tambem este argumento encontra-se em S. Thomaz, que diz: « Impossibile est aliqua contraria et dissonantia in unum ordinem concordare semper vel pluries, nisi alicuius gubernatione, ex qua omnibus et singulis tribuitur, ut ad certum finem tendant. Sed in mundo videmus res diversarum naturarum in unum ordinem concordare, non ut raro et a casu, sed ut semper vel in maiori parte. Oportet ergo esse aliquid, cujus providentia mundus gubernatur; et hoc dieimus Deum. (C. \emph{Gent}., 1. I, c. 13).}.

\pg{543} 104. A \textbf{Providencia extende-se} a \textbf{cada} uma \textbf{das coisas}.

a) Governando \textsc{Deus} as coisas, emquanto foram por Elle creadas, é claro que a sua Providencia se deve extender até onde se extende a sua causalidade. Ora a causalidade divina extende-se a cada uma das coisas, --- não só ás universaes mas tambem ás particulares, não só ás necessarias mas tambem ás contingentes, não só ás corporeas mas tambem ás incorporeas, não só ás substancias mas tambem aos accidentes, --- de modo que não ha coisa alguma, que não seja produzida por \textsc{Deus}. Logo a Providencia extende-se a cada uma das coisas, de modo que não ha e não pode haver coisa, que escape á disposição e ao governo de \textsc{Deus} \footnote{Diz S. Thomaz: « Cum omne agens agat propter finem, tantum se extendit ordinatio effectuum in finem, quantum se extendit causalitas primi agentis... Causalitas autem Dei, qui est primum agens, se extendit usque ad omnia entia, non solum quantum ad principia speciei, sed etiam quantum ad individualia principia, non solum incorruptibilium, sed etiam corruptibilium. Unde necesse est omnia, quae habent quocumque modo esse, ordinata esse a Deo in finem. Cum ergo nihil aliud sit Dei providentia quam ratio ordinis rerum in finem, necesse est omnia, inquantum participant esse, intantum subdi divinae providentiae » (\emph{Sum}. \emph{Th}., p. I, q. 22, a. 2).}.

b) Se a Providencia não se extendesse a cada uma das coisas, isto seria --- ou porque \textsc{Deus} \emph{não pode}, --- ou porque, embora possa, \emph{não quer}. Ora \textsc{DEus} --- \emph{pode}, porque é omnipotente e é capaz de regular e governar tudo com um acto simplieissimo da sua vontade, --- e \emph{quer}, porque um artifice infinitamente bom, como é \textsc{Deus}, não pode deixar de ordenar as suas obras e dirigil-as para o fim, que lhes estabeleceu e que ellas, abandonadas a si mesmas, não podem conseguir. Logo a Providencia extende-se a cada uma. das \textbf{egfífsas} \footnote{O cuidado, que, Deus tem de todas as coisas, ainda das que nos parecem mais vis, não é incompativel com a dignidade e grandeza de \textsc{Deus}. Todas as coisas são relativamente boas e perfeitas, porque todas participam da bondade e perfeição divina; e se \textsc{Deus} nada perdeu na sua dignidade, quando as creou, nada pode perder, quando as governa. Antes o contrario. A divina bondade não se revela nunca tão esplendidamente, como quando extende os seus cuidados aos seres mais pequeninos e imperceptiveis da terra. Diz Jrsus: « Respicite volatilia coeli, quoniam non serunt, neque metunt, neque congregant in horrea: et Pater vester coelestis pascit illa. Nonne vos magis pluris estis illis? » (\emph{Matth}., VI, 26; cf. \emph{ib}. X, 29-81; \emph{Luc}. XII, 21-29).}.

105. A Providencia tem um cuidado especial do homem.

e) A Providencia de \textsc{Deus}, sendo sapientissima e rectissima, deve ordenar e dirigir todas as creaturas em harmonia --- com a natureza, de que são dotadas, --- e com a ordem dos fins, \pg{544}  que as une e conserva unidas. Ora, o homem é mais nobre do que todas as creaturas racionaes, e, na ordem dos fins, é o fim \emph{proximo} d'essas creaturas; porque as creaturas irracionaes estão ordenadas ou dispostas para a perfeição do homem, e para o seu bem temporal e eterno. Logo a Providencia tem um euidado especial do homem \footnote{A Providencia é proporcionada á natureza e dignidade dos seres, que lhe estão sujeitos. Nem admira. Um bom pae de familia tem mais cuidado dos filhos do que dos servos, mais dos servos do que dos snimaes. Ora o homem é mais nobre do que todas as creaturas visiveis, e é o fim \emph{proximo} de todas ellas. Logo \textsc{Deus} deve ter e tem do homem um cuidado especial. --- HE? o que diz o mesmo Jesus nos logares citados.}.

b) O homem, por ser dotado de livre arbitrio, produz actos moralmente bons ou mãos, e por isso imputaveis e dignos de premio ou de castigo. Ora a indole de todo o governo deve corresponder á indole da coisa governada. Logo \textsc{Deus} deve governar o homem de modo que este se sirva rectamente do livre arbitrio, praticando o bem e fugindo do mal, e assim alcance o seu ultimo fim. --- De facto, \textsc{Deus} gravou no animo do homem a lei moral, que o dirija em todos os seus actos, e, para a devida sancção d'essa lei, prometteu premios e ameaçou castigos. Logo \textsc{Deus} tem um euidado especial do homem \footnote{Que \textsc{Deus} dirija e governe todas e cada uma das coisas, sobretudo o homem, com uma sapientissima e efficacissima providencia, prova-o tambem o unanime consenso do genero humano. A Religião, que se encontra em todos os povos, funda-se principalmente na convicção de que existe a Providencia divina. Por isso disse S. Agostinho: « Si Dei providentis non praesidet rebus humanis, nihil est de religione satagendum » (\emph{De util}. \emph{cred}., c. 16). Na verdade, se \textsc{Deus} abandonou todas as coisas creadas ao acaso, não têm razão de ser os juramentos, os sacrificios, os votos, as orações, --- Attesta o mesmo o instincto natural, que, na ameaça de um grande perigo, leva-nos a levantar os olhos, as mãos e \emph{a} voz para o ceu e a invocar o auxilio divino. --- Demonstram tambem esta verdade as blasphemias dos impios contra \textsc{Deus}; porque os impios não as pronunciariam, se não estivessem convencidos de que nada acontece no mundo, que \textsc{Deus} não queira ou não permitta.\par Concluimos, pois, --- que \textsc{Deus}, principio primeiro e fim ultimo de toda a creacüo, governa todas e cada uma das creaturas, quer materiaes quer espirituaes, não só em geral e emquanto ás especies, mas tambem em particular e emquanto aos individuos de cada especie; --- que os individuos materiaes e corruptiveis estão ordenados e servem para o bem e a conservação da especie, que se perpetua no meio de uma continua geração e corrupção dos individuos; --- que as creaturas racionaes, como é o homem, são\pgn{545}governadas com uma providencia especial, proporcionada á natureza o á dignidade d'ellas, e por isso por meio de leis, preceitos, conselhos, exhortações, promessas de premios e ameacas de castigos.\par *\par São muitos os adversarios da Providencia divina, não só entre os \emph{antigos}, mas tambem entre os \emph{modernos}. --- Na antiguidade, negaram a existencia da providencia --- os \emph{epicureus}, que chamavam a providencia um jugo insupportavel para o homem, e indecoroso para \textsc{Deus}, --- os \emph{estoicos}, ensinando que todas as coisas, mesmo as nossas vontades eo proprio \textsc{DEUS}, estão sujeitas ás leis do \emph{fato}, --- alguns sequazes de \emph{Platão}, que, embora admittissem uma providencia \emph{gera}! com relação aos generos, ás especies e ás causas universaes, negavam uma providencia \emph{particular} com relação aos individuos e ás causas particulares.\par Nos tempos modernos negaram a providencia --- \emph{indirectamente} os \emph{pessimistas} e os \emph{fatalistas}, --- \emph{directamente} os \emph{deistas}.\par Entre os \emph{pessimistas} destacam-se \emph{Schopenhauer}, \emph{Hartmann} e os \emph{Neo-Budhistas}. Na opinião de \emph{Schopenhauer}, o mundo \&\&a evolução de uma unica \emph{vontade} cega e maligna, que tudo move de uma maneira absurda. Na opiniio de \emph{Hartmann}, o principio da evolução universal é o \emph{ente-inconscio}, cuja maligna vontade é temperada por outro principio, isto é, pelo desejo da anniquilacao, e por isso o mundo, embora essencialmente mau e peior que o nada, comtudo não é absolutamente pessimo. Na opinião dos \emph{Neo-Budhistas}, a vida presente, cheia de miserias, é intolleravel, e os homens só podem encontrar a felicidade em \emph{Nirvana}, isto é, na total anniquilação da propria substancia. --- Os \emph{futalistas} sustentam que todas as coisas estão sujeitas a um \emph{fato} indeclinavel. Patrono do moderno \emph{fatalismo é Taine}.\par Os \emph{deistas}, embora admittam a existencia de Dus, rejeitam, em tudo ou em parte, o seu governo no mundo. Por isso, o seu systema, chamado \emph{deismo}, \emph{é absoluto} e \emph{relativo}. --- Os sequazes do deismo \emph{absoluto} rejeitam totalmente a providencia, porque incompativel com a justiça, bondade e santidade de \textsc{DEus}. Assim dizem \emph{Bayle}, \emph{Voltaire}, \emph{Rousseau}, \emph{Diderot}, \emph{d}' \emph{Alembert}. --- Os sequazes do deismo \emph{relativo} dizem --- que \textsc{Deus} estabeleceu leis universaes, pelas quaes todos os seres são governados, --- mas que a sua providencia não se extende a todos e a cada um dos seres creados, porque isso repugna á immutabilidade divina e á liberdade humana, --- e que por isso toda a religião revelada é impossivel e a oração, emquanto deprecativa, é inutil e injuriosa para \textsc{Deus}. Estas ideas encontram-se, mais ou menos explicitas, nas obras de \emph{Thiers}, \emph{Cousin}, \emph{J}. \emph{Simon}, \emph{Saisset}, \emph{Damiron}, \emph{G}. \emph{Séailles}.\par Tendo provado, no texto, a \emph{existencia} e a \emph{universalidade} da providencia divina, resta-nos demonstrar a falsidade do \emph{pessimismo} -e a inefficacia dos argumentos dos \emph{deístas}.\par O \emph{pessimismo} --- é absurdo nos seus principios, --- 6 contrario ao senso commum e á natureza, --- e leva ás mais pernieiosas consequencias. --- E' absurdo nos seus principios. Por quanto, a unica vontade cega e maligna de \emph{Schopenhauer}, o \emph{ente-inconscio} de \emph{Hartmann}, como o \emph{eu-puro} de \emph{Fichte},\par 35 Vor. II.\pgn{546}o \emph{absoluto} de \emph{Schelling}, a \emph{idea-ente} de \emph{Hegel}, encerram o mais grosseiro \emph{pantheismo}. Além d'isso, a vida não é puro soffrimento. A vida é, por si, boa: o soffrimento é \emph{accidental}. --- E' contrario ao senso commum; por quanto, todos os homens desejam a vida e temem a morte, e, commummente uma vida miserayel prefere-se a anniquilacüo da propria substancia. --- E' tambem contrario á natureza; visto que todos os homens, opprimidos pela dor, seniem um desejo innato de uma felicidade futura. --- Leva ás mais perniciosas consequencias; porque o homem, que soffre, se nega a existencia de uma Providencia superior e perde a esperança numa vida futura, não pode ver diante de si senão o desespero e o suicidio, e não pode desejar senão a extincção total do genero humano e a absorpção dos individuos pelo famoso \emph{Nirvana}, isto é, pelo \emph{nada}. Os \emph{pessimistas} admittem estas consequencias.\par Os \emph{deistas} adduzem muitas e varias difficuldades contra a Providencia divina. Expomos e resolvemos as principaes.\par a) Dizem: Se todas as coisas fossem governadas pela divina Providencia, não poderia produzirse nenhum effeito \emph{fortuito} ou \emph{casual}. Ora produzem-se muitos effeitos \emph{fortuitos} ou \emph{casuaes}. Logo nem todas as coisas estão sujeitas á divina Providencia.\par Resposta. Concedemos que haja effeitos \emph{furtwitos} ou \emph{casuaes} com relação ás causas \emph{segundas} e \emph{particulares}, mas negamos que os haja ou possa haver com relação á Causa \emph{primeira} e \emph{universal}. Na verdade, \emph{fortuito} ou \emph{casual} diz-se o effeito que se produz \emph{fóra da intenção do agente}. Ora nenhum effeito pode escapar e escapa á ordem e ao influxo da Causa primeira e universal, que tudo conhece e cuja acção não pode ser estorvada por nenhum outro agente; e por isso, com relação a \textsc{Deus}, nada pode chamar-se fortuito ou casual, pois tudo acontece por divina disposição. Pelo contrario, o effeito de uma causa segunda e particular pode ser estorvado pela acção contraria de outros agentes; e por isso, pode produzir-se um effeito, que é fortuito relativamente á causa particular, porque produz-se fóra da ordem e da intenção d'essa causa, ainda que não se produza fóra da ordem e da intenção da causa universal, porque todos os actos das causas particulares estão contidos na ordem da causa universal, e tudo acontece conforme a disposição: d'essa causa universal. (\emph{Sum}. \emph{Th}., p. T, q. 22, a. 2 ad 1).\par \emph{b}) Continuam: Um governador sapiente e bom afasta, por quanto é possivel, das coisas, que rege, todo o mal. Ora no mundo ha muitos males. Logo o mundo não é governado por Deus.\par Resposta. Um governador sapiente e bom afasta das coisas, que rege, o mal da \emph{communidade}, mas não deve afastar o mal do \emph{individuo}, quando este concorre para o bem da mesma communidade. Com effeito, o que constitue o maximo bem do mundo é a ordem ou disposição de todos os seres para o seu fim commum. Este bem não pode, ás vezes, alcançar-se senão com o prejuizo ou damno de um ser particular, que é como um membro da grande totalidade cosmica. N'este caso, o governador sapiente, que olha para o bem da sua communidade, permitte o mal do ser particular, porque e emquanto é ordenado ou concorre para o bem commum. Um governador bom e sapiente só é obrigado a afastar de um ser particular o defeito ou o mal, quando só tem ao seu cuidado esse ser particular. Ora \textsc{DEus} é o\pgn{547}governador \emph{universal} de todas as creaturas, e por isso deve \emph{permitir} o defeito ou o mal de qualquer ser particular, para que não soffra nem seja estorvado o bem perfeito e total do universo. Porquanto, se \textsc{DEus} não permittisse nenhum mal, muito bem deixaria de se dar ou do existir. Por isso dizia S. Agostinho (\emph{Enchir}., c. XI): « Deus omnipotens nullo modo sineret malum aliquod esse in operibus suis, nisi usque adeo esset omnipotens et bonus, ut bene faceret etiam de malo» (\emph{Cf}. \emph{Sum}. \emph{Th}, \emph{ib}., ad 2).\par \emph{c}) Replicam: Se existisse a Providencia divina, não existiriam no mundo os males \emph{physicos}, como a \emph{dór}, a \emph{pobreza}, as \emph{tempestades}, os \emph{terremotos}, etc. Ora estes males existem para a infelicidade do genero humano.\par Resposta. \textsc{Deus}, como dissemos, não quer \emph{directamente} os males \emph{physicos}; pois o mal \emph{physico} é a privação de uma entidade ou perfeição physica, e o termo directo da vontade divina não pode ser uma \emph{privação}. Logo \textsc{Deus} quer esses males \emph{indirectamente}, emquanto os \emph{permitte}. --- Esta permissão não só não exclue, mas até suppõe a providencia divina. Por quanto, o mal physico \emph{particular} --- não só serve para a consecução do bem physico \emph{geral}, --- mas tambem procura ao homem um bem de uma ordem mais elevada, assim a dor purifica e fortifica á nossa alma e é occasião de muitos meritos para a vida eterna.\par \emph{d}) Insistem: No mundo encontra-se o mal \emph{moral}, que é \emph{o peccado}. Ora o peccado é um mal tão grande, que por um governador sapiente e santo, como'é \textsc{Deus}, não pode, de maneira nenhuma, ser permittido ou tolerado.\par Resposta. O peccado, embora seja um mal tão grande, que chega a ter uma malicia quasi infinita, porque offensa de uma Bondade infinita, comtudo pode ser permittido. Por quanto, o poder peccar, ainda que não constitua a essencia da liberdade, comtudo é uma consequencia da liberdade, de que é dotada a creatura racional. Ora é conveniente que a Providencia divina governe as creaturas por um modo proporcionado á essencia d'ellas e que as deixe operar em conformidade com a natureza de cada uma, e por isso é conveniente que deixe ao homem o poder de praticar o mal, isto é, de commetter o peccado. --- Além d'isso, Deus não poderia supprimir no homem a faculdade:de peccar, sem destruir, ao mesmo tempo, a liberdade, Ora repugna que a Providencia divina possa destruir as qualidades, de que dotou as creaturas. --- Finalmente, esta \emph{permissão} do mal \emph{moral} não repugna nem á \emph{santidade}, nem \& \emph{sabedoria}, nem \& \emph{bondade} de \textsc{Deus}. Não repugna á \emph{santidade} ; porque esta exige que \textsc{Deus} \emph{odeie} o peccado, e não que \emph{impeça} todo o peccado, e porque \textsc{Deus} concede ao homem muitos auxilios para que possa evitar o peccado e exercer a virtude. Não repugna á \emph{sabedoria}; porque o peccado não destroe a ordem \emph{universal}, que consiste na \emph{gloria divina}, mas só destroe a ordem \emph{particular}, que consiste na subordinação de uma creatura ao seu Creador, e esta destruição da ordem \emph{particular} serve para a conservação e perfeição da ordem \emph{universal}, porque tambem os peccadores devem dar gloria a \textsc{Deus}, exaltando a sua \emph{Justiça}, se não quizerem exaltar a sua \emph{misericordia}. Não repugna \& \emph{bondade}; porque a liberdade, de que os peccadores abusam, é um dom excellente da munificencia divina, é principio de \emph{merito} para a vida presente e de um \emph{premio eterno} na vida futura. O proprio \emph{Rousseau} diz: « Como?! para impedir que o homem seja mau, devia\pgn{548}elle ser rebaixado até á condição dos brutos? Não, Deus da minha alma, eu nunca me queixarei por me teres feito á tua imagem (\emph{Emile}, l. IV, c. 61).\par €) Acerescentam: Embora Deus permitta o peccado, comtudo não devia ser indifferente para com os bons e para com os maus. Ora parece e é indifferente; porque os peccadores levam sempre uma vida feliz, e os justos uma vida cheia de miserias e de amarguras. Esta desegualdade de tratamento mostra que não ha uma \emph{providencia}; pois, se houvesse, seria injusta. --- E' o raciocinio de \emph{Bayle}, \emph{Quinet}, \emph{Havet}, \emph{Michelet}.\par Resposta. E' falso que os peceadores levem sempre uma vida feliz e os justos uma vida miseravel; muitas vezes acontece o contrario. Por isso \& pergunta não deve formular-se nestes termos: porque é que os justos soffrem e os peccadores gosam? mas nos seguintes: porque é que os justos e os peccadores estão, promiscuamente, sujeitos aos bens e aos males de vida presente, quando os justos deveriam estar isentos de todos os males, e os peccadores privados de todos os bens? --- Formulada nestes termos, a pergunta pode ter muitas respostas. Antes de tudo, esta sujeição dos justos e dos peccadores aos bens e aos males da vida presente é uma consequencia da condição humana, exposta á acção diversa e contraria das causas segundas, Deus poderia, com certeza, fazer com que os justos não recebessem algum mal e os peccadores não recebessem algum bem; mas é conveniente que o não faça, que não impeça a actividade das causas segundas, sobretudo das livres, e não perturbe a constancia das leis geraes. --- Alem d'isso, se \textsc{Deus} desse, immediatamente, aos justos o premio e aos peccadores o castigo, ficaria destruida toda a ordem moral; porque, então, os homens praticariam o bem não por amor da virtude, mas pela esperança de uma felicidade temporal, e evitariam o mal unicamente pelo medo do castigo. Acerescente-se que o bem e o mal d'esta vida não podem constituir a recompensa adequada da virtude e do vicio. --- Ha mais: essa desegualdade de tratamento, que ás vezes se encontra no mundo, poderia subministrar aos adversarios uma arma efficaz contra a providencia, se esta se extendesse e se limitasse á vida \emph{presente}. Mas não é assim. Porquanto \textsc{Deus} reserva, na vida \emph{futura}, \& recompensa \emph{adequada} para a virtude, e o castigo \emph{adequado} para o vicio. Por isso a desegualdade de tratamento na vida presente é o signal mais certo da existencia da providencia. --- De resto, os \emph{males}, de que são affligidos os bons, servem para os purificar dos seus defeitos, para os afastar dos perigos e seducções do mundo, e para lhes subministrar occasião de meritos, --- e os \emph{bens}, que os peccadores gosam ou possuem, são a recompensa de algumas boas obras, que elles não deixam de praticar, e são avisos da bondade divina, que não os quer eternamente perdidos.}.

106. Propriedades da Providencia divina. --- A Providencia divina é \emph{universal}, \emph{suave}, \emph{infallivel}, \emph{immediata}.

\pg{545} a) E' \emph{universal}. Por quanto, como dissemos, a \textbf{extensão} da Providencia de \textsc{Deus} está na razão directa da sua causalidade, e por isso extende-se a tudo o que tem razão de ente.

b) Eº \emph{suave}. \textsc{Deus}, sendo o Auctor da propria: natureza, \textbf{governa} as coisas conforme a inclinação de cada uma; e por \pg{546} \pg{547}  isso cada uma tende, não de um modo violento, mas espontaneo, para o seu fim, para o qual a dirige o impulso divino.

\pg{548} c) E' \emph{infallivel}. Sendo \textsc{Deus} o Agente absoluto e a Causa primeira de todas as coisas, só pode acontecer o que Dus decretou e no modo que \textsc{Deus} decretou, sem prejuizo do livre arbitrio, de que algumas creaturas são dotadas \footnote{« In hoe est \emph{immobilis} et \emph{certus} divinae providentiae ordo, quod ea quae ab ipso providentur, cuncta eveniunt eo modo quo ipse providet, sive \emph{necessario}, sive \emph{contingenter} » (\emph{Sum}. \emph{Th}., p. I, q. 82, a. 4 ad 2). --- E se perguntarmos a S. Thomaz a razão por que tudo acontece conforme os decretos da providencia divina, responde-nos o Angelico, dizendo que toda a inclinação dos seres, quer \emph{natural} quer \emph{voluntaria}, não é outra coisa senão uma\pgn{549}impressão do primeiro motor, como a direcção de uma flecha para o alvo não é senão o effeito da impulsão do archeiro, e que por isso os actos de todos os agentes, quer naturaes quer voluntarios, tendem.e chegam quasi espontaneamente ao termo, que a divina Providencia lhes determinou. Eis as suas palavras: « Omnis inclinatio alicujus rei vel naturalis vel voluntaria, nihil est aliud quam quaedam impressio a primo movente; sicut inclinatio sagittae ad signum determinatum nihil aliud est quam quaedam impressio a sagittante. Unde omnia quae aguntur vel \emph{naturaliter} vel \emph{voluntarie}, quasi \emph{propria sponte} perveniunt in id, ad quod divinitus ordinantur. Et ideo dicitur Deus \emph{omnia disponere suaviter} » (\emph{Sum}. \emph{Th}., p. I, q. 108, a. 8).}.

\pg{549} d) E' \emph{immediata}. Por quanto, a \emph{ordem} da Providencia foi estabelecida e predeterminada por \textsc{Deus}, desde a eternidade, sem o concurso das causas segundas; embora, na \emph{execução} d'essa ordem, \textsc{Deus} possa servir-se e se sirva d'essas causas. --- \textsc{Deus} serve-se das causas segundas, não porque lhe falte o poder, mas porque é t\&o bom, que deseja communicar a dignidade da causalidade tambem ás creaturas \footnote{Diz S. Thomaz: « Ad providentiam duo pertinent: scilicet \emph{ratio ordinis} verum provisarum in finem, et \emph{executio} hujus ordinis, quae \emph{gubernatio} dicitur. --- Quantum igitur ad \emph{primum} horum, Deus \emph{mmediate} omnibus providet; quia in suo intellectu habet \emph{rationem} omnium etiam minimorum, et quascumque causas aliquibus effectibus praefecit, dedit eis virtutem ad illos effectus producendos. Unde oportet quod \emph{ordinem} illorum effectuum \emph{in sua ratione praehabuerit}. --- Quantum autem ad \emph{secundum}, sunt aliqua \emph{media} divinae providentiae; quia inferiora gubernat per superiora, non propter defectum suae virtutis, sed propter abundantiam suae bonitatis, ut dignitatem causalitatis etiam creaturis communicet » (\emph{Sum}. \emph{Th}., p. I, q. 22. a. 9). --- Servindo-se \textsc{DEus} das causas segundas na \emph{execução} da ordem, que Elle, só por si, estabeleceu, é claro que os defeitos e as imperfeições, que se encontram no mundo, não se podem attribuir a \textsc{DEus}, mas unicamente a essas causas segundas. --- \textsc{Deus}, mesmo na \emph{execução} da sua ordem, pode prescindir e prescinde, algumas vezes, do concurso das causas segundas e opera immediatamente por si e fóra da ordem d'essas causas. E' quando faz um \emph{milagre}. Mas, neste caso, \textsc{DEUS} opera de um modo \emph{ewtraordénario}. --- Do \emph{milagre} \&ratàmos, com o devido desenvolvimento, na \emph{Cosmologia}.}.

107. A \textbf{Providencia} de \textbf{Deus não exclue} a \textbf{cooperação do homem}. --- A Providencia de \textsc{Deus} excluiria a cooperação do homem, se \textsc{Deus} fizesse, só por si, as nossas acções. Mas não é assim. \textsc{Deus} não faz, só por si, as nossas acções, mas presta o seu concurso aos que operam, e concede aos que pedem os auxilios necessarios para a consecução do fim. Logo a Providencia de \textsc{Deus} não exelue a cooperação do homem. --- D'ahi a utilidade e á necessidade da \emph{oração} \footnote{A \emph{providencia} de \textsc{Deus} não exelue, mas auxilia os cuidados humanos; assim como o seu \emph{concurso} não destroe, mas auxilia os actos da nossa vontade. Por isso, querendo alcançar um fim, devemos empregar todos os\pgn{550}meios licitos, que estiverem ao nosso alcance, e depois pedir a \textsc{Deus} pelo bom exito das nossas emprezas. A experiencia mostra que, muitas vezes, homens ignorantes e fracos, mas apoiados em \textsc{Deus}, levam ao fim negocios, que homens espertos e poderosos, mas confiados nas proprias forças, não poderam concluir. --- A \emph{immutabilidade} da providencia não exclue a utilidade da oração, como não impõe a necessidades ás coisas, por ella governadas. A \emph{oração}, mesmo quando \emph{efficas}, não muda a disposição eterna da providencia; porque o proprio acto, com que \textsc{Deus} decreta a concessão de uma graça a quem a pede, pertence á mesma ordem da providencia. Diz S. Thomaz: « Orationes apud Deum efficaces sunt, non tamen ordincm immutabilem divinae providentiae solvunt, quia et hoc ipsum quod tali petenti hoc concedatur sub ordine divinae providentiae cadit ». (\emph{C}. \emph{Gent}., 1. III, c. 96).}.

\pg{550} 108. \textbf{Corollario}. --- Creador e Conservador de todas as coisas. \textsc{Deus} tem sobre ellas um dominio absoluto, inalienavel. Ao dominio de \textsc{Deus} deve corresponder, da parte das creaturas, uma sujeição affectuosa, completa. Nesta sujeição \textsc{Deus} encontra a sua gloria, e a creatura uma felicidade eterna \footnote{Os theologos apontam outro titulo para o dominio de \textsc{Deus} sobre todos os homens: o titulo da \emph{Redempcüo}. Os homens eram escravos do \emph{peccado}, do \emph{demonio} e do \emph{inferno}. D'esta triplice escravidão resgatou-nos o \textsc{HowrEM-Dzvus}, levado unicamente pelo seu amor, e resgatou-nos offerecendo um preço infinito --- o seu \textsc{SANGUE}! Os resgatados da escravidão de um tyranno passam naturalmente para o dominio de quem deu o preco do resgate. Somos todos de Jrsus. Conquistado pelo amor, o dominio de \textsc{JESUS} exerce-se pelo amor e só exige uma sujeição de amor. E' o reino do Coração de Jesus! A este dominio e a esta sujeição referem-se as seguintes palavras de S. Paulo: u \emph{Non estis vestri}; \emph{empti enim estis pretio magno}. \emph{Glorificate et portate} DEUM \emph{in corpore vestro} » (I. \emph{Cor}., VI, 19-90).}.

\pg{551} 1. \textbf{Moral}. --- \emph{Moral} é a \emph{sciencia que}, \emph{por meio das regras de moralidade}, \emph{dirige os actos humanos ao nosso ultimo fim}. --- A \emph{Moral} --- é \emph{sciencia}, porque é um conjuncto ordenado de conclusões e principios relativos a um objecto determinado; --- \emph{que dirige os actos humanos}, isto é, os actos que derivam da nossa \emph{inteligencia} e \emph{livre vontade}; --- \emph{ao ultimo fim do homem}, porque a \emph{Moral} não se occupa dos fins \emph{proximos} e \emph{secundarios} dos nossos actos, mas do fim \emph{ultimo}; --- \emph{por meio das regras de moralidade}, porque os actos humanos só podem alcançar o nosso ultimo fim, quando \emph{moralmente} bons e honestos \footnote{A \emph{Moral} é \& sciencia dos \emph{costumes} (\emph{mores}) --- \emph{Costume}, em sentido \emph{rigoroso}, é a repetição frequente e uniforme de certos actos, \emph{causada pelo livre arbitrio}. D'ahi a denominação de \emph{moral} applicada ao \emph{acto}, cuja repetição produz o \emph{costume}. Por isso acto \emph{moral} é o que deriva da vontade \emph{livre} do homem. --- A existencia da \emph{liberdade} dos nossos actos é, pois, uma condição necessaria para a existencia da \emph{Moral}. é\par D'aqui se segue que a \emph{Moral} dos \emph{pantheistas}, \emph{fatalistas}, \emph{deterministas}, \emph{positivistas}, e, em geral, de todos os que não reconhecem no homem o \emph{livre arbitrio}, tudo poderá ser, menos a sciencia dos \emph{costumes}.\par a) Os \emph{pantheistas} admittem que o homem, com os actos que pratica, e o mundo, com as substancias que o compõem, são emanagóes \emph{necessarias} da substancia divina. Sendo assim, não ha acções livres no homem, e, não havendo acções livres, não é possivel uma direcção \emph{moral}.\par b) Os \emph{fatalistas}, sustentando que todos os acontecimentos, que se realizum no homem e no mundo, estão. sujeitos a uma necessidade fat\&l e. inevitavel, não reconhecem acções \emph{livres}, e por isso não podem deixar de desvirtuar a natureza da \emph{Moral}.\par 0) Os \emph{deterministas}, negando de \emph{facto} a \emph{liberdade} da nossa vontade, como Provámos na \emph{Anthropologia}, devem tambem negar o objecto da \emph{Moral}, ---dã) Os \emph{positivistas} pertencem á classe dos \emph{deterministas}. Na verdade, os \emph{positivistas} admittem --- que a sciencia \emph{moral} é uma parte da \emph{physica social}, que é tambem uma especie da \emph{p/iysica organica}, sujeita a leis \emph{biologicas}, --- que todos os phenomenos humanos são o producto do organismo sob a acção do meio, --- que entre o homem e o animal não existe nenhuma differença\pgn{552}\emph{essencial}. D'estes principios deriva necessariamente a negação do \emph{livre arbitrio}, e por isso a negação da verdadeira Philosophia \emph{moral}. --- Não negamos que os \emph{positivistas} tambem estabelecem uma sciencia \emph{moral}; mas para elles a \emph{Moral} é a \emph{sciencia das condições do aperfeiçoamento das sociedades humanas}, de modo que a \emph{Moral}, no sentido d'esses escriptores, tem uma esphera muito mais extensa do que a que lhe marcou o bom senso, abrangendo não só a educação, a instrucção, a penalidade, mas tambem as bellas artes, as bellas lettras, isto é, abrangendo tudo o que não devia abranger e não abrangendo o que devia abranger, que é a \emph{livre} actividade do homem em ordem ao seu fim. --- Mas, ao menos neste ponto, os \emph{positivistas} são coherentes. Se a alma humana é uma coisa \emph{ficticia}, se \textsc{DEUS} é uma coisa \emph{imaginaria}, ou, pelo menos, \emph{incognoscivel}, como é possivel formar a verdadeira sciencia \emph{moral}, que dirija os actos \emph{livres} do homem para \textsc{Deus}, que é o nosso ultimo fim? --- Teremos muitas occasiões de fallar nas doutrinas \emph{positivistas}, e sobretudo na 2.º parte, em que daremos um resumo da \emph{physica social} ou \emph{sociologia}.}.

\pg{552} 2. \textbf{Objecto da Moral}. --- O \emph{objecto} da \emph{Moral} é constituido pelos actos \emph{humanos}, emquanto, pelas regras da \emph{moralidade}, se hão de dirigir para o nosso \emph{ultimo} fim. --- Por isso a \emph{Moral} --- é uma sciencia \emph{pratica}, porque applica os conhecimentos ás operações, --- e distingue-se das outras duas partes da Philosophia, que são a \emph{Logica} e a \emph{Metaphysica}, porque a \emph{Moral} regula os actos da \emph{vontade}, ao passo que a \emph{Logica} dirige os actos de \emph{intelligencia}, e a \emph{Metaphysica} estuda os entes \emph{reaes} \footnote{A. \emph{Moral} distingue-se das outras duas partes da Philosophia, que são a \emph{Logica} e \emph{Metaphysica}. Na verdade, a Philosophia é altissima sabedoria. Ora é proprio da sabedoria estabelecer a ordem entre as coisas. Às coisas, que se podem ordenar, --- ou são os actos da nossa \emph{tintelligencia}, --- ou são os actos da nossa \emph{vontade}, --- ou são os objectos que compõem o universo. Os actos da nossa \emph{intelligencia} constituem o ente \emph{logico}, que é o objecto da \emph{Logica}, pois esta dirige a nossa intelligencia para a \emph{verdade}; --- os actos da \emph{vontade} formam o ente \emph{moral}, que é o objecto da \emph{Moral}, porque a \emph{Moral} dirige a nossa vontade para o \emph{bem}; --- os objectos, que compõem o universo, constituem o ente \emph{real}, que é o objecto da \emph{Metaphysica}; porque a \emph{Metaphysica} trata dos entes que existem independentemente dos nossos actos. Se, pois, o objecto da \emph{Moral} é distincto do objecto da \emph{Logica} e da \emph{Metaphysica}, é evidente que á primeira ó uma sciencia distincta d'estas duas ultimas.\par A \emph{Moral} não só se distingue da \emph{Logica} e da \emph{Metaphysica}, mas tambem da \emph{Theologia moral}. Na verdade, o \emph{sujeito} da. \emph{Theologia moral} é o \emph{christão}, o da philosophia \emph{moral} é o \emph{homem}; --- os \emph{principios} da \emph{Theologia moral} são revelados ou deduzidos das verdades reveladas, os da \emph{philosophia moral} são descobertos pela luz natural da razão; --- a \emph{origem} dos preceitos da \emph{Theologia moral é divina}, porque derivam da Revelação sobrenatural, ao passo que a origem dos preceitos da philosophia \emph{moral é natural}, porque esses preceitos estão baseados na lei natural. --- D'onde se vê o erro de \emph{Kant}, \emph{De} Wette, \emph{Bruch}, que identificam a philosophia \emph{mora}! com a 'Theologia \emph{moral}.\par Todavia a Theologia \emph{moral} auxilia e illumina a philosophia \emph{moral}, afastando-a dos perigos e guiando-a para um caminho seguro. A historia da philosophia attesta que o homem, abandonado ás suas forças naturaes, tem\pgn{553}descoberto algumas verdades, mas só depois de muito tempo e misturadas com muitos erros. Siryam de exemplo os escriptos dos maiores entre os philosophos pagãos --- \emph{Socrates}, \emph{Platão}, \emph{Aristoteles}. \emph{Kant} diz: « A maravilhosa) religião do Christianismo, na sua grande singeleza, enriquece a philosophia, \emph{moral} com ideas moraes mais exactas e mais puras do que as que a mesma \textbar{} philosophia tinha subministrado até essa epocha ». (\emph{Crit}. \emph{Judic}., pag. 442). \textbar{} (1) D'onde se vé a utilidade, que do estudo da philosophia \emph{moral} deriva para o individuo e parfa sociedade. Por quanto, a felicidade do individuo e da sociedade consiste na consecução dos seus fins respectivos. Ora é a \emph{Moral} que aponta aos individuos e á sociedade o \emph{fim}, que se deve consc guir, e os \emph{meios}, que se devem empregar. --- E não só a \emph{Moral} é util para a vida \emph{pratica} do homem e da sociedade, mas tambem para o estudo das outras sciencias praticas, sobretudo do \emph{direito}, da \emph{politica} e da \emph{economia politica}. E' util para o estudo do \emph{direito}; porque as leis \emph{positivas}, de que se occupa o \emph{direito}, nunca poderão ser conhecidas perfeitamente, se não forem relacionadas com a lei \emph{natural}, em que se devem fundar e que forma o objecto da \emph{Moral}. E' util para o estudo da \emph{politica}; porque é a \emph{Moral} que determina o fim do homem e da sociedade, e os deveres dos governantes e dos governados. E' util para a \emph{economia politica}; porque a \emph{Moral} ensina a subordinação do bem \emph{util}, que é o objecto de \emph{economia politica}, ao bem \emph{honesto}, e, definindo as relações entre patrões e operarios, entre amos e servos, lança os fundamentos do bem estar social..}.

\pg{553} 3. \textbf{Excellencia da Moral}. --- A \emph{Moral} não só é mais excellente que todas as sciencias racionaes e experimentaes, mas tambem é mais nobre que a \emph{Logica} e a \emph{Metaphysica}. --- Eº mais nobre que a \emph{Logica}; porque á \emph{Logica} ensina a raciocinar, e á \emph{Moral} a bem viver. --- E' mais nobre que a \emph{Metaphysica}; porque, se esta descobre á nossa intelligencia a natureza do mundo e da nossa alma, a existencia de um \textsc{Deus} infinitamente sabio, poderoso e bom; a \emph{Moral}, applieando á pratica esses conhecimentos, dirige a nossa alma a consecução da felicidade eterna, que consiste na posse do proprio \textsc{Deus} \footnote{\pgn{554}movimento da epocha actual, perderemos o tempo em questões inuteis ou applioaremos sem tacto e criterio as leis geraes e os principios universaes ás circumstancias individuaes e concretas.\par 5) Os \emph{racionalistas}, como \emph{Kant}, prescindem inteiramente do auxilio da divina Revelação. --- Tambem estes se enganam. A razão humana, abandonada a si mesma, tem sempre errado, sobre tudo em assumptos moraes e religiosos; e por isso, se quizer encontrar á verdade, deve acceitar, com reconhecimento e veneração, o auxilio de Revelação divina, não para provar as suas proposições, mas para conhecer o caminho da verdade.}.

4. Methodo da Moral. --- O \emph{metodo} da \emph{Moral} --- a) não póde ser exclusivamente \emph{synthetico}, --- 5) nem exclusivamente \emph{analytico}, --- c) mas deve ser \emph{synthetico-analytico}.

\textbar{} 8) Não póde ser exclusivamente \emph{synthetico}, isto é, \emph{racional}; porque a \emph{Moral} deve considerar os actos \emph{humanos} não em \emph{abstracto} mas em \emph{concreto}, para se conhecer o modo de applicar os principios absolutos e universaes da \emph{moralidade} ás diversas circumstancias da vida humana \footnote{Exaggeram a força da \emph{razão}, no estudo da \emph{Moral}, os sequazes do \emph{direito natural puro}, e os \emph{racionalistas}.\par \emph{a}) Os sequazes do \emph{direito natural puro}, como \emph{Rousseau}, sem se importar da \emph{historia} e da \emph{experiencia}, pretendem deduzir da natureza humana, considerada em \emph{abstracto}, todas as regras de \emph{moral}. --- Mas enganam-se. Se no estudo da \emph{Moral} não tivermos presentes a historia do nosso espirito e o}.

\pg{554} b) Não póde ser exclusivamente \emph{analytico}, isto é, \emph{experimental}; porque, para dirigir os actos humanos para o seu fim ultimo, é necessario conhecer esse fim, e o fim de um ser conhece-se pelo estudo da sua essencia, e o estudo das essencias exige uma operação da \emph{razão}, que das propriedades e operações se eleve, pelo raciocinio, ao descobrimento do principio, de que derivam \footnote{Attenuam ou rejeitam a \emph{força} da \emph{razão}, no estudo da \emph{Moral}, os \emph{tradicionalistas}, os \emph{historicos}, os \emph{positivistas}, os \emph{scepticos}.\par a) Os \emph{tradicionalistas} dizem que a razão individual não póde conhecer, com \emph{certeza}, nenhuma verdade, sobre tudo de ordem \emph{moral}, e que por isso todas as verdades devem ser ensinadas pela auctoridade, divina ou humana. --- Não nos detemos aqui em refutar o \emph{tradicionalismo}; porque já foi refutado na \emph{Anthropologia}.\par b) Os \emph{Aistoricos} reconhecem a \emph{historia} como a unica fonte das regras de \emph{moralidade}. --- O exagero é manifesto. A \emph{Astoria}, embora seja de grande auxilio no tratado da \emph{Moral}, não póde comtudo ser a \emph{unica} fonte das regras moraes; porque, se não for estudada á luz de principios sãos, leva facilmente a muitos erros. Não são os principios que se hão de accommodar aos costumes, mas são estes que se devem conformar com aquelles.\par 0) Os \emph{positivistas} affirmam que devemos estudar a \emph{Moral}, independentemente da \emph{Metaphysica}, isto é, de todo o \emph{supposto transcendental}, apoiando-nos unica e exclusivamente na \emph{experiencia}. A vida \emph{intellectual} e \emph{moral} do homem, segundo elles, está sujeita, como todas as coisas, a uma \emph{evolução} constante; e por isso a \emph{Moral} é uma parte da \emph{Mechanica} e tira os seus principios da \emph{Physica}, da \emph{Biologia}, da- \emph{Psychologia} e da \emph{Sociologia}, que explicam as leis universaes da evolução. D'este modo o munus da \emph{Moral} consiste em appliear essas leis aos phenomenos da vida moral, mostrando --- não o que os homens devem fazer, pois os homens, segundo os \emph{positivistas}, são destituidos de liberdade, --- mas o que \emph{farão}. --- O erro d'estes escriptores é evidente. A ordem \emph{especulativa} não pode deixar de ser o fundamento da ordem \emph{moral}. Be não conhecermos ou se desvirtuarmos as grandes verdades, que se ensinam na \emph{Metaphysica} e que são o principio de \emph{contradicção} e de \emph{causalidade}, a existencia de \textsc{Deus}, a espiritualidade, liberdade e immortalidade da alma humana, o fim ultimo da nossa vida, não é possivel compor um tratado de \emph{Moral}, sem errarmos, em cada passo, ácerca da \emph{lei natural}, da sua obrigação e sancção, ácerca da consciencia, da honestidade, do direito. A \emph{Moral} dos \emph{positivistas} não passa de uma \emph{historia natural}. De resto, na \emph{Anthropologia} refutàmos o \emph{determinismo}, que é a base de \emph{Moral} dos \emph{positivistas} e á negação\pgn{555} da liberdade, e mais adeante veremos que a nossa vida \emph{moral} não está nem póde estar sujeita a uma evolução fatal e irresistivel.\par d) Os \emph{scepticos} negam que as leis moraes sejam universaes e immutaveis para todos os homens, em todÜs os tempos e em todos os logares. --- Mais tarde refutaremos esta opinião.}.

\pg{555} c) Mas deve ser \emph{analytico-synthetico}, isto é, \emph{experimental} e \emph{racional}. --- Deve ser \emph{analytico} ou \emph{esperimental}; porque a \emph{Moral} dirige os actos humanos, que, por serem \emph{contingentes}, só podem ser conhecidos pela experiencia. --- Deve ser \emph{synthetico} ou \emph{racional}; porque a \emph{Moral} dirige os actos humanos por meio de principios universaes e absolutos, que só podem ser percebidos pela razão \footnote{O methodo \emph{analytico-synthetico} foi seguido por \emph{Socrates}, \emph{Platão}, \emph{Aristoteles}, \emph{S}. \emph{Agostinho}, \emph{S}. \emph{Thomas} e por todos os escholasticos e escriptores de bom senso, que, nas suas especulações, se servem dos principios da razão e da experiencia. !}.

5. \textbf{Divisão da Moral}. --- A \emph{Moral} divide-se em duas partes: \emph{geral} e \emph{especial}. --- A Moral \emph{geral} trata dos \emph{actos humanos}, considerando-os no seu \emph{fim}, \emph{essencia}, \emph{moralidade} e \emph{regras}. A Moral \emph{especial} tem por objecto os \emph{direitos} e os \emph{deveres}, que são proprios do homem emquanto é \emph{individuo} e \emph{membro da sociedade}. --- A Moral \emph{geral} costuma chamar-se \emph{Ethica}; --- a Moral \emph{especial} denomina-se commummente \emph{Direito natural} \footnote{Advirta-se que a palavra \emph{ethica} (\emph{etos-costume}) tem a mesma significação que a palavra \emph{moral} (\emph{mos-costume}), e que por isso não é muito rigoroso dar a uma \emph{parte} da divisão o nome que significa o \emph{todo}. Mas é questão de palavras, e accommodamo-nos aos mais auctorizados escriptores.\par Os \emph{philosophos} e os \emph{jurístas} discutem se o \emph{direito natural} é, ou não, inseparavel da \emph{ethica}. As opiniões divergem; uns identificam as duas partes, outros separam-nas. --- A questão não póde ser resolvida aqui; visto que, para esta solução, é necessario conhecer as definições da \emph{ethica} e do \emph{direito natural}, e as relações que existem entre uma parte e outra, --- definições e relações, que só serão conhecidas no decurso d'este tratado.}.

%%>>> parte-3

\part*{PARTE PRIMEIRA\\[0.3em]\large ETHICA}
\addcontentsline{toc}{part}{PARTE PRIMEIRA --- ETHICA}

\pg{556} 6. \textbf{Ethica}. --- \emph{Etica} é a parte da \emph{Mora}! que trata da \emph{moralidade} dos actos \emph{humanos}. --- Os actos \emph{humanos} são dotados de \emph{moralidade}, quando estão dispostos para o seu \emph{fim} \footnote{A sciencia tambem tira do seu objecto \emph{formal} a propria \emph{especie} e \emph{definico}; porque o objecto \emph{formal} é o \emph{termo} para o qual tende, e o \emph{termo} especifica e distingue os seres que tendem para elle. Ora o objecto \emph{formal} da \emph{Ethica} é constituido pela \emph{moralidade} dos actos \emph{humanos} (os proprios actos \emph{humanos} são apenas o objecto \emph{material} da \emph{Ethica}) Logo a \emph{Ethica} deve fundar a sua definição na \emph{moralidade} dos actos humanos. =}.

7. \textbf{Objecto} da \textbf{Ethica}. --- O \emph{objecto} da \emph{Ethica} é, pois, constituido pela \emph{moralidade} dos actos \emph{humanos}. --- Como se vê, esta parte da \emph{Moral} estuda o acto humano \emph{em geral}, considerando-o --- no \emph{fim}, para o qual tende, --- na \emph{essencia}, que o constitue, --- na sua \emph{relação} para o mesmo \emph{fim}, --- e nas \emph{regras}, que o devem dirigir \footnote{Ainda que a Ethica se occupe propriamente da \emph{moralidade} dos actos humanos, comtudo não póde deixar de tratar da \emph{essencia}, que constitue os mesmos actos, --- do \emph{fim}, para que elles tendem, --- e das \emph{regras}, que os dirigem. Na verdade, a \emph{moralidade} dos actos humanos não poderá ser conhecida distinctamente, se não conhecermos --- a sua \emph{essencia} (pois o ser \emph{pluysico} é o fundamento do ser \emph{moral}), --- o /\emph{im}, para que os actos tendem (porque a \emph{moralidade} consiste em ordenar ou dispór os actos para o seu fim), --- e as \emph{regras}, que os dirigem (porque toda a ordem suppõe as suas regras).}.

8. \textbf{Divisão} da \textbf{Ethica}. --- A \emph{Hthica} divide-se em \emph{quatro}. capitulos; --- no \emph{primeiro}, trata do \emph{fim} dos actos humanos; --- no \emph{segundo}, da \emph{essencia} dos mesmos actos; --- no \emph{terceiro}, da sua \emph{moralidade}; --- no \emph{quarto}, das suas \emph{regras} \footnote{Tratamos em primeiro logar do \emph{fim} dos actos humanos. Por quanto, a relação, que existe entre o \emph{fim} e os \emph{actos moraes}, é semelhante á que existe entre um \emph{principio} especulativo e as suas \emph{consequencias}, Ora o principio especulativo é a regra e a medida das conclusões. Logo tambem o \emph{fim} é a regra e a medida dos actos moraes. A regra e a medida devem ser conhecidas antes do proprio objecto, que se deve subordinar a essa regra e medida. Logo o conhecimento do \emph{fim} dos actos humanos deve ser anterior ao conhecimento dos proprios actos.}.

\pg{557} FIM DOS \textbf{ACTOS} HUMANOS

\textsc{SuMmMARIO}: --- Existencia do fim ultimo dos actos humanos. --- Natureza do fim ultimo dos actos humanos. --- Fim proximo dos actos humanos.

%%>>> art-48

\section*{Artigo I.\\ Existencia do fim ultimo dos actos humanos.}
\addcontentsline{toc}{section}{Artigo I. --- Existencia do fim ultimo dos actos humanos.}

9. \textbf{Os actos} humanos tendem sempre para um fim. --- Todo o acto, que deriva de uma faculdade, não póde deixar de tender para o objecto \emph{proprio} da mesma faculdade; assim a percepção, que deriva da intelligencia, tende para o \emph{intelligivel}, que é o objecto \emph{proprio} da mesma intelligencia. Ora os actos humanos, derivando da vontade deliberada, não podem deixar de

\emph{z} tender para o \emph{bem}, que é o objecto da mesma vontade e que por isso é o \emph{termo} ou \emph{fim} da propria tendencia. Logo os actos humanos tendem sempre para um \emph{fim} \footnote{Fallamos dos actos //\emph{umanos}, isto é, dos actos que são proprios do homem e por isso derivam da intelligencia e da vontade, --- e não dos actos do \emph{homem}, isto é, dos actos que derivam do homem emquanto é dotado de vida vegetativa e sensitiva e que por isso não emanam da intelligencia e da vontade. Ora é evidente que os actos /\emph{iwmanos} devem ter um \emph{fim}; não só porque a intelligencia opéra sempre para um \emph{fim} (\emph{o acaso} é negação da intelligencia), mas tambem porque a vontade tem por objecto o \emph{bem}, que é \emph{fim}. --- Ha, pois, differença entre os actos /\emph{iumanos} e as operações dos seres irracionaes. Estes tendem para um /\emph{im}, mas não o conhecem; e, se o conhecem (como os animaes), não o conhecem como \emph{tal}, não percebem \emph{a relação} entre o fim e os meios, e por isso não podem regular os seus actos, nem escolher os meios mais efficazes. --- D'aqui se vé o erro, ou, pelo menos, a inexactidão de \emph{Spencer}, quando no seu livro --- \emph{Les bases de la orale evolutioniste}, applica indifferentemente ao homem e ao animal a palavra \emph{conducta}.\par Dissemos que o \emph{bem é fim} da vontade, porque a vontade tende necessariamente para o \emph{bem}. --- O \emph{bem} distingue-se em \emph{util}, \emph{honesto} e \emph{agradavel}. O bem \emph{util} é o que se deseja, não por si, mas para se alcançar outro bem; --- o \emph{honesto} é o que se deseja por si, como conforme á razão; --- o \emph{agradavel} é o que, possuido, satisfaz a faculdade appetitiva. O bem /\emph{ionesto} e o \emph{agradavel} têm razão de \emph{fim}; o \emph{util} tem só razão de \emph{meio}. Por isso o homem, quando opéra, ou é movido pela \emph{honestidade} ou pelo \emph{praser} (\emph{Sum}. \emph{TA}., 1-2, q. 1, a. 1).}.

10. 0 fim dos \textbf{actos} humanos é ultimo ou proximo. --- O fim dos actos humanos é sempre um \emph{em}. Este \emph{bem} --- ou se deseja \pg{558}  por si mesmo e é o motivo por que se desejam os outros bens, --- ou não se deseja por si mesmo, mas emquanto é \emph{meio} para alcançar um outro bem. No primeiro caso, o \emph{bem} é fim \emph{ultimo}, porque a sua consecução importa o descanço de todas as tendencias; --- no segundo caso, o \emph{bem} é fim \emph{proximo}, porque está subordinado ao fim \emph{ultimo}, de que recebe a força para attrahir a vontade, e por isso não sacia completamente as aspirações da alma. Logo o fim dos actos humanos é \emph{ultimo} ou \emph{proximo} \footnote{O fim \emph{proximo} dos actos humanos póde ser e é \emph{multiplice}, mas o fim \emph{ultimo} é um só. Na verdade, embora o homem seja dotado de muitas e diversas faculdades, cada uma das quaes tem o seu \emph{fim especial}, todavia entre essas diversas faculdades existe uma harmonia tão perfeita, que todas as tendencias inferiores estão subordinadas a uma tendencia superior, a qual é a perfeição propria e distinctiva da nossa natureza, E, como essa perfeição, é uma só, o fim \emph{ultimo} dos actos humanos é \emph{um} só. --- Os fins \emph{particulares} das faculdades inferiores são \emph{meios}, \emph{condições} ou \emph{complementos} secundarios e subordinados ao fim \emph{ultimo} da nossa natureza (\emph{Sum}. \emph{Th}., 1-2, q. 1, a. 5).}.

11. \textbf{Os actos} humanos tendem \textbf{sempre para} um fim ultimo. --- Os actos humanos tendem sempre para um \emph{fim}. Este fim --- ou se deseja por si e socega completamente as aspirações da alma, --- ou se deseja por amor de um fim mais elevado e não torna completamente feliz quem o possue. --- Se se deseja por si, o fim é \emph{ultimo}, e per isso está provada a conclusão. --- Se se deseja por amor de outro fim, o fim é \emph{proximo}, e por isso continuamos a investigação, até chegarmos a um fim \emph{ultimo}, que se deseje por si e por amor do qual se desejem todos os outros fins; pois, sendo o fim \emph{ultimo} a razão por que se desejam os fins \emph{proximos}, se não desejassemos aquelle, n\&o poderiamos desejar estes. Logo os actos humanos tendem sempre para um fim \emph{ultimo} \footnote{O fim \emph{ultimo} é a razão das outras \emph{voligües}. Na verdade, tudo que o homem deseja, deseja-o emquanto é \emph{bem}. Mas o \emph{bem}, que não é o summo e perfeitissimo bem e por isso não sacia completamente a vontade, deseja-se como um bem menos perfeito. Ora o bem \emph{menos} perfeito deseja-se emquanto se relaciona com um bem \emph{absolutamente} perfeito. Logo o fim \emph{ultimo} é a razão por que desejamos os fins \emph{secundarios}, que são os bens \emph{particulares}.\par Todavia o homem não tende sempre \emph{explicitamente} para o fim \emph{ultimo}, mas só \emph{implicita} ou \emph{virtualmente}. --- Não tende sempre \emph{explicitamente}; porque a experiencia attesta que a vontade muitas vezes deseja os bens \emph{particulares}, sem que a intelligencia pense no fim \emph{ultimo}. --- Mas tende sempre \emph{implicitamente}; porque só desejamos os fins \emph{proximos}, que são os bens \emph{particulares}, emquanto desejamos o fim \emph{ultimo}, que é o bem \emph{universal} (\emph{Sum}. \emph{Th}., \emph{ib}., a. 6).}.

12. 0 fim ultimo dos \textbf{actos} humanos é o \textbf{mesmo} para \textbf{todos os} homens. --- O fim \emph{ultimo} dos actos humanos é o que, saciando \pg{559}  todas as tendencias da alma, completa a perfeição do homem. Ora o objecto, que completa a perfeição de um ente, deve ser um e o mesmo para todos os que são dotados de natureza especificamente egual á do mesmo ente. Logo o fim \emph{ultimo} dos actos humanos é o mesmo para todos os homens \footnote{O fim \emph{ultimo} dos actos \emph{humanos} é o mesmo para todos os homens, considerado \emph{objectivamente}, isto é, em si mesmo; porque, considerado \emph{subjectivamente}, isto é, no juizo dos homens, não é um e o mesmo para todos, mas varia, segundo a diversidade do objecto em que o fazem consistir; porque alguns fazem consistir o fim \emph{ultimo} nas \emph{riquezas}, outros nos \emph{prazeres}, outros na \emph{gloria}, outros no \emph{poder}, etc. (\emph{Sum}. \emph{Th}, \emph{ib}., a. 7).}.

18. O fim ultimo dos actos humanos \textbf{existe} realmente. --- O desejo da felicidade é tão irresistivel, constante e universal, que só póde derivar da natureza, e por isso de Deus, que é o Auctor da natureza. Ora, se não existisse realmente um fim \emph{ultimo}, centro e termo de todas as nossas tendencias, teriamos recebido de \textsc{Deus} um desejo vão e illusorio, principio de inquietações e tormentos; o que, além de impio, é absurdo. Logo o fim \emph{ultimo} dos actos humanos existe realmente \footnote{Todos os philosophos admittiram e admittem a \emph{objectividade} do fim \emph{ultimo} dos actos humanos, porque o desejo da felicidade, por não ser livre mas necessario, deriva da natureza, e por isso tem fundamento na realidade. --- \emph{A}. divergencia das opiniões começa, quando querem determinar o \emph{objecto}, em que consiste o \emph{ultimo} fim.}.

%%>>> art-49

\section*{Artigo II.\\ Natureza do fim ultimo dos actos humanos.}
\addcontentsline{toc}{section}{Artigo II. --- Natureza do fim ultimo dos actos humanos.}

14. O fim ultimo dos actos humanos é o bem universal. --- O fim \emph{ultimo} dos actos humanos é um bem que sacia completamente todas as tendencias e aspirações da vontade; de modo que, na posse d'elle, o homem encontra a sua perfeita \emph{felicidade}. Ora este bem deve ser \emph{universal}, isto: é, deve conter em si tudo o que é bem e excluir tudo o que é mal; aliás a sua posse não poderia formar a \emph{felicidade}, que é um \emph{estado perfeito}, \emph{constituido pela posse de todos os bens}. Logo o fim \emph{ultimo} dos actos humanos é o bem \emph{universal} \footnote{(\emph{Sum}. \emph{Th}., \emph{ib}., q. 2, a. 7), --- Ninguem se admire se aqui repetimos muitas ideas expostas na Secção IV da \emph{Antropologia}, quando tratámos do \emph{destino} do homem. A ordem \emph{logica} das materias exige esta repetição. Porquanto, o \emph{fim} do homem a o fundamento não só da \emph{Immortalidade} da sua alma, mas tambem das \emph{regras} dos seus actos. --- Queira o leitor consultar esse ponto da \emph{Anthropologia}, porque ahi encontrará algumas reflexões e difficuldades, que omittimos neste logar.}.

\pg{560} 15. O fim ultimo dos \textbf{actos humanos} não \textbf{póde ser} um \textbf{bem creado}. --- O fim \emph{ultimo} dos actos humanos é o bem \emph{universal}. Ora todo o bem \emph{creado}, por ser finito ou limitado, é sempre \emph{particular}. Logo o fim \emph{ultimo} dos actos humanos não póde ser um bem \emph{creado} \footnote{(\emph{Sum}. \emph{Th}., \emph{ib}., q. 2, a. 8). --- O fim \emph{ultimo} dos actos humanos é o \emph{objecto}, em que todas as tendencias da nossa alma devem encontrar o seu descanco, a sua \emph{felicidade}. --- O \emph{objecto} da felicidade costuma chamar-se felicidade \emph{objectiva}; e \& \emph{posse} d'este objecto dá-se o nome de felicidade \emph{subjectiva}.\par Os philosophos pagãos, não tendo o beneficio da Revelação, pensaram que a felicidade se encontrasse sobre a terra, e por isso a fizeram consistir na posse dos bens creados. E, como os bens \emph{creados} --- ou são bens \emph{extrinsecos} (de \emph{fortuna}), ou são bens \emph{intrinsecos} (do \emph{corpo} ou da \emph{alma}), alguns constituiram o objecto da felicidade nos bens \emph{eowtrinsecos} (\emph{riquezas}, \emph{honras}; \emph{poder}), --- outros, nos bens do \emph{corpo}, --- outros nos da \emph{alma}, --- outros, na \emph{reunião} dos bens da \emph{alma} e do \emph{corpo}. --- Os modernos incredulos, que desprezam o beneficio da Revelação, imitam os escriptores pagãos. Assim, os \emph{materialistas}, os \emph{sensualistas}, os \emph{positivistas} dizem que o \emph{praser corporeo} é o bem supremo, o fim \emph{ultimo} do homem; --- os \emph{Arausistas} e, em geral, os \emph{pantheistas} fazem consistir o objecto da felicidade humana na evolução completa de todas as nossas faculdades; --- os \emph{Aantistas}, no \emph{progresso constante} e \emph{indefinido para a felicidade perfeita e para a santidade infinita}; outros, como Wundt, \emph{Schleiermacher}, na perfeição do genero humano.\par Refutaremos estas theorias em outras tantas proposições.\par a) O fim \emph{ultimo} do homem ou o \emph{objecto} da nossa felicidade não consiste nos bens \emph{extrinsecos} ou de \emph{fortuna}. Na verdade, os bens de \emph{fortuna} são \emph{insufficientes}, não tornam completamente feliz o sujeito que os possue, porque não encerram todos os bens e não excluem todos os males. Ora o fim \emph{ultimo} deve saciar todos os desejos da alma. --- Os bens de \emph{fortuna} são o apanagio de poucos e dependem de circumstancias fortuitas, caprichosas ou alheias á vontade do homem. Ora o fim \emph{ultimo} deve ser accessivel a todos (porque todos nascerám para a felicidade): é independente das circumstancias fortuitas, caprichosas e alheias á nossa vontade; porque, tendo sido determinado pela natureza como alvo de todas as nossas tendencias, a sua consecução deve estar ao alcance de cada um. --- Os bens de \emph{fortuna} perdem-se, e por isso a sua posse está unida com o medo e a incerteza. Ora o \emph{objecto} da felicidade deve ser \emph{inamissivel}; aliás não poderia tornar completamente feliz quem o possue (\emph{Sum}. \emph{Th}., \emph{ib}., q. 2, a. 1).\par b) O fim \emph{ultimo} do homem não consiste nos bens do \emph{corpo}. Porquanto, os bens do \emph{corpo} são \emph{particulares}, porque não encerram os bens da \emph{alma}, --- não dependem da nossa vontade, --- são a origem de muitos males e inquietações, --- são \emph{instaveis}, pois podem perder-se e perdem-se com a maxima facilidade. --- Estes principios applicam-se a cada um dos bens do \emph{corpo}, mas sobretudo aos prazeres \emph{corporeos}, que são communs tambem aos animaes \emph{irracionaes}, e que por isso não são o fim \emph{ultimo} do homem, que é \emph{racional}; pois o fim de cada ente deve ser proporcionado á sua natureza (\emph{Sum}. \emph{Th}., \emph{ib}., a. 5, 6).\par c) O fim \emph{ultimo} do homem não consiste nos bens da \emph{alma}. Os bens da \emph{alma} são principalmente a \emph{sciencia} ea \emph{virtude}. Ora nem a \emph{sciencia}, nem á\pgn{561}\emph{virtude} é o fim ultimo do homem. --- Não a \emph{sciencia}; porque o fim \emph{ultimo} deve ser \emph{accessivel} a todos e saciar completamente as tendencias da alma, ao passo que a \emph{sciencia} é dote de poucas pessoas, e, longe de saciar os desejos da alma, torna-os cada vez mais vivos. --- Não póde ser a \emph{virtude}; visto que a \emph{virtude é meio} para a consecução da felicidade, mas não é a propria felicidade, pois não socega completamente a vontade, podendo perder-se e não excluindo todos os males e inquietações (\emph{Sum}. \emph{Th}, \emph{ib}., a. 7).\par 4) O fim \emph{ultimo} do homem não consiste no conjuncto de todos os bens \emph{creados}. Porquanto, embora o homem podesse possuir todos os bens creados --- de \emph{fortuna}, de \emph{corpo} e de \emph{alma}, comtudo não seria completamente feliz. Na verdade, o conjuncto de todos os bens creados é sempre uma coisa finita e sujeita a mudanças, nem exclue todos os males.\par €) O fim \emph{ultimo} do homem não consiste na evolução \emph{completa} de \emph{todas} as suas faculdades. Na verdade:\par \emph{a}) Esta opinião dos \emph{Zrausistas} baséa-se em que o fim \emph{ultimo} do homem é o proprio homem. Ora repugna que o homem seja o fim \emph{ultimo} de si mesmo; porque o fim \emph{ultimo} deve conter todos os bens e excluir todos os males, e a experiencia attesta que o homem não contém todo o bem e que está sujeito a muitos males.\par \emph{b}) A evolução \emph{completa} de \emph{todas} as nossas faculdades é impossivel. Porquanto, sendo a energia da alma uma coisa limitada, uma faculdade não póde desenvolver-se sem prejuizo das outras. E' o que a experiencia mostra. --- Os \emph{hrausistas} dizem que a evolução é harmonica. Mas respondemos que a evolução, se é harmonica, não é \emph{completa} em nenhuma faculdade.\par \emph{c}) A evolução \emph{completa} de \emph{todas} as nossas faculdades, embora fosse possivel, comtudo não seria o fim \emph{ultimo} do homem. Na verdade, o fim \emph{ultimo} deve saciar completamente as tendencias do homem. Ora a evolução \emph{completa} de \emph{todas} as nossas faculdades não póde socegar completamente as nossas tendencias, porque é sempre uma coisa finita e pouco duradoura.\par \emph{d}) Se o fim \emph{ultimo} do homem consistisse na evolução \emph{completa} de \emph{todas} as suas faculdades, teriamos obrigação de estudar, não só a Philosophia, mas tambem a \emph{gymnastica}, \emph{á pintura}, \emph{á musica}, e todas as artes liberaes; o que \emph{Ahrens} explicitamente affirma (\emph{Cours de Droit Naturel}, part. génér., ch. I, 8 17). Ora o senso commum nunca se lembrou de collocar, entre os deveres do homem, o estudo das sciencias e das artes liberaes, porque tal estudo é impossivel para a maxima parte dos individuos, e porque essas sciencias e artes não são necessarias para uma vida boa e feliz.\par f) O fim \emph{ultimo} do homem não consiste no progresso constante e indefinido do individuo para a perfeição. Com effeito, o fim \emph{ultimo é repouso} do espirito; ao passo que o progresso indefinido é \emph{movimento} e por isso \emph{agitação}. --- O fim \emph{ultimo} deve ser um bem \emph{infinito}, porque só este póde saciar o\par 36 Vor. II.\par ”\pgn{562}espirito; ao passo que o progresso indefinido é sempre /\emph{inito}, porque toda a serie dos movimentos é essencialmente finita.\par g) O fim \emph{ultimo} do homem não consiste na perfeição do genero humano. Na verdade, o fim \emph{ultimo} é proprio de \emph{cada individuo}, de modo que a sua posse deve tornar feliz cada um dos homens; ora a perfeição do genero humano é propria da \emph{collecção} dos homens, e não póde tornar feliz cada um dos individuos. --- O fim \emph{ultimo} deve ser. um bem \emph{infinito}; ora a perfeição do genero humano é sempre finita. --- Notamos que esta opinião está baseada tambem no \emph{pantheismo}, cuja falsidade foi demonstrada.}.

16. O fim ultimo dos \textbf{actos humanos} é \textbf{Deus}. --- O fim \emph{ultimo} dos actos humanos, sendo o objecto da felicidade, deve ser proporcionado á natural tendencia da vontade para o bem, e a natural tendencia da vontade é proporcionada. ao conhecimento da intelligencia; porque o conhecimento é a regra e a medida do desejo. Ora a nossa intelligencia conhece o bem \emph{universal}, ou \emph{infinito}, que encerra todos os bens. Logo a vontade \pg{561} \pg{562}  tende para um bem \emph{universal}. O bem \emph{universal} é \textsc{Deus}. Logo o fim \emph{ultimo} dos actos humanos é \textsc{Deus} \footnote{Diz S. Thomaz: « Impossibile est beatitudinem hominis esse in aliquo bono creato. Beatitudo enim est bonum perfectum, quod totaliter quietat appetitum; alioquin non esset finis ultimus, si adhuc restaret aliquid appetendum. Objectum autem voluntatis, quae est appetitus humanus, est universale bonum, sieut objectum intellectus est universale verum. Ex quo patet, quod nihil potest quietare voluntatem hominis nisi bonum universale; quod non invenitur in aliquo creato, sed solum in Deo; quia omnis creatura habet bonitatem participatam. Unde solus Deus voluntatem hominis implere potest... In solo igitur Deo beatitudo hominis consistit » (1-2, q. 2, a. 8). . --- Os proprios philosophos pagãos, como \emph{Pythagoras}, \emph{Platão}, \emph{Aristoteles}, guiados pela luz da razão, chegaram a conhecer que \textsc{Deus}, e só Elle, é o fim \emph{ultimo} do homem, o objecto da nossa felicidade. --- E nada mais logico. O fim \emph{ultimo} deve consistir num bem que. se deseje por si e que por isso não se refira \emph{a} um bem mais perfeito e elevado; aliás não seria fim \emph{ultimo}. Ora só \textsc{Deus} é o bem que se deseja por si e que se não refere a um bem mais elevado e perfeito; porque todos os bens creados, sendo participações do bem infinito, referem-se ao proprio \textsc{DeEUs}.\par Objectam: A vontade humana, sendo uma faculdade \emph{finita}, póde e deve saciar-se na posse de um bem \emph{finito}.\par Resposta. A vontade humana, se é \emph{finita} na sua \emph{entitade subjectiva}, emquanto é uma faculdade propria de um ente finito, é \emph{infinita} na sua \emph{capacidade objectiva}, emquanto póde desejar e deseja um bem infinito.\par Continuam: Se o objecto da vontade humana fosse o bem \emph{infinito}, não haveria differenca entre o fim do Creador e o da creatura; porque o mesmo\pgn{563}bem infinito é tambem objecto da vontade de \textsc{DEus}. Ora é absurdo que o fim da creatura se identifique com o fim do Creador.\par Resposta. Entre a felicidade de \textsc{Deus} e a do homem existe sempre uma distancia infinita; porque o objecto da nossa vontade a o bem infinito, \emph{possuido por um modo finito}, ao passo que o objecto da vontade divina é o bem infinito, \emph{mas possuido por um modo infinito} (\emph{Sum}. \emph{Th}., \emph{ib}., q. 8, a. 1).\par Outras diffieuldades foram resolvidas na sec. IV da \emph{Anthr}. (Cap. II, art. 1).}.

17. O fim ultimo dos \textbf{actos humanos consegue-se pela} intelligencia e \textbf{pela vontade}.

a) O fim \emph{ultimo} só póde conseguir-se por aquellas faculdades que lhe são semelhantes, visto que deve existir uma proporção entre o fim e os meios. Ora \textsc{Deus}, que é o fim \emph{ultimo} dos actos humanos, é purissimo espirito, e por isso só póde conseguir-se pelas nossas faculdades \emph{espirituaes}. As faculdades \emph{espirituaes} são a \emph{intelligencia} e a \emph{vontade}. Logo o fim \emph{ultimo} dos actos \emph{humanos} consegue-se pela \emph{intelligencia} e pela \emph{vontade}.

b) O fim de um ente consegue-se sempre pelas faculdades do mesmo ente, que tendem para elle. Ora o fim \emph{ultimo} dos actos \emph{humanos} é o centro, para o qual tendem a nossa intelligencia \pg{563}  e a nossa vontade; porque o acto \emph{humano} deriva da intelligencia e da vontade. Logo o fim \emph{ultimo} dos actos humanos consegue-se pela \emph{intelligencia} e pela \emph{vontade} \footnote{A felicidade do homem, pois, consiste em \textsc{DEus} conhecido e amado. A intelligencia deseja a verdade, e \textsc{Deus} é a propria Verdade; a vontade deseja o bem, e \textsc{Deus} é o Bem infinito.\par Os philosophos e os theologos discutem ácerca do acto, pelo qual se alcança a felicidade. Uns dizem que é o acto da \emph{intelligencia}, que conhece a \textsc{Deus}; outros affirmam que é o acto da \emph{vontade}, que ama a \textsc{Deus}; outros sustentam que é o acto de ambas essas faculdades.\par Para evitar equivocos, devemos fazer uma distinceção entre o \emph{estado} da felicidade e a sua \emph{essencia}. O \emph{estado} da felicidade abrange tudo o que torna o homem bemaventurado, e por isso abrange o \emph{conhecimento} de \textsc{Deus}, o seu \emph{amor}, o \emph{goso} que deriva d'este conhecimento e d'este amor, a \emph{perpetuidade}, a \emph{immutabilidade}, etc. A \emph{essencia} da felicidade é o acto, por que se chega á posse de \textsc{Deus} e que por isso é o principio de todos os outros actos, que precedem ou acompanham a propria felicidade. O ponto da questão é a \emph{essencia} da felicidade, e por isso quer saber-se qual é o \emph{acto} por que alcancamos o nosso fim \emph{ultimo}, que é \textsc{Deus}. --- Dadas estas explicações, procuraremos resolver brevemente a questão.\par A \emph{essencia} da felicidade não consiste no acto da \emph{vontade}, mas no acto da \emph{intelligencia}, que conhece perfeitamente a \textsc{Deus}. --- Não consiste no acto da \emph{vontade}. Porquanto, o acto da \emph{vontade} não só é posterior ao acto da \emph{intelligencia} (porque o amor e o gozo seguem o conhecimento), mas, longe de \emph{apprehender} o objecto, \emph{tende} para elle. --- Mas consiste no acto da \emph{intelligencia}. Na verdade, a felicidade consiste na \emph{consecução} ou \emph{posse} de \textsc{Deus}. Ora \textsc{Deus} \emph{consegue-se} ou \emph{possue-se} pela \emph{imtelligencia}, porque é proprio d'esta faculdade \emph{apprehender} ou \emph{attrahir} a si o objecto, e tornal-o presente ao espirito.\par Note-se que o perfeito \emph{conhecimento} de \textsc{Deus}, em que consiste a \emph{essencia} da felicidade, póde ser \emph{directo} ou \emph{indirecto}. --- X? \emph{directo}, quando o \emph{meio}, por que conhecemos a \textsc{Deus}, é a sua Essencia infinita, informante a nossa inteligencia; por isso este conhecimento chama-se \emph{intuitivo}, porque vê-se o nosso \textsc{Deus}, como Elle é em si mesmo. --- E' \emph{indirecto}, quando o \emph{meio}, por que \textsc{Deus} é conhecido, é uma especie analogica ou semelhança ideal, que representa os attributos divinos, como estes se manifestam na creação, e que é mais ou menos perfeita, conforme a diversidade dos merecimentos, --- O conhecimento \emph{directo} de \textsc{DEus} excede as forças naturaes da nossa intelligencia, \emph{e} por isso é \emph{sobrenatural}; ao passo que o \emph{indirecto} não excede as forças naturaes da intelligencia, e por isso é \emph{natural}.\par Embora a Bondade de \textsc{Deus} nos tenha destinado ao conhecimento \emph{sobrenatural} da sua essencia infinita, e por isso a uma felicidade \emph{sobrenatural}; comtudo o conhecimento \emph{natural} da divina Essencia, por meio de uma especie\pgn{564}intelligivel, teria sido bastante para a nossa felicidade \emph{natural}, isto é, para a felicidade, que é exigida pela nossa \emph{natureza}. :\par Poderia dizer-se que o homem não teria encontrado perfeita felicidade na visão \emph{indirecta} de \textsc{Deus}, visto que elle podia ainda desejar a visão \emph{directa} de \textsc{DEus}. Mas respondemos que o homem, no estado \emph{natural}, --- ou não teria suspeitado da possibilidade de uma visão \emph{directa} de \textsc{DEUS}, --- ou, se a tivesse conhecido, ter-se-hia convencido de que ella não era devida á sua natureza nem proporcionada ás suas forças, e assim não a teria desejado com desejo efficaz, nem a sua ausencia teria diminuido a felicidade \emph{natural}.\par Ao perfeito conhecimento de \textsc{Deus} corresponde um acto de amor perfeito, e este amor, unindo a nossa vontade com \textsc{Deus}, produz na alma uma consolação tão pura e constante, que nada pode inquietar ou destruir. Contemplar a \textsc{Deus}, amar a \textsc{DEus}, gosar a \textsc{Deus}: eis o fim \emph{ultimo} do homem. (Cf. \emph{Sum}. \emph{Th}., \emph{ib}., q. 9, aa. 2, 8, 4).}.

\pg{564} 18. \emph{A} \textbf{perfeita consecução} do fim ultimo dos \textbf{actos} humanos \textbf{está reservada para} a \textbf{vida futura}.

a) A \emph{perfeita} consecução do fim \emph{ultimo} dos actos humanos deve tornar a nossa alma completamente feliz. Ora a experiencia attesta que a nossa alma, na vida presente, não é completamente feliz, antes é repleta de muitas miserias e inquietações. Logo a \emph{perfeita} consecução do fim \emph{ultimo} dos actos humanos está reservada. para a vida \emph{futura}.

-b) A \emph{perfeita} consecução do fim \emph{ultimo} dos actos humanos consiste no \emph{perfeito} conhecimento e por isso no \emph{perfeito} amor de \textsc{Deus}. Ora o homem, sobre esta terra, não póde ter um \emph{perfeito} conhecimento nem um \emph{perfeito} amor de \textsc{Deus}; porque a especie \emph{analogica}, que a intelligencia, considerando as creaturas, forma de \textsc{Deus}, não leva a comprehensão de toda a ordem creada em si e nas suas relações para com a Creador, e porque o amor, que se funda nesse conhecimento, é fraco, mutavel, destruetivel. Logo a \emph{perfeita} consecução do fim \emph{ultimo} dos actos humanos está reservada para a vida \emph{futura} \footnote{\emph{Sum}. \emph{Th}., \emph{ib}., q. 5, a. 8). --- A idea \emph{analogica}, por meio da qual o homem, no estado \emph{natural}, conhece a \textsc{Deus}, representa, por um modo \emph{finito} mas relativamente perfeito, os attributos do proprio \textsc{DEus}, que se manifestam na \emph{creagão}. 'l'al conhecimento é o \emph{ideal} da \emph{sciencia natural}, porque encerra ua unidade perfeita de todos os principios a grande variedade de todas as suas consequencias e applicações. Na verdade, o perfeito conhecimento do Ser divino importa o conhecimento de toda a ordem creada, porque \textsc{DEus} é a causa \emph{efficiente}, \emph{exemplar} e \emph{final} de todas as coisas. Esta idea, que basta para a nossa felicidade \emph{natural}, é infusa por \textsc{Deus} na intelligencia humana e é mais ou menos elevada, conforme a diversidade dos merecimentos de cada um dos homens.\par Repetimos que tudo o que dissemos ácerca do estado \emph{natural} do homem é \emph{hypothetico}; porque esse estado \emph{natural}, se é \emph{possivel}, não existe na \emph{realidade}\pgn{565}, pois o homem foi elevado ao estado \emph{sobrenatural}, constituido \emph{essencialmente} pela visão \emph{directa} ou \emph{intuitiva} de \textsc{DEus}.\par Objectam: A philosophia só considera a ordem \emph{natural}, Logo não deve oceupar-se da ordem \emph{sobrenatural}.\par Resposta. A philosophia deve occupar-se do homem \emph{real}, o não do homem \emph{hypothetico}. Ora o homem \emph{real} foi elevado \& ordem \emph{sobrenatural}; e por isso a philosophia não póde prescindir inteiramente da ordem \emph{sobrenatural}. --- Uma coisa é que a philosophia deve fundar as suas demonstrações na \emph{razão natural} e não na \emph{revelação divina}, o que concedemos; outra coisa é que a philosophia prescinda inteiramente da \emph{revelação divina}, o que negamos.}.

%%>>> art-50

\section*{Artigo III.\\ Fim proximo dos actos humanos.}
\addcontentsline{toc}{section}{Artigo III. --- Fim proximo dos actos humanos.}

\pg{565} 19. O fim proximo dos actos humanos existe realmente. --- O fim \emph{proximo} dos actos humanos é o bem, que deve ser alcançado pelo homem nesta vida e que lhe deve servir de \emph{meio} para a consecução do fim \emph{ultimo}. Ora, se este fim \emph{proximo} não existisse realmente, deveria dizer-se que a vida do homem sobre a terra não tem razão de ser, e que o fim \emph{ultimo} póde alcançar-se sem o emprego dos \emph{meios} convenientes: o que é absurdo. Logo o fim \emph{proximo} dos actos humanos existe realmente \footnote{\textsc{DEus}, sendo a propria Sabedoria, não podia deixar de fixar um escopo ou fim \emph{provimo}, que o homem devesse conseguir aqui sobre a terra. --- Nem se diga que o fim \emph{ultimo} substitue o \emph{proximo}; porque não póde imaginar-se um fim \emph{ultimo} sem um fim \emph{intermedio} ou \emph{proximo}. --- De resto, todos os philosophos estão de accordo neste ponto. As suas divergencias comecam, quando explicam a \emph{natureza} d'esse fim \emph{proximo}, a qual não póde deixar de ser diversa, conforme a diversidade do fim \emph{ultimo}; porque o fim \emph{provimo é ele} em relação ao fim \emph{ultimo}, e o \emph{meio} deve ser proporcionado ao \emph{fim}.}.

20. O fim proximo dos \textbf{actos} humanos consiste na \textbf{sua disposição para} o fim ultimo.

a) O homem tem o dever de conseguir o seu fim \emph{ultimo}. Ora este dever deve cumprir-se na vida presente; porque a vida futura está reservada, não para a \emph{conquista}, mas para a \emph{posse} da felicidade. Logo o fim \emph{proximo} dos actos humanos consiste na sua disposição para o fim \emph{ultimo}.

b) O fim \emph{proximo} dos actos humanos é um \emph{meio} em relação ao seu fim \emph{ultimo}; aliás o homem não teria unidade de \emph{natureza} nem de \emph{fim}. Ora toda a perfeição do \emph{meio} consiste na aptidão que tem de levar á consecução do \emph{fim}. Logo o fim \emph{proximo} dos actos humanos consiste na sua disposição para o fim \emph{ultimo} \footnote{A vida presente e a vida futura constituem a duração da existencia\pgn{566}da nossa alma. Ora uma só existencia exige um só fim. Esta unidade de \emph{fim} só póde conceber-se, se o fim \emph{prowimo} dos nossos actos estiver disposto e ordenado para o fim \emph{ultimo}. Porquanto, se não existisse esta subordinação, o fim \emph{provimo} dos nossos actos seria tambem \emph{ultimo} (\emph{porque} não seria subordinado a outro superior), e assim teriamos dois fins \emph{ultimos} do homem. Ora repugna que existam dois fins \emph{ultimos}. Na verdade, o fim \emph{ultimo} ê o bem absolutamente perfeito, que sacia todas as tendencias da alma. Mas, se existissem dois fins \emph{ultimos}, nenhum d'elles seria o bem absolutamente perfeito; visto que nenhum d'elles saciaria todas as tendencias da alma, pois esta, depois de estar na posse de um bem, podia desejar tambem o outro. Logo não podem existir dois fins \emph{ultimos}.}. :

\pg{566} 21. O fim proximo dos actos humanos consiste no conhecimento e no \textbf{amor} de \textbf{Deus}.

a) O fim \emph{proximo} dos actos humanos consiste na sua disposição para o fim \emph{ultimo}. Ora o fim \emph{ultimo} consiste no perfeito conhecimento e amor de \textsc{Deus}. Logo o fim \emph{proximo} consiste na disposição para o perfeito conhecimento e amor de \textsc{Deus}. Mas semelhante disposição exige que o homem comece a fazer na vida presente o que mais perfeitamente fará na vida futura; porque o fim \emph{proximo} 6 O \emph{começo} do fim \emph{ultimo}. Logo o fim \emph{proximo} dos actos humanos consiste no conhecimento e no amor de \textsc{Deus}.

b) O fim, que \textsc{Deus} determinou ao homem, não só para a vida futura, mas tambem para a presente, foi, e não podia deixar de ser, a sua gloria. Ora o homem procura a gloria de \textsc{Deus}, emquanto o conhece e o ama. Logo o fim \emph{proximo} dos actos humanos consiste no conhecimento e no amor de \textsc{Deus} \footnote{O homem póde dispor a sua vida, em relação ao fim \emph{ultimo}, por dois modos contrarios --- ou amando e respeitando a \textsc{DEus} pelo cumprimento dos seus preceitos, --- ou desprezando a \textsc{Deus} pela violação da sua lei. Ora é absurdo pensar que uma vida boa leve para uma desgraça eterna, e que uma vida criminosa seja preparação para a felicidade eterna. --- E' claro que todos os que constituem o nosso fim \emph{ultimo} em algum bem creado, não estão de accordo com a nossa doutrina ácerca do fim \emph{prowvimo} dos actos humanos.}.

29. A consecução do fim proximo dos \textbf{actos} humanos produz uma felicidade imperfeita. --- A felicidade \emph{imperfeita}, possivel nesta vida, é uma participação e um começo da felicidade \emph{perfeita}, que nos está reservada na vida futura; porque, embora não contenha todos os bens e não exelua todos os males, todavia é luz para a intelligencia e consolação para a vontade. Ora esta felicidade \emph{imperfeita} é o resultado da consecução do fim \emph{proximo} dos actos humanos; porque a intelligencia encontra a sua mais nobre satisfação no conhecimento de \textsc{Deus}, e a vontade encontra a sua mais intima alegria no amor da \textsc{BonDADE} \pg{567}  \textsc{INFINITA}. Logo a consecução do, fim \emph{proximo} dos actos humanos produz uma felicidade \emph{imperfeita} \footnote{A \emph{esperança} de possuir um bem infinito e por isso de alcançar uma felicidade perfeita é já um começo da propria felicidade, é uma felicidade, embora imperfeita. --- E' a historia de todos os corações. Os que amam a \textsc{Deus}, passam uma vida tranquilla, socegada e serena, como o céo sem nuvens. Os que desprezam a \textsc{DEus}, não observando os preceitos que Elle nos impoz, vivem sempre agitados, lacerados pelos remorsos, e, embora possuam riquezas, glorias, prazeres, não encontram um momento de allivio. A razão d'isto foi já dada por \emph{S}. \emph{Agostinho} naquellas palavras: u \emph{Senhor}, \emph{creastes-nos para} Vós, \emph{e o nosso coração estará sempre inquieto emquanto não descancar em} Vós n (Cf. \emph{Sum}. \emph{TA}., 1-2, q. 5, a. 4).\par Objectam: O homem é composto de alma e de corpo. Logo para a sua felicidade, aqui sobre a terra, é necessario o repouso --- não só das suas faculdades \emph{intellectuaes}, pelo conhecimento e amor de \textsc{Deus}, --- mas tambem das suas faculdades \emph{sensitivas}, pelo gozo dos bens sensiveis.\par Resposta. A felicidade, que o homem póde obter neste mundo, é a \emph{imperfeita}. Tal felicidade deve consistir na consecução de um bem, que seja \emph{summo} relativamente ao estado presente e que seja \emph{proporcionado} á nossa natureza. O bem \emph{summo} e \emph{proporcionado} a cada natureza é o que corresponde á mais elevada das suas tendencias. Ora a mais elevada tendencia do homem é a \emph{racional}, isto é, a que deriva da livre vontade com o previo conhecimento da intelligencia. As outras tendencias inferiores devem estar subordinadas a esta suprema; de modo que devem ser contrariadas ou podem ser satisfeitas, conforme a sua satisfação se oppõe, ou não, aquella tendencia \emph{racional} para \textsc{DEus}. Nem admira se, na vida presente, não ficam satisfeitas \emph{todas} as tendencias do homem; porque a felicidade, que é possivel sobre a terra, é a \emph{imperfeita}, e por isso póde estar acompanhada de dores e contrariedades..--- Não negamos que a posse dos bens creados da alma e do corpo concorre, quando não é contraria á ordem moral, para o bem estar do homem sobre a terra; porém este concurso não \emph{constitue} a felicidade do homem, a qual consiste no conhecimento e no amor \textsc{Deus}, mas é um \emph{meio} para o homem se elevar a esse conhecimento e amor (\emph{Sum}. \emph{Th}., 1-2, q. 4, a. 1).\par Continuam: A felicidade exclue todo o medo da perda do bem. Ora o homem, aqui sobre a terra, póde perder a sua felicidade. Logo n\&o ha verdadeira felicidade do homem sobre a terra.\par Resposta. A proposição \emph{maior} é verdadeira, se se refere á felicidade \emph{perfeita}; mas é falsa, se se refere á felicidade \emph{imperfeita}. --- E' \emph{verdadeira}, se se refere \& felicidade \emph{perfeita}; por quanto a alma, tendo conseguido o seu \emph{ultimo} fim, conhecerá a \textsc{Deus} por um modo muito elevado e distincto, e, vendo-o tão bom e tão perfeito, não poderá deixar de o amar. --- E' \emph{falsa}, se se refere á felicidade \emph{imperfeita}; porque, no estado actual, a nossa intelli- ' gencia conhece a \textsc{Deus}, mas muito imperfeitamente, e á nossa vontade encontra outros bens, cuja posse é menos ardua que a de \textsc{Deus}, e por isso o homem póde deixar e muitas vezes deixa de amar a \textsc{DEus}.}.

23. \textbf{Corollarios}. --- Das doutrinas expostas resultam os seguintes corollarios: :

I. O homem póde ser relativamente feliz no meio das dores e contrariedades da vida; porque póde servir e amar a \textsc{Deus}.

\pg{568} II. Os males \emph{physicos} devem evitar-se ou procurar-se, conforme nos afastam do amor de \textsc{Deus}, ou nos levam a elle; pois estes-males, na ordem \emph{moral}, são por si coisas 2\emph{ndifferentes}.

III. O homem deve preferir todos os males physicos e a propria morte á offensa de \textsc{Deus}; porque esses males e a morte não são contrarios á ordem \emph{moral}, como lhe é contrario o peceado \footnote{Os \emph{incredulos}, os \emph{positivistas} riem-se da ingenuidade com que escrevemos estes corollarios, alias muito conformes com o o \emph{senso communm}. Dormem, infelizmente, nas desordens das paixões e não pensam que o seu despertar pode ser o principio de um desengano irremediavelmente funesto!}.

%%>>> cap-8

\chapter*{CAPITULO SEGUNDO\\[0.3em]\large ESSENCIA DOS ACTOS HUMANOS}
\addcontentsline{toc}{chapter}{CAPITULO SEGUNDO --- ESSENCIA DOS ACTOS HUMANOS}

SuwwaRr0: --- Elementos e divisio do acto humano. --- Diversas especies do

acto elicito. --- Impedimentos da voluntariedade dos actos humanos.

%%>>> art-51

\section*{Artigo I.\\ Elementos e divisão dos actos humanos.}
\addcontentsline{toc}{section}{Artigo I. --- Elementos e divisão dos actos humanos.}

24. Elementos do acto humano. --- Os elementos, que constituem o acto \emph{humano}, emquanto este é objecto da \emph{Moral}, são \emph{ires}, a saber: a \emph{percepção} da \emph{intelligencia}, o \emph{consentimento} da \emph{vontade} e a \emph{liberdade}. Com effeito, o nosso acto é e diz-se \emph{Humano}, quando deriva da \emph{intelligencia} e da \emph{livre vontade} \footnote{Os actos, que o homem exerce, dividem-se em duas classes, emquanto --- ou são \emph{communs} tambem aos animaes irracionaes (taes são os actos que derivam da \emph{sensibilidade}), --- ou são exclusivamente \emph{proprios} do homem (taes são os actos que derivam da ;\emph{intelligencia} e da \emph{vontade}). Os da primeira classe chamam-se actos \emph{do tomem}; os da segunda, actos \emph{Awmeanos}. --- Os actos \emph{humanos} excedem os \emph{do homem} --- ou só emquanto á \emph{substancia}, como os actos : \emph{necessarios} da \emph{intelligencia} (relativamente aos \emph{primeiros principios}) e da \emph{vontade} (relativamente á \emph{felicidade}), --- ou só emquanto ao \emph{modo}, como os actos de \emph{comer}, \emph{beber}, etc., que se fazem com deliberação, --- ou emquanto á \emph{substancia} e ao \emph{modo}, como são os actos \emph{livres} da \emph{intelligencia} e da \emph{vontade}.\par Os actos que são \emph{humanos} só emquanto á \emph{substancia}, por não estarem sujeitos ao livre arbitrio, não podem -ser regulados pela \emph{Moral}. Por isso os actos \emph{humanos}, que constituem o objecto d'esta sciencia, são os que são proprios do homem emquanto ao \emph{modo} de operar, quer excedam, ou não, emquanto á \emph{substancia}, os actos dos seres irracionaes,}.

\pg{569} 25. O \textbf{acto} humano \textbf{exige} a percepção da intelligencia.

a) O acto \emph{humano} é o que deriva do homem, emquanto \emph{tal}, isto é, emquanto se distingue de todos os outros seres sensiveis. Ora o homem distingue-se de todos os outros seres sensiveis pela \emph{intelligencia}, de que é dotado. Logo o acto \emph{humano} exige a \emph{percepção} da \emph{intelligencia}.

b) O acto \emph{Humano}, como demonstraremos d'aqui a pouco, exige o \emph{consentimento} da \emph{vontade}. Ora a \emph{vontade} não póde dar O seu \emph{consentimento} a uma coisa, que se não conhece; porque o desconhecido não póde ser alvo do desejo. Logo o acto \emph{humano} exige a \emph{percepção} da \emph{intelligencia} \footnote{Na \emph{Anthropologia} (Sec. IT, cap. III art. III) tratámos, com o devido desenvolvimento, das \emph{relações entre a intelligencia e a vontade}.}.

26. O \textbf{acto} humano exige o consentimento da vontade.

a) O acto \emph{humano} é o que deriva do homem, emquanto \emph{tal}, isto é, emquanto se distingue de todos os outros seres sensiveis. Ora o homem distingue-se de todos os outros seres sensiveis pela \emph{vontade}, de que é dotado. Logo o acto \emph{humano} exige o \emph{consentimento} da \emph{vontade}.

b) O acto \emph{humano}, tendendo para o \emph{bem}, apontado pela razão, deve proceder da faculdade, que tem por objecto o \emph{bem}. Ora a faculdade, que tem por objecto o \emph{bem}, é a \emph{vontade}. Logo o acto \emph{humano} exige o \emph{consentimento} da \emph{vontade}, ou, com mais exactidão, é o proprio \emph{consentimento} da \emph{vontade}. --- Por isso o acto \emph{humano} costuma chamar-se tambem \emph{voluntário} \footnote{Por isso á causa \emph{immediata} dos actos humanos é constituida pela \emph{intelligencia} e pela \emph{vontade}. Mas ha differenca; porque o acto humano deriva \emph{substancialmente} da \emph{vontade}, e.só \emph{radicalmente} da \emph{intelligencia}; porque é a vontade que tende para o bem, embora esta faculdade seja movida pela intelligencia, que \emph{o} aponta. --- Dissemos que a intelligencia e a vontade são a causa \emph{inmediata}, ou \emph{prosima}, do acto humano; porque o \emph{appetite sensitivo}, podendo actuar na intelligencia e na vontade, póde ser causa \emph{mediata}, ou \emph{remota}, do mesmo acto (\emph{Sum}. \emph{Th}., 1-2, q. 75, a. 2).}.

27. O \textbf{acto} humano exige a liberdade. --- Acto \emph{humano}, emquanto forma o objecto d'esta parte da Philosophia, é o que, pelas regras da \emph{moralidade}, deve ser dirigido para o nosso \emph{ultimo} fim. Ora esta direcção seria impossivel, se o homem não fosse dotado de \emph{liberdade}, pela qual um ser é senhor e arbitro dos seus actos. Logo o acto \emph{humano} exige a \emph{liberdade}. --- A \emph{liberdade} É a razão da \emph{imputabilidade} dos nossos actos \footnote{Para a \emph{essencia} do acto humano, que é o objecto da \emph{Moral}, não basiam a percepção da intelligencia e o consentimento da vontade, mas é necessaria tambem a \emph{liberdade}, isto é, requer-se que a vontade seja \emph{livre}. Sem\pgn{570}esta liberdade, é impossivel estabelecer normas e marcar deveres. Quando a vontade, consultando a razão, escolhe uma coisa e rejeita outra, então produz um acto /\emph{lumano}. --- A nossa vontade, como vimos, é livre na sua tendencia para os bens \emph{particulares}, mas não o é na sua tendencia para o bem \emph{universal}, para a felicidade. D'ahi a divisão do \emph{voluntario} em \emph{livre} e \emph{necessario}. Por isso todo o acto \emph{humano é voluntario}, mas nem todo o acto \emph{voluntario é humano}; pois o acto \emph{humano} deve proceder da \emph{vontade livre}.\par Tambem não nos demoraremos em demonstrar a \emph{existencia do livre arbitrio} no homem; porque esta demonstração foi feita na \emph{Anthropologia}.}.

\pg{570} \emph{Poa} \textbf{ao NUES ORT} \emph{Ru} RUE asa ini

570 MORAL --- PARTE PRIMEIRA --- ETHICA

. 28. \textbf{Divisão do acto humano}. --- O acto \emph{humano} divide-se principalmente --- a) em \emph{elicito} e \emph{imperado}, --- \emph{b}) \emph{perfeito} e \emph{imperfeito}, --- c) \emph{directo} e \emph{indirecto}, --- \emph{d}) \emph{positivo} e \emph{negativo}, --- e) \emph{actual} e \emph{virtual}.

a) \textbf{Elicito} e imperado. --- .\emph{Elicito} é o que emana \emph{immediatamente} da propria vontade; tal é o acto de \emph{amar}, de \emph{odiar}, ete. --- \emph{Imperado} é o que deriva das outras faculdades, sob o imperio ou moção da vontade, tal é o acto de \emph{fallar}, \emph{passear}, ete. \footnote{A vontade póde considerar-se: ou como \emph{faculdade especial}, ou como \emph{universal motor} das forças da alma. Considerada como \emph{faculdade especial}, produz os actos \emph{elicitos}; --- considerada como \emph{universal motor}, produz os \emph{actos imperados}. --- O acto da \emph{vontade} é sempre \emph{elicito}; ainda que, num certo sentido, possa ás vezes ser \emph{imperado}, emquanto um acto da vontade é determinado por outro. --- O acto \emph{imperado} suppõe o \emph{eZicito}.\par Note-se que, embora o \emph{imperio} seja, em rigor, proprio da \emph{raso} (porque o \emph{imperio} é a acção de quem ordena e dispõe o agente para o effeito, por meio de uma indicação e intimação, e só a razão póde dispor e ordenar uma coisa para outra, e só ella póde denunciar ou manifestar uma coisa), comtudo costuma tambem attribuir-se á \emph{vontade}, emquanto esta move a razão a produzir o acto do \emph{imperio}, e as outras faculdades a obedecer ao imperio da razão. Por isso o \emph{imperio} é um acto da razão, mas este acto exige um acto previo da vontade, e d'esta recebe a força de intimar ás outras faculdades que produzam o acto respectivo. --- O \emph{imperio} exerce-se sobre os actos do appetite sensitivo (se estes não tiverem prevenido o exercicio da razão) e sobre os actos dos membros externos, e pode excercer-se tambem sobre os actos da vontade e da intelligencia, emquanto a razão pode ordenar que a vontade queira alguma coisa e que a intelligncia se applique a um certo objecto ou dé o seu consentimento a uma verdade não evidente. --- O \emph{imperio} e o \emph{acto imperado} constituem \emph{moralmente} um \emph{unico} acto; porque \emph{unico} é o acto do agente \emph{principal} e do agente \emph{instrumental}, e o agente \emph{principal} é a razão que impera e o \emph{instrumental} é a faculdade que executa esse imperio. Dissemos \emph{moralmente}; porque, \emph{physicamente}, o acto \emph{imperante} e o acto \emph{imperado} são \emph{dois} actos. --- O acto \emph{imperado} attribue-se e pertence \emph{formalmente} á faculdade \emph{imperante}, ainda que \emph{material} ou \emph{immediatamente} derive de uma outra faculdade, porque um effeito attribue-se e pertence ao agente \emph{principal} e não ao agente \emph{instrumental} (\emph{Sum}. \emph{Th}, 1-2, q. 17).}.

b) \textbf{Perfeito} e imperfeito. --- \emph{Perfeito} é o que deriva da vontade com \emph{perfeito} ou \emph{completo} conhecimento da intelligencia.

\pg{571} --- \emph{Imperfeito} é o que deriva da vontade com \emph{imperfeito} ou \emph{incompleto} conhecimento da intelligencia \footnote{A vontade segue a intelligencia, de modo que o acto d'aquella está proporcionado ao d'esta. Por isso a \emph{perfeição} ou \emph{imperfeição} do acto humano depende da perfeição ou da imperfeição do acto da intelligencia.}.

c) Directo e indirecto. --- \emph{Directo} é o que a vontade quer e deseja \emph{em si mesmo}; tal é a morte que o homicida de proposito quiz dar a um individuo. --- \emph{Indirecto} é o que a vontade quer e deseja, não em si mesmo, mas na \emph{causa}, na qual está e se conhece contido; tal é a \emph{morte} causada por uma pedra que se langou imprudentemente, quando um individuo passava, prevendo-se que o podia matar \footnote{Diz S. Thomaz (\emph{Sum}. \emph{Th}., 1-2, q. 77, a. 7): « Aliquid potest esse "\emph{boluntarium} --- vel \emph{secundum se}, sicut quando voluntas \emph{directe} in ipsum fertur, --- vel secundum \emph{suam causam}, quando voluntas fertur in causam, et non in effectum, ut patet in eo, qui voluntarie inebriatur, ex hoc enim quasi voluntarium ei imputatur quod per ebrietatem committit v.\par Pergunta-se: Quaes são as condições necessarias para que um effeito mau, produzido por uma acção, possa ser imputado como \emph{voluntario} ?\par Resposta. As condições necessarias para que um effeito \emph{mau} possa ser imputado ao agente, como coisa \emph{voluntaria}, são as seguintes: --- 1.º) a causa ou acção, de que deriva o effeito, deve estar no \emph{poder} do agente; --- 2.9) o effeito deve ter sido previsto, ao menos de um modo confuso; --- 8.º) o effeito que nasce da \emph{omissão} não é \emph{voluntario}, se não houve \emph{obrigação} de o impedir; --- 4.9) o effeito mau, que nasce da \emph{acção}, não póde ser imputado como \emph{voluntario}, --- se a acção é por si boa ou indifferente, --- se da mesma acção procede tambem immediatamente um effeito \emph{bom}, --- se existe uma causa grave, que permitte o effeito \emph{mau}.\par Na pergunta e na resposta tratámos só do effeito \emph{mau}; porque o effeito \emph{bom}, para ser imputado ao agente, como coisa \emph{voluntaria}, deve ser desejado e procurado \emph{positivamente} em si mesmo.}.

d) Positivo e negativo. --- E' \emph{positivo}, quando a vontade quer operar e realmente opéra; tal é o acto de escrever. --- E' \emph{negativo}, quando a vontade quer deixar de operar; tal é o acto de deixar de soecorrer o proximo \footnote{Alguns modernos opinam que o acto \emph{positivo} corresponde ao acto \emph{directo} e o \emph{negativo} ao \emph{indirecto}. --- Mas esta opinião é falsa; porque, embora todo o acto \emph{negativo} seja \emph{indirecto}, todavia nem todo o acto \emph{indirecto é negativo}; pois ha muitos actos \emph{indirectos}, que são \emph{positivos}. (Cf. \emph{Sum}. \emph{Th}., 1-2, q. 6,a. 8 ad 1; \emph{De Malo}, q. 2, a. 1 ad 2).}.

e) Actual e virtual. --- Actual é o que deriva do influxo \emph{presente} da vontade. --- \emph{Virtual} é o que depende da vontade em virtude de uma determinação anterior que de algum modo continua a influir no proprio acto \footnote{Alguns auctores accrescentam a estas duas especies o acto \emph{habitual}, que, embora tenha existido e não tenha sido desapprovado, todavia não conserva\pgn{572} nenhuma influencia sobre o acto. Mas é claro que o acto /\emph{iabitual} não póde contar-se entre os actos \emph{internos} da nossa vontade, mas só póde'verificar-se nos actos \emph{externos} ou nos seus effeitos.}.

%%>>> art-52

\section*{Artigo II.\\ Diversas especies do acto elicito.}
\addcontentsline{toc}{section}{Artigo II. --- Diversas especies do acto elicito.}

\pg{572} 29. \textbf{Diversas especies} do \textbf{acto} elicito. --- As diversas especies do acto \emph{e/icito} são \emph{seis}, a saber: \emph{amor}, \emph{intenção}, \emph{gozo}, \emph{consentimento}, \emph{eleição}, \emph{uso}. --- As ires primeiras referem-se ao \emph{fim} (bem \emph{honesto}); as tres ultimas, aos \emph{meios} (bem \emph{util}) \footnote{Para bem entendermos a razão da multiplicidade e variedade dos actos \emph{elicitos} da vontade, recordamos que entre esta faculdade e a intelligencia existe mutua dependencia e harmonia. --- A vontade move a intelligencia, quanto ao \emph{exercicio}, porque a determina ao acto. A intelligencia move a vontade quanto á \emph{especificação}, porque lhe apresenta o objecto da tendencia (Cf. \emph{Anthr}., Sec. IT, cap. III, art. IIl). Por isso aos diversos actos \emph{elicitos} da vontade devem corresponder diversos actos da intelligencia.}.

30. \textbf{Amor}. --- Amor é o acto pelo qual a vontade se compraz no bem, que a intelligencia lhe apontou como amavel por si mesmo e honesto, \emph{sem se importar de o alcançar}. --- Por isso este \emph{amor} costuma chamar-se tambem \emph{ve/leidade} \footnote{O primeiro acto da intelligencia é a \emph{apprehensüo} do bem. A esta \emph{apprehensdo} corresponde o primeiro acto da vontade, que é o \emph{amor} do mesmo \emph{bem}; porque a vontade não pôde tender para o bem, se este não lhe for apresentado. Como se vê, a intelligencia move a vontade, não quanto ao \emph{exercicio}, mas quanto á \emph{especificação} (\emph{Of}. \emph{Sum}. \emph{Th}., 1-2, qq. 26, 27, 98).}.

31. Intenção. --- \emph{Intenção} é o acto pelo qual a vontade tende para o bem, que a intelligencia lhe propõe, \emph{para} o \emph{alcançar}. --- A \emph{intenção} póde referir-se não só ao fim \emph{ultimo}, mas tambem aos fins \emph{proximos} ou \emph{secundarios}; pois que tambem estes podem ser o termo da tendencia da vontade \footnote{Se o \emph{amor é effica}:, \& vontade move a intelligencia para esta investigar se aquelle bem 6 /\emph{ionesto} e \emph{possivel}. A intelligencia então profere um \emph{Juizo}, pelo qual affirma que o bem é \emph{honesto} e \emph{possivel}. A este \emph{juizo} segue-se a \emph{intenção}, pela qual a vontade deseja com mais ardor o bem e determina alcancalo pelo emprego dos meios convenientes. Dissemos que a \emph{intenção} póde referir-se tambem aos fins \emph{proximos} ; porque o fim \emph{proximo}, posto que seja \emph{meio} em relação ao flm \emph{ultimo}, todavia é verdadeiro \emph{fin}: em relação aos meios, que se empregam para a sua consecução (\emph{Sum}. \emph{Th}., 1-9, q. 1).}.

32. \textbf{Gozo}. --- \emph{Gozo} é o acto pelo qual a vontade docemente deseanca no bem possuido. --- O gozo \emph{perfeito} encontra-se exclusivamente na posse do bem \emph{universal}, porque só este é o \emph{fim ultimo} de todas as tendencias da nossa alma \footnote{O terceiro acto da vontade em relação ao /\emph{im} é o \emph{gozo}. A intelligencia\pgn{573} attrahe e possue o bem; a vontade goza o bem possuído. O \emph{gozo} é o \emph{ultimo} acto da vontade na ordem da \emph{execução}, porque se alcança, depois do emprego dos meios, --- mas é o \emph{primeiro} na ordem da \emph{intenção}, porque é elle que leva a vontade ao emprego dos proprios meios. --- Embora o objecto \emph{adequado} do \emph{gozo} seja todo o bem \emph{agradavel}, porque este, possuido, produz \emph{gozo}; todavia o seu objecto \emph{principal} é o bem do fim \emph{ultimo}, porque só este, possuido, sacia completamente as tendencias da alma (\emph{Sum}. \emph{T}\textasciicircum{}., 1-2, q. 11).\par Além d'estes tres actos, que á vontade exerce em relação ao \emph{fim}, não ha outros \emph{especificamente} diversos. Porquanto a vontade póde referir-se ao fim por tres diversos modos --- ou emquanto o fim se ama, e temos 0 \emph{amor}, --- ou emquanto se deseja efficazmente, e temos \emph{a intenção}, --- ou emquanto é centro de descanco, e temos o \emph{gozo}.}.

\pg{573} 33. \textbf{Consentimento}. --- \emph{Consentimento} é o acto pelo qual a vontade approva os \emph{meios}, que a intelligencia lhe apontou como aptos e uteis para a consecução do fim. --- Este \emph{consentimento} suppõe o dominio dos proprios actos, e por isso não póde encontrar-se nos animaes irracionaes, que são destuidos de livre arbitrio \footnote{Na ordem da \emph{execução}, o \emph{consentimento} succede \emph{immediatamente á intenção}. Porquanto a vontade, querendo tender \emph{efficazmente} para o fim, move a intelligencia a investigar os meios proporcionados e a julgar ácerca da sua aptidão. A este \emph{conselho} da intelligenza a vontade dá o seu \emph{consentimento}, approvando os meios propostos. --- Este consentimento é absolutamente livre; porque os meios são bens \emph{particulares}, que levam ao bem \emph{universal}, e por isso a vontade tende para elles \emph{livremente} (\emph{Sum}. \emph{Th}., 1-2, q. 15).}.

34. \textbf{Eleição}. --- \emph{Hleição} é o acto pelo qual a vontade, entre os varios meios, escolhe os que a intelligencia julgou mais aptos e mais uteis para a consecução do fim. --- Tambem a \emph{eleição} exige a \emph{liberdade}; e é por isso que os animaes irracionaes empregam só os meios, que a natureza lhes assignou \footnote{Ás vezes os meios, que a inteligencia julgou aptos para a consecução do fim e a vontade approvou, são muitos e diversos, e não podem empregar-se todos ao mesmo tempo. Então a vontade move a intelligencia a julgar qual é, entre todos, o mais apto e efficaz. A intelligencia investiga, julga e propõe á vontade o meio mais apto. A vontade então \emph{escolhe} esse meio. --- Quando, pois, os meios aptos para a consecução do fim são muitos e diversos, o \emph{consentimento} e a \emph{eleição} são \emph{realmente} distinctos entre si. Quando, porém, ha apenas um meio apto para o fim, então aquelles dois actos só \emph{logicamente} differem entre si (\emph{Sum}. \emph{Th}, 1-2, q. 13).}.

35. Uso. --- \emph{Uso} é a determinação pela qual a vontade move \emph{actualmente} as respectivas faculdades a empregar os meios, que foram escolhidos para a consecução do fim proposto \footnote{Em virtude da \emph{eleição}, a intelligencia \emph{intima} ou \emph{declara} o que as outras faculdades devem fazer, para que a vontade alcance o que \emph{escolheu}. Este acto da intelligencia chama-se \emph{émperiío}; porque, por esta \emph{intimação} ou \emph{declaração}, a intelligencia \emph{ordena} as faculdades para os seus actos. Ao /\emph{mperio} da intelligenecia corresponde o acto da vontade, que se chama \emph{uso} e\pgn{574}consiste em mover \emph{actualmente} as faculdades ao emprego do meio escolhido para a consecução do fim proposto (\emph{Sum}, \emph{Th}., 1-2, q. 16).}.

%%>>> art-53

\section*{Artigo III.}
\addcontentsline{toc}{section}{Artigo III.}

\pg{574} \textbf{Impedimentos da voluntariedade}

\textbf{dos actos humanos}.

36. Impedimentos da voluntariedade dos \textbf{actos} humanos. --- Os impedimentos da \emph{voluntariedade} dos actos humanos são \emph{quatro}, a saber: \emph{ignorancia}, \emph{concupiscencia}, \emph{medo} e \emph{violencia} \footnote{Para a \emph{voluntariedade} do acto humano são necessarios dois elementos: o conhecimento da intelligencia e o influxo da livre vontade; de modo que, se faltar um d'estes dois elementos, o acto não é nem póde dizer-se \emph{voluntario}. Portanto toda a causa, que impedir o conhecimento da intelligencia ou o influxo da livre vontade, inpedirá tambem a \emph{voluntariedade} do acto humano. A causa, que póde impedir o conhecimento da intelligencia, é a \emph{ignorancia}. A causa, que póde impedir o influxo da vontade livre, é \emph{intrinseca} ou \emph{extrinseca}. \emph{A intrinseca} é o \emph{appetite sensitivo}, que, --- se leva a vontade a operar pelo «\emph{mor} ou \emph{desejo} do bem sensivel, é \emph{concupiscencia}, --- e, se leva a vontade a operar pelo \emph{receio} do mal sensivel, é \emph{medo}. A causa \emph{extrinseca} é a \emph{violencia}. Logo as causas, que podem impedir a \emph{voluntariedade} dos actos humanos, são, principalmente, a \emph{ignorancia}, \emph{á concupiscencia} o \emph{medo} e a \emph{violencia}.}.

37. ignorancia. --- \emph{Ignorancia} é a ausencia do conhecimento, que se devia possuir. --- Pode ser \emph{antecedente} e \emph{concomitante}. E' \emph{antecedente}, quando \emph{precede} a determinação da vontade e tão \emph{efficazmente} influe no acto, que este não se teria realizado, se tivesse havido conhecimento. E' \emph{concomitante}, quando \emph{acompanha} o acto da vontade, mas não influe nelle, de modo que o acto não teria deixado de se realizar, embora tivesse havido conhecimento \footnote{A ignorancia divide-se em \emph{antecedente} e \emph{concomitante} em relação ao \emph{acto}, emquanto o \emph{precede} ou o \emph{acompanha}. --- Alguns auctores enumeram tambem a ignorancia \emph{consequente}, mas esta não se refere propriamente ao \emph{acto}, mas sim á \emph{vontade}; porque a ignorancia \emph{cansequente} é effeito de um acto da vontade, a qual, \emph{directa} ou \emph{indirectamente}, não quer possuir um determinado conhecimento.}. :

38. \emph{A} ignorancia concomitante não impede a voluntariedade do \textbf{acto} humano, mas sim a \textbf{antecedente}.

a) A ignorancia concomitante não impede a voluntariedade do acto humano. --- Toda a causa, cujo effeito não contraria a vontade, não impede a voluntariedade do acto humano. Ora a ignorancia \emph{concomitante} é uma causa, cujo effeito não contraria a vontade; porque o acto teria sido realizado, embora \pg{575}  tivesse havido conhecimento. Logo a ignorancia \emph{concomitante} não impede a voluntariedade do acto humano \footnote{Apresentamos um exemplo da ignorancia \emph{concomitante}. Pedro vai procurar o seu inimigo Paulo, para o matar. Um dia, julgando matar um veado, Pedro mata Paulo. Obrou com ignorancia \emph{concomitante}, porque ella \emph{acompanhou} o acto. Por isso este acto de Pedro é \emph{não-voluntario}, porque a vontade nao póde querer uma coisa, que ignora, --- mas não é \emph{nvoluntario}, isto é, \emph{não impede nem contraria} a voluntariedade do acto humano, porque não só não repugna á vontade, mas está muito conforme com os seus desejos. Logo a ignorancia. \emph{concomitante} não impede a voluntariedade do acto.\par O que dissemos da ignorancia \emph{concomitante} deve tambem dizer-se da \emph{consequente}. Na verdade, quem quer uma causa, quer tambem os effeitos, que necessariamente derivam d'ella. Ora a vontade, directa ou indirectamente, quer a ignorancia \emph{consequente}. Logo quer tambem os actos, que ella produz em consequencia d'essa ignorancia.}.

b) A ignorancia antecedente impede a voluntariedade do acto humano. --- Toda a causa, cujo effeito contraria a vontade, impede a voluntariedade do acto humano. Ora a ignorancia \emph{antecedente} é uma causa, cujo effeito contraria a vontade; porque,

\emph{Z} sem ella, a vontade, não teria realizado o acto. Logo a ignorancia \emph{antecedente} impede a voluntariedade do acto humano \footnote{Note-se que a ignorancia \emph{antecedente}, para impedir a voluntariedade do acto humano, deve ser \emph{invencivel}, de modo que não tenha podido acabar ---- ou porque todos os esforços empregados não deram o resultado desejado, --- ou porque não occorreu nenhuma duvida ou suspeita ácerca da obrigação de se instruir; pois, se a ignorancia foi \emph{vencivel}, isto é, se ella não cessou quando podia e devia cessar, não é contraria á vontade, e por isso não impede a voluntariedade do acto humano. --- Mas esta ignorancia \emph{vencivel}, se não \emph{impede} a voluntariedade, \emph{diminue-a}, e o individuo é mais ou menos culpado, conforme o grau da negligencia commettida (\emph{Sum}. TÀ., 1-2, q. 6, a. 8).}.

39. Concupiscencia. --- \emph{Concupiscencia} é o movimento do appetite sensitivo para um bem \emph{agradavel}, mas \emph{ausente}. --- Pode ser \emph{antecedente} e \emph{consequente}. E' \emph{antecedente}, quando, na ordem da causalidade, não só precede todo o acto da vontade, mas até solieita esta faculdade ao consentimento. E' \emph{consequente}, quando a vontade excita o appetite sensitivo para o bem agradavel \footnote{A \emph{concupiscencia}, e, em geral, a \emph{paixão}, embora não possa influir \emph{directamente} na vontade (porque uma faculdade \emph{organica} não póde operar numa faculdade \emph{espiritual}) todavia póde influir nella \emph{indirectamente}. Este influxo \emph{indirecto} da \emph{concupiscencia} na vontade póde dar-se por dois modos: --- 1.º) ou porque o acto do appetite sensitivo é muito intenso, e esta intensidade diminue ou impede a tendencia da vontade; pois, derivando as nossas faculdades da alma, se uma augmenta em força, \& outra deve diminuir e até suspender o seu acto; --- 2.") ou porque o appetite sensitivo inflamma a imaginação, e esta altera o juizo da razão, e a alteração do juizo da razão leva a vontade para o bem sensivel.}.

\pg{576} 40. \emph{A} concupiscencia \textbf{antecedente} póde impedir, e póde augmentar a voluntariedade do \textbf{acto} humano.

a) \textbf{Póde} impedir. --- A oconeupisceencia \emph{antecedente} poderá impedir a voluntariedade do acto humano, se impedir o conhecimento da intelligencia; pois o acto \emph{voluntario} exige o conhecimento d'esta faculdade. Ora a coneupiscencia \emph{antecedente} póde impedir, e muitas vezes impede o conhecimento da intelligencia. Logo a coneupiscencia \emph{antecedente} póde impedir a voluntariedade do acto humano \footnote{Quando, pois, a paixão é tão vehemente que impede o uso da razão, o acto não é \emph{voluntario} nem \emph{involuntario}, mas é propriamente \emph{não-voluntario}. Este principio vale para os casos em que a paixão não foi \emph{voluntaria} no seu principio; porque, se o foi, o acto deve considerarse voluntario \emph{indirectamente} ou na sua \emph{causa} (\emph{Sum}. \emph{Th}., 1-2, q. 6, a. 7 ad 8; q. 77, a a. 6, 7).}.

b) \textbf{Póde} augmentar, --- A coneupiscencia \emph{antecedente} poderá augmentar a voluntariedade dos actos humanos, se excitar a vontade a tender com maior intensidade de affecto para o bem sensivel. Ora a concupiscencia \emph{antecedente} póde excitar, e muitas vezes excita a vontade a tender com maior intensidade de affecto para o bem sensivel. Logo a coneupiscencia \emph{antecedente} póde augmentar a voluntariedade do acto humano \footnote{A \emph{concupiscencia}, quando augmenta a voluntariedade do acto humano, deixa á intelligencia o poder sufficiente para reflectir sobre a natureza do acto; porque, se lhe não deixasse esta reflexão, não só não augmentaria a voluntariedade, mas destruil-a-hia, pois o \emph{acto-voluntario} exige essencialmente a reflexão da intelligencia. --- Mas, para não cahirmos em equivoco, devemos advertir que a concupiscencia \emph{antecedente} augmenta a voluntariedade do acto humano, --- não emquanto o torna mais proprio da vontade e mais immanente nesta faculdade, porque neste caso o acto, dependendo da paixão, que é uma coisa extrinseca, é \emph{menos voluntario} do que se dependesse exclusivamente da razão e da vontade, e por isso dizemos que a paixão diminue o \emph{voluntario}, --- mas emquanto torna a vontade mais propensa ou inclinada para o bem sensivel. (\emph{Sum}. \emph{T}., 1-2, q. 6, a. 1).\par Fallamos só da concupiscencia \emph{antecedente}; pois a \emph{consequente}, sendo effeito da vontade, augmenta a voluntariedade, ou, melhor, é um signal que demonstra a intensidade do acto da vontade (\emph{Sum}. \emph{Th}., 1-2, q. 77, a. 6).}.

41. Medo. --- \emph{Medo} é uma trepidação do animo pela apprehensão de um mal \emph{Zmmznente}. --- Costuma dividir-se em \emph{grave} e \emph{leve}. E' \emph{grave}, quando o mal é \emph{grande} e \emph{certo}. E' \emph{leve}, quando o mal é \emph{pequeno} e \emph{incerto} \footnote{O \emph{medo} é \emph{absoluta} ou \emph{relativamente grave}. E' \emph{absolutamente} grave, quando é considerado como tal por homens fortes, tal é o medo da \emph{morte}. E' \emph{relativamente} grave, quando só o é para as pessoas fracas ou meticulosas.}.

49. O medo, embora \textbf{grave}, não impede, ordinariamente, a \textbf{voluntariedade} do \textbf{acto} humano. --- O \emph{medo}, embora \emph{grave}, não impedirá \pg{577}  a voluntariedade do acto humano, se a vontade, sob a impressão do \emph{medo}, poder ainda deliberar e preferir uma coisa á outra. Ora a vontade, mesmo sob a impressão de um medo \emph{grave}, póde ainda deliberar e preferir uma coisa á outra, embora com algum sacrificio. Logo o \emph{medo}, embora \emph{grave}, não impede a voluntariedade do acto humano.

Dissemos --- \emph{ordinariamente}; porque o \emph{medo}, se impedir ou perturbar o uso da razão, impedirá tambem a voluntariedade do acto humano \footnote{Quando o \emph{medo é grave}, o eapirito encontra-se entre dois bens, dos quaes um é ou parece \emph{maior}, e outro é ou parece \emph{menor}. Assim, na imminencia de um naufragio, deveriam salvar-se a vida, que é um bem \emph{maior}, e o thesouro, que é um bem \emph{menor}. Mas ambos estes bens não podem salvar-se, e por isso lanca-se ao mar o thesouro. O homem, neste caso, quer absoluta e simplesmente o fim, que é a salvação da sua vida, e por isso quer absoluta e simplesmente o meio, que é lançar ao mar o thesouro, para aliviar o navio. Logo o acto de lançar ao mar o thesouro, motivado pelo medo \emph{grave}, \emph{é voluntario}. Logo o medo, embora \emph{grave}, não impede ordinariamente a voluntariedade do acto humano. --- Todavia, como o homem não teria lançado ao mar o seu thesouro, se tivesse encontrado outro meio para salvar a vida, o acto, motivado pelo medo \emph{grave}, chama-se \emph{énvoluntario hypotheticamente} (\emph{Sum}. \emph{Th}., 1-2, q. 6, a. 6. --- Embora o acto praticado pelo \emph{medo grave} seja \emph{por si} voluntario, comtudo as leis positivas, tendo em vista o bem commum, declararam irritos ou rescindiveis alguns contractos, que o medo grave extorquiu; assim pode ser nullo o matrimonio, nulla a promessa ou o pagamento do dote, nulla a jurisdicção dada ou recebida, nulla a eleição dos prelados, nullos os votos solemnes.\par D'onde se vé a diversidade da influencia que o \emph{medo} e a \emph{concupiscencia} exercem sobre o acto humano. O acto, motivado pelo \emph{medo}, pratica-se com uma certa repugnancia da vontade, e só para evitar um mal maior; mas o acto, motivado pela \emph{concupiscencia}, pratica-se com todo o transporte e complacencia da vontade, porque o seu objecto é o bem agradavel.}.

43. Violencia. --- \emph{Violencia} é á moção de um agente exterior, contraria á inclinação da vontade do sujeito, que a recebe. --- A \emph{violencia}, pois, é diversa da \emph{necessidade}, que leva a vontade a tender \emph{úrresistivelmente} para o bem \emph{universal} \footnote{A. vontade, na sua tendencia \emph{necessaria} para o bem \emph{universal}, opéra e quer; ao passo que, na \emph{violencia}, não opera nem quer.\par 37 Vor. II.}.

44. A violencia não póde impedir a voluntariedade dos \textbf{actos} elicitos, mas só a dos imperados.

a) Não póde impedir a voluntariedade dos actos elicitos. --- A \emph{violencia} não poderá impedir a voluntariedade dos actos \emph{e/citos}, se estes n\&o estiverem sujeitos ao poder de um principio exterior. Ora os actos \emph{e/citos} não estão sujeitos ao poder de um principio exterior; porque derivam \emph{Immediatamente} da vontade \pg{578}  e alli ficam. Logo a \emph{violencia} não póde impedir a voluntariedade dos actos \emph{elicitos} \footnote{Um acto que fosse \emph{elicito} e \emph{violento} seria uma contradieção. Na verdade, por ser \emph{elécito}, seria \emph{conforme} á inclinação da vontade, --- e, por ser \emph{violento}, seria \emph{contrario} á inclinação da vontade. Ora um acto \emph{conforme} e \emph{contrario} á inclinação de uma faculdade é uma contradicção.}.

b) \textbf{Póde} impedir a voluntariedade dos actos imperados. --- A \emph{violencia} poderá impedir a voluntariedade dos actos \emph{imperados}, se estes estiverem sujeitos ao poder de um principio exterior. Ora os actos \emph{imperados} podem estar sujeitos ao poder de um principio exterior, apezar da resistencia da vontade; porque estes actos podem derivar \emph{immediatamente} de uma faculdade organica e externa, e por isso sujeita á acção dos agentes externos. Logo a \emph{violencia} póde impedir a voluntariedade dos actos \emph{émperados} \footnote{Neste caso temos uma lucta; e os actos são propriamente \emph{énvoluntarios}, por contrarios á tendencia da vontade (\emph{Sum}. \emph{T}7., 1-2, q. 6, aa. 4, 5).\par Para evitarmos repetições, pedimos ao leitor que torne a examinar o que escrevemos na \emph{Anthropologia} (Sec. II, cap. III, art. III) ácerca dos diversos factos \emph{physiologicos} e \emph{pathologicos}, que podem impedir e impedem a voluntariedade do acto humano.}.

%%>>> cap-9

\chapter*{CAPITULO TERCEIRO\\[0.3em]\large MORALIDADE DOS ACTOS HUMANOS}
\addcontentsline{toc}{chapter}{CAPITULO TERCEIRO --- MORALIDADE DOS ACTOS HUMANOS}

\textsc{SuuMaRIO}: --- Moralidade, sua existencia e fundamento. --- Fontes da mora-

lidade dos actos humanos. --- Consequencias da moralidade dos actos

humanos. --- A moralidade dos actos humanos e os habitos.

%%>>> art-54

\section*{Artigo I.\\ Moralidade, sua existencia e fundamento.}
\addcontentsline{toc}{section}{Artigo I. --- Moralidade, sua existencia e fundamento.}

45. Moralidade. --- \emph{Moralidade} é a propriedade, pela qual os actos \emph{humanos} são dotados, ou destituidos, de \emph{honestidade}, de modo que convenham, ou não convenham, á natureza humana

£ e possam, ou não, levar á consecução do fim \emph{ultimo}. --- Por isso o acto \emph{moral} é o \emph{genero}, cujas \emph{especies} são o acto \emph{honesto} e o acto \emph{não-honesto}. O acto \emph{honesto} chama-se \emph{bem moral}; o acto \emph{não-honesto} diz-se \emph{mal moral} \footnote{A \emph{moralidade} é uma propriedade do acto \emph{humano}, isto é, do acto que\pgn{579}procede da intelligencia e da livre vontade. Por isso o acto, que não é \emph{livre}, não póde ser \emph{moral}. --- A \emph{moralidade} do acto \emph{humano} é, como dissemos, o objecto \emph{formal} d'esta parte da philosophia, e o proprio acto \emph{umano} é o seu objecto \emph{material}. D'ahi a razão por que primeiramente tratámos da \emph{essencia} do acto \emph{humano}, e agora da sua \emph{moralidade}.}.

\pg{579} 46. Divisão da moralidade. --- A \emph{moralidade} divide-se em \emph{extrinseca} e \emph{intrinseca}. --- \emph{Extrinseca} é á que convém ao acto unicamente \emph{em consequencia} de uma prohibição ou de um preceito do legitimo superior. \emph{Intrinseca} é a que convém ao acto em virtude da sua natureza, independentemente de toda e qualquer prohibição ou preceito do superior \footnote{Como se vé, a moralidade \emph{ewtrinseca é accidental} ao acto humano, ao passo que a \emph{énírinseca} é \emph{essencial} ao mesmo acto.}.

4T. A moralidade \textbf{existe}.

a) A \emph{moralidade} existirá, se existirem actos \emph{moralmente bons} ou \emph{honestos} e actos \emph{moralmente maus} ou \emph{não-honestos}. Ora existem effectivamente actos \emph{moralmente} bons (como são o \emph{temor de} Dus, o \emph{respeito aos paes}, o \emph{amor da patria}, ete.), e actos \emph{moralmente} maus (como são a \emph{blasphemia}, o \emph{parricidio}, \emph{á traição da patria}, ete.). Logo a \emph{moralidade} existe.

b) Todos os homens, em todos os tempos e em todos os logares, admittiram que alguns actos são \emph{Honestos} e outros \emph{não-honestos}. Este consenso unanime, universal e constante, podia provir ou da evidencia da coisa ou da moção da vontade. Ora não podia provir da moção da vontade; pois que esta faculdade fica naturalmente contrariada com a existencia de actos \emph{honestos}, que se devem praticar, e de actos \emph{müo-honestos}, que se devem evitar. Logo deriva da evidencia da coisa. Ora a evidencia é criterio da verdade. Logo existem actos \emph{honestos} e actos \emph{não-honestos}, e por isso a \emph{moralidade} existe \footnote{Cf. \emph{C}. \emph{Gentes}, l. IIl, c. 129. --- Negam ou devem logicamente negar a existencia da verdadeira \emph{moralidade} todos os que negam a \emph{liberdade} da vontade humana, como os \emph{materialistas}, os \emph{pantheistas}, os \emph{fatalistas}, os \emph{deterministas}. Porquanto, seja qual for a natureza da \emph{moralidade}, é certo que, sem liberdade, não póde haver \emph{moralidade}; alias tambem as acções do animal irracional seriam \emph{moraes}: o que é absurdo.}.

48. A moralidade extrinseca existe. --- A moralidade \emph{extrinseca} existirá, se existirem actos que são \emph{honestos} ou \emph{müo-honestos} unicamente em \emph{consequencia} de uma prohibição ou de um pre-. ceito do legitimo superior. Ora a experiencia attesta que semelhantes actos existem. Logo a moralidade \emph{extrinseca} existe \footnote{Esta proposição é admittida por todos os escriptores, que reconhecem e sustentam a \emph{existencia} da \emph{moralidade}.}.

\pg{580} 49. A \textbf{moralidade} intrinseca \textbf{existe}.

a) A moralidade \emph{intrinseca} existirá, se existirem actos \emph{Z}0\emph{nestes} e \emph{não-honestos}, que, sempre e em toda parte, independentemente de toda e qualquer vontade, convêm, ou não, á nossa natureza racional; porque, nestes casos, a \emph{conveniencia} ou \emph{não conveniencia} destes actos só póde derivar da sua essencia, que só é immutavel. Ora existem semelhantes actos; assim o \emph{temor de} Dus, o \emph{respeito aos paes}, o \emph{amor do proximo}, ete., são actos \emph{naturalmente} honestos, independentemente de toda a vontade, e não deixariam de ser taes, embora uma vontade mal avisada os prohibisse. Logo a moralidade \emph{intrinseca} existe.

b) Todos os povos, em todos os tempos e em todos os logares, admittiram que a differença existente entre alguns actos \emph{honestos} e alguns \emph{não-honestos} (como entre o amor e o odio aos paes, á patria) é tão \emph{intrinseca} e tão \emph{essencial}, que, em nenhuma circumstancia, póde deixar de existir; e esta persuasão, tão universal, constante e unanime, só póde derivar da natureza, aliás teriamos um effeito sem causa proporcionada. Ora o que deriva da natureza não está sujeito a erro. Logo a- moralidade \emph{Intrinseca} existe \footnote{A moralidade \emph{intrinseca} é o fundamento da \emph{exstrinseca}. Com effeito, quando um subdito transgride a lei do principe, commette uma acção má. E porque é má essa acção? Porque a desobediencia ao principe é uma coisa \emph{essencialmente md}. --- D'aqui se segue que, se negarmos \& moralidade \emph{intrinseca}, devemos tambem negar a moralidade \emph{ewtrinseca},; e com isso destruir-se-ha a propria sciencia \emph{moral}.\par E' necessario, pois, dizer que a differenca que existe entre alguns actos \emph{moralmente} bons e outros \emph{moralmente} maus, é \emph{objectiva} e \emph{absoluta}, isto á, fundada na propria essencia das coisas e independente de toda e qualquer causa \emph{subjectiva} ou \emph{variavel}. E\par Os systemas, que negam a moralidade \emph{intrinseca}, são muitos e diversos, mas todos podem reduzir-se a \emph{quatro}, que são o \emph{materialismo dos epicureus}, o \emph{utilitarismo dos empiricos}, \emph{dos economistas e dos positivistas}, o \emph{autonomismo dos racionalistas}, e o \emph{voluntarismo dos juristas}. --- Na verdade, todos os \emph{systemas}, que negam uma differenca \emph{absoluta} e \emph{objectiva} entre o \emph{ben} e o \emph{mal moral}, e só reconhecem uma moralidade \emph{subjectiva} e \emph{relativa}, podem fundar esta moralidade e esta differenca --- ou no sentimento do \emph{praser}, ou no instincto do \emph{interesse}, \emph{individual} ou \emph{social}, ou nos \emph{principios} da razão sem caracter \emph{absoluto}, ou na \emph{livre vontade} dos homens. --- Daremos o resumo e faremos a critica de cada um d'estes systemas, procurando dar maior desenvolvimento aos dos \emph{positivistas}, que são mais seguidos.\par I. Materialismo. --- O \emph{materialismo}, devido a \emph{Epicuro}, formou a \emph{Moral} sem \textsc{Deus}, sem principios absolutos, inteiramente egoista, fundada sobre o bem \emph{sensivel}, que é o prazer. O \emph{prazer}, segundo este systema, divide-se em \emph{movel} e \emph{permanente}. O \emph{movel} é o que dura pouco tempo e cujo effeito produz\pgn{581}mais tarde soffrimento; ao passo que o \emph{permanente}, que é mais intimo e doce, consiste no perfeito repouso ou na isenção de toda a dor ou magoa. A \emph{Moral} deve proporcionar o meio para a acquisição d'este repouso perfeito. Tal meio é a fidelidade ás quatro virtudes cardeaes, que são: a \emph{prudencia} (que leva a evitar os perigos das inquietações), a \emph{coragem} (que nos defende do medo vão e illusorio), a \emph{justiça} (que nos indica a respeitar o pacto social, fundado no interesse) e a \emph{temperanca} (que, moderando as necessidades e os desejos, diminue a somma das dores e augmenta a dos prazeres). Por isso \& \emph{Moral} de Epicuro é a arte de procurar o prazer e evitar a dor, pelo modo mais facil, seguro e duradouro; visto que todo o acto, que produz \emph{prazer}, \emph{é bom}, e todo o acto, que produz \emph{dor}, \emph{é mau}. --- \emph{Epicuro}, foi seguido --- no sec. XVIII, por \emph{Lamettrie}, \emph{Helvetius}, \emph{Diderot}. \emph{Holbach} e \emph{Volney}, --- e, no sec. XIX, por \emph{Haechkel}, \emph{Vogt}, \emph{Molleschot}, e por todos os que professam o \emph{materialismo} na sua forma mais grosseira.\par A Moral \emph{materialista} é falsa, não só pelas razões com que temos demonstrado, em varios pontos do presente trabalho, os absurdos do \emph{materialismo}, mas tambem pelas seguintes:\par \emph{a}) A verdadeira \emph{Moral} não póde existir, se o homem não for dotado de \emph{liberdade}; porque é absurda a pretenção de dirigir a actividade, quando esta, por não ser livre, não é capaz de direcção. Ora a Moral \emph{materialista} nega que o homem é dotado de \emph{liberdade}; porque reduz todos os nossos actos a movimentos meramente mechanicos. Logo a moral \emph{materialista} é falsa.\par \emph{b}) A moral \emph{materialista} colloca o bem supremo do homem no bem \emph{sensivel}, \emph{e} o mal supremo no mal \emph{sensivel}. Ora o bem \emph{supremo} do homem é o bem \emph{universal} e \emph{infinito} --- \textsc{Deus}. Logo a moral \emph{materialista} é falsa.\par \emph{c}) Se o systema dos \emph{materialistas} fosse verdadeiro, o homem, no caso de um conílieto entre o prazer \emph{sensivel} e o dever, deveria pronunciar-se sempre em favor do prazer e contra o dever; e por isso o heroe, que morre pela patria, seria um insensato, ao passo que um cobarde, que atraicoa os seus para salvar a propria vida, seria um homem sabio e prudente. Ora estas consequencias repugnam. Logo a moral \emph{materialista} é falsa.\par \emph{d}) Se o systema dos \emph{materialistas} fosse verdadeiro, a medida do \emph{egoismo} seria, sempre e em toda parte, a medida da perfeição; de modo que, na ordem moral, quem mais exclusivamente procurasse, sempre e em toda parte, o seu maior prazer, seria o mais perfeito. Ora a humanidade não acceita uma theoria, que seria a apotheose do egoismo e produziria uma funesta e irremediavel desordem social. Logo á moral \emph{materialista} é falsa.\par IT, Utilitarismo dos empiricos, dos economistas e dos positivistas. --- Este systemd, nos seus principios, não differe muito do systema anterior; porque o fim \emph{ultimo} dos nossos actos, tanto para os \emph{materialistas} como para os \emph{utilitaristas}, é um bem finito e creado, constituido pelos \emph{prazeres} ou pelo \emph{interesse}, ou pela \emph{gloria}, ou por qualquer outra vaidade. --- Alguns auctores dividem o \emph{utilitarismo} em \emph{individual} e \emph{social}, conforme os seus defensores fundam a \emph{moralidade} dos actos humanos no interesse do \emph{individuo} ou no da \emph{sociedade}, isto é, conforme os seus defensores são \emph{egoistas} ou \emph{altruistas}. E nada mais logico. --- ''odavia, aqui fundamos a divisão d'este systema nas diversas\pgn{582}phases da sua evolução. A razão d'isto é que, embora o \emph{utilitarismo} possa ser \emph{egoista} e \emph{altruísta}, comtudo as formas, que o exprimem, não têm, \emph{em geral}, um caracter definitivamente marcado, mas são todas, mais ou menos, \emph{ego-altruistas}. Dissemos --- em \emph{geral}; porque \emph{Comte} parece exclusivamente \emph{altruísta}, \emph{e Bentham} exclusivamente \emph{egoista}.\par A) A primeira phase do \emph{utilitarismo} é representada pelos sequazes do \emph{empirismo}, por \emph{Locke} (1632-1704), pelo \emph{Conde de Shaftesburg} (1671-1718), pelo \emph{Marquez de Saint-Lambert} (1011-1808) e por Adam \emph{Smith} (1723-1790).\par a) \emph{Locke}, levado pelos principios do \emph{sensualismo}, baseou as suas theorias moraes na \emph{felicidade terrena} e no \emph{interesse}. A idea do \emph{direito} reduz-se á da \emph{utilidade}, individual e social. O primeiro direito, que é egual em todos, é o de procurar sobretudo o interesse proprio. O fundamento de todos os direitos civis e politicos é a \emph{egualdade na liberdade} ; porque todos os homens nascem eguaes, livres e independentes uns dos outros. Para a paz e harmonia d'estas liberdades é necessario um \emph{contracto social}.\par b) O \emph{Conde de Shaftesburg}, amigo e discipulo de \emph{Loche}, no seu livro u \emph{Etudes sur la vertu et le mérite} », começa por dizer que \emph{bem} é tudo o que é conforme ás \emph{tendencias} da natureza, e \emph{mal} é o que lhes é contrario. As \emph{tendencias} da natureza são de duas especies: tendencias de \emph{egoismo individual}, que levam o homem a procurar o seu interesse particular, e tendencias de \emph{sympathia social}, que têm por objecto o interesse commum e social. As segundas oceupam o primeiro logar; mas não devemos sacrificar as primeiras. O verdadeiro principio da \emph{moralidade} é a \emph{perfeita harmonia entre as inclinações egoistas} ou \emph{individuaes e as inclinações sympathicas ou sociaes}. Por isso um acto é \emph{bom}, quando é conforme ao interesse commum da sociedade, e é \emph{mau}, quando é contrario a esse interesse.\par 0) O \emph{Marques de Saint-Lambert}, na sua obra « \emph{Principes de moeurs ches toutes les nations ou Catéchisme universal} », diz que a origem da \emph{Moral} se deve reduzir ao instincto da conservação ou do interesse. Este instincto soffre modificações, conforme o clima e as exigencias da vida social. D'ahi as opiniões e as tradições, que, depois de uma longa experiencia dos seus bons resultados, se tornaram principios da \emph{Moral}.\par 4) Adam \emph{Smith}, sensualista, reconhece no sentimento da \emph{sympathia universal} o principio e a regra da moralidade dos nossos actos; e por isso um acto é \emph{bom} ou \emph{mau}, \emph{honesto} ou \emph{náüo-honesto}, conforme excita sentimentos de \emph{sympathia} ou de \emph{antipathia}. A sympathia, passiva e fatal, varia com os caracteres, paixões e prejuizos, e por conseguinte deve ser regulada pela \emph{sympathia} de um \emph{espectador imparcial}. --- E' esta a \emph{Moral sentimental}.\par B) A segunda phase do \emph{utilitarismo} é representada pelos \emph{economistas-materialistas}, que reconhecem como chefe a \emph{Jeremias Bentham}.\par \emph{Bentham} diz que o unico principio dos nossos actos é o \emph{interesse pessoal} --- o interesse da propria felicidade actual e sensivel. Na verdade, todos os homens são essencial e fatalmente \emph{egoistas}; ninguem poderá fazer bem aos outros, se nisto não encontrar o seu maior interesse, o seu maior prazer. A \emph{Moral}, portanto, consiste em \emph{regular} o \emph{egoismo}, em \emph{calcular} a \emph{quantidade} dos diversos prazeres; porque o maior bem é o maior prazer. --- Para calcular a \emph{quantidade} dos prazeres, é necessaria uma \emph{aritlimetica moral}, pela\pgn{583}qual os prazeres se considerem nas suas cirecumstancias de \emph{intensidade}, \emph{duração}, \emph{fecundidade}, \emph{isenção das dores}, \emph{proximidade}, \emph{certeza}, \emph{extensão das suas consequencias sociaes}. Todo o acto moral, que se não reduza a uma quantidade superior de prazeres, é absurdo. « O mais abominavel prazer de um malfeitor é bom em si, e só se torna mau pelos soffrimentos, que póde causar ». Por isso o \emph{dever} é á obrigação, que a todo o homem incumbe de procurar o seu prazer ou o seu interesse. O interesse \emph{particular} não póde ser contrario ao interesse \emph{commum}. Entre estas duas ordens de interesse existe uma identidade natural, sanccionada pelas leis humanas; porque o bem particular e o commum derivam das mesmas causas, o bem dos outros é o bem nosso, por causa da sympathia, que une todos os homens. Por isso o \emph{utilitarismo particular} tende para o \emph{utilitarismo social}. --- Como se vê, \emph{Bentham} segue a theoria de \emph{Epicuro}. Limita-se apenas a alargar a esphera da actividade humana, mas sempre em beneficio do \emph{individuo}; porque o individuo isolado não póde ter todas as satisfações, que a sociedade proporciona.\par €) A terceira phase de \emph{utilitarismo} é representada pelos \emph{positivistas}, sequazes do celebre eseriptor francez --- \emph{Augusto Comte}.\par a) \emph{Augusto Comte} (1198-1857), applicanda á \emph{Moral} os principios do methodo \emph{positivista}, que consiste em rejeitar tudo o que escapa á experiencia sensivel, --- nega, como coisas chimericas, a idea de um fim \emph{ultimo} ultramundano, de \emph{liberdade} e \emph{responsabilidade moral}, e por isso de \emph{deveres absolutos}, --- e affirma que a \emph{Moral} deve fundar-se no imperio do instincto, que nos inclina para os outros e que se chama \emph{altruísmo}. O homem tem instinctos \emph{egoistas} e \emph{altruistas}. Se os primeiros são necessarios, não o são menos os segundos. Por quanto, a humanidade só póde viver em sociedade, e esta só continua e progride pelo sacrificio e pela dedicação; de modo que a lei da existencia da humanidade é uma lei que se resume nesta formula: Vive \emph{para os outros}. Por isso a moralidade consiste propriamente na preponderancia dos instinctos \emph{altruistas} sobre os \emph{egoístas}, --- preponderancia, que resulta \emph{fatalmente} da educação e da sciencia. D'ahi se segue que o acto é \emph{bom}, quando deriva do \emph{a/truismo}.\par O \emph{positivismo} de \emph{Comte} foi seguido por duas escholas: pela eschola franceza, representada principalmente por \emph{Littré} e \emph{Taine}, e pela eschola ingleza, representada sobretudo por \emph{Stuart Mill}, \emph{Herbert Spencer} e \emph{Alevandre Bain}. --- Ambas estas escholas --- convêm em que a \emph{Moral} é a sciencia dos meios para transformar fatalmente o \emph{egoismo} em \emph{altruísmo}, no \emph{interesse} commum do individuo e da sociedade, --- mas differem emquanto ao modo por que explicam a formação do \emph{altruísmo}. Porquanto, a escola franceza basóa-se principalmente na \emph{physiologia}, para descobrir, no nosso organismo, a origem e a lei da evolução do \emph{altruísmo}; e a eschola ingleza, fundando-se sobretudo na \emph{psychologia} e no systema do \emph{darwinismo}, procura demonstrar a evolução psychologica dos nossos sentimentos, de \emph{egoistas} em \emph{altruistas}, sob a influencia do meio social. Mas, para todos, a \emph{Moral é relativa} e \emph{variavel}, porque relativas e variaveis são as tendencias e as evoluções do organismo. --- Damos o resumo da doutrina dos mencionados escriptores.\par bp) \emph{Littré} resume a sua \emph{Moral} nas seguintes palavras: « Ha dois prineipios de Moral: o \emph{egoísmo} e o \emph{altruísmo}. O primeiro, representando as ne-:\pgn{584}cessidades, é um principio de conservação do individuo; o segundo, representando a sexualidade (\emph{sexualité}) é um principio de expansão fôra do individuo. Ambos elles partem do mesmo ponto, que é a acção da substancia viva no cerebro por meio dos nervos. Assim considerados, o \emph{egoismo} e o \emph{altruismo} não passam de germens, desenvolvidos em seguida pela experiencia, pelo raciocinio e pelo tempo. E' por este modo que se formam as sciencias \emph{moraes}, sempre \emph{relativas}, das differentes epochas e das diversas nações, que progridem sempre á medida que a noção da humanidade, desembaracandose, comprime o \emph{egoismo} e dilata o \emph{altruismo} ».\par No sentido de \emph{Littré}, pois, a origem dos phenomenos moraes deve procurar-se na estructura e no organismo da substancia viva, porque as tendencias \emph{egoístas} e \emph{altruistas} são apenas uma transformação de duas tendencias essenciaes, de duas necessidades primitivas e irreductiveis da substancia viva. --- A primeira tendencia ou necessidade é a da \emph{nutrição}, e cria o \emph{egoismo} ou o amor de si, que é o instincto da conservação. D'este instincto, que é mais ou menos complexo, conforme a maior ou menor complexidade dos orgãos, nasce o cuidado de procurar os meios de conservação, o desejo do poder, da propriedade e de todos os bens, que tornam agradavel a vida. --- Mas o ser vivo não vive só para si; não é só individuo, é tambem \emph{especie}. D'ahi um novo instincto, emanante do intimo do organismo, --- o \emph{altruísmo}, que é o amor da especie com a necessidade da reproducção, da expansão e da beneficencia em todas as suas formas.\par As necessidades da \emph{nutrição} e da \emph{reproducção}, \emph{pluysiologicas} no começo, transformam-se em necessidades \emph{moraes}, pela complexidade crescente dos orgãos e pela progressiva evolução das sensações, das imagens, dos sentimentos, das aptidões correspondentes.\par Para se entender esta transformação, é necessario fazer uma distineção entre a \emph{producção} da Moral e a sua \emph{elaboração}. \emph{A elaboração} da Moral é obra do \emph{cerebro}; ao passo que a sua \emph{producção} é obra de \emph{toda} a \emph{substancia viva}, desde os pés até á cabeça. Assim como o estomago não é o \emph{creador}, mas o \emph{elaborador} do alimento, e é necessario que lhe forneçam o que elle deve digerir, assim tambem o cerebro \emph{elabora} as necessidades, que a substancia viva lhe leva para as transformar em sentimentos moraes. --- A substancia \emph{cinzenta} do cerebro, --- recebendo as necessidades \emph{egoistas}, vindas de todos os pontos do organismo, realiza uma \emph{reacção}, e produz os sentimentos e as paixões relativas á conservação do individuo, ao seu desenvolvimento e perfeição, --- e, recebendo as necessidades \emph{altruistas} da substancia viva, produz, por uma outra reacção, os sentimentos da familia, da. dedicação para com os seus, da sociabilidade, do amor da patria e da humanidade. --- Estes sentimentos, \emph{egoistas} e \emph{altruistas}, pelo nexo intimo que existe entre as cellulas cerebraes, desenvolvem-se num circulo de acções e reacções, numa lucta constante, e nisto consiste a \emph{vida moral}. A. victoria, porém, será do \emph{altruísmo}; porque as necessidades de um organismo mais complexo vencem as de um organismo menos complexo, e as funcções da reproducção, da expansão, do amor, de que deriva o \emph{altruísmo}, são mais complexas que as 'funcções da nutrição, que é o principio do \emph{egoismo}. Por isso o termo, para\pgn{585}o qual á humanidade tende por uma lei necessaria e irresistivel de evolução, é á \emph{fraternidade universal}.\par Dos principios expostos vê-se o que afinal denotam, no systema de \emph{Littré}, as ideas de \emph{lei}, de \emph{liberdade}, do \emph{Justiça}. --- A \emph{lei} é á \emph{consistencia das combinações sentimentaes}, ou mais claramente, é o \emph{instincto}. --- A \emph{liberdade} consiste em \emph{augmentar o numero dos motivos no espirito do individuo}, \emph{para que o conflicto}, \emph{que se dá entre elles}, \emph{o esclareça e o salve da omnipotencia de um motivo unico}. --- A \emph{justiça} reduz-se ao \emph{principio de identidade}, cuja intuição, porém, é \emph{irreductivel} ou \emph{incomprehensivel}. A idea de \emph{justiça} é o simples facto de uma evolução superior, e torna-se \emph{dever} emquanto assume a forma de uma \emph{inclinacdo intellectual}; porque, assim como a \emph{demonstração} nos obriga a consentir numa verdade, assim o \emph{dever} nos obriga á egualdade.\par €) \emph{Taine} propõe, nos seus livros (\emph{Les Philosophes francais du XIX siécle} --- \emph{ZLintelligence}), uma theoria \emph{moral}; mas esta theoria reduz-se a uma simples descripção de alguns parallelismos, que existem entre as sensações, as imagens, as ideas e os factos moraes. E' diffüicil achar ahi uma explicação, muito mais uma demonstração. Por isso é escusado insistir no assumpto.\par d) \emph{Stuart Mill}, imitando, ao mesmo tempo, \emph{Bentham} e \emph{Comte}, \emph{é econonista-positivista}. \emph{Bentham} tinha collocado o bem e o interesse na \emph{quantidade} do prazer. \emph{Stuart Mill}, para evitar as consequencias immoraes d'esse principio, introduz um novo elemento no prazer --- a \emph{qualidade}. \emph{A qualidade} do prazer á superior á sua \emph{quantidade}. O homem deve procurar e preferir sempre os prazeres, \emph{que valem mais}; porque « é melhor ser um Socrates descontente do que um animal satisfeito ». Os prazeres, que valem mais e devem ser preferidos, são os de \emph{sympathia}, que derivam do \emph{egoismo}, mas estão ligados com a felicidade dos outros. « Fazei aos outros o que querieis que se fizesse a vós: esta regra d'ouro de Jesus de Nazareth é a mais \emph{perfeita eeperiencia da Moral utilitaria} ». --- O que leva o homem a escolher sempre os prazeres de \emph{sympathia} é a \emph{associação das ideas}. A. educação, a leitura, o Juizo ácerca da conducta dos outros, formam no homem \emph{associações} poderosas, por que ligamos a alguns actos a idea de louvor e de admiração, e a outros a idea de censura e de desprezo, e que por isso produzem prazer, quando praticamos os da primeira especie, e causam remorso, quando produzimos os da segunda especie. --- E' assim que se explicam os phenomenos moraes, o imperio da consciencia, o prazer ou o remorso depois de termos : praticado o acto. E' assim que se prefere a morte \& ignominia e o sacrificio á cobardia. Tudo é objecto de caleulo e ex pliea-se pela \emph{associação}. A pratica constante do bem, embora no principio inspirada pelo interesse e pelo prazer, torna-se agradavel em si mesma; é então que o homem se torna \emph{virtuoso}. Assim como o avarento chega a amar o dinheiro pelo dinheiro; assim tambem o homem generoso chega a encontrar um prazer superior e virtuoso nos actos, que mais o honram deante da sociedade.\par e) \emph{Herbert Spencer} principia dizendo que a velha Moral, á qual se attribuia uma supposta origem divina, perdeu toda a sua auctoridade, porque era demasiadamente severa e não se accommodava ás modernas exigencias da humanidade. Mas é necessaria uma regra para a nossa vida. Deve, pois, fundar-se uma nova Moral, que, prescindindo de um legislador supremo,\pgn{586}seja mais facil e indulgente. Neste intuito, o escriptor inglez, applicando o \emph{darwinismo á Moral}, apresenta a seguinte theoria.\par A. \emph{conducta} é o complexo dos actos adaptados para o fim, --- ou, melhor, é a adaptação dos actos para o fim. O seu desenvolvimento é proporcionado ao desenvolvimento do organismo. Assim, imperfeitissima é a \emph{conducta} dos infusorios, que, sem direcção determinada, erram algum tempo, e morrem, quando as circumstancias do meio deixam de ser favoraveis. Menos imperfeita é \& \emph{conducta} dos vertebrados, dotados de sentidos mais desenvolvidos para a consecução do seu fim; mas, empregando poucos e simplicissimos meios, não tem vida longa. Mais perfeita é a \emph{conducta} dos mammiferos. Perfeitissima é a do homem, principalmente civilisado, que sabe adaptar meios mais numerosos e exactos, chegando a possuir os sentimentos de \emph{moralidade} e a formar a idea do \emph{bem} e do \emph{mal}.\par Como chega o homem a possuir os sentimentos de \emph{moralidade} e a idea do \emph{bem} e do \emph{mal}? --- Por meio de tres principios: do \emph{interesse pessoal}, do \emph{instincto social} e da \emph{hereditariedade}\par O homem \emph{primitivo} vivia num estado absolutamente selvagem. Então não havia moralidade, nem propriedade, nem mutua benevolencia, nem familia, nem religião. --- Mas este estado não podia durar. Na desordenada satisfação das proprias paixões, \emph{o} homem \emph{primitivo} muitas vezes provocava conflictos, e, ficando vencido, teve medo do mais forte e fez-se \emph{violencia}. Outras vezes, tendo accumulado grande quantidade de alimento, devoraya tudo num dia, e, como no dia seguinte experimentava os estimulos da fome, começou a ter a idea de \emph{prudencia}, que Ihe aconselhava a \emph{economia} e oppoz \emph{resistencia} aos proprios appetites. Assim o primeiro sentimento foi o \emph{egoismo}. O \emph{interesse pessoal}, pois, foi o primeiro factor da \emph{moralidade}.\par Mas o homem é sociavel por sua natureza. Como as abelhas e as formigas, os homens primitivos formaram \emph{énstinctivamente} uma sociedade. Formada a sociedade, o homem ficou sujeito a novos deveres. Porquanto, comprehendendo que o bem particular está indissoluvelmente ligado ao bem commum, concluiu que o individuo deve contribuir para o bem de todos e que tudo o que é \emph{util} para a sociedade é \emph{bem}, e tudo o que lhe é \emph{nocivo é mal}. A \emph{utilidade publica}: eis o principio da \emph{moralidade}, o criterio da distinecção entre o bem e o mal. Assim do \emph{egoísmo} nasceu necessaria e fatalmente o \emph{altruísmo}, com todas as suas formas mais elevadas, que são a generosidade, o sacrificio, a dedicação, etc. O \emph{egoismo} e o \emph{altruísmo} desenvolveram-se quasi ao mesmo tempo. O \emph{altruísmo}, no homem, torna-se \emph{consciente}, e abrange a familia (\emph{virtudes domesticas}) a sociedade (\emph{virtudes sociaes}), \emph{a} humanidade inteira (\emph{pAilantropia} e \emph{caridade}). O \emph{altruísmo} não só nasce do \emph{egoismo}, mas robustece-se com elle, até que o vence ou parece que o vence; porque a victoria do \emph{aliruismo} é tambem a do \emph{egoismo}, pois o bem da communidade é o bem de cada individuo. O \emph{egoismo} e o \emph{altruismo} harmonizar-se-hão numa sociedade futura, para a qual tendemos, pois são co-essenciaes, e formarão o \emph{ego-altruismo}. Do seu perfeito accordo dependerà a felicidade da humanidade \emph{ideal}, que a humanidade presente esta preparando. --- D'este modo a moralidade, se tem o começo no \emph{interesse pessoal}, tem o complemento no \emph{instincto social}.\par Mas não basta. Os juizos dos homens primitivos ácerca da \emph{utilidade} de\par ,\pgn{587}certas acções não foram isolados. Cada geração renovou as experiencias, e, ampliando-as e organisando-as, fixou-as tenazmente no cerebro e transmittiu-as, por \emph{hereditariedade}, sob a forma de \emph{modificações organicas}. Por isso as nossas \emph{intuições moraes}, sendo o resultado das experiencias de \emph{utilidade}, accumuladas pelos nossos antepassados e transmittidas por geração, têm todos os caracteres do \emph{instincto}, são independentes de toda a experiencia consciente. « As intuições moraes são o resultado das accumuladas experiencias de utilidade; organisadas gradualmente, passando de geração em geração, tornaram-se independentes da experiencia consciente. Assim como a intuição do \emph{espaço}, possuida por todo o ser vivo, é o resultado das experiencias organisadas e fortalecidas dos antepassados, que lhe transmittiram a propria organisação nervosa, lentamente desenvolvida, e não precisa de experiencia pessoal para ser definitiva e completa, mas praticamente é uma forma do pensamento, que parece inteiramente independente da experiencia; assim tambem as experiencias de \emph{utilidade}, organisadas e fortalecidas por muitas gerações humanas passadas, produziram modificações nervosas correspondentes, que continuamente transmittidas e accumuladas desenvolveram em nós algumas faculdades de intuição moral, algumas emoções correspondentes á conducta bón ou má, que não têm nenhuma base apparente nas experiencias individuaes » (\emph{Carta a Stuart Mill}).\par Proposto o criterio e o modo por que deve regular-se a \emph{conducta} do homem, vejamos o seu \emph{fim}. Este \emph{fim} é o bem, o prazer, a felicidade de cada um e de todos. E, como é o fim que especifica os meios, a conducta será \emph{bód}, se procurar prazeres ou sensações agradaveis, e será \emph{víciosa}, se não poder afastar as dores ou as sensações desagradaveis. A conducta chegará a ter sancção moral, quando as actividades, tornando-se cada vez menos militantes e mais industriaes, não precisarem de injustiças, nem de opposições mutuas, mas consistirem na cooperação, no auxilio reciproco e no desenvolvimento. A consciencia individual e social será bemaventurada, quando tender para um grau mais alto de felicidade, resultante da diminuição de afflicções ou do augmento de prazeres.\par f) \emph{Alexandre Bain} não diverge de \emph{Spencer} a não ser por uma analyse mais minuciosa do modo por que, em conformidade com os principios do \emph{darwinismo}, as gerações transmittem as experiencias e as modificações organicas, favoraveis ao principio de harmonia social.\par *\par Expostas as diversas phases do \emph{utilitarismo}, é necessario passar á sua\par \textasciicircum{} \textasciicircum{} critica, --- Para maior clareza, consideraremos primeiramente o \emph{utilitarismo} na sua \emph{essencia} e na sua \emph{divisão} em \emph{particular} e \emph{publico}, \emph{e} depois nas varias formas, que o expõem e defendem.\par Pnor. I. --- A utilidade não é o fundamento da moralidade.\par \emph{a}) Be a \emph{utilidade} fosse o fundamento da \emph{moralidade}, de modo que uma acção fosse /\emph{tonesta} por ser \emph{util}, e outra fosse \emph{ndão-honesta} por ser \emph{nociva}, seguir-se-hia --- que o \emph{parricidio}, a \emph{traição dos amigos e da patria}, e outros actos semelhantes seriam /\emph{ionestos}, se nos causassem alguma \emph{utilidade} ou \emph{interesse}, --- e que o \emph{amor da patria e do proximo}, \emph{á justiça}, a \emph{generosidade},\pgn{588}seriam actos \emph{não-honestos}, se nos trouxessem alguma \emph{perda} ou \emph{afllicção}. Ora estas consequencias repugnam á consciencia e ao senso commum dos homens. Logo a \emph{utilidade} não é o fundamento da \emph{moralidade}.\par \emph{b}) O bem \emph{util}, desejando-se, não por si mesmo, mas por amor de outro bem, tem razão de \emph{meio}; ao passo que o bem /\emph{onesto}, desejando-se por si mesmo, tem razão de \emph{fine}. Ora o \emph{meio} é essencialmente diverso do \emph{fim}. Logo o bem \emph{util} é essencialmente diverso do bem \emph{honesto}, e por isso nem tudo o que é \emph{util é honesto}. Logo a \emph{utilidade} não é o fundamento da \emph{moralidade}.\par \emph{c}) O fundamento da \emph{moralidade} deve ser uma coisa \emph{absoluta}, \emph{invariavel}, \emph{necessaria}; porque existem actos, que são absoluta, invariavel e necessariamente \emph{honestos}, independentemente de toda e qualquer vontade e circumstancia. Ora a \emph{utilidade} é uma coisa \emph{relativa}, \emph{variavel} e \emph{contingente}. Logo a \emph{utilidade} não é o fundamento da \emph{moralidade}.\par Pror. II. --- \emph{A} utilidade particular não é o fundamento da moralidade.\par \emph{a}) Se a utilidade \emph{particular} fosse o fundamento da \emph{moralidade}, todo o acto desinteressado e inspirado pela mais sublime abnegação, como a morte por causa da salvação da patria, seria \emph{immoral} e \emph{não-honesto}; como, pelo contrario, todo o acto baseado no \emph{egoismo}, como a traição da patria com o fim de enriquecer, seria \emph{moral} e \emph{honesto}. Ora estas consequencias são absurdas. Logo a utilidade \emph{particular} não é o fundamento da \emph{moralidade}.\par 5) Todos os homens julgam não haver coisa mais contraria á /\emph{onestidade} do que operar pela utilidade propria, contrariando a dos outros. Logo a utilidade \emph{particular} não é o fundamento da \emph{moralidade}.\par \emph{c}) Se a utilidade \emph{particular} fosse o fundamento da \emph{moralidade}, o homem seria o fim \emph{ultimo} dos seus actos e da sua existencia. Ora esta consequencia é absurda, como demonstrámos no Capitulo IL. Logo a utilidade \emph{particular} não é o fundamento da \emph{moralidade}.\par Pnor. III. --- \emph{A} utilidade publica não é o fundamento da moralidade.\par \emph{a}) Se a utilidade \emph{publica} fosse o fundamento da \emph{moralidade}, as acções, que não se referem ao bem publico ou que são exercidas por um individuo afastado da sociedade, seriam \emph{immoraes} e contrarias á \emph{honestidade}. Ora estas consequencias são falsas. Logo a utilidade \emph{publica} não póde ser o fundamento da \emph{moralidade}.\par 5) O fundamento da \emph{moralidade} deve ser uma coisa \emph{invariavel}, em todos os tempos e em todos os logares. Ora nada mais variavel do que a utilidade \emph{publica}, porque essa utilidade varia com a diversidade dos tempos e dos logares, e porque os juizos dos cidadãos ácerca da utilidade \emph{publica} são, ás vezes, contradictorios, e o que a uns parece \emph{util} para a patria, a outros parece \emph{nocivo}. Logo a utilidade \emph{publica} não é o fundamento da \emph{moralidade}.\par c) Se a utilidade \emph{publica} fosse o fundamento da \emph{moralidade}, todo o acto \emph{util} para o publico seria \emph{Zionesto}. Ora esta consequencia é desmentida pela consciencia humana. Assim, quando T'emistocles, para assegurar á sua patria a supremacia sobre as outras republicas da Grecia, fez a proposta de queimar as esquadras dos alliados, surtas no porto de Athenas, levantou-se Aristides e exclamou: « O projecto é \emph{util}, mas é \emph{injusto} ». Logo a utilidade \emph{publica} não é o fundamento da \emph{moralidade}.\par Tendo refutado o \emph{utilitarismo}, \emph{particular} e \emph{publico}, procuraremos agora\pgn{589}refutar as suas diversas formas, tendo sobretudo em vista o que ellas apresentam de especial ou de proprio.\par A) O \emph{utilitarismo} dos \emph{dinpiricos} refuta-se com as razões, com que na \emph{Anthropologia} provámos a falsidade do \emph{sensualismo}. Limitamo-nos a censurar a opinião de A. \emph{Smith}.\par A opinião de \emph{Adam Smith} é falsa pelas seguintes razões:\par 4) Não podem identificar-se duas coisas, das quaes uma emana de outra; aliás o effeito e a causa seriam uma e á mesma coisa, o que é absurdo. Ora a \emph{sympathia}, isto é, a aptidão que um acto tem de despertar benevolencia deriva da \emph{honestidade} do proprio acto; pois um acto é \emph{sympatico}, porque é \emph{honesto}. Logo a \emph{sympathia} não é o fundamento da \emph{moralidade}.\par \emph{b}) Além d'isso, a \emph{moralidade} é uma coisa \emph{absoluta} e \emph{invariavel}, de modo que um acto, que é \emph{honesto} para um, deve-o ser para todos. Ora, se a \emph{sympathia} fosse o fundamento da \emph{moralidade}, um e o mesmo acto seria, ao mesmo tempo, \emph{honesto} e \emph{ndo-honesto}, \emph{bom} e \emph{mau}, porque a \emph{sympathia} varia com as disposições de cada um, e por isso um acto, que é \emph{sympathico} para um, é \emph{antipathico} para outro. Ora isto repugna ao senso commum. Logo a \emph{sympathia} não póde ser o fundamento da \emph{moralidade}. --- E' verdade que \emph{Smith}, no caso da \emph{sympathia} ser \emph{erronea}, appella para um \emph{espectador imparcial}; mas este \emph{espectador imparcial}, em ultima analyse, não pode deixar de ser a \emph{razão}, e esta é e será sempre a regra da \emph{moralidade}.\par B) O systema de \emph{Bentham} contém muitas falsidades, --- E' falso que o unico principio dos actos humanos seja necessariamente o \emph{interesse} pessoal e sensivel; porque o homem, sendo dotado de uma alma espiritual, póde tender €, se não fór sensualista, tenderá para bens superiores e espirituaes e procurará a felicidade dos outros, mesmo com sacrificio do seu bem \emph{sensivel}. --- E falso que possa sempre construir-se uma \emph{arithmetica oral}, para se decidir se de um certo acto resultará um prazer de maior ou menor \emph{quanlitade}; porque prazeres ha que se não podem reduzir á mesma medida nem estar sujeitos ao mesmo criterio de calculo. --- E' falso que o \emph{dever} do homem consista em procurar o seu interesse ou prazer sensivel, e que por isso o homem mais honesto e virtuoso é o que procura o seu maior bem e interesse; porque então Nero e Sardanapalo deveriam ser mais dignos da nossa estima e admiração do que tantos heroes, que sacrificaram as commodidades e a vida pelo bem da sua patria. --- E' falso que o interesse \emph{particular} não possa ser contrario ao interesse \emph{commu}, e que entre um e outro exista necessariumente a mais perfeita harmonia; porque, embora o bem \emph{commum} possa produzir e produza muitas vezes o bem \emph{particular}, todavia é certo que o individuo póde, compromettendo o interesse e o bem \emph{commun}, obter grandes bens e interesses particulares. Por isso a harmonia, que deve existir entre o interesse particular e o commum, deve fundar-se num principio absoluto e mais elevado do que o interesse \emph{particular}.\par 0) Os systemas dos \emph{positivistas} são todos falsos.\par a) E' falso o systema de \emph{Comte} pelas seguintes razões:\par 4) Baséa-se em falso supposto. Na verdade, o systema parte do principio de que se devem rejeitar, como coisas chimericas, as ideas dos seres que escapam á experiencia sensivel, como são as ideas de \emph{liberdade}, de \emph{responsabilidade}\pgn{590}\emph{ moral}, de \emph{deveres absolutos}, de \textsc{Deus}. Ora este principio, como vimos, é falso e reduz o homem á condição dos animaes irracionaes.\par \emph{b}) O \emph{positivismo} funda o seu \emph{altruísmo}, isto é, os sentimentos de desinteresse e de benevolencia na \emph{sexualidade}. Ora uma vulgar experiencia attesta que esta paixão desenvolve o mais feroz egoismo, e que os animaes, que apresentam os mais bellos exmplos de \emph{altruísmo}, são os destituidos de \emph{sexualidade}. Assim, nos formigueiros, o cidadão mais \emph{alirwista} é um pobre \emph{neutro}, que não pertence a nenhum dos dois sexos.\par \emph{c}) \emph{Comte} reduz todos os deveres do homem a esta maxima: Vive \emph{para os outros}. Ora esta maxima, embora pareça generosa, é comtudo inacceitavel ; não só porque leva o homem a prescindir sempre dos seus direitos e deveres individuaes, o que repugna á dignidade humana, mas tambem porque, se abrange os deveres da humanidade, não abrange os da justiça.\par \emph{d}) A maxima de \emph{Comte}: Vive \emph{para os outros}, se fosse posta em pratica, seria causa de dissolução da sociedade. Na verdade, se cada um devesse sacrificar-se pelo bem do outro, ninguem poderia acceitar o sacrificio dos outros. Sendo assim, a lei do sacrificio universal daria em resultado a suppressão de todo o acto moral e social.\par e) Finalmente, a maxima de \emph{Comte} não póde impor-se a ninguem, como regra de conducta. O \emph{positivismo} exclue os \emph{deveres absolutos}. Pois. bem, se o individuo não estiver convencido de que um dever imprescindivel o obriga a \emph{viver para os uiros}, sacrificarà elle por ventura o interesse proprio ao interesse dos outros? Certamente que não. --- \emph{Comte} responde que o \emph{altruísmo}, pela \emph{educação} e pela \emph{sciencia}, terá \emph{fatalmente} preponderancia sobre o \emph{egoismo}. Mas esta resposta é inutil. A \emph{educação} e a \emph{sciencia}, se não estiverem subordinadas a principios transcendentes ou sobresensiveis, darão o resultado opposto, isto é, farão com que o \emph{egoismo} vença o \emph{altruísmo}. E' o que a experiencia de todos os dias attesta.\par b) E' falso tambem o systema de \emph{Littré} pelas seguintes razões:\par \emph{a}) E' \emph{insufficiente} e \emph{superficial}. --- E' \emph{insufficiente}; porque não explica muitos e diversos factos moraes, como são o exercicio da nossa liberdade, as ideas e os sentimentos religiosos, o amor desinteressado da verdade, do bello e do bem honesto. --- E' \emph{superficial}; porque limita-se a descrever o parallelismo, que existe entre alguns instinctos dos animaes e os nossos sentimentos moraes, --- parallelismo, que só prova a perfeita harmonia, estabelecida por \textsc{Deus} entre os diversos seres creados.\par \emph{b}) E' \emph{illogico} e \emph{falso}. --- E' \emph{illogico}; porquanto, \emph{Littré} considera como \emph{causa} de um phenomeno o que é apenas uma \emph{condição} ou \emph{circumstancia}, --- E' \emph{falso}; porque, querendo explicar tudo por meio dos dois instinctos da \emph{conservação} e da \emph{reproduccüo}, reduz a factos meramente sensiveis os actos da intelligencia e da vontade, que, como demonstrámos na \emph{Antropologia} (Sec. III), são \emph{espiérituaes}.\par \emph{c}) \emph{Littré} diz que as necessidades \emph{egoístas} e \emph{altruistas} da substancia viva são levadas ao cerebro pelos nervos, e que o cerebro, reagindo, produz uma paixão, um sentimento. --- Antes de tudo, notamos que uma \emph{necessidade} é uma \emph{negação}, é a \emph{ausencia} ou \emph{privação} de alguma coisa, é um \emph{nada}. Ora á curiosa a viagem de uma \emph{negação}, de uma \emph{ausencia}, que, apoiada aos nervos,\pgn{591}chega até ao cerebro!! --- Além d'isso, nenhum homem, que não tenha perdido o senso commum, póde admittir que a \emph{paixão} e o \emph{sentimento} se reduzem a uma acção do objecto e a uma reacção do cerebro (Cf. \emph{Anthr}., Sec. T1, cap. I e II) --- Finalmente, se as \emph{ideas} e os \emph{sentimentos} fossem necessidades \emph{egoistas} e \emph{altruistas}, elaboradas pelo cerebro, seguir-se-hia que todos os homens deveriam possuir as mesmas \emph{ideas} o os mesmos \emph{sentimentos}, porque todos tem as mesmas necessidades \emph{egoistas} e \emph{altruistas}. E, todavia, ha tão grande diversidade de ideas e sentimentos nos homens!\par \emph{d}) \emph{Littré} reduz a \emph{lei moral} ao \emph{instincto}. Ora esta reducção é simplesmente immoral. Se a \emph{lei moral} é um \emph{instincto}, um \emph{impulso da natureza}, é claro que a \emph{moralidade} consiste em seguir o \emph{instincto}. E todavia, o homem deve resistir a alguns instinctos, aliás será tido por \emph{immoral}, e, se póde seguir outros, deve fazel-o com \emph{moderação racional}. Para que elogiar tanto a amizade, a generosidade, a abnegação, o sacrificio, se estas coisas são apenas a satisfacção de uma necessidade?!\par e) A definição de \emph{liberdade}, dada por \emph{Littré}, contém o \emph{determinismo}. Ora este systema é falso, como demonstrámos na \emph{Anthropologia} (Sec. IT). --- Não menos falsa é a idea que \emph{Littré} faz da \emph{Justiça}. O eseriptor francez reduz a \emph{Justiça} ao \emph{principio da identidade}: \emph{A---A}. Sem extranharmos que um positivista admitta um principio \emph{absoluto} e \emph{metaphysico}, como é o da \emph{identidade}, \&dvertimos que, para justificar esta equação A---A, é necessario que tudo o que está contido no segundo A esteja tambem contido no primeiro. Mas, se a \emph{justiça} consiste numa simples equação de \emph{A----A}, os deveres dos outros, representados por um \emph{A}, são eguaes aos nossos representados por outro A. E, sendo assim, a obrigação, que os direitos dos outros produzirüo em nós, será egual á que em nós póde produzir o quadrado da /\emph{Aypothenusa} ou o \emph{binomio} de \emph{Newton}. Se \emph{Littré} tivesse unido \& idea de \emph{identidade} a idea de \emph{dever} ou de \emph{obrigação}, então, e só então, teria dado a verdadeira definição de \emph{justiça} (Cf. Bonniot, \emph{Les Malheurs de la Philospohie}, pag. 52 e seg.).\par 0) O systema de \emph{Taine} não merece uma refutação especial, depois do que temos exposto contra \emph{Littré}. Quem teye a ousadia de affirmar que « a virtude € o vicio são productos naturaes, como o assucar e o vitriolo », mostra não ter disposição para tratar d'uma sciencia tão delicada, como é a \emph{Moral}.\par d) O systema de \emph{Stuart Mill} é inacceitavel pelas seguintes razàes:\par 4) O escriptor inglez faz derivar do \emph{egoismo} a, \emph{sympathia}, do amor proprio o amor do proximo. Ora é impossivel que do amor do interesse e do prazer derive o desinteresse e o sacrificio. --- Nào póde negar-se que, em tempo de paz e em cireumstancias favoraveis, o individuo encontra o seu interesse pessoal na subordinação ás leis; mas seria absurdo pensar que, no tempo de guerra e num momento difficil, o amor do interesse pessoal e do prazer aconselhem a praticar, ás occultas e sem esperança de recompensa, actos heroicos de abnegação e de sacrificio. Semelhantes actos só podem derivar do amor de um bem espiritual e supremo.\par \emph{b}) \emph{Stuart Mill} afirma que o homem, pela pratica constante do bem, se torna virtuoso, e que chega a encontrar o seu maior prazer nos actos que mais o honram, assim como o avarento chega a amar o dinheiro pelo dinheiro. --- Esta explicação não é satisfactoria, Embora o homem chegasse\pgn{592}a ser virtuoso, a sua virtude, sem um principio superior de garantia e de seguranca, estaria exposta a todos os caprichos e desordens das paixões. --- Além d'isso, a virtude, se derivasse exclusivamente do egoismo e do prazer, que são coisas communs a todos, perderia todo o seu valor, não mereceria os louvores e a gloria dos outros, e assim deixaria de ser amada e praticada; assim como o avarento deixaria de amar e amontoar o dinheiro, se este perdesse o seu valor e deixasse de ser um meio de adquirir e de gosar.\par \emph{c}) \emph{A Moral} de \emph{Stuart Mill} é insufficiente. Na verdade, que obrigação tem o homem de preferir a \emph{qualidade} dos prazeres á sua \emph{quantidade}? Nenhuma. Uma alma nobre e delicada escolherá, talvez, os prazeres mais nobres e delicados, mas terá ella o direito de prohibir aos outros a escolha dos prazeres \emph{quantitativos} ? Não, com certeza. E, se não existe tal obrigação, poderá chamar-se mau o individuo, que escolhe os prazeres \emph{sensiveis}? Não. A. \emph{Moral} de \emph{Stuart Mill} é, pois, destituida de base.\par \emph{d}) Além d'isso, a \emph{associação} das ideas é um meio insufficiente para o estabelecimento da \emph{Moral}. Na verdade, o \emph{egoismo}, que se manifesta pelo amor desordenado dos prazeres, do interesse e da gloria, é, ás vezes, tão forte e violento, que, se não houver um obstaculo efficacissimo, arrastará, no seu impeto, a intelligencia, a vontade e todas as faculdades humanas para a consecução do seu fim. Ora este obstaculo efficacissimo não póde ser a \emph{associação} das ideas; porque o homem, dominado pela paixão, despreza os louvores e as censuras dos outros homens.\par e) A \emph{Moral} de \emph{Spencer} é falsa pelas seguintes razões:\par \emph{a}) Contém muitos suppostos, cuja falsidade foi já demonstrada no decurso d'este trabalho. O escriptor inglez suppõe --- que \textsc{Deus} não existe, --- que.o homem é a evolução de uma especie inferior, e por isso entre nós e o animal não existe uma differença \emph{especifica} e \emph{essencial}, --- que o estado primitivo do homem foi o selvagem, --- que as faculdades humanas se transformam, e de \emph{vegetativas} e \emph{sensitivas} se tornam \emph{intellectuaes}, --- que o \emph{pensamento} e a \emph{consciencia} são funcções da actividade nervosa, --- que por isso a \emph{alma} do homem não é espiritual, nem immortal. Ora todas estas supposições são falsas. E' falso que não existe um Creador e Conservador supremo de todas as coisas (\emph{Theodicea}, cap. I). E' falso que o homem é apenas a evolução de uma especie inferior (\emph{Anthr}., Sec. IV, cap. I). E' falso que o estado primitivo do homem foi o selvagem (\emph{Anthr}., Secção IV, cap. I, art. IT). E' falso que as faculdades \emph{intellectuaes} são o resultado da transformação das faculdades \emph{vegetativas} e \emph{sensitivas}; porque as faculdades :\emph{ntellectwaes} são \emph{espirituaes}, ao passo que as \emph{vegetativas} e as \emph{sensitivas} são \emph{materiaes}, \emph{e} repugna que o \emph{espiritual} derive do \emph{material} (\emph{Anthr}., Sec. II). E' falso que o \emph{pensamento} e a \emph{consciencia} são funcções da actividade nervosa (\emph{Anthr}., Sec. I, cap. II, art. IV). E' falso que a nossa vontade está sujeita ao \emph{determinismo}, \emph{i} que é uma negação do \emph{livre} «\emph{arbitrio} (\emph{Anthr}., Sec. II, cap. III, art. 11. E' falso que á nossa alma não é \emph{espiritual} (\emph{Anthr}., Sec. I, cap. II, art, LT), nem \emph{immortal} (\emph{Anthr}. Sec. IV, cap. IL, art. II).\par \emph{b}) \emph{Spencer} diz que o caracter \emph{moral} dos individuos é transmittido pela \emph{hereditariedade}. Ora isto é falso. Pela \emph{hereditariedade}, os paes podem transmittir aos filhos o temperamento, umas tendencias para um determinado\pgn{593}estado pathologico, umas predisposições inherentes ao sangue; mas não transmittem nem podem transmittir o que elles adquiriram pelo seu proprio esforço, nem o que depende da sua vontade, como a sciencia, a virtude, etc. Quem não sabe que de paes sabios nascem filhos ignorantes, e de paes virtuosos nascem filhos maus? --- E' verdade que os \emph{evolucionistas}, para. justificarem aquella transmissão, condemnada pela experiencia e pela razão, dizem que a \emph{moralidade} \& uma coisa gravada no cerebro, pois \emph{Gall} descobriu a bossa da moralidade. Mas a justificação é mais absurda que o principio. O que está unido ao cerebro é uma entidade material, tem grandeza e figura. Ora ninguem dirá que a idea de \emph{moralidade}, de \emph{justiça}, de \emph{bem} e de \emph{mal}, é uma coisa grande ou pequena, branca ou vermelha, quadrada ou circular. De resto, o systema de \emph{Gall} foi refutado na \emph{Anthr}. (Sec. I, cap. IL).\par €) \emph{Spencer} diz que o \emph{egoismo}, isto é, o amor de si mesmo, e o \emph{altruísmo}, isto é, 0 amor dos outros, são o \emph{fundamento da moralidade}, são \emph{co-essenciaes}, e do perfeito accordo entre elles depende a felicidade da humanidade \emph{ideal}. --- Antes de tudo, notamos que o \emph{egoismo} não se identifica com o \emph{amor de si mesmo}; porque o \emph{egoísmo} é sempre um vicio e nunca é um dever, e por isso o \emph{amor de si mesmo} não deve ser \emph{egoista}. --- Em seguida, perguntamos: Qual é o principio que póde estabelecer um perfeito accordo entre o \emph{egoismo} e o \emph{altruísmo}, entre o interesse \emph{particular} e o \emph{commum}? E' o \emph{meio social}: responde \emph{Spencer}. Mas estamos num circulo vicioso. O meio social, podendo formar-se só depois de se ter manifestado um perfeito accordo entre o \emph{egoismo} e o \emph{aliruismo}, \emph{é effeito} e não \emph{causa} do mesmo accordo. Qual será, pois, esse principio? Será o proprio \emph{egoismo}? não: porque o \emph{egoismo} tende sempre para o que é util ou agradavel ao individuo, e o que é util ou agradavel ao individuo é geralmente nocivo ou desagradavel á communidade. Será o \emph{altruísmo}? tambem não; porque o que convém á communidade é, muitas vezes, contrario no interesse do individuo. Logo o \emph{egoismo} e o \emph{altruismo} não podem ser os unicos factores do seu perfeito accordo, em que consistiria a felicidade da humanidade \emph{ideal}. --- Diremos, pois, a \emph{Spencer} que o amor honesto de si mesmo e o amor honesto dos outros se harmonizam perfeitamente, mas só no amor de Deus..E' o amor de \textsc{Deus} que nos leva a consecução da nossa felicidade eterna, pelo sacrificio das nossas paixões desordenadas, e que nos sujeita ao serviço do proximo e ao bem da sociedade. Este é unico fundamento capaz de estabelecer uma perfeita harmonia entre o interesse \emph{particular} e o \emph{commum}. Mas esta harmonia é apenas um fim \emph{procimo} e \emph{secundario}; porque o fim \emph{ultimo} e \emph{principal} do homem é Dus.\par \emph{d}) \emph{Spencer} admitte que a \emph{moralidade} é um effeito, uma consequencia da \textbar{} sociedade; porque esta é a mãe de todos os sentimentos moraes. --- Ora isto é falso. A \emph{moralidade}, longe de ser um effeito ou uma consequencia da sociedade, é o seu principio e fundamento. Por quanto, não póde haver sociedade, domestica ou civil, sem a auctoridade de um superior e sem a obediencia dos subditos. Ora o que é que leva os subditos a obedecerem ao superior? Será a \emph{sympathia}? não; porque ella fica vencida pelo \emph{egoismo}. --- Será a \emph{força material}? tambem não; porque a \emph{força material} só póde levar o subdito a obdecer ao superior, quando este a possue, e não quando não a tem, e porque deve haver um principio, que conserve a ordem e a disciplina\par 38 Vor. II.\pgn{594}no exercito. --- Serão os castigos e os premios? tambem não; porque os castigos e os premios suppõem a moralidade. Se, pois, quizermos assignar ao effeito uma causa proporcionada, devemos concluir que o principio, que leva o subdito a obedecer ao superior, é o sentimento \emph{moral} do \emph{dever}, é a \emph{moralidade}. Logo a \emph{moralidade} é o fundamento e o principio da sociedade, e por isso não é o seu effeito, a sua consequencia.\par e) \emph{Spencer} diz --- que a \emph{hereditariedade}, se transmitte a idea do bem e do mal, não transmitte a razão por que uma acção é boa, e outra é má, --- que os filhos, não conhecendo as causas da differenca entre o bem e o mal, deram ao bem e ao mal um caracter \emph{absoluto} e \emph{necessario} e erradamente disseram que deve fazer-se o bem, porque é bem, e que deve evitar-se o mal, porque é mal, --- que da mesma ignorancia nasceram as ideas de obrigação, de dever e de sancção, fundadas na existencia de \textsc{Deus}, --- e que por isso a moral \emph{absoluta} e \emph{religiosa}, sendo effeito da illusão, do prejuizo e da ignorancia, deve ser destruida pela sciencia. --- Poderiamos responder a \emph{Spencer} que a moral \emph{relativa} e \emph{athea} deve ser destruida, porque se funda na illusão, no prejuizo e na ignorancia. Mas limitamo-nos a notar que a opinião de \emph{Spencer} é destruida pela propria lei do \emph{evolucionismo}, em que ella se funda. Na verdade, segundo o escriptor inglez, a \emph{moralidade} nos antepassados era \emph{consciente}, pois elles conheciam a razão pela qual uma acção é boa e outra é má, e era verdadeira, porque correspondia á realidade; ao passo que, nos descendentes, é \emph{inconsciente}, é uma illusão, um engano. Sendo assim, é evidente que os antepassados eram mais perfeitos que nós, que o \emph{homo primigenius} era mais sabio que o \emph{omo sapiens}, e que por isso, na evolução das especies, não houve progresso, mas retrocesso. D'este modo cahe por terra o \emph{evolucionismo}, que admitte o progresso continuo e necessario; pois, neste caso, a \emph{selecção}, longe de accumular em volta da especie a perfeição das especies anteriores, teria eliminado uma perfeição existente.\par \emph{f}) \emph{Moral} de \emph{Spencer} é destituida de \emph{fim}. Porquanto o \emph{praser}, que \emph{Spencer} propõe como escopo da conducta humana, sendo passageiro, amissivel, enganador, não póde ser o \emph{fim ultimo}, como demonstrámos no Cap. I. --- E verdade que \emph{Spencer} propõe á humanidade \emph{actual}, como \emph{fim}, á humanidade \emph{futura}, ideal e perfeita. Mas, além de ser injusto e impossivel sacrificar a humanidade de hoje á de amanhã, essa humanidade \emph{ideal} e \emph{perfeita} é uma verdadeira chimera, em virtude do proprio principio da \emph{evolução}. Na verdade, se tudo está sujeito á lei fatal da evolução, se tudo progride necessariamente, não póde haver \emph{ideal} nem \emph{fim}. Nem fallem no \emph{equilibrio}, na adaptação perfeita do homem ao meio; porque, se tudo está sujeito á lei da evolução, tambem o \emph{equilibrio} e o \emph{meio} devem mudar. O proprio \emph{Spencer} reconhece “esta verdade, dizendo que « a adaptação da natureza humana ás condições da existencia nunca poderá attingir uma perfeição completa, e que ao periodo da \emph{evolução} succederá um periodo inverso --- o da \emph{dissolução} ». A felicidade, pois, é, segundo os evolucionistas, o fim \emph{necessario} dos actos humanos, mas é um fim que nunca será alcançado, é a agua a fugir sempre dos labios de Tantalos sequiosos! --- Concluimos que a Moral de \emph{Spencer}, não assignando ao homem um fim \emph{ultimo}, é destituida de base.\par \emph{g}) A \emph{Moral} de \emph{Spencer} não só é destituida de //\emph{m}, mas tambem de \emph{leis},\pgn{595}que dirijam a actividade livre do homem, marcando-lhe os direitos e os deveres. O escriptor inglez reduz as leis \emph{moraes} ás leis \emph{physicas}. Esta reducção repugna; porque o nome de \emph{lei} convém \emph{primeira} e \emph{propriamente} ás leis \emph{moraes}, e só \emph{secundaria} e \emph{analogicamente} ás leis \emph{physicas}; ora é absurdo reduzir o primeiro ao segundo, o proprio ao analogico. --- Além d'isso, tal reducção equivale a uma destruição de toda a Moral. Com effeito, se a lei \emph{moral} não se distingue das leis physiologicas, biologicas, ete., \emph{a} moralidade será \emph{organica}, será o fructo de experiencias accumuladas e transmittidas por \emph{hereditariedade}, \emph{e} haverá um orgão \emph{moral}, como ha um orgão \emph{visual}, \emph{auditivo}, etc. Mas, sendo assim, a \emph{Moral} desapparece. Uma lei, que nos impelle e leva necessariamente ao acto, não é uma lei que \emph{dirige}. E de que modo poderia ella dirigir a nossa conducta? Se o homem herda a virtude, como herda a saude, deve estar sujeito ao vicio, como o está á doença. Como não póde mudar de temperamento, de constituição, de orgãos, de physionomia; tambem não póde mudar de conducta. --- Não existindo verdadeiras leis \emph{moraes}, é claro que não póde existir verdadeira \emph{obrigação}, nem verdadeiro \emph{direito}. Por quanto, \emph{Spencer} diz que a \emph{obrigação} e o \emph{direito} são apenas experiencias desenvolvidas e transmittidas, são faculdades organicas, que nos impellem para os seus actos respectivos. Mas, se não ha verdadeira \emph{obrigação}, as tendencias \emph{individuaes} e \emph{egoistas} podem vencer as tendencias \emph{sociaes} e \emph{altruistas}, \emph{e} o equilibrio será quebrado e a sociedade dissolvida. Os \emph{evolucionistas} appellarao, talvez, para o \emph{ordem}; mas a \emph{ordem}, sem \emph{leis moraes} e por isso sem um Legislador soberano, é uma palavra và, e irrisoria. --- Esta Moral de \emph{Spencer} poderá ser uma \emph{physica}, uma \emph{chimica}, uma \emph{historia natural}; mas nunca será verdadeira \emph{Moral}, a \emph{Moral} que dirige a livre actividade do homem para o seu fim.\par \emph{A}) Se a \emph{Moral} de \emph{Spencer} não tem \emph{fim} nem \emph{leis} verdadeiras, nem póde ter uma \emph{sancção}. Se o8 bons e os maus não recebem, aqui sobre a terra, a recompensa ou o castigo, que lhes é devido, é claro que deve existir uma vida futura, em que \textsc{Deus} dará a cada um o que lhe pertence. Mas os \emph{evolucionistas}, supprimindo a vida futura, a existencia de \textsc{DEus}, a espiritualidade da nossa alma, não podem estabelecer uma \emph{sancção} da conducta humana. E aonde encontrar essa sancção?.. Na justiça humana? Não; porque a justiça humana é fallivel, está sujeita á corrupção, e não distingue sempre o bom do mau, o innocente do criminoso... Nos \emph{remorsos}? Tambem não; porque o remorso é uma censura, não é um castigo; é ministro mais da misericordia de \textsc{Deus} do que da sua justiça, e é por isso que o remorso diminue com o augmento dos crimes, o que não deveria acontecer, se elle fosse uma \emph{sanccüo}. Logo a \emph{Moral} dos \emph{evolucionistas} não tem verdadeira \emph{sancção}.\par f) O systema de \emph{Bain} é tambem inaboe vol e falso sobretudo pelas razões que apresentámos contra \emph{Spencer}.\par III. Autonomismo dos racionalistas. --- Este systema admitte que a razão deve formar a \emph{Moral}, \emph{independentemente} de todo e qualquer principio absoluto e objectivo, que regule a nossa actividade. --- O \emph{autonomismo} é uma derivação da \emph{Critica da} razão \emph{pura} de \emph{Kant}, em que o escriptor allemão pretende provar que a razão é independente da metaphysica e que os primeiros principios da razão são apenas formas «\emph{a-priori}, identificadas com a\pgn{596}propria razão e destituidas de todo o caracter objectivo e absoluto. Todavia \emph{Kant}, na sua \emph{Critica da} razão \emph{pratica}, não applicou, como veremos, á Moral os principios emittidos na \emph{Critica da razão pura}, porque reconheceu que os principios absolutos são indispensaveis para a ordem moral.\par Os defensores do \emph{autonomismo} dizem que a Moral não se deve fundar \emph{no sentimento}, nem na \emph{utilidade}, nem na \emph{força exterior}, mas unicamente em alguns \emph{principios racionaes}. Estes principios racionaes não são metaphysicos ou absolutos, nem subsistem em si mesmos e fóra de nós, mas encontram-se na analyse das \emph{relações logicas} da nossa natureza e das nossas faculdades. --- Entre as nossas faculdades deve existir uma ordem, uma subordinação, fundada na sua natureza e tendencia instinctiva. Ora, em nós, ha duas especies de faculdades, que, por causa da diversidade do seu objecto e do seu modo de operar, se chamam \emph{inferiores} e \emph{superiores}. Como é evidente, a razão julga que as faculdades \emph{inferiores} devem estar subordinadas ás \emph{superiores}. E' o primeiro principio da \emph{Moral}. --- Entre as faculdades \emph{superiores} encontramos a razão e a vontade. A relação entre estas faculdades exige que a vontade, na sua tendencia para o bem, seja esclarecida e regulada pela razão. E' o segundo principio da \emph{Moral}. --- Eis, pois, os dois principios fundamentaes da \emph{Moral}: subordinação perfeita das faculdades inferiores ás superiores, e da vontade á razão. --- E' assim que a ordem das relações \emph{logicas} funda a ordem das relações \emph{moraes}. Por isso á razão é \emph{autonoma}; ella encontra em si mesma a lei moral; ella é a propria lei moral; ella manda na vontade e em toda a nossa existencia.\par Este systema foi exposto e defendido principalmente por \emph{Vacherot} no seu livro --- \emph{Essais sur la Philosophie critique}. D'elle extrahimos o seguinte trecho: « Que importa á Moral se o universo se desenvolve, ou não, segundo um plano concebido pelo pensamento divino, se o mundo tende, ou não, para um escopo fixado pela mão divina? Quer o mundo seja a obra de um Dus bom ou de um genio mau, quer seja governado por uma providencia ou abandonado á fatalidade, o homem não deixa de ter a sua natureza, o seu fim, a sua lei, o seu direito e o seu dever, --- coisas que só podem ser determinadas pela psychologia e pela moral » (pag. 320).\par O \emph{autonomismo} dos racionalistas é falso pelas seguintes razões:\par a) Este systema rejeita todo e qualquer principio absoluto e objectivo. Ora esta rejeição é absurda; porque, como demonstrámos varias vezes, os primeiros principios da razão não são apenas formas a \emph{priori} da nossa intelligencia, mas são verdades absolutas e necessarias, que se impõem á nossa intelligencia e têm o seu ultimo fundamento na essencia divina.\par \emph{b}) O \emph{autonomáismo}, fazendo consistir a Moral na subordinação perfeita das faculdades inferiores ás superiores, admitte uma serie de fins \emph{relativos}, de bens \emph{particulares}, e não reconhece um fim \emph{absoluto}, um bem \emph{universal}, que domine e regule essa serie. Ora isto é absurdo; porque fins \emph{relativos} não podem existir sem um fim \emph{absoluto}, \emph{e} bens \emph{particulares} suppõem necessariamente um bem \emph{universal}. --- Como se vé, este systema admitte a Moral \emph{utilitaria}, porque, como dissemos, os fins \emph{relativos}, os bens \emph{particulares} têm razão de \emph{meios}, ou de bens \emph{uteis}.\par \emph{c}) O \emph{autonomismo} destroe a noção de \emph{dever}. A noção de \emph{dever} importa\pgn{597}necessariamente a noção de \emph{obrigação}. Esta \emph{obrigação} deve ter um \emph{fundamento}. Ora no systema, que eriticamos, a \emph{obrigação} não tem fundamento algum. E, na verdade, onde existe tal fundamento? Na razão, respondem. Mas isto é impossivel. Porquanto a razão só póde imperar na vontade, emquanto lhe indica os motivos capazes de se Ihe imporem; ora os motivos capazes de se imporem á vontade são unicamente os que se referem ao bem \emph{absoluto} e \emph{universal}; mas o \emph{autonomismo} só admitte bens \emph{relativos} e \emph{particulares}. Logo o \emph{autonomismo} destroe a noção de \emph{dever}.\par \emph{d}) O \emph{autonomismo} não só destroe a noção de \emph{dever}, mas tambem a de \emph{direito}. O \emph{direito} dos outros cria em nós deveres correspondentes. Mas, se a razão não descobre em nós nenhum principio de \emph{dever}, que se imponha aos nossos actos, tambem não póde descobrir esse principio nos outros. Sem a noção de um \emph{fim absoluto} e de um \emph{legislador supremo}, o \emph{direito} é destituido de \emph{fundamento}, e por isso de \emph{inviolabilidade}, isto é, deixa de ser \emph{direito}. Logo o \emph{autonomismo} dos racionalistas destroe a noção de \emph{direito}.\par IV. Voluntarismo dos juristas. --- Este systema admitte que a \emph{Honestidade} ou \emph{no-honestidade} dos actos humanos depende da livre vontade dos homens. Seus principaes representantes são \emph{Hobbes} e \emph{Saint}- \emph{Lambert}.\par a) \emph{Hobbes} diz que os homens, no principio, viviam no estado selvagem, abandonados aos seus caprichos e levados a operar unicamente pelo prazer e pela dór. Então os mais fortes obrigaram os mais fracos a formar uma sociedade, impondo-lhes algumas leis, que regulassem a sua actividade. D'este modo começaram a julgar \emph{honesto} o que era conforme com estas leis, e \emph{não-honesto} o que lhes era contrario. Estas leis estavam baseadas na \emph{utilidade}. --- Esta ultima proposição mostra a razão por que o systema de \emph{Hobbes} costuma reduzir-se tambem ao \emph{utilitarismo} ; porque, na verdade, para este escriptor, uma coisa 6 \emph{Aonesta}, quando é \emph{agradavel}.\par b) \emph{Saint-Lambert}, embora pertença á eschola dos \emph{utilitaristas empíricos}, ensina tambem o \emph{voluntarismo}, porque funda \emph{a honestidade} ou \emph{não-honestidade} dos nossos actos na \emph{opinião} dos povos; de modo que um acto é \emph{honesto} ou \emph{não-honesto}, conforme é tido como tal pela \emph{opinião publica}.\par Refutaremos estes dois escriptores, provando quo a \emph{moralidade} dos actos humanos não depende da \emph{lei civil}, nem da \emph{opinião} dos povos.\par á) A \emph{moralidade dos actos humanos não depende da lei civil}.\par \emph{a}) O fundamento da \emph{moralidade} dos actos humanos deve ser \emph{absoluto} e \emph{objectivo}; porque \emph{absoluta} e \emph{objectiva} é a \emph{moralidade} (\emph{intrinseca}) dos actos humanos. Ora a \emph{lei civil} é variavel e relativa ás opiniões dos povos. Logo a \emph{moralidade} dos nossos actos não depende da \emph{lei civil}.\par \emph{b}) Se a \emph{moralidade} dos actos humanos dependesse da \emph{lei civil}, tudo o que a \emph{ler civil} determinasse seria por isso honesto e santo. Ora esta consequencia é absurda; aliás não haveria nem poderia haver iniquidade, injustiça, tyrannia, que se não podesse tornar bôa e honesta. Logo á \emph{moralidade} dos actos humanos não depende da \emph{lei civil}.\par \emph{c}) Às leis humanas podem ser \emph{justas} e podem tambem ser \emph{injustas}, e neste ultimo caso deixam de ser leis e não obrigam. Ora isto demonstra que a \emph{lei civil} não é a norma unica da bondade e da justiça, mas que, acima\pgn{598}das leis civis, ha uma norma suprema, com que ellas se devem conformar, para que sejam \emph{justas}. Logo a \emph{moralidade} não depende da \emph{lei civil}.\par b) A \emph{moralidade dos actos humanos não depende da opinião}. \emph{dos povos}.\par \emph{a}) Se a \emph{moralidade} dos actos humanos dependesse da \emph{opinião} dos povos, deveria haver uma causa por que a \emph{opinião} julga alguns actos \emph{honestos} e outro \emph{não-honestos}; aliás dar-se-hia um effeito sem causa. Sendo assim, a \emph{moralidade} não depende da \emph{opinião}, mas sim da \emph{causa}, que leva a \emph{opinião} a julgar da \emph{moralidade} dos actos humanos. Logo a \emph{moralidade} dos actos humanos não depende da \emph{opinião} dos povos.\par \emph{b}) Além d'isso, á \emph{moralidade} não póde depender de um principio, que deve ser apreciado e julgado por ella; pois repugna que o superior dependa do inferior. Ora a \emph{opinião} dos povos é julgada pela \emph{moralidade}, e não viceversa; porque approvamos uma \emph{opinião}, que se conforma com a \emph{honestidade}, e reprovamos a que se lhe não conforma. Seria por ventura o \emph{parricidio} um acto \emph{bom} e \emph{honesto}, se a \emph{opinião} assim o julgasse? Não, com certeza. Logo a \emph{moralidade} dos actos humanos não depende da \emph{opinião} dos povos.\par c) Finalmente, o juizo ácerca da \emph{honestidade} de alguns actos (como a \emph{fidelidade}, o \emph{respeito aos paes}, o \emph{amor da patria}) e da \emph{ndo-honestidade} de outros (como a \emph{traição da patria}, o \emph{perjurio}, \emph{á blasphemia}, etc.) foi \emph{e} ê um e o mesmo em todos os povos. Ora este juizo não póde depender da \emph{opinião}, que é mutavel e diversa, conforme a diversidade dos povos, das suas inclinações e interesses; pois repugna que um effeito constante e universal dependa de uma causa variavel e particular. Logo a \emph{moralidade} dos actos humanos não depende da \emph{opinião} dos povos.}.

\pg{581} 50. A \textbf{moralidade intrinseca funda-se proximamente na conveniencia} ou \textbf{não-conveniencia} do \textbf{acto humano com} a \textbf{nossa natureza racional}.

a) A moralidade \emph{intrinseca} do acto humano consiste na sua \pg{582}  \emph{objectiva honestidade} ou \emph{não-honestidade}, independentemente de toda e qualquer prohibição ou preceito. Ora o acto humano \pg{583}  é \emph{objectivamente honesto} ou \emph{não-honesto}, quando convém ou não convém á nossa natureza \emph{racional}; porque \emph{honesto} é o que \pg{584}  convém á natureza \emph{especifica} de um ente. Logo a moralidade \emph{intrinseca}, funda-se \emph{proximamente} na conveniencia ou n\&o-conveniencia do acto humano com a nossa natureza \emph{racional}.

b) O fundamento da moralidade \emph{intrinseca} deve satisfazer ás seguintes condições: --- deve ser \emph{objectivo}, porque a moralidade \pg{585} \pg{586}  \emph{intrinseca} existe na realidade, independentemente do nosso modo de pensar e operar; --- deve ser \emph{absoluto}, porque a \pg{587}  differença entre o bem e o mal é invariavel e necessaria; --- deve ser \emph{fonte} de todas as regras de honestidade, como é evidente \pg{588} . Ora a conveniencia ou não-conveniencia do acto humano com a nossa natureza racional satisfaz a essas condições. Porquanto \pg{589}  aquella conveniencia ou não-conveniencia --- é \emph{objectiva}, porque não depende do nosso pensamento nem da nossa vontade \pg{590} ; --- é \emph{absoluta}, porque a nossa natureza racional é invariavel; --- é \emph{fonte} de todas as regras de honestidade, porque \pg{591}  o homem deve fazer ou deixar de fazer tudo o que convém ou não-convém á sua natureza racional. Logo a moralidade \pg{592}  \emph{intrinseca} funda-se \emph{proximamente} na conveniencia ou n\&o-conveniencia do acto humano com a nossa natureza racional.

Dissemos --- na sua conveniencia com a nossa natureza \emph{racional}, para significarmos que o acto honesto não é propriamente \pg{593} \pg{594}  o que convém á nossa natureza \emph{physica}, como é a força, a belleza, a indole, o talento, etc.; mas é tudo o que convém \pg{595}  á nossa natureza \emph{racional}, como \emph{tal}, --- quer o acto tenda para um bem \emph{espiritual}, que aperfeiçõa exclusivamente a alma, \pg{596}  como o acto da sciencia, da justiça, da caridade, do perdão das injurias e offensas recebidas, ete. ete., --- quer tenda par:

um bem \emph{material}, que aperfeiçõa o composto, comquanto elle se deseje e pratique por um modo \emph{ordenado}, como é o acto de repousar, de tomar alimento, de passear, com o fim de conservar \pg{597} \pg{598}  as forças e manter a propria vida no serviço de \textsc{Deus}, cumprindo d'este modo a sua soberana vontade \footnote{Esta tambem é a doutrina de S. Thomaz, que, quando quer dar a razão por que um acto é \emph{intrinsecamente} bom ou mau, recorre sempre á \emph{natureza racional}, com a qual alguns actos convém, e outros não convêm. (Cf. \emph{C}. \emph{Gent}., 1. ILI, c. 129; --- \emph{Sum}. \emph{Th}., 1-2, q. 54, a. 8; q. 94, a. 2). --- Por isso o fundamento \emph{provimo} da moralidade do acto humano não consiste --- a) nem na conformidade do mesmo acto com o \emph{fim ultimo} do homem, --- \emph{b}) nem na sua conformidade com a nossa razão.\par \emph{a}) Não consiste na conformidade do acto com o nosso \emph{fim ultimo}. Por quanto, o acto honesto leva á consecução do fim \emph{ultimo}, porque é honesto, mas não é honesto, porque leva á consecução do fim \emph{ultimo}. A \emph{honestidade} é uma propriedade \emph{intrinseca} do acto humano; ao passo que a sua \emph{relação} com o fim \emph{ultimo} é uma coisa \emph{extrinseca} ao proprio acto. Esta relação é \emph{effeito} da honestidade, e não \emph{causa}.\par \emph{b}) Não consiste na conformidade do acto com a nossa \emph{razão}. Porquanto, á nossa razão, assim como não é a medida da verdade \emph{ontologica} das coisas, assim tambem não o póde ser da bondade \emph{objectiva} dos actos. A razão conhece a honestidade, mas não a forma. --- De mais, o fundamento da \emph{noralidade} deve ser \emph{objectivo} e \emph{absoluto}; e a razão é \emph{subjectiva} e \emph{mutavel},\par Dizem: Se o fundamento da \emph{moralidade} consistisse na conformidade do acto com a \emph{natureza racional}, seria tambem verdadeira a opinião dos que collocam esse fundamento na conformidade do acto com a \emph{razão}; porque ser conforme com a natureza \emph{racional} é ser conforme com a razão.\par Resposta. Embora a razão indique o acto que nos é conveniente, todavia\pgn{599}esse acto é conveniente, não pela sua conformidade com o juizo da razão, mas pela sua proporção com a natureza racional; assim como, quando julgamos que o calor é para nós conveniente, esse calor não é conveniente porque assim a razão o julga, mas porque é proporcionado á nossa natureza.\par Continuam: Se o fundamento da \emph{moralidade} consistisse na conveniencia do acto com a nossa natureza racional, a Moral seria \emph{independente}.\par Resposta. A Moral seria \emph{independente}, se a conveniencia do acto com a nossa natureza racional fosse o fundamento \emph{remoto} e \emph{supremo}. Ora não é assim; porque essa conveniencia é apenas fundamento \emph{provimo} e \emph{subordinado} da \emph{moralidade} do acto humano, --- fundamento, que tira toda a sua força de um principio supremo --- de \textsc{Deus}.}.

\pg{599} 51. A moralidade intrinseca \textbf{funda-se remotamente na conveniencia} ou não-conveniencia do acto humano com a \textbf{natureza} divina. --- A moralidade \emph{Intrinseca} funda-se \emph{proximamente} na conveniencia ou n\&o-conveniencia do acto humano com a nossa natureza \emph{racional}. Ora toda a natureza creada, como é a do homem, funda-se na infinita natureza de \textsc{Deus}, como na sua causa \emph{exemplar}; de modo que uma coisa convém ou não convém a uma natureza creada, pelo facto de convir ou não convir á natureza divina. Logo a moralidade \emph{intrinseca} funda-se \emph{remotamente} na conveniencia ou não conveniencia do acto humano com a natureza divina \footnote{O fundamento \emph{ultimo} da \emph{moralidade} deve ser uma \emph{natureza} infinitamente perfeita e causa exemplar de todos os seres; de modo que todo o acto, que se conformar com este soberano archetypo de perfeição, será honesto, e todo acto, que se lhe não conformar, será não-honesto. Ora esta natureza não é a da creatura; porque esta é \emph{relativa}, \emph{finita} e \emph{contingente}; e todo o ente \emph{relativo}, \emph{finito} e \emph{contingente} tem a razão sufficiente da sua existencia no Ente \emph{absoluto}, \emph{infinito} e \emph{necessario}. Logo deve ser a natureza de \textsc{DEUS}. --- Por isso provámos que, embora o fundamento \emph{proximo} da \emph{moralidade} dos nossos actos seja a nossa natureza \emph{racional}, comtudo o seu fundamento \emph{ultimo} é á propria \emph{natureza} de \textsc{DEUS}, emquanto esta é modelo ou exemplar de toda a natureza creada, sobretudo da natureza \emph{racional}.\par Dissemos que o fundamento \emph{ultimo} da \emph{moralidade} é a \emph{natureza divina}, para excluirmos as opiniões, que encontram esse fundamento na livre \emph{von}- (\emph{ade} de \textsc{Deus}, ou na sua \emph{lei eterna}. Estas opiniões são falsas.\par \emph{a}) E' falsa á primeira opinião, que colloca o fundamento \emph{ultimo} da \emph{moralidade} na livre \emph{vontade} de \textsc{DEus}, e que foi defendida sobretudo por \emph{Pufendorf}. Na verdade, alguns actos humanos, por si mesmos e sem nenhuma relação com a \emph{livre vontade} de \textsc{DEus}, levam ao fim ultimo; e outros, tambem por si mesmos e sem relação com a \emph{livre vontade} de \textsc{DEUS}, nos afastam d'aquelle fim; porquanto, as essencias das coisas não dependem, como vimos, da \emph{livre vontade} de \textsc{Deus}. --- Além d'isso, se a moralidade \emph{ntrinseca} dependesse da \emph{livre vontade de} \textsc{Deus}, Elle poderia fazer com que um acto essencialmente mau fosse bom, e um acto essencialmente bom fosse mau. Ora isto é absurdo.\par \emph{b}) E' falsa a segunda opinião, que constitue o fundamento \emph{ultimo} da \emph{moralidade} na \emph{lei} eterna \emph{de} \textsc{Deus}. Na verdade, se assim fosse, um acto seria\pgn{600}honesto ou não-honesto, porque assim o julgou a sabedoria de \textsc{Deus}. Ora dá-se o contrario. Um acto não é bom ou mau, porque \textsc{Deus} assim o julga; mas \textsc{Deus} o julga bom ou mau, porque é tal em si mesmo; porque a essencia das coisas não têm o seu fundamento \emph{ultimo} na \emph{intelligencia} de \textsc{DEUS}. --- Além disso, a lei eterna de \textsc{Deus} suppõe os actos moralmente bons ou maus, e só lhes accrescenta a razão de preceito ou de prohibicao; porque, embora haja muitos actos que são maus por serem prohibidos, muitos outros ha que são prohibidos por serem maus.\par Os defensores da \emph{moral independente} dizem: A moralidade \emph{objectiva}, se não prescindir inteiramente de \textsc{Deus}, não será \emph{intrinsecu}, mas só \emph{extrinseca}, porque está sujeita a \textsc{Deus}. Ora a moralidade \emph{objectiva é intrinseca}.\par Resposta. A moralidade \emph{objectiva} do nosso acto deixaria de ser \emph{intrinseca}, se se fundasse exclusivamente no preceito divino. Mas não é assim. A moralidade não se funda exclusivamente no preceito divino, mas tambem, e sobre tudo, na essencia divina, que é o exemplar de todo o bem.\par *\par Depois de termos exposto a verdadeira theoria Acerca do fundamento \emph{ultimo} da «\emph{oralidade} do acto humano, resta-nos examinar algumas opiniões, que são diversas da nossa, e que por isso, embora fundem a moralidade num principio \emph{absoluto} e \emph{objectivo}, comtudo são falsas ou insufficientes. As principaes opiniões, diversas da nossa, podem reduzir-se a \emph{sete}; porque --- 1) alguns collocam o fundamento \emph{ultimo} da \emph{moralidade} na \emph{lei da evolução indefinida}, --- 2) outros, na \emph{lei da natureza} ou \emph{da razão}, --- 8) outros, na \emph{lei do respeito devido d pessoa ou d especie humana}, --- 4) outros, no \emph{imperativo categorico}, --- 5) outros, na \emph{idea de justiça}, --- 6) outros, na \emph{idea transcendental do bem}, --- \emph{V}) outros, na \emph{idea das relações essenciaes das coisas}. As quatro primeiras opiniões são falsas, as tres ultimas são incompletas ou confusas --- De cada uma faremos uma breve critica.\par A, Lei da evolução indefinida. --- Os \emph{pantheistas}, partindo do principio de que existe um só ente, eterno e necessario, que constante e fatalmente se determina, produzindo a serie dos diversos seres, concluem que todos os actos do homem, como tambem de todos os outros entes, são o producto de uma lei de evolução fatal, universal e indefinida, e que por isso esta lei é o unico e exclusivo fundamento da \emph{moralidade}.\par Não insistimos aqui na refutação d'este systema, cujos principios foram expostos e censurados na \emph{Cosmologia}. Só dizemos que o \emph{pantheismo}, negando a \emph{liberdade}, o principio de \emph{dever}, a distincção entre o \emph{bem} e o \emph{mal}, não só não aponta o fundamento da \emph{moralidade}, mas destroe toda a \emph{moralidade}. --- Feita esta summaria exposição e refutação do \emph{pantheismo}, resta-nos dizer alguma coisa ácerca das differentes formas, que este systema revestiu sobretudo nos ultimos tempos. Limitar-nos-hemos ás theorias de \emph{Spinoza}, \emph{Fichte}, \emph{Schelling}, \emph{Hegel}, \emph{Schopenhauer}, \emph{Hartmann}.\par a) \emph{Spinoza} diz o seguinte: « O ser eterno e infinito, que se chama \textsc{Deus} ou natureza, opéra com a mesma necessidade com que existe, sem liberdade, sem fim. Por isso a perfeição e a imperfeição, relativas a um fim, são noções\pgn{601}vãs, sem realidade; o bem e o mal são modos de pensar e de comparar as coisas. A convicção de que uma necessidade eterna liga a infelicidade á iyrannia das paixões e a felicidade ao desenvolvimento da razão, obriga fatalmente o homem a conformar a sua vida com os dictames da propria razão. Esta obrigação é a \emph{lei moral}. A. felicidade consiste em conceber todas ' as coisas sob a idea de eternidade e em communicar ás consolações da sciencia o caracter de eternidade. Comprehender a necessidade, por que tudo emana de \textsc{DEus}, é amar o proprio \textsc{Deus}, com amor intellectual ».\par A \emph{Moral} de \emph{Spinosa} é á do \emph{determinismo universal}. Nega a \emph{liberdade}, \emph{á finalidade}, a differenga \emph{intrinseca} entre o bem e o mal, identifica o mundo com \textsc{Deus}, nega a existencia pessoal e immortal do homem. Ora todas estas affirmações, como demonstrámos, são falsas e absurdas.\par b) \emph{Fichte} applica á Moral a concepção metaphysica do seu \emph{pantheismo}. Segundo este escriptor, o \emph{absoluto é o eu}. O eu \emph{absoluto} não é o que existe actualmente e que nos apparece sob uma forma individual, mas é o que ha de ser, o que tende para a universalidade e para a infinidade. A \emph{Moral} consiste em ter consciencia d'este eu verdadeiramente livre, e, por meio d'esta consciencia, desenvolver-se constantemente, crear o infinito, e, em seguida, levar o desconhecido, que nos limita, a um estado conhecido e consciente, e por esta forma realizar o \emph{ideal}. Este \emph{deal} chama-se \textsc{Deus}, que não é uma pessoa individua, mas é a suprema realidade, é o eu que lucta, que se desenvolve, que se aperfeicós até ao infinito.\par O systema de \emph{Fichte} é absurdo não só pelas razões que apresentamos na \emph{Cosmologia}, mas tambem porque destroe todos os fundamentos da \emph{moralidade}, que são a existencia de um Dus pessoal, a espiritualidade, a liberdade e a immortalidade da nossa alma, a existencia do mundo, a differenca \emph{intrinseca} entre o bem e o mal, etc.\par c) \emph{Schelling} admitte que o \emph{absoluto} se desenvolve em duas ordens oppostas: na ordem \emph{real} e na ordem \emph{ideal}. O \emph{absoluto} póde ser conhecido por uma intuição \emph{intellectual}, e esta é o principio de toda a \emph{sciencia}. Acima da Sciencia está a \emph{arte}, que é uma revelação mais profunda do \emph{absoluto}, em que a natureza e o pensamento se identificam. Superior á \emph{arte} é a \emph{moralidade} ou \emph{religião}. Esta consiste na resignação com que um ente deixa absorver as : suas forças individuaes pela vontade universal do \emph{absoluto}.\par Este systema tambem é absurdo como o de \emph{Fichte}.\par \emph{á}) \emph{Hegel} diz que o espirito, que é o termo da evolução da natureza, percorre tambem na sua evolução um triplice \emph{rhythmo}: a \emph{arte}, a \emph{religião} e a \emph{sciencia}. \emph{A arte} é a representação \emph{exterior} do \emph{absoluto}. \emph{A religião} é a representação \emph{interior} do mesmo \emph{absoluto}. Ao passo que a «\emph{ríe} se esforçava para unir á \emph{natureza} com \textsc{Deus}, a \emph{religio} procurava separala de \textsc{Deus}. D'ahi as differentes formas de religião: --- umas, como as orientaes, sacrificavam o natural ao divino, o finito ao infinito, --- outras, como o \emph{brahamanismo} e o \emph{budhismo}, julgavam o \emph{finito} uma \emph{illusão}, --- outras, como o \emph{monotheismo} dos semitas, pensavam que o \emph{finito é nada}, e que \textsc{DEus} é (\emph{udo}, --- outras, pelo contrario, como as dos Gregos e dos Romanos, sacrificavam o infinito ao finito. Este antagonismo devia acabar pela \emph{synthese}: d'ahi o \emph{Deus-homem} do\pgn{602}Christianismo. --- Mas a \emph{religião}, por não ser uma fórma perfeita do espirito, deve ser substituida pela \emph{sciencia}; é esta que deve dominar.\par Como se vé, o systema de \emph{Hegel} é uma mistura de \emph{pantheismo}, de \emph{determinismo}, de \emph{atheismo}. Se o fundamento da \emph{moralidade} é á \emph{sciencia}, o mais douto será o mais honesto: ou, melhor, tudo será honesto e divino, pois o principio da evolução do \emph{pantheismo} é o \emph{absoluto}. Ora isto é absurdo.\par e) \emph{Schopenhauer} (1188-1860) propoz uma fórma de \emph{pantheismo}, contraria á de \emph{Hegel} e favoravel á de \emph{Kant}. Rejeitando a intuição intellectual do \emph{absoluto}, \& \emph{metaphysica} e a \emph{theodicea}, \emph{Schopenhauer} pretende basear-se unicamente na experiencia \emph{total} e \emph{systematica}. A base da Moral é, para elle, a \emph{vontade universal}, \emph{indestructivel}, dotada de \emph{liberdade absoluta} mas \emph{invariavel}. O \emph{dever absoluto} e \& \emph{lei de liberdade} são contradicgóes. --- Esta \emph{vontade}, desenvolvendo-se, torna-se \emph{conhecimento}, pois que conhece a sua natureza e o mundo. O mundo apparece \emph{essencialmente mau}, porque a vida é soffrimento. O verdadeiro mal é, pois, o desejo da vida, é e \emph{egoísmo}. --- Para o homem ser \emph{honesto}, é necessario destruir o \emph{egoismo}, o desejo da vida, e sobre estas ruinas levantar o \emph{monismo} ou a \emph{unidade do universal}. Este \emph{monismo} traduz-se por duas tendencias moraes, que são a \emph{caridade universal} e a destruição do \emph{egoismo} pelo \emph{ascetismo} e pela \emph{castidade}.\par O systema de \emph{Schopenhauer} contém muitas falsidades. --- E' falso que a vontade da humanidade seja \emph{invariavel}; porque a experiencia, em que este escriptor quer unicamente basear as suas asserções, attesta o contrario. --- E' falso que a unica realidade é a \emph{vontade}, e que esta é a mesma em todos os homens e, pela evolução, se torna \emph{conhecimento}. --- E' falso que desejar é soffrer, e que a vida é o soffrimento. --- E' falso que a \emph{moralidade} consiste na caridade para com todos os homens, considerados na identidade de uma vontade universal, e na destruição do desejo da vida.\par f) \emph{Hartmann}, corrigindo o seu mestre \emph{Schopenhauer}, inventou a philosophia do \emph{inconsciente}. O principio supremo, que para \emph{Hegel} é a \emph{idea} e para \emph{Schopenhauer} é a \emph{vontade}, é dotado, segundo \emph{Hartmann}, de dois attributos parallelos: de \emph{vontade infinita} e de \emph{razão infinita} mas \emph{éncosciente}; e por isso o principio supremo é a \emph{vontade intelligente}, mas \emph{incosciente}. Estes dois attributos, estando em continua lucta, produzem o mundo com todos os phenomenos, que são as diversas vontades \emph{particulares}. Estas diversas vontades hão de desapparecer, quando os homens conhecerem que o desejo da vida individual é uma loucura, pois a vida é um continuo soffrimento. O \emph{dever}, pois, é a propria \emph{rasdo}, emquanto se oppõe a este desejo da vida. O fim da humanidade é o \emph{suicidio collectivo}, é á victoria do \emph{universal inconsciente}.\par Como é evidente, o systema de \emph{Hartmann} é uma nova edição do systema de \emph{Schopenhauer}, nao menos esteril, falso e arbitrario que o do mestre.\par IL Lei da natureza ou da razão. --- Os \emph{Stoicos}, embora rejeitassem a noção de um \textsc{Deus} pessoal e a immortalidade da nossa alma, comtudo admittiam alguns principios universaes e immutaveis, reguladores da razão e da natureza. O bem, segundo elles, consiste na rectidão immutayvel da vontade, e esta rectidão consiste na conformidade com as \emph{leis essenciaes da natureza ou da rasão}. D'este modo uma acção é boa, quando se conforma com essas leis, e é má, quando lhes é contraria. Por isso os bens exteriores, como a\par \emph{L}\pgn{603}vida, as riquezas, os prazeres, as honras, são bens falsos; e a morte, as doenças, os soffrimentos, os desprezos, não são verdadeiros males. A virtude consiste na independencia completa de todo o influxo externo no reino perfeito da razão. Para viver conforme á razão, é necessario \emph{so/frer} todos os males e \emph{abster-se} de tudo o que nos póde inquietar.\par Este principio dos \emph{Stoicos} não póde legitimar um só dever moral. Na verdade, para essa legitimação, seria necessario proyar que temos obrigação de desprezar todos os bens e males da terra, e de sacrificar os nossos interesses ao imperio da razão, sem nos importarmos de outras vantagens superiores nem do fim do nosso sacrificio. Mas, se \textsc{Deus} não existe e se a nossa alma não é immortal, esta obrigação não existe, e não se vê a razão por que o homem deve resistir ás tendencias da sua natureza, desprezando os bsz*da terra e a satisfação das proprias paixões. --- E' claro que este systema é muito semelhante ao \emph{autonomismo} dos racionalistas, e refuta-se com as mesmas razões com que este foi refutado.\par III. Lei do respeito devido á \textbf{pessoa} ou á especie humana. --- Alguns sequazes do \emph{Stoicismo} substituiram á \emph{lei da natureza}, ou \emph{da} razão, a lei do \emph{respeito}, \emph{devido d pessoa ou d especie humana}, dizendo que um acto é \emph{bom}, se traduz esse respeito, e é \emph{mau}, no caso contrario.\par Nem esta opinião é verdadeira. O respeito devido a si mesmo, o respeito devido á pessoa e á especie humana, são verdadeiros deveres, mas não podem constituir o fundamento da \emph{moralidade} dos nossos actos. Por quanto, se prescindirmos de \textsc{Deus}, não é possivel impor ao homem o respeito para comsigo nem para com os seus semelhantes.\par IV. Imperativo categorico. --- E' o celebre systema de \emph{Kant}. Os fundamentos da \emph{Moral} do escriptor allemão não são os \emph{principios}, mas sim os \emph{factos moraes} da nossa consciencia. Estes factos não se demonstram ; analyzam-se, não para se determinar por um modo racional a sua realidade objectiva, mas para se conhecerem os seus caracteres subjectivos; pois a fé \emph{moral} deve substituir a sciencia. --- Ha duas especies de \emph{factos moraes}: \emph{maxvimas} e \emph{imperativos}. \emph{Às maacimas} são regras de conducta, essencialmente individuaes e particulares. Os \emph{imperativos} (ordens) são regras universaes, applicaveis a todo o genero humano. --- Os \emph{imperativos} dividem-se em \emph{Aypotheticos e categoricos}. --- Os \emph{hypotheticos}, que se chamam tambem \emph{materiaes}, são regras que se seguem por causa da \emph{materia} do preceito, como: \emph{Homem}, \emph{trabalha}, \emph{para não soffres a fome}. Estes \emph{imperativos} só têm força de preceito na /\emph{nypothese} de que se queira o effeito. --- Os \emph{categoricos}, que se dizem tambem \emph{absolutos} ou \emph{formaes}, são regras que se devem seguir por causa de um \emph{principio universal}, que constitue a \emph{forma} do preceito. Estes \emph{imperativos} têm sempre a força de \emph{preceito}.\par Todos os imperativos \emph{categoricos} podem reduzir-se a este preceito absoluto: \emph{Opéra por um principio de ordem universal}. Este preceito, que é \emph{formal}, \emph{universal}, \emph{necessario} e \emph{inmutavel}, constitue a lei completa e unica da vontade humana. Por quanto, a lei moral não póde conter nenhum elemento empirico ou experimental, mas deve ser \emph{universal}, \emph{necessaria}, \emph{immvutavel}. Ora só aquelle imperativo \emph{categorico} é dotado de taes caracteres. Logo a lei moral é uma \emph{forma a-priori}, producto da razão \emph{pura}, mas emquanto\pgn{604}é \emph{pratica}, pois manda na vontade. Esta forma identifica-se com a razão \emph{purapratica}, como \emph{funcção} ou \emph{condição necessaria dos nossos actos moraes}.\par D'este \emph{imperativo categorico} derivam duas consequencias logicas:\par 1.º) A razão \emph{pratica} é a unica fonte de toda a honestidade e de toda a virtude. Ella produz o uso perfeito da liberdade. O uso da liberdade é perfeito, quando a razão e a vontade são independentes de todas as influencias externas, materiaes, mutaveis e particulares, \emph{e} quando o homem se regula por esta \emph{forma pura}, identificada com a nossa razão.\par 2.º) Dos caracteres da lei \emph{pratica} podem deduzir-se, pela analyse, a \emph{liberdade}, a \emph{existencia} de \textsc{Deus} e a \emph{immortalidade} da \emph{alma}. --- Póde deduzir-se a \emph{liberdade}; porque o \emph{imperativo categorico}, contendo um preceito e uma obrigação, suppõe uma vontade \emph{livre}. Póde deduzir-se a \emph{existencia} de \textsc{Deus}; porque o nosso bem apparece ligado com o bem universal, e este nexo suppõe uma causa suprema e omnipotente, que é \textsc{DEus}. Póde deduzir-se a \emph{espiritualidade} da \emph{alma}; porque a lucta, que existe entre o bem honesto e o nosso bem particular, mostra que o estado actual é provisorio e que o estado definitivo 80 comecará depois da morte; o que suppõe a \emph{immortalidade} da \emph{alma}. --- Estas consequencias, porém, não são demonstrações \emph{absolutas} e \emph{objectivas}, mas derivam logicamente de um \emph{facto} da consciencia, --- facto, que é apenas uma \emph{forma subjectiva} da nossa razão.\par O systema de \emph{Kant} é falso pelas seguintes razões:\par a) A ordem \emph{moral} não póde conceber-se, se as noções de \emph{bem}, de \emph{fim}, de bem \emph{supremo}, de fim \emph{ultimo}, não tiverem um valor \emph{real} e \emph{objectivo}, independentemente da nossa razão e vontade. Ora, no systema de \emph{Kant}, essas noções são meramente \emph{subjectivas}, são objectos de fé e não de \emph{sciencia}. Logo, no systema de \emph{Kant}, é impossivel a concepção da ordem \emph{moral}. --- Nem se diga que este philosopho baséa a ordem \emph{moral} na \emph{existencia} de \textsc{DEUS} e na \emph{immortalidade da alma}. Com effeito, no systema de \emph{Kant}, a necessidade de admittir estas duas verdades deriva da necessidade de conciliar a \emph{forma pura} e \emph{universal} da lei moral com a sua \emph{materia}, que é o bem particular. Ora esta necessidade não é \emph{absoluta} e \emph{demonstrativa}; porque, segundo \emph{Kant}, a lei moral é \emph{formal} e manda absolutamente, sem se importar do escopo ou da felicidade a conseguir. Se essas duas verdades são \emph{relativas} e \emph{hypotheticas}, não são a base da ordem \emph{moral}. Alem d'isso, essas verdades são, no systema de \emph{Kant}, \emph{consequencias logicas} da \emph{moralidade}, e não o seu \emph{principio}.\par \emph{bj} O \emph{imperativo categorico}, que se resume nesta formula: « ÉE' \emph{necessario operar conforme uma regra universal}, \emph{independente de todo o elemento variavel e particular}, é um principio gratuito e não justificado perante a razão. Na verdade, uma regra \emph{universal}, que possa servir a todos os homens, devo basear-se nas relações \emph{essenciaes}, que existem entre todos os homens. Mas, so ordem \emph{objectiva} é, segundo \emph{Kant}, duvidosa, é claro que essas relações não podem ser conhecidas. --- Além d'isso, é falso que o acto \emph{moral}, para ser honesto, deva ser absolutamente desinteressado; porque, se a esperança de uma recompensa estiver subordinada ao amor supremo do bem, o acto humano não póde deixar de ser dotado de honestidade.\par \emph{c}) Segundo \emph{Kant}, a lei moral, sendo uma \emph{forma a-priori}, é o producto de\pgn{605}uma simples necessidade psychologica e subjectiva, e \emph{o} cumprimento perfeito da lei moral identifica-se com a perfeita autonomia ou independencia da razão e da vontade. Sendo assim, é evidente que toda a Moral se reduz a uma simples subordinação das faculdades inferiores ás superiores, e que o juiz d'esta subordinação é o proprio eu. D'este modo, sendo a razão e a vontade regras de si mesmas, todo o acto, por ser objecto da vontade, é legitimo e honesto, e por isso o fundamento da \emph{moralidade} é a \emph{autonomia} ou \emph{liberdade} da vontade. Ora estes principios e estas consequencias são outros tantos absurdos. E' absurdo que a lei moral se identifica com a nossa razão e a vontade; porque o legislador é necessariamente superior e não egual a quem está sujeito á lei, aliás não póde conceber-se a idea de \emph{obrigação}. E' absurdo, como demonstrimos contra os \emph{awfonomistas}, que a Moral consiste na subordinação das faculdades inferiores ás superiores. E') absurdo que um acto é bom pelo facto de ser o termo da vontade; porque a bondade é uma propriedade objectiva das coisas. E' absurdo que a liberdade constitue a moralidade; pois a \emph{liberdade} é apenas uma \emph{condição} de moralidade.\par \emph{d}) \emph{O imperativo categorico} é, para \emph{Kant}, um juizo \emph{synthetico}- \emph{a-priori}; pois a razão \emph{pratica} une a idea de \emph{lei moral}, que é uma \emph{forma a-priori}, com o acto humano, que é um elemento experimental, não contido na essencia da lei \emph{moral}. Ora os juizos \emph{syntheticos a-priori}, como vimos, são absurdos.\par \emph{V}. Lei da justiça, --- \emph{Julio Simon} admitte a \emph{justiça} como fundamento de \emph{moralidade}. A idea de \emph{justiça} é simples, natural, innata, necessaria, absoluta, immutavel, soberana. E' conhecida por todos os homens e applicavel a todos. A \emph{justiça} existe, não em si mesma, porque é um attributo, mas num ser necessario, eterno, immutavel, porque esse attributo é necessario, eterno, immutavel. A \emph{justiça}, pois, é um attributo da substancia divina. A idea abstracta da justiça é \emph{innata}; é uma das formas sob as quaes \textsc{Deus} se manifesta ao nosso espirito. --- Por isso o principio da \emph{obrigação moral} é a idea \emph{abstracta} da \emph{justiça} absoluta e infinita, e este principio influe efficazmente na nossa vontade, emquanto existe nella um amor correspondente.\par Esta theoria de \emph{J}. \emph{Simon}, embora elevada e eminentemente espiritualista, é comtudo insufficiente e inacceitavel pelas seguintes razões:\par \emph{a}) Suppõe a idea \emph{innata} ou a \emph{intuição} da justiça divina. Esta supposição é falsa. Porquanto, a idea primitiva, que temos da \emph{justiça}, é uma idea geral e deduzida da analyse dos actos humanos, --- e a idea, que temos da \emph{justiça divina}, é deduzida, composta e analoga.\par \emph{b}) O fundamento da \emph{moralidade} deve ser a fonte de todas as regras dos actos humanos; Ora uma idem \emph{abstracta}, como é a de \emph{justiça}, á insufficiente para determinar a ordem dos nossos deveres. Logo a idea \emph{abstracta} de \emph{justiça} não póde ser o fundamento \emph{ultimo} da \emph{moralidade}.\par \emph{c}) O fundamento \emph{ultimo} da \emph{moralidade} deve ser tambem principio efficaz de \emph{obrigação}. Ora a idea \emph{abstracta} de \emph{justiça} não pôde, só por si, determinar efficazmente a nossa vontade á pratica dos seus deveres. Logo a idea \emph{abstracta} de \emph{justiça} não póde ser o fundamento \emph{ultimo} da \emph{moralidade}.\par VI. \textbf{Idea} transcendental do bem. --- Outros escriptores dizem que o fundamento \emph{ultimo} da \emph{moralidade} dos nossos actos é a idea \emph{transcendental} do\pgn{606}\emph{bem}; de modo que um acto é bom, quando exprime esta idea, e é mau, quando a não exprime.\par Tambem esta opinião é insufficiente. Porquanto, a, idea \emph{transcendental} do \emph{bem}, sendo muito vaga, precisa de ser completada e determinada por outro elemento proprio e especial. Além d'isso, o \emph{bem}, tomado em toda a sua extensão, marca a esphera em que se move necessariamente á nossa vontade, e por isso não póde dirigir esta faculdade.\par \textbf{VII}. \textbf{Idea} das relações essenciaes das coisas. --- Outros "UI de dif-\par ferentes escholas dizem que o fundamento \emph{ultimo} da \emph{moralidade} se encontra nas \emph{relações essenciaes} das coisas, de modo que um acto é bom, se conserva aquellas relações, e é mau, no caso contrario.\par Tambem esta opinião é inaeceitavel. As \emph{relações essenciaes} das coisas são de varias especies: \emph{logicas}, \emph{estheticas}, \emph{moraes}. As \emph{logicas} fundam a ordem \emph{logica}; as \emph{esthelicas}, á ordem \emph{esthelica}; as \emph{moraes}, a ordem \emph{moral}. Por isso esta opinião aponta o \emph{genero}, quando devia apontar a \emph{especie}; porque, na verdade, ninguem dirá que um artista offende a ordem \emph{moral}, se fizer uma obra deforme. --- Além d'isso, é necessario, para a \emph{moralidade} dos nossos actos, conservar as relações \emph{moraes}, que ligam essencialmente as creaturas ao Creador. --- (Cf. Castelein, \emph{Les principes de la Philosophie Morale}).\par Concluindo a critica dos systemas, que admittem a Moral \emph{independente} de \textsc{Deus}, fazemos nossas as seguintes palavras de \emph{Sc/Aérer} (Art. do \emph{Temps}, á de out. de 1884): « Saibanios ver as coisas como são. A Moral, a Moral verdadeira, boa, antiga, imperativa, precisa do \emph{absoluto}, aspira á transcendencia, não encontra o seu ponto de apoio senão em \textsc{Deus}. A consciencia é como o coração; precisa de um \emph{au delá}. O dever, que não é sublime, é nada; e a vida torna-se vã, se não está relacionada com a eternidade. A Moral, como á Religião, é o sobrenatural; pois « \emph{Moral}, \emph{que não é religiosa}, \emph{é nada}. O sobrenatural é a esphera natural da alma ». O testemunho não podia ser mais eloquente nem mais insuspeito, embora o escriptor, obedecendo aos seus instinctos de sceptico, duvide da possibilidade de alcançar o \emph{absoluto}.}.

\pg{600} 52. \textbf{Corollario}. --- Das doutrinas expostas resultam os seguintes importantes corollarios:

I. A \emph{regra} de \emph{toda á moralidade}, pela qual extremamos o \pg{601}  acto bom do acto mau, é a seguinte: \emph{O teu acto seja conforme com a tua natureza racional}. Com effeito, a regra da \emph{moralidade} deve ter os seguintes requisitos: --- deve ser \emph{universal}, isto é, \pg{602}  applicavel a cada um dos actos, aliás não seria a regra de \emph{toda} a moralidade; --- deve ser \emph{suprema}, de modo que \emph{todas} \pg{603}  as outras se reduzam a ella; --- deve ser \emph{evidente}, porque, se devesse ser demonstrada por outra regra, deixaria de ser \emph{suprema} \pg{604} ; --- deve fundar-se na propria natureza \emph{especifica}. Ora a regra sobredita tem estes requisitos. Na verdade, --- é \emph{universal}, isto é, applicavel a todo o acto, porque todo o acto, \pg{605}  que convém á nossa natureza \emph{racional}, é por si bom, e todo o acto, que não convém, é por si mau; --- é \emph{suprema}, e della dependem todas as outras regras, porque a ultima razão da \pg{606}  bondade ou malicia moral do acto é a sua conveniencia ou não-conveniencia com a nossa natureza racional; --- é \emph{evidente}, porque nada mais claro do que a obrigação de operarmos em conformidade com a nossa natureza \emph{racional}; --- funda-se na propria natureza \emph{especifica} do homem; pois esta, \emph{especificamente}, \emph{é racional}. Logo essa regra de \emph{moralidade} é verdadeira \footnote{As opiniões dos philosophos ácerca da \emph{regra} de \emph{moralidade}, são muitas e diversas, conforme a multiplicidade e diversidade das suas theorias Acerca do \emph{fundamento} da propria \emph{moralidade}; porque é o \emph{fundamento} que nos dá o criterio para julgarmos da qualidade dos actos humanos. Assim, para os \emph{materialistas}, a \emph{regra} de \emph{moralidade} é a seguinte: Vive \emph{para as tuas satisfações} ; para os \emph{utilitaristas}, é a seguinte: Vive \emph{para a tua utilidade ou para a dos outros}; etc. --- Para evitarmos repetições inuteis, limitar-nos-hemos a examinar apenas uma ou outra d'essas regras.\par a) \emph{Krause} dá a seguinte regra: \emph{O homem deve realizar no tempo infinito toda} a \emph{sua essencia e todas as suas propriedades eternas}, \emph{e}, \emph{em cada instante}, \emph{uma parte determinada da sua essencia} n. (Tiberghien, \emph{Esquisse de phil}. \emph{moral}).\par Esta formula contém muitos erros, consequencias do erro fundamental,\pgn{607}que é o \emph{pantheismo}, e por isso, longe de conter a regra da \emph{moralidade}, é á negação de toda a regra. Com effeito, não se concebe moralidade sem livre .arbitrio. Ora a tal realização da essencia, sendo fatal e necessaria, destroe toda a liberdade. --- Alem d'isso, a lei suppõe o superior, que a promulgue, e o subdito, que a execute. Mas, se existe uma só substancia, infinita e eterna, não póde haver subordinação do inferior ao superior.\par b) \emph{Kant} propoz a seguinte regra: \emph{Opéra por forma que o teu acto possa ser norma de todo} o \emph{agente livre}.\par Esta regra tambem não satisfaz. Porquanto, ella é obscura e indeterminada; visto que é necessario apontar o que é necessario, para que o acto possa ser norma de cada agente livre. --- Além d'isso, sendo tão variada a condição e o estado dos agentes livres, é difficil que. a norma, que serve para uns, possa servir para outros.\par 0) Para \emph{Jou/froy} é a seguinte: O \emph{teu acto}, \emph{para ser bom}, \emph{deve ser conforme com o teu fim}. A razão é porque o bem do ente é o seu fim.\par Esta regra tambem é imperfeita. Porquanto, o acto geralmente não é bom, porque leva ao fim; mas leva ao fim, porque é bom. --- Além d'isso, se fosse verdadeira essa regra, a bondade ou a malicia seria sempre \emph{extrinseca e accidental} ao acto, porque consistiria na \emph{relação} do acto com o fim. Ora a bondade ou a malicia é geralmente \emph{intrinseca} e \emph{essencial} ao acto; pois muitos actos são bons ou maus independentemente da consideração do fim.\par 4) Para alguns escriptores essa regra é a seguinte: A \emph{vontade deve fazer o que a intelligencia mostra ser mais perfeito}. A razão é porque o amor deve ser proporcionado ao acto do entendimento, e por isso a vontade deve escolher o que a intelligencia conhece ser mais perfeito.\par Esta formula é inacceitavel. Na verdade, ella suppõe que a vontade esteja em discordia com a intelligencia, quando não prefere o mais perfeito, Ora tal supposição é falsa; porque nem a intelligencia manda que a vontade escolha necessariamente o mais perfeito, nem a vontade, quando escolhe o menos perfeito, declara que este é superior ao bem, que a intelligencia julgou mais perfeito. --- Além d'isso, essa formula confunde os \emph{conselhos} com os \emph{preceitos}; porque ninguem é obrigado ao \emph{optimo}, nem \emph{absoluto} nem \emph{relativo}, comtanto que não se pratiquem coisas prohibidas, nem se omittam coisas preceituadas.\par e) Outros apresentaram a seguinte regra: \emph{Observa a ordem}.\par Esta regra não é \emph{suprema}. O homem deve \emph{observar a ordem}, mas só porque a \emph{rusão} mostra a justiça e a necessidade d'esta observancia.\par f) Para \emph{Balmes} é esta: \emph{O teu acto}, \emph{para ser bom}, \emph{deve derivar do amor divino}, \emph{explicito ou implicito}. A razão é porque a moralidade do homem deve imitar a moralidade de \textsc{Deus}. Ora a moralidade em \textsc{Deus} é o acto, com que Elle ama a sua bondade infinita. Logo todos os actos humanos devem derivar, explicita ou implicitamente, do amor de \textsc{Deus}.\par Esta regra tambem é inacceitavel. Porquanto, não ha comparação entre\pgn{608}a moralidade de \textsc{Deus} e a do homem; visto que, em \textsc{Deus}, não é possivel outro acto que não seja o amor da sua bondade, ao passo que no homem ha muitos outros actos moraes possiveis e obrigatorios, como o acto de fê, de esperança, etc. --- Além d'isso, essa regra não é \emph{suprema}; porque é necessario demonstrar que o homem é obrigado a operar sempre por amor de \textsc{Deus}.}.

\pg{607} IT. A faculdade perceptiva da \emph{moralidade} é só a \emph{intelligencia}. Na verdade, a moralidade é uma idea abstracta, necessaria e universal; pois consiste na conveniencia ou não-conveniencia do acto com a nossa natureza racional, e essa conveniencia ou não-conveniencia é uma idea abstracta, necessaria e universal.

\pg{608} Ora só a intelligencia é capaz de conceber e possuir ideas abstractas, necessarias e universaes. Logo a faculdade perceptiva da \emph{moralidade} é só a \emph{intelligencia} \footnote{Esta proposição é contra os \emph{materialistas} e os \emph{sensualistas}.\par \emph{a}) Os \emph{materialistas}, como \emph{Robinet}, \emph{Gall}, dizem que a faculdade perceptiva do bem e do mal é um \emph{sentido interno corporeo}, semelhante aos cinco sentidos externos. --- Esta opinião é absurda. A bondade e a malicia moral são qualidades abstractas e immateriaes, porque denotam uma relação entre o acto e a nossa natureza racional. Ora qualidades abstractas e immateriaes não podem ser percebidas por um sentido corporeo, porque o objecto d'este é constituido pelas qualidades materiaes e concretas.\par \emph{c}) Os \emph{sensualistas}, como \emph{Hutcheson}, \emph{Reid}, pensam que a faculdade perceptiva do bem e do mal é um \emph{sentido}, \emph{espiritual} ou \emph{moral}, superior aos sentidos corporeos e distincto da razão. --- Esta opinião tambem é falsa. Porquanto, uma faculdade perceptivo-espiritual, distincta da razão ou intelligencia, não existe. A razão \emph{pratica} não differe da razão \emph{especulativa}. --- Além d'isso, esta opinião leva a admittir uma moralidade \emph{subjectiva}, \emph{relativa e variavel}, porque esse sentido seria uma especie de \emph{instincto}, e todo o instincto, por si, é \emph{subjectivo}, \emph{relativo} e \emph{variavel}.\par A nossa intelligencia concebe a idea de \emph{bem} e de \emph{mal moral} pelo modo por que concebe as outras ideas universaes, --- pela \emph{abstracção} e pela \emph{comparação}; emquanto, considerando os bens e os males \emph{particulares}, compara-os entre si, percebe o elemento que os constitue, tornando-os bons ou maus, e d'este modo eleva-se á idea do bem e do mal \emph{universal}. Por isso a idea do \emph{bem} e do mal não é \emph{innata}, como pretendem alguns escriptores.}.

%%>>> art-55

\section*{Artigo II.\\ Fontes da moralidade dos actos humanos.}
\addcontentsline{toc}{section}{Artigo II. --- Fontes da moralidade dos actos humanos.}

58. \textbf{Fontes} da moralidade dos \textbf{actos humanos}. --- \emph{Fontes} da \emph{moralidade} dos actos humanos são os principios, que não só tornam os nossos actos moralmente bons ou maus, honestos ou não-honestos, mas tambem lhes communicam uma \emph{especie} particular de bondade ou de malicia, que os caracteriza e distingue uns dos outros. --- Estas fontes são \emph{tres}, a saber: o \emph{objecto}, o \emph{fim} e as \emph{circumstancias} \footnote{A \emph{moralidade}, que é uma qualidade inherente aos nossos actos, consiste, como dissemos, na conveniencia ou nàáo-conveniencia dos actos com a nossa natureza \emph{racional}. Tal conveniencia ou não-conveniencia deriva do triplice elemento ou principio, que apontamos no texto.}.

\pg{609} 94. \textbf{Objecto} do acto humano. --- \emph{Objecto} do acto humano é a coisa, que forma o termo \emph{immediato} ou \emph{directo} da tendencia da vontade humana \footnote{\emph{Objecto}, pois, do acto humano --- não é tudo o que forma o termo da tendencia da vontade; aliás o \emph{fim} e as \emph{circumstancias} identificar-se-hiam com o objecto dos nossos actos, --- mas é a propria \emph{coisa}, para a qual a vontade humana tende,}.

55. \textbf{Divisão} do objecto do \textbf{acto} humano. --- O \emph{objecto} do acto humano é \emph{material} ou \emph{formal}. --- O objecto \emph{material} é a propria coisa, que constitue o termo do acto humano, \emph{considerada na sua physica entidade}; assim o objecto \emph{material} da esmola é o \emph{dinheiro} que se dá ao pobre. O \emph{formal} é o \emph{aspecto}, sob o qual se attinge o objecto. \emph{material} e que é conforme, ou não, com a natureza racional; assim o objecto \emph{formal} da esmola é o \emph{allivio} do pobre \footnote{A divisão do \emph{objecto} em \emph{material} e \emph{formal} é de uma importancia capital, como em todas as sciencias, tambem na \emph{Moral}. Um e o mesmo objecto \emph{material} póde communicar ao nosso acto a bondade ou a malicia, conforme o \emph{aspecto} sob o qual se attinge ou o \emph{escopo} que se deseja alcançar. Assim podemos dar a um pobre uma \emph{moeda} para o \emph{alliviar} ou para o \emph{levar} a commetter um \emph{crime}; e deste odo, ao passo que, em ambos os casos, o objecto \emph{material} é o mesmo, o objecto \emph{formal} é muito diverso.}.

56. O \textbf{acto} humano tira do objecto formal a sua \textbf{moralidade especifica}. --- O acto humano é um movimento da nossa alma consciente e livre para um certo e determinado termo, que é o objecto \emph{formal}. Ora o termo, em que pára o movimento, determina e especifica o proprio movimento. Logo o acto humano tira do objecto \emph{formal} a sua moralidade \emph{especifica}. --- Por isso o acto terá uma certa e determinada \emph{especie} de bondade ou malicia moral, conforme a bondade ou malicia moral, que foi reconhecida e desejada no objecto \emph{formal}, e que consiste na conveniencia ou não-conveniencia d'este objecto com a nossa natureza racional. i :

Dissemos --- do objecto \emph{formal}; porque é este, e não o objecto \emph{material}, o termo proprio e especial do acto \footnote{Todo o acto é \emph{especificado} pelo seu objecto \emph{formal}; por isso o acto moral será dotado d'aquella \emph{especie} de bondade ou honestidade, de que se acha revestido o objecto \emph{formal}. (\emph{Sum}. \emph{Th}., 1-9, q. 18, a. 2). --- Advirta-se que a bondade ou honestidade se divide em \emph{objectiva} e \emph{subjectiva}. \emph{Objectiva} é a que está no proprio \emph{objecto}, emquanto este convém á nossa natureza racional. \emph{Subjectiva} é a que está no proprio \emph{acto} /\emph{wmano}; emquanto este, por tender para um objecto \emph{bom}, convém á natureza racional. A bondade \emph{objectiva} é sempre e necessariamente anterior \& bondade \emph{subjectiva}; pois aquella é causa e fundamento d'esta,\par 89 Vor. II.}.

\pg{610} 57. Fim do \textbf{acto} humano. --- \emph{Fim} do acto humano é o \emph{motivo} que leva o homem a operar \footnote{A definição, que apresentamos do \emph{fim} dos actos humanos, refere-se ao \emph{fim do agente}, \emph{e} não ao \emph{fim do acto}. O \emph{fim do agente} é o \emph{motivo} que leva a operar, e por isso é \emph{extrinseco} ao proprio acto; --- \emph{o fim do acto} é o \emph{effeito}, que o proprio acto naturalmente produz, e por isso é \emph{intrinseco} ao acto e identifica-se com o objecto \emph{formal}. --- Às vezes, o \emph{fim do agente} e o \emph{fim do acto} identificam-se, como quando fazemos esmola com o fim de alliviarmos um pobre. Outras vezes, estes dois fins são distinctos um do outro; assim, quando fazemos esmola para satisfazermos pelos nossos peccados, o fim do \emph{acto}, que é o allivio do pobre, é distincto do fim do \emph{agente}, que é a satisfação pelos peceados.}.

58. \textbf{Divisão} do fim do \textbf{acto} humano. --- O /\emph{im} do acto humano divide-se em \emph{proximo}, \emph{remoto} e \emph{ultimo}.

a) Fim \emph{proximo} é o motivo \emph{immediato}, que nos leva a praticar certos actos; assim a \emph{mortificação} das paixões póde ser o fim \emph{proximo} de alguns actos de penitencia.

b) Fim \emph{remoto} é o motivo \emph{mediato}, que nos leva a praticar certos actos, de modo que é um \emph{meo} para a aquisição de um fim ulterior; tal é a \emph{virtude} da \emph{caridade}, que podemos ter em vista quando, com alguns actos de penitencia, mortifieamos as nossas paixões desordenadas.

c) Fim \emph{ultimo} é o motivo \emph{supremo}, que nos leva a praticar certos actos, de modo que elle se deseja por si e é a causa por que se desejam as outras coisas; assim a \emph{felicidade} é fim \emph{ultimo}, pois deseja-se por si e é a causa por que desejamos a \emph{caridade}, \& \emph{mortificação} das nossas paixões, etc. \footnote{Um só acto, ainda que possa ter varios fins \emph{remotos}, dos quaes um é fim do outro, comtudo não póde ter senão um fim \emph{proximo}, que o especifica; porquanto, embora possam ser varios os fins \emph{proximos} do acto, todavia um só é o fim \emph{principal}, de que o proprio acto tira a sua \emph{especie substancial}, sendo os outros meramente \emph{accidentaes}. O mesmo deve dizer-se, e com mais razão, do fim que é \emph{ultimo} numa serie de fins; porque este é o primeiro na intenção, e por isso não póde deixar de especificar o acto.}.

59. O \textbf{acto} humano tira do fim a sua moralidade. --- Com effeito, o \emph{fim} é o proprio e verdadeiro bem da vontade; porque é para elle que esta faculdade tende como para o seu termo. Ora o acto de uma faculdade tira do bem, que é proprio da mesma faculdade, a sua especie ou natureza. Logo o acto humano tira do \emph{fim} á sua \emph{moralidade} \footnote{O \emph{fim} é sempre fonte de moralidade. Mas é necessaria uma advertencia. Quando o \emph{fim} se torna \emph{objecto} da vontade, o que acontece todas as vezes que o \emph{fim do agente} se identifica com o \emph{fim do acto}, então é fonte \emph{principal} e \emph{essencial} de moralidade e especifica o acto humano. Quando, \emph{oh} f 59 \textbf{mesa} (eue \textbar{} \textbar{} mb \emph{il n} RA ia asl\pgn{611}porém, o \emph{fim} é distincto do \emph{objecto}, nesse caso é fonte \emph{secundaria} e \emph{accidental} de moralidade (\emph{a príncital} e \emph{essencial} é o \emph{objecto}) e não especifica o acto.\par Esta ultima regra admitte excepções. Por quanto, se o \emph{objecto} fór moralmente \emph{indifferente}, o acto tira a sua especie do \emph{fim}, embora \emph{ectrínseco}; porque, então, esse fim é o proprio e principal termo da tendencia da vontade. Como tambem, quando \emph{o fim} não se limita a acompanhar o acto, mas leva o proprio agente a operar e a tender para um dado objecto, neste caso a moralidade do acto deriva principalmente do \emph{fim}, embora \emph{extrinseco}, porque é elle o principal termo da Xontade; assim, quem trabalha para a gloria do mundo, é mais vaidoso do que trabalhador (\emph{Sum}. \emph{TJ}., 1-2, q. 18, a. 4).}.

\pg{611} 60. Circumstancias do acto humano. --- \emph{Circumstancias} do acto humano são as condições extrinsecas e accidentaes, as quaes augmentam ou diminuem a sua \emph{moralidade}. --- As \emph{circumstancias} suppõem o acto constituido na sua moralidade substancial.

61. Especies das circumstancias do \textbf{acto} humano. --- As \emph{circumstancias} podem ser \emph{aggravantes} ou \emph{attenuantes}, \emph{antecedentes} ou \emph{concomitantes}, \emph{objectivas} ou \emph{subjectivas}.

a) Aggravantes ou attenuantes, --- \emph{Aggravantes} são as que \emph{augmentam} a \emph{moralidade}, isto é, á bondade ou a malicia \emph{moral} do acto humano. \emph{Attenuantes} são as que a \emph{diminuem} \footnote{Algumas \emph{circumstancias} são tão \emph{aggravantes}, que dão ao acto humano uma nova especie de bondade ou malicia moral; assim o roubo de um objecto \emph{sagrado} encerra, além da malicia do \emph{furto}, uma nova especie de malicia, que é a \emph{irreligião}.}.

b) Antecedentes ou concomitantes, --- \emph{Antecedentes} são as que \emph{precedem} o acto e exercem nelle uma influencia causal. \emph{Concomitantes} são as que \emph{acompanham} o acto, encontrando-o constituido na sua especie.

c) Objectivas ou subjectivas. --- \emph{Objectivas} são as que cercam, como \emph{accidentes}, o \emph{objecto} do acto humano. \emph{Subjectivas} são as que affectam o acto interior e a pessoa do agente \footnote{As \emph{circumstancias} estão contidas no celebre verso:\par \emph{Quis}, \emph{quid}, \emph{ubi}, \emph{quibus auxiliis}, \emph{cur}, \emph{quomodo}, \emph{quando} ?}.

62. O \textbf{acto} humano tira das circumstancias a \textbf{sua moralidade}. --- O acto humano tirará das \emph{circumstancias} a sua \emph{moralidade}, se ellas tiverem uma especial conveniencia ou não-conveniencia com a nossa natureza \emph{racional}, e se tal conveniencia ou não-conveniencia for, de algum modo, conhecida e approvada pelo agente. Ora algumas \emph{circumstancias} têm uma especial conveniencia ou não-conveniencia com a nossa natureza racional, e esta conveniencia ou não-conveniencia póde ser, de algum modo, conhecida e approvada pelo agente. Logo o acto humano tira das \emph{circumstancias} a sua \emph{moralidade} \footnote{O acto humano não tende para um objecto \emph{abstracto}, mas para um\pgn{612}objecto \emph{concreto}, isto é, cercado de varias \emph{circumstancias}. Logo o acto humano, assim como tira do objecto a sua moralidade \emph{especifica}, assim tambem tira das \emph{circumstancias} a sua moralidade \emph{accidental}. --- E se, ás vezes, o acto tira de alguma \emph{circumstancia} a propria moralidade \emph{especifica}, isto dá-se emquanto uma \emph{circumstancia aggravante}, por causa da sua conveniencia ou n\&o-conveniencia especial com a ordem racional, se torna condição principal do \emph{objecto}; tal é o roubo de um objecto \emph{sagrado}. Mas, em geral, \& \emph{circumstancia}, como \emph{tal}, não muda a especie do acto moral; só augmenta ou diminue a sua gravidade, ou produz nelle uma alteração accidental. (\emph{Sum}. \emph{Th}., 1-2, q. 18, aa. 3, 10, 11).}.

\pg{612} 63. \textbf{Corollario}. --- O acto humano será simples e absolutamente \emph{bom} na ordem moral, se fór bom o \emph{objecto}, bom o \emph{fim}, e bóas as \emph{circumstancias}. Se um destes elementos faltar, o acto será total ou parcialmente mau \footnote{Os tres elementos ou principios de \emph{moralidade} podem combinar-se por varios modos; e d'estas varias combinações resultam notaveis differencas para a bondade ou malicia dos actos.\par Enumeremos os casos principaes.\par 1.9) O acto póde ser bom em relação ao \emph{objecto}, ao \emph{fim}, e ás \emph{circumstancias}. Neste caso, o acto é inteiramente bom, e póde pertencer a differentes virtudes; tal é a esmola feita por amor de \textsc{Deus} a um pobre, com o fim de o ajudar para a boa criação dos filhos.\par 9,9) O acto póde ser mau em relação ao \emph{objecto}, mas bom em relação ao \emph{fim}. Neste caso, o acto é mau, e a sua malicia essencial póde ser \emph{altenuada} pela bondade do fim, mas não totalmente \emph{destruida}. A celebre theoria de \emph{Macchiavelli}, segundo a qual o \emph{fim justifica}. \emph{os meios}, embora \emph{maus}, é simplesmente immoral.\par 8.9) O acto póde ser bom em relação ao \emph{objecto}, mas mau em relação ao \emph{fim}. Neste caso, o acto é mau. --- Ás vezes, o mesmo acto póde ser motivado por dois fins, dos quaes um é bom e outro é mau. Em tal caso, o fim mau, se não influe efficazmente na substancia do acto e está subordinado ao fim bom, não destroe a bondade essencial do acto, embora lhe acerescente uma malicia especial; tal seria o caso de quem fizesse esmola, levado pela caridade e tambem, embora secundariamente, pela vaidade.\par 4,0) O acto póde ser bom em relação ao \emph{objecto} e ao \emph{fim}, mas mau em relação a alguma \emph{circumstancias}. Nesse caso, geralmente, o acto é bom na sua essencia, e só perde uma parte do seu merito. --- Dissemos --- \emph{geralmente}; porque, se a \emph{circumstancia} assume: a condição ou natureza do \emph{objecto}, o acto é essencialmente mau.\par Do que deixamos dito resulta que muitos actos, embora considerados \emph{em abstracto} e sô no seu \emph{objecto}, sejam \emph{indifferentes} (como os actos de comer, \emph{passear}, ete.), comtudo, considerados \emph{em concreto}, isto é, no \emph{individuo}, são \emph{bons} ou \emph{maus}. Na verdade, todo o acto humano ou é necessario, util, conveniente á natureza racional, ou é inutil, ocioso, n\&o-conveniente á natureza racional. No primeiro caso, o acto é moralmente bom; no segundo, não é. --- Além d'isso, sendo proprio da razão dirigir para o fim conveniente todos os seus actos deliberados, segue-se que o acto --- ou está disposto para o fim e é moralmente bom, --- ou não está, e é moralmente mau. --- Nem é impossivel\pgn{613} dirigir todos os actos para o fim conveniente; porque, para isso, basta uma intenção \emph{virtual} ou \emph{habitual} (\emph{Sum}. \emph{Th}., 1-2, q. 18, aa. 8, 9).\par Pergunta-se: O acto \emph{externo}, ou \emph{imperado}, augmenta a moralidade do acto \emph{interno}, ou \emph{elicito}?\par Respondemos que o acto \emph{externo}, \emph{por} si, nio augmenta a moralidade do acto \emph{interno}: Na verdade, o acto \emph{externo} recebe do \emph{interno} toda a moralidade, de que é dotado; pois o acto \emph{externo é humano,e} por isso \emph{moral}, emquanto é informado pelo acto \emph{interno}. Ora o acto não póde, \emph{por si}, augmentar a moralidade, qué elle recebe de outro acto. --- Dissemos --- \emph{por si}; porque, \emph{accidentalmente}, o acto \emph{externo} pode augmentar a moralidade do acto \emph{interno}, emquanto apresenta á vontade a occasiüo ou de repetir o acto \emph{inferno}, ou de o tornar mais intenso e dqimorado (\emph{Sum}. \emph{T}7., 1-2, q. 20, a. 4).}.

%%>>> art-56

\section*{Artigo III.}
\addcontentsline{toc}{section}{Artigo III.}

\pg{613} \textbf{Consequencias da moralidade}

\textbf{dos actos humanos}.

64. \textbf{Consequencias da} moralidade dos actos humanos. --- As principaes consequencias, que resultam da \emph{moralidade} dos actos humanos, são as seguintes: a .\emph{Imputabilidade}, o \emph{merito} ou o \emph{demerito}, o \emph{premio} ou o \emph{castigo} \footnote{O acto \emph{moral} é o que a livre vontade produz, depois de a intelligencia ter conhecido a sua honestidade ou não-honestidade. D'ahi se segue que o acto moral --- \emph{a}) \emph{é imputavel}, --- \emph{b}) tem razão de \emph{merito} ou de \emph{demerito}, --- \emph{c}) merece \emph{premio} ou \emph{castigo}.}.

65. Imputabilidade. --- \emph{Imputabilidade} é á propriedade pela qual a bondade ou a malicia \emph{mora}? do acto humano póde e deve ser attribuida ao agente, como ao seu principio e auctor. --- A \emph{imputabilidade} é causa de \emph{estima} ou de \emph{desprezo}, conforme a qualidade do acto humano \footnote{A \emph{imputabilidade}, de que tratamos, não é a \emph{physica}, mas sim a \emph{moral}. A imputabilidade \emph{p/iysica} é quando se attribue ao agente o acto, considerado na sua entidade \emph{p/usica}, isto é, emquanto deriva da intelligencia e da vontade deliberada. A \emph{moral} é quando. se attribue ao agente o acto, considerado na sua entidade \emph{moral}, isto é, emquanto é /\emph{ionesto} ou \emph{n\&o-honesto}. --- A imputabilidade \emph{p/jjsica} póde estar sem a \emph{moral}, como quando uma pessoa produz livremente um acto, mas não adverte, por ignorancia \emph{énvencivel,se} é honesto ou não; mas a \emph{moral} nao póde estar sem a \emph{phusica}, porque um acto \emph{moral}, para ser tal, deve derivar du intelligencia e da livre vontade. --- A" \emph{imputabilidade} corresponde a \emph{responsabilidade}, pela qual o homem deye considerar-se como auctor dos seus actos, e por isso sujeito a todas as consequencias,}.

66. O \textbf{acto} moral \textbf{deve ser} imputado \textbf{ao agente} racional e \textbf{livre}. --- "Todo o effeito deve imputar-se á causa de que depende e no modo por que depende. Ora o acto \emph{mora}? depende do agente racional e livre. Na verdade, esse acto realiza-se, só porque esse

\pg{614} \textasciicircum{} Em \textbf{OMNEM XR} EAR. nog

614 MORAL --- PARTE PRIMEIRA --- ETHICA agente quer realizal-o. Logo o acto \emph{moral} deve ser imputado a0 agente racional e livre \footnote{O fundamento, pois, da \emph{imputabilidade}, e por isso da \emph{responsabilidade}, é o \emph{livre arbitrio}. Ao homem deve \emph{imputar-se} o acto, que elle livremente \textbar{} quiz produzir. --- Por isso tudo o que tira, diminue ou augmenta o acto \emph{livre}, tira, diminue ou augmenta a sua \emph{imputabilidade}.\par Notamos que, para a \emph{malicia moral} do acto ser imputada, basta que seja \emph{advertida} pelo agente; ao passo que, para a \emph{bondade moral} do acto ser imputada, não basta que seja \emph{advertóda}, mas é preciso que seja de algum modo \emph{desejada}. A razão d'esta differenca é porque o agente, se \emph{adverte} na malicia de um acto e todavia não deixa de o fazer, é responsavel, pois quer \emph{livremente} o mal; ao passo que, embora advirta na bondade de um acto, não póde dizer-se que é responsavel por elle, pois póde não querel-o \emph{positivamente}, mas simplesmente toleral-o ou permittil-o.\par Não só os \emph{actos}, mas tambem as suas \emph{livres omissões}, com seus \emph{effeitos}, E nos são imputaveis. --- Para que os effeitos \emph{maus}, que derivam do acto ou da sua omissão, sejam imputavels, é necessario, como dissemos, que o agente ' não só os preveja, mas tambem tenha o poder e a obrigação de abster-se do acto ou da sua omissão (\emph{Sum}. \emph{Th}., 1-2, q. 21, a. 2).}.

67. \textbf{Merito} e \textbf{demerito}. --- \emph{Merito} é o acto, que foi praticado em vantagem de outrem e por isso exige por justiça um premio. --- \emph{JDemerito} é o acto que foi praticado em prejuizo de outrem e por isso exige por justiça um castigo \footnote{As definições apresentadas consideram o \emph{merito} e o \emph{demerito} em \emph{concreto}. Em \emph{abstracto}, o \emph{merito} é a \emph{evigencia do premio}, \emph{inherente ao acto}, \emph{pela vantagem}, \emph{que este} causou a \emph{outrem}, --- e o \emph{demerito} É a \emph{evigencia do castigo}, \emph{inherente} ao \emph{acto}, \emph{pelo prejuizo}, \emph{que este causou a outrem}.\par Vejamos os requisitos para que o acto exija \emph{premio} ou \emph{castigo}.\par a) O acto, para que exija \emph{premio}, deve ter os seguintes requisitos: --- 1.9) deve ser moralmente \emph{bom}, pois o acto \emph{mau} merece castigo; --- 2.º) deve ser \emph{livre}, não só porque o acto \emph{necessario} não é digno de louvor nem de premio, mas tambem porque o \emph{merito} é posterior á \emph{imputabilidade}, e esta, suppõe a \emph{liberdade}; --- 9.9) deve causar a outrem uma \emph{vantagem}, pois o premio só póde exigir-se d'aquelle, que recebeu do nosso acto uma vantagem; --- 4.9) deve ser explicita ou implicitamente acceito, como merecedor de premio, por aquelle em cuja vantagem foi praticado; aliás o premio não será devido por \emph{justiça}, mas por \emph{conveniencia}; --- 5.º) deve não ser \emph{devido} por causa de uma recompensa anteriormente determinada e recebida; pois, nesse caso, o acto não é \emph{meritorio}, não \emph{exige} recompensa ou premio, á não ser que o acto se pratique com tão extraordinario cuidado e celeridade, que produza ; uma vantagem superior á que se podia e devia esperar. \emph{r}\par bh) O acto, para que exija \emph{castigo}, deve ter os seguintes requisitos: --- 1.º) deve ser moralmente mau; --- 2.9) deve ser voluntario e livre; --- 3.º) deve causar a outrem um prejuizo. --- A razão d'estes requisitos é manifesta.}.

68. \textbf{Os actos humanos} podem ter uma \textbf{razão de merito} ou \textbf{de demerito relativamente aos} individuos, á \textbf{sociedade} e a \textbf{Deus}.

a) \textbf{Podem} ter uma razão de merito ou de demerito relativamente \textbar{} \pg{615}  aos individuos, --- Na verdade, o homem póde praticar um acto, de que resulte uma vantagem para outrem, a qual não seja uma restituição de coisa devida por contracto; como tambem póde praticar um acto que produza prejuizos a quem não tem obrigação de os receber. Sendo assim, o homem, no primeiro caso, tem direito ao premio, e, no segundo caso, deve sujeitar-se ao castigo. Logo os actos humanos têm razão de \emph{merito} ou de \emph{demerito} relativamente aos outros individuos \footnote{A conclusão é evidente. Um individuo póde produzir um acto \emph{meri}-\par , \emph{torio} relativamente a outro individuo. Assim o filho que venera o pae, o servo que obedece ao amo, o discipulo que respeita o mestre, produzem actos que exigem uma recompensa proporcionada. Em taes casos, a acceitação do acto por parte do pae, do amo ou do mestre, é \emph{expressa} ou \emph{explicita}. Quando não é \emph{explicita}, ordinariamente é \emph{zmplicita}, emquanto a lei natural impõe a obrigação de recompensarmos, ao menos pela gratidão, aquelles de quem recebemos alguma vantagem, --- obrigação que equivale a uma acceitacüo ?\emph{mplicita}, da parte dos que recebem o obsequio.}. . b) \textbf{Podem} ter uma razão de merito ou de demerito \textbf{relativamente} á \textbf{sociedade}, --- Com effeito, \emph{o} homem póde praticar actos, de que resulte vantagem ou damno para a sociedade ou para os seus membros. Ora é justo que a sociedade recompense ou castigue os actos, de que Ihe resulta vantagem ou damno \footnote{Para que um acto tenha razão de \emph{merito} ou de \emph{demerito} velativamente á sociedade, não é preciso que elle seja praticado \emph{directamente} para o bem ou para o mal da communidade; mas basta que o seja \emph{indirectamente}, emquanto procura o bem ou o mal de algum cidadão em particular, \emph{e} mesmo do proprio sujeito agente, porque o bem ou o mal do cidadão particular concorre para o bem ou para o mal de toda a sociedade.}.

c) \textbf{Podem} ter uma razão de merito ou de demerito relativamente a \textbf{Deus}. --- Na verdade, \textsc{Deus} estabeleceu uma ordem moral, promettendo ao homem uma felicidade eterna, se a guardasse, e ameaçando um castigo eterno, se a transgredisse. Logo o homem, se observar a lei, dará gloria a \textsc{Deus}, e alcançará \emph{meritos} para a felicidade eterna, --- e, se a violar, desprezará a Drvus, e os seus actos merecerão \emph{castigo} \footnote{Os actos \emph{moraes} têm razão de \emph{merito} ou de \emph{demerito} deante do \textsc{Deus}. Porquanto alguns actos, como a obediencia aos paes, o amor do proximo, etc., têm todos os requisitos necessarios para o \emph{merito}; poisque --- 1.º) são moralmente \emph{bons}, emquanto são conformes com a nossa natureza racional; --- 2.º) são \emph{livres}, porque a vontade é regulada pela lei, mas não constrangida \emph{qhysicamente}; --- 8.") dão gloria a \textsc{Deus}, pois quem guarda a lei, honra o legislador; --- 4.") sao acceitos por \textsc{Deus}, emquanto merecedores de premio; porque \textsc{Deus} prometteu, pela propria voz da natureza, uma felicidade eterna a quem guardasse a ordem ou a lei moral; --- 5.º) não são devidos por um premio antecedentemente promettido e recebido. Ora os actos, que satisfazem\pgn{616}a estes requisitos, são \emph{meritorios} deante de \textsc{Deus}, fim ultimo do homem e supremo ordenador de todas as coisas. --- Pelo contrario, outros actos, como a blasphemia, etc., têm razão de \emph{demerito} deante de \textsc{Deus}; porque são moralmente maus, são livres, tiram a \textsc{Deus} a gloria devida, e por isso exigem castigo (\emph{Sum}. \emph{TA}., 1-2, q. 21, aa. 3, 4).\par Dizem: O acto, que é devido a outrem, não póde ser \emph{meritorio} deante d'elle. Ora os actos honestos são devidos a \textsc{Deus}. Logo não podem ser \emph{meritorios} deante de \textsc{Deus}.\par Resposta. Não póde ser \emph{meritorio} o acto que é devido por causa de uma recompensa anteriormente determinada e recebida, mas póde ser \emph{meritorio} o acto devido pela obrigação imposta pela \emph{ler}. Ora os actos honestos são \emph{meritorios} deante de \textsc{Deus}, não por causa da recompensa anteriormente promettida por \textsc{Deus} e recebida pelo homem, mas por causa da obrigação imposta ao homem pela lei divina. --- Todavia deve notar-se que muitos actos honestos não são impostos pela lei divina, mas dependem exclusivamente da bóa vontade humana, e merecem recompensa.\par Insistem: Um acto, que é devido por \emph{lei}, não póde ser \emph{meritorio}.\par Resposta. Negamos que o acto, devido por \emph{lei}, não seja \emph{meritorio}. Com effeito, \emph{o} homem, emquanto, pela \emph{sua livre} e \emph{bôa vontade}, faz o que deve fazer, merece; aliás o acto de \emph{justiça} não seria \emph{meritorio}, o que é falso.\par Continuam: O acto, que é vantajoso para quem o realiza, e não para outrem, não póde ser \emph{meritorio} deante d'este. Ora o acto honesto é vantajoso para o homem, e não para \textsc{Deus}. Logo não é \emph{meritorio} deante de \textsc{Deus}.\par Resposta. Para o acto ser \emph{meritorio} deante de outrem, basta que este receba algum bem \emph{ezatrinseco}. Ora \textsc{DEus} recebe dos nossos actos honestos um bem \emph{extrinseco}, que é a sua \emph{gloria}, e por isso deve recompensal-os.\par Replicam: Se os nossos actos honestos tivessem razão de \emph{merito} deante de \textsc{Deus}, \textsc{Deus} seria nosso devedor. Ora isto repugna.\par Resposta. Se os nossos actos honestos fossem \emph{meri?torios} deante de \textsc{DEUS}, Dus seria nosso devedor, não \emph{absolutamente}, mas \emph{hypotheticamente}, emquanto \textsc{Deus}, na sua bondade, ordenou que as nossas bóas acções, feitas com o seu auxilio, fossem \emph{meritorias} de premio. Ora isto não repugna; pois, em tal caso, \textsc{DEus} não é devedor do homem, mas da propria fidelidade e justiça.}.

-

\pg{616} 69. \textbf{Premio} e \textbf{castigo}. --- \emph{Premio} é o bem com que se recompensa ou deve recompensar-se a vantagem proveniente do acto \emph{moralmente meritorio}. --- \emph{Castigo} é o mal com que se deve punir quem, pela sua acção, consciente, voluntaria e injusta, causou prejuizos a outrem \footnote{O \emph{premio} e o \emph{castigo} são consequencias do \emph{merito} e do \emph{demerito}; porque todo o \emph{merito} encerra o \emph{direito}, embora remoto, ao \emph{premio}, como todo o \emph{demerito} envolve a divida do \emph{castigo}.}.

70. O \textbf{premio} e o \textbf{castigo} não só \textbf{conservam} e \textbf{promovem}, \textbf{mas tambem completam} e \textbf{restabelecem} a \textbf{ordem moral}.

a) Conservam e promovem a ordem moral --- Na verdade, o \emph{premio} e o \emph{castigo} são motivos muito efficazes para que a vontade, conservando a sua liberdade, pratique o bem e fuja do mal. Ora é nesta livre pratica do bem e fuga do mal que \pg{617}  consiste a conservação e promoção da ordem moral. Logo o premio e o castigo conservam e promovem a ordem moral.

b) Completam e restabelecem a ordem moral. --- Com effeito, a ordem moral --- é completa, quando ao \emph{merito} se retribue o \emph{premio}, pois todo o merito exige uma recompensa, --- e é restabelecida, quando ao \emph{demerito} se inflige o \emph{castigo}, pois a idea de \emph{castigo} está necessariamente ligada com a idea de reparação da ordem violada. Logo o premio e o castigo completam e restabelecem a ordem moral \footnote{O fim do \emph{premio} e do \emph{castigo} é duplice: --- 1." conservar e promover a ordem moral; --- 2.º) completal-a e reparal-a por meio da retribuição devida ao \emph{merito} ou ao \emph{demerito}.\par Considerados em relação ao primeiro fim, o \emph{premio} e o \emph{castigo} são \emph{anteriores} ao exercicio da vontade; mas, considerados em relação ao segundo fim, são posteriores ao mesmo exercicio. --- A relação entre a retribuição e esse duplice fim é \emph{essencial}, e por isso tem o seu fundamento, a sua razão de ser, na intenção do Auctor da natureza. D'onde se vé o absurdo do \emph{atheismo social}, que pretende estabelecer o fundamento \emph{supremo} da \emph{justiça penal} no direito da sociedade; quando a razão e a fó nos attestam que o poder de punir os crimes só póde derivar de \textsc{Deus}. --- Mais adeante voltaremos ao assumpto, e com o devido desenvolvimento.\par Sendo \textsc{Deus} o auctor e o ordenador do mundo, não só do mundo \emph{physico} mas tambem do \emph{moral}, só Elle póde providenciar pela conservação da ordem com meios convenientes e estabelecer premios ou castigos, e a auctoridade humana só póle fazel-o \emph{por delegação dibina}. --- O \emph{premio} e o \emph{castigo}, emquanto \emph{conservam} e \emph{promovem} a ordem \emph{moral}, --- referem-se á \emph{sabedoria} de \textsc{DEus}; --- e, emquanto \emph{completam} e \emph{reparam} a mesma ordem, referem-se á \emph{justiça} do mesmo \textsc{Deus}. --- \textsc{Deus} quer e inflige o \emph{castigo}, emquanto este representa a \emph{sancção da ordem moral} contra a \emph{violação} da mesma ordem, e emquanto é uma victoriosa \emph{affirmação} da inviolabilidade objectiva da mesma ordem contra a rebelde \emph{negação} da vontade humana. E' assim que \textsc{Deus} restabelece, pelo \emph{castigo}, a ordem violada.\par \textsc{Deus} applica o \emph{premio} e o \emph{castigo} por um modo \emph{natural}, emquanto ás obras dignas de \emph{premio} correspondem \emph{cfífeitos bons}, e ás obras dignas de \emph{castigo} correspondem \emph{ejfeitos maus}. --- Estes effeitos, \emph{bons} ou \emph{maus}, referem-se á ordem \emph{individual}, \emph{social} ou \emph{universal}, conforme o acto tem razão de «\emph{merito} ou de \emph{demerito} em relação á ordem \emph{individual}, \emph{social} ou \emph{universal}, Assim na ordem \emph{individual}, o \emph{merito} encontra o \emph{premio} na paz da alma, na saude, etc., --- e o \emph{demerito} tem o \emph{castigo} no remorso, nas doenças da alma e do corpo\pgn{618}etc. --- Na ordem \emph{social}, o \emph{merito} tem o \emph{premio} na estima dos outros, na gloria, na posse pacifica da liberdade e dos bens, etc., --- e o \emph{demerito} encontra o \emph{castigo} no desprezo dos autros, na infamia, no odio das pessoas honestas. --- Na ordem \emph{universal}, o \emph{premio} do \emph{merito} consiste na \emph{consecução} da felicidade eterna, --- e o \emph{castigo} do \emph{demerito}, na sua \emph{privação}.}.

%%>>> art-57

\section*{Artigo IV.\\ A moralidade dos actos humanos e os habitos.}
\addcontentsline{toc}{section}{Artigo IV. --- A moralidade dos actos humanos e os habitos.}

71. \textbf{Habito}. --- \emph{Zabito} é a qualidade estavel ou permanente, que dispõe ou inclina a faculdade a operar mais de um modo \pg{618}  do que de outro, bem ou mal, conforme o acto convém ou repugna á natureza \emph{racional} do homem \footnote{Na \emph{Anthropologia} \&rat\&mos do \emph{habito}, declarando o seu \emph{sujeito}, as suas \emph{especies} e os seus \emph{effeitos}. --- Por isso aqui limitamo-nos a recordar que o \emph{habito} compete exclusivamente ás faculdades \emph{intellectuaes}, que por si são de algum modo indeterminadas em relação ao acto. --- S. Thomaz trata, com muito desenvolvimento, d'este assumpto na \emph{Sum}. \emph{Th}., 1-2, q. 49-54.}.

72. \textbf{Especies} do habito. --- O \emph{Habito}, considerado em relação ao \emph{acto}, para o qual dispõe, póde ser \emph{dom} ou \emph{mau}. --- O habito \emph{bom} é a qualidade estavel ou permanente, que dispõe ou inclina a faculdade a produzir actos \emph{convenientes} á natureza racional do homem. O habito \emph{bom} diz-se \emph{virtude}. --- O habito \emph{mau} é a qualidade estavel ou permanente, que dispõe ou inclina a faculdade a produzir actos \emph{repugnantes} á natureza \emph{racional} do homem. O habito \emph{mau} chama-se \emph{vicio} \footnote{O \emph{habito}, inclinando a faculdade para o acto como para o seu \emph{termo}, deve estar subordinado á divisão do proprio acto. Ora este é \emph{bom} ou \emph{mau}, conforme conyém ou repugna á natureza racional do homem.}.

73. \textbf{Virtude}. --- \emph{Virtude} é, pois, o habito que aperfeiçoa e inclina a faculdade a produzir actos \emph{convenientes} á natureza \emph{racional} do homem \footnote{Para \emph{Aristoteles}, a \emph{virtude} é « \emph{o habito que torna bom o agente e o acto} », ou « a \emph{disposição do perfeito para o optimo} ». --- Estas definições são ambas verdadeiras. E' verdadeira a primeira; pois é a \emph{virtude} que torna bom o agente, emquanto o dispõe para o fim, e torna bom o acto, emquanto este ou é fim ou é meio para o fim. --- E' verdadeira a segunda; porque a \emph{virtude} é uma \emph{disposição} estavel ou permanente, pela qual o agente é \emph{perfeito}, emquanto está disposto para o fim, que é uma coisa \emph{optima}, pois é a suprema felicidade do homem. (Cf. \emph{Sum}. \emph{Th}., 1-2, q. 55-67; 9-9, q. 41 sqq.).\par Os philosophos, antigos.e modernos, apresentaram outras definições da \emph{virtude}. Omittindo as definições, que têm sido apresentadas pelos materialistas e que se fundam na negação da verdadeira \emph{moralidade}, citaremos as de \emph{Platão}, \emph{M«lebranche} e \emph{Kant}.\par \emph{a}) \emph{Platão} disse que a \emph{virtude} consiste na \emph{harmonia das faculdades}. --- Esta definição é inacceitavel; porque é muita vaga e indeterminada e apenas exprime a \emph{condição} ou o \emph{effeito} da virtude, mas não a sua \emph{noturesa}.\par \emph{b}) \emph{Malebranche} ensinou que a \emph{virtude} consiste no \emph{amor da ordem}. --- Esta definição é defeituosa. A \emph{virtude} abrange varios actos da vontade, que são o \emph{desejo}, o \emph{amor}, a \emph{decisão}, \emph{á execução}; o por isso a sua definição não póde limitar-se a um só desses actos, que é o \emph{amor} da ordem.\par \emph{c}) \emph{Kant}, seguido por \emph{Cousin}, diz que a virtude consiste na \emph{difficuldade} com \emph{que a razão lucta contra as paíeoes}, --- Esta definição tambem é falsa.\pgn{619}Na verdade, a \emph{virtude} não pode referir-se exclusivamente ás paixões, mas tambem, e por um modo mais elevado, ás operações externas. Além disso, o acto, para que seja virtuoso, basta que tenda para o bem /\emph{onesto}, sendo por isso as paixões apenas uma \emph{occasião} para a perfeição da virtude.}.

\pg{619} ,

74. \textbf{Especies de virtude}. --- A \emph{virtude}, considerada em relação ao seu \emph{fim}, póde ser \emph{intellectual} ou \emph{moral}. --- E \emph{intellectual}, quando o seu \emph{fim} é aperfeiçoar a \emph{ntelligencia}, de modo que esta possa facil e ordenadamente aleangar a \emph{verdade}, que é o seu bem. --- E' \emph{moral}, quando o seu \emph{fim} é aperfeiçoar a \emph{vontade}, de modo que esta possa, com facilidade, ordem e agrado, alcan- 'car o bem \emph{moral}, que é o verdadeiro bem do homem \footnote{Embora as virtudes \emph{intellectuaes} sejam mais nobres que as \emph{moraes} (pois o objecto d'aquellas é mais. universal do que o d'estas), comtudo os effeitos, que as virtudes \emph{moraes} produzem no homem, s\&o mais beneficos e estimaveis do que os que derivam das \emph{intellectuaes}. Na verdade, a virtude \emph{intellectual} torna o homem \emph{bom}, emquanto lhe confero a aptidão para saber ou fazer \emph{bem} uma coisa, assim o «\emph{bito} da \emph{philosophia} faz com que o homem seja um \emph{bom} philosopho; ao passo que a virtude \emph{moral} torna o homem \emph{absoluta} e \emph{inteiramente} bom, emquanto a sua vontade é boa ou honesta. --- D'onde se vé o erro de \emph{Socrates} e de \emph{Platão}, que identificavam a \emph{sciencia} com a \emph{virtude moral}; pois a \emph{verdade}, que é o objecto da \emph{sciencia}, distingue-se da \emph{bondade}, que é o objecto da \emph{virtude moral}. --- Lembramos que as viitudes \emph{intellectuaes} se dividem em \emph{especulativas} e \emph{praticas}. As \emph{especulativas} são a \emph{sabedoria}, a \emph{sciencia} e o \emph{entendimento}, As \emph{praticas} são a \emph{arte} e a \emph{prudencia}.\par As virtudes \emph{moraes} residem, como em seu \emph{sujeito}, na \emph{vontade} ou \emph{numa faculdade subordinada d vontade e movida por ella}. Porquanto a virtude \emph{moral} não só dá o \emph{poder} de operar bem, mas faz com que produzamos \emph{actualmente} obras boas. Ora só as virtudes, que residem na \emph{vontade} ou numa outra faculdade \emph{dependente} d'ella, podem levar-nos a produzir actualmente. obras boas; pois que, para a pratica do bem, é preciso que a \emph{vontade} se incline para elle. --- Além d'isso, os habitos \emph{moraes} são produzidos pelos actos \emph{moraes}. Ora o acto \emph{moral} ou \emph{honesto} só pôde derivar da \emph{vontade} ou de alguma faculdade dependente d'ella. Logo tambem o habito \emph{moral} só póde residir na \emph{vontade} ou nalguma faculdade que depende d'ella.\par As virtudes \emph{moraes}, embora residam na \emph{vontade}, como em seu \emph{sujeito}, devem comtudo ser reguladas ou dirigidas pela \emph{razão}. Na verdade, as virtudes moraes inclinam a vontade ou as faculdades dependentes d'ella para o acto \emph{bom} ou \emph{recto}. Ora o acto é \emph{bom} ou \emph{recto}, quando a razão o julga tal, isto é, quando se conforma com a regra que a razão impõe, de modo que o mesmo acto corresponda exactamente á regra, sem se inclinar para um extremo nem para outro, evitando do mesmo modo o \emph{excesso} e o \emph{defeito}. Por isso, toda a virtude \emph{moral}, embora se eleve á mais alta perfeição e pratique as mais heroicas emprezas, deve conservar sempre o caracter de moderação e de justa proporção, deve estar sempre em harmonia com os criterios da razão. D'ahi o celebre axioma: \emph{Vértus consistit in medio}. --- Mais tarde veremos como as virtudes \emph{moraes} consistem num meio termo entre dois extremos, egualmente viciosos.}.

75. \textbf{Numero das virtudes moraes}. --- As principaes virtudes \pg{620}  \emph{moraes}, a que todas as outras se reduzem, são \emph{quatro}, a saber: \emph{prudencia}, \emph{justiça}, \emph{temperança} e \emph{fortaleza} \footnote{O numero das principaes virtudes \emph{moraes} é indicado pelo \emph{sujeito}, que ellas \emph{immediatamente} aperfeiçoam. Na verdade, as virtudes \emph{moraes} aperfeiçoam as nossas faculdades em relação aos actos honestos. Logo devem ser tantas, quantas são as principaes faculdades, que concorrem para a producção dos actos honestos. Ora as principaes faculdades, que concorrem para a producção d'esses actos, são quatro, --- duas \emph{espirituaes}, que são a \emph{intelligencia} e a \emph{vontade}, e duas \emph{materiaes}, que são o appetite \emph{concupiscivel} e o \emph{érascivel}. Logo são quatro as principaes virtudes \emph{moraes}: uma aperfeiçoa \emph{directamente} a. \emph{intelligencia}, e é a \emph{prudencia}, --- outra, a \emph{vontade}, e é a \emph{justiça}, --- á terceira, o appetite \emph{concupiscivel}, \emph{e} é a \emph{temperanca}, --- a quarta, o appetite \emph{érascivel}. e é a \emph{fortaleza}.\par Estas virtudes chamam-se \emph{cardeaes}, porque são a base em que assentam todas as outras virtudes, que aperfeicoam a vontade humana, para que esta produza o acto \emph{honesto}, e assim seja moralmente bóa.\par Todas as outras virtudes \emph{moraes} reduzem-se a uma ou outra d'estas quatro \emph{cardeaes} por tres modos, --- ou como partes \emph{integrantes}, emquanto se unem á \emph{cardeal}, para que o acto d'ella seja perfeito, --- ou como partes \emph{subjectivas}, emquanto estão subordinadas á \emph{cardeal}, como \emph{especies} ao \emph{genero}, --- ou como partes \emph{potenciaes}, emquanto exprimem uma parte ou um elemento da virtude \emph{cardeal}, como a \emph{potencia} exprime uma parte do poder da \emph{essencia}.}.

a) Prudencia, --- \emph{Prudencia} é a virtude \emph{moral}, que dirige e inclina a razão \emph{prática} a escolher os meios mais aptos para a producção dos actos \emph{honestos}. --- A \emph{prudencia}, embora, pelo facto de dirigir a \emph{intelligencia} na investigação da \emph{verdade}, se classifique entre as virtudes \emph{intellectuaes}, todavia, emquanto imprime nos actos humanos uma rectidão moral e suppõe a vontade bem disposta para o fim, participa tambem da natureza da virtude \emph{moral} \footnote{Digamos alguma coisa ácerca das \emph{funcções} da \emph{prudencia}, das suas partes \emph{integrantes} e \emph{subjectivas}, e dos \emph{vicios} que lhe são oppostos.\par a) As funcções da \emph{prudencia} são tres: --- \emph{consultar com cuidado}, examinando com diligencia os meios e as cireumstancias, --- \emph{julgar com rectidão}, determinando os meios, que se hão de empregar --- \emph{imperar com efficacia}.\par ph) As partes \emph{integrantes da prudencia} são as oito seguintes: 1.º) \emph{memoria} do passado; 2.º) \emph{intelligencia} ou conhecimento do presente e dos principios moraes; 3.º) \emph{providencia}, pela qual consideramos o fim da obra, e os effeitos futuros; 4.º) \emph{sagacidade} para encontrarmos o meio conveniente; 5.º) \emph{docilidade} para tomarmos conselho; 6.3) \emph{razão} ou promptidao de raciocinio; 7) \emph{circumspecção}, pela qual attendemos a todas as circumstancias; 8.º) \emph{precaução}, pela qual prevenimos os males e os impedimentos.\par 0) As partes \emph{subjectivas} ou \emph{especies da}, \emph{prudencia} são duas: --- prudencia \emph{pessoal} ou \emph{monastica}, pela qual o homem se rege a si mesmo, --- e prudencia \emph{governativa}, pela qual o homem rege-os outros. Esta ultima divide-se em \emph{domestica}, \emph{civil}, \emph{pastoral}, etc.\par d) As partes \emph{potenciaes} da \emph{prudencia} são tres: --- a \emph{eubulia}, que é o\pgn{621}habito do bom conselho, --- a \emph{synesis}, que é o habito de resolver bem as consultas, --- e a \emph{gnoma}, que é o habito de interpretar a mente do legislador.\par e) Como toda a virtude está geralmente entre dois vicios, á \emph{prudencia} oppõe-se, --- por \emph{defeito}, a \emph{imprudencia}, a que se reduzem a \emph{precipitação}, \emph{á inconsideração}, \emph{á inconstancia} e 8 \emph{negligencia}, --- e, por \emph{excesso}, ou melhor, sob a apparencia de \emph{prudencia}, a \emph{prudencia da carne}, \emph{á astucia}, \emph{á excessiva solicitude} pelas coisas temporaes (\emph{Sum}. \emph{Th}., 2-2, q. 47 sqq.).}.

\pg{621} yp----. EO A ao J------------------------------ ------ ------ ---------

CAPITULO HI --- MORALIDADE DOS ACTOS' HUMANOS 621

b) Justiça. --- \emph{Justiça} é a virtude \emph{moral}, que inclina a \emph{vontade} á dar a cada um o que lhe pertence por direito. --- Esta virtude reside, como em \emph{sujeito}, na \emph{vontade} \footnote{Examinemos as partes \emph{integrantes}, \emph{subjectivas} e \emph{potenciaes} da \emph{justiça}.\par a) As partes \emph{integrantes} são duas: --- o habito de \emph{fazer o bem}, que-ê devido a outrem, --- o habito de \emph{fugir do mal}, que prejudica o proximo.\par b) As partes \emph{subjectivas} da \emph{justiça} são tres: 1.3) a justiça \emph{geral} ou \emph{legal}, que leva os subditos e os superiores a tudo o que é devido á sociedade em ordem ao bem commum; --- 2.3) a justiça \emph{distributiva}, que leva os governantes a distribuir os empregos e os beneficios em razão dos meritos de cada um; --- 8.º) a justiça \emph{commutativa}, que obriga o individuo a dar a outro individuo o que lhe deve por verdadeiro e rigoroso direito, de modo que haja egualdade entre o que se dà e o que se deve. --- Estas especies de \emph{justiça} não são \emph{wnivocas}, mas \emph{analogas}. Por quanto a verdadeira natureza da \emph{justiça} encontra-se exclusivamente na justiça \emph{commutativa}.\par 0) As partes \emph{potenciaes} da \emph{justiça} são as virtudes, que se referem a outrem, mas não têm a verdadeira natureza da \emph{justiça}, --- ou porque não podem satisfazer \emph{adequadamente} a divida, --- ou porque restituem o que não é rigorosamente devido. --- A! primeira classe d'estas virtudes (que não podem satisfazer \emph{adequadamente} a divida) pertencem a \emph{religido} para com \textsc{Deus}, a \emph{piedade} para com os parentes e a patria, e o \emph{obsequio} para com os superiores. --- A! segunda classe d'estas virtudes (que restituem o que não é rigorosamente devido) pertencem a \emph{veracidade}, a \emph{gratidão}, a \emph{affabilidade}, \& \emph{liberalidade}, \emph{á equidade} ou \emph{epicheia}, etc. (\emph{Sum}. \emph{Th}., 2-2, q. 58 sqq.).}.

c) Temperanga. --- \emph{T'emperanca} é a virtude \emph{moral}, que inclina o homem a regular as paixões do appetite \emph{concupiscivel}, conforme as preseripeóes da razão. --- A \emph{temperança} reside no appetite \emph{concupiscivel}, emquanto este é movido ou regulado pela nossa vontade \footnote{Vejamos as partes \emph{integrantes}, \emph{subjectivas} e \emph{potenciaes} da \emph{temperanga}.\par a) As partes \emph{integrantes} são: --- a \emph{honestidade}, que é o amor do proprio decoro, --- e o \emph{pudor}, que é a fuga da torpeza contraria á \emph{temperança}.\par bp) As partes \emph{subjectivas} são: --- a \emph{abstinencia}, que modera o uso da comida, --- a \emph{sobriedade}, que regula o modo da bebida, --- a \emph{castidade}, ou \emph{pudicicia}.\par €) As partes \emph{potenciaes} da \emph{temperanca} são: --- a \emph{continencia}, que reprime os movimentos da concupiscencia, para que a vontade não fique vencida, --- a \emph{humildade}, que regula o desejo da propria grandeza, --- a \emph{modestia}, que ordena os movimentos e os gestos do corpo, --- a \emph{mansidão}, que reprime os movimentos da ira, --- a \emph{clemencia}, que leva a dar um castigo inferior á culpa, --- a \emph{estudiosidade}, que modera o desejo de saber, --- e a \emph{eutrapelia}, que conserva a moderação no recreio (Sum. \emph{Th}., 2-2, q. 141 sqq.).}.

\pg{622} d) Fortaleza. --- \emph{Fortaleza} é a virtude \emph{moral}, que inclina o homem a soffrer com paciencia os trabalhos e a emprehender com valor coisas arduas e perigosas. --- Esta virtude reside no appetite \emph{irascivel}, emquanto está subordinado á vontade \footnote{As virtudes \emph{moraes}, que se reduzem \& \emph{fortaleza}, são: --- a \emph{magnanimidade}, que inclina a grandes e gloriosas emprezas, --- a \emph{magnilicencia}, que leva o animo a grandes despezas, se á razão o julgar conveniente, --- a \emph{paciencia}, que dá força ao animo para elle soffrer os males presentes, --- a \emph{perseverança}, que firma a vontade contra as difficuldades provenientes da diuturnidade de um mal ou de uma empreza, --- \& \emph{constancia}, que firma contra os impedimentos externos. --- Estas virtudes, --- se se referirem aos \emph{perigos de morte}, pertencem á \emph{fortaleza} como partes \emph{integrantes}, porque, sem elles, não ha verdadeira \emph{fortaleza}; --- e, se se referirem a objectos menos difficeis, são partes \emph{potenciaes} da fortaleza (\emph{Sum}. \emph{T}., 2-2, q. 123 sqq.).\par *\par Digamos alguma coisa ácerca da \emph{dignidade} das virtudes \emph{moraes}, seu \emph{nexo}, \emph{origem} e \emph{perda}, \emph{augmento} e \emph{diminuição}.\par a) Dignidade das virtudes. --- Na ordem da \emph{dignidade}, a primeira virtude é a \emph{prudencia}, porque modera e regula todas as outras; em seguida é a \emph{Justiça}, que prescreve ao homem o que deve a si mesmo e o que deve aos outros; depois é á \emph{fortaleza}, pela qual o homem, para operar o bem, despreza os perigos e a propria morte; finalmente é a \emph{temperanca}, pela qual o homem prefere a honestidade a todos os prazeres corporeos.\par b) Nexo das virtudes. --- Entre as virtudes moraes \emph{perfeitas} existe um nexo, tão intimo e tão estreito, que uma não póde estar sem as outras, que são necessarias para o ordinario e regular modo de vida. Na verdade, uma virtude é e chama-se \emph{perfeita}, quando o homem, em todas as occasióes, firme e constantemente escolhe e executa o bem /\emph{Honesto}. Ora uma virtude \emph{perfeita} não póde estar sem as outras, que são necessarias para o ordinario e regular modo de vida. Com effeito, para que a virtude seja \emph{perfeita}, é necessaria uma. \emph{prudencia perfeita}, porque a \emph{prudencia} é a regra moderadora de todas as virtudes. Ora a \emph{prudencia} não pode ser \emph{perfeita} sem as outras virtudes \emph{cardeaes}; porque a \emph{prudencia} suppõe a ordenada disposição do \emph{appetite} intellectual e sensitivo, não só em relação a uma virtude, mas em relação a todas as virtudes, pois todos julgam conforme as suas affeições ou desejos. Logo entre todas ás virtudes \emph{moraes} existe um nexo. --- Note-se que fallamos do \emph{Habito} das virtudes \emph{moraes}, e não do \emph{acto}.\par C) Origem e perda das virtudes. --- As virtudes \emph{moraes adquirem-se}, ng ordem \emph{natural}, pela frequencia ou repetição dos actos, e \emph{perdem-se} pelo vicio contrario ou pela cessação do acto.\par a) As virtudes \emph{adqwuirem-se} pela frequencia ou repetição dos actos. Porquanto, á virtude é um \emph{Habito}, isto é, uma inclinação constante para a producção dos mesmos actos, accrescentada \& faculdade. Ora esta inclinação só se adquire, na ordem natural, pela repetição dos actos.\par \emph{b}) As virtudes \emph{perdem-se} pelo vicio contrario, ou pela cessação do acto. --- Perdem-se pelo \emph{vicio contrario}. Na verdade, a mesma faculdade não póde ser sujeito, ao mesmo tempo e sob o mesmo aspecto, de duas qualidades oppostas; e por isso, se adquire um vicio, não póde deixar de perder a\pgn{623}virtude contraria; assim o homem que commette uma injustiça, deixa de ser justo. --- Perdem-se pela \emph{cessação do acto}; não porque o \emph{habito}, para \& sua conservação, depende do \emph{acto}, mas porque as paixões, não sendo mortificadas pelos actos das virtudes contrarias, acabam por subjugar a vontade.\par d) Augmento e diminuição das virtudes. --- O \emph{augmento} e a \emph{diminuição} das virtudes consistem na maior ou menor perfeição e facilidade com que se produzem os seus actos. Este augmento e esta diminuição fundam-se na perfectibilidade e mutabilidade das nossas faculdades.}.

\pg{623} 76. \emph{As} virtudes \textbf{moraes} são \textbf{necessarias para} o \textbf{homem operar} perfeitamente em \textbf{materia} de honestidade. --- Para o homem operar perfeitamente em materia de \emph{honestidade}, É necessario que os seus actos sejam produzidos com \emph{uniformidade}, \emph{facilidade} e \emph{agrado}. Ora estas tres perfeições dos actos honestos só podem derivar das virtudes \emph{moraes}; pois a virtude \emph{moral é habito}, e todo o \emph{habito} tem por fim tornar o acto \emph{uniforme}, \emph{facil} e \emph{agradavel}. Logo as virtudes \emph{moraes} são necessarias para o homem operar perfeitamente em materia de \emph{honestidade} \footnote{Dizem; Os habitos são necessarios, porque facilitam o acto. Ora, para a producção dos actos virtuosos, não precisamos de facilidade alguma; porque dependem da eleição, e muito facil é o que está em o nosso poder. Logo o homem, para operar perfeitamente, não precisa de habitos.\par Resposta. Concedendo a proposição \emph{maior}, negamos a \emph{menor}. Na verdade, as paixões podem por tal forma contrariar os actos virtuosos, que estes, embora dependam de. nós, todavia não se possam executar sem difficuldade. --- Além d'isso, o acto, embora não seja contrariado pelas paixões e dependa da nossa eleição, todavia, sem os \emph{habitos}, não póde realizar-se com facilidade, uniformidade e agrado; pois este modo perfeitissimo de operar não acompanha necessariamente a faculdade.\par Continuam: O \emph{Habito}, determinando e inclinando a vontade para um certo acto, tira a liberdade. Logo, longe de aperfeiçoar, deteriora a faculdade.\par Resposta. Negamos o antecedente. Porquanto, toda a qualidade deixa intacta a essencia do sujeito, em que é recebida. Ora o /\emph{rabito} é recebido na vontade, que é essencialmente livre. Logo.o \emph{Habito} não tira a liberdade, mas só a \emph{inclina} para um certo e determinado acto.}.

77. \emph{As} virtudes \textbf{moraes} augmentam a bondade dos \textbf{actos honestos}. ,--- As virtudes \emph{moraes} augmentarão a bondade dos actos \emph{Honestos}, se augmentarem o affecto da vontade; pois a bondade dos actos honestos não póde deixar de ser proporcionada ao affecto da vontade, da qual elles derivam. Ora as virtudes \emph{moraes} augmentam o affecto da vontade; porque é proprio do \emph{habito} tornar os actos da faculdade, em que residem, mais intensos e mais perfeitos na sua especie. Logo as virtudes \emph{moraes} augmentam a bondade dos actos \emph{honestos} \footnote{Nem se diga que o /\emph{abito}, se augmenta o affecto da vontade, diminue a sua indifferenca e liberdade e por isso o merito do acto. Porquanto, esta diminuição é effeito, pelo menos \emph{indirecto}, da propria vontade, pois deriva\pgn{624}da virtude que se adquiriu com a livre repetição de muitos actos, e por isso não póde diminuir o merito do acto.}.

\pg{624} 78. Vicio. --- Vicio é o habito moral que deteriora e inclina a vontade a praticar actos \emph{repugnantes} á natureza \emph{racional} do homem. --- H' o contrario da virtude \emph{moral} \footnote{No \emph{vicio} devemos considerar a entidade \emph{positiva}, \emph{e} a \emph{malicia}, ou \emph{privação} da rectidão; porque o vicio consiste num bem de ordem inferior, que não é conveniente á nossa natureza \emph{racional} (\emph{Of}. \emph{Sum}. \emph{Th}., 1-2, q. X1 sqq).}.

79. Vicios \textbf{capitaes}. --- Os \emph{vicios} são muitos e diversos, mas todos derivam de \emph{sete}, que por isso se dizem \emph{capitaes}, e que são: \emph{soberba}, \emph{luxuria}, \emph{gula}, \emph{avareza}, \emph{preguiça}, \emph{inveja} e \emph{ira}.

a) Soberba. --- \emph{Soverba} é o vicio, que inclina a vontade a desejar \emph{desordenadamente} a propria estima ou gloria.

b) Luxuria. --- \emph{Luxuria} é o vicio, que inclina a vontade a desejar \emph{desordenadamente} os prazeres sensuaes.

c) Gula. --- \emph{Gula} é o vicio, que inclina a vontade a desejar \emph{desordenadamente} os prazeres sensuaes.

d) Avareza, --- Avareza é o vicio, que inclina a vontade a desejar \emph{desordenadamente} as riquezas.

e) \textbf{Preguiça}, --- \emph{Preguiça} É O vicio, que leva a vontade a entristecer-se \emph{desordenadamente} do bem proprio, só porque a sua consecução dá trabalho ou incommodo.

f) Inveja. --- \emph{Inveja} é o vicio, que leva a vontade a entristecer-se \emph{desordenadamente} do bem alheio, só porque este se considera como mal proprio.

£) Ira. --- \emph{Ira} é o vicio, que leva a vontade não só a entristecer-se \emph{desordenadamente} do bem alheio, mas tambem a desejar a vingança \footnote{A razão por que se enumeram \emph{sete} vicios \emph{capitaes} é a seguinte. O \emph{vicio} consiste --- ou no desejo desordenado do bem, --- ou na fuga desordenada do mesmo bem, motivada por um trabalho ou incommodo que lhe é inherente. Ora --- o bem, que o appetite deseja desordenadamente, --- ou é um bem da \emph{alma}, como é a \emph{gloria} propria, e temos a \emph{soberba}, --- ou é um bem do \emph{corpo}, e temos a \emph{lucuria} e á \emph{gula}, --- ou é um bem \emph{externo}, como é a \emph{riqueza}, e temos a \emph{avareza}; --- e o bem, de que o appetite foge desordenadamente por causa do trabalho ou incommodo, que á sua consecução póde causar, --- ou é um bem \emph{proprio}, e temos á \emph{preguiça}, --- ou é um bem do \emph{proximo}, e temos a \emph{ira} e a \emph{inveja}, conforme a tristeza pelo bem \emph{prazimo} é, ou não, acompanhada do desejo da \emph{vingança}. --- O numero dos vicios \emph{capitaes} excede o das virtudes \emph{Moraes}, porque a \emph{virtude} consiste no \emph{meio termo}, de que podemos afastar-nos ou por \emph{excesso} ou por \emph{defeito}.\par Os \emph{vicios} não estão ligados entre si, como o estão as \emph{virtudes}. Porquanto, toda a virtude consiste em seguir a regra da razão, que a \emph{prudencia} prescreve; e por isso todas as virtudes suppõem a \emph{prudencia} e por causa d'esta conservam uma união entre si. Mas o \emph{vicio} não consiste propriamente em se\pgn{625}afastar da regra da razão (porque o afastamento é uma \emph{privação}, e esta não póde ser o termo das nossas tendencias), mas consiste em desejar um bem inferior, que não é conveniente á nossa natureza racional, --- desejo que importa um afastamento da regra da razão, isto é, uma \emph{privação} de honestidade. E como os bens, que formam o objecto do vicio, são diversos e muitas vezes contrarios entre si, deve dizer-se que não existe nexo algum entre os vicios; pois os habitos distinguem-se e especificam-se pelos seus objectos. --- Todavia devemos dizer que os vicios \emph{capitaes} estão ás vezes por tal forma unidos entre si, que um deriva de outro, --- não emquanto o objecto de um vicio deriva do objecto de outro vicio, --- mas emquanto um vicio inclina para outro vicio; assim a \emph{gula} inclina para a \emph{luxuria}, \emph{á avareza} para a \emph{ira}. E' neste sentido que o \emph{orgulho} se chama o pae de todos os vicios.\par Do que deixamos escripto vê-se a falsidade da opinião dos \emph{Stoicos}, segundo a qual todos os vicios seriam \emph{eguaes}. Porquanto, um vicio póde, mais do que um outro vicio, afastar-se do meio termo onde está a virtude, e por isso ter mais malicia do que um outro vicio; assim como o erro, que mais se afasta da verdade, é mais grave do que um outro erro, que se afasta menos. --- Os \emph{Stoicos}, dizendo que todos os vicios são \emph{eguaes}, baseavam-se nas duas seguintes razões: 1.º) todo o vicio consiste na \emph{privação} da honestidade, ora a \emph{privação} não admitte graus; --- 2.º) a virtude consiste no \emph{indivisel}; logo tambem no \emph{indivisel} deve consistir o vicio, que é o contrario da virtude. Mas estas razões não são concludentes. Não é concludente a primeira; pois vimos que o \emph{vicio} não consiste na mera \emph{privação} da honestidade. Não é concludente a segunda; porque, embora houvesse proporção entre as virtudes e os vicios, estes deveriam ser deseguaes, porque tambem as virtudes são deseguaes, não só em relação ao seu \emph{objecto}, mas tambem em relação á sua \emph{intensidade}.}.

\pg{625} 80. Os vicios augmentam a malicia dos \textbf{actos humanos}. --- Os \emph{vicios} augmentarão a malicia dos actos humanos, se augmentarem o affecto da vontade; pois deve existir uma correspondencia entre os effeitos e a sua causa. Ora os \emph{vicios} augmentam o affecto da vontade; porque todos os habitos tornam os actos mais intensos e completos na sua especie. Logo os \emph{vicios} augmentam a malicia dos actos humanos \footnote{Devemos advertir que, quando o homem, levado por um sincero arrependimento, detesta os seus vicios, estes já não são \emph{voluntarios}; e por isso a intensidade maior, que os mesmos vicios dão aos actos moralmente maus, deixa de ser imputavel ao agente.\par As \emph{virtudes} e os \emph{vicios} encontram poderosos auxiliares nas \emph{paixões}, conforme a direcção que estas receberem da vontade; porque as \emph{paixões} augmentam a intensidade do \emph{voluntario}, tanto na decisão como na execução. Quando o bem, para o qual tende a \emph{paicüo}, é contrario ao verdadeiro bem da nossa natureza racional, então a paixão é auxiliar do vicio; como, pelo contrario, quando o bem, para o qual inclina a paixão, é conforme com a nossa natureza racional, nesse caso a paixão é auxiliar da virtude. Por isso a vontade deve moderar as paixões, dirigindo-as para o bem conveniente \& nossa natureza racional. A victoria contra as paixões exige trabalho e sacrificio; mas nunca houve victoria mais util, nem mais gloriosa.\par 40 Vor. II.}.

{\small\noindent\textsc{Summario}: --- Lei moral, sua divisão e existencia. --- Consciencia moral, sua divisão e direcção.\par}

\pg{626} \textbf{REGRAS} DOS ACTOS HUMANOS

%%>>> art-58

\section*{Artigo I.}
\addcontentsline{toc}{section}{Artigo I.}

\textbf{Lei moral}, \textbf{sua divisão} e \textbf{existencia},

81. \textbf{Regras} dos \textbf{actos} humanos. --- \emph{Regras} dos actos humanos são as normas \emph{praticas}, que impõem á vontade uma necessidade \emph{moral} ou \emph{obrigação}. --- Estas regras são duas: a \emph{let moral} e a \emph{consciencia moral}. A \emph{lei moral} é regra \emph{remota} e \emph{externa}; a \emph{consciencia moral} é regra \emph{proxima} e \emph{interna}, pois é a razão \emph{pratica} emquanto applica a \emph{Ze}? aos casos particulares.

Neste artigo tratamos da \emph{lei}; no seguinte, da \emph{consciencia} \footnote{\emph{A} norma ou a regra de \emph{moralidade}, de que fallamos no cap. III d'este tratado, é apenas \emph{t/ieorica}, pois limita-se a \emph{indicar} quaes são os actos intrinsecamente bons ou maus; ao passo que as normas ou regras dos actos humanos, de que nos estamos occupando, são \emph{praticas}, porque impõem uma \emph{obrigação}. --- Sendo a \emph{consciencia moral} uma \emph{applicação} da \emph{lei moral}, póde dizer-se que a regra ou norma \emph{pratica} dos actos humanos é uma só --- a lei \emph{moral}.}.

82. Lei moral. --- \emph{Let moral é um preceito racional para o bem commum}, \emph{promulgado por quem governa a communidade}.

a) A \emph{lei moral} é um \emph{preceito}; porque não se limita a mero conselho ou exhortação, mas recebe da \emph{vontade} do superior a efficacia ou força de impor uma \emph{obrigação}.

b) A \emph{lei} é um preceito \emph{racional}; não só porque toda a \emph{let} deve ser conforme á \emph{razão}, mas tambem porque é proprio da \emph{lei} ordenar os actos dos subditos para o bem commum, e só a \emph{razão} é capaz de ordenar ou dispor os meios para o fim.

c) A \emph{lei} é para o \emph{bem commum}; porque o fim da \emph{lei} é, e não póde deixar de ser, o bem da communidade.

d) A \emph{lei} a um preceito \emph{promulgado}; porque a vontade do superior não póde obrigar os seus subditos, se não lhes for sufficientemente manifestada.

e) A \emph{lei} deve ser promulgada \emph{por quem governa a communidade}; porque só este, devendo procurar o bem commum, pôde e deve dirigir e obrigar os seus subditos \footnote{E' a definição de lei, dada pelo angelico Doutor: « Lex est quaedam ordinatio rationis ad bonum commune et ab eo, qui curam communitatis\pgn{627}habet, promulgata » (\emph{Sum}. \emph{T}., 1-2, q. 90, a. 4). O mesmo Santo explica a definição nos quatro artigos d'essa questão.\par As seguintes observações são muito necessarias.\par a) A \emph{lei} é a \emph{regra} dos nossos actos, mas é uma regra \emph{efficaz}, isto é, que impõe uma \emph{obrigação}, Emquanto é \emph{regra}, a \emph{lei} deriva da \emph{razão} do superior: porque só á \emph{razão}, conhecendo a relação entre os meios e o fim, póde dirigir os actos livres para o bem commum. Emquanto é regra \emph{efficaz}, isto ê, emquanto impõe uma \emph{obrigação}, a \emph{lei} deriva da \emph{vontade} do superior; porque só o acto d'esta faculdade póde induzir os subditos a pór em pratica certos e determinados meios para a consecução do fim commum. S. Thomaz diz que « o imperio é um acto da razão, que presuppõe o acto da vontade ». (\emph{Sum}. \emph{Th}., 1-9, q. 17, a. 1. --- Sendo uma \emph{norma efficas}, a \emph{lei} é composta de dois elementos: um \emph{commum} e outro \emph{proprio}. O \emph{commum} é constituido pelo acto da \emph{rasdo}, e consiste em \emph{ordenar}, \emph{ensinar}, \emph{regular}, etc., e é elemento \emph{commun}, porque tambem são \emph{normas} o \emph{conselho}, a \emph{exhortação}, ete.; --- o \emph{proprio} é constituido pelo acto da \emph{vontade}, e consiste em impor, como \emph{obrigacdo}, a norma apontada pela razão, e diz-se elemento \emph{proprio} ou \emph{especifico}, pois é por elle que a \emph{lei} é constituida na sua especie, distinguindo-se do do \emph{conselho}, da \emph{exhortação}, ete.\par b) A \emph{obrigação}, imposta pela \emph{lei} moral, é uma \emph{necessidade ab}; --- * \textbar{} \emph{O} que incumbe ao subdito de sujeitar os seus actos á vontad que deriva do nexo \emph{necessario} de certos e determinados meios \emph{necessario}. --- A \emph{obrigação} --- 6 uma necessidade \emph{absoluta}, porque replica da parte do subdito; --- é uma necessidade \emph{moral}, pois não liberdade \emph{physica}, mas a liberdade \emph{moral}, emquanto o subdito deve escolhét e pôr em execução os meios que estão /;\emph{ecessariamente} ligados com o fim, que se deve alcançar. --- A \emph{obrigação}, portanto, é o effeito \emph{proximo} da lei; o bem dos subditos e da communidade é o seu effeito \emph{remoto} e \emph{final}.\par 0) A \emph{promulgação} é a intimação absoluta e externa da lei, em harmonia com as formalidades legitimamente estabelecidas. Esta intimação é necessaria; porque a \emph{lei} é norma, e toda a norma deve ser applieada ao agente; aliás este não poderá conformar com ella os seus actos. --- A \emph{promulgação} é diversa da \emph{notificação}, que é a manifestação da lei a todos os subditos em geral e a cada um em particular. Na verdade, a \emph{promulgação} é um complemento necessario da lei, é um acto do legislador, e deve ser feita segundo as formalidades legitimamente estabelecidas, pelas quaes se conheça a intenção que o legislador tem de obrigar os subditos; ao passo que a \emph{notificação} suppõe a lei já constituida, póde ser feita por qualquer empregado, e por todo e qualquer meio. Por isso os magistrados ou governadores notificam\pgn{628}as leis, mas não as promulgam; as leis existem e vigoram onde não foram notificadas, embora os subditos, pela sua ignorancia, não sejam responsaveis dos actos praticados contra ellas.\par á) A \emph{lei} póde considerar-se num triplice estado, a saber --- 1.º) na mente do legislador, --- 2.º) na mente dos subditos, aos quaes foi promulgada, --- 9.9) no signal externo por que se intima aos subditos. --- A lei, no primeiro estado, é \emph{activa}; no segundo e no terceiro, é \emph{passiva}. A lei, para ser completa, basta que exista na mente do \emph{legislador} e dos \emph{subditos}; o \emph{signal externo} não é necessario para toda a lei. --- A \emph{lei}, considerada \emph{activamente}, isto é, na mente do \emph{legislador}, suppõe varios actos; o \emph{juizo}, pelo qual o superior conhece que uma certa e determinada norma é util e necessaria para a communidade; --- o acto da \emph{vontade}, pelo qual o superior quer obrigar os subditos a observar aquella norma; --- o acto da \emph{razão}, pelo qual o superior refere a propria vontade aos subditos, e que se chama \emph{mperio}. O \emph{imperio}, pois, embora seja acto da razão suppõe essencialmente o acto da vontade, de que tira toda a sua efficacia.\par e) Erra, pois, \emph{Montesquieu}, quando faz consistir \emph{a lei moral nas relações essenciaes}, que \emph{nascem da natureza humana} (\emph{De lesprit des lois}, l. I, o. 2). Porquanto, embora a \emph{lei moral} deva necessariamente corresponder ás relações essenciaes das coisas, comtudo aquella definição não exprime, como devia, a relação ás operações, nem a relação do. inferior ao superior, nem a idea de \emph{obrigação}. --- Erram tambem, ácerca da verdadeira natureza da \emph{lei moral}, os \emph{pantheistas}, os \emph{materialistas}, que ou não admittem em nós o \emph{livre arbitrio}, ou dizem que o homem não está subordinado a nenhuma auctoridade. Por isso \emph{Littré}, como vimos, reduziu a \emph{lei moral} ao \emph{instincto}.}.

\pg{627} 83. \textbf{Especies da} lei moral. --- A lei \emph{moral}, conforme os differentes aspectos sob os quaes se considera, é \emph{eterna} ou \emph{temporal}, --- \emph{natural} ou \emph{positiva}, --- \emph{divina} ou \emph{humana}.

a) Lei eterna ou temporal --- \emph{Eterna} é a que existe, \emph{desde} a \emph{eternidade}, na mente do legislador. \emph{Temporal} é a que o legislador concebeu e formou no \emph{tempo}.

b) Lei natural ou positiva. --- \emph{Natural} é a que se promulga pela luz \emph{natural} da razão. \emph{Positiva} é a que se promulga por uma expressão, oral ou escripta, do legislador.

\pg{628} c) Lei divina ou humana. --- \emph{Divina} é a que deriva de \textsc{Deus}. \emph{Humana} á a que \emph{immediatamente} deriva do homem, mas em virtude de poder, que \textsc{Deus} lhe communicou. --- A lei \emph{humana}, que é sempre \emph{temporal} e \emph{positiva}, divide-se em \emph{ecclesiastica} ou

4 \emph{civil}, conforme emana do poder \emph{eccleszastico} ou \emph{civil}. :

Diremos da lei \emph{eterna}, \emph{natural} e \emph{positiva} \footnote{A \emph{lei moral}, em razão dos seus \emph{e/feitos}, divide-se em \emph{preceptiva}, \emph{prohibitiva}, \emph{permissiva} e \emph{punitiva}. --- Lei \emph{preceptiva}, que se diz tambem \emph{affirmativa}, 6 \& que obriga a certos e determinados actos: transgride-se pela \emph{omissão}. --- Lei \emph{prohibitiva}, que se diz tambem \emph{negativa}, é a que prohibe a realização de certos e determinados actos: transgride-se, praticando esses actos. --- Lei \emph{permissiva} é a que \emph{positivamente} permitte a pratica de certos e determinados actos: impõe a obrigação de não impedir o acto permittido. --- Lei \emph{punitiva} é a que applica o castigo ao transgressor.}.

84. \textbf{Lei eterna}. --- Lei \emph{eterna é o acto da razão} e \emph{vontade de} \textsc{Deus}, \emph{emquanto}, \emph{desde a eternidade}, \emph{prescreve ds creaturas vacionaes a conservação da ordem moral} e \emph{prohibe a sua perturbação}. --- Ordem \emph{moral} é, como dissemos, a harmonia entre os actos \emph{moraes} do homem e a sua natureza \emph{racional} \footnote{\textsc{Deus}, tendo estabelecido a ordem que existe no mundo, exige que todas as creaturas a observem. D'ahi a lei \emph{eterna}. Mas esta lei é diversa, conforme a \emph{ordem}, que se deve observar. Se a ordem é \emph{physica}, a lei tambem\pgn{629}6 \emph{physica}, e \textsc{DEUS} applica, no tempo, esta lei aos movimentos \emph{necessarios} das creaturas. Se a ordem é \emph{moral}, a lei tambem é \emph{moral}, e \textsc{Deus} appliea, no tempo, esta lei aos movimentos \emph{livres} das creaturas. --- Aqui tratamos da lei eterna \emph{moral}, isto é, da lei eterna, emquanto prescreve a observancia da ordem \emph{moral} e prohibe a sua violac\&o.\par A lei \emph{eterna}, embora não possa dizer-se \emph{necessaria}, emquanto suppóe o acto da vontade de \textsc{DEus}, que livremente criou o mundo, comtudo deve dizer-se \emph{necessaria}, emquanto, supposto que \textsc{Deus} quizesse criar o mundo, não podia deixar de prescrever uma lei ás suas creaturas e de executar por meio da sua providencia infinita.}.

\pg{629} 85. A lei \textbf{eterna existe}.

a) A lei \emph{eterna} é o acto da razão e vontade de \textsc{Deus}, emquanto, desde a eternidade, presereve a conservação da ordem \emph{moral} e prohibe a sua perturbação. Ora semelhante acto existe em \textsc{Deus}. Na verdade, \textsc{Deus}, governando com sapiente providencia todas as coisas, deve ter na sua mente, desde a eternidade, a regra da ordem \emph{moral}, que deve dirigir cada creatura para o seu fim, --- e, sendo infinitamente santo e justo, não póde deixar de exigir, desde a eternidade, que tal ordem seja inviolavelmente observada por cada creatura racional. Logo a lei \emph{eterna} existe.

b) A lei \emph{mora}? é o dictame \emph{efficaz} da razão \emph{pratica} do superior, a quem pertence o governo da communidade. Ora Deus, governando todas as creaturas, sobretudo as racionaes, deve ter em si, desde a eternidade, o dictame \emph{efficaz} da Sua razão \emph{pratica}, pelo qual mova cada uma das creaturas para o seu fim. Logo a lei \emph{eterna} existe \footnote{S. Thomaz trata admiravelmente, como é seu costume, da \emph{lei eterna} na \emph{Sum}. \emph{Th}., 1-2, q. 91, a. 1; q. 98. --- Embora a lei eterna pertenca á providencia, e esta áquella, comtudo a lei e a providencia são coisas distinctas. Por quanto --- \emph{a}) a lei eterna é a razão primeira e original do governo das coisas e refere-se á providencia, como o principio universal se refere a conclusão; --- 4) a lei tem a força de obrigar, a providencia prescinde d'essa força. Diz S. Thomaz « Providentia in Deo proprie non nominat legem aeternam, sed aliquid ad legem aeternam \emph{consequens}. Lex enim aeternà consideranda est in Deo, sicut accipiuntur in nobis principia operabilium naturaliter nota, ex quibus procedimus in consiliando et eligendo: quod est prudentiae, sive providentiae. Unde hoc modo se habet lex intellectus nostri ad prudentiam, sicut principium in demonstratione. Et similiter etiam in Deo lex aeterna non est providentia, \emph{sed providentiae quasi principium}. Unde et convenienterlegi aeternae attribuitur actus providentiae; sicut et omnis effectus demonstrationis principiis indemonstrabilibus attribuitur » (\emph{De verit}, q: 5, a. 1 ad 6).\par \emph{Thomasio} (1655-1128) rejeita a doutrina Acerca. da \emph{lei eterna}, comparando-a á opinião dos pagãos Acerca de \emph{materia eterna}. \emph{Stahl} diz que a \emph{lei eterna} não existe, porque seria uma restricção da \emph{liberdade de} \textsc{DEUS} nos seus\pgn{630}actos \emph{transeuntes}. --- Mas ambos erram. Erra \emph{TÀomasio}; porque, devendo a lei estar na razão e na vontade do superior anteriormente á sua applicação, é evidente que a lei, pela qual \textsc{Deus} governa o mundo, deve ter existido, desde a eternidade, na sua razão e vontade infinita. Erra \emph{Stahl}; porque, sendo a lei \emph{eterna} destinada para o governo das creaturas, que \textsc{Deus}, desde a eternidade, conhece como existentes no tempo, em virtude do \emph{livre} decreto de creacüo, não se percebe como essa lei possa restringir a \emph{liberdade} de \textsc{Deus} nos seus actos \emph{transeuntes}.\par Objectam: A \emph{promulgação é essencial} a toda a lei, e por isso a \emph{lei eterna}, não sendo \emph{promulgada}, não é verdadeira lei.\par Resposta. A \emph{promulgação} da lei póde considerar-se no seu \emph{principio}, que é o superior, e no seu \emph{termo}, que é constituido pelos subditos. Considerada no seu \emph{principio}, \emph{á promulgação} é necessaria, emquanto ella é um decreto que fixa irrevogavelmente o acto da razão e da vontade do legislador. Considerada no seu \emph{termo}, a \emph{promulgação} não é \emph{necessaria} para a \emph{essencia} da lei, mas só para o \emph{e/feito} da propria lei, que é a \emph{obrigação}. Ora a lei \emph{eterna} tem a \emph{promulgação}, considerada no seu \emph{principio}; porque em \textsc{Deus} existe, desde a eternidade, o decreto immutavel de \emph{promulgar} no tempo o acto eterno da sua vontade, ao qual devem sujeitar-se as creaturas.}.

\pg{630} 86. \textbf{Lei natural}. --- \emph{Let natural} é o acto da razão e vontade de Dzvus, que prescreve a observancia da ordem moral e prohibe a sua violação, e que se \emph{manifesta ás creaturas na luz natural da razão}. --- Diremos alguma coisa ácerca do seu \emph{objecto}, \emph{existencia}, \emph{propriedades}, \emph{principio} e \emph{sancção} \footnote{Os dictames, pois, da lei \emph{eterna}, que se referem ás nossas operações, emquanto são conhecidos pela razão, consituem a \emph{lei natural}. Por isso, á nossa razão, assim como, emquanto é \emph{theorica}, conhece, pela sua propria luz, os primeiros principios especulativos e d'elles deduz as conclusões, assim tambem, emquanto é \emph{pratica}, conhece os primeiros principios da honestidade, e por elles conhece o bem que deve fazer e o mal que deve evitar.\par Sendo a \emph{let natural} uma participação da \emph{let eterna}, é impossivel que o acto, que repugna á recta razão, não seja uma transgressão da lei divina. --- E\textasciicircum{}, pois, absurda a divisão do \emph{peccado} em \emph{theologico} (que é a transgressão moral da lei divina) e \emph{philosophico} (que seria um acto exclusivamente contrario á recta razão). O peccado \emph{philosophico} é tambem \emph{theologico}; e esta identidade é attestada pela propria consciencia, a qual receia o castigo de um poder superior, quando commette uma falta contraria á recta razão.\par Veja-se a \emph{Sum}. \emph{Th}., 1-2, q. 91, a. 2; q. 94.}.

87. \textbf{Objecto da lei natural}. --- O \emph{objecto} da lei \emph{natura}? é constituido por todas as verdades \emph{praticas}, conhecidas pela luz natural da razão. --- Estas verdades reduzem-se a tres categorias: umas exprimem os \emph{primeiros principios moraes}, e são immediatamente evidentes para todos; outras são consequencias \emph{émmediatas} d'esses primeiros principios, e conhecem-se por uma simples reflexão; outras são consequencias \emph{remotas} dos primeiros principios, e sô se conhecem pelo raciocinio. --- As verdades das duas primeiras categorias são o objecto \emph{primario} \pg{631}  da lei \emph{natural}; as da terceira categoria constituem o objecto \emph{secundario} \footnote{E', pois, evidente que a lei \emph{eterna}, --- emquanto se \emph{reflecte} na \emph{creatura racional}, tem a mesma extensão que a lei \emph{nataral}, --- mas, emquanto se considera em si mesma, isto é, emquanto é a regra que dirige todas as creaturas para os seus fins, tem uma extensão maior que a da lei \emph{natural}.}.

88. A lei \textbf{natural existe}.

a) \textsc{Deus}, para governar e dirigir todas as creaturas para o fim \emph{ultimo}, deve applicar a cada d'ellas a sua lei \emph{eterna}; pois não ha governo sem a applicação das leis respectivas. Esta applicação da lei \emph{eterna} deve ser proporcionada á natureza das creaturas; e por isso, se essa lei ás creaturas irracionaes é applieada, emquanto recebem uma inclinação necessaria para o seu fim, ás creaturas \emph{racionaes} deve ser applicada, emquanto a conhecem pela luz natural da sua razão. Ora a lei \emph{eterna}, emquanto se manifesta ás creaturas na luz natural da razão, constitue a \emph{lei natural}. Logo a lei \emph{natural} existe.

b) A consciencia intima attesta --- que o homem, pela luz natural da sua razão, não só conhece a distineção entre 0 bem e o mal moral, mas se sente obrigado a obedecer, queira ou não, a uma lei superior, que manda alguns actos honestos e prohibe outros n\&o-honestos, --- e que, se obedece a esta auctoridade, experimenta paz e socego de espirito, e, se a despreza, prova a censura do remorso e o medo de castigo. Ora esta lei superior ao homem, que manda a pratica do bem e a fuga do mal, e que se conhece pela luz da razão é a \emph{let natural}. Logo a \emph{lei natural} existe.

c) Todos os povos, em todos os tempos, admittiram a existencia da \emph{lei natural}, como um reflexo e manifestação da \emph{let eterna}; pois conheciam, pela luz natural da razão, que estavam obrigados a praticar alguns actos, e a evitar outros, sob pena de incorrer na inimizade e nos castigos do Ser supremo, --- conhecimento, que se traduziu no empenho de aplaear, sobre tudo por meio do sacrificio, a. Divindade irritada. Este consenso \emph{unanime}, \emph{universal} e \emph{constante}, derivou --- ou da moção da vontade, --- ou da evidencia da coisa. Ora não derivou da moção da vontade; pois a lei \emph{natural} é gravosa e molesta para a vontade. Logo derivou da evidencia da coisa, e por isso da verdade; porque a evidencia \emph{objectiva} é criterio de verdade. Logo a lei \emph{natural} existe \footnote{A existencia da lei \emph{natural} foi admittida por todos os povos, não só christios, mas tambem pagãos. Os mais insuspeitos escriptores, antigos\pgn{632}e modernos, como \emph{Aristoteles} (\emph{Rhet}., I, 18), \emph{Sophocles} (\emph{Antig}., v. 454), \emph{Voltaire} (\emph{La loi naturelle}), \emph{Rousseau} (\emph{L}' \emph{Emile}, l. V, exprimem a convicção de todos os povos ácerca de uma \emph{regra para toda} a \emph{especie humana}, que é \emph{anterior á opinião}, \emph{e de que tiram a força e a que se reduzem todas as outras regras} (\emph{Rousseau}). Mas ninguem falou da lei \emph{natural} com maior exactidão e eloquencia do que \emph{Cicero}. O grande orador romano diz assim: « Est quidem vera lex recta ratio, naturae congruens, diffusa in omnes, constans, sempiterna, quae vocet ad officium jubendo, vetando a fraude deterreat, quae tamen neque probos frustra jubet aut vetat, nec improbos jubendo aut vetando moyet. Huie legi nec abrogari fas est, neque derogari ex hac aliquid licet, neque tota abrogari potest, nec vero aut per senatum, aut per populum solvi hae lege possumus, neque est quaerendus explanator aut interpres ejus alius, neo erit alia lex Romae, alia Athenis, alia nunc, alia posthac, sed et omnes \emph{gentes} et \emph{omni tempore} una lex et \emph{sempiterna} et \emph{immutabilis continebit}, \emph{wunusque erit communis quasi magister et imperator omnium} \textsc{Deus}; \emph{ille legis huius inventor}, \emph{disceptator}, \emph{lator}, cui qui non parebit, ipse se fugiet, ac naturam hominis aspernatus hoc ipso luet \emph{maximas poenas}, etiam si cetera supplicia quae putantur effugerit » (\emph{Frag}. \emph{de Rep}., l. III).\par Os \emph{atheus} (\emph{praticos}), os \emph{materialistas}, os \emph{pantheistas}, os \emph{racionalistas}, os \emph{fatalistas}, os \emph{positivistas}, e outros infelizes da mesma raga, negam a existencia da lei \emph{natural}, e attribuem os dictames de \emph{moralidade}, que a razão manifesta a cada um, á \emph{educação}, ou aos \emph{enganos} dos sacerdotes, ou ás \emph{astucias} dos legisladores. --- Mas essas opiniões são absurdas. A \emph{educação}, os \emph{enganos} dos sacerdotes, as \emph{astucias} dos legisladores, são causas \emph{particulares} e \emph{diversas}, conforme a diversidade dos povos, e por isso incapazes de produzir os dictames \emph{universaes} e \emph{constantes} da lei \emph{natural}; aliás o effeito seria superior á causa. Um effeito \emph{universal} e \emph{constante} só póde derivar de uma causa \emph{universal} e \emph{constante}, como é a \emph{natureza humana}, que é sempre uma e a mesma em todos os homens.\par *\par Como o effeito proprio da lei é a \emph{obrigação}, é claro que a \emph{obrigação}, que a lei \emph{natural} impõe, deriva de \textsc{Deus}; porque \textsc{Deus} é o auctor da lei \emph{natural}. Esta \emph{obrigação} deriva do \emph{direito} que \textsc{DEus} tem sobre as creaturas, --- \emph{direito}, a que, por parte das creaturas, corresponde a devida obediencia.\par Alguns escriptores desvirtuam a natureza e a origem da \emph{obrigação}, que peza sobre a nossa consciencia e a que devemos sujeitar o nosso livre arbitrio. Limitamo-nos a citar as opiniões de \emph{TAÀomasio}, de \emph{Kant}, de \emph{Spencer}.\par a) \emph{Thomasio}, ensinando que os dictames da \emph{lei natural} não são \emph{preceitos} mas \emph{conselhos}, não reconhece verdadeira \emph{obrigação}, imposta por essa lei.\par Esta opinião é falsa. A consciencia da humanidade pensa que os dictames da lei \emph{natural} são verdadeiros \emph{preceitos} e nao meros \emph{conselhos}. Porquanto, cada um de nós sente o \emph{dever} de obedecer a esses dictames, ao qual não\pgn{633}póde nem deve subtrahir-se. Ora o \emph{dever} corresponde necessariamente a um \emph{preceito}, \emph{e} não a, um simples \emph{conselho}.\par b) \emph{Kant} desvirtuou a idea da \emph{obrigação}, imposta pela lei \emph{natural}, inventando a sua theoria acerca da \emph{autonomia} da razão. O philosopho allemão admitte que existe uma lei \emph{moral} universal, dotada de verdadeira força de \emph{obrigar} os homens, mas nega que esta lei seja participação e reflexo de uma lei superior. Para elle, a lei \emph{moral} ou \emph{natural} consiste no \emph{dictame meramente formal} e \emph{absoluto da} razão \emph{pratica}. A razão \emph{pratica é autonoma} ou \emph{nomothetica}, e a sua dignidade consiste em não estar sujeita senão á lei, que ella prescreveu a si mesma. Esta lei, pelo facto de emanar da razão, merece toda a veneração. A forma absoluta, por que a lei emana da razão pratica, é e diz-se \emph{imperativo categorico}.\par O systema de \emph{Kant} deve rejeitar-se pelas seguintes razões:\par a) \emph{Basêa-se num falso fundamento}. Porquanto \emph{Kant} suppõe que o \emph{homem}, \emph{emquanto ente racional e livre}, é o \emph{fim de si proprio}, não só \emph{relativamente} aos outros seres, mas \emph{absolutamente}. Ora este supposto é absurdo; porque o fim \emph{ultimo} do homem, como provámos, é \textsc{Deus}.\par \emph{b}) \emph{E}' \emph{absurdo na} sua \emph{essencia}. A essencia d'este systema consiste na \emph{autonomia} ou \emph{independencia} da \emph{razão humana}. Ora esta autonomia repugna. Na verdade, ninguem póde prescrever uma lei a si proprio; porque a lei suppõe o subdito e o superior, e ninguem é subdito de si mesmo; Ora se a \emph{obrigação} da lei natural derivasse da nossa razão, o homem legislaria a si mesmo. --- Além d'isso, se a lei \emph{natural} derivasse da nossa razão, esta faculdade teria estabelecido a ordem \emph{moral}; porque a ordem \emph{moral} funda-se na lei \emph{natural}. Ora á ordem \emph{moral}, ou a \emph{moralidade}, funda-se, \emph{remotamente}, no proprio \textsc{Deus}, como principio e fim de todas as coisas. --- Ha mais, a consciencia attesta que a nossa razão \emph{descobre} a lei \emph{natural}, mas não a \emph{fórma}; porque essa lei manifesta-se-nos como promulgada por uma auctoridade superior, á qual o homem não póde desobedecer, sob pena de expiar a sua transgressão. --- Finalmente, a nossa razão, sendo fallivel e mutavel, póde e deve ser dirigida para a consecução do seu fim. Ora um ente, que precisa de direcção, não póde dirigir-se a si mesmo, mas deve ser dirigido por outro ente, que, em ultima analyse, seja infinitamente santo e perfeito.\par \emph{c}) \emph{Encerra muitas contradicções}. Na verdade, \emph{Kant} admitte que a razão humana é \emph{absoluta} e illimitada, porque o seu dictame merece \emph{por si} uma veneração illimitada e absoluta, e tambem admitte que ella (a razão humana) é limitada e relativa, porque é creada e por isso derivada e contingente. --- Além disso, segundo \emph{Kant}, a razão humana, ou vontade racional é nomothetica ou autonoma e, ao mesmo tempo, é \emph{subdita} e, muitas vezes, \emph{rebelde}.\par \emph{d}) \emph{Leva ao panteismo}. Na verdade, uma razão independente é só propria do Ente absoluto e necessario, que é \textsc{Deus}. (Veja-se o que dissemos ácerca do \emph{imperativo categorico} de \emph{Kant} no cap. IIL).\par b) \emph{Herbert Spencer} explica, pela sua theoria da \emph{evolução}, a origem e a natureza da \emph{obrigação}, que todos sentimos na nossa consciencia. O eseriptor inglez comeca por dizer que ha um duplice genero de \emph{sancções}, que limitam\pgn{634}á nossa actividade: umas \emph{intrinsecas}, outras \emph{ectrinsecas}. As \emph{intrinsecas} são os varios effeitos maus que acompanham os nossos actos, como são as \emph{dores}, os \emph{dumnos}, etc. As \emph{ewtrinsecas} são constituidas pelo medo \emph{dos companheiros}, \emph{dos principes}, etc. A estas sancções começou gradualmente a ligar-se a idea de \emph{coacção} ou \emph{obrigação}; de modo que o homem sente a obrigação de evitar um acto, emquanto este acto produz dór ou incommodo. A idea de \emph{coacção} ou \emph{obrigação}, pela \emph{evolução}, diminuirá pouco a pouco, até desapparecer completamente. Depois d'este desapparecimento total, o homem operará exclusivamente pelo \emph{prazer}, que acompanha sempre o bem proprio e o commum.\par Esta theoria deve rejeitar-se pelas seguintes razões:\par a) Suppõe que o homem é a \emph{evolução} de uma especie inferior, e por isso nega a distincção essencial entre o homem e o animal, e rejeita a \emph{espiritualidade} e a \emph{liberdade} da nossa alma. Ora tudo isto, como vimos, é absurdo.\par \emph{b}) Longe de explicar a \emph{obrigação}, nega-a e destroe-a; porque ninguem dirá que um homem.se guia pelo sentimento da \emph{obrigação}, quando opéra pelo \emph{medo} das penas temporaes.\par c) A \emph{obrigagüo}, no systema de \emph{Spencer}, não tem nenhuma força ou efficacia. A \emph{obrigação} suppõe duas vontades não só distinctas, mas tambem subordinadas, de modo que uma imponha a \emph{obrigação} e a outra a execute. Ora, no systema de \emph{Spencer}, o homem não tem superior, e por isso não póde estar \textbf{sujeito}. (Veja-se o que escrevemos contra \emph{Spencer} no cap. III).}.

\pg{632} 89. \textbf{Propriedades da lei natural}. --- As principaes \emph{propriedades} da lei \emph{natural} são a \emph{universalidade} e a \emph{immutabilidade}.

a) A lei natural é universal. --- A lei \emph{natural} será \emph{universal}, Se Os seus mais geraes preceitos não poderem ser ignorados, nem 2\emph{nvencivelmente}, pelos homens dotados do uso perfeito da \pg{633}  razão. Ora os mais geraes preceitos da lei \emph{natural} não podem ser ignorados, nem \emph{Znvencivelmente}, pelos homens no uso perfeito da razão; não só porque ou são immediatamente evidentes, ou facilmente se deduzem dos primeiros principios, \pg{634}  mas tambem porque \textsc{Deus} não alcançaria o seu escopo de guiar, pela lei \emph{natural}, os homens para o seu fim, se a mesma lei não fosse por elles conhecida, pelo menos emquanto aos preceitos mais geraes. Logo a lei \emph{natural é universal} \footnote{A \emph{universalidade} da lei \emph{natural} diz respeito á sua \emph{promulgação}. --- Para esta promulgação são necessarias e bastam, da parte do legislador, as seguintes condições: --- 1.º) que os \emph{wuniversalissimos} principios da lei, que constituem os preceitos da \emph{primeira} categoria (como é este: \emph{deve fazer-se O bem}, etc.) e encerram virtualmente toda a lei, não possam de \emph{modo algum} ser \emph{ignorados} por quem tem o uso da razão desenvolvido; --- 2.) que o homem seja dotado de tal faculdade de raciocinio, que possa, com facilidade e sem esforço, deduzir dos principios geraes as conclusões praticas, absolutamente necessarias para a recta disposição da vida.\par Estas duas condições são necessarias e sufficientes; porque não é preciso que todos os homens possam deduzir, com a mesma facilidade, dos primeiros principios de moralidade as mais remotas e complexas conclusões, --- O que seria contrario á natural desegualdade dos homens. Ora é certo que todos os homens, que têm o uso perfeito ou desenvolvido da razão, não podem ignorar os primeiros preceitos da lei moral, como tambem não podem ter uma ignorancia \emph{invencivel} dos preceitos, que immediata e facilmente se deduzem dos primeiros principios (como são os preceitos do \emph{decalogo}); embora muitos homens possam ter uma ignorancia \emph{invencivel} dos preceitos da terceira categoria (como é a malicia do \emph{duello}) porque a deducção d' estas conclusões exige uma consideração mais subtil e depende, em grande parte, das cireumstancias das coisas e das pessoas.\par Os \emph{tradicionalistas} rigorosos combatem a nossa conclusão, dizendo que a razão humana, abandonada ás suas forças, não póde conhecer nenhuma verdade da ordem metaphysica e moral. --- Não insistimos na refutação\pgn{635}d'esta opinião, cuja falsidade foi demonstrada tambem na \emph{Antropología}. --- Os exemplos de povos, que teriam ignorado os preceitos moraes da segunda categoria (que constituem o \emph{decalogo}) nada provam contra a nossa these. Esses preceitos não eram ignorados na sua substancia, mas eram falsamente applieados aos casos particulares, e essa falsa applicação dependia da corrupção dos costumes. E, ainda que se admittisse uma ignorancia total, ella era \emph{vencivel}; e por isso é sempre verdade que o homem não póde ter uma ignorancia \emph{invencivel} dos preceitos da segunda categoria.}.

\pg{635} b) A lei natural é immutavel --- A lei \emph{natural é immutavel} não só \emph{intrinseca} mas tambem \emph{extrinsecamente}. --- E' immutavel \emph{intrinsecamente}, isto é, em relação ao seu \emph{objecto}; porque ella presereve o que é conforme com a natureza \emph{racional} do homem, e prohibe o que lhe é contrario. Ora a natureza do homem é \emph{immutavel}. --- Eº immutavel \emph{ectrinsecamente}, isto é, em relação ao \emph{legislador}, que é \textsc{Deus}; porque a Vontade de \textsc{Deus}, sendo : infinitamente santa e justa, não póde deixar de prescrever sempre o bem e de prohibir sempre o mal \footnote{Para a solução das difficuldades, que se fazem contra a \emph{immutabilidade} da lei \emph{natural}, deve attender-se ás seguintes considerações:\par 1.3) A' lei \emph{natural} podem adorescentar-se muitas coisas pela lei \emph{positiva}, divina ou humana. Porquanto esta addição, deixando intacta a essencia da lei natural não a muda \emph{intrinsecamente}, mas só a determina, explica, etc.\par 2.º) Quando o \emph{objecto} da lei muda, deixa de ser objecto da mesma lei, mas esta fica intacta e não muda. Assim \textsc{Deus}, quando mandou a Abrahão que immolasse seu filho Isaac, não alterou o \emph{quinto} mandamento do decalogo, mas fez com que a immolacüo de Isaac não fosse comprehendida no \emph{objecto} d'esse mandamento, que, prohibindo a morte \emph{injusta} do proximo, não podia prohibir a morte de Isaac, que era \emph{justa}, porque foi preceituada por \textsc{Deus}, Senhor da vida e da morte de todas as creaturas.\par A? \emph{immutabilidade} da lei \emph{natural} são contrarias --- a opinião dos sequazes do \emph{positivismo}, que admittem a Moral \emph{variavel} e \emph{relativa}, --- e a theoria de \emph{Pufendorf}, que colloca na vontade divina o fundamento da \emph{moralidade}, embora confesse que a lei \emph{natural}, uma vez livremente estabelecida, dura immutavel. --- Estas theorias tambem foram refutadas no cap. III d'este tratado.}.

90. \textbf{Primeiro} principio da lei \textbf{natural}. --- O primeiro principio da lei \emph{natural}, de que derivam e a que se reduzem todos os preceitos moraes, é o seguinte: Pratica o bem e evita o mal. --- Na verdade, o primeiro principio da lei \emph{natural} deve ter os seguintes requisitos: --- deve ser \emph{immediatamente} evidente; --- deve ser \emph{universalissimo}, contendo virtualmente todos os preceitos da lei natural; --- deve ser \emph{supremo}, exprimindo a razão ultima e commum de todos os preceitos. Ora o principio apontado satisfaz a estes requisitos; porque --- é \emph{mmediatamente} evidente; --- é \emph{wuniversalissimo}, pois todos os preceitos têm por objecto o que é intrinsecamente bom ou mau; --- é \emph{supremo}, \pg{636}  pois a razão ultima, pela qual nos são impostos os preceitos naturaes, encontra-se na intrinseca bondade ou malicia das coisas preceituadas ou prohibidas \footnote{O bem, que o primeiro principio da lei \emph{natural} manda praticar, não é \emph{todo} o bem, mas é o bem \emph{necessario}, é o bem cuja omissão é positivamente peccaminosa; pois ha muitas coisas boas, que não são \emph{preceituadas}, mas simplesmente \emph{aconselhadas}.\par O primeiro principio da \emph{le}? \emph{natural} é o complemento do outro principio, que se chama \emph{norma} ou \emph{regra} da \emph{moralidade}; porque á \emph{norma} da \emph{moralidade} Serve para extremar e distinguir o bem do mal, e o primeiro principio da lei \emph{natural} manda que se pratique o bem, que foi indicado pela \emph{norma} da moralidade e que se refere \emph{necessariamente} ao nosso fim ultimo. Como se vê, á norma da \emph{moralidade} tem mais extensão do que o primeiro principio da lei \emph{natural}; pois este principio só abrange o bem \emph{preceituado}, e \& \emph{norma} da \emph{moralidade} abrange tudo o que é bem moral, quer \emph{preceituado}, quer não.\par Os philosophos tém apresentado varias e diversas opiniões ácerca do primeiro principio da lei \emph{natural}, conforme a diversidade dos seus systemas ácerca do \emph{fundamento da moralidade}. As principaes d'estas opiniões foram expostas, quando tratámos da \emph{norma} ou \emph{regra} da \emph{moralidade}, com a qual identificaram o primeiro principio da lei \emph{natural}.}.

91. \textbf{Sancção da} lei natural. --- \emph{Sancção} da lei \emph{natural} é a somma dos premios e dos castigos, que \textsc{Deus} estabeleceu para os que observam ou violam a lei \emph{natural}. --- A \emph{sancção é sufficiente} ou \emph{insufficiente}. E'º \emph{sufficiente}, se é proporcionada ao merito ou ao demerito, e se tem força para levar a vontade á pratica do bem e á fuga do mal. E' \emph{insuficiente}, no caso contrario \footnote{A sancção é \emph{activa} ou \emph{passiva}. A sanccüo \emph{activa} é o decreto pelo qual o legislador estabelece premios para os que observarem a lei, e castigos para os que a violarem. A sancção \emph{passiva} é constituida por esses mesmos premios ou castigos. --- Aqui tratamos da sancção \emph{passiva}, que suppõe necessariamente a sancção \emph{activa}.\par Alguns auctores pretendem que a \emph{sancção}, para ser \emph{justa}, deve \emph{garantir} sempre a observancia da lei. --- Esta pretenção não é rasoavel; porque não se harmoniza nem com a liberdade \emph{p/iysiea} dos subditos, que a lei deixa sempre intacta, nem com os factos, que mostram que a \emph{sancção} não garante sempre a observancia da lei. Para que, pois, \emph{a sancção} seja justa, não é preciso que ella garanta sempre a observancia da lei, mas basta que seja por si sufficiente para garantir essa observancia, de modo que o motivo, que ella apresenta em favor da observancia da lei, seja por si superior aos motivos, que podem levar á sua violação.\par \emph{A}}.

92. \emph{A} lei \textbf{natural deve} ter uma \textbf{sancção}. --- \textsc{Deus} é um legislador infinitamente \emph{sabio}, \emph{justo} e \emph{santo}. Ora \textsc{Deus} não teria estes attributos, se não tivesse estabelecido uma sancção para a lei \emph{natural}. Não seria \emph{sabio}; porque um legislador \emph{sabio} deve providenciar pelo cumprimento da lei, e não ha providencia \pg{637}  mais efficaz que a \emph{sancgão}; --- não seria \emph{justo}; pois a \emph{justiça} exige que ao merito se dó o premio e ao demerito o castigo; --- não seria. \emph{santo}; pois repugna á \emph{santidade} que o destino do homem bom seja egual ao destino do malvado. Logo a lei \emph{natural} deve ter uma \emph{sancção} \footnote{Nada mais razoavel do que a \emph{sancção} da lei \emph{natural}. \textsc{DEus} promulgou a lei \emph{natural}, para que o homem, seguindo-a, aleancasse o bem. Se, pois, o homem seguir a lei, encontrará o bem, e este terá razão de \emph{premio}; e, se a não seguir, encontrará o mal, e este será \emph{castigo} (\emph{C}. \emph{Gent}., 1. TIT, c. 140).\par Todavia alguns negam a \emph{sancção} da lei \emph{natural}; porque, dizem, os nossos peceados não offendem a \textsc{Deus}, e porque repugna que a Bondade de \textsc{DEus} se vingue e se deleite dos tormentos dos peccadores. --- Mas estas razões nada provam. --- Não prova a primeira; porque o peccador, embora não tire a \textsc{Deus} nenhuma perfeição e não lhe possa causar um mal intrinseco, comtudo, \emph{quanto depende de} si, despreza a \textsc{Deus}, tira-lhe a gloria \emph{ewtrinseca}, e por isso deve ser punido. --- Não prova a segunda; porque não é o desejo da vinganca ou o sentimento da crueldade, mas é o proprio amor da Justiça, da santidade e da bondade, que leva \textsc{Deus} a punir os peceadores, porque o castigo não só conserva e promove, mas tambem completa e restabelece, como dissemos, a ordem moral, violada pelo peccado. Por isso ó a propria Bondade de \textsc{Deus} que exige a sancção da lei \emph{natural}.}.

93. A lei natural tem \textbf{nesta} vida uma \textbf{sancção insufficiente}.

a) Tem nesta vida uma sancção. --- Na verdade, a obediencia á lei natural recebe a sua recompensa na paz da consciencia, na saude do corpo, na estima dos homens, nos bens sociaes, ete. --- como, pelo contrario, a sua lrangressão recebe o seu castigo no remorso, nas doencas, no desprezo dos homens, nas penas, que lhe inflige a sociedade, ete. \footnote{As principaes sancções da lei \emph{natural} na \emph{vida presente} são quatro, a saber: a sanccüo da \emph{consciência}, a sancção \emph{natural}, a sanccüo \emph{social}, a Sancção da \emph{opinião publica}. --- A sanccüo da \emph{consciencia} comprehende as alegrias e as afflicções da consciencia. Depois de um acto moralmente bom, o homem sente ordinariamente uma ineffavel satisfação e alegria da alma, superior a todas as satisfações dos sentidos e alegrias do mundo. Depois de um acto moralmente mau, o homem prova, pelo contrario, uma tristeza humilhante, um sentimento vivissimo da sua queda, --- tristeza e sentimento, que constituem o \emph{remorso}. --- A sancção \emph{natural} consiste na saude, que é uma consequencia ordinaria da virtude, \emph{e} nas doenças, que em regra derivam de uma vida desordenada. --- A sancção \emph{social} abrange as recompensas e os castigos, que a sociedade estabelece para os que praticam o bem ou o mal. --- A sanecüo da \emph{opinião publica} consiste na estima ou desestima, que os homens geralmente tributam \& virtude ou ao vicio.}.

b) \textbf{Esta} sancgão é insufficiente, --- Na verdade, a experiencia

"\emph{n} attesta que as recompensas não são sempre proporcionadas aos meritos e os castigos não são sempre proporcionados aos demeritos, e que essas recompensas e esses castigos não têm \pg{638}  força sufficiente para conter no seu dever a nossa vontade, assaltada por tantas e tão perigosas tentações \footnote{A sancção da lei \emph{natural} na vida presente é insufficiente. Na verdade, a sancção, para ser justa, deve ser principalmente \emph{universal} (isto é, não deve deixar nenhum acto virtuoso sem premio, e nenhum acto vicioso sem castigo) e \emph{proporcionada} (isto é, o premio e o castigo devem ser proporcionados ao merito ou ao demerito do agente) Ora a sancção da lei \emph{natural}, na vida presente, não é nem \emph{universal} nem \emph{proporcionada}.\par \emph{a}) Não é \emph{universal}. Porquanto ha muitos actos bons, que, sobre a terra, não recebem premio, e muitos actos maus que não recebem castigo. Na verdade, a sancção \emph{natural} não se verifica sempre, porque ha pessoas boas muito doentes, e pessoas más cheias de saude; --- a sancção \emph{social} não se applica a todos os casos, pois á sociedade não se importa com as virtudes nem com os vicios particulares, não castiga todos os reus, e muitas vezes condemna os innocentes; --- a sancção da \emph{opinião publica} não se extende a todas as virtudes nem a todos os vicios, é indulgente para com o vicio elegante, approva muitos erros moraes, condemna a pobreza como vicio, etc.\par \emph{b}) Não é \emph{proporcionada}. Com effeito, a sanccüo da \emph{consciencia}, embora seja \emph{universal}, não corresponde sempre ao grau da bondade ou da malicia; assim uma consciencia virtuosa e delicada soffre ás vezes, por uma leve imperfeição, mais do que uma consciencia, habituada no vicio, pelos crimes mais graves e hediondos; --- como tambem não são proporcionadas a sancção \emph{natural}, a sancção \emph{social} e a sancção da \emph{opinião publica}; porque, como vimos, nem a \emph{natureza}, nem a \emph{sociedade}, nem a \emph{opinião publica}, distribuem os premios e os castigos em conformidade com os meritos e com os demeritos.}.

94. A lei \textbf{natural deve} ter uma sancção sufficiente na vida futura. --- A lei \emph{natural} deve ter uma sancção \emph{sufficiente}; porque \textsc{Deus}, exigindo a fiel observancia da sua lei, deve ter garantido esta observancia por um meio conveniente e efficaz, e este meio só póde consistir numa sancção \emph{su/ficiente}. Ora esta sancção \emph{sufficiente} não existe, como provámos, na vida presente. Logo deve existir na vida futura. Logo a lei \emph{natural} deve ter uma sancção sufliciente na vida futura \footnote{Esta these é uma consequencia logica das proposições antecedentes. \textsc{Deus}, promulgando a lei \emph{natural}, deve ter estabelecido uma sancção \emph{sufficiente} ou \emph{perfeita}. Esta sanecüo \emph{sufficiente} ou \emph{perfeita} não se encontra na vida presente. Logo deve encontrar-se na vida futura. --- E' esta uma verdade admittida por todos os que têm feito bom uso da razão; e os proprios incredulos e racionalistas, como \emph{Rousseau}, \emph{Cousin}, deduziram da sancção \emph{insufficiente} ou \emph{imperfeita}, que a lei \emph{natural} tem nesta vida a existencia de uma vida futura e por isso a immortalidade da alma humana. --- D'onde\pgn{639}se vê que a immortalidade da alma e a existencia de \textsc{DEus} são verdades não só metaphysica, mas tambem moralmente necessarias.\par Os \emph{positivistas} dizem que a lei \emph{natural} tem uma sanccüo sufficiente sobre a terra, e que esta sancção consiste nos prazeres e nas vantagens, que derivam de uma conducta sabia e virtuosa, e nas satisfações intimas e puras, que acompanham a pratica do bem.\par A resposta a esta diffieuldade está dada. Os prazeres e as vantagens não estão sempre em proporção com o merito, e muitas vezes as afflicções e as perdas augmentam com o augmento da virtude. A satisfação da alma (que, na eschola \emph{positivista}, é muito superficial e illusoria, pois não passa de um sentimento \emph{Aereditario}) não é recompensa sufficiente para sacrificios heroicos; assim como o remorso não é castigo adequado para os criminosos, pois o remorso diminue com o augmento do vicio. --- E' preciso, pois, procurar a sancção sufficiente da lei \emph{natural} na vida futura, onde a felicidade se una necessaria e adequadamente com a virtude, e a infelicidade com o vicio. Os \emph{positivistas} ou admittem esta conclusão, ou destroem o proprio fundamento de toda a ordem moral.}.

95. A \textbf{sancção da} lei natural, \textbf{na vida} futura, \textbf{consiste na posse eterna} ou na \textbf{perda eterna} da felicidade.

a) Consiste na posse ou na perda da felicidade. --- \textsc{Deus} criou os homens para a felicidade, e quer que a consigam pela observancia da lei natural. Ora á observancia da lei é uma \pg{639}  condição indispensavel para a consecução da felicidade; pois \textsc{Deus} não póde premiar, indistinetamente, os bons e os maus. Logo os homens, observando a lei de \textsc{Deus}, conseguem a felicidade, e, transgredindo a lei, perdem-na. Logo a sancção da lei \emph{natural}, na vida futura, consiste na posse ou na perda da felicidade \footnote{O fim \emph{ultimo} do homem é a felicidade perfeita na vida futura. A vida presente é um caminho. As acções moralmente boas são outros tantos passos, por que nos approximamos da felicidade; como, pelo contrario, as acções moralmente más são outros tantos passos, por que nos afastamos da mesma felicidade; pois, consistindo a felicidade perfeita na posse do Summo Bem, repugna que o Summo Bem se consiga por meio do mal moral, visto que os meios devem ser proporcionados ao fim. Por isso, acabando o caminho com a morte, os que percorreram o caminho que leva para a felicidade, devem necessariamente conseguil-a, \emph{e} os que d'ella se afastaram, não a conseguem. Ora a posse e a perda da felicidade, que são consequencias das obras, têm razão de \emph{premio} ou de \emph{castigo}. Logo a sancção da lei \emph{natural}, na vida futura, consiste na posse ou na perda da felicidade.\par Todavia isto não basta para a justa e completa sancção da lei \emph{natural}.\par Como os que conseguem o seu ultimo fim, não são eguaes nos meritos e nas virtudes, mas um é mais perfeito do que outro, é necessario dizer que, embora todos os virtuosos consigam o seu fim, comtudo nem todos recebem egual premio, mas cada um é elevado ao grau de felicidade, que está em harmonia com as suas obras boas. --- O mesmo criterio deve applicar-se aos impios que morrem no peccado. Estes infelizes não são todos egualmente culpados, mas os crimes de um excedem, em numero e gravidade, os crimes de outro, e por isso, devendo a sancção ser \emph{proporcionada}, é justo que o castigo de um seja mais intenso que o do outro. Ora, se o castigo dos impios se limitasse apenas á perda do fim \emph{ultimo} ou da felicidade, a sancção não seria \emph{proporcionada}. Logo é necessario dizer que os impios, além da perda do fim \emph{ultimo}, estarão sujeitos a outros castigos positivos, diversos conforme\pgn{640}a diversidade das suas culpas. E nada mais natural. O impio, peccando, colloca o seu fim \emph{ultimo} não no Creador, mas nas creaturas, e assim despreza a \textsc{Deus} pelo amor desordenado das creaturas. Ora, devendo o castigo ser proporcionado \& culpa, é justo que o peccador seja privado da posse de \textsc{Deus} e seja punido pelas creaturas.\par Um professor \emph{krausista}, da Hespanha, \emph{Gonzáles Serrano}, diz que a \emph{sancção} da lei, como tambem a sua observancia, consiste na \emph{infinita continuidade} da nossa \emph{vida}. --- Esta opinião é falsa, não só porque se basea no \emph{pantheismo}, mas tambem porque a observancia da lei pertence á vida presente, ea sancção da lei não consiste na \emph{infinita continuídade} da vida, mas sim nos premios ou nos castigos, que se hão de dar em conformidade com as obras.}.

\pg{640} b) A posse ou a perda da felicidade é eterna, --- Deve ser \emph{eterna} a \emph{posse} da felicidade; pois a felicidade, que é o fim ultimo do homem; é a \emph{perfeita}, e a felicidade não é perfeita, se não é \emph{eterna}, --- Deve ser \emph{eterna} a Sua \emph{perda}; porque, não podendo o homem merecer depois da morte, a perda do fim \emph{ultimo} é definitiva e eternamente irremediavel \footnote{Os incredulos têm-se sempre levantado furiosamente contra a \emph{eternidade das penas}. --- Citemos alguns trechos dos mais famigerados. --- \emph{Peszani} diz que « a theologia, ensinando o dogma da eternidade do infermo, tem commettido um, crime de lesa humanidade » (\emph{Dieu}, \emph{V homme}). --- \emph{Larroque} julga que a \emph{eternidade das penas} é « uma horrivel blasphemia, pela qual se desvirtua a justiça de \textsc{Deus}, ao mesmo tempo que se desconhece a sua santidade e bondade » (\emph{Examen}. \emph{critique de la religion chrétienne}, t. T). --- \emph{Maury} diz: « A falta de reconciliação de todas as creaturas com \textsc{Deus}, e a eternidade das penas inuteis para os culpados, --- eternidade que transforma a justiça e a expiação numa pura vingança, é uma idea barbara que o christianismo concebeu segundo as falsas noções do judaismo ácerca da justiça divina » (\emph{Encycl}. \emph{Didot}, art. \emph{Enfer}) --- \emph{Strauss} afirma que o dogma da \emph{eternidade das penas} « corrompe o sentimento e repugna á bondade e santidade de \textsc{DEus} ». --- Do mesmo modo se exprimem \emph{Reynaud} (\emph{Terre et Ciel}), \emph{Leroux} (\emph{De} /\emph{Umanité}), \emph{Callet} (\emph{L}? \emph{Enfer}), \emph{Franch} (\emph{Dict}. \emph{philos}., avt. \emph{Penalité}), \emph{Ahrens} (\emph{Cours de droit}).\par Como se vê, os incredulos attribuem ao Christianismo a invenção do dogma da \emph{eternidade das penas}, e dizem que este dogma é contrario á razão e aos attributos de \textsc{Deus}. Por isso, primeiramente demonstraremos que o dogma da \emph{eternidade das penas} foi sempre admittido por todos os povos e não repugna á razão, e depois resolveremos as objecções, sobre tudo as que se basêam nos attributos de \textsc{Deus}.\par A) \emph{A} eternidade das penas foi sempre admittida por todos os povos. --- Os historiadores, os poetas, os philosophos attestam que não houve povo algum, antigo ou moderno, que não admittisse o dogma da eternidade das penas. Em relação aos \emph{chaldeus}, \emph{persas}, \emph{assyrios}, \emph{egypcios}, temos a auctoridade de \emph{Tacito} e de \emph{Deodoro Siculo}. Em relação aos \emph{indios}, \emph{celtas}, \emph{gawlezes}, e outros povos antigos, temos o testemunho de \emph{Stacio}, de \emph{Strabdo}, de \emph{Valerio Maximo}. --- \emph{Homero} descreve os suplicios de \emph{Tantalo}, de \emph{Theseu}, etc. (\emph{Odyssea}, liv. X-XI) --- \emph{Virgilio} traduz a crença religiosa do seu tempo nestas tremendas expressões: « \emph{sedet}, \emph{aeternumque sedebiti nfelic Theseus} » (\emph{Eneid}., VI, 618),\par -\pgn{641}'« \emph{rostroque immanis vultur obunco immortale jecur tondens}... \emph{nec fibris requies datur ulla venatis} » (VI, 591). --- O proprio \emph{Lucrecio}, embora materialista, não deixa de attestar a convicção dos povos na eternidade das penas nestas palavras: « \emph{Aeternas quoniam poenas in morte timendum est} » (\emph{De nat}. \emph{rerum}). --- \emph{Ovidio} falla d'um caminho sombrio, silencioso, que leva ao Inferno (\emph{Metamorph}., liv. IV). --- \emph{Platão} é o mais explicito: « Os que morrem culpados de grandes crimes, de crimes irremediaveis, como o sacrilegio e o assassinio, julgados pelo destino, são encerrados no Tartaro, d'onde \emph{nunca} hão de sahir » (\emph{Phedon}, 113). --- Citados estes testemunhos, raciocinamos assim: Esta crenea, firme, universal e constante num dogma, que não só não lisongeia os interesses e as paixões, mas é contrario a todos os interesses e paixões, só pôde derivar --- ou de uma revelação primitiva, feita ao homem por \textsc{Deus}, --- ou de uma deducção facilima da nossa razão, a qual, sendo especificaments» egual em todos os homens, descobre a mesma verdade com a mesma certeza. Num e noutro caso, a crença na eternidade das penas é uma verdade; porque \textsc{Deus} não se engana e não engana, e porque uma affirmação certa, universal, constante, contraria ás paixões, só póde derivar da evidencia \emph{objectiva}, que é criterio de verdade.\par B) O dogma da eternidade das penas é conforme aos dictames da \textbf{razão}. --- O dogma da \emph{eternidade} das penas será conforme aos dictames da razão, se as \emph{penas eternas} forem \emph{justas} e \emph{necessarias}. Ora são \emph{justas} e \emph{necessarias}.\par a) As \emph{penas eternas} são \emph{justas}. Uma pena, para que seja \emph{justa}, deve ser proporcionada ao crime. Ora a pena \emph{eterna} é proporcionada ao peccado. Na verdade, uma offensa é mais ou menos grave, conforme a distancia que existe entre o offensor e a pessoa offendida. Ora entre o peccador e \textsc{Deus} \textbf{existe} uma distancia infinita; porque \textsc{Deus} é a magestade infinita e o peccador é um ente finito, e a distancia entre o infinito e o finito é infinita. Logo o peceado contém uma gravidade ou malicia infinita. Logo as penas, que o peccador merece, devem ser infinitas. Ora as penas, que se infligem ao peccador, não podem ser infinitas na \emph{intensidade}; porque nenhuma coisa creada, como é a pena, póde ser \emph{intensivamente} infinita, Logo devem ser infinitas na \emph{duração}, isto é, devem ser \emph{eternas}. --- Os adversarios poderiam responder que a privação da posse de \textsc{Deus}, embora \emph{temporal}, é uma pena infinita; porque \textsc{Deus} é um bem infinito. Mas esta resposta seria inutil. Por quanto, a privação \emph{temporal} da posse de \textsc{Deus}, se é uma pena infinita em relação ao \emph{objecto}, que se perde e que é \textsc{Deus}, não é uma pena infinita em relação ao \emph{sujeito}, \emph{a} que se inflige; porque a creatura, não comprehendendo o Bem infinito, não póde conceber uma pena infinita, se fôr privada da posse d' Elle.\par b) As \emph{penas eternas} são \emph{necessarias}. Uma \emph{pena} será \emph{necessaria}, se for a unica sancção \emph{sufficiente} da lei. Ora só uma \emph{pena eterna} póde ser a sancção \emph{sufficiente} da lei natural. Na verdade, se considerarmos a natureza do homem inclinada para o mal, a vehemencia das paixões, os attractivos dos bens sensiveis, as difficuldades da virtude, veremos que penas \emph{passageiras} não são motivo sufficiente para conter, sempre e em todas as circumstancias, os homens na observancia da lei. --- Esta verdade encontra a confirmação na propria conducta dos impios, que, para se abandonarem mais livremente ás desordens, procuram convencer-se de que as penas da vida futura hão de\par 41 Vor. II.\pgn{642}acabar e de que as almas dos impios serão um dia admittidas á bemaventuranca. \emph{Lucrecio} exprime esta idea nos seguintes versos (\emph{De nat}. \emph{ter}., I. 108).\par \&... Nam si certum finem esse viderent\par Aerumnarum homines, aliqua ratione valerent\par Religionibus atque minis obsistere vatum.\par Nune ratio nulla est restandi, nulla facultas,\par Aeternas quoniam \emph{poenas} in morte timendum . est ».\par C) Examinemos as principaes objecções dos incredulos.\par a) Dizem: Um juiz justo não dá um castigo não proporcionado á culpa. Ora a pena \emph{eterna} não é proporcionada ao peccado, que é uma coisa \emph{temporal}, e por vezes \emph{momentanea}. Logo Dreus, sendo um juiz infinitamente justo, não dà, nem póde dar uma \emph{pena eterna}.\par Resposta. A proposição \emph{menor} do syllogismo é falsa. Com effeito, a gravidade da culpa mede-se --- não pela sua \emph{duração}, aliás os maiores crimes deveriam ter um castigo de poucos momentos, porque, em geral, commettem-se no espaço de poucos momentos, --- mas pela \emph{distancia}, que existe entre a pessoa offendida e o offensor, assim um insulto feito ao rei tem mais gravidade que um insulto feito a um egual. Por isso, tendo o peccado uma malicia infinita, deve ser punido com uma pena infinita.\par b) Insistem: O fim do \emph{castigo} é a emenda do culpado. Ora a pena \emph{eterna} não emenda o culpado. Logo a pena \emph{eterna} não existe.\par Resposta. O \emph{castigo} tem dois fins: um \emph{principal} e outro \emph{secundario}. O \emph{principal} consiste em reparar a ordem moral, violada pela culpa: d'ahi o castigo \emph{vingativo} ou \emph{ewpiatorio}. O \emph{secundario} consiste em afastar o homem da culpa; d'ahi o castigo \emph{medicinal} ou \emph{exemplar}, conforme tem por fim a emenda do culpado, ou a obediencia á lei da parte dos outros homens. Ora a \emph{pena eterna} alcança sempre o seu fim \emph{principal}. Porquanto, se o peccado é uma proclamação \emph{pratica} de que a creatura póde desobedecer ao Creador, a pena \emph{eterna} é uma affirmação \emph{pratica} de que a creatura deve estar sujeita ao Creador. Pelo \emph{peccado}, á creatura levanta-se; pelo castigo, abate-se. --- A pena \emph{eterna} --- não só alcança sempre o fim \emph{principal}, --- mas alcança tambem o fim \emph{secundario}, tanto em relação aos outros, que afasta do peccado, quanto em relação ao proprio peccador, que o medo d'esse castigo leva a detestar as culpas e a reformar a vida.\par o) Continuam: A pena \emph{eterna}, se póde levar ao arrependimento o homem, emquanto este é vivo, não o leva, depois de morto. Logo a pena \emph{eterna} não alcança o seu fim, que é a emenda do culpado.\par Resposta. O \emph{castigo} deve alcançar sempre o seu fim \emph{principal}, aliás não teria razão de ser, e tambem deve alcançar sempre o seu fim \emph{secundario}, quando o \emph{secundario} é compativel com o \emph{principal}. Mas, quando o fim \emph{secundario} é incompativel com o \emph{principal}, neste caso basta que o castigo alcance o seu fim \emph{principal}. Ora a pena \emph{eterna}, que se inflige ao impio, morto no seu peoccado, alcança o seu fim \emph{principal}, que é a reparação da ordem moral, embora não alcance um dos seus fins \emph{secundarios}, que é a emenda do peccador, porque este é incapaz de emenda.\par 4d) Continuam: O dogma da \emph{eternidade das penas} suppõe que o tempo\pgn{643}de merecer é exclusivamente a vida presente. Ora esta supposicüo é falsa; porque o homem pode, em qualquer tempo da sua existencia, praticar o bem, e assim adquirir meritos e alcançar o fim ultimo.\par Resposta. Admittimos o \emph{antecedente}, mas negamos o \emph{consequente}. Na verdade, a Providencia de \textsc{Deus}, sendo infinitamente sabia, devia prescrever um tempo determinado, em que os homens podessem, com os seus actos livres, merecer o fim ultimo; aliás os homens teriam, com toda a facilidade e sem receio, commettido o peccado, seguros de que podiam, sempre que quizessem, tornar á amizade de \textsc{Deus}. Ora a ordem das coisas exige que este tempo determinado seja. só a vida présente; não só porque só na vida presente existe propriamente o /\emph{iomem}, isto é, o composto de alma e de corpo, --- mas tambem porque, aliás, os homens desprezariam impunemente a lei de \textsc{Deus} aqui sobre a terra, reservando a sua conversão para a vida futura. Por isso a propria natureza da \emph{sancção} da lei natural exige que o tempo util para merecer seja limitado á vida presente.\par e) Dizem: \emph{O dogma da eternidade} das penas é contrario á bondade de \textsc{Deus}. Logo as penas não podem ser \emph{eternas}.\par Resposta. Não póde ser contraria a um attributo divino uma coisa, que é exigida por outros attributos; pois, em \textsc{Deus}, os attributos não são \emph{realmente} distinctos entre si. Ora a \emph{eternidade das penas} é exigida pela \emph{Justiça} e \emph{Sabedoria} de \textsc{Deus}; pois demonstrámos que o castigo eterno é \emph{justo} e \emph{necessario} para a conservação da ordem moral. Logo o dogma da \emph{eternidade das penas} não é contrario á \emph{Bondade} de \textsc{Deus}. --- E não só não é contrario á \emph{Bondade} divina, mas é exigido por ella. Com effeito, á \emph{Bondade divina} consiste em propor meios efficazes para a destruição do peccado; pois o peccado é diametralmente opposto á Bondade divina. Ora \textsc{Deus} não teria proposto meios efficazes para a destruição do peccado, se, um dia, justos e peccadores conseguissem a felicidade eterna, que uns, os justos, alcançaram com tanto trabalho, e outros, os peccadores, rejeitaram com tanto desprezo no tempo da graça e da misericordia.\par f) Replicam: A \emph{eternidade das penas}, longe de tutelar e promover a virtude e a honestidade, é a destruição de toda a virtude e de toda a honestidade; porque uma pessoa, que ama a virtude e a honestidade por causa da recompensa, deixa por isso de ser virtuosa e honesta.\par Resposta. A difficuldade encerra a theoria de \emph{Kant}, seguida por \emph{Krause}, segundo a qual é necessario praticar o bem pelo bem, sem nenhum desejo da recompensa; aliás o acto é vicioso. --- Mas esta theoria é falsa. Porquanto, o homem, que segue a virtude e foge do vicia em attenção á posse ou á perda eterna do seu fim \emph{ultimo}, não só não se afasta de \textsc{Deus}, mas tende directamente para Elle; porque o nosso fim \emph{ultimo} não é o homem (como disse \emph{Kant}), mas é o proprio \textsc{Deus}. Por isso o homem, que ama a \textsc{Deus}, não póde deixar de desejar ardententemente a sua posse. D'onde se vê que o amor que temos a Deus, emquanto Elle é o nosso ultimo fim, póde chamar-se \emph{gratuito}; porquanto quem não deseja outra recompensa senão o bem, que elle ama, ama \emph{gratuitamente}; pois nada quer fóra do bem desejado. Se a pratica do bem pelo desejo da recompensa nasce do amor de \textsc{Deus}, é claro que essa pratica não póde deixar de ser virtuosa e honesta.}.

96. \textbf{Lei positiva}. --- Lei \emph{positiva} é um preceito racional para o bem commum, promulgado \emph{livremente} por quem governa a \pg{641} \pg{642}  communidade e addicionado á lei natural. --- “A lei \emph{positiva} póde ser \emph{divina} ou \emph{humana}, conforme é promulgada por \textsc{Deus}, \pg{643}  ou pelo homem, embora em virtude da auctoridade recebida de D\&us. --- A lei \emph{positivo-humana} subdivide-se em \emph{ecclesiastica} \pg{644}  ou \emph{civil} conforme deriva da auctoridade ecclesiastica ou da auctoridade civil \footnote{Em toda a lei devemos distinguir dois elementos: o \emph{material} e o \emph{formal}. O elemento \emph{material} é constituido pelas acções ou omissões, que a lei prescreve; o \emph{formal} é a propria \emph{obrigação}. --- A lei \emph{positiva} differe da \emph{natural} em relação a esse duplice elemento. Differe em relação ao elemento \emph{formal}; porque a \emph{obrigação}, imposta pela lei \emph{natural}, deriva da vontade divina \emph{necessariamente} (embora essa necessidade seja \emph{Aypothetica}, porque suppõe a criação), ao passo que a \emph{obrigação}, imposta pela lei \emph{positiva}, deriva da vontade \emph{livre} do superior. --- Differe em relação ao elemento \emph{material}; porque a lei \emph{natural} só prohibe o que é intrinsecamente mau, e só prescreve o bem \emph{indispensavel} ou \emph{necessario}, ao passo que a lei \emph{positiva} póde tambem prescrever coisas, que são indifferentes, ou que, embora honestas, não são comtudo \emph{necessarias}. --- D'estas differenças fundamentaes resultam outras, e são as seguintes: --- 1.9) \emph{A} lei \emph{natural} promulga-se pela luz de razão; a \emph{positiva}, por uma expressão, oral ou escripta, do superior. --- 2.º) A lei \emph{natural é universal}, e extende-se a todos os homens; a \emph{positiva é particular} em relação ao tempo, aos logares, e aos paizes. --- 8.º) A lei \emph{natural} é essencialmente divina, porque deriva de \textsc{Deus}; a \emph{positiva} póde tambem derivar \emph{immediatamente} do homem, em virtude da auctoridade que \textsc{Deus} lhe communicou.}.

97. A lei \textbf{positivo-divina} é \textbf{necessaria}. --- O homem, para que possa aleangar o ultimo fim, deve conhecer, com toda a certeza, o que deve fazer e o que deve evitar nas varias cireumstancias da sua vida. Ora a lei \emph{natural}, embora seja \emph{por si} sufficiente para dirigir os nossos passos para o ultimo fim, não é sufficiente no estado actual da humanidade; porque a luz de razão, a qual nos manifesta a lei natural, é por tal modo eseurecida pelos prejuizos e pelas paixões, que não só se acha, quasi sempre, perplexa e incerta na applicação dos principios d'essa lei aos casos particulares e contingentes, mas chega, muitas vezes, a perverter ou desvirtuar os mesmos principios. Logo é necessaria, além da lei \emph{natural}, uma outra lei \emph{divina}, promulgada por um modo \emph{positivo}, a qual possa dirigir seguramente os homens para o ultimo fim \footnote{E' claro que a necessidade, de que fallamos, é \emph{moral}, e não \emph{absoluta}. A razão humana é \emph{por si} sufficiente para ordenar a nossa vida. Mas essa razão não devemos consideral-a \emph{em abstracto}, mas \emph{em concreto}, isto é, como se encontra actualmente nos homens. Considerada \emph{em concreto}, a razão precisa de um auxilio, de um conforto superior; e esse auxilio, longe de deprimir á nossa razão, ennobrece-a, porque a preserva do erro. --- E claro tambem que essa necessidade \emph{moral} se refere á ordem \emph{natural}; porque, na hypothese da elevação do homem á ordem \emph{sobrenatural}, á necessidade de uma lei \emph{positivo-divina é absoluta}, porque á nossa razão, abandonada a si mesma, não pode, de modo algum, conhecer o que pertence a essa ordem.\par Nem se diga que a lei \emph{positivo-divina} pode ser substituida pela lei \emph{positivo-humana}, a qual pode applicar os principios da lei \emph{natural} aos casos\pgn{645}particulares da nossa vida. Por quanto, o homem não pode alcançar o seu ultimo fim, se não tiver uma perfeita rectidão, não só nos actos \emph{externos}, mas tambem, e sobretudo, nos actos \emph{internos}, e por isso precisa de uma lei, que julgue e regule \emph{todos} esses actos. Ora, se a lei \emph{positivo-humana} julga e regula os actos \emph{externos} do homem, não pode julgar e regular os seus actos \emph{internos}, que não apparecem. Logo é necessaria uma lei \emph{positiva}, que julgue e regule \emph{todos} os actos do homem --- \emph{externos} e \emph{internos}; e esta lei só pode ser promulgada por \textsc{Deus}. (\emph{Sum}. \emph{Th}., 1-2, q. 91, a. 8).}.

\pg{645} 98. A lei \textbf{positivo-humana} é \textbf{necessaria}. --- A lei \emph{positivo-humana} será necessaria, se a lei \emph{natural} não fôr por si mesma sufficiente para todos os casos. Ora a lei \emph{natural} não é sufficiente para todos os casos, --- não só porque ella é evidente apenas nos mais geraes principios, e por isso é necessario que a legitima auctoridade deduza as conclusões e as proponha como regras de conducta, --- mas tambem porque não determina muitas coisas, que São necessarias para a communidade e que não estão contidas nos primeiros principios da honestidade. \textbf{Logo} a lei \emph{positivo-humana} é necessaria \footnote{A lei \emph{positivo-humana}, quer seja uma \emph{conclusão} quer seja uma \emph{determinação} da lei \emph{natural}, é necessaria. --- E' necessaria emquanto é \emph{conclusão} da lei natural. Na verdade, se a lei \emph{natural} é sempre evidente em relação aos mais geraes principios, não o é sempre em relação ás conclusões. Para evitar, pois, a diversidade das opiniões e a desharmonia da actividade social, é necessario que o superior deduza dos principios geraes as conclusões, que não são evidentes e que nos mesmos principios estão \emph{provavelmente} contidas, e as applique aos subditos, como regras dos seus actos. --- E' necessaria emquanto é \emph{determinação} da lei \emph{natural}. Com effeito, a lei \emph{natural}, se dicta que o crime merece castigo, não determina nem a qualidade do castigo, nem a forma do julgamento, nem outras circumstancias, que se verificam nos casos concretos e que não se deduzem da propria lei \emph{natural}; como tambem, se esta dicta que todos os membros da sociedade devem cooperar para o bem commum, não determina o modo, o tempo, o grau d'essa cooperação.}.

99. A lei positivo-humana tem o \textbf{fundamento na} lei \textbf{natural}.

a) A lei \emph{positivo-humana} terá o fundamento na lei \emph{natural}, se aquella derivar d'esta. Ora a lei \emph{positivo-humana} deriva da \emph{lei natural}; pois a lei \emph{positivo-humana} é apenas uma \emph{conclusão} ou uma \emph{determinação} da lei \emph{natural}. Logo a lei \emph{positivo-humana} tem o fundamento na lei \emph{natural} \footnote{A lei \emph{positivo-humana} tem o fundamento na lei \emph{natural}, porque aquola é uma \emph{conclusão} ou uma \emph{determinação} d'esta. --- Todavia, a lei \emph{positivo-lumana} não é uma conclus\&o \emph{necessaria} da lei \emph{natural}; mas apenas \emph{provavel}. --- Não é uma conclusão \emph{necessaria}. Por quanto, uma conclusão, que deriva \emph{necessariamente} dos principios da lei \emph{natural}, pertence á propria lei \emph{natural}, e é imposta por esta e não pela lei \emph{positivo-humana}. Assim, quando dizemos: é \emph{necessario evitar} o \emph{peccado}; ora a \emph{blasphemia é peccado}; \emph{logo é necessario evitar} a \emph{blasphemia}, esta conclusão, por ser \emph{necessaria}, pertence\pgn{646} á lei natural. Se, portanto, o legislador humano promulgar uma lei, que está contida \emph{necessariamente} na lei \emph{natural}, essa lei \emph{positiva} receberá a sua força especialmente da propria lei \emph{natural}. --- Mas é uma conclusão apenas \emph{provavel}. Porquanto, ha muitas conclusões praticas que não estão ou se não vêm \emph{necessariamente} contidas nos principios geraes da honestidade, e por isso são apenas \emph{provaveis}. Sendo apenas \emph{provaveis}, não obrigam por si mesmas, e, sendo necessarias ou uteis para a communidade, é preciso que a legitima auctoridade as imponha, e toda a força, de que ellas são dotadas, depende da vontade do legislador. --- As leis \emph{positivo-humanas}, que são \emph{determinações} da lei \emph{natural}, embora tenham nesta o fundamento, comtudo, por não derivarem \emph{necessariamente} dos primeiros principios da propria lei \emph{natural}, não recebem a força da lei \emph{natural}, mas exclusivamente da legitima auctoridade humana.}.

\pg{646} b) A lei \emph{positivo-humana} terá o seu fundamento na lei \emph{natural}, se o principio, que nos leva a obedecer á lei \emph{positivohumana}, pertencer á lei \emph{natural}. Ora o principio, que nos leva a obedecer á lei \emph{positivo-humana} e que é o seguinte: « Deve \textbf{obedecer-se} ás \textbf{leis} justas dos superiores », pertence á lei \emph{natural}. Logo a, lei \emph{positivo-humana} tem o fundamento na lei \emph{natural} \footnote{A \emph{obrigação}, pois, que nos é imposta pela lei \emph{positivo-humana}, se \emph{proxima} e \emph{immediatamente} deriva da vontade do homem, \emph{remota} e \emph{mediatamente} deriva da vontade de \textsc{Deus}, que, sendo o Auctor da lei \emph{natural}, manda a obediencia aos superiores. Por isso quem transgride as ordens dos superiores, transgride as ordens de \textsc{Deus}.}.

100. \textbf{Requisitos da} lei positivo-humana. --- A lei \emph{posztvo-humana} deve ser \emph{honesta}, \emph{justa} e \emph{possivel}. --- Deve ser \emph{honesta}, de modo que não seja contraria a uma outra lei \emph{superior}, \emph{natural} ou \emph{positiva}. --- Deve ser \emph{justa} --- em relação ao \emph{fim}, que deve ser o bem da communidade, --- em relação ao \emph{auctor}, que deve ser o superior legitimo, --- em relação á \emph{fórma}, de modo que a divisão dos deveres seja proporcionada ás condições de cada individuo. --- Deve ser \emph{possivel}, de modo que não seja demasiadamente difficil ou gravosa \footnote{Quando a lei \emph{positivo-humana} é contraria á lei \emph{natural} ou \& lei \emph{positivo-divina}, não produz \emph{obrigação} no subdito; porque, derivando toda a obrigação da vontade de \textsc{Deus}, repugna que \textsc{Deus}, pela lei \emph{postitivo-humana}, nos obrigue a fazer o que Elle nos prohibe pela lei \emph{natural} ou \emph{positivo-divina}. Nem o subdito póde obedecer á lei \emph{positivo-humana}, quando não é \emph{honesta}; porque devemos antes obedecer a DrEus do que aos homens. --- Quando a lei \emph{positivo-humana é injusta}, porque quem a promulgou não tinha auctoridade, podemos e, muitas vezes, devemos observal-a, não porque ella nos obrigue, mas porque a sua transgressão póde produzir males para a sociedade. --- Se a lei é \emph{injusta}, o subdito nem póde ser obrigado a expiar a pena estabelecida contra os seus transgressores; pois, n'esse caso, tambem a pena é \emph{injusta}, e não obriga.\par A lei \emph{positivo-humana} está sujeita a \emph{mutabilidade}; porque podem variar\pgn{647}as circumstancias, que a determinaram. --- A mudança é \emph{total} ou \emph{parcial}. E' \emph{total}, quando a lei acaba inteiramente --- ou porque expressa ou tacitamente é \emph{abrogada} pelo legislador, --- ou porque cessa o \emph{fim adequado}, por que foi promulgada, --- ou porque foi substituida por um costume \emph{legitimo}, contrario á propria lei. A mudança é \emph{parcial}, quando o seu effeito, que é a \emph{obrigação}, fica suspenso em certos casos, por \emph{derogagão} ou por \emph{dispensa}.}.

%%>>> art-59

\section*{Artigo II.\\ Consciencia moral, sua divisão e direcção.}
\addcontentsline{toc}{section}{Artigo II. --- Consciencia moral, sua divisão e direcção.}

\pg{647} 101. \textbf{Consciencia moral}. --- \emph{Consciencia moral} é um juizo da razão \emph{pratica}, que, applicando os principios da lei moral a um certo e determinado acto, dieta se este é licito ou illicito \footnote{A applicação dos principios geraes aos casos concretos, que é feita pela nossa razão \emph{pratica} e constitue a consciencia \emph{moral}, forma um syllogismo, expresso ou tacito. A premissa \emph{maior} d'este syllogismo contém um principio geral da lei, dictando que uma coisa é, ou não, licita, por ex.: a \emph{mentira é illicita}; --- a premissa \emph{menor} é um juizo, indicando que um certo e determinado acto está, ou não, comprehendido no \emph{sujeito} d'aquelle principio geral, por ex.: \emph{esta resposta é mentira}; --- a \emph{conclusão} é tambem um juizo, exprimindo que tal acto é, ou não, licito, por ex.: \emph{logo esta resposta é illicita}. --- Este ultimo juizo, contido na \emph{conclusão}, é a \emph{consciencia}. Por isso a \emph{consciencia moral} não é propriamente uma \emph{faculdade}, como o é a consciencia \emph{psychologica}, mas é um acto da razão \emph{pratica}, pois consiste na \emph{applicação} dos principios moraes aos casos particulares. --- O nome de \emph{consciencia} dá-se tambem ao \emph{habito} natural dos principios moraes, embora este tenha o especial nome de \emph{synderesis} (\emph{Of}. \emph{Sum}. \emph{Th}., p. I, q. 19, a. 18).\par O que deixamos eseripto é indicado pela propria palavra --- \emph{consciencia}, a qual é composta de \emph{sciencia} e de \emph{com}; poisque a \emph{consciencia} é a \emph{sciencia} dos principios da lei \emph{moral}, comparada \emph{com} algum acto. --- Por isso a \emph{consciencia} é regra \emph{proxima} dos nossos actos, e a \emph{lei} é regra \emph{remota}. Com effeito, a lei não póde dirigir os nossos actos, se lhes não for applicada; ora esta applicação faz-se pela \emph{consciencia}. Todavia a consciencia não regula os actos humanos senão em virtude da lei; porque a consciencia não forma, mas , descobre e intima a lei ao homem. ;}.

102. \textbf{Divisão} da consciencia \textbf{moral}. --- A consciencia \emph{moral} divide-se em \emph{antecedente} e \emph{consequente}, --- \emph{recta} e \emph{erronea}, --- \emph{certa}, \emph{duvidosa} e \emph{provavel}, --- \emph{delicada}, \emph{laxa} e \emph{escrupulosa} \footnote{A consciencia \emph{moral} póde considerar-se --- em relação ao proprio acto, a que se deve applicar a lei, e é \emph{antecedente} ou \emph{consequente}; --- em relação á verdade \emph{objectiva} do juizo, e é \emph{recta} ou \emph{erronea}; --- em relação ao assentimanto da razão, e é \emph{certa}, \emph{duvidosa} ou \emph{provavel}, --- em relação ao estado \emph{habitual} do espirito, e é \emph{delicada}, \emph{laza} ou \emph{escrupulosa},\par T\par D\par r]}.

a) Consciencia antecedente e consequente, \emph{Antecedente} é a, que se refere a um acto \emph{futuro} e dicta elle se deve praticar \pg{648}  ou omittir. --- \emph{Consequente} é a que se refere a um acto \emph{passado} e dieta que elle foi moralmente bom ou mau \footnote{Aqui tratamos da consciencia \emph{antecedente}, pois é esta a regra \emph{procima} dos nossos actos, e é esta que verdadeiramente merece o nome de \emph{consciencia moral}. --- Tal divisão da consciencia \emph{moral} baséa-se nas diversas funcções, que ella exerce. Na verdade, a consciencia \emph{attesta}, \emph{liga} ou \emph{instiga}, \emph{approva} ou censura (pelo \emph{remorso}). Com effeito, a consciencia consiste, como dissemos, na applicação da sciencia ou do conhecimento aos nossos actos. Esta applicação faz-se por tres modos: --- ou emquanto conhecemos que praticámos ou deixámos de praticar um acto, e neste caso a consciencia \emph{a!festa} (é a consciencia \emph{consequente}, embora não seja propriamente \emph{moral}, mas \emph{psychologica}); --- ou emquanto julgamos que um acto é licito ou illicito, e neste caso a consciencia \emph{liga} ou \emph{instiga} (6 a consciencia \emph{moral-antecedente}, que \emph{liga} a vontade, emquanto prescreve o bem \emph{necessario} e prohibe o mal, e \emph{instiga} a mesma vontade a praticar o bem, que, embora não seja necessario, é comtudo conveniente); --- ou emquanto decidimos que um acto praticado é bom ou mau, e então a consciencia \emph{approva} ou \emph{censura} (pelos \emph{remorsos}), conforme a qualidade do acto (é a consciencia \emph{moral-consequente}).}.

b) Consciencia recta e erronea, --- \emph{Recta} é a que dicta ser lieito ou illieito um acto, que na verdade é lieito ou illicito. --- \emph{Erronea} é a que dieta ser licito ou illicito um acto, que na verdade não é licito ou illieito. --- O erro é \emph{vencivel} ou \emph{invencivel}, conforme póde ou não póde corrigir-se \footnote{Na consciencia \emph{recta}, as premissas do raciocinio, em que se realiza a applicação da lei a um acto particular, devem ser verdadeiras, e a consequencia, que se deduz das premissas, deve ser legitima. --- A consciencia póde ser erronea em relação ao \emph{direito} ou ao \emph{facto}. O erro é \emph{vencivel} ou \emph{invencivel}, conforme póde ou não póde ser vencido pelos meios ordinarios. ;}.

c) \textbf{Consciencia} certa, duvidosa e provavel. --- E' \emph{certa}, quando adherimos \emph{firmemente} á verdade, sem receio de errar. --- E' \emph{duvidosa}, quando não nos decidimos a emittir um juizo; por isso é mera \emph{privação}. --- Eº \emph{provavel}, quando, baseados numa razão \emph{provavel}, emittimos um juizo com receio de errar \footnote{As seguintes observações são de grande importancia.\par a) A consciencia \emph{certa} pôde fundar-se --- ou em motivos que excluem toda a possibilidade da duvida, e temos a certeza \emph{physica}, \emph{metaphysica}, ou moral no sentido rigoroso, --- ou em motivos que excluem a duvida, não \emph{absolutamente}, mas relativamente a uma certa e determinada coisa, e temos a certeza moral no sentido mais lato. --- A certeza \emph{moral} póde ser \emph{directa} ou \emph{reflexa}. E' \emph{directa}, quando se funda nos principios \emph{proprios} e \emph{intrinsecos} da coisa; assim temos certeza \emph{directa} de que \emph{se deve prestar culto} a \textsc{DEUS}, pelo facto da creatura ser \emph{intrinseca e essencialmente dependente de} \textsc{Deus}. E' \emph{reflexa}, quando se funda em principios, que são \emph{extrinsecos} á propria coisa e se referem á \emph{obrigação}, e que indicam se esta existe, ou não; assim, ignorando se um acto, prescripto pelo superior, é por si bom ou mau, recorremos ao principio \emph{reflexo}: « \emph{na duvida}, \emph{deve obedecer-se} ao \emph{superior} n.\par \emph{b}) A consciencia é \emph{duvidosa}, quando está suspensa relativamente á\pgn{649}honestidade do acto. --- A \emph{duvida} pôde ser \emph{especulativa} ou \emph{pratica}, \emph{positiva} ou \emph{negativa}, de \emph{direito} ou de \emph{facto}. --- E' \emph{especulativa}, quando tem por objecto um acto considerado em si, sem nenhuma relação ao agente; --- é \emph{pratica}, quando se não sabe com certeza se o acto, que se deve praticar ou omittir, consideradas todas as circumstancias, é moralmente bom ou mau. Esta duvida \emph{pratica} pertence propriamente á consciencia. --- \emph{E}? \emph{positiva}, quando tem razões, em que possa basear-se; --- é \emph{negativa}, quando não as tem. --- E' de \emph{direito}, quando alguem duvida ou da existencia da lei, ou da extensão da lei a um determinado acto; --- é de \emph{facto}, quando se não sabe com certeza se um certo e determinado acto foi, ou não, praticado. --- Advertimos que a consciencia \emph{duvidosa} só impropriamente se chama \emph{consciencia}; pois quem duvida não emitte \emph{juizo}, e á \emph{consciencia} é essencialmente um \emph{juizo}.\par \emph{c}) A consciencia \emph{provavel} é a que, fundada numa razão \emph{provavel}, emitte um juizo, mas com receio d'errar. --- E' necessario considerar que esta consciencia se diz \emph{provavel}, emquanto ao motivo ou á razão \emph{directa}, em que se funda e que é só \emph{provavel}; porque, emquanto ao motivo ou á razão \emph{reflexa}, que a leya a pronunciar um juizo \emph{pratico} ácerca da bondade ou malicia de um certo e determinado acto, á consciencia não é \emph{provavel}, mas é \emph{certa}. Em menos palavras: a consciencia \emph{provavel} --- \emph{especulativamente é provavel}, --- mas \emph{praticamente é certa}. Assim o subdito póde ter uma razão apenas provavel, que o leva a julgar da malicia de um acto prescripto pelo superior; mas, considerando o principio \emph{reflexo}: « \emph{na duvida}, \emph{deve obedecer-se}. ao \emph{superior} n, forma um dictame \emph{pratico} e \emph{certo} ácerca da bondade d'aquelle acto. --- Ha, pois, differenca entre a consciencia \emph{duvidosa} e a \emph{provavel}.}.

\pg{649} d) Consciencia delicada, laxa e escrupolosa. --- E' \emph{delicada}, quando o homem costuma, em todos os seus actos, julgar com todo o euidado ácerca da bondade ou malicia d'elles, e obedecer a esse juizo. --- E' /\emph{axa}, quando o homem, sem grave e prudente razão, julga não haver malicia, ao menos grave, num certo e determinado acto. --- E' \emph{escrupolosa}, quando o homem, por leves fazoes, julga que ha peccado onde não ha, ou que o peccagf é mais grave do que é na realidade \footnote{Alguns auctores citam tambem a consciencia \emph{perplexa}, que se dá quando se julga illicita tanto a realização como a omissão d'um acto. --- Outros, porém, considerando que a consciencia \emph{perplexa} se refere ao \emph{acto} e não ao \emph{habito}, reduzem-na á consciencia \emph{duvidosa}. :}.

1020 homem não póde desobedecer á consciencia \textbf{certa}, \textbf{quando recta} \&4 invencivelmente \textbf{erronea}.

a) Não póde desobedecer \& consciencia certa, quando recta, --- Quem desobedece á consciencia \emph{certa}, quando \emph{vecta}, quer o que elle sabe ser illicito, e adhere de proposito a um objecto \emph{prohibido}, como tal. Ora isto é moralmente mau. Logo o homem não póde desobedecer á consciencia \emph{certa}, quando \emph{recta} \footnote{Cf. \emph{Sum}. \emph{Th}., 1-2, q. 19, aa. 8-4. --- A consciencia, para que seja a regra dos nossos actos, é \emph{necessario} e \emph{basta} que seja \emph{moralmente certa}.\par \emph{a}) E' \emph{necessario} que á consciencia seja \emph{moralmente certa}, tanto para que\pgn{650}a realização ou a omissão de um acto seja \emph{licita}, como para que essa realização ou omissão seja \emph{obrigatoria}. --- E' \emph{necessario} que a consciencia seja moralmente \emph{certa}, para que a realização ou a omissão de um acto seia \emph{licita}. Por quanto, \emph{o} homem não póde praticar um acto, se estiver \textbar{} \emph{Xalmente certo} de que esse acto é objecto de alguma lei \emph{prohibitiva}, Amo tambem não póde deixar de praticar um acto, se estiver \emph{moralmegeê} certo de que esse acto é objecto d'uma lei \emph{preceptiva}; porque, se não tiver esta \emph{certeza moral} e todavia praticar ou deixar de praticar um acto, despreza a lei e o legislador, e este desprezo é peccado. --- E' \emph{necessario} que a consciencia seja \emph{moralmente certa}, para que a realização ou a omissão de um acto seja \emph{obrigatoria}. Por quanto o homem, para que seja \emph{obrigado} a praticar ou a omittir um acto, deve estar \emph{moralmente certo} de que existe uma lei que preceitua a sua realização ou omissão; visto que um direito \emph{certo}, como é o que o homem tem sobre o exercicio das suas faculdades, não póde ser destruido, alterado ou suspenso, por uma excepção \emph{incerta}. D'ahi o grande principio \emph{reflexo}: a \emph{lei duvidosa} ndo \emph{obriga}, isto é, a lei, cuja existencia, depois de um diligente exame, fica incerta, não impõe nem póde impor uma obrigação \emph{certa}, porque, nesse caso, fulta-lhe a sufficiente promulgação, indispensavel para ella produzir o seu effeito.\par \emph{b}) \emph{Basta} que a consciencia seja \emph{moralmente certa}. A consciencia é \emph{moralmente certa}, quando se basóa num motivo, que merece o assentimento de pessoas prudentes, isto é, quando exclue uma duvida \emph{prudente}. Ora basta esta certeza; porque a certeza deve ser proporcionada á materia, de que se trata, e a materia \emph{moral} só admitte, geralmente, a certeza \emph{moral}.}.

\pg{650} b) \textbf{Não póde desobedecer} á \textbf{consciencia certa}, \textbf{quando invencivelmente erronea}, --- Com effeito, quem desobedece á consciencia \emph{certa}, quando \emph{Znvencivelmente} erronea, quer uma coisa que elle actualmente julga prohibida e má. Ora isto é illieito. Logo o homem não póde desobedecer á consciencia \emph{certa}, quando \emph{invencivelmente} erronea \footnote{Cf. \emph{Sum}. \emph{Th}., \emph{ib}., a. 5. --- Quando, pois, a consciencia \emph{invencivelmente} erronea é \emph{permissiva}, \emph{podemos}, e, quando é \emph{preceptiva} ou \emph{prohibitiva}, \emph{devemos} obedecer-lhe. Quem se conforma com esta consciencia, commette um acto \emph{material} ou \emph{objectivamente} mau, mas \emph{formal} ou \emph{subjectivamente} bom; porque a violação da lei, não sendo \emph{voluntaria} nem em si nem na sua causa (aliás a consciencia não seria \emph{énvencivelmente} erronea), não é imputavel ao agente. Quem, pelo contrario, não se conforma com esta consciencia, \emph{preceptiva} ou \emph{prohibitiva}, pratica um acto \emph{material} ou \emph{objectivamente} bom, mas \emph{formal} ou \emph{subjectivamente} mau. --- A obrigação ou de praticar o acto, que \& consciencia \emph{invencivelmente} erronea prescreve, ou de omittir o que ella prohibe, não deriva da lei particular, àcereà da qual versa o erro \emph{invencivel}, mas sim do principio da lei natural, que obriga a seguir o dictame da consciencia, \emph{moralmente} certo, chamado \emph{lei reflexa da natureza}.}.

104. O \textbf{homem} não póde \textbf{desobedecer}, nem \textbf{obedecer} á \textbf{consciencia vencivelmente erronea}. --- Na verdade, a consciencia \emph{vencivelmente} erronea é um juizo pratico falso, com alguma suspeita do erro. Ora o homem, que tem esta consciencia, faz sempre mal, se \pg{651}  opéra. Faz mal, se opéra contra a consciencia; porque faz o que julga ser illicito, e por isso quer o mal. Faz mal, se opera em conformidade com a consciencia; porque, não empregando esforços sufficientes para subordinar o seu acto á lei, quer livre e conscientemente expor-se ao perigo de pecear e por isso quer o mal. Logo o homem não póde desobedecer, nem obedecer á consciencia \emph{vencivelmente} erronea \footnote{Se o homem não póde desobedecer nem obedecer á consciencia \emph{vencivelmente} erronea, quer seja \emph{preceptiva} quer \emph{prolibitiva}, é claro que elle, antes do acto, deve vencer o erro.}.

105. O homem não pode obedecer á consciencia duvidosa, mas pode obedecer á consciencia \textbf{provavel}.

a) O homem não póde obedecer á consciencia duvidosa, --- O homem não póde querer o mal. Ora quem obedece á consciencia \emph{duvidosa}, pelo facto de se expor, livre e conscientemente, ao perigo de pecear, quer o mal. Logo o homem não póde obedecer á consciencia \emph{duvidosa} \footnote{O homem não póde obedecer á consciencia \emph{duvidosa}, quando a duvida é \emph{pratica}, isto é, quando a duvida se refere á honestidade de um certo e determinado acto; pois quem tenciona praticar um acto, quer elle seja bom quer seja mau, tem uma vontade moralmente má. --- Todavia, se a duvida for \emph{especulativa} e \emph{invencivel}, então \& consciencia, que, baseada num principio \emph{re/fexo}, julga que um certo e determinado acto é licito, \emph{póde} (se é \emph{permissiva}) e \emph{deve} (se é \emph{preceptiva} ou \emph{prohibitiva}) ser attendida, pois está \emph{moralmente} certa, porque a regra da vontade não é constituida pelo juizo \emph{especulativo}, mas pelo \emph{pratico}. --- A duvida \emph{especulativa} póde estar com o juizo \emph{pratico} e \emph{certo}; porque a verdade do juizo \emph{especulativo} e a do \emph{pratico} derivam de principios e razões diversas. Assim, embora o subdito tenha uma \emph{duvida especulativa} ácerca da justiça de uma guerra, feita pelo soberano, todavia, na \emph{pratica}, póde julgar \emph{com certeza} de que pôde licitamente entrar na guerra, emquanto esta se lhe não manifestar certamente injusta, e este juizo \emph{pratico} e \emph{certo} funda-se no principio \emph{reflexo} --- « \emph{na duvida}, \emph{deve obedecer-se ao principe} n.}.

b) O homem póde obedecer á consciencia provavel. --- O homem póde praticar um acto, de cuja honestidade está \emph{moralmente} certo. Ora quem obedece á consciencia \emph{provavel}, está \emph{moralmente} certo da honestidade do acto; porque então a lei, sendo \emph{duvidosa}, não prohibe o acto, e por isso este é licito. Logo o homem póde obedecer á consciencia \emph{provavel} \footnote{O homem póde obedecer á consciencia \emph{provavel}; porque, como dissemos, ella é \emph{moralmente certa} em relação á honestidade de um acto. Com effeito, quando a consciencia é \emph{provavel}, isto é, quando ha razões \emph{provaveis} (que merecem a \emph{approvacdo} de pessoas prudentes), contrarias á existencia da lei, é signal que a propria lei não é \emph{certa}, mas \emph{duvidosa}. Sendo \emph{duvidosa}, a lei não obriga; e estamos \emph{moralmente certos} de que o acto, por não ser\pgn{652}prohibido pela lei, é honesto. Logo podemos obedecer á consciencia \emph{provavel}. --- Note-se que as razões \emph{provaveis,que} negam ou affirmam a existencia da lei, podem ser \emph{mais} ou \emph{menos provaveis}. Como nenhuma razão \emph{provavel} attinge o grau da certeza, é evidente que, tratando-se da \emph{existencia da lei}, podemos deixar uma opinião \emph{mais provavel}, que nega a existencia da lei, para seguirmos a opinião \emph{menos provavel}, que affirma a existencia da lei, comtanto que se funde em solidas razões; porque a opinião, que é \emph{mais provavel}, não exelue o perigo do erro, e por isso póde ser verdadeira a opinião menos \emph{provavel} e póde ser falsa a que o é \emph{mais}. Dissemos --- tratando-se da \emph{existencia da lei}; porque quando o homem deve, por justiça ou por caridade, evitar algum damno ou incommodo, deve preferir a opinião \emph{mais provavel} á que o é menos, isto é, deve seguir a opinião \emph{mais segura}, que mais afasta o perigo; tal deve ser a regra de conducta de um medico relativamente aos remedios, que applica aos doentes.\par Digamos poucas palavras ácerca da consciencia \emph{delicada}, \emph{laxa} e \emph{escrupulosa}. --- Quem tem uma consciencia \emph{delicada}, dé graças ao bom Dus, principio e fonte de toda a perfeição. --- Quem tem uma consciencia \emph{laxa}, pense que, para operar com rectidão, é \emph{necessaria} a certeza «\emph{oral} cerca da honestidade do acto. --- Quem tem a consciencia \emph{escrupulosa}, pense que, para evitar o peceado, basta a certeza \emph{moral}.}.

\pg{652} 106. \textbf{Corollario}. --- Sendo a \emph{consciencia} a regra \emph{proxima} dos nossos actos, e dependendo sobretudo d'ella a conveniente ordenação da nossa vida para \textsc{Deus}, a sua recta \emph{formação} deve ser o escopo de todos os nossos esforços \footnote{Os meios aptos para a recta \emph{formação} da \emph{consciencia} são, principalmente, a reflexão sobre os avisos e os exemplos de pessoas bóas e prudentes o confronto dos nossos actos com a lei moral, a pratica das virtudes, a oração humilde e perseverante.\par --- 9 ---}.

%%>>> parte-4

\part*{PARTE SEGUNDA\\[0.3em]\large DIREITO NATURAL}
\addcontentsline{toc}{part}{PARTE SEGUNDA --- DIREITO NATURAL}

\pg{653} PROLOGO

107. \textbf{Direito natural}. --- \emph{Direito natural} é a parte da philosophia \emph{moral}, que trata dos \emph{deveres} e dos \emph{direitos}, que pertencem ou podem pertencer ao homem em virtude da lei \emph{natural} \footnote{A palavra \emph{jus} costuma tomar-se pelos juristas e philosophos em varios sentidos. --- Umas vezes, a palavra \emph{jus} denota a propria \emph{lei}, e exprime a razão do licito e do illicito; porque os nossos actos são licitos ou illicitos, emquanto são permittidos ou prohibidos pela lei do superior. D'ahi o nome de \emph{jus} dado \& collecção da leis. --- Outras vezes, denota o \emph{objecto} da lei, e significa tudo o que a lei permitte ou manda, tudo o que é justo. --- Outras vozes, exprime o \emph{effeito} da lei; e, como a lei não só impõe \emph{obrigações}, mas tambem dá \emph{poderes} para fazer ou exigir alguma coisa, \emph{jus} significa as \emph{obrigações} e os \emph{poderes moraes} do homem. As \emph{obrigações}, que a lei impõe, chamam-se \emph{deveres}; --- os \emph{poderes}, que a lei concede, conservam o nome de \emph{direitos}. E' neste terceiro sentido que empregamos a palavra \emph{direito}.\par O \emph{direito} pôde ser \emph{natural} ou \emph{positivo}, conforme os \emph{deveres} e os \emph{direitos} derivam da lei \emph{natural} ou da \emph{positiva}. Aqui tratamos do direito \emph{natural}, isto é, dos \emph{deveres} e dos \emph{direitos}, que pertencem ou podem pertencer ao homem em virtude da lei \emph{natural}. --- D'aqui póde ver-se que os adversarios da lei \emph{natural} não reconhecem nem podem. reconhecer o \emph{direito natural}, na verdadeira accepção da palavra. Mas d'isto nos occuparemos mais tarde. --- Dissemos que o \emph{direito natural} trata dos deveres e dos direitos, que pertencem ou \emph{podem pertencer} ao homem, para indicarmos, contra a opinião de \emph{Rousseau}, que o homem, no decurso da vida, póde vir a ter deveres e direitos \emph{naturaes}, diversos dos que tem no seu nascimento.}.

108. \textbf{Dever}. --- \emph{Dever} é a \emph{necessidade moral de praticar ou omittir alguns actos}, \emph{imposta ao homem pela lei}. --- O \emph{dever} --- é uma \emph{necessidade}, porque é \emph{determinação para uma só coisa}; --- é uma necessidade \emph{moral}, porque é uma determinação, que não violenta mas obriga a vontade a praticar ou omittir \emph{livremente} uma coisa; --- é uma necessidade \emph{imposta} pela \emph{lei}, porque a lei não se limita a \emph{aconselhar}, mas \emph{obriga} o homem \footnote{No \emph{dever} devemos considerar a \emph{entidade}, o \emph{sujeito} e o \emph{termo}.\par a) A \emph{entidade} do \emph{dever} é uma coisa \emph{moral}, não \emph{physica}; pois é uma \emph{necessidade moral}, a qual consiste numa determinação \emph{moral} e provém de um\pgn{654}juizo, pelo qual conhecemos que a obediencia ao superior está intimamente ligada com a nossa felicidade eterna.\par \emph{b}) O \emph{sujeto} do \emph{dever} só póde ser o ente \emph{racional e subordinado a um superior}; porque só um ente dotado de \emph{razão} é capaz de \emph{necessidade moral}, e porque esta só póde ser imposta pelo \emph{superior}. Por isso \textsc{Deus}, sendo absolutamente independente, não tem \emph{deveres} a cumprir, Se, pois, dizemos, em algumas circumstancias, que \textsc{Deus} \emph{deve} fazer uma coisa, esta obrigação deriva da sua infinita sabedoria. Se \textsc{Deus} deve alguma coisa, deve-a a si mesmo.\par \emph{c}) O \emph{termo} do \emph{dever}, isto é, o \emph{termo}, a que se refere o \emph{sujeito} do dever, é tambem exclusivamente o ser \emph{racional}; porque ao \emph{dever} corresponde o \emph{direito}, e só o ser \emph{racional} é, como veremos, o sujeito do \emph{direito}. Por isso é um erro estabelecer \emph{deveres} do homem para com os animaes \emph{irracionaes}. --- O homem, pois, só póde ter deveres para com \textsc{Deus}, para \emph{comsigo}, \emph{e} para com os \emph{outros homens}. Na verdade, a perfeição dos nossos actos consiste nas suas relações para com os seus fins; e por isso só os seres, que têm a razão de \emph{fim}, podem ser \emph{termos} dos nossos actos. Ora a razão de \emph{fim} encontra-se em \textsc{Deus}, no proprio agente e nos outros homens. Encontra-se em \textsc{Deus}; porque \textsc{Deus} não só é o primeiro principio, mas tambem o fim ultimo do homem e de todas as creaturas. Encontra-se no proprio agente; porque \textsc{Deus} produziu a8 creaturas para a utilidade do homem. Encontra-se nos outros homens; pois foi tambem para a utilidade dos outros que \textsc{Deus} produziu as creaturas. --- Os \emph{deveres} para com \textsc{DEUS} são os mais nobres e universaes, e o seu cumprimento é mais importante e necessario; porque, sendo \textsc{DEUS} o nosso primeiro principio e o nosso ultimo fim, as nossas relações com Elle são as mais intimas e essenciaes, que se podem imaginar. Os deveres para comnosco excedem, em excellencia, os que temos para com os outros; porque o vinculo, que nos liga aos outros, é de \emph{semelhança}, e o que liga o homem a si mesmo é de \emph{identidade}.}.

\pg{654} 109. \textbf{Divisão} do \textbf{dever}. --- O \emph{dever} divide-se principalmente --- em \emph{natural} e \emph{positivo}, \emph{primitivo} e \emph{derivado}, \emph{juridico} e \emph{moral}, \emph{affirmativo} e \emph{negativo}, \emph{dispensavel} e \emph{indispensavel}.

a) \textbf{Dever} natural e positivo, --- \emph{Natural} é o que é imposto ao homem pela lei \emph{natural}. --- \emph{Positivo} é o que é imposto pela lei \emph{positiva}. --- Aqui tratamos dos deveres \emph{naturaes} do homem \footnote{A lei \emph{positiva}, como vimos na \emph{Ethica}, tira toda a sua força da lei \emph{natural}, \emph{e} pôde ser \emph{divina} ou \emph{humana}, conforme é promulgada por \textsc{Deus}, ou pelo homem, mas em virtude da auctoridade recebida de \textsc{Deus}.}.

b) \textbf{Dever} primitivo e derivado. --- \emph{Primitivo} é o que \emph{immediata} e \emph{proximamente} se funda na \emph{propria natureza} do homem s tal é o dever de \emph{amar a} \textsc{Deus}, de \emph{respeitar os paes}, ete, --- \emph{Derivado} é o que \emph{immediata} e \emph{proximamente} se basóa nalgum \emph{facto contingente}; tal é o dever que resulta de um \emph{contracto} \footnote{Esta divisão indica a diversa \emph{origem} do dever. O dever \emph{derivado}, se \emph{immediata} e \emph{prosimamente} se funda num \emph{facto contingente}, \emph{mediata} e remotamente se basêa na propria natureza do homem; porque os factos só servem\pgn{655}para determinar e applicar os deveres, que, por um modo indeterminado, derivam da natureza. E) sob este aspecto que os deveres \emph{derivados} estão comprehendidos no \emph{direito natural}. --- O dever \emph{primitivo} chama-se tambem \emph{connatural}; e o \emph{derivado} diz-se tambem \emph{adventicio}.}.

c) \textbf{Dever} juridico e moral. --- \emph{Jurídico} é o que obriga por titulo de rigorosa \emph{justiça}, isto é, que obriga a dar a cada um \pg{655}  o que lhe pertence; tal é a restituição de um deposito. --- \emph{Moral} é o que obriga por titulo de outra virtude, distineta da \emph{Justiça}; tal é o dever da \emph{misericordia} para com os pobres \footnote{Alguns auctores negam esta divisão, e dizem que todo o dever é \emph{Juridico}. Mas erram. Na verdade, dever \emph{jurídico} é o que se não póde transgredir, sem offender a \emph{justiça}. Logo haverá \emph{deveres não juridicos}, mas simplesmente \emph{moraes}, se \& sua transgressão envolver offensa não á \emph{justiça}, mas a qualquer outra virtude \emph{moral}. Ora ha effectivamente deveres, que envolvem offensa não á \emph{justiça}, mas a outras virtudes \emph{moraes}; taes são os actos meramente \emph{internos}, os de mera beneficencia ou benevolencia, etc. Logo existem deveres \emph{não-juridicos}, isto é, simplesmente \emph{moraes}. --- Os deveres \emph{juridicos} dizem-se \emph{legaes} ou \emph{perfeitos}; os \emph{moraes}, \emph{imperfeitos}.\par Devemos advertir duas coisas: 1.º) A lei \emph{civil} póde obrigar o homem ao cumprimento dos. deveres \emph{jurídicos}, porque isto é necessario para a ordem publica; mas não póde obrigar ao cumprimento dos- deveres \emph{moraes}, pela razão opposta. --- 2.º) Entre os deveres \emph{moraes}, uns obrigam por forma, que, sem o seu cumprimento, não póde conservar-se a honestidade moral, tal é o dever de conformar a palavra com o pensamento, o exterior com o interior; --- outros são apenas necessarios para a maior perfeição, tal é dever da \emph{affabilidade}, da \emph{liberalidade}, ete.}.

d) \textbf{Dever} affirmativo e negativo. --- \emph{Affirmativo} é a necessidade de praticar um acto \emph{intrinsecamente} honesto; tal é o preceito de \emph{respeitar os paes}. --- \emph{Negativo} é a necessidade de omittir um acto \emph{intrinsecamente} deshonesto; tal é o dever de não \emph{blasphemar} \footnote{Esta divisão do dever é de uma grande importancia. Os deveres \emph{affirmativos} obrigam, ma esta obrigação póde soffrer excepções; ao passo que os \emph{negativos} obrigam sempre e em cada instante, seja qual for a necessidade ou o perigo. Na verdade, o objecto dos deveres \emph{affirmativos}, que por . si é bom e até necessario de um modo geral, póde, pelas circumstancias, tornar-se indifferente e até mau; assim temos o dever de auxiliar o proximo, que está em perigo de morte, mas este dever acaba ou transforma-se num conselho, se a vida dos outros só se póde salvar com o sacrificio da vida propria. --- Pelo contrario, o \emph{objecto} dos deveres \emph{negativos é intrinsecamente} mau, e a malicia \emph{intrinseca} do objecto não póde ser destruida pelas melhores intenções; assim, nunca é permittido \emph{blasphemar}, \emph{odiar os paes}, etc, Sustentar o contrario, com os \emph{atheus}, \emph{materialistas}, \emph{positivistas}, é destruir toda a ordem moral. --- Além d'isso, os deveres \emph{effirmativos} referem-se ao exercicio da \emph{virtude}; os \emph{negativos}, á fuga do \emph{vicio}. Ora é claro que, se não é sempre necessario praticar a virtude, é sempre necessario fugir do vicio. --- Finalmente, o homem, pelo cumprimento dos deveres \emph{negativos} não se afasta do fim ultimo, mas persevera no amor supremo do bem moral; e, pelo cumprimento dos deveres \emph{affirmativos}, tende positivamente para \textsc{DEus} por um certo e determinado modo. Ora, se é absolutamente necessario que o homem nunca\pgn{656}se afaste do seu fim ultimo, e nunca deixe de amar o Bem supremo, não é sempre necessario que elle tenda sempre e positivamente para \textsc{DEUS} por um certo e determinado modo.}.

\pg{656} e) \textbf{Dever} dispensavele indispensavel. --- E' \emph{dispensavel}, quando o subdito, por graves motivos e com a auctoridade do superior, póde deixar de o cumprir. --- E' \emph{indispensavel}, no caso contrario \footnote{Os deveres \emph{primitivos}, quando \emph{negativos}, não admittem dispensa; porque, como dissemos, nunca é permittido praticar o que é intrinsecamente mau, seja qual for o perigo. --- Os \emph{deveres primitivos}, quando \emph{affirmativos}, admittem dispensa; porque, como dissemos, a lei \emph{positiva} não obriga em cada instante; e muitas vezes o homem não só póde mas deve deixar de cumprir um dever \emph{affirmativo}, com tanto que com esta omissão não se transgrida a lei \emph{negativa}, nem se \emph{despreze} a \textsc{DEus}, ete. --- Os deveres \emph{derivados} admittem \textbar{} dispensa, com tanto que a sua omissão não seja prejudicial ao direito dos outros, nem ao bem publico.}.

110. \textbf{Todo} o \textbf{dever}, immediata ou \textbf{mediatamente}, \textbf{deriva} de \textsc{Deus}.

a) Todo o \emph{dever} é uma obrigação \emph{moral}, imposta ás creaturas dotadas de razão e de livre arbitrio. Ora só \textsc{Deus} póde obrigar as creaturas dotadas de razão e de livre arbitrio; porque só Elle é o soberano Senhor de todas as creaturas, e o essencial Ordenador do mundo moral. Logo todo o dever, \emph{immediata} ou \emph{mediatamente}, deriva de \textsc{Deus} \footnote{A obrigação \emph{moral} deriva do nexo essencial, que existe entre um certo e determinado acto e o fim ultimo. Ora só \textsc{DEus} póde fazer com que a consecução do fim ultimo. Logo só \textsc{Deus} póde impor um \emph{dever}. " um certo e determinado acto seja ligado, como condição indispensavel, com}.

b) Todo o \emph{dever} é imposto por uma lei. Ora a lei é \emph{natura}? ou \emph{positiva}. Se a lei é \emph{natural}, o \emph{dever}, imposto por ella, deriva \emph{immediatamente} de \textsc{Deus}; porque \textsc{Deus} é o auctor 2\emph{mmediato} da lei \emph{natural}; --- e, se a lei é \emph{positiva}, o \emph{dever}, imposto por ella, deriva \emph{mediatamente} de \textsc{Deus}; pois a lei \emph{positiva}, como demonstrámos, tem o fundamento na lei \emph{natural}. Logo todo o dever, \emph{immediata} ou \emph{mediatamente}, deriva de \textsc{Deus} \footnote{Não negamos que os \emph{deveres} nos são impostos tambem pela lei \emph{posi}- \textbar{} \emph{tivo-humana}; mas, como a lei \emph{positivo-humana} recebe da lei \emph{natural}, e por isso de \textsc{Deus}, a força de obrigar, sustentamos que todo o \emph{dever} deriva de \textsc{Deus}, \emph{immediata} ou \emph{mediatamente}. --- Esta proposição é combatida pelos \emph{materialistas}, \emph{positivistas}, e por todos os que reconhecem na lei humana a \emph{i} fonte de todos os \emph{deveres} e \emph{direitos}.}.

111. Direito. --- \emph{Direito} é um \emph{poder moral e énviolavel de possuir}, \emph{fazer ou exigir alguma coisa}, \emph{concedido pela lei}. --- O \emph{direito} --- é um \emph{poder}, porque é a faculdade de possuir, fazer / ou exigir alguma coisa; --- é um poder \emph{moral}, porque o \emph{direito} pertence á ordem \emph{moral}, e dirige-se á intelligencia e á von- 4

tade; --- é um poder \emph{énviolavel}, porque impõe aos outros a \pg{657}  \emph{obrigação} ou o \emph{dever} de o não lesar; --- é um poder concedido pela \emph{lei}; porque o \emph{direito} é o effeito proprio da \emph{lei}. --- O \emph{direito} é correlativo ao \emph{dever} \footnote{Todo o \emph{direito} refere-se essencialmente a. um \emph{dever}. Na verdade, não póde conceber-se, num sujeito, um \emph{poder moral} e \emph{inviolavel}, sem se conceber, ao mesmo tempo, nos outros homens, uma obrigação \emph{moral} de omittir ou de fazer o que for necessario para o exercicio d'esse poder; aliás a lei daria a um individuo a faculdade de operar e aos outros a de impedir o acto: o que é contradictorio. --- O \emph{direito} e o \emph{dever} correspondente constituem a \emph{relação juridica}. Por isso as \emph{relações jurídicas} são tantas, quantos são os \emph{direitos} com os seus \emph{deveres} respectivos.\par Em toda a \emph{relação juridica} devemos distinguir quatro elementos: o \emph{sujeito}, o \emph{termo}, \& \emph{materia} e o \emph{titulo} do \emph{direito}.\par \emph{a}) \emph{Sujeito do direito}. --- O \emph{sujeito} do \emph{direito} é o ente que o possue. O \emph{sujeito} só póde ser um ente \emph{racional}, isto é, uma \emph{pessoa} (\emph{physica} ou \emph{moral}). Na verdade, o \emph{direito} é um poder \emph{moral}; ora á \emph{moralidade} só é propria da natureza racional. Além d'isso, o \emph{direito}, na creatura, é uma consequencia do \emph{dever} primitivo e essencial de tender para o fim ultimo: ora só a creatura racional tem esse dever. --- D'aqui se vê o erro dos que, seguindo \emph{Damiron}, \emph{Ahrens}, os \emph{pantheistas} e os \emph{evolucionistas}, admittem os \emph{deveres} do homem para com os animaes irracionaes. Porquanto, se o homem tivesse esses deveres, o animal deveria possuir verdadeiros \emph{direitos}; ora o animal, como vimos, não póde ser o \emph{sujeito} do \emph{direito}. E' verdade que \emph{Ahrens}, em confirmação d'esse erro, diz que o homem tem obrigação de não ser cruel, sem necessidade, para com os animaes; mas respondemos que essa obrigação não nasce \emph{dos-direitos} dos animaes para com o homem, mas sim da \emph{lei divina}, que prohibe o abuso das creaturas, e os sentimentos de erueldade.\par \emph{b}) \emph{Termo do direito}. --- O \emph{termo} do \emph{direito} deve ser tambem um ente racional; pois o \emph{termo} do \emph{direito} tem o \emph{dever} correspondente, e o \emph{dever} é uma obrigação \emph{moral}, de que só póde ser capaz um ente dotado de liberdade e por isso de razão. --- Não é necessario que o \emph{termo} do \emph{direito} exista \emph{actualmente}; basta que seja \emph{possivel} ou \emph{Aypothetico}.\par \emph{c}) \emph{Materia da direito}. --- A \emph{materia} do \emph{direito} é tudo o que é \emph{justo}, isto é, tudo o que é conforme com o fim e com a ordem estabelecida por \textsc{Deus}, tudo o que concorre para estabelecer ou conservar a proporção e o equilibrio nas relações humanas, dando a cada um o que lhe pertence. --- Podem ser \emph{materia} do \emph{direito} as creaturas irracionaes, os actos das nossas faculdades, e os actos dependentes do livre arbitrio dos outros; porquanto, uma coisa, para que possa ser objecto do \emph{poder humano}, deve ter uma natureza inferior e subordinada á do ente racional. Por isso a pessoa \emph{humana}, como tal, não póde ser \emph{materia} de direito. D'ahi a illegitimidade e injustiça da escravatura.\par \emph{d}) \emph{Titulo do direito}. --- O \emph{titulo}, que é o fundamento em que descanca o \emph{direito}, é a razão \emph{objectiva}, \emph{de que provcimamente nasce a efficacia} moral \emph{de um certo e determinado direito sobre} as \emph{intelligencias e as vontades dos seres vacionaes}. Por isso \emph{todo} o titulo do direito é a \emph{apresentação de uma verdade pratica}; pois só a \emph{verdade}, claramente conhecida, póde impor-se á intelligencia; e só uma verdade \emph{pratica}, ou um \emph{dictame da razão}, póde impor-se á vontade. O \emph{titulo} contém dois elementos: um principio \emph{geral}, que é o dictame da razão \emph{pratica}, e a applicação d'esse dictame a um caso particular. ---\par 49 Vor. II.\pgn{658}O \emph{titulo}, pois, do direito é como que um instrumento espiritual, pelo qual um direito particular se demonstra relacionado com toda a ordem racional, sanccionada e imposta por \textsc{Deus}, e que exerce o seu influxo \emph{moral} no \emph{termo} do mesmo direito. D'onde se vé que a efficacia do \emph{titulo} do direito é a propria efficacia da lei \emph{natural}, a qual prescreve a observancia da ordem racional. --- \emph{Num sentido mais lato}, porém, o \emph{titulo} do direito é o \emph{facto particular}, por que o \emph{dictame} da razão se applica a um caso particular e que basta para a determinação de um certo direito; pois, sendo evidente o principio \emph{geral}, só é preciso demonstrar a sua \emph{applicação} a um certo e determinado caso. Assim, sendo evidente o principio \emph{geral} de que tudo o que se compra pertence ao comprador, é claro que, para se provar o \emph{direito} do comprador sobre a coisa comprada, basta provar o facto da compra. Este \emph{facto} é o \emph{titulo} do direito.\par Pergunta-se: Será o \emph{direito} anterior ao \emph{dever}, ou vice-versa?\par Resposta. A questão tem varias soluções, conforme os seus diversos aspectos. Considerado em \textsc{Deus}, o \emph{direito} é anterior ao \emph{dever}; porque \textsc{Deus} só tem \emph{direitos}, e não tem \emph{deveres}. --- Considerado nas creaturas, o \emph{direito}, ás vezes, é posterior ao \emph{dever}; e, as vezes, é anterior. E' \emph{posterior} ao \emph{dever}, quando ambos se encontram no mesmo sujeito; pois cada uma das creaturas tem o \emph{dever} de servir a \textsc{Deus} e de tender para o seu fim ultimo, e este \emph{dever} é o principio de todos os \emph{direitos}. E' \emph{anterior} ao \emph{dever}, quando o \emph{direito} e. o \emph{dever} se encontram em sujeitos distinctos; porque o \emph{direito} de um individuo impõe um \emph{dever} aos outros individuos.}.

\pg{658} 112. \textbf{Divisão} de direito. --- O \emph{direito} divide-se em \emph{natural} e \emph{positivo}, \emph{innato} e \emph{adquirido}, \emph{jurídico} e \emph{moral}, \emph{pessoal} e \emph{real}, \emph{alienavel} e \emph{inalienavel}.

a) Direito natural e positivo. --- \emph{Natural} é o que deriva \emph{immediatamente} da \emph{lei natural}. A somma dos direitos \emph{naturaes} constitue a ordem \emph{juridico-natural}. --- \emph{Positivo} é o que deriva \emph{immediatamente} da lei \emph{positiva}. A somma dos direitos \emph{positivos} constitue a ordem \emph{juridico-positiva}. --- Aqui tratamos dos direitos \emph{maturaes} \footnote{O direito \emph{positivo} póde ser \emph{divino} ou \emph{humano}, conforme deriva da lei \emph{divina} ou \emph{humana}. O direito \emph{humano} póde ser \emph{ecclesiastico} ou \emph{civil}, conforme a lei, de que deriva, é \emph{ecclesiastica} ou \emph{civil},}.

b) Direito innato e adquirido. --- \emph{Innato} é o que se basêa na propria natureza do homem; tal é o direito de aperfeiçoar o espirito. --- \emph{Adquirido} é o que se basêa nalgum facto contingente; tal é o direito sobre uma easa, proveniente da compra \footnote{A. existencia dos direitos \emph{innatos} é incontestavel. O homem, devendo tender para o seu fim, tem o direito de procurar, sem offensa dos outros, os meios adequados para a consecução do mesmo fim; e este direito nasce com o homem, isto é, compete ao homem em virtude da sua \emph{existencia actual} ; tal é o direito de conservar a vida, de aquirir um dominio, de formar sociedades, etc. --- Além d'esses direitos \emph{innatos}, devemos tambem reconhecer outros, que se dizem \emph{adquiridos} e que são determinações dos direitos \emph{innatos}. Assim, temos o direito \emph{innato} de adquirir o que é util para a nossa vida ; mas a determinação d'esse direito, relativamente a uma ou outra coisa,\pgn{659}depende de um facto \emph{contingente}. Por isso os deveres \emph{adquiridos} basêam-se \emph{remotamente} na propria \emph{natureza}, e \emph{procimamente} num facto contingente.}.

\pg{659} c) Direito juridico e moral. --- \emph{Jurídico} é o que se refere directamente a um objecto de rigorosa justiça; tal é o direito á boa fama. --- \emph{Moral} é o que se refere a um objecto de qualquer. outra virtude, distincta da \emph{justiça} tomada no seu mais rigoroso sentido; tal é o direito á \emph{gratidão} \footnote{Os direitos \emph{jurídicos}, que se chamam tambem \emph{perfeitos}, são correlativos aos deveres \emph{juriídicos}, ou \emph{perfeitos}; e os direitos \emph{moraes}, ou \emph{imperfeitos}, correspondem aos deveres \emph{moraes}, ou \emph{imperfeitos}.}.

d) Direito pessoal e real. --- \emph{Pessoal} é o que deriva da dignidade da propria pessoa; tal é o direito de conservar a vida. --- \emph{Real} é o que deriva da posse de uma coisa; tal é o direito de propriedade sobre uma easa \footnote{O direito \emph{real} não existe na \emph{coisa} (pois o sujeito do \emph{direito} é sempre a \emph{pessoa}), mas existe na \emph{pessoa} em razão \emph{da coisa}. Por isso quem adquire a coisa, adquire tambem o \emph{direito}; assim como quem perde a coisa, não póde deixar de perder tambem o direito.}.

e) Direito alienavel e inalienavel. --- E' \emph{aZienavel}, quando, por não ser uma condição indispensavel para o fim ultimo ou para o cumprimento de uma obrigação moral, póde ser transmittido ou renunciado. --- E' \emph{znalienavel}, no caso contrario \footnote{A norma objectiva do \emph{exercicio} e da \emph{alienação} (temporaria ou perpetua) do \emph{direito}, \emph{innato} ou \emph{adquirido}, encontra-se na sua relação com o fim ultimo ou com alguma obrigação moral. --- Erra, pois \emph{Ahrens}, quando diz que o exercicio dos direitos \emph{innatos} é sempre necessario; porque, se assim fosse, o \emph{direito} transformar-se-hia num \emph{dever}.}.

113. Propriedades do direito, --- As propriedades do \emph{direito} são \emph{tres}, a Saber: \emph{coacção}, \emph{collisao} e \emph{limitação}.

a) Coacção. --- \emph{Coacção} é o poder \emph{moral}, que assiste ao homem de empregar a força \emph{physica} para o livre exercicio do seu \emph{direito}. --- A \emph{coacção} não constitue o \emph{direito}, nem o acompanha sempre; embora, se poder empregar-se, sirva para a defeza do direito. --- Todavia, na sociedade, o exercicio da \emph{coacção} não pertence ao individuo, mas á auctoridade publica \footnote{O \emph{direito} é, por si mesmo, \emph{coactivo}. Na verdade, o fim do \emph{direito} é garantir a liberdade e a independencia dos individuos e da sociedade. Ora o \emph{direito} não poderia alcançar esse fim sem a \emph{coacção}, isto é, sem o poder \emph{moral} de empregar a força \emph{p/wysica}, para obrigar os outros ao cumprimento dos \emph{deveres} respectivos; pois os homens, sendo dotados de liberdade, podem violar os \emph{direitos} alheios. --- Todavia a \emph{essencia} do \emph{direito} não é constituida pela \emph{coacção}, como ensinou \emph{Kant}; porque ha muitos direitos, que não são \emph{coactivos} (como o direito de não sermos julgados maus pelos outros, sem razão), e porque a coacção concebe-se posterior ao proprio direito. --- Se a essencia do direito não é constituida pela coacção, muito menos póde ser\pgn{660}constituida pela \emph{actual possibilidade de constranger} os outros ao respeito do direito. Na verdade, o \emph{direito} é um poder \emph{moral}, não \emph{p/ysico}. Ora, se consistisse na \emph{actual possibilidade} de \emph{constranger} os outros, o \emph{direito} deveria comparar-se com a \emph{força physica}, e os mais fortes nunca violariam os \emph{direitos}. Ora isto é absurdo; aliás o passageiro, que não tem a força necessaria para obrigar o ladrão ao respeito da sua propriedade, perderia o \emph{direito} \& mesma propriedade. E' verdade que, nesses casos, póde recorrer-se á coacção \emph{legal}; mas isso mesmo demonstra que o \emph{direito} e a coacção não se identificam. --- Devemos advertir que a coacção existe só em favor dos direitos \emph{jurídicos}; porque só a \emph{violação} d' esses direitos constitue uma injuria e perturba a ordem publica. :}.

\pg{660} b) Collisão. --- \emph{Collisão} é a lucta entre dois \emph{direitos}, que tém a mesma \emph{materia} e não podem ser ambos satisfeitos. --- O \emph{direito} de uma ordem deve ceder ao da ordem mais elevada \footnote{Para a comprehensão do assumpto, advertimos que dois direitos, que entram em \emph{collisão} e que por isso são oppostos, não tém ambos a verdadeira natureza de \emph{direito} (alias a razão, de que os \emph{direitos} derivam, contradizer-se-hia); mas o direito da ordem menos elevada deixa, nesse caso, de ser \emph{direito} ou fica suspenso. --- A \emph{collisão}, pois, é uma \emph{subordinação}.\par Na \emph{collisão} dos \emph{direitos} é necessario observar os seguintes criterios. --- 1.9) Quando ha \emph{collisão} entre direitos \emph{deseguaes}, o verdadeiro \emph{direito} é o. que tem por objecto uma coisa \emph{mais importante} ou \emph{universal}. --- 2.9) Quando ha \emph{collisão} entre dois \emph{direitos}, um dos quaes tem um titulo \emph{evidente} e outro um titulo \emph{incerto}, o verdadeiro é o que tem o titulo \emph{evidente}. --- Estes criterios regulam tambem a \emph{collisüo} dos \emph{deveres} ; pois a collisão dos \emph{direitos} envolve necessariamente a collisão dos \emph{deveres}. Assim, em eguaes circumstancias, os deveres \emph{naturaes} devem ser preferidos aos \emph{positivos}, os \emph{negativos} aos \emph{affirmativos}, os \emph{particulares} aos \emph{communs}, os de \emph{justiça} aos de \emph{caridade}.}.

c) Limitação. --- Todo o direito \emph{humano} tem limites. Para se julgar até onde se extende o \emph{direito}, deve examinar-se a natureza do seu \emph{objecto} e os deveres do seu \emph{sujeito} \footnote{O homem tem direitos, porque deve alcançar o seu fim. Por isso tudo o que se não relaciona com o fim, não é \emph{direito}. --- O \emph{direito} nunca póde extender-se ao que está em coniradieção com o \emph{dever}, que obriga \emph{actualmente} o sujeito; pois o que está em contradieção com tal dever, não é honesto; e não póde haver \emph{direito} para praticar o que não é honesto.}.

114. \textbf{Todo} o direito \textbf{deriva}, \textbf{immediata} ou \textbf{mediatamente}, \textbf{de} \textsc{Deus}.

a) \textsc{Duas} coisas, essencialmente relativas, derivam do mesmo principio. Ora o \emph{direito} e o \emph{dever} são essencialmente relativos; pois todo o \emph{direito} suppõe um \emph{dever}, e vice-versa. Logo o \emph{direito} e o \emph{dever} derivam do mesmo princi pio. Ora todo o \emph{dever} deriva, \emph{Immediata} ou \emph{mediatamente}, de Dus. Logo tambem todo o \emph{direito} deriva, \emph{immediata} ou \emph{mediatamente}, de \textsc{Deus} \footnote{Nada mais absurdo do que pensar que, emquanto o \emph{dever} deriva da ordem moral ou da vontade de \textsc{Deus}, o \emph{direito} respectivo derive de um facto contingente ou da vontade humana. Quem impõe um \emph{dever}, deye ter o \emph{direito} respectivo; --- e quem dá um \emph{direito}, não póde deixar de impor um\pgn{661}\emph{dever}; pois, repetimos, não se concebe \emph{direito} sem \emph{dever}, nem \emph{dever} sem \emph{direito}. Um só, pois, é o principio do \emph{dever} e do \emph{direito}.}.

\pg{661} b) Todo o \emph{direito} é concedido pela \emph{lei natural} ou pela lei \emph{positiva}. Ora, como demonstrámos, a lei \emph{natural} deriva \emph{immediatamente} de \textsc{Deus}; e a lei \emph{positiva}, \emph{mediatamente}. Logo todo o \emph{direito} deriva, \emph{immediata} ou \emph{mediatamente}, de \textsc{Deus} \footnote{Deve attender-se aos seguintes corollarios:\par \emph{a}) Um direito humano, que por si evidentemente repugna a um direito divino, não póde ser verdadeiro \emph{direito}; pois é absolutamente incapaz de ter a força obrigatoria, proveniente da auctoridade divina. --- Quando o direito humano é invalido ou illegal, mas não se oppõe directamente ao direito divino, merece um certo respeito, não por causa da sua legitimidade, mas por causa da auctoridade, de que deriva.\par \emph{b}) A causa \emph{efficiente} do \emph{direito} é sô a \emph{auctoridade divina}, \emph{immediata} ou \emph{mediate}. Por isso, se os contractos, que se fazem entre os homens, tem a força de obrigar, essa obrigação deriva de \textsc{Deus}. Os homens determinam ou subministram apenas as condições do \emph{direito}; dadas as condições indispensaveis, o \emph{direito} nasce da propria ordem moral, e por isso de Dkvus, auctor de toda a ordem, physica e moral..\par \emph{c}) O homem, no uso do \emph{direito}, transgride, ás vezes, uma lei divina. Mas, então, se o uso do direito é mau, não é mau o \emph{direito}. Por isso, se o uso \emph{desordenado} do direito deriva do homem, o proprio direito deriva de \textsc{Deus}.}.

115. A ordem juridica \textbf{existe}.

a) A ordem juridico-positiva existe, --- A ordem \emph{juridico-positiva} existirá, se existirem \emph{deveres} e \emph{direito jurídicos}, e \emph{leis positivas} que impõem esses deveres e concedem esses direitos; porque a ordem \emph{juridica} é o conjuncto dos \emph{deveres} e \emph{direitos juridicos}, e das suas \emph{leis} respectivas, Ora existem \emph{deveres} e \emph{direitos juridicos}, provenientes da lei \emph{positiva}. Logo a ordem \emph{juridico-positiva} existe \footnote{Esta proposição é admittida por todos, embora nem todos concedam á lei \emph{positiva} os caracteres e os effeitos, que nós lhe reconhecemos.}.

b) A ordem juridico-natural existe, --- A ordem \emph{juridico-natural} existirá, se existirem \emph{deveres} e \emph{direitos jurídicos}, e leis \emph{naturaes} que impõem esses deveres e concedem esses direitos. Ora existem \emph{deveres} e \emph{direitos jurídicos}, como é o \emph{dever} e o \emph{direito} á liberdade pessoal, á propriedade, etc.; e estes \emph{deveres} e \emph{direitos juridicos} derivam da lei \emph{natural}, independentemente da lei \emph{positiva}. Logo a ordem \emph{jwuridico-natural} existe \footnote{Esta proposição é negada por todos os que negam a existencia da lei \emph{natural}, e por isso attribuem a origem de todos os \emph{direitos} ou ao poder civil, ou a um contracto feito entre os homens, ou á lei publica, ou ao costume. D'aqui a pouco faremos uma exposição d'esses systemas erroneos. Agora limitamo-nos a dar, em confirmação da nossa these, estas outras provas:\par \emph{a}) Existirá a ordem \emph{juridico-natural}, se o individuo e a sociedade possuirem \emph{direitos}, independentes de toda e qualquer lei \emph{positive}. Ora o individuo\pgn{662} e a sociedade possuem effectivamente taes \emph{direitos}. --- Po:sue-os o individuo. Porquanto, o homem tem \emph{direito} sobre o exercicio das suas faculdades, sobre a sua liberdade, etc. Ora estes \emph{direitos} não derivam da lei \emph{positiva}, mas unicamente da \emph{vontade do Creador}, que deu ao homem esses bens, para que d'elles se servisse em utilidade propria. --- Possue-os a sociedade. Porquanto, a sociedade existe, como veremos, pela \emph{vontade} de \textsc{Deus}. Ora \textsc{Deus} não podia querer a existencia da sociedade, sem lhe conceder tambem os meios para a sua existencia e para a consecução do seu fim; e 1 \textsc{Deus} não teria concedido \& sociedade esses meios, se lhe não tivesse dado o \textbar{} \emph{direito} de exigir dos membros o que é necessario para o bem commum. Logo a ordem \emph{juridico-natural} existe.\par \emph{b}) A lei, que impõe um \emph{dever}, deve conceder tambem o \emph{direito} de o cumprir. Ora á lei \emph{natural} impõe ao homem muitos \emph{deveres}. Logo concede-lhe tambem muitos \emph{direitos}; aliás \textsc{DEUS} teria mandado ao homem que conseguisse o fim, sem lhe dar os meios para essa consecução. Logo existem muitos j \emph{direitos}, provenientes da lei \emph{natural}. Logo a ordem \emph{juridico-natural} existe.\par \emph{c}) Se não houvesse algum direito \emph{natural}, nem poderia haver algum direito \emph{positivo}. Por quanto, a auctoridade civil, quando promulgou a primeira lei, devia possuir o \emph{direito} de preceituar, assim como os subditos deviam ter o \emph{dever} de obedecer. Ora pergunta-se: D'onde derivou esse \emph{direito}? --- Derivou de um contracto? Mas ninguem póde obrigar o homem á observancia d'esse contracto, se não se admitte o principio, anterior a todo o contracto e proveniente da lei \emph{natural}, de que é necessario observar os contractos. --- Derivou do costume ou da persuasão dos povos? Mas este costume ou esta persuasão póde ter força de obrigar só em virtude do direito \emph{natural}. Logo não poderia haver direito \emph{positivo}, se não houvesse direito \emph{natural}.\par \emph{d}) Se não existisse o direito \emph{natural}, mas só o \emph{positivo}, não só não haveria algum acto \emph{intrinsecamente} injusto, mas tambem todo o acto mau seria bom, se a lei \emph{positiva} o preceituasse. Ora estas consequencias são absurdas, como vimos na \emph{Ethica}. Logo a ordem \emph{juridico-natural} existe.\par Objectam: Todo o direito é \emph{coactivo}. Ora o direito, fóra da sociedade civil, não é \emph{coactivo}. Logo, fóra da sociedade civil, não ha direito. !\par Resposta. E' verdade que todo o direito é \emph{coactivo}; mas á \emph{coacção} não e constitue a \emph{essencia} do direito, é apenas uma sua \emph{propriedade}, emquanto quem possue o \emph{direito} tem o poder de empregar a força \emph{physica}, conforme as circumstancias. Este poder de empregar a força existe mesmo fóra da sociedade, embora seja insuficiente e muitas vezes ineffiaz, --- insufficiencia e inefficacia, que só mostram a necessidade da sociedade civil para a defeza dos direitos naturaes do homem,}.

\pg{662} 116. A \textbf{ordem} juridica é \textbf{uma parte} da \textbf{ordem moral}.

a) A ordem \emph{juridica} será uma parte da ordem \emph{moral}, se o objecto da ordem \emph{juridica} for uma parte do da ordem \emph{moral}. Ora o objecto da ordem \emph{juridica}, que é constituido pelos actos de \emph{justiça}, \emph{é parte} do objecto da ordem \emph{moral}, que é constituido pelos actos bons e honestos; porque todo o acto de \emph{justiça} é * por si bom e honesto, embora haja muitos actos bons e honestos, que não pertencem á virtude da \emph{justiça}. Logo a ordem \emph{juridica} é uma parte da ordem moral.

\pg{663} b) A ordem \emph{juridica} será uma parte da ordem \emph{moral}, se a violação de um \emph{direito} envolver a violação de uma virtude \emph{moral}. Ora a violação de um direito envolve necessariamente a violação de uma virtude \emph{moral}, que é a \emph{justiça}; pois esta manda que se dê a cada um o que lhe pertence. Logo a ordem \emph{juridica} é uma parte da ordem \emph{moral} \footnote{Sendo a ordem \emph{juridica} uma parte da ordem \emph{moral}, é facil reconhecer a verdade das seguintes relações, que existem entre a \emph{Ethica} e o \emph{Direito natural}.\par I. A \emph{Ethica} é o \emph{fundamento} do \emph{Direito natural}. Por quanto, os \emph{deveres} e os \emph{direitos}, de que se occupa o \emph{Direito}, encontram a sua origem e a sua força na \emph{justiça}, que é uma virtude \emph{moral}, de que trata a \emph{Ethica}.\par II. Os actos meramente \emph{internos} estão subordinados á \emph{Ethica}, \emph{e} não ao \emph{Direito}. Os actos \emph{externos} não só estão sujeitos á \emph{Ethica}, emquanto se conformam, ou não, com a natureza \emph{racional}, mas tambem ao \emph{Direito}, emquanto se relacionam com a ordem \emph{social}. j\par IIL A \emph{Ethica}, apreciando um acto, considera sobre tudo a intenção do agente; ao passo que o \emph{Direito} suppõe sempre que o acto \emph{interno} corresponde ao \emph{externo}, emquanto se não mostrar o contrario.\par O que deixamos escripto serve para resolver o celebre problema ácerca da distinegão entre a \emph{Ethica} e o \emph{Direito natural}. --- Ha tres opiniões principaes a este respeito: uma de \emph{Kant}, outra de \emph{Ahrens}, e a terceira dos \emph{Escholasticos}. Expomos em resumo cada opinião, fazendo a critica respectiva.\par a) \emph{Kant} diz que o \emph{Direito natural} deve prescindir e separar-se inteiramente da \emph{Ethica}. Porquanto devem distinguir-se duas especies de actos humanos: uns \emph{internos}, outros \emph{externos}. Os \emph{internos} desenvolvem-se na esphera da consciencia e estão sujeitos ás leis internas; os \emph{externos} consistem nas mutuas relações entre os cidadãos, e estão \emph{unica} e \emph{exclusivamente} sujeitos ás leis externas ou civis, prescindindo de toda a moralidade e de qualquer intenção do agente. E, como dos actos \emph{internos} se occupa a \emph{Ethica} e dos \emph{externos} trata o \emph{Direito}, é claro que o \emph{Direito} deve separar-se da \emph{Ethica}.\par A opinião de \emph{Kant} é absurda. Na verdade, provámos que a ordem \emph{juridica}, de que trata o \emph{Direito}, é parte da ordem \emph{moral}, de que se occupa á \emph{Ethica}. --- Além d'isso, é falso que o \emph{Direito}, por se occupar dos actos \emph{ewternos}, possa e deva prescindir \emph{inteiramente} dos \emph{internos}. Por quanto, os actos \emph{externos}, que o \emph{Direito} considera, são actos \emph{humanos}, isto é, provenientes da razão e da \emph{vontade livre} do homem, e por isso dotados de \emph{moralidade interna}. E' esta \emph{moralidade}, que o \emph{Direito} considera, quando a um reu, que assassinou o seu bemfeitor com premeditacao, inflige uma pena mais severa do que a pena, que inflige a quem matou um individuo por inadvertencia ou por um subitaneo movimento de ira. --- Finalmente, como poderia o legislador prescindir inteiramente da \emph{moralidade} e por isso da \emph{Ethica}, se a sua missão é defender a liberdade que os subditos têm de viver honestamente, e promover á virtude e os bons costumes nos cidadãos? --- Logo o \emph{Direito} não deve prescindir nem separar-se inteiramente da \emph{Ethica}.\par bp) \emph{Ahrens}, embora reconheça que estas duas sciencias --- a \emph{Ethica} e o \emph{Direito} --- estão unidas, comtudo sustenta que são \emph{realmente} distinctas, que são entidades completamente diversas; pois imagina que só esta distincção\pgn{664}\emph{real} e \emph{completa} póde obstar a que a auctoridade civil se intrometta nas coisas internas do homem.\par Esta opinião tambem é falsa. Por quanto, nem o Direito \emph{natural} nem o \emph{positivo} prescindem da \emph{Ethica}. Não prescinde da \emph{Ethica} o \emph{Direito natural}; pois este, derivando de \textsc{Deus}, é perfeitissimo e illimitado, e por isso não póde deixar de mandar ao homem que pratique o que é \emph{honesto e} de um \emph{modo honesto}; ora praticar o que é \emph{honesto}, \emph{e} de um \emph{modo honesto}, é preceito tambem da \emph{Ethica}. Não prescinde da \emph{Ethica} o \emph{Direito positivo}; pois este, embora imperfeito e limitado a certos e determinados objectos, não prescinde, nem póde prescindir, da \emph{moralidade interna}, que forma o objecto da \emph{Ethica}, não só porque preceitua actos \emph{Zmanos}, mas tambem porque, na applicação das penas e na irritação dos contractos dolosos ou feitos por medo, considera essa moralidade e o fim do agente. --- A razão, que \emph{AArens} apresenta para sustentar a distincção \emph{real} entre a \emph{Etica} e o \emph{Direito natural}, é futil. Na verdade, aqui trata-se do Direito \emph{natural}, que deriva de \textsc{Deus}, e ninguem póde obstar a que \textsc{Deus} attinja \emph{directamente}, pelas suas leis, \& nossa moralidade \emph{interna}. Mas, embora se tratasse do Direito \emph{civil}, ninguem poderia obstar a que a auctoridade civil attingisse \emph{indirectamente} a moralidade \emph{interna} dos subditos, preceituando actos \emph{internos}, emquanto necessarios para que o acto \emph{externo} seja \emph{Àwmano} e assim corresponda ao fim da lei; pois que o legislador governa homens, não governa corpos humanos. Por isso nem o Direito \emph{natural} nem o \emph{civil} podem prescindir da \emph{E/Aica}.\par €) Os \emph{Escholasticos} dizem que deve admittir-se uma distineção entre a \emph{Ethica} e o \emph{Direito}, mas que esta distineção é apenas \emph{logica} ou de \emph{razão}, pois a \emph{Ethica} e o \emph{Direito} constituem a mesma sciencia \emph{moral}.\par Esta opinião é a verdadeira, e é a que seguimos. --- Com effeito, é necessatio admittir alguma distincção entre a \emph{Ethica} e o \emph{Direito natural}, Porquanto, a \emph{Ethica} considera o homem \emph{em absoluto}, isto é, emquanto é uma pessoa capaz de honestidade; ao passo que o \emph{Direito} considera o homem nas suas \emph{relações} para com outras pessoas. Além d'isso, a Ethica considera a \emph{forma} dos actos humanos, emquanto prescreve as condições necessarias para a sua honestidade; ao passo que o \emph{Direito} considera a \emph{materia} dos mesmos actos, emquanto inyestiga os \emph{deveres} e os \emph{direitos} do homem para cem os outros. --- Todavia esta distincção não excede a que existe entre as varius partes da mesma sciencia. A \emph{justiça}, que é o objecto do \emph{Direito}, é \emph{parte} do bem \emph{moral} ou \emph{honesto}, que é o objecto da \emph{Ethica}. D'este modo, a \emph{Ethica} está para com o \emph{Direito natural} na proporção em que a Metaphysica \emph{geral} está para com a \emph{especiul}. Por isso alguns auctores dão á primeira parte da \emph{Moral} o nome de \emph{Ethica geral} ou \emph{Nomologia}, e \& segunda parte, que é o \emph{Direito natural}, o nome de \emph{Ethica particular}. ou \emph{Deontologia}.\par Para conclusão d'este \emph{prologo}, resta-nos examinar os varios systemas ácerca da \emph{origem} do \emph{Direito}. --- Antes da era christã, o \emph{Direito}, entre os Romanos, que eram senhores do mundo, encobria, sob uma apparencia de liberdade, o mais feroz despotismo. Alli todos eram \emph{escravos}, ao menos do Estado. O Estado, a que se attribuia uma omnipotencia absoluta, era o principio\pgn{665} de todo o direito. --- Este erro e este despotismo foram destruidos pela doutrina eminentemente verdadeira e caridosa do Christianismo. Os escravos tornaram a conquistar a dignidade de \emph{pessoa}; a mulher voltou a ser a companheira do homem. O Estado renunciou á sua omnipotencia, e reconheceu que dependia de um \textsc{Deus} e que, além dos direitos civis e politicos, havia outros direitos mais elevados --- os direitos naturaes e divinos, proclamando a \textsc{Deus} como fonte principal de todos os direitos humanos. D'este modo a Egreja e o Estado, desenvolvendo-se harmonicamente na sua esphera de acção, formavam, na obediencia ás leis divinas, a felicidade espiritual e temporal dos seus subditos. --- Esta união, tão justa e tao util, comecada no seculo IV, foi destruida, infelizmente, no seculo XVI, pelo \emph{protestantismo}. A auctoridade \emph{religiosa} foi devolvida aos principes temporaes; a \emph{Ethica} e o \emph{Direito} perderam a sua origem divina; e o Estado tornou a ser fonte principal, senão unica, de todos os direitos e de todos os deveres.\par As doutrinas do \emph{protestantismo} crearam \emph{cinco} escholas para a explicação da \emph{origem} do \emph{direito}: --- a \emph{utilitarista} de \emph{Hobbes}, de \emph{Bentham} e dos outros materialistas, --- a \emph{individualista} de \emph{Rousseau}, de \emph{Kant}, de \emph{Rosmini}, --- a \emph{pantheista} de \emph{Schelling}, de \emph{Hegel}, de \emph{Krause}, --- á \emph{positivista} de \emph{Stuart Mill}, de \emph{Spencer}, --- a \emph{historica} de \emph{Savigny}. Estas escholas, embora partam de principios diversos, chegam todas á mesma conclusão; pois todas ensinam que a principal fonte dos direitos e dos deveres é o Estado ou a Sociedade civil. --- De cada eschola faremos uma breve exposição com a critica respectiva.\par A) Eschola utilitarista, --- A eschola \emph{utilitarista} é representada sobretudo por \emph{Hobbes} e por \emph{Bentham}. ;\par a) \emph{Hobbes} diz o seguinte: Não ha distincção \emph{intrinseca} entre o bem e o mal. O estado primitivo do homem foi o selvagem. Nesse estado, os homens, errando nos bosques sem noção alguma de honestidade, combatiam-se mutuamente, até que veio a lei civil, imposta pelos fortes aos fracos, pelos reis aos subditos. A lei civil é o \emph{direito}. : :\par O systema de \emph{Hobbes} é falso pelas seguintes razões: --- 1.º) Nega uma. \emph{distincção intrinseca}. entre o bem e o mal. Ora esta negação é absurda. --- 2.º) Suppõe que o estado primitivo do homem foi o selvagem. Ora esta supposicüo é falsa. --- 3.º) Colloca a essencia do \emph{direito} na força \emph{physica}. Ora isto é absurdo. O \emph{direito} é uma força \emph{moral}; pois vigora não só quando encontra garantia na força \emph{p/ysica}, mas tambem, e sobretudo, quando é opprimido. --- 4.º) Destroe toda a idea de \emph{obrigação}. Porquanto, se o homem, como os materialistas sustentam, é absolutamente independente, não ha nenhum motivo que o possa levar a obedecer á lei, sobretudo se for dotado de muita força \emph{p/iysica}, porque então o seu \emph{direito} é mais forte e eficaz. --- 5.º) Leva directamente ao \emph{despotismo} e \& \emph{anarchia}. . b) \emph{Bentham}, tendo collocado na \emph{utilidade} o primeiro principio do bem \emph{honesto}, concluiu que o \emph{direito} é a propria lei civil, promulgada pelo Estado para a \emph{utilidade} de todos.\par O systema de \emph{Bentham} é falso pelas seguintes razões: --- 1.º) Funda todo o \emph{direito} na \emph{utilidade} ou no \emph{interesse}, dizendo que o homem tem o \emph{direito} de fazer tudo o que lhe é \emph{util}. Ora este fundamento é falso; pois demonstrámos na \emph{Ethica} que o \emph{direito} e a \emph{moralidade} não se referem aos bens corporeos\pgn{666}e agradaveis, mas nos bens espirituaes e honestos. --- 2.º) Destroe toda a idea de \emph{obrigação}. Porquanto, se o \emph{direito} consiste em procurar o que é mais \emph{util} e \emph{vantajoso} para o \emph{individuo}, não ha motivo nenhum para \emph{obrigar} o homem a observar as leis, quando esta observancia é contraria aos seus interesses e ás suas satisfações. --- 3.º) Leva ás mesmas consequencias, a que leva o systema de \emph{Hobbes}.\par B) \textbf{Eschola} individualista. --- Esta eschola é representada especialmente por \emph{J}. \emph{J}. \emph{Rousseau}, \emph{Kant}, \emph{Rosmini}.\par a) \emph{Rousseau}, no seu \emph{Cantracto social}, professa a mesma theoria de \emph{Hobbes} ácerca do estado primitivo do homem e ácerca da origem da \emph{moralidade} e do \emph{direito}. Todavia diz que os homens não se reuniram em sociedade pela força de algum poderoso, como pensou \emph{Hobbes}, mas pela sua livre vontade, ede commum accordo crearam a \emph{moralidade} e o \emph{direito}, que deviam servir de norma aos cidadãos. Posto tal fundamento, este escriptor affirma que \& principal fonte de todos os direitos e deveres é o Estado, emquanto, no exercicio das suas funcções, representa a \emph{vontade geral dos cidadãos}. Por isso a \emph{lei} é a \emph{vontade geral dos cidadãos}, e o \emph{direito} é o poder concedido pela \emph{ler}.\par O systema de \emph{Rousseau} é falso pelas seguintes razões: --- 1.º) Suppõe que o estado primitivo do homem foi o selvagem. Ora esta supposição é falsa. --- 2.º) Suppõe que os homens se reuniram em sociedade, em virtude de um \emph{contracto}, livremente concluido entre elles. Ora esta supposicüo é tambem falsa; pois, como demonstraremos mais tarde, não foi a livre vontade do homem que formou a sociedade, mas sim a propria natureza. --- 3.º) Destroe toda a idea de \emph{obrigação}. Porquanto, fazendo \emph{Rousseau} consistir o \emph{direito} numa simples \emph{legalidade externa}, produzida exclusivamente pela \emph{vontade geral} do povo e sem relação a \textsc{Deus}, não ha força nenhuma, a não ser a \emph{f/iysica}, que \emph{obrigue} os cidadãos á observancia das leis do Estado. --- 4.º) Erradamente sustenta que a \emph{lei} é a vontade \emph{geral} dos cidadãos. Porquanto, a \emph{ler}, por ser um \emph{preceito}, deriva da vontade do legislador, e, por ser um preceito \emph{racional}, deve ser conforme á \emph{razão}. Se a \emph{lei} consistisse exclusivamente na \emph{vontade} do legislador, o despotismo seria legitimo e a tyrannia seria um ' direito. --- 5.º) Leva ás mais perniciosas consequencias. Na verdade, se a força do \emph{direito} consiste na mera força material ou na coacção \emph{physica}, como pensa \emph{Rousseau}, o socialismo e a anarchia deixam de ser coisas injustas e illegitimas, com tanto que possuam mais força material ou physica do que 0s partidos conservadores da ordem moral.\par b) \emph{Kant}, na sua \emph{Critica da razão pratica}, corrompeu não só a \emph{Etica} mas tambem o \emph{Direito}. O escriptor allemão, assim como basóa a ordem \emph{moral} na liberdade \emph{interna} do ente racional, assim tambem funda a ordem \emph{juridica} na liberdade \emph{externa}. A liberdade \emph{interna} é, segundo \emph{Kant}, a \emph{autonomia} da razão, de que deriva, como do seu principio, toda a obrigação moral; por isso todos os deveres \emph{moraes}, seja qual for o seu objecto, são \emph{internos}. A liberdade \emph{externa} é consequencia e complemento da \emph{interna}, é um postulado da razão; mas, para a sua subsistencia, precisa de uma legislação \emph{coactiva}, pela qual a liberdade de cada um possa coexistir com a liberdade dos outros. D'ahi a grande diversidade entre a legislação \emph{interna} ou \emph{moral} e a legislação \emph{externa} ou \emph{juridica}; pois nada é mais contrario á legislação \emph{moral} do que\pgn{667}a obrigação \emph{coactiva}. Por isso a differença entre a ordem \emph{moral} e a \emph{juridica} funda-se principalmente na diversidade da legislação. --- D'onde se segue que o \emph{direito} é o \emph{conjuncto das condições}, \emph{pelas quaes a liberdade de cada um póde coexistir harmonicamente com a liberdade dos outros}, \emph{segundo uma lei geral}. E, como todo o verdadeiro \emph{direito} é por si \emph{coactivo}, o mesmo \emph{direito}, tomado no seu mais rigoroso sentido, emquanto exclue todo o elemento \emph{ethico} ou \emph{moral}, póde conceber-se como \emph{a simwultanea possibilidade da} mutua \emph{coacção de todos} com a \emph{liberdade de cada um}, \emph{em conformidade} com as \emph{leis geraes}. D'ahi o \emph{principio geral} do \emph{direito}: « \emph{E}" \emph{conforme ao direito todo o acto}, \emph{em que} a \emph{liberdade de cada um póde existir harmonicamente} com a \emph{liberdade dos outros}, \emph{segundo uma lei geral} ». D'ahi tambem a seguinte \emph{lei juridica}: « \emph{Opéra de modo que o livre uso da tua liberdade possa coexistir com a liberdade de todos os outros}, \emph{segundo uma lei geral} n. --- Estabelecida a \emph{liberdade externe}, como o proprio fundamento philosophico de todo o \emph{direito}, \emph{Kant} conclue que a \emph{liberdade} é o unico direito \emph{innato} do homem e que ella é o principio não só da \emph{egualdade innala} dos homens, da \emph{personalidade Juridica}, de todos os outros direitos, mas tambem da constituição da \emph{sociedade civil}; porque foi um \emph{contracto social}, feito para conservar a cada um a sua esphera de \emph{liberdade externa}, que deu origem \& \emph{sociedade civil}. Por isso o \emph{direito} suppõe o \emph{Estado}, e d'elle deriva exclusivamente. --- A opinião de \emph{Kant}, foi seguida tambem por \emph{Fichte}.\par O systema de \emph{Kant} é falso pelas seguintes razões: --- 1.º) Affirma que o homem é independente de toda a lei divina, que é senhor e legislador de si mesmo, não só como \emph{individuo}, mas tambem como \emph{membro da sociedade}. Ora esta affirmacüo é falsa. O homem não póde ser independente da lei divina, porque não é independente de \textsc{Deus} na creação, na conservação, e na sua finalidade. --- 2.*) Separa a ordem \emph{juridica da ordem moral}, quando, como demoristraremos, aquella é apenas uma parte d'esta. --- 3.º) Destroe a \emph{obrigação} ou o \emph{dever moral}. Na verdade, a obrigação \emph{moral} suppõe sempre dois seres: um superior, que manda, e um subdito, que obedece. Mas, no systema de \emph{Kant}, o homem é, ao mesmo tempo e sob o mesmo aspecto, superior e subdito. --- 4.º). Como o \emph{direito}, segundo \emph{Kant}, é meramente \emph{negativo} (pois consiste na conciliação da liberdade de cada um com a liberdade dos outros), segue-se que o homem teria direito a commetter as maiores immoralidades e crimes, com tanto que não offendesse a liberdade dos outros; pois isso seria conforme com a lei geral da \emph{liberdade}. Nem a esta licença poderia obstar a lei \emph{moral}; porque, no systema de \emph{Kant}, o \emph{direito} deve prescindir inteiramente da \emph{Etica} e de toda á lei \emph{moral}. A lei \emph{geral}, de que falla \emph{Xant}, não é a lei \emph{moral}, mas é a lei da propria \emph{liberdade}. --- 5.º) Leva ao despotismo e á tyrannia. Na verdade, affirmando que o \emph{direito} não se funda na \emph{liberdade individual}, mas num. \emph{contracto social}, feito pela vontade de todos, e que a essencia do \emph{direito} consiste na coacção, é claro que este eseriptor proclama o reinado da força, sacrificando a liberdade \emph{real} do individuo, sem que haja uma norma objectiva, que defenda os \emph{direitos} de cada um contra a maioria ou contra a sociedade, quando esta os desconhecer e calcar, como incompativeis com a sua liberdade. E' assim que, \emph{em nome e pela vontade do povo soberano}, este é reduzido á escravidão e á miseria! (Veja-se o que dissemos, na primeira\pgn{668}parte d' este tratado, contra o systema de \emph{Kant} ácerca do fundamento da \emph{moralidade}). á:\par c) \emph{Rosmini} fundou um systema, em que á soberania de \textsc{Deus} substituiu a soberania do \emph{ente ideal}. Eis o resumo da sua doutrina: O \emph{dever} deriva da lei moral. A lei é uma \emph{noção}, uma \emph{idea} da nossa intelligencia. A primeira \emph{lei} é, pois, a idea do ser \emph{ideal} ou \emph{universal}. « E' claro que o ser, cuja intuição primitiva é illimitada e meramente possivel, o \emph{ser ideal}, tomado universalmente, é para nós a medida de todos os entes particulares; pois é o meio por que os percebemos como são em si mesmos, isto é, com os varios graus da sua entidade. A regra suprema e ultima, por que medimos a entidade dos objectos, é, pois, o mesmo ser por que os concebemos. Este ser ideal, produzindo na intelligencia a idea dos objectos, subministra as regras particulares para a sua apreciação e estima. Estas mesmas regras, consideradas em relação á nossa vontade, que produz a estima pratica, transformam-se em outras tantas leis moraes. Por isso, na Ethica natural, a primeira e \& suprema lei é o ser \emph{ideal}, luz da nossa intelligencia; --- o primeiro \emph{imperativo} é o seguinte: A \emph{tua vontade deve conformar-se com o ser ideal}, \emph{isto é}, \emph{com a luz da} razão » (\emph{Storia comparativa e critica dei Sistemi di morale}, c. T, a. 4). A \emph{obrigação moral} nasce de dois principios: da ordem objectiva das coisas, e do respeito que o homem deve a si mesmo (\emph{Storia comparativa}, etc., c. VII, a. 4; --- \emph{Filosofia del diritto}, Sistema morale, VIL).\par O systema de \emph{Rosmini} é falso pelas seguintes razões: --- 1.º) Destroe a idea de \emph{dever}. Na verdade, o homem não póde estar sujeito a um \emph{dever}, se não existir uma \emph{pessoa}, que tenha o \emph{direito} respectivo. Ora, no systema de \emph{Rosmini}, o sujeito do \emph{direito} é um ser \emph{impessoal}, é o \emph{ser ideal}. Logo não existe \emph{direito} e por isso não existe \emph{dever}. E como poderiamos ter o \emph{dever} de obedecer a um principe \emph{impessoal}, a um legislador \emph{ideal}, se o proprio \emph{Rosmini} confessa que « o conceito de \emph{direito} suppõe uma \emph{pessoa}? » (\emph{Filos}. \emph{del diritto}, pag. 129). E' verdade que o \emph{ser ideal} de \emph{Rosmini} se identifica com o proprio ser de \textsc{Deus}; mus esta identidade, embora logicamente derive dos principios d'esse escriptor, é negada por elle. --- 2.3) Favorece o \emph{racionalismo}. Na verdade, a \emph{obrigação moral}, no systema de \emph{Rosmini}, não deriva da lei de \textsc{Deus}, mas sim da ordem objectiva das coisas e do respeito que o homem deve a si mesmo. D'este modo, o peccado não consiste na offensa feita a \textsc{Deus}, mas na offensa feita a si proprio, ou ao \emph{ser ideal}, cujo \emph{imperativo} se transgride. D'ahi o \emph{peccado philosophico}, condemnado pela razão e pela Egreja. --- 8.º) Leva a pernieiosas consequencias. Porquanto, o peccado, segundo \emph{Rosmini}, consiste na offensa ao respeito, quo o homem deve a si mesmo. Bem. Mas o homem, podendo renunciar ao seu direito, póde fazer com que uma acção deshonesta seja honesta. Poderia dizer-se que a ordem \emph{objectiva} das coisas se oppõe a isso; mas quem diz que o homem deve conformar-se com a ordem objectiva das coisas? E' a luz da razão? Mas esta luz, como dissemos, é \emph{subjectiva}, e não póde \emph{obrigar} o homem; pois repugna que o mesmo individuo seja superior e subdito de si mesmo. --- Como se vê, o \&ystema de \emph{Rosmini}, embora não parta dos mesmos principios, comtudo chega ao mesmo resultado, a que chega o de \emph{Kant}.\par C) Eschola pantheista. --- Ao \emph{subjectivismo} de \emph{Kant} succedeu, por natural\pgn{669}reacção, o \emph{objectivismo} dos \emph{pantheistas}. Os principaes representantes da eschola \emph{pantheista} são \emph{Schelling}, \emph{Hegel}, \emph{Krause}.\par a) \emph{Schelling} diz: Existe o \emph{absoluto}, que é a identidade do sujeito com o objecto. O \emph{absoluto} manifesta-se no mundo da \emph{natureza} e na \emph{historia}. Manifestando-se no mundo da \emph{naíuresa}, o \emph{absoluto} forma as especies, os generos e os reinos; manifestando-se na \emph{historia}, forma á familia, a cidade e o Estado. O Estado é o organismo objectivo da liberdade ou a harmonia da necessidade com a liberdade, e tem por condição a \emph{ordem Juridica}. O \emph{Estado} e á \emph{ordem juridica} não são resultantes das livres vontades dos individuos, mas são a forma \emph{immediata} por que se manifestam a \emph{razão} e a \emph{vontade universal}. Por isso o Estado é um organismo \emph{absoluto}.\par Este systema é absurdo, não só pelas razões com que combatemos, na \emph{Cosmologia}, o \emph{absoluto} de \emph{Schelling}, mas tambem porque, destruindo toda a base da dignidade pessoal e do livre arbitrio, reduz o \emph{Direito} a uma cega fatalidade, que pesa sobre os cidadãos.\par b) \emph{Hegel} desenvolveu e completou o trabalho de \emph{Schelling}. O systema de \emph{Hegel} é o seguinte: O \emph{absoluto} est\& sujeito a evoluções \emph{necessarias}. O mundo não é a propria vida do \emph{absoluto}, desenvolvida por um movimento \emph{real}, como disse \emph{Schelling}, mas é a propria analyse \emph{logica} do \emph{pensamento absoluto}. Este \emph{pensamento absoluto} é substancial, mas não é pessoal nem adhere a um sujeito pensante. D'onde se segue que as regras absolutas do \emph{pensamento} são tambem as do \emph{ser}, e que tudo o que é ou que se faz, é ou faz-se por necessidade \emph{logica}. Por isso a philosophia-divide-se em tres paries: --- a \emph{primeira} parte, que é a \emph{logica}, é a sciencia da \emph{idea} em si e por si; --- a segunda, que é a philosophia da \emph{natureza}, é a sciencia da idea, que se manifesta fóra de si mesma; --- a terceira, que é a philosophia do \emph{espirito}, é a sciencia da idea, que faz reflexão sobre si e d'este modo adquire a consciencia de si mesma.\par A philosophia do \emph{espirito} subdivide-se em tres partes, conforme o triplice grau da sua evolução. Porquanto, o \emph{espirito} desenvolve-se ou revela-se --- primeiramente no estado \emph{suljectivo}, emquanto se considera em si e por si, --- em seguida, no estado \emph{objectivo}, emquanto se traduz na realidade ou no mundo, --- finalmente, no estado \emph{absoluto}, emquanto, recolhendo-se e differenciando-se de tudo o que está fóra d'elle, toma posse completa de si mesmo. --- A philosophia do \emph{direito} está contida na philosophia do \emph{espirito objectivo}. O principio, de que deriva o \emph{direito}, é a vontade universal e objectiva da idea \emph{absoluta}, --- vontade substancial e racional, mas impessoal. Esta vontade é o fim de si mesma, isto é, o seu fim interno é a liberdade; o seu objecto é o mundo externo, em que ella se realiza. Mas, nesta realização externa, a liberdade assume a forma de \emph{necessidade}. Por isso o \emph{direito} é a \emph{liberdade realizada}.\par O \emph{direito}, assim entendido, póde considerar-se sob um triplice aspecto. Porquanto, a livre vontade póde manifestar-se como \emph{individua} ou \emph{pessoal}, e a sua relação para com as coisas constitue a \emph{propriedade}, e para com as outras pessoas constitue o \emph{contracto}; o direito, neste caso, é formal e abstracto. Em seguida, a livre vontade reflecte sobre si mesma; d'ahi resulta o direito da vontade \emph{subjectiva} --- a \emph{moralidade}. Finalmente, a vontade torna-se\pgn{670}\emph{social}, traduzindo-se na vida historica e nas tradições dos povos; d'ahi o \emph{ethos} (idea da liberdade transformada em mundo real). --- N'este ultimo grau, a livre vontade manifesta-se sob tres formas successivas: --- primeiramente, na \emph{fumilia}, cujo fim é a unidade pelo amor, --- em seguida, na \emph{sociedade civil}, que liga as familias numa unidade mais elevada; --- finalmente, no \emph{Jistado}, em que o \emph{e/hos} attinge o supremo grau de perfeição. --- O Estado, pela sua vez, desenvolve-se no direito interno, no direito externo e na historia do mundo. O Estado é a realidade da idea ethica, a realidade da vontade substancial, a propria razão, o fim absoluto, em que a liberdade attinge o mais alto grau de direito. O Estado tem o supremo direito sobre cada individuo. O Estado é a vontade divina, é o proprio \textsc{Deus} presente, é o universo espiritual em que a razão divina se realizou.\par O \emph{direito publico} é triplice: \emph{interno}, \emph{externo}, e \emph{absoluto} ou \emph{direito} da \emph{historia universal}, conforme a idea do estado --- ou se refere a um determinado organismo individual, considerado em si mesmo, --- ou consiste na relação de um estado para com os outros, --- ou é a idea universal do estado, considerado como genero o assim dotado de poder absoluto sobre todos os estados particulares, isto é, é o espirito, que se realiza no processo da historia do mundo, e por isso é o unico e absoluto juiz. Os espiritos proprios dos varios povos encontram a suprema unidade no espirito do mundo. --- A Egreja o a Religião estão identificadas com o Estado, ou subordinadas a elle; porque o espirito, envolvido sob o véo da religião, é subjectivo e objectivamente indefinido, e, só no Estado adquire uma objectividade externa e real. Por isso a Religião não póde governar. --- \emph{Hegel} foi seguido por \emph{Karl} Marce e por muitos outros, sobretudo na Allemanha.\par O systema de \emph{Hegel} é falso, não só pelas razões que adduzimos na \emph{Cosmologia}, mas tambem pelas seguintes: --- 1.3) Funda-se no \emph{pantheismo}, negando a existencia de um \textsc{Deus} pessoal, Creador e Conservador do mundo, o \emph{livre arbitrio}, a \emph{imputabilidade moral}, \emph{á immortalidade} de cada alma humana, e fazendo derivar tudo do \emph{pensamento}, que não tem \emph{sujeito} nem \emph{objecto}, que é \emph{nada}. Ora tudo isto é absurdo. --- 2.º) Encerra contradicções. Porquanto \emph{Hegel}, rejeitando o \emph{idealismo} subjectivo e espiritual de \emph{Kant}, esforca-se por recorrer a um principio unico, universal e absoluto de todas as coisas. Mas, longe de procurar esse principio em \textsc{Deus}, Creador e Governador de todas as coisas, que, pela sua lei \emph{eterna}, é a fonte immutavel e objectiva de toda a moralidade e de todo o direito, procurou-o e disse tel-o encontrado numa supposta divindade, não real mas ficticin e imaginaria, que se desenvolve até adquirir a consciencia de si mesma no homem, e que é realmente presente, como vontade universal e como alma universal no Estado e na historia. D'este modo o \emph{objectivo de Hegel}, sendo apenas a evolução do \emph{pensamento} (sem sujeito e sem objecto), nao differe, na realidade, do \emph{subjectivo} de \emph{Kant}. --- 8.3) Destroe todo o fundamento da justiça e da moralidade social; porque tudo está numa continua evolução, e todas as nações têm apenas um valor historico e relativo, determinado pelas phases do desenvolvimento do espirito universal --- 4,3) Destroe a liberdade individual e os direitos individuaes, reconhecendo o Estado, como unica fonte absoluta de todo o direito, como principio de todo o dever e regra de toda a moralidade, completamente autonomo\pgn{671}. Por isso, tudo o que o Estado fizer ou mandar, é justo, é honesto e tem força do \emph{direito}, só porque foi o Estado quem o fez ou mandou fazer. D'este modo, o \emph{direito} consiste na força e na felicidade dos acontecimentos. Ora tudo isto é falso, como fica demonstrado.\par o) \emph{Krause} propoz o seguinte systema: \textsc{Deus} é a essencia unica, infinita e total; é indeterminado e infinito; contém em si todos os seres finitos, como o todo contém as partes. Este ser realiza a sua essencia, variando continuamente no tempo os seus estados e as suas determinações. A realização, num tempo infinito, da essencia de \textsc{Deus} é a \emph{vida}. A vida de \textsc{Deus} contém a vida de todos os seres. --- A essencia é o \emph{bem}. Realizar a \emph{essencia} é fazer o \emph{bem}. Para que o ente realize a essencia ou uma parte da unica essencia, deve ter \emph{liberdade}, que por isso é a propriedade pela qual o ente se determina para a realização da essencia. --- A \emph{moralidade}, pois, é \emph{a} propriedade, pela qual o ser finito produz livremente o bem, porque é bem; e a \emph{maxima} d'esta moral é a seguinte: \emph{Fase o bem}, \emph{porque é bem}, isto é, porque é uma parte da essencia divina, que manifesta ou realiza a sua vida no tempo.\par Mas o bem do homem consiste na realização de tudo o que está contido na sua essencia, isto é, no desenvolvimento de todas as suas faculdades e das diversas relações, que póde contrahir. As varias espheras, em que o homem póde desenvolver-se, são as seguintes: a \emph{religião}, \emph{á moralidade}, \emph{á sciencia}, \emph{á arte}, a \emph{educação}, o \emph{commercio} e o \emph{direito}. Para a realização d'esses fins, o homem deve ter promptos os \emph{meios} ou as \emph{condições} indispensaveis. Estas \emph{condições} podem ser subministradas pela natureza, como a agua, o fogo, etc., --- ou podem depender da vontade humana, como são os auxilios physicos, intellectuaes, moraes, etc. --- O homem póde exigir as \emph{condições}, que são indispensaveis para o desenvolvimento da sua vida e que dependem da vontade humana. D'ahi a origem do \emph{direito}. O \emph{direito}, pois, é o \emph{conjuncto das condições}, \emph{dependentes da vontade humana e necessarias para a} consecução \emph{do fim}, \emph{que foi determinado} ao \emph{homem pela sua natureza racional e que é o desenvolvimento da sua vida ou a realização do que está contido} na \emph{sua essencia}. --- Existe um só \emph{direito} --- o direito immanente de \textsc{Deus}, abrangendo todos os direitos particulares dos seres finitos. Para tornar effectivos os direitos do homem e da humanidade, isto é, para garantir as condições, o: são necessarias para o desenvolvimento do homem e da humanidade e-que dependem da vontade humana, foi instituido o \emph{Estado}, que tem um poder \emph{absoluto} sobre todas as sociedades particulares.\par \emph{Krause} foi seguido por \emph{A/Arens}, \emph{Tiberghien}, e por outros escriptores allemães, hespanhoes (\emph{Sans del Rio}, \emph{Giner de los Rios}, \emph{Salmerón}) e portuguezes.\par O systema de \emph{Krause} deve rejeitar-se, porque, baseado no \emph{pantheismo}, apresenta uma definição, que: desvirtua completamente a noção de \emph{direito} e é um conjuncto de falsidades. \emph{Xrause} diz que o \emph{direito} é o \emph{conjuncto das condições}, \emph{dependentes da vontade humana e necessarias para a consecução do fim}, \emph{que foi determinado} ao \emph{homem pela sua natureza racional e que é o desenvolvimento da sua vida} ou a \emph{realização do que está contido} na \emph{sua essencia}. As falsidades d'esta definição são tantas, quantas são as palavras. --- 1.0) E' falso que o fim do homem é o desenvolvimento do que está contido na sua essencia; porque, como vimos, o nosso fim ultimo é um bem absoluto e\pgn{672}infinito, que 6 \textsc{Deus}. --- 2.0) E' falso ser a nossa natureza racional, que nos assigna o fim ultimo; porque o poder do estabelecer uni fim, e de impor o dever de o conseguir, não póde derivar do proprio ser que tende para o fim, mas deriva da intelligencia e da vontade do superior, embora este tenha dese conformar com a natureza do mesmo ser. --- 3.9) E' falsa a supposição de que o homem possa desenvolver a sua essencia nas sete espheras, apontadas por \emph{Krause}; porque, sendo as forças do homem limitadas, o exercicio de uma faculdade impede o das outras, e quem quizesse desenvolver-se em todas as espheras, acabaria por se não desenvolver em nenhuma. --- 4,0) E falso que as \emph{condições}, indispensuveis para o desenvolvimento, \emph{dependem da vontade humana}; porque tudo o que é necessario para a consecução de um fim necessario, não depende da vontade humana, mas é subministrado pela propria natureza. --- 5.9) E falso que o \emph{direito}, que é uma faculdade ou um poder \emph{moral}, se identifica com as condições, que muitas vezes acompanham a realização do mesmo \emph{direito}.\par D) \textbf{Eschola} positivista. --- A eschola \emph{positivista} tem alguma analogia --- com a \emph{pantheista}, emquanto ao processo de \emph{evolução}, --- e com a \emph{utilitarista}, emquanto prescinde inteiramente de toda a idea do absoluto e da Divindade. Seus principaes representantes, nesta questão, são \emph{H}. \emph{Spencer} e \emph{Stuart Mill}.\par a) \emph{H}. \emph{Spencer} funda o systema \emph{juridico} na sua doutrina acerca da \emph{evolução}. Diz assim: Existe uma lei, cujo fim é a preservação da especie --- a lei da sobrevivencia do mais forte. Em virtude d'essa lei, os individuos, que mais se adaptam ás condições da existencia, recebem mais vantagem e utilidade do que os individuos, que menos se adaptam a essas condições. Esta lei suppõe que todo o individuo deve receber a sua parte de bens e de males, proporcionada á sua natureza e á sua conducta. --- A idea de \emph{justiça} contém dois elementos: um \emph{positivo}, e outro \emph{negativo}. O \emph{positivo} é o reconhecimento do direito, que todo o homem tem á sua livre actividade e ás vantagens que d'ella resultam; o \emph{negativo} é o reconhecimento dos limites, impostos pela presença dos outros homens, que gozam direitos analogos. D'ahi a formula da \emph{justiça}, que, reunindo os dois elementos, affirma que a cada um deve conceder-se a livre actividade, mas limitada pela presença dos outros homens, que tem as mesmas exigencias. O elemento \emph{positivo} exprime as condições da vida em geral; o \emph{negativo} determina estas condições, na hypothese da pluralidade dos homens. Por isso a formula é a seguinte: \emph{Cada um póde fazer o que quiser}, \emph{com tanto que ndo lese a liberdade dos outros}. --- A posse desta \emph{liberdade limitada} é o \emph{direito}. :\par O systema de \emph{Spencer} é falso pelas seguintes razões: --- 1.º) Basea-se num falso supposto, que é a \emph{evolução}; pois esta \emph{evolução}, como demonstrámos, é contraria á razão e á experiencia. --- Os adversarios dizem que essa evolução é uma hypothese. Embora. Mas, sendo uma hypothese, deve ser rejeitada pelos \emph{positivistas}, que só admittem o \emph{positivo}, isto é, só o que foi verificado pela experiencia externa e sensivel. --- 9.3) Faz consistir o \emph{direito} na força \emph{physica}; porque sustenta que deve respeitar-se a actividade de cada um, e que esta actividade é proporcionada ao fatal desenvolvimento das melhores e mais perfeitas aptidões, e por isso vencerá sempre o mais forte. Nem se diga que o mais forte tem o \emph{dever} de respeitar a liberdade do mais\pgn{673}fraco; porque esse mesmo \emph{dever} incumbe ao mais fraco em relação ao mais forte. --- 3.º) Destroe toda a idea de \emph{obrigação moral}. Na verdade, para que deveria o homem respeitar a livre actividade dos outros?.. Se tem a força bastante para se apossar de um determinado bem, elle póde exercer essa força; porque, no systema de \emph{Spencer}, não ha um principio mais alto e absoluto, que approve a legitima actividade e condemne o abuso do poder e da força. --- 4.º) Leva ás mais funestas consequencias. Porquanto, se o homem tem o direito de fazer o que não offende a liberdade dos outros, todo o crime, que não offende a liberdade dos outros, será licito e honesto.\par b) \emph{Stuart Mill} applica o seu \emph{utilitarismo positivista} ao \emph{Direito}, como o tinha applicado á \emph{Ethica}. Todo o esforço do escriptor inglez consiste em dar-nos a idea da \emph{justiça} e do \emph{direito}. --- O que é a \emph{justiça}? \emph{Mill} não \& define directamente. Limita-se a descrever o \emph{sentimento da justiça}, que é o desejo ou instincto da vingança e da defeza pessoal, o qual, pela faculdade que o homem tem de communicar os seus sentimentos de sympathia e pela concepção humana de um \emph{egoismo intelligente}, abrange todas as creaturas. --- O que é o \emph{direito}? O \emph{direito} consiste em ter \emph{titulos}, capazes de levar a sociedade a garantir a posse de alguma coisa e a punir quem se oppõe á mesma posse. A razão, por que a sociedade deve garantir esta posse, é a \emph{utilidade geral}. --- Eis a que se reduz toda a doutrina de \emph{Stuart Mill}.\par Este systema é falso pelas seguintes razões: --- 1.º) Desvirtua a idea de \emph{justiça}; pois esta não consiste num sentimento de vingança e de defeza pessoal, mas consiste em dar a cada um o que lhe pertence. --- 2.2) Não explica a idea de \emph{direito}. Na verdade, os \emph{titulos}, que levam a sociedade a garantir a posse de uma coisa e constituem o \emph{direito}, d'onde derivam? Da sociedade? Não; porque são anteriores a ella. Logo este systema, não explicando a origem dos \emph{titulos}, não explica a idea de \emph{direito}. --- 3.º) Se este systema não explica a idea de \emph{direito}, nem póde explicar a idea de \emph{obrigação}, correspondente a esse \emph{direito}. Na verdade, quem póde \emph{obrigar} os outros a respeitar o \emph{direito}? Respondem que é a \emph{utilidade}. Mas perguntamos: Esta \emph{utilidade}, que \emph{obriga} ao respeito do \emph{direito}, é a utilidade do \emph{sujeito} ou do \emph{termo} da obrigação? Se é a \emph{utilidade} do \emph{sujeito} da obrigação, temos a consagração do \emph{egoismo}, com prejuizo de quem possue o \emph{direito}. Se é a do \emph{termo} da obrigação, o \emph{direito} fica assim exposto a todas as violações; porque, neste systema, não ha uma lei superior, que obrigue o homem a preferir á utilidade propria a utilidade dos outros.\par E) Eschola historica. --- A \emph{eschola historica}, fundada por \emph{Savigny} (1779-1861) para remediar os males causados pelo \emph{transcendentalismo}, nega o direito \emph{natural} e só reconhece o direito \emph{positivo}. O direito \emph{positivo}, segundo \emph{Savigny}, deriva da vontade livre dos povos, a qual se manifesta pelas leis escriptas e pelos patrios costumes. Por isso não deve procurar-se um direito \emph{commum} a todos os povos; cada povo forma naturalmente um direito \emph{especial}, accommodado ás suas cireumstancias. Por isso só a \emph{historia} deve guiar-nos na investigação dos direitos humanos. --- Este systema foi seguido por \emph{Stahl}, \emph{Ihering}, \emph{Bluntschli}.\par O systema, proposto pela \emph{eschola historica}, é falso pelas seguintes razões: --- 1.3) Destroe a idea de \emph{obrigação moral}. Na verdade, se se prescindir da\par 48 Vor. IL.\pgn{674}lei \emph{natural}, não ha, nem póde haver, um motivo capaz de obrigar a vontade humana a sujeitar-se á lei \emph{positiva}, a não ser a força \emph{physica} ou \emph{material}; e assim chega a estabelecer-se a sociedade civil sobre o fundamento da força physica ou material --- 2.*) Encerra o mais despotico e tyrannieo \emph{naturalismo}; porque sustenta que as leis civis, embora contrarias á Moral e ao Direito natural, são verdadeiras leis e têm a força de obrigar, e que ninguem, nesses casos, póde resistir ao poder civil. Ora é simplesmente impio e deshonesto que a sociedade civil prescinda inteiramente de \textsc{Deus}, Senhor e Governador de todas as creaturas. --- 3.º) Erradamente ensina que é só a \emph{historia} que nos deve guiar na investigação dos \emph{direitos} humanos. Porquanto, a \emph{historia}, se póde indicar o que é ou o que foi, não diz o que \emph{deve} ser, nem o que é conforme ao ideal da \emph{justiça}, e por isso não póde levar ao conhecimento dos \emph{direitos}, de que o homem é dotado para a consecução do seu fim. --- 4.º) Leva ao \emph{fatalismo} dos pantheistas. Porquanto, os seus defensores affirmam que a \emph{vontade} do povo, que é a fonte de toda a lei, se desenvolve por um processo natural, á semelhança do organismo physico. Ora isto é \emph{fatalismo}.\par Concluindo a exposição d'estes systemas cerca da \emph{origem} do \emph{direito}, vemos que todos prescindem da lei \emph{natural}, que foi promulgada por \textsc{DEus} ao homem pela propria luz da razão, e que é fonte e norma do direito \emph{positivo}, ao qual impõe regras e limites, salvando assim os direitos dos individuos e a existencia da sociedade, e obstando a que o direito se transforme no capricho do mais forte. --- Foi com o fim de tutelar a liberdade e a dignidade do homem e de firmar as bases da sociedade, que esses falsos systemas foram condemnados pela Egreja Catholica, Mestra infallivel da verdade.}.

117. \textbf{Divisão do Direito natural}. --- Dividimos o \emph{Direito natural} em duas secções. --- Na \emph{primeira} secção tratamos dos \emph{deveres} \pg{664} \pg{665}  e \emph{direitos}, que pertencem ás relações \emph{individuaes} do \textbf{homem}, isto é, que pertencem ao homem, emquanto \emph{individuo}; e por \pg{666}  isso consideramos os \emph{deveres} do homem para com \textsc{Deus}, os \emph{deveres} do homem para \emph{comsigo}, e os \emph{deveres} e \emph{direitos} do \pg{667}  \textbf{homem} para com o \emph{proximo}, isto é, para com os outros individuos da mesma especie humana. Na \emph{segunda} secção tratamos \pg{668}  dos \emph{deveres} e dos \emph{direitos}, que pertencem ás relações \emph{sociaes} do homem, isto é, que pertencem ao homem, emquanto é \emph{membro} \pg{669}  \emph{da sociedade}, quer esta se considere em si, quer se considere em relação a outra sociedade \emph{civil} ou \emph{religiosa}; e por isso \pg{670}  consideramos --- os deveres e os direitos \emph{domesticos}, isto é, os \emph{deveres} e os \emph{direitos} que são proprios do homem emquanto é \pg{671}  membro da sociedade \emph{conjugal}, ou \emph{parental}, ou da \emph{dominical}, --- 08 deveres e os direitos \emph{civis}, isto é, que são proprios do \pg{672}  homem, emquanto é membro da sociedade \emph{civil}, --- os deveres e os direitos \emph{internacionaes}, isto é, que são proprios de uma \pg{673}  sociedade civil, chamada vulgarmente \emph{nação}, com relação a outra sociedade, tanto civil quanto religiosa. --- A \emph{primeira} \pg{674}  secção constitue o \emph{Direito individual}; a \emph{segunda} secção constitue o \emph{Direito social}.

\part*{SECÇÃO PRIMEIRA\\[0.3em]\large DIREITO INDIVIDUAL}
\addcontentsline{toc}{part}{SECÇÃO PRIMEIRA --- DIREITO INDIVIDUAL}

\pg{675} 118. \textbf{Divisão} do Direito individual. --- Dividimos o Direito \emph{individual} em \emph{tres} capitulos. No \emph{primeiro} capitulo tratamos dos \emph{deveres} do homem para com \textsc{Deus}; no \emph{segundo}, dos \emph{deveres} do homem para \emph{comsigo}; no \emph{terceiro}, dos \emph{deveres} e \emph{direitos} do homem para com o \emph{proximo}.

A

%%>>> cap-10

\chapter*{CAPITULO PRIMEIRO\\[0.3em]\large DEVERES DO HOMEM PARA COM DEUS}
\addcontentsline{toc}{chapter}{CAPITULO PRIMEIRO --- DEVERES DO HOMEM PARA COM DEUS}

{\small\noindent\textsc{Summario}: --- Deveres theoricos do homem para com \textsc{Deus}. --- Deveres praticos do homem para com \textsc{Deus}.\par}

%%>>> art-60

\section*{Artigo I.\\ Deveres theoricos do homem para com Deus.}
\addcontentsline{toc}{section}{Artigo I. --- Deveres theoricos do homem para com Deus.}

119. \textbf{Deveres} do homem \textbf{para} com \textbf{Deus}. --- \emph{Deveres} do homem para com \textsc{Deus} são as obrigações, que temos de prestar a \textsc{Deus} o culto que Lhe é devido. --- Como a virtude, que nos inclina a prestar a \textsc{Deus} o culto que Lhe é devido, chama-se \emph{Religião}, póde dizer-se que todos os \emph{deveres} do homem para com Dzvus se resumem na virtude da \emph{Religião} \footnote{Num sentido \emph{mais lato}, todos os deveres do homem são deveres para com \textsc{DEus}; não só porque são impostos por Dus, mas tambem porque nos levam ao fim ultimo, que é \textsc{DEus}, nem se podem transgredir, sem que ao mesmo tempo se transgrida um dever para com Elle. Todavia, num sentido \emph{mais rigoroso}, \emph{deveres} para com \textsc{DEUS} são os que \emph{directa} e \emph{immediatamente} nos referem a \textsc{Deus}. E' d'estes deveres que nos occupamos.\par A \emph{Religião} toma-se em varios sentidos. --- No sentido proprio, significa o \emph{culto devido} a \textsc{DEus}. --- Em sentido lato, derivado do primeiro, significa a \emph{virtude moral}, que nos inclina a prestar a \textsc{DEus} o culto, que Lhe é devido. --- Em sentido mais lato, significa o \emph{conjuncto} --- tanto das \emph{verdades}, que exprimem as relações entre \textsc{DEUs} e o homem, --- como dos \emph{deveres}, que se fundam naquellas relações. --- 'l'omada no \emph{segundo} sentido, a \emph{Religião} é a \emph{parte principal} da ordem moral ou dos \emph{deveres} impostos pela lei natural. Tomada\pgn{676}no \emph{terceiro} sentido, a \emph{Religião} é o \emph{fundamento} de toda a ordem moral; porque Dzvus é o \emph{fim ultimo} da ordem moral, --- a sua \emph{Essencia} é a \emph{norma ultima}, de que a ordem moral tira a sua necessidade, --- a sua Vontade é a causa da \emph{obrigação} e da \emph{sancgão} da ordem moral, e a fonte de todos os direitos.\par Na virtude da \emph{Religião} devemos distinguir o objecto \emph{formal}, o objecto \emph{material} e \emph{directo}, \emph{e} o objecto \emph{indirecto}. --- O objecto \emph{formal} da Religião é a 1 \emph{honestidade}, \emph{que se encerra nos actos religiosos}; porque nada mais \emph{honesto} ou \emph{conveniente}, do que dar a \textsc{DEus} o culto, que Lhe é devido. O seu objecto \emph{material} e \emph{directo} é o \emph{culto}, que se presta a \textsc{Deus}. O seu objecto \emph{indirecto} é o proprio \textsc{Deus}, emquanto é o \emph{Primeiro Principio de todas as coisas}. --- D'onde se vê que a \emph{Religião} não é uma virtude \emph{theologal}, mas \emph{moral}; porque a virtude \emph{theologal} tem por objecto \emph{directo} o proprio \textsc{Deus}, tal é a virtude da Fé, da \emph{Esperança} e da \emph{Caridade}. --- Todavia a \emph{Religião}, emquanto nos inclina a prestar a \textsc{Deus} o culto, que lhe é. devido, comprehende todos os \emph{deveres} do homem para com \textsc{Deus} e por isso tambem os que se referem ás virtudes da Fé, \emph{Esperança} e \emph{Caridade}. Nesse caso, os actos d'estas tres virtudes não são actos \emph{elícitos} da propria \emph{Religião} (porque uma virtude \emph{moral} não póde produzir os actos proprios de uma virtude \emph{theologal}), mas são actos \emph{imperados} por ella, emquanto a Religião leva essas virtudes \emph{theologaes} a praticar os * actos respectivos. D'este modo os actos das virtudes \emph{theologaes} derivam \emph{immediatamente} das proprias virtudes; \emph{med?atamente}, da \emph{Religião}.\par Muitos auctores modernos, escrevendo tratados de \emph{Direito natural}, omittem os \emph{deveres} do homem para com \textsc{DEus}. As razões d'esta omissão são varias. --- Uns dizem que os \emph{deveres} do homem para com \textsc{Deus} são apenas \emph{moraes}, mas não são \emph{jurídicos}; porque \textsc{DEus} não póde \emph{coagir} as creaturas. --- Outros affirmam que esses \emph{deveres} são inuteis; porque \textsc{Deus} não precisa dos nossos obsequios. --- Outros finalmente sustentam que a creatura não póde nem deve ter deveres \emph{religiosos} para com \textsc{DEus}; porque a Religião é \emph{parte} da \emph{justiça}, a qual preceitua que a divida seja \emph{adequadamente} satisfeita, \emph{e} o homem não póde satisfazer \emph{adequadamente} as suas dividas para com \textsc{Deus},\par Todas estas razões são falsas ou inconcludentes. --- E' falsa a \emph{primeira} razão. Os \emph{deveres} do homem para com \textsc{Deus} são rigorosamente \emph{jurídicos} ; \textbar{} porque \textsc{Deus}, sendo o Creador e o Senhor de todas as coisas, tem verdadeiro \textbar{} \emph{direito} a exigir. a observancia da sua lei. Nem se diga com \emph{Ahirens} que 5 \textsc{Deus}, pelo facto de não poder \emph{coagír} as creaturas, não tem \emph{direitos}; porque a respondemos que \textsc{Deus}, embora não tivesse a força \emph{coactiva}, não deixaria de \textbar{}, ter \emph{direitos}, pois \& \emph{coacção}, como vimos, não constitue a essencia do \emph{direito}. --- E' falsa a \emph{segunda} razão. Porquanto, embora \textsc{Deus} não \emph{precise} dos nossos Obsequiod, nem por isso esses obsequios são inuteis, porque são exigidos pela \emph{justiça}, \emph{a} qual preceitua a observancia das relações entre \textsc{Deus} e o \textbar{} homem. Em poucas palavras: O culto, que prestamos a Deus, se não é util\pgn{677}e necessario para Elle, é util e necessario para o \textbf{homem}. --- Não é concludente a \emph{terceira} razão. Porquanto, embora o homem não possa satisfazer \emph{adequadamente}. as suas dividas para com \textsc{Deus}, nem por isso fica dispensado de dar o que póde, o que está ao alcance das suas forças. "\par E' evidente que o homem, com relação a \textsc{Deus}, sô tem \emph{deveres} e não \emph{direitos}; porque o homem é essencialmente servo de \textsc{Deus}, e o Creador nada póde dever ás suas creaturas (\emph{Of}. \emph{Sum}. \emph{Th}., 2-2, q. 81 e seg.).}.

120. \textbf{Divisão} dos \textbf{deveres} do homem \textbf{para} com \textbf{Deus}. --- O homem não póde prestar á \textsc{Deus} o \emph{culto}, que Lhe é devido, se não tiver um verdadeiro e conveniente \emph{conhecimento} da essencia e das infinitas perfeições do mesmo \textsc{Deus}. D'ahi resulta a divisão dos \emph{deveres} do homem para com \textsc{Deus} em \emph{theoricos} e \textbar{} \emph{praticos}. Os \emph{theoricos} referem-se ao \emph{conhecimento} de \textsc{Deus}, e , por isso são cumpridos directamente pela \emph{éntelligencia}; --- OS \pg{676} \pg{677}  \emph{praticos} referem-se ao \emph{culto} de \textsc{Deus}, e por isso são cumpridos directamente pela \emph{vontade}. --- Neste artigo tratamos dos deveres \emph{theoricos} do homem para com \textsc{Deus}; no seguinte trataremos dos deveres \emph{praticos} \footnote{São, pois, duas as partes da \emph{Religião}: uma \emph{theorica}, outra \emph{pratica}. A \emph{theorica} consiste no conhecimento, embora inadequado, da infinita excellencia de \textsc{DEus}; --- a \emph{pratica} consiste no \emph{culto}, que devemos prestar a \textsc{Deus}, pela sua infinita excellencia. --- Quem erra na parte \emph{theorica}, isto é, no \emph{conhecimento} de \textsc{Deus}, emquanto Lhe nega uma propriedade que possue, ou Lhe attribue uma propriedade que não possue, chama-se \emph{blasphemo}, --- e quem nega a \textsc{Deus} o \emph{culto}, que Lhe é devido, diz-se ?\emph{mpio}.}. \textbar{}

121. \textbf{Deveres theoricos} do \textbf{homem para com Deus}. --- Deveres \emph{theoricos} do homem para com \textsc{Deus} são, como dissemos, os que se referem ao \emph{conhecimento} de Dxzus, isto é, ás verdades ácerca da essencia, das perfeições e das obras de \textsc{Deus}. --- Estas verdades podem ser \emph{racionaes} ou \emph{supraracionaes}, conforme excedem, ou não, a capacidade da limitada razão humana. --- Os deveres \emph{theoricos} do homem para com \textsc{Deus} resumem-se na obrigação de adquirir um conveniente \emph{conhecimento} das verdades divinas \emph{racionaes}, e de acceitar as verdades \emph{supraracionaes}, se \textsc{Deus} as revelar por um modo sobrenatural- \footnote{\textsc{DEus} é unica, simplicissima, infinita Verdede, identificada com a sua Essencia ineffavel. --- A nossa intelligencia, sendo finita e por isso não podendo abranger com um só olhar o infinito, procura formar da unica e simplieissima Verdade muitos e differentes conceptos, cada um dos quaes reproduz, embora por um modo infinitamente inadequado, um ou outro attributo da Essencia divina.\par Esta maneira de conhecer a Essencia infinita de \textsc{DEus} não é arbitraria nem falsa; porque os diversos conceptos, que formamos da Essencia divina, tém o seu fundamento na realidade, isto ó, as diversas verdades, que formamos Acerca de Drs, existem effectivamente na Essencia divina, não inadequadas e accidentaes, como na intelligencia humana, mas perfeitissimas e unidas numa unica Verdade, substancial e infinita. --- Em \textsc{Deus}, pois, podem distinguir-se muitas e differentes verdades, distribuidas em duas classes. Uma classe abrange as verdades, que não excedem a capacidade ou comprehens\&o da intelligencia creada, como são existencia e a unidade de \textsc{DEUS}, a creação do mundo, a espiritualidade e immortalidade da alma humana, etc. ; --- a outra classe comprehende as verdades, que excedem a capacidade ou comprehensão de \emph{toda} a inteligencia creada, como são a Unidade e a Trindade de \textsc{Deus}, a Encarnação do Filho de \textsc{Deus}, etc. --- As verdades, que\pgn{678}excedem a capacidade ou comprehensão de \emph{toda} a intelligencia creada, chamam-se \emph{mysterios}. \emph{Mysterio}, pois, é uma proposição em que não podemos perceber a razão \emph{intrinseca}, que une o predicado ao sujeito.\par Os \emph{mysterios} são uma pedra de escandalo para os incredulos. E todavia nada mais logico, nada mais natural do que. a existencia dos \emph{mysterios}. A natureza sensivel está cheia de mysterios! A essencia da attracção, do som, da luz, do calor, da electricidade, não são outros tantos mysterios para a intelligencia creada? Sim, certamente. E, se a nossa intelligencia encontra \emph{muysterios} nas coisas \emph{accessiveis} aos nossos sentidos, não deverá ella encontral-os em \textsc{Deus}, que é invisivel aos nossos olhos corporeos e que, emquanto é objecto da mesma intelligencia creada, está infinitamente distante? --- Esta razão é sufficiente para provar a existencia dos \emph{musterios} em \textsc{Deus}. Mas, como o assumpto é importantissimo, accrescentamos duas outras razões, mais intrinsecas e convincentes.\par \textsc{Deus} é infinito, e a intelligencia creada é finita. Ora repugna que o \emph{finito} comprehenda e meça o \emph{infinito}, ou que o \emph{infinito} esteja contido no \emph{finito}. Logo devem existir verdades em \textsc{Deus}, que excedem a capacidade ou comprehensüo de toda a intelligencia creada. Logo existem os \emph{mysterios}. --- Com effeito, para os incredulos sustentarem que em \textsc{Deus} não existem \emph{mysterios}, isto é, verdades que excedem a capacidade de \emph{toda} a intelligencia creada, deveriam provar ou que \textsc{Deus} é finito, ou que a intelligencia creada é infinita. Não se concebe outro meio termo. Ora um \textsc{Deus} \emph{finito} é contradicção, assim como é contradicção uma \emph{coisa creada infinita}.\par Além d'isso, a intelligencia, para que possa conhecer todas as propriedades, que se referem a um objecto, deve comprehender perfeitamente a \emph{essencia} do mesmo; porque a essencia é a razão, o fundamento de tudo o que o objecto possue. Ora a \emph{Essencia} de \textsc{Deus} é incomprehensivel para a intelligencia creada; porque as creaturas, que são o meio por que nos elevamos ao Creador, sendo effeitos finitos e inadequados, não podem manifestar a infinita grandeza de \textsc{DEus}. Logo existem \emph{mysterios}, isto é, existem em \textsc{Deus} verdades, que excedem a capacidade de \emph{toda a} intelligencia creada.\par Os \emph{mysterios} excedem a nossa intelligencia limitada, mas não a destroem; pois tudo o que é \emph{superior} a uma faculdade, não lhe é por ei \emph{contrario}. --- Se as verdades \emph{supraracionaes} e as \emph{racionaes} fossem contrarias em nós, deveriam ser contrarias tambem em \textsc{Deus}. Ora isto repugna; porque estas duas especies de verdades, consideradas em \textsc{Deus}, não são distinctas, mas identificam-se com a simplicissima Verdade infinita. --- Só o erro é contrario á verdade; e ninguem provou nem póde provar que uma verdade \emph{suwpraracional} é contraria a uma verdade \emph{racional}; visto que, se isto se provasse, a verdade \emph{supraracional} \&ornar-se-hia \emph{racional}, pois a intelligencia comprehenderia a \emph{relação} ou o \emph{nexo} entre o sujeito e o predicado (\emph{Sum}. \emph{c}. \emph{Gentes}, 1, 1, c. 8).}.

\pg{678} 122. O \textbf{homem deve adquirir} um \textbf{conveniente conhecimento das verdades divinas racionaes}.

a) O homem, derivando de \textsc{Deus}, como do seu primeiro principio, e tendendo para \textsc{Deus}, como para o seu ultimo fim, deve unir-se-Lhe pelos actos da \emph{intelligencia} e da \emph{vontade}; porque um Ente espiritual e purissimo, como é \textsc{Deus}, só póde ser attingido por faculdades espirituaes, como são a intelligencia \pg{679}  e á vontade. Ora o homem une-se com \textsc{Deus} pelos actos da \emph{intelligencia}, emquanto adquire um conveniente \emph{conhecimento} das verdades divinas \emph{racionaes}, isto é, das verdades que estão ao alcance da nossa razão. Logo o homem deve adquirir um conveniente \emph{conhecimento} das verdades divinas \emph{racionaes}.

b) O homem deve, pelos actos da sua \emph{vontade}, prestar a \textsc{Deus} o \emph{culto}, que é \emph{devido} a tão grande Senhor. Ora é impossivel prestar a \textsc{Deus} o culto, que Lhe é devido, se não se adquirir um conveniente \emph{conhecimento} das verdades divinas \emph{racionaes}; pois ninguem póde honrar convenientemente uma pessoa, se não conhecer a sua dignidade. Logo o homem deve adquirir um: conveniente \emph{conhecimento} das verdades divinas \emph{racionaes}.

Dissemos --- \emph{conveniente} conhecimento; pois basta que o homem adquira o conhecimento das verdades divinas, necessario para o cumprimento dos deveres do seu estado \footnote{Examinemos duas objecções:\par Dizem: A parte \emph{teorica} da Religião, consistindo no conhecimento de um objecto tão elevado e abstracto, como é \textsc{Deus}, é de. diffieil execução, e por isso não póde ser prescripta pela lei \emph{natural}, que abrange todos os homens, doutos e ignorantes.\par Resposta. A parte \emph{theorica} de Religião, que é imposta a todos pela lei \emph{natural}, não consiste no conhecimento das verdades, que são mais difficeis e abstractas, relativas a \textsc{DEUS}, mas sim no conhecimento das verdades divinas, que são mais obvias e elementares, e que são necessarias e sufficientes para o homem ter uma idea exacta de \textsc{Deus} e para Lhe prestar o culto conveniente, devido a tão grande Senhor.\par Insistem: A parte \emph{theorica} da Religião, por se referir á \emph{intelligencia}, não póde estar comprehendida na lei \emph{natural}; que só obriga a \emph{vontade}.\par Resposta. A parte \emph{theorica} da Religião refere-se á \emph{intelligencia} e \& \emph{vontade}. Refere-se \& \emph{intelligencia}, emquanto o conhecimento de \textsc{Deus} é um acto \emph{elicito} ou \emph{proprio} d'essa faculdade; refere-se á \emph{vontade}, emquanto tal conhecimento é um acto \emph{imperado}, isto é, preceituado á propria \emph{intelligencia} pela \emph{vontade}. --- Se \& parte \emph{theorica} da Religião se refere á \emph{vontade}, é claro que póde estar comprehendida na lei \emph{natural}, isto é, póde ser materia de preceito. Por isso quem, por uma negligencia injustificada, deixa de adquirir um conveniente conhecimento de \textsc{Deus}, transgride a lei natural. --- D'onde se vê o erro dé \emph{Hobbes}, quando diz que os \emph{athews} não são injustos, porque peecam com a intelligencia, não com a vontade.}.

193.-0 homem \textbf{deve} ecceitar \textbf{as verdades supraracionaes}, \textbf{se Deus as revelar}.

a) O homem deve prestar a \textsc{Deus} o \emph{culto}, que é devido a tão grande Senhor. Ora não prestariamos a \textsc{Deus} o culto devido, se não acceitassemos as verdades \emph{supraracionaes}, que \textsc{Deus} revelasse por um modo sobrenatural; porque, nesse caso, semelhantes verdades seriam a norma --- não só do que deveriamos \pg{680}  conhecer, --- mas tambem do que deveriamos fazer para o cumprimento dos nossos deveres para com \textsc{Deus}, e por isso para a consecução do fim ultimo. Logo o homem deve acceitar as verdades \emph{supraracionaes}, se Dzvus as revelar.

b) O homem terá o dever de aeceitar as verdades \emph{supraracionaes}, que \textsc{Deus} revelar, se o mesmo \textsc{Deus} tiver o \emph{direito} de exigir a nossa fé. Ora \textsc{Deus} tem esse \emph{direito}; porque, sendo Elle a propria Verdade, não se engana nem póde enganar, e, sendo Senhor absoluto do homem e de todas as suas faculdades, póde impor-se á intelligencia humana. Logo o homem deve aeceitar as verdades \emph{supraracionaes}, se \textsc{Deus} as revelar \footnote{Como esta these é combatida pelos incredulos, digamos mais alguma coisa Acerca da \emph{Revelação sobrenatural}, seu \emph{objecto}, \emph{possibilidade}, \emph{necessidade} e \emph{existencia}, deixando para o fim a refutação das principaes difficuldades.\par a) Revelação sobrenatural. --- Revelação \emph{sobrenatural} é a manifestação de uma verdade, ou de um systema de verdades, feita por Dxzus ao ente finito e racional por um modo \emph{sobrenatural}, isto é, superior a todas as forças creadas. --- \emph{Cousin}, os \emph{racionalistas} da \emph{Allemanha}, \emph{Arhens} admittem uma revelação divina, como fundamento da sciencia ou da religião; mas essa \emph{revelação} é inteiramente diversa da revelação \emph{sobrenatural}, de que nos occupamos. Porquanto, --- para \emph{Cousin} e para os \emph{racionalistas} da \emph{Allemanha}, a \emph{revelação} é a espontanea evolução da intelligencia, ou a manifestação \emph{natural} de \textsc{Deus} atravez dos seres finitos, --- e, para Ahrens, a \emph{revelação} é a manifestecao da natureza divina no mundo e sobre tudo no espirito humano.\par (b) Objecto da Revelação sobrenatural. --- O \emph{objecto} da Revelação sobrenatural é constituido por todas as verdades, que \textsc{Deus} manifestou ou póde manifestar ás intelligencias creadas por um modo superior a todas as forças finitas. Estas verdades podem ser \emph{racionaes} e \emph{supraracionaes} (ainda que á Revelação \emph{sobrenatural}, tomada no sentido mais rigoroso, se refira exclusivamente ás verdades \emph{swupraracionaes}) --- Por isso, quando \textsc{Deus} manifesta \emph{sobrenaturalmente} uma verdade \emph{supraracional}, a Revelação é \emph{sobrenatural} emquanto ao \emph{modo} e emquanto ao \emph{objecto} revelado; --- e quando manifesta \emph{sobrenaturalmente} uma verdade \emph{racional}, \emph{a} Revelação só é \emph{sobrenatural} emquanto ao \emph{modo} de manifestar a verdade.\par o) Possibilidade da Revelação sobrenatural. --- Uma coisa é \emph{possivel}, quando não repugna. Mas a Revelação \emph{sobrenatural} não repugna. Porquanto, esta repugnancia dar-se-hia --- ou da parte de \textsc{Deus}, --- ou da parte da \emph{creatura} . \emph{racional}, --- ou da parte da \emph{verdade revelada}. --- Ora a repugnancia não se dá da parte de \textsc{DEus}; porque \textsc{Deus} é infinitamente sabio e conhece todas as verdades, é omnipotente e tem os meios aptos para as manifestar, é acto purissimo e desde a eternidade determinou a revelação de certas verdades. --- Não se dá da parte da \emph{creatura racional}; porque esta, se póde sujeitar-se ao magisterio de outra creatura, muito mais póde sujeitar-se ao magisterio de \textsc{Deus}. --- Não se dá da parte da \emph{verdade revelada}; porque a verdade, quer exceda quer não exceda a capacidade da intelligencia creada, é sempre o objecto \emph{adequado} da mesma intelligencia. Nem se diga que as verdades \emph{supraracionaes} são \emph{incomprehensiveis}; porque, se são \emph{éncomprehensiteis}\pgn{681}emquanto não \emph{vemos} o nexo entre o sujeito e o predicado, são \emph{comprehensiveis} emquanto percebemos o \emph{sentido}, que \emph{exprimem}.\par á) Necessidade da Revelação sobrenatural. --- A Revelação \emph{sobrenatural} póde dizer-se \emph{necessaria} --- ou emquanto ás verdades \emph{supraracionaes}, --- ou emquanto ás verdades \emph{racionaes}. D'ahi uma duplice proposição.\par a) A Revelação \emph{sobrenatural é absolutamente} necessaria na \emph{hypothese} da elevação da creatura racional a um fim \emph{sobrenatural}. Com effeito, nessa hypothese, o fim da creatura racional, sendo superior ás exigencias da /sua natureza, deve manifestar-sc-lhe por um modo \emph{sobrenatural}; aliás o homem não poderá dirigir para elle os seus desejos e as suas obras; pois ninguem tende para o desconhecido. Se é \emph{absolutamente} necessaria a Revelação do fim \emph{sobrenatural}, é tambem \emph{absolutamente} necessaria a Revelação dos \emph{meios} proporcionados para esse fim; não só porque a determinação dos meios depende da causa, que determinou o fim, quando este é \emph{sobrenatural}, --- mas tambem porque os meios, proporcionados para a consecução de um fim \emph{sobrenatural}, devem ser \emph{sobrenaturaes} e por isso só dependem da Sabedoria e da Vontade de \textsc{Deus}. --- Dissemos que a Revelação, neste caso, é \emph{absolutamente} necessaria; porque a intelligencia creada não póde por si conhecer, pelo menos com \emph{certeza}, a existencia do fim \emph{sobrenatural}.\par \emph{b}) A Revelação \emph{sobrenatural é moralmente} necessaria emquanto-ás verdades \emph{racionaes}, relativas á Religião e á Moral. Na verdade, sem a Revelação, só poucos homens e depois de muito tempo teriam apprehendido algumas verdades religiosas e moraes, misturadas com muitos erros; quando é certo que \emph{todos} os homens, desde o principio da sua existencia, devem conhecer, sem mistura de erro, essas verdades; pois é d'ellas que depende a felicidade do homem. --- Dizemos que a Revelação \emph{sobrenatural} das verdades \emph{racionaes} é sô \emph{moralmente} necessaria, porque a intelligencia creada, embora, considerada em si, possa chegar ao conhecimento de todas as verdades \emph{religiosas} e moraes da ordem \emph{racional}, comtudo, considerada no homem, no qual ella existe cercada de difficuldades, é \emph{moralmente} (não \emph{physicamente}) impotente para conhecer todas essas verdades. A esta impotencia \emph{moral} corresponde a necessidade \emph{moral} da Revelação (\emph{Sum}. \emph{c}. \emph{Gentes}, l. I, c. 4).\par e) Existencia da Revelação sobrenatural. --- A Revelação \emph{sobrenatural} não só é \emph{possivel} e \emph{necessaria}, \emph{mas tambem existe}, \emph{e} á sua depositaria authentica é a Egreja Catholica. --- Não pertence á Philosophia a demonstração d'esta grande e consoladora verdade. Limitamo-nos a dizer que os \emph{motivos} de \emph{credibilidade}, que demonstram tal existencia, são tantos e tão valiosos, que um espirito sensato não póde deixar de dizer a \textsc{DEUS} com um celebre pensador: u \emph{Senhor}, \emph{se} a \emph{nossa} fé é \emph{falsa}, \emph{fomos enganados por Vós}; \emph{porque ella foi authenticada por milagres}, \emph{que} só Vós \emph{podeis operar} » (Ricardo de S. Victor, \emph{De Trinitate}, lib. I, c. 2).\par f) Difficuldades dos incredulos. --- Examinemos as principaes difficuldades dos incredulos contra a Revelação \emph{sobrenatural}.\par Dizem: A Revelação \emph{sobrenatural} repugna da parte de \textsc{Deus}, da parte do homem, e da parte das verdades reveladas. --- Repugna da parte de Dus; porque \textsc{Deus} não póde nem deve communicar ás creaturas os seus segredos. --- Repugna da parte do homem; porque a Revelação humilha e escraviza\pgn{682}a nossa razão. --- Repugna da parte das verdades reveladas; porque ellas, não tendo \emph{evidencia}, deixam de ser \emph{verdades}.\par Resposta. A Revelação \emph{sobrenatural} não repugna da parte de \textsc{DEus}, nem da parte do homem, nem da parte das verdades reveladas. --- Não repugna da parte de \textsc{Deus}; porque \textsc{Deus}, que concedeu ao homem a faculdade de communicar aos outros os proprios pensamentos, não póde deixar de ter o poder de communicar ás creaturas os seus segredos, e porque essa communicação, longe de abater a Grandeza e a Magestade de \textsc{Deus}, revela brilhantemente a sua infinita Bondade e Sabedoria. --- Não repugna da parte do \emph{homem}; porque, se é proprio do erro humilhar e escravizar á nossa razão, é proprio da verdade ennobrecel-a e tornal-a livre. --- Não repugna da parte das \emph{verdades reveladas}; porque muitas, isto é, as \emph{racionaes}, têm uma evidencia \emph{intrinseca}, \emph{e} muitas outras, isto é, as \emph{supraracionaes}, têm uma evidencia \emph{extrinseca}, porque se baséam na infallivel Verdade, que é \textsc{DEUs}, e por isso não podem deixar de ser verdades. Notamos que as verdades \emph{supraracionaes}, se tem uma evidencia \emph{extrinseca} relativamente á intelligencia creada, têm uma evidencia \emph{intrinseca} relativamente á intelligencia divina.\par Continuam: A Revelação \emph{sobrenatural} não é \emph{necessaria} relativamente ás verdades \emph{supraracionaes}, nem relativamente ás verdades \emph{racionaes}. Não é necessaria relativamente ás verdades \emph{supraracionaes}; porque essa necessidade basear-se-hia na existencia de um fim \emph{sobrenatural}, \emph{e} repugna que um ente seja destinado para um fim superior ás suas tendencias naturaes. Não é necessaria relativamente ás verdades \emph{racionaces}; porque essa necessidade arguiria imperfeição e deficiencia na creatura mais elevada entre todos os seres, que é o homem.\par Resposta. A Revelação \emph{sobrenatural} é necessaria em relação ás verdades \emph{supraracionaes}, e em relação ás verdades \emph{racionaes}. --- E' necessaria em relação ás verdades \emph{supraracionaes}, se \textsc{Deus} tiver destinado o homem para um fim \emph{sobrenatural}; e esta destinação não repugna, porque tal fim, embora \emph{superior} ás exigencias da hossa natureza, comtudo não lhe é \emph{contrario}, mas ennobrece-a, elevando-a de um modo maravilhoso. --- E' necessaria, embora \emph{moralmente}, em relação ás verdades \emph{racionaes}; porque, sem ella, só \emph{poucos} homens, e depois de \emph{muitos} annos, teriam chegado a conhecer uma ou outra verdade religiosa, misturada com muitos erros. Esta verdade foi reconhecida pelos proprios philosophos pagãos. \emph{Pithagoras} dizia: « O homem só deve fazer o que é do agrado de \textsc{DEus}, mas elle não poderá conhecer o que é do agrado de \textsc{Deus}, se o proprio \textsc{DEUs}, ou um genio superior, não ]h'o ensinar, e se uma luz divina não o esclarecer » (\emph{Janblich}. \emph{in vit}. \emph{Pithag}.). \emph{Socrates} eerescentaya: « A pietade é uma virtude summamente desejavel; mas ninguem, á excepção de \textsc{Deus}, póde ensinal-a.» (\emph{Plato}, \emph{Politic}.)\par Replicam: O homem não pôde acreditar numa verdade, que se diz \emph{revelada}, se não souber com \emph{certeza} de que \textsc{Deus} fallou. Mas, quando o homem sabe \emph{com certeza} de que \textsc{Deus} fallou, já não tem fé, mas \emph{sciencia}; porque a fé exclue toda a \emph{certeza}.\par Resposta. Concordamos em que o homem não póde acreditar numa verdade, que se diz \emph{revelada}, se não souber com \emph{certeza} de que \textsc{Deus} a revelou; mas negamos que o homem, quando sabe \emph{com certeza} de que \textsc{Deus} fallou, já\pgn{683}não tem fé mas \emph{sciencia}. Por quanto, esta certeza, sendo \emph{extrinseca}, isto é, baseada em argumentos \emph{extrinsecos} á propria Verdade, que \textsc{Deus} manifestou, é compativel com a fé, a qual só exclue a \emph{certeza intrinseca}, isto é, baseada em argumentos \emph{intrinsecos} á mesma verdade. --- D'ahi se segue que o acto de fé, por se não basear em argumentos \emph{intrinsecos}, não é \emph{necessario}, mas \emph{livre}, e por isso é \emph{meritorio}.}.

194. \textbf{Corollario}. --- Das precedentes considerações resulta que a \emph{autonomia} ou \emph{liberdade} de \emph{consciencia}, que é o chamado direito \pg{681} \pg{682}  do \emph{livre pensamento} em materia de Religião, e o \emph{indifferentismo religioso}, que É a opinião dos que julgam todas as \pg{683}  religiões egualmente boas e agradaveis a \textsc{Deus} e por isso capazes de levar o homem á consecução do seu fim ultimo, são erros injuriosos para \textsc{Deus} e funestos para o homem \footnote{Digamos alguma coisa ácerca da \emph{autonomia} ou \emph{liberdade} de \emph{consciencia} e ácerca do \emph{éndifferentismo religioso}.\par a) Liberdade de consciencia. --- \emph{Liberdade de consciencia} é o supposto direito, attribuido ao homem pelos \emph{racionalistas}, de abraçar e professar a Religião, que elle, só guiado pela luz da razão, reconheça por verdadeira.\par Esta theoria é falsa pelas seguintes razões: --- 1.º) Suppõe que a Religião, que o homem deve professar, é a \emph{natural}. Ora isto é falso. A Religião meramente \emph{natural} nunca. existiu; porque \textsc{Deus}, desde o principio do mundo, fallou ao. homem, revelando-lhe muitas verdades, \emph{theoricas} e \emph{praticas}, \emph{racionaes} e \emph{supraracionaes}. O philosopho não póde prescindir d'esta revelação. --- 2.º) \textsc{Deus}, sendo a Verdade infinita, tem todo o direito de se impor á intelligencia humana, que foi creada para a verdade. Ora este direito seria destituido de todo o valor, se o homem podesse rejeitar a Religião revelada por \textsc{DEUs}, para seguir a que elle quizesse.\par Objectam: A \emph{liberdade de pensamento} é um dos direitos da nossa razão. Logo homem póde, mesmo em materia de Religião, pensar como quizer.\par Resposta. Se por \emph{liberdade do pensamento} entendem que a nossa razão nao póde ser constrangida por nenhuma \emph{força extrinseca} a adherir a uma ou outra proposição, concedemos que esta \emph{liberdade} é um direito da razão. Se por \emph{liberdade do pensamento} entendem que a razão é independente do influxo da vontade, e por isso do influxo da lei moral, negamos que essa \emph{liberdade} é um direito da nossa razão, porque a lei moral regula a nossa vontade e, por meio d'esta, todas as outras faculdades do homem.\par 3) Indifferentismo religioso. --- \emph{Indifferentismo religioso} é á opinião dos que proclamam boas e egualmente agradaveis a \textsc{Deus} todas as religiões, sustentando que todas ellas levam o homem ao seu fim ultimo.\par Esta opinião é absurda pelas seguintes razões: --- 1.*) Se em toda e qualquer Religião o homem podesse conseguir o. seu fim ultimo, \textsc{Deus} teria egual amor á verdade e ao erro, á virtude e ao vicio; pois Elle daria a mesma recompensa tanto aos que seguem uma Religião de erros e de vicios, como aos que seguem uma Religião cheia de virtude e de verdade. Ora esta consequencia repugna. Logo o homem só póde salvar-se na verdadeira Religião, que é a revelada por \textsc{DEus}. --- 2.3) Para agradarmos a \textsc{Deus} e para alcancarmos a vida eterna, não basta honrar e venerar a \textsc{Deus} com um culto, que agrada a nós, mas é necessario tributar-lhe o culto e a veneração, que agrada a Elle; porque a determinação da \emph{qualidade} do serviço, e do «\emph{odo} por que se deve prestar, pertence ao senhor e não ao servo. Ora a \textsc{Deus}, infinitamente santo, não póde agradar todo e qualquer culto, mas só o que é inteiramente santo e o que foi determinado por Elle.\par Objectam: Se não se admittir o \emph{ndifferentismo religioso}, devemos dizer\pgn{684}que a maxima parte dos homens estão fóra do caminho da salvação eterna. Ora isto é contrario á Providencia divina.\par Resposta. Estão fóra do caminho da salvação eterna só os que \emph{não querem} reconhecer a verdadeira Religião, \emph{desprezando os auxílios}, que para isso recebem de \textsc{Deus}. --- Nem se diga que é difficil distinguir, entre tantas, a Religião verdadeira; porque a belleza e a harmonia dos dogmas, a santidade e a pureza da vida, a gloria dos milagres, a fortaleza dos martyres, --- todo este admiravel conjuncto de perfeições, que só se encontram na Religião Catholica, fundada por N.S. Jesus \emph{OunisTo}, é motivo mais que sufficiente para uma intelligencia serena e para um \emph{coração} bem feito \emph{abraçar}, amar e praticar esta augusta Religião.}.

%%>>> art-61

\section*{Artigo II.\\ Deveres praticos do homem para com Heus.}
\addcontentsline{toc}{section}{Artigo II. --- Deveres praticos do homem para com Heus.}

\pg{684} 125. \textbf{Deveres} praticos do homem para com \textbf{Deus}. --- Deveres \emph{praticos} do. homem para com \textsc{Deus} são, como dissemos, os que se referem ao \emph{culto}, que é devido a \textsc{Deus}. --- \emph{Culto} é o \emph{reconhecimento da excellencia de} \textsc{Deus} e \emph{da nossa sujeição para com Elle}. Divide-se em \emph{interno} e \emph{externo}, conforme se presta no intimo da nossa alma pelos actos da \emph{vontade}, ou se manifesta exteriormente por signaes \emph{sensiveis}. O \emph{culto} póde ser determinado pela \emph{natureza} ou pela \emph{Revelação}. --- O homem deve prestar . á \textsc{Deus} o culto, \emph{interno} e \emph{externo}, não só o que é determinado pela \emph{natureza}, mas tambem o que o póde ser pela \emph{Revelação} \footnote{O culto devido a \textsc{DEus} em virtude da sua infinita excellencia, isto é, emquanto \textsc{Deus} é o principio de todo o bem e o soberano Senhor de todas as coisas, chama-se \emph{latría}. Por isso o culto de \emph{latria} compete exclusivamente a \textsc{Deus}. --- Este culto diz-se \emph{devido} a \textsc{Deus}, não só porque \textsc{Deus} é o Ser dotado de infinita excellencia, mas tambem porque tudo o que o homem póde prestar a \textsc{Deus}, é devido.}.

126. 0 homem deve prestar a \textbf{Deus} o culto interno.

a) O homem, derivando de \textsc{Deus}, como do seu primeiro principio, e, tendendo para \textsc{Deus}, como para o seu ultimo fim, deve unir-se-lhe pelos actos da \emph{ntelligencia} e da \emph{vontade}; porque um Ente espiritual, como é \textsc{Deus}, só é attingido por faculdades espirituaes, como são a intelligencia e a vontade. Ora o homem une-se com \textsc{Deus} pelos actos da \emph{vontade}, emquanto Lhe presta o culto \emph{interno}. Logo o homem deve prestar a \textsc{Deus} o culto \emph{interno}.

d) O homem deve, pelos actos livres da \emph{vontade}, observar e reproduzir em si o que a ordem moral exige. Ora a ordem \pg{685}  moral exige que --- o servo preste reverencia e obsequio ao senhor, --- o beneficiado seja reconhecido ao seu bemfeitor, --- o ser finito reconheça a excellencia do Ser infinito; e \textsc{Deus} é o senhor e bemfeitor do homem, é o Ser dotado de excellencia infinita. Logo o homem deve, pelos actos da vontade, prestar a \textsc{Deus} o culto devido. Os actos da vontade são \emph{internos}. Logo o homem deve prestar a \textsc{Deus} o culto \emph{interno} \footnote{Os \emph{atheus}, os \emph{materialistas} e os \emph{pantheistas}, que ou negam a \textsc{Deus}, ou querem confundil-o com a universalidade das coisas, não reconhecem no homem, como não podia deixar de ser, o \emph{dever} de prestar a \textsc{Deus} um culto, quer \emph{interno} quer \emph{externo}. --- Mas este erro é uma excepção, pois não houve, nem ha povo algum, ainda que ignorante e barbaro, que não tivesse ou não tenha prestado algum culto \& Divindade.}.

127. \textbf{Principaes actos} do culto interno. --- Os principaes actos do culto \emph{interno} são os seguintes: a fé, o \emph{amor}, \emph{á esperança}, a \emph{oração}, o \emph{agradecimento}, \emph{á obediencia}, a \emph{satisfação}.

a) Fé, --- Fé é a virtude, pela qual, baseando-nos na infinita verdade e sinceridade de \textsc{Deus}, acreditamos em tudo o que Elle revelou. --- A fé é imperada pela vontade, e por.isso é meritoria.

b) Amor, --- Amor é a virtude pela qual, baseando-nos na infinita bondade e perfeição de \textsc{Deus}, O estimamos sobre todas as coisas. --- O \emph{amor} é ou ha de ser o principio de todos os actos do \emph{culto}, interno e externo, que o homem deve prestar ao seu \textsc{Deus} \footnote{O homem deve \emph{amar} a \textsc{Deus} sobre todas as coisas. Porquanto, a recta ordem exige que amemos os bens em proporção da sua dignidade e da estima, que merecem; pois o bem é por si amavel. Ora \textsc{Deus} é o Bem infinito, fonte e fim ultimo de todos os bens creados, e por isso infinitamente amavel. Logo a recta ordem exige que amemos a \textsc{Deus} sobre todas as coisas, e todas as coisas por \textsc{DEus}. --- Não é preciso que este amor seja summo na \emph{intensidade}; basta que o seja na \emph{apreciação}. --- Note-se que o Bem infinito póde considerar-se \emph{em si} e \emph{em relação a nós}. Considerado \emph{em si}, o Bem infinito é objecto do \emph{amor} de \emph{benevolencia}; --- considerado em \emph{relação a nós}, é objecto do amor de \emph{concupiscencia}.\par \emph{Kant} nega que o homem \emph{possa amar} a \textsc{Deus}, pois sustenta que \textsc{Deus} é invisivel, \emph{e} é um dos entes \emph{transcendentaes}, que estão fóra de toda a esphera da nossa investigação.\par O erro do philosopho allemão é evidente. O amor consiste na tendencia da vontade para o bem conhecido; de modo que todo o bem, que é ou póde ser conhecido, é ou póde ser um objecto capaz de attrahir o amor da vontade. Ora \textsc{DEus} é um Bem infinito, que não só póde ser, mas que é effectivamente conhecido pela intelligencia humana. Logo o homem \emph{póde} amar a \textsc{Deus}. --- Este amor do homem para com \textsc{Deus}, se não é de \emph{comprehensão}, porque \textsc{Deus} não póde ser amado pela creatura tanto, quanto \textsc{Deus} é amavel, --- é todavia de \emph{adhesão}, emquanto a nossa vontade \emph{adhere} a \textsc{Deus}.\par As razões de \emph{Kant} em abono da sua theoria são falsas. --- E' falso que\pgn{686}\textsc{Deus} é \emph{invisivel}; porque \textsc{Deus}, se é invisivel em si e aos olhos do nosso corpo, não é invisivel nas suas obras e aos olhos da nossa alma, isto é, \& nossa intelligencia. --- E' falso que \textsc{Deus} está fóra da esphera da nossa investigação; porque a nossa razão, baseada na experiencia, descobre e \emph{demonstra} muitas verdades ácerca de \textsc{Deus} e dos seus attributos.}.

\pg{686} c) \textbf{Esperança}. --- \emph{Esperança} é a virtude pela qual, baseando-nos na infinita liberalidade e fidelidade de \textsc{Deus}, confiamos em que Elle nos dará a felicidade eterna e os meios necessarios para a sua consecução \footnote{O homem deve \emph{esperar} em \textsc{Deus}. Na verdade, o homem, devendo tender para a felicidade eterna e empregar os meios proporcionados, e não podendo, só coni as suas forças, realizar essa tendencia e esse uso, deve ter confiança em \textsc{Deus}, que não sabe negar as suas graças a quem recorre 4 Sua infinita misericordia. --- O objecto \emph{formal}, pois, da \emph{esperança} é a infinita \emph{liberalidade} e \emph{fidelidade} de \textsc{Deus}; o objecto \emph{final} é a \emph{felicidade eterna}, a qual consiste na posse do proprio \textsc{Deus}.}.

d) Oração. --- \emph{Oração} é o acto pelo qual a nossa alma se eleva a \textsc{Deus}, e Lhe pede as graças convenientes \footnote{O homem é obrigado, pelo proprio \emph{Direito natural}, \& \emph{oração}. As trevas que encobrem a nossa intelligencia, as inclinações desordenadas da vontade, a vehemencia das paixões, mil perigos e contrariedades, tornam muito diffieil se não impossivel, a pratica do bem \emph{Honesto}. Sendo assim, e devendo alcançar a felicidade eterna, \emph{o} homem não póde deixar de recorrer á fonte inexhaurivel: de todas as graças --- a \textsc{Deus}. --- Este recurso a \textsc{Deus} é tão natural, que, no meio das angustias, sahe-nos da alma espontaneamente, e é causa de allivio e de consolação. --- Todos os povos, ainda os mais selvagens, dirigem as suas orações á Divindade, da qual esperam todos os bens.\par Os \emph{d\textasciicircum{}istas} e alguns \emph{racionalistas}, que não acreditam na \emph{Providencia divina} e na \emph{Revelação sobrenatural}, negam a efficacia da \emph{oração}, dizendo que \textsc{Deus} não se importa com as nossas supplicas, o que estas são inuteis, porque ás leis da natureza são absolutamente necessarias e immutaveis.\par Esta negação é absurda. --- Antes de tudo é falso que \textsc{Deus} se não importa com as nossas supplieas. \textsc{DEus} tem um cuidado particular das creaturas racionaes, e, tendo, pela voz da razão naturale pela voz da Revelação, mandado ao homem de recorrer a Elle por meio da \emph{oração}, não póde deixar de attender as nossas supplicas. --- Em segundo logar, é falso que as leis immutaveis e necessarias da natureza tornem inteiramente inuteis as nossas supplicas. Porquanto, as leis da natureza não são absolutamente immutaveis e necessarias, e \textsc{DEUs}, que as fez, póde, em attenção ás nossas supplicas, suspendel-as. Mas, embora as leis da natureza fossem immutaveis e necessarias, nem por isso seria inutil a \emph{oração}; não só porque dentro da actual ordem das coisas podem dar-se diversas mudanças, sem que se alterem as leis geraes da natureza, --- mas tambem porque, tendo \textsc{Deus} previsto desde a eternidade as orações dos homens, poude estabelecer leis, cujos effeitos correspondessem aos desejos dos supplicantes. --- O celebre \emph{Leibnitz} dizia;\pgn{687}« Não só os cuidados e os trabalhos são uteis ao homem, mas tambem as \emph{orações} ; porque \textsc{DEus} as viu desde a eternidade, antes de ordenar a existencia e o movimento dos seres. Por isso os que não fazem as diligencias necessarias para pôr em ordem as suas coisas, allegando o vão pretexto da necessidade dos acontecimentos, incorrem no sophisma conhecido pelos antigos sob o nome de \emph{sophisma do preguiçoso} » (\emph{Theod}., t. IT, pag. 416).}.

e) \textbf{Agradecimento}. --- \emph{Agradecimento} é o acto pelo qual o homem traduz a sua gratidão para com \textsc{Deus} por tantos beneficios, recebidos da sua infinita Bondade. i

\pg{687} f) Obediencia. --- \emph{Obediencia} é o acto pelo qual a vontade creada.se sujeita aos preceitos, impostos pela vontade de \textsc{Deus}, Creador e Senhor de todas as creaturas \footnote{Não é menos evidente o \emph{dever} que o homem tem de \emph{obedecer} a \textsc{Deus}; não só porque o homem é servo e \textsc{DEus} é o Soberano Senhor, mas tambem porque a vontade do homem, sendo fallivel, deve estar sujeita á vontade de \textsc{Deus}, que é infallivel e modelo de toda a santidade. --- Esta obediencia extende-se a todos os preceitos divinos, \emph{naturaes} e \emph{positivos}.}.

£) Satisfação. --- \emph{Satisfação} é o acto pelo qual a nossa vontade repara, pelo arrependimento e pelo proposito, as offensas ' feitas á infinita Magestade e Bondade de \textsc{Deus} \footnote{O homem, que transgrediu os preceitos de \textsc{Deus}, fica immediatamente privado da amizade de \textsc{DEus} e condemnado ao castigo proporcionado ás suas culpas. Se quizer sahir d'esse estado, deve reconciliar-se com \textsc{Deus} por meio do \emph{arrependimento}. O \emph{arrependimento}, pois, é uma dor intima da alma, pela qual detestamos as offensas feitas a \textsc{DEus}, e promettemos a emenda. --- O \emph{arrependimento} póde ser \emph{perfeito} ou \emph{imperfeito}, conforme nasce do \emph{amor de Deus} ou de \emph{temor do castigo}.\par Dissemos que os actos enumerados do culto \emph{interno} são os \emph{principaes}, porque, além d'esses, ha muitos outros, como o \emph{temor filial}, a \emph{gratidão}, etc. --- "Todos estes actos, como é evidente, derivam da \emph{vontade} por um modo \emph{elicito} ou \emph{iémperado}, isto é, \emph{immediata} ou \emph{mediatamente}.}.

128. O homem \textbf{deve} prestar a \textbf{Deus} o culto externo.

a) O homem deve prestar a \textsc{Deus} o culto devido, não com a alma só, mas com todo o \emph{composto}; pois todo o \emph{composto}, derivando de \textsc{Deus}, está sujeito a \textsc{Deus}, e por isso tem o dever de exprimir e attestar a sua sujeição para com \textsc{Deus}. Ora o \emph{culto}, prestado a \textsc{Deus} pelo \emph{composto humano}, não póde deixar de ser tambem \emph{externo}; porque deriva do \emph{corpo} informado pela \emph{alma}. Logo o homem deve prestar a \textsc{Deus} o culto \emph{externo}.

b) O homem deverá prestar a \textsc{Deus} o culto \emph{externo}, se este fór necessario para excitar e completar o \emph{interno}; pois quem tem a obrigação de alcançar um fim, tem tambem a obrigação de empregar os meios convenientes. Ora o culto \emph{externo} é necessario para \emph{excitar} e \emph{completar} o \emph{interno}; porque é proprio do homem elevar-se ao espiritual pelo sensivel, e traduzir por actos \emph{externos} o intimo affecto da alma, quando é sincero. Logo o homem deve prestar a \textsc{Deus} o culto \emph{externo} \footnote{Todos os povos, em todos os tempos, tem prestado a \textsc{DEus} o culto\pgn{688}\emph{externo}, publico e particular. Este facto, tão constante e universal, não póde ser attribuito a causas particulares, mas deve basear-se na propria natureza racional, a qual, sendo uma e a mesma em todos os tempos e em todos os logares, impõe a todos os homens os mesmos preceitos.\par \emph{Kant}, \emph{Ahrens} rejeitam o culto \emph{externo}, porque não prescripto pela \emph{lei natural}, e porque \textsc{Deus}, purissimo espirito, não precisa de culto \emph{externo}.\par Respondemos que as razões, adduzidas por estes escriptores, são todas falsas. --- E' falso que o culto \emph{externo} não é prescripto pela \emph{lei natural}. Porquanto, a lei \emph{natural}, embora por si não determine a \emph{especie} do culto \emph{externo}, todavia prescreve que o culto \emph{inferno} seja acompanhado pelo \emph{externo}. As razões d'esta prescripgão encontram-se nos \emph{direitos} de \textsc{DEus} sobre o homem (emquanto \textsc{Deus}, tendo dado ao homem \emph{alma} e \emph{corpo}, exige que o homem o sirva com a totalidade da \emph{essencia}) \emph{e} na propria \emph{natureza humana} (emquanto o homem traduz naturalmente o \emph{espiritual} no \emph{sensivel}, e d'este sobe para aquelle. --- E' falso que o culto \emph{externo} repugna á purissima e espiritual Essencia de \textsc{Deus}; porque o culto \emph{externo}, sendo informado pelo \emph{interno}, é o acto \emph{proprio} e \emph{natural} do homem, e, como tal não póde repugnar á Essencia purissima de \textsc{DEus}. O preceito, que temos de adorar a \textsc{Deus} em \emph{espirito}, não quer dizer que o homem deve adorar a \textsc{Deus} exclusivamente com a alma, mas quer dizer que a adoração \emph{externa} derive do \emph{espirito}, isto é, que não seja \emph{Aypocrita} mas \emph{sincera}. --- E' falso que o culto \emph{externo} deve rejeitar-se, porque \textsc{DEus} não precisa d'elle. O homem deve, prestar a \textsc{Deus} o culto \emph{externo}, não porque \textsc{Deus} precisa d'elle (pois, se assim fosse, nem Lhe deveriamos prestar o culto \emph{interno}), mas porque a \emph{justiça} exige que o homem, que recebeu de \textsc{Deus} \emph{a/ma.e corpo}, sirva ao seu Creador com actos provenientes de ambos estes elementos.}.

\pg{688} 129. \textbf{Principaes actos} do culto externo. --- Os principaes actos do culto \emph{externo} são a \emph{adoração externa}, o \emph{louvor}, O \emph{juramento}, a \emph{oblação}, O \emph{sacrificio}.

a) Adoração externa, --- Adoração \emph{externa} é o acto pelo qual nos prostramos deante de \textsc{Deus}, para d'este modo attestarmos a sua excellencia e o seu dominio sobre todas as creaturas, e a nossa inferioridade e submissão para com Elle.

b) \textbf{Louvor}. --- \emph{Louvor} é o acto pelo qual celebramos, com a lingua, as grandezas de \textsc{Deus} e as suas infinitas misericordias.

c) Juramento. --- \emph{Juramento} é o acto pelo qual invocamos o Nome Santo de \textsc{Deus} em confirmação de um nosso dieto ou de uma nossa promessa \footnote{O \emph{juramento} é a invocação do Santo Nome de \textsc{Deus} em testumunho, ou para dar credito ao que se diz, ou para confirmar alguma promessa. --- O \emph{juramento}, invocando a \textsc{DEus} como testemunha da verdade, deve ser feito em \emph{abono da propria verdade}, por causa \emph{grave}, e com summa. \emph{reverencia}.}.

d) Oblagão. --- \emph{Oblação} é o acto pelo qual offerecemos a \textsc{Deus} alguma coisa em testemunho da sua excellencia e soberania.

e) Sacrificio. --- \emph{Sacrificio} é o acto pelo qual se destroe, physica ou moralmente, uma coisa offerecida a \textsc{Deus}, em testemunho \pg{689}  do dominio soberano e absoluto, que Elle tem sobre todas as creaturas \footnote{O \emph{sacrificio} é a expressão mais naturale legitima do supremo dominio de \textsc{Deus} sobre o homem, e da nossa sujeição para com Elle. Porquanto, o homem deve offerecer a \textsc{DEus} a vida, que d'Elle recebeu. Mas, como isto não convém, escolhemos um animal ou uma outra coisa, que nos represente, e, offerecendo e destruindo, moral ou physicamente, esse animal ou essa coisa em honra de \textsc{Deus}, entendemos offerecer e destruir a nossa propria vida em testemunho do soberano dominio de \textsc{Deus} sobre a nossa existencia. --- Na hypothese da humanidade decahida, o \emph{sacrificio} offerece-se tambem em expiação do peccado, para aplacar a Justiça divina.\par Aos actos do culto \emph{externo}, enumerados no texto, alguns accrescentam o \emph{voto}, que é a promessa deliberada de um bem maior, feita a \textsc{Deus}. Mas o \emph{voto} póde pertencer tambem ao culto \emph{interno}; porque póde ser feito no intimo da alma, sem que transpareca exteriormente.}.

130. O homem \textbf{deve} prestar a \textbf{Deus} não só o culto que foi \textbf{imposto pela} natureza, mas tambem o que póde \textbf{ser determinado pela Revelação}. --- O homem deve prestar a \textsc{Deus} o culto, interno e externo, que \textsc{Deus} determina; porque a determinação do serviço não pertence ao subdito mas ao superior. Ora \textsc{Deus} póde determinar o culto, que o homem Lhe deve prestar, ou pela luz natural da razão, ou pela luz sobrenatural da Revelação. Logo o homem deve prestar a \textsc{Deus} não só o culto, interno e externo, que foi determinado pela natureza, mas tambem o que póde ser determinado pela Revelação \footnote{A verdade d'esta proposição é evidente. Além dos actos do culto interno e externo, prescriptos pela \emph{lei natural}, póde haver outros actos, tambem internos e externos, preseriptos pela Revelação divina, e o homem deve sujeitar-se a essa prescripcüo; porque, repetimos, o homem é servo de \textsc{Deus}, e não pertence ao servo a escolha da \emph{qualidade} ou do \emph{modo} do servico, que deve prestar ao seu senhor.\par Depois de termos tratado da virtude da \emph{Religião}, resta-nos dizer alguma coisa ácerca dos \emph{vicios}, que lhe são oppostos. Estes vicios são \emph{seis}, a saber: a \emph{indifferença religiosa}, a \emph{incredulidade}, \emph{á impiedade}, o \emph{atheismo}, a. \emph{superstição} e o \emph{fanatismo}. Os primeiros quatro vicios oppõem-se á Religião por \emph{defeito}; os ultimos dois, por \emph{excesso}.\par a) Indifferença religiosa. --- A \emph{indifferença religiosa} é o vicio dos que dizem, com as \emph{palavras} ou com os \emph{factos}, que \textsc{DeUS} não recebe as nossas homenagens, porque não precisa d'ellas, e que o homem póde ser bom e honesto, embora não seja religioso. --- Estas affirmações são falsas. E' falso, como demonstrámos, que \textsc{Deus} não recebe as nossas homenagens, pelo facto de não precisar d'ellas. E' falso tambem que o homem póde ser bom e honesto, sem ser religioso; porque quem falta aos seus \emph{deveres} para com \textsc{Deus}, é \emph{injusto} e \emph{ingrato}, \emph{e} uma pessoa \emph{injusta} e \emph{ingrata} não é boa nem honesta.\par b) Incredulidade. --- A \emph{incredulidade} é o vicio dos que dizem que só deve\par 44 Vor. II.\pgn{690}admittir-se o que estiver ao alcance da razão humana, e que por isso deve rejeitar-se a Revelação sobrenatural.\par 0) Impiedade. --- A \emph{impiedade} é o vicio dos que, embora admittam a existencia de \textsc{Deus}, comtudo têm a ousadia de o ultrajar e combater.\par d) Atheismo. --- O \emph{atheismo} é o vicio dos que não admittem ou dizem não admittir a existencia de um Creador e Conservador do universo.\par e) Superstição. --- A \emph{superstição} é o vicio dos que adoram a creatura em logar do Creador, ou misturam coisas perniciosas ou superfluas no culto do verdadeiro \textsc{Deus}. --- Todavia, não é \emph{superstição} o culto, que os catholicos prestam ás imagens de Nosso Senhor e dos Santos, pois o \emph{termo} d'esse culto não é um pedaço de pedra, de pau, ou de tela, mas é a Pessoa de Nosso Senhor e a dos Santos, que as imagens representam. Assim a estima, que temos a um retrato de um amigo, não tem por termo um pedaço de papel, mas a propria pessoa do amigo.\par £) Fánatismo. --- O \emph{fanatismo} é o vicio dos que propagam a \emph{superstição} por meio do ferro e do fogo, ou defendem a verdadeira Religião com meios perniciosos e condemnados. ; J\par Estes vicios encontram-se quasi todos na \emph{religio positivista}, fundada por \emph{Comte}. Eis o resumo dos erros do escriptor francez a este respeito: « Existe um \textsc{Deus}, que é o \emph{Grande Ser}, a \emph{Humanidade}. Fazem parte do Grande Ser todos os seres aptos para a assimilação, e por isso não sÓ os homens, mas tambem os animaes. Ao Grande Ser devem acerescentar-se o \emph{Grande Fetiche}, que é a nossa terra com o \emph{systema solar}, e o \emph{Grande Meio}, que é o \emph{espaço}. E' a Trindade positivista ». Esta Religião tem o seu signal da Cruz, cuja formula é a seguinte: « \emph{O amor por principio} --- a \emph{ordem por base} --- o \emph{progresso por fim} ». --- E não basta. \emph{Comte} reformou o \emph{almanak}; instituiu sacramentos; fundou um sacerdocio com os tres graus de \emph{aspirantes}, \emph{vigavios} e \emph{sacerdotes}, subordinados a um \emph{Grande Sacerdote}. --- E não é tudo. A nova religião devia ter os seus santos. A primeira pessoa, canonizada por \emph{Comte}, foi a \emph{Grande Sacerdotiza}, --- \emph{Clotilde de} Vaus. Morta esta, o philosopho\pgn{691}prestou-lhe culto, venerando as suas reliquias e passando varias horas do dia em adoração deante da sua imagem.\par Tudo isto seria ridiculo, se não fosse impio e pagão. E houve e ha infelizes que desprezam a Religião Catholica, tão bella, tão pura, tão elevada, para se abaterem deante da louca impiedade, que se chama \emph{religião positivista}. Simplesmente triste!}.

{\small\noindent\textsc{Summario}: --- Deveres do homem relativos á sua alma. --- Deveres do ho- mem relativos ao seu corpo.\par}

\pg{690} \textbf{DEVERES} DO HOMEM PARA COMSIGO

%%>>> art-62

\section*{Artigo I.\\ Deveres do homem relativos á sua alma.}
\addcontentsline{toc}{section}{Artigo I. --- Deveres do homem relativos á sua alma.}

131. \textbf{Deveres} do \textbf{homem para} comsigo. --- \emph{Deveres} do homem para \emph{comsigo} são as obrigações que elle tem de praticar alguns actos, que tendem \emph{directa} e \emph{proximamente} para o seu aperfeiçoamento. --- O principio, de que derivam e a que se reduzem \pg{691}  todos os deveres do homem para \emph{comsigo}, é o preceito do \emph{amor ordenado de si mesmo}, imposto pela lei \emph{natural} \footnote{O principio, de que derivam e a que se reduzem todos os deveres do homem para \emph{comsigo}, é o \emph{amor ordenado de si mesmo}. É, pois, necessario notar que o principio d'esses deveres não é o preceito do \emph{amor de si mesmo}; porque o \emph{amor de si mesmo} é uma \emph{necessidade}, não é um \emph{dever}; --- mas é o preceito do amor \emph{ordenado} de si mesmo; pois o amor de si mesmo, emquanto deve ser \emph{ordenado}, é materia de \emph{preceito} e constitue um \emph{dever}.\par Mas quando é que o amor de si mesmo é \emph{ordenado}? --- Para que a resposta seja clara, devemos notar que tres amores devem Occupar o coração do homem: o amor de \textsc{Deus}, o amor de si mesmo, o amor do proximo. O amor de \textsc{Deus} deve ser supremo, absoluto e superior a todos os outros amores; porque \textsc{Deus} é o Bem supremo, absoluto e principio de todos os bens. Os outros bens devem ser amados, não só depois de \textsc{Deus}, mas tambem em \textsc{Deus} e por \textsc{DEus}; porque, como dissemos, todo o bem creado é uma participação e uma semelhança da infigita Bondade de \textsc{Deus}. Por isso o \emph{amor de si mesmo} será \emph{ordenado}, quando estiver sujeito ao amor de \textsc{Deus}, quando o homem se amar a si mesmo em \textsc{Deus} e por \textsc{Deus}, cuja imagem e semelhança está gravada em a nossa alma, simples e espiritual. O amor do \emph{proximo} deve ser semelhante ao amor de si \emph{mesmo}; porque tambem o \emph{proximo} deriva de \textsc{Deus} e tem gravada em si a imagem e a semelhança do Creador. --- D'onde se segue que o amor, que o homem deve ter a si \emph{mesmo} e aos outros, se funda \emph{iminediata} e \emph{prozimamente} na dignidade da \emph{pessoa humana}, emquanto é semelhante ao proprio \textsc{Deus}; «\emph{mediata} e \emph{remotamente}, no Bem infinito.\par Os \emph{pantheistas}, que fazem do homem uma \emph{evolução}, emanação ou \emph{phenomeno} do proprio \textsc{Deus}, e os \emph{materialistas}, que reduzem a pessoa humana a um conjuncto de atomos destituidos de finalidade e só movidos pela força mechaniea, não podem reconhecer, como não reconhecem, deveres do homem para \emph{comsigo}; porque o dever importa \emph{inferioridade} e \emph{livre arbitrio}, e para os \emph{pantheistas} não ha \emph{inferioridade}, nem para os \emph{materialistas} ha \emph{livre arbitrio}.\par Alguns escriptores negam a existencia dos \emph{deveres} do homem para \emph{comsigo}. Porquanto, dizem elles, sendo o \emph{dever} correlativo ao \emph{direito}, se o homem tivesse \emph{deveres} para \emph{comsigo}, devia tambem ter \emph{direitos} para \emph{comsigo}; e d'este modo uma e a mesma pessoa seria o sujeito do \emph{direito} e do \emph{dever respectivo}. Ora o sujeito do \emph{direito} deve ser distincto do sujeito do \emph{dever respectivo}. Logo o homem não póde ter \emph{deveres} para \emph{comsigo}.\par Estes escriptores erram. Porquanto suppõem que todo o dever é \emph{jurídico}. Ora isto é falso; porque ha tambem deveres \emph{moraes}, como são os que o homem tem para \emph{comsigo}. Na verdade, para que o homem possa ter e tenha deveres para \emph{comsigo}, basta que nelle existam differentes faculdades, capazes de serem ordenadas, e uma força, que tenha o \emph{dever} de as ordenar e dirigir. Ora existem effectivamente no homem differentes faculdades, capazes de serem ordenadas, e existe nelle uma força, que é a vontade, e que tem o\pgn{692}\emph{dever}, imposto pela lei natural, de regular e dirigir todas as outras faculdades, conforme os dictames da razão. --- Estes \emph{deveres} --- não são \emph{juridicos}, porque os deveres \emph{jurídicos} fundam-se na \emph{justiça}, tomada no seu mais rigoroso sentido, e a \emph{justiça} exige effectivamente um \emph{direito} correlativo ao \emph{dever} e uma distinecüo \emph{real} entre o sujeito do \emph{dever} e o do \emph{direito}; --- mas são \emph{moraes}, porque são impostos pela lei \emph{natural}, que é norma de \emph{moralidade}. Por isso não existem \emph{direitos} do homem para \emph{comsigo}.}.

\pg{692} 132. \textbf{Divisão} dos deveres do homem para comsigo. --- Os mais importantes deveres do homem para \emph{comsigo} referem-se ao aperfeiçoamento da sua \emph{alma}, que é a parte mais nobre do \emph{composto humano}, e á conservação do seu \emph{corpo}, que é vivificado pela alma. --- Por isso ha duas especies de deveres: uns relativos á \emph{alma}; outros, ao \emph{corpo}. --- Neste artigo tratamos dos deveres relativos á \emph{alma}; no seguinte, dos deveres relativos ao \emph{corpo} \footnote{Dizendo que o homem tem \emph{deveres} relativos á conservação do \emph{corpo}, consideramos o \emph{corpo} não como separado da alma, mas como unido substancialmente com ella e dotado de vida \emph{vegetativa} e \emph{sensitiva}. ;}.

133. \textbf{Deveres} do homem relativos á sua alma. --- O aperfeiçoamento da nossa \emph{alma} consiste no aperfeiçoamento das faculdades \emph{espirituaes}, que derivam d'ella e que .são a \emph{intelligencia} e a \emph{vontade}. Os deveres, pois, do homem, relativos á sua \emph{alma}, consistem na obrigação que elle tem de aperfeiçoar --- a sua \emph{intelligencia}, apprehendendo as verdades necessarias para o exacto cumprimento dos seus deveres, --- e a sua \emph{vontade}, praticando a \emph{virtude} e dirigindo-se para o seu fim \emph{ultimo}, que é a felicidade eterna \footnote{O aperfeiçoamento da alma consiste no aperfeiçoamento das faculdades \emph{espirituaes}, que d'ella derivam. Porquanto, o aperfeiçoamento da alma consiste em produzir, com facilidade, agrado e uniformade, os actos necessarios para o conhecimento da verdade e para a pratica do bem. Ora a alma não produz aquelles actos pela sua \emph{essencia}, mas sim pelas suas faculdades \emph{espirituaes}, que são a \emph{intelligencia} e a \emph{vontade}. Logo o aperfeiçoamento da alma consiste no aperfeiçoamento da \emph{intelligencia} e da \emph{vontade}.}.

134. O homem tem o \textbf{dever} de aperfeiçoar a sua intelligencia. --- O homem terá o dever de aperfeiçoar a sua \emph{intelligencia}, se tiver o dever de praticar certos e determinados actos, que se não podem praticar devidamente sem o conhecimento de certas e determinadas \emph{verdades}; pois quem tem o dever de conseguir um fim, tem tambem o dever de empregar os meios indispensaveis. Ora o homem tem o dever de praticar certos e determinados actos, como são os da \emph{Religião}, e estes actos não se podem praticar devidamente sem o conhecimento de algumas \emph{verdades}, como são a \emph{existencia} de \textsc{Deus}, á \emph{espiritualidade} \pg{693}  da nossa alma, a sua \emph{origem} e \emph{destino}, ete. Logo o homem tem o dever de aperfeiçoar a sua \emph{Intelligencia} \footnote{O homem, pois, é obrigado a adquirir os conhecimentos necessarios para a pratica dos seus deveres, não só dos que são communs a todos os homens, mas tambem dos que são proprios do seu estado. A esta classe de conhecimentos pertencem os que se referem a DEus e aos deveres religiosos, e os que dizem respeito ao proprio estado e condição social.\par Os conhecimentos, absolutamente necessarios para o homem, estão compendiados num pequeno livro, chamado \emph{catecismo}. --- \emph{Jouffroy} diz: « Existe um livrinho, que as creanças aprendem e recitam na Egreja. Lêde-o: é o \emph{catecismo}. Achareis nelle uma solução para todos os problemas da ordem moral, --- para todos, sem excepção. Perguntai ao christão a origem da especie humana, elle sabe-o; perguntai-lhe o seu destino, elle sabe-o. Perguntai a uma pobre criança a razão por que existe sobre a terra e o que fará depois da sua morte, e ella dar-vos-ha uma sublime resposta. Origem do mundo, origem das raças, destino do homem nesta e noutra vida, relações do homem para com \textsc{Deus}, deveres do homem para com os seus semelhantes, direitos do homem sobre os seres creados, tudo esta criança conhece; e, mais tarde, não ignorará o direito natural, o direito politico, o direito internacional. Eis o que para mim é uma grande Religião; reconheço-a neste signal, que ella não deixa sem resposta nenhum dos problemas, que agitam o genero humano ». --- \emph{T'roplong} accrescenta: « Depois de ter lido muito, estudado muito, vivido muito, quando se approxima a hora da morte, conhece-se que a unica coisa verdadeira é o \emph{catecismo} n.\par Pelo mesmo motivo, por que deve abraçar a verdade, o homem deve evitar o erro, que sabe insinuar-se por muitos e differentes modos.\par Ao \emph{dever}, que o homem tem de aperfeiçoar a sua \emph{ntelligencia} pelo amor da \emph{verdade} e pela fuga do \emph{erro}, oppõem-se as modernas theorias, que, sob o especioso nome de \emph{liberdade}, abatem o homem, tornando-o escravo de pernieiosos erros e de funestas desordens. 'Tal é a theoria ácerca da \emph{liberdade} de \emph{pensamento}, de \emph{consciencia}, de \emph{ler}, de \emph{ouvir}, ete.}.

135. O homem tem o \textbf{dever} de \textbf{aperfeiçoar} a \textbf{sua vontade}. --- O homem terá o dever de aperfeicoar a sua vontade, se tiver o dever de praticar a \emph{virtude} e de tender para o \emph{fim ultimo}; pois é nessa pratica e nessa tendencia, que deve consistir e consiste a perfeição da nossa vontade. Ora o homem tem effectivamente o dever, imposto pela lei \emph{natural}, de praticar a \emph{virtude} e de tender, por tal meio, para o seu \emph{fim ultimo}. Logo o homem tem o dever de aperfeiçoar a sua \emph{vontade} \footnote{O homem, pois, esta obrigado a amar a \textsc{Deus} sobre todas as coisas, e a amar as outras coisas em \textsc{Deus} e por \textsc{Deus}, --- amor de que derivam todas as outras virtudes. --- E, como este amor encontra difficuldades, sobretudo nas paixões do \emph{appetite}, concupiscivel e irascivel, \emph{o} homem deve mortificar as paixões, sobretudo o orgulho e a sensualidade. Os meios principaes, proporcionados\pgn{694} para essa mortificação, são a pratica da \emph{humildade} e da \emph{modestia}, e o \emph{trabalho} (pois o \emph{ocio} é o pae de todos os vicios). --- D'este modo haverá harmonia entre a vontade e a intelligencia, --- harmonia, em que consiste o esplendor, a belleza propria da nossa vida moral. i\par Do que temos dito vê-se que o dever, que o homem tem de aperfeiçoar a sua intelligencia e a sua vontade, é uma consequencia da lei \emph{natural}. Porquanto a lei \emph{natural} dirige o homem para \textsc{Deus}, que é o summo e perfeitissimo Bem, e por isso o centro de toda a felicidade. Ora o homem dirige-se ou tende para \textsc{DEus}, emquanto O conhece com a \emph{intelligencia e} O ama com a \emph{vontade}. Logo é a propria lei \emph{natural}, que impõe ao homem o dever de aperfeiçoar a sua intelligencia e a sua vontade.}.

136. 0 \textbf{homem} tem o \textbf{dever} de empregar \textbf{os meios proporcionados para} o \textbf{aperfeiçoamento} da sua intelligencia e da \textbf{sua vontade}. --- A proposição é evidente. --- Quem tem o dever de conseguir \pg{694}  um fim, tem tambem o dever de empregar os \emph{meios} proporcionados para a consecução d'esse fim; pois, se assim não fosse, um sujeito, ao mesmo tempo e sob o mesmo aspecto, teria e deixaria de ter o dever de conseguir um fim. Ora o homem, como demonstrámos, tem o dever de aperfeiçoar a sua intelligencia e a sua vontade. Logo o homem tem tambem o dever de empregar os \emph{meios} proporcionados para o aperfeiçoamento da sua intelligencia e da sua vontade \footnote{Nada é mais claro \emph{em theoria} do que esta verdade, e todavia nada é mais difficil \emph{em pratica}. Todos estão mais ou menos convencidos de que é necessario aperfeiçoar a propria intelligencia e vontade; mas poucos empregam os meios proporcionados para esse fim. Uns não empregam meio algum; outros só empregam meios faceis e de pouca efficacia; outros empregam meios contrarios a esse aperfeiçoamento, e assim abatem a intelligencia até ao erro e inclinam a vontade até á desordem das paixões. A consecução do fim ultimo é uma felicidade tão grande, que vale a pena fazer algum sacrificio, luctando contra os muitos e varios obstaculos, que impedem o exercicio da virtude, e subordinando as faculdades appetitivas ás faculdades intellectuaes. 4}.

137. \textbf{Corollario}. --- Do que deixamos exposto vê-se que o homem, relativamente ás coisas creadas, deve servir-se ou abster-se d'ellas, conforme concorrem, ou não, para o aperfeiçoamento da sua intelligencia e da sua vontade \footnote{E' esta uma regra pratica da altissima sabedoria. O homem não foi creado para os bens caducos da terra, mas foi creado para servir e amar \textsc{Deus} nesta vida, e para conseguir a felicidade eterna, a qual consiste na posse do proprio \textsc{Deus}. O Bem Infinito é, pois, o fim \emph{ultimo} do homem; os bens finitos só podem ser \emph{meios} para a consecução do fim \emph{ultimo}. Ora toda a bondade do \emph{meio} consiste na sua proporção com o fim. --- D'ahi resultam estes corollarios: --- 1.9) Se os bens finitos concorrem para a consecução do nosso fim, podemos servir-nos d'elles; --- 2.9) Se os bens finitos impedem essa consecução, devemos abster-nos d'elles; --- 3.9) Se os bens finitos tornam mais diffieil o exercicio da virtude e por isso a posse da felicidade, convem tambem deixal-os. --- D'onde se vê que a pratica dos \emph{conselhos evangelicos} não só se harmoniza com a /é, mas tambem com a \emph{sciencia}, embora se não harmonize muito com as ruins paixões da humanidade decahida!}.

%%>>> art-63

\section*{Artigo II.\\ Deveres do homem relativos ao seu corpo.}
\addcontentsline{toc}{section}{Artigo II. --- Deveres do homem relativos ao seu corpo.}

\pg{695} 138. \textbf{Deveres} do homem \textbf{relativos ao} seu corpo. --- \emph{Deveres} do homem relativos ao \emph{corpo} são, como dissemos, os que se referem á conservação do elemento \emph{material} e \emph{sensivel} do composto humano. --- Estes deveres são \emph{negativos} e \emph{positivos}. Os \emph{negativos} consistem sobretudo em não procurar \emph{directamente} a propria morte pelo \emph{suicidio}. Os \emph{positivos} consistem em empregar os meios convenientes para a conservação da saude e da vida \footnote{Não é nossa intenção enumerar cada um dos deveres do homem, relativos ao \emph{corpo}. \emph{SO} enumeraremos os principaes, e, entre estes, os que têm sido objecto de discussão entre os escriptores.\par Os antigos \emph{Stoicos} julgaram que o homem sabio deve descuidar-se da conservação da sua saude e da vida e procurar \emph{directamente} a propria morte, quando esta morte é util aos amigos, á patria, ou é o unico remedio para os males da vida presente. Esta velha opinião, resuscitada pelos incredulos do seculo passado, é muito seguida em nossos dias, chegando a considerar-se o \emph{suicidio} como um acto heroico. --- Os protestantes e os sensualistas, pelo contrario, exageram o dever de conservar á saude e a vida, chegando a condemnar, como peceado, a mortificação do corpo, embora moderada e pratieada para obedecer ás leis da Egreja, ou para satisfazer pelas proprias culpas, ou para reprimir as paixões desordenadas.\par Heprovamos estes erros e exageros. D'ahi a razão das proposições, que vamos estabelecer e provar no texto.}.

139. O homem tem o dever de não procurar \textbf{directamente} a propria \textbf{morte pelo} suicidio.

a) O homem terá o dever de não procurar \emph{directamente} a propria morte pelo \emph{suicidio}, se não tiver o \emph{dominio} da sua vida, mas só o \emph{uso}; pois o direito de destruir uma coisa não pertence ao depositario, mas unicamente ao dono. Ora o homem só tem o \emph{uso} da propria vida; porque o \emph{dominio} sobre a creatura \emph{racional} só pertence ao Creador. Logo o homem tem o dever de não procurar \emph{directamente} á propria morte pelo \emph{suicidio} \footnote{Pergunta-se: Como se prova que \textsc{DEus} concedeu ao homem só o uso da propria vida, e não o \emph{dominio}?\par Resposta. E' a voz da natureza, a qual nos diz que \textsc{DEUS} concedeu ao homem só o uso da vida, e não o \emph{dominio}. Na verdade, o homem tem um \emph{instincto natural} para a conservação da sua vida. Ora tal instincto leva-nos a descobrir a intenção e a vontade de \textsc{Deus}, que nol-o deu; assim como da tendencia natural para a felicidade, que experimentamos em a nossa alma, concluimos que \textsc{Deus} nos destinou para essa felicidade, e que nol-a concederá, se empregarmos os meios proporcionados. Logo do instincto \emph{natural},\pgn{696}que temos para a nossa conservação, deduzimos que \textsc{Deus} quer esta conservação. Se \textsc{Deus} quer a conservação da nossa vida, é claro que Elle só nos concedeu o uso da propria vida, e não o \emph{dominio}; porque \textsc{Deus}, se nos tivesse concedido tambem o \emph{dominio}, seria indifferente com relação á conservação ou \& destruição da nossa vida. --- Além d'isso, o homem é \emph{essencialmente} servo de \textsc{DEus}, e deve trabalhar nesta vida para a consecução da felicidade immortal. Ora o servo não póde estabelecer a duração do serviço, mas deve esperar até que a voz do seu senhor o chame para o descanço, para a recompensa. Por isso, o \emph{suicida} é um servo infiel, e merece o castigo, que se dá a um soldado, quando caprichosamente abandona, antes do tempo estabelecido, o posto que lhe foi confiado.}.

\pg{696} b) O homem, que procura \emph{directamente} a propria morte, offende a \textsc{Deus}, a \emph{si mesmo} e a \emph{sociedade}. --- Offende a \textsc{Deus}; porque, dispondo da sua vida, eujo \emph{dominio}, como vimos, pertence a \textsc{Deus}, lesa um direito do proprio \textsc{Deus}. --- Offende a \emph{si mesmo}; porque se priva de um grande bem, que é a existencia, e se condemna a uma infelicidade eterna. --- Offende a \emph{sociedade}; porque injustamente a priva de um membro, e é occasião de escandalo e causa de damno para os outros. Logo o homem tem o dever de não procurar \emph{directamente} a propria morte pelo \emph{suicidio}. --- Dissemos --- \emph{directamente}; para indicarmos que o homem não póde praticar um acto unicamente com o fim de se matar \footnote{Cf. \emph{Sum}. \emph{Th}., 2-2, q. 64, a. 5. --- Os damnos, que o \emph{suicidio} produz nas familias e na sociedade, são incalculaveis. Uma familia precisa do trabalho e do amparo do pae; este dá-se a morte; a familia cobre-se de lucto e de miserias. A sociedade precisa de homens, que possam cooperar para o bem communm; o \emph{suicidio} inutiliza tantas esperanças. --- E este mal é excessivamente contagioso. A repetição dos casos torna menos horrorosa a idea da morte; a descripção e, por vezes, a glorificação do \emph{suicidio}, feita, pelos jornaes, exaltam a phantasia de modo, que se attenta contra a propria vida, como se se cumprisse a mais nobre e sublime das emprezas!\par Os defensores do \emph{suicidio} apresentam os seguintes argumentos: --- 1.º) o \emph{suicida} não se offende a si mesmo, porque a maxima diz: « \emph{volenti non fit injuria} »; --- 2.9) é rasoavel deixar a vida, quando é pesada ou inutil; --- 8.9) quando se n\&o póde viver com honra, o \emph{suicidio} é um acto heroico, digno de um homem sabio.\par Estes argumentos são todos falsos. --- E' falso o primeiro; porque á maxima --- \emph{volenti non fit injuria} --- só vale, quando se trata de bens \emph{alienaveis}, mas não vale, quando se trata de bens \emph{inalienaveis}, isto é, de bens de que não podemos dispor; ora a vida é um d'estes bens, --- E falso o segundo; porque a nossa vida, embora pareça pesada ou inutil, depende sempre de Drs, e não podemos dispor della. --- E' falso o terceiro; porque o \emph{suicidio}, longe de ser um acto de heroismo, digno de um homem sabio, é um acto de cobardia, proprio de um homem fraco, que não sabe soffrer.\par A causa ordinaria do \emph{suicidio}, fóra dos casos de alienação mental, é a \emph{irreligido}. Porquanto, se o homem perder de vista a vida futura e se não\pgn{697}tiver a esperança de um premio eterno, nada é mais natural do que procurar no \emph{suicidio} o remedio para os males, que actualmente o atormentam. A experiencia demonstra que a idea do \emph{suicidio} começa, quando e onde acaba a idea da fé. --- Por isso são dignos de compaixão os que, para impedirem os continuos attentados contra a propria existencia, invocam meios humanos, excluindo um meio divino, verdadeiramente efficaz, que é a pratica dos deveres religiosos. A cada mal o remedio proporcionado.\par O que dizemos ácerca do \emph{suicidio} deve applicar-se ao dever que o homem tem de não procurar \emph{directamente} a \emph{mutilação} do seu corpo.\par Note-se que, declarando \emph{illicito} e \emph{injusto} o \emph{suicidio}, queremos fallar do suicidio, commettido unicamente pela livre vontade do homem, e não do \emph{suicidio}, que fosse ordenado por \textsc{Deus}; porque, se \textsc{DEus} o manda, o \emph{suicidio} perde toda a sua immoralidade, o torna-se licito e necessario.}.

\pg{697} 140. O \textbf{homem póde}, por um motivo \textbf{justo}, \textbf{procurar indirectamente} a propria \textbf{morte}. --- O homem póde fazer o que o Direito \emph{natural} não só não reprova, mas até louva e recompensa. Ora o Direito \emph{natural} --- não só não condemna o homem, que, levado por üm motivo justo, isto é, no intuito de alcançar um bem mais nobre do que a propria vida, procura \emph{indirectamente} ou \emph{permitte} a propria morte; --- mas até chega a prodigalizar-lhe louvores e recompensas; assim todos elogiam a abnegação do soldado, que se expõe aos golpes dos inimigos para a salvação da patria. Logo o homem póde, por um motivo justo, procurar \emph{indirectamente} a” propria morte. --- Dissemos --- \emph{indirectamente}; porque, neste caso, o homem não quer propriamente a sua morte, mas um bem muito elevado e digno, cuja consecução ceustar-lhe-ha a morte \footnote{E' necessario advertir que o homem, que, levado por um justo motivo, procura \emph{indirectamente} a morte, não se mata a si mesmo, não quer propria ou directamente a propria morte, mas só a \emph{permitte} com o fim de alcançar um bem mais elevado que a vida, e esta permissão não deriva do capricho, mas sim do amor da \emph{justiça}, da \emph{caridade} ou de qualquer outra virtude. Ora isto não é contrario ao Direito \emph{natural}, que só prohibe o suicidio \emph{directo} e não justificado por um motivo mais elevado, como seria uma ordem formal de \textsc{Deus}. D'ahi a idea de \emph{sacrificio}.\par Quanto mais deforme e desprezivel é o \emph{suicidio}, tanto mais bello e glorioso é o \emph{sacrificio}. A pratica dos proprios deveres póde levar o homem a desprezar a sua vida corporea, e a derramar o seu sangue sem medida e sem remorso. O soldado, que faz do seu peito um escudo para a defeza da patria, --- o christão, que, pela confissão da sua fé, se offerece aos tormentos, ao martyrio, --- o homem dedicado, que, para salvar um naufrago, affronta o furor da tempestade e despreza o perigo imminente, --- eis os verdadeiros heroes, dignos da estima e da admiração do mundo. E' por isso que a Pessoa adoravel de N. S. Jesus CHmisro, que deu todo o seu sangue precioso para a salvação do mundo, é e será sempre o objecto da estima e do amor para todo o coração bem formado!}. :

\pg{698} 141. O \textbf{homem deve procurar positivamente} a \textbf{conservação da sua vida} e da \textbf{sua saude}. --- A vida e a saude são meios, que \textsc{Deus} concedeu ao homem, para este poder cumprir os proprios deveres e assim alcançar a felicidade eterna. Ora estes meios precisam, para a sua conservação, de um cuidado \emph{positivo}, da parte do homem. Logo o homem deve procurar \emph{positivamente} a conservação da sua vida e da sua saude. --- Dissemos --- \emph{positivamente}; porque, para a conservação da vida e da saude, não basta que o homem nada faça que lhes seja contrario, mas é preciso que empregue os meios convenientes \footnote{Deve notar-se que o homem, para a conservação da sua vida e saude, não tem o dever de empregar meios \emph{eowtraordinarios}. Porquanto, sendo \textsc{Deus} um legislador de infinita prudencia, devemos suppor, e com razão, que Elle não quer impor-nos um preceito \emph{positivo} com tão grande incommodo nosso, que exceda proporcionalmente os bens da vida e da saude. Ora se, para a conservação da nossa saude, devessemos empregar meios \emph{ewtraordinarios}, este cuidado seria uma coisa insupportayel e absorveria todas as outras occupações, ainda as da ordem mais elevada, que é a espiritual. --- Por isso devemos afastar os perigos proximos, que ameaçam a saude e a vida, empregando os meios \emph{ordinarios}.}.

142. Ao dever da conservação da vida e da \textbf{saude} não \textbf{se oppõe} a \textbf{moderada} mortificacdo do corpo. --- O homem tem o dever de procurar \emph{positivamente} a conservação da sua vida e saude, emquanto esta conservação é um meio para o homem cumprir os proprios deveres e assim alcançar a felicidade eterna. Ora a moderada mortificação do corpo é tambem um \emph{meio}, e dos mais effieazes, para o homem cumprir os proprios deveres E porque os peccados nascem geralmente da rebellião do appetite contra a razão. Logo ao dever da conservação da vida e da saude não se oppõe a moderada mortificação do corpo \footnote{O \emph{protestantismo}, nascido da desordem das paixões, inculeou uma doutrina conforme á sua origem. D'ahi o odio contra a mortificação christa, prégada pela Egreja Catholica. Não insistimos neste ponto. A experiencia attesta que a moderada mortificação do corpo, longe de destruir a saude, robustece-a e protege-a até á mais longa velhice; como, pelo contrario, nada mais proprio para arrüinar a saude e abreviar a existencia do que uma vida molle e desordenada. --- Mas, supponhamos que a moderada mortificação do corpo chegue a alterar a saude e abreviar de alguns annos a existencia. Será ella prohibida, como contraria ao dever que temos de conservar a vida? Não, com certeza. Porquanto, se é permittido e louvavel enfraquecer a saude e abreviar a vida por amor da \emph{sciencia}, que todavia é um bem secundario e não se relaciona com \emph{a} vida eterna, muito mais o será, tratando-se da acquisição da \emph{virtude}, que é um bem essencial e um meio indispensavel para a consecução da felicidade. Os que combatem, pois, a mortificação chrisià, são \emph{materialistas} em \emph{theoria} ou em \emph{pratica}. 3}.

{\small\noindent\textsc{Summario}: --- Direitos do homem para com o proximo. --- Deveres do homem para com o proximo.\par}

\pg{699} \textbf{DIREITOS} E \textbf{DEVERES DO} HOMEM

PARA COM O PROXIMO

%%>>> art-64

\section*{Artigo I.\\ Direitos do homem para com o proximo.}
\addcontentsline{toc}{section}{Artigo I. --- Direitos do homem para com o proximo.}

143. \textbf{Direitos} do homem \textbf{para} com o proximo. --- \emph{Direitos} do homem para com o \emph{proximo} são os poderes moraes, que elle tem de exigir dos seus \emph{semelhantes} o que.lhe é devido. --- Estes \emph{direitos} são \emph{jurídicos} ou \emph{moraes}, conforme se baseam na idea da \emph{justiça} ou de alguma outra virtude.

144. \textbf{Direitos} juridicos do homem \textbf{para} com o proximo. --- Os direitos \emph{jurídicos} do homem para com o proximo podem ser \emph{innatos} ou \emph{adquiridos}, conforme se baséam na propria natureza do homem ou nalgum facto contingente. --- Os principaes direitos 2\emph{nnatos} são os seguintes: o direito á \emph{liberdade religiosa}, á \emph{dignidade pessoal}, \emph{á vida}, o direito de \emph{legitima defeza}, de \emph{propriedade}, de \emph{independencia} e de \emph{associação}. Estes \emph{direitos} nascem dos \emph{deveres} do homem para com \textsc{Deus} e para \emph{comsigo} \footnote{Os direitos \emph{innatos} do homem nascem dos \emph{deveres}, que elle tem para com \textsc{Deus} e para \emph{comsigo}. --- Na verdade, o homem, tendo sido creado por \textsc{Deus}, tem o \emph{dever} de prestar livremente a \textsc{Deus} o culto devido: d'ahi o \emph{direito á liberdade religiosa}. Além d'isso, \emph{o} homem, sendo um ser nobilissimo, dotado de alma \emph{espiritual}, capaz de aperfeiçoamento, e destinado para um fim proprio, tem o \emph{dever} de aperfeicoar a propria alma e de tender para o seu fim ultimo; e por conseguinte póde exigir que os outros respeitem nelle o ser nobilissimo, que, derivando de \textsc{Deus}, deve voltar para \textsc{Deus}; d'ahi o direito á \emph{dignidade pessoal}. --- E não basta. O homem recebeu de \textsc{Deus} a vida, que devemos conservar para a consecução do fim imposto pelo Creador; d'ahi o direito á \emph{vida}, de que deriva o direito de \emph{legitima defeza} e o de \emph{propriedade}. --- Estes direitos, que resultam dos deveres do homem para com \textsc{Deus} e para comsigo, exigem, como natural complemento, dois outros direitos: o de \emph{independencia} e o de \emph{associação}; pois quem tem o direito ao fim, tem tambem o direito aos respectivos meios indispensaveis.\par Alguns escriptores não reconhecem no homem direitos 2\emph{naatos}, provenientes da lei \emph{natural}, e dizem que todos os nossos direitos são \emph{adquiridos}, --- dependentes exclusivamente da lei \emph{positiva}.\par A falsidade d'esta opinião já foi demonstrada. A lei \emph{natural}, como dissemos\pgn{700}, impõe ao homem o \emph{dever} de conseguir o fim ultimo pela observancia da ordem moral. Se a lei \emph{natural} impõe ao homem o dever de conseguir o fim ultimo pela observancia da ordem \emph{moral}, deve dar-lhe necessariamente os respectivos direitos de exigir tudo o que for necessario para essa observancia e por isso para a consecução do mesmo fim ultimo; pois repugna que a lei imponha ao subdito um \emph{dever}, e não lhe conceda o \emph{direito} de exigir o que é necessario para o cumprimento d'esse dever.\par Os adversarios dizem que o homem não tem direitos \emph{innatos}, porque não os conhece desde o \emph{nascimento}. --- Esta razão é và. Uma coisa é o \emph{direito}, e outra coisa é o seu \emph{exercicio}. Se as crianças, como tambem os loucos, não \emph{exercem} o direito, porque o não conhecem, nem por isso deixam de o possuir, nem os outros deixam de ter o \emph{dever} de o respeitar.\par Os direitos \emph{innatos} tem tres propriedades, pois que são \emph{inalienaveis}, \emph{eguaes} em todos os homens, e \emph{limitados}.\par a) Os direitos \emph{innatos} são \emph{inalienaveis}. M verdade, os direitos \emph{innatos}\par são meios indispensaveis, concedidos pela lei \emph{natural}, para o cumprimento dos nossos deveres e para a realização da ordem moral. Quem, pois, renuncia a esses direitos, transgride os seus deveres e perturba a ordem. moral. --- Como se vé, a \emph{inalienabilidade} dos direitos \emph{innatos} funda-se no nexo necessario, que existe entre elles e o nosso fim ultimo.\par \emph{b}) Os direitos \emph{innatos} são \emph{eguaes} em todos os homens. Na verdade, estes. direitos basêam-se exclusivamente na propria natureza do homem, e, como esta natureza é uma e a mesma em todos os homens, tambem estes direitos devem ser \emph{eguaes}.--- Esta \emph{egualdade} só existe nos direitos \emph{innatos}, e não nos direitos \emph{adquiridos}; porque estes fundam-se em factos contingentes, e como estes factos são deseguaes (pois os individuos são deseguaes nas aptidões, na força, na actividade, \emph{etc}.), tambem os direitos \emph{adquiridos} devem ser \emph{deseguaes}. --- A \emph{egualdade absoluta} dos direitos \emph{adquiridos} em todos os homens, proclamada pelos revolucionarios francezes do seculo passado e pelos socialistas modernos, além de injusta, porque não corresponde aos meritos de cada um, é contraria --- a toda ordem, porque esta exige a variedade na unidade, --- e ao progresso social, porque o motor principal da actividade humana é o desejo da desegualdade ou da superioridade.\par c) Os direitos \emph{innatos} são \emph{limitados}. Os direitos \emph{innatos}, derivando da propria lei \emph{natural}, têm á extensão, que ella lhes concede. Ora a lei \emph{natural} limita os direitos \emph{innatos}, quando o exercicio d'elles involve a transgressão de um dever, ou a violação de um outro direito. --- Os limites, pois, dos direitos \emph{innatos} derivam primeiramente dos deveres, que são impostos ao mesmo individuo, e, em segundo logar, dos direitos mais elevados dos outros individuos ou da sociedade. Estes limites podem ser determinados não só pelo Direito \emph{natural}, mas tambem pelo Direito \emph{positivo}, embora este deva sempre basear-se naquelle.}.

a) \textbf{Direito} á liberdade religiosa, --- Direito á \emph{liberdade religiosa} é o poder moral de crer e praticar tudo o que a verdadeira Religião ensina. --- Este direito existe no homem. Porquanto, o homem, tendo o \emph{dever}, imposto pela lei \emph{natural}, de crer e praticar tudo o que a verdadeira Religião ensina, deve ter tambem o \emph{direito}, concedido pela mesma lei, relativo áquella \pg{700} \pg{701}  crença e áquella pratica. Logo o homem tem direito á \emph{liberdade religiosa} \footnote{Os \emph{incredulos} e os \emph{racionalistas} reconhecem no homem o direito á \emph{liberdade religiosa}, mas este direito, para elles, consiste no poder de praticar toda a religião, que á propria razão julgar verdadeira. --- Esta liberdade, como demonstr\&mos, oppõe-se aos nossos deveres para com \textsc{Deus}, e por isso não póde ser um \emph{direito}. --- A liberdade deve existir com a verdade, e não com o erro; porque é só a verdade que nos torna verdadeiramente livres.}.

b) Direito á dignidade pessoal. --- Direito á \emph{dignidade pessoal} é o poder, commum a todo o homem, de exigir que seja reconhecido como ente racional, dotado de fim proprio, e não como simples meio para o fim dos outros. --- Este direito existe no homem. Porquanto, o homem, sendo um ente racional e tendo o \emph{dever} de conseguir o seu fim ultimo, deve ter tambem o \emph{direito} de exigir que os outros o reconheçam e respeitem como ente racional, creado para conseguir esse fim. Logo o homem tem direito á \emph{dignidade pessoal} \footnote{O direito á \emph{dignidade pessoal} é desconhecido e violado, quando se considera o homem não como um ente racional, dotado de fim proprio, mas como um simples \emph{meto} para os fins temporaes dos outros homens. D'ahi a injustiça da \emph{escravatura}. --- Digamos alguma coisa ácerca da \emph{escravatura}, suas principaes \emph{causas}, e \emph{abolição}.\par \emph{a}) \emph{Escravatura}. --- A \emph{escravatura} é a instituição, que, violando por um modo permanente a \emph{dignidade pessoal}, sujeita o homem ao dominio absoluto e universal de outro homem. O escravo não é \emph{pessoa}, é uma \emph{coisa}; tudo o que elle é, alma e corpo, pertence ao seu dono. Este é a unica lei do escravo, a sua patria, a sua cidade, o fim da sua vida, a regra do justo e do injusto. O escravo, por ser \emph{coisa} e não \emph{pessoa}, não tem \emph{direitos}; para elle não ha familia, nem matrimonio, nem paternidade, nem filhos, nem irmãos; é como um animal irracional, sujeito a todas as violencias e crueldade do dono.\par \emph{b}) \emph{Causas da escravatura}. --- As principaes causas da escravatura foram, sobretudo na antiguidade, o \emph{desprezo pelo homem} e \emph{o.desprezo pelo trabalho}. --- A primeira causa foi o \emph{desprezo pelo homem}; pois os antigos pensavam geralmente que os homens não pertencem todos á mesma especie, e que os escravos constituem uma raça vil, inferior a dos homens livres. Nem os maiores talentos evitaram este tão grande erro. \emph{Homero} diz na \emph{Odissea} que Jupiter tirou aos escravos metade do entendimento; \emph{Platão} acerescenta que o animo dos escravos não é são nem inteiro; \emph{Aristoteles} ensina, na \emph{Politica}, que alguns homens nascem para escravos. --- A segunda causa foi o \emph{desprezo pelo trabalho}. O homem esforçou-se sempre por se subtrahir ao jugo do trabalho, impondo-o \&os seres mais fracos. A mulher e os filhos foram os primeiros a trabalhar na vida domestica; --- em seguida os paes e os filhos cahiram sob o dominio do mais forte, e esta sujeição tornou-se um estado permanente. Assim o trabalho acabou por ser considerado indigno do homem livre, e só proprio dos seres mais fracos e menos nobres.\par \emph{c}) \emph{Abolição da escravatura}. --- A \emph{abolição da escravatura} constituiu sempre o escopo do ensino e dos esforços do Christianismo. Para alcançar o seu\pgn{702}generoso intento, a Religião Christã começou por ensinar --- a egualdade \emph{natural} entre todos os homens, judeus ou gregos, homens ou mulheres, livres ou escravos, --- e a necessidade do trabalho, que, além de ser uma \emph{expiação} e uma \emph{necessidade}, exerce benefica influencia sobre a nossa vida moral. E este ensino tornou-se pratico. A Egreja não só reconheceu a egualdade \emph{natural} entre todos os homens, tratando a todos do mesmo modo e elevando ás mais sublimes dignidades os livres e os escravos, mas tambem empregando os seus esforços junto dos senhores, para que íratassem os seus escravos como irmãos e lhes reconhecessem todos os \emph{direitos}, que são proprios do ente racional e immortal. --- Todos conhecem a influencia poderosa e benefica do Summo Pontifice Leão XIII na abolição da \emph{escravatura} no Brazil, e á sua cruzada em favor dos escravos da Africa.\par O direito \& \emph{dignidade pessoal} transgride-se tambem em muitas organizações do trabalho, quando se considera o trabalhador como uma continuação da maquina e se lhe negam os direitos essenciaes do homem. Assim o excesso nas horas do trabalho, o abuso no trabalho das mulheres e das crianças, a falta de cuidado relativamente á saude e á moralidade dos operarios, a obrigação do trabalho em dias santos, são outras tantas infracções do direito á \emph{dignidade pessoal}. --- Nem se objecte que existe um contracto voluntario entre os operarios e os patrões; porque o homem não póde renuneiar aos seus direitos \emph{innatos}, e por isso o contracto é nullo.\par Se a condição de \emph{escravo}, como foi praticada pelo paganismo, é contraria ao direito que todo o homem tem á \emph{dignidade pessoal}, a condição de \emph{servo} ou \emph{criado} não é contraria a esse direito, como veremos mais tarde.}.

\pg{702} c) \textbf{Direito} á vida, --- Direito Á \emph{vida} é o que o homem tem á conservação da sua existencia. --- Este direito existe no homem. Porquanto, o homem, recebendo de \textsc{Deus} a vida, recebeu tambem o \emph{dever} de a conservar, empregando os meios proporcionados e necessarios. Se o homem tem o \emph{dever} de conservar a vida, deve ter tambem o \emph{direito} respectivo; pois este direito é o meio proporcionado e necessario para o cumprimento d'aquelle dever. Logo o homem tem direito á \emph{vida} \footnote{O direito é \emph{vida é innato}, mas não é tão absoluto, que não possa ser suspenso deante de um direito superior. Assim como o \emph{dever} de conservar a vida, pelo facto de ser um \emph{meio} para o cumprimento da ordem moral, póde ceder deante de um dever superior, que é incompativel com elle, e cujo cumprimento é necessario para a realização da ordem moral (o que acontece com os martyres, e com os que expõem e sacrificam a propria vida para a salvação dos outros); assim tambem o direito \emph{d vida} póde e deve ceder deante de um direito mais elevado. --- As regras, que servem para o conhecimento da superioridade ou inferioridade dos \emph{direitos}, foram apresentadas no \emph{Prologo} d'este tratado.\par Este direito \emph{d vida} foi, na antiguidade, negado não só \emph{em pratica}, porque a vida dos \emph{escravos} dependia do capricho do dono, mas tambem \emph{em theoria}, porque \emph{Platão}, na sua \emph{Republica}, diz que devem ser condemnados á morte os que têm um corpo mal disposto, e \emph{Aristoteles}, na sua \emph{Politica}, ensina que, em alguns casos, o infanticidio deve ser prescripto e que o aborto é licito. --- Mas a Egreja Catholica, baseada na grande verdade de que só \textsc{Deus} é\pgn{703}Senhor da nossa vida e da nossa morte, luctou sempre em defeza da inviolabilidade da vida humana, condemnando o suicidio, o homicidio, o infanticidio, o aborto (Kurth, \emph{Les origines de la civilisation moderne}).}.

\pg{703} d) \textbf{Direito} de legitima defeza, --- Direito de \emph{legitima defeza} é o poder moral, que assiste o todo o homem, de empregar a força physica para repellir uma aggressão actual e injusta, quando não ha outro meio de defeza. --- Este direito existe no homem. Porquanto, quem possue um direito, possue tambem o direito de o defender com a força physica contra uma aggressão actual e injusta, quando não ha outro meio; pois o direito é naturalmente \emph{coactivo}. Ora o homem possue o direito á \emph{vida}. Logo possue tambem o direito de a defender, empregando a força physica, contra a actual e injusta aggressão, quando não ha outro meio \footnote{Todo o direito \emph{jurídico}, como é o \emph{innato}, é por si \emph{coactivo}. A coacção, como dissemos, deve ser geralmente exercida pela auctoridade publica; mas, quando não póde ser exercida pela auctoridade publica e o direito está exposto a uma injusta *aggressão, é claro que póde ser exercida pelo proprio sujeito do direito; aliás não haveria segurança na sociedade.\par O direito de \emph{legitima difesa} póde extender-se até á morte do injusto aggressor, quando isto é absolutamente necessario para a defeza da propria vida. Porquanto, neste caso, ha uma collisão entre o direito do injusto aggressor e o direito do innocente aggredido. Ora, sendo o direito do innocente aggredido moralmente mais forte que o direito do injusto aggressor, é claro que este deve ceder áquelle; pois são incompativeis. --- \emph{A/rens} não reconhece o direito de \emph{legitima defeza}, pois diz que o homicidio é mal moral, e este nunca é permittido. Mas erra; poisque o homicidio, no caso de \emph{legitima defesa}, deixa de ser mal moral, e por isso não é prohibido.\par D'onde se vé que o direito de \emph{legitima defesa} presuppõe as seguintes condições: --- 1.º) A aggressüo deve ser \emph{injusta}, isto é; contraria á lei. --- 2,3) A aggressao deve ser \emph{actual}, isto é, deve constituir um perigo \emph{certo} e \emph{imminente} para a \emph{pessoa} do aggredido; porque, se o nosso direito podesse ser defendido pelo auctoridade publica, ou se o perigo podesse evitar-se por outro meio, que não seja o damno e a morte do aggressor, cessa então o direito de \emph{legitima defeza}. --- 3.º) A defeza deve ser \emph{proporcionada} á aggressão. Porquanto o fim do direito de \emph{legitima defeza} é a conservação da vida do innocente. Ora é o fim que determina os meios. Logo os meios, que o aggredido deve empregar, não devem exceder os limites necessarios e sufficientes para a repulsão do ataque.\par O direito de \emph{legitima defesa}, assim limitado, póde exercer-se não só contra o aggressor da vida, mas tambem contra o aggressor da integridade dos nossos membros, do pudor e dos bens de fortuna; porque o respeito do direito do aggredido é um bem \emph{universal}, que se relaciona com a paz e a ordem commum, e por isso prevalece ao bem individual do aggressor.\par Todavia este direito não se extende ás aggressóes contra a propria \emph{fama} e \emph{honra}. Na verdade, a defeza, como dissemos, deve ser \emph{proporcionada} \& aggressão. Ora o uso da força physica não é o meio \emph{proporcionado} para a\pgn{704}repulsão dos ataques contra a \emph{fama} e a \emph{honra}; pois não é com a morte ou com o mal do aggressor que se repara a \emph{fama} e a \emph{honra}, mas só com a demonstração da falsidade do ataque contra esses bens. Além d'isso, o direito, que o homem tem á propria \emph{fama} e \emph{honra}, é inferior ao direito, que o aggressor da mesma \emph{fama} e \emph{honra} tem á propria \emph{vida}.\par Do que temos exposto resulta a illegitimidade do \emph{duello}.\par \emph{Duello} é um combate singular e privado, combinado por mutuo accordo entre duas pessoas, para resolver differencas particulares, com perigo de morte ou de lesão corporal, tendo precedido a determinação das armas, da hora, do sitio do combate, e de outras circumstancias e condições. --- O \emph{duello}, que é um resto da barbaridade da \emph{edade media}, é illicito pelas seguintes razões: --- 1.º) despreza a \textsc{DEus}, que não concedeu ao homem particular o dominio sobre a vida propria nem sobre a do proximo; --- 2.º) é uma violação do dever, que temos de conservar a vida; --- 3.º) é uma violação dos direitos do proximo, cujo bem temporal e eterno elle injustamente compromette; --- 4.3) despreza os direitos da sociedade, usurpando as funcções da auctoridade publica, criando odios, levando a viuvez e a orphandade ao seio das familias, etc.\par \emph{Benjamin Franklin} diz a respeito do \emph{duello} o seguinte: « O homem offende continuamente a honra de \textsc{DEus}, e \textsc{DEus} o conserva; o homem offende um outro homem, e este julga, no seu orgulho, que a offensa só póde ser expiada com' a morte! Um principe, quando condemna á morte um subdito, que o ultrajou, chama-se tyranno; e um particular, exercendo as funcções de juiz na propria causa e por auctoridade propria, não só condemna á morte o offensor, mas elle mesmo é o carrasco! v. --- O proprio \emph{Rousseau} escreve: « Deverá a honra de um homem honesto ser exposta ao capricho do primeiro louco furioso, que elle póde encontrar? Dir-se-ha que o \emph{dwello} prova a coragem do homem. E bastará isto para apagar a vergonha e o remorso de todo e qualquer facto? Mas, então, um malvado, para ser bom, não haverá a fazer outra coisa se não bater-se em .duello. As palavras de um mentiroso tornar-se-hão verdades, desde o momento que possam ser sustentadas pela espada; e, se tu fores accusado de ter assassinado um homem, bastará, para provar que não é \textbf{verdade}, quo mates um outro homem, D'este modo a virtude e o vicio, a honra e a' infamia, a verdade e a mentira, tudo depende do \textbf{exito} do um duello. Uma 5 elle. Uma sala d'armas é a séde de toda a justiça; o unico\par direito é a força; a unica razão é o homicidio; toda a reparação devida aos offendidos consiste em lhes dar a morte; e toda a defeza é egualmente lavada no sangue do offensor e do offendido... Se os lobos soubessem raciocinar, teriam maximas differentes!... O duello é um esforço sem virtude, um crime sem prazer, uma honra sem razão \emph{m(Nouvelle Hêl}.).\par Os defensores do \emph{duello} fazem as seguintes objecções:\par a) Dizem: O \emph{duello} é uma consequencia do direito da \emph{legitima defeza}, de que todo o homem é dotado. Logo o \emph{duello} é licito.\par Resposta. E' falso que o \emph{duello} é uma consequencia do direito de \emph{legitima}\pgn{705}\emph{defesa}. Porquanto o direito da \emph{legitima defesa}, para o seu exercicio, suppõe uma aggressão \emph{actual}, exige que não haja opportunidade de recorrer á auctoridade publica e que a defeza seja proporcionada á aggress\&o. Ora o \emph{duello} --- suppõe uma offensa \emph{passada}, --- representa um desprezo da auctoridade publica, de que usurpa os direitos e a que não se recorre, quando este recurso é possivel, --- não é \emph{proporcionado} á offensa, não só porque esta póde e deve reparar-se por outro meio, diverso do \emph{duello}, mas tambem porque a \emph{honra} não se repara com a morte do offendido nem com a morte do offensor.\par 5) Continuam: O \emph{duello} é uma necessidade para um homem honrado, quando é injustamente calumniado ou offendido.\par Resposta. As affrontas, como dissemos, reparam-se por outro meio, diverso do \emph{duello}. Se o \emph{duello} reparasse as affrontas, ninguem seria mais honrado do que um espadachim, que, apezar dos seus crimes, soubesse manejar com felicidade as armas, assim como ninguem seria mais infame do que um homem de bem, que, apezar da sua honestidade, não soubesse pegar numa espada. E' verdadeiramente um raciocinio de \emph{lobos} o dos adversarios!}.

e) \textbf{Direito} de propriedade, --- Direito de \emph{propriedade} é o poder moral de adquirir bens materiaes externos, e de usar e 1 dispor d'elles livre e exclusivamente. --- Este direito existe no homem. Porquanto o homem tem o \emph{dever} de conservar a \pg{704} \pg{705}  propria vida. Ora este \emph{dever} involve o \emph{direito} de adquirir bens materiaes externos e de usar e dispor d'elles livre e exclusivamente; porque esse \emph{direito} é um meio indispensavel para a nossa existencia presente e futura. Logo o homem tem o direito de \emph{propriedade} \footnote{A mais importante questão, que agita actualmente os individuos e a sociedade, é a \emph{questão social}. --- Reservando para mais tarde a solução d'este problema capital, aqui limitamo-nos a langar as bases d'essa solução, desenvolvendo a theoria ácerca do direito de \emph{propriedade}.\par a) Direito de propriedade. --- Direito de \emph{propriedade} é o poder moral de adquirir os bens materiaes externos e de usar e dispor d'elles livre e exclusivamente. --- Dissemos --- \emph{exclusivamente}, para indicar que o proprietario póde, no uso e disposição dos seus bens, excluir os outros individuos, embora o seu direito possa ter e tenha limitações moraes e \emph{jurídicas}.\par b) Materia do direito de propriedade. --- Como o direito de \emph{propriedade} tem por fim a satisfação das nossas necessidades por meio dos bens materiaes externos, e exclue os outros do uso e disposição do que é nosso, a \emph{materia} d'esse direito não é constituida por bens não sujeitos á apropriação, como são O ar, a \emph{agua do mar}, etc. mas é constituida por bens exgoteveis pelo uso e capazes de serem apropriados por um individuo, de modo que os outros sejam excluidos, como são as \emph{terras}, as \emph{producções}, etc.\par 0) Divisão do direito de propriedade. --- O direito de \emph{propriedade} divide-se principalmente em \emph{transitorio} e \emph{estavel}, \emph{perfeito} e \emph{imperfeito}, \emph{substancial} e \emph{util}, \emph{pessoal} e \emph{collectivo}.\par \emph{a}) Direito \emph{transitorio} e \emph{estavel}. --- E' \emph{transitorio}, quando se refere a bens, que duram pouco tempo, porque se destroem ou consomem pelo \emph{uso}; tal é o direito que temos aos \emph{vestidos}, aos \emph{alimentos}, etc. --- E' \emph{estavel}, quando se refere a bens, que duram muito tempo, porque não se destroem ou consomem pelo \emph{uso}; tal é o direito que temos a uma \emph{casa}, a um \emph{terreno}, etc.\par \emph{b}) Direito \emph{perfeito} e \emph{imperfeito}. --- E' \emph{perfeito}, quando o proprietario póde fazer da coisa possuida o uso que elle quizer, e dispor d'ella conforme lhe agradar. --- E' \emph{imperfeito}, no caso contrario.\par \emph{c}) Direito \emph{substancial} e \emph{util}, --- E' \emph{substancial}, ou \emph{directo}, quando se refere\par 45 Vor. II.\pgn{706}unicamente á \emph{substancia} da coisa. --- E' \emph{util}, quando se refere unicamente ao \emph{uso} ou ao \emph{usufructo} da coisa. --- Esta divisão é só propria do direito \emph{imperfeito}.\par \emph{d}) Direito \emph{pessoal}! e \emph{collectivo}. --- E' \emph{pessoal}, quando é possuido por uma só pessoa \emph{physica}. --- E' \emph{collectivo}, quando é possuido por uma pessoa \emph{moral}; tal é a \emph{familia}, o \emph{estado}, a Egreja, e, em geral, toda a collecção de individuos humanos.\par d) Existencia do direito de propriedade. --- Todos admittem que o individuo tem o direito de \emph{propriedade particular}, relativa aos bens \emph{transitorios} (porque, o homem, tendo direito á \emph{vida}, deve ter tambem o direito relativo aos meios, necessarios e indispensaveis para a sua conservação, que são os \emph{alimentos}, os \emph{vestidos}, etc.); mas nem todos admittem que o individuo possue verdadeiro direito de \emph{propriedade particular}, relativa aos bens \emph{estaveis} ou \emph{immoveis}. Mas este direito existe, como o mostram as seguintes razões.\par 4) O homem, devendo conservar a sua vida, deve procurar o qué é necessario para ess\&.conservagüo. Ora, para procurar o que é necessario para a conservação da sua vida, deve empregar um esforço e um trabalho constante e cuidadoso. Mas não poderia empregar esse esforco e esse trabalho, se não tivesse o direito de \emph{propriedade particular}, relativa aos bens \emph{estaveis}; porque ninguem enceta longos e pezados trabalhos, sem a certeza de que elle, e só elle, ha de perceber todo o fructo, --- certeza que inyolve a posse \emph{estavel} e \emph{esclusiva} dos bens materiaes. Logo o homem tem o direito de \emph{propriedade particular}, relativa aos bens \emph{estaveis},\par 5) O homem possue o direito \emph{innato} de servir-se dos bens materiaes e exteriores para a sua utilidade; pois esses bens foram creados para a conservação da sua vida. Ora o homem não poderia exercer ordenada e pacificamente tal direito, se não podesse adquirir e fazer seus os bens \emph{estaveis}. Porquanto, o homem não só deve prover ás suas necessidades presentes, mas tambem ás \emph{futuras}, que hão de pesar sobre elle no tempo do \emph{nverno}, da \emph{carestia}, da \emph{doença}, da \emph{velhice}. Mas esta providencia seria impossivel e a nossa vida seria constantemente atormentada pela incerteza do futuro, se o homem não podesse adquirir os bens \emph{estáveis}, e se não podesse excluir do uso d'elles os outros individuos. Logo.\par c) Além d'isso, o homem, tendo o \emph{direito} de formar a sociedade \emph{domestica}, tem tambem o \emph{dever} de procurar a educação corporal e espiritual dos seus filhos, e de esforcar-se para que estes, depois da sua morte, continuem a viver honrada e convenientemente. Ora, para isso, é indispensavel que o individuo possua bens \emph{proprios} e \emph{estaveis}, com o direito de excluir do uso d'elles os outros individuos, para que os possa transmittir aos seus legitimos descendentes. Logo o homem tem o direito de \emph{propriedade particular}, relativa aos bens \emph{estaveis}.\par \emph{d}) A historia de todos os povos confirma a verdade da nossa proposição. Com effeito, em todos os povos, mesmo nos mais barbaros e selvagens, encontra-se, numa forma mais ou menos perfeita, a propriedade \emph{particular} e \emph{estavel}. Ora este facto, tão geral e constante, só póde explicar-se na theoria, que reconhece no homem o direito da mesma \emph{propriedade}.\par Examinemos tres objecções contra o direito de \emph{propriedade particular}.\par 1.º) Dizem: A terra é o patrimonio de todos os homens. Logo quem\pgn{707}apropria a si uma parte d'ella, commette um roubo em prejuizo do genero humano e destroe a ordem estabelecida por \textsc{DEus}.\par Resposta. A terra póde dizer-se patrimonio de todos os homens --- em sentido \emph{negativo}, emquanto nenhuma parte d'ella foi assignada a esta ou \&quella pessoa, --- mas não em sentido \emph{positivo}, emquanto todos podem apropriar a si uma coisa, que uma outra pessoa fez sua pela occupação e pelo trabalho. Na verdade, \textsc{Deus}, tendo creado a terra inculta, criou ao mesmo tempo o direito de \emph{propriedade particular}; visto que, sem esse direito, faltaria todo o estimulo para a cultura da mesma terra. Por isso a propriedade \emph{particular}, longe de ser um roubo e destruir a ordem estabelecida por \textsc{Deus}, representa o exercicio de um direito natural e completa essa ordem.\par 2. 3) Continuam: O Estado tem o dever de prover a todas as necessidades do individuo, e por isso o direito de \emph{propriedade} só pagó. pertencer ao Estado e não ao individuo.\par Resposta. O individuo existe antes do Estado; pois o Estado é o conjuncto de individuos. Existindo antes do Estado, o individuo tem o \emph{dever}, proveniente da lei natural, de prover ás suas necessidades, presentes e futuras, e por isso tem o direito de \emph{propriedade particular}, derivado da mesma lei. Tal direito, por ser \emph{natural}, não póde ser destruido pelo Estado, mas deve ser garantido no seu exercicio pelo proprio Estado.\par 3.º) \emph{Bentham} objecta: Todo o direito é essencialmente \emph{coactévo}. Ora o individuo é destituido de força \emph{coactiva}. Logo o individuo não póde ter o direito de \emph{propriedade}.\par Resposta. O direito é \emph{coactivo}, não emquanto é constituido essencialmente pela \emph{coacção actual}, ou pela força \emph{physica}, mas emquanto póde ser defendido pela força physica, se o seu \emph{sujeito} poder e quizer. A força \emph{coactiva} é uma \emph{propriedade} do direito; de modo que o direito subsistiria, embora o individuo nunca podesse empregar a força physica para obrigar os outros ao respeito do mesmo direito.\par e) Origem do direito de propriedade. --- Esta questão está já resolvida na proposição antecedente. Mas, tratando-se de um assumpto tão importante, não é inutil a insistencia. --- Ha varias opiniões ácerca da \emph{origem} do direito de \emph{propriedade}. --- Uns, como \emph{Hobbes}, \emph{Montesquieu}, \emph{Kant}, \emph{Bentham}, dizem que o direito de \emph{propriedade} deriva, unica e exclusivamente, da \emph{let civil}. --- Outros, como \emph{Grocio}, \emph{Puffendorf}, \emph{Heinecio}, sustentam que esse direito deriva de um \emph{pacto primitivo}, pelo qual os homens, querendo eliminar os inconvenientes, que resultavam da propriedade \emph{collectiva} dos bens, estabeleceram a divisão e a propriedade \emph{particular} e \emph{estavel}. --- Outros, como \emph{Locke}, \emph{George}, dizem que o direito de \emph{propriedade} deriva do direito \emph{innato}, que o homem tem ao \emph{fructo do seu trabalho}. --- Outros, finalmente, ensinam que o direito de \emph{propriedade}, isto é, o poder moral de adquirir e fazer seu um bem \emph{estavel}, deriva da propria \emph{natureza}, ou da \emph{lei natural}. --- De cada uma d'estas opiniões faremos a critica nas proposições seguintes.\par a) O direito de \emph{propriedade} não deriva da \emph{lei civil}. Com effeito, a \emph{lei civil} n\&o póde explicar sufficientemente o facto, universal e constante, da propriedade \emph{particular}, não só entre os povos civilisados, mas tambem entre os selvagens, pelo menos em relação aos bens moveis. --- Além d'isso, o dominio\pgn{708} ou a propriedade \emph{particular} é geralmente necessaria a cada individuo e a cada familia. Ora os individuos e as familias são anteriores á sociedade e por isso á \emph{lei civil}. Logo tambem o dominio ou a propriedade \emph{particular} é anterior á \emph{lei civil}, \emph{e} por isso não deriva d'ella.\par \emph{b}) O direito de \emph{propriedade} não deriva de um \emph{pacto primitivo}. Porquanto, este \emph{pacto primitivo} é uma simples invenção; porque nenhum povo o recorda nas suas tradições. --- Além d'isso, esta theoria suppõe que todos os bens \emph{estaveis}, no principio, pertenciam \emph{positivamente} a toda a \emph{collecção} dos homens; de.modo que não eram proprios de nenhum individuo. Ora esta supposição é falsa; porque, nesse caso, teriamos um dos seguintes absurdos: --- ou todo o genero humano devia estar sujeito \emph{a} uma unica auctoridade suprema, que desse a cada um o necessario para a sustentação da vida, --- ou ninguem podia servir-se de coisa alguma, sem o consentimento de todos os outros. --- Finalmente, embora esse \emph{pacto primitivo} tivesse existido, os homens do nosso tempo poderiam rescindil-o, e assim o direito de \emph{propriedade} deixaria de existir. Nem digam que a ordem e a paz social nos obrigam a respeitar esse pacto; porque, então, admittem, em ultima analyse, --- que os homens são obrigados pela \emph{lei natural} a respeitar a propriedade particular e estavel, --- e que propriedade é tão necessaria, que se não póde abrogar, isto é, admittem que o direito de \emph{propriedade} se funda na propria lei \emph{natural}.\par \emph{c}) O direito de \emph{propriedade} não deriva, originaria e exclusivamente, do direito do homem ao \emph{fructo do seu trabalho}. Porquanto, o trabalho suppõe --- a \emph{terra}, que se cultiva, --- o objecto, em que se exerce a industria, --- os instrumentos e os utensilios, de que o homem se serve para o trabalho, etc. Ora, devendo ser \emph{nossa} a terra que se cultiva, \emph{nosso} o objecto em que se exerce a industria, \emph{nossos} os instrumentos e utensilios de que nos servimos, é claro que o \emph{fructo} do trabalho presuppõe o direito de \emph{propriedade}, e que por isso este direito não deriva exclusivamente do direito que o homem tem ao \emph{fructo do seu trabalho}. --- Além d'isso, se por \emph{trabalho} se entende o exercicio da actividade humana, emquanto transforma ou modifica as coisas materiaes, este trabalho não é meio sufficiente para a acquisição de \emph{todos} os bens materiaes; porque muitos d'esses bens, não estando sujeitos á modificação ou transformação, como são os \emph{bosques}, os \emph{prados}, etc. não podem ser apropriados. Se por \emph{trabalho} se entende toda e qualquer actividade, que se exerce com esforço e com fadiga, este trabalho, sendo uma norma demasiadamente vaga, subjectiva e mutavel, longe de ser um titulo sufficiente para a acquisição dos bens materiaes, é um meio para semear discordias no seio das familias e da sociedade.\par \emph{d}) O direito de \emph{propriedade} deriva, originaria e immediatamente, da propria \emph{natureza}, ou da lei \emph{natural}. Esta proposição prova-se com os mesmos argumentos, que adduzimos em favor da \emph{existencia} do direito de \emph{propriedade}. Porquanto, esses argumentos provam que o homem, tendo o \emph{dever}, imposto pela \emph{natureza}, ou pela lei \emph{natural}, de conservar e aperfeiçoar a propria vida e a dos seus filhos, deve ter tambem o \emph{direito}, proveniente da mesma \emph{natureza}, ou lei \emph{natural}, de empregar os meios indispensaveis para o cumprimento d'esse dever. Ora um d'esses meios indispensaveis é o direito de \emph{propriedade}, isto ó, o direito de possuir, livre'e exclusivamente, alguns bens materiaes e\pgn{709}externos. Cf. \emph{Sum}. \emph{Th}, 2-2, q. 66, aa. 1, 2. --- E' verdade que S. Thomaz faz derivar do \emph{direito das gentes} o direito da \emph{propriedade}; mas o \emph{direito das gentes}, na mente do S. Doutor, não é o direito \emph{meramente positivo}, introduzido por um paeto internacional, mas é uma conclusão deduzida do direito \emph{natural} e recebida pelo senso commum. Eis as suas palavras: « Ad \emph{jus gentium} pertinent ea, quae derivantur ex \emph{lege naturae}, sicut conclusiones ex principiis, ut justae emptiones, venditiones et alia hujusmodi, sine quibus homines ad invicem convivere non possunt; quod est de \emph{lege naturae}, quia homo est naturaliter animal sociabile » (\emph{Sum}. \emph{Th}, 1-2, q. 95, a. 4). --- A Encyeliea \emph{Rerum novarum} do S. Padre Leão XIII, 16 Maio 1891, estuda e resolve a \emph{questão social} por um modo exhauriente á luz da fé e da razão.\par e) Determinação do direito de propriedade. --- O direito de \emph{propriedade} póde considerar-se em \emph{abstracto} e \emph{em concreto}. Considerado \emph{em abstracto}, é o poder moral de adquirir um dominio particular e estavel, e este poder é concedido pela lei \emph{natural}; --- considerado \emph{em concreto}, é o poder de usar e dispor dos bens materiaes e externos, que lhe pertencem. --- O direito de propriedade em \emph{concreto} deriva do direito de propriedade em \emph{abstracto}, por meio de um \emph{facto}, emquanto este \emph{determina} ou limita a certos bens o poder que temos de adquirir e possuir. --- A ordem logica das ideas exige que tratemos agora do \emph{facto primitivo}, que \emph{determinou} ou tornou \emph{concreto} o direito de \emph{propriedade} em \emph{abstracto}. --- Este facto \emph{primitivo} n\&o poude ser um \emph{pacto}, nem um \emph{decreto}, nem o \emph{trabalho} só, mas foi a \emph{occupacdo}.\par \emph{a}) O facto determinativo do direito de \emph{propriedade} n\&o foi um \emph{pacto}. Por quanto, tal \emph{pacto} poderia conceber-se nalguns casos, quando, multiplicado o numero dos homens, tivessem distribuido entre si os bens das familias ou iribus. Mas esta distribuição supporia já o direito de \emph{propriedade}, existente nas familias ou nas tribus, e determinado por um facto precedente.\par \emph{b}) Não foi um \emph{decreto} d'algum principe. Porquanto, tal \emph{decreto} deveria suppór que, no principio, todas as coisas eram \emph{positivamente} communs, ou que os primeiros principes tinham jurisdicc\&o sobre toda a terra. Ora ambas estas supposições são gratuitas e inacceitaveis.\par c) O facto \emph{determinativo} do direito de \emph{propriedade} não foi exclusivamente o \emph{trabalho}. Porquanto, ha coisas, que n\&o precisam de ser transformadas pelo \emph{trabalho}, porque são preparadas pela propria natureza para uso do homem, como as \emph{fructas}, os \emph{animaes}, etc., e por isso, relativamente a estas coisas, o \emph{trabalho} não póde ser o \emph{facto determinativo} do direito de \emph{propriedade}. --- Além d'isso, quando alguem se apossa de uma coisa com a intenção de a transformar com o trabalho, esta coisa pertence-lhe desde o primeiro momento da occupação; aliás todos poderiam, sem injustiça, impedir o exercicio do \emph{trabalho}, e por isso o direito de \emph{propriedade}.\par \emph{d}) O facto \emph{determinativo} do direito de \emph{propriedade} foi a \emph{occupagdo}. --- \emph{Occupação} é a spprehensao das coisas, que não pertencem a ninguem, com a intenção de as reter como proprias. Para a \emph{occupação}, pois, é necessario que a coisa \emph{occupada} seja susceptivel de apprehensão e de posse exclusiva, e que o \emph{occupante}, além da \emph{intenção} de fazer sua uma coisa, occupe \emph{e/fectivamente} a mesma coisa, de modo que todos possam conhecer que lhe pertence. Ora a occupação é um facto sufficiente para determinar \emph{em} concreto o direito de\pgn{710}\emph{propriedade}. Na verdade, a \emph{occupação} não lesa o direito dos outros (pois trata-se de coisas, que não pertencem a ninguem), e manifesta claramente a vontade, que uma pessoa tem, de reter uma coisa como sua, para o seu uso exclusivo. --- E' por isso que todos os povos, mesmo selvagens, têm sempre considerado a \emph{occupação}, como titulo sufficiente de \emph{propriedade} em relação ás coisas, que não pertencem a ninguem,\par \emph{Ahrens} objecta: Se o direito de \emph{propriedade} se adquire pela \emph{occupação}, um só individuo poderia fazer seu o mundo inteiro. Ora isto é absurdo.\par Resposta. O individuo, que occupasse o mundo inteiro, não o tornava seu \emph{juridicamente}, mas \emph{arbitrariamente}. Não o tornava seu \emph{juridicamente}; porque o homem só tem direito ao que é necessario para a sua sustentação; ora para essa sustentação não é necessaria a posse do mundo inteiro. Mas tornava-o seu \emph{arbitrariamente}, e assim nunca teria o \emph{direito} de propriedade; pois esse direito não se funda no arbitrio proprio, mas na lei natural.\par Ao facto principal, que determinou o direito de \emph{propriedade} e que foi a \emph{occupação}, seguiu-se um outro facto, pelo qual os homens adquiriram ou augmentaram as suas propriedades. Este facto é a \emph{accessão}. --- Accessão, pois, é o augmento que uma coisa nossa adquire --- ou por beneficio da natureza, --- ou pela industria humana, --- ou pela obra unida da natureza e da arte. D'ahi a divisão da \emph{accessão} em \emph{natural}, \emph{artificial} e \emph{mista}.\par f) Transmissão do direito de propriedade. --- O direito de \emph{propriedade}, referindo-se á bens materiaes e externos, de que podemos dispor livre e exclusivamente, póde transmittir-se. Esta transmissão faz-se principalmente por \emph{successão ereditaria} ou por \emph{contracto}.\par a) \emph{Successão hereditaria}, --- \emph{Successão hereditaria} é o facto pelo qual uma pessoa entra na posse dos bens, que uma outra pessoa lhe deixou por uma disposição, declarada por \emph{testamento} ou por qualquer outro acto. --- Aqui provaremos --- que o proprietario póde transmittir os seus bens por meio do \emph{testamento}, --- e que os filhos succedem aos paes no direito da \emph{propriedade}, embora estes não tenham deixado \emph{testamento}.\par --- O proprietario póde transmittir os seus bens por meio do \emph{testamento}. Por quanto, o proprietario tem o direito sobre os seus bens, e por isso póde dispor d'elles, como quizer, transformando-os, alienando-os, ou mesmo destruindo-os. Ora, se o proprietario tem esse direito durante \& sua vida, deve tel-o tambem no tempo da sua morte; pois não ha razão nenhuma, que exija uma excepção para aquelle momento. Logo o proprietario póde transmittir os seus bens por meio do \emph{testamento}. --- E' claro que este direito póde ser legitimamente limitado pela lei \emph{natural} e \emph{positiva}.\par --- Os filhos succedem aos paes no direito de \emph{propriedade}, embora estes não tenham deixado \emph{testamento}, isto é, embora tenham fallecido \emph{a5-intestato}. Por quanto, o elemento do genero humano não é propriamente o individuo, mas sim a familia, a qual constitue uma verdadeira sociedade, cuja cabeça e director é o pae; e por isso o direito de \emph{propriedade}, embora pertença ao pae, reflecte-se sobre toda a familia. Ficando, pois, a familia, quando morre o pae, é claro que o direito de \emph{propriedade} deve passar para ella, á qual de algum modo já pertencia e em cuja vantagem era exercido pelo pae. --- Além d'isso, o pae adquire as propriedades sobretudo para os filhos, a cujas\pgn{711}necessidades deve prover. Logo, embora o pae morra sem deixar testamento (\emph{ab-intestato}), presuppõe-se, em virtude da lei natural, que tivesse intenção de deixar os bens aos seus filhos. --- Finalmente, o pae gravou .nos bens, que adquiriu, o cunho da propria personalidade. Perseverando, pois, a personalidade do pae nos filhos, que s\&o uma parte do pae, os bens d'este devem pertencer aquelles pelo proprio direito \emph{natural}.\par \emph{b}) \emph{Contracto}. --- \emph{Contracto} é um consentimento livre e expresso, pelo qual duas ou mais pessoas conferem e acceitam algum direito. --- Para a \emph{legitimidade do contracto} são necessarias as seguintes condições: conhecimento do \emph{objecto} do contracto, liberdade nos contrahentes, coexistencia do mutuo consentimento, possibilidade e honestidade do objecto. --- O \emph{contracto} divide-se principalmente-em \emph{egual} e \emph{desegual}. \emph{Egual}, que se chama tambem \emph{oneroso} e \emph{bilateral}, é o que obriga ambos os contrahentes, embora o objecto da obrigação possa ser diverso; --- \emph{desegual}, que se diz tambem \emph{gratuito} e \emph{unilateral}, . 6 o que obriga um só dos contrahentes, porque o outro não faz senão acceitar a offerta. --- O homem póde, por um contracto \emph{oneroso} ou \emph{gratuito}, transmittir a outro o direito de \emph{propriedade} sobre os seus bens; porque esse direito extende-se, como dissemos, a toda e qualquer disposição relativa aos bens possuidos. --- Todavia, o direito de transmittir os proprios bens, por meio do \emph{contracto}, póde ser e é limitado pela lei \emph{natural} e \emph{positiva}.\par E' necessario advertir que a \emph{occwupação}, a, \emph{accessdio}, a \emph{herança}, o \emph{contracto}, sendo factos contingentes, dão logar a direitos \emph{adquiridos}. Porquanto, o direito de \emph{propriedade} póde, como dissemos, considerar-se em \emph{abstracto} e em \emph{concreto}. --- Considerado em \emph{abstracto}, refere-se ao poder de adquirir e de dispor, livre e exclusivamente, das coisas adquiridas; e este direito é \emph{innato}. --- Considerado em \emph{concreto}, refere-se á coisa, que nos pertence; e este direito é \emph{adquirido}, porque se basêa num facto contingente.\par h) Limitação do direito de propriedade. --- O direito de \emph{propriedade}, embora derive da lei \emph{natural}, não é comtudo tão absoluto e independente, que não admitta \emph{limites}; pois não ha direito, entre as creaturas, que não seja limitado. --- Estes limites são \emph{moraes} e \emph{jurídicos}. Os limites \emph{moraes}, que não podem ser impostos pela coacção, nascem não só da natureza do direito de \emph{propriedade}, mas tambem dos deveres do homem para com \textsc{Deus}, para comsigo, e para com o proximo. Por isso o homem não póde abusar dos seus bens, nem destruil-os inutilmente, mas deve empregal-os para a subsistencia propria e da sua familia, deve soccorrer aos pobres e contribuir para a felicidade social. --- Os limites \emph{jurídicos}, que podem ser impostos pela coacção, derivam dos \emph{direitos} do proximo e da sociedade, e têm por fim a garantia d'estes mesmos direitos, contra os abusos do exercicio do direito de \emph{propriedade}. Por isso o Estado póde exigir os impostos, póde fazer expropriações forçosas, póde, numa palavra, limitar o direito de \emph{propriedade}, como e quando o exigir o bem social ou a conservação da sociedade.\par O direito de \emph{propriedade} é contestado por dois systemas, que são o \emph{communismo} e o \emph{socialismo}. --- Digamos alguma coisa ácerca d'estes systemas,\pgn{712}que se alastram por toda a parte, minando os fundamentos da sociedade e ameaçando os mais funestos cataclysmos.\par I. Communismo. --- \emph{Communismo} é o systema que rejeita o direito de propriedade \emph{particular}, como contrario á lei \emph{natural}, e admitte que todos os bens devem ser \emph{communs} a todos os homens. --- Foi defendido principalmente pelos \emph{Manicheus}, \emph{Albigenses}, \emph{Anabaptistas}, etc.\par II. Divisão do communismo. --- O \emph{communismo} divide-se em \emph{negativo} e \emph{positivo}. --- O \emph{negativo} rejeita todo o direito de \emph{propriedade} e sustenta que todos os bens devem ser \emph{negativamente} communs, de modo que todo o individuo possa servir-se do que lhe agradar. O \emph{positivo} ensina que todos os bens devem ser proprios não dos individuos, mas de alguma communidade. --- O communismo \emph{positivo} subdivide-se em \emph{absoluto} e \emph{moderado}. O \emph{absoluto} sustenta que \emph{todos} os bens, quer \emph{productivos} quer de \emph{consumo}, devem ser proprios de alguma communidade, e que a producção e distribuição d'esses bens devem ser feitas de modo que tudo, quanto é possivel, seja commum, como os \emph{jantares}, os \emph{dormitorios}, etc. \emph{O moderado} afirma que só os bens \emph{productivos} (\emph{vinhas}, \emph{olivaes}, etc.) devem ser communs. --- O communismo \emph{moderado} divide-se ainda em \emph{anarchismo} e \emph{socialismo}. O \emph{anarchismo} affirma que todos os bens \emph{productivos} devem passar para o dominio inalienavel dos \emph{municipios} ou das \emph{sociedades operaréas}, independentes entre si, e que tal estado de coisas se deve realisar por meio da força. O \emph{socialismo} sustenta que todos os bens \emph{productivos} devem passar para o dominio de toda a \emph{sociedade civil}, e que tal estado de coisas se deve realisar por meio do suffragio popular. --- Digamos alguma coisa ácerca do \emph{socialismo}.\par III. Socialismo. --- \emph{Socialismo} é o systema economico-politico, que ensina --- que todos os bens \emph{productivos} devem passar para a sociedade civil --- e que a produecáo e distribuição dos mesmos bens deve ser feita pela mesma sociedade, constituida em forma \emph{democratica}. --- O \emph{socialismo} costuma chamar-se tambem \emph{collectivismo}.\par IV. Divisão do socialismo. --- O \emph{socialismo} divide-se em \emph{absoluto} e \emph{moderado}. --- O socialismo \emph{absoluto} sustenta que \emph{todos} os bens \emph{productivos}, \emph{immoveis} e \emph{moveis}, devem passar para o dominio da sociedade civil. Este systema foi sustentado principalmente por \emph{Saint-Simon} (1160-1825), \emph{Carlos Fourier} (1712-1881), \emph{Proudhon} (1809-1865), que dizia: « A \emph{propriedade é um roubo n}, \emph{Karl Mare} (1818-1883), que o desenvolve no livro « \emph{O Capital} ». --- O socialismo \emph{moderado}, que se chama tambem \emph{agrario}, ensina que, entre os bens \emph{productivos}, só os \emph{immoveis} devem passar para o dominio da sociedade civil. Este systema, esbocado anteriormente por \emph{Laveleye} na Belgica e por \emph{Stuart Mill} na Inglaterra, foi aperfeiçoado por \emph{George} na America, por \emph{Flürscheim} na Allemanha, e por \emph{Hertzka} na Austria.\par \emph{V}. Doutrina do socialismo absoluto. --- A doutrina do \emph{socialismo absoluto} póde resumir-se nos seguintes pontos:\par O estado da sociedade actual não poderia ser mais desgracado. Favorecida pelo influxo do individualismo liberal, a \emph{anarchia da producção} domina soberana. Todos podem fazer, na ordem economica, o que querem. D'ahi o ardor da \emph{concorrencia} e á lucta universal pela \emph{existencia}, em que os mais fracos perecem continuamente, e as riquezas vão-se accumulando cada vez mais\pgn{713}nas mãos de poucos. D'este modo, desapparecidas as classes medias, ficàm apenas duas classes extremas, cheias de mutuo odio: uma composta por poucos riquissimos, que nadam na abundancia, outra constituida por uma multidão de proletarios, que vivem na miseria. --- A esta anarchia de producção accrescentam-se outras causas d'este immenso desequilibrio social. Ás vezes, na mesma cidade, abrem-se muitas fabricas, cujo numero é superior ás exigencias e á exportação, e produzem-se muitos generos, que se não podem vender, D'ahi a ruina de muitas fabricas, as crises economicas, e por isso o licenciamento de operarios, lançados na incerteza da vida e na miseria.\par Esta condição de coisas não só é desgraçada, mas é tambem \emph{injusta}. Os \emph{capitalistas}, tornando-se ricos emquanto possuem exclusivamente os meios de producção, como são as terras, as fabricas, os caminhos de ferro, etc., despojam injustamente os operarios de uma parte do fructo do trabalho.\par Esta injustiça apparece claramente, se considerarmos a \emph{theoria do valor}.\par O \emph{valor} de uma coisa é duplice: \emph{usual} e \emph{comnutativo}. O valor é \emph{usual}, quando a coisa serve para o uso \emph{pessoal} do homem; --- 6 \emph{commutativo}, quando a coisa serve para o uso \emph{communm}, isto é, quando póde ser trocada por outras mercadorias. O valor \emph{usual} é constituido pelas \emph{propriedades} physicas e chimicas da coisa; --- o \emph{commutativo} é constituido exclusivamente pelo \emph{trabalho humano}, de modo que o valor \emph{comnmutativo} de uma coisa póde e deve medir-se pelo trabalho, que o homem desenvolveu na producção da propria coisa. --- O valor \emph{commutativo} da coisa é o trabalho humano crystallisado. Na verdade, podem commutar-se generos diversissimos entre si. Logo o valor \emph{commutativo} deve ser um elemento \emph{communm} a todos elles. Este elemento não póde ser constituido pelas propriedades physicas e chimicas; porque estas propriedades são causa de \emph{di/ferenca} entre as coisas, e constituem o valor \emph{usual}. Logo deve ser o \emph{producto} do \emph{trabalho humano}, não d'este ou d'aquelle trabalho \emph{em concreto}, mas do trabalho \emph{em geral} ou \emph{em abstracto}. Logo uma coisa tem o seu valor \emph{commutativo}, emquanto contém o \emph{trabalho humano}, como tal. O trabalho humano, embora \emph{abstracto}, mede-se pelo tempo, que é ordinariamente necessario para a sua producção.\par Esta theoria do \emph{valor} mostra o modo por que os capitalistas vão augmentando de dia para dia a sua fortuna. As forças, que o operario alluga ao capitalista, tem o seu valor \emph{commutativo} e \emph{usual}. O valor \emph{commoutativo} das forças mede-se pelo trabalho, que é necessario para a conservação e restauração d'ellas e que é egual á somma das coisas, de que o operario precisa para a sustentação; --- o valor \emph{usual} das forças consiste em que augmentam o valor \emph{commutativo} das coisas, nas quaes se exercem, de modo que o valor \emph{commutativo} das coisas chega a ser superior ao valor \emph{commutativo} das forças. Todavia o operario, pela mais flagrante injustiça, não recebe a recompensa d'este excesso. Supponhamos que o operario precise de cinco tostóes para a sua sustentação. Esta quantia representa o valor \emph{commwutativo} das forças do operario, ou a somma que o capitalista deve ao operario como recompensa de um dia de trabalho. Mas, se essas forças, que o operario alluga pelo espaço de um dia ao capitalista, chegarem a dar ás coisas, em que se exercem, um valor de dez tostóes, o capitalista não póde, sem injustiça, ficar com os outros cinco tostóes, que representam o fructo do operario; porque, se é apreciado\pgn{714}o valor commutativo das forças humanas, não é apreciado o seu valor \emph{usual}, e o capitalista commette um roubó em damno do operario.\par Semelhante processo, que diminue cada vez mais o numero dos capitalistas e augmenta o dos proletarios, deve crear finalmente, como resultado fatal mas logico, um estado de violencia e de guerra, em quo os capitalistas serão expropriados por todo o povo.\par Como remedio para tão grandes males, é necessario instituir uma nova ordem social, em que sejam devolvidos para o dominio inalienavel do Estado todos os bens productivos, como terras, casas, machinas, etc. O Estado deve ser democratico; nada de privilegios; tudo deve ser egual. O povo faz as leis; escolhe, por meio do suffragio universal, os magistrados, e julga-os nos seus actos. O officio dos magistrados deve consistir em \emph{ordenar a producção publica}; e por isso deve determinar, depois de um diligente exame, a qualidade e quantidade --- não só de cada genero de producção, --- mas tambem do trabalho, que deve ser egual para todos na duração. Ninguem poderá eximir-se do trabalho commum. Todos terão os mesmos direitos e os mesmos deveres. --- Não sô a \emph{producção}, mas tambem a \emph{circulação} e \& \emph{distribuição} dos bens serão confiados aos magistrados. Estes --- ou darão a cada um parte dos productos, conforme o trabalho feito ou o estado de indigencia, --- ou passarão uma certidão do trabalho feito ou do estado de indigencia, de modo que os operarios possam reclamar o que lhes for devido ou necessario. --- A propriedade particular, pois, longe de ser destruida, fica estabelecida na mais firme base, que é o trabalho (Cathrein, \emph{Philosophia Moral}).\par VI. Doutrina do socialismo moderado. --- O socialismo \emph{moderado}, ou \emph{agrario}, sustenta, como dissemos, que todos os bens \emph{productivos} e \emph{immoveis} devem passar para o dominio da sociedade civil, ou do Estado. Os defensores d'este systema, pretendendo demonstrar que a propriedade \emph{particular} das terras é contraria á natureza, dizem --- que, no principio, a posse das terras era commum, --- que o \emph{trabalho} é o unico titulo justo da propriedade \emph{particular} e a terra não é producto do trabalho, --- que todos têm direito á \emph{vida}, e por isso todos devem ter direito a uma porção egual de terra.\par VIX. Critica do communismo. --- O \emph{communismo}, embora não repugne á natureza humana, considerada em \emph{abstracto}, isto é, subordinada á recta razão e isenta das paixões desordenadas, repugna á nossa natureza, considerada em \emph{concreto}, isto é, cercada de fraqueza e de miserias. Porquanto:\par a) \emph{Baséa-se num falso principio}. --- Os \emph{communistas} partem do principio de que a propriedade \emph{particular} repugna á lei \emph{natural}. Ora este principio é falso. Porquanto, não repugna á lei natural o que não repugna --- nem ás relações do homem para com \textsc{DEus}, --- nem ás suas relações para com os bens materiaes, --- nem ás suas relações para com o proximo. Ora a propriedade particular --- não repugna ás relações do homem para com \textsc{Deus}, --- nem ás suas relações para com os bens materiaes, --- nem ás suas relações para com o proximo. --- Não repugna ás suas relações para com \textsc{DEus}; o que é evidente. --- Não repugna ás suas relações para com os bens materiaes; não só porque estes foram creados para uso do homem, mas tambem porque são indifferentes em relação aos seus possuidores. --- Não repugna ás suas relações para com o proximo; porquanto --- nem a propriedade particular torna por si\pgn{715}impossivel o exercicio dos direitos \emph{innatos} dos outros, --- nem os bens materiaes são \emph{positivamente} communs a todos os homens, --- nem a desegualdade, que deriva da propriedade particular, é contraria ao Direito natural, pois essa desegualdade é causa de progresso na ordem physica e moral.\par 5) \emph{Repugna d tendencia natural do homem}. --- Existe no homem uma tendencia natural --- para o livre desenvolvimento das suas forças, --- para a posse livre e exclusiva dos bens materiaes, --- para a propria dignidade, independencia e liberdade. Ora o \emph{communismo} repugna a esta triplice tendencia do homem. Repugna á tendencia para o livre desenvolvimento das forças; pois obriga egualmente todos os homens aos mesmos trabalhos. Repugna á tendencia para a posse livre e exclusiva dos bens materiaes; o que é evidente. Repugna á tendencia para a dignidade, independencia e liberdade pessoal; porque obriga, mesmo com a força, todos os homens a trabalhar para a utilidade dos outros, ruduzindo-os á condição de servos.\par \emph{c}) \emph{Não dá uma compensação sufficiente}. --- O \emph{communismo} diz que o Estado deve distribuir aos cidadãos, como compensação, uma parte do fructo, proporcionada ás suas necessidades e ao seu trabalho. Esta compensação poderia admittir-se, se os superiores e os subditos procedessem sempre com justiça e com consciencia. Ora um procedimento justo e consciencioso é tão difficil, que se não póde basear sobre elle o interesse da sociedade.\par \textbf{VIII}. Critica do socialismo absoluto. --- Este socialismo baséa-se em falsos fundamentos, e, no estado actual da sociedade, é impossivel e absurdo.\par a) \emph{O socialismo absoluto baséa-se em falsos fundamentos}. --- Porquanto, o socialismo \emph{absoluto} baséóa-se --- a) na negação de \textsc{Deus}, da espiritualidade da alma humana, da existencia da vida futura, --- 0) na theoria \emph{materialista} da evolução, --- \emph{c}) na egualdade de direitos em todos os homens, --- \emph{d}) na theoria do \emph{valor} de \emph{Marz}, --- e) na supposicio de que a propriedade particular deve necessariamente dividir a sociedade em duas classes inimigas: uma composta de poucos riquissimos, e outra constituida por uma multidão immensa de mercenarios. Ora estes fundamentos são todos falsos.\par a) E' falso que \textsc{Deus} não existe, que a nossa alma não é espiritual nem immortal, que não existe a vida futura, e que por isso o fim unico e supremo da vida humana é o goso material. Não nos demoramos na refutação d'estes erros, pois que foram sufficientemente refutados nos respectivos logares.\par \emph{b}) E' falsa a theoria \emph{materialista} da \emph{evolução}. Na verdade, esta theoria póde ter dois sentidos, --- ou que tudo está sujeito a uma perpetua mudança, e por isso cada edade forma novas ideas, relativas á moral, á politica, \& Religião e tambem á economia social, --- ou que a evolução economica tende sempre para a \emph{concentração}, de modo que, assim como muitas coisas, que eram administradas pelos individuos, passaram depois para as companhias e para o Estado, assim tambem a producção e a distribuição dos bens devam centralisar-se fatalmente na mãos da sociedade civil. Tanto no primeiro sentido, como no segundo, esta theoria é falsa. --- E' falsa no primeiro sentido, que é o de \emph{Marc}, discipulo de \emph{Hegel}; porque, como demonstrámos, existem verdades necessarias e immutaveis, que o proprio \emph{Marx}, não sabemo; com quanta coherencia, suppõe nos seus raciocinios. --- E' falsa no segundo sentido; nao só porque, como veremos, é impossivel que todos os\pgn{716}bens productivos passem para o dominio da sociedade, mas tambem porque a centralisação não exclue necessariamente a propriedade particular.\par c) E' falsa a theoria da \emph{egualdade} de direitos em \emph{todos} os homens. Na verdade, os direitos, se são eguaes em todos os homens, \emph{considerados em abstracto}, isto é, na sua natureza especifica, não são eguaes em \emph{todos} os homens, \emph{considerados em concreto}, isto é, nas suas condições individuaes. Esta diversidade não só se extende aos direitos \emph{adquiridos}, cuja desegualdade depende da desegualdade dos talentos, da actividade, da força, da prudencia, etc. mas tambem aos direitos \emph{naturaes}, cuja desegualdade depende da propria natureza (assim os direitos do pae não são eguaes aos direitos do filho, nem os direitos dos amos são eguaes aos dos servos). Por isso a egualdade ' \emph{absoluta} dos direitos envolve a destruição da familia e da sociedade.\par \emph{d}) E' falsa a theoria do \emph{valor}, inventada por \emph{Marx}. Porquanto, Marx suppõe que o valor \emph{commutativo} de uma coisa é constituido, unica e exclusivamente, pelo trabalho; de modo que o valor está na razão \emph{directa} do trabalho. Ora esta supposição é falsa; porque entre o valor e o trabalho ha muitas vezes grande desproporção; assim de duas vinhas, em que se gastou o mesmo trabalho, uma, bem situada, dá um fructo mais abundante e mais apreciado do que a outra, plantada num terreno ingrato. O elemento, pois, que constitue o valor \emph{commutativo} não é o trabalho, mas é \emph{principalmente} a \emph{utilidade} ou o valor \emph{usual}. Dissemos --- \emph{principalmente}; porque não negamos que a \emph{raridade} concorre tambem para o valor da coisa. --- O argumento, que \emph{Mara} adduzem favor da sua opinião, é falso. Porquanto, duas mercadorias, que se trocam, devem ter um elemento commum. Ora esse elemento commum não póde deixar de ser sobretudo a sua \emph{utilidade}; porque o valor de uma coisa é maior ou menor, conforme o grau da necessidade, em que o homem se encontra, e da aptidão, que a coisa tem para remediar aquella necessidade. Esta aptidão, por sua vez, depende das qualidades chimicas e physicas da coisa, das quaes por isso não podemos nem devemos prescindir na determinação do valor da propria coisa.\par \emph{e}) Concedemos que a propriedade \emph{particular} divide a sociedade em duas classes: uma de ricos e outra de pobres; mas negamos que estas devem ser oppostas, porquanto a desegualdade, se estiver contida nos limites marcados pelas leis da justiça e da caridade, concorre para o exercicio das virtudes Sociaes e para a cultura das sciencias e das bellas artes.\par b) \emph{O socialismo absoluto}, \emph{considerado o estado actual da sociedade}, \emph{á impossivel e absurdo}. Na verdade:\par a) E' difficil, para não dizer impossivel, extremar praticamente os bens \emph{productivos} dos de \emph{consumo}. Com effeito, muitas coisas ha que, pela diversidade do seu fim, podem contar-se tanto entre os bens \emph{productivos}, como entre os de \emph{consumo}. Logo --- ou devem passar para o dominio da communidade todos os bens, --- ou, se ficarem entregues aos particulares, deve empregar-se uma constante vigilancia para que não sejam fontes de producção.\par \emph{b}) E' tambem difficil a determinação da \emph{qualidade} e da \emph{quantidade} das coisas productiveis, exigida pelos socialistas, com o fim de abolir a \emph{anarchia} da \emph{producção}. Porquanto, na opinião dos socialistas, cada familia deve manifestar aos officiaes da sociedade a lista de todas as coisas necessarias\pgn{717}para a sua conservação; para que, reunidas todas estas listas, um comicio dos delegados do povo determine a somma das coisas, que se devem produzir. Ora esta operação é immensamente complexa, e exige uma multidão de escrivães e de escripturas. --- Se disserem que o governo, sem inquirição, póde determinar a quantidade da producção, responderemos que esta determinação não será exacta, e haverá muita producção desnecessaria.\par c) Muito mais difficil é a organisação do \emph{trabalho publico}. Na verdade, para que a assembleia suprema do Estado possa impor a cada districto, cidade, villa ou aldéa, uma determinada quantidade de trabalho, deve conhecer o numero dos homens aptos para o trabalho durante um certo espaço de tempo; porque é absurdo impor a toda a aldóa, villa, cidade ou districto, a mesma quantidade de trabalho, sem ter em conta o numero dos operarios presentes. Posto isto, pergunta-se: Poderá um cidadão do estado socialista mudar de domicilio, ou não? Se não póde, sem previa licença dos magistrados, a liberdade individual fica muito reduzida. Se póde, cada um abandonará uma terra ingrata, onde nada o prende (pois, na theoria socialista, tudo é commum, casas, terras, etc.) e escolherá os logares mais salubres e os terrenos mais ferteis, e assim a organisacüo do trabalho será impossivel.\par \emph{d}) Supponhamos, porém, que esteja averiguado o numero dos operarios, que habitualmente residem num logar. E' necessario \emph{distribuir} entre elles o trabalho, de modo que um seja applicado aos trabalhos das terras, outro aos da cidade, --- um cuide da limpeza das ruas e outro da assistencia dos doentes, etc. --- Ora pergunta-se: A escolha da qualidade do trabalho depende da vontade dos magistrados ou do arbitrio dos cidadãos? Se depende do arbitrio dos cidadãos, cada um escolherá o trabalho mais facil e decoroso, deixando para os outros o mais ingrato e abjecto; d'ahi a desordem. Se depende da vontade dos magistrados, os livres cidadãos tornar-se-hão verdadeiros escravos; d'onde a revolução social.\par e) A difficuldade augmenta, se considerarmos a \emph{distribuição} das coisas, que foram produzidas em commum. Porquanto a \emph{distribuição} deve ter uma norma. Esta norma póde basear-se num dos seguintes criterios: --- ou no \emph{numero} das pessoas, --- ou no \emph{tempo} do trabalho, --- ou na \emph{qualidade} e \emph{quantidade} do trabalho --- ou na \emph{applicação}, --- ou na \emph{éndigencia}. Ora cada um d'estes criterios é inacceitavel. --- E' inaeceitavel o criterio, que basêa a \emph{distribuição} no \emph{numero} das pessoas; porque, sendo diversas as aptidões, as forças, o cuidado, é injusto dar a cada um a mesma quantidade do fructo accumulado. --- E' inacceitavel o criterio do \emph{tempo} do trabalho; porque, então, deveriam receber a mesma recompensa os que foram diligentes ou emprehenderam obras difficeis e perigosas e os que foram preguiçosos ou fizeram trabalhos faceis e seguros, o que é summamente injusto. --- E inacceitavel o criterio da \emph{qualidade} e \emph{quantidade} do trabalho; não só porque a determinação d'essa \emph{qualidade} e \emph{quantidade} é muito difficil e póde depender de uma norma subjectiva e injusta, mas tambem porque é injusto não dar ou dar uma muito pequena recompensa ao operario, que, apesar da sug diligencia e applicação, não tem força nem talento para o trabalho. --- E' inacceitavel o criterio da \emph{diligencia}; porque, para a averiguação da \emph{diligencia} dos operarios, seria necessario collocar ao lado de cada operario um fiscal\pgn{718}consciencioso. Mas onde encontrar esses fiscaes? --- E' inacceitavel finalmente o criterio da \emph{indigencia}; porque, para se conhecer o estado de pobreza de cada familia, seria necessario constituir um conselho. Mas quem não và que este conselho não poderia cumprir com o seu dever, sem suscitar contradieções, discordias, invejas, capazes de revoltar a sociedade inteira?\par \emph{f}) Mas não é tudo. A justa distribuição do trabalho e dos fructos suppõe --- nos magistrados uma rara prudencia, uma probidade excepcional, e uma abnegação generosa, --- e nos subditos uma obediencia cega, uma resignação humilde, uma constante mortificação do proprio egoismo. Mas estas virtudes, tão raras mesmo entre os christios (embora a nossa Religião proponha tantos e tão poderosos motivos para a pratica da virtude e para a fuga do vicio) serão possiveis no systema socialista, que nega a existencia de \textsc{Deus}, a espiritualidade e immortalidade da alma humana, a differença entre a virtude e o vicio, a existencia de uma vida futura? Não, com certeza.\par \emph{A}) Ha mais. Se se traduzisse em pratica a theoria socialista, faltaria todo o estimulo para o trabalho. Porquanto, o que leva o homem ao trabalho é o desejo de procurar o que é necessario para a sua sustentação presente e futura, e é a certeza de gosar exclusivamente o fructo do seu trabalho. Ora estes dois estimulos para o trabalho não existem na theoria socialista. --- Nào existe o primeiro estimulo, que é o desejo de procurar o que é necessario para a propria sustentação presente e futura; pois, nesse systema, o homem, quando tem saude, tem direito á recompensa do seu trabalho quotidiano, e, quando está doente, fica entregue aos cuidados da sociedade. --- Não existe o segundo estimulo, que é a certeza de gosar exclusivamente o fructo do proprio trabalho; pois, para esta certeza, deveria cada um estar convencido de que os outros empregaram os. mesmos esforços e de que a recompensa corresponde ao trabalho: ora tal convicção é impossivel.\par \emph{i}) O systema socialista não só não promove o trabalho, mas chega a impedir o progresso das artes e das sciencias. Porquanto, se todos os cidadãos fossem obrigados ao trabalho, faltaria o tempo necessario para o estudo, sobretudo das materias ou sciencias mais difficeis, que exigem uma applicação diuturna e constante. --- Os socialistas dizem que se póde empregar no estudo o tempo, que cresce do trabalho. Mas esta resposta não é sufficiente. Porquanto, o homem --- ou recebe unicamente a recompensa do trabalho, que elle emprehende por ordem do Estado, e não a do trabalho, que lhe dá o estudo das sciencias e das artes, --- ou recebe tambem esta recompensa. Na primeira hypothese, faltará o estimulo para o estudo das sciencias e das bellas artes, e assim não haverá progresso. Na segunda, teremos desegualdade nos bens e nos direitos, e assim ficará destruido o ponto capital do socialismo, que é a perfeita egualdade entre todos os cidadãos.\par . IX. Objecções dos defensores do socialismo absoluto. --- Os defensores do socialismo \emph{absoluto} fazem as seguintes objecções.\par a) Dizem: O socialismo tem por fim melhorar a condição do operario; pois a miseria das classes indigentes não se deve attribuir tanto á falta de productos, como á sua má repartição, que se funda na propriedade particular e no capital. Ora um systema tão philantropico deve ser acceite.\par Resposta. O socialismo, longe de melhorar a condição do operario, tornal-a-hia\pgn{719} mais miseravel. Porquanto, no systema da propriedade \emph{particular}, o operario póde dispor, livre e exclusivamente, do salario, e, se fizer algumas economias, póde comprar terrenos e assim melhorar a sua condição; e este caso dá-se com todos os operarios honestos e laboriosos, que enriquecem, adquirindo fundos, que não são outra coisa senão o salario transformado. Ora, no systema socialista, que supprime a propriedade \emph{particular}, o operario não póde melhorar a sua condição, não póde de trabalhador tornar-se proprietario; pois que só póde gastar o seu salario em objectos de consumo, em alimentos, em vestidos, em divertimentos, fontes perennes de desmoralisação e de miseria. Eis o amor dos socialistas pelos trabalhadores!\par \emph{b}) Continuam: O capital é fructo de injustiça. Porquanto, o augmento do capital deriva do augmento do valor, que a materia bruta recebeu por meio do trabalho, e por isso esse augmento não pertence ao capitalista, mas sim ao trabalhador.\par Resposta. O capital não é fructo de injustiça. O seu augmento não depende só do trabalho; porque o trabalho suppõe a materia bruta, e esta adquire-se com o capital já existente. Por isso, se o operario tem direito ao fructo do seu trabalho, tambem o proprietario deve ter parte nos beneficios da industria. E, pois, injusto exigir que o proprietario entregue aos operarios todo o augmento de valor, que a materia bruta adquiriu pelo trabalho, tanto mais que o salario é uma coisa presente e certa, e o beneficio da industria é futuro, incerto e sujeito a riscos e perdas. --- Em poucas palavras: O capital não é fructo de injustiça, porque o proprietario comprou a materia bruta e deu ao operario o salario proporcionado ao trabalho.\par c) Insistem. O capital é essencialmente immoral; porque leva o proprietario a diminuir o mais possivel o salario e por isso reduz á miseria uma innumeravel multidão de pessoas.\par Resposta. Este perigo, que é real, póde e deve ser afastado por meios legitimos e proporcionados. A auctoridade religiosa e a civil devem intervir nesta importantissima questão social, impedindo a rebelli\&o dos operarios. e o despotismo dos capitalistas. O Christianismo resolveu, ha muitos seculos, a grande \emph{questão social}. \emph{Emilio de Laveleye}, depois de examinar as theorias de \emph{Karl Marx}, diz o seguinte: « Como o Christianismo, considerado apenas sob o ponto de vista de uma forma social, é superior a todos estes systemas, onde umas vezes falta a justa apreciação da realidade, outras vezes a verdadeira caridade! No Evangelho reina em toda a parte uma grande ternura pelos desherdados, e, ao mesmo tempo, um sentimento sublime de justiça social. A verdade essencial, que resulta de todos os ensinamentos de Christo, é que nenhum melhoramento é possivel, se primeiramente não se tornar melhor o proprio homem. A renovação moral: eis a fonte de todo o progresso verdadeiro. Não é pela critica das doutrinas economicas, por mais subtil que seja, nem por uma nova forma de associação, que hão de curar-se os males da sociedade actual » (\emph{Le socialisme contemporain}, pag. 46). --- Mais tarde voltaremos ao assumpto.\par X. Critica do socialismo moderado. --- O socialismo \emph{moderado}, ou \emph{agrario}, que ensina o devolvimento inteiro e exclusivo de todas as terras ao Estado, é falso. Na verdade, ha tres hypotheses possiveis para explicar o modo por\pgn{720}que todas as terras podem ser possuidas e administradas por conta do Estado. Porquanto, o Estado --- ou administra e cultiva por officiaes seus todas as terras, --- ou alluga a muitos colonos, por um preço determinado, todas as terras, --- ou deixa os actuaes proprietarios na posse dos fundos, e só exige, a titulo de tributo, o rendimente dos mesmos fundos. Ora estas tres hypotheses são \emph{praticamente} impossiveis.\par a) E' impossivel a primeira hypothese. Com effeito, a terra precisa de cultura assidua e diligente, e esta cultura só póde ser feita por pessoas, que têm residencia actual nessas terras (pois só ellas as podem conhecer) e que empregam todo o cuidado e diligencia para que o fructo seja abundante. Ora os operarios ambulantes, de quem o Estado teria de servir-se para a cultura das terras, devendo entregar todo o rendimento e só contentar-se com o seu salario, não conheceriam a natureza das terras nem teriam empenho em multiplicar o seu producto. --- De resto, esta hypothese é rejeitada pelos proprios defensores do socialismo \emph{moderado} ou \emph{agrario}.\par 5) E' impossivel a segunda hypothese. Porquanto, esta hypothese, além de acabar com o \emph{estado dos agricultores}, com as suas residencias fixas, impede o progresso da cultura. Com effeito, os colonos devem, muitas vezes, fazer grandes despesas e trabalhar durante largos annos, para que a terra dé finalmente um fructo abundante. Mas, os colonos, estando sujeitos aos caprichos dos governos e das facções politicas, nunca poderiam estar certos de que hão de ficar, durante longo tempo, na posse da terra allugada, e por isso nunca fariam grandes despesas nem se sujeitariam a gravosos trabalhos. D'ahi a paralysacao da agricultura.\par c) E' impossivel a terceira hypothese, sustentada por \emph{George} e sequazes. Na verdade, estes socialistas dizem que não se deve tirar a posse das terras aos actuaes proprietarios, mas que se devem impor graves tributos sobre as terras, de modo que todo o rendimento d'ellas seja devolvido ao Estado, com excepção da parte correspondente ao capital e ao trabalho, e de alguma outra coisa necessaria aos melhoramentos dos fundos e á manutenção das casas, etc. Estes tributos sobre as terras dispensam, na opinião de \emph{George}, todos os outros tributos e bastam para alliviar a miseria e para tornar feliz a sociedade. --- Tal theoria é impossivel na realidade. Ordinariamente o producto liquido das terras não excede 3 ?/,. Os agricultores, embora sejam proprietarios e empreguem todos os esforços, apenas tiram do seu trabalho o necessario para a honesta sustentação da sua familia. E poderiam elles sustent\&r-se, se augmentassem os tributos sobre os fundos até não haver precisão de novos tributos? Certamente que não. D'ahi o augmento da miseria nas familias e as desordens e revoluções na sociedade civil.\par XI. Objecções dos defensores do socialismo moderado. --- Os defensores do socialismo \emph{moderado}, ou \emph{agrario}, fazem muitas objecções. Limitamo-nos a expor e criticar as principaes.\par a) Dizem: A propriedade particular das terras é contraria á natureza; porque, no principio, a posse das terras era commum.\par Resposta. Náo é verdade que, no principio, \& posse das terras era commum; pois a historia diz que, entre os Hebreus, Assyrios, Egypceios, Chinezes, a propriedade era particular, e não commum. --- E, embora entre esses\pgn{721}povos a posse das terras tivesse sido commum, a propriedade particular nunca poderia dizer-se contraria á natureza. A propriedade particular é um fructo da civilisação, e esta não é contraria á natureza humana.\par 5) Continuam: O \emph{trabalho} é o unico titulo justo da propriedade particular; pois ninguem póde possuir senão o que produziu com o seu trabalho. Ora a terra não é producto do trabalho. Logo a terra não póde ser objecto de propriedade particular.\par Resposta. Essa razão prova muito, e por isso não prova nada. Se o homem só tivesse direito ao fructo do seu trabalho, a posse particular das pedras preciosas, do ouro, etc., seria injusta e o proprietario deveria entregar ao Estado uma grande parte do valor dos objectos, que não produziu com o seu trabalho; como tambem seria injusto o uso que fazemos das fruetas, da lenha, dos animaes e de muitas outras coisas, que foram preparadas pela natureza e que, não podendo ser modificadas pelo trabalho, não poderiam passar para a nossa propriedade.\par \emph{c}) Replicam: Quem cultivou uma porção de terreno inculto, tem direito somente aos fructos e não á propriedade do terreno. Logo é injusto que o agricultor conserve a propriedade particular dos terrenos.\par Resposta. Quem cultivou uma porção de terreno inculto, não tem somente direito aos fructos. Porquanto, o trabalho transformou e melhorou o solo, tornando-o fecundo não só durante um anno, mas perpetuamente (com tanto que o homem conserve esse melhoramento) e por isso augmentou o valor primitivo do terreno. Se o homem tornou fertil e melhor o terreno, este pertence-lhe, e ninguem Ih'o póde tirar, sem commetter uma grande injustiça.\par \emph{d}) Insistem: Todos os homens têm direito á \emph{vida}. Logo todos devem ter direito a uma porção \emph{egual} de terra.\par Resposta. E' verdade que todos os homens têm direito á \emph{vida}, mas é falso que todos têm direito a uma porção \emph{egual} de terra. Na verdade, essa porção \emph{egual} de terra não é o unico e indispensavel meio para a conservação da vida; pois muitas pessoas ha que não possuem terras, e todavia têm o necessario para a sua sustentação. --- Além d'isso, se todos tivessem direito \& egual porção de terra, porque todos tem egual direito á \emph{vida}, deveria dizer-se, pela mesma razão, que todos têm o mesmo direito á egual porção de bens \emph{moveis}; o que não é admittido nem pelos proprios defensores do socialismo \emph{moderado} (Cathrein, \emph{Philosophia Moral}).}.

f) \textbf{Direito} de independencia. --- Direito de \emph{independencia} é o poder moral que o homem tem --- de exercer, livre e legitimamente \pg{706} \pg{707} , a sua actividade, --- de aproveitar-se dos fructos e resultados da mesma actividade, --- e de exigir que sejam \pg{708}  respeitados todos os bens, que lhe são proprios. --- Este direito existe no homem. Por quanto o homem está sujeito a muitas \pg{709}  e diversas necessidades, --- organicas, intellectuaes e moraes, --- a que elle não póde deixar de satisfazer, se quizer alcançar os \pg{710}  fins, que Dzvus, creando-o, lhe determinou. Ora-o homem não

, , poderia satisfazer a estas muitas e diversas necessidades, se \pg{711}  não podesse. exercer, livre e legitimamente, a sua actividade. Logo o homem tem o \emph{dever}, e por isso o \emph{direito} de exercer, \pg{712}  livre e legitimamente, a sua actividade. --- Se o homem tem o direito de exercer, livre e legitimamente, a propria actividade, \pg{713}  deve ter tambem o direito de aproveitar-se dos fructos e resultados, que derivam da mesma actividade, Porquanto, sendo \pg{714}  justo que o effeito siga a propria causa, é tambem justo que os fruetos e resultados aproveitem a quem exerceu a actividade \pg{715} , de que elles derivam. Logo o homem tem o direito aos fruetos e resultados da propria actividade. --- Finalmente, a \pg{716}  egualdade \emph{essencial}, --- que existe entre todos os homens (pois todos os. homens são dotados de alma espiritual e immortal, \pg{717}  creada á imagem e semelhança de \textsc{Deus}) e que, como dissemos, ficaria destruida, se o homem fosse considerado simplesmente \pg{718}  como um \emph{meto} para os fins temporaes dos outros e dependesse, por um modo absoluto e arbitrario, do capricho do \pg{719}  mais forte, --- exige que todos os homens respeitem os varios e diversos bens, que nos são proprios, isto é, que nos estão unidos \pg{720}  por um vineulo, physico ou moral. Logo o homem tem o direito de exigir que todos os homens respeitem os varios e \pg{721}  diversos bens, que lhe são proprios. Logo \emph{o} homem tem o direito de \emph{independencia} \footnote{O direito de \emph{independencia} refere-se --- \emph{a}) ao respeito devido aos \emph{bens} proprios do homem, --- \emph{db}) e ao legitimo \emph{exercicio} da actividade humana.\par a) Os \emph{bens} proprios do homem, cujo respeito é imposto pelo direito de \emph{independencia}, são --- 1.º) os bens \emph{corporeos}, como a vida, a integridade do organismo, a saude, etc., --- 2.9) os bens \emph{internos} da \emph{alma}, \emph{intellectuaes} e \emph{moraes}, --- 9.9) os bens \emph{externos} e \emph{moraes}, como a \emph{fama}, a \emph{honra}, --- 4.9) os bens \emph{externos} e \emph{materiaes}. --- Relativamente aos bens \emph{corporeos}, o direito de \emph{independencia} prohibe os ataques, directos ou indirectos, contra a nossa vida, a integridade dos membros, a saude, o pudor, etc. Estes ataques são tambem prohibidos pelo direito á \emph{vida}. --- Relativamente aos bens \emph{internos} da \emph{alma}, o direito de \emph{independencia} exige que ninguem nos prive do uso da nossa razão pela \emph{embriaguez}, pelo \emph{hypnotismo}, ou nos ensine erros e falsidades, e tambem que\par 46 Vor. II.\pgn{722}não deprave a nossa vontade com maus exemplos, com a seducção, etc. Estes ataques são tambem prohibidos pelo direito á verdadeira \emph{liberdade religiosa}. --- Relativamente aos bens \emph{moraes} e \emph{externos}, o direito de \emph{independenciía} prohibe --- os ataques contra a \emph{fama}, que são o \emph{juizo temerario}, a \emph{calumnia}, a \emph{murmuração}, --- e os ataques contra a \emph{honra}, que são a \emph{injuria}, o \emph{escarneo}, a \emph{maldição}. --- Relativamente aos bens \emph{materiaes} e \emph{externos}, o direito de \emph{independencia} prohibe o \emph{roubo}, \emph{á fraude} (principalmente nos \emph{contractos}) a usura, a violação do \emph{domicilio}.\par 5) O legitimo \emph{exercicio} da actividade humana, garantido pelo direito de \emph{independencia}, consiste em que o homem, dentro dos limites da ordem moral e social, póde transferir-se de um logar para outro, escolher a profissüo e o estado, fazer contractos, dispor da sua actividade, etc.\par O direito de \emph{independencia} tem, como os outros direitos, os seus limites, marcados pela propria \emph{let natural}. --- Estes limites nascem --- dos \emph{deveres}, que temos para com \textsc{Deus}, para comnosco e para com o proximo, --- e dos \emph{direitos} mais elevados, de que os outros são dotados e cujo exercicio é incompativel com este direito de \emph{independencia}.}.

\pg{722} £) \textbf{Direito} de associação, --- Direito de \emph{associação} é o poder moral, que todo o homem tem de unir, por um modo constante, as suas forças com as dos outros, para a consecução de um fim communm, licito e honesto. --- Este. direito existe no homem. Porquanto, quem tem o dever de conseguir um fim, deve ter tambem o direito de empregar os meios necessarios para a consecução d'esse fim. Ora o homem tem o dever de desenvolver e aperfeiçoar as suas faculdades, e um dos meios mais necessarios para isso é a \emph{associação}, pela qual unimos as nossas forças com as dos nossos semelhantes. Logo o homem tem o direito de \emph{associação} \footnote{O direito de \emph{associação}, fundando-se na propria natureza \emph{social} do homem e na impossibilidade de desenvolver as suas faculdades fóra da sociedade, deriva da lei \emph{natural}, e não é mera concessão da \emph{lei civil}.\par Este direito tambem tem os seus limites, que derivam --- não só do proprio \emph{fim} do direito de associação, pois nao póde haver direito de \emph{associação} para fins contrarios á natureza racional do homem, e por isso á lei moral, --- mas tambem dos \emph{direitos}, mais elevados do proximo, que, na concorrencia, devem prevalecer.}.

145. \textbf{Direitos moraes} do \textbf{homem para} com o proximo. --- Direitos moraes do homem para com o proximo são os que se baséam na virtude de \emph{caridade}. --- Os principaes são o direito á \emph{humanidade} e o direito á \emph{beneficencia}.

a) Direito á humanidade, --- Direito á \emph{humanidade} é o poder moral, que o homem tem de exigir uma coisa, que é util a si proprio e não é prejudicial aos outros. --- Este direito existe no homem. Porquanto o homem, tendo o direito á sua conservação e ao seu aperfeiçoamento, deve ter o direito ao que

E \pg{723}  \emph{positivamente} concorre para esse fim e que não é prejudicial aos outros. Logo o homem tem direito á \emph{humanidade}.

b) Direito á beneficencia, --- Direito á \emph{beneficencia} é o poder moral que o homem tem de exigir, nas suas necessidades, o soccorro dos seus semelhantes. --- Este direito existe no homem. Porquanto o direito, que o necessitado tem á propria subsistencia ou ao bem da sua alma, excede o direito, que os outros têm a uma fortuna superabundante ou ás suas commodidades. Logo o homem tem direito á \emph{beneficencia} \footnote{A \emph{beneficencia}, que é \emph{particular} ou \emph{publica}, conforme provém dos individuos ou da sociedade civil, deve ser \emph{prudente}, \emph{generosa}, \emph{benigna}, \emph{ordenada}. --- Deve ser \emph{prudente}, isto é, não deve ser excessiva para quem a faz, nem inutil ou prejudicial para quem a recebe. --- Deve ser \emph{generosa}, isto é, não deve ser feita por fins secundarios e egoistas. --- Deve ser \emph{benigna}, isto é, deve ser affavel, de modo que não envergonhe o beneficiado. --- Deve ser \emph{ordenada}, preferindo os mais proximos, os que melhor hão de aproveitar-se da esmola, os mais necessitados.\par Os direitos \& \emph{humanidade} e \& \emph{beneficencia} --- não são \emph{jurídicos}, não se basêam sobre a idea da \emph{justiça}, porque o sujeito d'estes direitos não tem o poder de exigir dos outros uma coisa como \emph{propria}, isto é, como pertencente ao mesmo sujeito; --- mas são \emph{moraes}, porque se fundam sobre a idea da \emph{caridade}. --- Por isso quem transgride os direitos \emph{jurídicos} do proximo, chama-se \emph{injusto}; --- e quem transgride os direitos \emph{moraes}, diz-se \emph{deshumano}, \emph{duro} de \emph{coração}. Os direitos \emph{jurídicos} são \emph{negativos}; --- os \emph{moraes} são \emph{positivos}.}.

%%>>> art-65

\section*{Artigo II.\\ Deveres do homem para com o proximo.}
\addcontentsline{toc}{section}{Artigo II. --- Deveres do homem para com o proximo.}

146. \textbf{Deveres} do homem para com o proximo. --- \emph{Deveres} do homem para com o \emph{proximo} são as obrigações que o homem tem de respeitar os \emph{direitos} do proximo. --- Estes deveres são \emph{Jurídicos} ou \emph{moraes}, conforme derivam da virtude da \emph{justiça} ou da virtude da \emph{caridade}, e por isso correspondem aos direitos \emph{jurídicos} ou \emph{moraes} do proximo \footnote{Tendo sufficientemente desenvolvido a theoria ácerca dos \emph{direitos} do homem para com o \emph{proximo}, daremos aqui, muito em resumo, a theoria Acerca dos \emph{deveres}, correlativos a esses \emph{direitos}.}.

147. \textbf{Deveres} juridicos do homem \textbf{para} com o proximo. --- Os deveres \emph{jurídicos} do homem para com o proximo, sendo as obrigações de respeitar os direitos \emph{jurídicos} dos nossos semelhantes, são tantos, quantos são esses \emph{direitos}. --- Por isso o homem tem o \emph{dever} de respeitar o \emph{direito}, que o homem tem \pg{724}  4 \emph{liberdade religiosa}, \emph{á dignidade pessoal}, \emph{á vida}, \emph{á legitima defeza}, á \emph{propriedade}, \& \emph{independencia} e \& \emph{associação}. --- Estes \emph{deveres}, como se vê, fundam-se na virtude da \emph{justiça}, que manda dar a cada um o que lhe pertence, e são \emph{negativos}, porque obrigam a não fazer coisa, que possa prejudicar o proximo \footnote{A existencia dos deveres \emph{juridicos} é uma consequencia da existencia dos direitos \emph{jurídicos}. Limitar-nos-hemos, pois, a indicar a natureza de cada dever \emph{jurídico}, e a esphera em que se deve exercer.}.

a) \textbf{Dever} de respeitar o direito á liberdade \textbf{religiosa}, --- Este dever consiste em deixar ao proximo a liberdade de erer e praticar tudo o que a verdadeira Religião ensina e em não fazer coisa que possa lesar este direito \footnote{Lesam o direito que o homem tem á \emph{liberdade religiosa} --- todos os que propagam o erro e procuram impedir, por qualquer forma, a pratica da verdadeira Religião. --- E' claro que não se oppõem a este direito os que impedem o progresso do erro ou a pratica de uma religião falsa e funesta; pois não ha nem póde haver o direito de propagar o erro e de prejudicar a alma do proximo. --- Se se fizesse alguma reflexão neste ponto, não seriamos obrigados a ler e a ouvir tantos disparates em materia de Religião!}.

b) Dever de respeitar o direito á dignidade pessoal. --- Este dever consiste em reconhecer o homem como ente racional, dotado de fim proprio, e não como simples meio para o fim dos outros, e em não fazer coisa contraria a este direito \footnote{O homem deve considerar o seu proximo não como um simples \emph{meto} para os seus fins temporaes, mas como um ser nobilissimo, creado á imagem e semelhança de \textsc{Deus}. E, pois, illicito não só reduzir os nossos semelhantes á condição de \emph{escravos}, mas tambem impor-lhes um trabalho superior ás suas forças, ou impor-lhes o trabalho em dias santos, ou descuidar-se da sua saude, temporal e eterna. A responsabilidade dos patrões, que maltratam os operarios, é tremenda deante de \textsc{Deus}, que se chama o Pae dos orphãos e o Juiz das viuvas, e que julgará os poderosos com todo o poder da sua severidade!}.

c) Dever de respeitar o direito á vida, --- Este dever consiste em não attentar contra a \emph{vida} do proximo, nem contra a integridade do seu organismo, nem contra a sua saude \footnote{Este dever prohibe todo e qualquer ataque contra a \emph{vida} do proximo, contra a integridade dos seus membros, contra a sua saude. A vida do proximo é um dom de \textsc{Deus}, que se reservou o direito absoluto e exclusivo sobre ella, \emph{e} é um dom \emph{pessoal}, que \textsc{Deus} fez ao individuo humano. Quem, pois, lesar este direito, commette uma injustiça contra \textsc{Deus} e contra o proximo. Commette uma injustiça contra \textsc{DEus}; porque invade um dominio, que \textsc{Deus} reservou para si mesmo. Commette uma injustiça contra o proximo; porque, sem razão, tira-lhe um bem que lhe pertence.}.

d) Dever de respeitar o direito de legitima \textbf{defeza}. --- Este dever consiste em reconhecer no homem o direito de empregar \pg{725}  a força physica para repellir uma ageressão actual e injusta, e em não pôr obstaculo ao exercicio d'este direito \footnote{Este dever obriga, como dissemos, a não fazer nada contra o direito de \emph{legitima defeza}. Seria, pois, injusta uma pessoa, que, vendo um innocente nggredido por um malfeitor, auxiliasse o aggressor ou obstasse, de algum modo, á que o aggredido se podesse defender do ataque.}.

e) \textbf{Dever} de respeitar o direito de propriedade. --- Este dever consiste em reconhecer no proximo o direito de.adquirir os bens materiaes e externos, e de usar e dispor d'elles livre e exclusivamente, e em não praticar coisa contraria a tal direito \footnote{O dever de respeitar o direito de \emph{propriedade} consiste em não pôr obstaculos ao exercicio d'esse direito, e sobre tudo em não apropriar a si uma coisa alheia contra a \emph{justa} prohibição do dono. --- Dizemos --- contra a \emph{Justa} prohibição do dono, porque, se a prohibição for \emph{injusta}, \& apropriação da coisa alheia não offende o direito de \emph{propriedade}. Assim um pobre, que se encontra numa necessidade \emph{extrema}, póde apropriar a si uma coisa alheia, que baste para prover á sua extrema necessidade. Porquanto, o homem tem o \emph{direito} de procurar os meios \emph{necessarios} para a sua sustentação; e o direito, que o pobre tem á sua \emph{vida}, é superior ao direito, que o proprietario tem aos seus bens, e por isso, na \emph{collisüo}, deve prevalecer a este. Nesse caso uma prohibicio do proprietario seria \emph{injusta}. --- Isto, porém, suppõe que o proprietario não se encontre na mesma necessidade, em que se encontra o pobre; porque, se o proprietario tambem se encontrar nas mesmas circumstancias do pobre, este não póde apropriar a si a coisa alheia, visto que, então, o direito que o proprietario tem á \emph{vida} e \& \emph{propriedade} é superior ao direito que o pobre tem á \emph{vida}.}.

f) \textbf{Dever} de respeitar o direito de independencia, --- Este dever consiste em reconhecer no proximo o direito --- de exercer, livre e legitimamente, a sua actividade, --- de aproveitar-se dos fruetos e resultados da mesma actividade, --- e de exigir que sejam respeitados todos os bens, que lhe são proprios, e em não fazer coisa alguma que possa lesar este direito \footnote{Este dever prohibe, como dissemos, --- os ataques contra os bens \emph{corporeos} de proximo, e por isso o homicidio, a mutilação, o duello, etc., --- os ataques contra os bens \emph{internos} da alma, \emph{intellectuaes} e \emph{moraes}, e por isso é illicito*privar os outros do uso da razão pela \emph{embriaguez} ou pelo \emph{Aypnotismo}, inocular o erro nas intelligencias, sobretudo juvenís, dizer \emph{mentiras}, depravar a vontade dos outros com os maus exemplos, com a seducção, etoc., --- 08 ataques contra os bens moraes e \emph{externos}, e por isso é illicito o \emph{juizo temerario}, a \emph{calumnia}, a \emph{murmuração}, a \emph{injuria}, etc., --- os ataques contra os bens \emph{materiaes} e \emph{externos}, e por isso é prohibido o \emph{roubo}, a \emph{fraude} (especialmente nos \emph{contractos}) a usura, a violação do \emph{domicilio}, etc.;, --- os ataques contra o legitimo \emph{exercicio} da actividade humana, e por isso é prohibido impor-se; sem motivo justificado, á vontade do proximo, obstando a que este disponha da sua liberdade na escolha da profissão, do estado, etc.}.

£) \textbf{Dever de} respeitar o direito de associação, --- Este dever \pg{726}  consiste em reconhecer no proximo o direito que elle tem de unir, de um modo constante, as suas forças com as dos outros, --- e em não impedir o exercicio d'este direito \footnote{Este dever transgride-se, quando, injustamente, se impede a fundação das associações, ou quando estas se dissolvem, depois de fundadas.}.

148. \textbf{Deveres moraes} do homem \textbf{para} com o \textbf{proximo}. --- Os deveres \emph{moraes} do homem para com o proximo, sendo as obrigações de respeitar os direitos \emph{moraes} dos nossos semelhantes, são tantos, quantos são esses direitos. --- Por isso o homem está sujeito aos deveres de \emph{humanidade} e de \emph{beneficencia}. --- Estes \emph{deveres} fundam-se na virtude da \emph{caridade}, a qual manda que \emph{amemos o proximo como} a \emph{nós mesmos}, e são \emph{positivos}, pois nos obrigam a soecorrer ao proximo com os nossos bens \footnote{Os deveres \emph{moraes} do homem para com o proximo derivam do grande preceito da \emph{caridade}, imposto por \textsc{Deus} ao homem e contido na celebre formula: Ama \emph{o teu proximo como a ti Mesmo}. O adverbio --- \emph{como} --- não significa \emph{egualdade}, mas \emph{semelhança}; pois não somos obrigados a ter ao proximo um amor \emph{egual} ao que temos a nós mesmos (o que seria contra a ordem natural) mas devemos ter-lhe um amor \emph{semelhante} ao amor, com que amamos a nós mesmos, emquanto devemos amar a nós mesmos e ao proximo em \textsc{Deus} e por \textsc{DEus}, desejando a todos a consecução do fim ultimo, que é um e o mesmo para todos os homens, e consiste na posse do Bem infinito.\par A razão ou o fundamento do amor do proximo é a participação da mesma natureza racional, pela qual todos somos filhos de \textsc{Deus} e herdeiros da mesma gloria; porque a semelhança é causa de amor.\par O amor, que devemos ter ao nosso proximo, deve ser de \emph{benevolencia}, de modo que desejemos e procuremos o bem e a felicidade d'elle e não a nossa satisfação e o nosso proveito. --- O amor de \emph{benevolencia}, quando é \emph{reciproco}, chama-se \emph{amizade}; que tende sempre para a \emph{egualdade} e \emph{commanicacdo} de bens. D'hai a maxima: \emph{Amicitia pares aut invenit aut facit}.\par O amor do proximo deve ser \emph{ordenado}. Por isso os que nos são mais unidos pelo parentesco ou por outro laco devem ser preferidos aos outros. --- Mas ninguem póde ser exeluido do nosso coração, nem o nosso maior inimigo; pois tambem este conyém comnosco na participação da mesma natureza racional, que é, como dissemos, a causa e o fundamento do amor para com o proximo. --- O preceito de amar os proprios inimigos é o claro distinetivo do Christianismo, pois era quasi inteiramente desconhecido dos pagãos. \emph{Socrates} dizia: « E' merecedor dos maiores elogios quem previne os inimigos nas offensas e os amigos nos beneficios »v. No mesmo sentido fallam \emph{Seneca} e \emph{Cicero}. Veio N. S. Jesus CunisTo, e, purificando de todo o egoismo o principal sentimento da nossa alma e elevando o homem aos mais sublimes ideaes, deu ao mundo o grande preceito: « Àmai os vossos inimigos, fazei bem aos que vos odeiam --- \emph{JDiligite inimicos vestros}, \emph{benefacite lis qui oderunt} vos » (Matth, V, 44).\par Se onosso amor deve extender-se.aos inimigos, muito mais deve abranger a nossa \emph{patria}. O amor da \emph{patria} deve ser \emph{sincero}, \emph{desinteressado}, \emph{ordenado}, \emph{activo}; e, para ter estes attributos, deve fundar-se na /\emph{onestidade} ou na\pgn{727}\emph{virtude}. \emph{Silvio Pellico} diz: « Se alguem violar os altares, a santidade cofiju4 gal, a decencia, a probidade, e depois bradar: patria, patria! não soredilaig\par nelle. E' um hypocrita do patriotismo. E' um pessimo cidadão ».\par D'onde se vé --- que o amor de \textsc{Deus} envolve necessariamente o amor do proximo; pois quem ama o original, não póde deixar de amar o retrato, --- e que o verdadeiro amor do proximo suppõe necessariamente o amor de \textsc{Deus}; pois quem ama o proximo por \textsc{Deus} e para \textsc{DEus}, não póde deixar de amar o proprio \textsc{Deus}.}.

\pg{727} á) \textbf{Deveres} de humanidade. --- Os deveres de \emph{humanidade} consistem em praticar, em beneficio do proximo, um acto, que é util a quem o recebe, e não é prejudicial a quem o faz; tal é o acto de quem se faz guia a um -cego \footnote{Os deveres da \emph{humanidade} resumem-se naquelle principio da jurisprudencia romana: « Estás obrigado a fazer o que não prejudica a ti e é util a outrem --- \emph{Quod tibi non nocet et alteri prodest}; \emph{ad id obligatus es} n. --- \emph{Cicero} applicou a estes deveres os seguintes versos de Ennio:\par \emph{Homo qui erranti comiter monstrat viam}\par \emph{Quasi lumen de suo lumine accendat facit}.\par \emph{Nihilo minus ipsi lucet cum illi accenderit}.}.

b) \textbf{Deveres} de beneficencia, --- Os deveres de \emph{beneficencia} consistem na obrigação que temos de soccorrer o nosso proximo nas suas necessidades, temporaes e espirituaes \footnote{Os homens, pois, cujo rendimento é superior ás suas necessidades, têm o dever rigoroso de soccorrer os pobres. --- E não basta. O proximo, além das necessidades corporaes, póde ter necessidades espirituaes, superiores aquellas na essencia e na importancia. E, pois, nosso dever prestar auxilio espiritual aos nossos semelhantes, livrando-os dos perigos de perdição e de morte eterna. --- Estes deveres foram compendiados nos versos seguintes:\par \emph{Vésito}, \emph{poto}, \emph{cibo}, \emph{redimo}, \emph{tego}, \emph{colligo}, \emph{condo},\par \emph{Consule}, \emph{carpe}, \emph{doce}, \emph{solare}, \emph{remitte}, \emph{fer}, ora.\par Como é facil ver, as obras de \emph{misericordia} s\&o deveres \emph{inoraes}, impostos pela lei natural Nada mais bello e consolador do que praticar essas obras por amor de \textsc{Deus}, reconhecendo e agasalhando na creatura o proprio Oreador! Este é o grande principio, que leva tantas almas g generosas a encontrar, em beneficio do proximo, os mais heroicos sacrificios e a propria morte! A Verdade increada disse: « \emph{Bemaventwurados os misericordiosos}, \emph{porque alcançarão misericordia} ! ».}.

\pg{728} \textbf{DIREITO SOCIAL}

INTRODUCÇÃO

149. \textbf{Sociedade}. --- \emph{Sociedade} é a \emph{união de duas ou mais pessoas}, \emph{que}, \emph{pelos seus actos}, \emph{cooperam para} um \emph{fim commum}.

a) A \emph{sociedade} é a \emph{união} de duas ou mais pessoas; porque, para a formação da \emph{sociedade}, é preciso que as pessoas se unam. --- Esta união deve ser \emph{moral} e \emph{constante}. Deve ser \emph{moral}, isto é, formada por laços moraes ou espirituaes, que são os da intelligencia e da vontade, emquanto todos conhecem e querem o mesmo fim. Deve ser \emph{constante}; porque uma reunião \emph{passageira} de pessoas não merece o nome de \emph{sociedade} \footnote{A \emph{sociedade} não é um simples aggregado de pessoas, aliás as pessoas\par , que estão num templo, num theatro, etc. formariam uma sociedade; mas é uma união \emph{moral}, uma união de intelligencias e de vontades, de modo que os socios podem viver em diversos paizes, sem comtudo deixar de pertencer \& sociedade, e de formar uma só entidade moral.}.

b) A \emph{sociedade} é a união de pessoas, que cooperam para um \emph{fim commum}; porque uma união \emph{moral} e \emph{constante} consiste na cooperação de todos para um fim \emph{commum}. O fim \emph{commum} deve ser necessariamente \emph{honesto} \footnote{O \emph{fim}, por ser \emph{bem}, attrahe. --- O \emph{fim}, que attrahe varias pessoas e produz nellas a unidade de pensamentos, de desejos.e de esforços, deve ser \emph{commum};, isto é, deve ser proprio da \emph{collectividade}, e não dos individuos; aliás haveria um effeito sem causa proporcionada. --- O fim, que reune varias pessoas em \emph{sociedade}, deve ser \emph{honesto}, porque uma união. \emph{moral} e \emph{constante} de intelligencias e de vontades só póde basear-se no bem \emph{honesto}, conveniente á natureza racional. Por isso uma conspiração de malfeitores não merece o nome nem a dignidade de \emph{sociedade}.}.

c) A \emph{sociedade} É a união de pessoas, que, \emph{pelos} seus \emph{actos}, cooperam para um fim commum; visto que os socios devem unir as suas forças --- \emph{intellectuaes}, \emph{moraes} e \emph{physicas} --- para a consecução do \emph{fim commum}, e por isso devem combinar-se na escolha e no uso dos \emph{meios} aptos para o mesmo fim \footnote{Para a formação da \emph{sociedade} não basta que as pessoas conheçam e desejem o fim \emph{commum} e \emph{honesto}, mas é necessario que unam as suas forças e empreguem os meios convenientes. Sem o emprego d'esses meios, a \emph{sociedade} seria uma coisa ideal e subjectiva, e nunca uma união externa e real, e por isso nunca seria verdadeira \emph{sociedade}.}.

\pg{729} 150. \textbf{Elementos da sociedade}. --- Os elementos, que constituem a essencia da \emph{sociedade}, são dois: a \emph{pluralidade} das pessoas e a sua \emph{união}. \emph{A pluralidade} das pessoas é a \emph{materia} da sociedade; a sua \emph{união} é á \emph{forma}. --- Os principios, de que deriva a \emph{união}, são o \emph{fim ultimo} e a \emph{auctoridade}. A \emph{auctoridade} é o \emph{direito de dirigir eficazmente os membros da sociedade para a consecução do fim commum}. --- O \emph{fim} não só é o principio da \emph{união social}, mas é tambem o seu \emph{objecto formal}; por isso \emph{especifica} a sociedade, e é a \emph{razão} e a \emph{regra} de todos os seus \emph{deveres} e \emph{direitos} \footnote{Desenvolvamos a theoria, dada muito em resumo no texto --- \emph{a}) ácerca dos dois elementos constitutivos da \emph{sociedade}, que são a \emph{pluralidade} das pessoas e a sua \emph{união}, --- \emph{D}) ácerca dos dois principios, de que deriva a \emph{união}, e que são o \emph{fim commum} e a \emph{auctoridade}, --- \emph{c}) ácerca do mesmo \emph{fim commun}, emquanto, por ser o \emph{objecto formal} da \emph{união}, \emph{especifica} a \emph{sociedade} e é a razão e a \emph{regra} de todos os seus \emph{deveres} e \emph{direitos}.\par a) Os elementos constitutivos de toda a \emph{sociedade} são dois: um indeterminado e determinavel, que é a \emph{pluralidade} de pessoas, --- outro determinado e determinante, que é a \emph{união} moral] das mesmas pessoas em ordem ao \emph{fim commun}. A \emph{pluralidade} das pessoas, que é a \emph{materia} da sociedade, recebe da \emph{união social} uma nova unidade entitativa, que é distincta da unidade entitativa de cada individuo e pela qual a sociedade, embora multiplice nos seus membros, forma um só \emph{corpo social} e uma só \emph{pessoa moral}.\par A denominação de \emph{corpo}, dada á \emph{sociedade}, mostra a \emph{analogia}, que existe entre o \emph{corpo organico}, ou \emph{vivo}, e a \emph{sociedade}. Porquanto, o corpo \emph{organico} é composto de muitos e diversos elementos, que estão unidos por uma força commum e intrinseca, que se chama \emph{alma}; --- a sociedade é composta de varios membros, que estão unidos por um laco commum de cohesao moral e cooperação, distincto de cada uma das partes. O corpo organico, embora sujeito a perpetuas mudanças, permanece sempre identico; --- a sociedade, quando numerosa, embora soffra continuas perdas, todavia não acaba.\par A denominação de \emph{pessoa moral}, dada á \emph{sociedade}, mostra que esta tem, relativamente á \emph{pessoa humana}, não só uma analogia \emph{generica} (emquanto a sociedade é um organismo \emph{moral}, e o homem é um organismo \emph{natural}), mas tambem uma analogia \emph{especifica}. Por quanto, assim como a pessoa \emph{natural} é dotada de razão e de actividade propria; assim tambem a pessoa \emph{moral}, que é a sociedade, é dotada de \emph{razão} (pois se compõe de membros \emph{racionaes}) e de \emph{actividade propria}, distincta da actividade de cada um dos socios. Assim como a pessoa \emph{natural} é sujeito de \emph{deveres} e de \emph{direitos}, e produz operações \emph{moraes}, de que é responsavel; assim tambem a pessoa \emph{moral}, que é a \emph{sociedade}, é sujeito de \emph{deveres} e de \emph{direitos}, e produz acções, que lhe são inteiramente imputaveis.\par Todavia devemos evitar os exaggeros dos \emph{positivistas}. Estes julgam que a comparação da \emph{sociedade} com um \emph{organismo} é uma descoberta feita por elles neste seculo de progresso. Se tivessem lido as Epistolas de S. Paulo, as obras dos Santos Padres e Doutores da Egreja e dos Escholasticos, antigos e modernos, ter-se-hiam convencido de que a descoberta é muito velha. A verdade é que os \emph{positivistas} desvirtuaram aquella doutrina. Porquanto, do facto de\pgn{730}haver \emph{analogia} entre o corpo \emph{organico} e a \emph{sociedade}, deduziram que entre elles ha \emph{egualdade}, e que por isso, assim como o corpo \emph{organico} está sujeito ás leis \emph{biologicas}, que são \emph{fataes} e \emph{necessarias}, assim tambem a sociedade está sujeita a leis \emph{fataes} e \emph{necessarias}. Basta o simples senso commum para mostrar --- que organismo \emph{moral} não é \emph{egual} ao organismo \emph{physico} ou \emph{natural}, --- e que por isso a sociedade, sendo um organismo \emph{moral}, composto de seres \emph{racionaes} e \emph{livres}, deve ser regulado por leis diversas das que regulam o organismo \emph{physico}, composto de cellulas materiaes, destituidas de razão e de livre arbitrio. --- Esta theoria dos \emph{positivistas} é uma copia ou reproducção da theoria de \emph{Heckel}, que citàmos na \emph{Cosmologia}.\par b) A \emph{união} ou \emph{conspiração} dos membros para o mesmo fim é, como dissemos, a \emph{forma} da sociedade; porque é ella que determina a \emph{pluralidade} das pessoas ao estado e á natureza da sociedade. Concebida esta \emph{união} entre as pessoas, concebe-se a \emph{sociedade}. --- A "\emph{nido}, porém, suppõe um \emph{principio}, de que deriva e por que é constituida; pois o effeito deve ter uma causa proporcionada. Este principio é duplice: o \emph{fim commum} e a \emph{auctoridade}. O \emph{fim} é principio da «\emph{nido}; porque é o conhecimento do fim e o desejo da sua consecugao que leva os homens a unir as suas forças: --- \emph{intellectuaes}, \emph{moraes} e \emph{physicas}. \emph{O fim commum} é principio \emph{ideal} e \emph{objectivo} da união social. Mas, se considerarmos a sociedade, não numa ordem \emph{idea}! e \emph{abstracta}, mas sim na ordem \emph{real} e \emph{concreta}, em que os homens são dotados de livre arbitrio, e divergem nos juizos e nas forças, veremos que o \emph{fim} não é por si sufficiente para formar e conservar essa união constante e ordenada entre os membros. E, pois, necessario outro principio \emph{real} e \emph{subjectivo}, que vivifique todos os membros da sociedade, como o principio vital vivifica todo o organismo natural, e que, por um modo constante e efficaz, una e dirija a livre vontade dos membros para a consecução do fim commum. Este outro principio é a \emph{auctoridade}. Por isso estabelecemos a seguinte\par Proposição. --- A \emph{auctoridade é essencialmente necessaria na sociedade}.\par \emph{a}) \emph{A auctoridade} será essencialmente necessaria, se, sem ella, não for possivel a \emph{união moral} e \emph{constante} dos membros; pois esta \emph{união}, por sera \emph{forma} da sociedade, é essencialmente necessaria. Ora, sem «\emph{uctoridade}, não é possivel a união \emph{moral} e \emph{constante} entre os membros da sociedade. Porquanto, os homens, divergindo entre si nos \emph{juizos} ou \emph{apreciações} ácerca dos meios \emph{proporcionados} para a consecução do fim, e, sendo dotados de \emph{livre arbitrio}, que os leva a fazer o que lhes agrada e não o que é necessario para o bem de todos, nunca poderão formar nem conservar constantemente a união de intelligencias, de vontades e de forças, se não estiverem sujeitos a um principio efficaz de obrigação, que é a \emph{auctoridade}. --- E, embora os homens podessem estar unidos moral e constantemente pelo desejo e amor do fim commum, nem por isso era menos necessaria a \emph{auctoridade}. Com effeito, a cooperação social é multiplice e varia nas suas funcções. E, pois, necessario distribuir essas funcções, conforme a qualidade dos membros; pois nem todos são capazes de fazer tudo. Ora esta distribuição, não estando por si determinada\pgn{731}, deve ser determinada por um principio efficaz, que possa obrigar os membros. Este principio é a \emph{a«uctoridade}. Logo a \emph{auctoridade} é essencialmente necessaria na sociedade.\par \emph{b}) A sociedade tem «\emph{analogia} com o corpo \emph{organico}. Ora, em todo o organismo, os muitos e diversos membros, as muitas e diversas faculdades precisam de ser informadas por um principio intrinseco e vital, que as conserve unidas entre si e as dirija para o fim de todo o organismo. Logo tambem o. \emph{corpo social} deve ter um principio intrinseco e efficaz, que conserve unidos os muitos e diversos membros e dirija as suas actividades para um \emph{fim commam}. Este principio intrinseco e efficaz é a \emph{auctoridade}.\par \emph{c}) A experiencia attesta que nunca tem existido sociedade sem \emph{auctoridade}. Ora este facto, tão universal e constante, mostra que a \emph{auctoridade} é essencialmente necessaria á sociedade; aliás, alguma vez, teria existido uma \emph{sociedade sem auctoridade}. --- A força do argumento augmenta, se considerarmos que a \emph{auctoridade} não se conforma muito com a inclinação que o homem tem de executar sempre as tendencias da sua vontade.\par Da doutrina exposta resultam os seguintes corollarios:\par I. Na ordem \emph{real}, basta a \emph{auctoridade} para constituir a \emph{sociedade} e para a tornar \emph{pessoa moral}. Sem \emph{auctoridade}, \& \emph{sociedade} não é \emph{pessoa moral}, porque não tem em si o principio da sua essencia, unidade e actividade.\par 11. Embora, na ordem \emph{metaphysica}, \emph{á forma} da sociedade seja constituida pela \emph{união moral} e \emph{constante} das pessoas em ordem ao fim commum ; comtudo, na ordem \emph{real}, a \emph{forma} da sociedade é a propria \emph{auctoridade}, porque é esta o principio \emph{concreto}, que constitue essa \emph{união}; ora o elemento, de que deriva a unidade de entidade e de actividade, é e chama-se \emph{forma}.\par III. A \emph{auctoridade} póde considerar-se \emph{em st} e na sua \emph{evistencia concreta}. Considerada em si, é o \emph{direito de obrigar os membros da sociedade para cooperarem para o fim commwum}, não pertence a nenhum membro, mas é propria de todo o corpo social (pois representa o seu principio vital), \emph{e} por isso é por si indifferente para existir neste ou naquelle membro. Considerada na sua \emph{evistencia concreta}, é o \emph{direito particular}, \emph{por que reside numa pessoa determinada}. Em virtude d'este \emph{direito particular}, \emph{á auctoridade}, que considerada \emph{em si} ou no seu \emph{fim} pertencia a toda a sociedade, torna-se propria de uma pessoa particular, physica ou moral.\par o) Façamos algumas observações ácerca do \emph{fim} da sociedade.\par \emph{a}) O \emph{fim commum} é o objecto \emph{formal} da sociedade; porque é por elle que existe, e é para elle que tende a \emph{união moral} e \emph{constante} das pessoas.\par \emph{b}) O \emph{fim especifica} a sociedade. Porquanto, o \emph{fim}, embora não pertença á essencia da sociedade, comtudo \emph{especifica} a \emph{união} moral e \emph{constante} das pessoas ; porque o /\emph{im} é o termo d'esta união, e toda a tendencia é \emph{especificada} pelo \emph{termo}. Se o \emph{fim} especifica a \emph{união}, deve tambem especificar a propria sociedade; visto que essa \emph{união} 6a \emph{forma} da sociedade, e o \emph{composto}\pgn{732}é determinado na sua \emph{especie} pela \emph{forma}. --- Por isso uma sociedade differirà \emph{especificamente} de outra, se o \emph{fim} d'ambas for \emph{especificamente} diverso.\par \emph{c}) O \emph{fim} é \& \emph{razão} e a \emph{regra} de todos os \emph{direitos} e \emph{deveres} na sociedade. E' a \emph{razão}; pois quem tem direito ao \emph{fim}, tem tambem direito aos \emph{meios} necessarios. Ora os \emph{deveres} e os \emph{direitos} são \emph{meios} para a consecução do fim. --- E'a \emph{regra}; porque os \emph{meios} devem ser \emph{proporcionados} ao \emph{fim}. Por isso o \emph{fim} limita tambem o exercicio da propria \emph{auctoridade}.}.

\pg{730} 151. \textbf{Divisão da sociedade}. --- A \emph{sociedade} divide-se --- a) em \emph{universal} e \emph{particular}, --- \emph{b}) \emph{simples} e \emph{composta}, --- c) \emph{necessaria} e \emph{voluntaria}, --- \emph{d}) \emph{equal} e \emph{desegual}, --- e) \emph{completa} e \emph{incompleta}, --- \emph{f}) \emph{perfeita} e \emph{imperfeita}.

a) \textbf{Sociedade} universal e particular, --- \emph{Universal} é a que abrange todos os homens, actualmente existentes, unidos pelos deveres, que lhes são impostos --- pela sua natureza, em que todos convém, --- pelo fim ultimo, que é commum a todos, --- e pela lei natural, a que todos estão subordinados. --- \emph{Particular} é a que se limita a uma parte dos homens, unidos moralmente \pg{731} \pg{732}  para a consecução de um fim especial e determinado. --- Esta divisão funda-se na \emph{extensão} da sociedade \footnote{A sociedade \emph{universal} é formada por todos os homens, actualmente existentes. Na verdade, todos os homens estão subordinados á mesma auctoridade do legislador supremo, que governa pela voz da consciencia, e todos devem tender, pelo cumprimento dos seus deveres, para o fim \emph{commum}, que é a perfeita felicidade. Estes deveres mutuos constituem o vinculo social, que une todos os homens. --- Mais perfeita do que esta é a sociedade religiosa, que, além de ser \emph{universal}, \emph{é sobrenatural}, e produz obras superiores a todas as. forças da natureza. --- E' claro que o \emph{Direito natural} se occupa principalmente da sociedade \emph{particular}.}.

.b) Sociedade simples e composta. --- \emph{Simples} é a que se compõe \emph{immediatamente} de pessoas \emph{physicas}, ou de \emph{individuos}; tal é a sociedade \emph{conjugal}. --- \emph{Composta} é a que resulta de pessoas \emph{moraes}, isto é, de outras sociedades menores; tal é a sociedade \emph{civil}. --- Esta divisão basêa-se nos \emph{elementos} da sociedade \footnote{A sociedade \emph{composta} póde admittir muitos e diversos graus, chegando a identificar-se com a sociedade \emph{universal}.}.

c) Sociedade necessaria e voluntaria. --- \emph{Necessaria} é a que se funda numa determinação da natureza ou numa ordem positiva de \textsc{Deus}; taes são a sociedade civil e a Egreja. --- \emph{Voluntaria} é a que, na sua existencia, depende do livre arbitrio dos homens; tal é uma sociedade \emph{commercial}. --- Esta divisão funda-se na diversidade da \emph{origem} da sociedade \footnote{A sociedade \emph{necessaria} existe por determinação da \emph{lei divina}, natural ou positiva, que estabelece tambem o fim, a natureza, e a constituição interna. Por isso a auctoridade social e os homens não podem mudar nem destruir nenhum dos seus elementos \emph{substanciaes}. --- A sociedade \emph{voluntaria} existe pelo arbitrio dos homens, que livremente determinam o fim particular, a constituição e a forma. Por isso os homens podem, por commum consentimento, mudar e alterar os seus elementos, embora substanciaes.}.

d) \textbf{Sociedade} egual e desegual. --- E' \emph{egual}, quando todos os seus meinbros possuem eguaes direitos, relativamente a administração e ao governo. --- E \emph{desegual}, quando esses direitos pertencem a um só ou a alguns dos membros; tal é a sociedade \emph{domestica}. --- Esta divisão funda-se nos \emph{direitos} dos socios \footnote{Na sociedade \emph{egual}, a \emph{auctoridade} reside na \emph{collecção} dos membros; --- na \emph{desegual}, reside num ou em alguns, a quem todos obedecem. Em regra,-as sociedades \emph{voluntarias} são \emph{eguaes}; as \emph{necessarias} são \emph{deseguaes}.}.

\pg{733} e) Sociedade completa e incompleta. --- \emph{Completa} é a que, pelos meios que lhe são proprios, procura e attinge um bem \emph{completo} ou \emph{total} no seu \emph{genero}; tal é a sociedade \emph{domestica}, a \emph{civil}, e a \emph{religiosa}. --- \emph{Incompleta} é a que se limita a procurar e attingir um bem \emph{incompleto} ou \emph{particular} no proprio \emph{genero}; tal é uma sociedade \emph{litteraria}, \emph{industrial}, etc. --- Esta divisão funda-se na \emph{comprehensão} do \emph{fim} da sociedade \footnote{E' necessario insistir nesta distincção. O bem, que é o \emph{fim} da sociedade, póde ser \emph{incompleto}, isto é, póde ser \emph{parte} de outro bem e por isso subordinado a outro do mesmo \emph{genero}, ou póde ser \emph{completo}, \emph{total e supremo} no seu \emph{genero}. No primeiro caso, a sociedade é \emph{incompleta}, ou \emph{particular}; --- mo segundo caso, a sociedade diz-se \emph{completa}, ou \emph{total}.\par Existem tres sociedades \emph{completas}, --- a \emph{domestica}, \emph{á civil} e a \emph{religiosa}; pois estas, na sua propria esphera, procuram o bem \emph{completo} e \emph{total} no seu genero. Assim a sociedade \emph{domestica} procura tudo o que se refere á cducac\&o e desenvolvimento da criança; --- a sociedade \emph{civil} procura tudo o que concorre para a felicidade \emph{temporal} do homem; --- a sociedade \emph{religiosa} procura tudo o que se refere á felicidade \emph{eterna} do mesmo homem. --- Alem d'estas tres, não ha outras sociedades \emph{completas}. Porquanto, o homem primeiramente deve nascer e desenvolver-se, e, em seguida, deve procurar a sua felicidade temporal e eterna. Ora o homem nasce e desenvolve-se na sociedade \emph{domestica}, --- procura a felicidade \emph{temporal} na sociedade \emph{civil}, --- e consegue a felicidade \emph{eterna} na sociedade \emph{religiosa}.}.

f) Sociedade perfeita e imperfeita, --- \emph{Perfeita} é a que possue em si todos os meios necessarios para a consecução do seu fim, e que por isso não se refere a uma outra sociedade, em que encontre a perfeição ou o complemento, mas é \emph{independente} na \emph{sua esphera}; tal é a sociedade civil, a Egreja --- \emph{Imperfeita} é a que, embora \emph{completa}, comtudo não possue em si todos os meios necessarios para a consecução do fim, e que por isso se refere a outra, de que recebe alguns d'esses meios e de que por isso está \emph{dependente}; tal é a sociedade \emph{domestica}. --- Esta divisão funda-se nos \emph{meios}, de que a sociedade é dotada \footnote{A sociedade \emph{perfeita}, para ser tal, não precisa de possuir \emph{actualmente todos} os meios indispensaveis para a consecução do seu fim; basta que os possua \emph{virtualmente}, de modo que possa exigil-os de outra sociedade, sem que esta lh'os possa recusar. --- Pelo contrario, uma sociedade que não possue, nem \emph{actual} nem \emph{virtualmente}, \emph{todos} os meios necessarios para a consecução do fim, mas, em relação aos proprios meios, depende de uma autra sociedade, que lh'os póde recusar, é certamente \emph{imperfeita}.}.

152. O homem é naturalmente sociavel.

a) O homem será \emph{naturalmente} sociavel, se encontrar, unica e exclusivamente na união com os seus semelhantes, a satisfação das suas necessidades \emph{corporaes}, \emph{intellectuaes} e \emph{moraes}. Ora o homem encontra, unica e exclusivamente na união com \pg{734}  os seus semelhantes, a satisfação de todas essas necessidades. --- Encontra a satisfação das suas necessidades \emph{corporaes}; pois, quando nasce, não possue os meios necessarios para a sua conservação e desenvolvimento, nem os póde adquirir só pelos seus esforços. --- Encontra a satisfação das suas necessidades \emph{intellectuaes}; pois a sciencia e as artes só se apprehendem depois de muito tempo e estudo, sob a disciplina de pessoas sabias. --- Encontra a satisfação das suas necessidades \emph{moraes}; pois a virtude tem o seu estimulo no exemplo e conselho dos bons. Logo o homem é \emph{naturalmente} sociavel \footnote{Desenvolvamos as partes, que constituem a proposição \emph{menor} do syllogismo, exposto no texto.\par a) O homem só na sociedade póde encontrar a satisfação das suas necessidades \emph{corporaes}. Porquanto, ao passo que o animal,logo no seu nascimento, recebe da natureza o alimento, o vestido, os meios de defeza, que são as unhas, as pontas, os dentes, ou, ao menos, a velocidade para fugir dos males e perigos, o homem não recebeu da natureza nenhum d'esses meios, mas 80 a razão. Ora o homem não póde, só por si, desenvolver a sua razão de modo que esta possa ihdicar-lhe tudo o que a necessario para a sustentação e conservação da sua vida. Logo o homem só na sociedade póde encontrar a satisfação das suas necessidades \emph{corporaes}.\par 5) O homem só na sociedade póde encontrar a satisfação das sua necessidades \emph{intellectuaes}. Porquanto, as sciencias e as artes são necessarias para o desenvolvimento \emph{intellectual} do homem. Mas o homem, só por si, n\&o póde cultivar, com a relativa \emph{proficiencia}, as sciencias o as artes; pois toda a vida do homem não basta para a \emph{descoberta} das verdades relativas a uma só Sciencia ou a uma só arte. Por isso é preciso que sabios experimentados communiquem aos outros o fructo dos seus trabalhos e das suas lucubrações. Este ensino só é possivel na sociedade. Logo o homem só na sociedade póde encontrar a satisfação das suas necessidades \emph{intellectuaes}.\par \emph{c}) O homem só na sociedade póde encontrar a satisfação das suas necessidades moraes. Porquanto a pratica da virtude só na sociedade encontra o seu estimulo, que consiste no exemplo e na exhortação das pessoas boas, na direcção das leis \emph{positivas}, divinas e humanas, etc. --- Além d'isso, ha muitas virtudes, como a justiça, á misericordia, a caridade, a esmola, que só na sociedade podem praticar-se e desenvolver-se. --- Logo o homem só na sociedade póde encontrar a satisfação das suas necessidades moraes.}.

b) O homem será \emph{naturalmente} sociavel, se a natureza o tiver dotado de faculdades, que têm uma relação \emph{intrinseca} e \emph{necessaria} com a vida social; pois repugna que a natureza não relacione um ser com uma coisa, com a qual se relacionam essencialmente as faculdades do proprio ser. Ora a natureza dotou o homem das faculdades de \emph{fa/lar} e de \emph{ouvir}, e estas têm uma relação \emph{intrinseca} e \emph{necessaria} com a vida social ; pois é absurdo dizer que o homem tem o dom da \emph{linguagem}, só \pg{735}  para fallar comsigo mesmo, e possue o dote de \emph{ouvir}, só para escutar a sua voz. Logo o homem é \emph{naturalmente} sociavel \footnote{Não só a faculdade de \emph{ouvir} e de \emph{fallar} mostra que a vida do homem é essencialmente uma vida de \emph{relação}, mas tambem a \emph{face} e os \emph{olhos}. Porquanto a face e os olhos são espelhos da alma; não ha affeição, porquanto intima, não ha sentimento, por mais occulto, que se não revele e traduza na face e nos olhos, de modo que todos os outros podem conhecer o que uma pessoa pensa, sente, quer. Ora esta linguagem muda demonstra que \emph{o} homem foi creado para a sociedade, para dividir com os outros os seus affectos, as suas alegrias, as suas magoas.}.

c) A experiencia attesta que os homens, sempre e em toda a parte, viveram unidos em sociedade. Este facto, tão constante, universal e uniforme, deve ter uma causa constante, universal e uniforme. Ora tudo no homem é diverso e mutavel, á excepção da \emph{natureza}; porque só a natureza é universal, constante e identica em todos os homens. Logo o homem é \emph{naturalmente} sociavel \footnote{Cf. S. Thom., \emph{De regimine principum}, lib. I. --- \emph{Quatrefages} escreve: « O homem é um ser essencialmente sociavel. Um sabio da Grecia disse: « Se alguem subisse sosinho ao ceu, e sosinho escutasse as harmonias do mundo, não gosaria estas maravilhas ». Assim, encontramos por toda a parte a especie humana reunida em sociedades, mais ou menos numerosas. Sempre, com raras excepções que se explicam ordinariamente por uma diversão violenta, estas sociedades contam um numero mais ou menos consideravel de familias e merecem, pelo menos, o nome de povoações ». O mesmo diz \emph{Cousin} (\emph{Cours d'hist}. \emph{de phil}. \emph{mor}., \emph{leg}. VI).}.

153. O homem é determinado pela \textbf{natureza} a ser membro \textbf{da sociedade} domestica, civil e religiosa. --- O homem é \emph{naturalmente} sociavel; porque só na sociedade póde encontrar a sua felicidade temporal e eterna. Logo o homem é determinado pela natureza a ser membro da sociedade, em que elle encontre o seu bem \emph{absolutamente total} ou \emph{completo}, que forma a sua felicidade temporal e eterna. Ora a sociedade, em que o homem encontra o seu bem \emph{absolutamente total} ou \emph{completo}, não é uma só, mas é triplice: a \emph{domestica}, \emph{á civil} e a \emph{religiosa}. Logo o homem é determinado pela natureza a ser membro da sociedade \emph{domestica}, \emph{civil} e \emph{religiosa} \footnote{Nenhuma sociedade é \emph{absolutamente completa} e \emph{perfeita}, porque nenhuma attinge um bem \emph{absolutamente total} e nenhuma possue \emph{todos} os meios, necessarios para a felicidade temporal e eterna dos seus membros. D'ahi a insufficiencia de uma só sociedade, embora \emph{completa} e \emph{perfeita no seu genero}, e a necessidade da triplice sociedade: \emph{domestica}, \emph{civil} e \emph{religiosa}, determinada pela natureza. --- E' necessaria a sociedade \emph{domestica}; porque o homem nasce falto de todos os meios e recursos, necessarios para a sustentação e conservação da vida, e é a propria natureza que entrega a\pgn{736}criança nas mãos carinhosas dos paes, para que estes lhe dêem agasalho, soffram por ella e lhe procurem os meios para a vida. --- Mas a sociedade \emph{domestica} não basta. As artes, as sciencias, a seguranca da vida e da propriedade, são coisas que não se encontram na \emph{familia}, E' necessaria, pois, a sociedade \emph{civil}. E, como a natureza nos dotou de faculdades de \emph{relação}, e de tendencias de sociabilidade, devemos concluir que é a propria natureza que nos inclina para a sociedade \emph{civil} --- Mas a sociedade \emph{civil}, se procura a felicidade \emph{temporal}, não procura a \emph{eterna}. E' necessaria, pois, a sociedade \emph{religiosa}, cujo fim é a felicidade \emph{eterna} dos homens. E, como a natureza inclina o homem a procurar a felicidade \emph{eterna} e a professar a Religião, é necessario concluir que é á propria natureza que \emph{determina} o homem a fazer parte da sociedade \emph{religiosa}. --- Concluimos, pois, que a natureza não só fez o homem \emph{sociavel}, mas tambem o inclinou a fazer parte da triplice sociedade --- \emph{domestica}, \emph{civil} e \emph{religiosa}.\par Esta proposição é confirmada pela experiencia. No principio do mundo, encontramos a \emph{familia}. Crescidas em numero, as familias unem-se e começam a formar a sociedade \emph{civil}, mais ou. menos perfeita e desenvolvida, até se formar um corpo moral, solido e vigoroso. Dentro da sociedade \emph{civil}, ou ao seu lado, vemos a sociedade \emph{religiosa}, com os seus dogmas, ritos, sacrificios. A causa d'este facto, tão \emph{universal} e \emph{constante}, deve ser \emph{universal} e \emph{constante}, e por isso só póde ser a \emph{natureza}, pois só esta é \emph{universal}, porque ' 6 \emph{identica} em todos os homens, e é \emph{constante}, porque é \emph{immutavel}.\par O Direito \emph{natural} não trata \emph{directamente} da sociedade \emph{religiosa}, mas só \emph{indirectamente}. --- Não trata \emph{directamente}, Na verdade, a sociedade \emph{religiosa} foi elevada por N. S. Jesus \textsc{CuHRrISTO} ao estado \emph{sobrenatural}, e por isso foi dotada de \emph{fim} e de \emph{meios}, que excedem as exigencias \emph{naturaes} de uma sociedade. Sendo assim, é claro que não pertence ao Direito \emph{natural}, que se guia unicamente pela razão, mas sim \& sagrada 'theologia, que se basóa na “\emph{Revelação divina}, manifestar a essencia, o fim, os dotes da sociedade \emph{religiosa}, que, no \emph{concreto}, é a Egreja Catholica. --- Mas o Direito \emph{natural}, se não trata da sociedade \emph{religiosa directamente}, trata della \emph{indirectamente}. Porquanto, o Direito \emph{natural}, devendo considerar \& sociedade \emph{civil} em relação ás outras sociedades, no intuito de determinar os \emph{deveres} e os \emph{direitos internacionaes}, não póde deixar de a considerar tambem em relação á sociedade \emph{religiosa}; pois ha e deve haver relações entre as duas sociedades, E, tratando \emph{indirectamente} da sociedade \emph{religiosa}, o Direito \emph{natural} deve reconhecel-a como Ella é, isto é, como N. S. Jesus Onnisro a fundou; e, só depois d'este reconhecimento, baseados na luz natural da razão devemos investigar as relações e os deveres da sociedade \emph{civil} para com a \emph{religiosa}, com o fim de procurarmos a felicidade \emph{temporal} e \emph{eterna} dos homens.\par *\par Do exposto vê-se que a razão da \emph{sociabilidade} do homem se encontra na sua \emph{perfectibilidade}. O animal nasce \emph{perfeito}, e por isso o seu estado é o \emph{selvagem}; o homem nasce \emph{perfectivel}, e por isso o seu estado é o \emph{social}. --- Afastando-se d'esta grande verdade, baseada na propria experiencia, alguns escriptores sustentam que o estado \emph{primitivo} e \emph{natural} do homem foi\pgn{737}o \emph{selvagem} e que, só em virtude de uma \emph{convenção} ou \emph{contracto}, se uniram em \emph{sociedade}. Este erro foi defendido sobretudo por \emph{Hobbes} e \emph{Rousseau}, seguidos depois por \emph{Spencer} e por todos os discipulos de \emph{Darwin}, que fazem descender o homem de um anthropoide. --- Façamos uma breve exposição e critica dos systemas de \emph{Hobbes} e de \emph{Rousseau}, que são a base da \emph{revolução social e politica} dos nossos dias.\par a) \emph{Hobbes} (1588-1679) expõe as suas ideas ácerca da \emph{origem} da sociedade nos tres livros: \emph{De cive} --- \emph{De corpore politico} --- \emph{Leviathan}. Eis o resumo: O homem, no seu estado \emph{natural} e \emph{originario}, tinha duas tendencias: o egoismo ou desejo illimitado de adquirir e de gosar, --- e a solicitude de afastar \& morte e de conservar a vida. Do egoismo resultava a guerra de todos contra todos; pois a natureza tinha dado a todos o direito sobre todas as coisas. Da solicitude de conservar a vida derivava o desejo de sahir d'esse estado e de procurar companheiros. Por isso a sociedade; na sua origem, foi uma convenção de paz mutua; não foi o resultado da benevolencia, mas sim do medo e da necessidade. Constituida a sociedade, todas as vontades devem formar uma só vontade, quer seja a de um individuo, quer seja a de um conselho; e esta vontade \emph{publica} deve ter sobre a vontade de cada um o mesmo poder, que todo o individuo tinha sobre si mesmo, antes da formação da sociedade. Este poder deve ser summo e absoluto, e a elle deve estar subordinada a propria \emph{consciencia moral} dos cidadãos.\par O systema de \emph{Hobbes} é falso pelas seguintes razões: --- 1.º) Basêa-se no \emph{materialismo}, suppondo que o homem é apenas um animal aperfeiçoado, destituido de alma \emph{espiritual}. Ora estas asserções são falsas e absurdas. --- 2.º) Suppõe que o estado \emph{primitivo} do homem foi \emph{anti-social} e \emph{selvagem}, e que deste estado passou para o de \emph{civilisação}. Ora esta supposição, como vimos na \emph{Anthropologia}, é falsa; porque, se o estado primitivo do homem tivesse sido o \emph{selvagem}, elle nunca teria chegado ao estado de civilisação. O proprio \emph{Rénan} o confessa nas seguintes palavras: « Não ha um unico exemplo de povo selvagem, que se tenha, só \emph{por si}, elevado á civilisação. E' necessario, portanto, suppor que as raças civilisadas não atravessaram o estado selvagem, mas tiveram em si mesmas, desde o principio, o germen ' dos progressos futuros.» (\emph{Historia comparada das linguas semiticas}, pag. 268). --- 8.º) E' falso que o estado primitivo e \emph{natural} da humanidade foi a guerra de todos contra todos --- \emph{bellum omnium contra omnes}; não só porque o estado de guerra é violento, e tudo o que é violento não dura e por isso não é natural --- mas tambem porque a guerra é apenas um \emph{meio} para alcançar a paz, a qual é o estado natural do homem. --- Nem se diga que ha luctas e guerras entre todos os povos, e que por consequencia a guerra é o estado \emph{natural} do homem. Porquanto, a guerra não deriva da natureza (pois esta inclina ao amor mutuo) mas sim da \emph{corrupção} da natureza. --- 4,2) E' falso que o poder publico da \emph{sociedade} é tão absoluto, que a elle deva estar subordinada a propria \emph{consciencia moral}. Porquanto, acima da auctoridade \emph{publica} está a auctoridade \emph{divina}, de que a outra tira a sua origem e a sua força, e por isso a auctoridade \emph{publica} não póde preceituar á nossa consciencia moral uma coisa contraria á lei divina.\par b) \emph{Rousseau} (1712-1778), nos seus livros --- \emph{Discours sur Vorígine et les}\par 47 Vor. II.\pgn{738}\emph{fondements de l'inégalité parmi les hommes} --- \emph{Du contrat social ou principe du droit politique} --- expõe a seguinte theoria: O primeiro estado do homem foi o selvagem, dirigido unicamente pelo \emph{instincto}. As suas funcções eram meramente animaes, e os seus desejos não excediam as suas necessidades physicas. \textbar{} O homem, nesse estado, era feliz; não experimentava nenhum dos males, que o affligem agora na sociedade. --- A este primeiro periodo, de duração \textbar{} incerta, seguiu-se um outro, em que \emph{o} homem começou a desenvolver as suas faculdades (pois, segundo \emph{Rousseau}, o homem embora não diffira do \textbar{} animal em relação á intelligenocia, differe em relação á liberdade e perfectibilidade). Desenvolvendo-se gradualmente a razão com as outras faculdades, começou a communicação entre os homens, inventou-se a \emph{palavra} e fundou-se a familia. Cada familia tornou-se uma pequena sociedade, cujos laços eram unicamente a affeição reciproca e a liberdade, e que, se devia o seu prin- \textbar{} cipio á natureza (pois a \emph{familia} é a unica sociedade \emph{natural}), devia a sua conservação a um \emph{contracto}. Este periodo, em que começou a existir a sociedade e a apparecer a moralidade, devia ser a mais feliz epocha do mundo; porque o homem conservaya um meio termo entre a indolencia do estado primitivo e a petulante actividade do amor proprio. --- No terceiro periodo, os homens, mais desenvolvidos, inventaram as artes mechanicas e adquiriram a propriedade. Mas, começando a manifestar-se desegualdade nos talentos e nos caracteres, começou a accentuar-se a desegualdade entre os homens, que no principio eram eguaes. D'esta desegualdade nasceram todos os vicios, todas as paixões, todas as desordens, as prepotencias, as guerras. Para remediar a estas desordens, acabar com o estado de guerra, e obter a paz e a segurança, determinaram os homens formar uma sociedade, que defendesse e protegesse, com toda a força commum, a pessoa e os bens de cada um, e em que todo o individuo se unisse por tal modo com os seus semelhantes, que sô obedecesse a si proprio e ficasse livre, como no estado primitivo. Este escopo alcançava-se, emquanto cada um punha em \emph{communm} 2 sua pessoa e o seu poder; e assim \emph{todos} mandavam em cada um, e o individuo recebia o equivalente de tudo o que perdia, e, além d'isso, alcançava maior força para conservar tudo o que tinha. --- O \emph{contracto social}, pois, consistiu em que todo o individuo collocou a sua pessoa e o seu poder debaixo da suprema direcção da \emph{vontade geral}. Este acto de associação produziu immediatamente um corpo moral e collectivo, que é composto de tantos membros quantos são os contrahentes, e que recebe d'este mesmo acto a unidade, a personalidade, a vida e a vontade. A \emph{vontade geral} é a \emph{soberania}. --- Assim, á força foi substituida a justiça, ao impulso physico a voz do dever, ao appetite o direito, ao instincto a razão; e o homem, se perdeu a liberdade natural e o direito illimitado a tudo o que lhe agradava, ganhou a liberdade civil e a propriedade de tudo o que possuia.\par O systema de \emph{Rousseau} é falso pelas seguintes razões: --- 1.23) Suppõe que o homem, sendo naturalmente \emph{extra-social}, vivia, no principio, num estado selvagem, sem palavra, sem artes, sem religião, guiando-se unicamente pelo ?\emph{nstincto}. Ora isto é falso e gratuito; pois, como o proprio \emph{Rousseau} confessa, o homem nunca se encontrou no estado de \emph{natureza pura}, muito menos no\pgn{739}estado \emph{selvagem}. --- 2.º) Affirma que o homem, que é essencial ou naturalmente racional, era, no estado primitivo, destituido de razão, e só produzia operações sensitivas, e que tal estado era \emph{natural}, pois teria sido muito conforme com a natureza. Ora isto é absurdo. --- 3.3) Sustenta que o homem primitivo, se era destituido de intelligencia e do uso da razão, tinha comtudo a liberdade. Ora isto é contradictorio; porque a liberdade suppõe necessariamente o uso da razão. --- 4.º) Ensina --- que o homem, no estado selvagem, carecia de todos os bens, os quaes só se encontram na sociedade, --- e que, todavia, gosava uma \emph{felicidade real}. Ora isto tambem é contradictorio; pois a \emph{felicidade real} não pôde existir, quando faltam os bens indispensaveis e necessarios, como são o desenvolvimento das faculdades, a defesa da vida e da propriedade, etc. --- 5.º) Diz que os homens, depois do \emph{contracto social}, ficaram tão livres e independentes, como o eram no estado primitivo. Ora isto é inconcebivel. A liberdade é uma propriedade da vontade, e por isso quem renuncia á sua vontade, renuncia tambem á sua liberdade. Logo os que fizeram o \emph{contracto social} cederam, com os outros direitos, tambem a propria liberdade. --- 6.º) Destroe as bases essenciaes da sociedade, a constituição da familia, da propriedade, etc., subordinando tudo ao capricho da maioria. --- Mais tarde examinaremos este systema sob o aspecto \emph{politico}.}.

154. \textbf{Divisão} do \textbf{Direito social}. --- Dividimos o Direito \emph{social} em tres capitulos. No \emph{primeiro} capitulo tratamos dos \emph{deveres} e \pg{736} \pg{737}  \emph{direitos domesticos}, e por isso dos \emph{deveres} e \emph{direitos} na sociedade \emph{conjugal}, na \emph{parental} e na \emph{heril} (porque a sociedade \emph{domestica} \pg{738}  compõe-se de tres sociedades \emph{simples}: a \emph{conjugal}, a \emph{parental} e a \emph{heril}); no \emph{segundo} capitulo tratamos dos \emph{deveres} e \emph{direitos civis} ; \textbar{} no \emph{terceiro}, dos \emph{deveres} e \emph{direitos internacionaes}. \textbar{}

{\small\noindent\textsc{Summario}: --- Deveres e direitos na sociedade conjugal. --- Deveres e direi- - tos na sociedade parental. --- Deveres e direitos na sociedade heril.\par}

\pg{739} \textbf{DEVERES} E \textbf{DIREITOS DOMESTICOS}

%%>>> art-66

\section*{Artigo I.\\ Deveres e direitos na sociedade conjugal.}
\addcontentsline{toc}{section}{Artigo I. --- Deveres e direitos na sociedade conjugal.}

155. \textbf{Sociedade domestica}. --- Sociedade \emph{domestica}, ou \emph{familia}, é a união de pessoas, que vivem sob o mesmo tecto e cooperam, pelos \emph{seus actos}, para um fim commum, que é o bem de todos os membros. --- Os membros da sociedade \emph{domestica}, unidos por laços que derivam immediatamente da lei natural, estão subordinados a um superior commum, que é o \emph{pae} \footnote{Na sociedade \emph{domestica}, como em toda a sociedade, devemos distinguir dois elementos: a \emph{materia} e a \emph{forma}. \emph{A materia} é constituida pelos membros da familia; a \emph{forma} é a \emph{união}, que existe entre elles e que é produzida pela \emph{auctoridade}, a qual reside no pae.}.

156. \textbf{Partes} da \textbf{sociedade} domestica. --- A sociedade \emph{domestica é composta} de tres partes, ou sociedades \emph{simples}, que são: a sociedade \emph{conjugal}, a \emph{parental} e a \emph{heril}. --- Neste artigo tratamos \pg{740}  da sociedade \emph{conjugal}; nos dois seguintes trataremos da sociedade \emph{parental} e da \emph{heril} \footnote{A sociedade \emph{domestica} foi instituida pelo Auctor da natureza para a conservação da vida. Ora a vida conserva-se por dois modos: --- ou conservando o \emph{individuo}, e para esta conservação são necessarios os bens exteriores, que subministrem alimentos, vestidos, etc., \emph{e} são tambem hecessarios os servO8; --- ou conservando a \emph{especie}, e para isso é necessaria a união do homem e da mulher. Na sociedade \emph{domestica}, pois, nascem varias relações, --- relações entre o homem e a mulher, entre paes e filhos, entre amos e servos. D'este modo a \emph{familia} divide-se em tres sociedades \emph{simples}, que a constituem e que são a \emph{conjugal}, a \emph{parental} e a \emph{heril}.}.

157. \textbf{Sociedade} conjugal. --- Sociedade \emph{conjugal}, ou \emph{matrimonio}, é a união legitima e indissoluvel de um só homem com uma só mulher, para a geração e educação da prole e para auxilio de ambos. --- O \emph{matrimonio} é um \emph{contracto natural}, \emph{distincto} dos outros contractos, tem indole essencialmente \emph{religiosa}, \emph{é necessario} para o \emph{genero humano}, mas não para o \emph{individuo}. --- O sujeito da \emph{auctoridade}, na sociedade \emph{conjugal}, é o \emph{marido} \footnote{Apresentamos a razão por que --- \emph{a}) o \emph{matrimonio} é um \emph{contracto natuval}, \emph{distincto} dos outros contractos, --- 5») tem indole essencialmente \emph{religiosa}, --- €) é necessario para o \emph{genero humano}, mas não para o \emph{individuo}, --- \emph{d}) o sujeito da \emph{auctoridade}, na sociedade \emph{conjugal}, é o \emph{marido}.\par a) O \emph{matrimonio} é um \emph{contracto}. \emph{Contracto}, como dissemos, é um consentimento livre e expresso, pelo qual duas ou mais pessoas conferem e acceitam algum direito. Ora o \emph{matrimonio} é o consentimento livre e expresso, pelo qual os conjuges conferem e aeceitam direitos; donde resultam os \emph{direitos} e os \emph{deveres}, necessarios para a consecução dos fins d'essa união. Logo o \emph{matrimonio} é um \emph{contracto}. --- E' um contracto \emph{natural}; porque a \emph{essencia}, as \emph{propriedades} e os \emph{fins} do \emph{matrimonio} são determinados pela propria \emph{natureza} ou pela lei \emph{natural}. Por isso o \emph{matrimonio é anterior} á lei \emph{civil}, e independente d'ella. --- E' um contracto \emph{distincto} dos outros; porque os \emph{fins} do \emph{matrimonio são distinctos} dos fins dos outros contractos, e os contractos distinguem-se, conforme a distineção dos seus fins.\par b) O \emph{matrimonio} tem uma indole essencialmente \emph{religiosa}. Porquanto, os fins do \emph{matrimonio} têm uma relação estreita e immediata com a Religião. Na \emph{geração}, \textsc{Deus} tira do \emph{nada} a alma dos filhos, criando-a á sua imagem e semelhança. Na \emph{educação} da prole, \textsc{DEus} é o objecto, para o qual tendem os primeiros e os mais importantes \emph{deveres} do homem. No \emph{mutuo auxilio}, que os esposos devem prestar, \textsc{DEus} é o Pae providentissimo, que a todos dispensa as suas graças e a todos concede força para o cumprimento dos seus deveres. --- A experiencia mostra --- que, em todos os povos, o \emph{matrimonio} tem sido sempre celebrado com ceremonias religiosas, --- e que, quanto mais ia perdendo da sua indole religiosa, tanto mais ia perdendo da sua pureza e dignidade. Basta citar, entre tantos, o exemplo do povo romano. --- \emph{Affonso Martinez} diz a este respeito: « A philosophia do Direito, investigando e definindo a natureza desta união (do \emph{matrimonio}) bem póde começar estabelecendo um facto historico universal: que, em todos os tempos o em\pgn{741}todos os povos, os homens têm procurado celebrar o matrimonio com certas ceremonias religiosas, como para invocar a protecção de \textsc{Deus} e attrahir a sua benção sobre a nova familia... Só espiritos superficiaes podem ver um \emph{simples contracto} num acto tão elevado e duradouro, que não só decide da vida dos esposos, mas tambem do destino dos outros seres, e que está intimamente ligado com a propria existencia da sociedade » (\emph{Estudios sobre la Filos}. \emph{del Derecho}).\par Em conformidade com este sentimento universal e constante do genero humano, N. S. Jesus Onnisro elevou á dignidade de \emph{Sacramento} o \emph{contracto matrimonial}; de modo que, entre os christãos, o \emph{Sacramento} não se distingue do \emph{contracto}, mas todo o \emph{contracto matrimonial}, que seja feito \emph{legalmente}, \emph{é Sacramento}. E, como os \emph{Sacramentos} estão confiados á Egreja, segue-se que o Estado não tem poder algum sobre a \emph{substancia} do \emph{Matrimonio}, celebrado entre os christãos, mas 80 tem o direito de regular os seus effeitos civis, para garantir os direitos de cada um.\par Do que deixamos dito vê-se que o chamado \emph{matrimonio civil}, entre os christãos, é um attentado contra os direitos de \textsc{Deus}, auctor da natureza e da graça. Aos que dizem que o \emph{Matrimonio} é um \emph{contracto} e que todo o \emph{contracto} está sujeito ao poder civil, respondemos que o \emph{Matrimonio} não é um \emph{contracto}, como os outros, mas é um contracto \emph{elevado d dignidade de Sacramento}, e, como tal, está sujeito ao poder da Egreja.\par 0) O \emph{matrimonio} é necessario para o \emph{genero humano}, mas não para o \emph{individuo}. --- E' necessario para o \emph{genero humano}; porque, sem o \emph{matrimonio}, acabaria o mundo. --- Não é necessario para o \emph{individuo}; visto que, embora muitos se abstenham, a maior parte contrahe-o, e assim n\&o ha perigo de que se extinga o genero humano. E esta liberdade é rasoavel; pois seria uma tyrannia impor a obrigação de contrahir matrimonio a quem não tem genio, forças, recursos para satisfazer a tão graves encargos.\par O \emph{celibato}, pois, não só não é prohibido pela lei \emph{natural}, mas, quando se observa por motivo justo e elevado (como fazem as pessoas, que se consagram ao culto de \textsc{Deus} e ao bem do proximo), é mais excellente do que o matrimonio, e constitue uma das mais bellas virtudes, que possam exornar o coração humano. --- Esta verdade foi conhecida pelos proprios pagãos. Assim \emph{Pithagoras} prescreveu aos seus discipulos uma continencia absoluta; os Romanos impunham a virgindade ás Vestaes; \emph{Theophrasto}, successor de \emph{Aristoteles}, dizia que o \emph{matrimonio} n\&o convém ao sabio. (\emph{Sum}. \emph{Th}., \emph{Suppl}., q. 41, a. 2).\par Sendo injusta a lei que obrigasse os homens ao \emph{matrimonio}, não o seria menos a que o prohibisse. D'onde se vé a falsidade da theoria do economista inglez --- \emph{Thomas Roberto Malthus}. Este escriptor, no seu livro --- \emph{Ensaio sobre a população}, ensinou que, ao passo que a população humana augmenta em \emph{progressão geometrica} (1. 2. 4. 8.16.92...) os meios de subsistencia só augmentam em \emph{progressão arithmetica} (1.9.8.4.5. 6...) . E, pois, preciso que á lei prohiba o matrimonio dos proletarios e supprima as \emph{Obras pias}, que subministram dotes ás jovens solteiras. E não basta. O inglez propõe\pgn{742}outros meios, tão immoraes como funestos para a sociedade. --- Esta theoria deve rejeitar-se. Porquanto, nunca tem sido averiguada a tal proporção entre o augmento da população e o dos meios de subsistencia. --- Além disso, uma lei geral, que prohibisse o \emph{matrimonio} aos pobres, seria um attentado contra o direito \emph{natural}, que o homem tem á liberdade. Os pobres, embora faltos de bens de fortuna, podem trabalhar, e, pelo trabalho, satisfazer aos encargos do matrimonio, --- Finalmente, os meios que \emph{Malthus} aconselha para a limitação da prole em cada familia, são, como dissemos, não só immoraes, mas tambem funestos para a sociedade. A Franca, onde estes conselhos são postos em vigor mais que em qualquer outro paiz, vai diminuindo constantemente de população. Está averiguado que as victimas, que a immoralidade faz nesse paiz, em cada anno, excedem as que fez a guerra de 1870!\par 4) O \emph{sujeito} da \emph{auctoridade}, na sociedade \emph{conjugal}, é o \emph{marido}. Na sociedade \emph{conjugal}, como em toda a sociedade, deve haver uma \emph{auctoridade}; pois esta é um elemento indispensavel da sociedade. A \emph{auctoridade}, na sociedade \emph{conjugal}, deve residir no \emph{marido}; porquanto o homem, sendo dotado de maior talento, força, actividade e resistencia do que a mulher, é mais apto do que esta para tratar dos negocios mais importantes da familia e para dirigir todas as coisas para os fins da sociedade \emph{conjugal}. --- Esta auctoridade emana \emph{inmediatamente} de \textsc{Deus}; porque o que é exigido pela \emph{natureza} das coisas, não póde deixar de emanar immediatamente do Auctor da propria natureza. --- A. \emph{auctoridade} deve esforcar-se por obter aquella união de pareceres e de operações, que é necessaria para a consecução dos \emph{fins} da sociedade \emph{conjugal}. Por isso o seu exercicio deve ter por principal motor o affecto mutuo, e não deve extender-se fóra dos fins do matrimonio, nem deve lesar os direitos \emph{innatos} da mulher, como são os direitos á verdadeira \emph{liberdade de consciencia}, \emph{á dignidade}, etc.}.

158. Fim \textbf{da sociedade conjugal}. --- O \emph{fim} da sociedade \emph{conjugal} é duplice: \emph{proximo} e \emph{remoto}. O \emph{proximo} divide-se em \emph{principal} e \emph{secundario}. --- O fim \emph{proximo} e \emph{principal} da sociedade \emph{conjugal} é \emph{geração} e a \emph{educação} da \emph{prole}. --- O fim \emph{proximo} e \pg{741} \pg{742}  \emph{secundario} é o \emph{mutuo auxilio} dos \emph{contrahentes}. --- O fim \emph{remoto} da sociedade \emph{conjugal} é a salvação eterna da prole \footnote{Os fins \emph{proximos} do \emph{matrimonio} são os tres apontados. --- O primeiro fim é a \emph{geração} da prole. --- O segundo é a \emph{educação} da mesma prole; pois a criança deye desenvolver-se na intelligencia pelo conhecimento da verdade, e na vontade pelo amor e pratica do bem; e, para isso, é preciso que os paes comecem a semear nella os principios da verdade e do bem, que, desabrochando com o tempo, dêem frucios de virtude e de perfeição. --- O terceiro é o \emph{mutuo auxilio} dos conjuges, porque o homem precisa dos cuidados da mulher, e esta precisa do trabalho do homem. (\emph{Sum}. \emph{T}7., \emph{Suppl}., q. 41, a. 1.; \emph{C}. \emph{Gentes}, 1. III, c. 128).}.

159. \textbf{Propriedades} da \textbf{sociedade} conjugal. --- As principaes \emph{propriedades} da \emph{sociedade conjugal} são duas, a saber: a \emph{unidade} e a, \emph{indissolubilidade}.

a) A unidade é uma propriedade da sociedade conjugal. --- A \emph{unidade}, isto é, a união de um só homem com uma só mulher, será uma \emph{propriedade} da sociedade \emph{conjugal}, se ella fôr necessaria para a mesma sociedade conseguir os seus fins \emph{proximos}. Ora a \emph{unidade} é necessaria para a sociedade \emph{conjugal} conseguir os seus fins \emph{proximos}, que são a \emph{geração} e a \emph{educação} da \pg{743}  prole, e o mutuo auxilio dos esposos: o que é evidente. Logo a \emph{unidade} é uma \emph{propriedade} da sociedade \emph{conjugal} \footnote{A" \emph{unidade} da sociedade \emph{conjugal} oppõem-se --- a \emph{poligamia}, que é a união de um homem com duas ou mais mulheres, --- e a \emph{poliandria}, que é \& união de uma mulher com dois ou mais homens. --- \emph{A poligamia}, que é proscripta não só em todos os paizes catholicos, mas tambem em quasi todos os protestantes, não se oppõe ao primeiro fim do \emph{matrimonio}, que é a \emph{geração}; mas oppõe-se aos outros fins, que são a boa \emph{educação} da prole, a paz entre os esposos, etc. --- A \emph{poliandria} oppõe-se a \emph{todos} os fins do \emph{matrimonio}. --- N. S. Jesus \textsc{CnrISTO} restituiu ao Matrimonio o caracter \emph{monogamico}, que tinha sido o \emph{primitivo}. « Ab \emph{initio autem non fuit sic} » (\emph{Matth}., XIX, 8).\par Os \emph{positivistas} e os \emph{evolucionistas} dizem que a primeira forma do matrimonio foi a \emph{promiscuidade}, passando-se depois para a \emph{poliandria}, em seguida para a \emph{poligamia}, finalmente para a \emph{monogamia}. --- Esta theoria, sustentada por \emph{Bachofen} e \emph{Giraud-Tewlon}, é gratuita e falsa. --- E' gratuita; pois, como diz \emph{Starcke}, escriptor darwinista e por isso insuspeito, esta theoria é destituida de base e não tem provas positivas. --- E' falsa; não só porque o estado de \emph{promiscuidade} teria impedido a propagação do genero humano, mas tambem porque não se póde explicar o modo por que os homens se elevaram, exclusivamente pelos seus esforcos, do estado brutal e selvagem aos estados superiores, sabendo-se por experiencia que são necessarios immensos esforços para civilisar os selvagens e leval-os á \emph{monogamia}.}.

b) A indissolubilidade é uma propriedade da sociedade conjugal. --- \emph{A indissolubilidade} será uma \emph{propriedade} da sociedade \emph{conjugal}, se ella fôr necessaria para a mesma sociedade conseguir os seus fins \emph{proximos}. Ora. a \emph{indissolubilidade} é necessaria para a sociedade \emph{conjugal} conseguir os seus fins \emph{proximos}; porque Os esposos, se não estiverem certos de que a sua união é \emph{indissoluvel}, não procuram, com todo \emph{o} empenho e sacrificio, os meios eflicazes para a boa \emph{educação} dos filhos, não vivem em paz e harmonia, nem encontram o mutuo auxilio, que os levou a eontrahir o \emph{matrimonio}. Logo a \emph{indissolubilidade} é uma \emph{propriedade} da sociedade \emph{conjugal} \footnote{Tudo o que enfraquece a intima união de vontades, que é o fundamento da sociedade \emph{conjugal}, é contrario ao estado e á perfeição da mesma sociedade. Os encargos, que pesam sobre os esposos, são tantos e tão difficeis, que só a certeza de que a união é indissoluvel póde contel-os nos seus deveres, leval-os aos mais duros sacrificios, sujeital-os aos mais rudes trabalhos. Desde que no coração dos esposos entrar o medo de que a união possa acabar um dia, acabará a paz e a tranquillidade, começarão as mutuas desconfianças; e assim, paralysadas as forças e os trabalhos, a sociedade \emph{domestica} será um logar de desordem e de inquietações.\par \emph{A}? \emph{indissolubilidade} do \emph{matrimonio} oppõe-se o \emph{divorcio}, que é a separação dos esposos, pela qual, quebrado o vinculo conjugal, adquirem a liberdade de contrahir outro matrimonio. O \emph{divorcio} é a destruição da paz e da harmonia na familia, é a origem das paixões e das desordens, é a completa desgraça dos filhos, é a humilhação e a ruina da mulher. E' por isso que o Christianismo\pgn{744}, defensor desvelado da virtude, da justiça e da fraqueza opprimida, proclamou sempre a \emph{indissolubilidade} do matrimonio. \emph{Kwrth} escreve: « O matrimonio foi elevado á dignidade de Sacramento, e proclamado indissoluvel e inviolavel; e por isso o divorcio, que o quebrava, e o concubinato, que o manchayva, foram proscriptos com egual severidade. Protegida no seu pudor de mulher e nos seus direitos de mãe por uma legislação immensamente delicada, a esposa foi, pela primeira vez, a nobre e pura companheira do homem, sujeita, é verdade, mas livre, e levando perante \textsc{Deus} uma alma do mesmo valor que a d'elle » (\emph{Les orig}. \emph{de la civil}. \emph{mod}.).\par Não podemos negar que, em alguns casos, a \emph{éndissolubilidade} do matrimonio dá logar a graves inconvenientes. Mas tambem devemos dizer --- que esses inconvenientes podem remediar-se com a separação \emph{pessoal} (ficando illeso o vinculo), --- e que elles são coisas pequenas em comparação do mal immenso, que a instituição do \emph{divorcio} produz na familia e na sociedade. --- Os resultados do \emph{divorcio} vêm-se sobretudo na França, onde a população diminue com o augmento da desmoralisação. --- Concluimos que o \emph{divorcio}, se favorece as paixões desordenadas, destroe a familia e a sociedade, e que o bem geral deve prevalecer ao bem particular.}.

\pg{744} 160. \textbf{Deveres} e \textbf{direitos} dos conjuges. --- Os \emph{deveres} e os \emph{direitos} dos conjuges nascem, como sempre, da necessidade de conseguir os fins do \emph{matrimonio}. --- D'estes \emph{deveres} e \emph{direitos}, uns são communs a ambos os conjuges, outros são \emph{proprios} de cada um.

a) \textbf{Deveres} e direitos communs. --- Os principaes \emph{deveres} e \emph{direitos}, \emph{communs} a ambos os conjuges, são \emph{tres}: \emph{fidelidade inviolada}, \emph{amor constante} e \emph{mutuo auxilio}. Porquanto, a propria natureza do \emph{matrimonio} impõe a cada um dos conjuges os \emph{deveres} de guardar uma \emph{fidelidade inviolada} em relação ao outro, de o \emph{amar constantemente}, e de o \emph{auxiliar}; e por isso dá ao outro conjuge os \emph{direitos} correspondentes a estes \emph{deveres}.

b) \textbf{Deveres} e direitos proprios. --- Apresentamos os principaes. --- O \emph{marido} tem o \emph{dever} de dirigir, alimentar, defender a mulher, de tratar os mais importantes negocios, de escolher o domicilio, etc.; e tem o \emph{direito} á obediencia e á submissão da mulher. --- A \emph{mulher} tem o \emph{dever} de obedecer ao marido, de seguil-o relativamente ao domicilio, de tratar dos negocios domesticos, etc.; e tem o \emph{direito} ao respeito e á estima do marido \footnote{Os \emph{deveres} e os \emph{direitos} dos membros de uma sociedade fundam-se nos fins da mesma sociedade. D'onde se vê a razão por que o \emph{Direito natural}, antes de determinar os \emph{deveres} e os \emph{direitos} dos conjuges, trata da \emph{essencia} e dos \emph{fins} da sociedade \emph{conjugal}.\par Nada dissemos ácerca da \emph{capacidade juridica} dos conjuges, em relação aos bens materiaes. D'este assumpto occupa-se principalmente o Direito \emph{positivo}. O Direito \emph{natural} limita-se a dizer o seguinte: --- ambos os conjuges devem concorrer, com os seus bens, para a realisação dos fins do matrimonio; --- o marido póde dispor livremente dos bens que lhe pertencem (com tanto que não prejudique os interesses da casa); --- o mesmo marido tem o \emph{direito} de\pgn{745}administrar e dispor do rendimento dos bens da mulher (com tanto que respeite o direito de \emph{propriedade}); --- a mulher precisa, para todo o acto juridico, da auctorisacao do marido.\par Resta-nos dizer alguma coisa ácerca da \emph{capacidade} e da \emph{incapacidade} para contrahir o \emph{matrimonio}. -- \emph{Capacidade} é a \emph{aptidão} necessaria para contrahir o matrimonio e conseguir os seus fins. \emph{Incapacidade} é a \emph{inaptidão} para o matrimonio. --- A \emph{éncapacidade} é \emph{absoluta} ou \emph{relativa}, conforme é a inaptidão para contrahir o matrimonio com toda e qualquer pessoa, ou só com algumas pessoas. --- A incapacidade \emph{absoluta é physica} ou \emph{moral}, conforme deriva de impedimentos \emph{physicos}, ou \emph{moraes} (como a loucura, ete.) --- São \emph{absolutamente} incapazes de contrahir o matrimonio os que não podem conseguir os \emph{fins} do matrimonio, como são as crianças, os loucos, etc. --- Os parentes são \emph{relativamente} incapazes de contrahir o matrimonio com os seus ascendentes ou descendentes em linha recta, os irmãos com as suas irmãs. A prohibicao do matrimonio entre parentes proximos funda-se em razões moraes e physiologicas. --- As \emph{incapacidades} constituem os chamados \emph{impedimentos} do matrimonio, que podem ser \emph{dirimentes} ou simplesmente \emph{émpedientes}. --- Entre os christios, só á Egreja compete o direito de estabelecer os \emph{impedimentos}, quer \emph{dirimentes} quer \emph{émpedientes}, do Matrimonio; porque á Egreja, e só a Ella, está confiada por \textsc{Deus} a administração dos Sacramentos.}.

%%>>> art-67

\section*{Artigo II.\\ Deveres e direitos na sociedade parental.}
\addcontentsline{toc}{section}{Artigo II. --- Deveres e direitos na sociedade parental.}

%%>>> t261

\pg{745} \tese{261}{Sociedade parental.}
Sociedade \emph{parental} é a união entre os paes e os filhos, instituida para a \emph{educação} dos mesmos \emph{filhos}. --- Esta sociedade tambem se compõe de \emph{materia} e \emph{forma}. A \emph{materia} é constituida pelas \emph{pessoas} dos paes e dos filhos. A \emph{forma} é a \emph{união} entre essas pessoas. --- A \emph{união} é produzida pela \emph{auctoridade}, que reside no \emph{pae}, e, na falta d'este, na \emph{mãe} \footnote{Os \emph{communistas} negam a \emph{existencia} da sociedade \emph{parental}, e tam aos paes toda a auctoridade sobre seus filhos, dizendo que estes hão de estar immediatamente sujeitos á auctoridade \emph{civil}, e d'esta devem receber a educação. Seguem a opinião dos \emph{communistas} os que sustentam as escholas leigas e obrigatorias, fundadas pela auctoridade \emph{civil} para a educação das crianças. --- Esta opinião é falsa. A sociedade \emph{parental} é verdadeira \emph{sociedade}. Por quanto, a \emph{sociedade} é a união de duas ou mais pessoas, que, pelos seus actos, cooperam para a consecução de um fim commum. Ora a sociedade \emph{parental} é união entre paes e filhos, que, pelos seus actos, cooperam para a consecução de um fim commum --- a conservação e a educação da prole. Logo á sociedade \emph{parental} é verdadeira \emph{sociedade}. --- Dizem que a sociedade \emph{civil} é a grande familia, \& qual por isso pertencem todos os homens. Mas não se lembram de que a sociedade \emph{civil} não é uma \emph{fumilia} propriamente dicta, mas é uma sociedade mais elevada, instituida não para a destruição da familia, mas para a sua defesa e conservação.\par A sociedade \emph{parental} é uma sociedade \emph{desegual}. Na verdade, os filhos\pgn{746}são \emph{inferiores} aos paes --- em razão da \emph{geração}, pela qual receberam d'elles a existencia, --- em razão da \emph{cohabitação}, porque os filhos, quando nascem, encontram a auctoridade já estabelecida no pae e, depois d'este, na máe, --- em razão da \emph{educação}, que recebem dos seus paes. --- A \emph{auctoridade}, pois, que constitue a \emph{união} na sociedade \emph{parental}, deve residir no \emph{pae}, como chefe da sociedade \emph{domestica} e como \emph{superior} da mulher, e, depois d'elle, na mãe.\par \emph{Ahrens} affirma que o direito do \emph{patrío poder} deriva exclusivamente da \emph{educação}. --- Como se vé, este escriptor confunde o \emph{fim} do \emph{patrio poder} com a sua \emph{rais}. E' certo que o \emph{fim} do \emph{patrio poder} é a \emph{educação} da prole, mas é certo tambem que a \emph{rais} d'esse poder é a natureza da sociedade \emph{parental}.}.

%%>>> t262

\pg{746} \tese{262}{Fim da sociedade parental.}
O \emph{fim} da sociedade \emph{parental} é, como dissemos, a \emph{educação} dos \emph{filhos} \footnote{Os filhos nascem destituidos de todos os meios para a sua educação não só \emph{physica}, mas tambem \emph{intellectual} e \emph{moral}. Para esta educação, são necessarios muitos esforços, cuidados e sacrificios. Estes esforços, cuidados e sacrificios exigem um grande espirito de abnegação e um grande affecto da parte de quem deve subministrar a mesma educação. Este espirito de abnegação e este affecto só se encontram nos paes. O /\emph{im}, pois, da sociedade \emph{parental} é a \emph{educação} dos \emph{filhos}.}.

%%>>> t263

\tese{263}{A sociedade parental é necessaria.}
Uma sociedade é \emph{necessaria}, quando é determinada immediatamente pela natureza. Ora a sociedade \emph{parental} foi determinada immediatamente pela natureza. Porquanto a natureza, exeitando no coração dos paes sentimentos de affecto e de desvelo para com os seus filhos, que, quando nascem, precisam de tudo e não podem procurar-se coisa alguma, une a indigeneia d'estes com o amor e a actividade d'aquelles para a consecução de um fim, --- a educação dos mesmos filhos. Logo a sociedade \emph{parental é necessaria} \footnote{\emph{Puffendorf}, embora reconheça que a união dos paes com os filhos constitue uma verdadeira sociedade, comtudo affirma que esta sociedade não é necessaria, isto é, não deriva da \emph{natureza}, mas é \emph{voluntaria}, isto é, deriva de um \emph{consentimento} implieito entre paes e filhos. --- Repetimos que a sociedade \emph{parental á necessaria}, --- não só porque o seu fim, que é a educação dos filhos, é \emph{necessario}, --- mas tambem porque \textsc{DEus}, Auctor da natureza, deu ao coração dos paes e dos filhos taes impulsos e sentimentos, que levam uns e outros a constituir e conservar a sociedade \emph{parental}.}.

%%>>> t264

\tese{264}{A sociedade parental deve subministrar aos filhos a educação physica, intellectual e moral.}
A natureza entrega aos cuidados dos paes o filho, para elles o educarem em conformidade com as necessidades, a que está sujeito. Ora o filho, sendo \emph{homem}, é dotado de faculdades \emph{physicas}, \emph{intellectuaes} e \emph{moraes}, que precisam da respectiva educação e desenvolvimento. Logo a sociedade \emph{parental} deve subministrar aos filhos a educação \emph{physica}, \emph{intellectual} e \emph{moral} \footnote{Os \emph{communistas}, seguidos por muitos defensores do \emph{despotismo} moderno, chamado (por ironia) \emph{liberalismo}, querem, como dissemos, que a\pgn{747}educação das crianças, sobre tudo a \emph{intellectual} e \emph{moral}, seja confiada ao Estado, que deve instituir escholas \emph{obrigatorias}, em que, todavia, não se falle de Religião nem de moral. --- Esta theoria é \emph{injusta} e \emph{absurda}. E' \emph{injusta}; porque despoja os paes do direito, que a natureza lhes concedeu, de educar os filhos, tratando assim os cidadãos como outros tantos escravos do Estado. E' \emph{absurda}; porque, ao passo que apregoam a \emph{liberdade de consciencia}, tiram aos paes a liberdade de educar seus filhos como melhor lhes parece, para os obrigar a receber a educação dos mercenarios nomeados pelo Estado. --- O Estado, pois, se tem direito de abrir escholas, não o tem de obrigar os paes a enviar a ellas seus filhos. --- E' bom que os filhos saibam ler e escrever, mas é melllor que sejam justos é virtuosos. A illustração superficial, sem o temor de \textsc{Deus}, só póde produzir \emph{socialistas} e \emph{materialistas}. Basta o exemplo da \emph{Alemana}. --- Em poucas palavras. A \emph{familia} é anterior á sociedade \emph{civil}. A natureza entrega \emph{Z/mmediatamente} os filhos aos paes, e não á sociedade \emph{civil}. Logo a \emph{educação} dos filhos é um direito \emph{énnato} dos paes, que o Estado deve proteger, e não destruir.}.

%%>>> t265

\pg{747} \tese{265}{Deveres e direitos dos paes.}
Enumeramos os principaes \emph{deveres} e \emph{direitos} dos paes para com os filhos.

a) Deveres dos paes para com os filhos. --- Os paes têm o \emph{dever} de dar aos seus filhos a educação \emph{physica}, \emph{intellectual} e \emph{moral}. --- Em relação á educação \emph{physica}, os paes devem alimentar os filhos; cuidar da sua saude; prover á sua subsistencia futura, ou deixando-lhes os proprios bens, ou acostumando-os ao trabalho. --- Em relação á educação \emph{intellectual}, os paes devem ensinar aos filhos as principaes verdades, que são necessarias para o exacto cumprimento dos seus deveres para com \textsc{Deus}, para comsigo e para com os outros homens, inferiores, eguaes ou superiores. --- Em relação á educação \emph{moral}, os paes devem, com os exemplos e com os conselhos, inspirar no coração dos filhos o amor de \textsc{Deus} e do proximo, a inclinação para a virtude e a fuga do vicio, empregando tambem os meios da correcção e do castigo \footnote{E' claro que os principaes \emph{deveres} dos paes se, referem á educação \emph{intellectual} e \emph{moral} dos filhos. Para isso é indispensavel o elemento \emph{religioso} ; pois, como diz \emph{Claudio Janet}, « o ideal da vida da familia, que corresponde á verdadeira natureza humana, e cuja noção a recta razão conservou, apesar dos impulsos para o mal, devidos ao peccado original, foi fixado plenamente pelo Evangelho e pelos ensinos da Egreja. A graça de Jesus Onmisro deu á fraqueza humana o modo de praticar esta doutrina ». \emph{M}. \emph{Davas} demonstra esta verdade no artigo --- \emph{La constitution de la famille}, ete., publicado na revista --- \emph{La Réforme sociale}, 15 de Julho de 1886. :\par \emph{Rousseau} ensina, no seu \emph{Emilio}, que aos filhos não se deve fallar em \textsc{Deus}, na Religião, na alma, emquanto não tiveram attingido, pelo menos, a etade de dezoito annos. --- Esta doutrina é \emph{impia e inepta}. --- E' \emph{impia}. Porquanto, o homem, desde o uso da razão, tem a obrigação de cumprir os deveres, a que está sujeito, principalmente os religiosos. --- E' \emph{inepta}. Porquanto\pgn{748}, o proprio \emph{Rousseau} diz que os jovens devem esforçar-se por adquirir as virtudes necessarias. Mas esse esforço não é possivel, sem o conhecimento da existencia de \textsc{Deus}, da nobreza da alma, da importancia da salvação eterna, --- conhecimento, que é indispensavel para a repressão e mortificação das paixões desordenadas, que agitam o coração do joven, sobretudo na perigosa edade de doze annos, como o proprio \emph{Rousseau} confessa.}.

b) Direito dos paes para com os filhos. --- Os \emph{direitos} dos paes são \emph{perpetuos} e \emph{temporarios}. --- Os \emph{perpetuos} resumem-se no \emph{direito}, que os paes tém ao amor, a gratidão, á reverencia e (se precisarem) ao auxilio dos filhos. Estes \emph{direitos} são \emph{perpetuos}, porque duram sempre. Os \emph{temporarios} consistem no \emph{direito}, que os paes têm de dar aos seus filhos a educação \emph{physica}, \emph{intellectual} e \emph{moral}. Por isso os paes podem dirigir e obrigar os filhos, --- julgar dos seus actos, --- leval-os á submissão pelo emprego de penas proporcionadas. Estes direitos são \emph{temporarios}, porque acabam quando os filhos attingem a \emph{madoridade} e abandonam a casa paterna \footnote{A somma dos \emph{direitos}, que os paes têm para com os filhos, constitue o \emph{patrio poder}, que por isso é o \emph{conjuncto dos direitos}, \emph{que existem} no \emph{pae}, \emph{e}, \emph{na falta d'este}, \emph{na mãe}, \emph{sobre} as \emph{pessoas e os bens de seus filhos}. --- \emph{O patrio poder} tem os seus limites, provenientes da propria \emph{natureza} da sociedade \emph{parental} e dos seus \emph{fins}, que são a educação \emph{physica}, \emph{intellectual} e \emph{moral} dos filhos. Por isso os paes não podem fazer coisa alguma contraria ao bem moral dos filhos, ou á sua saude e integridade; e os filhos devem ser protegidos pela auctoridade publica contra os abusos do \emph{patrio poder}. --- Aqui devemos notar que a abolição do poder, que os paes, na Roma pagã, tinham sobre a vida e a morte dos filhos, é devida exclusivamente ao Christianismo. (Troplong, \emph{De Vinfluence du Christianisme sur le droit civil des Romains}).}.

%%>>> t266

\pg{748} \tese{266}{Deveres e direitos dos filhos.}
Enumeramos os principaes \emph{deveres} e \emph{direitos} dos filhos para com os paes.

a) \textbf{Deveres} dos filhos para com os paes, --- Os \emph{deveres} dos filhos para com os paes, correspondendo aos \emph{direitos} d'estes para com aquelles, são \emph{perpetuos} e \emph{temporarios}. --- \emph{Perpetuos} são os que os filhos têm de amar os seus paes, ser-lhes agradecidos, prestar-lhes reverencia, auxilial-os e allivial-os nas suas necessidades e afflicções. Estes \emph{deveres} são \emph{perpetuos}, porque duram sempre. --- \emph{Temporarios} são os que os filhos têm de se sujeitar á direcção dos paes, aos seus juizos e castigos. Estes deveres são \emph{temporarios}, pois acabam, quando os filhos attingem a \emph{maioridade}, com tanto que não vivam na casa paterna \footnote{Se os filhos, depois de attingir a maioridade, continuarem a viver na casa paterna, são obrigados a estar sujeitos á auctoridade do pae, e esta sujeição é exigida pela \emph{ordem domestica}; pois não ha ordem sem a devida sujeição á auctoridade constituida.}.

b) \textbf{Direitos} dos filhos para com os paes. --- Os \emph{direitos} dos filhos para com os paes, correspondendo aos \emph{deveres} d'estes para com aquelles, consistem no poder \emph{moral}, que os filhos têm de exigir dos paes a educação \emph{physica}, \emph{intellectual} e \emph{moral} \footnote{Os filhos não podem empregar a força \emph{physica} para levar os paes ao cumprimento dos deveres, que estes têm para com elles, mas, nas coisas de maior importancia, podem recorrer á auctoridade publica.}.

%%>>> art-68

\section*{Artigo III.\\ Deveres e direitos na sociedade heril.}
\addcontentsline{toc}{section}{Artigo III. --- Deveres e direitos na sociedade heril.}

{\small\noindent\textsc{Summario}: --- Sociedade civil, suas propriedades e elementos. --- Deveres e direitos dos soberanos. --- Deveres e direitos dos subditos.\par}

%%>>> t267

\pg{749} \tese{267}{Sociedade heril.}
Sociedade \emph{heril} é a união moral do eriado e do amo, formada para procurar o bem util á familia do proprio amo. --- A sociedade \emph{heril} é, pois, parte integrante ou complemento da sociedade \emph{domestica} \footnote{A sociedade \emph{Aeril} funda-se na diversidade das condições, em que os homens se encontram. De uma parte, a familia precisa ordinariamente de quem faça os serviços necessarios para o seu conveniente desenvolvimento; --- de outra parte, muitas pessoas, não possuindo os meios necessarios para \& sua vida corporal e espiritual, podem fazer parte de uma outra familia, em que encontrem tudo o que lhes é necessario para a consecução do seu fim, proximo e ultimo. --- A diversidade das condições funda-se, por sua vez, na desegualdade natural, que existe entre os homens, em relação ao talento, ás aptidões, e, por isso, aos bens de fortuna. É, pois, a natureza que leva os pobres e os ricos a constituir a sociedade \emph{Aeril}.\par O \emph{criado} differe do \emph{jornaleiro} e do \emph{escravo}. --- Differe do \emph{jornaleiro}; pois o \emph{criado} é parte integrante da sociedade \emph{domestica}, vive sob o mesmo tecto, e trabalha para o bem da familia, ao passo que o \emph{jornaleiro} vive sob a dependencia do amo só em relação á obra, a que se obrigou e por que recebe o salario determinado. --- Differe do \emph{escravo}; pois o \emph{criado} conserva a dignidade e os direitos essenciaes do homem, e voluntariamente se obriga a um genero de servicos, ao passo que o \emph{escravo} foi e é considerado d'uma condição inferior á humana, sem direitos, sem liberdade, sem dignidade. --- O paganismo não fazia distincção entre o \emph{criado} e o \emph{escravo}; o \emph{criado} acabou sempre por ser escravo. A Religião Christa distinguiu o escravo do \emph{criado}; conservou este, mas aboliu aquelle. Todos somos filhos de \textsc{Deus}, e por isso todos somos livres! Eis a razão por que a Egreja, vencendo a resistencia, que,oppunham o interesse e o egoismo dos principes e dos povos, tem sempre luetado pela abolição da \emph{escravatura}.}.

%%>>> t268

\tese{268}{Fim da sociedade heril.}
O \emph{fim} da sociedade \emph{Aer} é a consecução do bem util á familia do amo \footnote{O \emph{fim} da sociedade \emph{Aeril} é a consecução do bem util á familia do amo. Na verdade, se considerarmos o fim dos servicos e por isso o fim dos deveres do eriado, veremos que tudo converge sempre para a utilidade do amo e da sua familia. --- E' verdade que o criado, offerecendo e prestando os seus serviços, tem em vista exclusivamente o seu bem; mas este bem é o fim do\pgn{750}\emph{agente}, e não o fim do \emph{acto}; pois, como dissemos, o fim do \emph{acto} é sempre a \emph{utilidade} do amo ou da familia.}.

%%>>> t269

\pg{750} \tese{269}{A sociedade heril não repugna á lei natural.}
Uma coisa, que não se oppõe aos essenciaes direitos do homem e á sua tendencia para o fim ultimo, não repugna á lei \emph{natural}. Ora a sociedade \emph{heril} não se oppõe aos essenciaes direitos do criado, nem á sua tendencia para o fim ultimo; pois o criado, embora deva dirigir os seus actos externos para a utilidade do amo, comtudo conserva a sua dignidade e póde tender para o fim ultimo. Logo a sociedade \emph{Zeri}/ não repugna á lei \emph{natural} \footnote{A sociedade \emph{Aeril}, como está geralmente constituida, não é contraria á lei \emph{natural}, como pretendem os \emph{racionalistas}; porque o amo não tem \emph{direito} sobre a \emph{pessoa} do criado, de modo que possa dispor d'elle como quizer, mas só tem direito sobre o seu \emph{trabalho}. Prohibir que o homem ponha \emph{voluntariamente} a sua actividade ao serviço de um outro é lesar o direito \emph{énnato}, que o homem tem á sua \emph{independencia}. --- O homem póde obrigar-se a um serviço \emph{temporario} ou \emph{perpetuo}; embora este ultimo caso offereça muitos perigos e occasiões de desordens na sociedade domestica e na civil.\par A. sociedade \emph{Zieril} não só não é contraria á lei \emph{natural}, mas, quando informada pelo verdadeiro espirito christão, é uma eschola pratica de \emph{moralidade} e da \emph{paz social}. --- E' uma eschola de \emph{moralidade}; porque os criados, educados na piedade e na justiça, levam depois a moralidade ao seio das familias, que podem constituir. --- E' uma eschola de paz \emph{social}; porque os amos verdadeiramente christãos amam os seus criados com amor paternal e cuidam do bem e da felicidade das familias dos mesmos criados, estabelecendo-se assim uma união intima e sincera entre as Pessoas e as familias. --- Hoje, infelizmente, são poucas as familias, em que reina esta harmonia entre servos e amos, porque são poucas as familias, em que reina o espirito christão.}.

%%>>> t270

\tese{270}{Deveres e direitos dos amos.}
Enumeramos os principaes \emph{deveres} e \emph{direitos} dos amos para com os criados.

a) Deveres dos amos para com os oriados, --- Os \emph{deveres} dos amos para com os criados são: pagar-lhes a soldada pontualmente; --- tratal-os com a bondade e consideração devida ao homem, por mais humilde que seja a sua condição; --- instruil-os Sobre os seus deveres moraes, e proporeionar-lhes os meios para que os cumpram, --- não lhes impor serviços superiores ás suas forças, --- cuidar da sua saude, quando doentes, etc. \footnote{A. estes deveres os amos podem faltar por muitos modos, mas sobretudo não cuidando da alma dos criados, impondo-lhes serviços incompativeis com o cumprimento dos deveres religiosos. \textsc{DEus} disse: « \emph{Si est tibi servus fidelis}, \emph{sit tibi quasi anima tua}: \emph{quasi fratrem}, \emph{sic} cum \emph{tracta} » (Eccl., 33).}.

b) Direitos dos amos para com os criados. --- Os amos têm \emph{direito} ao serviço, amor, respeito e fidelidade dos criados.

%%>>> t271

\tese{271}{Deveres e direitos dos criados.}
Enumeramos os principaes \emph{deveres} e \emph{direitos} dos criados para com os amos.

a) \textbf{Deveres} dos criados para com os amos, --- Os \emph{deveres} dos criados são: --- prestar aos amos com zelo os serviços a que se obrigaram; --- ter-lhes amor, respeito e fidelidade \footnote{A sociedade \emph{Aeril} constitue uma relação de desegualdade entre o amo, em que reside á auctoridade \emph{domestica}, e o criado, que está sujeito a esta auctoridade. Por isso o criado tem o dever de obedecer ao amo, como ao proprio \textsc{DEus}, de quem procede toda a auctoridade, em todas as coisas, que não são contrarias aos seus direitos essenciaes, relativos 4 Religião, \& moralidade, á saude, etc.\par *}.

b) Direitos dos criados para com os amos. --- Os criados têm \emph{direito} á sua soldada, --- ao cuidado, respeito e consideração dos amos, --- a instrucção moral, á liberdade religiosa, ete. \footnote{O que temos exposto ácerca das relações entre o amo e o criado deve appliear-se ás relações entre o \emph{patrão} e o \emph{operario}. --- Este argumento é muito discutido, sobretudo em nossos dias, e devemos tratal-o com algum desenvolvimento. Fallaremos, pois, no \emph{contracto do trabalho}, no \emph{salario}, nos \emph{syndicatos profissionaes}, na \emph{greve}.\par a) Contracto do trabalho. --- E' o \emph{contracto}, \emph{pelo qual} o \emph{operario se obriga} a \emph{prestar} a \emph{sua obra por um determinado salario}. Este contracto tem uma indole especial, em quanto a obra humana não deve considerar-se como uma mercadoria qualquer, mas como uma coisa que participa da dignidade humana, e por isso deve ser regulado por leis especiaes. --- D'ahi os deveres dos \emph{operarios} e dos \emph{patroes}.\par 4) \emph{Deveres dos operarios}. --- Os operarios devem --- 1.9) prestar por inteiro e com fidelidade o que foi livremente ajustado, isto é, fazer a obra determinada no tempo marcado e com a diligencia, que ó-exigida pelo pacto, ou pela lei, ou pelo uso; --- 2.9) abster-se de causar damno ao patrão, ou servindo-se das coisas sem. parcimonia, ou destruindo as machinas, ou produzindo mercadorias avariadas; --- 3.º) não offender a pessoa do patrão, com recriminações, ameaças; --- 4.9) não abandonar o trabalho antes do tempo determinado, a não ser que haja uma causa grave.\par \emph{b}) \emph{Deveres dos patrões}. --- Uns derivam da \emph{justiça}, outros da \emph{caridade}. --- O patrão deve \emph{por justiça} --- 1.9) dar ao operario o justo salario no tempo determinado; --- 2.9) respeitar no operario a dignidade da pessoa humana, ennobrecida pelo caracter de christão; e por isso é obrigado --- a dar ao operario o tempo necessario para este satisfazer aos seus deveres religiosos, --- a afastar d'elle todas as occasióes proximas do peccado, --- a não desvial-o do cumprimento dos deveres domesticos; --- 8.9) não impor um trabalho superior ás forças, ou incompativel com a edade ou com o sexo do operario; --- 4,9) remover todos os perigos para a vida e a saude do operario, e compensar os damnos, que possam \emph{casualmente} derivar do trabalho. --- O patrão deve \emph{por caridade} --- 1.º) acolher com benevolencia e ouvir de bom grado os operarios ou os delegados d'elles; --- 2.) favorecer as varias instituições, que têm por fim melhorar a condição dos operarios, como são as escholas para as creanças, as associações cooperativas, etc., sem, comtudo, faltar \& delicadeza para com a liberdade e dignidade dos operarios.\par b) Salario. --- Dissemos que um dos deveres de \emph{justiça}, para o patrão, é\pgn{752}dar ao operario o \emph{justo salario}. \emph{Salario} é a retribuição do serviço, que qualquer individuo presta a outro, dia por dia, ou hora por hora. Mas quando é que o \emph{salario é justo}?\par \emph{a}) \emph{Diversas respostas d pergunta}. --- A eschola \emph{liberal} ensina que o salario \emph{justo} é o que o patrão offerece e o operario acceita. Mas a resposta não é acceitavel; porque o operario pode ser \emph{constrangido} pela miseria, ou pela concorrencia, a acceitar um salario inferior ao trabalho e insufficiente para a honesta sustentação da vida. --- A eschola \emph{socialista} pretende que ao operario é devido, e por \emph{justiça}, o fructo \emph{inteiro} do seu trabalho, deduzidas \emph{somente} as despesas que se fazem com a compra dos generos, ou da materia bruta, dos instrumentos, e com os que presidem ao trabalho dos operarios (o trabalho \emph{intellectual} não vale mais que o trabalho \emph{manual}). Tambem esta opinião é inacceitavel, como vimos na refutação do \emph{socialismo}, --- porque está baseada num supposto falso (\emph{a propriedade é um roubo}), --- é injusta para com os patrões, que empregam e compromettem o seu capital, --- e é prejudicial ao proprio operario, que não poderá encontrar trabalho, se os patrões não tiverem dinheiro para construir officinas, comprar machinas e outros instrumentos necessarios, e tentar novas empresas. --- A doutrina \emph{catholica}, exposta na Encyclica \emph{Rerum novarum} do S. Padre Leão XIII, 16 Maio 1891, resume-se nas seguintes proposições: --- \emph{a}) \emph{em todo} e \emph{qualquer caso}, \emph{o} patrão deve, por \emph{justiça}, dar ao operario o salario, que corresponde, no juizo de homens prudentes, ao valor economico do trabalho feito; --- 5) no \emph{estado normal da industria e do commercio}, \emph{o} patrão deve, por \emph{justiça}, dar o salario, que é necessario e sufficiente para a honesta sustentação do operario, comtanto que este trabalho com a devida solercia e actividade; --- c) o salarió ajustado com o operario capaz deve ser tal, que, unido com o salario ganho pela mulher e filhos por um trabalho proporcionado ás suas forças, chegue para a honesta sustentação \emph{da familia inteira}; e isto por \emph{justiça} na opinião de alguns escriptores, e só por \emph{caridade} na opinião de outros.\par \emph{b}) \emph{Remedios contra a insufficiencia do salario}. --- Para augmentar o\pgn{753}salario, insufficiente para a honesta sustentação do operario e da sua familia, foram excogitados varios meios, que, se não se baseam na \emph{justiça}, baseam-se na \emph{caridade}. Os principaes sao --- \emph{a}) os \emph{premios} concedidos aos operarios, que trabalham mais ou melhor; --- 5) á \emph{participação dos lucros}, sobretudo quando esses lucros excedem as esperanças e os calculos do patrão; --- c) os \emph{subsidios} concedidos aos operarios, emquanto o patrão destina uma parte dos lucros para favorecer os institutos, fundados em beneficio dos operarios, ou para garantir a vida e a sustentação dos mesmos operarios contra os casos fortuitos; --- \emph{d}) a \emph{moderada intervenção dos poderes publicos}, não emquanto devam fixar o salario, mas emquanto devem dar bom exemplo, assignando aos operarios, que trabalham nas obras do Governo, das Provincias, dos Municipios, um salario sufficiente para uma honesta sustentação dos mesmos operarios e das suas familias; --- \emph{e}) as \emph{associações dos operarios}, chamadas \emph{syndicatos profissionaes}, porque um operario, unido com outros, pode mais facilmente alcançar o que só por si não pode. --- Mas d'este ultimo meio digamos mais alguma coisa.\par o) Syndicatos profissionaes. --- Estes \emph{syndicatos} são uteis e até necessarios para os operarios, que, se ficassem entregues a si mesmos, seriam facilmente vietimas de uma concorrencia desenfreada, nem poderiam fazer valer os seus direitos, menosprezados, muitas vezes, por patrões prepotentes. --- Mas estes syndieatos devem ser --- \emph{mistos} (compostos de operarios e patrões) --- ou \emph{separados} (constituidos só por operarios, ou só por patrões)? Ha divergencia na resposta. Alguns querem os syndicatos \emph{mistos}, porque 1.9 promovem não só a mutua confiança e concordia entre patrões e operarios, mas tambem a prosperidade industrial e por isso o melhoramento economico de ambas as classes, \emph{e} porque 2.º) assim afastam efficazmente o perigo do socialismo. Outros, embora reconheçam estas vantagens, que derivam dos syndicatos \emph{mistos}, comtudo inclinam-se para os \emph{separados}, --- porque 1.º) os operarios, geralmente muito agarrados á propria independencia, dignidade, opinião, difficilmente se entendem com os patrões, geralmente auctoritarios, ambiciosos, soberbos; --- 2.º) destroe-se, pela base, á \emph{concorrencia}, emquanto os operarios, de commum accordo, pedem um salario proporcionado e sufficiente, e os patrões, tambem de commum accordo, estabelecem o justo preço dos productos e, por isso, o justo salario do operario; e, no caso de conflicto, pode recorrer-se ao arbitrio dos \emph{deputados}; --- 8.º) promove-se mais facilmente a \emph{educação dos operarios}, sobretudo se \emph{catholicos}, emquanto, em virtude dos seus estatutos, levam uma vida laboriosa e morigerada, evitam os viciosos, e, animados pelos bons exemplos e santas exhortações, cumprem fielmente os seus deveres religiosos e levam outros so amor da religião e da vida christá, --- Ambas estas opiniões são respeitaveis. --- Na \emph{pratica}, --- quando é possivel a formação dos syndicatos \emph{mistos} (como nas \emph{pequenas fabricas}), faga-se, comtanto que os patrões deixem aos operarios a liberdade de expor os seus juizos, e de fundar e dirigir varias associações, que têm por fim o bem da classe; --- quando não é possivel a formação dos syndicatos \emph{mistos} (como nas \emph{grandes fabricas}) façam-se os \emph{separados}, comtanto\par 48 Vor. IL\pgn{754}que se afaste o espirito de sedição, e se promova a religião e a moralidade. --- Nos paizes, onde ha catholicos e protestantes ou schismaticos, são apenas \emph{tolerados} os syndicatos \emph{interconfessionaes}, em vista dos perigos, a que estão ou podem estar expostos os catholicos.\par d) Greve. --- Dissemos que os operarios não devem abandonar o trabalho antes do tempo determinado, a não ser que haja uma causa grave. Quando assim se abandona o trabalho, dá-se a \emph{greve}. --- A \emph{greve} é a \emph{suspensão do trabalho}, \emph{feita de commum accordo por operarios associados}, \emph{com o fim de conseguirem do patrão a satisfação das suas pretenções}. As \emph{pretenções} podem ser um salario mais elevado, um trabalho menos duro, a reintegração de um operario associado, etc. --- Com relação á \emph{greve}, estabelecemos os seguintes principios. --- 1.9) E' \emph{iliécita á greve}, a não ser que haja uma causa grave, e sempre depois de uma deliberação madura e depois de se terem exgottado todos os meios para uma solução pacifica; porquanto a \emph{greve} traz sempre comsigo funestas consequencias, que não devem permittir-se, a não ser que d'ella tambem resultem, para o operario, alguns beneficios, a que este tem direito e que não poderia conseguir pelo emprego de outro meio. --- 2.º) Embora a \emph{greve} seja, nalguns casos, \emph{legitima}, todavia é \emph{sempre} prohibido o emprego de meios illegitimos e violentos, que causam prejuizo aos patrões, ou obrigam os outros operarios a absterem-se do trabalho (a não ser que estes operarios pertençam á mesma associação, a \emph{greve} tenha sido deliberada na forma devida, e os meios de coacção sejam \emph{moraes}, como á persuasão, a ameaça da expulsão do syndicato, etc.) --- 9.9) Quando a \emph{greve} é \emph{illegitima}, ou \emph{por si}, ou \emph{pelos meios empregados}, o Estado, que tem o direito e a obrigação de impedir tudo o que possa perturbar a ordem publica, deve intervir e fazer cessar a greve. --- O que temos dito refere-se á greve \emph{economica}; porque a greve \emph{politica}, que foi inventada nestes ultimos tempos pelo socialismo e que consiste na suspensão de todos os serviços publicos e de toda a vida da nação para fazar pressão sobre o governo, é uma forma de resistencia nem legal nem legitima, mas eminentemente anarchica.}.

\textbf{DEVERES} E \textbf{DIREITOS} CIVIS

%%>>> art-69

\section*{Artigo I.\\ Sociedade civil, suas propriedades e elementos.}
\addcontentsline{toc}{section}{Artigo I. --- Sociedade civil, suas propriedades e elementos.}

%%>>> t272

\tese{272}{Sociedade civil.}
Sociedade \emph{civil} é a união permanente de muitas familias, que, sob o mesmo regime, cooperam e tendem para um \emph{fim commum}, que é a \emph{felicidade temporal}. --- A sociedade \emph{civil} differe --- não só da \emph{domestica}, pois esta \pg{751} \pg{752} \pg{753}  se compõe \emph{immediatamente} de \emph{individuos} e foi instituida para a \emph{conservação} do genero humano, e aquella se compõe \emph{Immediatamente} de \emph{familias} e foi instituida para a \emph{perfeição} ou \emph{desenvolvimento} do mesmo genero humano, --- mas tambem da \pg{754}  \emph{religiosa}, que \emph{directamente} tende para a \emph{felicidade eterna}, e só \emph{indirectamente} procura a \emph{temporal} \footnote{A sociedade \emph{civil} derivou da \emph{domestica}. Esta derivação realisou-se por dois modos: --- por expansão natural de uma só familia, emquanto os filhos, emancipados, constituiram novas familias e permaneceram nas terras de seu pae, que, deste modo, se tornou patriarcha, concentrando em si a paternidade, o principado e o sacerdocio; --- e por mutuo accomdo entre as familias, estabelecidas na mesma terra, emquanto os filhos, emancipados, constituiram novas familias, estabeleceram-se em novas terras e uniram-se em sociedade, formando um governo e escolhendo um rei.}.

%%>>> t273

\tese{273}{Fim da sociedade civil.}
O \emph{fim} da sociedade \emph{civil} é a \emph{felicidade temporal} dos seus membros. Esta \emph{felicidade temporal} consiste no \emph{bem commum e externo}, \emph{coordenado com o bem interno e subordinado ao fim ultimo}. --- O \emph{bem}, em que consiste a \emph{felicidade temporal} da sociedade, deve ser --- \emph{commum}, porque convém a todos os membros da sociedade, --- \emph{externo}, porque a sociedade cuida da ordem externa, --- \emph{coordenado com o bem interno}, que é a \emph{virtude}; porque o homem não póde ser feliz, \pg{755}  mesmo temporalmente, se não fôr virtuoso, --- \emph{subordinado} ao \emph{fim ultimo}, porque a vida \emph{social} é uma \emph{parte} da vida humana, a qual está essencialmente \emph{subordinada} ao \emph{fim ultimo} \footnote{A questão Acerca da \emph{fim} da sociedade \emph{civil} é da maxima importancia. Por quanto, se não se determinar exactamente o fim que a sociedade quer alcançar, não se poderão regular com exactidão os actos dos seus membros, nem conhecer os deveres e os direitos que pertencem á auctoridade civil, nem julgar da natureza e excellencia da propria sociedade.\par O \emph{fim} da sociedade \emph{civil} é, como dissemos, o \emph{bem commum e externo}, \emph{coordenado com o bem interno e subordinado ao fim ultimo}. Desenvolvamos mais esta theoria. --- O \emph{fim} da sociedade \emph{civil} --- é um bem commum a todos os membros, porque se alcança com os esforços communs e porque pertence a todos; --- é um bem \emph{externo}, porque os socios empregam meios \emph{externos}, e com meios \emph{externos} só póde conseguir-se um fim ou bem \emph{externo}; --- \emph{coordenado com} o \emph{bem interno}, porque a sociedade não é um \emph{fim} para o homem, mas é um \emph{meio} para elle conseguir, com maior segurança e facilidade, o bem \emph{interno}, que constitue o fim da vida presente e que é a \emph{virtude}: ora, se o \emph{fim}. da sociedade \emph{civil} não fosse \emph{coordenado} com o bem \emph{interno}; a propria sociedade não seria util ao homem, mas muito prejudicial, e o bem \emph{temporal}, que fosse contrario ao \emph{bem interno}, não seria um verdadeiro bem, mas um verdadeiro mal; --- \emph{subordinado ao fim ultimo}, que é a \emph{felicidade} eterna, porque, se o fim da sociedade fosse contrario ao fim \emph{ultimo}, a mesma sociedade seria prejudicial ao homem. O bem \emph{commum} e \emph{externo} é o fim \emph{proximo} e \emph{directo} da sociedade \emph{civil}; --- o bem \emph{interno} dos socios é o fim \emph{proximo} e \emph{indirecto}; --- o fim \emph{ultimo}, que é a felicidade \emph{eterna}, é o fim remoto.\par D'esta theoria resulta que a sociedade civil, --- para alcançar o seu fim \emph{proximo} e \emph{directo}, deve tutelar os direitos dos cidadãos, conservar a ordem publica e a paz entre as familias, obrigar os homens ao cumprimento dos seus deveres sociaes, promover as sciencias, as artes, a industria, o commercio, etc.; --- para alcançar o seu fim \emph{proximo e indirecto}, que é o bem \emph{interno} dos socios, deve proteger e recompensar a virtude, impedir e castigar os vicios e os escandalos, não deve praticar nenhuma acção nem promulgar nenhuma lei, que sejam contrarias á Religião e á moralidade, mas deve fazer tudo o que, nos limites da sua esphera, concorre para a felicidade \emph{eterna} dos homens. --- A sociedade civil, procurando o bem \emph{interno} dos homens, procura implicitamente a felicidade \emph{eterna}, para a qual foram creados.\par Examinemos duas objecções contra a theoria exposta.\par Dizem: Se o fim da sociedade civil é o bem \emph{externo}, mas coordenado com o bem \emph{interno}, essa, sociedade influirá na ordem \emph{moral}; o que é absurdo.\par Resposta. A sociedade \emph{civil} influe na ordem \emph{moral} não \emph{directamente}, mas \emph{indirectamente}. Não influe \emph{directamente}; porque a ordem \emph{moral}, derivando de \textsc{Deus}, está \emph{directamente} sujeita \& auctoridade, que para isso foi instituida + pelo proprio \textsc{Deus} e que é a \emph{religiosa}. Mas influe \emph{indirectamente} ; porque a sociedade civil, devendo procurar o bem \emph{externo}, que se harmonise com o bem \emph{interno} dos homens, \emph{indirectamente} procura e facilita o bem \emph{interno}.\par Replicam. A sociedade civil, se procurasse, embora \emph{indirectamente}, o bem \emph{interno} dos homens, poderia preceituar actos \emph{internos}; o que é absurdo.\par Resposta. A sociedade \emph{civil} póde preceituar actos \emph{internos}, quando estes estão. necessariamente ligados com os \emph{externos}, que a \textbf{mesma} sociedade preceitua\pgn{756}; mas não póde preceituar actos \emph{internos}, quando estes não estão ligados com os \emph{externos}, ou quando os actos \emph{externos} não dependem da mesma sociedade \emph{civil}. Assim a sociedade civil, tendo o direito de exigir dos soldados o juramento de fidelidade, tem o direito de exigir que o mesmo juramento tenha as condições \emph{internas}, que o constituem.\par *\par Muitos philosophos e juristas modernos discordam da nossa doutrina, e estabelecem outros \emph{fins} á sociedade civil. Limitamo-nos a citar as opiniões de \emph{Kant}, \emph{Krause}, \emph{Spencer} e \emph{Hegel}.\par a) \emph{Kant} diz que, sendo a liberdade \emph{externa} illimitada e commum \emph{a} todos os homens, e tendo ella por seu fim o mesmo objecto, deve nascer entre os homens um mutuo conflicto, que a razão não póde approvar. D'ahi o preceito racional: A \emph{liberdade de cada um deve limitar-se de modo}, \emph{que possa coexistir com a liberdade dos outros}. Este preceito chama-se \emph{principio da coewistencia}, e d'elle emanam todos os outros preceitos juridicos. O Estado, ou a sociedade civil, é um \emph{postulado} deste \emph{principio de coexistencia}; porque, não estando os homens naturalmente dispostos a obedecer áquelle principio, foi necessario formar a sociedade civil, para que podesse obrigar, com a sua \&uctoridade, os cidadãos e defender a mutua coexistencia. Por isso a sociedade civil não tem outro fim, senão uma egual limitação da liberdade de cada um. --- \emph{Kant} é seguido por quasi todos os modernos racionalistas.\par O systema de \emph{Kant} é falso não só pelas razões, que adduzimos contra a sua theoria da \emph{origem do direito}, mas tambem pelas seguintes:\par a) E' falso que a liberdade externa é \emph{illimitada}. Por quanto, se a liberdade \emph{ideal} e \emph{abstracta} se concebe sem limites, a liberdade \emph{real} e \emph{concreta} tem sempre limites, que derivam do seu \emph{objecto}, dos \emph{deveres} do mesmo agente, e dos \emph{direitos} dos outros homens.\par \emph{b}) Se o \emph{fim} da auctoridade civil consistisse unicamente em obstar a que uma pessoa offenda a liberdade dos outros, esse fim seria mais \emph{negativo} do que \emph{positivo}; e assim a sociedade não teria um fim commum para alcançar, nem poderia haver união de intelligencias e de vontades entre os cidadãos, nem haveria um freio para a auctoridade, nem haveria sociedade.\par c) Se todos os preceitos jurídicos derivassem do \emph{principio da coecistencia}, a auetoridade civil não poderia prohibir nenhum dos actos, que são favoraveis á liberdade de uns, sem serem contrarios á liberdade dos autros. Assim, seria licito o \emph{suicidio}, \& \emph{blasphemia}, a \emph{irreligido}, o \emph{sacrilegio}, o \emph{escandolo}, e toda e qualquer acção impia e desonesta, com tanto que não offendesse a liberdade dos outros ou se fizesse com o consentimento dos interessados. Ora estas consequencias são simplesmente immoraes.\par b) \emph{Krause}, seguido por \emph{A/irens}, desenvolveu a theoria de \emph{Kant}, mas em . harmonia com o seu systema pantheista. Disse que o fim do Estado consiste em proteger o \emph{direito} dos cidadãos, tomando o \emph{direito} --- não em sentido \emph{negativo}, emquanto significa o que não se oppõe á liberdade dos outros, --- mas em sentido \emph{positivo}, emquanto significa as \emph{condições necessarias para} a consecução \emph{do fim}, \emph{que foi assignado} ao \emph{homem pela natureza racional}, e que consiste no desenvolvimento das faculdades humanas, sem relação ao ultimo\pgn{757}fim e sem subordinação ás regras da moralidade. Este desenvolvimento realisa-se por associações, independentes do Estado e destinadas para differentes fins, como são a sociedade \emph{religiosa}, a \emph{scientifica}, \emph{á commercial}, etc. O fim do Estado consiste em obstar a que uma associação perturbe a outra e em subministrar-lhes os meios para aquelle desenvolvimento.\par Esta theoria foi refutada em diversos pontos d'este compendio, quando demonstrámos que o fim do homem sobre a terra não é o desenvolvimento das suas faculdades e que o fim da sociedade não póde prescindir --- nem do bem \emph{interno}, que é a moralidade ou a virtude, --- nem do fim \emph{ultimo}.\par 0) \emph{Spencer}, partindo do conceito da \emph{evolução}, segundo a qual os seres vivos se vão aperfeiçoando por meio da \emph{lucta pela existencia}, sustenta que o Estado não deve obstar a que os organismos mais fortes triumphem e os mais fracos pereçam, mas deve fazer com que todos os seres operem livremente, com tanto que não impeçam a liberdade dos outros.\par Este systema é falso pelas seguintes razões.\par \emph{a}) Basóa-se na hypothese da \emph{evolução}. Ora esta hypothese, applicada sobretudo ao homem, é, como vimos, falsa e absurda.\par \emph{b}) A lei sociologica da sobrevivencia dos organismos mais aptos é a lei do mais forte, o por isso é a negação de toda a ordem moral, que impõe o dever de respeitar o direito dos outros, quer sejam fortes quer sejam fracos. --- E' por isso que os \emph{evolucionistos} bradam contra as obras de caridade, que julgam contrarias á obra benefica e purificadora da natureza!\par \emph{c}) Embora a lei da \emph{evolução} fosse certa e podesse applicar-se ao homem, comtudo a simples protecção do \emph{direito}, no sentido de \emph{Kant}, não seria por si sufficiente para dar a victoria aos organismos mais fortes; pois vemos que pessoas vigorosas, mas destituidas de meios, constituem familias pobres e doentias, ao passo que pessoas fracas, mas dotadas de meios e condições convenientes, fundam familias sás e florescentes.\par d) \emph{Hegel} ensina que o \emph{fim} do Estado, ou da sociedade civil, não é o bem commum de seus membros, mas é o proprio Estado, isto é, a grandeza do seu poder, a riqueza progressiva e a harmonia da sua organisac\&o, a perfeição cada vez mais completa da vida individual por meio do organismo social. O estado é o \textsc{Deus} \emph{presente}, pois é o infinito, e por isso é o \emph{fim} immovele absoluto de si mesmo; é omnipotente; tem todos os direitos, de que é unica fonte. A acção do Estado é absolutamente independente da lei moral.\par Não gastaremos palavras na refutação d'esse systema; pois elle é uma consequencia do \emph{pantheismo}, que já foi refutado. Só diremos que essa dou- \&rina contém o mais absoluto despotismo, a mais feroz tyrannia.}.

%%>>> t274

\pg{756} \tese{274}{Propriedades da sociedade civil.}
A sociedade \emph{civil é necessaria}, \emph{completa}, \emph{perfeita}, \emph{organica}, \emph{desegual}.

a) A sociedade civil é necessaria. --- Uma sociedade é \emph{mecessaria}, quando é exigida pela propria \emph{natureza}. Ora a sociedade \emph{civil} é exigida pela propria \emph{natureza}; pois o homem é essencialmente perfectivel, e não póde aperfeiçoar-se senão na sociedade \emph{civil}. Logo a sociedade \emph{civil é necessaria} \footnote{Esta proposição foi sufficientemente desenvolvida, quando iratàmos da \emph{sociabilidade} do homem. O homem, fóra da sociedade civil, n\&o póde conservar-se, ou, pelo menos, não póde desenvolver as suas faculdades e tendencias \emph{naturaes}. Logo a sociedade \emph{civil} é exigida pela \emph{natureza}, \emph{é necessaria}, \emph{e o} homem é \emph{natural} e \emph{necessariamente} cidadão, embora o não queira\pgn{758}ser. --- Esta \emph{necessidade}, porém, póde admittir excepções com relação a um ou outro \emph{individuo}, que póde viver fóra da sociedade.}.

b) A sociedade civil é completa, --- Uma sociedade é \emph{completa}, quando, pelos meios que lhe são proprios, procura e attinge um bem \emph{completo} ou \emph{total} no seu \emph{genero}. Ora a sociedade \emph{civil} procura e attinge um bem \emph{completo} ou \emph{total} no seu \emph{genero}, que é a felicidade \emph{temporal} dos homens. Logo é \emph{completa} \footnote{O fim da sociedade civil é a felicidade \emph{temporal} dos seus membros, de modo que a sociedade deve procurar e procura tudo o que, directa ou indirectamente, concorre para essa felicidade, não se limitando a uma só especie de bens temporaes. D'este modo o bem, que a sociedade civil procura, é \emph{completo} e \emph{total} no seu \emph{genero}. Por isso \& sociedade civil é \emph{completa} ; pois é o fim que especifica uma sociedade.\par Esta tendencia \emph{natural}, que o homem tem para a \emph{sociedade}, não bastou para constituir as varias e particulares sociedades civis; mas, como já notámos, foi necessario, no principio, o livre consentimento, implicito ou explicito, pelo qual os homens e as familias se obrigassem a procurar um fim commum, pelo emprego de meios communs. --- "Tal consentimento é muito diverso do \emph{contracto social}, inventado por \emph{Rousseau}; porquanto esse consentimento --- não cria a moralidade nem o direito, --- não é effeito de uma convenção, mas de uma determinação da natureza, --- não suppõe o estado selvagem do homem, mas só a existencia de varias familias, --- não é necessariamente \emph{explicito}. --- Nem se diga que, se \emph{a} formação da sociedade civil se realisou pelo livre consentimento dos homens, a mesma sociedade civil é \emph{livre} e não \emph{necessaria}. Porquanto uma coisa é a sociedade em \emph{abstracto}, e outra coisa é a sociedade em \emph{concreto}. A sociedade em \emph{abstracto é necessaria}, porque é a propria natureza que determina o homem a viver em sociedade; mas a sociedade em \emph{concreto é livre}, porque a formação das diversas sociedades civis não poude depender da natureza.}.

c) \emph{A} sociedade civil é perfeita. --- Uma sociedade é \emph{perfeita}, quando possue todos os meios necessarios para a consecução do seu fim e por isso, na esphera do mesmo fim, é independente

\emph{A} de outra sociedade. Ora a sociedade \emph{civi}? possue todos os meios necessarios para a consecução da felicidade \emph{temporal}, e, na esphera d'este fim, é independente de outra sociedade, de modo que é \emph{suprema} no seu \emph{genero}. Logo é \emph{perfeita} \footnote{Alguns escriptores negam que a sociedade civil é \emph{perfeita}; porque ella, dizem, depende, em muitas coisas, da Egreja. --- Mas erram. Uma sociedade, para ser \emph{perfeita}, deve ser independente, \emph{na sua esphera}, de qualquer outra sociedade. Ora a sociedade civil, na esphera do seu fim \emph{provimo} e \emph{directo}, que é o bem \emph{commuwm} e \emph{externo} dos seus membros, é independente da Egreja, e por isso é \emph{perfeita}. --- E' verdade que, muitas vezes, o Estado deve estar subordinado á Egreja, e isto dá-se, quando os interesses do Estado vem a contacto com os da Egreja; mas tal subordinação é apenas \emph{indirecta}, e não tira ao Estado a qualidade de sociedade \emph{perfeita}. Este ponto, que é muito importante e muito discutido, será tratado mais desenvolvidamente no capitulo III d'esta secção,}.

d) A sociadade civil é organica, --- Uma sociedade é \emph{organica}, quando se compõe de outras sociedades menores, que, embora não sejam \emph{perfeitas}, comtudo têm uma esphera de acção propria e possuem direitos proprios. Ora a sociedade \emph{civil} compõe-se de \emph{familias} e de varias \emph{associações}, que, embora sejam sociedades \emph{imperfeitas}, comtudo conservam uma certa autonomia, isto é, têm uma esphera propria de acção e possuem direitos proprios. Logo a sociedade \emph{civil é organica} \footnote{Os philosophos e juristas discutem se a sociedade é um composto \emph{organico,ou} se é um composto \emph{mechanico}. --- O composto \emph{organico} differe do \emph{mechanico}; porque este resulta de partes homogeneas, dispostas artificiosamente para a execução do movimento, a que são determinadas por uma força \emph{ectrinseca}, --- e aquelle resulta de partes heterogeneas, dotadas de especial estructura e de força \emph{intrinseca}, pela qual operam para a consecução do seu fim particular, embora em harmonia com o fim commum. --- Posto isto, dizemos que a sociedade imita o composto \emph{organico}, e não O \emph{mechanico}. --- Não imita o composto \emph{mecAhanico}. Porquanto, como demonstraremos contra \emph{Rousseau}, a sociedade não resulta apenas de individuos, que renunciaram a todos os seus direitos para os devolver ao corpo social. --- Mas imita o composto \emph{organico}. Na verdade, a sociedade civil é composta de pessoas e de sociedades menores, que são as familias e as diversas associações com os seus fins especiaes. Ora as pessoas têm muitos direitos inalienaveis, e as associações e as familias, embora estejam sujeitas ao supremo poder do Estado, possuem um fim \emph{intrinseco}, leis \emph{intrinsecas}, e direitos \emph{intrinsecos}, de que não podem ser despojadas, sem grande injustiça.\par O que deixamos dito serve-de criterio para podermos julgar do \emph{centralismo}, apregoado pelos \emph{socialistas} e \emph{liberaes}, para o qual as sociedades inferiores, privadas de todo o movimento espontaneo, não passam de meras machinas, destinadas a executar o movimento, que uma força extrinseca lhes imprime. Este systema enfraquece as forças dos membros e por isso a propria sociedade, cuja vida é a vida de cada uma das partes, subordinadas a uma força suprema, que lhe dê unidade e direcção.}.

e) \emph{A} sociedade civil é desegual. --- Uma sociedade é \emph{desegual}, quando os seus membros não possuem todos, os mesmos direitos. Ora os membros da sociedade civil não possuem todos, os mesmos direitos; porque esta sociedade é \emph{organica}, e as funcções e os direitos de uns não são eguaes ás funcções e aos direitos dos outros. Logo a sociedade \emph{civil é desegual} \footnote{A sociedade civil, sendo composta de \emph{familias} e de outras \emph{associações}, cujos membros não possuem todos os mesmos direitos, deve ser necessariamente \emph{desegual}. \emph{A egualdade absoluta}, proclamada pela revolução franceza, suppõe, como. observa Taine (\emph{Les origines de la France contemporaine}), que todos os homens nascem de maior edade, sem paes, sem passado, sem tradições, sem deveres, sem patria. --- Esta \emph{desegualdade} da sociedade civil não se oppõe \& \emph{egualdade civil}, que é propria dos individuos, segundo a qual todos os membros da sociedade têm os mesmos direitos \emph{innatos}, sagrados e inviolaveis, e todos podem adquirir novos direitos.}.

%%>>> t275

\pg{757} \pg{758} \pg{759} \pg{760} \tese{275}{Elementos da sociedade civil.}
Os elementos, que, na ordem \emph{real}, compõem a sociedade \emph{civil}, são dois: um \emph{material}, e outro \emph{formal}. --- O elemento \emph{material} é constituido \emph{immediatamente} pelas \emph{familias}, \emph{e mediatamente} pelos \emph{individuos}. O elemento \emph{formal}, de que deriva a \emph{união} entre as \emph{familias}, é a \emph{auctoridade}. --- Sendo, pois, a \emph{auctoridade} a alma da sociedade T \emph{civil}, devemos oecupar-nos d'ella, indicando a sua \emph{natureza}, o ] seu \emph{fim}, a sua \emph{origem}, o seu \emph{sujeito} e as suas \emph{funcções} \footnote{A sociedade é um corpo, não \emph{physico}, mas \emph{moral}. Ora, em todo o corpo, dotado de unidade, devemos encontrar as partes, que se unem, e o n principio, que as une. As partes, que se unem, são proximamente as \emph{familias} \textbar{} e \emph{remotamente} os \emph{individuos}; --- o principio, que as une, é a \emph{auctoridade}.\par Dizendo que a \emph{materia} da sociedade civil é constituida \emph{provimantente} pelas \emph{familias}, e remotamente pelos \emph{individuos}, afirmamos que os \emph{individuos} constituem a sociedade civil, emquanto já fazem parte da sociedade \emph{dome}- \textbar{} \emph{stica}, porque a sociedade civil formou-se primitivamente pelo desenvolvimento Successivo das \emph{familias}, e não dos \emph{individuos}. Alem d'isso, os homens nascem na familia, e, antes da propria emancipação, dependem da auctoridade paterna, não gosam dos direitos politicos, e por tanto não pertencem ao Estado, senão emquanto pertencem á familia. Advertimos que nem todos cs individuos se podem propriamente chamar \emph{cidadãos}, mas só os \emph{paes de familia}, isto é, todos os que são taes \emph{actualmente}, ou, pelo menos, representam \emph{virtualmente} a familia, em quanto têm uma proxima aptidão natural para a constituir; só estes estão proximamente habilitados Para exercer cargos Publicos; só estes constituem \emph{Immediatamente} a sociedade, ao passo que a \emph{mulher} e os \emph{filhos} só pertencem á sociedade, emquanto \emph{pertencem á familia}.\par D'ahi se vê o inconveniente da \emph{emancipação politica da mulher}. Os sequazes do \emph{liberalismo}, como \emph{John Stuart Mill}, \emph{Ed}, \emph{Laboulaye}, e outros, sustentam que á mulher se devem conceder os \emph{direitos políticos}, isto é, o direito de eleger os deputados, de exercer os cargos de juiz, de advogado, etc. --- Mas essas pretenções são inadmissiveis. Só os homens e os paes de familia são, como dissemos, os verdadeiros cidadãos, e não as mulheres, que, pela lei natural, estão sujeitas á auctoridade do proprio marido, chefe da familia, e que só por meio da familia pertencem á sociedade. Ora, se ás mulheres se concedessem os \emph{direitos políticos}, dar-se-hiam frequentes collisões entre esses direitos da mulher e os direitos do marido com grave prejuizo da auctoridade do marido e da paz domestica. Alem d' isso, o dever da mulher, imposto pela Providencia, é o governo interno da casa e, sobretudo, a educação dos filhos. Ora a mulher não poderia satisfazer convenientemente a esse dever sacrosanto, se, ao mesmo tempo, devesse exercer os seus deveres politicos; o que é evidente. Finalmente, a \emph{decencía} e o \emph{pudor} da mulher obstam a que ella se mostre frequentemente ao publico, assista aos comicios, \& falle ás multidões, etc., tanto mais que o seu talento, mais accessivel aos affectos do que Às razões, não é, em geral, capaz de comprehender convenientemente as coisas politicas e muito menos de regulal-as efficazmente.\par \emph{Rousseau} e \emph{Beccaría}, seguidos pelos racionalistas, fazendo derivar a sociedade não da natureza mas da livre vontade dos membros, dizem que a materia \emph{provima} da mesma sociedade é constituida pelos \emph{individuos} e não\par ,\pgn{761}pelas \emph{familias}. --- Esta opinião é falsa, não só porque se funda na theoria do estado \emph{selvagem} do homem primitivo e compara a sociedade a um composto \emph{mechanico}, mas tambem porque é contraria ao modo de operar da natureza, a qual, progredindo do imperfeito para o perfeito, começa a formar a sociedade no matrimonio, desenvolve-a na familia, e completa-a no Estado.\par Para a constituição do organismo politico concorre, de algum modo, tambem \emph{o territorio}. --- O \emph{territorio} n\&o pertence, rigorosamente fallando, á \emph{essencia} da sociedade. Por quanto, para a \emph{essencia} da sociedade são necessarios, e sufficientes, dois elementos: varias familias e uma auctoridade suprema, que as una e governe. Ora estes dois elementos podem subsistir sem territorio. Assim o povo d'Israel, vivendo no deserto, formava uma verdadeira sociedade civil, sob o governo de Moyses, revestido de uma auctoridade suprema. --- Mas \emph{o territorio}, se não pertence \& \emph{essencia} da sociedade, é um \emph{complemento}, uma \emph{condição necessaria} para a sua perfeição. Na verdade, a historia demonstra que um povo, uma tribu, uma familia, logo que se eleve a um certo grau de civilisagão, escolhe uma mansão fixa e determinada. Alem d'isso, a agricultura, os opificios e, sobretudo, as artes líberaes e as sciencias não se podem cultivar, se o povo não residir estayelmente num determinado territorio; e a propria evolução physica, moral e politica depende por tal forma do territorio, que da diversidade d'este depende a diversidade d'aquella. Por isso, para que se possam comprehender a indole e as necessidades particulares de um povo, é necessario considerar as suas relações para com o territorio. D'onde se vê a falsidade da opinião, que no organismo politico reconhece apenas a \emph{mera união de homens livres}, e por isso não reconhece no soberano \emph{nenhuma jurisdiccdo local} no territorio, e se sauda o Rei dos Belgas, não sauda o Hei da Belgica. --- Mas é necessario evitar um duplice erro. O primeiro é dos que afirmam quo o territorio pertence exclusivamente a \emph{todo} o povo, de modo que \emph{todo} o povo tem o direito de \emph{propriedade} (dominio \emph{directo}) e os particulares só têm o direito do \emph{usufructo} (dominio \emph{util}), O segundo é dos que, confundindo o direito publico com o privado, sustentam que todos os direitos dos principes derivam do originario dominio sobre os fundos; e assim o poder de jurisdicção não differe do direito de propriedade e os subditos mesmos são \emph{propriedades} do principe.}. \textbar{}

%%>>> t276

\pg{761} \tese{276}{Auctoridade civil.}
Auctoridade \emph{civil} é o direito de dirigir efficazmente os membros da sociedade civil para a consecução do fim commum, que é a \emph{felicidade temporal}. --- Esta auctoridade é independente e suprema no seu genero \footnote{\emph{A}' auctoridade civil, que é o \emph{direito} de dirigir os cidadãos para o fim commum da sociedade, corresponde a sujeição civil, que é o \emph{dever} de seguir o impulso da auctoridade, para a consecução do fim. --- A auctoridade civil não é um \emph{dominio}; pois o \emph{dominio} consiste no poder de dispor \emph{livremente} de uma \emph{coisa} para a utilidade \emph{propria}, --- ao passo que a \emph{auctoridade} refere-se ás \emph{pessoas}, mandando o que é \emph{justo} e tendo em vista o bem da \emph{communidade}, e não o do superior. ;\par Fallando da sociedade em \emph{geral}, demonstrámos que a \emph{auctoridade} é essencialmente necessaria em toda a sociedade. As razões, que adduzimos, têm maior vigor, applicadas á sociedade \emph{civil}; porque a união das intelligencias e das vontades, que é o fim da auctoridade, encontra os maiores obstaculos na lucta dos interesses e das aspirações dos membros d'esta sociedade.}.

%%>>> t277

\pg{762} \tese{277}{Fim da auctoridade civil.}
O \emph{fim} da auctoridade \emph{civil} é procurar a \emph{felicidade temporal} dos homens, e por isso procurar o \emph{bem commum} e \emph{externo}, \emph{coordenado com o bem interno} e \emph{subordinado ao fim ultimo}. --- O \emph{fim} da sociedade \emph{civil} determina a \emph{extensão} da propria auctoridade \footnote{O fim da auctoridade \emph{civil} identifica-se com o fim da sociedade \emph{civil}. E, como o fim da sociedade civil é a \emph{felicidade temporal}, e esta consiste no \emph{bem} commum \emph{e externo}, \emph{coordenado com o bem interno e subordinado ao fim ultimo}, assim dizemos que o fim da auctoridade civil é procurar a \emph{felicidade temporal}, e por isso o \emph{bem externo e contmum}, \emph{coordenado com o bem interno e subordinado} ao \emph{fim ultimo},\par O fim da sociedade civil determina a \emph{extensão} da propria auctoridade. Por isso as opiniões ácerca da \emph{extensão} da autoridade civil são tantas, quantas são as opiniões ácerca do \emph{fim} da sociedade civil. As principaes são as dos \emph{individualistas} e dos \emph{pantheistas}, apoiados pelos \emph{socialistas}.\par a) Os \emph{individualistas}, que, seguindo \emph{Rousseau}, \emph{Kant}, \emph{Humbold}, fazem consistir o fim da sociedade civil unicamente na protecção dos direitos dos cidadãos, dizem que a auctoridade civil se deve limitar á defesa e á protecção dos individuos contra os ataques dos seus inimigos, e por isso só póde prohibir os actos, que offendem os direitos dos outros. --- Esta theoria, como vimos, é falsa; pois ensina --- que o direito é um mero producto da vontade humana, --- que a sociedade civil é um simples aggregado de homens, dotados de eguaes direitos e deveres, --- que a soberania reside inalienavelmente nos cidadãos, --- que a liberdade absoluta dos individuos póde levar e leva effectivamente á felicidade.\par \emph{b}) Os \emph{pantheistas}, sequazes de \emph{Hegel}, consideram o Estado como a realisação da idea, do absoluto, e por isso como o fim de si mesmo. D'onde se segue que o Estado e a auctoridade podem fazer tudo o que quizerem. Esta opinião, como vimos, é absurda. --- \emph{Krause} e \emph{Ahrens}, admittindo que o fim da sociedade civil é a cultura da humanidade, dizem que a auctoridade civil é o juiz supremo e absoluto d'esta cultura e do que deve exigir-se dos individuos para a sua consecução, e d'este modo attribuem á mesma auctoridade poderes excessivos. Tambem esta theoria é inacceitavel; pois o fim da sociedade é, como vimos, diverso do fim apontado por estes escriptores. --- \emph{Bluntschli}, na sua obra --- \emph{Theoria geral do Estado}, ensina que o fim da sociedade civil é o desenvolvimento das faculdades da nação, é o aperfeiçoamento da sua vida por uma marcha progressiva, que não se opponha. aos destinos da humanidade. Esta theoria, que afasta do Estado a idea de \textsc{Deus}, proclama o despotismo das maiorias e, favorecendo excessivamente a auctoridade, offende e destroe a legitima liberdade e os direitos dos cidadãos.\par c) Os \emph{socialistas}, partindo do principio de que o fim da sociedade civil consiste em procurar \emph{directa} e \emph{immediatamente} o bem \emph{particular} e \emph{individual} dos cidadãos, concedem á auctoridade civil o direito de repartir a propriedade, organisar a producção em todas as espheras, distribuir os seus rendimentos, restringindo assim os legitimos direitos dos individuos. O Estado, para os socialistas, deve ser \emph{atAew}. --- Esta theoria é falsa; porque, concedendo ao Estado o direito de dispor de todos os direitos, do todas as propriedades e de todos os rendimentos, importa a destruição das familias e a morte da verdadeira liberdade.}.

%%>>> t278

\pg{763} \tese{278}{Origem da auctoridade civil.}
A auctoridade \emph{civi} deriva \emph{immediatamente} de \textsc{Deus}. Por quanto, \textsc{Deus}, pela lei natural, prescreve a observancia da ordem. Ora a ordem, na multidão, não se concebe sem uma auctoridade suprema. Logo \textsc{Deus}, pela propria lei natural, prescreve uma auctoridade suprema. Ora tudo o que \textsc{Deus} prescreve pela lei natural deriva, como vimos, \emph{immediatamente} do proprio \textsc{Deus}. Logo a auctoridade civil procede \emph{Immediatamente} de \textsc{Deus} \footnote{Esta proposição é uma consequencia dos dois principios acima estabelecidos: --- que a sociedade \emph{civil é necessaria}, isto é, instituida pela natureza e por isso por \textsc{DEus}, auctor da natureza, --- e que a auctoridade é essencialmente necessaria na sociedade. Porquanto, postos esses dois principios, argumentamos do modo seguinte: O ser, que é o immediato auctor e instituidor de uma sociedade, é por isso mesmo o immediato auctor e instituidor dos elementos, que a natureza da mesma sociedade exige. Ora \textsc{Deus} é o immediato auctor e instituidor da sociedade civil (porque esta sociedade é \emph{natural} ou \emph{necessaria}, e o que deriva da natureza, deriva do proprio Deus) e a natureza da sociedade civil exige uma auctoridade suprema. Logo \textsc{Deus} é tambem o immediato auctor e instituidor da autoridade civil. --- \textsc{Deus}, pois, é a causa \emph{proxima} e \emph{particular} (\emph{e} não só \emph{remota} e \emph{universal}) da auctoridade civil. --- \textsc{Deus}, porém, não concede a auctoridade civil por um modo \emph{sobrenatural} ou \emph{positivo}, mas concede-a por um modo inteiramente \emph{natural}, emquanto a auctoridade ó um elemento essencial da sociedade civil. --- Nas sociedades \emph{livres}, \textsc{DEus} é tambem causa da auctoridade, mas é causa exclusivamente \emph{primeira} e \emph{universal}.\par A esta, doutrina ácerca da \emph{origem} da auctoridade, que é a doutrina de todos os povos (pois todos os povos cercaram sempre a auctoridade de prestigio e de veneração, como proveniente de Deus), oppõem-se os \emph{protestantes}, os \emph{racionalistas}, os \emph{liberaes}, e todos os que consideram a sociedade como o simples resultado de uma convenção arbitraria. Mas ninguem, neste ponto, excedeu \emph{Rousseau}. Este escriptor, no \emph{Contracto social}, considera a auctoridade civil como o resultado da livre vontade dos socios.\par O systema de \emph{Rousseau} foi sufficientemente exposto. Aqui limitamo-nos a repetir a parte que diz respeito á \emph{auctoridade civil} e á sua \emph{origem}. Diz assim: « A liberdade é essencial ao homem, á inalienavel. Quando, pois, se institue a sociedade, deve ficar illesa a liberdade de cada um, de modo que, obedecendo á auctoridade, obedeça a si mesmo. Para isso é necessario que todo o cidadão ceda absolutamente \emph{todos} os seus direitos á, \emph{sociedade}, e não a uma pessoa particular. E, como os direitos são eguaes em todos os homens, segue-se que todos elles concorrem, com partes eguaes, para a constituição da sociedade e da auctoridade. Com esta cessão, os cidadãos nada perdem. Porquanto, cada um, sendo membro da sociedade, tem direito a uma parte egual da somma, composta dos direitos de todos e confiada a todo o corpo; e, por esta \emph{reciprocidade}, recebem tanto, quanto deram. --- D'esta alienação de direitos e de vontades particulares nasce o \emph{direito} e a \emph{vontade geral} de toda a sociedade. Nesta \emph{vontade geral} reside a \emph{soberania}, que, embora seja formada pelas vontades de todos, comtudo, emquanto é uma unidade collectiva, é superior a todos. A soberania é \emph{inalienavel} da collecção dos cida-\par x\pgn{764}dãos, é \emph{indivisivel}, \emph{é illimitada}, e não é capaz de ser \emph{representada}. O povo 6, pois, essencialmente \emph{soberano}. Os que exercem a auctoridade são simples magistrados do povo; dependem d'elle; podem ser demittidos e expulsos, sem que o povo se possa chamar \emph{rebelde}, pois usa do seu direito ».\par Esta theoria deve rejeitar-se tambem pelas seguintes razões:\par a) E' falso que a liberdade humana seja \emph{inalienavel}. Porquanto, se é \emph{inalienavel} a liberdade \emph{physica}, não é \emph{Znalienavel} a liberdade \emph{moral}, de que falla \emph{Rousseau} e que é o poder moral e inviolavel de regular livremente os proprios actos externos; porque a lei natural sujeita os filhos aos paes, o homem a. \textsc{Deus}, e confere direitos, de que derivam diversas relações.\par \emph{b}) \emph{Rousseau} admitte uma egualdade \emph{absoluta} entre todos os homens. Ora isto é falso; pois que os homens, se, considerados em \emph{abstracto}, são eguaes e gosam de eguaes direitos, --- considerados em \emph{concreto}, ou na \emph{realidade}, são deseguaes e têm direitos deseguaes.\par c) O escriptor francez --- por um lado sustenta que a liberdade individual é inalienavel, --- e por outro affirma que os cidadãos devem devolver irrevogavelmente todos os seus direitos á sociedade. Ora estas duas affirmações são contradictorias. --- E' verdade que \emph{Rousseau} diz que os homens, mesmo depois de terem devolvido a sua liberdade á sociedade, ficam tão livres como eram anteriormente, mas, como vimos, um absurdo não se remedeia com outro absurdo.\par \emph{d}) No systema de \emph{Rousseau}, toda a ordem moral e juridica deriva exclusivamente dos homens, e a lei civil é a propria vontade dos cidadãos, sem nenhuma relação a \textsc{DEus}. Ora isto é falso; porque vimos que a ordem moral e juridica, immediata ou mediatamente, deriva, \textsc{Deus}, e que a lei creada\par deve conformar-se com a lei eterna de \textsc{Deus}.\par e) E' falso que a auctoridade civil seja á \emph{somma das auctoridades individuaes}; não só porque a auctoridade civil, devendo regular o exercicio dos direitos individuaes, pertence a uma ordem superior á collecção de todos os direitos dos cidadãos; mas tambem porque ella póde naturalmente praticar actos que não podem ser praticados pela collecção dos direitos individuaes, como são declarar a guerra, punir com a pena de morte os malfeitores, etc., porque o homem, não tendo o direito de se matar a si mesmo, não póde conferir a outros o poder de o matarem.\par \emph{f}) A theoria do \emph{povo soberano} é a destruição da sociedade. Por quanto, se a vontade do povo é a fonte de todos os deveres e de todos os direitos, se é o unico principio e juiz da moralidade, da justiça, da religião, segue-se que tudo o que o povo manda é santo e tudo o que o povo prohibe é deshonesto, --- que a revolução, o despotismo, a tyrannia, o socialismo e o anarchismo são direitos, que o povo legitimamente exerce. Ora estas consequencias envolyem a destruição da sociedade.}.

%%>>> t279

\tese{279}{Sujeito da auctoridade civil.}
\emph{Sujeito} da auctoridade \emph{civil} é a pessoa, que a possue e exerce. --- O \emph{sujeito} da auctoridade civil não é determinado pela natureza, mas por algum \emph{facto humano}, como é a \emph{successdo}, a \emph{eleição}, \emph{á victoria} numa guerra \emph{justa}, ete. Determinado o \emph{sujeito}, este recebe \emph{immediatamente} de \textsc{Deus} a auctoridade. --- A pessoa, que é o \emph{sujeito} da auctoridade, póde ser \emph{physica} ou \emph{moral}. D'ahi as varias e \pg{764} \pg{765}  diversas formas de governo, que são a \emph{monarchica}, \& \emph{aristocratica}, e á \emph{democratica} \footnote{A auctoridade civil deve residir necessariamente num \emph{sujeito}. O \emph{sujeito} não é determinado pela natureza; pois esta, se exige que na sociedade exista uma auoctoridade, não exige que o sujeito seja uma certa e deéterminada pessoa. O \emph{sujeito} da auctoridade deve ser, pois, determinado por um \emph{facto humano}. --- Este facto, ou \emph{titulo}, póde considerar-se na primitiva constituigüo da sociedade, e nas sociedades posteriormente constituidas.\par Considerado na primitiva constituição da sociedade, o \emph{titulo} é o \emph{patrio poder}; pois, sendo a sociedade \emph{civil} um desenvolvimento da sociedade \emph{domestica}, o \emph{sujeito} da auctoridade \emph{civil} devia ser o mesmo que o da \emph{paterna}. Neste caso, a auctoridade \emph{paterna} não se confundia com a \emph{civil}, nem era absorvida por esta, mas era o \emph{facto}, que designava o \emph{sujeito} da auctoridade.\par Considerado nas sociedades posteriormente constituidas, o \emph{titulo} é multiplice; pois póde ser a \emph{successão hereditaria}, a \emph{eleição}, \emph{á victoria} numa guerra \emph{justa}, e a \emph{prescripção}. --- A \emph{successão hereditaria} dá-se, quando o principe transmitte ao seu herdeiro, com todos os outros bens e direitos, a auctoridade suprema. Este modo de transmittir \& auctoridade evita muitas discordias e guerras, e robustece as instituições. --- A \emph{eleição} é necessaria, quando a auctoridade não se transmitte por \emph{successão}. Neste caso, o povo, ou parte do povo, deve escolher uma pessoa dotada de todas as qualidades necessarias, e esta recebe de \textsc{DEus} a auctoridade suprema. --- A \emph{victoria} em guerra \emph{justa} é tambem um \emph{titulo} legitimo. Mas a guerra deve ser \emph{justa}; aliás a oceupação será um roubo. Além disso, é preciso que \emph{a} mudança do regimen ou a perda da independencia da parte do povo vencido: seja reclamada pela segurança do vencedor ou dos outros povos; aliás a reparação ou o castigo será superior ao crime do povo vencido. --- A \emph{prescripção} pode ser um \emph{titulo} legitimo para adquirir o poder supremo. A questão refere-se ao \emph{usurpador}. E' necessario fazer varias considerações a este respeito. O \emph{usurpador} tem sempre o \emph{dever} de restituir a auctoridade injustamente possuida ao legitimo soberano; e até póde ser obrigado a essa restituição pelo povo, comtanto que haja esperança de um exito favoravel, e a insurreição não seja prejudicial ao paiz. O povo tem obrigação de obedecer ás ordens justas do \emph{usurpador}; não só porque esta obediencia é necessaria para o bem do paiz, mas tambem porque o usurpador possue, embora illegitimamente, a auctoridade suprema, e esta conserva toda a sua força e todo o seu direito relativamente aos sub- ditos, que não obedecem ao usurpador, mas á auctoridade. Se, com o andar dos annos, a causa do usurpador por tal modo se ligar com o bem e a tranquillidade do povo, que o passado não possa repristinar-se sem uma revolução e effusüó de sangue, o usurpador torna-se legitimo, o povo não pode revoltar-se contra elle; por quanto, neste caso, é verdadeira a maxima: \emph{Salus publica suprema lex}, e a ella devem obedecer não só os subditos, mas tambem os principes, se não quizerem. preferir ao interesse publico o seu interesse particular.\par Determinado o \emph{sujeito} por algum dos \emph{factos} indicados, elle recebe \emph{immediatamente} de \textsc{Deus} a auctoridade suprema, não por modos extraordinarios e \emph{positivos}, mas polos meios com que Elle, emquanto Auctor da natureza, : governa a sociedade. --- Como se vê, rejeitamos a opinião dos que affirmam que \textsc{DEus} concedeu a auctoridade suprema a todo o povo, como ao seu \emph{sujeito}\pgn{766}, \emph{e} que o povo, não potendo exercel-a, a transfere no principe; de modo que \textsc{Deus} daria o poder ao principe só \emph{mediatamente}, isto é, só por meio do povo. Não podemos acceitar esta opinião. Na verdade, o elemento, que reune a multidão, é anterior á propria multidão e não depende d'ella. Ora a auctoridade suprema é o elemento, que reune em sociedade a multidão. Logo a auctoridade suprema é anterior \& multidão, e por isso não reside na multidão, como no seu \emph{sujeito}, \emph{e} não deriva d'ella. --- Além d'isso, a auctoridade é, como dissemos, o elemento \emph{formal} da sociedade, e a multidão é o elemento \emph{material}. Ora o elemento \emph{formal} não pode depender do elemento \emph{material}; porque o superior não pode depender do inferior. Logo a auctoridade não depende da multidão, não reside na multidão, como no seu \emph{sujeito}. --- A nossa opinião encontra uma confirmação na Encyclica \emph{Diuturnum} do S. Padre Leão XIII, que diz: « Quo sane delectu (multitudinis) \emph{designatur} princeps, \emph{non conferuntur jura principatus}: \emph{neque mandatur}. \emph{imperéum}, sed \emph{statuitur} a quo sit regendum ». Concluimos que a auctoridade é \emph{para} a multidão, mas não está \emph{na} multidão.\par Conforme o \emph{sujeito} da auctoridade civil é uma pessoa \emph{physica} ou \emph{moral}, nascem as diversas \emph{formas} do governo, que por isso são as \emph{differentes organisacües da auctoridade suprema em cada sociedade civil}. --- Se a auctoridade reside numa pessoa \emph{physica}, a forma é \emph{monarchica}; --- se reside numa pessoa \emph{moral}, a forma é \emph{polyarchica}, que se subdivide em \emph{arístocratica} e \emph{democratica}, conforme a auctoridade reside numa corporação de pessoas nobres e distinctas, ou em todo o povo. D'este modo, a \emph{monarchia}, a \emph{aristocracia} e a \emph{democracia} representam as tres \emph{formas simples} da soberania.\par Estas formas \emph{simples} podem combinar-se e dar outras formas \emph{compostas} ou \emph{mistas}. --- Todas as \emph{formas} de governo são \emph{legitimas}, comtanto que sejam \emph{justamente} estabelecidas, levem a sociedade para a consecução do seu fim, reconheçam e respeitem os principios relativos á origem divina, ao fim, ás faculdades e aos limites da auctoridade. Os vicios, que se notam na maior parte dos governos, dependem da maldade dos homens e não do vicio da forma.\par Resta-nos dizer alguma coisa ácerca de cada forma \emph{particular}.\par \emph{A monarchia} póde ser \emph{absoluta} ou \emph{temperada}. --- E' \emph{absoluta}, quando a auctoridade, que reside no monarcha, é limitada exclusivamente pela lei \emph{natural}. E \emph{temperada}, quando a auctoridade, que reside no monarcha, é limitada por um estatuto ou por costumes estabelecidos. --- A monarchia \emph{absoluta}, se o imperante tiver as qualidades necessarias, dá maior força e unidade ao regimen; mas, se o monarcha não tiver essas qualidades, é a ruina do paiz. A monarchia \emph{temperada} é muito util, se os que governam com o monarcha são pessoas honestas, competentes e amigas do bem do paiz. --- A \emph{monarchia}, quer \emph{absoluta} quer \emph{moderada}, póde ser \emph{Aereditaria} ou \emph{electiva}. --- O abuso da \emph{monarchia} é a \emph{tyrannia}.\par \emph{A aristocracia}, como \emph{classe} e como \emph{forma} de governo, offereco muitas vantagens, mas tambem póde ser occasião de discordias e miserias. --- Offerece muitas vantagens, porque os.aristocratas, sendo distinctos pelo talento, pela fortuna ou por nobres qualidades, podem melhor conhecer os males e os remedios, e com maior desinteresse e justiça exercer os cargos publicos. --- Mas póde ser occasiüo de discordias e de miserias, se os aristocratas,\pgn{767}esquecidos da sua missão, tratarem o povo com dureza e crueldade, fizerem monopolio dos rendimentos do Estado, e assim provocarem reacções e guerras. --- A \emph{aristocracia} tambem póde ser \emph{Aereditaria} ou \emph{electiva}. --- O abuso da \emph{aristocracia} é a \emph{oligarchia}.\par A. \emph{democracia} pôde ser \emph{directa} ou \emph{indirecta}. E' \emph{directa}, quando \emph{todos} os paes de familia concorrem por si para fazer leis no Conselho geral; o que se verifioava nos antigos municipios da Italia. E' \emph{indirecta}, ou \emph{representativa}, quando os cidadãos não fazem por si as leis, mas escolhem os legisladores; o que se verifica na Republica Franceza. --- A \emph{democracia}, para ser boa e util forma de governo, deve satisfazer ás seguintes condições: 1.º) a massa do povo, sobre tudo os operarios, deve ser muito religiosa e moral, --- 2.3) na sociedade devem predominar os agricultores, os chefes de familia, as pessoas mais distinctas pela sua virtude, prudencia e abnegação. Sem estas condições, a \emph{democracia} é perigosa e fatal para o paiz, arrastando-o para os horrores do \emph{anarchismo}. --- O abuso da \emph{democracia} é a \emph{demagogia}.\par A \emph{forma} de governo, que actualmente predomina e que se reduz á \emph{monarchia temperada} (pois é uma mistura de \emph{monarchia}, \emph{aristocracia} e \emph{democracia}), é a \emph{representativa}. Esta forma, que suppõe um pacto fundamental entre o povo e o principe, chamado \emph{carta}, attribue a auctoridade --- ao \emph{Rei}, que, conjunctamente com os ministros, tem o poder \emph{executivo}, --- e á duas \emph{camaras}, uma de \emph{Senadores} ou \emph{Pares} e outra de \emph{Deputados}, cujo direito e dever consistem em discutir e elaborar as leis. --- O \emph{fundamento} d'este machinismo politico é a \emph{soberania} do povo, que, não podendo por si fazer leis, escolhe quem as faga em seu nome; --- o seu \emph{fim} é tornar impossivel o abuso do poder; --- o \emph{meio} é a divisão das funcções da auctoridade, pois esta divisão é a causa de mutuas limitações. E' por isso que ao collegio dos \emph{Deputados} se oppõe o collegio dos \emph{Pares}; e, para que ambos estes collegios não abusem da sua auctoridade, as leis não são validas sem a sancção real. O poder \emph{executivo} é contido nos seus limites pela \emph{responsabilidade}, que pesa sobre os ministros, que podem ser julgados e punidos pelo poder \emph{legislativo}. Mas, como esta funcção do poder \emph{legislativo} podia ser perigosa, o Rei é \emph{inviolavel}. Este nada póde fazer sem o consentimento dos ministros; por isso o Rei \emph{reina}, \emph{mas não governa}.\par Tal \emph{forma} de governo é imperfeita. Na verdade, a auctoridade suprema, estando dividida, não é muito capaz de manter a unidade e ordem social. Embora imperfeita, esta forma tem as suas vantagens, e póde ser a mais conveniente para o povo, que já se accommodou a ella, e, sendo legitimamente introduzida, tem direito á obediencia dos cidadãos. Mas, então, para que se torne um principio duradouro de bem social, deve ser purificada dos principaes vicios, que a deturpam e que são --- o principio do \emph{povo soberano}, --- o principio de que a opinião publica deve ser o criterio do governo e da administração, --- o dogma de que cada facção politica deve, por sua vez, ser chamada ao governo, o que é motivo de discordias.}.

179. \textbf{Funcções} da \textbf{auctoridade} civil. --- As principaes \emph{funcções}, que a auctoridade \emph{civil} exerce para a consecução do fim da \pg{767}  sociedade, podem reduzir-se a tres, a saber: a funeção \emph{legislativa}, \emph{á judicial} e a \emph{executiva} \footnote{Todas as \emph{funcções}, que a auctoridade civil póde realisar para o cumprimento da sua missão, foram reduzidas por Montesquieu (\emph{L'esprit de lois}, t. L l. XI) ás tres apontadas no texto. O direito de as exercer chama-se poder \emph{legislativo}, \emph{executivo} e \emph{judicial}. --- Esta divisão não agrada a todos;\pgn{768}mas julgamos que, por ser muito commum e por dividir convenientemente o objecto nas suas partes, é acceitavel. Porquanto, a auctoridade tem os seguintes deveres. Primeiramente deve prescrever aos cidadãos as normas dos seus actos em ordem ao fim; d'ahi a necessidade do poder \emph{legislativo}. Além d'isso, deve dirimir as controversias entre os particulares, ou entre os particulares e o governo, ácerca dos direitos, que a lei concede ou defende, e punir as violações da propria lei; d'ahi a necessidade do poder \emph{judicial}. Finalmente, deve applicar ao caso concreto o que foi deliberado pelo poder \emph{legislativo} e definido pelo \emph{judicial}; d'ahi a necessidade do poder \emph{executivo}.\par \emph{Montesquieu} afirma que as tres partes da auctoridade devem ser separadas, de modo que não se reunam no mesmo sujeito.\par O jurisconsulto francez, que é seguido por todos os liberaes, confunde duas coisas, muito distinctas, que são: o \emph{exercicio} d'estas funcções e a \emph{raiz} ou o \emph{principio} de que derivam. --- Se considerarmos o \emph{exercicio} d'essas funcções, é certo que devem ser exercidas por \emph{distinctos sujeitos}; não só porque, para esse \emph{ewercicio}, se requerem muitas e diversas qualidades de intelligencia e de vontade, que difficilmente se encontram num só sujeito, --- mas tambem porque a concentração dos tres poderes num só sujeito dá logar a abusos e despotismos. --- Se considerarmos a \emph{rais} ou o \emph{principio}, de que derivam e que é a \emph{auctoridade}, é certo que exigem \emph{um só sujeito}. Porquanto, \& auctoridade, embora seja dotada de varias faculdades, todavia, em \emph{concreto}, é uma só; pois é o direito de governar a multidão. E a auctoridade, em \emph{concreto}, é uma só, porque informa um só sujeito; pois uma coisa \emph{abstracta}, como é a \emph{auctoridade}, torna-se \emph{concreta} pelo sujeito, em que reside. Logo as tres funcções da auctoridade devem encontrar-se reunidas numa só Pessoa, physica ou moral, que é a pessoa do supremo reitor da sociedade. Esta \emph{unidade} é necessaria; porque, aliás, o movimento social não seria harmonico.}.

a) Funcção legislativa, --- Funoec\&o \emph{legislativa} é o poder de fazer ou promulgar \emph{leis}. --- Esta funeção compete á auctoridade \emph{civil}. Por quanto, a auctoridade tem o dever e o direito de dirigir os actos dos cidadãos para a consecução do fim ou bem commum. Ora esta direcção dá-se por meio de \emph{leis}. Logo a funec\&o \emph{legislativa} compete á auctoridade \emph{civil}. --- As leis, como vimos, devem ser \emph{honestas}, \emph{justas} e \emph{possiveis} \footnote{A funcção \emph{legislativa} é a mais nobre de todas, porque prescreve as normas dos actos, e ordena toda a vida social. --- O \emph{objecto} d'esta funcção é constituido por tudo o que é necessario e util para a consecução do fim social, que é o bem \emph{commum e externo}, \emph{coordenado com o bem interno e subordinado ao fim ultimo}. Por isso o poder legislativo não se extende aos actos meramente internos, --- nem póde prescrever coisas contrarias \& lei divina, natural ou positiva; como tambem não póde alterar a forma do governo, nem os elementos essenciaes do Estado (pois esta alteração deve ser feita pelo principe e pelo povo, isto é, pelo poder \emph{constituinte}) --- As \emph{propriededes} das leis \emph{positivas} são tres, a saber: \emph{honestidade}, \emph{justiça} e \emph{possibilidade}. --- As leis dividem-se em grupos, e constituem as diversas partes\pgn{769}do \emph{Direito}, \emph{e} assim temos o Direito \emph{privado}, o \emph{publico}, o \emph{internacional}, etc. --- Nem todas as disposições, que emanam do poder \emph{legislativo}, têm a forma solemne de \emph{leis}. Os caracteres, que distinguem a \emph{lei} de uma simples \emph{disposição}, são tres: a sua \emph{importancia}, a \emph{generalidade} e a \emph{permanencia}.}.

b) \textbf{Funeção} judicial. --- Funeção \emph{judicia}? é o poder de julgar ácerca da qualidade de um \emph{direito}, e ácerca da qualidade de \pg{769}  um \emph{crime} e da \emph{pena} que se lhe deve applicar. --- Esta funeção compete á auctoridade \emph{civil}. Porquanto, a auctoridade tem o dever de defender os direitos dos individuos e da sociedade. Ora esta defesa exige um \emph{juizo}, pelo qual a auctoridade conheca a qualidade do \emph{direito}, ou a qualidade do \emph{crime}, que violou o proprio direito, e da \emph{pena}, que deve reparar essa violação. Logo a funecüo \emph{judicial} pertence á auctoridade \emph{civil} \footnote{As funegóes \emph{judiciaes} consistem em \emph{decidir as contendas}, \emph{que nascem das pretencües encontradas de duas ou mais pessoas}, \emph{ácerca da mesma coisa}, \emph{e em restabelecer a ordem juridico-moral}, \emph{perturbada por uma} “\emph{violação criminosa}. --- Por isso o \emph{poder judicial} divide-se em \emph{civil} e \emph{criminal}. O poder \emph{civil} tem por fim tutelar os direitos dos cidadãos e defendel-os dos ataques dos outros. O poder \emph{criminal} consiste em julgar um acto nas suas relações com a lei moral, applicando-lhe, se elle for mau, o castigo conveniente, para reparar a ordem perturbada e evitar ulteriores transgressões. Por isso a pena é, ao mesmo tempo, \emph{vindicativa} e \emph{defensiva}.\par Aqui levanta-se a questão ácerca da justiça da \emph{pena de morte}. \emph{Beccaria}, \emph{Bentham} e \emph{Ahrens}, seguidos por muitos escriptores modernos, negam que a auctoridade civil tem o direito de infligir a pena capital. --- Mas estes escriptores erram. Porquanto, commettem-se crimes horrorosos (como o \emph{parricidio}, o \emph{Ahomicidio premeditado}, ete.), que não podem ser convenientemente punidos senão com a privação da vida; pois um mal summo deve ser punido com a perda de um bem summo (como é a vida); visto que deve existir uma proporção entre o crime e a pena. --- Além d'isso, o Estado tem o direito de empregar os meios mais efficazes para a defesa dos direitos seus e dos cidadãos. Ora o meio mais efficaz para essa defesa é a \emph{pena de morte}; visto que os homens depravados geralmente deixam de praticar um crime, só quando sabem que este crime se expia com a propria morte. --- A experiencia confirma a nossa these. Quando, num paiz, fica abolida a pena de morte, os crimes, e sobretudo os homicidios premeditados, augmentam; e, quando essa pena fica restabelecida, os crimes diminuem. Foi o que aconteceu, ha poucos annos, na Suissa.\par Examinemos as objecções dos adversarios da pena de morte.\par a) \emph{Beccaria} (\emph{Dei delitti e delle pene}) e \emph{Bentham} (\emph{Oeuvres}, \emph{t}. IT) fazem as seguintes objecções. --- La) \emph{A} auctoridade civil, sendo a somma dos direitos dos cidadãos, não se extende á pena de morte; porque ninguem tem direito sobre a propria vida. --- IL.a) A \emph{pena de morte} não tem efficacia para retrahir os homens do mal; pois a morte não é considerada pelos criminosos como o peior dos males. --- IIL2) A \emph{pena de morte} é uma especie de homicidio, e leva os espectadores mais ao crime do que á virtude. --- IV.2) Para a sociedade é mais util a vida do criminoso, condemnado aos trabalhos, do que a morte violenta, que a priva de um membro.\par Estas razões sao falsas. --- E' falsa a primeira; porque a auctoridade civil não é á somma dos direitos dos cidadãos, mas é uma participação da aucto-\par 49 Vor. II.\pgn{770}ridade, que \textsc{Deus} tem sobre a nossa vida e sobre a nossa morte. --- E' falsa a segunda; pois os condemnados á \emph{morte} pedem para que essa pena seja commutada com os trabalhos por toda a vida, e ficam immensamente satisfeitos, quando alcançam essa commutação, que lhes permitte esperar sempre pela liberdade. --- E' falsa a terceira; porque a \emph{pena de morte} é justa, e ineute um terror salutar nos circumstantes e um odio efficaz ao crime. --- E' falsa a quarta; porque a sociedade tira um proveito immenso da \emph{morte} do eriminoso, emquanto essa morte retrahe dos crimes os outros, e assim a morte de um criminoso é a salvação de muitos innocentes.\par \emph{b}) \emph{Ahrens} faz as seguintes objecções. --- L?) Matar o homem é acto essencialmente mau. --- IL?) O homem é \emph{fim} e não \emph{meto}. Ora, na applicação da \emph{pena de morte}, \emph{é meio}; porque serve para afastar dos crimes os outros. IIL2) O direito á \emph{vida}, sendo \emph{innato} e \emph{natural}, é absolutamente inamissivel.\par Estas razões tambem são falsas. --- E' falsa a primeira; porque, se é mal matar um homem por auctoridade \emph{particular}, não é mal matal-o por auctoridade \emph{publica}; aliás as guerras seriam todas injustas e contrarias á razão, o que o proprio \emph{AZrens} não admitte. --- E' falsa a segunda; porque o fim do homem é a \emph{ordem moral}, e a vida presente é \emph{meto} em relação a essa ordem, e por isso a \emph{pena de morte}, quando subordina a vida temporal do homem á ordem moral, não transforma o homem de \emph{fim} em \emph{meio}; aliás a\par \emph{n} vida do individuo nunca poderia pospor-se á utilidade publica. --- E' falsa a terceira; porque, se o direito á \emph{vida}, que o criminoso possue, é \emph{énnato} ou \emph{natural}, tambem o direito, que a auctoridade civil tem de prover á sua seguranca e por isso de applicar a \emph{pena de morte}, \emph{é innato} e \emph{natural}; ora, na collisão dos dois direitos, deve prevalecer o direito da auctoridade civil, que é naturalmente o mais forte.\par c) Outros escriptores fazem a seguinte objecção: O fim da pena é a emenda do reu. Ora a \emph{pena de morte} não consegue tal fim. Logo essa pena, não alcançando o seu fim, é contraria á razão e por isso é illicita.\par Resposta. A pena tem dois fins: um \emph{principal}, que é reparar a ordem violada e incutir medo aos mal intencionados, para os afastar do crime; --- outro \emph{secundario}, que é a emenda do reu (pois o bem \emph{publico} deve preferir-se ao \emph{particular}). Ora a \emph{pena de morte}, embora não consiga \emph{sempre} o fim \emph{secundario}, consegue sempre o fim \emph{principal}; e isto é sufficiente para que essa pena seja conforme com a razão e por isso licita. --- Dissemos --- embora não consiga \emph{sempre} o fim \emph{secundario}; porque, muitas vezes, os criminosos, á vista da morte que os espera e da eternidade que se approxima, detestam seus crimes; o que talvez nunca teriam feito, se lhes tivessem poupado a vida.}.

c) Funeção executiva. --- Funcção \emph{executiva} é o poder de fazer \emph{executar} as leis, regulando, em harmonia com ellas, os actos dos cidadãos. --- Esta funeção compete á auctoridade \emph{civil}. Por quanto, para a consecução do fim social, não basta dirigir as intelligencias por meio das \emph{leis} e dos \emph{juízos}; mas é preciso dirigir as vontades, varias e geralmente rebeldes, dos cidadãos e as suas forças physicas por meio da coacção juridica. Ora esta direcção das vontades e das forças dos cidadãos \pg{771}  exige a funeção \emph{executiva}. Logo a funeção \emph{executiva} compete á auctoridade \emph{civil}. --- Os principaes dotes da funeção \emph{executiva} são: \emph{legalidade}, \emph{força} e \emph{prudencia} \footnote{Ás funceções \emph{legislativas} e \emph{judiciaes} succedem as funcções \emph{executivas}. Não basta fazer leis, dirimir contendas e julgar crimes; é preciso applicar em concreto o que foi deliberado pelo poder \emph{legislativo} e definido pelo \emph{judicial}. A somma das funcções \emph{executivas} constitue a \emph{administração publica}, e o conjuncto das leis, que as regulam, constitue o \emph{Direito administrativo}. --- A \emph{materia} do \emph{Direito administrativo} costuma ser dividida pelos \emph{modernos} (Persico, \emph{Principii di diritto amministrativo}) em quatro partes, que são --- a \emph{jerarchia} ou \emph{organisação administrativa}, --- \& \emph{propriedade publica} ou \emph{fazenda}, --- a \emph{tutela}, isto é, o exercito, a policia e o contencioso, --- a \emph{administração social} ou concurso do Estado na civilisação economica, intellectual e moral. Porquanto, a auctoridade, para o exercicio das suas funcções \emph{executivas}, --- precisa de um grupo de pessoas, \emph{capazes} e \emph{conscienciosas}, que a representem em todos os pontos do paiz (d'ahi \& necessidade da \emph{jerarchéa} ou \emph{organisação administrativa}); --- precisa de recursos e rendimentos para as despesas e necessidades do Estado (d'ahi a necessidade e o direito da \emph{propriedade} ou \emph{fazenda publica}; --- precisa de força para coagir juridicamente os insubordinados (d'ahi a necessidade do \emph{exercito}, da \emph{policia} e do \emph{contencioso}); --- precisa do concurso do Estado para promover a prosperidade temporal da sociedade (d'ahi a necessidade da \emph{administração social}).}.

%%>>> art-70

\section*{Artigo II.\\ Deveres e direitos dos soberanos.}
\addcontentsline{toc}{section}{Artigo II. --- Deveres e direitos dos soberanos.}

\pg{766} \pg{768} \pg{770} 181. \textbf{Origem} dos \textbf{deveres} e dos direitos dos \textbf{soberanos}. --- Os \emph{deveres} e os \emph{direitos}, que são proprios dos. \emph{soberanos}, têm a sua \emph{origem} no proprio \emph{fim} da sociedade civil. Por isso os soberanos \emph{devem} e \emph{podem} fazer o que, de algum modo, concorre para o bem \emph{commum} e \emph{externo}, comtanto que este bem esteja \emph{coordenado com o bem interno}, que é a \emph{virtude}, e \emph{subordinado}

\emph{A ao fim ultimo}, que é a \emph{felicidade eterna} \footnote{A sociedade civil resulta do \emph{soberano} e dos \emph{subditos}; porque estes constituem a \emph{multiddo}, e aquelle representa a \emph{auctoridade}, e já vimos que a \emph{multidão} e á \emph{auctoridade} são os elementos constituitivos da sociedade civil. Por isso os \emph{deveres} e os \emph{direitos sociaes} reduzem-se aos \emph{deveres} e aos \emph{direitos} dos \emph{soberanos} e dos \emph{subditos}. Estes \emph{deveres} e \emph{direitos} têm a sua \emph{origem} no \emph{fim} da propria sociedade; pois, como dissemos, o \emph{fim} é a rasãoe \& \emph{regra} dos \emph{meios}, e os \emph{deveres} e os \emph{direitos} são \emph{meios} por que a sociedade consegue a felicidade temporal dos cidadãos. --- E' esta a razão, por que o \emph{Direito natural}, para determinar os \emph{deveres} e os \emph{direitos sociaes}, deve primeiramente estudar a natureza ou \emph{constituição}, e, sobretudo, o \emph{fim} da sociedade civil.}. :

182. \textbf{Especies} de \textbf{deveres} dos \textbf{soberanos}. --- Os soberanos tém \emph{deveres} para com \textsc{Deus}, para com a sociedade em \emph{geral}, para \pg{772}  com as diversas \emph{classes sociaes}, para com as \emph{familias}, para com os \emph{individuos}, e, em harmonia com estes deveres, são obrigados a exercer as diversas funcções da sua auctoridade.

a) \textbf{Deveres} para com \textbf{Deus}. --- Os soberanos, além dos deveres \emph{religiosos} a que estão sujeitos como pessoas \emph{particulares}, têm a obrigação de fazer com que a sociedade preste a \textsc{Deus} um culto \emph{publico} e \emph{legitimo}. (pois \textsc{Deus} é o Auctor, o Senhor, o Governador e o Fim ultimo da sociedade), e de prohibir todos os actos \emph{externos}, que envolvam o desprezo da Magestade Divina \footnote{\textsc{Deus} é o \emph{Auctor} e o \emph{Senhor} da sociedade civil. Por isso a sociedade, como tal, deve reconhecer esta dependencia, e por conseguinte prestar a \textsc{Deus} um culto \emph{publico} e \emph{social}. --- Além d'isso, a sociedade, como tal, recebe continuamente de \textsc{DEus} muitos e grandes beneficios; como são a prosperidade temporal, a victoria contra os inimigos, etc.; e precisa do auxilio divino para o bom exito dos seus actos. Logo a sociedade, como tal, deve ser reconhecida 2 \textsc{DEus}, offerecer-lhe os seus votos, isto é, deve, publica e socialmente, mosirar a \textsc{DEus} a propria gratidão, e as proprias necessidades. Os Romanos, embora pagãos, faziam supplicas aos deuses para alcançar victorias, e, depois de as terem alcançado, davam graças e offereciam sacrificios. E não só em Roma, mas em todos os outros povos vigorou sempre o culto \emph{publico} e \emph{social} da Divindade, e sempre existiram sacerdotes, sacrificios, templos, etc. --- D'onde se vê o erro dos que, apesar de reconhecerem no \emph{individuo} o dever de prestar a \textsc{Deus} o culto \emph{publico}, não o reconhecem comtudo na \emph{sociedade}; porque julgam que a sociedade não é uma instituição \emph{natural} e proveniente de \textsc{DEus}, mas uma instituição \emph{arbitraria} dos homens. E' o \emph{indifferentismo politico em materia de religião}.\par Tendo a sociedade o dever de prestar a \textsc{Deus} um culto \emph{publico}, é claro que devem ser prohibidos todos os actos \emph{externos}, que, \emph{directa} ou \emph{indirectamente}, offendem a Magestade Divina, como são a negação de qualquer attributo de \textsc{Deus}, ou do seu dominio sobre os homens, as publicas violações dos deveres fundamentaes, que a lei natural impõe, com a força da evidencia, a todos os homens. Por isso a \emph{liberdade de imprensa}, sobretudo em materia de \emph{Religião}, é uma licença impis, que deve ser enfreada. \emph{PAilostrato} narra que, em Athenas, foram queimados publicamente uns livros, porque ensinavam o \emph{atheismo}. \emph{Tito Livio} diz que o mesmo se fez em Roma.\par O culto, que a sociedade deve prestar a \textsc{DEus}, não só deve ser \emph{publico}, mas tambem \emph{legitimo}, em harmonia com a vontade de \textsc{Deus}; pois ninguem acceita um servico ou um culto, que não seja conforme com a propria vontade. Por isso, se \textsc{Deus} revelar, por via sobrenatural, \emph{o} modo por que deseja ser servido e venerado pela sociedade, esta tem o dever de se conformar com a vontade de \textsc{Deus}. E, como \textsc{Deus} se dignou revelar o modo por que deseja ser venerado, e como a depositaria d'esta revelação é a Egreja Catholica, segue-se que a sociedade civil deve prestar a \textsc{Deus} o culto publico sob a direcção da mesma Egreja Catholica. A\par D'onde se vê que a \emph{liberdade de culto}, que concede ao Estado o poder de escolher a Religião que elle quizer, é \emph{impia}, \emph{é injusta}, é \emph{perniciosa}. ---\pgn{773}E' \emph{impia}, porque envolve uma desobedieneia e insubordinação contra \textsc{DEus}. --- E' \emph{injusta}; porque não reconhece o \emph{direito} que a Egreja Catholica tem \textasciicircum{}de ser reconhecida como a unica e verdadeira Religião. --- E' \emph{perniciosa} não só ás almas, mas tambem á sociedade; pois a diversidade de culto ou de Religião trouxe sempre discordias e guerras no seio da sociedade. --- 'T'odavia concedemos que a \emph{liberdade de culto}, convenientemente reduzida ou limitada, possa, algumas vezes e nos paizes onde uma grande parte professa uma religião falsa, \emph{tolerar-se} conforme as condições dos tempos e dos paizes. Esta \emph{tolerancia} póde admittir-se nestes casos: --- 1.º) quando produz um bem maior; --- 2.9) quando póde evitar um mal mais funesto; --- 3.º) quando é \emph{physicamente} impossivel a prohibição de um culto falso. Nos paizes, porém, onde domina a verdadeira Religião, o Estado, se não póde impor por força o exercicio d'ella aos que a não professam, deve prohibir o exercicio publico e a propaganda das religiões falsas.}.

\pg{773} b) \textbf{Deveres para} com a \textbf{sociedade em geral}..--- Os soberanos, tendo por fim procurar o bem \emph{externo} e \emph{commum}, resultado de dois factores indispensaveis, que são a prosperidade \emph{materia}? e \emph{economica} e á prosperidade \emph{religiosa}, \emph{moral} e \emph{intellectual}, Vêm o \emph{dever} de fazer, ordenar e exigir tudo o que póde concorrer para esse fim, comtanto que não offendam os direitos 2\emph{nnatos} dos individuos, das familias ou das outras sociedades \footnote{O bem \emph{externo} e \emph{commum} da sociedade abrange a prosperidade \emph{material} ou \emph{economica}, e a prosperidade \emph{religiosa}, \emph{moral} e \emph{intellectual}.\par a) A prosperidade \emph{material} ou \emph{economica}, á qual consiste na abundancia de todos os bens materiaes, capazes de satisfazer a todas as necessidades dos homens, deve ser o objecto dos esforcos do soberano, ou da auctoridade civil. Estes bens são: --- 1.º) os varios productos da actividade e do trabalho dos cidadãos; --- 2.º) os differentes meios de satisfazer ás necessidades da vida, como os hospieios, os hospitaes, as vias de communicação, etc., --- 9.9) os meios e as instituições publicas, cujo fim é a ordem, a paz e a segurança; --- 4.9) tudo o que concorre para o bem estar publico e não póde ser conseguido pelos esforços dos individuos. Ora o soberano tem o dever de procurar esses bens ou favorecer os individuos que se esforçam para os conseguir, afastando os obstaculos que se oppõem ao desenvolvimento da producção e do \emph{commercio}, e protegendo e defendendo a actividade productora contra as causas, que \& possam prejudicar, e estimulando a vontade dos socios.\par O soberano, para poder cumprir com estes seus deveres, deve --- não só repartir os impostos de modo que recaiam proporcionalmente sobre todos os ramos da producção, e não sejam tão excessivos, que absorvam e oxgottem as fontes de producção, --- mas tambem deve conceder plena liberdade commercial, com tanto que a concorrencia estrangeira não prejudique a producção nacional. Esta intervenção do soberano não deve ser \emph{directa}, de modo que o Estado se torne proprietario e industrial com damno dos individuos e das familias; pois o Estado só póde ser \emph{productor} naquelles ramos de actividade, que, por serem de interesse geral, não podem ser desenvolvidos por particulares, e que, sem uma intervenção \emph{directa} de Estado, seriam condemnados a desapparecer. --- Além d'isso, o soberano deve escolher magistrados e empregados publicos dignos e competentes, observando nessa escolha a\pgn{774}justiça \emph{déstributiva}; pois que a escolha de pessoas indignas ou incapazes prejudica os interesses da sociedade.\par b) A prosperidade \emph{moral} e \emph{religiosa}, ou a \emph{moralidade} e a \emph{religiosidade}, são necessarias para a sociedade civil conseguir o seu fim temporal. --- E' necessaria a \emph{moralidade}; porque a experiencia attesta que, onde não existe o sentimento de honestidade, não existe a justiça, nem a paz, nem por isso o bem temporal. --- E' necessaria a \emph{religiosidade}; porque não ha moralidade sem religiosidade. Por isso \emph{Platão}, \emph{Aristoteles}, entre os antigos, \emph{Washington} e \emph{Rousseau}, entre os modernos, confessam que a Religião é o fundamento do Estado, e que o homem irreligioso nunca será bom cidadão nem bom subdito. \emph{Donoso Cortés} dizia que numa sociedade verdadeiramente catholica são impossiveis duas coisas: o despotismo e a revolução. \emph{Taine}, embora positivista, diz: « Quando se tem assistido de perto a este espectaculo, pôde avaliar-se a influencia do Christianismo nas nossas modernas sociedades, pois é elle que traz o pudor, a doçura, a humanidade, e que, entre nós, conserva a honradez, a boa fê e a justiça. Nem a razão philosophica, nem a cultura artistica e litteraria, nem a honra feudal, militar e cavalheiresca, nem os codigos, nem os governos, nem as administrações podem substituil-o nesta obra benefica » (\emph{Orig}. \emph{de la France contemporaine}).\par Sendo a \emph{moralidade} e a \emph{religiosidade} tao estreitamente ligadas com o bem temporal e publico dos povos, é claro que o soberano não só deve (deveres \emph{negativos}) impedir, prohibir e castigar todos os actos, que são dirigidos contra a Morale a Religião e constituem um escandalo publico, --- mas tambem deve (deveres \emph{positivos}) proteger e auxiliar tudo o que se refere á moralidade e \& religiosidade \emph{publica}. Este dever do soberano é exigido pelo proprio fim \emph{directo} da sociedade; pois, como vimos, a Moral e a Religião são meios indispensaveis para a felicidade temporal dos homens.\par Dizemos que a moralidade e a religiosidade \emph{publica} formam o fim \emph{directo} da auctoridade do soberano; porque o fim \emph{directo} da sociedade é o bem \emph{externo} e \emph{commun} dos homens. O bem \emph{interno} e \emph{particular}, que fórma o fim \emph{directo} da sociedade \emph{religiosa}, só forma o fim \emph{indirecto} da auctoridade \emph{civil}, emquanto esta, devendo promover e procurar \emph{directamente} a religiosidade e a moralidade \emph{publica}, promove tambem, ainda que \emph{indirectamente}, a religiosidade e a moralidade \emph{particular} de cada membro da sociedade.\par Concluimos com as palavras insuspeitas de \emph{MaccAiavelli}: « Os principes e as republieas, que desejam manter-se incorruptas, devem sobretudo manter incorruptas as ceremonias da religião e prestar-]hes a devida veneração. Porquanto o desprezo do culto divino é o maior indicio da ruina d'uma provincia » (\emph{La mente di un womo di stato}, 1. I, c. 18). ,\par Note-se quo o soberano, ou a auctoridade civil, tem a obrigação de defender e promover a Religião na sociedade --- não só para satisfazer aos deveres, que o proprio Estado tem para com \textsc{Deus}, --- mas tambem para cumprir os deveres que o mesmo soberano tem para com a sociedade, pois o soberano deve promover \emph{directamente} o bem \emph{communm} e \emph{externo} dos cidadãos e \emph{indivectamente} o bem \emph{interno}, e a religiosidade é um meio para a consecução do bem \emph{commun}: e \emph{externo}, e é um elemento do bem \emph{interno}.\par 0) A \emph{instrucção}, em sentido mais rigoroso, é o conjuncto de conhecimentos\pgn{775} religiosos, litterarios, scientificos e technicos, adquiridos por meio dos livros, dos professores, ou pela observação racional e directa dos phenomenos da vida physica ou moral. --- A. \emph{insiruccüo} refere-se á \emph{intelligencia}, como a \emph{educação á vontade} (pois á \emph{educação} tem por fim acostumar a vontade a vencer todos os instinctos e impulsos desordenados, que afastam o homem do cumprimento dos seus deveres). A \emph{educação} suppõe á \emph{instrucção}, pelo menos em tudo o que se refere á lei moral e religiosa, porque a vontade não póde querer uma coisa, que a intelligencia não conhece.\par Para avaliarmos a importancia da \emph{instrucção}, devemos distinguir os conhecimentos \emph{moraes} e \emph{religiosos} dos conhecimentos \emph{litterartos}, \emph{scientificos} e \emph{tecnhicos}. --- Os conhecimentos \emph{religiosos} e \emph{moraes} são, como vimos, absolutamente necessarios para que a sociedade exista e prospere, e para que os individuos alcancem o fim da vida presente, que é o \emph{bem interno}, \& \emph{virtude}. --- Os conhecimentos \emph{literarios}, \emph{scientificos} e \emph{technicos}, quando são bem regulados e baseados na Religião, contribuem para o bem social e individual, embora não sejam necessarios para o homem alcançar o seu fim, nem para a sociedade subsistir. Estes conhecimentos constituem a \emph{instrucção superior}.\par Em relação á \emph{instrucção superior}, o soberano deve fundar universidades, bibliothecas, museus, observatorios, etc. --- cuidar para que as doutrinas sejam conformes com a Religião e com a Moral, --- conformidade, que, nos paizes onde a Religião catholica é official, deve ser julgada pelos Prelados, etc. --- Todavia o Estado não deve attribuir-se o monopolio da \emph{instrucção superior}, o que é contrario ao direito dos paes.\par Além da \emph{énstruccüo superior}, ha uma outra \emph{instrucção}, chamada \emph{primaria} ou \emph{elementar}, que é formada pelas noções da leitura, escriptura e contabilidade, e que é apenas um \emph{meio} para adquirir os conhecimentos, que formam a \emph{instrucco superior}. Neste seculo, dá-se a instruccüo \emph{primaria} ou \emph{elementar} uma importancia, que ella não merece, chegando-se a dizer que ella, só por si e desacompanhada dos conhecimentos religiosos e scientificos, póde causar o bem estar individual e social, e assim tornar-se a panacea de todos os males sociaes. E' um grande exaggero e um grande principio de desordens e agitações, como dizem auctores insuspeitos --- \emph{Spence} e \emph{Prins}. O homem ignorante resigna-se com a sua sorte; o homem instruido revolta-se, quer uma posição, e, vendo frustrados os seus desejos, deixa-se arrastar para os mais horrorosos crimes. \emph{Garofulo}, fundando-se na estatistica criminal, demonstra que a instrucção não influe na diminuição da criminalidade, e faz notar que, com o augmento do salario e com a diffusão da instrucção na ultima metade d'este seculo, coincide um extraordinario augmento de crimes, revestidos das mais horrorosas circumstancias.\par Espostas estas noções, vejamos brevemente os deveres do soberano, ou da auctoridade civil, em relação á \emph{instrucção primaria}. O soberano deve fazer tudo o que não póde ser feito por iniciativa particular,e por isso deve fundar escholas, velar para que a instruecao litteraria seja acompanhada da instrucção moral e religiosa, prevenir e punir os abusos dos professores, etc. --- O soberano deve, porém, lembrar-se de que os paes têm o direito ?\emph{nnato} de dirigir á educação e a insiruecüo de seus filhos e de que a instrucção é tambem, e sobretudo, um bem particular e proprio do individuo. Por isso o\pgn{776}\emph{ensino obrigatorio} e o \emph{monopolio do ensino}, attribuidos ao Estado, são despotismos e tyrannias; pois que atacam os direitos dos paes e a liberdade individual. --- Aos partidarios do \emph{ensino obrigatorio}, que dizem que o Estado deve preceituar a instrucção, por esta tambem concorrer para o bem social, respondemos que o fim do Estado não póde ser um bem \emph{particular}, que o individuo possa alcançar, mas só deve ser um bem \emph{commum}, que só elle (o Estado) pôde conseguir; ora a instrucção alcança-se sem o Estado.}.

c) \textbf{Deveres} para com as classes sociaes, --- Os soberanos devem favorecer as diversas \emph{c/asses sociaes}, sobretudo a \emph{agricola} \pg{774} \pg{775}  e a \emph{industrial}, necessarias para o progresso e prosperidade social, tutelando os seus direitos, promovendo o seu desenvolvimento \pg{776} , afastando os perigos que as ameaçam, e conservando entre ellas o espirito da paz e da justiça \footnote{\emph{Classes sociaes} são os diversos conjunctos de homens, que exercem a mesma profissão e industria, ou que, por occuparem egual posição social, têm eguaes interesses. --- A existencia das \emph{diversas classes} é uma necessidade social. As causas d'esta necessidade são duas: --- a \emph{divisão do trabalho}, a qual deriva da limitação das forças do homem e da diversidade das suas aptidões e produz as varias classes, a \emph{agricola}, \emph{á industrial}, etc., --- e a necessidade de direcção e protecção nos assumptos sociaes e politicos, --- direcção e protecção que, devendo ser ou sendo exercida por pessoas de maior intelligencia, prestigio ou audacia, dá origem ás classes dirigentes e dirigidas. --- A existencia das classer sociaecs, se é um bem para o Estado, quando ellas se contém nos limites marcados pelo proprio fim, é um mal e uma causa de desordens, quando saiam da esphera da sua acção, pretendendo privilegios e regalias. --- A \emph{missão} ou o \emph{dever} da auctoridade em relação ás \emph{classes sociaes} consiste, como dissemos, --- em defender o proteger os direitos, que. ellas têm em razão do seu fim, --- promover o seu desenvolvimento por meio de leis e providencias accommodadas á indole de cada uma das classes, --- corrigir e evitar os perigos, que as podem deteriorar material e moralmente, obstando sobre tudo a que offendam os direitos das outras classes e da propria sociedade civil.\par As \emph{classes sociaes} são muitas e diversas, e todas muito importantes. 'Todavia, limitamos o nosso estudo ás \emph{prodwctoras}; porque a ellas se refere a \emph{questão social}, que agita tão fortemente as nações.\par As principaes classes \emph{productoras} são a \emph{agricola} e a \emph{industrial}.\par \&) A classe \emph{agricola} é muito importante, quer se considere sob o aspecto \emph{economico}, quer sob o aspecto \emph{moral}, quer sob o aspecto \emph{social}. E' muito importante sob o aspecto \emph{economico}; pois, como diz \emph{Cawwés}, á primeira questão economica é a da \emph{subsistencia}, e por isso a producção agricola domina todas as outras questões sociaes. E' muito importante sob o aspecto \emph{moral}; pois que, em geral, é mais religiosa, morigerada e simples que as outras. E' muito importante sob o aspecto \emph{social}; porque, predominando nella o espirito de tradição, é um importante elemento de conservação e de ordem, e, sendo composta de familias e individuos sãos e vigorosos, forma uma fonte viva de forças para as outras classes. --- Para que a classe \emph{agricola} consiga os seus fins, é necessario que o Estado --- promulgue leis opportunas para salvar os seus direitos e interesses legitimos, --- facilite a sua existencia e desenvolvimento, --- defenda a \emph{pequena propriedade}, para que não seja absorvida pela \emph{grande propriedade}, --- promova ou favoreca as \emph{associações} e o \emph{credito agricola}, para evitar as \emph{Aypothecas} e as «\emph{suras} e proporcionar aos agricultores os meios de vida e os instrumentos da lavoura.\par b) A classe \emph{industrial}, no sentido mais rigoroso, é a que tem por fim\pgn{777}transformar e modificar os productos ou materias primas, offerecidas pela terra, pará satisfazer ás necessidades do homem. --- Esta classe é necessaria, não só porque o homem não póde subsistir e progredir sem esta transformação e modificação dos productos naturaes, senão tambem porque o trabalho deve ser dividido, e uma especie d'esta divisão constitue a classe \emph{industrial}.\par A \emph{industria} divide-se principalmente em \emph{pequena} e \emph{grande}. --- A \emph{pequena industria} exige um capital mais modesto; o proprietario é, ao mesmo tempo, trabalhador e emprega um pequeno numero de operarios, que podem, com o tempo, tornar-se modestos industriaes. --- A \emph{grande industria}, pelo contrario, requer grandes capitaes, um grande exercito de operarios, que difficilmente chegam a ter uma industria propria, e exige a separação entre o trabalho material e a sua direcção, o emprego de grandes machinas, etc.\par A \emph{pequena industria} é muito vantajosa sob o aspecto \emph{moral} e \emph{social}. E' vantajosa sob o aspecto \emph{moral}; emquanto, não afastando o operario da sua familia, subtrahe-o a tantos males moraes. E' vantajosa sob o aspecto \emph{social}; porque, não havendo entre o operario e o proprietario uma grande differença de costumes e de posição, não ha entre elles antagonismos e odios.\par A \emph{grande industria}, porém, traz geralmente comsigo muitos inconvenientes moraes e \emph{sociaes}. --- Traz inconvenientes \emph{moraes}; porque, reunindo muitos operarios de ambos os sexos na mesma fabrica, dà logar \& muitas desordens e relaxa o vinculo entre as pessoas da mesma familia. --- Traz inconvenientes \emph{sociaes}; porque a differenca de posição entre o proprietario e os operarios, a difficuldade que o operario encontra para subir na jerarchia industrial, as crises da producção provenientes da concorrencia demasiada, e, como resultado d'estas crises, a diminuição do salario, a paralysação do trabalho, --- todo este conjuncto de causas deram origem ao \emph{pauperismo}, que afllige tantas familias e á \emph{questão social}, que ameaça a sociedade moderna.\par Para evitar os inconvenientes, a que dão logar as \emph{grandes industrias}, fundaram-se na Edade media gremios e associações, cujo regimen tinha um triplice fim: \emph{economico-technico}, \emph{benefico} e \emph{religioso}. Este regimen, considerado sob o ponto de vista social, « assegurava á classe operaria, em todas as suas categorias, o bem estar e o desafogo, e por conseguinte a educação, \& situação social. Por outra parte, o systema corporativo tirava ao individuo a possibilidade de se elevar muito acima dos outros. A liberdade absoluta, se cria fortunas colossaes, leva muitas vezes a especular com as forças do trabalho e por isso a opprimir centenares e milhares de individuos » (Janssen, \emph{ZL/Allemagne d la fin du moyen-dgé}, pag. 3892).\par Hoje não existe este regimen corporativo, e por isso acabaram as intimas relações entre patrões e operarios; não existem laços entre os officiaes, nem entre os companheiros do mesmo officio; o operario não tem esperança de encontrar soccorro no momento da necessidade e da velhice, nem apoio num caso de oppressão. --- Deram, é verdade, ao operario o direito politico do suffragio, mas este mesmo direito devia servir-Ihe de occasião para se tornar escravo de especuladores, que, para subirem ao poder, não duvidam em excitar nas massas populares as mais perniciosas paixões e lançar no seio das familias a semente das mais funestas discordias.\par Um dos principaes meios; pois, para resolver a \emph{questão social} e evitar uma\pgn{778}catastrophe, é a reconstituição dos gremios e das associações operarias. As bases d'estes gremios devem ser as seguintes: --- 1.º) o elemento religioso e catholico, porque toda a associação, sem este elemento, é perigosa; --- 2.9) a união entre patrões e operarios, e portanto a protecção que as classes superiores devem dispensar ás inferiores na mesma associação; --- 9.9) a representação dos operarios na administração dos interesses communs, profissionaes e economicos; --- 4.º) a \textbf{capacidade profissional}, \textbf{promovida}, pelo ensino profissional, que restitue ao operario a sua dignidade e o seu patrimonio; --- 5.9) a jerarehia profissional, que honre o officio e diminua a injusta concorrencia, feita por incapazes a operarios distinctos e honestos.\par O Estado, pois, deve favorecer e proteger o restabelecimento e a conservação do regimen corporativo, promulgando as leis necessarias, para evitar a collisao dos direitos entre os varios gremios e harmonisar o interesse particular dos mesmos gremios com o interesse geral da sociedade.\par Mas não acabam aqui os deveres do Estado, ou da auctoridade civil, para com as classes \emph{éndustriaes}. E' necessario regular o \emph{trabalho}, eliminando todos os abusos, de que a corrupção e o interesse são \emph{origem} continua e perniciosa. Sujeitar a um trabalho excessivo as crianças e as mulheres, é destruir-Ihes a saude, \textbf{deprayar-l]hes} a alma, preparar gerações atrophiadas na alma e no corpo, relaxar ou dissolver os sagrados laços da familia. Sujeitar o homem adulto a um trabalho demasiadamente prolongado, mesmo nos dias santos, é desprezar por completo a dignidade humana. --- O Estado deve regular o \emph{trabalho} de modo, que fiquem garantidos os direitos do individuo e da familia e o bem geral da sociedade. Porquanto a missao do Estado consiste em defender os direitos, de que são dotados os individuos e as familias, \emph{e} em procurar a tranquillidade e o bem estar social, e, para cumprir com esta missão, deve servir-se da auctoridade, que \textsc{Deus} lhe confiou, regulando a actividade dos subditos, e eliminando todos os abusos. --- O Estado, pois, --- deve proteger as \emph{crianças}, fixando a edade em que podem começar a trabalhar, determinando as horas, e prohibindo o que é contrario ao desenvolvimento physico e moral das mesmas crianças; --- deve proteger as \emph{mulheres}, solteiras e casadas, afastando tudo o que as possa prejudiear na saude e na moralidade, e dando-lhes o tempo necessario para cumprirem com os seus deveres domesticos, prohibindo o trabalho excessivo e nocturno; --- deve proteger os \emph{adultos}, determinando a duração do trabalho que seja compativel com as forças e saude dos operarios, e, sobretudo, prohibindo o trabalho nos dias santos, --- prohibição utilissima não só sob o aspecto moral e religioso, mas turabem sob o aspecto economico e physiologico. --- O Estado, nestas medidas, deve respeitar, com o maior escrupulo, os direitos dos individuos e das familias, não intervindo se não quando os individuos e as familias não poderem fazer respeitar os direitos seus ou das pessoas que lhes estão confiadas.}.

d) \textbf{Deveres para} com as \textbf{familias}, --- Os soberanos devem proteger, defender e garantir os direitos das familias. E, como \pg{777} \pg{778}  a familia se compõe de tres sociedades, que são a \emph{conjugal}, a \emph{parental} e a \emph{heril}, \emph{o} dever dos soberanos é proteger, defender e garantir os direitos de cada uma d'estas sociedades, para que possam conseguir os seus fins respectivos \footnote{Os soberanos, promovendo \emph{directamente} a felicidade \emph{publica} da sociedade civil, promovem \emph{indirectamente} a felicidade \emph{particular} da familia, emquanto esta participa do bem da mesma sociedade civil. --- Para podermos determinar os deveres dos soberanos para com as familias, deyemos lembrar-nos\pgn{779}de que a familia é anterior ao Estado, e possue direitos, que derivam immediatamente da lei natural. Os soberanos, pois, --- 1.9) devem proteger a sociedade domestica, os seus direitos, e prohibir tudo o que a possa perturbar; --- 2.º) devem tomar as devidas disposições, relativas á parte economica da familia, ou aos seus bens; --- 3.º) devem proteger e defender os direitos proprios da familia contra os abusos de algum dos seus membros, ou do seu chefe. --- Todavia os soberanos não podem legislar ácerca da \emph{essencia} da familia, nem ordenar o que é contrario aos fins, propriedades e direitos da sociedade \emph{conjugal}, nem fazer o que se oppõe aos fins e aos direitos da sociedade \emph{parental}.}.

\pg{779} e) \textbf{Deveres para} com os individuos. --- Os soberanos têm o dever de tutelar e garantir os \emph{direitos} do individuo, contra todos os ataques, \emph{directos} e \emph{indirectos}, e de procurar, embora \emph{indirectamente}, o bem estar do proprio \emph{end?viduo}, etc. \footnote{Os deveres dos soberanos, ou da auctoridade civil, em relação ao \emph{individuo}, são os seguintes: --- 1.º) defender e proteger os \emph{direitos}, de que a lei \emph{natural} dotou todo o ente racional, não só contra os ataques \emph{directos} mas tambem contra os \emph{indirectos}; --- 2.º) reprimir as manifestações abusivas da liberdade humana, que não podem constituir um direito; --- 3.º) probibir e castigar todos os vicios e todas as manifestações irreligiosas, que prejudicam o bem estar da sociedade, como são os abusos da imprensa, a usura, etc.; --- 4.9) limitar os direitos do individuo, quando entram em collisão com os direitos mais elevados dos outros individuos ou da sociedade, e produzem perturbação; --- 5.º) exigir dos cidadãos tudo o que for necessario para a conservação e progresso da sociedade. --- Todavia o Estado --- não póde exigir nada que não se relacione com o bem \emph{commum}, --- nem desprezar certos direitos do individuo, que são essencialmente superiores aos do Estado, --- nem procurar \emph{directamente} o bem \emph{particular} do individuo, --- nem absorver e anniquilar a iniciativa e a actividade individual nas emprezas de caracter particular e geral. Por isso os soberanos --- 1.º) devem permittir certos males, quando os prejuizos, que o remedio póde causar á sociedade, são superiores aos que os proprios males produzem; --- 2.º) não podem impor contribuições sem verdadeiros motivos; --- 9.9) devem repartir os impostos segundo a justiça \emph{distributiva}, etc.}.

183. \textbf{Especies de direitos dos soberanos}. --- Os soberanos possuem o \emph{direito} de exercer as tres funegóes da auoctoridade, que são a \emph{legislativa}, a \emph{judicial}, e á \emph{executiva}, e de cumprir todos os \emph{deveres}, a que estão sujeitos. --- Além d'isso, têm direito á \emph{veneração}, \emph{obediencia}: e \emph{fidelidade} dos subditos \footnote{Os soberanos têm direito --- á \emph{veneração} dos subditos, porque a dignidade dos soberanos é uma participação da infinita dignidade de \textsc{Deus}; --- á \emph{obediencia}, porque, aliás, á auctoridade não poderia cumprir a sua missão e deixaria de ser o que é; --- á \emph{fidelidade}, porque os soberanos possuem um titulo justo para exercer a suprema auctoridade. O direito, que os soberanos tem á \emph{obediencia} dos subditos, é \emph{jurídico}. --- Os soberanos, que reconhecem no povo, e não em \textsc{DEUs}, a \emph{origem} e a \emph{causa} da sua auctoridade, tiram todo o prestigio á sua alta dignidade, e dão aos subditos o ensejo de menoscabar e lesar o \emph{direito} que têm \& \emph{veneração}, \emph{obediencia} e \emph{fidelidade}.}.

%%>>> art-71

\section*{Artigo III.\\ Deveres e direitos dos subditos.}
\addcontentsline{toc}{section}{Artigo III. --- Deveres e direitos dos subditos.}

\pg{780} 184. \textbf{Origem} dos \textbf{deveres} e dos \textbf{direitos} dos subditos. --- Os \emph{deveres} e os \emph{direitos}, que são proprios dos \emph{subditos}, têm \emph{origem} no \textbf{mesmo} \emph{fim} da sociedade civil. Por isso os subditos \emph{devem} e \emph{podem} fazer o que, de algum modo, concorre para a felicidade \emph{temporal} da sociedade \footnote{O homem, emquanto membro da sociedade civil, refere-se a um triplice termo: a \emph{cada um dos cidadãos}, \& \emph{sociedade inteira}, \& \emph{pessoa do soberano}. --- Em relação a \emph{cada um dos cidadãos}, o homem deve amar, com especial affecto, os seus compairiotas, --- respeitar os seus direitos, provenientes da lei, natural ou positiva, --- Em relação á \emph{sociedade inteira}, o homem devo amar a patria mais do que os outros paizes do mundo, mas com amor razoavel desinteressado, generoso, procurando ser-lhe util com a palavra, com o exemplo, com o trabalho e, sobretudo, com o sacrificio. --- Em relação á \emph{pessoa do soberano}, o homem deve-lhe veneração, \emph{obediencia} e \emph{fidelidade}. --- Aqui só tratamos, com algum desenvolvimento, dos deveres dos subditos para com os soberanos; porque o subdito refere-se principal e directamente \& auctoridade, e porque, noutros logares, já tivemos occasiüo de fallar nos deveres do homem para com os semelhantes e para com a patria.\par 4) Nada mais grandioso e efficaz do que considerar a magestade humana como um reflexo e uma participação da Magestade Divina! Só \textsc{Deus} é o Soberano essencial e absoluto do homem, e por isso só a \textsc{DEus} pertence a auctoridade sobre o homem. Mas, o homem precisa de uma direcção sensivel. ti, pois, necessario que a auctoridade divina se manifeste sensivelmente ao homem, enearnando-se, por assim dizer, na pessoa do soberano. D'este modo, o soberano apresenta-se em nome e por ordem de \textsc{Deus}, e por isso,0s homens devem venerar na pessoa do soberano terreno a propria Magestade de \textsc{Deus}. Considerado sob este aspecto, o \emph{soberano}, embora como pessoa particular tenha defeitos e vicios, comtudo, como representante de Deus, merece sempre o respeito e á veneração do povo.}.

185. \textbf{Especies} de \textbf{deveres} dos subditos. --- Os \emph{deveres} dos \emph{subditos} para com os soberanos reduzem-se a tres, que são a \emph{veneração}, a \emph{obediencia} e a \emph{fidelidade}.

a) \textbf{Dever} de veneração. --- O soberano, sendo revestido de auctoridade, que é uma participação da auctoridade divina, representa a propria Divindade, e por isso a sua pessoa é inviolavel e merece a \emph{veneração} dos subditos \footnote{Derivando de \textsc{Deus} a auctoridade, de que os soberanos estão revesti-}.

b) \textbf{Dever} de obediencia. --- O soberano, tendo recebido de \textsc{Deus} a auctoridade suprema, tem o \emph{direito} de mandar nos subditos. Por isso, correspondendo a cada \emph{direito} um \emph{dever}, os subditos têm o \emph{dever de obedecer} ás ordens do soberano \footnote{\pgn{781}dos, é claro que toda a lei justa, promulgada por elles, por tirar a sua força da propria, auctoridade, deriva de algum modo de \textsc{Deus}. Por isso quem desobsdece a essas leis, desobedece ao proprio \textsc{Deus}.\par Pergunta-se: Estarüo os subditos obrigados a obedecer ao soberano, que abusa do poder e exerce sobre elles um dominio tyrannico?\par Resposta. O tyranno póde ser de \emph{governo} ou de \emph{usurpação}. O tyranno de \emph{governo} é o que possue \emph{legitimamente} a auctoridade, de que abusa; --- o de \emph{usurpação} é o que \emph{illegitimamente} occupa a auctoridade. --- Quando o tyranno é de \emph{usurpação}, mas faz bom uso da auctoridade, os subditos devem obedecer-lhe, pois assim o exige a ordem publica. --- Quando o tyranno é de \emph{governo}, o subdito póde fazer-lhe uma \emph{resistencia passiva}, a qual consiste em não obedecer ás suas ordens injustas, sem comtudo conspirar contra elle; mas não póde fazer-lhe uma \emph{resistencia activa}, \& qual consiste em destituir o soberano da sua auctoridade, não só porque os subditos não têm o direito de destituir o soberano (pois a auctoridade vem de \textsc{Deus} e não do povo), mas tambem porque abrir-se-hia a porta a continuas guerras civis.}.

\pg{781} c) \textbf{Dever} de fidelidade, --- Os subditos devem cooperar para o bem commum da sociedade civil. Ora não só não cooperariam para o bem commum, mas conspirariam contra o seu paiz, se não fossem \emph{fieis} ao seu soberano \footnote{Os subditos não podem attentar contra as formas do governo, nem conspirar e excitar reBellióes contra o legitimo soberano, sem se tornarém culpados de impiedade e de parricidio. --- A felicidade de um paiz não depende da forma do governo, mas sim da qualidade dos imperantes. Além d'isso, os. males, de que a revolução é causa, excedem toda e qualquer vanta- \textbar{} gem, que possa derivar da mudança do soberano ou da forma de governo.}.

186. \textbf{Especies de direitos} dos subditos. --- Os \emph{subditos} podem considerar-se ou emquanto constituem a \emph{sociedade civil}, as \emph{classes sociaes} e a \emph{familia}, --- ou emquanto são \emph{pessoas particulares}. Os \emph{direitos}, pois, dos \emph{subditos} comprehendem os direitos proprios da \emph{sociedade em geral}, das \emph{classes sociaes}, das \emph{familias} e das \emph{pessoas particulares} ou \emph{individuos} \footnote{A todo o \emph{dever} corresponde um \emph{direito}. Se o soberano tem \emph{deveres} para com os \emph{subditos}, quer estes se considerem emquanto constituem a sociedade civil, ou as classes sociaes, oü as familias, quer se considerem emquanto individuos, é evidente que os \emph{subditos} devem ter e effectivamente têm \emph{direitos} para com o soberano. --- Não nos demoramos nestes \emph{direitos} dos \emph{subditos}; porque a sua \emph{maturesa} e \emph{extensão} foram declaradas, quando tratavamos dos correspondentes \emph{deveres} dos \emph{soberanos}.}.

a) \textbf{Direitos proprios} da \textbf{sociedade em geral}, --- A \emph{sociedade em geral} tem o direito de exigir que o soberano procure o bem \emph{externo} e \emph{commum}, e por isso promova a prosperidade \emph{material} e \emph{economica}, e á prosperidade \emph{religiosa}, \emph{moral} e \emph{intellectual} \footnote{O bem \emph{commun} e \emph{externo} é á causa que formou e conserva a sociedade civil, composta da multidão e da auctoridade. --- A multidão, se tem o \emph{dever} de obedecer ao soberano em vista do fim \emph{commum}, tem tambem o \emph{direito} de exigir que o soberano exerça o poder para a consecução d'esse fim.}.

\pg{782} b) Direitos proprios das classes, --- As \emph{classes sociaes} têm o \emph{direito} de exigir que o soberano as favoreça, tutelando e promovendo os seus interesses, afastando os perigos, e conservando entre ellas o espirito de paz e de justiça \footnote{Em geral, uma classe social, que tem um fim honesto, meios licitos, e cujo bem não é contrario ao bem da sociedade civil, tem o direito ao respeito e á protecção do soberano; porque os homens têm o direito \emph{énnato} da \emph{associação}, com tanto que esta não contenha elementos viciados e não seja contraria ao bem social --- D'onde se vê --- 1.º) que as Ordens Religiosas, embora se considerem como associações meramente naturaes, tem direito á protecção do Estado, porque o seu fim e os meios são coisas boas e honestas, e summamente vantajosas para a sociedade civil; --- 2.9) que as \emph{associações secretas}, como as \emph{maçonicas}, cujo fim é contrario á Religião e á Patria, devem ser prohibidas pela auctoridade publica.}.

c) Direitos proprios das familias, --- As \emph{familias} têm o \emph{direito} de exigir que o soberano respeite e defenda os seus direitos, quer estes se refiram á sociedade \emph{conjugal}, quer á \emph{parental}, quer á \emph{Aeri} \footnote{As \emph{familias}, sendo a origem e a base da sociedade civil, são anteriores á mesma sociedade e possuem \emph{direitos} provenientes da propria lei natural e não da auctoridade do Estado. Por isso a \emph{familia} póde exigir que o soberano não a destrua e opprima, mas que a defenda e auxilie.}.

d) Direitos proprios dos individuos. --- Os \emph{individuos} têm o direito de exigir que o soberano respeite e defenda os seus direitos, e procure, embora \emph{indirectamente}, a sua felicidade \footnote{O \emph{individuo} tem muitos e diversos direitos \emph{innatos}, e \emph{adquiridos}. Estes \emph{direitos} são limitados, e a limitação tem a sua razão, como dissemos, ou nos \emph{deveres}, a que está sujeito o proprio individuo, --- ou nos \emph{direitos}\pgn{783}mais elevados, que se encontram nos outros individuos ou na sociedade. O soberano, pois, deve tutelar e limitar esses \emph{direitos} dos individuos. --- Hoje, infelizmente, julga-se que o direito humano é absoluto e illimitado. D'ahi a licença desenfreada de pensar, dizer e fazer tudo o que as paixões desordenadas dictam, com grande damno dos individuos, da familia e da sociedade. D'ahi as queixas contra a auctoridade civil; quando prohibe ou castiga o abuso da liberdade. A liberdade, que degenera em licença, é a peior das escravidóes, porque, sujeitando o homem ao imperio das paixões ignobeis, abate-o até á condição dos seres irracionaes.}.

%%>>> cap-11

\chapter*{CAPITULO TERCEIRO\\[0.3em]\large DIREITOS E DEVERES INTERNACIONAES}
\addcontentsline{toc}{chapter}{CAPITULO TERCEIRO --- DIREITOS E DEVERES INTERNACIONAES}

{\small\noindent\textsc{Summario}: --- Direitos e deveres mutuos das sociedades civis. --- Deveres da sociedade civil para com a Egreja.\par}

%%>>> art-72

\section*{Artigo I.\\ Direitos e deveres mutuos das sociedades civis.}
\addcontentsline{toc}{section}{Artigo I. --- Direitos e deveres mutuos das sociedades civis.}

187. \textbf{Direito} internacional. --- Direito \emph{internacional} é a somma dos \emph{direitos} e dos \emph{deveres}, que liga uma sociedade civil com as outras sociedades civis, ou com a sociedade religiosa, que \pg{783}  é a Egreja. --- D'ahi a razão dos dois artigos, em que se divide este \emph{terceiro} capitulo \footnote{O direito \emph{internacional} baséa-se na propria \emph{lei natural}, emquanto esta regula as mutuas relações entre as diversas sociedades. Na verdade, \textsc{Deus}, creando o homem sociavel, creou tambem as diversas sociedades, que são o meio de satisfazer a esta sociabilidade; e, assim como dictou leis ao homem, para que este consiga o seu fim, tambem dictou leis ás sociedades, para que realisem a ordem estabelecida entre ellas. Se o Direito \emph{internacional} não fosse regulado pela \emph{let natural}, devia defender-se um d'estes absurdos: --- ou que as sociedades, nas mutuas relações, não estão sujeitas a nenhuma ordem ou plano, --- ou que, existindo esta ordem ou plano, \textsc{Deus} não promulgou leis para a sua realisação.\par O Direito \emph{internacional} costuma ás vezes chamar-se \emph{Direito das gentes}. Mas ha differença entre elles: porquanto o \emph{Direito das gentes}, \emph{actualmente}, significa o Direito \emph{positivo} que, por \emph{costume} ou por \emph{pactos}, se tem introduzido entre todos os povos; ao passo que o Direito \emph{internacional} é uma parte do Direito \emph{natural}, porque se basêa nas leis que \emph{immediatamente} derivam da natureza racional do homem. --- Dissemos --- \emph{actualmente}; porque os jurisconsultos romanos entendiam por \emph{Direito das gentes} o \emph{Direito} que vigora entre todos os homens e deriva da razão natural.\par O Direito \emph{internacional} costuma geralmente limitar-se ás mutuas relações entre as diversas sociedades civis, e não ás relações entre a sociedade civil e a religiosa. Collocamos, comtudo, no Direito \emph{internacional} as relações entre a sociedade civil e a Egreja, porque tambem a Egreja é uma sociedade \emph{perfeita}, \emph{e} porque, embora Ella seja uma sociedade sobrenatural, comtudo as suas relações com o Estado são determinadas á luz dos principios da razão natural.}.

188. \textbf{Sujeito} dos direitos e \textbf{deveres internacionaes}. --- O \emph{sujeito} dos direitos e dos deveres \emph{internacionaes} não é constituido \emph{immediatamente} pela pessoa \emph{physica} dos subditos nem pela dos soberanos, mas pela sociedade, que forma uma verdadeira pessoa \emph{juridica} \footnote{Todavia, devendo as pessoas \emph{moraes} ou \emph{jurídicas} ter um representante, que exerça os direitos, de que são dotadas, segue-se que o sujeito da auctoridade suprema é o orgão do Direito \emph{internacional} e o representante da sociedade. --- Como se vê, tratamos do Direito internacional \emph{publico}, que regula as mutuas relações das \emph{sociedades}, consideradas como pessoas juridicas, e não do Direito internacional \emph{privado}, que regula as relações juridicas entre os \emph{individuos} das diversas sociedades.}.

\pg{784} 189. \textbf{Especies} de direitos mutuos das sociedades civis. --- Os \emph{direitos} mutuos, que as sociedades civis têm entre si, são de duas especies: \emph{juridicos} e \emph{moraes}, conforme se baséam na virtude da \emph{justiça} ou na da \emph{caridade}. --- Os direitos \emph{juridicos} são \emph{innatos} ou \emph{adquiridos}, conforme derivam da propria natureza da \emph{sociedade}, ou de alguma cireumstancia ou facto accidental \footnote{O homem, pessoa \emph{p/ysica}, é dotado de muitos e diversos direitos.\par n D'estes direitos uns são \emph{innatos}, outros são \emph{adquiridos}. Os \emph{innatos} pertencem ao homem, como tal, isto é, emquanto é dotado de natureza humana, e por isso encontram-se em todos os homens, e dizem-se \emph{inznatos}, porque nascem com o homem. Os \emph{adquiridos} pertencem ao homem, não pelo facto d'elle possuir a natureza humana, mas por alguma circumstancia particular ou pessoal, em que elle se encontra. --- A sociedade \emph{civil}, pessoa \emph{moral}, é tambem dotada dos seus \emph{direitos}. Estes podem ser \emph{innatos} ou \emph{adquiridos}. Os \emph{innatos} pertencem á sociedade, como tal, isto é, emquanto possue a natureza de sociedade, e por isso encontram-se em todas as sociedades \emph{perfeitas}. Os \emph{adquiridos} pertencem á sociedade, não porque é sociedade, mas por causa de algum facto, costume ou circumstancia particular, e por isso não pertencem a todas as sociedades, embora perfeitas, e são diversos, conforme a diversidade das condições e das circumstancias da sociedade.}.

190. \textbf{Direitos} juridico-innatos das \textbf{sociedades} civis. --- Os direitos \emph{Juridico-inmatos} das sociedades civis são tres: o direito de \emph{conservação}, de \emph{aperfeiçoamento} e de \emph{livre actividade} \footnote{O homem, recebendo de \textsc{Deus} a vida, para que a empregue em se aperfeiçoar e alcançar o fim ultimo, que é a felicidade eterna, recebe da propria lei \emph{natural} o direito de \emph{conservação}, de \emph{aperfeiçoamento}, e de \emph{livre actividade}. Do mesmo modo, a natureza, querendo a existencia de varias sociedades civis, perfeitas e independentes, quer que ellas se conservem e tendam livremente para a perfeita consecução do seu fim, que é a felicidade \emph{temporal}, e por isso dá a cada uma d'ellas o triplice direito: de \emph{conservação}, de \emph{aperfeiçoamento} e de \emph{livre actividade}.}.

a) \textbf{Direito} de conservação. --- Este \emph{direito} abrange --- o direito de conservar a propria independencia, em que consiste a vida da sociedade perfeita, --- o direito de não soffrer nenhuma diminuição nos membros, de que consta, ou nos bens, que possue, --- e o direito de conservar a sua paz e prosperidade \footnote{No homem, o direito de, conservação abrange o direito de conservar a vida, a integridade dos membros e dos bens externos, e a saude da alma e do corpo. Na sociedade, o direito de conservação refere-se á \emph{independencia}, aos \emph{membros} e \emph{bens externos}, e á \emph{prosperidade material}. --- A sociedade, como tal, póde e deve possuir propriedades ou bens materiaes, de que tenha verdadeiro \emph{dominio}, e de que possa dispor livremente para o bem commum:\par Pergunta-se: Terá toda a nação direito \& \emph{independencia} ?\par Antes de respondermos a esta pergunta, devemos considerar que a palavra nação se toma em dois sentidos: um \emph{ethnografico}, outro \emph{jurídico}. No sentido \emph{ethnografico}, nação é um aggregado de familias, que tem a mesma\pgn{785}origem, fallam a mesma lingua e vivem num territorio geographicamente determinado, que se chama patria, isto é, \emph{terra dos paes}; --- no sentido \emph{juridico}, nação significa \emph{sociedade independente} ou \emph{Estado}. Esta distincção é muito sensivel; porque uma só nação, no sentido \emph{ethnografico}, póde ser um \emph{Estado} (como o \emph{Japão}), mas póde tambem constituir varios \emph{Estados} (como um dia a nação \emph{Hespanhola}, á qual pertenciam o Reino de \emph{Hespanha} e as republicas da \emph{America meridional}) ou póde estar unida com outras nações extrangeiras (como a nação \emph{Hungara}, que, no Imperio \emph{Austro-Hungaro}, vive com a nação \emph{allemá}, com a \emph{slava}, com a \emph{rumena} e com a \emph{italiana}).\par Esta distinecüo dà occasião a uma duplice resposta.\par Se a palavra \emph{nação} se tomar em sentido \emph{jurídico}, é claro que ella tem direito á \emph{independencia}; porque todo o Estado, grande ou pequeno, é sociedade \emph{perfeita}, e por isso essencialmente \emph{independente},\par Se a palavra \emph{nação} se tomar exclusivamente em sentido \emph{ethnografico}, ella não têm direito algum \& \emph{independencia} nem á \emph{unidade}; e por isso á usurpação praticada, em nome do \emph{principio} de \emph{nacionalidade}, contra um soberano legitimo, não passa de um verdadeiro latrocinio. --- Para darem uma apparencia de \emph{legalidade} ao que é uma aberta violação da justiça e do direito \emph{natural}, os modernos usurpadores recorreram ao celebre expediente dos \emph{plebiscitos}. Mas o recurso é tão pueril, como flagrante é a sua illegalidade. Os \emph{plebiscitos} não representam a vontade dos povos; e, embora a representassem, não poderiam justificar os roubos praticados contra o principe legitimo; poisque a auctoridade, de que o soberano está investido, não deriva do povo, mas immediatamente do proprio Dus.\par Todavia, póde e deve desejar-se --- que os povos da mesma \emph{nacionalidade}, sem offensa de nenhum Estado nem dos legitimos soberanos, formem um só Estado, federal ou unitario, --- ou que, embora estejam unidos com povos de diversas nacionalidades e sujeitos á mesma auctoridade, conseryem, dentro da unidade social, os seus caracteres, costumes, tradições, lingua, etc. --- Os povos da mesma \emph{nacionalidade} podem \emph{legitimamente} formar um \emph{Estado} em consequencia --- de \emph{cessão voluntaria}, --- ou de \emph{casamento} entre os soberanos (assim o reino de \emph{Castella} e o de \emph{Aragão} uniram-se pelo casamento de \emph{Isabel a Catholica} com \emph{Fernando}), --- ou de \emph{conquista} depois de uma guerra \emph{Justa}, --- ou de \emph{prescripção}.\par A \emph{prescripcüo}, como dissemos, verifica-se em relação ao \emph{usurpador}. Com -o andar do tempo, póde succeder que o legitimo soberano perca o seu direito ao throno, de que foi injustamente privado; e 9 isto acontece porque o seu direito se acha em collisão com o outro mais forte, que á nação tem de não ser lacerada por discordias intestinas ou guerras civis; pois o bem particular deve ceder ao bem commum. --- Ha só um caso em que a \emph{prescripção} não é applieavel, e este caso refere-se á soberania temporal do Romano Pontifice sobre os seus Estados. Porquanto, na collisão do direito, que o Papa tem á soberania, e do direito, que o Estado tem á sua paz, está sempre superior o direito do Papa; visto que essa soberania é o direito da Egreja, que, sob todos os aspectos, ó immensamente superior ao do Estado. --- Além d'isso, o poder politico do Papa não é para a vantagem material das suas terras, mas para a garantia da liberdade do seu poder jurisdicional sobre toda a\par 50 Vor. II.\pgn{786}Egreja, e por isso, emquanto permanece a necessidade d'esta garantia, não ha e não pode haver prescripção em contrario. A este caso não pode applicar-se o direito commum; porque o Principado do Papa não tem o mesmo escopo que o dos outros reis.}.

b) Direito de aperfeiçoamento. --- Este direito consiste em que a sociedade póde empregar, sem que ninguem a possa \textbf{impedir} \pg{785} \pg{786} , meios cada vez mais efficazes, para alcançar, com maior segurança e perfeição, o proprio fim \footnote{Muitos politicos modernos, imitando \emph{Montesquieu}, quando vêm que uma nação entra, pelo emprego de meios licitos, no caminho do progresso e da prosperidade, ensinam que, para a manutenção do equilibrio internacional, as outras nações podem e devem impedir, mesmo pela \emph{guerra}, esse progresso e prosperidade, que póde constituir uma ameaca continua contra a paz universal --- Nada mais injusto. Cada paiz tem o direito de se aperfeigoar, e se, para a consecução do seu fim, emprega meios licitos, não póde, sem grande injustiça, ser contrariado nos seus planos.}.

c) Direito de livre actividade. --- Este direito consiste em que a sociedade póde tender para o seu fim pelo modo, que julgar melhor, e pelo emprego dos meios, que lhe parecerem mais aptos, sem ser obrigada a receber de outra sociedade a lei ou regra, que deve seguir no exercicio da sua actividade \footnote{Uma nação tem, pois, direito a que as outras nações não se intromettam na propria administração politica, nem impeçam ao soberano o livre exercicio dos seus deveres. Porquanto, toda a nação é sociedade \emph{perfeita}, isto é, independente das outras na sua existencia e nas suas operações. --- Esta theoria, porém, deve applicar-se ás nações, cujo regimen interno é regular e não constitue um perigo imminente para as outras nações, . (8) E' triplice, pois, a fonte dos direitos \emph{juridico-adquiridos}; --- o uso ou o \emph{costume} dos povos, a \emph{acquisição do dominio}, os \emph{tratados internacionaes}.\par a) O uso ou o \emph{costume} dos povos póde ser fonte de \emph{direitos}. Assim como, numa só sociedade, o \emph{uso} ou o \emph{costume} introduz um direito, assim tambem, entre varias sociedades, ou nações, o mesmo \emph{uso} ou \emph{costume} póde dar origem a alguns direitos; tal é o direito de \emph{inviolabilidade} e de \emph{independencia}, de que os embaixadores gosam. --- Estes direitos encontram-se em todas as nações, porque estão em harmonia com a natureza humana e facilmente se deduzem dos direitos \emph{innatos}. :\par b) \emph{A acquisicdo do dominio} é tambem origem de direitos \emph{juridico-adquiridos}. --- A razão é manifesta. Quem adquire o \emph{dominio} sobre uma coisa, adquire direitos, de que não póde ser despojado sem grande injustiça. --- Advertimos que a sociedade póde adquirir e possuir bens e propriedades em territorio pertencente a uma outra nação; mas, neste caso, o dominio não é absolutamente independente, pois é semelhante ao dos particulares e sujeito á auctoridade da nação, em cujo territorio estão situados os bens. Advertimos tambem que uma nação póde \emph{legitimamente} adquirir o \emph{dominio}}.

191. \textbf{Direitos} juridico-adquiridos das \textbf{sociedades} civis. --- Os principaes direitos \emph{juridico-adquiridos}, que as sociedades civis têm ou podem ter entre si, reduzem-se a uma triplice classe; pois uns desivam do \emph{uso} ou \emph{costume}, --- outros da \emph{acquésigao do dominio}, --- outros dos \emph{tratados internacionaes} \footnote{\pgn{787}de uma porção de mar, que banha as costas ou que avança pela terra dentro (\emph{golfo}), porque esta oceupação póde ser exigida pela segurança e utilidade da propria nação, --- mas não póde adquirir o dominio do \emph{alto Mar}, pois este não póde tornar-se \emph{util} por industria pessoal, e póde satisfazer egualmente ás necessidades e exigencias de todos.\par 0) Os \emph{tratados internacionaes} dão tambem origem a direitos \emph{Juridicos}. --- Estes tratados, para serem licitos e validos, devem ser feitos pelo soberano e em nome da sociedade, que elle representa, devem ser honestos e justos em relação ao \emph{objecto}, \emph{ás circumstancias} e ao \emph{fim}. --- Os tratados \emph{internacionaes} podem ter por \emph{objecto} --- ou uma coisa, que já pertence ao Direito \emph{natural}, --- ou uma coisa, que se accrescenta ao Direito \emph{natural}. Os que têm por \emph{objecto} uma coisa, que se accrescenta ao Direito \emph{natural}, produzem uma obrigação que anteriormente não vigorava; e dividem-se em \emph{eguaes} ou \emph{deseguaes}, conforme as coisas, que as partes contrahentes dão e recebem, são \emph{eguaes} ou \emph{deseguaes}. --- Os tratados podem ter por \emph{termo principal} e \emph{directo} a pessoa do soberano, ou a propria nação. No primeiro caso, o tratado diz-se \emph{pessoal}; no segundo, \emph{real}. Sendo \emph{pessoal}, o tratado dissolve-se com a morte do soberano; --- sendo \emph{real}, não se dissolve com essa morte. --- Os tratados \emph{dissolvem-se} --- 1.9) por mutuo consentimento das partes; --- 2.9) por morte de uma das partes contrahentes, e por isso por morte \emph{physica} do soberano (se o tratado é \emph{pessoal}) ou por morte \emph{civil} da nação (se o tratado é real); --- 9.9) por acabar o tempo determinado em que devia vigorar o tratado; --- 4.9) pela absoluta e perpetua impotencia, que uma parte tem de satisfazer aos seus deveres.\par *\par A violação dos direitos \emph{jurídicos} não só é uma \emph{offensa}, mas é uma \emph{injuria} ou \emph{injustiça}, praticada contra a nação, que os possuia. --- Esta \emph{injuria} deve ser reparada. --- Quando a nação, que violou os direitos \emph{jurídicos} de outra nação, não dá, por si e espontaneamente, uma reparação, ou porque julga ter procedido rectamente, ou porque, embora conheça a necessidade da reparação, não a quer dar, então, se não existir uma \emph{auctoridade commum} a todas as nações, que possa impor as suas decisões e fazer respeitar o direito, ha dois meios para reparar a ordem \emph{juridica} violada. Estes dois meios são a \emph{guerra} e a \emph{arbitragem}.\par a) Guerra. --- \emph{Guerra}, tomada no sentido rigoroso, é o \emph{estado de duas} ou \emph{mais nações que luctam entre si}, \emph{fazendo uso da força}, \emph{com o fim de tutelar o direito}; --- ou mais-brevemente: é \emph{uma violenta defesa da ordem}. --- E' \emph{offensiva} ou \emph{defensiva}, conforme se emprega a força --- para vingar uma injuria recebida, --- ou para repellir os ataques dos inimigos.\par A guerra é \emph{licita}. Porquanto as pessoas, revestidas da auctoridade suprema, têm a obrigação de proteger os seus subüitos não só contra os perigos, que nascem do seio da propria nação a que pertencem, mas tambem contra os ataques, que partem das outras nações e que são mais funestos. --- Esta \emph{liceidade} torna-se mais evidente, se considerarmos as duas especies de guerra. E' licita a guerra \emph{defensiva}; porquanto, na guerra \emph{defensiva},\pgn{788}uma nação repelle a injusta aggressao da outra nação; ora, se um particular póde repellir a força com a força, muito mais o poderá uma nação, aggredida por outra. E' licita á guerra \emph{offensiva}; porquanto a ordem juridica, violada pela injuria, deve ser restabelecida; ora muitas vezes não póde ser restabelecida senão pela guerra. --- Nem se diga que a guerra traz muitos males; pois estes são menores do que aquelles que se seguiriam, se a guerra nunca fosse licita.\par Para que a \emph{guerra} seja \emph{licita}, são necessarias as seguintes condições: --- 1.8) a auctoridade, que declara ou acceita a guerra, deve ser a do \emph{soberano}, porque só este tem o dever e o direito de defender o bem publico do seu povo; --- 2.º) a causa deve ser \emph{justa}, o que .é evidente; --- 3.º) a intenção, com que se faz a guerra, deve ser \emph{recta}, de modo que se tenha em vista a promoção do bem e a repressão do mal.\par O modo por que se deve fazer a \emph{guerra} é o seguinte. Antes de comecar a guerra contra um Estado, a nação offendida deve pedir-lhe uma reparação ou satisfação adequada, e, só no caso de recusa, póde empregar a força. Depois de começada a guerra, é permittido fazer ao inimigo todo o damno, necessario para alcançar a victoria, comtanto que não seja intrinsecamente mau (assim n\&o s\&o permittidas as mentiras, a perfidia, a morte \emph{directa} dos innocentes, etc.) --- Aocabada a guerra, o vencedor póde exigir do vencido, como justa compensação dos damnos soffridos na guerra, uma condigna satisfação da injuria e uma garantia segura contra futuras aggressóes. --- Esta garantia póde ser a annexacüo de alguma provincia, e, nalguns casos, a morte politica do paiz vencido, a qual consiste na perda da independencia. D'ahi o \emph{direito de conquista}. --- Como se vê, o \emph{direito de conquista} basóa-se sobre a \emph{justiça} da guerra e sobre a necessidade de uma \emph{garantia} contra possiveis aggressóes da nação vencida. Por isso deve reprovar-se a theoria --- de \emph{Hegel}, para quem \emph{toda} a conquista é \emph{legitima}, porque representaria a fatale progressiva evolução da humanidade, --- e de \emph{Cousin}, para quem toda a conquista é \emph{justa}, porque o vencedor seria sempre melhor que o vencido. E' a famosa \emph{theoria dos factos consummados}, que saneciona o direito da forga brutal, o direito dos ladróes e dos assassinos.\par \textbar{}) Arbitragem. --- \emph{Arbitragem} é a sentença, que, numa questão entre duas ou mais nações, é dictada por um juiz escolhido livremente por ellas. --- A \emph{arbitragem} tem a immensa utilidade de poupar as grandes despezas e os funestos horrores da \emph{guerra}, e, ao mesmo tempo, de fazer justiça a quem a merece, --- A instituição, pois, de um tribunal \emph{internacional} e \emph{permanente} de \emph{pas} seria um beneficio inealeulavel para as pobres nações, que gemem sob o peso dos impostos, lançados para a manutenção da \emph{pas armada}. Os membros d'este tribunal deviam ser os mais illustres e honestos personagens de cada nação. O seu presidente devia ser naturalmente o Summo Pontifice\pgn{789}Romano, que, no pensamento de todos, é a mais alta personificação da Religi\&o, da justiça e da paz. Os proprios protestantes concordam comnosco neste ponto. Em 21 de margo de 1887, o jornal inglez --- \emph{Saint James's Gazette} --- publicou um artigo, demonstrando que é immensamente util a constituição de um tribunal de \emph{arbitragem} ou de \emph{pas} para as questões \emph{internacionaes}, o que, para desempenhar esse papel, ninguem está em melhores condições de imparcialidade e de auctoridade do que o Papa. Em 25 do mesmo mes e anno, outro jornal inglez --- \emph{Pall Mall Gasette} --- publicou outro artigo no mesmo sentido, acabando por dizer que talvez o Papa chegará a ser o centro da paz do mundo. --- A idea da instituição d'esse tribunal de \emph{pas}, presidido pelo Santo Padre, começou a desenvolver-se nestes ultimos tempos, sobretudo por causa da celebre \emph{arbitragem}, exercida pelo grande Pontifice \emph{Leão XIII}, entre a \emph{Allemanha} e a \emph{Hespanha}, na questão das \emph{Carolinas}. A escolha do Papa, como arbitro, indicada por um \emph{Bismarck}, mostra a todos os povos que no Vigario de N. S. Jesus \textsc{CHRISTO} encontrarão sempre um juiz imparcial e justo, e reconhece que a missão do Pontificado Romano é missão de paz e de amor.\par Dissemos que á \emph{guerra} e a \emph{arbitragem} são meios para reparar a ordem \emph{juridica} violada, \emph{se ndo existir uma auctoridade commum a todas} as nações, para nos referirmos ao projecto de uma sociedade \emph{universal}, defendido por alguns juristas. E' o projecto da \emph{etAnarchia}, isto é, da sociedade universal, que seria formada por todas as nações, se estas se unissem sob uma só auctoridade, salva a independencia de cada uma.\par Esta sociedade \emph{universal} seria admiravel e não poderia deixar de dar excellentes resultados; mas a sua actuação é muito difficil, considerando os antagonismos e os odios, que dividem os varios povos do mundo. Só num caso, a \emph{ethnarchia} poderia constituir-se: se todos os povos do mundo abracassem, como é seu dever, a verdadeira Religião, que é a Catholica. Com effeito, tal feliz acontecimento, sujeitando todos os homens ao Summo Pontifice Romano, reduziria os espiritos a uma unidade de crenças e de affectos, e assim tornaria possivel a união entre todas as nações, que, mesmo nas coisas temporaes, dependeriam da auctoridade, sempre justa, desinteressada e benefica, do mesmo Vigario de N. S. Jesus CHRISTO.\par A \emph{ethnarchia}, de que fallamos, não deve confundir-se com a idea pantheista ou semi-pantheista, segundo a qual, assim como o espirito tem um só corpo em cada homem, assim tambem a humanidade deve ter um só corpo politico, que é o \emph{Estado} no seu mais alto grau de perfeição. \emph{BluntscAli} diz: « O Estado perfeito é egual á humanidade corporalmente visivel. O Estado ou o reino universal ó o ideal da humanidade progressiva » (A \emph{doutrina do Estado moderno}, vol. 1, pag. 26 e 27). O escriptor allemão, como é evidente, transforma em identidade uma simples analogia, que existe entre o Estado e o corpo humano, fazendo da humanidade um todo homogeno e por isso excluindo a pluralidade e variedade das nações, em que está dividida a humanidade e que é necessaria para o conceito da verdadeira \emph{ethnarchia}. \emph{Krause}, coherente com o seu systema pantheista, tambem falla de um unico Estado, \& que um dia ha de reduzir-se a humanidade,}.

\pg{787} 192. \textbf{Direitos moraes} das \textbf{sociedades civis}. --- Os direitos \emph{moraes}, que as sociedades civis têm entre si, podem reduzir-se aos direitos á \emph{humanidade} e á \emph{beneficencia}.

\pg{788} [93] [5]

a) \textbf{Direito} á humanidade. --- Direito á \emph{Humanidade} é o poder moral, que toda a sociedade tem de exigir das outras um auxilio, que é util a si e não é prejudicial a ellas. --- Tal direito existe na sociedade; pois esta, tendo o direito á propria conservação e aperfeiçoamento, deve ter o direito ao que \emph{positivamente} concorre para esse fim e que não é prejudicial ás outras sociedades.

\pg{789} b) \textbf{Direito} á \textbf{beneficencia}, --- Direito á \emph{beneficencia} é, o \textbf{poder} moral, que toda a sociedade tem de exigir, nas suas hecessidades, o socorro das outras sociedades, embora estas soffram

\pg{790} PARTE SEGUNDA --- DIREITO NATURAL algum incommodo. --- Tal direito existe na sociedade, pois é uma consequencia do direito, que assiste a toda a sociedade em relação á sua subsistencia e que é superior ao direito, que as outras sociedades têm ás suas commodidades \footnote{Os obsequios de \emph{humanidade} distinguem-se dos de \emph{beneficencia}; pois os de \emph{humanidade} fazem-se sem incommodo ou damno proprio, como é ensinar o caminho a um cego, dar um bom conselho a quem o pede, etc., ao passo que os obsequios de \emph{beneficencia} trazem algum incommodo ou damno, como é soccorrer aos outros com o dinheiro ou com o trabalho. Os obsequios da \emph{beneficencia} são mais nobres e elevados que os de \emph{humanidade}.\par Os direitos, que uma sociedade tem á /\emph{iwmanidade} e \& \emph{beneficencia} das outras sociedades, são meramente \emph{moraes}, porque não se baséam ma \emph{justiça} mas só na \emph{caridade}; visto que a \emph{humanidade} e á \emph{beneficencia} dos \emph{outros} não são coisas nossas. Por isso, estes direitos cessam, quando uma sociedade não póde cumprir os \emph{deveres}, correspondentes a estes direitos \emph{moraes}, sem grande damno seu; pois, nestas circumstancias, a \emph{caridade}, devendo ser \emph{ordenada}, manda que o bem proprio seja preferido ao dos outros. --- Todavia, em caso de extrema \emph{necessidade}, o direito, que uma sociedade tem á \emph{liumanidade} e \emph{beneficencia}, transforma-se de \emph{moral} em \emph{juridico} e por isso impõe ás outras sociedades, que se não encontram na \textbf{mfesma} necessidade, um dever \emph{jurídico}, de que não podem esquivar-se sem grande injustiça; porque, em caso de \emph{extrema necessidade}, torna-se \emph{propriedade} da sociedade e do proximo o que é necessario para a sua conservação.\par D'onde se segue que uma nação, aggredida injustamente por outra, tem o \emph{direito} de pedir auxilio ás outras nações, \emph{e} nenhum Estado póde obstar a que se preste este auxilio. Porquanto, cada nação tem o direito \& \emph{conservação} e por isso tem o direito de empregar todos os meios para esse fim, comtanto que não offenda os direitos das outras nações. Ora pedir auxilio ás outras nações é sempre util e muitas vezes necessario, e não offende os seus direitos. Logo a nação, que é injustamente aggredida, póde pedir auxilio á outra nação, e esta, podendo, tem o dever e portanto o direito de prestar \textbf{esse} auxilio. --- D'onde se vé que o famoso principio de \emph{não-intervenção} é \emph{immoral}, \emph{injusto} e \emph{funesto}. E) \emph{immoral}; porque se oppõe ao amor do proximo, que obriga a prestar auxilio, não só ao individuo, mas tambem á sociedade. E \emph{injusto}; porque offende os direitos --- não só da nação, que pede auxilio, pois a priva de um meio util e talvez necessario para a sua conservação, --- mas tambem da nação, que o presta, porque lhe nega o exercicio de um direito \emph{e} o cumprimento de um dever. E' \emph{funesto}; porque protege os ataques da força material contra as nações que só têm a força do direito.}.

193. \textbf{Especies} de \textbf{deveres} mutuos \textbf{das sociedades} civis. --- Os \emph{deveres} mutuos, que as sociedades civis têm entre si, são de duas especies: \emph{jurídicos} e \emph{moraes}. Na verdade, correspondendo a todo o \emph{direito} um \emph{dever}, é claro que, tendo cada sociedade direitos \emph{jurídicos} e \emph{moraes} em relação ás outras, estas devem ter deveres \emph{jurídicos} e \emph{moraes} em relação áquella \footnote{Como já dissemos, o dever, que a lei impõe a uma pessoa, ou corresponde a um verdadeiro direito, existente noutra pessoa, ou não, Se corresponde\pgn{791}, o dever é \emph{juridico}; se não corresponde, é apenas moral. Quem não cumpre os deveres \emph{jurídicos}, faz uma \emph{injuria} e offende a \emph{justiça} (assim quem não restitue um deposito, diz-se \emph{injusto}); mas quem não cumpre os deveres \emph{moraes}, \emph{offende} a pessoa que possue o respectivo direito \emph{moral}, mas não commette uma \emph{injuria}, pois não transgride a lei da \emph{justiça}, mas só a da \emph{caridade} (assim quem nega a esmola a um pobre, é \emph{deshumano}, mas não \emph{injusto}). Por isso, se o sujeito que deve cumprir um dever \emph{jurídico} póde ser obrigado pela coacção, o sujeito que deve cumprir um dever \emph{moral} não póde ser obrigado, pois trata-se de dar uma coisa que é propria \emph{d'elle} e não dos \emph{outros}.}.

\pg{791} 194. \textbf{Deveres} jurídicos das sociedades civis. --- Os deveres \emph{juridicos}, que as sociedades civis têm entre si, são os que correspondem aos mutuos direitos \emph{jurídicos}, e por isso são os deveres de respeitar --- os direitos de \emph{conservação}, de \emph{aperfeiçoamento}, de \emph{livre actividade} (direitos \emph{nnatos}), --- e os direitos, que derivam do \emph{uso} ou \emph{costume} dos povos, da \emph{acquisigdo do dominio}, e dos \emph{tratados internacionaes} (direitos \emph{adquiridos}) \footnote{Não é preciso que demonstremos a razão ou a base dos deveres \emph{juridicos}, que uma sociedade civil tem em relação ás outras; porque a existencia d'esses deveres foi demonstrada, quando tratámos da existencia dos direitos \emph{juridicos}, que as sociedades tem entre si, pois a todo o \emph{direito} corresponde necessariamente um \emph{dever}. :}.

a) Dever de respeitar o direito de conservação. --- Este \emph{dever} consiste em que uma nação deve respeitar --- a independencia das outras nações, --- os individuos, de que ellas se compõem, --- e os bens materiaes, que possuem, e não deve fazer coisa contraria á paz e á prosperidade das outras \footnote{E, pois; injusto --- conspirar ou attentar contra a independencia de uma nação, --- lesar os seus subditos na pessoa, na honra ou nos bens materiaes, --- excitar e manter, no seu seio, discordias e rebellióes, etc.}.

b) Dever de respeitar o direito de aperfeiçoamento. --- Este \emph{dever} consiste em que uma nação não póde obstar a que uma outra nação avance, pelo emprego de meios honestos e licitos, no caminho da civilisação e do progresso \footnote{É, pois, injusto obstar a que uma \emph{nação} se entregue \& industria e ao commercio, ou se allie com outra nação e d'ella receba auxilio, ou promova o maior bem e interesse dos seus subditos, etc.}.

c) \textbf{Dever} de respeitar o direito de livre actividade. --- Este \emph{dever} consiste em que uma nação não póde contrariar, sem motivo justificado, a liberdade de movimentos numa outra, nem impor-se-lhe no legitimo exercicio das suas funceções \footnote{E, pois, injusto --- pretender dirigir os negocios das outras nações, --- privar os outros soberanos da liberdade necessaria para o bom andamento da administração publica, etc.}.

d) Dever de respeitar os direitos, que derivam do uso ou costume, da acquisição do dominio e dos tratados internacionaes., --- \pg{792}  Este dever consiste em que uma nação não póde fazer coisa alguma, contraria aos direitos, que as outras nações possuem em virtude --- ou do \emph{uso} e \emph{costume} dos povos, --- ou da \emph{acquisição do dominio}, --- ou de \emph{tratados internacionaes} \footnote{É, pois, injusto --- violar o domicilio ou offender de qualquer modo os \emph{embaizadores} das outras nações, --- arrebatar por violencia ou por fraude uma coisa, cujo dominio foi adquirido por outra nação, --- não satisfazer a empenhos assumidos por meio de tratados \emph{internacionaes}, etc.}.

195. \textbf{Deveres moraes das sociedades} civis. --- Os deveres \emph{moraes}, que as sociedades civis. tém entre si, reduzem-se aos deveres de \emph{humanidade} e de \emph{beneficencia}.

a) \textbf{Deveres} de humanidade. --- Os deveres de \emph{humanidade} consistem em que uma sociedade tem a obrigação de praticar, em favor de outra sociedade, obsequios, que são uteis a quem os recebe e não são prejudiciaes a quem os faz. --- Estes deveres existem; porquanto uma sociedade, reconhecendo nas outras sociedades uma natureza semelhante á sua, deve ter-lhes um amor semelhante ao que tem a si mesma, e por isso deve auxilial-as nas suas criticas cireumstancias.

b) \textbf{Deveres} de beneficencia, --- Os deveres de \emph{beneficencia} consistem em que uma sociedade tem a obrigação de soccorrer, com algum incommodo proprio, as outras sociedades nas desgraças, que as acommettem. --- Estes deveres existem; porque a lei da \emph{caridade} pede que se preste auxilio ao proximo nas suas necessidades, embora com algum incommodo proprio \footnote{As sociedades civis têm entre si deveres não só \emph{jurídicos}, mas tambem moraes. Estes deveres reduzem-se todos ao grande preceito de mutua caridade, imposto pela lei \emph{natural} ás sociedades. Porquanto, se cada homem deve amar os outros homens, porque a natureza de um é semelhante a dos outros; pelo mesmo motivo, uma nação deve amar as outras nações, pois a nação é constituida por homens, elevados a uma condição mais nobre, e dotados;emquanto compõem o corpo social, d'uma natureza mais culta e poderosa.\par Os deveres de mutua \emph{caridade} só obrigam em certos casos, Na \emph{pratica}, devem observar-se os seguintes criterios: --- 1.9) A obrigação existe só relativamente aos Estados, que têm verdadeira necessidade, e que não podem auxiliar-se a si mesmos. --- 2.º) Quando um Estado se encontra em verdadeira necessidade, os outros Estados, se poderem \emph{prestar-lhe} auxilio sem incommodo ou damno proprio, devem fazel-o, pois assim manda a caridade. Se não poderem sem incommodo ou damno, devem comparar este seu incommodo ou damno com o damno que o tal Estado \emph{soffrerà}, se não for \emph{auxiliado}; e esta comparação mostrará se, nesse caso, a lei de \emph{caridade} obriga, ou não. --- 9.9) Se varios Estados se encontrarem em \emph{eguaes} criticas \emph{circumstancias}, e o Estado não poder prestar auxilio a todos, deve observar-se a \emph{ordem devida}, preferindo os amigos aos inimigos, os dignos aos indignos, etc.\par As nações, para satisfazerem aos deveres de \emph{Denevolencia}, a que estão\pgn{793}sujeitas, não só devem exercer o dever e o direito de \emph{intervenção}, mas tambem devem exercer entre si mutuo \emph{commercio}, que consiste na. compra, venda ou permuta de coisas necessarias ou uteis; pois o \emph{commercio}, além de favorecer os interesses economicos e materiaes, promoye a cultura intellectual e moral das nações. --- Este dever, que uma nação tem de entreter relações \emph{commerciaes} com as outras, é \emph{geralmente} um dever \emph{moral} e não \emph{jurídico}; porque não se trata de restituir uma coisa \emph{alheia}, mas sim de dar ou permutar uma coisa \emph{propria}. Por isso o Estado póde --- julgar se deve, ou não, abrir relações commerciaes com outros Estados, --- prohibir ou permittir a exportação e a importação das mercadorias, --- impor os direitos, etc., conforme o exigir o interesse proprio e o de todos os cidadãos. --- Dissemos --- \emph{geralmente}; porque este dever de dar ou permutar as proprias producções com as de outro Estado pôde de \emph{moral} tornar-se \emph{jurídico}, e por isso ser dotado de força \emph{coactiva}, quando um Estado se encontra em extrema ou quasi extrema necessidade, ou quando os Estados se obrigaram por um contracto á permuta das proprias producções.\par Para conservar as mutuas relações de \emph{beneficencia} e de \emph{humanidade}, e tutelar seus interesses, os Estados têm o direito de communicar entre si por meio de \emph{embaixadores} ou \emph{ministros}. Estes, representando o seu paiz, devem ter toda a liberdade de se corresponder com os proprios governos, e não podem estar sujeitos á auctoridade do paiz, junto do qual estão acreditados. Por isso, segundo o Direito \emph{internacional}, os \emph{embaixadores} ou \emph{ministros} são \emph{inviolaveis} e \emph{independentes}.}.

%%>>> art-73

\section*{Artigo II.\\ Deveres da sociedade civil para com a Egreja.}
\addcontentsline{toc}{section}{Artigo II. --- Deveres da sociedade civil para com a Egreja.}

\pg{793} 196. \textbf{Egreja}. --- \emph{Egreja} é a sociedade de homens, unidos entre si pela profissão da mesma verdadeira Fé e pela participação dos mesmos Sacramentos, sob o regimen dos legitimos Pastores e do Romano Pontifice, Vigario de N. S. Jesus Caunisro \footnote{São tres, pois, as condições necessarias para que um homem pertença á Egreja: --- a \emph{profissão da verdadeira} fé, a \emph{communhao dos Sacramentos}, a \emph{subordinação aos legitimos pastores e sobretudo ao Romano Pontifice}. --- Por isso estão fóra da Egreja --- os Judeus, os Turcos, os Pagãos, que nunca tiveram a verdadeira fé, os herejes e os apostatas, que a abandonaram; --- os catechumenos, que nunca foram admittidos aos Sacramentos, e os excommungados, que foram privados da sua participação; --- os schismaticos, que, embora tenham a Fé e os Sacramentos, não estão sujeitos ao Papa.\par O homem é naturalmente inclinado para a sociedade \emph{religiosa}; e esta satisfaz ás mais sublimes e irresistiveis tendencias do coração humano. --- Nosso Senhor Jesus CunisTro, que veio ao mundo para restabelecer todas as coisas e elevar o homem a um estado sobrenatural, isto é, superior a todas as exigencias e aspirações da alma, fundou a \emph{Egreja}, que possuisse, por um modo \emph{eminente}, todas as perfeições proprias de uma sociedade \emph{religiosa}, e\pgn{794}assim satisfizesse superabundantemente ás inclinações naturaes e aos deveres sacrosantos, que nos unem com Ds.\par Repetimos que estas e as seguintes noções \&cerca da Egreja Catholica nos são fornecidas pela Theologia \emph{revelada}. Postos os principios, manifestados pela Revelação sobrenatural, o \emph{Direito natural} tira as consequencias, em harmonia com o senso commum e a boa logica.}.

\pg{794} 197. Fim \textbf{da Egreja}. --- O \emph{fim} da Egreja é dirigir os homens, pelo caminho da \emph{virtude}, para a \emph{felicidade eterna}, a qual consiste na \emph{visão directa} e \emph{intuitiva} de \textsc{Deus}. --- A \emph{virtude} é o fim \emph{proximo} da Eoreja, pois consegue-se nesta vida; --- a \emph{felicidade eterna} é o seu fim \emph{ultimo}, pois consegue-se noutra vida, e é o \emph{ultimo} termo das tendencias da alma humana \footnote{O fim \emph{ultimo} é a \emph{felicidade eterna}. Esta felicidade, pelos merecimentos de N. S. Jesus OnnisTo, consiste na visão \emph{directa} ou \emph{intuitiva} da propria Essencia Divina, e por isso é um fim \emph{sobrenatural}. Sendo \emph{sobrenatural} o \emph{fim}, devem ser tambem \emph{sobrenaturaes} os \emph{meios}. Por isso a \emph{virtude}, que é o \emph{fim} da vida presente e o \emph{meio} para a futura, deve ser \emph{sobrenatural}, isto é, deve derivar de um principio \emph{sobrenatural}, que é á \emph{graça} do Redemptor.}.

198. \textbf{Propriedades} da \textbf{Egreja}. --- A Egreja é sociedade \emph{divina}, \emph{espiritual}, \emph{sobrenatural}, \emph{juridica}, \emph{perfeita} e \emph{suprema}. \footnote{Não é nosso proposito enumerar \emph{todas} as \emph{propriedades} da Egreja. Limitamo-nos a referir as principaes que convém á Egreja, emquanto \emph{sociedade}. --- O caracter \emph{social} da Egreja vê-se pela sua definição. j}.

a) A \textbf{Egreja} é sociedade divina, --- Uma sociedade é \emph{divina} no sentido rigoroso da palavra, se foi, por um modo \emph{Immediato} e \emph{positivo}, fundada pelo proprio \textsc{Deus}. Ora a Egreja foi, por um modo \emph{Zmmedi?ato} e \emph{positivo}, fundada por N. S. Jesus \textsc{CnrISTO}, que é \textsc{Deus}. Logo a Egreja é sociedade \emph{divina} \footnote{A sociedade \emph{civil} e a \emph{domestica} foram fundadas \emph{immediatamente} por \textsc{DEus}, mas por um modo \emph{natural}, isto é, emquanto \textsc{Deus} é Auctor da \emph{natureza}. A Egreja, porém, foi fundada \emph{immediatamente} por \textsc{DEus}, mas por um modo \emph{positivo}, isto é, foi fundada por \textsc{Deus}, emquanto é Auctor da \emph{graça}. --- Por isso a Egreja é sociedade \emph{divina}, no mais rigoroso sentido da palavra; porquanto foi o proprio \textsc{Deus} quem ordenou a \emph{forma} do regimen, escolheu o \emph{sujeito} da auctoridade, indicou e presereveu os \emph{meios} e promulgou as \emph{leis}, que a regulam. --- A divindade do Chistianismo, e por isso da Egreja Catholica, foi confirmada por innumeraveis prophecias e milagres, contidos no Antigo e no Novo Testamento. Ao philosopho basta \emph{repetir} o grande argumento, exposto por dois genios --- \emph{Santo Agostinho} e \emph{Dante}. A conversão do mundo ao Christianismo ou foi effeito de milagres, ou não. Se foi, o Christianismo é \emph{divino}, porque só \textsc{Deus} póde fazer milagres. Se não foi, a conversão do mundo ao Christianismo realisada, sem milagres, por doze pescadores, pobres e ignorantes, no meio de tantos obstaculos e perseguições, é o maior de todos os milagres, e por isso o Christianismo é \emph{divino}.}.

b) A \textbf{Egreja} é sociedade espiritual. --- Uma sociedade é \emph{espiritual}, quando o seu fim é \emph{espiritual}; pois é o \emph{fm} que \emph{especifica} \pg{795}  a sociedade. Ora o fim da Egreja é \emph{espiritual}; porque foi instituida para a consecução da \emph{felicidade eterna}, que é um bem \emph{espiritual}. Logo a Egreja é sociedade \emph{espiritual} \footnote{D'onde se vê que a Egreja e o Estado são sociedades \emph{especificamente} diversas; visto. que o fim do Estado é a felicidade \emph{temporal} dos cidadãos, e o fim da Egreja é a felicidade \emph{espiritual} dos fieis.}.

c) \emph{A} Egreja é sociedade sobrenatural --- Uma sociedade é \emph{sobrenatural}, se o seu fim é \emph{sobrenatural}. Ora o fim da Egreja é \emph{sobrenatural}; porque a visão \emph{intuitiva} e \emph{directa} de Dzvus, em que consiste a felicidade eterna, excede as exigencias da na--- tureza creada. Logo a Egreja é sociedade \emph{sobrenatural} \footnote{A Egreja é \emph{sobrenatural} por muitas e diversas razões: --- pela sua \emph{origem}, pois deriva de N. S. Jesus \textsc{CHRISTO}, --- pela sua \emph{constituição}, que foi dada pelo mesmo Jesus Onnisro, --- pelos \emph{meos}, que excedem os limites da força humana, --- e, sobretudo, pelo seu \emph{fim}, que é a visão \emph{directa} e \emph{intuitiva} de \textsc{Deus}. --- A visão \emph{directa} e \emph{intuitiva} de \textsc{Deus} é um fim \emph{sobrenatural}; pois, como vimos na \emph{An/Aropologia}, \emph{o} homem, entregue ás suas proprias forças, só póde elevar-se ao conhecimento de \textsc{Deus} por meio das creaturas e por isso por um modo \emph{indirecto} e \emph{abstracto}. --- A visão \emph{directa} e \emph{intuitiva} de \textsc{Deus} é um bem tão grande, uma felicidade tão perfeita, que N. S. Jesus \textsc{CHRISTO}, para nol-a alcançar, derramou todo o seu Sangue precioso!}.

d) A \textbf{Egreja} é sociedade juridica, --- Uma sociedade é \emph{juridica}, quando os \emph{deveres}, que os socios tém para com ella, são \emph{jurídicos}, isto é, provenientes de verdadeiros \emph{direitos}, de que a propria sociedade é dotada. Ora os \emph{deveres}, que os homens têm para com a Egreja, são \emph{jurídicos}; pois os homens tem o dever de entrar na Egreja e de obedecer-lhe, e este dever corresponde ao direito, que N. S. Jesus Cungisro deu á Egreja de chamar e dirigir todos os homens. Logo a Egreja é sociedade \emph{juridica} \footnote{A Egreja é o reino visivel de \textsc{Deus}. O Rei é N. S. Jesus CHRISTO, representado pelo Pontifice Romano. Por isso a auctoridade do Papa, ou da Egreja, é a propria auctoridade de Jrsus Cnnisro. D'onde se segue que os deveres dos homens para com a Egreja são da mesma natureza ou especie que os seus deveres para com \textsc{DEus}. Ora os deveres do homem para com \textsc{Deus} são rigorosamente \emph{juridicos}; porque se fundam na \emph{justiça} e correspondem ao \emph{direito absoluto} de \textsc{Deus} sobre todas as creaturas, que d'Elle receberam tudo o que tem. Logo tambem os deveres do homem para com a Egreja são \emph{jurídicos}. --- A distinecio dos deveres em \emph{jurídicos} e \emph{moraes} admitte-se só nas mutuas relações entre os homens. Porquanto, fóra dos casos de \emph{justiça}, o homem não póde pretender de outro homem nenhuma coisa, como \emph{devida}. Se ha obrigação de a dar, esta obrigação é apenas \emph{moral}, pois é imposta por uma virtude distincta da \emph{justiça}, como é a \emph{beneficência}.}.

e) A \textbf{Egreja} é sociedade perfeita. --- Uma sociedade é \emph{perfeita}, quando possue em si todos os meios necessarios para a consecução \pg{796}  do seu fim, de modo que não depende de outra sociedade superior. Ora a Egreja possue em si todos os meios necessarios para a consecução do seu fim, e por isso não depende de outra sociedade superior. Logo a Egreja é sociedade \emph{perfeita} \footnote{Uma sociedade de ordem \emph{sobrenatural}, como é a Egreja, é necessariamente \emph{perfeita}. Porquanto, uma sociedade é \emph{perfeita}, quando possue em si todos os meios necessarios para a consecução do seu fim. Ora a Egreja não pode deixar de possuir todos esses meios; visto que, se os n\&o possuisse, não os poderia receber das outras sociedades, pois os meios, que a Egreja emprega para o seu fim, devem ser \emph{sobrenaturaes}, e as outras sociedades, tendo um fim \emph{natural}, só possuem meios \emph{naturaes}. Logo a Egreja é uma sociedade necessariamente \emph{perfeita}.\par Nem se diga que a Egreja, sendo composta de homens, precisa de meios \emph{materiaes} e por isso depende do Estado, em que se encontram esses meios. Porquanto os meios \emph{materiaes}, embora \emph{indispensaveis}, são meramente \emph{subsidiarios,e} não são \emph{intrinsecamente proporcionados} e \emph{operativos} em relação ao fim da Egreja, como o são os Sacramentos. Ora a \emph{perfeição} ou \emph{imperfeição} de uma sociedade não depende da posse ou da falta dos meios meramente \emph{subsidiarios}, mas sim dos meios \emph{intrinsecamente proporcionados} e \emph{operativos} em relação ao fim. --- Além d'isso, a Egreja, embora não possua \emph{formalmente} os meios \emph{materiaes}, possue-os \emph{virtualmente}; pois o Estado tem, como veremos, o dever, a que não póde faltar, de lh'os subministrar.}.

f) \emph{A} \textbf{Egreja} é sociedade suprema. --- Uma sociedade é \emph{suprema}, quando o seu fim é \emph{supremo}, isto é, superior a todos os outros fins. Ora o fim da Egreja, sendo a \emph{felicidade eterna} e \emph{sobrenatural}, \emph{é supremo}. Logo a Egreja é sociedade \emph{suprema} \footnote{O fim da Egreja não é só \emph{supremo} no seu \emph{genero}, como é o fim de toda a sociedade \emph{perfeita}, mas é \emph{supremo} em relação a \emph{todos os generos}, isto é, \emph{absolutamente}; porque a salvação eterna das almas está acima de todos os interesses temporaes. Logo a Egreja, tambem por este \emph{lado}, é superior a todas as sociedades humanas. --- Esta superioridade hão \emph{deve ser só theorica}; deve ser tambem \emph{pratica}, de modo que as outras sociedades subordinem os seus actos e interesses aos actos e interesses da Egreja.}.

199. \textbf{Constituição} da \textbf{Egreja}. --- A Egreja, emquanto é uma \emph{sociedade}, resulta de duas partes, que são a \emph{multidão dos fieis} e a \emph{legitima auctoridade}. --- A \emph{auctoridade}, na Egreja, é duplice: . de \emph{ordem} e de \emph{jurisdicção}. A primeira refere-se a oblagão do S. Sacrificio e á administração dos Sacramentos; --- a segunda, ao governo e direcção dos actos externos e internos dos fieis. --- \emph{A jurisdicção} reside, como em fonte, no Romano Pontifice, que é o supremo Pastor de toda a Egroeja, e d Elle reflue, por instituição divina, para os Bispos. --- A \emph{forma} de governo na Egreja é \emph{monarchica} \footnote{A Egreja, como toda a sociedade, é composta de dois elementos: um \emph{material}, que é a multidão dos \emph{fieis}, --- outro \emph{formal}, que é a \emph{auctoridade}. --- A \emph{auctoridade} da Egreja é de \emph{ordem} e de \emph{jurisdicção}. Porquanto a Egreja,\pgn{797}sendo Religião e Reino, deve dirigir e guiar os fieis para a felicidade eterna, a qual se consegue pela graça de \textsc{Deus}, que nos é communicada pelos Sacramentos. Logo a auctoridade da Egreja deve ser de \emph{ordem} para a administração dos Sacramentos, e deve ser de \emph{jurisdicção} para a direcção social dos feis. --- D'estas duas auctoridades, a que se reduzem todos os poderes da Egreja e que, por serem sobrenaturaes, são independentes de toda a auctoridade natural, uma --- a de \emph{ordem} --- é \emph{immovel}, a outra --- a de \emph{jurisdicção} --- é \emph{movel}. --- A \emph{jurisdicção} pôde ser \emph{interna} ou \emph{externa}, conforme se refere ao \emph{foro interno}, que é a consciencia, ou ao \emph{foro externo}, que é. a publica direcção dos fieis. Aqui fallamos da jurisdiegao \emph{externa}. --- A. jurisdicção ecclesiastica é jerarchicamente ordenada. Porquanto ella está concentrada, em toda a sua plenitude, no Romano Pontifice, e d'Este desce para os Bispos, e isto por vontade expressa do Fundador da Egreja. Os simples Sacerdotes exercem jurisdicção por delegação do Papa e dos Bispos. --- \emph{A forma} da Egreja é, pois, a \emph{monarchica}; visto que á suprema auctoridade está concen- \&rada numa só pessoa, que é o Romano Pontifice. Todavia póde dizer-se que esta monarchia é \emph{temperada} pela \emph{aristocracia} e pela \emph{democracia}: --- pela \emph{aristocracia}, emquanto os Bispos, reunidos em Concilio, exercem com o Summo Pontifice a suprema auctoridade, e são verdadeiros Principes nas suas dioceses; --- pela \emph{democracia}, emquanto \emph{todos} podem subir aos mais elevados eargos da Egreja.}.

%%>>> t280

\pg{797} \tese{280}{Poderes da Egreja.}
A \emph{jurisdicção}, sendo o direito que a Esteja tem de governar os fieis, abrange quatro poderes, que são o \emph{doutrinal}, o \emph{legislativo}, O \emph{judicial} e o \emph{executivo}.

a) \emph{A} Egreja tem o poder doutrinal. --- A Egreja é uma Religião, constituida em forma de sociedade \emph{perfeita}. Se, emquanto é \emph{Religião}, tem o \emph{officio} de ensinar; emquanto é sociedade \emph{perfeita}, tem o \emph{poder} de obrigar os fieis a prestarem assentimento ás verdades, que ensina. Logo a Egreja tem o poder \emph{doutrinal} \footnote{O \emph{ensino} é uma das principaes funcções da \emph{jurisdicção} da Egreja. --- Este poder de ensinar deriva do fim da propria Egreja. Porquanto o fim da Egreja refere-se sobretudo á verdadeira Fê, que é a forma, que constitue os \emph{fieis} e lhes dá a denominação. A Fó é o principio de salvação, o fundamento e a raiz de toda a justificação. Ora a Fé requer e suppõe o \emph{ensino}; e este ensino deve ser ministrado por quem recebeu do Divino Fundador da Egreja a respectiva missão. --- O sujeito, que recebeu a missão de ensinar a toda a Egreja, Bispos e fieis, é o Romano Pontifice, legitimo successor de S. Pedro. Os Bispos exercem o magisterio \emph{particular} nas suas dioceses, e sô exercem o magisterio \emph{universal} em toda a Egreja, quando reunidos em Concilio, em união com o Romano Pontifice e sob a sua dependencia. --- O Romano Pontifice, quando, exercendo o seu munus de Mestre universal da Egreja, define coisas pertencentes \& Fé e á Moral, é \emph{infallivel}. \emph{A infallibilidade} é um privilegio necessario para a direcção dos fieis e para a conservagüo da Egreja, que é essencialmente o Reino da \emph{verdade}.}.

t) \emph{A} \textbf{Egreja} tem o poder legislativo. --- Toda a sociedade \emph{perfeita} tem o poder \emph{legislativo}; porque este é um meio necessario \pg{798}  para dirigir a multidão para o bem commum. Ora a Egreja é sociedade \emph{perfeita}, dotada de todos os meios necessarios para o seu fim. Logo a Egreja tem o poder \emph{legislativo} \footnote{A \emph{jurisdicção} é essencialmente um poder que dirige e move a multidão ao fim. Esta direcção e este movimento suppõem as regras respectivas, que são as \emph{leis}. A Egreja, devendo dirigir os homens para o fim \emph{sobrenatural}, que é a visão \emph{intuitiva} de \textsc{DEus}, deve prescrever aos fieis as normas do movimento, --- normas, que não devem ser \emph{naturaes}, mas \emph{positivas} (porque as normas \emph{naturaes} não podem servir para a consecução de um fim \emph{sobrenatural}). É, pois, evidente que a Egreja tem o poder \emph{legislativo}.\par Nem se diga que, para essa direcção, basta a \emph{lei evangelica}; porque \& \emph{lei evangelica}, sendo muito geral, deve ser determinada e applicada aos casos especiaes, por meio de leis \emph{positivas}. --- O poder \emph{legislativo} é de ordem \emph{sobrenatural}, como a \emph{jurisdicção}, de que elle deriva ; porque, sendo \emph{sobrenatural} o fim da Egreja, devem ser \emph{sobrenaturaes} as \emph{leis}, que dirigem os homens a esse fim (pois os \emph{meios} revestem a natureza do \emph{fim}). --- O poder \emph{legislativo} reside no Romano Pontifice em relação 4 Egreja \emph{universal}, \emph{e} nos Bispos em relação ás Egrejas \emph{particulares}. Os Bispos podem fazer leis para a Egreja \emph{universal}, mas só quando estão reunidos em Concilio Ecumenico. --- A \emph{materia} do poder \emph{legislativo} da Egreja refere-se á ordem \emph{espiritual} e \emph{externa}, abrangendo o \emph{culto divino}, o \emph{uso dos Sacramentos}, os \emph{costumes} dos \emph{povos}, etc. Por isso a Egreja não regula \emph{directamente} os actos \emph{internos} dos fieis, mas só \emph{indirectamente}, emquanto, preceituando um acto \emph{externo}, preceitua, ao mesmo tempo, o acto \emph{interno}, sem o qual o \emph{externo} não tem valor. Os actos \emph{internos} são preceituados \emph{directamente} pelo poder \emph{doutrinal}.}.

c) A Egreja tem o poder judicial. --- Toda a sociedade \emph{perfeita} tem o poder \emph{judicial}, isto é, o poder de \emph{Julgar} ácerca dos \emph{direitos}, e ácerca dos \emph{crimes} e das \emph{penas}, que se lhes devem applicar; porque, aliás, o poder \emph{legislativo} seria inutil. Ora a Egreja é sociedade \emph{perfeita}. Logo a Egreja tem o poder \emph{judicial} \footnote{Não póde conceber-se um Estado, em que, além do poder \emph{legislativo}, não haja juizos que interpretem as \emph{leis} e as appliquem aos casos particulares, e, em harmonia com ellas, decidam as contendas que se possam levantar entre cidadãos, e condemnem os actos contrarios ao fim social. Se isto se verifica nas sociedades humanas, é impossivel que se não verifique na Egreja, obra immediata e positivo da Sabedoria Increada. --- O poder \emph{judicial} da Egreja extende-se tambem á parte \emph{doutrinal}; porque a Egreja é essencialmente o \emph{Reino da verdade}, e por isso é necessario que haja juizos, que illuminem os fieis, interpretando os dogmas, especulativos e praticos, comparando com elles as crenças e os costumes com o fim de julgar ácerca da honestidade ou deshonestidade das acções, resolvendo as controversias e as duvidas em materia de Fé, condemnando heresias, etc.}.

d) A \textbf{Egreja} tem o poder executivo. --- Toda a sociedade \emph{perfeita} tem o poder \emph{executivo}, isto é, o poder de levar, por meio da \emph{coacção juridica}, as vontades dos subditos ao cumprimento da lei; aliás a sociedade não alcançaria o seu fim.

Ora a Egreja é sociedade \emph{perfeita}. Logo a Egreja tem o poder \emph{executivo} \footnote{A Egreja tem o poder de servir-se da \emph{coacção juridica} para obrigar os homens á observancia da lei e reprimir as insubordinações dos impios, que perturbam a ordem publica. Sem este poder, os outros poderes seriam simplesmente inuteis e a Egreja não poderia cumprir a sua missão. --- A coacção \emph{juridica}, porém, não se deve limitar ás penas \emph{espirituaes} e ao uso da força \emph{espiritual}, mas póde e deve abranger as penas \emph{temporaes} e o uso da força \emph{material}; não só porque, aliás, não se poderiam cohibir os actos \emph{externos} do homem, mas tambem porque, consistindo geralmente o crime na rebellião da sensibilidade contra a razão, é justo que a auctoridade possa reagir contra a mesma sensibilidade, pela applicação de penas temporaes, para assim ficar restabelecido o equilibrio da ordem.\par Examinemos algumas objecções contra esta doutrina.\par \emph{a}) Dizem: A Egreja, sendo sociedade \emph{espiritual}, só deve empreger meios \emph{espirituaes}. Logo não tem o direito de servir-se da \emph{força physica}.\par Resposta. A Egreja é sociedade \emph{espiritual}, emquanto tende para um fim \emph{espiritual}, mas não emquanto é composta de \emph{espiritos}, porque é composta de \emph{homens}. Sendo composta de \emph{homens}, tem direito a servir-se dos meios proporcionados para levar o \emph{homem} ao cumprimento dos seus deveres. Ora os meios, proporcionados para levar o homem ao cumprimento dos seus deveres, devem ser \emph{materiaes}, embora se devam applicar com a intenção de alcançar um fim \emph{espiritual}; porque o homem resulta de \emph{materia} e de \emph{espirito}.\par \emph{b}) Continuam: A Egreja não é dotada da \emph{força material}. Logo não póde empregar a \emph{coacção physica}.\par Resposta. A Egreja, se não possue a \emph{força material formalmente}, possue-a \emph{virtualmente}; emquanto o Estado tem o dever de lh'a prestar todas as vezes que a Egreja a requisitar; como veremos d'aqui a pouco.\par c) Insistem: A applicação das \emph{penas temporaes} --- repugna á indole mansa e caridosa da Egreja, --- leva o homem a operar pelo medo da pena e por isso imperfeitamente, --- e só serve pára fazer hypocritas.\par Resposta. --- I) E' falso que a applicação das \emph{penas temporaes} repugne á indole mansa e caridosa da Egreja; porque a Egreja deve procurar o bem de seus filhos, e este bem alcança-se pela applicação das dictas penas, emquanto ellas corrigem o culpado e preservam os outros. --- II). A applicação das \emph{penas temporaes}, embora levasse o homem a operar pelo medo e por isso imperfeitamente, comtudo não seria nociva nem inutil, porque é melhor operar o bem por medo e imperfeitamente do que operar deshonestamente. --- III) E' falso que a applicação das \emph{penas temporaes} só sirva para fazer hypoeritas; pois a coacção, \emph{por si}, tende a tornar honestos os subditos. E, embora, \emph{accidentalmente}, a \emph{coacção} fizesse hypocritas, nem por isso deixaria de ser util para a sociedade; porque serviria para impedir escandalos, que são \textbf{peiores} e \textbf{mais} funestos do que as hypocrisias.}.

%%>>> t281

\tese{281}{Especies de direitos da Egreja.}
\emph{Direitos} da Egreja são os \emph{poderes moraes} e \emph{inviolaveis}, que Ella tem de \emph{possuir}, \emph{fazer} ou \emph{exigir} alguma coisa, concedidos pelo seu Divino Fundador. --- Estes \emph{direitos} são de duas especies: \emph{internos} e \emph{externos}. \emph{Internos} são os que a Egreja tem para com os seus \emph{subditos}, e \pg{799} \pg{800}  por isso correspondem ás suas relações \emph{internas}; --- \emph{externos} são os que a Egreja tem para com as sociedades civis, ou nações, e por isso correspondem ás suas relações \emph{externas} \footnote{Depois de termos tratado da \emph{natureza} da Egreja, e dos seus \emph{poderes}, resta-nos agora tratar dos seus \emph{direitos}. --- Estes \emph{direitos} derivam de uma lei \emph{positiva} de \textsc{Deus}. --- Aqui limitar-nos-hemos a indicar os \emph{principaes}.}.

%%>>> t282

\tese{282}{Direitos internos da Egreja.}
Os principaes direitos \emph{internos} da Egreja são: o direito \emph{territorial}, o \emph{constitutivo}, o de \emph{escolher} os seus \emph{ministros}, o de \emph{dirigir} e \emph{tutelar o ensino religioso}, o de \emph{possuir bens temporaes}, o da \emph{livre communicação}.

a) \textbf{Direito} territorial. --- O direito \emph{territorial}! consiste no poder que a Egreja tem de propagar-se e estabelecer-se em todo o mundo, para dirigir todos os homens para a vida eterna. --- Tal poder assiste á Egreja; pois esta é o Reino de N. S. Jesus Cunisro, que é, de \emph{direito}, constituido por \emph{todos} os povos, e porque foi expressamente fundada para a salvação de todos os homens \footnote{O direito \emph{territorial}, de que a Egreja é dotada em relação ao mundo inteiro, --- não é direito de \emph{propriedade}, porque ninguem se lembrou de dizer que a Egreja é \emph{proprietaria} do mundo inteiro; --- nem é direito de \emph{Jurisdicção politica}, porque N. S. Jesus \textsc{CHRISTO} não concedeu \& Egreja o dominio \emph{politico} do mundo; --- mas é direito de \emph{jurisdicção religiosa}, porque refere-se ao dominio \emph{espiritual}, que a Egreja recebeu do seu Divino Fundador sobre todos os povos. --- Foi o nosso Redemptor quem deu á Egreja este dominio; e nenhum principe ou poder humano póde oppor-se á vontade d'Aquelle, que é o \emph{Rei dos reis e o Senhor dos dominadores}. --- A Egreja, pois, não é hospede em nenhuma parte do mundo, nem é um poder estrangeiro, que exerce a sua jurisdicção sobre subditos \emph{alheios}. Em toda a parte a Egreja está em sua casa, com maior direito do que o proprio soberano temporal (porque o mundo pertence mais a \textsc{Deus} do que aos principes); e o poder, que ella exerce sobre os fieis, é um poder \emph{ordinario}, um poder que exerce Sobre subditos \emph{proprios}. --- A existencia de um Reino de \textsc{DEUS} não póde ser contrária á autonomia do Estado. Porquanto, os interesses, de que se occupa o Estado e que são os \emph{materiaes}, são essencialmente diversos dos interesses, de que se occupa a Egreja e que são os \emph{espirituaes}; e por isso o Estado, se não sahir da esphera da sua acção, não só não encontrará opposição da parte da Egreja, mas, pelo contrario, encontrará auxilio e conforto; pois a Egreja recorda continuamente aos seus filhos o preceito do Divino Mestre: \emph{Dai a Cesar o que} é de \emph{Cesar} (mas não o que Cesar se \emph{attribue}) --- D'onde se vé que a Egreja tem o direito de annunciar, pelos seus missionarios, a Fé de Jesus Cnmisro nas terras dos hereges, dos schismaticos, dos pagãos, independentemente do consentimento e mesmo contra o consentimento dos soberanos temporaes; porque o direito da Egreja é o direito do proprio \textsc{Deus}, a quem todos os homens, embora soberanos, devem obedecer.}.

b) Direito constitutivo, --- Direito \emph{constitutivo} é o poder que a Egreja tem --- de regular o culto divino, interno e externo, \pg{801}  sobretudo a administração dos Sacramentos e a celebração do S. Sacrificio, --- de estabelecer a \emph{jerarchia} para o regimen dos fieis, --- de fundar, approvar e regular as \emph{ordens religiosas}, ete. --- Este direito pertence á Egreja; pois é necessario para a consecução do seu fim, que é a santificação e salvação eterna das almas \footnote{Este direito \emph{constitutivo}, que attribuimos á Egreja, --- não se refere á sua constituição \emph{primitiva} e \emph{fundamental}, nem aos elementos, que constituem a sua essencia, porque este direito é exclusiva e incommunicavelmente proprio de N. S. Jesus \emph{CHRISTO}; --- mas refere-se á constituição \emph{secundaria} e \emph{explicativa}. Com effeito, a Egreja é Religião e é sociedade. --- Emquanto é Religião, tem o direito de regular o culto divino, a administração dos Sacramentos e do S. Sacrificio; de levantar templos, determinar a liturgia, estabelecer os dias santos, e tudo, em fim, que se refere á gloria de \textsc{Deus} e á salvação das almas. --- Emquanto é sociedade, a Egreja tem o direito de instituir a \emph{jerarchia}, prescrever as normas para o regimen dos fieis, fundar as \emph{ordens religiosas}, determinando a sua organisacaão social, etc.}.

c) Direito de escolher os seus ministros. --- Este direito pertence á Hgreja; porque, se lhe não pertencesse, Ella não poderia conservar-se, nem operar para a consecução do seu fim. --- O direito de escolher os Bispos reside no Romano Pontifice, e o direito de escolher os Parochos reside nos Bispos \footnote{O Romano Pontifice tem jurisdicção \emph{immediata} e \emph{ordinaria} sobre a Egreja universal, e por isso pertence a Elle a nomeação dos Bispos, que tenham jurisdicção sobre uma determinada diocese; assim como pertence ao Rei, que tem jurisdicg\&o sobre todo o Reino, designar e nomear os governadores dos districtos. Pela mesma razão, os Bispos tem o direito de nomear os Parochos nas diversas freguezias das suas dioceses. --- A intervenção do povo ou dos soberanos, admittida nalgum tempo e nalguns logares, n\&o \emph{escolhia} os pastores, mas dava apenas um testemunho da honestidade e da capacidade dos candidatos. Hoje essa mesma intervenção seria motivo de perpetuas luctas, de damno para as almas e de escandalo para a Egreja.}.

d) Direito de dirigir e tutelar o ensino religioso. --- Este direito consiste no poder que a Egreja tem de ministrar o ensino religioso a todos os homens, e de velar para que nada se diga e propague contra a verdadeira doutrina de N. S. Jesus CunisTo. --- Tal direito pertence á Egreja; pois esta recebeu to seu Fundador a missão de ensinar a todos os homens \footnote{Aqui não fallamos do direito, que a Egreja tem de prégar livremente o Evangelho, mas do direito que Ella tem sobre o ensino em materia de Religião, que se ministra aos fieis. O ensino póde ser o que se ministra aos elerigos, e o que se ministra aos leigos. Vejamos os direitos da Egreja sobre um e outro ensino.\par a) Em relação ao ensino dos \emph{clerigos}, o direito da Egreja é evidente. A. formação dos ministros pertence á sociedade, a que devem servir. Além d'isso, a sciencia \emph{theologica} é propriedade da Egreja; porque se funda na\par 51 : . Vor. II.\pgn{802}divina Revelação, e a depositaria da Revelação é a Egreja. --- \emph{E}, pois, absurdo o estabelecimento de cadeiras de Theologia, com inteira independencia dos Bispos, sujeitas sómente ao ministerio da instruccüo publica.\par \emph{b}) Em relação ao ensino dos \emph{leigos}, o direito da Egreja não é menos evidente. O \emph{leigo} não deve ser theologo, mas deve ser \emph{c/Arístüo}. Para ser \emph{christão}, elle deve conhecer as coisas principaes relativas á nossa Santa Religião. Mas o ensino das coisas relativas á nossa Santa Religião pertence á Egreja; porque Ella, e só Ella, recebeu a missão de ensinar a toda a gente. Este ensino é dado pelo Episcopado em união com o Summo Pontifice. --- O direito ácerca do \emph{ensino} abrange tambem, embora \emph{indirectamente}, o direito de julgar das outras sciencias, condemnando as proposições, que são contrarias á doutrina Catholica e que por isso são necessariamente falsas.}.

e) Direito de possuir bens temporaes. --- Este direito pertence á Egreja. Porquanto, a Egreja, sendo sociedade religiosa e visivel, deve sustentar os seus ministros, levantar templos, comprar vasos e paramentos, etc.; e o cumprimento d'esses deveres importa grandes despezas, que não se podem fazer, se a Egreja não possuir bens temporaes. Ora quem tem o dever de alcançar um fim, tem direito aos meios indispensaveis para esse fim \footnote{A Egreja tem o direito de \emph{possuir} bens \emph{temporaes}, e este direito deriva de N. S, Jesus Cumirsro. Por isso os bens da Egreja chamam-se \emph{sagrados}, e quem os rouba commette um \emph{sacrilegio}. --- A Egreja, tendo o direito de possuir bens \emph{materiaes}, tem tambem o direito de os administrar. * Pergunta-se: O que é mais conveniente para o clero --- possuir bens materiaes, ou receber dos governos o estipendio?\par Resposta. Os governos modernos têm geralmente confiseado os bens da Egreja, e limitam-se a dar ao Clero um estipendio relativamente insignificante. Esta medida é \emph{injusta}, \emph{degradante} e \emph{perigosa}. --- E' \emph{injusta}, porque priva de seus bens a legitima proprietaria, que é a Egreja. --- E' \emph{degradante}, porque torna os ministros do altar empregados civis e escravos dos governos. --- E' \emph{perigosa}, porque uma revolução, uma crise financeira, ou a maldade de um governo póde fazer com que se suspenda o estipendio, com damno dos interessados e sobre tudo da Egreja e das almas.}.

f) \textbf{Direito} de livre communicação. --- O direito de \emph{livre communicação} consiste em que os Bispos podem communicar livremente com os fieis da sua diocese, \emph{e} o Papa com os Bispos e com os fieis de todo o mundo. --- Este direito existe-nos Bispos e no Summo Pontifice; porque, sendo verdadeiros pastores, o seu poder sobre os subditos é \emph{immediato} e \emph{ordinario} \footnote{A proposição é tão evidente, que, para ser acatada, basta o senso commum. --- O Bispo é \emph{mestre}, \emph{é pastor}, \emph{é pae}, é \emph{principe} do povo de \textsc{Deus}. Ora quem poderá dizer que o mestre não tem o direito de communicar livremente com os seus discipulos, o pastor com as suas ovelhas, o pae com os seus filhos, o principe com os seus subditos? --- O Summo Pontifice é o Bispo dos Bispos, é o Pae e o Mestre de todos os fieis, é o Pastor supremo das ovelhas e dos cordeiros. Ora quem lhe negará o direito de communicar livremente com os Bispos e com os fieis? --- Este direito do S. Padre é\pgn{803}\emph{divino}, porque o dever, de que deriva, é divino, e por isso não póde ser impedido nem limitado por nenhuma lei humana. --- O S8. Padre communica com os Bispos e com os fieis do mundo por meio dos \emph{Nuncios} ou \emph{Legados Apostolicos}; que, representando o Summo Pontifice, têm o direito de communicar livremente com os Bispos, Parochos e com os fieis do paiz, onde está acreditado, exercendo sobre elles a propria auctoridade do Papa. --- D'onde se vê que os que tolhem: ou diffieultam a communicação do Papa com os Bispos e com os fieis exercem um despotismo injusto e contrario á legitima liberdade. Tratam o Papa como se Elle fosse uma potencia estrangeira, e não querem convencer-se de que o pae nunca é estranho para os seus filhos, nem o pastor para as suas ovelhas! Têm medo de que o Papa invada as abtribuições dos governos, e não se lembram de que a Egreja, que dá o reino celeste, não ambiciona os thronos dos principes da terra!}.

%%>>> t283

\pg{802} \pg{803} \tese{283}{Direitos externos da Egreja.}
Os principaes direitos \emph{externos} da Egreja são: o direito á \emph{independencia}, \emph{á superioridade}, \emph{á immunidade}, \emph{á protecção} e á \emph{assistencia} do Estado.

a) Direito á independencia. --- Este direito pertence á Egreja . em relação á sociedade civil ou Estado. Porquanto, uma sociedade é \emph{independente} de outra, quando o fim d'aquella é \emph{independente} do fim d'esta; pois as sociedades estão entre si na razão \emph{directa} dos seus fins. Ora o fim da Egreja, sendo supremo, sobrenatural e divino, é independente do fim da sociedade

\emph{z} civil, que é secundario, natural e humano \footnote{A \emph{independencia} da Egreja é um corollario da these, em que demonstrámos que a Egreja é sociedade \emph{perfeita}. Até que se não prove que a Religião é apenas um negocio humano e politico, e que o espirito deve ser governado pelo corpo, é, e será sempre verdade, que a Egreja não póde depender do Estado. --- Esta \emph{independencia} é absolutamente necessaria e indispensayel para a Egreja poder cumprir a grande missão, que N. S. JESUS Cumisro lhe confiou. --- Para a defesa e segurança da independencia da Eigreja, é necessario, nas actuaes circumstancias, que o Romano Pontifice não esteja sujeito a nenhum poder, interno ou externo; porque, aliás, poderia ser contrariado no excercicio publico do seu ministerio, ou, pelo menos, nasceriam duvidas ácerca da plena liberdade dos seus actos, com grande perturbação das consciencias. Ora esta independencia de todo o poder, interno e externo, só póde obter-se, se o Summo Pontifice tiver o \emph{dominio temporal}, de modo que seja verdadeiro Principe e não subdito. --- Este \emph{dominio temporal} tinha sido concedido por N. S. Jesus \textsc{CHRISTO} ao seu Vigario mas foi usurpado pela maçonaria triumphante, sob o pretexto da unidade \emph{da Italia}, mas com o verdadeiro fim de tornar difficil e impossivel\par ieio do ministerio ecclesiastico o de destruir, se fosse destructivel, a\par de Jesus Cumisro.\par = \emph{MAN}}.

b) Direito á sugerioridade. --- Direito á \emph{superioridade} consiste em que o Estado, no regimen social dos povos, deve estar \emph{subordinado} á Egreja. Porquanto, o Estado, tendo o dever de estar subordinado a \textsc{Deus}, tem tambem o dever de estar subordinado \pg{804}  á Egreja, que representa o proprio \textsc{Deus}, e que, no seu Augusto Nome, dita leis aos povos e aos principes. --- Esta subordinagão do Estado á Egreja não é \emph{directa}, mas \emph{indirecta} \footnote{A coisa é evidente. O Estado, no seu governo, não póde eximir-se da obediencia, que deve a \textsc{DEus} e á sua santa lei. Ora a missão de indicar a vontade de \textsc{DEus} e de proclamar a sua lei foi confiada á Egreja. Logo o Estado deve obedecer ou estar subordinado á Egreja. --- Além d'isso, as sociedades estão entre si na razão directa dos seus fins --- \emph{Societates sunt} ut \emph{fines} earum; de modo que uma sociedade é superior ou inferior a outra, conforme o fim d'aquella é superior ou inferior ao fim d'esta. Ora o fim do Estado, que é a felicidade \emph{temporal}, deve estar subordinado ao fim da Egreja, que é a felicidade \emph{eterna}; porque a felicidade \emph{temporal} deve servir para a \emph{eterna}, ou pelo menos, não deve contrarial-a. --- Finalmente, tres opiniões podem conceber-se ácerca das relações entre a Egreja e o Estado. Uma, que sujeita a Egreja ao Estado; --- outra, que declara ambas as sociedades independentes entre si e parallelas; --- a terceira, que ensina a subordinação do Estado á Egreja. A primeira é absurda; porque sujeita o reino de \textsc{DEus} ao reino do homem, o bem da alma ao bem do corpo. A segunda é tambem insustentavel; porque suppõe dois fins \emph{ultimos} (o que é contradiecüo; pois, se o homem é um só, um só deve ser o seu fim \emph{ultimo}, a que devem estar subordinados todos os outros fins) A terceira é admittida por todos os que ainda n\&o renunciaram ao senso commum.\par Dos principios expostos descendem dois corollarios: --- T) Em caso de conflicto entre a Egreja e o Estado, o fiel deve obedecer áquella e não a este; porque o direito do superior prevalece ao do inferior, o direito de \textsc{Deus} ao do homem. --- II) Quando o Estado faz uma lei contrária a \textsc{Deus} ou á Egreja, essa lei é nulla; porque, em tal caso, \textsc{DEus} não dá \& lei a força necessaria para obrigar os subditos. --- A Egreja, tendo recebido do proprio \textsc{Deus} a missão de conservar e de defender o deposito da Fê e da moral, tem o direito de declarar tal nullidade.\par Todavia esta subordinação do Estado á Egreja não é \emph{directa}, mas é simplesmente \emph{indirecta}. --- Não é \emph{directa}. Na verdade, a subordinação é \emph{directa}, quando a sociedade inferior, nos seus actos relativos ao fim, está por tal modo sujeita á superior, que esta póde irrital-os ou tornal-os validos, conforme lhe agradar. Ora a Egreja, tendo direito só sobre as coisas, que se referem de algum modo ao seu fim, não póde ter nem tem direito sobre as coisas \emph{temporaes}, que de nenhum modo se referem ao seu fim; e por isso, em relação a estas, o Estado, sendo sociedade \emph{perfeita}, tem o direito de operar livremente sem dependencia da Egreja. Logo a subordinação do Estado á Egreja não é \emph{directa}. --- Mas é \emph{indirecta}. Porquanto, a subordinação é \emph{indirecta}, quando a sociedade inferior está sujeita á superior naquelles seus actos, que se prendem e relacionam com o fim da mesma sociedade superior. Ora o Estado está sujeito á Egreja naquelles seus actos, que se prendem e relacionam com o fim da mesma Egreja; porquanto a felicidade\pgn{805}\emph{temporal}, que é o fim do Estado, estando subordinada á felicidade eterna, que é o fim de Egreja, deve procurar-se de modo, que não só não impeça, mas até favoreça, por uma forma \emph{positiva}, a consecução da mesma felicidade \emph{eterna}. Logo a subordinação do Estado á Egreja é \emph{indirecta}.}.

c) Direito á immunidade. --- O direito á \emph{immunidade} consiste em que os ministros da Egreja e as coisas sagradas estão isentas da jurisdicção civil. --- Este direito pertence á Egreja. Porquanto, a Egreja é livre e independente do Estado. Ora: não \pg{805}  póde chamar-se livre e independente uma sociedade, cujos ministros e propriedades estão sujeitas a um poder estranho \footnote{A \emph{immunidade ecclesiastica} é de direito \emph{divino}; pois que foi o proprio Redemptor quem concedeu á sua Egreja esta isenção do poder civil.\par A. \emph{immunidade} divide-se em \emph{pessoal} e \emph{real}, conforme se refere á \emph{pessoa} dos ministros, ou ás \emph{coisas}, que são os \emph{templos}, os \emph{bens materiaes}, etc.\par a) A \emph{immunidade pessoal} significa que os ministros da Egreja, por especial privilegio de \textsc{Deus}, não estão sujeitos ao poder civil, não só nas coisas \emph{espirituaes} (em que os proprios leigos estão isentos do poder civil) mas tambem mas coisas \emph{temporaes}. Por isso os Olerigos não podem ser obrigados aos offioios, que repugnam \& dignidade e ao caracter clerical (como é a \emph{milicia}, e, embora estejam sujeitos a observar as leis civis, relativas ao bem publico, comtudo, se transgredirem essas leis, não podem ser julgados nem punidos pelo poder \emph{civil}, mas só pelo poder \emph{ecclesiastico}. --- A razão d'esta \emph{immunidade} é que, estando os leigos para com os Sacerdotes, como os filhos estão para com o pae, os discipulos para com o mestre, as ovelhas para com o pastor, é justo que o pae não seja julgado e punido pelos filhos, nem o mestre pelos discipulos, nem o pastor pelas ovelhas. --- A \emph{immunidade ecclesiastica} é total e absoluta no Romano Pontifice, que é o pastor supremo de todos, e não é julgado por ninguem, pois é superior a todos os tribunaes, civis e ecclesiasticos. --- Este privilegio foi reconhecido pelo primeiros Imperadores christãos --- \emph{Constantino}, \emph{Theodosio}, \emph{Valentiniano}, \emph{Arcadio}, e por todos os soberanos que se servem da auctoridade, que \textsc{Deus} lhes concedeu, para o bem dos povos.\par \emph{b}) \emph{A iémmunidade real} abrange os \emph{templos} e os \emph{bens ecclesiasticos}. O Estado não tem direito algum sobre os \emph{templos}, que estão estreitamente ligados com a Religião, nem póde lançar impostos sobre os bens \emph{ecclesiasticos}. Todos os povos, ainda que idolatras, estiveram sempre convencidos de que os bens destinados ao decoro dos templos e á sustentação dos ministros da Egrej\& devem ser isentos dos impostos communs. Assim Pharaó, no Egypto, isentou os Sacerdotes dà obrigação de vender as suas propriedades ao Estado; --- Artaxerxes, Rei da Persia, ordenou, depois da reedificação do templo de Jerusalem, que os Sacerdotes e os outros Ministros fossem isentos de toda a contribuição; --- Julio Cesar conta que os \emph{Druidos}, na Gallia, não pagavam impostos, nem estavam sujeitos ao serviço militar... Assim decretaram tambem os verdadeiros Principes Christãos. --- Nem esta isenção dos impostos, concedida ao Clero, era prejudicial ao Estado; porquanto o Clero, além de distribuir uma grande parte dos seus rendimentos entre os pobres, auxiliaya, com generosas offertas, o Estado nas suas necessidades,}.

d) Direito á protecção e é assistencia do Estado. --- Este direito pertence á Egreja. Porquanto, o Estado é \emph{indirectamente} subordinado á Egreja. Ora esta subordinação \emph{indirecta} exige que o Estado não só não faça coisa alguma em prejuizo da Egreja, mas tambem Lhe dispense toda a \emph{protecção} e a \emph{assistencia} \pg{806} , que são necessarias ou uteis para Ella poder conseguir o seu fim altissimo \footnote{Este direito da Egreja em relação ao Estado nasce do outro direito de \emph{superioridade}, que compete á Egreja. O fim \emph{subordinado} não só não deve contrariar a consecução do fim superior (subordinação \emph{negativa}), --- mas deve facilital-a pelo emprego de meios \emph{positivos} (subordinação \emph{positiva}).\par O dever, que o Estado tem de dispensar a sua \emph{protecção} e \emph{assistencia} á Egreja, deriva tambem dos deveres, que elle tem para com \textsc{DEus}, para com os seus \emph{subditos}, e para \emph{comsigo}. --- Deriva dos deveres, quo o Estado tem para com \textsc{Deus}. Porquanto, o Rei é ministro de \textsc{Deus}; pois \textsc{DEus} communicou-Ihe a sua suprema auctoridade, para que, em seu Augusto Nome, governasse os povos; ora o ministro deve fazer e promover os interesses do seu senhor, e os interesses de \textsc{Deus} são os que estão confiados aos cuidados da Egreja, que é o Reino de \textsc{Deus}. --- Deriva dos deveres, que o Estado tem para com os seus \emph{subditos}. Na verdade, o Estado, devendo proteger os direitos e favorecer os interesses dos subditos, deve sobretudo proteger e favorecer a Religião, que é o supremo de todos os direitos e de todos os deveres,e por isso a Egreja. --- Deriva dos deveres, que o Estado tem para \emph{comsigo}. Porquanto, o Estado deve procurar a sua prosperidade. A prosperidade de um Estado consiste sobretudo na justiça dos seus empregados, na obediencia. dos subditos, na paz e harmonia entre todas as classes, Ora, pertencendo á Egreja a missão de levar os homens ao amor e á pratica da justiça, da obediencia e da paz, segue-se que o Estado, favorecendo e protegendo a Egreja, favorece e promove a propria felicidade.\par Do que deixamos dito vê-se o absurdo da \emph{separação} da Egreja do Estado, significada pela celebre formula --- \emph{Egreja livre no Estado livre}. --- Não póde negar-se que, em alguns casos, esta separação é uma necessidade, como nos Estados, que perseguem a Egreja (como a \emph{Russia}), Mas são excepções. Considerando as coisas em si mesmas, nos designios de \textsc{Deus}, dizemos que o Estado não deve separar-se da Egreja. Porquanto, não podem separar-se duas sociedades, quando o fim de uma não póde separar-se do fim da outra. Ora o fim do Estado não póde separar-se do fim da Egreja; pois a felicidade \emph{temporal} não póde separar-se da felicidade \emph{eterna}; visto que a vida presente é um \emph{meio} em relação \& vida futura, e o meio não póde separar-se do fim. --- Além d'isso, esta separação seria causa de continuas luctas e de grandes inquietações para as consciencias dos subditos. Porquanto, sendo uma e a mesma pessoa que está sujeita 4 Egreja e ao Estado, segue-se que, se o Estado mandar uma coisa contraria é Egreja, a condição do \emph{cidadão} estará em conflieto com a condição do \emph{crente}, E, então, --- ou a pessoa obedece ao Estado, e fica mal com a sua consciencia e com \textsc{Deus}, --- ou obedece \& Egreja e fica mal com o Estado, que fará uso da força material para a levar á sujeição. --- Finalmente,.a separação da Egreja do Estado, significando um governo sem \textsc{Deus} e sem Religião, é funesta para o proprio soberano e para o Estado. E' funesta para o soberano; porque, não reconhecendo-.a Deus,\pgn{807}cuja vontade dá força ás leis, o soberano perde o titulo mais forte que póde ter ao respeito e á obediencia dos povos. E' funesta para o Estado; porque a Religião é necessaria, mesmo para o interesse material do povo. \emph{Macchiavelli} diz: « Onde existe a Religião, ha todo o bem; onde ella falta, ha todo o mal. E' impossivel que o soberano seja reverenciado por quem despreza a \textsc{Deus}. Os Romanos foram grandes emquanto tiveram o temor de \textsc{DEus} ».\par A verdadeira doutrina, pois, é a seguinte: A sociedade religiosa e a civil, devendo, cada uma na sua ordem ou esphera, procurar o bem do homem e a gloria de \textsc{Deus}, devem auxiliar-se mutuamente. --- A Egreja deve auxiliar o Estado, formando os bons costumes dos cidadãos, inculcando a reverencia aos poderes constituidos e a obediencia á legitima auctoridade, promovendo as virtudes moraes, que são a base de toda a ordem. --- O Estado tem para com a Egreja deveres \emph{negativos} e \emph{positivos}, que derivam da sua subordinação \emph{indirecta} para com Ella e que serão d'aqui a pouco indicados.\par Para a harmonia de ambas as auctoridades, \emph{ecclesiastica} e \emph{civil}, a Egreja e o Estado costumam celebrar tratados, a que se dá o nome de \emph{concordatas}, e pelos quaes a Egreja concede ao Estado algumas faculdades ou privilegios sobre coisas relativas á disciplina ecclesiastica. --- Nas \emph{concordatas}, a Egreja faz alguma \emph{concessão} (p. e. que o Estado apresente os Bispos), ou renuncia a qualquer direito seu (p. e. á immunidade ecclesiastica); e por isso dá o que não tem o dever de dar nem pertence ao direito do Principe. Pelo contrario, quando o Estado se compromette solemnemente a defender a liberdade e os direitos da Egreja, não dá senão o que \emph{devia} dar, e ao seu dever, imposto pela lei divina, \emph{accrescenta} a obrigação de um pacto solemne, de modo que se torna devido tambem pelo pacto o que antes era devido pelo direito divino. Por isso as \emph{concordatas} differem dos pactos internacionaes, estipulados entre sociedades juridicamente eguaes, pelos quaes tanto uma como a outra parte contrahente dá o seu e recebe o \emph{alheio}; porque a Egreja, se concede ou renuncia aos \emph{seus} direitos, não recebe do Estado nada que jà não lhe seja devido. --- As \emph{concordatas} obrigam ambas as partes contrahentes. Todavia os escriptores catholicos não estão de accordo ácerca da natureza. dessa obrigação. Alguns dizem que as \emph{concordatas} têm a força de contracto \emph{bilateral} e por isso obriga por \emph{justiça} não só o Estado, mas tambem a Egreja. Outros affirmam que as \emph{concordatas} obrigam por \emph{justiça} o Estado, que tem o dever de obedecer á Egreja, e só por uma certa \emph{decencia} ou \emph{fidelidade} obriga a propria Egreja. --- A \emph{interpretação} das \emph{concordatas} pertence só á Egreja; porque só a Egreja recebeu de \textsc{Deus} o \emph{supremo} poder \emph{legislativo}, e por isso o \emph{interpretativo}, sobre as coisas \emph{espirituaes} e \emph{mistas}, que constituem a materia ordinaria d'estes tratados. --- A Egreja pode \emph{abrogar} as \emph{concordatas} em todos os casos, nos quaes é permittido rescindir um contracto: --- quando o Estado falta ás suas promessas, --- quando a \emph{materia} se muda por tal modo, que se torna impossivel ou illicita, --- quando o exige o \emph{bem publico} da Egreja (é condição implicita nestes pactos), --- quando houve \emph{erro substancial}, ou \emph{engano}, ou as concessões foram extorquidas pela \emph{força} ou pelo \emph{medo}, porque, embora o acto, extorquido pela \emph{força}\pgn{808}ou pelo \emph{medo}, não seja irrito por si, comtudo o Papa, que não pode recorrer a um juiz superior, pode desligar-se, por si mesmo, do um vinculo injustamente estreitado.\par Os direitos da Egreja, que foram até aqui enumerados, dizem respeito propriamente aos Estados Catholicos, que, vivendo na Egreja, estão sujeitos á sua espiritual auctoridade. --- Em relação aos Estados \emph{infieis}, a Egreja tem o \emph{direito} --- de prégar o Evangelho, sem que nenhum principe possa ' oppor-se á missão recebida do proprio \textsc{Deus} (embora Ella não possa constranger ninguem a abraçar a Fé), --- e de exercer os seus poderes de \emph{ordem} e de \emph{jurisdicção}. --- Em relação aos Estados \emph{heterodoxos} (\emph{protestantes} ou \emph{schismaticos}), \emph{a} Egreja tem o direito de conservar e constituir a sua jerarchia, de educar o Clero, governar os fieis, possuir e administrar os bens, etc. (Cf. P. Liberatore, \emph{Del diritto pubblico ecclesiastico}, cap. 1V, art. VI).}.

%%>>> t284

\tese{284}{Deveres da sociedade civil para com a Egreja.}
Os \emph{deveres}, que incumbem á sociedade civil, ou Estado, para com a Egreja e que se fundam na superioridade dº Esta em relação áquelle, são de duas especies: \emph{negativos}. e \emph{positivos}.

a) \textbf{Deveres} negativos. --- Os deveres \emph{negativos} consistem em que a sociedade civil deve respeitar todos os \emph{direitos} da Egreja, \emph{internos} e \emph{externos}, nem deve fazer coisa alguma, que seja contrária á \emph{existencia} e á \emph{acção} da mesma Egreja.

b) \textbf{Deveres} positivos, --- Os deveres \emph{positivos} consistem em que a sociedade civil deve auxiliar e favorecer a Egreja, quando for necessario e Ella o pedir. --- Todos estes deveres, \emph{negativos} e \emph{positivos}, são \emph{jurídicos}, porque correspondem aos \emph{direitos}, de que a Egreja foi dotada por Nosso Senhor Jesus \textsc{CarISTO} \footnote{Não nos demoramos- na descripção dos \emph{deveres} da sociedade civil para com a Egreja; pois a sua \emph{natureza} e \emph{extensão} foram explicadas sufficientemente, quando tratámos dos correspondentes \emph{direitos} da Egreja. --- Só diremos que os soberanos terrenos, se imitassem, na sua obediencia para com a Egreja, a \emph{Carlos Magno}, que se assignava: \emph{Karolus Dei gratia rex}, \emph{devotus} sanctae \emph{Ecclesiae defensor}, \emph{atque adjutor in omnibus Apostolicae Sedis}, alcancariam, como elle, gloria nas conquistas e tornariam prospera a sua nação, convencendo-se, mais uma vez, da grande verdade, proclamada por \emph{Jontesquieu} (\emph{De esprit des lois}, l. XXIV): « Coisa admiravel! A Religião Christa, que parece não ter outro fim senão a bemaventurança da vida futura, faz tambem a nossa felicidade na presente ». Soberanos e subditos veriam despontar uma era de paz e de prosperidade, desde o momento em que, unidos numa só fé e num só affecto, reconhecessem, na theoria e na pratica, 2 soberania benefica e gloriosa de Nosso Senhor \textsc{JESUS CHRISTO},}.

DOS TRATADOS CONTIDOS NESTE VOLUME

\textbf{SYNOPSE DA ANTHROPOLOGIA}

\textsc{ANTHROPOLOGIA}. --- \emph{Antropologia} é a sciencia que trata das causas su-

H premas \emph{do homem}. O \emph{homem} é o \emph{animal racional}. --- O \emph{objecto}, pois, da \emph{Anthropologia} é o \emph{homem}, considerado nas suas causas supremas, que são a sua \emph{essencia} (de que derivam as \emph{faculdades} e as \emph{operações}), á sua \emph{origem} e o seu \emph{destino}. --- Por isso á \emph{Anthropologia} é distincta da \emph{Physiologia} e das outras sciencias, que se occupam das causas \emph{proximas} do organismo humano.

Divisão DA ANTrHROPOLOGIA. --- A \emph{Anthropologia} divide-se em \emph{quatro} secções. A 1.º secção occupa-se da \emph{essencia} do homem; --- a \emph{segunda}, das suas \emph{faculdades}; --- s terceira, das suas \emph{operações}; --- a quarta, da sua \emph{origem} e do seu \emph{destino}.

\part*{SECÇÃO I\\[0.3em]\large ESSENCIA DO HOMEM}
\addcontentsline{toc}{part}{SECÇÃO I --- ESSENCIA DO HOMEM}

\pg{807} \pg{808} \pg{809} Divisio DA SECÇÃO PRIMEIRA. --- A secção I divide-se em \emph{tres} capitulos: --- no 1.º trata-se do \emph{corpo humano}; --- no 2.º, da \emph{alma humana}; --- no 8.º, da \emph{união} da alma com o corpo.

A. CoRPO HUMANO.

a) \emph{Descripção do corpo humano}. --- As principaes regiões do corpo humano são: --- a \emph{cabeça}, que se divide em craneo e \emph{face}; --- o \emph{thoraz}, que tem tres cavidades secundarias, uma \emph{media}, que aloja o coração, etc., e duas \emph{lateraes}, contendo os \emph{pulmões}; --- o \emph{abdomen}, que encerra os principaes orgãos do apparelho \emph{digestivo}. --- O \emph{thorax} e o \emph{abdomen} constituem o tronco, a que se unem, como appendices, os membros \emph{thoracicos} e \emph{abdominaes}.

b) \emph{Orgãos da vida vegetativa}. --- Os orgãos da vida \emph{vegetativa}, que servem para a \emph{nutrição}, constituem tres \emph{apparelhos}: o \emph{digestivo}, \emph{o circulatorio}, e o \emph{respiratorio}. i

a) Os orgãos, que constituem o apparelho \emph{digestivo}, são: a \emph{bocca}, a \emph{pharinge}, o \emph{esophago}, o \emph{estomago}, o \emph{intestino delgado}, o \emph{intestino grosso}.

\emph{b}) Os orgãos do apparelho \emph{circulatorio} são: o córagüo, as \emph{arterias}, os \emph{vasos capillares} e as \emph{veias}.

\emph{c}) Os orgãos do apparelho \emph{respiratorio} são; a \emph{larynge}, a \emph{trachéa}, os \emph{bronchios}, os \emph{pulmões} e o \emph{thorac}.

\pg{810} €) \emph{Orgãos da sensibilidade externa}. --- Os orgãos da \emph{sensibilidade externa} são os que formam o apparelho dos \emph{cinco} sentidos \emph{externos}.

a) Os orgãos do apparelho \emph{visual} dividem-se em \emph{essenciaes} e \emph{accessorios}. --- Os \emph{essenciaes} são: o \emph{globo ocular}, que é formado por \emph{membranas} e \emph{meios transparentes}, \emph{e} o \emph{nervo optico}. --- Os \emph{accessorios} são: as \emph{orbitas}, as \emph{palpebras} e os \emph{musculos oculares}.

\emph{b}) O apparelho \emph{auditivo} divide-se em tres partes, que são: o \emph{ouvido externo}, composto do \emph{pavilhão} e \emph{canal auditivo externo}, --- o \emph{ouvido medio}, formado pela \emph{caixa do tympano}, --- o \emph{ouvido interno}, composto de tres cavidades, que são o \emph{vestibulo}, os \emph{canaes semi-circulares} e o \emph{caracol}.

c) O apparelho \emph{o/factivo} resume-se numa membrana mucosa, chamada \emph{pituitaria}, que guarnece internamente as \emph{fossas nasaes}, e que recebe o nervo \emph{olfactico}. :

\emph{d}) O orgão do apparelho \emph{gustativo} é principalmente a \emph{lingua}. O nervo \emph{gustativo} parece ser o \emph{glosso-pharingeo}.

€) A sede principal do \emph{tacto} existe na superficie palmar das extremidades dos dedos, na parte anterior da bocca e na ponta da lingua.

á) \emph{Orgãos da sensibilidade interna e da locomoção}.

4) Os orgãos da sensibilidade \emph{interna} são: --- o \emph{encephalo}, dividido em tres secções, que são o \emph{cerebro}, o \emph{cerebello} e o \emph{bolbo rachidiano}, --- a \emph{medulla espinal}, --- e os nervos.

5) Os orgãos da \emph{locomoção} são os \emph{ossos} (\emph{longos}, \emph{chatos} ou \emph{curtos}) e os \emph{musculos} (\emph{lisos} ou \emph{estreados}).

B. ALMA HUMANA,

a) Alma \emph{humana}, \emph{sua existencia e substancialidade}. --- A \emph{alma humana é o principio primeiro pelo qual o homem vegeta}, \emph{sente} e \emph{entende}. --- E' um ente \emph{real}, porque as nossas operações sao \emph{reaes}, e todo o effeito \emph{real} exige uma causa \emph{real}. --- E' \emph{substancia}, porque o \emph{accidente} sô póde ser principio \emph{secundario} de operações; mas é \emph{incompleta} na \emph{especie}, porque é \emph{parte} do homem. --- E' \emph{distincta do corpo}; porque o corpo está sujeito a perpetuas mudanças nos elementos, o que não se verifica na alma, e porque, aliás, todos os corpos seriam vivos, pois a alma é principio de \emph{vida}.

b) A \emph{alma humana é simples}. --- \emph{Simples} é o ente que não é composto de \emph{partes}, nem \emph{essenciíaes} (como o é o homem), nem. \emph{integrantes} (como o é toda a substancia extensa) Ora a alma humana não é composta de partes \emph{essenciaes}, nem de partes \emph{integrantes}. Não é composta de partes \emph{essenciaes}; aliás o homem não perceberia, como effectivamente percebe, os objectos materiaes com um acto \emph{unico} e \emph{completo}, mas percebel-os-hia com tantos actos incompletos, quantas fossem as partes da alma, o que é falso. Nào é composta de partes \emph{integrantes}; porque não é composta de partes \emph{essenciaes}, porque não é \emph{corpo}. Logo a alma humana é \emph{simples};

0) A \emph{alma humana é espiritual}. --- Um ente é \emph{espiritual}, quando é \emph{intrinsecamente} independente da materia nas \emph{operações} e por isso na \emph{existencia}. Ora a alma humana é \emph{intrinsecamente} independente da materia nas suas \emph{operações}. Porquanto, a nossa alma percebe objectos \emph{immateriaes} ou \emph{espiritwdes}, como são \textsc{Deus}, á \emph{virtude}, a \emph{Justiça}, etc. Mas se a alma fosse dependente da. materia, não poderia perceber semelhantes objectos; porque um principio material ou dependente da materia só póde perceber objectos materiaes e extensos. --- A \emph{espiritualidade} da alma demonstra-se tambem pelas operações da sua \emph{vontade}, a qual tende para objectos \emph{espirituaes} ou \emph{immateriaes}, e todo o ser tende para o bem que ihe é \emph{semelhante}: --- O \emph{materia}:

\pg{811} \emph{lismo}, que, para negar a \emph{espiritualidade} da nossa alma, affirma que o pensamento é um \emph{producto}, uma \emph{secreção}, um \emph{movimento} ou uma \emph{funcção} do cerebro, é um systema condemnado não só pela razão, mas tambem pelas mais recentes descobertas da \emph{anatomia} e da \emph{physiologia}.

0. União DA ALMA COM O CORPO E SUAS CONSEQUENCIAS.

a) \emph{União da alma com o corpo}. --- A alma está unida ao corpo, como a \emph{forma substancial á materia}. Com effeito, uma realidade é forma \emph{substancial} de um ser, quando o colloca numa determinada especie, distinguindo-o de todos os seres das outras especies. Ora a alma colloca o nosso corpo na especie de \emph{Humano}, distinguindo-o dos outros corpos. Logo a alma está unida ao corpo, como á forma \emph{substancial} á materia.

São falsos, pois, os systemas de Platão, de Mallebranche (causas \emph{occasionaes}, de Leibnitz (\emph{harmonia preestabelecida}), de Leclerc (\emph{mediador plastico}), de Locke (\emph{influxo physico}), e de Rosmini; porque estes systemas admittem apenas uma união \emph{accidental} entre a alma e o corpo, quando esta união é \emph{substancial}, porque d'ella resulta uma só \emph{pessoa} e uma sô natureza.

b) \emph{Unidade da alma humana}. --- A alma, por ser forma \emph{substancial} do corpo, é uma só em cada homem. Esta unica alma é a \emph{intellectual}, que é tambem principio da vida \emph{vegetativa} e \emph{sensitiva} no homem; pois a forma superior contem a força das inferiores. --- Além da alma \emph{intellectual}, não existe outra forma \emph{substancial} no homem.

c) \emph{Sede da alma humana}. --- A alma humana, sendo forma \emph{substancial} do corpo, está toda em todo o corpo e toda em cada parte d'elle, quanto á \emph{essencia}, mas não quanto á \emph{energia}.

\part*{SECÇÃO II\\[0.3em]\large FACULDADES DO HOMEM}
\addcontentsline{toc}{part}{SECÇÃO II --- FACULDADES DO HOMEM}

Divisão DA SECÇÃO SEGUNDA. --- A secção II divide-se em \emph{tres} capitulos: --- no 1.º trata-se das faculdades humanas em \emph{geral}; --- no 2.º, das faculdades \emph{sensitivas} do homem; --- no 8.º, das suas faculdades \emph{intellectuaes}.

A. AS FACULDADES HUMANAS EM GERAL.

a) As \emph{faculdades em si mesmas}. --- \emph{Faculdade} é o principio \emph{proximo} e \emph{immediato} das operações. --- O principio \emph{remoto} e \emph{mediato} das operações é a \emph{essencia} do homem, --- As faculdades conhecem-se e distinguem-se pelas Suas \emph{operações}, como as operações pelos seus \emph{objectos}. --- As \emph{faculdades} do homem são muitas e diversas, e reduzem-se a \emph{tres} classes supremas, que são a \emph{vegetativa}, \emph{á sensitiva} e a \emph{intellectual}.

b) As \emph{faculdades em relação á alma}. --- A alma é o \emph{principio} de que derivam \emph{todas} as faculdades, mas só é \emph{sujeito} das faculdades \emph{intellectuaes}. --- Distinguem-se \emph{realmente} da \emph{essencia} da alma, como os ramos do tronco.

0) As \emph{faculdades em relução ás operações}. --- As faculdades distinguem-se realmente das suas operações, como os ramos dos fructos.

B. FACULDADES SENSITIVAS \emph{DO} HOMEM.

a) \emph{Sentidos externos}. --- Sentidos \emph{externos} são os que existem na superfivie \emph{externa} do organismo humano, e por cujo intermedio o homem percebe os seres \emph{materiaes} e \emph{externos}. --- São \emph{cinco}, a saber: \emph{vista}, \emph{ouvido}, \emph{olfacto}, \emph{gosto} e \emph{tacto}. Não podem ser mais do que estes enumerados.

b) \emph{Sentidos internos}. --- Sentidos \emph{internos} são os que existem no interior do organismo, e-por cujo intermedio o homem percebe: --- \emph{directa} e imméz \pg{812}  \emph{diatamente} as proprias \emph{sensações} e as \emph{imagens} dos objectos \emph{externos}, e \emph{indirecta} e \emph{mediatamente} os mesmos objectos externos. --- São \emph{quatro}, a saber: o \emph{sentido commum}, pelo qual percebemos as sensações dos sentidos \emph{externos} e as suas \emph{differenças}, --- á \emph{imaginação}, pela qual conservamos e reproduzimos as \emph{imagens} dos objectos, mesmo na \emph{ausencia} d'elles, --- a \emph{estimativa}, pela qual distinguimos o que é util do que é nocivo, --- a \emph{memoria sensitiva}, pela qual percebemos as \emph{proprias} percepções sensitivas \emph{passadas}.

o) \emph{Appetite sensitivo}. --- Appetite \emph{sensitivo} é a faculdade, pela qual o homem tende para o bem \emph{sensivel}. --- Divide-se em .\emph{concupiscivel} e \emph{irascivel}, conforme tende para o bem e foge do mal sensivel, ou tende para um bem \emph{arduo} mas \emph{util} para remover os obstaculos, que se oppõem á consecução do bem ou \& fuga do mal. São duas faculdades distinctas.

d) \emph{Faculdade locomotiva}. --- Faculdade \emph{locomotiva} é aquella pela qual o homem se transporta de um logar para outro. --- No homem \emph{a} faculdade \emph{locomotiva} está sujeita ao imperio da vontade.

e) \emph{Faculdade vocal}, --- Faculdade \emph{vocal} é aquella pela qual o homem exprime e communica os seus \emph{conhecimentos} e \emph{paixões} por meio da vos.

O. FACULDADES INTELLECTUAES DO HOMEM.

a) \emph{Intelligencia}. --- \emph{Intelligencia} é a faculdade pela qual o homem percebe a \emph{essencia} das coisas. --- O seu objecto divide-se em \emph{adequado e proporcionado}. \emph{Adequado} é a \emph{essencia} das coisas, considerada no seu elemento mais fundamental, que é o \emph{ente}. \emph{Proporcionado} é a \emph{essencia} das coisas \emph{materiaes}, \emph{abstrahida das condições materiaes}. --- E' uma faculdade \emph{espiritual}, porque attinge o \emph{espiritual} ou \emph{immaterial}. --- Embora \emph{espiritual}, depende das faculdades \emph{sensitivas}, que lhe subministram o \emph{objecto} da percepção. --- E' faculdade \emph{potencial}, que deve ser determinada ao acto pelo objecto. E como o objecto só pode actuar na intelligencia por meio de uma sua imagem \emph{immaterial}, é necessario admittir uma força \emph{abstractiva}, que elabore aquella imagem. Esta força é a \emph{iéntelligencia agente}. --- A intelligencia diz-se \emph{razão}, emquanto o homem deduz dos principios as consequencias, --- \emph{consciencia}, em quanto reflecte sobre si, --- \emph{memoria intellectual}, emquanto recorda as \emph{proprias} operações \emph{intellectuaes passadas}, como \emph{proprias} e como \emph{passadas}.

b) Vontade. --- \emph{Vontade} é a faculdade pela qual o homem tende para o bem \emph{Immaterial}. --- O bem, para o qual a vontade tende, pode ser \emph{universal} e \emph{particular}. --- E' uma faculdade \emph{espiritual}, Não é livre na sua tendencia para o bem \emph{universal}, mas é livre na sua tendencia para o bem \emph{particular}. --- O \emph{fatalismo} e o \emph{determinismo} são absurdos. Ha factos \emph{physiologicos} e \emph{pathologicos}, que ás vezes tiram toda a \emph{voluntariedade} e outras vezes apenas a diminuem,

\part*{SECÇÃO III\\[0.3em]\large OPERAÇÕES DO HOMEM}
\addcontentsline{toc}{part}{SECÇÃO III --- OPERAÇÕES DO HOMEM}

Divisão DA SECÇÃO TERCEIRA. --- A secção III divide-se em dois capitulos: --- no 1.º trata-se das operações \emph{sensitivas} do homem; --- no 2.º, das \emph{intellectuaes}.

A. OPERAÇÕES SENSITIVAS DO HOMEM.

a) \emph{Percepção sensitiva}. --- Percepção \emph{sensitiva} é uma operação pela qual o homem conhece \emph{directamente} os entes \emph{materiaes} na sua \emph{singularidade}. --- E' uma operação do \emph{composto}, e verdadeiro \emph{conhecimento}. --- Divide-se em \emph{externa} e \emph{interna}.

b) \emph{Percepção sensitiva externa}. --- E' o acto proprio dos sentidos \emph{externos}, pelo qual o homem conhece, de um modo \emph{material} ou \emph{concreto}, os entes \emph{materiaes} \pg{813}  e \emph{externos}. --- Exige a união do objecto com a faculdade. O objecto deve unir-se á faculdade por uma sua imagem. A faculdade, determinada pela imagem, percebe o objecto externo, como é em si mesmo. A percepção dos objectos externos é \emph{directa} ou \emph{immediata}. \emph{A séde} d'esta percepção é o \emph{orgão animado externo}.

\emph{o}) \emph{Percepção sensitiva interna}. --- E' o acto proprio dos sentidos \emph{internos}, pelo qual o homem conhece, de um modo \emph{material} ou \emph{concreto}, seus actos e estados \emph{sensitivos internos}. --- Os sentidos \emph{internos} precisam de uma imagem para o seu acto. Essa imagem é produzida pela sensação \emph{externa}. Os sentidos \emph{internos}, informados pela imagem, percebem o objecto. O \emph{cerebro} é a \emph{séde} da percepção sensitiva \emph{interna}.

d) \emph{Appetite sensitivo}. --- Appetite sensitivo é a operação pela qual o homem tende para o bem \emph{sensivel}. --- Divide-se em \emph{concupiscivel} e \emph{irascivel}. --- As operações \emph{sensitivas}, elevadas a um grau intenso de força, dizem-se \emph{paixões}. --- Estas são \emph{onse}, a saber: \emph{amor} e \emph{odio}, \emph{desejo} e \emph{abominação}, \emph{alegria} e \emph{tristeza}, \emph{esperança} e \emph{desespero}, \emph{audacia} e \emph{medo}, e \emph{ira}. As primeiras \emph{seis} derivam do appetite \emph{concupiscivel}, as ultimas \emph{cinco} do \emph{érascivel}.

B. OPERAQÓES INTELLECTUAES DO HOMEM.

a) \emph{Conhecimento intellectual}, \emph{sua origem}. --- Conhecimento \emph{intellectual} é a operação pela qual a intelligencia percebe a \emph{essencia} das coisas. --- Os nossos conhecimentos não são \emph{énaatos}, não derivam exclusivamente da \emph{tradição}, nem são attingidos na visão intuitiva de Deus, não têm por principio só os \emph{sentidos} nem só a \emph{intelligencia}, mas procedem da \emph{intelligencia}, auxiliada pelos \emph{sentidos}.

b) \emph{Modo por que se produz o conhecimento intellectual}. --- Os sentidos \emph{externos} apresentam o \emph{objecto} do conhecimento \emph{intellectual}. A \emph{imaginação} forma a especie sensivel do objecto. A intelligencia \emph{agente} torna a especie sensivel capaz de operar na intelligencia \emph{possivel}. Esta, determinada pela especie, produz o \emph{verbo mental}, que exprime a \emph{essencia} das coisas.

c) \emph{Modo por que se desenvolve o conhecimento intellectual}. --- O desenvolvimento refere-se aos \emph{objectos}, que conhecemos, e ao \emph{modo}, por que os conhecemos. --- A nossa intelligencia conhece --- a \emph{essencia} das substancias \emph{materiaes} por um modo \emph{directo}, --- as substancias \emph{materiaes} na sua \emph{singularidade} por um modo \emph{indirecto}, --- á \emph{existencia} e a \emph{essencia} da alma humana por um modo \emph{indirecto}, --- \textsc{Deus} tambem por um modo \emph{indirecto}.

\emph{a}) \emph{Modo por que se exprime o conhecimento intellectual}. --- E' a \emph{linguagem}. A relação entre a linguagem e as ideas é \emph{convencional}. A linguagem suppõe as ideas; não é \emph{necessaria}, mas é \emph{utilissima}. O homem não inventou, mas podia inventar a linguagem.

e) \emph{Volicües}. --- São os actos, pelos quaes a vontade tende para o bem \emph{immaterial}, percebido pela \emph{intelligencia}. O seu \emph{objecto} é o mesmo que o da vontade: o \emph{bem}, que é \emph{universal} ou \emph{particular}.

\part*{SECÇÃO IV\\[0.3em]\large ORIGEM E DESTINO DO HOMEM}
\addcontentsline{toc}{part}{SECÇÃO IV --- ORIGEM E DESTINO DO HOMEM}

Divisão DA SECÇÃO \textsc{QUARTA}. --- A. secção IV divide-se em dois capitulos: --- no 1.º trata-se da \emph{origem} do homem; --- no 2.9, do seu \emph{destino}.

\emph{A}. ORIGEM DO HOMEM.

á) \emph{Origem do corpo humano}. --- O corpo humano não derivou do concurso de causas fortuitas, --- nem da evolução de uma especie inferior, --- \pg{814}  mas, na sua origem primitiva, foi formado immediatamente por \textsc{Deus}, e, no estado actual, propaga-se por geração.

b) \emph{Origem da alma humana}. --- A alma humana não é produzida pela evolução de uma especie inferior, nem por geração, nem por emanação da substancia divina, --- mas é creada \emph{immediatamente} por \textsc{DEUS}, no momento da sua união com o corpo.

c) \emph{Origem do composto humano}. --- O composto humano não derivou do concurso de causas fortuitas, nem da evolução de uma especie inferior, --- mas, na sua origem primitiva, foi formado \emph{immediatamente} por \textsc{DEUS}, e, no estado actual, deriva de outro homem e de \textsc{Deus}, mas por modo diver:o. --- Todos os homens pertencem á mesma especie, e as divergencias são accidentaes. Por isso todos poderam derivar da mesma familia ou par, porque a especie é fecunda até ao indefinido. --- O homem constitue um reino especial.

B. DESTINO DO HOMEM.

a) \emph{Fim do homem}. --- O homem foi creado para um fim, que é a felicidade. Esta encontra-se no bem \emph{universal}; por isso não nos bens finitos, mas no bem infinito, que é \textsc{Deus}. A felicidade, que consiste no conhecimento e no amor de \textsc{Deus}, não é perfeita na vida presente; só o será na futura.

b) \emph{Immortalidade da alma humana}. --- A alma humana é immortal, tanto \emph{intrinsecamente}, isto é, considerada em 8i, como \emph{extrinsecamente}, isto é, considerada em relação ao seu Creador. --- E' immortal \emph{intrinsecamente}; porque não pode acabar --- nem pela dissolução dos seus elementos, visto ser \emph{simples}, --- nem por destruição do organismo, visto ser \emph{espiritual}. --- E' immortal \emph{evtrinsecamente}; porque, na ordem actual das coisas, a destruição da alma repugna á Sabedoria, á Bondade, á Justiça, \& Santidade de \textsc{Deus}. . . 0) \emph{Operações da alma separada}. --- A alma, separada do corpo, é sempre \emph{individua}, conserva um desejo innato de se reunir ao \emph{seu} corpo, produz algumas operações, --- não as \emph{organicas}, --- mas as \emph{intellectuaes}. --- Conhece, por meio das especies intelligiveis, as causas \emph{immateriaes} do mundo corporeo, os objectos \emph{materiaes} na sua \emph{singularidade}, e as substancias \emph{espirituaes}, que são a sua essencia, os espiritos, \textsc{DEus}. E' immutavel na vontade e é dotada de livre arbitrio. --- Creada por \textsc{Deus}, a nossa alma volta para Dus.

---------------

\textbf{SYNOPSE DA THEODICEA}

TuEopicEA. --- \emph{Theodicea} é a parte da Metaphysica, que trata de \textsc{Deus}, \emph{emquanto é conhecido pela lus da razão}. --- D'ahi a \emph{ewcellencia} da "'heodicea e a sua \emph{utilidade} para o \emph{individuo} e para a \emph{sociedade}.

Divisão pA \textsc{TuxopicEA}. --- A Theodicea divide-se em \emph{quatro} capitulos: --- no l.? trata-se da \emph{existencia} de \textsc{DEUS}; --- no 2.9, da sua \emph{essencia} e dos seus \emph{attributos}; no 9.9, das suas \emph{faculdades} e \emph{operações immanentes}; --- no 4.0, das suas \emph{operações transeuntes}.

I. --- EXISTENCIA DE DEUS

á) DEMONSTIRABILIDADE DA EXISTENCIA DE \textsc{DEUS}. --- A \emph{existencia} de \textsc{DEus} \emph{pode} e \emph{deve} ser demonstrada. \emph{Pode}; porque a razão deduz da existencia dos \pg{815}  effeitos a existencia da Causa proporcionada, e os seres finitos são effeitos, que exigem uma Causa infinita. \emph{Deve}; porque a \emph{existencia} de \textsc{Deus} não é immediatamente evidente para nós. --- A \emph{existencia} de \textsc{Deus} não se demonstra --- nem \emph{a-priori} (isto é, pela \emph{causa}), porque \textsc{DEUS} não depende de nenhuma causa, --- nem \emph{a-simultaneo} (isto é, pela sua \emph{idea}), porque da \emph{idea} de um ente não pode deduzir-se a sua existencia \emph{real}; --- mas demonstra-se \emph{a-posteriori} (\emph{isto} é, pelos \emph{e/feitos}).

b) DEMONSTRAÇÃO DA EXISTENCIA DE \textsc{DEUS}. --- A \emph{existencia} de \textsc{Deus} demonstra-se por argumentos \emph{metuphysicos}, \emph{physicos} e moraes.

a) \emph{Argumentos metaphysicos}. --- Os argumentos \emph{metaphysicos} são \emph{cinco}, e baseam-se em que, --- 1.º) se no mundo existe o movimento ou a mudança, deve existir um motor immovel, pois, se tudo fosse movido, teriamos um effeito sem causa; --- 2.º) se existem as causas \emph{segundas}, deve existir a causa \emph{primeira}, porque as causas \emph{segundas} são effeitos; --- 8.9) se existem seres \emph{contingentes}, deve existir um ser \emph{necessario}, porque o \emph{contingente} não tem em si a razão da sua existencia; --- 4.9) se existem varios e diversos graus de perfeição, deve existir um ser infinitamente perfeito, porque o \emph{relativo} suppõe o \emph{absoluto}; --- 5.9) se existem verdades \emph{absolutas} e \emph{necessarias}, deve existir tambem uma intelligencia \emph{absoluta} e \emph{necessaria}, porque nada existe sem \emph{razão sufficiente}.

\emph{D}) \emph{Argumentos physicos}. --- Os argumentos \emph{p/iysicos} baseam-se em que, se existe a \emph{ordem} no mundo e, em particular, nas operações dos agentes \emph{irracionaes}, deve existir um ordenador infinitamente sabio e poderoso; porque a \emph{ordem} é um \emph{effeito}, que deve attribuir-se a uma causa, e esta causa não pode deixar de ser uma sabedoria e um poder infinito. \textbar{}

€) \emph{Argumentos moraes}. --- Os argumentos \emph{moraes} fundam-se em que, --- 1.º) se existe uma lei moral, deve existir um legislador supremo, pois não ha \emph{lei} sem \emph{legislador}; --- 2.9) se \emph{todos} os povos admittiram \emph{sempre} a existencia de \textsc{Deus}, \textsc{Deus} existe, porque tal consenso só pode provir da verdade.

0) A EXISTENCIA DE \textsc{DEUS} E O ATHEISMO. --- \emph{Atheismo} é o erro que nega a \emph{existencia} de \textsc{Deus}. --- E' \emph{pratico} ou \emph{theoretico}, conforme nega essa existencia com os \emph{factos} ou com a \emph{razão}. O \emph{theoretico é negativo} ou \emph{positivo}, conforme se nega a existencia de Deus, ou porque não o conhecem, ou porque estão convencidos de que Elle não existe. --- O atheismo \emph{pratico} é possivel; o \emph{theoretico-negativo} não é possivel nos adultos; o \emph{theoretico-positivo} é impossivel. --- O \emph{atheismo} é a destruição da sciencia, da sociedade, e da propria natureza humana,

IL --- ESSENCIA E ATTRIBUTOS DE DEUS

\&) EssENciA DE \textsc{DEus}, SEU CONSTITUTIVO E ATTRIBUTOS. --- A \emph{Essencia} de \textsc{Deus} é o que faz com que Elle seja \textsc{Deus}, e não outra coisa. --- E' \emph{physica} ou \emph{metaphysica}. \emph{A physica} é constituida por todas as perfeições divinas e por isso consiste na \emph{infinidade}. A. \emph{metaphysica} é a perfeição, que se \emph{concebe} como a \emph{vais} de todas as perfeições divinas, e por isso consiste na \emph{subsistencia} do \emph{proprio ser}. --- Os \emph{attributos} são \emph{quiescentes} ou \emph{operativos}. Os \emph{quiescentes} dividem-se em \emph{negativos} e \emph{positivos}, conforme \emph{negam} de \textsc{Deus} as imperfeições, ou \emph{affirmam} as perfeições creadas.

b) AttRIBUTOS DE \textsc{Deus} NEGATIVOS. --- São a \emph{simplicidade}, a \emph{immutabilidade}, a. \emph{eternidade}, a \emph{immensidade} e á \emph{unidade}. --- \textsc{Deus} é \emph{simples}; porque o ente \emph{composto é finito} e \emph{dependente}, e \textsc{Deus} \emph{infinito} e \emph{independente}. ---

;

é \textsc{Deus} é immnutavel; porque todo o ente mutavel é composto de \emph{potencia}. e de \pg{816}  \emph{acto}, \emph{e} \textsc{Deus} é \emph{acto purissimo}. Esta immutabilidade é \emph{physica} e \emph{moral}. --- \textsc{Deus} é \emph{eterno}; porque não tem \emph{principio}, nem \emph{fim}, nem está sujeito á \emph{successão}. --- \textsc{Deus} é \emph{immenso}, porque, sendo o supremo Creador e Conservador, está presente \emph{actualmente} em todas as coisas \emph{existentes} e \emph{virtualmente} em todas as coisas \emph{possiveis}. --- \textsc{Deus} é \emph{uno}; porque a Divindade, sendo subsistente em si e por si, não pode ser recebida em varios sujeitos, e por isso não pode multiplicar-se e é necessariamente \emph{unica}.

c) ATTRIBUTOS DE \textsc{DEUS} PosrrIVOS. --- São a \emph{verdade}, a \emph{bondade} e a \emph{belleza}. --- \textsc{Deus} não só é \emph{verdadeiro}, mas é a propria \emph{verdade}, porque a sua \emph{intelligencia} é uma e á mesma coisa com o seu \emph{ser}. --- \textsc{Deus} não só é \emph{bom}, porque é \emph{perfeito}, mas é a propria \emph{bondade}, porque é a sua propria \emph{perfeição}. --- \textsc{Deus} é \emph{bello}, \emph{e} é \& propria \emph{belleza}, porque as suas infinitas perfeições constituem a unica e simplicissima Essencia divina. A belleza divina é a causa \emph{efficiente}, \emph{exemplar} e \emph{final} de toda a belleza creada.

IIL --- FACULDADES E OPERACÓES IMMANENTES DE DEUS

4) INtELLIGENCIA DE \textsc{DEUS}, SUA NATUREZA E OBJECTO, --- A \emph{intelligencia} de Dus identifica-se com a sua essencia, e é acto purissimo e infinito. \textsc{Deus} conhece, pela sua essencia, o proprio ser, e todas as outras coisas, existentes e possiveis. --- Esta sciencia é \emph{causa} das coisas creadas, emquanto está unida, para o effeito, com o acto da vontade, e é \emph{substancial}, \emph{infinita}, \emph{intuitiva}, immutavel. O seu termo \emph{immanente} é o Verbo.

b) FACULDADE E OPERAÇÃO VOLITIVA IMMANENTE. --- A \emph{vontade} de \textsc{Deus} identifica-se com a sua essencia, e é acto purissimo, e infinitamente perfeita. \textsc{Deus} ama, pela sua bondade, o \emph{proprio ser} como \emph{fim} e por isso \emph{necessariamente} e as \emph{outras coisas} como \emph{meios} e por isso \emph{livremente}. O amor divino é causa das coisas creadas e é \emph{substancial}, \emph{immutavel}, \emph{efficacissimo}, \emph{benevolo}. O seu termo immanente é o \emph{Espirito}.

€) VIDA DIVINA, SUA EXISTENCIA E NATUREZA. --- \textsc{Deus} não só é dotado de \emph{vida}, porque opera de um modo \emph{immanente}, mas é a propria \emph{vida}, porque \emph{vita viventibus est esse} e \textsc{DEus} é o proprio ser. A vida de \textsc{Deus} é \emph{perfeitissima} e \emph{sempiterna}. Todas as coisas são vida em \textsc{Deus}.

IV. --- OPERACÓES TRANSEUNTES DE DEUS

a) CREAÇÃO, SUA NATUREZA, PRINCIPIO E FIM. --- \emph{Creação} é o acto, pelo qual

, \textsc{Deus} tirou do \emph{nada} todas as coisas. E' só proprio de \textsc{DEus}. \textsc{Deus} foi livre na creação e creou tudo para a sua gloria.

b) CONSERVAÇÃO, SUA DIVISÃO E NECESSIDADE. --- \emph{Conservação} é o acto da divina vontade, pelo qual as coisas, uma vez creadas, perseveram na existencia. \textsc{Deus} conserva todas as coisas de um modo \emph{positivo}, emquanto lhes continua o seu influxo, e as conserva pela continuação do acto creador, e n\&o anniquila coisa alguma creada.

o) CONCURSO, SUA NATUREZA E EXTENSÃO. --- \emph{Concurso} é o acto, pelo qual \textsc{Deus}, como Causa primeira, move ou auxilia as forças \emph{physicas} das causas \emph{segundas}, e assim influe em todos os actos naturaes d'essas causas. --- \textsc{DEus} concorre \emph{immediatamente} --- em todas as operações das creaturas, --- e em todas as operações da vontade humana. --- O concurso de \textsc{Deus} é tambem \emph{previo} aos actos das creaturas, e deixa intacta a liberdade da nossa vontade.

d) PROYIDENCIA, SUA EXISTENCIA, EXTENSÃO E PROPRIEDADES. --- \emph{Providencia}

\pg{817} al \textsc{Deus} \emph{ordena} e \emph{dirige} todas as coisas para o fim. --- Esta é o acto, pelo qu

Extende-se a cada uma das coisas; tem um providencia existe em \textsc{Deus}.

\emph{riversal}, \emph{suave}, \emph{infallivel,-immediata}. Mas cuidado especial do homem; e é \emph{ur}

ção do homem; d'ahi a necessidade da oração. não exclue a coopera

---------------------

\textbf{SYNOPSE DA MORAL}

Monar. --- \emph{Moral} é a sciencia que, por meio das regras de \emph{moralidade}, dirige os actos \emph{humanos} ao nosso \emph{ultimo fim}. --- O \emph{objecto}, pois, da \emph{Moral} é constituido pelos actos \emph{humanos}, emquanto, pelas regras da \emph{moralidade}, se hão de dirigir para o nosso \emph{ultimo fim}. :

Divisão va Momar. --- A \emph{Moral} divide-se em duas partes: \emph{geral} e \emph{especial}. --- \emph{A geral} trata dos \emph{actos humanos}, considerando-os no seu \emph{fim}, \emph{essencia}, \emph{moralidade} e \emph{regras}. \emph{A especial} occupa-se dos \emph{direitos} e dos \emph{deveres}, que são proprios do homem, emquanto é \emph{individuo} e \emph{membro da sociedade}. --- A Moral \emph{geral} chama-se \emph{Ethica}; \& \emph{especial}, \emph{Direito natural}.

ETHICA

\textsc{EtHicA}. --- \emph{Ethica} é a parte da Philosophia \emph{moral}, que trata dos \emph{actos humanos}, seu \emph{fim}, \emph{essencia}, \emph{moralidade} e \emph{regras}.

Divisão pa EtHica. --- \emph{A Ethica} divide-se em \emph{quatro} capitulos: --- no 1.* trata-se do \emph{fim} dos actos \emph{Humanos}; --- no 2. da \emph{essencia} dos mesmos actos; --- no 8.º, da sua \emph{moralidade}; --- no 4.º, das suas \emph{regras}.

\emph{A}, FIM DOS ACTOS HUMANOS.

a) \emph{Existencia do fim ultimo dos actos humanos}. --- Os actos humanos (que derivam da \emph{intelligencia} e da \emph{vontade} do homem) tendem sempre para um \emph{fim}, e tambem para um fim \emph{ultimo}, isto é, que satisfaça completamente as aspirações da nossa alma. --- Este fim \emph{ultimo}, que é o mesmo para todos os homens, existe realmente. :

b) \emph{Natureza do fim ultimo dos actos humanos}. --- O fim \emph{ultimo} dos actos humanos é o bem \emph{universal}, pois só este pode satisfazer completamente as aspirações da alma. --- Este bem \emph{universal} não é um bem \emph{creado}, pois este é particular e finito, mas é o bem \emph{increado}, é \textsc{Deus}. --- O fim \emph{ultimo}, que é \textsc{Deus}, consegue-se pela \emph{inteligencia} e pela \emph{vontade}, mas a sua consecução \emph{perfeita} esta reservada para a \emph{vida futura}.

o) \emph{Fim proximo dos actos humanos}. --- O fim \emph{procimo} dos actos \emph{humanos}, que se deve alcançar na vida presente, consiste na disposição d'elles para o fim \emph{ultimo}, e por isso no \emph{conhecimento} e no amor de \textsc{DEUS}. --- A consecugüo do fim \emph{proximo} produz uma felicidade \emph{imperfeita}.

B. ESSENCIA DOS ACTOS HUMANOS.

a) \emph{Elementos e divisão do} acto \emph{humano}. --- O acto \emph{humano} exige \& \emph{percepção} da \emph{intelligencia}, o \emph{consentimento da vontade} e a \emph{liberdade}. E' \emph{elicito} ou imperado, \emph{directo} ou \emph{indirecto}, \emph{positivo} ou \emph{negativo}, \emph{actual} ou \emph{virtual}. 3 52 Vor. II.

\pg{818} b) \emph{Diversas especies do acto elicito}. --- As especies do acto \emph{elicito} são \emph{seis}, a saber: \emph{amor}, \emph{intenção}, \emph{goso}, \emph{consentimento}, \emph{eleição}, \emph{uso}.

\emph{o}) \emph{Impedimentos da voluntariedade dos actos humanos}. --- Os impedimentos da \emph{voluntariedade} dos actos \emph{humanos} costumam reduzir-se a \emph{quatro}, \& saber: \emph{ignorancia}, \emph{concupiscencia}, \emph{medo} e \emph{violencia}.

O. MORALIDADE DOS ACTOS HUMANOS.

a) \emph{Moralidade}, \emph{sua existencia e fundamento}. --- \emph{Moralidade} é a propriedade pela qual os actos /\emph{umanos} são dotados ou destituidos de \emph{honestidade}, de modo que possam, ou não, levar á consecução do fim \emph{ultimo}. --- A \emph{moralidade} é \emph{intrinseca} ou \emph{extrinseca}, conforme convem ao acto em virtude da sua natureza, ou só em consequencia de um preceito ou de uma prohibição do superior. --- A moralidade \emph{intrinseca} existe; pois alguns actos são \emph{honestos} e outros são \emph{deshonestos}, independentemente de toda a vontade humana. Esta moralidade \emph{intrinseca} funda-se --- \emph{proximamente} na conveniencia ou não-conveniencia do acto humano com a nossa natureza racional, --- e \emph{remotamente} na conveniencia ou não-conveniencia do mesmo acto com a Natureza Divina. --- Os systemas, que se referem á \emph{moralidade} e que são contrarios á verdade, são muitos e diversos. Uns negam a moralidade \emph{intrinseca}, isto é, a differenca \emph{absoluta} e \emph{objectiva} entre o \emph{bem} e o \emph{mal moral}, e são: o \emph{materialismo} dos \emph{epicureus}, --- o \emph{utilitarismo} dos \emph{empiricos}, dos \emph{economistas} e dos \emph{positivistas}, --- o \emph{autonomismo} dos \emph{racionalistas}, --- e o \emph{voluntarismo} dos \emph{Juristas}. Outros, embora admittam a moralidade \emph{intrinseca}, todavia basóam o seu fundamento \emph{ultimo} --- ou na \emph{lei da evolução indefinida}, --- ou na \emph{lei da natureza} ou da \emph{razão}, --- ou na \emph{lei do respeito devido á pessoa} ou \emph{d especie humana}, --- ou no \emph{imperativo categorico}, --- ou na \emph{idea da justiça}, --- ou na \emph{idea transcendental do bem}, --- ou na \emph{idea das relações essenciaes das coisas}.

b) \emph{Fontes da moralidade dos actos humanos}. --- As \emph{fontes}, que communicam aos actos humanos uma certa \emph{especie} de bondade ou de malicia, são tres, a saber: o \emph{objecto}, o \emph{fim}, e as \emph{circumstancias}.

o) \emph{Consequencias da moralidade dos actos humanos}. --- As principaes consequencias, que resultam da \emph{moralidade} dos actos humanos, são a \emph{imputabilidade}, o \emph{merito} ou o \emph{demerito}, o \emph{premio} ou o \emph{castigo}.

d) A \emph{moralidade dos actos humanas e os habitos}. --- \emph{Habito} é a qualidade estavel ou permanente, que inclina a faculdade a operar mais de um modo do que de outro, bem ou mal, conforme o acto convem ou repugna \& natureza \emph{racional} do homem. --- Divide-se em \emph{bom}, que é a \emph{virtude}, e em \emph{mau}, que é o \emph{vicio}. --- A \emph{virtude é intellectual} ou \emph{moral}, conforme aperfeiçoa a \emph{intelligencia} ou a \emph{vontade}.

D. REGRAS DOS ACTOS HUMANOS. -

a) \emph{Lei moral e suas especies}. --- Lei \emph{moral} é um preceito racional para o bem commum, promulgada por quem governa a communidade. Divide-se em \emph{eterna} e \emph{temporal}, \emph{natural} e \emph{positiva}, \emph{divina} e \emph{humana}. --- A lei \emph{eterna}, que é um acto da razão e da vontade de \textsc{DEus}, emquanto, desde a eternidade, prescreve ás creaturas racionaes a conservação da ordem moral e prohibe a sua perturbação, existe, e manifesta-se ás creaturas pela luz natural da razão. --- A lei \emph{eterna}, emquanto assim se manifesta, constitue a lei \emph{natural}, --- A lei \emph{natural} existe, como o attestam a razão, a consciencia o a experiencia. E' \emph{universal} e \emph{immutavel}, e deve ter e tem uma sanec\&o \emph{sufficiente} na vida \emph{futura}; e esta sancção consiste na posse ou perda eterna da felicidade. --- A lei \emph{positiva} é um preceito racional para o bem commum, \pg{819}  promulgado \emph{livremente} por quem governa a communidade. --- A lei \emph{positiva} é necessaria, funda-se na lei \emph{natural}, e deve ser \emph{honesta}, \emph{justa} e \emph{possivel}.

b) \emph{Consciencia moral}, \emph{sua divisão e direcção}. --- Consciencia \emph{moral} é um juizo da razão \emph{pratica}, que, applicando os principios da lei moral a um certo e determinado caso, dicta se este é licito ou illicito. --- Divide-se em --- \emph{antecedente} e \emph{consequente}, --- \emph{recta} e \emph{erronea}, --- \emph{certa}, \emph{duvidosa} e \emph{provavel}, --- \emph{delicada}, \emph{laca} e \emph{escrupulosa}. --- O homem --- 1.º) não pode desobedecer á consciencia \emph{certa}, quando \emph{recta} ou \emph{énvencivelmente erronea}; --- 2.º) não pode desobedecer, nem obedecer \& consciencia \emph{vencivelmente erronea}; --- 8.º) não pode obedecer á consciencia \emph{duvidosa}; --- 4.º) pode obedecer á consciencia \emph{provavel}.

\textbf{DIREITO} NATURAL

\textsc{DIREITO} NATURAL. --- Direito \emph{natural} é a parte da Philosophia \emph{moral}, que trata dos \emph{deveres} e dos \emph{direitos}, que pertencem ou podem pertencer ao homem em virtude da lei \emph{natural}. --- \emph{Dever} é a necessidade moral de praticar ou omittir alguns actos, imposta ao homem pela lei. Divide-se em natural e \emph{positivo}, \emph{primitivo} e \emph{derivado}, \emph{juridico} e \emph{moral}, \emph{affirmativo} e \emph{negativo}, \emph{dispensavel} e \emph{indispensavel}; e, immediata ou mediatamente, deriva de Drs. --- \emph{Direito} é o poder moral e inviolavel de possuir, fazer ou exigir alguma coisa, concedido pela lei. E' \emph{natural} ou \emph{positivo}, \emph{innato} ou \emph{adquirido}, \emph{juridico} ou \emph{moral}, \emph{pessoal} ou \emph{real}, \emph{alienavel} ou \emph{inalienavel}. Tem tres propriedades: a coacção, \& \emph{collisüo} e a \emph{limitação}, --- e, immediata ou mediatamente, deriva de \textsc{Deus}. A ordem \emph{juridica} existe e é uma parte da ordem \emph{moral}.

Divisão no \textsc{DIREITO} NATURAL. --- O Direito \emph{natural} divide-se em duas secções. A 1.º abrange o Direito \emph{individual}; \& 2.*, o Direito \emph{social}.

\part*{SECÇÃO I\\[0.3em]\large DIREITO INDIVIDUAL}
\addcontentsline{toc}{part}{SECÇÃO I --- DIREITO INDIVIDUAL}

Divisão po \textsc{DIREITO} INDIVIDUAL. --- O Direito \emph{individual} divide-se em \emph{tres} capitulos: --- no 1.º trata-se dos \emph{deveres} do homem para com \textsc{Deus}; no 2», dos \emph{deveres} do homem para \emph{comsigo}; no 3.º, dos \emph{deveres} e \emph{direitos} do homem para com o \emph{proximo}.

\emph{A}. DEVERES DO HOMEM PARA COM \textsc{DEUS}.

a) \emph{Deveres theoricos do homem pára com Deus}. --- Os \emph{deveres} do homem para com \textsc{Deus} são \emph{theoricos} ou \emph{praticos}, conforme se referem ao \emph{conhecimento} de \textsc{DEus}, ou ao seu \emph{culto}. --- Para cumprir com os seus deveres \emph{theoricos} para com \textsc{Deus}, o homem deve adquirir um conveniente conhecimento das verdades divinas racionaes, e tambem deve acceitar as verdades supraracionaes, se \textsc{Deus} as revelar. D'ahi o absurdo da \emph{liberdade de consciencia} e do \emph{indifferentismo religioso}.

b) \emph{Deveres praticos do homem para com Deus}. --- Para o cumprimento d'esses deveres, o homem deve prestar a \textsc{Deus} o culto \emph{interno} e \emph{externo}, não só o que foi imposto pela natureza, mas tambem o que pode ser determinado pela Revelação sobrenatural.

B, DEVERES DO HOMEM PARA COMSIGO.

a) \emph{Deveres do homem relativos} dá sua \emph{alma}. --- O homem tem o dever de ap»rfeigoar a sua \emph{intelligencia} pelo conhecimento da \emph{verdade}, e a sua \emph{vontade} pela pratica da \emph{virtude}.

b) \emph{Deveres do homem relativos ao seu corpo}. --- O homem tem o dever \pg{820}  de não procurar \emph{directamente} a propria morte pelo suicidio, embora possa

, , por um motivo justo, procural-a \emph{indirectamente}. Além d'isso, deve procurar \emph{positivamente} \& conservação da sua vida e da sua saude.

0. DigEITOS E DEVERES DO HOMEM PARA COM O PROXIMO.

a) \emph{Direitos do homem para com o proximo}. --- Os \emph{direitos} do homem para com o \emph{proximo} são \emph{jurídicos} ou \emph{moraes}, conforme se baséam na idea de \emph{justiça} ou na da \emph{caridade}. --- Os direitos \emph{jurídicos} são \emph{innatos} ou \emph{adquiridos}, conforme nascem com o homem, ou derivam de um facto contingente. Os principaes direitos \emph{innatos} são: o direito á \emph{liberdade religiosa}, \& \emph{dignidade pessoal}, \emph{á vida}, o direito de \emph{legitima defesa}, de \emph{propriedade} (d'onde a falsidade do \emph{communismo} e do \emph{socialismo}), de \emph{independencia} e de \emph{associação}. --- Os direitos moraes reduzem-se ao direito \& \emph{humanidade} e \& \emph{beneficencia}.

b) \emph{Deveres do homem para com o proximo}. --- Os \emph{deveres} do homem para com o \emph{proximo} são \emph{jurídicos} ou \emph{moraes}, conforme derivam da virtude da \emph{Justiça} ou da \emph{caridade}. --- Os deveres \emph{jurídicos} consistem em respeitar o \emph{direito} que os outros têm á \emph{liberdade religiosa}, \& \emph{dignidade pessoal}, \emph{á vida}, \emph{á legitima defeza}, \emph{á propriedade}, \& \emph{independencia} e \& \emph{associação}. --- Os deveres moraes são os que correspondem aos direitos \emph{moraes} do proximo, e abrangem os deveres de \emph{humanidade} e de \emph{beneficencia}.

\part*{SECÇÃO II\\[0.3em]\large DIREITO SOCIAL}
\addcontentsline{toc}{part}{SECÇÃO II --- DIREITO SOCIAL}

\textsc{IwrRopUcQÀo}. --- \emph{Sociedade} é a união de duas ou mais pessoas, que, pelos seus actos, cooperam para um fim commum. --- Os elementos essenciaes da \emph{sociedade} são dois: a \emph{pluralidade} das pessoas o a sua \emph{união}. --- A \emph{união} deriva do \emph{fim} e da \emph{auctoridade}. --- A \emph{auctoridade} é o direito de \emph{dirigir efficazmente} os membros da sociedade para \& consecução do bem commum. = O \emph{fim} não só é principio da união social, mas é tambem o seu \emph{objecto formal}; por isso \emph{especifica} a sociedade, e é a \emph{razão} e a \emph{regra} de todos os seus \emph{deveres} e \emph{direitos}. --- A \emph{sociedade} divide-se em \emph{universal} e \emph{particular}, \emph{simples} e \emph{composta}, \emph{necessaria} e \emph{voluntaria}, \emph{egual} e \emph{desegual}, \emph{completa} e \emph{incompleta}, \emph{perfeita} e \emph{imperfeita}. --- O homem é um ser naturalmente \emph{sociavel}, e foi determinado pela natureza para ser membro das sociedades \emph{domestica}, \emph{civil} e \emph{veligiosa}. --- O Direito \emph{natural} occupa-se \emph{directamente da sociedade domestica} e da \emph{civil}, e só \emph{indirectamente} trata da \emph{religiosa}, que, tendo sido elevada ao estado \emph{sobrenatural}, forma o objecto \emph{directo} da "l'heologia \emph{revelada}.

Divisão po Dimgrro socrat. --- O direito \emph{social} divide-se em \emph{tres} capitulos: --- no 1.º trata-se dos \emph{deveres} e \emph{direitos domesticos}, \emph{e} por isso dos \emph{deveres} e \emph{direitos} na sociedade \emph{conjugal}, na \emph{parental} e na \emph{heril}; --- no 2.º, dos \emph{deveres} e \emph{direitos civis}, --- no 3.º, dos \emph{direitos} e \emph{deveres internacionaes}.

A, DEvERES E DIREITOS DOMESTICOS,

a) \emph{Deveres e direitos na sociedade conjugal}. --- Sociedade \emph{conjugal}, ou \emph{matrimonio}, é a união legitima e indissoluvel de um só homem com uma só mulher, para a geração e educação da prole, e para auxilio de ambos. --- O \emph{matrimonio} é, pois, dotado --- de \emph{unidade}, a que se oppõem a \emph{poligamia} e a \emph{poliandria}, --- e de \emph{indissolubilidade}, a que se oppõe o \emph{divorcio}. --- Os \emph{deveres} e \emph{direitos} communs a ambos os conjuges sao: \emph{fidelidade inviolada}, \emph{amor constante} e \emph{auxilio mutuo}. --- Além d'isso, o \emph{marido} tem o \emph{dever} de dirigir e alimentar a mulher, tratar dos negocios mais importantes da casa, etc., --- e tem o \emph{direito} á obediencia e submissão da mulher. Esta, por sua vez, tem \pg{821}  o \emph{dever} de obedecer ao marido, tratar dos negocios domesticos, etc., e tem o. \emph{direito} ao respeito e á estima do marido.

b) \emph{Deveres e direitos na sociedade parental}. --- Sociedade \emph{parental} é a união entre paes e filhos, instituida para a \emph{educação} d'estes. --- Os paes têm o \emph{dever} de dar aos filhos a educação \emph{physica}, \emph{intellectual} e \emph{moral}, e têm \emph{direito} ao amor, \& gratidão, á reverencia e (se precisarem) ao auxilio dos filhos. --- Os \emph{filhos} têm o \emph{dever} de amar, respeitar, auxiliar os paes, e de obedecer-lhes, emquanto estiverem na casa paterna, --- e têm \emph{direito} a educação \emph{physica}, \emph{intellectual} e \emph{moral}.

o) \emph{Deveres e direitos na sociedade heril}. --- Sociedade \emph{heri} é a união moral do criado e do amo, formada para procurar a utilidade da familia do proprio amo. --- Os amos têm o \emph{dever} de tratar bem os criados, pagar-lhes a soldada, cuidar da sua saude, etc., --- e têm o \emph{direito} ao serviço, amor, respeito e fidelidade dos criados. --- Os \emph{criados} têm o \emph{dever} de prestar aos amos os serviços com zelo, fidelidade, etc.; --- e têm o \emph{direito} á soldada, ao cuidado e respeito dos amos, á instrucção, etc.

B. DEVERES E DIREITOS CIVIS.

a) \emph{Sociedade civil}, \emph{suas propriedades e elementos}. --- Sociedade \emph{civil} é a união de muitas familias, que, pelos seus actos, cooperam para um \emph{fim commum}, que é a. \emph{felicidade temporal}. --- Esta sociedade é \emph{necessaria}, \emph{completa}, \emph{perfeita}, \emph{organica}, \emph{desegual}, --- e compõe-se de dois elementos: um \emph{material}, constituido pelas \emph{familias}, outro \emph{formal}, que é a \emph{auctoridade}. --- A \emph{auctoridade civil}, que é a alma da sociedade, deriva \emph{immediatamente} de \textsc{Deus}, --- pode residir numa pessoa \emph{physica} ou \emph{moral} (d'onde as diversas \emph{formas} de governo) --- e exerce tres funcções: a \emph{legislativa}, a \emph{judicial} e a. \emph{executiva},

b) \emph{Deveres e direitos dos soberanos}. --- Os soberanos tem \emph{deveres} --- para com \textsc{Deus}, emquanto devem fazer com que a sociedade preste a \textsc{Deus} um culto \emph{publico} e \emph{legitimo}, --- para com a \emph{sociedade em geral}, --- para com as \emph{classes sociaes}, para com as \emph{familias} e para com os \emph{individuos}, cujos direitos devem tutelar; e têm \emph{direito} \& veneração, \emph{obediencia} e \emph{fidelidade} dos subditos.

o) \emph{Deveres e direitos dos subditos}. --- Os subditos têm o \emph{dever} de prestar ao soberano veneração, \emph{obediencia} e \emph{fidelidade}; e têm \emph{direitos} correspondentes aos \emph{deveres} do soberano para com elles.

C. DIREITOS E DEVERES INTERNACIONAES.

a) \emph{Direitos e deveres} mutuos \emph{das sociedades civis}. --- Os \emph{direitos} mutuos das sociedades são \emph{jurídicos} e moraes. Os \emph{jurídicos} são \emph{innatos} e \emph{adquiridos}. Os direitos \emph{juridico-innatos} são tres: o direito de conservação, de \emph{aperfeiçoamento} e de \emph{livre actividade}; --- os \emph{juridico-adquiridos} são os que derivam do \emph{uso} ou \emph{costume}, da \emph{acquisição do dominio}, e dos \emph{tratados internacionaes}. Os moraes reduzem-se aos direitos á \emph{humanidade} e á \emph{beneficencia}. --- Os deveres mutuos das sociedades são tambem \emph{jurídicos} e \emph{moraes}.

b) \emph{Deveres da sociedade civil para} com a \emph{Egreja}. --- \emph{Egreja} é a sociedade de homens, unidos entre si pela profissão da mesma verdadeira Fé e pela participação dos mesmos Sacramentos, sob o regimen dos legitimos Pastores e do Romano Pontifice, Vigario de N. S. Jesus Ounisro. E' sociedade divina, \emph{espiritual}, \emph{sobrenatural}, \emph{juridica}, \emph{perfeita} e \emph{suprema}. --- A auctoridade da Egreja é duplice: de \emph{ordem} e de \emph{jurisdicção}. Esta abrange quatro poderes, que são o \emph{doutrinal}, o \emph{legislativo}, o \emph{judicial} e o \emph{executivo}. --- Os \emph{direitos} da, Egreja dividem-se em \emph{internos} e \emph{externos}, conforme se referem aos fieis, ou ás sociedades civis. Os direitos \emph{internos} são: o direito \emph{territorial} \pg{822} , o constitutivo, o de escolher os seus \emph{ministros}, o de \emph{dirigir} e \emph{tutelar} o ensino \emph{religioso}, o de possuir \emph{bens} temporaes, o de \emph{livre communicacdo}. Os direitos \emph{externos} da Egreja são: o direito á independencia, á \emph{superioridade}, \& \emph{immunidade}, á protecção e \emph{assistencia} do Estado. --- Os \emph{deveres} do Estado para com a Egreja são \emph{negativos} e \emph{affirmativos}. Os \emph{negativos} consistem em respeitar todos os direitos da Egreja, internos e \emph{externos}. Os \emph{positivos} consistem em auxiliar e favorecer a Egreja por um modo \emph{positivo}. Estes \emph{deveres} são \emph{jurídicos}; porque correspondem aos direitos, de que a Egreja foi dotada por Nosso Senhor \textsc{JESUS CHnisTO}.

ES be \pg{823}  \textbf{INDICE} DAS \textbf{MATERIAS}

\textbf{METAPHYSICA} ;

%%>>> parte-5

\part*{PARTE TERCEIRA\\[0.3em]\large ANTHROPOLOGIA}
\addcontentsline{toc}{part}{PARTE TERCEIRA --- ANTHROPOLOGIA}

Pag. 5

SECÇÃO PRIMEIRA --- Essencia do homem

%%>>> cap-12

\chapter*{CAPITULO PRIMEIRO}
\addcontentsline{toc}{chapter}{CAPITULO PRIMEIRO}

\emph{Corpo humano} Descripção do corpo humano . n 10 Orgãos da vida vegetativa --- . : x \emph{Anas} \$5 RED Org\&os da sensibilidade externa : n 90 Orgãos da sensibilidade interna e da locomoção s n 26

%%>>> cap-13

\chapter*{CAPITULO SEGUNDO}
\addcontentsline{toc}{chapter}{CAPITULO SEGUNDO}

Alma \emph{humana} Alma humana, sua existencia e substancialidade . 8. e Simplicidade da alma humana » 95 Espiritualidade da alma humana E POS A. espiritualidade da alma humana e o materialismo ;. o» AS

%%>>> cap-14

\chapter*{CAPITULO TERCEIRO}
\addcontentsline{toc}{chapter}{CAPITULO TERCEIRO}

\emph{Unido da alma com o corpo} União da alma com o corpo e suas consequencias 41 \emph{no} 160) Unidade da alma humana Miu ES : : 2.5 We ET Séde da alma humana . : \emph{HL} . . EE.

SECÇÃO SEGUNDA --- Faculdades do homem

%%>>> cap-15

\chapter*{CAPITULO PRIMEIRO}
\addcontentsline{toc}{chapter}{CAPITULO PRIMEIRO}

\emph{Faculdades humanas em geral} As faculdades consideradas em si mesmas As faculdades consideradas em relação á alma As faculdades consideradas em relação ás operações

%%>>> cap-16

\chapter*{CAPITULO SEGUNDO}
\addcontentsline{toc}{chapter}{CAPITULO SEGUNDO}

\pg{824} \emph{Faculdades sensitivas do homem} Sentidos externos Sentidos internos : Appetite sensitivo, faculdade io uEODEN e 55\textasciicircum{}

%%>>> cap-17

\chapter*{CAPITULO TERCEIRO}
\addcontentsline{toc}{chapter}{CAPITULO TERCEIRO}

\emph{Faculdades intellectuaes do homem} Intelligencia, sua natureza e varios nomes Vontade, sua natureza e liberdade . A liberdade da vontade humana e o determinismo Relações entre a intelligencia e a vontade

SECÇÃO TERCEIRA --- Operações do homem

%%>>> cap-18

\chapter*{CAPITULO PRIMEIRO}
\addcontentsline{toc}{chapter}{CAPITULO PRIMEIRO}

\emph{Operações sensitivas do homem} Percepção sensitiva, seu objecto e divisão Percepção sensitiva externa Percepção sensitiva interna ; : Appetite sensitivo, sua divisão e paixões

%%>>> cap-19

\chapter*{CAPITULO SEGUNDO}
\addcontentsline{toc}{chapter}{CAPITULO SEGUNDO}

\emph{Operações intellectuaes do homem} Conhecimento intellectual, sua origem . : : Modo por que se produz o conhecimento intellectual : Modo por que se desenvolve o conhecimento intellectual Modo por que se exprime o conhecimento intellectual Volições, seu objecto e divisão

SECÇÃO \textsc{QUARTA} --- \textbf{Origem} e destino do \textbf{homem}

%%>>> cap-20

\chapter*{CAPITULO PRIMEIRO}
\addcontentsline{toc}{chapter}{CAPITULO PRIMEIRO}

\emph{Origem do homem} =

Origem do corpo humano Origem da alma humana Origem do composto humano .

%%>>> cap-21

\chapter*{CAPITULO SEGUNDO}
\addcontentsline{toc}{chapter}{CAPITULO SEGUNDO}

\emph{Destino do homem} Fim do homem . : 5 d ; ; Immortalidade da alma lins X . Condição e operações da alma separada do corpo.

%%>>> parte-6

\part*{PARTE QUARTA\\[0.3em]\large THEODICEA}
\addcontentsline{toc}{part}{PARTE QUARTA --- THEODICEA}

\pg{825} Prologo

%%>>> cap-22

\chapter*{CAPITULO PRIMEIRO}
\addcontentsline{toc}{chapter}{CAPITULO PRIMEIRO}

\emph{Ecistencia de Deus}

Anr. L. Demonstrabilidade da existencia de \textsc{Deus} Anr. Il... Demonstração da existencia de \textsc{Deus} Amr. III. A existencia de \textsc{Deus} e o atheismo

%%>>> cap-23

\chapter*{CAPITULO SEGUNDO}
\addcontentsline{toc}{chapter}{CAPITULO SEGUNDO}

\emph{Essencia} e \emph{attributos de Deus} AnT. Ll. Essencia de \textsc{Deus}, seu constitutivo e attributos Amr. II. Attributos de \textsc{Deus} negativos . Amr. III. Attributos de \textsc{Deus} positivos . j

%%>>> cap-24

\chapter*{CAPITULO TERCEIRO}
\addcontentsline{toc}{chapter}{CAPITULO TERCEIRO}

\emph{Faculdades e operações immanentes de Deus} Amr. l. Intelligencia de \textsc{DEus}, sua natureza e objecto Amr. II. Faculdade e operação volitiva immanente Amr. III. Vida divina, sua existencia o natureza . 7

%%>>> cap-25

\chapter*{CAPITULO QUARTO}
\addcontentsline{toc}{chapter}{CAPITULO QUARTO}

\emph{Operações transeuntes de Deus} ArT. I. --- Creação, sua natureza, principio e fim . ArrT. II. Conservação, sua divisão e necessidade . Amr. III. Concurso, sua natureza e extensão . Amr. IV. Providencia, sua existencia, extensão e proprisdades

\textbf{MORAL}

Proemio

%%>>> parte-7

\part*{PARTE PRIMEIRA\\[0.3em]\large ETHICA}
\addcontentsline{toc}{part}{PARTE PRIMEIRA --- ETHICA}

%%>>> cap-26

\chapter*{CAPITULO PRIMEIRO}
\addcontentsline{toc}{chapter}{CAPITULO PRIMEIRO}

\emph{Fim dos actos humanos} Art. I. Existencia do fim ultimo dos actos humanos Anr. II. Natureza do fim ultimo dos actos humanos . Anr, III. Fim proximo dos actos humanos .. : õ

%%>>> cap-27

\chapter*{CAPITULO SEGUNDO}
\addcontentsline{toc}{chapter}{CAPITULO SEGUNDO}

\pg{826} \pg{827} -\emph{Essencia dos actos humanos} Elementos e divisão dos actos humanos Diversas especies do acto elicito : Impedimentos da voluntariedade dos actos Jamidnos

OAPITULO TERCEIRO

\emph{Moralidade dos actos humanos} Moralidade, sua existencia o fundamento, Fontes da moralidade dos actos humanos ; Consequencias da moralidade dos actos humanos. A moralidade dos actos humanos e os habitos

%%>>> cap-28

\chapter*{CAPITULO QUARTO}
\addcontentsline{toc}{chapter}{CAPITULO QUARTO}

\emph{Regras dos actos humanos} Lei moral, sua divisão e existencia Consciencia moral, sua divisão e direcção

SECÇÃO PRIMEIRA --- \textbf{Direito} individual

%%>>> cap-29

\chapter*{CAPITULO PRIMEIRO}
\addcontentsline{toc}{chapter}{CAPITULO PRIMEIRO}

\emph{Deveres do homem para com Deus} Deveres theoricos do homem para com \textsc{Deus} Deveres praticos do homem para com \textsc{Deus} .

%%>>> cap-30

\chapter*{CAPITULO SEGUNDO}
\addcontentsline{toc}{chapter}{CAPITULO SEGUNDO}

\emph{Deveres do homem para comsigo} Deveres do homem relativos á sua alma Deveres do homem relativos ao seu corpo

CAPITULO TERCEIRO \emph{Direitos e deveres do homem para com o proximo}

Direitos do homem para com o proximo Deveres do homem para com o proximo
